\documentclass[11pt,a4paper]{article}
\usepackage{isabelle,isabellesym}

% further packages required for unusual symbols (see also
% isabellesym.sty), use only when needed

%\usepackage{amssymb}
  %for \<leadsto>, \<box>, \<diamond>, \<sqsupset>, \<mho>, \<Join>,
  %\<lhd>, \<lesssim>, \<greatersim>, \<lessapprox>, \<greaterapprox>,
  %\<triangleq>, \<yen>, \<lozenge>

%\usepackage{eurosym}
  %for \<euro>

%\usepackage[only,bigsqcap]{stmaryrd}
  %for \<Sqinter>

%\usepackage{eufrak}
  %for \<AA> ... \<ZZ>, \<aa> ... \<zz> (also included in amssymb)

%\usepackage{textcomp}
  %for \<onequarter>, \<onehalf>, \<threequarters>, \<degree>, \<cent>,
  %\<currency>

\usepackage{wasysym}  
  
% this should be the last package used
\usepackage{pdfsetup}

% urls in roman style, theory text in math-similar italics
\urlstyle{rm}
\isabellestyle{it}

% for uniform font size
%\renewcommand{\isastyle}{\isastyleminor}


\begin{document}

\title{Temporal Planning Semantics}
% \subtitle{Proof Document}
\author{Mohammad Abdulaziz \and Lukas Koller}
\date{}

\maketitle

This is a formalisation of the semantics of temporal planning as used in the Planning Domain
Definition Language. Based on that, we also develop a plan validator which is verified w.r.t.\ those
semantics. It is based on an earlier formalisation of the semantics of classical planning~\cite{classicalPlanningSemAFP}.
 This work was described in an ealier publication~\cite{TemporalSemantics}. 
\tableofcontents

\clearpage

% sane default for proof documents
\parindent 0pt\parskip 0.5ex

% generated text of all theories
%
\begin{isabellebody}%
\setisabellecontext{TEMPORAL{\isacharunderscore}{\kern0pt}PDDL{\isacharunderscore}{\kern0pt}Semantics}%
%
\isadelimdocument
%
\endisadelimdocument
%
\isatagdocument
%
\isamarkupsection{Temporal PDDL Semantics%
}
\isamarkuptrue%
%
\endisatagdocument
{\isafolddocument}%
%
\isadelimdocument
%
\endisadelimdocument
%
\isadelimtheory
%
\endisadelimtheory
%
\isatagtheory
\isacommand{theory}\isamarkupfalse%
\ TEMPORAL{\isacharunderscore}{\kern0pt}PDDL{\isacharunderscore}{\kern0pt}Semantics\isanewline
\isakeyword{imports}\isanewline
\ \ {\isachardoublequoteopen}Propositional{\isacharunderscore}{\kern0pt}Proof{\isacharunderscore}{\kern0pt}Systems{\isachardot}{\kern0pt}Formulas{\isachardoublequoteclose}\isanewline
\ \ {\isachardoublequoteopen}Propositional{\isacharunderscore}{\kern0pt}Proof{\isacharunderscore}{\kern0pt}Systems{\isachardot}{\kern0pt}Sema{\isachardoublequoteclose}\isanewline
\ \ {\isachardoublequoteopen}Propositional{\isacharunderscore}{\kern0pt}Proof{\isacharunderscore}{\kern0pt}Systems{\isachardot}{\kern0pt}Consistency{\isachardoublequoteclose}\isanewline
\ \ {\isachardoublequoteopen}Automatic{\isacharunderscore}{\kern0pt}Refinement{\isachardot}{\kern0pt}Misc{\isachardoublequoteclose}\isanewline
\ \ {\isachardoublequoteopen}Automatic{\isacharunderscore}{\kern0pt}Refinement{\isachardot}{\kern0pt}Refine{\isacharunderscore}{\kern0pt}Util{\isachardoublequoteclose}\isanewline
\ \ {\isachardoublequoteopen}HOL{\isachardot}{\kern0pt}Rat{\isachardoublequoteclose}\isanewline
\ \ \ \isanewline
\isakeyword{begin}%
\endisatagtheory
{\isafoldtheory}%
%
\isadelimtheory
\isanewline
%
\endisadelimtheory
\isacommand{no{\isacharunderscore}{\kern0pt}notation}\isamarkupfalse%
\ insert\ {\isacharparenleft}{\kern0pt}{\isachardoublequoteopen}{\isacharunderscore}{\kern0pt}\ {\isasymtriangleright}\ {\isacharunderscore}{\kern0pt}{\isachardoublequoteclose}\ {\isacharbrackleft}{\kern0pt}{\isadigit{5}}{\isadigit{6}}{\isacharcomma}{\kern0pt}{\isadigit{5}}{\isadigit{5}}{\isacharbrackright}{\kern0pt}\ {\isadigit{5}}{\isadigit{5}}{\isacharparenright}{\kern0pt}%
\isadelimdocument
%
\endisadelimdocument
%
\isatagdocument
%
\isamarkupsubsection{Utility Functions%
}
\isamarkuptrue%
%
\endisatagdocument
{\isafolddocument}%
%
\isadelimdocument
%
\endisadelimdocument
\isacommand{abbreviation}\isamarkupfalse%
\ strict{\isacharunderscore}{\kern0pt}sorted{\isacharcolon}{\kern0pt}{\isacharcolon}{\kern0pt}{\isachardoublequoteopen}{\isacharprime}{\kern0pt}a{\isacharcolon}{\kern0pt}{\isacharcolon}{\kern0pt}{\isacharbraceleft}{\kern0pt}linorder{\isacharbraceright}{\kern0pt}\ list\ {\isasymRightarrow}\ bool{\isachardoublequoteclose}\ \isakeyword{where}\ {\isachardoublequoteopen}strict{\isacharunderscore}{\kern0pt}sorted\ {\isasymequiv}\ sorted{\isacharunderscore}{\kern0pt}wrt\ {\isacharparenleft}{\kern0pt}{\isacharless}{\kern0pt}{\isacharparenright}{\kern0pt}{\isachardoublequoteclose}\isanewline
\isanewline
\isacommand{definition}\isamarkupfalse%
\ {\isachardoublequoteopen}index{\isacharunderscore}{\kern0pt}by\ f\ l\ {\isasymequiv}\ map{\isacharunderscore}{\kern0pt}of\ {\isacharparenleft}{\kern0pt}map\ {\isacharparenleft}{\kern0pt}{\isasymlambda}x{\isachardot}{\kern0pt}\ {\isacharparenleft}{\kern0pt}f\ x{\isacharcomma}{\kern0pt}x{\isacharparenright}{\kern0pt}{\isacharparenright}{\kern0pt}\ l{\isacharparenright}{\kern0pt}{\isachardoublequoteclose}\isanewline
\isanewline
\isacommand{lemma}\isamarkupfalse%
\ index{\isacharunderscore}{\kern0pt}by{\isacharunderscore}{\kern0pt}eq{\isacharunderscore}{\kern0pt}Some{\isacharunderscore}{\kern0pt}eq{\isacharbrackleft}{\kern0pt}simp{\isacharbrackright}{\kern0pt}{\isacharcolon}{\kern0pt}\isanewline
\ \ \isakeyword{assumes}\ {\isachardoublequoteopen}distinct\ {\isacharparenleft}{\kern0pt}map\ f\ l{\isacharparenright}{\kern0pt}{\isachardoublequoteclose}\isanewline
\ \ \isakeyword{shows}\ {\isachardoublequoteopen}index{\isacharunderscore}{\kern0pt}by\ f\ l\ n\ {\isacharequal}{\kern0pt}\ Some\ x\ {\isasymlongleftrightarrow}\ {\isacharparenleft}{\kern0pt}x{\isasymin}set\ l\ {\isasymand}\ f\ x\ {\isacharequal}{\kern0pt}\ n{\isacharparenright}{\kern0pt}{\isachardoublequoteclose}\isanewline
%
\isadelimproof
\ \ %
\endisadelimproof
%
\isatagproof
\isacommand{unfolding}\isamarkupfalse%
\ index{\isacharunderscore}{\kern0pt}by{\isacharunderscore}{\kern0pt}def\isanewline
\ \ \isacommand{using}\isamarkupfalse%
\ assms\isanewline
\ \ \isacommand{by}\isamarkupfalse%
\ {\isacharparenleft}{\kern0pt}auto\ simp{\isacharcolon}{\kern0pt}\ o{\isacharunderscore}{\kern0pt}def{\isacharparenright}{\kern0pt}%
\endisatagproof
{\isafoldproof}%
%
\isadelimproof
\isanewline
%
\endisadelimproof
\isanewline
\isacommand{lemma}\isamarkupfalse%
\ index{\isacharunderscore}{\kern0pt}by{\isacharunderscore}{\kern0pt}eq{\isacharunderscore}{\kern0pt}SomeD{\isacharcolon}{\kern0pt}\isanewline
\ \ \isakeyword{shows}\ {\isachardoublequoteopen}index{\isacharunderscore}{\kern0pt}by\ f\ l\ n\ {\isacharequal}{\kern0pt}\ Some\ x\ {\isasymLongrightarrow}\ {\isacharparenleft}{\kern0pt}x{\isasymin}set\ l\ {\isasymand}\ f\ x\ {\isacharequal}{\kern0pt}\ n{\isacharparenright}{\kern0pt}{\isachardoublequoteclose}\isanewline
%
\isadelimproof
\ \ %
\endisadelimproof
%
\isatagproof
\isacommand{unfolding}\isamarkupfalse%
\ index{\isacharunderscore}{\kern0pt}by{\isacharunderscore}{\kern0pt}def\isanewline
\ \ \isacommand{by}\isamarkupfalse%
\ {\isacharparenleft}{\kern0pt}auto\ dest{\isacharcolon}{\kern0pt}\ map{\isacharunderscore}{\kern0pt}of{\isacharunderscore}{\kern0pt}SomeD{\isacharparenright}{\kern0pt}%
\endisatagproof
{\isafoldproof}%
%
\isadelimproof
\isanewline
%
\endisadelimproof
\isanewline
\isanewline
\isacommand{lemma}\isamarkupfalse%
\ lookup{\isacharunderscore}{\kern0pt}zip{\isacharunderscore}{\kern0pt}idx{\isacharunderscore}{\kern0pt}eq{\isacharcolon}{\kern0pt}\isanewline
\ \ \isakeyword{assumes}\ {\isachardoublequoteopen}length\ params\ {\isacharequal}{\kern0pt}\ length\ args{\isachardoublequoteclose}\isanewline
\ \ \isakeyword{assumes}\ {\isachardoublequoteopen}i{\isacharless}{\kern0pt}length\ args{\isachardoublequoteclose}\isanewline
\ \ \isakeyword{assumes}\ {\isachardoublequoteopen}distinct\ params{\isachardoublequoteclose}\isanewline
\ \ \isakeyword{assumes}\ {\isachardoublequoteopen}k\ {\isacharequal}{\kern0pt}\ params\ {\isacharbang}{\kern0pt}\ i{\isachardoublequoteclose}\isanewline
\ \ \isakeyword{shows}\ {\isachardoublequoteopen}map{\isacharunderscore}{\kern0pt}of\ {\isacharparenleft}{\kern0pt}zip\ params\ args{\isacharparenright}{\kern0pt}\ k\ {\isacharequal}{\kern0pt}\ Some\ {\isacharparenleft}{\kern0pt}args\ {\isacharbang}{\kern0pt}\ i{\isacharparenright}{\kern0pt}{\isachardoublequoteclose}\isanewline
%
\isadelimproof
\ \ %
\endisadelimproof
%
\isatagproof
\isacommand{using}\isamarkupfalse%
\ assms\isanewline
\ \ \isacommand{by}\isamarkupfalse%
\ {\isacharparenleft}{\kern0pt}auto\ simp{\isacharcolon}{\kern0pt}\ in{\isacharunderscore}{\kern0pt}set{\isacharunderscore}{\kern0pt}conv{\isacharunderscore}{\kern0pt}nth{\isacharparenright}{\kern0pt}%
\endisatagproof
{\isafoldproof}%
%
\isadelimproof
\isanewline
%
\endisadelimproof
\isanewline
\isacommand{lemma}\isamarkupfalse%
\ rtrancl{\isacharunderscore}{\kern0pt}image{\isacharunderscore}{\kern0pt}idem{\isacharbrackleft}{\kern0pt}simp{\isacharbrackright}{\kern0pt}{\isacharcolon}{\kern0pt}\ {\isachardoublequoteopen}R\isactrlsup {\isacharasterisk}{\kern0pt}\ {\isacharbackquote}{\kern0pt}{\isacharbackquote}{\kern0pt}\ R\isactrlsup {\isacharasterisk}{\kern0pt}\ {\isacharbackquote}{\kern0pt}{\isacharbackquote}{\kern0pt}\ s\ {\isacharequal}{\kern0pt}\ R\isactrlsup {\isacharasterisk}{\kern0pt}\ {\isacharbackquote}{\kern0pt}{\isacharbackquote}{\kern0pt}\ s{\isachardoublequoteclose}\isanewline
%
\isadelimproof
\ \ %
\endisadelimproof
%
\isatagproof
\isacommand{by}\isamarkupfalse%
\ {\isacharparenleft}{\kern0pt}metis\ relcomp{\isacharunderscore}{\kern0pt}Image\ rtrancl{\isacharunderscore}{\kern0pt}idemp{\isacharunderscore}{\kern0pt}self{\isacharunderscore}{\kern0pt}comp{\isacharparenright}{\kern0pt}%
\endisatagproof
{\isafoldproof}%
%
\isadelimproof
%
\endisadelimproof
%
\isadelimdocument
%
\endisadelimdocument
%
\isatagdocument
%
\isamarkupsubsection{Abstract Syntax%
}
\isamarkuptrue%
%
\isamarkupsubsubsection{Generic Entities%
}
\isamarkuptrue%
%
\endisatagdocument
{\isafolddocument}%
%
\isadelimdocument
%
\endisadelimdocument
%
\begin{isamarkuptext}%
Defines the time number type%
\end{isamarkuptext}\isamarkuptrue%
\isacommand{type{\isacharunderscore}{\kern0pt}synonym}\isamarkupfalse%
\ time\ {\isacharequal}{\kern0pt}\ rat\isanewline
\isanewline
\isacommand{type{\isacharunderscore}{\kern0pt}synonym}\isamarkupfalse%
\ name\ {\isacharequal}{\kern0pt}\ string\isanewline
\isanewline
\isacommand{datatype}\isamarkupfalse%
\ predicate\ {\isacharequal}{\kern0pt}\ Pred\ {\isacharparenleft}{\kern0pt}name{\isacharcolon}{\kern0pt}\ name{\isacharparenright}{\kern0pt}\isanewline
\isanewline
\isacommand{datatype}\isamarkupfalse%
\ func\ {\isacharequal}{\kern0pt}\ Func\ {\isacharparenleft}{\kern0pt}name{\isacharcolon}{\kern0pt}\ name{\isacharparenright}{\kern0pt}%
\begin{isamarkuptext}%
Some of the AST entities are defined over a polymorphic \isa{{\isacharprime}{\kern0pt}val} type,
  which gets either instantiated by variables (for domains)
  or objects (for problems).%
\end{isamarkuptext}\isamarkuptrue%
%
\begin{isamarkuptext}%
An atom is either a predicate with arguments, or an equality statement.%
\end{isamarkuptext}\isamarkuptrue%
\isacommand{datatype}\isamarkupfalse%
\ {\isacharprime}{\kern0pt}ent\ atom\ {\isacharequal}{\kern0pt}\ predAtm\ {\isacharparenleft}{\kern0pt}predicate{\isacharcolon}{\kern0pt}\ predicate{\isacharparenright}{\kern0pt}\ {\isacharparenleft}{\kern0pt}arguments{\isacharcolon}{\kern0pt}\ {\isachardoublequoteopen}{\isacharprime}{\kern0pt}ent\ list{\isachardoublequoteclose}{\isacharparenright}{\kern0pt}\isanewline
\ \ \ \ \ \ \ \ \ \ \ \ \ \ \ \ \ \ \ {\isacharbar}{\kern0pt}\ eqAtm\ {\isacharparenleft}{\kern0pt}lhs{\isacharcolon}{\kern0pt}\ {\isacharprime}{\kern0pt}ent{\isacharparenright}{\kern0pt}\ {\isacharparenleft}{\kern0pt}rhs{\isacharcolon}{\kern0pt}\ {\isacharprime}{\kern0pt}ent{\isacharparenright}{\kern0pt}%
\begin{isamarkuptext}%
A type is a list of primitive type names.
  To model a primitive type, we use a singleton list.%
\end{isamarkuptext}\isamarkuptrue%
\isacommand{datatype}\isamarkupfalse%
\ type\ {\isacharequal}{\kern0pt}\ Either\ {\isacharparenleft}{\kern0pt}primitives{\isacharcolon}{\kern0pt}\ {\isachardoublequoteopen}name\ list{\isachardoublequoteclose}{\isacharparenright}{\kern0pt}%
\begin{isamarkuptext}%
An effect contains a list of values to be added, and a list of values
  to be removed.%
\end{isamarkuptext}\isamarkuptrue%
\isacommand{datatype}\isamarkupfalse%
\ {\isacharprime}{\kern0pt}ent\ ast{\isacharunderscore}{\kern0pt}effect\ {\isacharequal}{\kern0pt}\ Effect\ {\isacharparenleft}{\kern0pt}adds{\isacharcolon}{\kern0pt}\ {\isachardoublequoteopen}{\isacharparenleft}{\kern0pt}{\isacharprime}{\kern0pt}ent\ atom\ formula{\isacharparenright}{\kern0pt}\ list{\isachardoublequoteclose}{\isacharparenright}{\kern0pt}\isanewline
\ \ \ \ \ \ \ \ \ \ \ \ \ \ \ \ \ \ \ \ \ \ \ \ \ \ \ \ \ \ \ \ \ \ {\isacharparenleft}{\kern0pt}dels{\isacharcolon}{\kern0pt}\ {\isachardoublequoteopen}{\isacharparenleft}{\kern0pt}{\isacharprime}{\kern0pt}ent\ atom\ formula{\isacharparenright}{\kern0pt}\ list{\isachardoublequoteclose}{\isacharparenright}{\kern0pt}%
\begin{isamarkuptext}%
Variables are identified by their names.%
\end{isamarkuptext}\isamarkuptrue%
\isacommand{datatype}\isamarkupfalse%
\ variable\ {\isacharequal}{\kern0pt}\ Var\ {\isacharparenleft}{\kern0pt}varname{\isacharcolon}{\kern0pt}\ name{\isacharparenright}{\kern0pt}%
\begin{isamarkuptext}%
Objects and constants are identified by their names%
\end{isamarkuptext}\isamarkuptrue%
\isacommand{datatype}\isamarkupfalse%
\ object\ {\isacharequal}{\kern0pt}\ Obj\ {\isacharparenleft}{\kern0pt}name{\isacharcolon}{\kern0pt}\ name{\isacharparenright}{\kern0pt}\ {\isacharbar}{\kern0pt}\ FuncEnt\ {\isacharparenleft}{\kern0pt}func{\isacharcolon}{\kern0pt}\ func{\isacharparenright}{\kern0pt}\ {\isacharparenleft}{\kern0pt}arguments{\isacharcolon}{\kern0pt}\ {\isachardoublequoteopen}object\ list{\isachardoublequoteclose}{\isacharparenright}{\kern0pt}\ {\isacharbar}{\kern0pt}\ TimeEnt\ time\isanewline
\isanewline
\isacommand{datatype}\isamarkupfalse%
\ {\isachardoublequoteopen}term{\isachardoublequoteclose}\ {\isacharequal}{\kern0pt}\ VAR\ variable\ {\isacharbar}{\kern0pt}\ CONST\ object\isanewline
\isanewline
\isacommand{hide{\isacharunderscore}{\kern0pt}const}\isamarkupfalse%
\ {\isacharparenleft}{\kern0pt}\isakeyword{open}{\isacharparenright}{\kern0pt}\ VAR\ CONST\ %
\isamarkupcmt{Refer to constructors by qualified names only%
}%
\isadelimdocument
%
\endisadelimdocument
%
\isatagdocument
%
\isamarkupsubsubsection{Domains%
}
\isamarkuptrue%
%
\endisatagdocument
{\isafolddocument}%
%
\isadelimdocument
%
\endisadelimdocument
%
\begin{isamarkuptext}%
A time specifier specifies the time point of conditions and effects in a durative action%
\end{isamarkuptext}\isamarkuptrue%
\isacommand{datatype}\isamarkupfalse%
\ temporal{\isacharunderscore}{\kern0pt}annotation\ {\isacharequal}{\kern0pt}\ At{\isacharunderscore}{\kern0pt}Start\ {\isacharbar}{\kern0pt}\ Over{\isacharunderscore}{\kern0pt}All\ {\isacharbar}{\kern0pt}\ At{\isacharunderscore}{\kern0pt}End\isanewline
\isanewline
\isacommand{datatype}\isamarkupfalse%
\ duration{\isacharunderscore}{\kern0pt}op\ {\isacharequal}{\kern0pt}\ EQ\ {\isacharbar}{\kern0pt}\ LEQ\ {\isacharbar}{\kern0pt}\ GEQ\isanewline
\isanewline
\isacommand{datatype}\isamarkupfalse%
\ {\isacharprime}{\kern0pt}ent\ duration{\isacharunderscore}{\kern0pt}constraint\ {\isacharequal}{\kern0pt}\ \isanewline
\ \ No{\isacharunderscore}{\kern0pt}Const\ \isanewline
{\isacharbar}{\kern0pt}\ Time{\isacharunderscore}{\kern0pt}Const\ {\isacharparenleft}{\kern0pt}d{\isacharunderscore}{\kern0pt}op{\isacharcolon}{\kern0pt}\ duration{\isacharunderscore}{\kern0pt}op{\isacharparenright}{\kern0pt}\ time\ \isanewline
{\isacharbar}{\kern0pt}\ Func{\isacharunderscore}{\kern0pt}Const\ {\isacharparenleft}{\kern0pt}d{\isacharunderscore}{\kern0pt}op{\isacharcolon}{\kern0pt}\ duration{\isacharunderscore}{\kern0pt}op{\isacharparenright}{\kern0pt}\ {\isacharparenleft}{\kern0pt}func{\isacharcolon}{\kern0pt}\ func{\isacharparenright}{\kern0pt}\ {\isacharparenleft}{\kern0pt}arguments{\isacharcolon}{\kern0pt}\ {\isachardoublequoteopen}{\isacharprime}{\kern0pt}ent\ list{\isachardoublequoteclose}{\isacharparenright}{\kern0pt}%
\begin{isamarkuptext}%
An action schema has a name, a typed parameter list, a precondition,
  and an effect. In case of a durative action schema there is additionally a 
  duration constraint, and the conditions and effects are timed with a time specifier%
\end{isamarkuptext}\isamarkuptrue%
\isacommand{datatype}\isamarkupfalse%
\ ast{\isacharunderscore}{\kern0pt}action{\isacharunderscore}{\kern0pt}schema\ {\isacharequal}{\kern0pt}\isanewline
\ \ Simple{\isacharunderscore}{\kern0pt}Action{\isacharunderscore}{\kern0pt}Schema\isanewline
\ \ \ \ {\isacharparenleft}{\kern0pt}name{\isacharcolon}{\kern0pt}\ name{\isacharparenright}{\kern0pt}\isanewline
\ \ \ \ {\isacharparenleft}{\kern0pt}parameters{\isacharcolon}{\kern0pt}\ {\isachardoublequoteopen}{\isacharparenleft}{\kern0pt}variable\ {\isasymtimes}\ type{\isacharparenright}{\kern0pt}\ list{\isachardoublequoteclose}{\isacharparenright}{\kern0pt}\isanewline
\ \ \ \ {\isacharparenleft}{\kern0pt}precondition{\isacharcolon}{\kern0pt}\ {\isachardoublequoteopen}term\ atom\ formula{\isachardoublequoteclose}{\isacharparenright}{\kern0pt}\isanewline
\ \ \ \ {\isacharparenleft}{\kern0pt}effect{\isacharcolon}{\kern0pt}\ {\isachardoublequoteopen}term\ ast{\isacharunderscore}{\kern0pt}effect{\isachardoublequoteclose}{\isacharparenright}{\kern0pt}\isanewline
{\isacharbar}{\kern0pt}\ Durative{\isacharunderscore}{\kern0pt}Action{\isacharunderscore}{\kern0pt}Schema\isanewline
\ \ \ \ {\isacharparenleft}{\kern0pt}name{\isacharcolon}{\kern0pt}\ name{\isacharparenright}{\kern0pt}\isanewline
\ \ \ \ {\isacharparenleft}{\kern0pt}parameters{\isacharcolon}{\kern0pt}\ {\isachardoublequoteopen}{\isacharparenleft}{\kern0pt}variable\ {\isasymtimes}\ type{\isacharparenright}{\kern0pt}\ list{\isachardoublequoteclose}{\isacharparenright}{\kern0pt}\ \ \ \ \ \ \ \ \ \ \ \ \ \ \ \ \ \ \ \ \ \ \ \ \ \ \ \ \ \isanewline
\ \ \ \ {\isacharparenleft}{\kern0pt}duration{\isacharunderscore}{\kern0pt}constraint{\isacharcolon}{\kern0pt}\ {\isachardoublequoteopen}term\ duration{\isacharunderscore}{\kern0pt}constraint{\isachardoublequoteclose}{\isacharparenright}{\kern0pt}\isanewline
\ \ \ \ {\isacharparenleft}{\kern0pt}condition{\isacharcolon}{\kern0pt}\ {\isachardoublequoteopen}{\isacharparenleft}{\kern0pt}temporal{\isacharunderscore}{\kern0pt}annotation\ {\isasymtimes}\ {\isacharparenleft}{\kern0pt}term\ atom\ formula{\isacharparenright}{\kern0pt}{\isacharparenright}{\kern0pt}\ list{\isachardoublequoteclose}{\isacharparenright}{\kern0pt}\isanewline
\ \ \ \ {\isacharparenleft}{\kern0pt}durative{\isacharunderscore}{\kern0pt}effect{\isacharcolon}{\kern0pt}\ {\isachardoublequoteopen}{\isacharparenleft}{\kern0pt}temporal{\isacharunderscore}{\kern0pt}annotation\ {\isasymtimes}\ {\isacharparenleft}{\kern0pt}term\ ast{\isacharunderscore}{\kern0pt}effect{\isacharparenright}{\kern0pt}{\isacharparenright}{\kern0pt}\ list{\isachardoublequoteclose}{\isacharparenright}{\kern0pt}%
\begin{isamarkuptext}%
A predicate declaration contains the predicate's name and its
  argument types.%
\end{isamarkuptext}\isamarkuptrue%
\isacommand{datatype}\isamarkupfalse%
\ predicate{\isacharunderscore}{\kern0pt}decl\ {\isacharequal}{\kern0pt}\ PredDecl\ {\isacharparenleft}{\kern0pt}pred{\isacharcolon}{\kern0pt}\ predicate{\isacharparenright}{\kern0pt}\ {\isacharparenleft}{\kern0pt}argTs{\isacharcolon}{\kern0pt}\ {\isachardoublequoteopen}type\ list{\isachardoublequoteclose}{\isacharparenright}{\kern0pt}\isanewline
\isanewline
\isacommand{datatype}\isamarkupfalse%
\ function{\isacharunderscore}{\kern0pt}decl\ {\isacharequal}{\kern0pt}\ FuncDecl\ {\isacharparenleft}{\kern0pt}func{\isacharcolon}{\kern0pt}\ func{\isacharparenright}{\kern0pt}\ {\isacharparenleft}{\kern0pt}argTs{\isacharcolon}{\kern0pt}\ {\isachardoublequoteopen}type\ list{\isachardoublequoteclose}{\isacharparenright}{\kern0pt}%
\begin{isamarkuptext}%
A domain contains the declarations of primitive types, predicates,
  and action schemas.%
\end{isamarkuptext}\isamarkuptrue%
\isacommand{datatype}\isamarkupfalse%
\ ast{\isacharunderscore}{\kern0pt}domain\ {\isacharequal}{\kern0pt}\ Domain\isanewline
\ \ {\isacharparenleft}{\kern0pt}types{\isacharcolon}{\kern0pt}\ {\isachardoublequoteopen}{\isacharparenleft}{\kern0pt}name\ {\isasymtimes}\ name{\isacharparenright}{\kern0pt}\ list{\isachardoublequoteclose}{\isacharparenright}{\kern0pt}\ %
\isamarkupcmt{\isa{{\isacharparenleft}{\kern0pt}type{\isacharcomma}{\kern0pt}\ supertype{\isacharparenright}{\kern0pt}} declarations.%
}\isanewline
\ \ {\isacharparenleft}{\kern0pt}predicates{\isacharcolon}{\kern0pt}\ {\isachardoublequoteopen}predicate{\isacharunderscore}{\kern0pt}decl\ list{\isachardoublequoteclose}{\isacharparenright}{\kern0pt}\isanewline
\ \ {\isacharparenleft}{\kern0pt}{\isachardoublequoteopen}functions{\isachardoublequoteclose}{\isacharcolon}{\kern0pt}\ {\isachardoublequoteopen}function{\isacharunderscore}{\kern0pt}decl\ list{\isachardoublequoteclose}{\isacharparenright}{\kern0pt}\isanewline
\ \ {\isacharparenleft}{\kern0pt}{\isachardoublequoteopen}consts{\isachardoublequoteclose}{\isacharcolon}{\kern0pt}\ {\isachardoublequoteopen}{\isacharparenleft}{\kern0pt}object\ {\isasymtimes}\ type{\isacharparenright}{\kern0pt}\ list{\isachardoublequoteclose}{\isacharparenright}{\kern0pt}\isanewline
\ \ {\isacharparenleft}{\kern0pt}actions{\isacharcolon}{\kern0pt}\ {\isachardoublequoteopen}ast{\isacharunderscore}{\kern0pt}action{\isacharunderscore}{\kern0pt}schema\ list{\isachardoublequoteclose}{\isacharparenright}{\kern0pt}%
\isadelimdocument
%
\endisadelimdocument
%
\isatagdocument
%
\isamarkupsubsubsection{Problems%
}
\isamarkuptrue%
%
\endisatagdocument
{\isafolddocument}%
%
\isadelimdocument
%
\endisadelimdocument
%
\begin{isamarkuptext}%
A fact is a predicate applied to objects.%
\end{isamarkuptext}\isamarkuptrue%
\isacommand{type{\isacharunderscore}{\kern0pt}synonym}\isamarkupfalse%
\ fact\ {\isacharequal}{\kern0pt}\ {\isachardoublequoteopen}predicate\ {\isasymtimes}\ object\ list{\isachardoublequoteclose}%
\begin{isamarkuptext}%
A problem consists of a domain, a list of objects,
  a description of the initial state, and a description of the goal state.%
\end{isamarkuptext}\isamarkuptrue%
\isacommand{datatype}\isamarkupfalse%
\ ast{\isacharunderscore}{\kern0pt}problem\ {\isacharequal}{\kern0pt}\ Problem\isanewline
\ \ {\isacharparenleft}{\kern0pt}domain{\isacharcolon}{\kern0pt}\ ast{\isacharunderscore}{\kern0pt}domain{\isacharparenright}{\kern0pt}\isanewline
\ \ {\isacharparenleft}{\kern0pt}objects{\isacharcolon}{\kern0pt}\ {\isachardoublequoteopen}{\isacharparenleft}{\kern0pt}object\ {\isasymtimes}\ type{\isacharparenright}{\kern0pt}\ list{\isachardoublequoteclose}{\isacharparenright}{\kern0pt}\isanewline
\ \ {\isacharparenleft}{\kern0pt}init{\isacharcolon}{\kern0pt}\ {\isachardoublequoteopen}object\ atom\ formula\ list{\isachardoublequoteclose}{\isacharparenright}{\kern0pt}\isanewline
\ \ {\isacharparenleft}{\kern0pt}goal{\isacharcolon}{\kern0pt}\ {\isachardoublequoteopen}object\ atom\ formula{\isachardoublequoteclose}{\isacharparenright}{\kern0pt}%
\isadelimdocument
%
\endisadelimdocument
%
\isatagdocument
%
\isamarkupsubsubsection{Plans%
}
\isamarkuptrue%
%
\endisatagdocument
{\isafolddocument}%
%
\isadelimdocument
%
\endisadelimdocument
\isacommand{datatype}\isamarkupfalse%
\ plan{\isacharunderscore}{\kern0pt}action\ {\isacharequal}{\kern0pt}\ \isanewline
\ \ Simple{\isacharunderscore}{\kern0pt}Plan{\isacharunderscore}{\kern0pt}Action\ {\isacharparenleft}{\kern0pt}name{\isacharcolon}{\kern0pt}\ name{\isacharparenright}{\kern0pt}\ {\isacharparenleft}{\kern0pt}arguments{\isacharcolon}{\kern0pt}\ {\isachardoublequoteopen}object\ list{\isachardoublequoteclose}{\isacharparenright}{\kern0pt}\isanewline
{\isacharbar}{\kern0pt}\ Durative{\isacharunderscore}{\kern0pt}Plan{\isacharunderscore}{\kern0pt}Action\ {\isacharparenleft}{\kern0pt}name{\isacharcolon}{\kern0pt}\ name{\isacharparenright}{\kern0pt}\ {\isacharparenleft}{\kern0pt}arguments{\isacharcolon}{\kern0pt}\ {\isachardoublequoteopen}object\ list{\isachardoublequoteclose}{\isacharparenright}{\kern0pt}\ {\isacharparenleft}{\kern0pt}duration{\isacharcolon}{\kern0pt}\ time{\isacharparenright}{\kern0pt}\isanewline
\isanewline
\isacommand{type{\isacharunderscore}{\kern0pt}synonym}\isamarkupfalse%
\ plan\ {\isacharequal}{\kern0pt}\ {\isachardoublequoteopen}{\isacharparenleft}{\kern0pt}time\ {\isasymtimes}\ plan{\isacharunderscore}{\kern0pt}action{\isacharparenright}{\kern0pt}\ list{\isachardoublequoteclose}%
\isadelimdocument
%
\endisadelimdocument
%
\isatagdocument
%
\isamarkupsubsubsection{Ground Actions%
}
\isamarkuptrue%
%
\endisatagdocument
{\isafolddocument}%
%
\isadelimdocument
%
\endisadelimdocument
%
\begin{isamarkuptext}%
The following datatype represents an action scheme that has been
  instantiated by replacing the arguments with concrete objects,
  also called ground action.%
\end{isamarkuptext}\isamarkuptrue%
\isacommand{datatype}\isamarkupfalse%
\ ground{\isacharunderscore}{\kern0pt}action\ {\isacharequal}{\kern0pt}\ Ground{\isacharunderscore}{\kern0pt}Action\isanewline
\ \ {\isacharparenleft}{\kern0pt}precondition{\isacharcolon}{\kern0pt}\ {\isachardoublequoteopen}{\isacharparenleft}{\kern0pt}object\ atom{\isacharparenright}{\kern0pt}\ formula{\isachardoublequoteclose}{\isacharparenright}{\kern0pt}\isanewline
\ \ {\isacharparenleft}{\kern0pt}effect{\isacharcolon}{\kern0pt}\ {\isachardoublequoteopen}object\ ast{\isacharunderscore}{\kern0pt}effect{\isachardoublequoteclose}{\isacharparenright}{\kern0pt}%
\begin{isamarkuptext}%
A happening represents a collection of grounded actions that are all
   executed simultaneously at a specified time point.%
\end{isamarkuptext}\isamarkuptrue%
\isacommand{type{\isacharunderscore}{\kern0pt}synonym}\isamarkupfalse%
\ happening\ {\isacharequal}{\kern0pt}\ {\isachardoublequoteopen}time\ {\isasymtimes}\ {\isacharparenleft}{\kern0pt}ground{\isacharunderscore}{\kern0pt}action\ list{\isacharparenright}{\kern0pt}{\isachardoublequoteclose}%
\isadelimdocument
%
\endisadelimdocument
%
\isatagdocument
%
\isamarkupsubsection{Closed-World Assumption, Equality, and Negation%
}
\isamarkuptrue%
%
\endisatagdocument
{\isafolddocument}%
%
\isadelimdocument
%
\endisadelimdocument
%
\begin{isamarkuptext}%
Discriminator for atomic predicate formulas.%
\end{isamarkuptext}\isamarkuptrue%
\ \ \isacommand{fun}\isamarkupfalse%
\ is{\isacharunderscore}{\kern0pt}predAtom\ \isakeyword{where}\isanewline
\ \ \ \ {\isachardoublequoteopen}is{\isacharunderscore}{\kern0pt}predAtom\ {\isacharparenleft}{\kern0pt}Atom\ {\isacharparenleft}{\kern0pt}predAtm\ {\isacharunderscore}{\kern0pt}\ {\isacharunderscore}{\kern0pt}{\isacharparenright}{\kern0pt}{\isacharparenright}{\kern0pt}\ {\isacharequal}{\kern0pt}\ True{\isachardoublequoteclose}\ {\isacharbar}{\kern0pt}\ {\isachardoublequoteopen}is{\isacharunderscore}{\kern0pt}predAtom\ {\isacharunderscore}{\kern0pt}\ {\isacharequal}{\kern0pt}\ False{\isachardoublequoteclose}\isanewline
\isanewline
\ \ \isacommand{fun}\isamarkupfalse%
\ is{\isacharunderscore}{\kern0pt}eqAtom\ \isakeyword{where}\isanewline
\ \ \ \ {\isachardoublequoteopen}is{\isacharunderscore}{\kern0pt}eqAtom\ {\isacharparenleft}{\kern0pt}Atom\ {\isacharparenleft}{\kern0pt}eqAtm\ {\isacharunderscore}{\kern0pt}\ {\isacharunderscore}{\kern0pt}{\isacharparenright}{\kern0pt}{\isacharparenright}{\kern0pt}\ {\isacharequal}{\kern0pt}\ True{\isachardoublequoteclose}\ {\isacharbar}{\kern0pt}\ {\isachardoublequoteopen}is{\isacharunderscore}{\kern0pt}eqAtom\ {\isacharunderscore}{\kern0pt}\ {\isacharequal}{\kern0pt}\ False{\isachardoublequoteclose}%
\begin{isamarkuptext}%
The world model is a set of (atomic) formulas%
\end{isamarkuptext}\isamarkuptrue%
\ \ \isacommand{type{\isacharunderscore}{\kern0pt}synonym}\isamarkupfalse%
\ world{\isacharunderscore}{\kern0pt}model\ {\isacharequal}{\kern0pt}\ {\isachardoublequoteopen}object\ atom\ formula\ set{\isachardoublequoteclose}%
\begin{isamarkuptext}%
It is basic, if it only contains atoms%
\end{isamarkuptext}\isamarkuptrue%
\ \ \isacommand{definition}\isamarkupfalse%
\ {\isachardoublequoteopen}wm{\isacharunderscore}{\kern0pt}basic\ M\ {\isasymequiv}\ {\isasymforall}a{\isasymin}M{\isachardot}{\kern0pt}\ is{\isacharunderscore}{\kern0pt}predAtom\ a{\isachardoublequoteclose}%
\begin{isamarkuptext}%
A valuation extracted from the atoms of the world model%
\end{isamarkuptext}\isamarkuptrue%
\ \ \isacommand{definition}\isamarkupfalse%
\ valuation\ {\isacharcolon}{\kern0pt}{\isacharcolon}{\kern0pt}\ {\isachardoublequoteopen}world{\isacharunderscore}{\kern0pt}model\ {\isasymRightarrow}\ object\ atom\ valuation{\isachardoublequoteclose}\isanewline
\ \ \ \ \isakeyword{where}\ {\isachardoublequoteopen}valuation\ M\ {\isasymequiv}\ {\isasymlambda}predAtm\ p\ xs\ {\isasymRightarrow}\ Atom\ {\isacharparenleft}{\kern0pt}predAtm\ p\ xs{\isacharparenright}{\kern0pt}\ {\isasymin}\ M\ {\isacharbar}{\kern0pt}\ eqAtm\ a\ b\ {\isasymRightarrow}\ a{\isacharequal}{\kern0pt}b{\isachardoublequoteclose}%
\begin{isamarkuptext}%
Augment a world model by adding negated versions of all atoms
    not contained in it, as well as interpretations of equality.%
\end{isamarkuptext}\isamarkuptrue%
\ \ \isacommand{definition}\isamarkupfalse%
\ close{\isacharunderscore}{\kern0pt}world\ {\isacharcolon}{\kern0pt}{\isacharcolon}{\kern0pt}\ {\isachardoublequoteopen}world{\isacharunderscore}{\kern0pt}model\ {\isasymRightarrow}\ world{\isacharunderscore}{\kern0pt}model{\isachardoublequoteclose}\ \isakeyword{where}\isanewline
\ \ \ \ {\isachardoublequoteopen}close{\isacharunderscore}{\kern0pt}world\ M\ {\isacharequal}{\kern0pt}\isanewline
\ \ \ \ \ \ M\ {\isasymunion}\ {\isacharbraceleft}{\kern0pt}\isactrlbold {\isasymnot}{\isacharparenleft}{\kern0pt}Atom\ {\isacharparenleft}{\kern0pt}predAtm\ p\ as{\isacharparenright}{\kern0pt}{\isacharparenright}{\kern0pt}\ {\isacharbar}{\kern0pt}\ p\ as{\isachardot}{\kern0pt}\ Atom\ {\isacharparenleft}{\kern0pt}predAtm\ p\ as{\isacharparenright}{\kern0pt}\ {\isasymnotin}\ M{\isacharbraceright}{\kern0pt}\isanewline
\ \ \ \ \ \ {\isasymunion}\ {\isacharbraceleft}{\kern0pt}Atom\ {\isacharparenleft}{\kern0pt}eqAtm\ a\ a{\isacharparenright}{\kern0pt}\ {\isacharbar}{\kern0pt}\ a{\isachardot}{\kern0pt}\ True{\isacharbraceright}{\kern0pt}\ {\isasymunion}\ {\isacharbraceleft}{\kern0pt}\isactrlbold {\isasymnot}{\isacharparenleft}{\kern0pt}Atom\ {\isacharparenleft}{\kern0pt}eqAtm\ a\ b{\isacharparenright}{\kern0pt}{\isacharparenright}{\kern0pt}\ {\isacharbar}{\kern0pt}\ a\ b{\isachardot}{\kern0pt}\ a{\isasymnoteq}b{\isacharbraceright}{\kern0pt}{\isachardoublequoteclose}\isanewline
\isanewline
\ \ \isacommand{definition}\isamarkupfalse%
\ {\isachardoublequoteopen}close{\isacharunderscore}{\kern0pt}neg\ M\ {\isasymequiv}\ M\ {\isasymunion}\ {\isacharbraceleft}{\kern0pt}\isactrlbold {\isasymnot}{\isacharparenleft}{\kern0pt}Atom\ a{\isacharparenright}{\kern0pt}\ {\isacharbar}{\kern0pt}\ a{\isachardot}{\kern0pt}\ Atom\ a\ {\isasymnotin}\ M{\isacharbraceright}{\kern0pt}{\isachardoublequoteclose}\isanewline
\ \ \isacommand{lemma}\isamarkupfalse%
\ {\isachardoublequoteopen}wm{\isacharunderscore}{\kern0pt}basic\ M\ {\isasymLongrightarrow}\ close{\isacharunderscore}{\kern0pt}world\ M\ {\isacharequal}{\kern0pt}\ close{\isacharunderscore}{\kern0pt}neg\ {\isacharparenleft}{\kern0pt}M\ {\isasymunion}\ {\isacharbraceleft}{\kern0pt}Atom\ {\isacharparenleft}{\kern0pt}eqAtm\ a\ a{\isacharparenright}{\kern0pt}\ {\isacharbar}{\kern0pt}\ a{\isachardot}{\kern0pt}\ True{\isacharbraceright}{\kern0pt}{\isacharparenright}{\kern0pt}{\isachardoublequoteclose}\isanewline
%
\isadelimproof
\ \ \ \ %
\endisadelimproof
%
\isatagproof
\isacommand{unfolding}\isamarkupfalse%
\ close{\isacharunderscore}{\kern0pt}world{\isacharunderscore}{\kern0pt}def\ close{\isacharunderscore}{\kern0pt}neg{\isacharunderscore}{\kern0pt}def\ wm{\isacharunderscore}{\kern0pt}basic{\isacharunderscore}{\kern0pt}def\isanewline
\ \ \ \ \isacommand{apply}\isamarkupfalse%
\ clarsimp\isanewline
\ \ \ \ \isacommand{apply}\isamarkupfalse%
\ {\isacharparenleft}{\kern0pt}auto\ {\isadigit{0}}\ {\isadigit{3}}{\isacharparenright}{\kern0pt}\isanewline
\ \ \ \ \isacommand{by}\isamarkupfalse%
\ {\isacharparenleft}{\kern0pt}metis\ atom{\isachardot}{\kern0pt}exhaust{\isacharparenright}{\kern0pt}%
\endisatagproof
{\isafoldproof}%
%
\isadelimproof
\isanewline
%
\endisadelimproof
\isanewline
\ \ \isacommand{abbreviation}\isamarkupfalse%
\ cw{\isacharunderscore}{\kern0pt}entailment\ {\isacharparenleft}{\kern0pt}\isakeyword{infix}\ {\isachardoublequoteopen}\isactrlsup c{\isasymTTurnstile}\isactrlsub {\isacharequal}{\kern0pt}{\isachardoublequoteclose}\ {\isadigit{5}}{\isadigit{3}}{\isacharparenright}{\kern0pt}\ \isakeyword{where}\isanewline
\ \ \ \ {\isachardoublequoteopen}M\ \isactrlsup c{\isasymTTurnstile}\isactrlsub {\isacharequal}{\kern0pt}\ {\isasymphi}\ {\isasymequiv}\ close{\isacharunderscore}{\kern0pt}world\ M\ {\isasymTTurnstile}\ {\isasymphi}{\isachardoublequoteclose}\isanewline
\isanewline
\ \ \isacommand{lemma}\isamarkupfalse%
\isanewline
\ \ \ \ close{\isacharunderscore}{\kern0pt}world{\isacharunderscore}{\kern0pt}extensive{\isacharcolon}{\kern0pt}\ {\isachardoublequoteopen}M\ {\isasymsubseteq}\ close{\isacharunderscore}{\kern0pt}world\ M{\isachardoublequoteclose}\ \isakeyword{and}\isanewline
\ \ \ \ close{\isacharunderscore}{\kern0pt}world{\isacharunderscore}{\kern0pt}idem{\isacharbrackleft}{\kern0pt}simp{\isacharbrackright}{\kern0pt}{\isacharcolon}{\kern0pt}\ {\isachardoublequoteopen}close{\isacharunderscore}{\kern0pt}world\ {\isacharparenleft}{\kern0pt}close{\isacharunderscore}{\kern0pt}world\ M{\isacharparenright}{\kern0pt}\ {\isacharequal}{\kern0pt}\ close{\isacharunderscore}{\kern0pt}world\ M{\isachardoublequoteclose}\isanewline
%
\isadelimproof
\ \ \ \ %
\endisadelimproof
%
\isatagproof
\isacommand{by}\isamarkupfalse%
\ {\isacharparenleft}{\kern0pt}auto\ simp{\isacharcolon}{\kern0pt}\ close{\isacharunderscore}{\kern0pt}world{\isacharunderscore}{\kern0pt}def{\isacharparenright}{\kern0pt}%
\endisatagproof
{\isafoldproof}%
%
\isadelimproof
\isanewline
%
\endisadelimproof
\isanewline
\ \ \isacommand{lemma}\isamarkupfalse%
\ in{\isacharunderscore}{\kern0pt}close{\isacharunderscore}{\kern0pt}world{\isacharunderscore}{\kern0pt}conv{\isacharcolon}{\kern0pt}\isanewline
\ \ \ \ {\isachardoublequoteopen}{\isasymphi}\ {\isasymin}\ close{\isacharunderscore}{\kern0pt}world\ M\ {\isasymlongleftrightarrow}\ {\isacharparenleft}{\kern0pt}\isanewline
\ \ \ \ \ \ \ \ {\isasymphi}{\isasymin}M\isanewline
\ \ \ \ \ \ {\isasymor}\ {\isacharparenleft}{\kern0pt}{\isasymexists}p\ as{\isachardot}{\kern0pt}\ {\isasymphi}{\isacharequal}{\kern0pt}\isactrlbold {\isasymnot}{\isacharparenleft}{\kern0pt}Atom\ {\isacharparenleft}{\kern0pt}predAtm\ p\ as{\isacharparenright}{\kern0pt}{\isacharparenright}{\kern0pt}\ {\isasymand}\ Atom\ {\isacharparenleft}{\kern0pt}predAtm\ p\ as{\isacharparenright}{\kern0pt}\ {\isasymnotin}\ M{\isacharparenright}{\kern0pt}\isanewline
\ \ \ \ \ \ {\isasymor}\ {\isacharparenleft}{\kern0pt}{\isasymexists}a{\isachardot}{\kern0pt}\ {\isasymphi}{\isacharequal}{\kern0pt}Atom\ {\isacharparenleft}{\kern0pt}eqAtm\ a\ a{\isacharparenright}{\kern0pt}{\isacharparenright}{\kern0pt}\isanewline
\ \ \ \ \ \ {\isasymor}\ {\isacharparenleft}{\kern0pt}{\isasymexists}a\ b{\isachardot}{\kern0pt}\ {\isasymphi}{\isacharequal}{\kern0pt}\isactrlbold {\isasymnot}{\isacharparenleft}{\kern0pt}Atom\ {\isacharparenleft}{\kern0pt}eqAtm\ a\ b{\isacharparenright}{\kern0pt}{\isacharparenright}{\kern0pt}\ {\isasymand}\ a{\isasymnoteq}b{\isacharparenright}{\kern0pt}\isanewline
\ \ \ \ {\isacharparenright}{\kern0pt}{\isachardoublequoteclose}\isanewline
%
\isadelimproof
\ \ \ \ %
\endisadelimproof
%
\isatagproof
\isacommand{by}\isamarkupfalse%
\ {\isacharparenleft}{\kern0pt}auto\ simp{\isacharcolon}{\kern0pt}\ close{\isacharunderscore}{\kern0pt}world{\isacharunderscore}{\kern0pt}def{\isacharparenright}{\kern0pt}%
\endisatagproof
{\isafoldproof}%
%
\isadelimproof
\isanewline
%
\endisadelimproof
\isanewline
\ \ \isacommand{lemma}\isamarkupfalse%
\ valuation{\isacharunderscore}{\kern0pt}aux{\isacharunderscore}{\kern0pt}{\isadigit{1}}{\isacharcolon}{\kern0pt}\isanewline
\ \ \ \ \isakeyword{fixes}\ M\ {\isacharcolon}{\kern0pt}{\isacharcolon}{\kern0pt}\ world{\isacharunderscore}{\kern0pt}model\ \isakeyword{and}\ {\isasymphi}\ {\isacharcolon}{\kern0pt}{\isacharcolon}{\kern0pt}\ {\isachardoublequoteopen}object\ atom\ formula{\isachardoublequoteclose}\isanewline
\ \ \ \ \isakeyword{defines}\ {\isachardoublequoteopen}C\ {\isasymequiv}\ close{\isacharunderscore}{\kern0pt}world\ M{\isachardoublequoteclose}\isanewline
\ \ \ \ \isakeyword{assumes}\ A{\isacharcolon}{\kern0pt}\ {\isachardoublequoteopen}{\isasymforall}{\isasymphi}{\isasymin}C{\isachardot}{\kern0pt}\ {\isasymA}\ {\isasymTurnstile}\ {\isasymphi}{\isachardoublequoteclose}\isanewline
\ \ \ \ \isakeyword{shows}\ {\isachardoublequoteopen}{\isasymA}\ {\isacharequal}{\kern0pt}\ valuation\ M{\isachardoublequoteclose}\isanewline
%
\isadelimproof
\ \ \ \ %
\endisadelimproof
%
\isatagproof
\isacommand{using}\isamarkupfalse%
\ A\ \isacommand{unfolding}\isamarkupfalse%
\ C{\isacharunderscore}{\kern0pt}def\isanewline
\ \ \ \ \isacommand{apply}\isamarkupfalse%
\ {\isacharparenleft}{\kern0pt}auto\ simp{\isacharcolon}{\kern0pt}\ in{\isacharunderscore}{\kern0pt}close{\isacharunderscore}{\kern0pt}world{\isacharunderscore}{\kern0pt}conv\ valuation{\isacharunderscore}{\kern0pt}def\ Ball{\isacharunderscore}{\kern0pt}def\ intro{\isacharbang}{\kern0pt}{\isacharcolon}{\kern0pt}\ ext\ split{\isacharcolon}{\kern0pt}\ atom{\isachardot}{\kern0pt}split{\isacharparenright}{\kern0pt}\isanewline
\ \ \ \ \isacommand{apply}\isamarkupfalse%
\ {\isacharparenleft}{\kern0pt}metis\ formula{\isacharunderscore}{\kern0pt}semantics{\isachardot}{\kern0pt}simps{\isacharparenleft}{\kern0pt}{\isadigit{1}}{\isacharparenright}{\kern0pt}\ formula{\isacharunderscore}{\kern0pt}semantics{\isachardot}{\kern0pt}simps{\isacharparenleft}{\kern0pt}{\isadigit{3}}{\isacharparenright}{\kern0pt}{\isacharparenright}{\kern0pt}\isanewline
\ \ \ \ \isacommand{apply}\isamarkupfalse%
\ {\isacharparenleft}{\kern0pt}metis\ formula{\isacharunderscore}{\kern0pt}semantics{\isachardot}{\kern0pt}simps{\isacharparenleft}{\kern0pt}{\isadigit{1}}{\isacharparenright}{\kern0pt}\ formula{\isacharunderscore}{\kern0pt}semantics{\isachardot}{\kern0pt}simps{\isacharparenleft}{\kern0pt}{\isadigit{3}}{\isacharparenright}{\kern0pt}{\isacharparenright}{\kern0pt}\isanewline
\ \ \ \ \isacommand{apply}\isamarkupfalse%
\ {\isacharparenleft}{\kern0pt}metis\ atom{\isachardot}{\kern0pt}collapse{\isacharparenleft}{\kern0pt}{\isadigit{2}}{\isacharparenright}{\kern0pt}\ formula{\isacharunderscore}{\kern0pt}semantics{\isachardot}{\kern0pt}simps{\isacharparenleft}{\kern0pt}{\isadigit{1}}{\isacharparenright}{\kern0pt}\ is{\isacharunderscore}{\kern0pt}predAtm{\isacharunderscore}{\kern0pt}def{\isacharparenright}{\kern0pt}\isanewline
\ \ \ \ \isacommand{done}\isamarkupfalse%
%
\endisatagproof
{\isafoldproof}%
%
\isadelimproof
\isanewline
%
\endisadelimproof
\isanewline
\ \ \isacommand{lemma}\isamarkupfalse%
\ valuation{\isacharunderscore}{\kern0pt}aux{\isacharunderscore}{\kern0pt}{\isadigit{2}}{\isacharcolon}{\kern0pt}\isanewline
\ \ \ \ \isakeyword{assumes}\ {\isachardoublequoteopen}wm{\isacharunderscore}{\kern0pt}basic\ M{\isachardoublequoteclose}\isanewline
\ \ \ \ \isakeyword{shows}\ {\isachardoublequoteopen}{\isacharparenleft}{\kern0pt}{\isasymforall}G{\isasymin}close{\isacharunderscore}{\kern0pt}world\ M{\isachardot}{\kern0pt}\ valuation\ M\ {\isasymTurnstile}\ G{\isacharparenright}{\kern0pt}{\isachardoublequoteclose}\isanewline
%
\isadelimproof
\ \ \ \ %
\endisadelimproof
%
\isatagproof
\isacommand{using}\isamarkupfalse%
\ assms\ \isacommand{unfolding}\isamarkupfalse%
\ wm{\isacharunderscore}{\kern0pt}basic{\isacharunderscore}{\kern0pt}def\isanewline
\ \ \ \ \isacommand{by}\isamarkupfalse%
\ {\isacharparenleft}{\kern0pt}force\ simp{\isacharcolon}{\kern0pt}\ in{\isacharunderscore}{\kern0pt}close{\isacharunderscore}{\kern0pt}world{\isacharunderscore}{\kern0pt}conv\ valuation{\isacharunderscore}{\kern0pt}def\ elim{\isacharcolon}{\kern0pt}\ is{\isacharunderscore}{\kern0pt}predAtom{\isachardot}{\kern0pt}elims\ is{\isacharunderscore}{\kern0pt}eqAtom{\isachardot}{\kern0pt}elims{\isacharparenright}{\kern0pt}%
\endisatagproof
{\isafoldproof}%
%
\isadelimproof
\isanewline
%
\endisadelimproof
\isanewline
\ \ \isacommand{lemma}\isamarkupfalse%
\ val{\isacharunderscore}{\kern0pt}imp{\isacharunderscore}{\kern0pt}close{\isacharunderscore}{\kern0pt}world{\isacharcolon}{\kern0pt}\ {\isachardoublequoteopen}valuation\ M\ {\isasymTurnstile}\ {\isasymphi}\ {\isasymLongrightarrow}\ M\ \isactrlsup c{\isasymTTurnstile}\isactrlsub {\isacharequal}{\kern0pt}\ {\isasymphi}{\isachardoublequoteclose}\isanewline
%
\isadelimproof
\ \ \ \ %
\endisadelimproof
%
\isatagproof
\isacommand{unfolding}\isamarkupfalse%
\ entailment{\isacharunderscore}{\kern0pt}def\isanewline
\ \ \ \ \isacommand{using}\isamarkupfalse%
\ valuation{\isacharunderscore}{\kern0pt}aux{\isacharunderscore}{\kern0pt}{\isadigit{1}}\isanewline
\ \ \ \ \isacommand{by}\isamarkupfalse%
\ blast%
\endisatagproof
{\isafoldproof}%
%
\isadelimproof
\isanewline
%
\endisadelimproof
\isanewline
\ \ \isacommand{lemma}\isamarkupfalse%
\ close{\isacharunderscore}{\kern0pt}world{\isacharunderscore}{\kern0pt}imp{\isacharunderscore}{\kern0pt}val{\isacharcolon}{\kern0pt}\isanewline
\ \ \ \ {\isachardoublequoteopen}wm{\isacharunderscore}{\kern0pt}basic\ M\ {\isasymLongrightarrow}\ M\ \isactrlsup c{\isasymTTurnstile}\isactrlsub {\isacharequal}{\kern0pt}\ {\isasymphi}\ {\isasymLongrightarrow}\ valuation\ M\ {\isasymTurnstile}\ {\isasymphi}{\isachardoublequoteclose}\isanewline
%
\isadelimproof
\ \ \ \ %
\endisadelimproof
%
\isatagproof
\isacommand{unfolding}\isamarkupfalse%
\ entailment{\isacharunderscore}{\kern0pt}def\ \isacommand{using}\isamarkupfalse%
\ valuation{\isacharunderscore}{\kern0pt}aux{\isacharunderscore}{\kern0pt}{\isadigit{2}}\ \isacommand{by}\isamarkupfalse%
\ blast%
\endisatagproof
{\isafoldproof}%
%
\isadelimproof
%
\endisadelimproof
%
\begin{isamarkuptext}%
Main theorem of this section:
    If a world model \isa{M} contains only atoms, its induced valuation
    satisfies a formula \isa{{\isasymphi}} if and only if the closure of \isa{M} entails \isa{{\isasymphi}}.

    Note that there are no syntactic restrictions on \isa{{\isasymphi}},
    in particular, \isa{{\isasymphi}} may contain negation.%
\end{isamarkuptext}\isamarkuptrue%
\ \ \isacommand{theorem}\isamarkupfalse%
\ valuation{\isacharunderscore}{\kern0pt}iff{\isacharunderscore}{\kern0pt}close{\isacharunderscore}{\kern0pt}world{\isacharcolon}{\kern0pt}\isanewline
\ \ \ \ \isakeyword{assumes}\ {\isachardoublequoteopen}wm{\isacharunderscore}{\kern0pt}basic\ M{\isachardoublequoteclose}\isanewline
\ \ \ \ \isakeyword{shows}\ {\isachardoublequoteopen}valuation\ M\ {\isasymTurnstile}\ {\isasymphi}\ {\isasymlongleftrightarrow}\ M\ \isactrlsup c{\isasymTTurnstile}\isactrlsub {\isacharequal}{\kern0pt}\ {\isasymphi}{\isachardoublequoteclose}\isanewline
%
\isadelimproof
\ \ \ \ %
\endisadelimproof
%
\isatagproof
\isacommand{using}\isamarkupfalse%
\ assms\ val{\isacharunderscore}{\kern0pt}imp{\isacharunderscore}{\kern0pt}close{\isacharunderscore}{\kern0pt}world\ close{\isacharunderscore}{\kern0pt}world{\isacharunderscore}{\kern0pt}imp{\isacharunderscore}{\kern0pt}val\ \isacommand{by}\isamarkupfalse%
\ blast%
\endisatagproof
{\isafoldproof}%
%
\isadelimproof
%
\endisadelimproof
%
\isadelimdocument
%
\endisadelimdocument
%
\isatagdocument
%
\isamarkupsubsubsection{Proper Generalization%
}
\isamarkuptrue%
%
\endisatagdocument
{\isafolddocument}%
%
\isadelimdocument
%
\endisadelimdocument
%
\begin{isamarkuptext}%
Adding negation and equality is a proper generalization of the
  case without negation and equality%
\end{isamarkuptext}\isamarkuptrue%
\isacommand{fun}\isamarkupfalse%
\ is{\isacharunderscore}{\kern0pt}STRIPS{\isacharunderscore}{\kern0pt}fmla\ {\isacharcolon}{\kern0pt}{\isacharcolon}{\kern0pt}\ {\isachardoublequoteopen}{\isacharprime}{\kern0pt}ent\ atom\ formula\ {\isasymRightarrow}\ bool{\isachardoublequoteclose}\ \isakeyword{where}\isanewline
\ \ {\isachardoublequoteopen}is{\isacharunderscore}{\kern0pt}STRIPS{\isacharunderscore}{\kern0pt}fmla\ {\isacharparenleft}{\kern0pt}Atom\ {\isacharparenleft}{\kern0pt}predAtm\ {\isacharunderscore}{\kern0pt}\ {\isacharunderscore}{\kern0pt}{\isacharparenright}{\kern0pt}{\isacharparenright}{\kern0pt}\ {\isasymlongleftrightarrow}\ True{\isachardoublequoteclose}\isanewline
{\isacharbar}{\kern0pt}\ {\isachardoublequoteopen}is{\isacharunderscore}{\kern0pt}STRIPS{\isacharunderscore}{\kern0pt}fmla\ {\isacharparenleft}{\kern0pt}{\isasymbottom}{\isacharparenright}{\kern0pt}\ {\isasymlongleftrightarrow}\ True{\isachardoublequoteclose}\isanewline
{\isacharbar}{\kern0pt}\ {\isachardoublequoteopen}is{\isacharunderscore}{\kern0pt}STRIPS{\isacharunderscore}{\kern0pt}fmla\ {\isacharparenleft}{\kern0pt}{\isasymphi}\isactrlsub {\isadigit{1}}\ \isactrlbold {\isasymand}\ {\isasymphi}\isactrlsub {\isadigit{2}}{\isacharparenright}{\kern0pt}\ {\isasymlongleftrightarrow}\ is{\isacharunderscore}{\kern0pt}STRIPS{\isacharunderscore}{\kern0pt}fmla\ {\isasymphi}\isactrlsub {\isadigit{1}}\ {\isasymand}\ is{\isacharunderscore}{\kern0pt}STRIPS{\isacharunderscore}{\kern0pt}fmla\ {\isasymphi}\isactrlsub {\isadigit{2}}{\isachardoublequoteclose}\isanewline
{\isacharbar}{\kern0pt}\ {\isachardoublequoteopen}is{\isacharunderscore}{\kern0pt}STRIPS{\isacharunderscore}{\kern0pt}fmla\ {\isacharparenleft}{\kern0pt}{\isasymphi}\isactrlsub {\isadigit{1}}\ \isactrlbold {\isasymor}\ {\isasymphi}\isactrlsub {\isadigit{2}}{\isacharparenright}{\kern0pt}\ {\isasymlongleftrightarrow}\ is{\isacharunderscore}{\kern0pt}STRIPS{\isacharunderscore}{\kern0pt}fmla\ {\isasymphi}\isactrlsub {\isadigit{1}}\ {\isasymand}\ is{\isacharunderscore}{\kern0pt}STRIPS{\isacharunderscore}{\kern0pt}fmla\ {\isasymphi}\isactrlsub {\isadigit{2}}{\isachardoublequoteclose}\isanewline
{\isacharbar}{\kern0pt}\ {\isachardoublequoteopen}is{\isacharunderscore}{\kern0pt}STRIPS{\isacharunderscore}{\kern0pt}fmla\ {\isacharparenleft}{\kern0pt}\isactrlbold {\isasymnot}{\isasymbottom}{\isacharparenright}{\kern0pt}\ {\isasymlongleftrightarrow}\ True{\isachardoublequoteclose}\isanewline
{\isacharbar}{\kern0pt}\ {\isachardoublequoteopen}is{\isacharunderscore}{\kern0pt}STRIPS{\isacharunderscore}{\kern0pt}fmla\ {\isacharunderscore}{\kern0pt}\ {\isasymlongleftrightarrow}\ False{\isachardoublequoteclose}\isanewline
\isanewline
\isacommand{lemma}\isamarkupfalse%
\ aux{\isadigit{1}}{\isacharcolon}{\kern0pt}\ {\isachardoublequoteopen}{\isasymlbrakk}wm{\isacharunderscore}{\kern0pt}basic\ M{\isacharsemicolon}{\kern0pt}\ is{\isacharunderscore}{\kern0pt}STRIPS{\isacharunderscore}{\kern0pt}fmla\ {\isasymphi}{\isacharsemicolon}{\kern0pt}\ valuation\ M\ {\isasymTurnstile}\ {\isasymphi}{\isacharsemicolon}{\kern0pt}\ {\isasymforall}G{\isasymin}M{\isachardot}{\kern0pt}\ {\isasymA}\ {\isasymTurnstile}\ G{\isasymrbrakk}\ {\isasymLongrightarrow}\ {\isasymA}\ {\isasymTurnstile}\ {\isasymphi}{\isachardoublequoteclose}\isanewline
%
\isadelimproof
\ \ %
\endisadelimproof
%
\isatagproof
\isacommand{apply}\isamarkupfalse%
{\isacharparenleft}{\kern0pt}induction\ {\isasymphi}\ rule{\isacharcolon}{\kern0pt}\ is{\isacharunderscore}{\kern0pt}STRIPS{\isacharunderscore}{\kern0pt}fmla{\isachardot}{\kern0pt}induct{\isacharparenright}{\kern0pt}\isanewline
\ \ \isacommand{by}\isamarkupfalse%
\ {\isacharparenleft}{\kern0pt}auto\ simp{\isacharcolon}{\kern0pt}\ valuation{\isacharunderscore}{\kern0pt}def{\isacharparenright}{\kern0pt}%
\endisatagproof
{\isafoldproof}%
%
\isadelimproof
\isanewline
%
\endisadelimproof
\isanewline
\isacommand{lemma}\isamarkupfalse%
\ aux{\isadigit{2}}{\isacharcolon}{\kern0pt}\ {\isachardoublequoteopen}{\isasymlbrakk}wm{\isacharunderscore}{\kern0pt}basic\ M{\isacharsemicolon}{\kern0pt}\ is{\isacharunderscore}{\kern0pt}STRIPS{\isacharunderscore}{\kern0pt}fmla\ {\isasymphi}{\isacharsemicolon}{\kern0pt}\ {\isasymforall}{\isasymA}{\isachardot}{\kern0pt}\ {\isacharparenleft}{\kern0pt}{\isasymforall}G{\isasymin}M{\isachardot}{\kern0pt}\ {\isasymA}\ {\isasymTurnstile}\ G{\isacharparenright}{\kern0pt}\ {\isasymlongrightarrow}\ {\isasymA}\ {\isasymTurnstile}\ {\isasymphi}{\isasymrbrakk}\ {\isasymLongrightarrow}\ valuation\ M\ {\isasymTurnstile}\ {\isasymphi}{\isachardoublequoteclose}\isanewline
%
\isadelimproof
\ \ %
\endisadelimproof
%
\isatagproof
\isacommand{apply}\isamarkupfalse%
{\isacharparenleft}{\kern0pt}induction\ {\isasymphi}\ rule{\isacharcolon}{\kern0pt}\ is{\isacharunderscore}{\kern0pt}STRIPS{\isacharunderscore}{\kern0pt}fmla{\isachardot}{\kern0pt}induct{\isacharparenright}{\kern0pt}\isanewline
\ \ \isacommand{apply}\isamarkupfalse%
\ simp{\isacharunderscore}{\kern0pt}all\isanewline
\ \ \isacommand{apply}\isamarkupfalse%
\ {\isacharparenleft}{\kern0pt}metis\ in{\isacharunderscore}{\kern0pt}close{\isacharunderscore}{\kern0pt}world{\isacharunderscore}{\kern0pt}conv\ valuation{\isacharunderscore}{\kern0pt}aux{\isacharunderscore}{\kern0pt}{\isadigit{2}}{\isacharparenright}{\kern0pt}\isanewline
\ \ \isacommand{using}\isamarkupfalse%
\ in{\isacharunderscore}{\kern0pt}close{\isacharunderscore}{\kern0pt}world{\isacharunderscore}{\kern0pt}conv\ valuation{\isacharunderscore}{\kern0pt}aux{\isacharunderscore}{\kern0pt}{\isadigit{2}}\ \isacommand{apply}\isamarkupfalse%
\ blast\isanewline
\ \ \isacommand{using}\isamarkupfalse%
\ in{\isacharunderscore}{\kern0pt}close{\isacharunderscore}{\kern0pt}world{\isacharunderscore}{\kern0pt}conv\ valuation{\isacharunderscore}{\kern0pt}aux{\isacharunderscore}{\kern0pt}{\isadigit{2}}\ \isacommand{by}\isamarkupfalse%
\ auto%
\endisatagproof
{\isafoldproof}%
%
\isadelimproof
\isanewline
%
\endisadelimproof
\isanewline
\isanewline
\isacommand{lemma}\isamarkupfalse%
\ valuation{\isacharunderscore}{\kern0pt}iff{\isacharunderscore}{\kern0pt}STRIPS{\isacharcolon}{\kern0pt}\isanewline
\ \ \isakeyword{assumes}\ {\isachardoublequoteopen}wm{\isacharunderscore}{\kern0pt}basic\ M{\isachardoublequoteclose}\isanewline
\ \ \isakeyword{assumes}\ {\isachardoublequoteopen}is{\isacharunderscore}{\kern0pt}STRIPS{\isacharunderscore}{\kern0pt}fmla\ {\isasymphi}{\isachardoublequoteclose}\isanewline
\ \ \isakeyword{shows}\ {\isachardoublequoteopen}valuation\ M\ {\isasymTurnstile}\ {\isasymphi}\ {\isasymlongleftrightarrow}\ M\ {\isasymTTurnstile}\ {\isasymphi}{\isachardoublequoteclose}\isanewline
%
\isadelimproof
%
\endisadelimproof
%
\isatagproof
\isacommand{proof}\isamarkupfalse%
\ {\isacharminus}{\kern0pt}\isanewline
\ \ \isacommand{have}\isamarkupfalse%
\ aux{\isadigit{1}}{\isacharcolon}{\kern0pt}\ {\isachardoublequoteopen}{\isasymAnd}{\isasymA}{\isachardot}{\kern0pt}\ {\isasymlbrakk}valuation\ M\ {\isasymTurnstile}\ {\isasymphi}{\isacharsemicolon}{\kern0pt}\ {\isasymforall}G{\isasymin}M{\isachardot}{\kern0pt}\ {\isasymA}\ {\isasymTurnstile}\ G{\isasymrbrakk}\ {\isasymLongrightarrow}\ {\isasymA}\ {\isasymTurnstile}\ {\isasymphi}{\isachardoublequoteclose}\isanewline
\ \ \ \ \isacommand{using}\isamarkupfalse%
\ assms\ \isacommand{apply}\isamarkupfalse%
{\isacharparenleft}{\kern0pt}induction\ {\isasymphi}\ rule{\isacharcolon}{\kern0pt}\ is{\isacharunderscore}{\kern0pt}STRIPS{\isacharunderscore}{\kern0pt}fmla{\isachardot}{\kern0pt}induct{\isacharparenright}{\kern0pt}\isanewline
\ \ \ \ \isacommand{by}\isamarkupfalse%
\ {\isacharparenleft}{\kern0pt}auto\ simp{\isacharcolon}{\kern0pt}\ valuation{\isacharunderscore}{\kern0pt}def{\isacharparenright}{\kern0pt}\isanewline
\ \ \isacommand{have}\isamarkupfalse%
\ aux{\isadigit{2}}{\isacharcolon}{\kern0pt}\ {\isachardoublequoteopen}{\isasymlbrakk}{\isasymforall}{\isasymA}{\isachardot}{\kern0pt}\ {\isacharparenleft}{\kern0pt}{\isasymforall}G{\isasymin}M{\isachardot}{\kern0pt}\ {\isasymA}\ {\isasymTurnstile}\ G{\isacharparenright}{\kern0pt}\ {\isasymlongrightarrow}\ {\isasymA}\ {\isasymTurnstile}\ {\isasymphi}{\isasymrbrakk}\ {\isasymLongrightarrow}\ valuation\ M\ {\isasymTurnstile}\ {\isasymphi}{\isachardoublequoteclose}\isanewline
\ \ \ \ \isacommand{using}\isamarkupfalse%
\ assms\isanewline
\ \ \ \ \isacommand{apply}\isamarkupfalse%
{\isacharparenleft}{\kern0pt}induction\ {\isasymphi}\ rule{\isacharcolon}{\kern0pt}\ is{\isacharunderscore}{\kern0pt}STRIPS{\isacharunderscore}{\kern0pt}fmla{\isachardot}{\kern0pt}induct{\isacharparenright}{\kern0pt}\isanewline
\ \ \ \ \isacommand{apply}\isamarkupfalse%
\ simp{\isacharunderscore}{\kern0pt}all\isanewline
\ \ \ \ \isacommand{apply}\isamarkupfalse%
\ {\isacharparenleft}{\kern0pt}metis\ in{\isacharunderscore}{\kern0pt}close{\isacharunderscore}{\kern0pt}world{\isacharunderscore}{\kern0pt}conv\ valuation{\isacharunderscore}{\kern0pt}aux{\isacharunderscore}{\kern0pt}{\isadigit{2}}{\isacharparenright}{\kern0pt}\isanewline
\ \ \ \ \isacommand{using}\isamarkupfalse%
\ in{\isacharunderscore}{\kern0pt}close{\isacharunderscore}{\kern0pt}world{\isacharunderscore}{\kern0pt}conv\ valuation{\isacharunderscore}{\kern0pt}aux{\isacharunderscore}{\kern0pt}{\isadigit{2}}\ \isacommand{apply}\isamarkupfalse%
\ blast\isanewline
\ \ \ \ \isacommand{using}\isamarkupfalse%
\ in{\isacharunderscore}{\kern0pt}close{\isacharunderscore}{\kern0pt}world{\isacharunderscore}{\kern0pt}conv\ valuation{\isacharunderscore}{\kern0pt}aux{\isacharunderscore}{\kern0pt}{\isadigit{2}}\ \isacommand{by}\isamarkupfalse%
\ auto\isanewline
\ \ \isacommand{show}\isamarkupfalse%
\ {\isacharquery}{\kern0pt}thesis\isanewline
\ \ \ \ \isacommand{by}\isamarkupfalse%
\ {\isacharparenleft}{\kern0pt}auto\ simp{\isacharcolon}{\kern0pt}\ entailment{\isacharunderscore}{\kern0pt}def\ intro{\isacharcolon}{\kern0pt}\ aux{\isadigit{1}}\ aux{\isadigit{2}}{\isacharparenright}{\kern0pt}\isanewline
\isacommand{qed}\isamarkupfalse%
%
\endisatagproof
{\isafoldproof}%
%
\isadelimproof
%
\endisadelimproof
%
\begin{isamarkuptext}%
Our extension to negation and equality is a proper generalization of the
  standard STRIPS semantics for formula without negation and equality%
\end{isamarkuptext}\isamarkuptrue%
\isacommand{theorem}\isamarkupfalse%
\ proper{\isacharunderscore}{\kern0pt}STRIPS{\isacharunderscore}{\kern0pt}generalization{\isacharcolon}{\kern0pt}\isanewline
\ \ {\isachardoublequoteopen}{\isasymlbrakk}wm{\isacharunderscore}{\kern0pt}basic\ M{\isacharsemicolon}{\kern0pt}\ is{\isacharunderscore}{\kern0pt}STRIPS{\isacharunderscore}{\kern0pt}fmla\ {\isasymphi}{\isasymrbrakk}\ {\isasymLongrightarrow}\ M\ \isactrlsup c{\isasymTTurnstile}\isactrlsub {\isacharequal}{\kern0pt}\ {\isasymphi}\ {\isasymlongleftrightarrow}\ M\ {\isasymTTurnstile}\ {\isasymphi}{\isachardoublequoteclose}\isanewline
%
\isadelimproof
\ \ %
\endisadelimproof
%
\isatagproof
\isacommand{by}\isamarkupfalse%
\ {\isacharparenleft}{\kern0pt}simp\ add{\isacharcolon}{\kern0pt}\ valuation{\isacharunderscore}{\kern0pt}iff{\isacharunderscore}{\kern0pt}close{\isacharunderscore}{\kern0pt}world{\isacharbrackleft}{\kern0pt}symmetric{\isacharbrackright}{\kern0pt}\ valuation{\isacharunderscore}{\kern0pt}iff{\isacharunderscore}{\kern0pt}STRIPS{\isacharparenright}{\kern0pt}%
\endisatagproof
{\isafoldproof}%
%
\isadelimproof
%
\endisadelimproof
%
\isadelimdocument
%
\endisadelimdocument
%
\isatagdocument
%
\isamarkupsubsection{Happening Execution Semantics%
}
\isamarkuptrue%
%
\endisatagdocument
{\isafolddocument}%
%
\isadelimdocument
%
\endisadelimdocument
%
\begin{isamarkuptext}%
For this section, we fix a domain \isa{D}, using Isabelle's
  locale mechanism.%
\end{isamarkuptext}\isamarkuptrue%
\isacommand{locale}\isamarkupfalse%
\ ast{\isacharunderscore}{\kern0pt}domain\ {\isacharequal}{\kern0pt}\isanewline
\ \ \isakeyword{fixes}\ D\ {\isacharcolon}{\kern0pt}{\isacharcolon}{\kern0pt}\ ast{\isacharunderscore}{\kern0pt}domain\isanewline
\isakeyword{begin}%
\begin{isamarkuptext}%
Check if two ground actions are interfering. This definition is taken from [FL03]%
\end{isamarkuptext}\isamarkuptrue%
\ \ \isacommand{definition}\isamarkupfalse%
\ acts{\isacharunderscore}{\kern0pt}non{\isacharunderscore}{\kern0pt}intrf\ {\isacharcolon}{\kern0pt}{\isacharcolon}{\kern0pt}\ {\isachardoublequoteopen}ground{\isacharunderscore}{\kern0pt}action\ {\isasymRightarrow}\ ground{\isacharunderscore}{\kern0pt}action\ {\isasymRightarrow}\ bool{\isachardoublequoteclose}\ \isakeyword{where}\isanewline
\ \ \ \ {\isachardoublequoteopen}acts{\isacharunderscore}{\kern0pt}non{\isacharunderscore}{\kern0pt}intrf\ a\ b\ {\isasymlongleftrightarrow}\ {\isacharparenleft}{\kern0pt}\isanewline
\ \ \ \ \ \ let\ \isanewline
\ \ \ \ \ \ \ \ add\isactrlsub a\ {\isacharequal}{\kern0pt}\ set{\isacharparenleft}{\kern0pt}adds{\isacharparenleft}{\kern0pt}effect\ a{\isacharparenright}{\kern0pt}{\isacharparenright}{\kern0pt}{\isacharsemicolon}{\kern0pt}\ del\isactrlsub a\ {\isacharequal}{\kern0pt}\ set{\isacharparenleft}{\kern0pt}dels{\isacharparenleft}{\kern0pt}effect\ a{\isacharparenright}{\kern0pt}{\isacharparenright}{\kern0pt}{\isacharsemicolon}{\kern0pt}\ pre\isactrlsub a\ {\isacharequal}{\kern0pt}\ Atom\ {\isacharbackquote}{\kern0pt}\ atoms\ {\isacharparenleft}{\kern0pt}precondition\ a{\isacharparenright}{\kern0pt}{\isacharsemicolon}{\kern0pt}\isanewline
\ \ \ \ \ \ \ \ add\isactrlsub b\ {\isacharequal}{\kern0pt}\ set{\isacharparenleft}{\kern0pt}adds{\isacharparenleft}{\kern0pt}effect\ b{\isacharparenright}{\kern0pt}{\isacharparenright}{\kern0pt}{\isacharsemicolon}{\kern0pt}\ del\isactrlsub b\ {\isacharequal}{\kern0pt}\ set{\isacharparenleft}{\kern0pt}dels{\isacharparenleft}{\kern0pt}effect\ b{\isacharparenright}{\kern0pt}{\isacharparenright}{\kern0pt}{\isacharsemicolon}{\kern0pt}\ pre\isactrlsub b\ {\isacharequal}{\kern0pt}\ Atom\ {\isacharbackquote}{\kern0pt}\ atoms\ {\isacharparenleft}{\kern0pt}precondition\ b{\isacharparenright}{\kern0pt}\ \isanewline
\ \ \ \ \ \ in\isanewline
\ \ \ \ \ \ pre\isactrlsub a\ {\isasyminter}\ {\isacharparenleft}{\kern0pt}add\isactrlsub b\ {\isasymunion}\ del\isactrlsub b{\isacharparenright}{\kern0pt}\ {\isacharequal}{\kern0pt}\ {\isacharbraceleft}{\kern0pt}{\isacharbraceright}{\kern0pt}\ {\isasymand}\isanewline
\ \ \ \ \ \ pre\isactrlsub b\ {\isasyminter}\ {\isacharparenleft}{\kern0pt}add\isactrlsub a\ {\isasymunion}\ del\isactrlsub a{\isacharparenright}{\kern0pt}\ {\isacharequal}{\kern0pt}\ {\isacharbraceleft}{\kern0pt}{\isacharbraceright}{\kern0pt}\ {\isasymand}\isanewline
\ \ \ \ \ \ add\isactrlsub a\ {\isasyminter}\ del\isactrlsub b\ {\isacharequal}{\kern0pt}\ {\isacharbraceleft}{\kern0pt}{\isacharbraceright}{\kern0pt}\ {\isasymand}\isanewline
\ \ \ \ \ \ add\isactrlsub b\ {\isasyminter}\ del\isactrlsub a\ {\isacharequal}{\kern0pt}\ {\isacharbraceleft}{\kern0pt}{\isacharbraceright}{\kern0pt}{\isacharparenright}{\kern0pt}{\isachardoublequoteclose}\isanewline
\isanewline
\ \ \isacommand{lemma}\isamarkupfalse%
\ acts{\isacharunderscore}{\kern0pt}non{\isacharunderscore}{\kern0pt}intrf{\isacharunderscore}{\kern0pt}symmetric{\isacharcolon}{\kern0pt}\ \isanewline
\ \ \ \ {\isachardoublequoteopen}{\isasymforall}a\ b{\isachardot}{\kern0pt}\ acts{\isacharunderscore}{\kern0pt}non{\isacharunderscore}{\kern0pt}intrf\ a\ b\ {\isasymlongleftrightarrow}\ acts{\isacharunderscore}{\kern0pt}non{\isacharunderscore}{\kern0pt}intrf\ b\ a{\isachardoublequoteclose}\isanewline
%
\isadelimproof
\ \ \ \ %
\endisadelimproof
%
\isatagproof
\isacommand{unfolding}\isamarkupfalse%
\ acts{\isacharunderscore}{\kern0pt}non{\isacharunderscore}{\kern0pt}intrf{\isacharunderscore}{\kern0pt}def\ \isacommand{by}\isamarkupfalse%
\ {\isacharparenleft}{\kern0pt}auto\ simp{\isacharcolon}{\kern0pt}\ Let{\isacharunderscore}{\kern0pt}def{\isacharparenright}{\kern0pt}%
\endisatagproof
{\isafoldproof}%
%
\isadelimproof
%
\endisadelimproof
%
\begin{isamarkuptext}%
It seems to be agreed upon that, in case of a contradictory effect,
    addition overrides deletion. We model this behaviour by first executing
    the deletions, and then the additions.%
\end{isamarkuptext}\isamarkuptrue%
\ \ \isacommand{fun}\isamarkupfalse%
\ apply{\isacharunderscore}{\kern0pt}happ\ {\isacharcolon}{\kern0pt}{\isacharcolon}{\kern0pt}\ {\isachardoublequoteopen}happening\ {\isasymRightarrow}\ world{\isacharunderscore}{\kern0pt}model\ {\isasymRightarrow}\ world{\isacharunderscore}{\kern0pt}model{\isachardoublequoteclose}\ \isakeyword{where}\isanewline
\ \ \ \ {\isachardoublequoteopen}apply{\isacharunderscore}{\kern0pt}happ\ {\isacharparenleft}{\kern0pt}t\isactrlsub i{\isacharcomma}{\kern0pt}\ A\isactrlsub i{\isacharparenright}{\kern0pt}\ M\ {\isacharequal}{\kern0pt}\ \isanewline
\ \ \ \ \ \ {\isacharparenleft}{\kern0pt}M\ {\isacharminus}{\kern0pt}\ {\isasymUnion}\ {\isacharparenleft}{\kern0pt}set\ {\isacharparenleft}{\kern0pt}map\ {\isacharparenleft}{\kern0pt}set\ o\ dels\ o\ effect{\isacharparenright}{\kern0pt}\ A\isactrlsub i{\isacharparenright}{\kern0pt}{\isacharparenright}{\kern0pt}{\isacharparenright}{\kern0pt}\ {\isasymunion}\ {\isasymUnion}\ {\isacharparenleft}{\kern0pt}set\ {\isacharparenleft}{\kern0pt}map\ {\isacharparenleft}{\kern0pt}set\ o\ adds\ o\ effect{\isacharparenright}{\kern0pt}\ A\isactrlsub i{\isacharparenright}{\kern0pt}{\isacharparenright}{\kern0pt}{\isachardoublequoteclose}%
\begin{isamarkuptext}%
Predicate to model that the given list happenings is
    executable, and transforms an initial world model \isa{M} into a final
    model \isa{M{\isacharprime}{\kern0pt}}.

    Note that this definition over the list structure is more convenient in HOL
    than to explicitly define an indexed sequence \isa{M\isactrlsub {\isadigit{0}}{\isasymdots}M\isactrlsub N} of intermediate world
    models, as done in [Lif87].%
\end{isamarkuptext}\isamarkuptrue%
\ \ \isacommand{fun}\isamarkupfalse%
\ valid{\isacharunderscore}{\kern0pt}happ{\isacharunderscore}{\kern0pt}seq\ {\isacharcolon}{\kern0pt}{\isacharcolon}{\kern0pt}\ {\isachardoublequoteopen}world{\isacharunderscore}{\kern0pt}model\ {\isasymRightarrow}\ happening\ list\ {\isasymRightarrow}\ world{\isacharunderscore}{\kern0pt}model\ {\isasymRightarrow}\ bool{\isachardoublequoteclose}\ \isakeyword{where}\isanewline
\ \ \ \ {\isachardoublequoteopen}valid{\isacharunderscore}{\kern0pt}happ{\isacharunderscore}{\kern0pt}seq\ M\ {\isacharbrackleft}{\kern0pt}{\isacharbrackright}{\kern0pt}\ M{\isacharprime}{\kern0pt}\ {\isasymlongleftrightarrow}\ {\isacharparenleft}{\kern0pt}M\ {\isacharequal}{\kern0pt}\ M{\isacharprime}{\kern0pt}{\isacharparenright}{\kern0pt}{\isachardoublequoteclose}\isanewline
\ \ {\isacharbar}{\kern0pt}\ {\isachardoublequoteopen}valid{\isacharunderscore}{\kern0pt}happ{\isacharunderscore}{\kern0pt}seq\ M\ {\isacharparenleft}{\kern0pt}{\isacharparenleft}{\kern0pt}t\isactrlsub i{\isacharcomma}{\kern0pt}A\isactrlsub i{\isacharparenright}{\kern0pt}{\isacharhash}{\kern0pt}hs{\isacharparenright}{\kern0pt}\ M{\isacharprime}{\kern0pt}\ {\isasymlongleftrightarrow}\ \isanewline
\ \ \ \ \ \ {\isacharparenleft}{\kern0pt}{\isasymforall}a{\isasymin}\ set\ A\isactrlsub i{\isachardot}{\kern0pt}\ M\ \isactrlsup c{\isasymTTurnstile}\isactrlsub {\isacharequal}{\kern0pt}\ precondition\ a{\isacharparenright}{\kern0pt}\isanewline
\ \ \ \ \ \ {\isasymand}\ {\isacharparenleft}{\kern0pt}pairwise\ acts{\isacharunderscore}{\kern0pt}non{\isacharunderscore}{\kern0pt}intrf\ {\isacharparenleft}{\kern0pt}set\ A\isactrlsub i{\isacharparenright}{\kern0pt}{\isacharparenright}{\kern0pt}\isanewline
\ \ \ \ \ \ {\isasymand}\ valid{\isacharunderscore}{\kern0pt}happ{\isacharunderscore}{\kern0pt}seq\ {\isacharparenleft}{\kern0pt}apply{\isacharunderscore}{\kern0pt}happ\ {\isacharparenleft}{\kern0pt}t\isactrlsub i{\isacharcomma}{\kern0pt}A\isactrlsub i{\isacharparenright}{\kern0pt}\ M{\isacharparenright}{\kern0pt}\ hs\ M{\isacharprime}{\kern0pt}{\isachardoublequoteclose}\isanewline
\isanewline
\ \ \isacommand{lemma}\isamarkupfalse%
\ valid{\isacharunderscore}{\kern0pt}happ{\isacharunderscore}{\kern0pt}seq{\isacharunderscore}{\kern0pt}M\isactrlsub j{\isacharunderscore}{\kern0pt}is{\isacharunderscore}{\kern0pt}M\isactrlsub i{\isacharcolon}{\kern0pt}\isanewline
\ \ \ \ \isakeyword{assumes}\ {\isachardoublequoteopen}{\isasymforall}h\ {\isasymin}\ set\ hs{\isachardot}{\kern0pt}\ apply{\isacharunderscore}{\kern0pt}happ\ h\ M\isactrlsub i\ {\isacharequal}{\kern0pt}\ M\isactrlsub i{\isachardoublequoteclose}\ \isanewline
\ \ \ \ \ \ \ \ \isakeyword{and}\ {\isachardoublequoteopen}valid{\isacharunderscore}{\kern0pt}happ{\isacharunderscore}{\kern0pt}seq\ M\isactrlsub i\ hs\ M\isactrlsub j{\isachardoublequoteclose}\isanewline
\ \ \ \ \isakeyword{shows}\ {\isachardoublequoteopen}M\isactrlsub i\ {\isacharequal}{\kern0pt}\ M\isactrlsub j{\isachardoublequoteclose}\isanewline
%
\isadelimproof
\ \ %
\endisadelimproof
%
\isatagproof
\isacommand{using}\isamarkupfalse%
\ assms\ \isacommand{by}\isamarkupfalse%
\ {\isacharparenleft}{\kern0pt}induction\ M\isactrlsub i\ hs\ M\isactrlsub j\ rule{\isacharcolon}{\kern0pt}\ valid{\isacharunderscore}{\kern0pt}happ{\isacharunderscore}{\kern0pt}seq{\isachardot}{\kern0pt}induct{\isacharparenright}{\kern0pt}\ auto%
\endisatagproof
{\isafoldproof}%
%
\isadelimproof
\isanewline
%
\endisadelimproof
\isanewline
\isacommand{end}\isamarkupfalse%
\ %
\isamarkupcmt{Context of \isa{ast{\isacharunderscore}{\kern0pt}domain}%
}%
\isadelimdocument
%
\endisadelimdocument
%
\isatagdocument
%
\isamarkupsubsection{Well-Formedness of PDDL%
}
\isamarkuptrue%
%
\endisatagdocument
{\isafolddocument}%
%
\isadelimdocument
%
\endisadelimdocument
\isacommand{fun}\isamarkupfalse%
\ ty{\isacharunderscore}{\kern0pt}term\ \isakeyword{where}\isanewline
\ \ {\isachardoublequoteopen}ty{\isacharunderscore}{\kern0pt}term\ varT\ objT\ {\isacharparenleft}{\kern0pt}term{\isachardot}{\kern0pt}VAR\ v{\isacharparenright}{\kern0pt}\ {\isacharequal}{\kern0pt}\ varT\ v{\isachardoublequoteclose}\isanewline
{\isacharbar}{\kern0pt}\ {\isachardoublequoteopen}ty{\isacharunderscore}{\kern0pt}term\ varT\ objT\ {\isacharparenleft}{\kern0pt}term{\isachardot}{\kern0pt}CONST\ c{\isacharparenright}{\kern0pt}\ {\isacharequal}{\kern0pt}\ objT\ c{\isachardoublequoteclose}\isanewline
\isanewline
\isanewline
\isacommand{lemma}\isamarkupfalse%
\ ty{\isacharunderscore}{\kern0pt}term{\isacharunderscore}{\kern0pt}mono{\isacharcolon}{\kern0pt}\ {\isachardoublequoteopen}varT\ {\isasymsubseteq}\isactrlsub m\ varT{\isacharprime}{\kern0pt}\ {\isasymLongrightarrow}\ objT\ {\isasymsubseteq}\isactrlsub m\ objT{\isacharprime}{\kern0pt}\ {\isasymLongrightarrow}\isanewline
\ \ ty{\isacharunderscore}{\kern0pt}term\ varT\ objT\ {\isasymsubseteq}\isactrlsub m\ ty{\isacharunderscore}{\kern0pt}term\ varT{\isacharprime}{\kern0pt}\ objT{\isacharprime}{\kern0pt}{\isachardoublequoteclose}\isanewline
%
\isadelimproof
\ \ %
\endisadelimproof
%
\isatagproof
\isacommand{apply}\isamarkupfalse%
\ {\isacharparenleft}{\kern0pt}rule\ map{\isacharunderscore}{\kern0pt}leI{\isacharparenright}{\kern0pt}\isanewline
\ \ \isacommand{subgoal}\isamarkupfalse%
\ \isakeyword{for}\ x\ v\isanewline
\ \ \ \ \isacommand{apply}\isamarkupfalse%
\ {\isacharparenleft}{\kern0pt}cases\ x{\isacharparenright}{\kern0pt}\isanewline
\ \ \ \ \isacommand{apply}\isamarkupfalse%
\ {\isacharparenleft}{\kern0pt}auto\ dest{\isacharcolon}{\kern0pt}\ map{\isacharunderscore}{\kern0pt}leD{\isacharparenright}{\kern0pt}\isanewline
\ \ \ \ \isacommand{done}\isamarkupfalse%
\isanewline
\ \ \isacommand{done}\isamarkupfalse%
%
\endisatagproof
{\isafoldproof}%
%
\isadelimproof
\isanewline
%
\endisadelimproof
\isanewline
\isanewline
\isacommand{context}\isamarkupfalse%
\ ast{\isacharunderscore}{\kern0pt}domain\ \isakeyword{begin}%
\begin{isamarkuptext}%
The signature is a partial function that maps the predicates
    of the domain to lists of argument types.%
\end{isamarkuptext}\isamarkuptrue%
\ \ \isacommand{definition}\isamarkupfalse%
\ sig\ {\isacharcolon}{\kern0pt}{\isacharcolon}{\kern0pt}\ {\isachardoublequoteopen}predicate\ {\isasymrightharpoonup}\ type\ list{\isachardoublequoteclose}\ \isakeyword{where}\isanewline
\ \ \ \ {\isachardoublequoteopen}sig\ {\isasymequiv}\ map{\isacharunderscore}{\kern0pt}of\ {\isacharparenleft}{\kern0pt}map\ {\isacharparenleft}{\kern0pt}{\isasymlambda}PredDecl\ p\ n\ {\isasymRightarrow}\ {\isacharparenleft}{\kern0pt}p{\isacharcomma}{\kern0pt}n{\isacharparenright}{\kern0pt}{\isacharparenright}{\kern0pt}\ {\isacharparenleft}{\kern0pt}predicates\ D{\isacharparenright}{\kern0pt}{\isacharparenright}{\kern0pt}{\isachardoublequoteclose}\isanewline
\isanewline
\ \ \isacommand{definition}\isamarkupfalse%
\ func{\isacharunderscore}{\kern0pt}sig\ {\isacharcolon}{\kern0pt}{\isacharcolon}{\kern0pt}\ {\isachardoublequoteopen}func\ {\isasymrightharpoonup}\ type\ list{\isachardoublequoteclose}\ \isakeyword{where}\isanewline
\ \ \ \ {\isachardoublequoteopen}func{\isacharunderscore}{\kern0pt}sig\ {\isasymequiv}\ map{\isacharunderscore}{\kern0pt}of\ {\isacharparenleft}{\kern0pt}map\ {\isacharparenleft}{\kern0pt}{\isasymlambda}FuncDecl\ f\ n\ {\isasymRightarrow}\ {\isacharparenleft}{\kern0pt}f{\isacharcomma}{\kern0pt}n{\isacharparenright}{\kern0pt}{\isacharparenright}{\kern0pt}\ {\isacharparenleft}{\kern0pt}functions\ D{\isacharparenright}{\kern0pt}{\isacharparenright}{\kern0pt}{\isachardoublequoteclose}%
\begin{isamarkuptext}%
We use a flat subtype hierarchy, where every type is a subtype
    of object, and there are no other subtype relations.

    Note that we do not need to restrict this relation to declared types,
    as we will explicitly ensure that all types used in the problem are
    declared.%
\end{isamarkuptext}\isamarkuptrue%
\ \ \isacommand{fun}\isamarkupfalse%
\ subtype{\isacharunderscore}{\kern0pt}edge\ \isakeyword{where}\isanewline
\ \ \ \ {\isachardoublequoteopen}subtype{\isacharunderscore}{\kern0pt}edge\ {\isacharparenleft}{\kern0pt}ty{\isacharcomma}{\kern0pt}superty{\isacharparenright}{\kern0pt}\ {\isacharequal}{\kern0pt}\ {\isacharparenleft}{\kern0pt}superty{\isacharcomma}{\kern0pt}ty{\isacharparenright}{\kern0pt}{\isachardoublequoteclose}\isanewline
\isanewline
\ \ \isacommand{definition}\isamarkupfalse%
\ {\isachardoublequoteopen}subtype{\isacharunderscore}{\kern0pt}rel\ {\isasymequiv}\ set\ {\isacharparenleft}{\kern0pt}map\ subtype{\isacharunderscore}{\kern0pt}edge\ {\isacharparenleft}{\kern0pt}types\ D{\isacharparenright}{\kern0pt}{\isacharparenright}{\kern0pt}{\isachardoublequoteclose}\isanewline
\isanewline
\ \ \isacommand{definition}\isamarkupfalse%
\ of{\isacharunderscore}{\kern0pt}type\ {\isacharcolon}{\kern0pt}{\isacharcolon}{\kern0pt}\ {\isachardoublequoteopen}type\ {\isasymRightarrow}\ type\ {\isasymRightarrow}\ bool{\isachardoublequoteclose}\ \isakeyword{where}\isanewline
\ \ \ \ {\isachardoublequoteopen}of{\isacharunderscore}{\kern0pt}type\ oT\ T\ {\isasymequiv}\ set\ {\isacharparenleft}{\kern0pt}primitives\ oT{\isacharparenright}{\kern0pt}\ {\isasymsubseteq}\ subtype{\isacharunderscore}{\kern0pt}rel\isactrlsup {\isacharasterisk}{\kern0pt}\ {\isacharbackquote}{\kern0pt}{\isacharbackquote}{\kern0pt}\ set\ {\isacharparenleft}{\kern0pt}primitives\ T{\isacharparenright}{\kern0pt}{\isachardoublequoteclose}%
\begin{isamarkuptext}%
This checks that every primitive on the LHS is contained in or a
    subtype of a primitive on the RHS%
\end{isamarkuptext}\isamarkuptrue%
%
\begin{isamarkuptext}%
For the next few definitions, we fix a partial function that maps
    a polymorphic entity type \isa{{\isacharprime}{\kern0pt}e} to types. An entity can be
    instantiated by variables or objects later.%
\end{isamarkuptext}\isamarkuptrue%
\ \ \isacommand{context}\isamarkupfalse%
\isanewline
\ \ \ \ \isakeyword{fixes}\ ty{\isacharunderscore}{\kern0pt}ent\ {\isacharcolon}{\kern0pt}{\isacharcolon}{\kern0pt}\ {\isachardoublequoteopen}{\isacharprime}{\kern0pt}ent\ {\isasymrightharpoonup}\ type{\isachardoublequoteclose}\ \ %
\isamarkupcmt{Entity's type, None if invalid%
}\isanewline
\ \ \isakeyword{begin}%
\begin{isamarkuptext}%
Checks whether an entity has a given type%
\end{isamarkuptext}\isamarkuptrue%
\ \ \ \ \isacommand{definition}\isamarkupfalse%
\ is{\isacharunderscore}{\kern0pt}of{\isacharunderscore}{\kern0pt}type\ {\isacharcolon}{\kern0pt}{\isacharcolon}{\kern0pt}\ {\isachardoublequoteopen}{\isacharprime}{\kern0pt}ent\ {\isasymRightarrow}\ type\ {\isasymRightarrow}\ bool{\isachardoublequoteclose}\ \isakeyword{where}\isanewline
\ \ \ \ \ \ {\isachardoublequoteopen}is{\isacharunderscore}{\kern0pt}of{\isacharunderscore}{\kern0pt}type\ v\ T\ {\isasymlongleftrightarrow}\ {\isacharparenleft}{\kern0pt}\isanewline
\ \ \ \ \ \ \ \ case\ ty{\isacharunderscore}{\kern0pt}ent\ v\ of\isanewline
\ \ \ \ \ \ \ \ \ \ Some\ vT\ {\isasymRightarrow}\ of{\isacharunderscore}{\kern0pt}type\ vT\ T\isanewline
\ \ \ \ \ \ \ \ {\isacharbar}{\kern0pt}\ None\ {\isasymRightarrow}\ False{\isacharparenright}{\kern0pt}{\isachardoublequoteclose}\isanewline
\isanewline
\ \ \ \ \isacommand{fun}\isamarkupfalse%
\ wf{\isacharunderscore}{\kern0pt}pred{\isacharunderscore}{\kern0pt}atom\ {\isacharcolon}{\kern0pt}{\isacharcolon}{\kern0pt}\ {\isachardoublequoteopen}predicate\ {\isasymtimes}\ {\isacharprime}{\kern0pt}ent\ list\ {\isasymRightarrow}\ bool{\isachardoublequoteclose}\ \isakeyword{where}\isanewline
\ \ \ \ \ \ {\isachardoublequoteopen}wf{\isacharunderscore}{\kern0pt}pred{\isacharunderscore}{\kern0pt}atom\ {\isacharparenleft}{\kern0pt}p{\isacharcomma}{\kern0pt}vs{\isacharparenright}{\kern0pt}\ {\isasymlongleftrightarrow}\ {\isacharparenleft}{\kern0pt}\isanewline
\ \ \ \ \ \ \ \ case\ sig\ p\ of\isanewline
\ \ \ \ \ \ \ \ \ \ None\ {\isasymRightarrow}\ False\isanewline
\ \ \ \ \ \ \ \ {\isacharbar}{\kern0pt}\ Some\ Ts\ {\isasymRightarrow}\ list{\isacharunderscore}{\kern0pt}all{\isadigit{2}}\ is{\isacharunderscore}{\kern0pt}of{\isacharunderscore}{\kern0pt}type\ vs\ Ts{\isacharparenright}{\kern0pt}{\isachardoublequoteclose}\isanewline
\isanewline
\ \ \ \ \isacommand{fun}\isamarkupfalse%
\ wf{\isacharunderscore}{\kern0pt}func{\isacharunderscore}{\kern0pt}args\ {\isacharcolon}{\kern0pt}{\isacharcolon}{\kern0pt}\ {\isachardoublequoteopen}func\ {\isasymtimes}\ {\isacharprime}{\kern0pt}ent\ list\ {\isasymRightarrow}\ bool{\isachardoublequoteclose}\ \isakeyword{where}\isanewline
\ \ \ \ \ \ {\isachardoublequoteopen}wf{\isacharunderscore}{\kern0pt}func{\isacharunderscore}{\kern0pt}args\ {\isacharparenleft}{\kern0pt}f{\isacharcomma}{\kern0pt}vs{\isacharparenright}{\kern0pt}\ {\isasymlongleftrightarrow}\ {\isacharparenleft}{\kern0pt}\isanewline
\ \ \ \ \ \ \ \ case\ func{\isacharunderscore}{\kern0pt}sig\ f\ of\isanewline
\ \ \ \ \ \ \ \ \ \ None\ {\isasymRightarrow}\ False\isanewline
\ \ \ \ \ \ \ \ {\isacharbar}{\kern0pt}\ Some\ Ts\ {\isasymRightarrow}\ list{\isacharunderscore}{\kern0pt}all{\isadigit{2}}\ is{\isacharunderscore}{\kern0pt}of{\isacharunderscore}{\kern0pt}type\ vs\ Ts{\isacharparenright}{\kern0pt}{\isachardoublequoteclose}%
\begin{isamarkuptext}%
Predicate-atoms are well-formed if their arguments match the
      signature, equalities are well-formed if the arguments are valid
      objects (have a type).

      TODO: We could check that types may actually overlap%
\end{isamarkuptext}\isamarkuptrue%
\ \ \ \ \isacommand{fun}\isamarkupfalse%
\ wf{\isacharunderscore}{\kern0pt}atom\ {\isacharcolon}{\kern0pt}{\isacharcolon}{\kern0pt}\ {\isachardoublequoteopen}{\isacharprime}{\kern0pt}ent\ atom\ {\isasymRightarrow}\ bool{\isachardoublequoteclose}\ \isakeyword{where}\isanewline
\ \ \ \ \ \ {\isachardoublequoteopen}wf{\isacharunderscore}{\kern0pt}atom\ {\isacharparenleft}{\kern0pt}predAtm\ p\ vs{\isacharparenright}{\kern0pt}\ {\isasymlongleftrightarrow}\ wf{\isacharunderscore}{\kern0pt}pred{\isacharunderscore}{\kern0pt}atom\ {\isacharparenleft}{\kern0pt}p{\isacharcomma}{\kern0pt}vs{\isacharparenright}{\kern0pt}{\isachardoublequoteclose}\isanewline
\ \ \ \ {\isacharbar}{\kern0pt}\ {\isachardoublequoteopen}wf{\isacharunderscore}{\kern0pt}atom\ {\isacharparenleft}{\kern0pt}eqAtm\ a\ b{\isacharparenright}{\kern0pt}\ {\isasymlongleftrightarrow}\ ty{\isacharunderscore}{\kern0pt}ent\ a\ {\isasymnoteq}\ None\ {\isasymand}\ ty{\isacharunderscore}{\kern0pt}ent\ b\ {\isasymnoteq}\ None{\isachardoublequoteclose}%
\begin{isamarkuptext}%
A formula is well-formed if it consists of valid atoms,
      and does not contain negations, except for the encoding \isa{\isactrlbold {\isasymnot}{\isasymbottom}} of true.%
\end{isamarkuptext}\isamarkuptrue%
\ \ \ \ \isacommand{fun}\isamarkupfalse%
\ wf{\isacharunderscore}{\kern0pt}fmla\ {\isacharcolon}{\kern0pt}{\isacharcolon}{\kern0pt}\ {\isachardoublequoteopen}{\isacharparenleft}{\kern0pt}{\isacharprime}{\kern0pt}ent\ atom{\isacharparenright}{\kern0pt}\ formula\ {\isasymRightarrow}\ bool{\isachardoublequoteclose}\ \isakeyword{where}\isanewline
\ \ \ \ \ \ {\isachardoublequoteopen}wf{\isacharunderscore}{\kern0pt}fmla\ {\isacharparenleft}{\kern0pt}Atom\ a{\isacharparenright}{\kern0pt}\ {\isasymlongleftrightarrow}\ wf{\isacharunderscore}{\kern0pt}atom\ a{\isachardoublequoteclose}\isanewline
\ \ \ \ {\isacharbar}{\kern0pt}\ {\isachardoublequoteopen}wf{\isacharunderscore}{\kern0pt}fmla\ {\isacharparenleft}{\kern0pt}{\isasymbottom}{\isacharparenright}{\kern0pt}\ {\isasymlongleftrightarrow}\ True{\isachardoublequoteclose}\isanewline
\ \ \ \ {\isacharbar}{\kern0pt}\ {\isachardoublequoteopen}wf{\isacharunderscore}{\kern0pt}fmla\ {\isacharparenleft}{\kern0pt}{\isasymphi}{\isadigit{1}}\ \isactrlbold {\isasymand}\ {\isasymphi}{\isadigit{2}}{\isacharparenright}{\kern0pt}\ {\isasymlongleftrightarrow}\ {\isacharparenleft}{\kern0pt}wf{\isacharunderscore}{\kern0pt}fmla\ {\isasymphi}{\isadigit{1}}\ {\isasymand}\ wf{\isacharunderscore}{\kern0pt}fmla\ {\isasymphi}{\isadigit{2}}{\isacharparenright}{\kern0pt}{\isachardoublequoteclose}\isanewline
\ \ \ \ {\isacharbar}{\kern0pt}\ {\isachardoublequoteopen}wf{\isacharunderscore}{\kern0pt}fmla\ {\isacharparenleft}{\kern0pt}{\isasymphi}{\isadigit{1}}\ \isactrlbold {\isasymor}\ {\isasymphi}{\isadigit{2}}{\isacharparenright}{\kern0pt}\ {\isasymlongleftrightarrow}\ {\isacharparenleft}{\kern0pt}wf{\isacharunderscore}{\kern0pt}fmla\ {\isasymphi}{\isadigit{1}}\ {\isasymand}\ wf{\isacharunderscore}{\kern0pt}fmla\ {\isasymphi}{\isadigit{2}}{\isacharparenright}{\kern0pt}{\isachardoublequoteclose}\isanewline
\ \ \ \ {\isacharbar}{\kern0pt}\ {\isachardoublequoteopen}wf{\isacharunderscore}{\kern0pt}fmla\ {\isacharparenleft}{\kern0pt}\isactrlbold {\isasymnot}{\isasymphi}{\isacharparenright}{\kern0pt}\ {\isasymlongleftrightarrow}\ wf{\isacharunderscore}{\kern0pt}fmla\ {\isasymphi}{\isachardoublequoteclose}\isanewline
\ \ \ \ {\isacharbar}{\kern0pt}\ {\isachardoublequoteopen}wf{\isacharunderscore}{\kern0pt}fmla\ {\isacharparenleft}{\kern0pt}{\isasymphi}{\isadigit{1}}\ \isactrlbold {\isasymrightarrow}\ {\isasymphi}{\isadigit{2}}{\isacharparenright}{\kern0pt}\ {\isasymlongleftrightarrow}\ {\isacharparenleft}{\kern0pt}wf{\isacharunderscore}{\kern0pt}fmla\ {\isasymphi}{\isadigit{1}}\ {\isasymand}\ wf{\isacharunderscore}{\kern0pt}fmla\ {\isasymphi}{\isadigit{2}}{\isacharparenright}{\kern0pt}{\isachardoublequoteclose}\isanewline
\isanewline
\ \ \ \ \isacommand{lemma}\isamarkupfalse%
\ {\isachardoublequoteopen}wf{\isacharunderscore}{\kern0pt}fmla\ {\isasymphi}\ {\isacharequal}{\kern0pt}\ {\isacharparenleft}{\kern0pt}{\isasymforall}a{\isasymin}atoms\ {\isasymphi}{\isachardot}{\kern0pt}\ wf{\isacharunderscore}{\kern0pt}atom\ a{\isacharparenright}{\kern0pt}{\isachardoublequoteclose}\isanewline
%
\isadelimproof
\ \ \ \ \ \ %
\endisadelimproof
%
\isatagproof
\isacommand{by}\isamarkupfalse%
\ {\isacharparenleft}{\kern0pt}induction\ {\isasymphi}{\isacharparenright}{\kern0pt}\ auto%
\endisatagproof
{\isafoldproof}%
%
\isadelimproof
%
\endisadelimproof
%
\begin{isamarkuptext}%
Special case for a well-formed atomic predicate formula%
\end{isamarkuptext}\isamarkuptrue%
\ \ \ \ \isacommand{fun}\isamarkupfalse%
\ wf{\isacharunderscore}{\kern0pt}fmla{\isacharunderscore}{\kern0pt}atom\ \isakeyword{where}\isanewline
\ \ \ \ \ \ {\isachardoublequoteopen}wf{\isacharunderscore}{\kern0pt}fmla{\isacharunderscore}{\kern0pt}atom\ {\isacharparenleft}{\kern0pt}Atom\ {\isacharparenleft}{\kern0pt}predAtm\ a\ vs{\isacharparenright}{\kern0pt}{\isacharparenright}{\kern0pt}\ {\isasymlongleftrightarrow}\ wf{\isacharunderscore}{\kern0pt}pred{\isacharunderscore}{\kern0pt}atom\ {\isacharparenleft}{\kern0pt}a{\isacharcomma}{\kern0pt}vs{\isacharparenright}{\kern0pt}{\isachardoublequoteclose}\isanewline
\ \ \ \ {\isacharbar}{\kern0pt}\ {\isachardoublequoteopen}wf{\isacharunderscore}{\kern0pt}fmla{\isacharunderscore}{\kern0pt}atom\ {\isacharunderscore}{\kern0pt}\ {\isasymlongleftrightarrow}\ False{\isachardoublequoteclose}\isanewline
\isanewline
\ \ \ \ \isacommand{lemma}\isamarkupfalse%
\ wf{\isacharunderscore}{\kern0pt}fmla{\isacharunderscore}{\kern0pt}atom{\isacharunderscore}{\kern0pt}alt{\isacharcolon}{\kern0pt}\ {\isachardoublequoteopen}wf{\isacharunderscore}{\kern0pt}fmla{\isacharunderscore}{\kern0pt}atom\ {\isasymphi}\ {\isasymlongleftrightarrow}\ is{\isacharunderscore}{\kern0pt}predAtom\ {\isasymphi}\ {\isasymand}\ wf{\isacharunderscore}{\kern0pt}fmla\ {\isasymphi}{\isachardoublequoteclose}\isanewline
%
\isadelimproof
\ \ \ \ \ \ %
\endisadelimproof
%
\isatagproof
\isacommand{by}\isamarkupfalse%
\ {\isacharparenleft}{\kern0pt}cases\ {\isasymphi}\ rule{\isacharcolon}{\kern0pt}\ wf{\isacharunderscore}{\kern0pt}fmla{\isacharunderscore}{\kern0pt}atom{\isachardot}{\kern0pt}cases{\isacharparenright}{\kern0pt}\ auto%
\endisatagproof
{\isafoldproof}%
%
\isadelimproof
%
\endisadelimproof
%
\begin{isamarkuptext}%
An effect is well-formed if the added and removed formulas
      are atomic%
\end{isamarkuptext}\isamarkuptrue%
\ \ \ \ \isacommand{fun}\isamarkupfalse%
\ wf{\isacharunderscore}{\kern0pt}effect\ \isakeyword{where}\isanewline
\ \ \ \ \ \ {\isachardoublequoteopen}wf{\isacharunderscore}{\kern0pt}effect\ {\isacharparenleft}{\kern0pt}Effect\ a\ d{\isacharparenright}{\kern0pt}\ {\isasymlongleftrightarrow}\ {\isacharparenleft}{\kern0pt}{\isasymforall}ae{\isasymin}set\ a{\isachardot}{\kern0pt}\ wf{\isacharunderscore}{\kern0pt}fmla{\isacharunderscore}{\kern0pt}atom\ ae{\isacharparenright}{\kern0pt}\ {\isasymand}\ {\isacharparenleft}{\kern0pt}{\isasymforall}de{\isasymin}set\ d{\isachardot}{\kern0pt}\ wf{\isacharunderscore}{\kern0pt}fmla{\isacharunderscore}{\kern0pt}atom\ de{\isacharparenright}{\kern0pt}{\isachardoublequoteclose}\isanewline
\isanewline
\ \ \ \ \isacommand{fun}\isamarkupfalse%
\ wf{\isacharunderscore}{\kern0pt}duration{\isacharunderscore}{\kern0pt}const\ {\isacharcolon}{\kern0pt}{\isacharcolon}{\kern0pt}\ {\isachardoublequoteopen}{\isacharprime}{\kern0pt}ent\ duration{\isacharunderscore}{\kern0pt}constraint\ {\isasymRightarrow}\ bool{\isachardoublequoteclose}\ \isakeyword{where}\isanewline
\ \ \ \ \ \ {\isachardoublequoteopen}wf{\isacharunderscore}{\kern0pt}duration{\isacharunderscore}{\kern0pt}const\ No{\isacharunderscore}{\kern0pt}Const\ {\isasymlongleftrightarrow}\ True{\isachardoublequoteclose}\isanewline
\ \ \ \ {\isacharbar}{\kern0pt}\ {\isachardoublequoteopen}wf{\isacharunderscore}{\kern0pt}duration{\isacharunderscore}{\kern0pt}const\ {\isacharparenleft}{\kern0pt}Time{\isacharunderscore}{\kern0pt}Const\ op\ d{\isacharparenright}{\kern0pt}\ {\isasymlongleftrightarrow}\ d\ {\isasymge}\ {\isadigit{0}}{\isachardoublequoteclose}\isanewline
\ \ \ \ {\isacharbar}{\kern0pt}\ {\isachardoublequoteopen}wf{\isacharunderscore}{\kern0pt}duration{\isacharunderscore}{\kern0pt}const\ {\isacharparenleft}{\kern0pt}Func{\isacharunderscore}{\kern0pt}Const\ op\ f\ vs{\isacharparenright}{\kern0pt}\ {\isasymlongleftrightarrow}\ \isanewline
\ \ \ \ \ \ \ \ {\isacharparenleft}{\kern0pt}case\ func{\isacharunderscore}{\kern0pt}sig\ f\ of\isanewline
\ \ \ \ \ \ \ \ \ \ \ \ None\ {\isasymRightarrow}\ False\isanewline
\ \ \ \ \ \ \ \ \ \ {\isacharbar}{\kern0pt}\ Some\ Ts\ {\isasymRightarrow}\ list{\isacharunderscore}{\kern0pt}all{\isadigit{2}}\ is{\isacharunderscore}{\kern0pt}of{\isacharunderscore}{\kern0pt}type\ vs\ Ts{\isacharparenright}{\kern0pt}{\isachardoublequoteclose}\isanewline
\isanewline
\ \ \isacommand{end}\isamarkupfalse%
\ %
\isamarkupcmt{Context fixing \isa{ty{\isacharunderscore}{\kern0pt}ent}%
}\isanewline
\isanewline
\isanewline
\ \ \isacommand{definition}\isamarkupfalse%
\ constT\ {\isacharcolon}{\kern0pt}{\isacharcolon}{\kern0pt}\ {\isachardoublequoteopen}object\ {\isasymrightharpoonup}\ type{\isachardoublequoteclose}\ \isakeyword{where}\isanewline
\ \ \ \ {\isachardoublequoteopen}constT\ {\isasymequiv}\ map{\isacharunderscore}{\kern0pt}of\ {\isacharparenleft}{\kern0pt}consts\ D{\isacharparenright}{\kern0pt}{\isachardoublequoteclose}%
\begin{isamarkuptext}%
An action schema is well-formed if the parameter names are distinct,
    and the precondition and effect is well-formed wrt.\ the parameters.%
\end{isamarkuptext}\isamarkuptrue%
\ \ \isacommand{fun}\isamarkupfalse%
\ wf{\isacharunderscore}{\kern0pt}action{\isacharunderscore}{\kern0pt}schema\ {\isacharcolon}{\kern0pt}{\isacharcolon}{\kern0pt}\ {\isachardoublequoteopen}ast{\isacharunderscore}{\kern0pt}action{\isacharunderscore}{\kern0pt}schema\ {\isasymRightarrow}\ bool{\isachardoublequoteclose}\ \isakeyword{where}\isanewline
\ \ \ \ {\isachardoublequoteopen}wf{\isacharunderscore}{\kern0pt}action{\isacharunderscore}{\kern0pt}schema\ {\isacharparenleft}{\kern0pt}Simple{\isacharunderscore}{\kern0pt}Action{\isacharunderscore}{\kern0pt}Schema\ n\ params\ pre\ eff{\isacharparenright}{\kern0pt}\ {\isasymlongleftrightarrow}\ {\isacharparenleft}{\kern0pt}\isanewline
\ \ \ \ \ \ let\ tyt\ {\isacharequal}{\kern0pt}\ ty{\isacharunderscore}{\kern0pt}term\ {\isacharparenleft}{\kern0pt}map{\isacharunderscore}{\kern0pt}of\ params{\isacharparenright}{\kern0pt}\ constT\ in\isanewline
\ \ \ \ \ \ \ \ distinct\ {\isacharparenleft}{\kern0pt}map\ fst\ params{\isacharparenright}{\kern0pt}\isanewline
\ \ \ \ \ \ {\isasymand}\ wf{\isacharunderscore}{\kern0pt}fmla\ tyt\ pre\isanewline
\ \ \ \ \ \ {\isasymand}\ wf{\isacharunderscore}{\kern0pt}effect\ tyt\ eff{\isacharparenright}{\kern0pt}{\isachardoublequoteclose}\isanewline
\ \ {\isacharbar}{\kern0pt}\ {\isachardoublequoteopen}wf{\isacharunderscore}{\kern0pt}action{\isacharunderscore}{\kern0pt}schema\ {\isacharparenleft}{\kern0pt}Durative{\isacharunderscore}{\kern0pt}Action{\isacharunderscore}{\kern0pt}Schema\ n\ params\ d\ cond\ eff{\isacharparenright}{\kern0pt}\ {\isasymlongleftrightarrow}\ {\isacharparenleft}{\kern0pt}\isanewline
\ \ \ \ \ \ let\ tyt\ {\isacharequal}{\kern0pt}\ ty{\isacharunderscore}{\kern0pt}term\ {\isacharparenleft}{\kern0pt}map{\isacharunderscore}{\kern0pt}of\ params{\isacharparenright}{\kern0pt}\ constT\ in\isanewline
\ \ \ \ \ \ \ \ distinct\ {\isacharparenleft}{\kern0pt}map\ fst\ params{\isacharparenright}{\kern0pt}\isanewline
\ \ \ \ \ \ {\isasymand}\ {\isacharparenleft}{\kern0pt}{\isasymforall}{\isacharparenleft}{\kern0pt}t{\isacharcomma}{\kern0pt}c{\isacharparenright}{\kern0pt}\ {\isasymin}\ set\ cond{\isachardot}{\kern0pt}\ wf{\isacharunderscore}{\kern0pt}fmla\ tyt\ c{\isacharparenright}{\kern0pt}\isanewline
\ \ \ \ \ \ {\isasymand}\ {\isacharparenleft}{\kern0pt}{\isasymforall}{\isacharparenleft}{\kern0pt}t{\isacharcomma}{\kern0pt}e{\isacharparenright}{\kern0pt}\ {\isasymin}\ set\ eff{\isachardot}{\kern0pt}\ wf{\isacharunderscore}{\kern0pt}effect\ tyt\ e\ {\isasymand}\ t\ {\isasymnoteq}\ Over{\isacharunderscore}{\kern0pt}All{\isacharparenright}{\kern0pt}\isanewline
\ \ \ \ \ \ {\isasymand}\ wf{\isacharunderscore}{\kern0pt}duration{\isacharunderscore}{\kern0pt}const\ tyt\ d{\isacharparenright}{\kern0pt}{\isachardoublequoteclose}%
\begin{isamarkuptext}%
A type is well-formed if it consists only of declared primitive types,
     and the type object.%
\end{isamarkuptext}\isamarkuptrue%
\ \ \isacommand{fun}\isamarkupfalse%
\ wf{\isacharunderscore}{\kern0pt}type\ \isakeyword{where}\isanewline
\ \ \ \ {\isachardoublequoteopen}wf{\isacharunderscore}{\kern0pt}type\ {\isacharparenleft}{\kern0pt}Either\ Ts{\isacharparenright}{\kern0pt}\ {\isasymlongleftrightarrow}\ set\ Ts\ {\isasymsubseteq}\ insert\ {\isacharprime}{\kern0pt}{\isacharprime}{\kern0pt}object{\isacharprime}{\kern0pt}{\isacharprime}{\kern0pt}\ {\isacharparenleft}{\kern0pt}fst{\isacharbackquote}{\kern0pt}set\ {\isacharparenleft}{\kern0pt}types\ D{\isacharparenright}{\kern0pt}{\isacharparenright}{\kern0pt}{\isachardoublequoteclose}%
\begin{isamarkuptext}%
A predicate is well-formed if its argument types are well-formed.%
\end{isamarkuptext}\isamarkuptrue%
\ \ \isacommand{fun}\isamarkupfalse%
\ wf{\isacharunderscore}{\kern0pt}predicate{\isacharunderscore}{\kern0pt}decl\ \isakeyword{where}\isanewline
\ \ \ \ {\isachardoublequoteopen}wf{\isacharunderscore}{\kern0pt}predicate{\isacharunderscore}{\kern0pt}decl\ {\isacharparenleft}{\kern0pt}PredDecl\ p\ Ts{\isacharparenright}{\kern0pt}\ {\isasymlongleftrightarrow}\ {\isacharparenleft}{\kern0pt}{\isasymforall}T{\isasymin}set\ Ts{\isachardot}{\kern0pt}\ wf{\isacharunderscore}{\kern0pt}type\ T{\isacharparenright}{\kern0pt}{\isachardoublequoteclose}%
\begin{isamarkuptext}%
A function is well-formed if its argument types are well-formed.%
\end{isamarkuptext}\isamarkuptrue%
\ \ \isacommand{fun}\isamarkupfalse%
\ wf{\isacharunderscore}{\kern0pt}function{\isacharunderscore}{\kern0pt}decl\ \isakeyword{where}\isanewline
\ \ \ \ {\isachardoublequoteopen}wf{\isacharunderscore}{\kern0pt}function{\isacharunderscore}{\kern0pt}decl\ {\isacharparenleft}{\kern0pt}FuncDecl\ f\ Ts{\isacharparenright}{\kern0pt}\ {\isasymlongleftrightarrow}\ {\isacharparenleft}{\kern0pt}{\isasymforall}T{\isasymin}set\ Ts{\isachardot}{\kern0pt}\ wf{\isacharunderscore}{\kern0pt}type\ T{\isacharparenright}{\kern0pt}{\isachardoublequoteclose}%
\begin{isamarkuptext}%
The types declaration is well-formed, if all supertypes are declared types (or object)%
\end{isamarkuptext}\isamarkuptrue%
\ \ \isacommand{definition}\isamarkupfalse%
\ {\isachardoublequoteopen}wf{\isacharunderscore}{\kern0pt}types\ {\isasymequiv}\ snd{\isacharbackquote}{\kern0pt}set\ {\isacharparenleft}{\kern0pt}types\ D{\isacharparenright}{\kern0pt}\ {\isasymsubseteq}\ insert\ {\isacharprime}{\kern0pt}{\isacharprime}{\kern0pt}object{\isacharprime}{\kern0pt}{\isacharprime}{\kern0pt}\ {\isacharparenleft}{\kern0pt}fst{\isacharbackquote}{\kern0pt}set\ {\isacharparenleft}{\kern0pt}types\ D{\isacharparenright}{\kern0pt}{\isacharparenright}{\kern0pt}{\isachardoublequoteclose}%
\begin{isamarkuptext}%
A domain is well-formed if

%
\begin{itemize}%
\item there are no duplicate declared predicate names,

\item all declared predicates are well-formed,

\item there are no duplicate declared function names,

\item all declared functions are well-formed,

\item there are no duplicate action names,

\item and all declared actions are well-formed%
\end{itemize}%
\end{isamarkuptext}\isamarkuptrue%
\ \ \isacommand{definition}\isamarkupfalse%
\ wf{\isacharunderscore}{\kern0pt}domain\ {\isacharcolon}{\kern0pt}{\isacharcolon}{\kern0pt}\ {\isachardoublequoteopen}bool{\isachardoublequoteclose}\ \isakeyword{where}\isanewline
\ \ \ \ {\isachardoublequoteopen}wf{\isacharunderscore}{\kern0pt}domain\ {\isasymequiv}\isanewline
\ \ \ \ \ \ wf{\isacharunderscore}{\kern0pt}types\isanewline
\ \ \ \ {\isasymand}\ distinct\ {\isacharparenleft}{\kern0pt}map\ {\isacharparenleft}{\kern0pt}predicate{\isacharunderscore}{\kern0pt}decl{\isachardot}{\kern0pt}pred{\isacharparenright}{\kern0pt}\ {\isacharparenleft}{\kern0pt}predicates\ D{\isacharparenright}{\kern0pt}{\isacharparenright}{\kern0pt}\isanewline
\ \ \ \ {\isasymand}\ {\isacharparenleft}{\kern0pt}{\isasymforall}p{\isasymin}set\ {\isacharparenleft}{\kern0pt}predicates\ D{\isacharparenright}{\kern0pt}{\isachardot}{\kern0pt}\ wf{\isacharunderscore}{\kern0pt}predicate{\isacharunderscore}{\kern0pt}decl\ p{\isacharparenright}{\kern0pt}\isanewline
\ \ \ \ {\isasymand}\ distinct\ {\isacharparenleft}{\kern0pt}map\ {\isacharparenleft}{\kern0pt}function{\isacharunderscore}{\kern0pt}decl{\isachardot}{\kern0pt}func{\isacharparenright}{\kern0pt}\ {\isacharparenleft}{\kern0pt}functions\ D{\isacharparenright}{\kern0pt}{\isacharparenright}{\kern0pt}\isanewline
\ \ \ \ {\isasymand}\ {\isacharparenleft}{\kern0pt}{\isasymforall}f{\isasymin}set\ {\isacharparenleft}{\kern0pt}functions\ D{\isacharparenright}{\kern0pt}{\isachardot}{\kern0pt}\ wf{\isacharunderscore}{\kern0pt}function{\isacharunderscore}{\kern0pt}decl\ f{\isacharparenright}{\kern0pt}\isanewline
\ \ \ \ {\isasymand}\ distinct\ {\isacharparenleft}{\kern0pt}map\ fst\ {\isacharparenleft}{\kern0pt}consts\ D{\isacharparenright}{\kern0pt}{\isacharparenright}{\kern0pt}\isanewline
\ \ \ \ {\isasymand}\ {\isacharparenleft}{\kern0pt}{\isasymforall}{\isacharparenleft}{\kern0pt}n{\isacharcomma}{\kern0pt}T{\isacharparenright}{\kern0pt}{\isasymin}set\ {\isacharparenleft}{\kern0pt}consts\ D{\isacharparenright}{\kern0pt}{\isachardot}{\kern0pt}\ wf{\isacharunderscore}{\kern0pt}type\ T{\isacharparenright}{\kern0pt}\isanewline
\ \ \ \ {\isasymand}\ distinct\ {\isacharparenleft}{\kern0pt}map\ ast{\isacharunderscore}{\kern0pt}action{\isacharunderscore}{\kern0pt}schema{\isachardot}{\kern0pt}name\ {\isacharparenleft}{\kern0pt}actions\ D{\isacharparenright}{\kern0pt}{\isacharparenright}{\kern0pt}\isanewline
\ \ \ \ {\isasymand}\ {\isacharparenleft}{\kern0pt}{\isasymforall}a{\isasymin}set\ {\isacharparenleft}{\kern0pt}actions\ D{\isacharparenright}{\kern0pt}{\isachardot}{\kern0pt}\ wf{\isacharunderscore}{\kern0pt}action{\isacharunderscore}{\kern0pt}schema\ a{\isacharparenright}{\kern0pt}\isanewline
\ \ \ \ {\isachardoublequoteclose}\isanewline
\isanewline
\isacommand{end}\isamarkupfalse%
\ %
\isamarkupcmt{locale \isa{ast{\isacharunderscore}{\kern0pt}domain}%
}%
\begin{isamarkuptext}%
We fix a problem, and also include the definitions for the domain
  of this problem.%
\end{isamarkuptext}\isamarkuptrue%
\isacommand{locale}\isamarkupfalse%
\ ast{\isacharunderscore}{\kern0pt}problem\ {\isacharequal}{\kern0pt}\ ast{\isacharunderscore}{\kern0pt}domain\ {\isachardoublequoteopen}domain\ P{\isachardoublequoteclose}\isanewline
\ \ \isakeyword{for}\ P\ {\isacharcolon}{\kern0pt}{\isacharcolon}{\kern0pt}\ ast{\isacharunderscore}{\kern0pt}problem\isanewline
\isakeyword{begin}%
\begin{isamarkuptext}%
We refer to the problem domain as \isa{D}%
\end{isamarkuptext}\isamarkuptrue%
\ \ \isacommand{abbreviation}\isamarkupfalse%
\ {\isachardoublequoteopen}D\ {\isasymequiv}\ ast{\isacharunderscore}{\kern0pt}problem{\isachardot}{\kern0pt}domain\ P{\isachardoublequoteclose}\isanewline
\isanewline
\ \ \isacommand{definition}\isamarkupfalse%
\ objT\ {\isacharcolon}{\kern0pt}{\isacharcolon}{\kern0pt}\ {\isachardoublequoteopen}object\ {\isasymrightharpoonup}\ type{\isachardoublequoteclose}\ \isakeyword{where}\isanewline
\ \ \ \ {\isachardoublequoteopen}objT\ {\isasymequiv}\ map{\isacharunderscore}{\kern0pt}of\ {\isacharparenleft}{\kern0pt}objects\ P{\isacharparenright}{\kern0pt}\ {\isacharplus}{\kern0pt}{\isacharplus}{\kern0pt}\ constT{\isachardoublequoteclose}\isanewline
\isanewline
\ \ \isacommand{lemma}\isamarkupfalse%
\ objT{\isacharunderscore}{\kern0pt}alt{\isacharcolon}{\kern0pt}\ {\isachardoublequoteopen}objT\ {\isacharequal}{\kern0pt}\ map{\isacharunderscore}{\kern0pt}of\ {\isacharparenleft}{\kern0pt}consts\ D\ {\isacharat}{\kern0pt}\ objects\ P{\isacharparenright}{\kern0pt}{\isachardoublequoteclose}\isanewline
%
\isadelimproof
\ \ \ \ %
\endisadelimproof
%
\isatagproof
\isacommand{unfolding}\isamarkupfalse%
\ objT{\isacharunderscore}{\kern0pt}def\ constT{\isacharunderscore}{\kern0pt}def\isanewline
\ \ \ \ \isacommand{apply}\isamarkupfalse%
\ {\isacharparenleft}{\kern0pt}clarsimp{\isacharparenright}{\kern0pt}\isanewline
\ \ \ \ \isacommand{done}\isamarkupfalse%
%
\endisatagproof
{\isafoldproof}%
%
\isadelimproof
\isanewline
%
\endisadelimproof
\isanewline
\ \ \isacommand{definition}\isamarkupfalse%
\ wf{\isacharunderscore}{\kern0pt}fact\ {\isacharcolon}{\kern0pt}{\isacharcolon}{\kern0pt}\ {\isachardoublequoteopen}fact\ {\isasymRightarrow}\ bool{\isachardoublequoteclose}\ \isakeyword{where}\isanewline
\ \ \ \ {\isachardoublequoteopen}wf{\isacharunderscore}{\kern0pt}fact\ {\isacharequal}{\kern0pt}\ wf{\isacharunderscore}{\kern0pt}pred{\isacharunderscore}{\kern0pt}atom\ objT{\isachardoublequoteclose}\isanewline
\isanewline
\ \ \isacommand{fun}\isamarkupfalse%
\ wf{\isacharunderscore}{\kern0pt}func{\isacharunderscore}{\kern0pt}assign\ {\isacharcolon}{\kern0pt}{\isacharcolon}{\kern0pt}\ {\isachardoublequoteopen}object\ atom\ formula\ {\isasymRightarrow}\ bool{\isachardoublequoteclose}\ \isakeyword{where}\ \isanewline
\ \ \ \ {\isachardoublequoteopen}wf{\isacharunderscore}{\kern0pt}func{\isacharunderscore}{\kern0pt}assign\ {\isacharparenleft}{\kern0pt}Atom\ {\isacharparenleft}{\kern0pt}eqAtm\ {\isacharparenleft}{\kern0pt}FuncEnt\ f\ args{\isacharparenright}{\kern0pt}\ {\isacharparenleft}{\kern0pt}TimeEnt\ d{\isacharparenright}{\kern0pt}{\isacharparenright}{\kern0pt}{\isacharparenright}{\kern0pt}\ {\isasymlongleftrightarrow}\ wf{\isacharunderscore}{\kern0pt}func{\isacharunderscore}{\kern0pt}args\ objT\ {\isacharparenleft}{\kern0pt}f{\isacharcomma}{\kern0pt}args{\isacharparenright}{\kern0pt}{\isachardoublequoteclose}\isanewline
\ \ {\isacharbar}{\kern0pt}\ {\isachardoublequoteopen}wf{\isacharunderscore}{\kern0pt}func{\isacharunderscore}{\kern0pt}assign\ {\isacharunderscore}{\kern0pt}\ {\isasymlongleftrightarrow}\ False{\isachardoublequoteclose}\isanewline
\isanewline
\ \ \isacommand{lemma}\isamarkupfalse%
\ wf{\isacharunderscore}{\kern0pt}func{\isacharunderscore}{\kern0pt}assign{\isacharunderscore}{\kern0pt}is{\isacharunderscore}{\kern0pt}eqAtom{\isacharcolon}{\kern0pt}\ {\isachardoublequoteopen}wf{\isacharunderscore}{\kern0pt}func{\isacharunderscore}{\kern0pt}assign\ f\ {\isasymLongrightarrow}\ is{\isacharunderscore}{\kern0pt}eqAtom\ f{\isachardoublequoteclose}\isanewline
%
\isadelimproof
\ \ \ \ %
\endisadelimproof
%
\isatagproof
\isacommand{apply}\isamarkupfalse%
\ {\isacharparenleft}{\kern0pt}cases\ f{\isacharparenright}{\kern0pt}\isanewline
\ \ \ \ \isacommand{apply}\isamarkupfalse%
\ auto\isanewline
\ \ \ \ \isacommand{apply}\isamarkupfalse%
\ {\isacharparenleft}{\kern0pt}metis\ is{\isacharunderscore}{\kern0pt}eqAtom{\isachardot}{\kern0pt}simps{\isacharparenleft}{\kern0pt}{\isadigit{1}}{\isacharparenright}{\kern0pt}\ wf{\isacharunderscore}{\kern0pt}func{\isacharunderscore}{\kern0pt}assign{\isachardot}{\kern0pt}elims{\isacharparenleft}{\kern0pt}{\isadigit{2}}{\isacharparenright}{\kern0pt}{\isacharparenright}{\kern0pt}\isanewline
\ \ \ \ \isacommand{done}\isamarkupfalse%
%
\endisatagproof
{\isafoldproof}%
%
\isadelimproof
%
\endisadelimproof
%
\begin{isamarkuptext}%
This definition is needed for well-formedness of the initial model,
    and forward-references to the concept of world model.%
\end{isamarkuptext}\isamarkuptrue%
\ \ \isacommand{definition}\isamarkupfalse%
\ wf{\isacharunderscore}{\kern0pt}world{\isacharunderscore}{\kern0pt}model\ \isakeyword{where}\isanewline
\ \ \ \ {\isachardoublequoteopen}wf{\isacharunderscore}{\kern0pt}world{\isacharunderscore}{\kern0pt}model\ M\ {\isacharequal}{\kern0pt}\ {\isacharparenleft}{\kern0pt}{\isasymforall}f{\isasymin}M{\isachardot}{\kern0pt}\ wf{\isacharunderscore}{\kern0pt}fmla{\isacharunderscore}{\kern0pt}atom\ objT\ f{\isacharparenright}{\kern0pt}{\isachardoublequoteclose}\ \isanewline
\isanewline
\ \ \isanewline
\ \ \isacommand{definition}\isamarkupfalse%
\ wf{\isacharunderscore}{\kern0pt}problem\ \isakeyword{where}\isanewline
\ \ \ \ {\isachardoublequoteopen}wf{\isacharunderscore}{\kern0pt}problem\ {\isasymequiv}\isanewline
\ \ \ \ \ \ wf{\isacharunderscore}{\kern0pt}domain\isanewline
\ \ \ \ {\isasymand}\ distinct\ {\isacharparenleft}{\kern0pt}map\ fst\ {\isacharparenleft}{\kern0pt}objects\ P{\isacharparenright}{\kern0pt}\ {\isacharat}{\kern0pt}\ map\ fst\ {\isacharparenleft}{\kern0pt}consts\ D{\isacharparenright}{\kern0pt}{\isacharparenright}{\kern0pt}\isanewline
\ \ \ \ {\isasymand}\ {\isacharparenleft}{\kern0pt}{\isasymforall}{\isacharparenleft}{\kern0pt}n{\isacharcomma}{\kern0pt}T{\isacharparenright}{\kern0pt}{\isasymin}set\ {\isacharparenleft}{\kern0pt}objects\ P{\isacharparenright}{\kern0pt}{\isachardot}{\kern0pt}\ wf{\isacharunderscore}{\kern0pt}type\ T{\isacharparenright}{\kern0pt}\isanewline
\ \ \ \ {\isasymand}\ distinct\ {\isacharparenleft}{\kern0pt}init\ P{\isacharparenright}{\kern0pt}\isanewline
\ \ \ \ {\isasymand}\ {\isacharparenleft}{\kern0pt}{\isasymforall}f{\isasymin}set\ {\isacharparenleft}{\kern0pt}init\ P{\isacharparenright}{\kern0pt}{\isachardot}{\kern0pt}\ wf{\isacharunderscore}{\kern0pt}fmla{\isacharunderscore}{\kern0pt}atom\ objT\ f\ {\isasymor}\ wf{\isacharunderscore}{\kern0pt}func{\isacharunderscore}{\kern0pt}assign\ f{\isacharparenright}{\kern0pt}\isanewline
\ \ \ \ {\isasymand}\ wf{\isacharunderscore}{\kern0pt}fmla\ objT\ {\isacharparenleft}{\kern0pt}goal\ P{\isacharparenright}{\kern0pt}{\isachardoublequoteclose}\ \isanewline
\isanewline
\ \ \isacommand{fun}\isamarkupfalse%
\ wf{\isacharunderscore}{\kern0pt}effect{\isacharunderscore}{\kern0pt}inst\ {\isacharcolon}{\kern0pt}{\isacharcolon}{\kern0pt}\ {\isachardoublequoteopen}object\ ast{\isacharunderscore}{\kern0pt}effect\ {\isasymRightarrow}\ bool{\isachardoublequoteclose}\ \isakeyword{where}\isanewline
\ \ \ \ {\isachardoublequoteopen}wf{\isacharunderscore}{\kern0pt}effect{\isacharunderscore}{\kern0pt}inst\ {\isacharparenleft}{\kern0pt}Effect\ {\isacharparenleft}{\kern0pt}a{\isacharparenright}{\kern0pt}\ {\isacharparenleft}{\kern0pt}d{\isacharparenright}{\kern0pt}{\isacharparenright}{\kern0pt}\ {\isasymlongleftrightarrow}\ {\isacharparenleft}{\kern0pt}{\isasymforall}a{\isasymin}set\ a\ {\isasymunion}\ set\ d{\isachardot}{\kern0pt}\ wf{\isacharunderscore}{\kern0pt}fmla{\isacharunderscore}{\kern0pt}atom\ objT\ a{\isacharparenright}{\kern0pt}{\isachardoublequoteclose}\isanewline
\isanewline
\ \ \isacommand{lemma}\isamarkupfalse%
\ wf{\isacharunderscore}{\kern0pt}effect{\isacharunderscore}{\kern0pt}inst{\isacharunderscore}{\kern0pt}alt{\isacharcolon}{\kern0pt}\ {\isachardoublequoteopen}wf{\isacharunderscore}{\kern0pt}effect{\isacharunderscore}{\kern0pt}inst\ eff\ {\isacharequal}{\kern0pt}\ wf{\isacharunderscore}{\kern0pt}effect\ objT\ eff{\isachardoublequoteclose}\isanewline
%
\isadelimproof
\ \ \ \ %
\endisadelimproof
%
\isatagproof
\isacommand{by}\isamarkupfalse%
\ {\isacharparenleft}{\kern0pt}cases\ eff{\isacharparenright}{\kern0pt}\ auto%
\endisatagproof
{\isafoldproof}%
%
\isadelimproof
\isanewline
%
\endisadelimproof
\isanewline
\isacommand{end}\isamarkupfalse%
\ %
\isamarkupcmt{locale \isa{ast{\isacharunderscore}{\kern0pt}problem}%
}%
\begin{isamarkuptext}%
Locale to express a well-formed domain%
\end{isamarkuptext}\isamarkuptrue%
\isacommand{locale}\isamarkupfalse%
\ wf{\isacharunderscore}{\kern0pt}ast{\isacharunderscore}{\kern0pt}domain\ {\isacharequal}{\kern0pt}\ ast{\isacharunderscore}{\kern0pt}domain\ {\isacharplus}{\kern0pt}\isanewline
\ \ \isakeyword{assumes}\ wf{\isacharunderscore}{\kern0pt}domain{\isacharcolon}{\kern0pt}\ wf{\isacharunderscore}{\kern0pt}domain%
\begin{isamarkuptext}%
Locale to express a well-formed problem%
\end{isamarkuptext}\isamarkuptrue%
\isacommand{locale}\isamarkupfalse%
\ wf{\isacharunderscore}{\kern0pt}ast{\isacharunderscore}{\kern0pt}problem\ {\isacharequal}{\kern0pt}\ ast{\isacharunderscore}{\kern0pt}problem\ P\ \isakeyword{for}\ P\ {\isacharplus}{\kern0pt}\isanewline
\ \ \isakeyword{assumes}\ wf{\isacharunderscore}{\kern0pt}problem{\isacharcolon}{\kern0pt}\ wf{\isacharunderscore}{\kern0pt}problem\isanewline
\isakeyword{begin}\isanewline
\ \ \isacommand{sublocale}\isamarkupfalse%
\ wf{\isacharunderscore}{\kern0pt}ast{\isacharunderscore}{\kern0pt}domain\ {\isachardoublequoteopen}domain\ P{\isachardoublequoteclose}\isanewline
%
\isadelimproof
\ \ \ \ %
\endisadelimproof
%
\isatagproof
\isacommand{apply}\isamarkupfalse%
\ unfold{\isacharunderscore}{\kern0pt}locales\isanewline
\ \ \ \ \isacommand{using}\isamarkupfalse%
\ wf{\isacharunderscore}{\kern0pt}problem\isanewline
\ \ \ \ \isacommand{unfolding}\isamarkupfalse%
\ wf{\isacharunderscore}{\kern0pt}problem{\isacharunderscore}{\kern0pt}def\ \isacommand{by}\isamarkupfalse%
\ simp%
\endisatagproof
{\isafoldproof}%
%
\isadelimproof
\isanewline
%
\endisadelimproof
\isanewline
\isacommand{end}\isamarkupfalse%
\ %
\isamarkupcmt{locale \isa{wf{\isacharunderscore}{\kern0pt}ast{\isacharunderscore}{\kern0pt}problem}%
}%
\isadelimdocument
%
\endisadelimdocument
%
\isatagdocument
%
\isamarkupsubsection{PDDL Semantics%
}
\isamarkuptrue%
%
\endisatagdocument
{\isafolddocument}%
%
\isadelimdocument
%
\endisadelimdocument
\isacommand{context}\isamarkupfalse%
\ ast{\isacharunderscore}{\kern0pt}domain\ \isakeyword{begin}\isanewline
\isanewline
\ \ \isacommand{definition}\isamarkupfalse%
\ resolve{\isacharunderscore}{\kern0pt}action{\isacharunderscore}{\kern0pt}schema\ {\isacharcolon}{\kern0pt}{\isacharcolon}{\kern0pt}\ {\isachardoublequoteopen}name\ {\isasymrightharpoonup}\ ast{\isacharunderscore}{\kern0pt}action{\isacharunderscore}{\kern0pt}schema{\isachardoublequoteclose}\ \isakeyword{where}\isanewline
\ \ \ \ {\isachardoublequoteopen}resolve{\isacharunderscore}{\kern0pt}action{\isacharunderscore}{\kern0pt}schema\ n\ {\isacharequal}{\kern0pt}\ index{\isacharunderscore}{\kern0pt}by\ ast{\isacharunderscore}{\kern0pt}action{\isacharunderscore}{\kern0pt}schema{\isachardot}{\kern0pt}name\ {\isacharparenleft}{\kern0pt}actions\ D{\isacharparenright}{\kern0pt}\ n{\isachardoublequoteclose}\isanewline
\isanewline
\ \ \isacommand{fun}\isamarkupfalse%
\ subst{\isacharunderscore}{\kern0pt}term\ \isakeyword{where}\isanewline
\ \ \ \ {\isachardoublequoteopen}subst{\isacharunderscore}{\kern0pt}term\ psubst\ {\isacharparenleft}{\kern0pt}term{\isachardot}{\kern0pt}VAR\ x{\isacharparenright}{\kern0pt}\ {\isacharequal}{\kern0pt}\ psubst\ x{\isachardoublequoteclose}\isanewline
\ \ {\isacharbar}{\kern0pt}\ {\isachardoublequoteopen}subst{\isacharunderscore}{\kern0pt}term\ psubst\ {\isacharparenleft}{\kern0pt}term{\isachardot}{\kern0pt}CONST\ c{\isacharparenright}{\kern0pt}\ {\isacharequal}{\kern0pt}\ c{\isachardoublequoteclose}%
\begin{isamarkuptext}%
To instantiate an action schema, we first compute a substitution from
    parameters to objects, and then apply this substitution to the
    precondition and effect. The substitution is applied via the \isa{map{\isacharunderscore}{\kern0pt}xxx}
    functions generated by the datatype package.%
\end{isamarkuptext}\isamarkuptrue%
\ \ \isacommand{fun}\isamarkupfalse%
\ instantiate{\isacharunderscore}{\kern0pt}action{\isacharunderscore}{\kern0pt}schema\ {\isacharcolon}{\kern0pt}{\isacharcolon}{\kern0pt}\ {\isachardoublequoteopen}ast{\isacharunderscore}{\kern0pt}action{\isacharunderscore}{\kern0pt}schema\ {\isasymRightarrow}\ object\ list\ {\isasymRightarrow}\ ground{\isacharunderscore}{\kern0pt}action{\isachardoublequoteclose}\ \isakeyword{where}\isanewline
\ \ \ \ {\isachardoublequoteopen}instantiate{\isacharunderscore}{\kern0pt}action{\isacharunderscore}{\kern0pt}schema\ {\isacharparenleft}{\kern0pt}Simple{\isacharunderscore}{\kern0pt}Action{\isacharunderscore}{\kern0pt}Schema\ n\ params\ pre\ eff{\isacharparenright}{\kern0pt}\ args\ {\isacharequal}{\kern0pt}\ \isanewline
\ \ \ \ \ \ {\isacharparenleft}{\kern0pt}let\isanewline
\ \ \ \ \ \ \ \ tsubst\ {\isacharequal}{\kern0pt}\ subst{\isacharunderscore}{\kern0pt}term\ {\isacharparenleft}{\kern0pt}the\ o\ {\isacharparenleft}{\kern0pt}map{\isacharunderscore}{\kern0pt}of\ {\isacharparenleft}{\kern0pt}zip\ {\isacharparenleft}{\kern0pt}map\ fst\ params{\isacharparenright}{\kern0pt}\ args{\isacharparenright}{\kern0pt}{\isacharparenright}{\kern0pt}{\isacharparenright}{\kern0pt}{\isacharsemicolon}{\kern0pt}\isanewline
\ \ \ \ \ \ \ \ pre{\isacharunderscore}{\kern0pt}inst\ {\isacharequal}{\kern0pt}\ {\isacharparenleft}{\kern0pt}map{\isacharunderscore}{\kern0pt}formula\ o\ map{\isacharunderscore}{\kern0pt}atom{\isacharparenright}{\kern0pt}\ tsubst\ pre{\isacharsemicolon}{\kern0pt}\isanewline
\ \ \ \ \ \ \ \ eff{\isacharunderscore}{\kern0pt}inst\ {\isacharequal}{\kern0pt}\ {\isacharparenleft}{\kern0pt}map{\isacharunderscore}{\kern0pt}ast{\isacharunderscore}{\kern0pt}effect{\isacharparenright}{\kern0pt}\ tsubst\ eff\isanewline
\ \ \ \ \ \ in\isanewline
\ \ \ \ \ \ \ \ Ground{\isacharunderscore}{\kern0pt}Action\ pre{\isacharunderscore}{\kern0pt}inst\ eff{\isacharunderscore}{\kern0pt}inst\isanewline
\ \ \ \ \ \ {\isacharparenright}{\kern0pt}{\isachardoublequoteclose}%
\begin{isamarkuptext}%
Filter a given timed-list by a time specifier. 
  \texttt{(at start ...)}, \texttt{(at end ...)}, or \texttt{(over all ...)} conditions/effects%
\end{isamarkuptext}\isamarkuptrue%
\ \ \isacommand{definition}\isamarkupfalse%
\ filter{\isacharunderscore}{\kern0pt}time{\isacharunderscore}{\kern0pt}spec\ {\isacharcolon}{\kern0pt}{\isacharcolon}{\kern0pt}\ {\isachardoublequoteopen}temporal{\isacharunderscore}{\kern0pt}annotation\ {\isasymRightarrow}\ {\isacharparenleft}{\kern0pt}temporal{\isacharunderscore}{\kern0pt}annotation\ {\isasymtimes}\ {\isacharprime}{\kern0pt}a{\isacharparenright}{\kern0pt}\ list\ {\isasymRightarrow}\ {\isacharprime}{\kern0pt}a\ list{\isachardoublequoteclose}\ \isakeyword{where}\isanewline
\ \ \ \ {\isachardoublequoteopen}filter{\isacharunderscore}{\kern0pt}time{\isacharunderscore}{\kern0pt}spec\ t\ l\ {\isacharequal}{\kern0pt}\ map\ snd\ {\isacharparenleft}{\kern0pt}filter\ {\isacharparenleft}{\kern0pt}{\isacharparenleft}{\kern0pt}{\isacharparenleft}{\kern0pt}{\isacharequal}{\kern0pt}{\isacharparenright}{\kern0pt}\ t{\isacharparenright}{\kern0pt}\ o\ fst{\isacharparenright}{\kern0pt}\ l{\isacharparenright}{\kern0pt}{\isachardoublequoteclose}\isanewline
\isanewline
\ \ \isacommand{lemma}\isamarkupfalse%
\ filter{\isacharunderscore}{\kern0pt}time{\isacharunderscore}{\kern0pt}spec{\isacharunderscore}{\kern0pt}correct{\isacharcolon}{\kern0pt}\ {\isachardoublequoteopen}{\isasymforall}x\ {\isasymin}\ set\ {\isacharparenleft}{\kern0pt}filter{\isacharunderscore}{\kern0pt}time{\isacharunderscore}{\kern0pt}spec\ t\ xs{\isacharparenright}{\kern0pt}{\isachardot}{\kern0pt}\ {\isacharparenleft}{\kern0pt}t{\isacharcomma}{\kern0pt}x{\isacharparenright}{\kern0pt}\ {\isasymin}\ set\ xs{\isachardoublequoteclose}\isanewline
%
\isadelimproof
\ \ \ \ %
\endisadelimproof
%
\isatagproof
\isacommand{unfolding}\isamarkupfalse%
\ filter{\isacharunderscore}{\kern0pt}time{\isacharunderscore}{\kern0pt}spec{\isacharunderscore}{\kern0pt}def\isanewline
\ \ \ \ \isacommand{by}\isamarkupfalse%
\ {\isacharparenleft}{\kern0pt}induction\ xs\ arbitrary{\isacharcolon}{\kern0pt}\ t{\isacharparenright}{\kern0pt}\ auto%
\endisatagproof
{\isafoldproof}%
%
\isadelimproof
%
\endisadelimproof
%
\begin{isamarkuptext}%
auxiliary-function to conjunct multiple effects into one single effect. 
  Needed to construct effect after filtering by a time specifier%
\end{isamarkuptext}\isamarkuptrue%
\ \ \isacommand{definition}\isamarkupfalse%
\ conjunct{\isacharunderscore}{\kern0pt}effects\ {\isacharcolon}{\kern0pt}{\isacharcolon}{\kern0pt}\ {\isachardoublequoteopen}{\isacharparenleft}{\kern0pt}{\isacharprime}{\kern0pt}a\ ast{\isacharunderscore}{\kern0pt}effect{\isacharparenright}{\kern0pt}\ list\ {\isasymRightarrow}\ {\isacharprime}{\kern0pt}a\ ast{\isacharunderscore}{\kern0pt}effect{\isachardoublequoteclose}\ {\isacharparenleft}{\kern0pt}{\isachardoublequoteopen}{\isasymAnd}\isactrlsub e\isactrlsub f\isactrlsub f\ {\isacharunderscore}{\kern0pt}{\isachardoublequoteclose}\ {\isacharbrackleft}{\kern0pt}{\isadigit{4}}{\isadigit{0}}{\isacharbrackright}{\kern0pt}\ {\isadigit{4}}{\isadigit{0}}{\isacharparenright}{\kern0pt}\ \isakeyword{where}\isanewline
\ \ \ \ {\isachardoublequoteopen}conjunct{\isacharunderscore}{\kern0pt}effects\ l\ {\isacharequal}{\kern0pt}\ Effect\ {\isacharparenleft}{\kern0pt}{\isacharparenleft}{\kern0pt}concat\ o\ {\isacharparenleft}{\kern0pt}map\ adds{\isacharparenright}{\kern0pt}{\isacharparenright}{\kern0pt}\ l{\isacharparenright}{\kern0pt}\ {\isacharparenleft}{\kern0pt}{\isacharparenleft}{\kern0pt}concat\ o\ {\isacharparenleft}{\kern0pt}map\ dels{\isacharparenright}{\kern0pt}{\isacharparenright}{\kern0pt}\ l{\isacharparenright}{\kern0pt}{\isachardoublequoteclose}%
\begin{isamarkuptext}%
instantiate a given Durative Action Schema; 
  Note: if an action schema is well formed then the effect for the snap action for the 
  temporal annotation \texttt{(over all ...)} will be an empty effect (\texttt{Effect [] []})%
\end{isamarkuptext}\isamarkuptrue%
\ \ \isacommand{fun}\isamarkupfalse%
\ inst{\isacharunderscore}{\kern0pt}snap{\isacharunderscore}{\kern0pt}action\ {\isacharcolon}{\kern0pt}{\isacharcolon}{\kern0pt}\ {\isachardoublequoteopen}ast{\isacharunderscore}{\kern0pt}action{\isacharunderscore}{\kern0pt}schema\ {\isasymRightarrow}\ object\ list\ {\isasymRightarrow}\ temporal{\isacharunderscore}{\kern0pt}annotation\ {\isasymRightarrow}\ ground{\isacharunderscore}{\kern0pt}action{\isachardoublequoteclose}\ \isakeyword{where}\isanewline
\ \ \ \ {\isachardoublequoteopen}inst{\isacharunderscore}{\kern0pt}snap{\isacharunderscore}{\kern0pt}action\ {\isacharparenleft}{\kern0pt}Durative{\isacharunderscore}{\kern0pt}Action{\isacharunderscore}{\kern0pt}Schema\ n\ params\ d\ cond\ eff{\isacharparenright}{\kern0pt}\ args\ ta\ {\isacharequal}{\kern0pt}\ \isanewline
\ \ \ \ \ \ {\isacharparenleft}{\kern0pt}let\ \isanewline
\ \ \ \ \ \ \ \ tsubst\ {\isacharequal}{\kern0pt}\ subst{\isacharunderscore}{\kern0pt}term\ {\isacharparenleft}{\kern0pt}the\ o\ {\isacharparenleft}{\kern0pt}map{\isacharunderscore}{\kern0pt}of\ {\isacharparenleft}{\kern0pt}zip\ {\isacharparenleft}{\kern0pt}map\ fst\ params{\isacharparenright}{\kern0pt}\ args{\isacharparenright}{\kern0pt}{\isacharparenright}{\kern0pt}{\isacharparenright}{\kern0pt}{\isacharsemicolon}{\kern0pt}\isanewline
\ \ \ \ \ \ \ \ pre\isactrlsub i\ {\isacharequal}{\kern0pt}\ {\isacharparenleft}{\kern0pt}map{\isacharunderscore}{\kern0pt}formula\ o\ map{\isacharunderscore}{\kern0pt}atom{\isacharparenright}{\kern0pt}\ tsubst\ {\isacharparenleft}{\kern0pt}BigAnd\ {\isacharparenleft}{\kern0pt}filter{\isacharunderscore}{\kern0pt}time{\isacharunderscore}{\kern0pt}spec\ ta\ cond{\isacharparenright}{\kern0pt}{\isacharparenright}{\kern0pt}{\isacharsemicolon}{\kern0pt}\isanewline
\ \ \ \ \ \ \ \ eff\isactrlsub i\ {\isacharequal}{\kern0pt}\ {\isacharparenleft}{\kern0pt}map{\isacharunderscore}{\kern0pt}ast{\isacharunderscore}{\kern0pt}effect{\isacharparenright}{\kern0pt}\ tsubst\ {\isacharparenleft}{\kern0pt}{\isasymAnd}\isactrlsub e\isactrlsub f\isactrlsub f\ {\isacharparenleft}{\kern0pt}filter{\isacharunderscore}{\kern0pt}time{\isacharunderscore}{\kern0pt}spec\ ta\ eff{\isacharparenright}{\kern0pt}{\isacharparenright}{\kern0pt}\ \isanewline
\ \ \ \ \ \ in\ \isanewline
\ \ \ \ \ \ \ \ Ground{\isacharunderscore}{\kern0pt}Action\ pre\isactrlsub i\ eff\isactrlsub i\isanewline
\ \ \ \ \ \ {\isacharparenright}{\kern0pt}{\isachardoublequoteclose}\isanewline
\isanewline
\isacommand{end}\isamarkupfalse%
\ %
\isamarkupcmt{Context of \isa{ast{\isacharunderscore}{\kern0pt}domain}%
}\isanewline
\isanewline
\isacommand{context}\isamarkupfalse%
\ wf{\isacharunderscore}{\kern0pt}ast{\isacharunderscore}{\kern0pt}domain\ \isakeyword{begin}%
\begin{isamarkuptext}%
Resolving an action yields a well-founded action schema.%
\end{isamarkuptext}\isamarkuptrue%
\ \ \isacommand{lemma}\isamarkupfalse%
\ resolve{\isacharunderscore}{\kern0pt}action{\isacharunderscore}{\kern0pt}wf{\isacharcolon}{\kern0pt}\isanewline
\ \ \ \ \isakeyword{assumes}\ {\isachardoublequoteopen}resolve{\isacharunderscore}{\kern0pt}action{\isacharunderscore}{\kern0pt}schema\ n\ {\isacharequal}{\kern0pt}\ Some\ a{\isachardoublequoteclose}\isanewline
\ \ \ \ \isakeyword{shows}\ {\isachardoublequoteopen}wf{\isacharunderscore}{\kern0pt}action{\isacharunderscore}{\kern0pt}schema\ a{\isachardoublequoteclose}\isanewline
%
\isadelimproof
\ \ %
\endisadelimproof
%
\isatagproof
\isacommand{proof}\isamarkupfalse%
\ {\isacharminus}{\kern0pt}\isanewline
\ \ \ \ \isacommand{from}\isamarkupfalse%
\ wf{\isacharunderscore}{\kern0pt}domain\ \isacommand{have}\isamarkupfalse%
\isanewline
\ \ \ \ \ \ X{\isadigit{1}}{\isacharcolon}{\kern0pt}\ {\isachardoublequoteopen}distinct\ {\isacharparenleft}{\kern0pt}map\ ast{\isacharunderscore}{\kern0pt}action{\isacharunderscore}{\kern0pt}schema{\isachardot}{\kern0pt}name\ {\isacharparenleft}{\kern0pt}actions\ D{\isacharparenright}{\kern0pt}{\isacharparenright}{\kern0pt}{\isachardoublequoteclose}\isanewline
\ \ \ \ \ \ \isakeyword{and}\ X{\isadigit{2}}{\isacharcolon}{\kern0pt}\ {\isachardoublequoteopen}{\isasymforall}a{\isasymin}set\ {\isacharparenleft}{\kern0pt}actions\ D{\isacharparenright}{\kern0pt}{\isachardot}{\kern0pt}\ wf{\isacharunderscore}{\kern0pt}action{\isacharunderscore}{\kern0pt}schema\ a{\isachardoublequoteclose}\isanewline
\ \ \ \ \ \ \isacommand{unfolding}\isamarkupfalse%
\ wf{\isacharunderscore}{\kern0pt}domain{\isacharunderscore}{\kern0pt}def\ \isacommand{by}\isamarkupfalse%
\ auto\isanewline
\isanewline
\ \ \ \ \isacommand{show}\isamarkupfalse%
\ {\isacharquery}{\kern0pt}thesis\isanewline
\ \ \ \ \ \ \isacommand{using}\isamarkupfalse%
\ assms\ \isacommand{unfolding}\isamarkupfalse%
\ resolve{\isacharunderscore}{\kern0pt}action{\isacharunderscore}{\kern0pt}schema{\isacharunderscore}{\kern0pt}def\isanewline
\ \ \ \ \ \ \isacommand{by}\isamarkupfalse%
\ {\isacharparenleft}{\kern0pt}auto\ simp\ add{\isacharcolon}{\kern0pt}\ index{\isacharunderscore}{\kern0pt}by{\isacharunderscore}{\kern0pt}eq{\isacharunderscore}{\kern0pt}Some{\isacharunderscore}{\kern0pt}eq{\isacharbrackleft}{\kern0pt}OF\ X{\isadigit{1}}{\isacharbrackright}{\kern0pt}\ X{\isadigit{2}}{\isacharparenright}{\kern0pt}\isanewline
\ \ \isacommand{qed}\isamarkupfalse%
%
\endisatagproof
{\isafoldproof}%
%
\isadelimproof
\isanewline
%
\endisadelimproof
\isanewline
\isacommand{end}\isamarkupfalse%
\ %
\isamarkupcmt{Context of \isa{ast{\isacharunderscore}{\kern0pt}domain}%
}\isanewline
\isanewline
\isacommand{context}\isamarkupfalse%
\ ast{\isacharunderscore}{\kern0pt}problem\ \isakeyword{begin}%
\begin{isamarkuptext}%
Initial model%
\end{isamarkuptext}\isamarkuptrue%
\ \ \isacommand{definition}\isamarkupfalse%
\ I\ {\isacharcolon}{\kern0pt}{\isacharcolon}{\kern0pt}\ {\isachardoublequoteopen}world{\isacharunderscore}{\kern0pt}model{\isachardoublequoteclose}\ \isakeyword{where}\isanewline
\ \ \ \ {\isachardoublequoteopen}I\ {\isasymequiv}\ set\ {\isacharparenleft}{\kern0pt}filter\ is{\isacharunderscore}{\kern0pt}predAtom\ {\isacharparenleft}{\kern0pt}init\ P{\isacharparenright}{\kern0pt}{\isacharparenright}{\kern0pt}{\isachardoublequoteclose}%
\begin{isamarkuptext}%
Resolve a plan action and instantiate the referenced action schema.%
\end{isamarkuptext}\isamarkuptrue%
\ \ \isacommand{fun}\isamarkupfalse%
\ res{\isacharunderscore}{\kern0pt}inst\ {\isacharcolon}{\kern0pt}{\isacharcolon}{\kern0pt}\ {\isachardoublequoteopen}plan{\isacharunderscore}{\kern0pt}action\ {\isasymrightharpoonup}\ ground{\isacharunderscore}{\kern0pt}action{\isachardoublequoteclose}\ \isakeyword{where}\isanewline
\ \ \ \ {\isachardoublequoteopen}res{\isacharunderscore}{\kern0pt}inst\ {\isacharparenleft}{\kern0pt}Simple{\isacharunderscore}{\kern0pt}Plan{\isacharunderscore}{\kern0pt}Action\ n\ args{\isacharparenright}{\kern0pt}\ {\isacharequal}{\kern0pt}\isanewline
\ \ \ \ \ \ Some\ {\isacharparenleft}{\kern0pt}instantiate{\isacharunderscore}{\kern0pt}action{\isacharunderscore}{\kern0pt}schema\ {\isacharparenleft}{\kern0pt}the\ {\isacharparenleft}{\kern0pt}resolve{\isacharunderscore}{\kern0pt}action{\isacharunderscore}{\kern0pt}schema\ n{\isacharparenright}{\kern0pt}{\isacharparenright}{\kern0pt}\ args{\isacharparenright}{\kern0pt}{\isachardoublequoteclose}\isanewline
\ \ {\isacharbar}{\kern0pt}\ {\isachardoublequoteopen}res{\isacharunderscore}{\kern0pt}inst\ {\isacharparenleft}{\kern0pt}Durative{\isacharunderscore}{\kern0pt}Plan{\isacharunderscore}{\kern0pt}Action\ {\isacharunderscore}{\kern0pt}\ {\isacharunderscore}{\kern0pt}\ {\isacharunderscore}{\kern0pt}{\isacharparenright}{\kern0pt}\ {\isacharequal}{\kern0pt}\ None{\isachardoublequoteclose}\isanewline
\isanewline
\ \ \isanewline
\ \ \isacommand{fun}\isamarkupfalse%
\ res{\isacharunderscore}{\kern0pt}inst{\isacharunderscore}{\kern0pt}snap{\isacharunderscore}{\kern0pt}action\ {\isacharcolon}{\kern0pt}{\isacharcolon}{\kern0pt}\ {\isachardoublequoteopen}plan{\isacharunderscore}{\kern0pt}action\ {\isasymRightarrow}\ temporal{\isacharunderscore}{\kern0pt}annotation\ {\isasymrightharpoonup}\ ground{\isacharunderscore}{\kern0pt}action{\isachardoublequoteclose}\ \isakeyword{where}\isanewline
\ \ \ \ {\isachardoublequoteopen}res{\isacharunderscore}{\kern0pt}inst{\isacharunderscore}{\kern0pt}snap{\isacharunderscore}{\kern0pt}action\ {\isacharparenleft}{\kern0pt}Durative{\isacharunderscore}{\kern0pt}Plan{\isacharunderscore}{\kern0pt}Action\ n\ args\ d{\isacharparenright}{\kern0pt}\ ta\ {\isacharequal}{\kern0pt}\ \isanewline
\ \ \ \ \ \ Some\ {\isacharparenleft}{\kern0pt}inst{\isacharunderscore}{\kern0pt}snap{\isacharunderscore}{\kern0pt}action\ {\isacharparenleft}{\kern0pt}the\ {\isacharparenleft}{\kern0pt}resolve{\isacharunderscore}{\kern0pt}action{\isacharunderscore}{\kern0pt}schema\ n{\isacharparenright}{\kern0pt}{\isacharparenright}{\kern0pt}\ args\ ta{\isacharparenright}{\kern0pt}{\isachardoublequoteclose}\isanewline
\ \ {\isacharbar}{\kern0pt}\ {\isachardoublequoteopen}res{\isacharunderscore}{\kern0pt}inst{\isacharunderscore}{\kern0pt}snap{\isacharunderscore}{\kern0pt}action\ {\isacharparenleft}{\kern0pt}Simple{\isacharunderscore}{\kern0pt}Plan{\isacharunderscore}{\kern0pt}Action\ {\isacharunderscore}{\kern0pt}\ {\isacharunderscore}{\kern0pt}{\isacharparenright}{\kern0pt}\ {\isacharunderscore}{\kern0pt}\ {\isacharequal}{\kern0pt}\ None{\isachardoublequoteclose}%
\begin{isamarkuptext}%
Check whether object has specified type%
\end{isamarkuptext}\isamarkuptrue%
\ \ \isacommand{definition}\isamarkupfalse%
\ {\isachardoublequoteopen}is{\isacharunderscore}{\kern0pt}obj{\isacharunderscore}{\kern0pt}of{\isacharunderscore}{\kern0pt}type\ n\ T\ {\isasymequiv}\ case\ objT\ n\ of\isanewline
\ \ \ \ None\ {\isasymRightarrow}\ False\isanewline
\ \ {\isacharbar}{\kern0pt}\ Some\ oT\ {\isasymRightarrow}\ of{\isacharunderscore}{\kern0pt}type\ oT\ T{\isachardoublequoteclose}%
\begin{isamarkuptext}%
We can also use the generic \isa{is{\isacharunderscore}{\kern0pt}of{\isacharunderscore}{\kern0pt}type} function.%
\end{isamarkuptext}\isamarkuptrue%
\ \ \isacommand{lemma}\isamarkupfalse%
\ is{\isacharunderscore}{\kern0pt}obj{\isacharunderscore}{\kern0pt}of{\isacharunderscore}{\kern0pt}type{\isacharunderscore}{\kern0pt}alt{\isacharcolon}{\kern0pt}\ {\isachardoublequoteopen}is{\isacharunderscore}{\kern0pt}obj{\isacharunderscore}{\kern0pt}of{\isacharunderscore}{\kern0pt}type\ {\isacharequal}{\kern0pt}\ is{\isacharunderscore}{\kern0pt}of{\isacharunderscore}{\kern0pt}type\ objT{\isachardoublequoteclose}\isanewline
%
\isadelimproof
\ \ \ \ %
\endisadelimproof
%
\isatagproof
\isacommand{apply}\isamarkupfalse%
\ {\isacharparenleft}{\kern0pt}intro\ ext{\isacharparenright}{\kern0pt}\isanewline
\ \ \ \ \isacommand{unfolding}\isamarkupfalse%
\ is{\isacharunderscore}{\kern0pt}obj{\isacharunderscore}{\kern0pt}of{\isacharunderscore}{\kern0pt}type{\isacharunderscore}{\kern0pt}def\ is{\isacharunderscore}{\kern0pt}of{\isacharunderscore}{\kern0pt}type{\isacharunderscore}{\kern0pt}def\ \isacommand{by}\isamarkupfalse%
\ auto%
\endisatagproof
{\isafoldproof}%
%
\isadelimproof
%
\endisadelimproof
%
\begin{isamarkuptext}%
HOL encoding of matching an action's formal parameters against an
    argument list.
    The parameters of the action are encoded as a list of \isa{name{\isasymtimes}type} pairs,
    such that we map it to a list of types first. Then, the list
    relator \isa{list{\isacharunderscore}{\kern0pt}all{\isadigit{2}}} checks that arguments and types have the same
    length, and each matching pair of argument and type
    satisfies the predicate \isa{is{\isacharunderscore}{\kern0pt}obj{\isacharunderscore}{\kern0pt}of{\isacharunderscore}{\kern0pt}type}.%
\end{isamarkuptext}\isamarkuptrue%
\ \ \isacommand{definition}\isamarkupfalse%
\ {\isachardoublequoteopen}action{\isacharunderscore}{\kern0pt}params{\isacharunderscore}{\kern0pt}match\ a\ args\isanewline
\ \ \ \ {\isasymequiv}\ list{\isacharunderscore}{\kern0pt}all{\isadigit{2}}\ is{\isacharunderscore}{\kern0pt}obj{\isacharunderscore}{\kern0pt}of{\isacharunderscore}{\kern0pt}type\ args\ {\isacharparenleft}{\kern0pt}map\ snd\ {\isacharparenleft}{\kern0pt}parameters\ a{\isacharparenright}{\kern0pt}{\isacharparenright}{\kern0pt}{\isachardoublequoteclose}%
\begin{isamarkuptext}%
All functions are required to be constant and defined in the initial state. 
  Therefore all duration constraints must be satisfied by the initial state.%
\end{isamarkuptext}\isamarkuptrue%
\ \ \isacommand{definition}\isamarkupfalse%
\ func{\isacharunderscore}{\kern0pt}eval\ {\isacharcolon}{\kern0pt}{\isacharcolon}{\kern0pt}\ {\isachardoublequoteopen}object\ {\isasymrightharpoonup}\ rat{\isachardoublequoteclose}\ \isakeyword{where}\isanewline
\ \ \ \ {\isachardoublequoteopen}func{\isacharunderscore}{\kern0pt}eval\ {\isasymequiv}\ map{\isacharunderscore}{\kern0pt}of\ {\isacharparenleft}{\kern0pt}fold\ {\isacharparenleft}{\kern0pt}{\isasymlambda}a\ fev{\isachardot}{\kern0pt}\ case\ a\ of\ Atom\ {\isacharparenleft}{\kern0pt}eqAtm\ f\ {\isacharparenleft}{\kern0pt}TimeEnt\ t{\isacharparenright}{\kern0pt}{\isacharparenright}{\kern0pt}\ {\isasymRightarrow}\ {\isacharparenleft}{\kern0pt}f{\isacharcomma}{\kern0pt}t{\isacharparenright}{\kern0pt}{\isacharhash}{\kern0pt}fev\ {\isacharbar}{\kern0pt}\ {\isacharunderscore}{\kern0pt}\ {\isasymRightarrow}\ fev{\isacharparenright}{\kern0pt}\ {\isacharparenleft}{\kern0pt}init\ P{\isacharparenright}{\kern0pt}\ {\isacharbrackleft}{\kern0pt}{\isacharbrackright}{\kern0pt}{\isacharparenright}{\kern0pt}{\isachardoublequoteclose}\isanewline
\isanewline
\ \ \isacommand{fun}\isamarkupfalse%
\ duration{\isacharunderscore}{\kern0pt}matches\ {\isacharcolon}{\kern0pt}{\isacharcolon}{\kern0pt}\ {\isachardoublequoteopen}time\ {\isasymRightarrow}\ term\ duration{\isacharunderscore}{\kern0pt}constraint\ {\isasymRightarrow}\ {\isacharparenleft}{\kern0pt}variable\ {\isasymtimes}\ type{\isacharparenright}{\kern0pt}\ list\ {\isasymRightarrow}\ object\ list\ {\isasymRightarrow}\ bool{\isachardoublequoteclose}\ \isakeyword{where}\isanewline
\ \ \ \ {\isachardoublequoteopen}duration{\isacharunderscore}{\kern0pt}matches\ d\ No{\isacharunderscore}{\kern0pt}Const\ params\ args\ {\isasymlongleftrightarrow}\ d\ {\isasymge}\ {\isadigit{0}}{\isachardoublequoteclose}\isanewline
\ \ {\isacharbar}{\kern0pt}\ {\isachardoublequoteopen}duration{\isacharunderscore}{\kern0pt}matches\ d\ {\isacharparenleft}{\kern0pt}Time{\isacharunderscore}{\kern0pt}Const\ op\ d{\isacharprime}{\kern0pt}{\isacharparenright}{\kern0pt}\ params\ args\ {\isasymlongleftrightarrow}\ \isanewline
\ \ \ \ \ \ {\isacharparenleft}{\kern0pt}case\ op\ of\isanewline
\ \ \ \ \ \ \ \ LEQ\ {\isasymRightarrow}\ d\ {\isasymle}\ d{\isacharprime}{\kern0pt}\isanewline
\ \ \ \ \ \ {\isacharbar}{\kern0pt}\ EQ\ {\isasymRightarrow}\ d\ {\isacharequal}{\kern0pt}\ d{\isacharprime}{\kern0pt}\isanewline
\ \ \ \ \ \ {\isacharbar}{\kern0pt}\ GEQ\ {\isasymRightarrow}\ d\ {\isasymge}\ d{\isacharprime}{\kern0pt}{\isacharparenright}{\kern0pt}{\isachardoublequoteclose}\isanewline
\ \ {\isacharbar}{\kern0pt}\ {\isachardoublequoteopen}duration{\isacharunderscore}{\kern0pt}matches\ d\ {\isacharparenleft}{\kern0pt}Func{\isacharunderscore}{\kern0pt}Const\ op\ f\ vs{\isacharparenright}{\kern0pt}\ params\ args\ {\isasymlongleftrightarrow}\isanewline
\ \ \ \ \ \ {\isacharparenleft}{\kern0pt}let\ tsubst\ {\isacharequal}{\kern0pt}\ subst{\isacharunderscore}{\kern0pt}term\ {\isacharparenleft}{\kern0pt}the\ o\ {\isacharparenleft}{\kern0pt}map{\isacharunderscore}{\kern0pt}of\ {\isacharparenleft}{\kern0pt}zip\ {\isacharparenleft}{\kern0pt}map\ fst\ params{\isacharparenright}{\kern0pt}\ args{\isacharparenright}{\kern0pt}{\isacharparenright}{\kern0pt}{\isacharparenright}{\kern0pt}\ in\isanewline
\ \ \ \ \ \ \ case\ func{\isacharunderscore}{\kern0pt}eval\ {\isacharparenleft}{\kern0pt}FuncEnt\ f\ {\isacharparenleft}{\kern0pt}map\ tsubst\ vs{\isacharparenright}{\kern0pt}{\isacharparenright}{\kern0pt}\ of\ \isanewline
\ \ \ \ \ \ \ \ \ \ None\ {\isasymRightarrow}\ False\ \isanewline
\ \ \ \ \ \ \ \ {\isacharbar}{\kern0pt}\ Some\ d{\isacharprime}{\kern0pt}\ {\isasymRightarrow}\ {\isacharparenleft}{\kern0pt}case\ op\ of\isanewline
\ \ \ \ \ \ \ \ \ \ \ \ LEQ\ {\isasymRightarrow}\ d\ {\isasymle}\ d{\isacharprime}{\kern0pt}\isanewline
\ \ \ \ \ \ \ \ \ \ {\isacharbar}{\kern0pt}\ EQ\ {\isasymRightarrow}\ d\ {\isacharequal}{\kern0pt}\ d{\isacharprime}{\kern0pt}\isanewline
\ \ \ \ \ \ \ \ \ \ {\isacharbar}{\kern0pt}\ GEQ\ {\isasymRightarrow}\ d\ {\isasymge}\ d{\isacharprime}{\kern0pt}{\isacharparenright}{\kern0pt}{\isacharparenright}{\kern0pt}{\isachardoublequoteclose}%
\begin{isamarkuptext}%
At this point, we can define well-formedness of a plan action:
    The action must refer to a declared action schema, the arguments must
    be compatible with the formal parameters' types.%
\end{isamarkuptext}\isamarkuptrue%
\ \ \isacommand{fun}\isamarkupfalse%
\ wf{\isacharunderscore}{\kern0pt}plan{\isacharunderscore}{\kern0pt}action\ {\isacharcolon}{\kern0pt}{\isacharcolon}{\kern0pt}\ {\isachardoublequoteopen}plan{\isacharunderscore}{\kern0pt}action\ {\isasymRightarrow}\ bool{\isachardoublequoteclose}\ \isakeyword{where}\isanewline
\ \ \ \ {\isachardoublequoteopen}wf{\isacharunderscore}{\kern0pt}plan{\isacharunderscore}{\kern0pt}action\ {\isacharparenleft}{\kern0pt}Simple{\isacharunderscore}{\kern0pt}Plan{\isacharunderscore}{\kern0pt}Action\ n\ args{\isacharparenright}{\kern0pt}\ {\isacharequal}{\kern0pt}\ {\isacharparenleft}{\kern0pt}\isanewline
\ \ \ \ \ \ case\ resolve{\isacharunderscore}{\kern0pt}action{\isacharunderscore}{\kern0pt}schema\ n\ of\isanewline
\ \ \ \ \ \ \ \ None\ {\isasymRightarrow}\ False\isanewline
\ \ \ \ \ \ {\isacharbar}{\kern0pt}\ Some\ {\isacharparenleft}{\kern0pt}Simple{\isacharunderscore}{\kern0pt}Action{\isacharunderscore}{\kern0pt}Schema\ n\ params\ pre\ eff{\isacharparenright}{\kern0pt}\ {\isasymRightarrow}\ \isanewline
\ \ \ \ \ \ \ \ \ \ action{\isacharunderscore}{\kern0pt}params{\isacharunderscore}{\kern0pt}match\ {\isacharparenleft}{\kern0pt}Simple{\isacharunderscore}{\kern0pt}Action{\isacharunderscore}{\kern0pt}Schema\ n\ params\ pre\ eff{\isacharparenright}{\kern0pt}\ args\isanewline
\ \ \ \ \ \ {\isacharbar}{\kern0pt}\ Some\ {\isacharparenleft}{\kern0pt}Durative{\isacharunderscore}{\kern0pt}Action{\isacharunderscore}{\kern0pt}Schema\ n\ params\ d\ cond\ eff{\isacharparenright}{\kern0pt}\ {\isasymRightarrow}\ False{\isacharparenright}{\kern0pt}{\isachardoublequoteclose}\isanewline
\ \ \ \ {\isacharbar}{\kern0pt}\ {\isachardoublequoteopen}wf{\isacharunderscore}{\kern0pt}plan{\isacharunderscore}{\kern0pt}action\ {\isacharparenleft}{\kern0pt}Durative{\isacharunderscore}{\kern0pt}Plan{\isacharunderscore}{\kern0pt}Action\ n\ args\ d{\isacharparenright}{\kern0pt}\ {\isacharequal}{\kern0pt}\ {\isacharparenleft}{\kern0pt}\isanewline
\ \ \ \ \ \ case\ resolve{\isacharunderscore}{\kern0pt}action{\isacharunderscore}{\kern0pt}schema\ n\ of\isanewline
\ \ \ \ \ \ \ \ None\ {\isasymRightarrow}\ False\isanewline
\ \ \ \ \ \ {\isacharbar}{\kern0pt}\ Some\ {\isacharparenleft}{\kern0pt}Simple{\isacharunderscore}{\kern0pt}Action{\isacharunderscore}{\kern0pt}Schema\ n\ params\ pre\ eff{\isacharparenright}{\kern0pt}\ {\isasymRightarrow}\ False\isanewline
\ \ \ \ \ \ {\isacharbar}{\kern0pt}\ Some\ {\isacharparenleft}{\kern0pt}Durative{\isacharunderscore}{\kern0pt}Action{\isacharunderscore}{\kern0pt}Schema\ n\ params\ dconst\ cond\ eff{\isacharparenright}{\kern0pt}\ {\isasymRightarrow}\ \isanewline
\ \ \ \ \ \ \ \ \ \ action{\isacharunderscore}{\kern0pt}params{\isacharunderscore}{\kern0pt}match\ {\isacharparenleft}{\kern0pt}Durative{\isacharunderscore}{\kern0pt}Action{\isacharunderscore}{\kern0pt}Schema\ n\ params\ dconst\ cond\ eff{\isacharparenright}{\kern0pt}\ args\isanewline
\ \ \ \ \ \ \ \ \ \ {\isasymand}\ duration{\isacharunderscore}{\kern0pt}matches\ d\ dconst\ params\ args\ {\isasymand}\ d\ {\isasymge}\ {\isadigit{0}}{\isacharparenright}{\kern0pt}{\isachardoublequoteclose}%
\begin{isamarkuptext}%
A plan is wellformed if all plan actions are wellformed and all starting 
  time points are greater-equal 0.%
\end{isamarkuptext}\isamarkuptrue%
\ \ \isacommand{definition}\isamarkupfalse%
\ wf{\isacharunderscore}{\kern0pt}plan\ {\isacharcolon}{\kern0pt}{\isacharcolon}{\kern0pt}\ {\isachardoublequoteopen}plan\ {\isasymRightarrow}\ bool{\isachardoublequoteclose}\ \isakeyword{where}\isanewline
\ \ \ \ {\isachardoublequoteopen}wf{\isacharunderscore}{\kern0pt}plan\ {\isasympi}s\ {\isasymlongleftrightarrow}\ {\isacharparenleft}{\kern0pt}{\isasymforall}{\isacharparenleft}{\kern0pt}t{\isacharcomma}{\kern0pt}{\isasympi}{\isacharparenright}{\kern0pt}\ {\isasymin}\ set\ {\isasympi}s{\isachardot}{\kern0pt}\ wf{\isacharunderscore}{\kern0pt}plan{\isacharunderscore}{\kern0pt}action\ {\isasympi}\ {\isasymand}\ t\ {\isasymge}\ {\isadigit{0}}{\isacharparenright}{\kern0pt}{\isachardoublequoteclose}\isanewline
\isanewline
\ \ \isacommand{definition}\isamarkupfalse%
\ is{\isacharunderscore}{\kern0pt}act{\isacharunderscore}{\kern0pt}simple\ {\isacharcolon}{\kern0pt}{\isacharcolon}{\kern0pt}\ {\isachardoublequoteopen}plan{\isacharunderscore}{\kern0pt}action\ {\isasymRightarrow}\ bool{\isachardoublequoteclose}\ \isakeyword{where}\isanewline
\ \ \ \ {\isachardoublequoteopen}is{\isacharunderscore}{\kern0pt}act{\isacharunderscore}{\kern0pt}simple\ {\isasympi}\ {\isasymlongleftrightarrow}\ {\isacharparenleft}{\kern0pt}case\ {\isasympi}\ of\ {\isacharparenleft}{\kern0pt}Simple{\isacharunderscore}{\kern0pt}Plan{\isacharunderscore}{\kern0pt}Action\ {\isacharunderscore}{\kern0pt}\ {\isacharunderscore}{\kern0pt}{\isacharparenright}{\kern0pt}\ {\isasymRightarrow}\ True\ {\isacharbar}{\kern0pt}\ {\isacharunderscore}{\kern0pt}\ {\isasymRightarrow}\ False{\isacharparenright}{\kern0pt}{\isachardoublequoteclose}\isanewline
\isanewline
\ \ \isacommand{definition}\isamarkupfalse%
\ simple{\isacharunderscore}{\kern0pt}acts\ {\isacharcolon}{\kern0pt}{\isacharcolon}{\kern0pt}\ {\isachardoublequoteopen}plan\ {\isasymRightarrow}\ {\isacharparenleft}{\kern0pt}time\ {\isasymtimes}\ plan{\isacharunderscore}{\kern0pt}action{\isacharparenright}{\kern0pt}\ set{\isachardoublequoteclose}\ \isakeyword{where}\ \isanewline
\ \ \ \ {\isachardoublequoteopen}simple{\isacharunderscore}{\kern0pt}acts\ {\isasympi}s\ {\isacharequal}{\kern0pt}\ set\ {\isacharparenleft}{\kern0pt}filter\ {\isacharparenleft}{\kern0pt}is{\isacharunderscore}{\kern0pt}act{\isacharunderscore}{\kern0pt}simple\ o\ snd{\isacharparenright}{\kern0pt}\ {\isasympi}s{\isacharparenright}{\kern0pt}{\isachardoublequoteclose}\isanewline
\isanewline
\ \ \isacommand{lemma}\isamarkupfalse%
\ simple{\isacharunderscore}{\kern0pt}acts{\isacharunderscore}{\kern0pt}subset{\isacharcolon}{\kern0pt}\ {\isachardoublequoteopen}simple{\isacharunderscore}{\kern0pt}acts\ {\isasympi}s\ {\isasymsubseteq}\ set\ {\isasympi}s{\isachardoublequoteclose}\isanewline
%
\isadelimproof
\ \ \ \ %
\endisadelimproof
%
\isatagproof
\isacommand{unfolding}\isamarkupfalse%
\ simple{\isacharunderscore}{\kern0pt}acts{\isacharunderscore}{\kern0pt}def\ \isacommand{by}\isamarkupfalse%
\ auto%
\endisatagproof
{\isafoldproof}%
%
\isadelimproof
\isanewline
%
\endisadelimproof
\isanewline
\ \ \isacommand{definition}\isamarkupfalse%
\ durative{\isacharunderscore}{\kern0pt}acts\ {\isacharcolon}{\kern0pt}{\isacharcolon}{\kern0pt}\ {\isachardoublequoteopen}plan\ {\isasymRightarrow}\ {\isacharparenleft}{\kern0pt}time\ {\isasymtimes}\ plan{\isacharunderscore}{\kern0pt}action{\isacharparenright}{\kern0pt}\ set{\isachardoublequoteclose}\ \isakeyword{where}\ \isanewline
\ \ \ \ {\isachardoublequoteopen}durative{\isacharunderscore}{\kern0pt}acts\ {\isasympi}s\ {\isacharequal}{\kern0pt}\ set\ {\isacharparenleft}{\kern0pt}filter\ {\isacharparenleft}{\kern0pt}{\isacharparenleft}{\kern0pt}HOL{\isachardot}{\kern0pt}Not{\isacharparenright}{\kern0pt}\ o\ is{\isacharunderscore}{\kern0pt}act{\isacharunderscore}{\kern0pt}simple\ o\ snd{\isacharparenright}{\kern0pt}\ {\isasympi}s{\isacharparenright}{\kern0pt}{\isachardoublequoteclose}\isanewline
\isanewline
\ \ \isacommand{lemma}\isamarkupfalse%
\ dur{\isacharunderscore}{\kern0pt}acts{\isacharunderscore}{\kern0pt}subset{\isacharcolon}{\kern0pt}\ {\isachardoublequoteopen}durative{\isacharunderscore}{\kern0pt}acts\ {\isasympi}s\ {\isasymsubseteq}\ set\ {\isasympi}s{\isachardoublequoteclose}\isanewline
%
\isadelimproof
\ \ \ \ %
\endisadelimproof
%
\isatagproof
\isacommand{unfolding}\isamarkupfalse%
\ durative{\isacharunderscore}{\kern0pt}acts{\isacharunderscore}{\kern0pt}def\ \isacommand{by}\isamarkupfalse%
\ auto%
\endisatagproof
{\isafoldproof}%
%
\isadelimproof
%
\endisadelimproof
%
\begin{isamarkuptext}%
Definition of a happening time point. Happening time point specify time points at, 
  which the current state can change. Therefore every start and end of any plan action is 
  a happening time point.%
\end{isamarkuptext}\isamarkuptrue%
\ \ \isacommand{definition}\isamarkupfalse%
\ is{\isacharunderscore}{\kern0pt}htp\ {\isacharcolon}{\kern0pt}{\isacharcolon}{\kern0pt}\ {\isachardoublequoteopen}plan\ {\isasymRightarrow}\ time\ {\isasymRightarrow}\ bool{\isachardoublequoteclose}\ \isakeyword{where}\isanewline
\ \ \ \ {\isachardoublequoteopen}is{\isacharunderscore}{\kern0pt}htp\ {\isasympi}s\ t\isactrlsub i\ {\isasymlongleftrightarrow}\ {\isacharparenleft}{\kern0pt}{\isasymexists}{\isasympi}{\isachardot}{\kern0pt}\ {\isacharparenleft}{\kern0pt}t\isactrlsub i{\isacharcomma}{\kern0pt}{\isasympi}{\isacharparenright}{\kern0pt}\ {\isasymin}\ set\ {\isasympi}s{\isacharparenright}{\kern0pt}\ {\isasymor}\ {\isacharparenleft}{\kern0pt}{\isasymexists}{\isacharparenleft}{\kern0pt}t\isactrlsub {\isasympi}{\isacharcomma}{\kern0pt}{\isasympi}{\isacharparenright}{\kern0pt}\ {\isasymin}\ durative{\isacharunderscore}{\kern0pt}acts\ {\isasympi}s{\isachardot}{\kern0pt}\ t\isactrlsub i\ {\isacharequal}{\kern0pt}\ t\isactrlsub {\isasympi}\ {\isacharplus}{\kern0pt}\ {\isacharparenleft}{\kern0pt}duration\ {\isasympi}{\isacharparenright}{\kern0pt}{\isacharparenright}{\kern0pt}{\isachardoublequoteclose}%
\begin{isamarkuptext}%
Predicate for sequence of happening time points.%
\end{isamarkuptext}\isamarkuptrue%
\ \ \isacommand{definition}\isamarkupfalse%
\ htps{\isacharunderscore}{\kern0pt}seq\ {\isacharcolon}{\kern0pt}{\isacharcolon}{\kern0pt}\ {\isachardoublequoteopen}plan\ {\isasymRightarrow}\ time\ list\ {\isasymRightarrow}\ bool{\isachardoublequoteclose}\ \isakeyword{where}\isanewline
\ \ \ \ {\isachardoublequoteopen}htps{\isacharunderscore}{\kern0pt}seq\ {\isasympi}s\ htps\ {\isasymlongleftrightarrow}\ strict{\isacharunderscore}{\kern0pt}sorted\ htps\ {\isasymand}\ {\isacharparenleft}{\kern0pt}{\isasymforall}t{\isachardot}{\kern0pt}\ t\ {\isasymin}\ set\ htps\ {\isasymlongleftrightarrow}\ is{\isacharunderscore}{\kern0pt}htp\ {\isasympi}s\ t{\isacharparenright}{\kern0pt}{\isachardoublequoteclose}%
\begin{isamarkuptext}%
Predicate for consecutive happening time points: There is no happening 
  time point in between.%
\end{isamarkuptext}\isamarkuptrue%
\ \ \isacommand{definition}\isamarkupfalse%
\ consec{\isacharunderscore}{\kern0pt}htps\ {\isacharcolon}{\kern0pt}{\isacharcolon}{\kern0pt}\ {\isachardoublequoteopen}plan\ {\isasymRightarrow}\ time\ {\isasymRightarrow}\ time\ {\isasymRightarrow}\ bool{\isachardoublequoteclose}\ \isakeyword{where}\isanewline
\ \ \ \ {\isachardoublequoteopen}consec{\isacharunderscore}{\kern0pt}htps\ {\isasympi}s\ t\isactrlsub i\ t\isactrlsub j\ \isanewline
\ \ \ \ \ \ {\isasymlongleftrightarrow}\ {\isacharparenleft}{\kern0pt}t\isactrlsub i\ {\isacharless}{\kern0pt}\ t\isactrlsub j\ {\isasymand}\ is{\isacharunderscore}{\kern0pt}htp\ {\isasympi}s\ t\isactrlsub i\ {\isasymand}\ is{\isacharunderscore}{\kern0pt}htp\ {\isasympi}s\ t\isactrlsub j\ {\isasymand}\ {\isacharparenleft}{\kern0pt}{\isasymforall}t{\isachardot}{\kern0pt}\ is{\isacharunderscore}{\kern0pt}htp\ {\isasympi}s\ t\ {\isasymlongrightarrow}\ t\ {\isasymle}\ t\isactrlsub i\ {\isasymor}\ t\isactrlsub j\ {\isasymle}\ t{\isacharparenright}{\kern0pt}{\isacharparenright}{\kern0pt}{\isachardoublequoteclose}%
\begin{isamarkuptext}%
Predicate to check is a grounded action is a instantiation of a plan action 
  for an induced happening sequence.%
\end{isamarkuptext}\isamarkuptrue%
\ \ \isacommand{fun}\isamarkupfalse%
\ inst{\isacharunderscore}{\kern0pt}of{\isacharunderscore}{\kern0pt}plan{\isacharunderscore}{\kern0pt}action\ {\isacharcolon}{\kern0pt}{\isacharcolon}{\kern0pt}\ {\isachardoublequoteopen}plan\ {\isasymRightarrow}\ time\ {\isasymtimes}\ plan{\isacharunderscore}{\kern0pt}action\ {\isasymRightarrow}\ time\ {\isasymtimes}\ ground{\isacharunderscore}{\kern0pt}action\ {\isasymRightarrow}\ bool{\isachardoublequoteclose}\ \isakeyword{where}\isanewline
\ \ \ \ {\isachardoublequoteopen}inst{\isacharunderscore}{\kern0pt}of{\isacharunderscore}{\kern0pt}plan{\isacharunderscore}{\kern0pt}action\ {\isasympi}s\ {\isacharparenleft}{\kern0pt}t\isactrlsub {\isasympi}{\isacharcomma}{\kern0pt}{\isasympi}{\isacharparenright}{\kern0pt}\ {\isacharparenleft}{\kern0pt}t\isactrlsub a{\isacharcomma}{\kern0pt}a{\isacharparenright}{\kern0pt}\ {\isasymlongleftrightarrow}\ {\isacharparenleft}{\kern0pt}\isanewline
\ \ \ \ \ \ {\isacharparenleft}{\kern0pt}res{\isacharunderscore}{\kern0pt}inst\ {\isasympi}\ {\isacharequal}{\kern0pt}\ Some\ a\ {\isasymand}\ t\isactrlsub {\isasympi}\ {\isacharequal}{\kern0pt}\ t\isactrlsub a{\isacharparenright}{\kern0pt}\isanewline
\ \ \ \ \ \ {\isasymor}\ {\isacharparenleft}{\kern0pt}res{\isacharunderscore}{\kern0pt}inst{\isacharunderscore}{\kern0pt}snap{\isacharunderscore}{\kern0pt}action\ {\isasympi}\ At{\isacharunderscore}{\kern0pt}Start\ {\isacharequal}{\kern0pt}\ Some\ a\ {\isasymand}\ t\isactrlsub {\isasympi}\ {\isacharequal}{\kern0pt}\ t\isactrlsub a{\isacharparenright}{\kern0pt}\isanewline
\ \ \ \ \ \ {\isasymor}\ {\isacharparenleft}{\kern0pt}res{\isacharunderscore}{\kern0pt}inst{\isacharunderscore}{\kern0pt}snap{\isacharunderscore}{\kern0pt}action\ {\isasympi}\ At{\isacharunderscore}{\kern0pt}End\ {\isacharequal}{\kern0pt}\ Some\ a\ {\isasymand}\ t\isactrlsub {\isasympi}\ {\isacharplus}{\kern0pt}\ {\isacharparenleft}{\kern0pt}duration\ {\isasympi}{\isacharparenright}{\kern0pt}\ {\isacharequal}{\kern0pt}\ t\isactrlsub a{\isacharparenright}{\kern0pt}\isanewline
\ \ \ \ \ \ {\isasymor}\ {\isacharparenleft}{\kern0pt}res{\isacharunderscore}{\kern0pt}inst{\isacharunderscore}{\kern0pt}snap{\isacharunderscore}{\kern0pt}action\ {\isasympi}\ Over{\isacharunderscore}{\kern0pt}All\ {\isacharequal}{\kern0pt}\ Some\ a\ {\isasymand}\ \isanewline
\ \ \ \ \ \ \ \ \ \ {\isacharparenleft}{\kern0pt}{\isasymexists}t\isactrlsub i\ t\isactrlsub j{\isachardot}{\kern0pt}\ consec{\isacharunderscore}{\kern0pt}htps\ {\isasympi}s\ t\isactrlsub i\ t\isactrlsub j\ {\isasymand}\ t\isactrlsub {\isasympi}\ {\isasymle}\ t\isactrlsub i\ {\isasymand}\ t\isactrlsub j\ {\isasymle}\ t\isactrlsub {\isasympi}\ {\isacharplus}{\kern0pt}\ {\isacharparenleft}{\kern0pt}duration\ {\isasympi}{\isacharparenright}{\kern0pt}\ {\isasymand}\ t\isactrlsub i\ {\isacharless}{\kern0pt}\ t\isactrlsub a\ {\isasymand}\ t\isactrlsub a\ {\isacharless}{\kern0pt}\ t\isactrlsub j{\isacharparenright}{\kern0pt}{\isacharparenright}{\kern0pt}{\isacharparenright}{\kern0pt}{\isachardoublequoteclose}\isanewline
\isanewline
\ \ \isanewline
\ \ \isanewline
\isanewline
\ \ \isacommand{abbreviation}\isamarkupfalse%
\ happ{\isacharunderscore}{\kern0pt}member\ {\isacharparenleft}{\kern0pt}\isakeyword{infix}\ {\isachardoublequoteopen}{\isasymin}\isactrlsub h{\isachardoublequoteclose}\ {\isadigit{5}}{\isadigit{3}}{\isacharparenright}{\kern0pt}\ \isakeyword{where}\isanewline
\ \ \ \ {\isachardoublequoteopen}ta\ {\isasymin}\isactrlsub h\ hs\ {\isasymequiv}\ {\isacharparenleft}{\kern0pt}case\ ta\ of\ {\isacharparenleft}{\kern0pt}t{\isacharcomma}{\kern0pt}a{\isacharparenright}{\kern0pt}\ {\isasymRightarrow}\ {\isacharparenleft}{\kern0pt}{\isasymexists}as{\isachardot}{\kern0pt}\ {\isacharparenleft}{\kern0pt}t{\isacharcomma}{\kern0pt}as{\isacharparenright}{\kern0pt}\ {\isasymin}\ set\ hs\ {\isasymand}\ a\ {\isasymin}\ set\ as{\isacharparenright}{\kern0pt}{\isacharparenright}{\kern0pt}{\isachardoublequoteclose}%
\begin{isamarkuptext}%
Definition of an induced happening sequence. 

%
\begin{itemize}%
\item for every simple action $(t,\pi)$ there is an instantiation at time point $t$

\item for every durative action $(t,\pi)$

%
\begin{itemize}%
\item there is a ground action at time point $t$, 
which contains all conditions and effect of $\pi$ with time specifier \texttt{(at start ...)}

\item there is a ground action at time point $t + (\texttt{duration}\ \pi)$, 
which contains all conditions and effect of $\pi$ with time specifier \texttt{(at end ...)}

\item in between every consecutive happening-time-points $t_i$ $t_j$ during the execution of $\pi$
there is a ground action at some time point $t'$ (in between $t_i$ \& $t_j$), 
which contains all conditions of $\pi$ with time specifier \texttt{(over all ...)}, the effect of is empty%
\end{itemize}

\item additionally an induced happening sequence is strictly sorted 
and contains no other actions besides the ones specified above.%
\end{itemize}%
\end{isamarkuptext}\isamarkuptrue%
\ \ \isacommand{definition}\isamarkupfalse%
\ ind{\isacharunderscore}{\kern0pt}happ{\isacharunderscore}{\kern0pt}seq\ {\isacharcolon}{\kern0pt}{\isacharcolon}{\kern0pt}\ {\isachardoublequoteopen}plan\ {\isasymRightarrow}\ happening\ list\ {\isasymRightarrow}\ bool{\isachardoublequoteclose}\ \isakeyword{where}\isanewline
\ \ \ \ {\isachardoublequoteopen}ind{\isacharunderscore}{\kern0pt}happ{\isacharunderscore}{\kern0pt}seq\ {\isasympi}s\ hs\ {\isasymlongleftrightarrow}\ {\isacharparenleft}{\kern0pt}\isanewline
\ \ \ \ \ \ strict{\isacharunderscore}{\kern0pt}sorted\ {\isacharparenleft}{\kern0pt}map\ fst\ hs{\isacharparenright}{\kern0pt}\isanewline
\ \ \ \ \ \ {\isasymand}\ {\isacharparenleft}{\kern0pt}{\isasymforall}{\isacharparenleft}{\kern0pt}t\isactrlsub {\isasympi}{\isacharcomma}{\kern0pt}{\isasympi}{\isacharparenright}{\kern0pt}\ {\isasymin}\ simple{\isacharunderscore}{\kern0pt}acts\ {\isasympi}s{\isachardot}{\kern0pt}\ \isanewline
\ \ \ \ \ \ \ \ \ \ let\ g\isactrlsub a\ {\isacharequal}{\kern0pt}\ the\ {\isacharparenleft}{\kern0pt}res{\isacharunderscore}{\kern0pt}inst\ {\isasympi}{\isacharparenright}{\kern0pt}\ in\isanewline
\ \ \ \ \ \ \ \ \ \ {\isasymexists}A{\isachardot}{\kern0pt}\ {\isacharparenleft}{\kern0pt}t\isactrlsub {\isasympi}{\isacharcomma}{\kern0pt}A{\isacharparenright}{\kern0pt}\ {\isasymin}\ set\ hs\ {\isasymand}\ g\isactrlsub a\ {\isasymin}\ set\ A{\isacharparenright}{\kern0pt}\ \isanewline
\ \ \ \ \ \ {\isasymand}\ {\isacharparenleft}{\kern0pt}{\isasymforall}{\isacharparenleft}{\kern0pt}t\isactrlsub {\isasympi}{\isacharcomma}{\kern0pt}{\isasympi}{\isacharparenright}{\kern0pt}\ {\isasymin}\ durative{\isacharunderscore}{\kern0pt}acts\ {\isasympi}s{\isachardot}{\kern0pt}\ \isanewline
\ \ \ \ \ \ \ \ \ \ let\ {\isasympi}\isactrlsub s\isactrlsub t\isactrlsub a\isactrlsub r\isactrlsub t\ {\isacharequal}{\kern0pt}\ the\ {\isacharparenleft}{\kern0pt}res{\isacharunderscore}{\kern0pt}inst{\isacharunderscore}{\kern0pt}snap{\isacharunderscore}{\kern0pt}action\ {\isasympi}\ At{\isacharunderscore}{\kern0pt}Start{\isacharparenright}{\kern0pt}{\isacharsemicolon}{\kern0pt}\isanewline
\ \ \ \ \ \ \ \ \ \ \ \ \ \ {\isasympi}\isactrlsub e\isactrlsub n\isactrlsub d\ {\isacharequal}{\kern0pt}\ the\ {\isacharparenleft}{\kern0pt}res{\isacharunderscore}{\kern0pt}inst{\isacharunderscore}{\kern0pt}snap{\isacharunderscore}{\kern0pt}action\ {\isasympi}\ At{\isacharunderscore}{\kern0pt}End{\isacharparenright}{\kern0pt}{\isacharsemicolon}{\kern0pt}\isanewline
\ \ \ \ \ \ \ \ \ \ \ \ \ \ {\isasympi}\isactrlsub i\isactrlsub n\isactrlsub v\ {\isacharequal}{\kern0pt}\ the\ {\isacharparenleft}{\kern0pt}res{\isacharunderscore}{\kern0pt}inst{\isacharunderscore}{\kern0pt}snap{\isacharunderscore}{\kern0pt}action\ {\isasympi}\ Over{\isacharunderscore}{\kern0pt}All{\isacharparenright}{\kern0pt}\ in\isanewline
\ \ \ \ \ \ \ \ \ \ {\isacharparenleft}{\kern0pt}{\isasymexists}A{\isachardot}{\kern0pt}\ {\isacharparenleft}{\kern0pt}t\isactrlsub {\isasympi}{\isacharcomma}{\kern0pt}A{\isacharparenright}{\kern0pt}\ {\isasymin}\ set\ hs\ {\isasymand}\ {\isasympi}\isactrlsub s\isactrlsub t\isactrlsub a\isactrlsub r\isactrlsub t\ {\isasymin}\ set\ A{\isacharparenright}{\kern0pt}\ {\isasymand}\ \isanewline
\ \ \ \ \ \ \ \ \ \ {\isacharparenleft}{\kern0pt}{\isasymexists}A{\isachardot}{\kern0pt}\ {\isacharparenleft}{\kern0pt}t\isactrlsub {\isasympi}\ {\isacharplus}{\kern0pt}\ {\isacharparenleft}{\kern0pt}duration\ {\isasympi}{\isacharparenright}{\kern0pt}{\isacharcomma}{\kern0pt}A{\isacharparenright}{\kern0pt}\ {\isasymin}\ set\ hs\ {\isasymand}\ {\isasympi}\isactrlsub e\isactrlsub n\isactrlsub d\ {\isasymin}\ set\ A{\isacharparenright}{\kern0pt}\ {\isasymand}\ \isanewline
\ \ \ \ \ \ \ \ \ \ {\isacharparenleft}{\kern0pt}{\isasymforall}t\isactrlsub i\ t\isactrlsub j{\isachardot}{\kern0pt}\ {\isacharparenleft}{\kern0pt}consec{\isacharunderscore}{\kern0pt}htps\ {\isasympi}s\ t\isactrlsub i\ t\isactrlsub j\ {\isasymand}\ t\isactrlsub {\isasympi}\ {\isasymle}\ t\isactrlsub i\ {\isasymand}\ t\isactrlsub j\ {\isasymle}\ t\isactrlsub {\isasympi}\ {\isacharplus}{\kern0pt}\ {\isacharparenleft}{\kern0pt}duration\ {\isasympi}{\isacharparenright}{\kern0pt}{\isacharparenright}{\kern0pt}\ \isanewline
\ \ \ \ \ \ \ \ \ \ \ \ {\isasymlongrightarrow}\ {\isacharparenleft}{\kern0pt}{\isasymexists}{\isacharparenleft}{\kern0pt}t{\isacharprime}{\kern0pt}{\isacharcomma}{\kern0pt}A{\isacharparenright}{\kern0pt}\ {\isasymin}\ set\ hs{\isachardot}{\kern0pt}\ t\isactrlsub i\ {\isacharless}{\kern0pt}\ t{\isacharprime}{\kern0pt}\ {\isasymand}\ t{\isacharprime}{\kern0pt}\ {\isacharless}{\kern0pt}\ t\isactrlsub j\ {\isasymand}\ {\isasympi}\isactrlsub i\isactrlsub n\isactrlsub v\ {\isasymin}\ set\ A{\isacharparenright}{\kern0pt}{\isacharparenright}{\kern0pt}{\isacharparenright}{\kern0pt}\isanewline
\ \ \ \ \ \ {\isasymand}\ {\isacharparenleft}{\kern0pt}{\isasymforall}{\isacharparenleft}{\kern0pt}t\isactrlsub i{\isacharcomma}{\kern0pt}A\isactrlsub i{\isacharparenright}{\kern0pt}\ {\isasymin}\ set\ hs{\isachardot}{\kern0pt}\ A\isactrlsub i\ {\isasymnoteq}\ {\isacharbrackleft}{\kern0pt}{\isacharbrackright}{\kern0pt}\ {\isasymand}\isanewline
\ \ \ \ \ \ \ \ \ \ {\isacharparenleft}{\kern0pt}{\isasymforall}a\ {\isasymin}\ set\ A\isactrlsub i{\isachardot}{\kern0pt}\ {\isasymexists}{\isacharparenleft}{\kern0pt}t\isactrlsub {\isasympi}{\isacharcomma}{\kern0pt}{\isasympi}{\isacharparenright}{\kern0pt}\ {\isasymin}\ set\ {\isasympi}s{\isachardot}{\kern0pt}\ inst{\isacharunderscore}{\kern0pt}of{\isacharunderscore}{\kern0pt}plan{\isacharunderscore}{\kern0pt}action\ {\isasympi}s\ {\isacharparenleft}{\kern0pt}t\isactrlsub {\isasympi}{\isacharcomma}{\kern0pt}{\isasympi}{\isacharparenright}{\kern0pt}\ {\isacharparenleft}{\kern0pt}t\isactrlsub i{\isacharcomma}{\kern0pt}a{\isacharparenright}{\kern0pt}{\isacharparenright}{\kern0pt}{\isacharparenright}{\kern0pt}{\isacharparenright}{\kern0pt}{\isachardoublequoteclose}\isanewline
\isanewline
\isacommand{lemma}\isamarkupfalse%
\ ind{\isacharunderscore}{\kern0pt}happ{\isacharunderscore}{\kern0pt}seq{\isacharunderscore}{\kern0pt}props{\isacharcolon}{\kern0pt}\isanewline
\ \ \isakeyword{assumes}\ {\isachardoublequoteopen}ind{\isacharunderscore}{\kern0pt}happ{\isacharunderscore}{\kern0pt}seq\ {\isasympi}s\ hs{\isachardoublequoteclose}\isanewline
\ \ \isakeyword{shows}\ {\isachardoublequoteopen}strict{\isacharunderscore}{\kern0pt}sorted\ {\isacharparenleft}{\kern0pt}map\ fst\ hs{\isacharparenright}{\kern0pt}{\isachardoublequoteclose}\isanewline
\ \ \ \ \ \ \ \ {\isachardoublequoteopen}{\isacharparenleft}{\kern0pt}t\isactrlsub {\isasympi}{\isacharcomma}{\kern0pt}{\isasympi}{\isacharparenright}{\kern0pt}\ {\isasymin}\ simple{\isacharunderscore}{\kern0pt}acts\ {\isasympi}s\ {\isasymLongrightarrow}\ {\isasymexists}A{\isachardot}{\kern0pt}\ {\isacharparenleft}{\kern0pt}t\isactrlsub {\isasympi}{\isacharcomma}{\kern0pt}A{\isacharparenright}{\kern0pt}\ {\isasymin}\ set\ hs\ {\isasymand}\ the\ {\isacharparenleft}{\kern0pt}res{\isacharunderscore}{\kern0pt}inst\ {\isasympi}{\isacharparenright}{\kern0pt}\ {\isasymin}\ set\ A{\isachardoublequoteclose}\isanewline
\ \ \ \ \ \ \ \ {\isachardoublequoteopen}{\isacharparenleft}{\kern0pt}t\isactrlsub {\isasympi}{\isacharcomma}{\kern0pt}{\isasympi}{\isacharparenright}{\kern0pt}\ {\isasymin}\ durative{\isacharunderscore}{\kern0pt}acts\ {\isasympi}s\ {\isasymLongrightarrow}\ {\isasymexists}A{\isachardot}{\kern0pt}\ {\isacharparenleft}{\kern0pt}t\isactrlsub {\isasympi}{\isacharcomma}{\kern0pt}A{\isacharparenright}{\kern0pt}\ {\isasymin}\ set\ hs\ {\isasymand}\ the\ {\isacharparenleft}{\kern0pt}res{\isacharunderscore}{\kern0pt}inst{\isacharunderscore}{\kern0pt}snap{\isacharunderscore}{\kern0pt}action\ {\isasympi}\ At{\isacharunderscore}{\kern0pt}Start{\isacharparenright}{\kern0pt}\ {\isasymin}\ set\ A{\isachardoublequoteclose}\isanewline
\ \ \ \ \ \ \ \ {\isachardoublequoteopen}{\isacharparenleft}{\kern0pt}t\isactrlsub {\isasympi}{\isacharcomma}{\kern0pt}{\isasympi}{\isacharparenright}{\kern0pt}\ {\isasymin}\ durative{\isacharunderscore}{\kern0pt}acts\ {\isasympi}s\ {\isasymLongrightarrow}\ {\isasymexists}A{\isachardot}{\kern0pt}\ {\isacharparenleft}{\kern0pt}t\isactrlsub {\isasympi}\ {\isacharplus}{\kern0pt}\ {\isacharparenleft}{\kern0pt}duration\ {\isasympi}{\isacharparenright}{\kern0pt}{\isacharcomma}{\kern0pt}A{\isacharparenright}{\kern0pt}\ {\isasymin}\ set\ hs\ {\isasymand}\ the\ {\isacharparenleft}{\kern0pt}res{\isacharunderscore}{\kern0pt}inst{\isacharunderscore}{\kern0pt}snap{\isacharunderscore}{\kern0pt}action\ {\isasympi}\ At{\isacharunderscore}{\kern0pt}End{\isacharparenright}{\kern0pt}\ {\isasymin}\ set\ A{\isachardoublequoteclose}\isanewline
\ \ \ \ \ \ \ \ {\isachardoublequoteopen}{\isasymlbrakk}{\isacharparenleft}{\kern0pt}t\isactrlsub {\isasympi}{\isacharcomma}{\kern0pt}{\isasympi}{\isacharparenright}{\kern0pt}\ {\isasymin}\ durative{\isacharunderscore}{\kern0pt}acts\ {\isasympi}s{\isacharsemicolon}{\kern0pt}\ consec{\isacharunderscore}{\kern0pt}htps\ {\isasympi}s\ t\isactrlsub i\ t\isactrlsub j\ {\isacharsemicolon}{\kern0pt}\ t\isactrlsub {\isasympi}\ {\isasymle}\ t\isactrlsub i\ {\isacharsemicolon}{\kern0pt}\ t\isactrlsub j\ {\isasymle}\ t\isactrlsub {\isasympi}\ {\isacharplus}{\kern0pt}\ {\isacharparenleft}{\kern0pt}duration\ {\isasympi}{\isacharparenright}{\kern0pt}{\isasymrbrakk}\isanewline
\ \ \ \ \ \ \ \ \ \ {\isasymLongrightarrow}\ {\isasymexists}{\isacharparenleft}{\kern0pt}t{\isacharprime}{\kern0pt}{\isacharcomma}{\kern0pt}A{\isacharparenright}{\kern0pt}\ {\isasymin}\ set\ hs{\isachardot}{\kern0pt}\ t\isactrlsub i\ {\isacharless}{\kern0pt}\ t{\isacharprime}{\kern0pt}\ {\isasymand}\ t{\isacharprime}{\kern0pt}\ {\isacharless}{\kern0pt}\ t\isactrlsub j\ {\isasymand}\ the\ {\isacharparenleft}{\kern0pt}res{\isacharunderscore}{\kern0pt}inst{\isacharunderscore}{\kern0pt}snap{\isacharunderscore}{\kern0pt}action\ {\isasympi}\ Over{\isacharunderscore}{\kern0pt}All{\isacharparenright}{\kern0pt}\ {\isasymin}\ set\ A{\isachardoublequoteclose}\isanewline
\ \ \ \ \ \ \ \ {\isachardoublequoteopen}{\isacharparenleft}{\kern0pt}t\isactrlsub i{\isacharcomma}{\kern0pt}A\isactrlsub i{\isacharparenright}{\kern0pt}\ {\isasymin}\ set\ hs\ {\isasymLongrightarrow}\ A\isactrlsub i\ {\isasymnoteq}\ {\isacharbrackleft}{\kern0pt}{\isacharbrackright}{\kern0pt}{\isachardoublequoteclose}\isanewline
\ \ \ \ \ \ \ \ {\isachardoublequoteopen}{\isasymlbrakk}{\isacharparenleft}{\kern0pt}t\isactrlsub i{\isacharcomma}{\kern0pt}A\isactrlsub i{\isacharparenright}{\kern0pt}\ {\isasymin}\ set\ hs\ {\isacharsemicolon}{\kern0pt}\ a\ {\isasymin}\ set\ A\isactrlsub i{\isasymrbrakk}\ {\isasymLongrightarrow}\ {\isasymexists}{\isacharparenleft}{\kern0pt}t\isactrlsub {\isasympi}{\isacharcomma}{\kern0pt}{\isasympi}{\isacharparenright}{\kern0pt}\ {\isasymin}\ set\ {\isasympi}s{\isachardot}{\kern0pt}\ inst{\isacharunderscore}{\kern0pt}of{\isacharunderscore}{\kern0pt}plan{\isacharunderscore}{\kern0pt}action\ {\isasympi}s\ {\isacharparenleft}{\kern0pt}t\isactrlsub {\isasympi}{\isacharcomma}{\kern0pt}{\isasympi}{\isacharparenright}{\kern0pt}\ {\isacharparenleft}{\kern0pt}t\isactrlsub i{\isacharcomma}{\kern0pt}a{\isacharparenright}{\kern0pt}{\isachardoublequoteclose}\isanewline
%
\isadelimproof
\ \ %
\endisadelimproof
%
\isatagproof
\isacommand{using}\isamarkupfalse%
\ assms\isanewline
\ \ \isacommand{by}\isamarkupfalse%
{\isacharparenleft}{\kern0pt}auto\ simp{\isacharcolon}{\kern0pt}\ ind{\isacharunderscore}{\kern0pt}happ{\isacharunderscore}{\kern0pt}seq{\isacharunderscore}{\kern0pt}def{\isacharparenright}{\kern0pt}%
\endisatagproof
{\isafoldproof}%
%
\isadelimproof
%
\endisadelimproof
%
\begin{isamarkuptext}%
A sequence of plan actions form a path, if they are well-formed and
    their instantiations form a path.%
\end{isamarkuptext}\isamarkuptrue%
\ \ \isacommand{definition}\isamarkupfalse%
\ plan{\isacharunderscore}{\kern0pt}happ{\isacharunderscore}{\kern0pt}path\ {\isacharcolon}{\kern0pt}{\isacharcolon}{\kern0pt}\ {\isachardoublequoteopen}world{\isacharunderscore}{\kern0pt}model\ {\isasymRightarrow}\ plan\ {\isasymRightarrow}\ world{\isacharunderscore}{\kern0pt}model\ {\isasymRightarrow}\ bool{\isachardoublequoteclose}\ \isakeyword{where}\isanewline
\ \ \ \ {\isachardoublequoteopen}plan{\isacharunderscore}{\kern0pt}happ{\isacharunderscore}{\kern0pt}path\ M\ {\isasympi}s\ M{\isacharprime}{\kern0pt}\ {\isasymlongleftrightarrow}\ {\isacharparenleft}{\kern0pt}{\isasymexists}hs{\isachardot}{\kern0pt}\ ind{\isacharunderscore}{\kern0pt}happ{\isacharunderscore}{\kern0pt}seq\ {\isasympi}s\ hs\ {\isasymand}\ valid{\isacharunderscore}{\kern0pt}happ{\isacharunderscore}{\kern0pt}seq\ M\ hs\ M{\isacharprime}{\kern0pt}{\isacharparenright}{\kern0pt}{\isachardoublequoteclose}%
\begin{isamarkuptext}%
A plan is valid wrt.\ a given initial model, if it forms a path to a
    goal model%
\end{isamarkuptext}\isamarkuptrue%
\ \ \isacommand{definition}\isamarkupfalse%
\ valid{\isacharunderscore}{\kern0pt}plan{\isacharunderscore}{\kern0pt}from{\isadigit{2}}\ {\isacharcolon}{\kern0pt}{\isacharcolon}{\kern0pt}\ {\isachardoublequoteopen}world{\isacharunderscore}{\kern0pt}model\ {\isasymRightarrow}\ plan\ {\isasymRightarrow}\ bool{\isachardoublequoteclose}\ \isakeyword{where}\isanewline
\ \ \ \ {\isachardoublequoteopen}valid{\isacharunderscore}{\kern0pt}plan{\isacharunderscore}{\kern0pt}from{\isadigit{2}}\ M\ {\isasympi}s\ {\isasymlongleftrightarrow}\ {\isacharparenleft}{\kern0pt}wf{\isacharunderscore}{\kern0pt}plan\ {\isasympi}s\ {\isasymand}\ {\isacharparenleft}{\kern0pt}{\isasymexists}M{\isacharprime}{\kern0pt}{\isachardot}{\kern0pt}\ plan{\isacharunderscore}{\kern0pt}happ{\isacharunderscore}{\kern0pt}path\ M\ {\isasympi}s\ M{\isacharprime}{\kern0pt}\ {\isasymand}\ M{\isacharprime}{\kern0pt}\ \isactrlsup c{\isasymTTurnstile}\isactrlsub {\isacharequal}{\kern0pt}\ {\isacharparenleft}{\kern0pt}goal\ P{\isacharparenright}{\kern0pt}{\isacharparenright}{\kern0pt}{\isacharparenright}{\kern0pt}{\isachardoublequoteclose}%
\begin{isamarkuptext}%
Finally, a plan is valid if it is valid wrt.\ the initial world
    model \isa{I}%
\end{isamarkuptext}\isamarkuptrue%
\ \ \isacommand{definition}\isamarkupfalse%
\ valid{\isacharunderscore}{\kern0pt}plan{\isadigit{2}}\ {\isacharcolon}{\kern0pt}{\isacharcolon}{\kern0pt}\ {\isachardoublequoteopen}plan\ {\isasymRightarrow}\ bool{\isachardoublequoteclose}\isanewline
\ \ \ \ \isakeyword{where}\ {\isachardoublequoteopen}valid{\isacharunderscore}{\kern0pt}plan{\isadigit{2}}\ {\isasymequiv}\ valid{\isacharunderscore}{\kern0pt}plan{\isacharunderscore}{\kern0pt}from{\isadigit{2}}\ I{\isachardoublequoteclose}%
\begin{isamarkuptext}%
Concise definition used in paper:%
\end{isamarkuptext}\isamarkuptrue%
\ \ \isacommand{lemma}\isamarkupfalse%
\ {\isachardoublequoteopen}valid{\isacharunderscore}{\kern0pt}plan{\isadigit{2}}\ {\isasympi}s\ {\isasymequiv}\ wf{\isacharunderscore}{\kern0pt}plan\ {\isasympi}s\ {\isasymand}\ {\isacharparenleft}{\kern0pt}{\isasymexists}M{\isacharprime}{\kern0pt}{\isachardot}{\kern0pt}\ plan{\isacharunderscore}{\kern0pt}happ{\isacharunderscore}{\kern0pt}path\ I\ {\isasympi}s\ M{\isacharprime}{\kern0pt}\ {\isasymand}\ M{\isacharprime}{\kern0pt}\ \isactrlsup c{\isasymTTurnstile}\isactrlsub {\isacharequal}{\kern0pt}\ {\isacharparenleft}{\kern0pt}goal\ P{\isacharparenright}{\kern0pt}{\isacharparenright}{\kern0pt}{\isachardoublequoteclose}\isanewline
%
\isadelimproof
\ \ \ \ %
\endisadelimproof
%
\isatagproof
\isacommand{unfolding}\isamarkupfalse%
\ valid{\isacharunderscore}{\kern0pt}plan{\isadigit{2}}{\isacharunderscore}{\kern0pt}def\ valid{\isacharunderscore}{\kern0pt}plan{\isacharunderscore}{\kern0pt}from{\isadigit{2}}{\isacharunderscore}{\kern0pt}def\ \isacommand{by}\isamarkupfalse%
\ auto%
\endisatagproof
{\isafoldproof}%
%
\isadelimproof
\isanewline
%
\endisadelimproof
\isanewline
\isacommand{end}\isamarkupfalse%
\ %
\isamarkupcmt{Context of \isa{ast{\isacharunderscore}{\kern0pt}problem}%
}%
\isadelimdocument
%
\endisadelimdocument
%
\isatagdocument
%
\isamarkupsubsection{Preservation of Well-Formedness%
}
\isamarkuptrue%
%
\isamarkupsubsubsection{Well-Formed Action Instances%
}
\isamarkuptrue%
%
\endisatagdocument
{\isafolddocument}%
%
\isadelimdocument
%
\endisadelimdocument
%
\begin{isamarkuptext}%
The goal of this section is to establish that well-formedness of
  world models is preserved by execution of well-formed plan actions.%
\end{isamarkuptext}\isamarkuptrue%
\isacommand{context}\isamarkupfalse%
\ ast{\isacharunderscore}{\kern0pt}problem\ \isakeyword{begin}%
\begin{isamarkuptext}%
As plan actions are executed by first instantiating them, and then
    executing the action instance, it is natural to define a well-formedness
    concept for action instances.%
\end{isamarkuptext}\isamarkuptrue%
\ \ \isacommand{fun}\isamarkupfalse%
\ wf{\isacharunderscore}{\kern0pt}ground{\isacharunderscore}{\kern0pt}action\ {\isacharcolon}{\kern0pt}{\isacharcolon}{\kern0pt}\ {\isachardoublequoteopen}ground{\isacharunderscore}{\kern0pt}action\ {\isasymRightarrow}\ bool{\isachardoublequoteclose}\ \isakeyword{where}\isanewline
\ \ \ \ {\isachardoublequoteopen}wf{\isacharunderscore}{\kern0pt}ground{\isacharunderscore}{\kern0pt}action\ {\isacharparenleft}{\kern0pt}Ground{\isacharunderscore}{\kern0pt}Action\ pre\ eff{\isacharparenright}{\kern0pt}\ {\isasymlongleftrightarrow}\ {\isacharparenleft}{\kern0pt}wf{\isacharunderscore}{\kern0pt}fmla\ objT\ pre\ {\isasymand}\ wf{\isacharunderscore}{\kern0pt}effect\ objT\ eff{\isacharparenright}{\kern0pt}{\isachardoublequoteclose}%
\begin{isamarkuptext}%
We first prove that instantiating a well-formed action schema will yield
    a well-formed action instance.

    We begin with some auxiliary lemmas before the actual theorem.%
\end{isamarkuptext}\isamarkuptrue%
\ \ \isacommand{lemma}\isamarkupfalse%
\ {\isacharparenleft}{\kern0pt}\isakeyword{in}\ ast{\isacharunderscore}{\kern0pt}domain{\isacharparenright}{\kern0pt}\ of{\isacharunderscore}{\kern0pt}type{\isacharunderscore}{\kern0pt}refl{\isacharbrackleft}{\kern0pt}simp{\isacharcomma}{\kern0pt}\ intro{\isacharbang}{\kern0pt}{\isacharbrackright}{\kern0pt}{\isacharcolon}{\kern0pt}\ {\isachardoublequoteopen}of{\isacharunderscore}{\kern0pt}type\ T\ T{\isachardoublequoteclose}\isanewline
%
\isadelimproof
\ \ \ \ %
\endisadelimproof
%
\isatagproof
\isacommand{unfolding}\isamarkupfalse%
\ of{\isacharunderscore}{\kern0pt}type{\isacharunderscore}{\kern0pt}def\ \isacommand{by}\isamarkupfalse%
\ auto%
\endisatagproof
{\isafoldproof}%
%
\isadelimproof
\isanewline
%
\endisadelimproof
\isanewline
\ \ \isacommand{lemma}\isamarkupfalse%
\ {\isacharparenleft}{\kern0pt}\isakeyword{in}\ ast{\isacharunderscore}{\kern0pt}domain{\isacharparenright}{\kern0pt}\ of{\isacharunderscore}{\kern0pt}type{\isacharunderscore}{\kern0pt}trans{\isacharbrackleft}{\kern0pt}trans{\isacharbrackright}{\kern0pt}{\isacharcolon}{\kern0pt}\isanewline
\ \ \ \ {\isachardoublequoteopen}of{\isacharunderscore}{\kern0pt}type\ T{\isadigit{1}}\ T{\isadigit{2}}\ {\isasymLongrightarrow}\ of{\isacharunderscore}{\kern0pt}type\ T{\isadigit{2}}\ T{\isadigit{3}}\ {\isasymLongrightarrow}\ of{\isacharunderscore}{\kern0pt}type\ T{\isadigit{1}}\ T{\isadigit{3}}{\isachardoublequoteclose}\isanewline
%
\isadelimproof
\ \ \ \ %
\endisadelimproof
%
\isatagproof
\isacommand{unfolding}\isamarkupfalse%
\ of{\isacharunderscore}{\kern0pt}type{\isacharunderscore}{\kern0pt}def\isanewline
\ \ \ \ \isacommand{by}\isamarkupfalse%
\ clarsimp\ {\isacharparenleft}{\kern0pt}meson\ in{\isacharunderscore}{\kern0pt}mono\ rtrancl{\isacharunderscore}{\kern0pt}image{\isacharunderscore}{\kern0pt}unfold{\isacharunderscore}{\kern0pt}right\ rtrancl{\isacharunderscore}{\kern0pt}reachable{\isacharunderscore}{\kern0pt}induct{\isacharparenright}{\kern0pt}%
\endisatagproof
{\isafoldproof}%
%
\isadelimproof
\ \isanewline
%
\endisadelimproof
\isanewline
\ \ \isacommand{lemma}\isamarkupfalse%
\ is{\isacharunderscore}{\kern0pt}of{\isacharunderscore}{\kern0pt}type{\isacharunderscore}{\kern0pt}map{\isacharunderscore}{\kern0pt}ofE{\isacharcolon}{\kern0pt}\isanewline
\ \ \ \ \isakeyword{assumes}\ {\isachardoublequoteopen}is{\isacharunderscore}{\kern0pt}of{\isacharunderscore}{\kern0pt}type\ {\isacharparenleft}{\kern0pt}map{\isacharunderscore}{\kern0pt}of\ params{\isacharparenright}{\kern0pt}\ x\ T{\isachardoublequoteclose}\isanewline
\ \ \ \ \isakeyword{obtains}\ i\ xT\ \isakeyword{where}\ {\isachardoublequoteopen}i{\isacharless}{\kern0pt}length\ params{\isachardoublequoteclose}\ {\isachardoublequoteopen}params{\isacharbang}{\kern0pt}i\ {\isacharequal}{\kern0pt}\ {\isacharparenleft}{\kern0pt}x{\isacharcomma}{\kern0pt}xT{\isacharparenright}{\kern0pt}{\isachardoublequoteclose}\ {\isachardoublequoteopen}of{\isacharunderscore}{\kern0pt}type\ xT\ T{\isachardoublequoteclose}\isanewline
%
\isadelimproof
\ \ \ \ %
\endisadelimproof
%
\isatagproof
\isacommand{using}\isamarkupfalse%
\ assms\isanewline
\ \ \ \ \isacommand{unfolding}\isamarkupfalse%
\ is{\isacharunderscore}{\kern0pt}of{\isacharunderscore}{\kern0pt}type{\isacharunderscore}{\kern0pt}def\isanewline
\ \ \ \ \isacommand{by}\isamarkupfalse%
\ {\isacharparenleft}{\kern0pt}auto\ split{\isacharcolon}{\kern0pt}\ option{\isachardot}{\kern0pt}splits\ dest{\isacharbang}{\kern0pt}{\isacharcolon}{\kern0pt}\ map{\isacharunderscore}{\kern0pt}of{\isacharunderscore}{\kern0pt}SomeD\ simp{\isacharcolon}{\kern0pt}\ in{\isacharunderscore}{\kern0pt}set{\isacharunderscore}{\kern0pt}conv{\isacharunderscore}{\kern0pt}nth{\isacharparenright}{\kern0pt}%
\endisatagproof
{\isafoldproof}%
%
\isadelimproof
\isanewline
%
\endisadelimproof
\isanewline
\ \ \isacommand{lemma}\isamarkupfalse%
\ wf{\isacharunderscore}{\kern0pt}atom{\isacharunderscore}{\kern0pt}mono{\isacharcolon}{\kern0pt}\isanewline
\ \ \ \ \isakeyword{assumes}\ SS{\isacharcolon}{\kern0pt}\ {\isachardoublequoteopen}tys\ {\isasymsubseteq}\isactrlsub m\ tys{\isacharprime}{\kern0pt}{\isachardoublequoteclose}\isanewline
\ \ \ \ \isakeyword{assumes}\ WF{\isacharcolon}{\kern0pt}\ {\isachardoublequoteopen}wf{\isacharunderscore}{\kern0pt}atom\ tys\ a{\isachardoublequoteclose}\isanewline
\ \ \ \ \isakeyword{shows}\ {\isachardoublequoteopen}wf{\isacharunderscore}{\kern0pt}atom\ tys{\isacharprime}{\kern0pt}\ a{\isachardoublequoteclose}\isanewline
%
\isadelimproof
\ \ %
\endisadelimproof
%
\isatagproof
\isacommand{proof}\isamarkupfalse%
\ {\isacharminus}{\kern0pt}\isanewline
\ \ \ \ \isacommand{have}\isamarkupfalse%
\ {\isachardoublequoteopen}list{\isacharunderscore}{\kern0pt}all{\isadigit{2}}\ {\isacharparenleft}{\kern0pt}is{\isacharunderscore}{\kern0pt}of{\isacharunderscore}{\kern0pt}type\ tys{\isacharprime}{\kern0pt}{\isacharparenright}{\kern0pt}\ xs\ Ts{\isachardoublequoteclose}\ \isakeyword{if}\ {\isachardoublequoteopen}list{\isacharunderscore}{\kern0pt}all{\isadigit{2}}\ {\isacharparenleft}{\kern0pt}is{\isacharunderscore}{\kern0pt}of{\isacharunderscore}{\kern0pt}type\ tys{\isacharparenright}{\kern0pt}\ xs\ Ts{\isachardoublequoteclose}\ \isakeyword{for}\ xs\ Ts\isanewline
\ \ \ \ \ \ \isacommand{using}\isamarkupfalse%
\ that\isanewline
\ \ \ \ \ \ \isacommand{apply}\isamarkupfalse%
\ induction\isanewline
\ \ \ \ \ \ \isacommand{by}\isamarkupfalse%
\ {\isacharparenleft}{\kern0pt}auto\ simp{\isacharcolon}{\kern0pt}\ is{\isacharunderscore}{\kern0pt}of{\isacharunderscore}{\kern0pt}type{\isacharunderscore}{\kern0pt}def\ split{\isacharcolon}{\kern0pt}\ option{\isachardot}{\kern0pt}splits\ dest{\isacharcolon}{\kern0pt}\ map{\isacharunderscore}{\kern0pt}leD{\isacharbrackleft}{\kern0pt}OF\ SS{\isacharbrackright}{\kern0pt}{\isacharparenright}{\kern0pt}\isanewline
\ \ \ \ \isacommand{with}\isamarkupfalse%
\ WF\ \isacommand{show}\isamarkupfalse%
\ {\isacharquery}{\kern0pt}thesis\isanewline
\ \ \ \ \ \ \isacommand{by}\isamarkupfalse%
\ {\isacharparenleft}{\kern0pt}cases\ a{\isacharparenright}{\kern0pt}\ {\isacharparenleft}{\kern0pt}auto\ split{\isacharcolon}{\kern0pt}\ option{\isachardot}{\kern0pt}splits\ dest{\isacharcolon}{\kern0pt}\ map{\isacharunderscore}{\kern0pt}leD{\isacharbrackleft}{\kern0pt}OF\ SS{\isacharbrackright}{\kern0pt}{\isacharparenright}{\kern0pt}\isanewline
\ \ \isacommand{qed}\isamarkupfalse%
%
\endisatagproof
{\isafoldproof}%
%
\isadelimproof
\isanewline
%
\endisadelimproof
\isanewline
\ \ \isacommand{lemma}\isamarkupfalse%
\ wf{\isacharunderscore}{\kern0pt}fmla{\isacharunderscore}{\kern0pt}atom{\isacharunderscore}{\kern0pt}mono{\isacharcolon}{\kern0pt}\isanewline
\ \ \ \ \isakeyword{assumes}\ SS{\isacharcolon}{\kern0pt}\ {\isachardoublequoteopen}tys\ {\isasymsubseteq}\isactrlsub m\ tys{\isacharprime}{\kern0pt}{\isachardoublequoteclose}\isanewline
\ \ \ \ \isakeyword{assumes}\ WF{\isacharcolon}{\kern0pt}\ {\isachardoublequoteopen}wf{\isacharunderscore}{\kern0pt}fmla{\isacharunderscore}{\kern0pt}atom\ tys\ a{\isachardoublequoteclose}\isanewline
\ \ \ \ \isakeyword{shows}\ {\isachardoublequoteopen}wf{\isacharunderscore}{\kern0pt}fmla{\isacharunderscore}{\kern0pt}atom\ tys{\isacharprime}{\kern0pt}\ a{\isachardoublequoteclose}\isanewline
%
\isadelimproof
\ \ %
\endisadelimproof
%
\isatagproof
\isacommand{proof}\isamarkupfalse%
\ {\isacharminus}{\kern0pt}\isanewline
\ \ \ \ \isacommand{have}\isamarkupfalse%
\ {\isachardoublequoteopen}list{\isacharunderscore}{\kern0pt}all{\isadigit{2}}\ {\isacharparenleft}{\kern0pt}is{\isacharunderscore}{\kern0pt}of{\isacharunderscore}{\kern0pt}type\ tys{\isacharprime}{\kern0pt}{\isacharparenright}{\kern0pt}\ xs\ Ts{\isachardoublequoteclose}\ \isakeyword{if}\ {\isachardoublequoteopen}list{\isacharunderscore}{\kern0pt}all{\isadigit{2}}\ {\isacharparenleft}{\kern0pt}is{\isacharunderscore}{\kern0pt}of{\isacharunderscore}{\kern0pt}type\ tys{\isacharparenright}{\kern0pt}\ xs\ Ts{\isachardoublequoteclose}\ \isakeyword{for}\ xs\ Ts\isanewline
\ \ \ \ \ \ \isacommand{using}\isamarkupfalse%
\ that\isanewline
\ \ \ \ \ \ \isacommand{apply}\isamarkupfalse%
\ induction\isanewline
\ \ \ \ \ \ \isacommand{by}\isamarkupfalse%
\ {\isacharparenleft}{\kern0pt}auto\ simp{\isacharcolon}{\kern0pt}\ is{\isacharunderscore}{\kern0pt}of{\isacharunderscore}{\kern0pt}type{\isacharunderscore}{\kern0pt}def\ split{\isacharcolon}{\kern0pt}\ option{\isachardot}{\kern0pt}splits\ dest{\isacharcolon}{\kern0pt}\ map{\isacharunderscore}{\kern0pt}leD{\isacharbrackleft}{\kern0pt}OF\ SS{\isacharbrackright}{\kern0pt}{\isacharparenright}{\kern0pt}\isanewline
\ \ \ \ \isacommand{with}\isamarkupfalse%
\ WF\ \isacommand{show}\isamarkupfalse%
\ {\isacharquery}{\kern0pt}thesis\isanewline
\ \ \ \ \ \ \isacommand{by}\isamarkupfalse%
\ {\isacharparenleft}{\kern0pt}cases\ a\ rule{\isacharcolon}{\kern0pt}\ wf{\isacharunderscore}{\kern0pt}fmla{\isacharunderscore}{\kern0pt}atom{\isachardot}{\kern0pt}cases{\isacharparenright}{\kern0pt}\ {\isacharparenleft}{\kern0pt}auto\ split{\isacharcolon}{\kern0pt}\ option{\isachardot}{\kern0pt}splits\ dest{\isacharcolon}{\kern0pt}\ map{\isacharunderscore}{\kern0pt}leD{\isacharbrackleft}{\kern0pt}OF\ SS{\isacharbrackright}{\kern0pt}{\isacharparenright}{\kern0pt}\isanewline
\ \ \isacommand{qed}\isamarkupfalse%
%
\endisatagproof
{\isafoldproof}%
%
\isadelimproof
\isanewline
%
\endisadelimproof
\isanewline
\isanewline
\ \ \isacommand{lemma}\isamarkupfalse%
\ constT{\isacharunderscore}{\kern0pt}ss{\isacharunderscore}{\kern0pt}objT{\isacharcolon}{\kern0pt}\ {\isachardoublequoteopen}constT\ {\isasymsubseteq}\isactrlsub m\ objT{\isachardoublequoteclose}\isanewline
%
\isadelimproof
\ \ \ \ %
\endisadelimproof
%
\isatagproof
\isacommand{unfolding}\isamarkupfalse%
\ constT{\isacharunderscore}{\kern0pt}def\ objT{\isacharunderscore}{\kern0pt}def\isanewline
\ \ \ \ \isacommand{apply}\isamarkupfalse%
\ rule\isanewline
\ \ \ \ \isacommand{by}\isamarkupfalse%
\ {\isacharparenleft}{\kern0pt}auto\ simp{\isacharcolon}{\kern0pt}\ map{\isacharunderscore}{\kern0pt}add{\isacharunderscore}{\kern0pt}def\ split{\isacharcolon}{\kern0pt}\ option{\isachardot}{\kern0pt}split{\isacharparenright}{\kern0pt}%
\endisatagproof
{\isafoldproof}%
%
\isadelimproof
\isanewline
%
\endisadelimproof
\isanewline
\ \ \isacommand{lemma}\isamarkupfalse%
\ wf{\isacharunderscore}{\kern0pt}atom{\isacharunderscore}{\kern0pt}constT{\isacharunderscore}{\kern0pt}imp{\isacharunderscore}{\kern0pt}objT{\isacharcolon}{\kern0pt}\ {\isachardoublequoteopen}wf{\isacharunderscore}{\kern0pt}atom\ {\isacharparenleft}{\kern0pt}ty{\isacharunderscore}{\kern0pt}term\ Q\ constT{\isacharparenright}{\kern0pt}\ a\ {\isasymLongrightarrow}\ wf{\isacharunderscore}{\kern0pt}atom\ {\isacharparenleft}{\kern0pt}ty{\isacharunderscore}{\kern0pt}term\ Q\ objT{\isacharparenright}{\kern0pt}\ a{\isachardoublequoteclose}\isanewline
%
\isadelimproof
\ \ \ \ %
\endisadelimproof
%
\isatagproof
\isacommand{apply}\isamarkupfalse%
\ {\isacharparenleft}{\kern0pt}erule\ wf{\isacharunderscore}{\kern0pt}atom{\isacharunderscore}{\kern0pt}mono{\isacharbrackleft}{\kern0pt}rotated{\isacharbrackright}{\kern0pt}{\isacharparenright}{\kern0pt}\isanewline
\ \ \ \ \isacommand{apply}\isamarkupfalse%
\ {\isacharparenleft}{\kern0pt}rule\ ty{\isacharunderscore}{\kern0pt}term{\isacharunderscore}{\kern0pt}mono{\isacharparenright}{\kern0pt}\isanewline
\ \ \ \ \isacommand{by}\isamarkupfalse%
\ {\isacharparenleft}{\kern0pt}simp{\isacharunderscore}{\kern0pt}all\ add{\isacharcolon}{\kern0pt}\ constT{\isacharunderscore}{\kern0pt}ss{\isacharunderscore}{\kern0pt}objT{\isacharparenright}{\kern0pt}%
\endisatagproof
{\isafoldproof}%
%
\isadelimproof
\isanewline
%
\endisadelimproof
\isanewline
\ \ \isacommand{lemma}\isamarkupfalse%
\ wf{\isacharunderscore}{\kern0pt}fmla{\isacharunderscore}{\kern0pt}atom{\isacharunderscore}{\kern0pt}constT{\isacharunderscore}{\kern0pt}imp{\isacharunderscore}{\kern0pt}objT{\isacharcolon}{\kern0pt}\ {\isachardoublequoteopen}wf{\isacharunderscore}{\kern0pt}fmla{\isacharunderscore}{\kern0pt}atom\ {\isacharparenleft}{\kern0pt}ty{\isacharunderscore}{\kern0pt}term\ Q\ constT{\isacharparenright}{\kern0pt}\ a\ {\isasymLongrightarrow}\ wf{\isacharunderscore}{\kern0pt}fmla{\isacharunderscore}{\kern0pt}atom\ {\isacharparenleft}{\kern0pt}ty{\isacharunderscore}{\kern0pt}term\ Q\ objT{\isacharparenright}{\kern0pt}\ a{\isachardoublequoteclose}\isanewline
%
\isadelimproof
\ \ \ \ %
\endisadelimproof
%
\isatagproof
\isacommand{apply}\isamarkupfalse%
\ {\isacharparenleft}{\kern0pt}erule\ wf{\isacharunderscore}{\kern0pt}fmla{\isacharunderscore}{\kern0pt}atom{\isacharunderscore}{\kern0pt}mono{\isacharbrackleft}{\kern0pt}rotated{\isacharbrackright}{\kern0pt}{\isacharparenright}{\kern0pt}\isanewline
\ \ \ \ \isacommand{apply}\isamarkupfalse%
\ {\isacharparenleft}{\kern0pt}rule\ ty{\isacharunderscore}{\kern0pt}term{\isacharunderscore}{\kern0pt}mono{\isacharparenright}{\kern0pt}\isanewline
\ \ \ \ \isacommand{by}\isamarkupfalse%
\ {\isacharparenleft}{\kern0pt}simp{\isacharunderscore}{\kern0pt}all\ add{\isacharcolon}{\kern0pt}\ constT{\isacharunderscore}{\kern0pt}ss{\isacharunderscore}{\kern0pt}objT{\isacharparenright}{\kern0pt}%
\endisatagproof
{\isafoldproof}%
%
\isadelimproof
\isanewline
%
\endisadelimproof
\isanewline
\ \ \isacommand{context}\isamarkupfalse%
\isanewline
\ \ \ \ \isakeyword{fixes}\ Q\ \isakeyword{and}\ f\ {\isacharcolon}{\kern0pt}{\isacharcolon}{\kern0pt}\ {\isachardoublequoteopen}variable\ {\isasymRightarrow}\ object{\isachardoublequoteclose}\isanewline
\ \ \ \ \isakeyword{assumes}\ INST{\isacharcolon}{\kern0pt}\ {\isachardoublequoteopen}is{\isacharunderscore}{\kern0pt}of{\isacharunderscore}{\kern0pt}type\ Q\ x\ T\ {\isasymLongrightarrow}\ is{\isacharunderscore}{\kern0pt}of{\isacharunderscore}{\kern0pt}type\ objT\ {\isacharparenleft}{\kern0pt}f\ x{\isacharparenright}{\kern0pt}\ T{\isachardoublequoteclose}\isanewline
\ \ \isakeyword{begin}\isanewline
\isanewline
\ \ \ \ \isacommand{lemma}\isamarkupfalse%
\ is{\isacharunderscore}{\kern0pt}of{\isacharunderscore}{\kern0pt}type{\isacharunderscore}{\kern0pt}var{\isacharunderscore}{\kern0pt}conv{\isacharcolon}{\kern0pt}\ {\isachardoublequoteopen}is{\isacharunderscore}{\kern0pt}of{\isacharunderscore}{\kern0pt}type\ {\isacharparenleft}{\kern0pt}ty{\isacharunderscore}{\kern0pt}term\ Q\ objT{\isacharparenright}{\kern0pt}\ {\isacharparenleft}{\kern0pt}term{\isachardot}{\kern0pt}VAR\ x{\isacharparenright}{\kern0pt}\ T\ \ {\isasymlongleftrightarrow}\ is{\isacharunderscore}{\kern0pt}of{\isacharunderscore}{\kern0pt}type\ Q\ x\ T{\isachardoublequoteclose}\isanewline
%
\isadelimproof
\ \ \ \ \ \ %
\endisadelimproof
%
\isatagproof
\isacommand{unfolding}\isamarkupfalse%
\ is{\isacharunderscore}{\kern0pt}of{\isacharunderscore}{\kern0pt}type{\isacharunderscore}{\kern0pt}def\ \isacommand{by}\isamarkupfalse%
\ {\isacharparenleft}{\kern0pt}auto{\isacharparenright}{\kern0pt}%
\endisatagproof
{\isafoldproof}%
%
\isadelimproof
\isanewline
%
\endisadelimproof
\isanewline
\ \ \ \ \isacommand{lemma}\isamarkupfalse%
\ is{\isacharunderscore}{\kern0pt}of{\isacharunderscore}{\kern0pt}type{\isacharunderscore}{\kern0pt}const{\isacharunderscore}{\kern0pt}conv{\isacharcolon}{\kern0pt}\ {\isachardoublequoteopen}is{\isacharunderscore}{\kern0pt}of{\isacharunderscore}{\kern0pt}type\ {\isacharparenleft}{\kern0pt}ty{\isacharunderscore}{\kern0pt}term\ Q\ objT{\isacharparenright}{\kern0pt}\ {\isacharparenleft}{\kern0pt}term{\isachardot}{\kern0pt}CONST\ x{\isacharparenright}{\kern0pt}\ T\ \ {\isasymlongleftrightarrow}\ is{\isacharunderscore}{\kern0pt}of{\isacharunderscore}{\kern0pt}type\ objT\ x\ T{\isachardoublequoteclose}\isanewline
%
\isadelimproof
\ \ \ \ \ \ %
\endisadelimproof
%
\isatagproof
\isacommand{unfolding}\isamarkupfalse%
\ is{\isacharunderscore}{\kern0pt}of{\isacharunderscore}{\kern0pt}type{\isacharunderscore}{\kern0pt}def\isanewline
\ \ \ \ \ \ \isacommand{by}\isamarkupfalse%
\ {\isacharparenleft}{\kern0pt}auto\ split{\isacharcolon}{\kern0pt}\ option{\isachardot}{\kern0pt}split{\isacharparenright}{\kern0pt}%
\endisatagproof
{\isafoldproof}%
%
\isadelimproof
\isanewline
%
\endisadelimproof
\isanewline
\ \ \ \ \isacommand{lemma}\isamarkupfalse%
\ INST{\isacharprime}{\kern0pt}{\isacharcolon}{\kern0pt}\ {\isachardoublequoteopen}is{\isacharunderscore}{\kern0pt}of{\isacharunderscore}{\kern0pt}type\ {\isacharparenleft}{\kern0pt}ty{\isacharunderscore}{\kern0pt}term\ Q\ objT{\isacharparenright}{\kern0pt}\ x\ T\ {\isasymLongrightarrow}\ is{\isacharunderscore}{\kern0pt}of{\isacharunderscore}{\kern0pt}type\ objT\ {\isacharparenleft}{\kern0pt}subst{\isacharunderscore}{\kern0pt}term\ f\ x{\isacharparenright}{\kern0pt}\ T{\isachardoublequoteclose}\isanewline
%
\isadelimproof
\ \ \ \ \ \ %
\endisadelimproof
%
\isatagproof
\isacommand{apply}\isamarkupfalse%
\ {\isacharparenleft}{\kern0pt}cases\ x{\isacharparenright}{\kern0pt}\ \isacommand{using}\isamarkupfalse%
\ INST\ \isacommand{apply}\isamarkupfalse%
\ {\isacharparenleft}{\kern0pt}auto\ simp{\isacharcolon}{\kern0pt}\ is{\isacharunderscore}{\kern0pt}of{\isacharunderscore}{\kern0pt}type{\isacharunderscore}{\kern0pt}var{\isacharunderscore}{\kern0pt}conv\ is{\isacharunderscore}{\kern0pt}of{\isacharunderscore}{\kern0pt}type{\isacharunderscore}{\kern0pt}const{\isacharunderscore}{\kern0pt}conv{\isacharparenright}{\kern0pt}\isanewline
\ \ \ \ \ \ \isacommand{done}\isamarkupfalse%
%
\endisatagproof
{\isafoldproof}%
%
\isadelimproof
\isanewline
%
\endisadelimproof
\isanewline
\isanewline
\ \ \ \ \isacommand{lemma}\isamarkupfalse%
\ wf{\isacharunderscore}{\kern0pt}inst{\isacharunderscore}{\kern0pt}eq{\isacharunderscore}{\kern0pt}aux{\isacharcolon}{\kern0pt}\ {\isachardoublequoteopen}Q\ x\ {\isacharequal}{\kern0pt}\ Some\ T\ {\isasymLongrightarrow}\ objT\ {\isacharparenleft}{\kern0pt}f\ x{\isacharparenright}{\kern0pt}\ {\isasymnoteq}\ None{\isachardoublequoteclose}\isanewline
%
\isadelimproof
\ \ \ \ \ \ %
\endisadelimproof
%
\isatagproof
\isacommand{using}\isamarkupfalse%
\ INST{\isacharbrackleft}{\kern0pt}of\ x\ T{\isacharbrackright}{\kern0pt}\ \isacommand{unfolding}\isamarkupfalse%
\ is{\isacharunderscore}{\kern0pt}of{\isacharunderscore}{\kern0pt}type{\isacharunderscore}{\kern0pt}def\isanewline
\ \ \ \ \ \ \isacommand{by}\isamarkupfalse%
\ {\isacharparenleft}{\kern0pt}auto\ split{\isacharcolon}{\kern0pt}\ option{\isachardot}{\kern0pt}splits{\isacharparenright}{\kern0pt}%
\endisatagproof
{\isafoldproof}%
%
\isadelimproof
\isanewline
%
\endisadelimproof
\isanewline
\ \ \ \ \isacommand{lemma}\isamarkupfalse%
\ wf{\isacharunderscore}{\kern0pt}inst{\isacharunderscore}{\kern0pt}eq{\isacharunderscore}{\kern0pt}aux{\isacharprime}{\kern0pt}{\isacharcolon}{\kern0pt}\ {\isachardoublequoteopen}ty{\isacharunderscore}{\kern0pt}term\ Q\ objT\ x\ {\isacharequal}{\kern0pt}\ Some\ T\ {\isasymLongrightarrow}\ objT\ {\isacharparenleft}{\kern0pt}subst{\isacharunderscore}{\kern0pt}term\ f\ x{\isacharparenright}{\kern0pt}\ {\isasymnoteq}\ None{\isachardoublequoteclose}\isanewline
%
\isadelimproof
\ \ \ \ \ \ %
\endisadelimproof
%
\isatagproof
\isacommand{by}\isamarkupfalse%
\ {\isacharparenleft}{\kern0pt}cases\ x{\isacharparenright}{\kern0pt}\ {\isacharparenleft}{\kern0pt}auto\ simp{\isacharcolon}{\kern0pt}\ wf{\isacharunderscore}{\kern0pt}inst{\isacharunderscore}{\kern0pt}eq{\isacharunderscore}{\kern0pt}aux{\isacharparenright}{\kern0pt}%
\endisatagproof
{\isafoldproof}%
%
\isadelimproof
\isanewline
%
\endisadelimproof
\isanewline
\isanewline
\ \ \ \ \isacommand{lemma}\isamarkupfalse%
\ wf{\isacharunderscore}{\kern0pt}inst{\isacharunderscore}{\kern0pt}atom{\isacharcolon}{\kern0pt}\isanewline
\ \ \ \ \ \ \isakeyword{assumes}\ {\isachardoublequoteopen}wf{\isacharunderscore}{\kern0pt}atom\ {\isacharparenleft}{\kern0pt}ty{\isacharunderscore}{\kern0pt}term\ Q\ constT{\isacharparenright}{\kern0pt}\ a{\isachardoublequoteclose}\isanewline
\ \ \ \ \ \ \isakeyword{shows}\ {\isachardoublequoteopen}wf{\isacharunderscore}{\kern0pt}atom\ objT\ {\isacharparenleft}{\kern0pt}map{\isacharunderscore}{\kern0pt}atom\ {\isacharparenleft}{\kern0pt}subst{\isacharunderscore}{\kern0pt}term\ f{\isacharparenright}{\kern0pt}\ a{\isacharparenright}{\kern0pt}{\isachardoublequoteclose}\isanewline
%
\isadelimproof
\ \ \ \ %
\endisadelimproof
%
\isatagproof
\isacommand{proof}\isamarkupfalse%
\ {\isacharminus}{\kern0pt}\isanewline
\ \ \ \ \ \ \isacommand{have}\isamarkupfalse%
\ X{\isadigit{1}}{\isacharcolon}{\kern0pt}\ {\isachardoublequoteopen}list{\isacharunderscore}{\kern0pt}all{\isadigit{2}}\ {\isacharparenleft}{\kern0pt}is{\isacharunderscore}{\kern0pt}of{\isacharunderscore}{\kern0pt}type\ objT{\isacharparenright}{\kern0pt}\ {\isacharparenleft}{\kern0pt}map\ {\isacharparenleft}{\kern0pt}subst{\isacharunderscore}{\kern0pt}term\ f{\isacharparenright}{\kern0pt}\ xs{\isacharparenright}{\kern0pt}\ Ts{\isachardoublequoteclose}\ \isakeyword{if}\isanewline
\ \ \ \ \ \ \ \ {\isachardoublequoteopen}list{\isacharunderscore}{\kern0pt}all{\isadigit{2}}\ {\isacharparenleft}{\kern0pt}is{\isacharunderscore}{\kern0pt}of{\isacharunderscore}{\kern0pt}type\ {\isacharparenleft}{\kern0pt}ty{\isacharunderscore}{\kern0pt}term\ Q\ objT{\isacharparenright}{\kern0pt}{\isacharparenright}{\kern0pt}\ xs\ Ts{\isachardoublequoteclose}\ \isakeyword{for}\ xs\ Ts\isanewline
\ \ \ \ \ \ \ \ \isacommand{using}\isamarkupfalse%
\ that\isanewline
\ \ \ \ \ \ \ \ \isacommand{apply}\isamarkupfalse%
\ induction\isanewline
\ \ \ \ \ \ \ \ \isacommand{using}\isamarkupfalse%
\ INST{\isacharprime}{\kern0pt}\isanewline
\ \ \ \ \ \ \ \ \isacommand{by}\isamarkupfalse%
\ auto\isanewline
\ \ \ \ \ \ \isacommand{then}\isamarkupfalse%
\ \isacommand{show}\isamarkupfalse%
\ {\isacharquery}{\kern0pt}thesis\isanewline
\ \ \ \ \ \ \ \ \isacommand{using}\isamarkupfalse%
\ assms{\isacharbrackleft}{\kern0pt}THEN\ wf{\isacharunderscore}{\kern0pt}atom{\isacharunderscore}{\kern0pt}constT{\isacharunderscore}{\kern0pt}imp{\isacharunderscore}{\kern0pt}objT{\isacharbrackright}{\kern0pt}\ wf{\isacharunderscore}{\kern0pt}inst{\isacharunderscore}{\kern0pt}eq{\isacharunderscore}{\kern0pt}aux{\isacharprime}{\kern0pt}\isanewline
\ \ \ \ \ \ \ \ \isacommand{by}\isamarkupfalse%
\ {\isacharparenleft}{\kern0pt}cases\ a{\isacharsemicolon}{\kern0pt}\ auto\ split{\isacharcolon}{\kern0pt}\ option{\isachardot}{\kern0pt}splits{\isacharparenright}{\kern0pt}\isanewline
\isanewline
\ \ \ \ \isacommand{qed}\isamarkupfalse%
%
\endisatagproof
{\isafoldproof}%
%
\isadelimproof
\isanewline
%
\endisadelimproof
\isanewline
\ \ \ \ \isacommand{lemma}\isamarkupfalse%
\ wf{\isacharunderscore}{\kern0pt}inst{\isacharunderscore}{\kern0pt}formula{\isacharunderscore}{\kern0pt}atom{\isacharcolon}{\kern0pt}\isanewline
\ \ \ \ \ \ \isakeyword{assumes}\ {\isachardoublequoteopen}wf{\isacharunderscore}{\kern0pt}fmla{\isacharunderscore}{\kern0pt}atom\ {\isacharparenleft}{\kern0pt}ty{\isacharunderscore}{\kern0pt}term\ Q\ constT{\isacharparenright}{\kern0pt}\ a{\isachardoublequoteclose}\isanewline
\ \ \ \ \ \ \isakeyword{shows}\ {\isachardoublequoteopen}wf{\isacharunderscore}{\kern0pt}fmla{\isacharunderscore}{\kern0pt}atom\ objT\ {\isacharparenleft}{\kern0pt}{\isacharparenleft}{\kern0pt}map{\isacharunderscore}{\kern0pt}formula\ o\ map{\isacharunderscore}{\kern0pt}atom\ o\ subst{\isacharunderscore}{\kern0pt}term{\isacharparenright}{\kern0pt}\ f\ a{\isacharparenright}{\kern0pt}{\isachardoublequoteclose}\isanewline
%
\isadelimproof
\ \ \ \ \ \ %
\endisadelimproof
%
\isatagproof
\isacommand{using}\isamarkupfalse%
\ assms{\isacharbrackleft}{\kern0pt}THEN\ wf{\isacharunderscore}{\kern0pt}fmla{\isacharunderscore}{\kern0pt}atom{\isacharunderscore}{\kern0pt}constT{\isacharunderscore}{\kern0pt}imp{\isacharunderscore}{\kern0pt}objT{\isacharbrackright}{\kern0pt}\ wf{\isacharunderscore}{\kern0pt}inst{\isacharunderscore}{\kern0pt}atom\isanewline
\ \ \ \ \ \ \isacommand{apply}\isamarkupfalse%
\ {\isacharparenleft}{\kern0pt}cases\ a\ rule{\isacharcolon}{\kern0pt}\ wf{\isacharunderscore}{\kern0pt}fmla{\isacharunderscore}{\kern0pt}atom{\isachardot}{\kern0pt}cases{\isacharsemicolon}{\kern0pt}\ auto\ split{\isacharcolon}{\kern0pt}\ option{\isachardot}{\kern0pt}splits{\isacharparenright}{\kern0pt}\isanewline
\ \ \ \ \ \ \isacommand{by}\isamarkupfalse%
\ {\isacharparenleft}{\kern0pt}simp\ add{\isacharcolon}{\kern0pt}\ INST{\isacharprime}{\kern0pt}\ list{\isachardot}{\kern0pt}rel{\isacharunderscore}{\kern0pt}map{\isacharparenleft}{\kern0pt}{\isadigit{1}}{\isacharparenright}{\kern0pt}\ list{\isacharunderscore}{\kern0pt}all{\isadigit{2}}{\isacharunderscore}{\kern0pt}mono{\isacharparenright}{\kern0pt}%
\endisatagproof
{\isafoldproof}%
%
\isadelimproof
\isanewline
%
\endisadelimproof
\isanewline
\ \ \ \ \isacommand{lemma}\isamarkupfalse%
\ wf{\isacharunderscore}{\kern0pt}inst{\isacharunderscore}{\kern0pt}effect{\isacharcolon}{\kern0pt}\isanewline
\ \ \ \ \ \ \isakeyword{assumes}\ {\isachardoublequoteopen}wf{\isacharunderscore}{\kern0pt}effect\ {\isacharparenleft}{\kern0pt}ty{\isacharunderscore}{\kern0pt}term\ Q\ constT{\isacharparenright}{\kern0pt}\ {\isasymphi}{\isachardoublequoteclose}\isanewline
\ \ \ \ \ \ \isakeyword{shows}\ {\isachardoublequoteopen}wf{\isacharunderscore}{\kern0pt}effect\ objT\ {\isacharparenleft}{\kern0pt}{\isacharparenleft}{\kern0pt}map{\isacharunderscore}{\kern0pt}ast{\isacharunderscore}{\kern0pt}effect\ o\ subst{\isacharunderscore}{\kern0pt}term{\isacharparenright}{\kern0pt}\ f\ {\isasymphi}{\isacharparenright}{\kern0pt}{\isachardoublequoteclose}\isanewline
%
\isadelimproof
\ \ \ \ \ \ %
\endisadelimproof
%
\isatagproof
\isacommand{using}\isamarkupfalse%
\ assms\isanewline
\ \ \ \ \ \ \isacommand{proof}\isamarkupfalse%
\ {\isacharparenleft}{\kern0pt}induction\ {\isasymphi}{\isacharparenright}{\kern0pt}\isanewline
\ \ \ \ \ \ \ \ \isacommand{case}\isamarkupfalse%
\ {\isacharparenleft}{\kern0pt}Effect\ x{\isadigit{1}}a\ x{\isadigit{2}}a{\isacharparenright}{\kern0pt}\isanewline
\ \ \ \ \ \ \ \ \isacommand{then}\isamarkupfalse%
\ \isacommand{show}\isamarkupfalse%
\ {\isacharquery}{\kern0pt}case\ \isacommand{using}\isamarkupfalse%
\ wf{\isacharunderscore}{\kern0pt}inst{\isacharunderscore}{\kern0pt}formula{\isacharunderscore}{\kern0pt}atom\ \isacommand{by}\isamarkupfalse%
\ auto\isanewline
\ \ \ \ \ \ \isacommand{qed}\isamarkupfalse%
%
\endisatagproof
{\isafoldproof}%
%
\isadelimproof
\isanewline
%
\endisadelimproof
\isanewline
\ \ \ \ \isacommand{lemma}\isamarkupfalse%
\ wf{\isacharunderscore}{\kern0pt}inst{\isacharunderscore}{\kern0pt}formula{\isacharcolon}{\kern0pt}\isanewline
\ \ \ \ \ \ \isakeyword{assumes}\ {\isachardoublequoteopen}wf{\isacharunderscore}{\kern0pt}fmla\ {\isacharparenleft}{\kern0pt}ty{\isacharunderscore}{\kern0pt}term\ Q\ constT{\isacharparenright}{\kern0pt}\ {\isasymphi}{\isachardoublequoteclose}\isanewline
\ \ \ \ \ \ \isakeyword{shows}\ {\isachardoublequoteopen}wf{\isacharunderscore}{\kern0pt}fmla\ objT\ {\isacharparenleft}{\kern0pt}{\isacharparenleft}{\kern0pt}map{\isacharunderscore}{\kern0pt}formula\ o\ map{\isacharunderscore}{\kern0pt}atom\ o\ subst{\isacharunderscore}{\kern0pt}term{\isacharparenright}{\kern0pt}\ f\ {\isasymphi}{\isacharparenright}{\kern0pt}{\isachardoublequoteclose}\isanewline
%
\isadelimproof
\ \ \ \ \ \ %
\endisadelimproof
%
\isatagproof
\isacommand{using}\isamarkupfalse%
\ assms\isanewline
\ \ \ \ \ \ \isacommand{by}\isamarkupfalse%
\ {\isacharparenleft}{\kern0pt}induction\ {\isasymphi}{\isacharparenright}{\kern0pt}\ {\isacharparenleft}{\kern0pt}auto\ simp{\isacharcolon}{\kern0pt}\ wf{\isacharunderscore}{\kern0pt}inst{\isacharunderscore}{\kern0pt}atom\ dest{\isacharcolon}{\kern0pt}\ wf{\isacharunderscore}{\kern0pt}inst{\isacharunderscore}{\kern0pt}eq{\isacharunderscore}{\kern0pt}aux{\isacharparenright}{\kern0pt}%
\endisatagproof
{\isafoldproof}%
%
\isadelimproof
\isanewline
%
\endisadelimproof
\isanewline
\ \ \ \ \isacommand{lemma}\isamarkupfalse%
\ wf{\isacharunderscore}{\kern0pt}BigAnd{\isacharunderscore}{\kern0pt}aux{\isacharcolon}{\kern0pt}\ {\isachardoublequoteopen}{\isasymforall}x\ {\isasymin}\ set\ xs{\isachardot}{\kern0pt}\ wf{\isacharunderscore}{\kern0pt}fmla\ tyt\ x\ {\isasymLongrightarrow}\ wf{\isacharunderscore}{\kern0pt}fmla\ tyt\ {\isacharparenleft}{\kern0pt}BigAnd\ xs{\isacharparenright}{\kern0pt}{\isachardoublequoteclose}\isanewline
%
\isadelimproof
\ \ \ \ \ \ %
\endisadelimproof
%
\isatagproof
\isacommand{by}\isamarkupfalse%
\ {\isacharparenleft}{\kern0pt}induction\ xs{\isacharparenright}{\kern0pt}\ auto%
\endisatagproof
{\isafoldproof}%
%
\isadelimproof
\isanewline
%
\endisadelimproof
\isanewline
\ \ \ \ \isacommand{lemma}\isamarkupfalse%
\ wf{\isacharunderscore}{\kern0pt}BigAnd{\isacharunderscore}{\kern0pt}filter{\isacharunderscore}{\kern0pt}time{\isacharunderscore}{\kern0pt}spec{\isacharcolon}{\kern0pt}\isanewline
\ \ \ \ \ \ \isakeyword{fixes}\ n\ params\ d\ cond\ eff\isanewline
\ \ \ \ \ \ \isakeyword{defines}\ {\isachardoublequoteopen}a\ {\isasymequiv}\ Durative{\isacharunderscore}{\kern0pt}Action{\isacharunderscore}{\kern0pt}Schema\ n\ params\ d\ cond\ eff{\isachardoublequoteclose}\ \ \isanewline
\ \ \ \ \ \ \isakeyword{assumes}\ {\isachardoublequoteopen}wf{\isacharunderscore}{\kern0pt}action{\isacharunderscore}{\kern0pt}schema\ a{\isachardoublequoteclose}\isanewline
\ \ \ \ \ \ \isakeyword{shows}\ {\isachardoublequoteopen}wf{\isacharunderscore}{\kern0pt}fmla\ {\isacharparenleft}{\kern0pt}ty{\isacharunderscore}{\kern0pt}term\ {\isacharparenleft}{\kern0pt}map{\isacharunderscore}{\kern0pt}of\ params{\isacharparenright}{\kern0pt}\ constT{\isacharparenright}{\kern0pt}\ {\isacharparenleft}{\kern0pt}BigAnd\ {\isacharparenleft}{\kern0pt}filter{\isacharunderscore}{\kern0pt}time{\isacharunderscore}{\kern0pt}spec\ ta\ cond{\isacharparenright}{\kern0pt}{\isacharparenright}{\kern0pt}{\isachardoublequoteclose}\isanewline
%
\isadelimproof
\ \ \ \ \ \ %
\endisadelimproof
%
\isatagproof
\isacommand{using}\isamarkupfalse%
\ assms\ filter{\isacharunderscore}{\kern0pt}time{\isacharunderscore}{\kern0pt}spec{\isacharunderscore}{\kern0pt}correct\ wf{\isacharunderscore}{\kern0pt}BigAnd{\isacharunderscore}{\kern0pt}aux\isanewline
\ \ \ \ \ \ \isacommand{by}\isamarkupfalse%
\ {\isacharparenleft}{\kern0pt}fastforce\ simp{\isacharcolon}{\kern0pt}\ Let{\isacharunderscore}{\kern0pt}def{\isacharparenright}{\kern0pt}%
\endisatagproof
{\isafoldproof}%
%
\isadelimproof
\isanewline
%
\endisadelimproof
\isanewline
\ \ \ \ \isacommand{lemma}\isamarkupfalse%
\ wf{\isacharunderscore}{\kern0pt}conjunct{\isacharunderscore}{\kern0pt}effects{\isacharunderscore}{\kern0pt}aux{\isacharcolon}{\kern0pt}\ {\isachardoublequoteopen}{\isasymforall}x\ {\isasymin}\ set\ xs{\isachardot}{\kern0pt}\ wf{\isacharunderscore}{\kern0pt}effect\ tyt\ x\ {\isasymLongrightarrow}\ wf{\isacharunderscore}{\kern0pt}effect\ tyt\ {\isacharparenleft}{\kern0pt}{\isasymAnd}\isactrlsub e\isactrlsub f\isactrlsub f\ xs{\isacharparenright}{\kern0pt}{\isachardoublequoteclose}\isanewline
%
\isadelimproof
\ \ \ \ \ \ %
\endisadelimproof
%
\isatagproof
\isacommand{unfolding}\isamarkupfalse%
\ conjunct{\isacharunderscore}{\kern0pt}effects{\isacharunderscore}{\kern0pt}def\isanewline
\ \ \ \ \ \ \isacommand{apply}\isamarkupfalse%
\ {\isacharparenleft}{\kern0pt}induction\ xs{\isacharparenright}{\kern0pt}\ \isanewline
\ \ \ \ \ \ \isacommand{apply}\isamarkupfalse%
\ auto\isanewline
\ \ \ \ \ \ \isacommand{apply}\isamarkupfalse%
\ {\isacharparenleft}{\kern0pt}metis\ ast{\isacharunderscore}{\kern0pt}effect{\isachardot}{\kern0pt}collapse\ wf{\isacharunderscore}{\kern0pt}effect{\isachardot}{\kern0pt}simps{\isacharparenright}{\kern0pt}{\isacharplus}{\kern0pt}\isanewline
\ \ \ \ \ \ \isacommand{done}\isamarkupfalse%
%
\endisatagproof
{\isafoldproof}%
%
\isadelimproof
%
\endisadelimproof
\ \isanewline
\isanewline
\ \ \ \ \isacommand{lemma}\isamarkupfalse%
\ wf{\isacharunderscore}{\kern0pt}conj{\isacharunderscore}{\kern0pt}eff{\isacharunderscore}{\kern0pt}filter{\isacharunderscore}{\kern0pt}time{\isacharunderscore}{\kern0pt}spec{\isacharcolon}{\kern0pt}\isanewline
\ \ \ \ \ \ \isakeyword{fixes}\ n\ params\ d\ cond\ eff\isanewline
\ \ \ \ \ \ \isakeyword{defines}\ {\isachardoublequoteopen}a\ {\isasymequiv}\ Durative{\isacharunderscore}{\kern0pt}Action{\isacharunderscore}{\kern0pt}Schema\ n\ params\ d\ cond\ eff{\isachardoublequoteclose}\ \ \isanewline
\ \ \ \ \ \ \isakeyword{assumes}\ {\isachardoublequoteopen}wf{\isacharunderscore}{\kern0pt}action{\isacharunderscore}{\kern0pt}schema\ a{\isachardoublequoteclose}\isanewline
\ \ \ \ \ \ \isakeyword{shows}\ {\isachardoublequoteopen}wf{\isacharunderscore}{\kern0pt}effect\ {\isacharparenleft}{\kern0pt}ty{\isacharunderscore}{\kern0pt}term\ {\isacharparenleft}{\kern0pt}map{\isacharunderscore}{\kern0pt}of\ params{\isacharparenright}{\kern0pt}\ constT{\isacharparenright}{\kern0pt}\ {\isacharparenleft}{\kern0pt}{\isasymAnd}\isactrlsub e\isactrlsub f\isactrlsub f\ {\isacharparenleft}{\kern0pt}filter{\isacharunderscore}{\kern0pt}time{\isacharunderscore}{\kern0pt}spec\ ta\ eff{\isacharparenright}{\kern0pt}{\isacharparenright}{\kern0pt}{\isachardoublequoteclose}\isanewline
%
\isadelimproof
\ \ \ \ \ \ %
\endisadelimproof
%
\isatagproof
\isacommand{using}\isamarkupfalse%
\ assms\ filter{\isacharunderscore}{\kern0pt}time{\isacharunderscore}{\kern0pt}spec{\isacharunderscore}{\kern0pt}correct\ wf{\isacharunderscore}{\kern0pt}conjunct{\isacharunderscore}{\kern0pt}effects{\isacharunderscore}{\kern0pt}aux\isanewline
\ \ \ \ \ \ \isacommand{by}\isamarkupfalse%
\ {\isacharparenleft}{\kern0pt}fastforce\ simp{\isacharcolon}{\kern0pt}\ Let{\isacharunderscore}{\kern0pt}def{\isacharparenright}{\kern0pt}%
\endisatagproof
{\isafoldproof}%
%
\isadelimproof
\isanewline
%
\endisadelimproof
\isanewline
\isacommand{end}\isamarkupfalse%
%
\begin{isamarkuptext}%
Next two theorems show that, instantiating a well-formed action schema 
  with compatible arguments will yield a well-formed action instance.%
\end{isamarkuptext}\isamarkuptrue%
\ \ \isacommand{theorem}\isamarkupfalse%
\ wf{\isacharunderscore}{\kern0pt}inst{\isacharunderscore}{\kern0pt}action{\isacharunderscore}{\kern0pt}schema{\isacharcolon}{\kern0pt}\isanewline
\ \ \ \ \isakeyword{fixes}\ n\ params\ pre\ eff\isanewline
\ \ \ \ \isakeyword{defines}\ {\isachardoublequoteopen}a\ {\isasymequiv}\ Simple{\isacharunderscore}{\kern0pt}Action{\isacharunderscore}{\kern0pt}Schema\ n\ params\ pre\ eff{\isachardoublequoteclose}\ \ \isanewline
\ \ \ \ \isakeyword{assumes}\ {\isachardoublequoteopen}action{\isacharunderscore}{\kern0pt}params{\isacharunderscore}{\kern0pt}match\ a\ args{\isachardoublequoteclose}\isanewline
\ \ \ \ \isakeyword{assumes}\ {\isachardoublequoteopen}wf{\isacharunderscore}{\kern0pt}action{\isacharunderscore}{\kern0pt}schema\ a{\isachardoublequoteclose}\isanewline
\ \ \ \ \isakeyword{shows}\ {\isachardoublequoteopen}wf{\isacharunderscore}{\kern0pt}ground{\isacharunderscore}{\kern0pt}action\ {\isacharparenleft}{\kern0pt}instantiate{\isacharunderscore}{\kern0pt}action{\isacharunderscore}{\kern0pt}schema\ a\ args{\isacharparenright}{\kern0pt}{\isachardoublequoteclose}\isanewline
%
\isadelimproof
\ \ %
\endisadelimproof
%
\isatagproof
\isacommand{proof}\isamarkupfalse%
\ {\isacharminus}{\kern0pt}\isanewline
\ \ \ \ \isacommand{have}\isamarkupfalse%
\ INST{\isacharcolon}{\kern0pt}\isanewline
\ \ \ \ \ \ {\isachardoublequoteopen}is{\isacharunderscore}{\kern0pt}of{\isacharunderscore}{\kern0pt}type\ objT\ {\isacharparenleft}{\kern0pt}{\isacharparenleft}{\kern0pt}the\ {\isasymcirc}\ map{\isacharunderscore}{\kern0pt}of\ {\isacharparenleft}{\kern0pt}zip\ {\isacharparenleft}{\kern0pt}map\ fst\ params{\isacharparenright}{\kern0pt}\ args{\isacharparenright}{\kern0pt}{\isacharparenright}{\kern0pt}\ x{\isacharparenright}{\kern0pt}\ T{\isachardoublequoteclose}\isanewline
\ \ \ \ \ \ \isakeyword{if}\ {\isachardoublequoteopen}is{\isacharunderscore}{\kern0pt}of{\isacharunderscore}{\kern0pt}type\ {\isacharparenleft}{\kern0pt}map{\isacharunderscore}{\kern0pt}of\ params{\isacharparenright}{\kern0pt}\ x\ T{\isachardoublequoteclose}\ \isakeyword{for}\ x\ T\isanewline
\ \ \ \ \ \ \isacommand{using}\isamarkupfalse%
\ that\isanewline
\ \ \ \ \ \ \isacommand{apply}\isamarkupfalse%
\ {\isacharparenleft}{\kern0pt}rule\ is{\isacharunderscore}{\kern0pt}of{\isacharunderscore}{\kern0pt}type{\isacharunderscore}{\kern0pt}map{\isacharunderscore}{\kern0pt}ofE{\isacharparenright}{\kern0pt}\isanewline
\ \ \ \ \ \ \isacommand{using}\isamarkupfalse%
\ assms\isanewline
\ \ \ \ \ \ \isacommand{apply}\isamarkupfalse%
\ {\isacharparenleft}{\kern0pt}clarsimp\ simp{\isacharcolon}{\kern0pt}\ Let{\isacharunderscore}{\kern0pt}def{\isacharparenright}{\kern0pt}\isanewline
\ \ \ \ \ \ \isacommand{subgoal}\isamarkupfalse%
\ \isakeyword{for}\ i\ xT\isanewline
\ \ \ \ \ \ \ \ \isacommand{unfolding}\isamarkupfalse%
\ action{\isacharunderscore}{\kern0pt}params{\isacharunderscore}{\kern0pt}match{\isacharunderscore}{\kern0pt}def\isanewline
\ \ \ \ \ \ \ \ \isacommand{apply}\isamarkupfalse%
\ {\isacharparenleft}{\kern0pt}subst\ lookup{\isacharunderscore}{\kern0pt}zip{\isacharunderscore}{\kern0pt}idx{\isacharunderscore}{\kern0pt}eq{\isacharbrackleft}{\kern0pt}\isakeyword{where}\ i{\isacharequal}{\kern0pt}i{\isacharbrackright}{\kern0pt}{\isacharsemicolon}{\kern0pt}\isanewline
\ \ \ \ \ \ \ \ \ \ {\isacharparenleft}{\kern0pt}clarsimp\ simp{\isacharcolon}{\kern0pt}\ list{\isacharunderscore}{\kern0pt}all{\isadigit{2}}{\isacharunderscore}{\kern0pt}lengthD{\isacharparenright}{\kern0pt}{\isacharquery}{\kern0pt}{\isacharparenright}{\kern0pt}\isanewline
\ \ \ \ \ \ \ \ \isacommand{apply}\isamarkupfalse%
\ {\isacharparenleft}{\kern0pt}frule\ list{\isacharunderscore}{\kern0pt}all{\isadigit{2}}{\isacharunderscore}{\kern0pt}nthD{\isadigit{2}}{\isacharbrackleft}{\kern0pt}\isakeyword{where}\ p{\isacharequal}{\kern0pt}i{\isacharbrackright}{\kern0pt}{\isacharsemicolon}{\kern0pt}\ simp{\isacharquery}{\kern0pt}{\isacharparenright}{\kern0pt}\isanewline
\ \ \ \ \ \ \ \ \isacommand{apply}\isamarkupfalse%
\ {\isacharparenleft}{\kern0pt}auto\isanewline
\ \ \ \ \ \ \ \ \ \ \ \ \ \ \ \ simp{\isacharcolon}{\kern0pt}\ is{\isacharunderscore}{\kern0pt}obj{\isacharunderscore}{\kern0pt}of{\isacharunderscore}{\kern0pt}type{\isacharunderscore}{\kern0pt}alt\ is{\isacharunderscore}{\kern0pt}of{\isacharunderscore}{\kern0pt}type{\isacharunderscore}{\kern0pt}def\isanewline
\ \ \ \ \ \ \ \ \ \ \ \ \ \ \ \ intro{\isacharcolon}{\kern0pt}\ of{\isacharunderscore}{\kern0pt}type{\isacharunderscore}{\kern0pt}trans\isanewline
\ \ \ \ \ \ \ \ \ \ \ \ \ \ \ \ split{\isacharcolon}{\kern0pt}\ option{\isachardot}{\kern0pt}splits{\isacharparenright}{\kern0pt}\isanewline
\ \ \ \ \ \ \ \ \isacommand{done}\isamarkupfalse%
\isanewline
\ \ \ \ \ \ \isacommand{done}\isamarkupfalse%
\isanewline
\ \ \ \ \isacommand{then}\isamarkupfalse%
\ \isacommand{show}\isamarkupfalse%
\ {\isacharquery}{\kern0pt}thesis\isanewline
\ \ \ \ \ \ \isacommand{using}\isamarkupfalse%
\ assms\ wf{\isacharunderscore}{\kern0pt}inst{\isacharunderscore}{\kern0pt}formula\ wf{\isacharunderscore}{\kern0pt}inst{\isacharunderscore}{\kern0pt}effect\isanewline
\ \ \ \ \ \ \isacommand{by}\isamarkupfalse%
\ {\isacharparenleft}{\kern0pt}fastforce\ split{\isacharcolon}{\kern0pt}\ term{\isachardot}{\kern0pt}splits\ simp{\isacharcolon}{\kern0pt}\ Let{\isacharunderscore}{\kern0pt}def\ comp{\isacharunderscore}{\kern0pt}apply{\isacharbrackleft}{\kern0pt}abs{\isacharunderscore}{\kern0pt}def{\isacharbrackright}{\kern0pt}{\isacharparenright}{\kern0pt}\isanewline
\ \ \isacommand{qed}\isamarkupfalse%
%
\endisatagproof
{\isafoldproof}%
%
\isadelimproof
\isanewline
%
\endisadelimproof
\isanewline
\ \ \isacommand{theorem}\isamarkupfalse%
\ wf{\isacharunderscore}{\kern0pt}inst{\isacharunderscore}{\kern0pt}durative{\isacharunderscore}{\kern0pt}action{\isacharunderscore}{\kern0pt}schema{\isacharcolon}{\kern0pt}\isanewline
\ \ \ \ \isakeyword{fixes}\ n\ params\ d\ cond\ eff\isanewline
\ \ \ \ \isakeyword{defines}\ {\isachardoublequoteopen}a\ {\isasymequiv}\ Durative{\isacharunderscore}{\kern0pt}Action{\isacharunderscore}{\kern0pt}Schema\ n\ params\ d\ cond\ eff{\isachardoublequoteclose}\ \ \isanewline
\ \ \ \ \isakeyword{assumes}\ {\isachardoublequoteopen}action{\isacharunderscore}{\kern0pt}params{\isacharunderscore}{\kern0pt}match\ a\ args{\isachardoublequoteclose}\isanewline
\ \ \ \ \isakeyword{assumes}\ {\isachardoublequoteopen}wf{\isacharunderscore}{\kern0pt}action{\isacharunderscore}{\kern0pt}schema\ a{\isachardoublequoteclose}\isanewline
\ \ \ \ \isakeyword{shows}\ {\isachardoublequoteopen}wf{\isacharunderscore}{\kern0pt}ground{\isacharunderscore}{\kern0pt}action\ {\isacharparenleft}{\kern0pt}inst{\isacharunderscore}{\kern0pt}snap{\isacharunderscore}{\kern0pt}action\ a\ args\ ta{\isacharparenright}{\kern0pt}{\isachardoublequoteclose}\isanewline
%
\isadelimproof
\ \ %
\endisadelimproof
%
\isatagproof
\isacommand{proof}\isamarkupfalse%
\ {\isacharminus}{\kern0pt}\isanewline
\ \ \ \ \isacommand{have}\isamarkupfalse%
\ INST{\isacharcolon}{\kern0pt}\isanewline
\ \ \ \ \ \ {\isachardoublequoteopen}is{\isacharunderscore}{\kern0pt}of{\isacharunderscore}{\kern0pt}type\ objT\ {\isacharparenleft}{\kern0pt}{\isacharparenleft}{\kern0pt}the\ {\isasymcirc}\ map{\isacharunderscore}{\kern0pt}of\ {\isacharparenleft}{\kern0pt}zip\ {\isacharparenleft}{\kern0pt}map\ fst\ params{\isacharparenright}{\kern0pt}\ args{\isacharparenright}{\kern0pt}{\isacharparenright}{\kern0pt}\ x{\isacharparenright}{\kern0pt}\ T{\isachardoublequoteclose}\isanewline
\ \ \ \ \ \ \isakeyword{if}\ {\isachardoublequoteopen}is{\isacharunderscore}{\kern0pt}of{\isacharunderscore}{\kern0pt}type\ {\isacharparenleft}{\kern0pt}map{\isacharunderscore}{\kern0pt}of\ params{\isacharparenright}{\kern0pt}\ x\ T{\isachardoublequoteclose}\ \isakeyword{for}\ x\ T\isanewline
\ \ \ \ \ \ \isacommand{using}\isamarkupfalse%
\ that\isanewline
\ \ \ \ \ \ \isacommand{apply}\isamarkupfalse%
\ {\isacharparenleft}{\kern0pt}rule\ is{\isacharunderscore}{\kern0pt}of{\isacharunderscore}{\kern0pt}type{\isacharunderscore}{\kern0pt}map{\isacharunderscore}{\kern0pt}ofE{\isacharparenright}{\kern0pt}\isanewline
\ \ \ \ \ \ \isacommand{using}\isamarkupfalse%
\ assms\isanewline
\ \ \ \ \ \ \isacommand{apply}\isamarkupfalse%
\ {\isacharparenleft}{\kern0pt}clarsimp\ simp{\isacharcolon}{\kern0pt}\ Let{\isacharunderscore}{\kern0pt}def{\isacharparenright}{\kern0pt}\isanewline
\ \ \ \ \ \ \isacommand{subgoal}\isamarkupfalse%
\ \isakeyword{for}\ i\ xT\isanewline
\ \ \ \ \ \ \ \ \isacommand{unfolding}\isamarkupfalse%
\ action{\isacharunderscore}{\kern0pt}params{\isacharunderscore}{\kern0pt}match{\isacharunderscore}{\kern0pt}def\isanewline
\ \ \ \ \ \ \ \ \isacommand{apply}\isamarkupfalse%
\ {\isacharparenleft}{\kern0pt}subst\ lookup{\isacharunderscore}{\kern0pt}zip{\isacharunderscore}{\kern0pt}idx{\isacharunderscore}{\kern0pt}eq{\isacharbrackleft}{\kern0pt}\isakeyword{where}\ i{\isacharequal}{\kern0pt}i{\isacharbrackright}{\kern0pt}{\isacharsemicolon}{\kern0pt}\isanewline
\ \ \ \ \ \ \ \ \ \ {\isacharparenleft}{\kern0pt}clarsimp\ simp{\isacharcolon}{\kern0pt}\ list{\isacharunderscore}{\kern0pt}all{\isadigit{2}}{\isacharunderscore}{\kern0pt}lengthD{\isacharparenright}{\kern0pt}{\isacharquery}{\kern0pt}{\isacharparenright}{\kern0pt}\isanewline
\ \ \ \ \ \ \ \ \isacommand{apply}\isamarkupfalse%
\ {\isacharparenleft}{\kern0pt}frule\ list{\isacharunderscore}{\kern0pt}all{\isadigit{2}}{\isacharunderscore}{\kern0pt}nthD{\isadigit{2}}{\isacharbrackleft}{\kern0pt}\isakeyword{where}\ p{\isacharequal}{\kern0pt}i{\isacharbrackright}{\kern0pt}{\isacharsemicolon}{\kern0pt}\ simp{\isacharquery}{\kern0pt}{\isacharparenright}{\kern0pt}\isanewline
\ \ \ \ \ \ \ \ \isacommand{apply}\isamarkupfalse%
\ {\isacharparenleft}{\kern0pt}auto\isanewline
\ \ \ \ \ \ \ \ \ \ \ \ \ \ \ \ simp{\isacharcolon}{\kern0pt}\ is{\isacharunderscore}{\kern0pt}obj{\isacharunderscore}{\kern0pt}of{\isacharunderscore}{\kern0pt}type{\isacharunderscore}{\kern0pt}alt\ is{\isacharunderscore}{\kern0pt}of{\isacharunderscore}{\kern0pt}type{\isacharunderscore}{\kern0pt}def\isanewline
\ \ \ \ \ \ \ \ \ \ \ \ \ \ \ \ intro{\isacharcolon}{\kern0pt}\ of{\isacharunderscore}{\kern0pt}type{\isacharunderscore}{\kern0pt}trans\isanewline
\ \ \ \ \ \ \ \ \ \ \ \ \ \ \ \ split{\isacharcolon}{\kern0pt}\ option{\isachardot}{\kern0pt}splits{\isacharparenright}{\kern0pt}\isanewline
\ \ \ \ \ \ \ \ \isacommand{done}\isamarkupfalse%
\isanewline
\ \ \ \ \ \ \isacommand{done}\isamarkupfalse%
\isanewline
\ \ \ \ \isacommand{then}\isamarkupfalse%
\ \isacommand{show}\isamarkupfalse%
\ {\isacharquery}{\kern0pt}thesis\isanewline
\ \ \ \ \ \ \isacommand{using}\isamarkupfalse%
\ assms\isanewline
\ \ \ \ \ \ \ \ wf{\isacharunderscore}{\kern0pt}inst{\isacharunderscore}{\kern0pt}formula{\isacharbrackleft}{\kern0pt}OF\ {\isacharunderscore}{\kern0pt}\ wf{\isacharunderscore}{\kern0pt}BigAnd{\isacharunderscore}{\kern0pt}filter{\isacharunderscore}{\kern0pt}time{\isacharunderscore}{\kern0pt}spec{\isacharcomma}{\kern0pt}\ \isanewline
\ \ \ \ \ \ \ \ \ \ \isakeyword{where}\ f{\isadigit{1}}{\isacharequal}{\kern0pt}{\isachardoublequoteopen}{\isasymlambda}x{\isachardot}{\kern0pt}\ the\ {\isacharparenleft}{\kern0pt}map{\isacharunderscore}{\kern0pt}of\ {\isacharparenleft}{\kern0pt}zip\ {\isacharparenleft}{\kern0pt}map\ fst\ params{\isacharparenright}{\kern0pt}\ args{\isacharparenright}{\kern0pt}\ x{\isacharparenright}{\kern0pt}{\isachardoublequoteclose}\ \isanewline
\ \ \ \ \ \ \ \ \ \ \ \ \ \ \isakeyword{and}\ params{\isadigit{1}}{\isacharequal}{\kern0pt}{\isachardoublequoteopen}params{\isachardoublequoteclose}\ \isakeyword{and}\ Q{\isadigit{1}}{\isacharequal}{\kern0pt}{\isachardoublequoteopen}{\isacharparenleft}{\kern0pt}map{\isacharunderscore}{\kern0pt}of\ params{\isacharparenright}{\kern0pt}{\isachardoublequoteclose}{\isacharbrackright}{\kern0pt}\isanewline
\ \ \ \ \ \ \ \ wf{\isacharunderscore}{\kern0pt}inst{\isacharunderscore}{\kern0pt}effect{\isacharbrackleft}{\kern0pt}OF\ {\isacharunderscore}{\kern0pt}\ wf{\isacharunderscore}{\kern0pt}conj{\isacharunderscore}{\kern0pt}eff{\isacharunderscore}{\kern0pt}filter{\isacharunderscore}{\kern0pt}time{\isacharunderscore}{\kern0pt}spec{\isacharcomma}{\kern0pt}\ \isanewline
\ \ \ \ \ \ \ \ \isakeyword{where}\ f{\isadigit{1}}{\isacharequal}{\kern0pt}{\isachardoublequoteopen}{\isasymlambda}x{\isachardot}{\kern0pt}\ the\ {\isacharparenleft}{\kern0pt}map{\isacharunderscore}{\kern0pt}of\ {\isacharparenleft}{\kern0pt}zip\ {\isacharparenleft}{\kern0pt}map\ fst\ params{\isacharparenright}{\kern0pt}\ args{\isacharparenright}{\kern0pt}\ x{\isacharparenright}{\kern0pt}{\isachardoublequoteclose}\ \isanewline
\ \ \ \ \ \ \ \ \ \ \ \ \ \ \isakeyword{and}\ params{\isadigit{1}}{\isacharequal}{\kern0pt}{\isachardoublequoteopen}params{\isachardoublequoteclose}\ \isakeyword{and}\ Q{\isadigit{1}}{\isacharequal}{\kern0pt}{\isachardoublequoteopen}{\isacharparenleft}{\kern0pt}map{\isacharunderscore}{\kern0pt}of\ params{\isacharparenright}{\kern0pt}{\isachardoublequoteclose}{\isacharbrackright}{\kern0pt}\isanewline
\ \ \ \ \ \ \isacommand{apply}\isamarkupfalse%
\ {\isacharparenleft}{\kern0pt}auto\ split{\isacharcolon}{\kern0pt}\ term{\isachardot}{\kern0pt}splits\ simp{\isacharcolon}{\kern0pt}\ Let{\isacharunderscore}{\kern0pt}def\ comp{\isacharunderscore}{\kern0pt}apply{\isacharbrackleft}{\kern0pt}abs{\isacharunderscore}{\kern0pt}def{\isacharbrackright}{\kern0pt}{\isacharparenright}{\kern0pt}\isanewline
\ \ \ \ \ \ \isacommand{done}\isamarkupfalse%
\ \isanewline
\ \ \isacommand{qed}\isamarkupfalse%
%
\endisatagproof
{\isafoldproof}%
%
\isadelimproof
\isanewline
%
\endisadelimproof
\isanewline
\ \ \isacommand{lemma}\isamarkupfalse%
\ wf{\isacharunderscore}{\kern0pt}over{\isacharunderscore}{\kern0pt}all{\isacharunderscore}{\kern0pt}empty{\isacharunderscore}{\kern0pt}eff{\isacharcolon}{\kern0pt}\isanewline
\ \ \ \ \isakeyword{assumes}\ {\isachardoublequoteopen}wf{\isacharunderscore}{\kern0pt}plan{\isacharunderscore}{\kern0pt}action\ {\isasympi}\ {\isachardoublequoteclose}\ \isanewline
\ \ \ \ \ \ \ \ \isakeyword{and}\ {\isachardoublequoteopen}wf{\isacharunderscore}{\kern0pt}action{\isacharunderscore}{\kern0pt}schema\ a\isactrlsub s\isactrlsub c\isactrlsub h\isactrlsub e\isactrlsub m\isactrlsub a{\isachardoublequoteclose}\isanewline
\ \ \ \ \ \ \ \ \isakeyword{and}\ {\isachardoublequoteopen}Some\ a\isactrlsub s\isactrlsub c\isactrlsub h\isactrlsub e\isactrlsub m\isactrlsub a\ {\isacharequal}{\kern0pt}\ resolve{\isacharunderscore}{\kern0pt}action{\isacharunderscore}{\kern0pt}schema\ {\isacharparenleft}{\kern0pt}name\ {\isasympi}{\isacharparenright}{\kern0pt}{\isachardoublequoteclose}\isanewline
\ \ \ \ \ \ \ \ \isakeyword{and}\ {\isachardoublequoteopen}Some\ a\ {\isacharequal}{\kern0pt}\ res{\isacharunderscore}{\kern0pt}inst{\isacharunderscore}{\kern0pt}snap{\isacharunderscore}{\kern0pt}action\ {\isasympi}\ Over{\isacharunderscore}{\kern0pt}All{\isachardoublequoteclose}\isanewline
\ \ \ \ \ \ \isakeyword{shows}\ {\isachardoublequoteopen}effect\ a\ {\isacharequal}{\kern0pt}\ Effect\ {\isacharbrackleft}{\kern0pt}{\isacharbrackright}{\kern0pt}\ {\isacharbrackleft}{\kern0pt}{\isacharbrackright}{\kern0pt}{\isachardoublequoteclose}\isanewline
%
\isadelimproof
\ \ \ \ %
\endisadelimproof
%
\isatagproof
\isacommand{using}\isamarkupfalse%
\ assms\isanewline
\ \ \isacommand{proof}\isamarkupfalse%
\ {\isacharparenleft}{\kern0pt}cases\ a\isactrlsub s\isactrlsub c\isactrlsub h\isactrlsub e\isactrlsub m\isactrlsub a{\isacharparenright}{\kern0pt}\isanewline
\ \ \ \ \isacommand{case}\isamarkupfalse%
\ {\isacharparenleft}{\kern0pt}Simple{\isacharunderscore}{\kern0pt}Action{\isacharunderscore}{\kern0pt}Schema\ n\ params\ pre\ eff{\isacharparenright}{\kern0pt}\isanewline
\ \ \ \ \isacommand{then}\isamarkupfalse%
\ \isacommand{show}\isamarkupfalse%
\ {\isacharquery}{\kern0pt}thesis\isanewline
\ \ \ \ \ \ \isacommand{using}\isamarkupfalse%
\ assms\isanewline
\ \ \ \ \ \ \isacommand{by}\isamarkupfalse%
\ {\isacharparenleft}{\kern0pt}cases\ {\isasympi}{\isacharparenright}{\kern0pt}\ {\isacharparenleft}{\kern0pt}auto\ split{\isacharcolon}{\kern0pt}\ option{\isachardot}{\kern0pt}splits{\isacharparenright}{\kern0pt}\isanewline
\ \ \isacommand{next}\isamarkupfalse%
\isanewline
\ \ \ \ \isacommand{case}\isamarkupfalse%
\ {\isacharparenleft}{\kern0pt}Durative{\isacharunderscore}{\kern0pt}Action{\isacharunderscore}{\kern0pt}Schema\ n\ params\ d\ cond\ eff{\isacharparenright}{\kern0pt}\isanewline
\ \ \ \ \isacommand{then}\isamarkupfalse%
\ \isacommand{have}\isamarkupfalse%
\ {\isachardoublequoteopen}{\isasymforall}{\isacharparenleft}{\kern0pt}t{\isacharcomma}{\kern0pt}e{\isacharparenright}{\kern0pt}\ {\isasymin}\ set\ eff{\isachardot}{\kern0pt}\ t\ {\isasymnoteq}\ Over{\isacharunderscore}{\kern0pt}All{\isachardoublequoteclose}\isanewline
\ \ \ \ \ \ \isacommand{using}\isamarkupfalse%
\ {\isacartoucheopen}wf{\isacharunderscore}{\kern0pt}action{\isacharunderscore}{\kern0pt}schema\ a\isactrlsub s\isactrlsub c\isactrlsub h\isactrlsub e\isactrlsub m\isactrlsub a{\isacartoucheclose}\ \isacommand{by}\isamarkupfalse%
\ {\isacharparenleft}{\kern0pt}auto\ simp{\isacharcolon}{\kern0pt}\ Let{\isacharunderscore}{\kern0pt}def{\isacharparenright}{\kern0pt}\isanewline
\ \ \ \ \isacommand{then}\isamarkupfalse%
\ \isacommand{have}\isamarkupfalse%
\ {\isachardoublequoteopen}{\isacharparenleft}{\kern0pt}{\isasymAnd}\isactrlsub e\isactrlsub f\isactrlsub f\ {\isacharparenleft}{\kern0pt}filter{\isacharunderscore}{\kern0pt}time{\isacharunderscore}{\kern0pt}spec\ Over{\isacharunderscore}{\kern0pt}All\ eff{\isacharparenright}{\kern0pt}{\isacharparenright}{\kern0pt}\ {\isacharequal}{\kern0pt}\ Effect\ {\isacharbrackleft}{\kern0pt}{\isacharbrackright}{\kern0pt}\ {\isacharbrackleft}{\kern0pt}{\isacharbrackright}{\kern0pt}{\isachardoublequoteclose}\isanewline
\ \ \ \ \ \ \isacommand{unfolding}\isamarkupfalse%
\ filter{\isacharunderscore}{\kern0pt}time{\isacharunderscore}{\kern0pt}spec{\isacharunderscore}{\kern0pt}def\ conjunct{\isacharunderscore}{\kern0pt}effects{\isacharunderscore}{\kern0pt}def\ \isacommand{by}\isamarkupfalse%
\ auto\isanewline
\ \ \ \ \isacommand{then}\isamarkupfalse%
\ \isacommand{show}\isamarkupfalse%
\ {\isacharquery}{\kern0pt}thesis\ \isanewline
\ \ \ \ \ \ \isacommand{using}\isamarkupfalse%
\ assms\ Durative{\isacharunderscore}{\kern0pt}Action{\isacharunderscore}{\kern0pt}Schema\isanewline
\ \ \ \ \ \ \isacommand{by}\isamarkupfalse%
\ {\isacharparenleft}{\kern0pt}cases\ {\isasympi}{\isacharparenright}{\kern0pt}\ {\isacharparenleft}{\kern0pt}auto\ split{\isacharcolon}{\kern0pt}\ option{\isachardot}{\kern0pt}splits{\isacharparenright}{\kern0pt}\ \isanewline
\ \ \isacommand{qed}\isamarkupfalse%
%
\endisatagproof
{\isafoldproof}%
%
\isadelimproof
\isanewline
%
\endisadelimproof
\ \ \isanewline
\isacommand{end}\isamarkupfalse%
\ %
\isamarkupcmt{Context of \isa{ast{\isacharunderscore}{\kern0pt}problem}%
}%
\isadelimdocument
%
\endisadelimdocument
%
\isatagdocument
%
\isamarkupsubsubsection{Preservation%
}
\isamarkuptrue%
%
\endisatagdocument
{\isafolddocument}%
%
\isadelimdocument
%
\endisadelimdocument
\isacommand{context}\isamarkupfalse%
\ ast{\isacharunderscore}{\kern0pt}problem\ \isakeyword{begin}\isanewline
\isanewline
\ \ \isacommand{fun}\isamarkupfalse%
\ happ{\isacharunderscore}{\kern0pt}precond\ {\isacharcolon}{\kern0pt}{\isacharcolon}{\kern0pt}\ {\isachardoublequoteopen}happening\ {\isasymRightarrow}\ world{\isacharunderscore}{\kern0pt}model\ {\isasymRightarrow}\ bool{\isachardoublequoteclose}\ \isakeyword{where}\isanewline
\ \ \ \ {\isachardoublequoteopen}happ{\isacharunderscore}{\kern0pt}precond\ {\isacharparenleft}{\kern0pt}t{\isacharcomma}{\kern0pt}A{\isacharparenright}{\kern0pt}\ M\ {\isasymlongleftrightarrow}\ {\isacharparenleft}{\kern0pt}{\isasymforall}a\ {\isasymin}\ set\ A{\isachardot}{\kern0pt}\ M\ \isactrlsup c{\isasymTTurnstile}\isactrlsub {\isacharequal}{\kern0pt}\ precondition\ a{\isacharparenright}{\kern0pt}{\isachardoublequoteclose}\isanewline
\isanewline
\ \ \isacommand{fun}\isamarkupfalse%
\ happ{\isacharunderscore}{\kern0pt}non{\isacharunderscore}{\kern0pt}intrf\ {\isacharcolon}{\kern0pt}{\isacharcolon}{\kern0pt}\ {\isachardoublequoteopen}happening\ {\isasymRightarrow}\ bool{\isachardoublequoteclose}\ \isakeyword{where}\isanewline
\ \ \ \ {\isachardoublequoteopen}happ{\isacharunderscore}{\kern0pt}non{\isacharunderscore}{\kern0pt}intrf\ {\isacharparenleft}{\kern0pt}t{\isacharcomma}{\kern0pt}A{\isacharparenright}{\kern0pt}\ {\isasymlongleftrightarrow}\ {\isacharparenleft}{\kern0pt}{\isasymforall}a\ {\isasymin}\ set\ A{\isachardot}{\kern0pt}\ {\isasymforall}b\ {\isasymin}\ set\ A{\isachardot}{\kern0pt}\ a\ {\isasymnoteq}\ b\ {\isasymlongrightarrow}\ acts{\isacharunderscore}{\kern0pt}non{\isacharunderscore}{\kern0pt}intrf\ a\ b{\isacharparenright}{\kern0pt}{\isachardoublequoteclose}%
\begin{isamarkuptext}%
Shorthand for enabled happening: all preconditions satisfied, no interfering actions, 
    \& all ground actions are wellformed .%
\end{isamarkuptext}\isamarkuptrue%
\ \ \isacommand{fun}\isamarkupfalse%
\ happ{\isacharunderscore}{\kern0pt}enabled\ {\isacharcolon}{\kern0pt}{\isacharcolon}{\kern0pt}\ {\isachardoublequoteopen}happening\ {\isasymRightarrow}\ world{\isacharunderscore}{\kern0pt}model\ {\isasymRightarrow}\ bool{\isachardoublequoteclose}\ \isakeyword{where}\ \isanewline
\ \ \ \ {\isachardoublequoteopen}happ{\isacharunderscore}{\kern0pt}enabled\ {\isacharparenleft}{\kern0pt}t{\isacharcomma}{\kern0pt}A{\isacharparenright}{\kern0pt}\ M\ {\isasymlongleftrightarrow}\ {\isacharparenleft}{\kern0pt}happ{\isacharunderscore}{\kern0pt}precond\ {\isacharparenleft}{\kern0pt}t{\isacharcomma}{\kern0pt}A{\isacharparenright}{\kern0pt}\ M\ {\isasymand}\ happ{\isacharunderscore}{\kern0pt}non{\isacharunderscore}{\kern0pt}intrf\ {\isacharparenleft}{\kern0pt}t{\isacharcomma}{\kern0pt}A{\isacharparenright}{\kern0pt}\ {\isasymand}\ {\isacharparenleft}{\kern0pt}{\isasymforall}a\ {\isasymin}\ set\ A{\isachardot}{\kern0pt}\ wf{\isacharunderscore}{\kern0pt}ground{\isacharunderscore}{\kern0pt}action\ a{\isacharparenright}{\kern0pt}{\isacharparenright}{\kern0pt}{\isachardoublequoteclose}%
\begin{isamarkuptext}%
Shorthand for wellformed-ness of a happening sequence: all ground actions wellformed .%
\end{isamarkuptext}\isamarkuptrue%
\ \ \isacommand{fun}\isamarkupfalse%
\ wf{\isacharunderscore}{\kern0pt}happ{\isacharunderscore}{\kern0pt}seq\ {\isacharcolon}{\kern0pt}{\isacharcolon}{\kern0pt}\ {\isachardoublequoteopen}happening\ list\ {\isasymRightarrow}\ bool{\isachardoublequoteclose}\ \isakeyword{where}\isanewline
\ \ \ \ {\isachardoublequoteopen}wf{\isacharunderscore}{\kern0pt}happ{\isacharunderscore}{\kern0pt}seq\ hs\ {\isasymlongleftrightarrow}\ {\isacharparenleft}{\kern0pt}{\isasymforall}{\isacharparenleft}{\kern0pt}t{\isacharcomma}{\kern0pt}A{\isacharparenright}{\kern0pt}\ {\isasymin}\ set\ hs{\isachardot}{\kern0pt}\ {\isacharparenleft}{\kern0pt}{\isasymforall}a\ {\isasymin}\ set\ A{\isachardot}{\kern0pt}\ wf{\isacharunderscore}{\kern0pt}ground{\isacharunderscore}{\kern0pt}action\ a{\isacharparenright}{\kern0pt}{\isacharparenright}{\kern0pt}{\isachardoublequoteclose}\isanewline
\isanewline
\isacommand{end}\isamarkupfalse%
\ %
\isamarkupcmt{Context of \isa{ast{\isacharunderscore}{\kern0pt}problem}%
}\isanewline
\isanewline
\isacommand{context}\isamarkupfalse%
\ wf{\isacharunderscore}{\kern0pt}ast{\isacharunderscore}{\kern0pt}problem\ \isakeyword{begin}%
\begin{isamarkuptext}%
The initial world model is well-formed%
\end{isamarkuptext}\isamarkuptrue%
\ \ \isacommand{lemma}\isamarkupfalse%
\ wf{\isacharunderscore}{\kern0pt}I{\isacharcolon}{\kern0pt}\ {\isachardoublequoteopen}wf{\isacharunderscore}{\kern0pt}world{\isacharunderscore}{\kern0pt}model\ I{\isachardoublequoteclose}\isanewline
%
\isadelimproof
\ \ \ \ %
\endisadelimproof
%
\isatagproof
\isacommand{using}\isamarkupfalse%
\ wf{\isacharunderscore}{\kern0pt}problem\isanewline
\ \ \ \ \isacommand{unfolding}\isamarkupfalse%
\ I{\isacharunderscore}{\kern0pt}def\ wf{\isacharunderscore}{\kern0pt}world{\isacharunderscore}{\kern0pt}model{\isacharunderscore}{\kern0pt}def\ wf{\isacharunderscore}{\kern0pt}problem{\isacharunderscore}{\kern0pt}def\isanewline
\ \ \ \ \isacommand{apply}\isamarkupfalse%
{\isacharparenleft}{\kern0pt}safe{\isacharparenright}{\kern0pt}\ \isanewline
\ \ \ \ \isacommand{subgoal}\isamarkupfalse%
\ \isakeyword{for}\ f\ \isanewline
\ \ \ \ \ \ \isacommand{apply}\isamarkupfalse%
\ {\isacharparenleft}{\kern0pt}induction\ f{\isacharparenright}{\kern0pt}\ \isanewline
\ \ \ \ \ \ \isacommand{apply}\isamarkupfalse%
\ auto\isanewline
\ \ \ \ \ \ \isacommand{by}\isamarkupfalse%
\ {\isacharparenleft}{\kern0pt}metis\ is{\isacharunderscore}{\kern0pt}predAtom{\isachardot}{\kern0pt}elims{\isacharparenleft}{\kern0pt}{\isadigit{2}}{\isacharparenright}{\kern0pt}\ wf{\isacharunderscore}{\kern0pt}func{\isacharunderscore}{\kern0pt}assign{\isachardot}{\kern0pt}simps{\isacharparenleft}{\kern0pt}{\isadigit{2}}{\isacharparenright}{\kern0pt}{\isacharparenright}{\kern0pt}\isanewline
\ \ \ \ \isacommand{done}\isamarkupfalse%
%
\endisatagproof
{\isafoldproof}%
%
\isadelimproof
\isanewline
%
\endisadelimproof
\isanewline
\ \ \isanewline
\ \ \isacommand{lemma}\isamarkupfalse%
\ wf{\isacharunderscore}{\kern0pt}ground{\isacharunderscore}{\kern0pt}action{\isacharunderscore}{\kern0pt}wf{\isacharunderscore}{\kern0pt}fmla{\isacharunderscore}{\kern0pt}atom{\isacharcolon}{\kern0pt}\ \isanewline
\ \ \ \ {\isachardoublequoteopen}wf{\isacharunderscore}{\kern0pt}ground{\isacharunderscore}{\kern0pt}action\ a\ {\isasymLongrightarrow}\ {\isasymforall}x\ {\isasymin}\ set\ {\isacharparenleft}{\kern0pt}{\isacharparenleft}{\kern0pt}adds\ o\ effect{\isacharparenright}{\kern0pt}\ a{\isacharparenright}{\kern0pt}{\isachardot}{\kern0pt}\ wf{\isacharunderscore}{\kern0pt}fmla{\isacharunderscore}{\kern0pt}atom\ objT\ x{\isachardoublequoteclose}\isanewline
%
\isadelimproof
\ \ \ \ %
\endisadelimproof
%
\isatagproof
\isacommand{using}\isamarkupfalse%
\ wf{\isacharunderscore}{\kern0pt}effect{\isacharunderscore}{\kern0pt}inst{\isacharunderscore}{\kern0pt}alt{\isacharbrackleft}{\kern0pt}symmetric{\isacharbrackright}{\kern0pt}\ wf{\isacharunderscore}{\kern0pt}effect{\isachardot}{\kern0pt}elims{\isacharparenleft}{\kern0pt}{\isadigit{2}}{\isacharparenright}{\kern0pt}\isanewline
\ \ \ \ \isacommand{by}\isamarkupfalse%
\ {\isacharparenleft}{\kern0pt}cases\ a{\isacharparenright}{\kern0pt}\ force%
\endisatagproof
{\isafoldproof}%
%
\isadelimproof
%
\endisadelimproof
%
\begin{isamarkuptext}%
The Application of a happening preserves well-formedness%
\end{isamarkuptext}\isamarkuptrue%
\ \ \isacommand{theorem}\isamarkupfalse%
\ wf{\isacharunderscore}{\kern0pt}apply{\isacharunderscore}{\kern0pt}happ{\isacharcolon}{\kern0pt}\isanewline
\ \ \ \ \isakeyword{assumes}\ {\isachardoublequoteopen}happ{\isacharunderscore}{\kern0pt}enabled\ {\isacharparenleft}{\kern0pt}t{\isacharcomma}{\kern0pt}as{\isacharparenright}{\kern0pt}\ s{\isachardoublequoteclose}\ \isanewline
\ \ \ \ \ \ \ \ \isakeyword{and}\ {\isachardoublequoteopen}wf{\isacharunderscore}{\kern0pt}world{\isacharunderscore}{\kern0pt}model\ s{\isachardoublequoteclose}\isanewline
\ \ \ \ \isakeyword{shows}\ {\isachardoublequoteopen}wf{\isacharunderscore}{\kern0pt}world{\isacharunderscore}{\kern0pt}model\ {\isacharparenleft}{\kern0pt}apply{\isacharunderscore}{\kern0pt}happ\ {\isacharparenleft}{\kern0pt}t{\isacharcomma}{\kern0pt}as{\isacharparenright}{\kern0pt}\ s{\isacharparenright}{\kern0pt}{\isachardoublequoteclose}\isanewline
%
\isadelimproof
\ \ \ \ %
\endisadelimproof
%
\isatagproof
\isacommand{using}\isamarkupfalse%
\ assms\ wf{\isacharunderscore}{\kern0pt}ground{\isacharunderscore}{\kern0pt}action{\isacharunderscore}{\kern0pt}wf{\isacharunderscore}{\kern0pt}fmla{\isacharunderscore}{\kern0pt}atom\ \isacommand{by}\isamarkupfalse%
\ {\isacharparenleft}{\kern0pt}auto\ simp{\isacharcolon}{\kern0pt}\ wf{\isacharunderscore}{\kern0pt}world{\isacharunderscore}{\kern0pt}model{\isacharunderscore}{\kern0pt}def{\isacharparenright}{\kern0pt}%
\endisatagproof
{\isafoldproof}%
%
\isadelimproof
\isanewline
%
\endisadelimproof
\isanewline
\ \ \isacommand{theorem}\isamarkupfalse%
\ wf{\isacharunderscore}{\kern0pt}apply{\isacharunderscore}{\kern0pt}happ{\isacharunderscore}{\kern0pt}compact{\isacharunderscore}{\kern0pt}notation{\isacharcolon}{\kern0pt}\isanewline
\ \ \ \ {\isachardoublequoteopen}happ{\isacharunderscore}{\kern0pt}enabled\ {\isacharparenleft}{\kern0pt}t{\isacharcomma}{\kern0pt}as{\isacharparenright}{\kern0pt}\ s\ {\isasymLongrightarrow}\ wf{\isacharunderscore}{\kern0pt}world{\isacharunderscore}{\kern0pt}model\ s\ {\isasymLongrightarrow}\ wf{\isacharunderscore}{\kern0pt}world{\isacharunderscore}{\kern0pt}model\ {\isacharparenleft}{\kern0pt}apply{\isacharunderscore}{\kern0pt}happ\ {\isacharparenleft}{\kern0pt}t{\isacharcomma}{\kern0pt}as{\isacharparenright}{\kern0pt}\ s{\isacharparenright}{\kern0pt}{\isachardoublequoteclose}\isanewline
%
\isadelimproof
\ \ \ \ %
\endisadelimproof
%
\isatagproof
\isacommand{by}\isamarkupfalse%
\ {\isacharparenleft}{\kern0pt}rule\ wf{\isacharunderscore}{\kern0pt}apply{\isacharunderscore}{\kern0pt}happ{\isacharparenright}{\kern0pt}%
\endisatagproof
{\isafoldproof}%
%
\isadelimproof
\isanewline
%
\endisadelimproof
\isanewline
\ \ \isacommand{corollary}\isamarkupfalse%
\ wf{\isacharunderscore}{\kern0pt}valid{\isacharunderscore}{\kern0pt}happ{\isacharunderscore}{\kern0pt}seq{\isacharcolon}{\kern0pt}\ \ \ \ \ \ \ \ \ \ \ \ \ \ \ \ \ \ \ \ \ \ \ \ \ \ \ \ \ \ \ \ \ \ \ \ \ \ \ \isanewline
\ \ \ \ \isakeyword{assumes}\ {\isachardoublequoteopen}wf{\isacharunderscore}{\kern0pt}world{\isacharunderscore}{\kern0pt}model\ M{\isachardoublequoteclose}\ \isanewline
\ \ \ \ \ \ \ \ \isakeyword{and}\ {\isachardoublequoteopen}wf{\isacharunderscore}{\kern0pt}happ{\isacharunderscore}{\kern0pt}seq\ hs{\isachardoublequoteclose}\isanewline
\ \ \ \ \ \ \ \ \isakeyword{and}\ {\isachardoublequoteopen}valid{\isacharunderscore}{\kern0pt}happ{\isacharunderscore}{\kern0pt}seq\ M\ hs\ M{\isacharprime}{\kern0pt}{\isachardoublequoteclose}\isanewline
\ \ \ \ \isakeyword{shows}\ {\isachardoublequoteopen}wf{\isacharunderscore}{\kern0pt}world{\isacharunderscore}{\kern0pt}model\ M{\isacharprime}{\kern0pt}{\isachardoublequoteclose}\isanewline
%
\isadelimproof
\ \ \ \ %
\endisadelimproof
%
\isatagproof
\isacommand{using}\isamarkupfalse%
\ assms\ wf{\isacharunderscore}{\kern0pt}apply{\isacharunderscore}{\kern0pt}happ\ \isacommand{by}\isamarkupfalse%
\ {\isacharparenleft}{\kern0pt}induction\ hs\ arbitrary{\isacharcolon}{\kern0pt}\ M{\isacharparenright}{\kern0pt}\ {\isacharparenleft}{\kern0pt}auto\ split{\isacharcolon}{\kern0pt}\ prod{\isachardot}{\kern0pt}splits\ simp{\isacharcolon}{\kern0pt}\ pairwise{\isacharunderscore}{\kern0pt}def{\isacharparenright}{\kern0pt}%
\endisatagproof
{\isafoldproof}%
%
\isadelimproof
%
\endisadelimproof
%
\begin{isamarkuptext}%
In a wellformed plan the action parameters match for every plan action%
\end{isamarkuptext}\isamarkuptrue%
\ \ \isacommand{lemma}\isamarkupfalse%
\ wf{\isacharunderscore}{\kern0pt}plan{\isacharunderscore}{\kern0pt}action{\isacharunderscore}{\kern0pt}params{\isacharunderscore}{\kern0pt}match{\isacharcolon}{\kern0pt}\ \isanewline
\ \ \ \ \isakeyword{assumes}\ {\isachardoublequoteopen}wf{\isacharunderscore}{\kern0pt}plan\ {\isasympi}s{\isachardoublequoteclose}\ \isanewline
\ \ \ \ \ \ \ \ \isakeyword{and}\ {\isachardoublequoteopen}{\isacharparenleft}{\kern0pt}t{\isacharcomma}{\kern0pt}a{\isacharparenright}{\kern0pt}\ {\isasymin}\ set\ {\isasympi}s{\isachardoublequoteclose}\ \isanewline
\ \ \ \ \ \ \ \ \isakeyword{and}\ {\isachardoublequoteopen}a\isactrlsub s\isactrlsub c\isactrlsub h\isactrlsub e\isactrlsub m\isactrlsub a\ {\isacharequal}{\kern0pt}\ the\ {\isacharparenleft}{\kern0pt}resolve{\isacharunderscore}{\kern0pt}action{\isacharunderscore}{\kern0pt}schema\ {\isacharparenleft}{\kern0pt}name\ a{\isacharparenright}{\kern0pt}{\isacharparenright}{\kern0pt}{\isachardoublequoteclose}\isanewline
\ \ \ \ \isakeyword{shows}\ {\isachardoublequoteopen}action{\isacharunderscore}{\kern0pt}params{\isacharunderscore}{\kern0pt}match\ a\isactrlsub s\isactrlsub c\isactrlsub h\isactrlsub e\isactrlsub m\isactrlsub a\ {\isacharparenleft}{\kern0pt}arguments\ a{\isacharparenright}{\kern0pt}{\isachardoublequoteclose}\isanewline
%
\isadelimproof
\ \ \ \ %
\endisadelimproof
%
\isatagproof
\isacommand{using}\isamarkupfalse%
\ assms\ \isacommand{by}\isamarkupfalse%
\ {\isacharparenleft}{\kern0pt}cases\ a{\isacharparenright}{\kern0pt}\ {\isacharparenleft}{\kern0pt}auto\ simp{\isacharcolon}{\kern0pt}\ wf{\isacharunderscore}{\kern0pt}plan{\isacharunderscore}{\kern0pt}def\ split{\isacharcolon}{\kern0pt}\ ast{\isacharunderscore}{\kern0pt}action{\isacharunderscore}{\kern0pt}schema{\isachardot}{\kern0pt}splits\ option{\isachardot}{\kern0pt}splits{\isacharparenright}{\kern0pt}%
\endisatagproof
{\isafoldproof}%
%
\isadelimproof
\isanewline
%
\endisadelimproof
\isanewline
\ \ \isacommand{lemma}\isamarkupfalse%
\ simple{\isacharunderscore}{\kern0pt}plan{\isacharunderscore}{\kern0pt}action{\isacharunderscore}{\kern0pt}schema{\isacharunderscore}{\kern0pt}type{\isadigit{1}}{\isacharcolon}{\kern0pt}\isanewline
\ \ \ \ \isakeyword{assumes}\ {\isachardoublequoteopen}wf{\isacharunderscore}{\kern0pt}plan{\isacharunderscore}{\kern0pt}action\ {\isacharparenleft}{\kern0pt}Simple{\isacharunderscore}{\kern0pt}Plan{\isacharunderscore}{\kern0pt}Action\ n\ args{\isacharparenright}{\kern0pt}{\isachardoublequoteclose}\isanewline
\ \ \ \ \isakeyword{shows}\ {\isachardoublequoteopen}{\isasymexists}params\ pre\ eff{\isachardot}{\kern0pt}\ resolve{\isacharunderscore}{\kern0pt}action{\isacharunderscore}{\kern0pt}schema\ n\ {\isacharequal}{\kern0pt}\ Some\ {\isacharparenleft}{\kern0pt}Simple{\isacharunderscore}{\kern0pt}Action{\isacharunderscore}{\kern0pt}Schema\ n\ params\ pre\ eff{\isacharparenright}{\kern0pt}{\isachardoublequoteclose}\isanewline
%
\isadelimproof
\ \ %
\endisadelimproof
%
\isatagproof
\isacommand{proof}\isamarkupfalse%
\ {\isacharparenleft}{\kern0pt}cases\ {\isachardoublequoteopen}resolve{\isacharunderscore}{\kern0pt}action{\isacharunderscore}{\kern0pt}schema\ n{\isachardoublequoteclose}{\isacharparenright}{\kern0pt}\isanewline
\ \ \ \ \isacommand{case}\isamarkupfalse%
\ None\isanewline
\ \ \ \ \isacommand{then}\isamarkupfalse%
\ \isacommand{show}\isamarkupfalse%
\ {\isacharquery}{\kern0pt}thesis\ \isanewline
\ \ \ \ \ \ \isacommand{using}\isamarkupfalse%
\ assms\ \isacommand{by}\isamarkupfalse%
\ auto\isanewline
\ \ \isacommand{next}\isamarkupfalse%
\isanewline
\ \ \ \ \isacommand{case}\isamarkupfalse%
\ {\isacharparenleft}{\kern0pt}Some\ a{\isacharparenright}{\kern0pt}\isanewline
\ \ \ \ \isacommand{from}\isamarkupfalse%
\ wf{\isacharunderscore}{\kern0pt}domain\ \isacommand{have}\isamarkupfalse%
\ {\isacharbrackleft}{\kern0pt}simp{\isacharbrackright}{\kern0pt}{\isacharcolon}{\kern0pt}\ {\isachardoublequoteopen}distinct\ {\isacharparenleft}{\kern0pt}map\ ast{\isacharunderscore}{\kern0pt}action{\isacharunderscore}{\kern0pt}schema{\isachardot}{\kern0pt}name\ {\isacharparenleft}{\kern0pt}actions\ D{\isacharparenright}{\kern0pt}{\isacharparenright}{\kern0pt}{\isachardoublequoteclose}\isanewline
\ \ \ \ \ \ \isacommand{unfolding}\isamarkupfalse%
\ wf{\isacharunderscore}{\kern0pt}domain{\isacharunderscore}{\kern0pt}def\ \isacommand{by}\isamarkupfalse%
\ auto\isanewline
\ \ \ \ \isacommand{then}\isamarkupfalse%
\ \isacommand{show}\isamarkupfalse%
\ {\isacharquery}{\kern0pt}thesis\ \isanewline
\ \ \ \ \ \ \isacommand{using}\isamarkupfalse%
\ assms\ {\isacartoucheopen}resolve{\isacharunderscore}{\kern0pt}action{\isacharunderscore}{\kern0pt}schema\ n\ {\isacharequal}{\kern0pt}\ Some\ a{\isacartoucheclose}\ resolve{\isacharunderscore}{\kern0pt}action{\isacharunderscore}{\kern0pt}schema{\isacharunderscore}{\kern0pt}def\isanewline
\ \ \ \ \ \ \isacommand{by}\isamarkupfalse%
\ {\isacharparenleft}{\kern0pt}cases\ a{\isacharparenright}{\kern0pt}\ fastforce{\isacharplus}{\kern0pt}\isanewline
\ \ \isacommand{qed}\isamarkupfalse%
%
\endisatagproof
{\isafoldproof}%
%
\isadelimproof
\isanewline
%
\endisadelimproof
\isanewline
\ \ \isacommand{lemma}\isamarkupfalse%
\ durative{\isacharunderscore}{\kern0pt}plan{\isacharunderscore}{\kern0pt}action{\isacharunderscore}{\kern0pt}schema{\isacharunderscore}{\kern0pt}type{\isadigit{1}}{\isacharcolon}{\kern0pt}\isanewline
\ \ \ \ \isakeyword{assumes}\ {\isachardoublequoteopen}wf{\isacharunderscore}{\kern0pt}plan{\isacharunderscore}{\kern0pt}action\ {\isacharparenleft}{\kern0pt}Durative{\isacharunderscore}{\kern0pt}Plan{\isacharunderscore}{\kern0pt}Action\ n\ args\ t{\isacharprime}{\kern0pt}{\isacharparenright}{\kern0pt}{\isachardoublequoteclose}\isanewline
\ \ \ \ \isakeyword{shows}\ {\isachardoublequoteopen}{\isasymexists}params\ d\ cond\ eff{\isachardot}{\kern0pt}\ resolve{\isacharunderscore}{\kern0pt}action{\isacharunderscore}{\kern0pt}schema\ n\ \isanewline
\ \ \ \ \ \ \ \ \ \ \ \ {\isacharequal}{\kern0pt}\ Some\ {\isacharparenleft}{\kern0pt}Durative{\isacharunderscore}{\kern0pt}Action{\isacharunderscore}{\kern0pt}Schema\ n\ params\ d\ cond\ eff{\isacharparenright}{\kern0pt}{\isachardoublequoteclose}\isanewline
%
\isadelimproof
\ \ %
\endisadelimproof
%
\isatagproof
\isacommand{proof}\isamarkupfalse%
\ {\isacharparenleft}{\kern0pt}cases\ {\isachardoublequoteopen}resolve{\isacharunderscore}{\kern0pt}action{\isacharunderscore}{\kern0pt}schema\ n{\isachardoublequoteclose}{\isacharparenright}{\kern0pt}\isanewline
\ \ \ \ \isacommand{case}\isamarkupfalse%
\ None\isanewline
\ \ \ \ \isacommand{then}\isamarkupfalse%
\ \isacommand{show}\isamarkupfalse%
\ {\isacharquery}{\kern0pt}thesis\ \isanewline
\ \ \ \ \ \ \isacommand{using}\isamarkupfalse%
\ assms\ \isacommand{by}\isamarkupfalse%
\ auto\isanewline
\ \ \isacommand{next}\isamarkupfalse%
\isanewline
\ \ \ \ \isacommand{case}\isamarkupfalse%
\ {\isacharparenleft}{\kern0pt}Some\ a{\isacharparenright}{\kern0pt}\isanewline
\ \ \ \ \isacommand{from}\isamarkupfalse%
\ wf{\isacharunderscore}{\kern0pt}domain\ \isacommand{have}\isamarkupfalse%
\ {\isacharbrackleft}{\kern0pt}simp{\isacharbrackright}{\kern0pt}{\isacharcolon}{\kern0pt}\ {\isachardoublequoteopen}distinct\ {\isacharparenleft}{\kern0pt}map\ ast{\isacharunderscore}{\kern0pt}action{\isacharunderscore}{\kern0pt}schema{\isachardot}{\kern0pt}name\ {\isacharparenleft}{\kern0pt}actions\ D{\isacharparenright}{\kern0pt}{\isacharparenright}{\kern0pt}{\isachardoublequoteclose}\isanewline
\ \ \ \ \ \ \isacommand{unfolding}\isamarkupfalse%
\ wf{\isacharunderscore}{\kern0pt}domain{\isacharunderscore}{\kern0pt}def\ \isacommand{by}\isamarkupfalse%
\ auto\isanewline
\ \ \ \ \isacommand{then}\isamarkupfalse%
\ \isacommand{show}\isamarkupfalse%
\ {\isacharquery}{\kern0pt}thesis\ \isanewline
\ \ \ \ \ \ \isacommand{using}\isamarkupfalse%
\ assms\ {\isacartoucheopen}resolve{\isacharunderscore}{\kern0pt}action{\isacharunderscore}{\kern0pt}schema\ n\ {\isacharequal}{\kern0pt}\ Some\ a{\isacartoucheclose}\ resolve{\isacharunderscore}{\kern0pt}action{\isacharunderscore}{\kern0pt}schema{\isacharunderscore}{\kern0pt}def\isanewline
\ \ \ \ \ \ \isacommand{by}\isamarkupfalse%
\ {\isacharparenleft}{\kern0pt}cases\ a{\isacharparenright}{\kern0pt}\ fastforce{\isacharplus}{\kern0pt}\isanewline
\ \ \isacommand{qed}\isamarkupfalse%
%
\endisatagproof
{\isafoldproof}%
%
\isadelimproof
%
\endisadelimproof
%
\begin{isamarkuptext}%
Instantiating a wellformed plan action always produces a wellformed ground action%
\end{isamarkuptext}\isamarkuptrue%
\ \ \isacommand{lemma}\isamarkupfalse%
\ wf{\isacharunderscore}{\kern0pt}inst{\isacharunderscore}{\kern0pt}of{\isacharunderscore}{\kern0pt}plan{\isacharunderscore}{\kern0pt}action{\isacharcolon}{\kern0pt}\isanewline
\ \ \ \ \isakeyword{assumes}\ {\isachardoublequoteopen}wf{\isacharunderscore}{\kern0pt}plan\ {\isasympi}s{\isachardoublequoteclose}\ \isanewline
\ \ \ \ \ \ \ \ \isakeyword{and}\ {\isachardoublequoteopen}{\isacharparenleft}{\kern0pt}t\isactrlsub {\isasympi}{\isacharcomma}{\kern0pt}{\isasympi}{\isacharparenright}{\kern0pt}\ {\isasymin}\ set\ {\isasympi}s{\isachardoublequoteclose}\ \isanewline
\ \ \ \ \ \ \ \ \isakeyword{and}\ {\isachardoublequoteopen}inst{\isacharunderscore}{\kern0pt}of{\isacharunderscore}{\kern0pt}plan{\isacharunderscore}{\kern0pt}action\ {\isasympi}s\ {\isacharparenleft}{\kern0pt}t\isactrlsub {\isasympi}{\isacharcomma}{\kern0pt}{\isasympi}{\isacharparenright}{\kern0pt}\ {\isacharparenleft}{\kern0pt}t\isactrlsub i{\isacharcomma}{\kern0pt}a{\isacharparenright}{\kern0pt}{\isachardoublequoteclose}\isanewline
\ \ \ \ \isakeyword{shows}\ {\isachardoublequoteopen}wf{\isacharunderscore}{\kern0pt}ground{\isacharunderscore}{\kern0pt}action\ a{\isachardoublequoteclose}\isanewline
%
\isadelimproof
\ \ \ \ %
\endisadelimproof
%
\isatagproof
\isacommand{using}\isamarkupfalse%
\ assms\isanewline
\ \ \isacommand{proof}\isamarkupfalse%
\ {\isacharminus}{\kern0pt}\isanewline
\ \ \ \ \isacommand{have}\isamarkupfalse%
\ wf{\isacharunderscore}{\kern0pt}{\isasympi}{\isacharcolon}{\kern0pt}\ {\isachardoublequoteopen}wf{\isacharunderscore}{\kern0pt}plan{\isacharunderscore}{\kern0pt}action\ {\isasympi}{\isachardoublequoteclose}\isanewline
\ \ \ \ \ \ \isacommand{using}\isamarkupfalse%
\ assms\ wf{\isacharunderscore}{\kern0pt}plan{\isacharunderscore}{\kern0pt}def\ \isacommand{by}\isamarkupfalse%
\ blast\isanewline
\ \ \ \ \isacommand{obtain}\isamarkupfalse%
\ ta\ \isakeyword{where}\ {\isachardoublequoteopen}res{\isacharunderscore}{\kern0pt}inst\ {\isasympi}\ {\isacharequal}{\kern0pt}\ Some\ a\ {\isasymor}\ res{\isacharunderscore}{\kern0pt}inst{\isacharunderscore}{\kern0pt}snap{\isacharunderscore}{\kern0pt}action\ {\isasympi}\ ta\ {\isacharequal}{\kern0pt}\ Some\ a{\isachardoublequoteclose}\isanewline
\ \ \ \ \ \ \isacommand{using}\isamarkupfalse%
\ assms\ \isacommand{by}\isamarkupfalse%
\ auto\isanewline
\ \ \ \ \isacommand{then}\isamarkupfalse%
\ \isacommand{show}\isamarkupfalse%
\ {\isacharquery}{\kern0pt}thesis\isanewline
\ \ \ \ \isacommand{proof}\isamarkupfalse%
\ {\isacharparenleft}{\kern0pt}elim\ disjE{\isacharparenright}{\kern0pt}\isanewline
\ \ \ \ \ \ \isacommand{assume}\isamarkupfalse%
\ case{\isacharunderscore}{\kern0pt}smpl{\isacharunderscore}{\kern0pt}act{\isacharcolon}{\kern0pt}\ {\isachardoublequoteopen}res{\isacharunderscore}{\kern0pt}inst\ {\isasympi}\ {\isacharequal}{\kern0pt}\ Some\ a{\isachardoublequoteclose}\isanewline
\ \ \ \ \ \ \isacommand{then}\isamarkupfalse%
\ \isacommand{show}\isamarkupfalse%
\ {\isacharquery}{\kern0pt}thesis\ \isanewline
\ \ \ \ \ \ \ \ \isacommand{proof}\isamarkupfalse%
\ {\isacharparenleft}{\kern0pt}cases\ {\isasympi}{\isacharparenright}{\kern0pt}\isanewline
\ \ \ \ \ \ \ \ \ \ \isacommand{case}\isamarkupfalse%
\ {\isacharparenleft}{\kern0pt}Simple{\isacharunderscore}{\kern0pt}Plan{\isacharunderscore}{\kern0pt}Action\ n\ args{\isacharparenright}{\kern0pt}\isanewline
\ \ \ \ \ \ \ \ \ \ \isacommand{then}\isamarkupfalse%
\ \isacommand{obtain}\isamarkupfalse%
\ params\ pre\ eff\ \isakeyword{where}\ \isanewline
\ \ \ \ \ \ \ \ \ \ \ \ \ \ \ \ a\isactrlsub s\isactrlsub c\isactrlsub h\isactrlsub e\isactrlsub m\isactrlsub a{\isacharunderscore}{\kern0pt}def{\isacharcolon}{\kern0pt}\ {\isachardoublequoteopen}resolve{\isacharunderscore}{\kern0pt}action{\isacharunderscore}{\kern0pt}schema\ n\ {\isacharequal}{\kern0pt}\ Some\ {\isacharparenleft}{\kern0pt}Simple{\isacharunderscore}{\kern0pt}Action{\isacharunderscore}{\kern0pt}Schema\ n\ params\ pre\ eff{\isacharparenright}{\kern0pt}{\isachardoublequoteclose}\isanewline
\ \ \ \ \ \ \ \ \ \ \ \ \isacommand{using}\isamarkupfalse%
\ wf{\isacharunderscore}{\kern0pt}{\isasympi}\ simple{\isacharunderscore}{\kern0pt}plan{\isacharunderscore}{\kern0pt}action{\isacharunderscore}{\kern0pt}schema{\isacharunderscore}{\kern0pt}type{\isadigit{1}}\ Simple{\isacharunderscore}{\kern0pt}Plan{\isacharunderscore}{\kern0pt}Action\ \isacommand{by}\isamarkupfalse%
\ blast\isanewline
\ \ \ \ \ \ \ \ \ \ \isacommand{with}\isamarkupfalse%
\ wf{\isacharunderscore}{\kern0pt}domain\ \isacommand{have}\isamarkupfalse%
\ {\isachardoublequoteopen}wf{\isacharunderscore}{\kern0pt}action{\isacharunderscore}{\kern0pt}schema\ {\isacharparenleft}{\kern0pt}Simple{\isacharunderscore}{\kern0pt}Action{\isacharunderscore}{\kern0pt}Schema\ n\ params\ pre\ eff{\isacharparenright}{\kern0pt}{\isachardoublequoteclose}\isanewline
\ \ \ \ \ \ \ \ \ \ \ \ \isacommand{unfolding}\isamarkupfalse%
\ wf{\isacharunderscore}{\kern0pt}domain{\isacharunderscore}{\kern0pt}def\ \isacommand{using}\isamarkupfalse%
\ resolve{\isacharunderscore}{\kern0pt}action{\isacharunderscore}{\kern0pt}schema{\isacharunderscore}{\kern0pt}def\ \isacommand{by}\isamarkupfalse%
\ force\isanewline
\ \ \ \ \ \ \ \ \ \ \isacommand{then}\isamarkupfalse%
\ \isacommand{show}\isamarkupfalse%
\ {\isacharquery}{\kern0pt}thesis\isanewline
\ \ \ \ \ \ \ \ \ \ \ \ \isacommand{using}\isamarkupfalse%
\ assms\ wf{\isacharunderscore}{\kern0pt}{\isasympi}\ case{\isacharunderscore}{\kern0pt}smpl{\isacharunderscore}{\kern0pt}act\ a\isactrlsub s\isactrlsub c\isactrlsub h\isactrlsub e\isactrlsub m\isactrlsub a{\isacharunderscore}{\kern0pt}def\ Simple{\isacharunderscore}{\kern0pt}Plan{\isacharunderscore}{\kern0pt}Action\ wf{\isacharunderscore}{\kern0pt}inst{\isacharunderscore}{\kern0pt}action{\isacharunderscore}{\kern0pt}schema\ \isanewline
\ \ \ \ \ \ \ \ \ \ \ \ \isacommand{by}\isamarkupfalse%
\ force\ \isanewline
\ \ \ \ \ \ \ \ \isacommand{next}\isamarkupfalse%
\isanewline
\ \ \ \ \ \ \ \ \ \ \isacommand{case}\isamarkupfalse%
\ {\isacharparenleft}{\kern0pt}Durative{\isacharunderscore}{\kern0pt}Plan{\isacharunderscore}{\kern0pt}Action\ n\ args\ d{\isacharparenright}{\kern0pt}\isanewline
\ \ \ \ \ \ \ \ \ \ \isacommand{then}\isamarkupfalse%
\ \isacommand{show}\isamarkupfalse%
\ {\isacharquery}{\kern0pt}thesis\ \isacommand{using}\isamarkupfalse%
\ case{\isacharunderscore}{\kern0pt}smpl{\isacharunderscore}{\kern0pt}act\ \isacommand{by}\isamarkupfalse%
\ auto\isanewline
\ \ \ \ \ \ \ \ \isacommand{qed}\isamarkupfalse%
\isanewline
\ \ \ \ \isacommand{next}\isamarkupfalse%
\isanewline
\ \ \ \ \ \ \isacommand{assume}\isamarkupfalse%
\ case{\isacharunderscore}{\kern0pt}snp{\isacharunderscore}{\kern0pt}act{\isacharcolon}{\kern0pt}\ {\isachardoublequoteopen}res{\isacharunderscore}{\kern0pt}inst{\isacharunderscore}{\kern0pt}snap{\isacharunderscore}{\kern0pt}action\ {\isasympi}\ ta\ {\isacharequal}{\kern0pt}\ Some\ a{\isachardoublequoteclose}\isanewline
\ \ \ \ \ \ \isacommand{then}\isamarkupfalse%
\ \isacommand{show}\isamarkupfalse%
\ {\isacharquery}{\kern0pt}thesis\ \isanewline
\ \ \ \ \ \ \ \ \isacommand{proof}\isamarkupfalse%
\ {\isacharparenleft}{\kern0pt}cases\ {\isasympi}{\isacharparenright}{\kern0pt}\isanewline
\ \ \ \ \ \ \ \ \ \ \isacommand{case}\isamarkupfalse%
\ {\isacharparenleft}{\kern0pt}Simple{\isacharunderscore}{\kern0pt}Plan{\isacharunderscore}{\kern0pt}Action\ n\ args{\isacharparenright}{\kern0pt}\isanewline
\ \ \ \ \ \ \ \ \ \ \isacommand{then}\isamarkupfalse%
\ \isacommand{show}\isamarkupfalse%
\ {\isacharquery}{\kern0pt}thesis\ \isacommand{using}\isamarkupfalse%
\ case{\isacharunderscore}{\kern0pt}snp{\isacharunderscore}{\kern0pt}act\ \isacommand{by}\isamarkupfalse%
\ auto\isanewline
\ \ \ \ \ \ \ \ \isacommand{next}\isamarkupfalse%
\isanewline
\ \ \ \ \ \ \ \ \ \ \isacommand{case}\isamarkupfalse%
\ {\isacharparenleft}{\kern0pt}Durative{\isacharunderscore}{\kern0pt}Plan{\isacharunderscore}{\kern0pt}Action\ n\ args\ d{\isacharparenright}{\kern0pt}\isanewline
\ \ \ \ \ \ \ \ \ \ \isacommand{then}\isamarkupfalse%
\ \isacommand{obtain}\isamarkupfalse%
\ params\ dcond\ cond\ eff\ \isakeyword{where}\ \isanewline
\ \ \ \ \ \ \ \ \ \ \ \ a\isactrlsub s\isactrlsub c\isactrlsub h\isactrlsub e\isactrlsub m\isactrlsub a{\isacharunderscore}{\kern0pt}def{\isacharcolon}{\kern0pt}\ {\isachardoublequoteopen}resolve{\isacharunderscore}{\kern0pt}action{\isacharunderscore}{\kern0pt}schema\ n\ {\isacharequal}{\kern0pt}\ Some\ {\isacharparenleft}{\kern0pt}Durative{\isacharunderscore}{\kern0pt}Action{\isacharunderscore}{\kern0pt}Schema\ n\ params\ dcond\ cond\ eff{\isacharparenright}{\kern0pt}{\isachardoublequoteclose}\isanewline
\ \ \ \ \ \ \ \ \ \ \ \ \isacommand{using}\isamarkupfalse%
\ wf{\isacharunderscore}{\kern0pt}{\isasympi}\ durative{\isacharunderscore}{\kern0pt}plan{\isacharunderscore}{\kern0pt}action{\isacharunderscore}{\kern0pt}schema{\isacharunderscore}{\kern0pt}type{\isadigit{1}}\ \isacommand{by}\isamarkupfalse%
\ blast\isanewline
\ \ \ \ \ \ \ \ \ \ \isacommand{with}\isamarkupfalse%
\ wf{\isacharunderscore}{\kern0pt}domain\ \isacommand{have}\isamarkupfalse%
\ {\isachardoublequoteopen}wf{\isacharunderscore}{\kern0pt}action{\isacharunderscore}{\kern0pt}schema\ {\isacharparenleft}{\kern0pt}Durative{\isacharunderscore}{\kern0pt}Action{\isacharunderscore}{\kern0pt}Schema\ n\ params\ dcond\ cond\ eff{\isacharparenright}{\kern0pt}{\isachardoublequoteclose}\isanewline
\ \ \ \ \ \ \ \ \ \ \ \ \isacommand{unfolding}\isamarkupfalse%
\ wf{\isacharunderscore}{\kern0pt}domain{\isacharunderscore}{\kern0pt}def\ \isacommand{using}\isamarkupfalse%
\ resolve{\isacharunderscore}{\kern0pt}action{\isacharunderscore}{\kern0pt}schema{\isacharunderscore}{\kern0pt}def\ \isacommand{by}\isamarkupfalse%
\ force\ \ \ \ \ \ \ \ \ \ \ \ \ \ \ \ \ \ \ \ \ \ \ \ \ \ \ \ \ \ \ \ \ \ \ \ \ \ \ \ \ \ \ \isanewline
\ \ \ \ \ \ \ \ \ \ \isacommand{then}\isamarkupfalse%
\ \isacommand{show}\isamarkupfalse%
\ {\isacharquery}{\kern0pt}thesis\isanewline
\ \ \ \ \ \ \ \ \ \ \ \ \isacommand{using}\isamarkupfalse%
\ assms\ wf{\isacharunderscore}{\kern0pt}{\isasympi}\ case{\isacharunderscore}{\kern0pt}snp{\isacharunderscore}{\kern0pt}act\ a\isactrlsub s\isactrlsub c\isactrlsub h\isactrlsub e\isactrlsub m\isactrlsub a{\isacharunderscore}{\kern0pt}def\ Durative{\isacharunderscore}{\kern0pt}Plan{\isacharunderscore}{\kern0pt}Action\ wf{\isacharunderscore}{\kern0pt}inst{\isacharunderscore}{\kern0pt}durative{\isacharunderscore}{\kern0pt}action{\isacharunderscore}{\kern0pt}schema\isanewline
\ \ \ \ \ \ \ \ \ \ \ \ \isacommand{by}\isamarkupfalse%
\ force\ \isanewline
\ \ \ \ \ \ \ \ \isacommand{qed}\isamarkupfalse%
\isanewline
\ \ \ \ \ \ \isacommand{qed}\isamarkupfalse%
\isanewline
\ \ \isacommand{qed}\isamarkupfalse%
%
\endisatagproof
{\isafoldproof}%
%
\isadelimproof
%
\endisadelimproof
%
\begin{isamarkuptext}%
Every ground action in an induced happening sequence for a wellformed plan 
  is wellformed.%
\end{isamarkuptext}\isamarkuptrue%
\ \ \isacommand{lemma}\isamarkupfalse%
\ happ{\isacharunderscore}{\kern0pt}seq{\isacharunderscore}{\kern0pt}wf{\isacharunderscore}{\kern0pt}ground{\isacharunderscore}{\kern0pt}action{\isacharcolon}{\kern0pt}\ \isanewline
\ \ \ \ \isakeyword{assumes}\ {\isachardoublequoteopen}wf{\isacharunderscore}{\kern0pt}plan\ {\isasympi}s{\isachardoublequoteclose}\ \isanewline
\ \ \ \ \ \ \ \ \isakeyword{and}\ {\isachardoublequoteopen}ind{\isacharunderscore}{\kern0pt}happ{\isacharunderscore}{\kern0pt}seq\ {\isasympi}s\ hs{\isachardoublequoteclose}\ \isanewline
\ \ \ \ \ \ \ \ \isakeyword{and}\ {\isachardoublequoteopen}{\isacharparenleft}{\kern0pt}t\isactrlsub i{\isacharcomma}{\kern0pt}A\isactrlsub i{\isacharparenright}{\kern0pt}\ {\isasymin}\ set\ hs{\isachardoublequoteclose}\ \isanewline
\ \ \ \ \ \ \ \ \isakeyword{and}\ {\isachardoublequoteopen}a\ {\isasymin}\ set\ A\isactrlsub i{\isachardoublequoteclose}\isanewline
\ \ \ \ \isakeyword{shows}\ {\isachardoublequoteopen}wf{\isacharunderscore}{\kern0pt}ground{\isacharunderscore}{\kern0pt}action\ a{\isachardoublequoteclose}\isanewline
%
\isadelimproof
\ \ \ \ %
\endisadelimproof
%
\isatagproof
\isacommand{unfolding}\isamarkupfalse%
\ ind{\isacharunderscore}{\kern0pt}happ{\isacharunderscore}{\kern0pt}seq{\isacharunderscore}{\kern0pt}def\isanewline
\ \ \isacommand{proof}\isamarkupfalse%
\ {\isacharminus}{\kern0pt}\isanewline
\ \ \ \ \isacommand{have}\isamarkupfalse%
\ {\isachardoublequoteopen}{\isasymexists}{\isacharparenleft}{\kern0pt}t\isactrlsub {\isasympi}{\isacharcomma}{\kern0pt}{\isasympi}{\isacharparenright}{\kern0pt}\ {\isasymin}\ set\ {\isasympi}s{\isachardot}{\kern0pt}\ inst{\isacharunderscore}{\kern0pt}of{\isacharunderscore}{\kern0pt}plan{\isacharunderscore}{\kern0pt}action\ {\isasympi}s\ {\isacharparenleft}{\kern0pt}t\isactrlsub {\isasympi}{\isacharcomma}{\kern0pt}{\isasympi}{\isacharparenright}{\kern0pt}\ {\isacharparenleft}{\kern0pt}t\isactrlsub i{\isacharcomma}{\kern0pt}a{\isacharparenright}{\kern0pt}{\isachardoublequoteclose}\isanewline
\ \ \ \ \ \ \isacommand{using}\isamarkupfalse%
\ assms\ ind{\isacharunderscore}{\kern0pt}happ{\isacharunderscore}{\kern0pt}seq{\isacharunderscore}{\kern0pt}def\ \isacommand{by}\isamarkupfalse%
\ blast\isanewline
\ \ \ \ \isacommand{then}\isamarkupfalse%
\ \isacommand{obtain}\isamarkupfalse%
\ t\isactrlsub {\isasympi}\ {\isasympi}\ \isakeyword{where}\ {\isachardoublequoteopen}{\isacharparenleft}{\kern0pt}t\isactrlsub {\isasympi}{\isacharcomma}{\kern0pt}{\isasympi}{\isacharparenright}{\kern0pt}\ {\isasymin}\ set\ {\isasympi}s{\isachardoublequoteclose}\ \isakeyword{and}\ {\isachardoublequoteopen}inst{\isacharunderscore}{\kern0pt}of{\isacharunderscore}{\kern0pt}plan{\isacharunderscore}{\kern0pt}action\ {\isasympi}s\ {\isacharparenleft}{\kern0pt}t\isactrlsub {\isasympi}{\isacharcomma}{\kern0pt}{\isasympi}{\isacharparenright}{\kern0pt}\ {\isacharparenleft}{\kern0pt}t\isactrlsub i{\isacharcomma}{\kern0pt}a{\isacharparenright}{\kern0pt}{\isachardoublequoteclose}\isanewline
\ \ \ \ \ \ \isacommand{by}\isamarkupfalse%
\ auto\isanewline
\ \ \ \ \isacommand{then}\isamarkupfalse%
\ \isacommand{show}\isamarkupfalse%
\ {\isacharquery}{\kern0pt}thesis\ \isanewline
\ \ \ \ \ \ \isacommand{using}\isamarkupfalse%
\ assms\ wf{\isacharunderscore}{\kern0pt}inst{\isacharunderscore}{\kern0pt}of{\isacharunderscore}{\kern0pt}plan{\isacharunderscore}{\kern0pt}action\isanewline
\ \ \ \ \ \ \isacommand{by}\isamarkupfalse%
\ blast\isanewline
\ \ \isacommand{qed}\isamarkupfalse%
%
\endisatagproof
{\isafoldproof}%
%
\isadelimproof
%
\endisadelimproof
%
\begin{isamarkuptext}%
An induced happening sequence for a wellformed plan is always a wellformed 
  happening sequence.%
\end{isamarkuptext}\isamarkuptrue%
\ \ \isacommand{lemma}\isamarkupfalse%
\ inst{\isacharunderscore}{\kern0pt}wf{\isacharunderscore}{\kern0pt}plan{\isacharunderscore}{\kern0pt}wf{\isacharunderscore}{\kern0pt}happ{\isacharunderscore}{\kern0pt}seq{\isacharcolon}{\kern0pt}\ \isanewline
\ \ \ \ \isakeyword{assumes}\ {\isachardoublequoteopen}wf{\isacharunderscore}{\kern0pt}plan\ {\isasympi}s{\isachardoublequoteclose}\ \isanewline
\ \ \ \ \ \ \ \ \isakeyword{and}\ {\isachardoublequoteopen}ind{\isacharunderscore}{\kern0pt}happ{\isacharunderscore}{\kern0pt}seq\ {\isasympi}s\ hs{\isachardoublequoteclose}\isanewline
\ \ \ \ \isakeyword{shows}\ {\isachardoublequoteopen}wf{\isacharunderscore}{\kern0pt}happ{\isacharunderscore}{\kern0pt}seq\ hs{\isachardoublequoteclose}\isanewline
%
\isadelimproof
\ \ \ \ %
\endisadelimproof
%
\isatagproof
\isacommand{unfolding}\isamarkupfalse%
\ Let{\isacharunderscore}{\kern0pt}def\ simple{\isacharunderscore}{\kern0pt}acts{\isacharunderscore}{\kern0pt}def\ durative{\isacharunderscore}{\kern0pt}acts{\isacharunderscore}{\kern0pt}def\ is{\isacharunderscore}{\kern0pt}act{\isacharunderscore}{\kern0pt}simple{\isacharunderscore}{\kern0pt}def\isanewline
\ \ \ \ \isacommand{using}\isamarkupfalse%
\ assms\ happ{\isacharunderscore}{\kern0pt}seq{\isacharunderscore}{\kern0pt}wf{\isacharunderscore}{\kern0pt}ground{\isacharunderscore}{\kern0pt}action\ \isacommand{by}\isamarkupfalse%
\ auto%
\endisatagproof
{\isafoldproof}%
%
\isadelimproof
%
\endisadelimproof
%
\begin{isamarkuptext}%
Execution of a plan preserves well-formedness%
\end{isamarkuptext}\isamarkuptrue%
\ \ \isacommand{corollary}\isamarkupfalse%
\ wf{\isacharunderscore}{\kern0pt}plan{\isacharunderscore}{\kern0pt}happ{\isacharunderscore}{\kern0pt}path{\isacharcolon}{\kern0pt}\isanewline
\ \ \ \ \isakeyword{assumes}\ {\isachardoublequoteopen}wf{\isacharunderscore}{\kern0pt}world{\isacharunderscore}{\kern0pt}model\ M{\isachardoublequoteclose}\ \isakeyword{and}\ {\isachardoublequoteopen}wf{\isacharunderscore}{\kern0pt}plan\ {\isasympi}s{\isachardoublequoteclose}\ \isakeyword{and}\ {\isachardoublequoteopen}plan{\isacharunderscore}{\kern0pt}happ{\isacharunderscore}{\kern0pt}path\ M\ {\isasympi}s\ M{\isacharprime}{\kern0pt}{\isachardoublequoteclose}\isanewline
\ \ \ \ \isakeyword{shows}\ {\isachardoublequoteopen}wf{\isacharunderscore}{\kern0pt}world{\isacharunderscore}{\kern0pt}model\ M{\isacharprime}{\kern0pt}{\isachardoublequoteclose}\isanewline
%
\isadelimproof
\ \ \ \ %
\endisadelimproof
%
\isatagproof
\isacommand{using}\isamarkupfalse%
\ assms\ inst{\isacharunderscore}{\kern0pt}wf{\isacharunderscore}{\kern0pt}plan{\isacharunderscore}{\kern0pt}wf{\isacharunderscore}{\kern0pt}happ{\isacharunderscore}{\kern0pt}seq\ wf{\isacharunderscore}{\kern0pt}valid{\isacharunderscore}{\kern0pt}happ{\isacharunderscore}{\kern0pt}seq\isanewline
\ \ \ \ \isacommand{unfolding}\isamarkupfalse%
\ plan{\isacharunderscore}{\kern0pt}happ{\isacharunderscore}{\kern0pt}path{\isacharunderscore}{\kern0pt}def\ \isacommand{by}\isamarkupfalse%
\ {\isacharparenleft}{\kern0pt}induction\ {\isasympi}s\ arbitrary{\isacharcolon}{\kern0pt}\ M{\isacharparenright}{\kern0pt}\ blast{\isacharplus}{\kern0pt}%
\endisatagproof
{\isafoldproof}%
%
\isadelimproof
\isanewline
%
\endisadelimproof
\isanewline
\isacommand{end}\isamarkupfalse%
\ %
\isamarkupcmt{Context of \isa{wf{\isacharunderscore}{\kern0pt}ast{\isacharunderscore}{\kern0pt}problem}%
}%
\begin{isamarkuptext}%
Auxiliary function that are used to divide happening sequences.%
\end{isamarkuptext}\isamarkuptrue%
\ \ \isacommand{definition}\isamarkupfalse%
\ {\isachardoublequoteopen}leq\ {\isasymequiv}\ {\isasymlambda}t\isactrlsub i{\isachardot}{\kern0pt}\ {\isasymlambda}{\isacharparenleft}{\kern0pt}t\isactrlsub a{\isacharcomma}{\kern0pt}as{\isacharparenright}{\kern0pt}{\isachardot}{\kern0pt}\ t\isactrlsub a\ {\isasymle}\ t\isactrlsub i{\isachardoublequoteclose}\isanewline
\ \ \isanewline
\ \ \isacommand{definition}\isamarkupfalse%
\ {\isachardoublequoteopen}le\ {\isasymequiv}\ {\isasymlambda}t\isactrlsub i{\isachardot}{\kern0pt}\ {\isasymlambda}{\isacharparenleft}{\kern0pt}t\isactrlsub a{\isacharcomma}{\kern0pt}as{\isacharparenright}{\kern0pt}{\isachardot}{\kern0pt}\ t\isactrlsub a\ {\isacharless}{\kern0pt}\ t\isactrlsub i{\isachardoublequoteclose}\isanewline
\isanewline
\isacommand{lemma}\isamarkupfalse%
\ strict{\isacharunderscore}{\kern0pt}sorted{\isacharunderscore}{\kern0pt}app{\isacharunderscore}{\kern0pt}iff{\isacharcolon}{\kern0pt}\ {\isachardoublequoteopen}strict{\isacharunderscore}{\kern0pt}sorted\ a\ {\isasymand}\ strict{\isacharunderscore}{\kern0pt}sorted\ b\ {\isasymand}\ {\isacharparenleft}{\kern0pt}{\isasymforall}x\ {\isasymin}\ set\ a{\isachardot}{\kern0pt}\ {\isasymforall}y\ {\isasymin}\ set\ b{\isachardot}{\kern0pt}\ x\ {\isacharless}{\kern0pt}\ y{\isacharparenright}{\kern0pt}\isanewline
\ \ \ \ \ \ \ \ \ \ {\isasymlongleftrightarrow}\ strict{\isacharunderscore}{\kern0pt}sorted\ {\isacharparenleft}{\kern0pt}a\ {\isacharat}{\kern0pt}\ b{\isacharparenright}{\kern0pt}{\isachardoublequoteclose}\isanewline
%
\isadelimproof
\ \ %
\endisadelimproof
%
\isatagproof
\isacommand{by}\isamarkupfalse%
\ {\isacharparenleft}{\kern0pt}induction\ a\ arbitrary{\isacharcolon}{\kern0pt}\ b{\isacharparenright}{\kern0pt}\ auto%
\endisatagproof
{\isafoldproof}%
%
\isadelimproof
\isanewline
%
\endisadelimproof
\isanewline
\isacommand{lemma}\isamarkupfalse%
\ empty{\isacharunderscore}{\kern0pt}filter{\isacharunderscore}{\kern0pt}conv{\isacharcolon}{\kern0pt}\ {\isachardoublequoteopen}{\isacharbrackleft}{\kern0pt}{\isacharbrackright}{\kern0pt}\ {\isacharequal}{\kern0pt}\ filter\ P\ xs\ {\isasymlongleftrightarrow}\ {\isacharparenleft}{\kern0pt}{\isasymforall}x{\isasymin}set\ xs{\isachardot}{\kern0pt}\ {\isasymnot}\ P\ x{\isacharparenright}{\kern0pt}{\isachardoublequoteclose}\isanewline
%
\isadelimproof
\ \ %
\endisadelimproof
%
\isatagproof
\isacommand{by}\isamarkupfalse%
{\isacharparenleft}{\kern0pt}auto\ dest{\isacharcolon}{\kern0pt}\ sym\ simp\ add{\isacharcolon}{\kern0pt}\ filter{\isacharunderscore}{\kern0pt}empty{\isacharunderscore}{\kern0pt}conv{\isacharparenright}{\kern0pt}%
\endisatagproof
{\isafoldproof}%
%
\isadelimproof
%
\endisadelimproof
%
\isadelimdocument
%
\endisadelimdocument
%
\isatagdocument
%
\isamarkupsubsubsection{Completeness of Semantics%
}
\isamarkuptrue%
%
\endisatagdocument
{\isafolddocument}%
%
\isadelimdocument
%
\endisadelimdocument
\isacommand{context}\isamarkupfalse%
\ wf{\isacharunderscore}{\kern0pt}ast{\isacharunderscore}{\kern0pt}problem\isanewline
\isakeyword{begin}\isanewline
\isanewline
\isacommand{fun}\isamarkupfalse%
\ consume{\isadigit{2}}\ \isakeyword{where}\isanewline
{\isachardoublequoteopen}consume{\isadigit{2}}\ {\isacharbrackleft}{\kern0pt}{\isacharbrackright}{\kern0pt}\ {\isacharbrackleft}{\kern0pt}{\isacharbrackright}{\kern0pt}\ {\isacharequal}{\kern0pt}\ {\isadigit{0}}{\isachardoublequoteclose}\isanewline
{\isacharbar}{\kern0pt}\ \ {\isachardoublequoteopen}consume{\isadigit{2}}\ {\isacharparenleft}{\kern0pt}x{\isacharhash}{\kern0pt}xs{\isacharparenright}{\kern0pt}\ {\isacharparenleft}{\kern0pt}y{\isacharhash}{\kern0pt}ys{\isacharparenright}{\kern0pt}\ {\isacharequal}{\kern0pt}\ {\isadigit{0}}\ {\isacharplus}{\kern0pt}\ consume{\isadigit{2}}\ xs\ ys{\isachardoublequoteclose}\isanewline
\isanewline
\isacommand{lemma}\isamarkupfalse%
\ empty{\isacharunderscore}{\kern0pt}happ{\isacharcolon}{\kern0pt}\isanewline
\ {\isachardoublequoteopen}{\isacharparenleft}{\kern0pt}{\isasymAnd}a{\isachardot}{\kern0pt}\ a\ {\isasymin}\ set\ A\ {\isasymLongrightarrow}\ effect\ a\ {\isacharequal}{\kern0pt}\ Effect\ {\isacharbrackleft}{\kern0pt}{\isacharbrackright}{\kern0pt}\ {\isacharbrackleft}{\kern0pt}{\isacharbrackright}{\kern0pt}{\isacharparenright}{\kern0pt}\ {\isasymLongrightarrow}\ apply{\isacharunderscore}{\kern0pt}happ\ {\isacharparenleft}{\kern0pt}t{\isacharcomma}{\kern0pt}\ A{\isacharparenright}{\kern0pt}\ M\ {\isacharequal}{\kern0pt}\ M{\isachardoublequoteclose}\isanewline
%
\isadelimproof
\ \ %
\endisadelimproof
%
\isatagproof
\isacommand{by}\isamarkupfalse%
\ {\isacharparenleft}{\kern0pt}induction\ A{\isacharparenright}{\kern0pt}\ auto%
\endisatagproof
{\isafoldproof}%
%
\isadelimproof
\isanewline
%
\endisadelimproof
\isanewline
\isacommand{lemma}\isamarkupfalse%
\ valid{\isacharunderscore}{\kern0pt}happ{\isacharunderscore}{\kern0pt}seq{\isacharunderscore}{\kern0pt}append{\isacharunderscore}{\kern0pt}{\isadigit{1}}{\isacharcolon}{\kern0pt}\isanewline
\ \ {\isachardoublequoteopen}{\isasymlbrakk}{\isasymAnd}t\ A\ a{\isachardot}{\kern0pt}{\isasymlbrakk}{\isacharparenleft}{\kern0pt}t{\isacharcomma}{\kern0pt}\ A{\isacharparenright}{\kern0pt}\ {\isasymin}\ set\ hs{\isadigit{1}}{\isacharsemicolon}{\kern0pt}\ a\ {\isasymin}\ set\ A{\isasymrbrakk}\ {\isasymLongrightarrow}\ effect\ a\ {\isacharequal}{\kern0pt}\ Effect\ {\isacharbrackleft}{\kern0pt}{\isacharbrackright}{\kern0pt}\ {\isacharbrackleft}{\kern0pt}{\isacharbrackright}{\kern0pt}{\isacharsemicolon}{\kern0pt}\isanewline
\ \ \ \ \ \ \ \ valid{\isacharunderscore}{\kern0pt}happ{\isacharunderscore}{\kern0pt}seq\ M\ hs{\isadigit{2}}\ M{\isacharprime}{\kern0pt}{\isacharsemicolon}{\kern0pt}\isanewline
\ \ \ \ \ \ \ \ {\isacharparenleft}{\kern0pt}{\isasymAnd}t\ A\ a{\isachardot}{\kern0pt}\ {\isasymlbrakk}\ {\isacharparenleft}{\kern0pt}t{\isacharcomma}{\kern0pt}\ A{\isacharparenright}{\kern0pt}\ {\isasymin}\ set\ hs{\isadigit{1}}{\isacharsemicolon}{\kern0pt}\ a{\isasymin}\ set\ A{\isasymrbrakk}\ {\isasymLongrightarrow}\ M\ \isactrlsup c{\isasymTTurnstile}\isactrlsub {\isacharequal}{\kern0pt}\ precondition\ a{\isacharparenright}{\kern0pt}{\isacharsemicolon}{\kern0pt}\isanewline
\ \ \ \ \ \ \ \ {\isacharparenleft}{\kern0pt}{\isasymAnd}t\ A\ a{\isachardot}{\kern0pt}\ {\isacharparenleft}{\kern0pt}t{\isacharcomma}{\kern0pt}\ A{\isacharparenright}{\kern0pt}\ {\isasymin}\ set\ hs{\isadigit{1}}\ {\isasymLongrightarrow}\ pairwise\ acts{\isacharunderscore}{\kern0pt}non{\isacharunderscore}{\kern0pt}intrf\ {\isacharparenleft}{\kern0pt}set\ A{\isacharparenright}{\kern0pt}{\isacharparenright}{\kern0pt}\ {\isasymrbrakk}\ {\isasymLongrightarrow}\isanewline
\ \ \ \ \ \ \ \ valid{\isacharunderscore}{\kern0pt}happ{\isacharunderscore}{\kern0pt}seq\ M\ {\isacharparenleft}{\kern0pt}hs{\isadigit{1}}\ {\isacharat}{\kern0pt}\ hs{\isadigit{2}}{\isacharparenright}{\kern0pt}\ M{\isacharprime}{\kern0pt}{\isachardoublequoteclose}\isanewline
%
\isadelimproof
\ \ %
\endisadelimproof
%
\isatagproof
\isacommand{by}\isamarkupfalse%
\ {\isacharparenleft}{\kern0pt}induction\ hs{\isadigit{1}}\ arbitrary{\isacharcolon}{\kern0pt}\ M{\isacharparenright}{\kern0pt}\ {\isacharparenleft}{\kern0pt}fastforce\ simp\ del{\isacharcolon}{\kern0pt}\ apply{\isacharunderscore}{\kern0pt}happ{\isachardot}{\kern0pt}simps\ simp\ add{\isacharcolon}{\kern0pt}\ empty{\isacharunderscore}{\kern0pt}happ{\isacharparenright}{\kern0pt}{\isacharplus}{\kern0pt}%
\endisatagproof
{\isafoldproof}%
%
\isadelimproof
\isanewline
%
\endisadelimproof
\isanewline
\isacommand{lemma}\isamarkupfalse%
\ valid{\isacharunderscore}{\kern0pt}happ{\isacharunderscore}{\kern0pt}seq{\isacharunderscore}{\kern0pt}append{\isacharunderscore}{\kern0pt}{\isadigit{2}}{\isacharcolon}{\kern0pt}\isanewline
\ \ {\isachardoublequoteopen}{\isasymlbrakk}\ valid{\isacharunderscore}{\kern0pt}happ{\isacharunderscore}{\kern0pt}seq\ M\ {\isacharparenleft}{\kern0pt}hs{\isadigit{1}}\ {\isacharat}{\kern0pt}\ hs{\isadigit{2}}{\isacharparenright}{\kern0pt}\ M{\isacharprime}{\kern0pt}{\isacharsemicolon}{\kern0pt}\ {\isasymAnd}t\ A\ a{\isachardot}{\kern0pt}{\isasymlbrakk}{\isacharparenleft}{\kern0pt}t{\isacharcomma}{\kern0pt}\ A{\isacharparenright}{\kern0pt}\ {\isasymin}\ set\ hs{\isadigit{1}}{\isacharsemicolon}{\kern0pt}\ a\ {\isasymin}\ set\ A{\isasymrbrakk}\ {\isasymLongrightarrow}\ effect\ a\ {\isacharequal}{\kern0pt}\ Effect\ {\isacharbrackleft}{\kern0pt}{\isacharbrackright}{\kern0pt}\ {\isacharbrackleft}{\kern0pt}{\isacharbrackright}{\kern0pt}{\isasymrbrakk}\ {\isasymLongrightarrow}\isanewline
\ \ \ \ \ \ \ \ valid{\isacharunderscore}{\kern0pt}happ{\isacharunderscore}{\kern0pt}seq\ M\ hs{\isadigit{2}}\ M{\isacharprime}{\kern0pt}{\isachardoublequoteclose}\isanewline
%
\isadelimproof
\ \ %
\endisadelimproof
%
\isatagproof
\isacommand{by}\isamarkupfalse%
\ {\isacharparenleft}{\kern0pt}induction\ hs{\isadigit{1}}\ arbitrary{\isacharcolon}{\kern0pt}\ M{\isacharparenright}{\kern0pt}\ {\isacharparenleft}{\kern0pt}fastforce\ simp\ del{\isacharcolon}{\kern0pt}\ apply{\isacharunderscore}{\kern0pt}happ{\isachardot}{\kern0pt}simps\ simp\ add{\isacharcolon}{\kern0pt}\ empty{\isacharunderscore}{\kern0pt}happ{\isacharparenright}{\kern0pt}{\isacharplus}{\kern0pt}%
\endisatagproof
{\isafoldproof}%
%
\isadelimproof
\isanewline
%
\endisadelimproof
\isanewline
\isacommand{lemma}\isamarkupfalse%
\ valid{\isacharunderscore}{\kern0pt}happ{\isacharunderscore}{\kern0pt}seq{\isacharunderscore}{\kern0pt}append{\isacharunderscore}{\kern0pt}{\isadigit{2}}{\isacharprime}{\kern0pt}{\isacharcolon}{\kern0pt}\isanewline
\ \ {\isachardoublequoteopen}{\isasymlbrakk}\ valid{\isacharunderscore}{\kern0pt}happ{\isacharunderscore}{\kern0pt}seq\ M\ {\isacharparenleft}{\kern0pt}hs{\isadigit{1}}\ {\isacharat}{\kern0pt}\ hs{\isadigit{2}}{\isacharparenright}{\kern0pt}\ M{\isacharprime}{\kern0pt}{\isacharsemicolon}{\kern0pt}\ {\isasymAnd}t\ A\ a{\isachardot}{\kern0pt}{\isasymlbrakk}{\isacharparenleft}{\kern0pt}t{\isacharcomma}{\kern0pt}\ A{\isacharparenright}{\kern0pt}\ {\isasymin}\ set\ hs{\isadigit{2}}{\isacharsemicolon}{\kern0pt}\ a\ {\isasymin}\ set\ A{\isasymrbrakk}\ {\isasymLongrightarrow}\ effect\ a\ {\isacharequal}{\kern0pt}\ Effect\ {\isacharbrackleft}{\kern0pt}{\isacharbrackright}{\kern0pt}\ {\isacharbrackleft}{\kern0pt}{\isacharbrackright}{\kern0pt}{\isasymrbrakk}\ {\isasymLongrightarrow}\isanewline
\ \ \ \ \ \ \ \ valid{\isacharunderscore}{\kern0pt}happ{\isacharunderscore}{\kern0pt}seq\ M\ hs{\isadigit{1}}\ M{\isacharprime}{\kern0pt}{\isachardoublequoteclose}\isanewline
%
\isadelimproof
\ \ %
\endisadelimproof
%
\isatagproof
\isacommand{apply}\isamarkupfalse%
\ {\isacharparenleft}{\kern0pt}induction\ hs{\isadigit{1}}\ arbitrary{\isacharcolon}{\kern0pt}\ hs{\isadigit{2}}\ M\ M{\isacharprime}{\kern0pt}{\isacharparenright}{\kern0pt}\ \isanewline
\ \ \isacommand{apply}\isamarkupfalse%
\ {\isacharparenleft}{\kern0pt}auto\ simp\ del{\isacharcolon}{\kern0pt}\ apply{\isacharunderscore}{\kern0pt}happ{\isachardot}{\kern0pt}simps\ simp\ add{\isacharcolon}{\kern0pt}\ empty{\isacharunderscore}{\kern0pt}happ{\isacharparenright}{\kern0pt}\isanewline
\ \ \isacommand{by}\isamarkupfalse%
\ {\isacharparenleft}{\kern0pt}metis\ empty{\isacharunderscore}{\kern0pt}happ\ old{\isachardot}{\kern0pt}prod{\isachardot}{\kern0pt}exhaust\ valid{\isacharunderscore}{\kern0pt}happ{\isacharunderscore}{\kern0pt}seq{\isacharunderscore}{\kern0pt}M\isactrlsub j{\isacharunderscore}{\kern0pt}is{\isacharunderscore}{\kern0pt}M\isactrlsub i{\isacharparenright}{\kern0pt}{\isacharplus}{\kern0pt}%
\endisatagproof
{\isafoldproof}%
%
\isadelimproof
\isanewline
%
\endisadelimproof
\isanewline
\isacommand{lemma}\isamarkupfalse%
\ valid{\isacharunderscore}{\kern0pt}happ{\isacharunderscore}{\kern0pt}seq{\isacharunderscore}{\kern0pt}append{\isacharunderscore}{\kern0pt}{\isadigit{3}}{\isacharcolon}{\kern0pt}\isanewline
\ \ {\isachardoublequoteopen}{\isasymlbrakk}{\isacharparenleft}{\kern0pt}t{\isacharcomma}{\kern0pt}\ A{\isacharparenright}{\kern0pt}\ {\isasymin}\ set\ hs{\isadigit{1}}{\isacharsemicolon}{\kern0pt}\ a{\isasymin}\ set\ A{\isacharsemicolon}{\kern0pt}\ valid{\isacharunderscore}{\kern0pt}happ{\isacharunderscore}{\kern0pt}seq\ M\ {\isacharparenleft}{\kern0pt}hs{\isadigit{1}}\ {\isacharat}{\kern0pt}\ hs{\isadigit{2}}{\isacharparenright}{\kern0pt}\ M{\isacharprime}{\kern0pt}{\isacharsemicolon}{\kern0pt}\isanewline
\ \ \ \ {\isasymAnd}t\ A\ a{\isachardot}{\kern0pt}{\isasymlbrakk}{\isacharparenleft}{\kern0pt}t{\isacharcomma}{\kern0pt}\ A{\isacharparenright}{\kern0pt}\ {\isasymin}\ set\ hs{\isadigit{1}}{\isacharsemicolon}{\kern0pt}\ a\ {\isasymin}\ set\ A{\isasymrbrakk}\ {\isasymLongrightarrow}\ effect\ a\ {\isacharequal}{\kern0pt}\ Effect\ {\isacharbrackleft}{\kern0pt}{\isacharbrackright}{\kern0pt}\ {\isacharbrackleft}{\kern0pt}{\isacharbrackright}{\kern0pt}{\isasymrbrakk}\isanewline
\ \ \ \ {\isasymLongrightarrow}\ M\ \isactrlsup c{\isasymTTurnstile}\isactrlsub {\isacharequal}{\kern0pt}\ precondition\ a{\isachardoublequoteclose}\isanewline
%
\isadelimproof
\ \ %
\endisadelimproof
%
\isatagproof
\isacommand{by}\isamarkupfalse%
\ {\isacharparenleft}{\kern0pt}induction\ hs{\isadigit{1}}\ arbitrary{\isacharcolon}{\kern0pt}\ M{\isacharparenright}{\kern0pt}\ {\isacharparenleft}{\kern0pt}fastforce\ simp\ del{\isacharcolon}{\kern0pt}\ apply{\isacharunderscore}{\kern0pt}happ{\isachardot}{\kern0pt}simps\ simp\ add{\isacharcolon}{\kern0pt}\ empty{\isacharunderscore}{\kern0pt}happ{\isacharparenright}{\kern0pt}{\isacharplus}{\kern0pt}%
\endisatagproof
{\isafoldproof}%
%
\isadelimproof
\isanewline
%
\endisadelimproof
\isanewline
\isacommand{lemma}\isamarkupfalse%
\ valid{\isacharunderscore}{\kern0pt}happ{\isacharunderscore}{\kern0pt}seq{\isacharunderscore}{\kern0pt}append{\isacharunderscore}{\kern0pt}{\isadigit{4}}{\isacharcolon}{\kern0pt}\isanewline
\ \ {\isachardoublequoteopen}{\isasymlbrakk}valid{\isacharunderscore}{\kern0pt}happ{\isacharunderscore}{\kern0pt}seq\ M\ {\isacharparenleft}{\kern0pt}hs{\isadigit{1}}\ {\isacharat}{\kern0pt}\ hs{\isadigit{2}}{\isacharparenright}{\kern0pt}\ M{\isacharprime}{\kern0pt}{\isacharsemicolon}{\kern0pt}\ {\isasymAnd}t\ A\ a{\isachardot}{\kern0pt}{\isasymlbrakk}{\isacharparenleft}{\kern0pt}t{\isacharcomma}{\kern0pt}\ A{\isacharparenright}{\kern0pt}\ {\isasymin}\ set\ hs{\isadigit{1}}{\isacharsemicolon}{\kern0pt}\ a\ {\isasymin}\ set\ A{\isasymrbrakk}\ {\isasymLongrightarrow}\ effect\ a\ {\isacharequal}{\kern0pt}\ Effect\ {\isacharbrackleft}{\kern0pt}{\isacharbrackright}{\kern0pt}\ {\isacharbrackleft}{\kern0pt}{\isacharbrackright}{\kern0pt}{\isasymrbrakk}\ {\isasymLongrightarrow}\isanewline
\ \ \ \ {\isacharparenleft}{\kern0pt}{\isasymAnd}t\ A\ a{\isachardot}{\kern0pt}\ {\isacharparenleft}{\kern0pt}t{\isacharcomma}{\kern0pt}\ A{\isacharparenright}{\kern0pt}\ {\isasymin}\ set\ hs{\isadigit{1}}\ {\isasymLongrightarrow}\ pairwise\ acts{\isacharunderscore}{\kern0pt}non{\isacharunderscore}{\kern0pt}intrf\ {\isacharparenleft}{\kern0pt}set\ A{\isacharparenright}{\kern0pt}{\isacharparenright}{\kern0pt}{\isachardoublequoteclose}\isanewline
%
\isadelimproof
\ \ %
\endisadelimproof
%
\isatagproof
\isacommand{by}\isamarkupfalse%
\ {\isacharparenleft}{\kern0pt}induction\ hs{\isadigit{1}}\ arbitrary{\isacharcolon}{\kern0pt}\ M{\isacharparenright}{\kern0pt}\ {\isacharparenleft}{\kern0pt}fastforce\ simp\ del{\isacharcolon}{\kern0pt}\ apply{\isacharunderscore}{\kern0pt}happ{\isachardot}{\kern0pt}simps\ simp\ add{\isacharcolon}{\kern0pt}\ empty{\isacharunderscore}{\kern0pt}happ{\isacharparenright}{\kern0pt}{\isacharplus}{\kern0pt}%
\endisatagproof
{\isafoldproof}%
%
\isadelimproof
\isanewline
%
\endisadelimproof
\isanewline
\isacommand{lemma}\isamarkupfalse%
\ valid{\isacharunderscore}{\kern0pt}happ{\isacharunderscore}{\kern0pt}seq{\isacharunderscore}{\kern0pt}empty{\isacharunderscore}{\kern0pt}eff{\isacharcolon}{\kern0pt}\isanewline
\ \ {\isachardoublequoteopen}{\isasymlbrakk}{\isasymAnd}t\ A\ a{\isachardot}{\kern0pt}{\isasymlbrakk}{\isacharparenleft}{\kern0pt}t{\isacharcomma}{\kern0pt}\ A{\isacharparenright}{\kern0pt}\ {\isasymin}\ set\ hs{\isacharsemicolon}{\kern0pt}\ a\ {\isasymin}\ set\ A{\isasymrbrakk}\ {\isasymLongrightarrow}\ effect\ a\ {\isacharequal}{\kern0pt}\ Effect\ {\isacharbrackleft}{\kern0pt}{\isacharbrackright}{\kern0pt}\ {\isacharbrackleft}{\kern0pt}{\isacharbrackright}{\kern0pt}{\isacharsemicolon}{\kern0pt}\isanewline
\ \ \ \ \ \ \ \ {\isacharparenleft}{\kern0pt}{\isasymAnd}t\ A\ a{\isachardot}{\kern0pt}\ {\isasymlbrakk}\ {\isacharparenleft}{\kern0pt}t{\isacharcomma}{\kern0pt}\ A{\isacharparenright}{\kern0pt}\ {\isasymin}\ set\ hs{\isacharsemicolon}{\kern0pt}\ a{\isasymin}\ set\ A{\isasymrbrakk}\ {\isasymLongrightarrow}\ M\ \isactrlsup c{\isasymTTurnstile}\isactrlsub {\isacharequal}{\kern0pt}\ precondition\ a{\isacharparenright}{\kern0pt}{\isacharsemicolon}{\kern0pt}\isanewline
\ \ \ \ \ \ \ \ {\isacharparenleft}{\kern0pt}{\isasymAnd}t\ A\ a{\isachardot}{\kern0pt}\ {\isacharparenleft}{\kern0pt}t{\isacharcomma}{\kern0pt}\ A{\isacharparenright}{\kern0pt}\ {\isasymin}\ set\ hs\ {\isasymLongrightarrow}\ pairwise\ acts{\isacharunderscore}{\kern0pt}non{\isacharunderscore}{\kern0pt}intrf\ {\isacharparenleft}{\kern0pt}set\ A{\isacharparenright}{\kern0pt}{\isacharparenright}{\kern0pt}\ {\isasymrbrakk}\ {\isasymLongrightarrow}\isanewline
\ \ \ \ \ \ \ \ valid{\isacharunderscore}{\kern0pt}happ{\isacharunderscore}{\kern0pt}seq\ M\ hs\ M{\isachardoublequoteclose}\ \ \isanewline
%
\isadelimproof
\ \ %
\endisadelimproof
%
\isatagproof
\isacommand{by}\isamarkupfalse%
\ {\isacharparenleft}{\kern0pt}induction\ hs{\isacharparenright}{\kern0pt}\ {\isacharparenleft}{\kern0pt}fastforce\ simp\ del{\isacharcolon}{\kern0pt}\ apply{\isacharunderscore}{\kern0pt}happ{\isachardot}{\kern0pt}simps\ simp\ add{\isacharcolon}{\kern0pt}\ empty{\isacharunderscore}{\kern0pt}happ{\isacharparenright}{\kern0pt}{\isacharplus}{\kern0pt}%
\endisatagproof
{\isafoldproof}%
%
\isadelimproof
\isanewline
%
\endisadelimproof
\isanewline
\isacommand{lemma}\isamarkupfalse%
\ valid{\isacharunderscore}{\kern0pt}happ{\isacharunderscore}{\kern0pt}seq{\isacharunderscore}{\kern0pt}append{\isacharunderscore}{\kern0pt}{\isadigit{5}}{\isacharcolon}{\kern0pt}\isanewline
\ \ {\isachardoublequoteopen}{\isasymlbrakk}{\isasymAnd}t\ A\ a{\isachardot}{\kern0pt}{\isasymlbrakk}{\isacharparenleft}{\kern0pt}t{\isacharcomma}{\kern0pt}\ A{\isacharparenright}{\kern0pt}\ {\isasymin}\ set\ hs{\isadigit{2}}{\isacharsemicolon}{\kern0pt}\ a\ {\isasymin}\ set\ A{\isasymrbrakk}\ {\isasymLongrightarrow}\ effect\ a\ {\isacharequal}{\kern0pt}\ Effect\ {\isacharbrackleft}{\kern0pt}{\isacharbrackright}{\kern0pt}\ {\isacharbrackleft}{\kern0pt}{\isacharbrackright}{\kern0pt}{\isacharsemicolon}{\kern0pt}\isanewline
\ \ \ \ \ \ \ \ valid{\isacharunderscore}{\kern0pt}happ{\isacharunderscore}{\kern0pt}seq\ M\ hs{\isadigit{1}}\ M{\isacharprime}{\kern0pt}{\isacharsemicolon}{\kern0pt}\isanewline
\ \ \ \ \ \ \ \ {\isacharparenleft}{\kern0pt}{\isasymAnd}t\ A\ a{\isachardot}{\kern0pt}\ {\isasymlbrakk}\ {\isacharparenleft}{\kern0pt}t{\isacharcomma}{\kern0pt}\ A{\isacharparenright}{\kern0pt}\ {\isasymin}\ set\ hs{\isadigit{2}}{\isacharsemicolon}{\kern0pt}\ a{\isasymin}\ set\ A{\isasymrbrakk}\ {\isasymLongrightarrow}\ M{\isacharprime}{\kern0pt}\ \isactrlsup c{\isasymTTurnstile}\isactrlsub {\isacharequal}{\kern0pt}\ precondition\ a{\isacharparenright}{\kern0pt}{\isacharsemicolon}{\kern0pt}\isanewline
\ \ \ \ \ \ \ \ {\isacharparenleft}{\kern0pt}{\isasymAnd}t\ A\ a{\isachardot}{\kern0pt}\ {\isacharparenleft}{\kern0pt}t{\isacharcomma}{\kern0pt}\ A{\isacharparenright}{\kern0pt}\ {\isasymin}\ set\ hs{\isadigit{2}}\ {\isasymLongrightarrow}\ pairwise\ acts{\isacharunderscore}{\kern0pt}non{\isacharunderscore}{\kern0pt}intrf\ {\isacharparenleft}{\kern0pt}set\ A{\isacharparenright}{\kern0pt}{\isacharparenright}{\kern0pt}\ {\isasymrbrakk}\ {\isasymLongrightarrow}\isanewline
\ \ \ \ \ \ \ \ valid{\isacharunderscore}{\kern0pt}happ{\isacharunderscore}{\kern0pt}seq\ M\ {\isacharparenleft}{\kern0pt}hs{\isadigit{1}}\ {\isacharat}{\kern0pt}\ hs{\isadigit{2}}{\isacharparenright}{\kern0pt}\ M{\isacharprime}{\kern0pt}{\isachardoublequoteclose}\ \ \isanewline
%
\isadelimproof
\ \ %
\endisadelimproof
%
\isatagproof
\isacommand{by}\isamarkupfalse%
\ {\isacharparenleft}{\kern0pt}induction\ hs{\isadigit{1}}\ arbitrary{\isacharcolon}{\kern0pt}\ M{\isacharparenright}{\kern0pt}\isanewline
\ \ \ \ \ {\isacharparenleft}{\kern0pt}auto\ simp\ del{\isacharcolon}{\kern0pt}\ apply{\isacharunderscore}{\kern0pt}happ{\isachardot}{\kern0pt}simps\ simp\ add{\isacharcolon}{\kern0pt}\ empty{\isacharunderscore}{\kern0pt}happ\ intro{\isacharbang}{\kern0pt}{\isacharcolon}{\kern0pt}\ valid{\isacharunderscore}{\kern0pt}happ{\isacharunderscore}{\kern0pt}seq{\isacharunderscore}{\kern0pt}empty{\isacharunderscore}{\kern0pt}eff{\isacharparenright}{\kern0pt}%
\endisatagproof
{\isafoldproof}%
%
\isadelimproof
\isanewline
%
\endisadelimproof
\isanewline
\isacommand{lemma}\isamarkupfalse%
\ strict{\isacharunderscore}{\kern0pt}sorted{\isacharunderscore}{\kern0pt}cons{\isacharcolon}{\kern0pt}\ {\isachardoublequoteopen}{\isasymlbrakk}strict{\isacharunderscore}{\kern0pt}sorted\ {\isacharparenleft}{\kern0pt}map\ f\ {\isacharparenleft}{\kern0pt}x\ {\isacharhash}{\kern0pt}\ xs{\isacharparenright}{\kern0pt}{\isacharparenright}{\kern0pt}{\isacharsemicolon}{\kern0pt}\ x{\isacharprime}{\kern0pt}\ {\isasymin}\ set\ xs{\isasymrbrakk}\ {\isasymLongrightarrow}\ f\ x\ {\isasymnoteq}\ f\ x{\isacharprime}{\kern0pt}{\isachardoublequoteclose}\isanewline
%
\isadelimproof
\ \ %
\endisadelimproof
%
\isatagproof
\isacommand{by}\isamarkupfalse%
\ {\isacharparenleft}{\kern0pt}induction\ xs{\isacharparenright}{\kern0pt}\ auto%
\endisatagproof
{\isafoldproof}%
%
\isadelimproof
\isanewline
%
\endisadelimproof
\isanewline
\isacommand{lemma}\isamarkupfalse%
\ strict{\isacharunderscore}{\kern0pt}sorted{\isacharunderscore}{\kern0pt}order{\isacharcolon}{\kern0pt}\isanewline
\ \ \ \ \ \ {\isachardoublequoteopen}{\isasymlbrakk}{\isacharparenleft}{\kern0pt}{\isasymAnd}x{\isachardot}{\kern0pt}\ x\ {\isasymin}\ s\ {\isasymLongrightarrow}\ {\isasymexists}y\ {\isasymin}\ set\ {\isacharparenleft}{\kern0pt}x{\isadigit{1}}\ {\isacharhash}{\kern0pt}\ xs{\isadigit{1}}\ {\isacharat}{\kern0pt}\ x{\isadigit{2}}\ {\isacharhash}{\kern0pt}\ xs{\isadigit{2}}{\isacharparenright}{\kern0pt}{\isachardot}{\kern0pt}\ f\ y\ {\isacharequal}{\kern0pt}\ x{\isacharparenright}{\kern0pt}{\isacharsemicolon}{\kern0pt}\ strict{\isacharunderscore}{\kern0pt}sorted\ {\isacharparenleft}{\kern0pt}map\ f\ {\isacharparenleft}{\kern0pt}x{\isadigit{1}}\ {\isacharhash}{\kern0pt}\ xs{\isadigit{1}}\ {\isacharat}{\kern0pt}\ x{\isadigit{2}}\ {\isacharhash}{\kern0pt}\ xs{\isadigit{2}}{\isacharparenright}{\kern0pt}{\isacharparenright}{\kern0pt}{\isacharsemicolon}{\kern0pt}\isanewline
\ \ \ \ \ \ \ \ {\isacharparenleft}{\kern0pt}{\isasymAnd}y{\isachardot}{\kern0pt}\ y\ {\isasymin}\ set\ xs{\isadigit{1}}\ {\isasymLongrightarrow}\ f\ y\ {\isasymnotin}\ s{\isacharparenright}{\kern0pt}{\isacharsemicolon}{\kern0pt}\ x\ {\isasymin}\ s{\isacharsemicolon}{\kern0pt}\ x\ {\isacharless}{\kern0pt}\ f\ x{\isadigit{2}}{\isasymrbrakk}\ {\isasymLongrightarrow}\ x\ {\isacharequal}{\kern0pt}\ f\ x{\isadigit{1}}{\isachardoublequoteclose}\isanewline
%
\isadelimproof
\ \ %
\endisadelimproof
%
\isatagproof
\isacommand{by}\isamarkupfalse%
\ {\isacharparenleft}{\kern0pt}fastforce\ simp{\isacharcolon}{\kern0pt}\ image{\isacharunderscore}{\kern0pt}def\ strict{\isacharunderscore}{\kern0pt}sorted{\isacharunderscore}{\kern0pt}app{\isacharunderscore}{\kern0pt}iff{\isacharbrackleft}{\kern0pt}symmetric{\isacharbrackright}{\kern0pt}{\isacharparenright}{\kern0pt}%
\endisatagproof
{\isafoldproof}%
%
\isadelimproof
\isanewline
%
\endisadelimproof
\isanewline
\isacommand{lemma}\isamarkupfalse%
\ in{\isacharunderscore}{\kern0pt}suff{\isacharcolon}{\kern0pt}\isanewline
\ \ {\isachardoublequoteopen}{\isasymlbrakk}{\isasymAnd}x{\isachardot}{\kern0pt}\ x\ {\isasymin}\ set\ xs{\isadigit{1}}\ {\isasymLongrightarrow}\ x\ {\isasymnoteq}\ x{\isadigit{2}}{\isacharsemicolon}{\kern0pt}\ x{\isadigit{2}}\ {\isasymin}\ set\ {\isacharparenleft}{\kern0pt}xs{\isadigit{1}}\ {\isacharat}{\kern0pt}\ xs{\isadigit{2}}{\isacharparenright}{\kern0pt}{\isasymrbrakk}\ {\isasymLongrightarrow}\ x{\isadigit{2}}\ {\isasymin}\ set\ xs{\isadigit{2}}{\isachardoublequoteclose}\isanewline
%
\isadelimproof
\ \ %
\endisadelimproof
%
\isatagproof
\isacommand{by}\isamarkupfalse%
\ auto%
\endisatagproof
{\isafoldproof}%
%
\isadelimproof
\isanewline
%
\endisadelimproof
\isanewline
\isacommand{lemma}\isamarkupfalse%
\ strict{\isacharunderscore}{\kern0pt}sorted{\isacharunderscore}{\kern0pt}order{\isacharunderscore}{\kern0pt}in{\isacharunderscore}{\kern0pt}suff{\isacharcolon}{\kern0pt}\isanewline
\ \ {\isachardoublequoteopen}{\isasymlbrakk}strict{\isacharunderscore}{\kern0pt}sorted\ {\isacharparenleft}{\kern0pt}xs{\isadigit{1}}\ {\isacharat}{\kern0pt}\ xs{\isadigit{2}}{\isacharparenright}{\kern0pt}{\isacharsemicolon}{\kern0pt}\ x{\isadigit{1}}\ {\isasymin}\ set\ xs{\isadigit{1}}{\isacharsemicolon}{\kern0pt}\ x{\isadigit{2}}\ {\isasymin}\ set\ xs{\isadigit{2}}{\isasymrbrakk}\ {\isasymLongrightarrow}\ x{\isadigit{1}}\ {\isacharless}{\kern0pt}\ x{\isadigit{2}}{\isachardoublequoteclose}\isanewline
%
\isadelimproof
\ \ %
\endisadelimproof
%
\isatagproof
\isacommand{by}\isamarkupfalse%
\ {\isacharparenleft}{\kern0pt}induction\ xs{\isadigit{1}}{\isacharparenright}{\kern0pt}\ {\isacharparenleft}{\kern0pt}auto\ simp{\isacharcolon}{\kern0pt}\ image{\isacharunderscore}{\kern0pt}def\ strict{\isacharunderscore}{\kern0pt}sorted{\isacharunderscore}{\kern0pt}app{\isacharunderscore}{\kern0pt}iff{\isacharbrackleft}{\kern0pt}symmetric{\isacharbrackright}{\kern0pt}{\isacharparenright}{\kern0pt}%
\endisatagproof
{\isafoldproof}%
%
\isadelimproof
\isanewline
%
\endisadelimproof
\isanewline
\isacommand{lemma}\isamarkupfalse%
\ empty{\isacharunderscore}{\kern0pt}eff{\isacharunderscore}{\kern0pt}non{\isacharunderscore}{\kern0pt}intrf{\isacharcolon}{\kern0pt}\isanewline
\ \ {\isachardoublequoteopen}{\isasymlbrakk}{\isasymAnd}a{\isachardot}{\kern0pt}\ a\ {\isasymin}\ A\ {\isasymLongrightarrow}\ effect\ a\ {\isacharequal}{\kern0pt}\ Effect\ {\isacharbrackleft}{\kern0pt}{\isacharbrackright}{\kern0pt}\ {\isacharbrackleft}{\kern0pt}{\isacharbrackright}{\kern0pt}{\isasymrbrakk}\ {\isasymLongrightarrow}\ pairwise\ acts{\isacharunderscore}{\kern0pt}non{\isacharunderscore}{\kern0pt}intrf\ A{\isachardoublequoteclose}\isanewline
%
\isadelimproof
\ \ %
\endisadelimproof
%
\isatagproof
\isacommand{by}\isamarkupfalse%
\ {\isacharparenleft}{\kern0pt}auto\ simp{\isacharcolon}{\kern0pt}\ pairwise{\isacharunderscore}{\kern0pt}def\ acts{\isacharunderscore}{\kern0pt}non{\isacharunderscore}{\kern0pt}intrf{\isacharunderscore}{\kern0pt}def{\isacharparenright}{\kern0pt}%
\endisatagproof
{\isafoldproof}%
%
\isadelimproof
\isanewline
%
\endisadelimproof
\isanewline
\isacommand{lemma}\isamarkupfalse%
\ filter{\isacharunderscore}{\kern0pt}eq{\isacharunderscore}{\kern0pt}consD{\isacharcolon}{\kern0pt}\ {\isachardoublequoteopen}x\ {\isacharhash}{\kern0pt}\ l{\isacharprime}{\kern0pt}\ {\isacharequal}{\kern0pt}\ filter\ Q\ l\ {\isasymLongrightarrow}\ x{\isasymin}set\ l\ {\isasymand}\ Q\ x{\isachardoublequoteclose}\isanewline
%
\isadelimproof
\ \ %
\endisadelimproof
%
\isatagproof
\isacommand{by}\isamarkupfalse%
\ {\isacharparenleft}{\kern0pt}induction\ l{\isacharparenright}{\kern0pt}\ {\isacharparenleft}{\kern0pt}auto\ split{\isacharcolon}{\kern0pt}\ if{\isacharunderscore}{\kern0pt}splits{\isacharparenright}{\kern0pt}%
\endisatagproof
{\isafoldproof}%
%
\isadelimproof
\isanewline
%
\endisadelimproof
\isanewline
\isacommand{lemma}\isamarkupfalse%
\ valid{\isacharunderscore}{\kern0pt}hap{\isacharunderscore}{\kern0pt}seq{\isacharunderscore}{\kern0pt}construct{\isacharcolon}{\kern0pt}\isanewline
\ \ \isakeyword{assumes}\ {\isachardoublequoteopen}{\isasymAnd}t{\isachardot}{\kern0pt}\ t\ {\isasymin}\ htps\ {\isasymLongrightarrow}\ {\isasymexists}A{\isachardot}{\kern0pt}\ {\isacharparenleft}{\kern0pt}t{\isacharcomma}{\kern0pt}\ A{\isacharparenright}{\kern0pt}\ {\isasymin}\ set\ hs{\isachardoublequoteclose}\isanewline
\ \ \ \ \isakeyword{and}\ {\isachardoublequoteopen}{\isasymAnd}t{\isachardot}{\kern0pt}\ t\ {\isasymin}\ htps\ {\isasymLongrightarrow}\ {\isasymexists}A{\isachardot}{\kern0pt}\ {\isacharparenleft}{\kern0pt}t{\isacharcomma}{\kern0pt}\ A{\isacharparenright}{\kern0pt}\ {\isasymin}\ set\ hs{\isacharprime}{\kern0pt}{\isachardoublequoteclose}\isanewline
\ \ \ \ \isakeyword{and}\ {\isachardoublequoteopen}strict{\isacharunderscore}{\kern0pt}sorted\ {\isacharparenleft}{\kern0pt}map\ fst\ hs{\isacharparenright}{\kern0pt}{\isachardoublequoteclose}\isanewline
\ \ \ \ \isakeyword{and}\ {\isachardoublequoteopen}strict{\isacharunderscore}{\kern0pt}sorted\ {\isacharparenleft}{\kern0pt}map\ fst\ hs{\isacharprime}{\kern0pt}{\isacharparenright}{\kern0pt}{\isachardoublequoteclose}\isanewline
\ \ \ \ \isakeyword{and}\ {\isachardoublequoteopen}{\isasymAnd}t\ A\ A{\isacharprime}{\kern0pt}{\isachardot}{\kern0pt}\ {\isasymlbrakk}t\ {\isasymin}\ htps{\isacharsemicolon}{\kern0pt}\ {\isacharparenleft}{\kern0pt}t{\isacharcomma}{\kern0pt}\ A{\isacharparenright}{\kern0pt}\ {\isasymin}\ set\ hs{\isacharsemicolon}{\kern0pt}\ {\isacharparenleft}{\kern0pt}t{\isacharcomma}{\kern0pt}\ A{\isacharprime}{\kern0pt}{\isacharparenright}{\kern0pt}\ {\isasymin}\ set\ hs{\isacharprime}{\kern0pt}{\isasymrbrakk}\ {\isasymLongrightarrow}\ set\ A\ {\isacharequal}{\kern0pt}\ set\ A{\isacharprime}{\kern0pt}{\isachardoublequoteclose}\isanewline
\ \ \ \ \isakeyword{and}\ {\isachardoublequoteopen}{\isasymAnd}t\ A\ A{\isacharprime}{\kern0pt}\ a{\isachardot}{\kern0pt}\ {\isasymlbrakk}t\ {\isasymnotin}\ htps{\isacharsemicolon}{\kern0pt}\ {\isacharparenleft}{\kern0pt}t{\isacharcomma}{\kern0pt}\ A{\isacharparenright}{\kern0pt}\ {\isasymin}\ set\ hs{\isacharsemicolon}{\kern0pt}\ a\ {\isasymin}\ set\ A{\isasymrbrakk}\ {\isasymLongrightarrow}\ effect\ a\ {\isacharequal}{\kern0pt}\ Effect\ {\isacharbrackleft}{\kern0pt}{\isacharbrackright}{\kern0pt}\ {\isacharbrackleft}{\kern0pt}{\isacharbrackright}{\kern0pt}{\isachardoublequoteclose}\isanewline
\ \ \ \ \isakeyword{and}\ {\isachardoublequoteopen}{\isasymAnd}t\ A\ A{\isacharprime}{\kern0pt}\ a{\isachardot}{\kern0pt}\ {\isasymlbrakk}t\ {\isasymnotin}\ htps{\isacharsemicolon}{\kern0pt}\ {\isacharparenleft}{\kern0pt}t{\isacharcomma}{\kern0pt}\ A{\isacharparenright}{\kern0pt}\ {\isasymin}\ set\ hs{\isacharprime}{\kern0pt}{\isacharsemicolon}{\kern0pt}\ a\ {\isasymin}\ set\ A{\isasymrbrakk}\ {\isasymLongrightarrow}\ effect\ a\ {\isacharequal}{\kern0pt}\ Effect\ {\isacharbrackleft}{\kern0pt}{\isacharbrackright}{\kern0pt}\ {\isacharbrackleft}{\kern0pt}{\isacharbrackright}{\kern0pt}{\isachardoublequoteclose}\isanewline
\ \ \ \ \isakeyword{and}\ {\isachardoublequoteopen}{\isasymAnd}t\ A\ t\isactrlsub {\isadigit{1}}\ t\isactrlsub {\isadigit{2}}\ a{\isachardot}{\kern0pt}\isanewline
\ \ \ \ \ \ \ \ \ \ \ \ {\isasymlbrakk}t\ {\isasymnotin}\ htps{\isacharsemicolon}{\kern0pt}\ t\isactrlsub {\isadigit{1}}\ {\isasymin}\ htps{\isacharsemicolon}{\kern0pt}\ t\isactrlsub {\isadigit{2}}\ {\isasymin}\ htps{\isacharsemicolon}{\kern0pt}\ t\ {\isasymin}\ {\isacharbraceleft}{\kern0pt}t\isactrlsub {\isadigit{1}}{\isacharless}{\kern0pt}{\isachardot}{\kern0pt}{\isachardot}{\kern0pt}{\isacharless}{\kern0pt}\ t\isactrlsub {\isadigit{2}}{\isacharbraceright}{\kern0pt}{\isacharsemicolon}{\kern0pt}\ {\isacharparenleft}{\kern0pt}t{\isacharcomma}{\kern0pt}\ A{\isacharparenright}{\kern0pt}\ {\isasymin}\ set\ hs{\isacharsemicolon}{\kern0pt}\ a\ {\isasymin}\ set\ A{\isasymrbrakk}\ {\isasymLongrightarrow}\ \isanewline
\ \ \ \ \ \ \ \ \ \ \ \ \ \ {\isasymexists}t{\isacharprime}{\kern0pt}{\isasymin}{\isacharbraceleft}{\kern0pt}t\isactrlsub {\isadigit{1}}{\isacharless}{\kern0pt}{\isachardot}{\kern0pt}{\isachardot}{\kern0pt}{\isacharless}{\kern0pt}\ t\isactrlsub {\isadigit{2}}{\isacharbraceright}{\kern0pt}{\isachardot}{\kern0pt}{\isasymexists}A{\isacharprime}{\kern0pt}{\isachardot}{\kern0pt}\ {\isacharparenleft}{\kern0pt}t{\isacharprime}{\kern0pt}{\isacharcomma}{\kern0pt}A{\isacharprime}{\kern0pt}{\isacharparenright}{\kern0pt}\ {\isasymin}\ set\ hs{\isacharprime}{\kern0pt}\ {\isasymand}\ a\ {\isasymin}\ set\ A{\isacharprime}{\kern0pt}{\isachardoublequoteclose}\isanewline
\ \ \ \ \isakeyword{and}\ {\isachardoublequoteopen}{\isasymAnd}t\ A\ t\isactrlsub {\isadigit{1}}\ t\isactrlsub {\isadigit{2}}\ a{\isachardot}{\kern0pt}\isanewline
\ \ \ \ \ \ \ \ \ \ \ \ {\isasymlbrakk}t\ {\isasymnotin}\ htps{\isacharsemicolon}{\kern0pt}\ t\isactrlsub {\isadigit{1}}\ {\isasymin}\ htps{\isacharsemicolon}{\kern0pt}\ t\isactrlsub {\isadigit{2}}\ {\isasymin}\ htps{\isacharsemicolon}{\kern0pt}\ t\ {\isasymin}\ {\isacharbraceleft}{\kern0pt}t\isactrlsub {\isadigit{1}}{\isacharless}{\kern0pt}{\isachardot}{\kern0pt}{\isachardot}{\kern0pt}{\isacharless}{\kern0pt}\ t\isactrlsub {\isadigit{2}}{\isacharbraceright}{\kern0pt}{\isacharsemicolon}{\kern0pt}\ {\isacharparenleft}{\kern0pt}t{\isacharcomma}{\kern0pt}\ A{\isacharparenright}{\kern0pt}\ {\isasymin}\ set\ hs{\isacharprime}{\kern0pt}{\isacharsemicolon}{\kern0pt}\ a\ {\isasymin}\ set\ A{\isasymrbrakk}\ {\isasymLongrightarrow}\ \isanewline
\ \ \ \ \ \ \ \ \ \ \ \ \ \ {\isasymexists}t{\isacharprime}{\kern0pt}{\isasymin}{\isacharbraceleft}{\kern0pt}t\isactrlsub {\isadigit{1}}{\isacharless}{\kern0pt}{\isachardot}{\kern0pt}{\isachardot}{\kern0pt}{\isacharless}{\kern0pt}\ t\isactrlsub {\isadigit{2}}{\isacharbraceright}{\kern0pt}{\isachardot}{\kern0pt}{\isasymexists}A{\isacharprime}{\kern0pt}{\isachardot}{\kern0pt}\ {\isacharparenleft}{\kern0pt}t{\isacharprime}{\kern0pt}{\isacharcomma}{\kern0pt}A{\isacharprime}{\kern0pt}{\isacharparenright}{\kern0pt}\ {\isasymin}\ set\ hs\ {\isasymand}\ a\ {\isasymin}\ set\ A{\isacharprime}{\kern0pt}{\isachardoublequoteclose}\isanewline
\ \ \ \ \isakeyword{and}\ {\isachardoublequoteopen}valid{\isacharunderscore}{\kern0pt}happ{\isacharunderscore}{\kern0pt}seq\ M\ hs\ M{\isacharprime}{\kern0pt}{\isachardoublequoteclose}\isanewline
\ \ \ \ \isakeyword{and}\ {\isachardoublequoteopen}hs\ {\isasymnoteq}\ {\isacharbrackleft}{\kern0pt}{\isacharbrackright}{\kern0pt}\ {\isasymLongrightarrow}\ {\isacharparenleft}{\kern0pt}fst\ {\isacharparenleft}{\kern0pt}hd\ hs{\isacharparenright}{\kern0pt}{\isacharparenright}{\kern0pt}\ {\isasymin}\ htps{\isachardoublequoteclose}\isanewline
\ \ \ \ \isakeyword{and}\ {\isachardoublequoteopen}hs{\isacharprime}{\kern0pt}\ {\isasymnoteq}\ {\isacharbrackleft}{\kern0pt}{\isacharbrackright}{\kern0pt}\ {\isasymLongrightarrow}\ {\isacharparenleft}{\kern0pt}fst\ {\isacharparenleft}{\kern0pt}hd\ hs{\isacharprime}{\kern0pt}{\isacharparenright}{\kern0pt}{\isacharparenright}{\kern0pt}\ {\isasymin}\ htps{\isachardoublequoteclose}\isanewline
\ \ \ \ \isakeyword{and}\ {\isachardoublequoteopen}hs\ {\isasymnoteq}\ {\isacharbrackleft}{\kern0pt}{\isacharbrackright}{\kern0pt}\ {\isasymLongrightarrow}\ {\isacharparenleft}{\kern0pt}fst\ {\isacharparenleft}{\kern0pt}last\ hs{\isacharparenright}{\kern0pt}{\isacharparenright}{\kern0pt}\ {\isasymin}\ htps{\isachardoublequoteclose}\isanewline
\ \ \ \ \isakeyword{and}\ {\isachardoublequoteopen}hs{\isacharprime}{\kern0pt}\ {\isasymnoteq}\ {\isacharbrackleft}{\kern0pt}{\isacharbrackright}{\kern0pt}\ {\isasymLongrightarrow}\ {\isacharparenleft}{\kern0pt}fst\ {\isacharparenleft}{\kern0pt}last\ hs{\isacharprime}{\kern0pt}{\isacharparenright}{\kern0pt}{\isacharparenright}{\kern0pt}\ {\isasymin}\ htps{\isachardoublequoteclose}\isanewline
\ \ \isakeyword{shows}\ {\isachardoublequoteopen}valid{\isacharunderscore}{\kern0pt}happ{\isacharunderscore}{\kern0pt}seq\ M\ hs{\isacharprime}{\kern0pt}\ M{\isacharprime}{\kern0pt}{\isachardoublequoteclose}\isanewline
%
\isadelimproof
\ \ %
\endisadelimproof
%
\isatagproof
\isacommand{using}\isamarkupfalse%
\ assms\isanewline
\isacommand{proof}\isamarkupfalse%
{\isacharparenleft}{\kern0pt}induction\ {\isachardoublequoteopen}filter\ {\isacharparenleft}{\kern0pt}{\isasymlambda}h{\isachardot}{\kern0pt}\ fst\ h\ {\isasymin}\ htps{\isacharparenright}{\kern0pt}\ hs{\isachardoublequoteclose}\ {\isachardoublequoteopen}filter\ {\isacharparenleft}{\kern0pt}{\isasymlambda}h{\isachardot}{\kern0pt}\ fst\ h\ {\isasymin}\ htps{\isacharparenright}{\kern0pt}\ hs{\isacharprime}{\kern0pt}{\isachardoublequoteclose}\isanewline
\ \ \ \ arbitrary{\isacharcolon}{\kern0pt}\ M\ hs\ hs{\isacharprime}{\kern0pt}\ htps\ rule{\isacharcolon}{\kern0pt}\ consume{\isadigit{2}}{\isachardot}{\kern0pt}induct{\isacharparenright}{\kern0pt}\isanewline
\ \ \isacommand{case}\isamarkupfalse%
\ {\isacharparenleft}{\kern0pt}{\isadigit{2}}\ h\ xs\ h{\isacharprime}{\kern0pt}\ ys\ hs\ hs{\isacharprime}{\kern0pt}\ htps\ M{\isacharparenright}{\kern0pt}\isanewline
\ \ \isacommand{then}\isamarkupfalse%
\ \isacommand{obtain}\isamarkupfalse%
\ hs{\isadigit{1}}\ hs{\isadigit{2}}\ hs{\isadigit{1}}{\isacharprime}{\kern0pt}\ hs{\isadigit{2}}{\isacharprime}{\kern0pt}\isanewline
\ \ \ \ \isakeyword{where}\ {\isachardoublequoteopen}hs\ {\isacharequal}{\kern0pt}\ hs{\isadigit{1}}\ {\isacharat}{\kern0pt}\ h\ {\isacharhash}{\kern0pt}\ hs{\isadigit{2}}{\isachardoublequoteclose}\ {\isachardoublequoteopen}{\isacharparenleft}{\kern0pt}{\isasymforall}h{\isasymin}set\ hs{\isadigit{1}}{\isachardot}{\kern0pt}\ fst\ h\ {\isasymnotin}\ htps{\isacharparenright}{\kern0pt}{\isachardoublequoteclose}\isanewline
\ \ \ \ \ \ {\isachardoublequoteopen}fst\ h\ {\isasymin}\ htps{\isachardoublequoteclose}\ {\isachardoublequoteopen}xs\ {\isacharequal}{\kern0pt}\ filter\ {\isacharparenleft}{\kern0pt}{\isasymlambda}h{\isachardot}{\kern0pt}\ fst\ h\ {\isasymin}\ htps{\isacharparenright}{\kern0pt}\ hs{\isadigit{2}}{\isachardoublequoteclose}\isanewline
\isanewline
\ \ \ \ \ \ {\isachardoublequoteopen}hs{\isacharprime}{\kern0pt}\ {\isacharequal}{\kern0pt}\ hs{\isadigit{1}}{\isacharprime}{\kern0pt}\ {\isacharat}{\kern0pt}\ h{\isacharprime}{\kern0pt}\ {\isacharhash}{\kern0pt}\ hs{\isadigit{2}}{\isacharprime}{\kern0pt}{\isachardoublequoteclose}\ {\isachardoublequoteopen}{\isacharparenleft}{\kern0pt}{\isasymforall}h{\isasymin}set\ hs{\isadigit{1}}{\isacharprime}{\kern0pt}{\isachardot}{\kern0pt}\ fst\ h\ {\isasymnotin}\ htps{\isacharparenright}{\kern0pt}{\isachardoublequoteclose}\isanewline
\ \ \ \ \ \ {\isachardoublequoteopen}fst\ h{\isacharprime}{\kern0pt}\ {\isasymin}\ htps{\isachardoublequoteclose}\ {\isachardoublequoteopen}ys\ {\isacharequal}{\kern0pt}\ filter\ {\isacharparenleft}{\kern0pt}{\isasymlambda}h{\isachardot}{\kern0pt}\ fst\ h\ {\isasymin}\ htps{\isacharparenright}{\kern0pt}\ hs{\isadigit{2}}{\isacharprime}{\kern0pt}{\isachardoublequoteclose}\isanewline
\ \ \ \ \isacommand{by}\isamarkupfalse%
\ {\isacharparenleft}{\kern0pt}auto\ simp{\isacharcolon}{\kern0pt}\ Cons{\isacharunderscore}{\kern0pt}eq{\isacharunderscore}{\kern0pt}filter{\isacharunderscore}{\kern0pt}iff{\isacharparenright}{\kern0pt}\isanewline
\ \ \isacommand{hence}\isamarkupfalse%
\ {\isacharbrackleft}{\kern0pt}simp{\isacharbrackright}{\kern0pt}{\isacharcolon}{\kern0pt}\ {\isachardoublequoteopen}hs{\isadigit{1}}\ {\isacharequal}{\kern0pt}\ {\isacharbrackleft}{\kern0pt}{\isacharbrackright}{\kern0pt}{\isachardoublequoteclose}\ {\isachardoublequoteopen}hs{\isadigit{1}}{\isacharprime}{\kern0pt}\ {\isacharequal}{\kern0pt}\ {\isacharbrackleft}{\kern0pt}{\isacharbrackright}{\kern0pt}{\isachardoublequoteclose}\isanewline
\ \ \ \ \isacommand{using}\isamarkupfalse%
\ {\isacartoucheopen}hs{\isacharprime}{\kern0pt}\ {\isasymnoteq}\ {\isacharbrackleft}{\kern0pt}{\isacharbrackright}{\kern0pt}\ {\isasymLongrightarrow}\ fst\ {\isacharparenleft}{\kern0pt}hd\ hs{\isacharprime}{\kern0pt}{\isacharparenright}{\kern0pt}\ {\isasymin}\ htps{\isacartoucheclose}\ {\isacartoucheopen}hs\ {\isasymnoteq}\ {\isacharbrackleft}{\kern0pt}{\isacharbrackright}{\kern0pt}\ {\isasymLongrightarrow}\ fst\ {\isacharparenleft}{\kern0pt}hd\ hs{\isacharparenright}{\kern0pt}\ {\isasymin}\ htps{\isacartoucheclose}\isanewline
\ \ \ \ \ \isacommand{apply}\isamarkupfalse%
\ {\isacharminus}{\kern0pt}\isanewline
\ \ \ \ \isacommand{by}\isamarkupfalse%
\ {\isacharparenleft}{\kern0pt}rule\ list{\isachardot}{\kern0pt}exhaust{\isacharcomma}{\kern0pt}\ auto{\isacharparenright}{\kern0pt}{\isacharplus}{\kern0pt}\isanewline
\isanewline
\ \ \isacommand{note}\isamarkupfalse%
\ {\isacharbrackleft}{\kern0pt}simp{\isacharbrackright}{\kern0pt}\ {\isacharequal}{\kern0pt}\ {\isacartoucheopen}hs\ {\isacharequal}{\kern0pt}\ hs{\isadigit{1}}\ {\isacharat}{\kern0pt}\ h\ {\isacharhash}{\kern0pt}\ hs{\isadigit{2}}{\isacartoucheclose}\ {\isacartoucheopen}hs{\isacharprime}{\kern0pt}\ {\isacharequal}{\kern0pt}\ hs{\isadigit{1}}{\isacharprime}{\kern0pt}\ {\isacharat}{\kern0pt}\ h{\isacharprime}{\kern0pt}\ {\isacharhash}{\kern0pt}\ hs{\isadigit{2}}{\isacharprime}{\kern0pt}{\isacartoucheclose}\isanewline
\isanewline
\ \ \isacommand{obtain}\isamarkupfalse%
\ t\ t{\isacharprime}{\kern0pt}\ A\ A{\isacharprime}{\kern0pt}\ \isakeyword{where}\ {\isacharasterisk}{\kern0pt}{\isacharbrackleft}{\kern0pt}simp{\isacharbrackright}{\kern0pt}{\isacharcolon}{\kern0pt}\ {\isachardoublequoteopen}h\ {\isacharequal}{\kern0pt}\ {\isacharparenleft}{\kern0pt}t{\isacharcomma}{\kern0pt}\ A{\isacharparenright}{\kern0pt}{\isachardoublequoteclose}\ {\isachardoublequoteopen}h{\isacharprime}{\kern0pt}\ {\isacharequal}{\kern0pt}\ {\isacharparenleft}{\kern0pt}t{\isacharprime}{\kern0pt}{\isacharcomma}{\kern0pt}A{\isacharprime}{\kern0pt}{\isacharparenright}{\kern0pt}{\isachardoublequoteclose}\isanewline
\ \ \ \ \isacommand{by}\isamarkupfalse%
\ {\isacharparenleft}{\kern0pt}cases\ h{\isacharcomma}{\kern0pt}\ cases\ h{\isacharprime}{\kern0pt}{\isacharcomma}{\kern0pt}\ auto{\isacharparenright}{\kern0pt}\isanewline
\ \ \isacommand{hence}\isamarkupfalse%
\ in{\isacharunderscore}{\kern0pt}htp{\isacharcolon}{\kern0pt}\ {\isachardoublequoteopen}t{\isasymin}htps{\isachardoublequoteclose}\ {\isachardoublequoteopen}t{\isacharprime}{\kern0pt}{\isasymin}htps{\isachardoublequoteclose}\isanewline
\ \ \ \ \isacommand{using}\isamarkupfalse%
\ {\isacartoucheopen}fst\ h\ {\isasymin}\ htps{\isacartoucheclose}\ {\isacartoucheopen}fst\ h{\isacharprime}{\kern0pt}\ {\isasymin}\ htps{\isacartoucheclose}\isanewline
\ \ \ \ \isacommand{by}\isamarkupfalse%
\ auto\isanewline
\ \ \isacommand{moreover}\isamarkupfalse%
\ \isacommand{have}\isamarkupfalse%
\ {\isachardoublequoteopen}{\isasymAnd}t{\isachardot}{\kern0pt}\ t\ {\isasymin}\ htps\ {\isasymLongrightarrow}\ {\isasymexists}A{\isachardot}{\kern0pt}\ {\isacharparenleft}{\kern0pt}t{\isacharcomma}{\kern0pt}\ A{\isacharparenright}{\kern0pt}\ {\isasymin}\ set\ {\isacharparenleft}{\kern0pt}h\ {\isacharhash}{\kern0pt}\ hs{\isadigit{2}}{\isacharparenright}{\kern0pt}{\isachardoublequoteclose}\isanewline
\ \ \ \ \ \ \ \ \ \ \ \ \ \ \ \ {\isachardoublequoteopen}{\isasymAnd}t{\isachardot}{\kern0pt}\ t\ {\isasymin}\ htps\ {\isasymLongrightarrow}\ {\isasymexists}A{\isachardot}{\kern0pt}\ {\isacharparenleft}{\kern0pt}t{\isacharcomma}{\kern0pt}\ A{\isacharparenright}{\kern0pt}\ {\isasymin}\ set\ {\isacharparenleft}{\kern0pt}h{\isacharprime}{\kern0pt}\ {\isacharhash}{\kern0pt}\ hs{\isadigit{2}}{\isacharprime}{\kern0pt}{\isacharparenright}{\kern0pt}{\isachardoublequoteclose}\isanewline
\ \ \ \ \isacommand{using}\isamarkupfalse%
\ {\isadigit{2}}{\isacharparenleft}{\kern0pt}{\isadigit{4}}{\isacharcomma}{\kern0pt}{\isadigit{5}}{\isacharparenright}{\kern0pt}\ {\isacartoucheopen}{\isacharparenleft}{\kern0pt}{\isasymforall}h{\isasymin}set\ hs{\isadigit{1}}{\isachardot}{\kern0pt}\ fst\ h\ {\isasymnotin}\ htps{\isacharparenright}{\kern0pt}{\isacartoucheclose}\ {\isacartoucheopen}{\isacharparenleft}{\kern0pt}{\isasymforall}h{\isasymin}set\ hs{\isadigit{1}}{\isacharprime}{\kern0pt}{\isachardot}{\kern0pt}\ fst\ h\ {\isasymnotin}\ htps{\isacharparenright}{\kern0pt}{\isacartoucheclose}\isanewline
\ \ \ \ \isacommand{by}\isamarkupfalse%
\ fastforce{\isacharplus}{\kern0pt}\isanewline
\ \ \isacommand{moreover}\isamarkupfalse%
\ \isacommand{have}\isamarkupfalse%
\ {\isachardoublequoteopen}strict{\isacharunderscore}{\kern0pt}sorted\ {\isacharparenleft}{\kern0pt}map\ fst\ {\isacharparenleft}{\kern0pt}h{\isacharhash}{\kern0pt}hs{\isadigit{2}}{\isacharparenright}{\kern0pt}{\isacharparenright}{\kern0pt}{\isachardoublequoteclose}\ {\isachardoublequoteopen}strict{\isacharunderscore}{\kern0pt}sorted\ {\isacharparenleft}{\kern0pt}map\ fst\ {\isacharparenleft}{\kern0pt}h{\isacharprime}{\kern0pt}{\isacharhash}{\kern0pt}hs{\isadigit{2}}{\isacharprime}{\kern0pt}{\isacharparenright}{\kern0pt}{\isacharparenright}{\kern0pt}{\isachardoublequoteclose}\isanewline
\ \ \ \ \isacommand{using}\isamarkupfalse%
\ {\isadigit{2}}{\isacharparenleft}{\kern0pt}{\isadigit{6}}{\isacharcomma}{\kern0pt}{\isadigit{7}}{\isacharparenright}{\kern0pt}\ \isanewline
\ \ \ \ \isacommand{by}\isamarkupfalse%
\ {\isacharparenleft}{\kern0pt}auto\ simp{\isacharcolon}{\kern0pt}\ strict{\isacharunderscore}{\kern0pt}sorted{\isacharunderscore}{\kern0pt}app{\isacharunderscore}{\kern0pt}iff{\isacharbrackleft}{\kern0pt}symmetric{\isacharbrackright}{\kern0pt}{\isacharparenright}{\kern0pt}\isanewline
\ \ \isacommand{ultimately}\isamarkupfalse%
\ \isacommand{have}\isamarkupfalse%
\ {\isacharbrackleft}{\kern0pt}simp{\isacharbrackright}{\kern0pt}{\isacharcolon}{\kern0pt}\ {\isachardoublequoteopen}t\ {\isacharequal}{\kern0pt}\ t{\isacharprime}{\kern0pt}{\isachardoublequoteclose}\isanewline
\ \ \ \ \isacommand{by}\isamarkupfalse%
\ force\isanewline
\ \ \isacommand{hence}\isamarkupfalse%
\ A{\isacharunderscore}{\kern0pt}eq{\isacharcolon}{\kern0pt}\ {\isachardoublequoteopen}set\ A\ {\isacharequal}{\kern0pt}\ set\ A{\isacharprime}{\kern0pt}{\isachardoublequoteclose}\isanewline
\ \ \ \ \isacommand{using}\isamarkupfalse%
\ {\isadigit{2}}{\isacharparenleft}{\kern0pt}{\isadigit{8}}{\isacharparenright}{\kern0pt}\ in{\isacharunderscore}{\kern0pt}htp\isanewline
\ \ \ \ \isacommand{by}\isamarkupfalse%
\ auto\isanewline
\ \ \isacommand{moreover}\isamarkupfalse%
\ \isacommand{have}\isamarkupfalse%
\ {\isachardoublequoteopen}valid{\isacharunderscore}{\kern0pt}happ{\isacharunderscore}{\kern0pt}seq\ M\ {\isacharparenleft}{\kern0pt}h\ {\isacharhash}{\kern0pt}\ hs{\isadigit{2}}{\isacharparenright}{\kern0pt}\ M{\isacharprime}{\kern0pt}{\isachardoublequoteclose}\isanewline
\ \ \ \ \isacommand{using}\isamarkupfalse%
\ {\isadigit{2}}{\isacharparenleft}{\kern0pt}{\isadigit{9}}{\isacharparenright}{\kern0pt}\ {\isacartoucheopen}valid{\isacharunderscore}{\kern0pt}happ{\isacharunderscore}{\kern0pt}seq\ M\ hs\ M{\isacharprime}{\kern0pt}{\isacartoucheclose}\ {\isacartoucheopen}{\isacharparenleft}{\kern0pt}{\isasymforall}h{\isasymin}set\ hs{\isadigit{1}}{\isachardot}{\kern0pt}\ fst\ h\ {\isasymnotin}\ htps{\isacharparenright}{\kern0pt}{\isacartoucheclose}\isanewline
\ \ \ \ \isacommand{by}\isamarkupfalse%
\ {\isacharparenleft}{\kern0pt}force\ simp{\isacharcolon}{\kern0pt}\ intro{\isacharbang}{\kern0pt}{\isacharcolon}{\kern0pt}\ valid{\isacharunderscore}{\kern0pt}happ{\isacharunderscore}{\kern0pt}seq{\isacharunderscore}{\kern0pt}append{\isacharunderscore}{\kern0pt}{\isadigit{2}}{\isacharbrackleft}{\kern0pt}\isakeyword{where}\ {\isacharquery}{\kern0pt}hs{\isadigit{1}}{\isachardot}{\kern0pt}{\isadigit{0}}\ {\isacharequal}{\kern0pt}\ hs{\isadigit{1}}{\isacharbrackright}{\kern0pt}{\isacharparenright}{\kern0pt}{\isacharplus}{\kern0pt}\isanewline
\ \ \isacommand{ultimately}\isamarkupfalse%
\ \isacommand{have}\isamarkupfalse%
\ {\isachardoublequoteopen}valid{\isacharunderscore}{\kern0pt}happ{\isacharunderscore}{\kern0pt}seq\ M\ {\isacharbrackleft}{\kern0pt}h{\isacharbrackright}{\kern0pt}\ {\isacharparenleft}{\kern0pt}apply{\isacharunderscore}{\kern0pt}happ\ {\isacharparenleft}{\kern0pt}t{\isacharcomma}{\kern0pt}\ A{\isacharprime}{\kern0pt}{\isacharparenright}{\kern0pt}\ M{\isacharparenright}{\kern0pt}{\isachardoublequoteclose}\isanewline
\ \ \ \ \isacommand{using}\isamarkupfalse%
\ {\isacartoucheopen}h\ {\isacharequal}{\kern0pt}\ {\isacharparenleft}{\kern0pt}t{\isacharcomma}{\kern0pt}A{\isacharparenright}{\kern0pt}{\isacartoucheclose}\ A{\isacharunderscore}{\kern0pt}eq\isanewline
\ \ \ \ \isacommand{by}\isamarkupfalse%
\ auto\isanewline
\ \ \isacommand{hence}\isamarkupfalse%
\ {\isachardoublequoteopen}valid{\isacharunderscore}{\kern0pt}happ{\isacharunderscore}{\kern0pt}seq\ M\ {\isacharbrackleft}{\kern0pt}h{\isacharprime}{\kern0pt}{\isacharbrackright}{\kern0pt}\ {\isacharparenleft}{\kern0pt}apply{\isacharunderscore}{\kern0pt}happ\ {\isacharparenleft}{\kern0pt}t{\isacharcomma}{\kern0pt}\ A{\isacharprime}{\kern0pt}{\isacharparenright}{\kern0pt}\ M{\isacharparenright}{\kern0pt}{\isachardoublequoteclose}\isanewline
\ \ \ \ \isacommand{using}\isamarkupfalse%
\ {\isacartoucheopen}h\ {\isacharequal}{\kern0pt}\ {\isacharparenleft}{\kern0pt}t{\isacharcomma}{\kern0pt}A{\isacharparenright}{\kern0pt}{\isacartoucheclose}\ A{\isacharunderscore}{\kern0pt}eq\isanewline
\ \ \ \ \isacommand{by}\isamarkupfalse%
\ auto\isanewline
\isanewline
\ \ \isacommand{show}\isamarkupfalse%
\ {\isacharquery}{\kern0pt}case\isanewline
\ \ \isacommand{proof}\isamarkupfalse%
{\isacharparenleft}{\kern0pt}cases\ {\isachardoublequoteopen}xs{\isachardoublequoteclose}{\isacharparenright}{\kern0pt}\isanewline
\ \ \ \ \isacommand{case}\isamarkupfalse%
\ Nil\isanewline
\ \ \ \ \isacommand{hence}\isamarkupfalse%
\ {\isachardoublequoteopen}set\ {\isacharparenleft}{\kern0pt}map\ fst\ hs{\isadigit{2}}{\isacharparenright}{\kern0pt}\ {\isasyminter}\ htps\ {\isacharequal}{\kern0pt}\ {\isacharbraceleft}{\kern0pt}{\isacharbraceright}{\kern0pt}{\isachardoublequoteclose}\isanewline
\ \ \ \ \ \ \isacommand{using}\isamarkupfalse%
\ {\isacartoucheopen}xs\ {\isacharequal}{\kern0pt}\ filter\ {\isacharparenleft}{\kern0pt}{\isasymlambda}h{\isachardot}{\kern0pt}\ fst\ h\ {\isasymin}\ htps{\isacharparenright}{\kern0pt}\ hs{\isadigit{2}}{\isacartoucheclose}\isanewline
\ \ \ \ \ \ \isacommand{by}\isamarkupfalse%
\ {\isacharparenleft}{\kern0pt}auto\ simp\ add{\isacharcolon}{\kern0pt}\ filter{\isacharunderscore}{\kern0pt}empty{\isacharunderscore}{\kern0pt}conv{\isacharparenright}{\kern0pt}\isanewline
\ \ \ \ \isacommand{hence}\isamarkupfalse%
\ {\isachardoublequoteopen}M{\isacharprime}{\kern0pt}\ {\isacharequal}{\kern0pt}\ {\isacharparenleft}{\kern0pt}apply{\isacharunderscore}{\kern0pt}happ\ {\isacharparenleft}{\kern0pt}t{\isacharcomma}{\kern0pt}\ A{\isacharprime}{\kern0pt}{\isacharparenright}{\kern0pt}\ M{\isacharparenright}{\kern0pt}{\isachardoublequoteclose}\isanewline
\ \ \ \ \ \ \isacommand{using}\isamarkupfalse%
\ {\isacartoucheopen}valid{\isacharunderscore}{\kern0pt}happ{\isacharunderscore}{\kern0pt}seq\ M\ hs\ M{\isacharprime}{\kern0pt}{\isacartoucheclose}\ {\isacartoucheopen}valid{\isacharunderscore}{\kern0pt}happ{\isacharunderscore}{\kern0pt}seq\ M\ {\isacharbrackleft}{\kern0pt}h{\isacharbrackright}{\kern0pt}\ {\isacharparenleft}{\kern0pt}apply{\isacharunderscore}{\kern0pt}happ\ {\isacharparenleft}{\kern0pt}t{\isacharcomma}{\kern0pt}\ A{\isacharprime}{\kern0pt}{\isacharparenright}{\kern0pt}\ M{\isacharparenright}{\kern0pt}{\isacartoucheclose}\isanewline
\ \ \ \ \ \ \ \ \ \ \ \ \ {\isacartoucheopen}hs\ {\isasymnoteq}\ {\isacharbrackleft}{\kern0pt}{\isacharbrackright}{\kern0pt}\ {\isasymLongrightarrow}\ fst\ {\isacharparenleft}{\kern0pt}last\ hs{\isacharparenright}{\kern0pt}\ {\isasymin}\ htps{\isacartoucheclose}\isanewline
\ \ \ \ \ \ \isacommand{by}\isamarkupfalse%
\ {\isacharparenleft}{\kern0pt}auto\ split{\isacharcolon}{\kern0pt}\ if{\isacharunderscore}{\kern0pt}splits{\isacharparenright}{\kern0pt}\isanewline
\ \ \ \ \isacommand{moreover}\isamarkupfalse%
\ \isacommand{from}\isamarkupfalse%
\ {\isacartoucheopen}set\ {\isacharparenleft}{\kern0pt}map\ fst\ hs{\isadigit{2}}{\isacharparenright}{\kern0pt}\ {\isasyminter}\ htps\ {\isacharequal}{\kern0pt}\ {\isacharbraceleft}{\kern0pt}{\isacharbraceright}{\kern0pt}{\isacartoucheclose}\ \isacommand{have}\isamarkupfalse%
\ {\isachardoublequoteopen}set\ {\isacharparenleft}{\kern0pt}map\ fst\ hs{\isadigit{2}}{\isacharprime}{\kern0pt}{\isacharparenright}{\kern0pt}\ {\isasyminter}\ htps\ {\isacharequal}{\kern0pt}\ {\isacharbraceleft}{\kern0pt}{\isacharbraceright}{\kern0pt}{\isachardoublequoteclose}\isanewline
\ \ \ \ \ \ \isacommand{using}\isamarkupfalse%
\ {\isacartoucheopen}{\isasymAnd}t{\isachardot}{\kern0pt}\ t\ {\isasymin}\ htps\ {\isasymLongrightarrow}\ {\isasymexists}A{\isachardot}{\kern0pt}\ {\isacharparenleft}{\kern0pt}t{\isacharcomma}{\kern0pt}\ A{\isacharparenright}{\kern0pt}\ {\isasymin}\ set\ {\isacharparenleft}{\kern0pt}h\ {\isacharhash}{\kern0pt}\ hs{\isadigit{2}}{\isacharparenright}{\kern0pt}{\isacartoucheclose}\ {\isacartoucheopen}fst\ h{\isacharprime}{\kern0pt}\ {\isasymin}\ htps{\isacartoucheclose}\isanewline
\ \ \ \ \ \ \ \ strict{\isacharunderscore}{\kern0pt}sorted{\isacharunderscore}{\kern0pt}cons{\isacharbrackleft}{\kern0pt}OF\ {\isacartoucheopen}strict{\isacharunderscore}{\kern0pt}sorted\ {\isacharparenleft}{\kern0pt}map\ fst\ {\isacharparenleft}{\kern0pt}h{\isacharprime}{\kern0pt}{\isacharhash}{\kern0pt}hs{\isadigit{2}}{\isacharprime}{\kern0pt}{\isacharparenright}{\kern0pt}{\isacharparenright}{\kern0pt}{\isacartoucheclose}{\isacharbrackright}{\kern0pt}\isanewline
\ \ \ \ \ \ \isacommand{by}\isamarkupfalse%
\ {\isacharparenleft}{\kern0pt}simp\ add{\isacharcolon}{\kern0pt}\ disjoint{\isacharunderscore}{\kern0pt}iff{\isacharunderscore}{\kern0pt}not{\isacharunderscore}{\kern0pt}equal{\isacharparenright}{\kern0pt}\ {\isacharparenleft}{\kern0pt}metis\ fst{\isacharunderscore}{\kern0pt}conv{\isacharparenright}{\kern0pt}\isanewline
\ \ \ \ \isacommand{ultimately}\isamarkupfalse%
\ \isacommand{show}\isamarkupfalse%
\ {\isacharquery}{\kern0pt}thesis\isanewline
\ \ \ \ \ \ \isacommand{using}\isamarkupfalse%
\ {\isacartoucheopen}hs{\isacharprime}{\kern0pt}\ {\isasymnoteq}\ {\isacharbrackleft}{\kern0pt}{\isacharbrackright}{\kern0pt}\ {\isasymLongrightarrow}\ fst\ {\isacharparenleft}{\kern0pt}last\ hs{\isacharprime}{\kern0pt}{\isacharparenright}{\kern0pt}\ {\isasymin}\ htps{\isacartoucheclose}\ {\isacartoucheopen}valid{\isacharunderscore}{\kern0pt}happ{\isacharunderscore}{\kern0pt}seq\ M\ {\isacharbrackleft}{\kern0pt}h{\isacharprime}{\kern0pt}{\isacharbrackright}{\kern0pt}\ {\isacharparenleft}{\kern0pt}apply{\isacharunderscore}{\kern0pt}happ\ {\isacharparenleft}{\kern0pt}t{\isacharcomma}{\kern0pt}\ A{\isacharprime}{\kern0pt}{\isacharparenright}{\kern0pt}\ M{\isacharparenright}{\kern0pt}{\isacartoucheclose}\isanewline
\ \ \ \ \ \ \isacommand{by}\isamarkupfalse%
\ {\isacharparenleft}{\kern0pt}auto\ split{\isacharcolon}{\kern0pt}\ if{\isacharunderscore}{\kern0pt}splits{\isacharparenright}{\kern0pt}\isanewline
\ \ \isacommand{next}\isamarkupfalse%
\isanewline
\ \ \ \ \isacommand{case}\isamarkupfalse%
\ {\isacharparenleft}{\kern0pt}Cons\ h{\isadigit{2}}\ xs{\isacharprime}{\kern0pt}{\isacharparenright}{\kern0pt}\isanewline
\ \ \ \ \isacommand{then}\isamarkupfalse%
\ \isacommand{obtain}\isamarkupfalse%
\ hs{\isadigit{3}}\ hs{\isadigit{4}}\isanewline
\ \ \ \ \ \ \isakeyword{where}\ {\isachardoublequoteopen}hs{\isadigit{2}}\ {\isacharequal}{\kern0pt}\ hs{\isadigit{3}}\ {\isacharat}{\kern0pt}\ h{\isadigit{2}}\ {\isacharhash}{\kern0pt}\ hs{\isadigit{4}}{\isachardoublequoteclose}\ {\isachardoublequoteopen}{\isacharparenleft}{\kern0pt}{\isasymforall}h{\isasymin}set\ hs{\isadigit{3}}{\isachardot}{\kern0pt}\ fst\ h\ {\isasymnotin}\ htps{\isacharparenright}{\kern0pt}{\isachardoublequoteclose}\isanewline
\ \ \ \ \ \ \ \ {\isachardoublequoteopen}fst\ h{\isadigit{2}}\ {\isasymin}\ htps{\isachardoublequoteclose}\ {\isachardoublequoteopen}xs{\isacharprime}{\kern0pt}\ {\isacharequal}{\kern0pt}\ filter\ {\isacharparenleft}{\kern0pt}{\isasymlambda}h{\isachardot}{\kern0pt}\ fst\ h\ {\isasymin}\ htps{\isacharparenright}{\kern0pt}\ hs{\isadigit{4}}{\isachardoublequoteclose}\isanewline
\ \ \ \ \ \ \isacommand{using}\isamarkupfalse%
\ {\isacartoucheopen}xs\ {\isacharequal}{\kern0pt}\ filter\ {\isacharparenleft}{\kern0pt}{\isasymlambda}h{\isachardot}{\kern0pt}\ fst\ h\ {\isasymin}\ htps{\isacharparenright}{\kern0pt}\ hs{\isadigit{2}}{\isacartoucheclose}\isanewline
\ \ \ \ \ \ \isacommand{by}\isamarkupfalse%
\ {\isacharparenleft}{\kern0pt}auto\ simp{\isacharcolon}{\kern0pt}\ filter{\isacharunderscore}{\kern0pt}eq{\isacharunderscore}{\kern0pt}Cons{\isacharunderscore}{\kern0pt}iff{\isacharparenright}{\kern0pt}\isanewline
\ \ \ \ \isacommand{moreover}\isamarkupfalse%
\ \isacommand{hence}\isamarkupfalse%
\ {\isachardoublequoteopen}fst\ h{\isadigit{2}}\ {\isasymnoteq}\ fst\ h{\isachardoublequoteclose}\isanewline
\ \ \ \ \ \ \isacommand{using}\isamarkupfalse%
\ {\isacartoucheopen}strict{\isacharunderscore}{\kern0pt}sorted\ {\isacharparenleft}{\kern0pt}map\ fst\ {\isacharparenleft}{\kern0pt}h\ {\isacharhash}{\kern0pt}\ hs{\isadigit{2}}{\isacharparenright}{\kern0pt}{\isacharparenright}{\kern0pt}{\isacartoucheclose}\isanewline
\ \ \ \ \ \ \isacommand{by}\isamarkupfalse%
\ auto\isanewline
\ \ \ \ \isacommand{hence}\isamarkupfalse%
\ {\isachardoublequoteopen}fst\ h{\isadigit{2}}\ {\isasymnoteq}\ fst\ h{\isacharprime}{\kern0pt}{\isachardoublequoteclose}\isanewline
\ \ \ \ \ \ \isacommand{using}\isamarkupfalse%
\ {\isacartoucheopen}t\ {\isacharequal}{\kern0pt}\ t{\isacharprime}{\kern0pt}{\isacartoucheclose}\isanewline
\ \ \ \ \ \ \isacommand{by}\isamarkupfalse%
\ auto\isanewline
\ \ \ \ \isacommand{ultimately}\isamarkupfalse%
\ \isacommand{obtain}\isamarkupfalse%
\ A{\isadigit{2}}{\isacharprime}{\kern0pt}\ \isakeyword{where}\ {\isachardoublequoteopen}{\isacharparenleft}{\kern0pt}fst\ h{\isadigit{2}}{\isacharcomma}{\kern0pt}\ A{\isadigit{2}}{\isacharprime}{\kern0pt}{\isacharparenright}{\kern0pt}\ {\isasymin}\ set\ hs{\isadigit{2}}{\isacharprime}{\kern0pt}{\isachardoublequoteclose}\isanewline
\ \ \ \ \ \ \isacommand{using}\isamarkupfalse%
\ {\isacartoucheopen}{\isasymAnd}t{\isachardot}{\kern0pt}\ t\ {\isasymin}\ htps\ {\isasymLongrightarrow}\ {\isasymexists}A{\isachardot}{\kern0pt}\ {\isacharparenleft}{\kern0pt}t{\isacharcomma}{\kern0pt}\ A{\isacharparenright}{\kern0pt}\ {\isasymin}\ set\ {\isacharparenleft}{\kern0pt}h{\isacharprime}{\kern0pt}\ {\isacharhash}{\kern0pt}\ hs{\isadigit{2}}{\isacharprime}{\kern0pt}{\isacharparenright}{\kern0pt}{\isacartoucheclose}\isanewline
\ \ \ \ \ \ \isacommand{by}\isamarkupfalse%
\ {\isacharparenleft}{\kern0pt}cases\ h{\isadigit{2}}{\isacharcomma}{\kern0pt}\ cases\ h{\isacharprime}{\kern0pt}{\isacharcomma}{\kern0pt}\ auto\ simp\ add{\isacharcolon}{\kern0pt}\ disjoint{\isacharunderscore}{\kern0pt}iff{\isacharunderscore}{\kern0pt}not{\isacharunderscore}{\kern0pt}equal{\isacharparenright}{\kern0pt}\isanewline
\ \ \ \ \isacommand{then}\isamarkupfalse%
\ \isacommand{obtain}\isamarkupfalse%
\ hs{\isadigit{3}}{\isacharprime}{\kern0pt}\ hs{\isadigit{4}}{\isacharprime}{\kern0pt}\ \isakeyword{where}\ {\isacharbrackleft}{\kern0pt}simp{\isacharbrackright}{\kern0pt}{\isacharcolon}{\kern0pt}\ {\isachardoublequoteopen}hs{\isadigit{2}}{\isacharprime}{\kern0pt}\ {\isacharequal}{\kern0pt}\ hs{\isadigit{3}}{\isacharprime}{\kern0pt}\ {\isacharat}{\kern0pt}\ {\isacharparenleft}{\kern0pt}fst\ h{\isadigit{2}}{\isacharcomma}{\kern0pt}\ A{\isadigit{2}}{\isacharprime}{\kern0pt}{\isacharparenright}{\kern0pt}\ {\isacharhash}{\kern0pt}\ hs{\isadigit{4}}{\isacharprime}{\kern0pt}{\isachardoublequoteclose}\isanewline
\ \ \ \ \ \ \isacommand{by}\isamarkupfalse%
\ {\isacharparenleft}{\kern0pt}auto\ dest{\isacharcolon}{\kern0pt}\ split{\isacharunderscore}{\kern0pt}list{\isacharunderscore}{\kern0pt}first{\isacharparenright}{\kern0pt}\isanewline
\isanewline
\ \ \ \ \isacommand{note}\isamarkupfalse%
{\isacharbrackleft}{\kern0pt}simp{\isacharbrackright}{\kern0pt}\ {\isacharequal}{\kern0pt}\ {\isacartoucheopen}hs{\isadigit{2}}\ {\isacharequal}{\kern0pt}\ hs{\isadigit{3}}\ {\isacharat}{\kern0pt}\ h{\isadigit{2}}\ {\isacharhash}{\kern0pt}\ hs{\isadigit{4}}{\isacartoucheclose}\ strict{\isacharunderscore}{\kern0pt}sorted{\isacharunderscore}{\kern0pt}app{\isacharunderscore}{\kern0pt}iff{\isacharbrackleft}{\kern0pt}symmetric{\isacharbrackright}{\kern0pt}\isanewline
\isanewline
\ \ \ \ \isacommand{have}\isamarkupfalse%
\ fst{\isacharunderscore}{\kern0pt}h{\isadigit{3}}{\isacharprime}{\kern0pt}{\isacharunderscore}{\kern0pt}n{\isacharunderscore}{\kern0pt}htp{\isacharcolon}{\kern0pt}\ {\isachardoublequoteopen}fst\ h{\isadigit{3}}{\isacharprime}{\kern0pt}\ {\isasymnotin}\ htps{\isachardoublequoteclose}\ \isakeyword{if}\ {\isachardoublequoteopen}h{\isadigit{3}}{\isacharprime}{\kern0pt}{\isasymin}set\ hs{\isadigit{3}}{\isacharprime}{\kern0pt}{\isachardoublequoteclose}\ \isakeyword{for}\ h{\isadigit{3}}{\isacharprime}{\kern0pt}\isanewline
\ \ \ \ \isacommand{proof}\isamarkupfalse%
{\isacharparenleft}{\kern0pt}rule\ ccontr{\isacharcomma}{\kern0pt}\ safe{\isacharparenright}{\kern0pt}\isanewline
\ \ \ \ \ \ \isacommand{assume}\isamarkupfalse%
\ {\isachardoublequoteopen}fst\ h{\isadigit{3}}{\isacharprime}{\kern0pt}\ {\isasymin}\ htps{\isachardoublequoteclose}\isanewline
\ \ \ \ \ \ \isacommand{hence}\isamarkupfalse%
\ {\isachardoublequoteopen}fst\ h{\isadigit{3}}{\isacharprime}{\kern0pt}\ {\isacharless}{\kern0pt}\ fst\ h{\isadigit{2}}{\isachardoublequoteclose}\isanewline
\ \ \ \ \ \ \ \ \isacommand{using}\isamarkupfalse%
\ that\ {\isacartoucheopen}strict{\isacharunderscore}{\kern0pt}sorted\ {\isacharparenleft}{\kern0pt}map\ fst\ hs{\isacharprime}{\kern0pt}{\isacharparenright}{\kern0pt}{\isacartoucheclose}\isanewline
\ \ \ \ \ \ \ \ \isacommand{by}\isamarkupfalse%
\ auto\isanewline
\ \ \ \ \ \ \isacommand{moreover}\isamarkupfalse%
\ \isacommand{have}\isamarkupfalse%
\ {\isachardoublequoteopen}x\ {\isasymin}\ htps\ {\isasymLongrightarrow}\ {\isasymexists}y{\isasymin}set\ {\isacharparenleft}{\kern0pt}h\ {\isacharhash}{\kern0pt}\ hs{\isadigit{3}}\ {\isacharat}{\kern0pt}\ h{\isadigit{2}}\ {\isacharhash}{\kern0pt}\ hs{\isadigit{4}}{\isacharparenright}{\kern0pt}{\isachardot}{\kern0pt}\ fst\ y\ {\isacharequal}{\kern0pt}\ x{\isachardoublequoteclose}\ \isakeyword{for}\ x\isanewline
\ \ \ \ \ \ \ \ \isacommand{using}\isamarkupfalse%
\ {\isacartoucheopen}{\isasymAnd}t{\isachardot}{\kern0pt}\ t\ {\isasymin}\ htps\ {\isasymLongrightarrow}\ {\isasymexists}A{\isachardot}{\kern0pt}\ {\isacharparenleft}{\kern0pt}t{\isacharcomma}{\kern0pt}\ A{\isacharparenright}{\kern0pt}\ {\isasymin}\ set\ {\isacharparenleft}{\kern0pt}h\ {\isacharhash}{\kern0pt}\ hs{\isadigit{2}}{\isacharparenright}{\kern0pt}{\isacartoucheclose}\isanewline
\ \ \ \ \ \ \ \ \isacommand{by}\isamarkupfalse%
\ force\isanewline
\ \ \ \ \ \ \isacommand{ultimately}\isamarkupfalse%
\ \isacommand{have}\isamarkupfalse%
\ {\isachardoublequoteopen}fst\ h{\isadigit{3}}{\isacharprime}{\kern0pt}\ {\isacharequal}{\kern0pt}\ fst\ h{\isachardoublequoteclose}\isanewline
\ \ \ \ \ \ \ \ \isacommand{using}\isamarkupfalse%
\ {\isacartoucheopen}fst\ h{\isadigit{3}}{\isacharprime}{\kern0pt}\ {\isasymin}\ htps{\isacartoucheclose}\ {\isacartoucheopen}{\isasymforall}h{\isasymin}set\ hs{\isadigit{3}}{\isachardot}{\kern0pt}\ fst\ h\ {\isasymnotin}\ htps{\isacartoucheclose}\ {\isacartoucheopen}strict{\isacharunderscore}{\kern0pt}sorted\ {\isacharparenleft}{\kern0pt}map\ fst\ hs{\isacharparenright}{\kern0pt}{\isacartoucheclose}\isanewline
\ \ \ \ \ \ \ \ \isacommand{by}\isamarkupfalse%
\ {\isacharparenleft}{\kern0pt}auto\ simp{\isacharcolon}{\kern0pt}\ simp\ del{\isacharcolon}{\kern0pt}\ {\isacharasterisk}{\kern0pt}\isanewline
\ \ \ \ \ \ \ \ \ \ \ \ \ \ \ \ \ intro{\isacharbang}{\kern0pt}{\isacharcolon}{\kern0pt}\ strict{\isacharunderscore}{\kern0pt}sorted{\isacharunderscore}{\kern0pt}order{\isacharbrackleft}{\kern0pt}\isakeyword{where}\ s\ {\isacharequal}{\kern0pt}\ htps\ \isakeyword{and}\ {\isacharquery}{\kern0pt}xs{\isadigit{1}}{\isachardot}{\kern0pt}{\isadigit{0}}\ {\isacharequal}{\kern0pt}\ {\isachardoublequoteopen}hs{\isadigit{3}}{\isachardoublequoteclose}\ \isakeyword{and}\ {\isacharquery}{\kern0pt}x{\isadigit{2}}{\isachardot}{\kern0pt}{\isadigit{0}}\ {\isacharequal}{\kern0pt}\ {\isachardoublequoteopen}h{\isadigit{2}}{\isachardoublequoteclose}\ \isakeyword{and}\ {\isacharquery}{\kern0pt}xs{\isadigit{2}}{\isachardot}{\kern0pt}{\isadigit{0}}\ {\isacharequal}{\kern0pt}\ {\isachardoublequoteopen}hs{\isadigit{4}}{\isachardoublequoteclose}{\isacharbrackright}{\kern0pt}{\isacharparenright}{\kern0pt}{\isacharplus}{\kern0pt}\isanewline
\ \ \ \ \ \ \isacommand{thus}\isamarkupfalse%
\ False\isanewline
\ \ \ \ \ \ \ \ \isacommand{using}\isamarkupfalse%
\ that\ {\isacartoucheopen}strict{\isacharunderscore}{\kern0pt}sorted\ {\isacharparenleft}{\kern0pt}map\ fst\ hs{\isacharprime}{\kern0pt}{\isacharparenright}{\kern0pt}{\isacartoucheclose}\isanewline
\ \ \ \ \ \ \ \ \isacommand{by}\isamarkupfalse%
\ auto\isanewline
\ \ \ \ \isacommand{qed}\isamarkupfalse%
\isanewline
\isanewline
\ \ \ \ \isacommand{define}\isamarkupfalse%
\ hs{\isacharunderscore}{\kern0pt}IH\ \isakeyword{where}\ {\isachardoublequoteopen}hs{\isacharunderscore}{\kern0pt}IH\ {\isasymequiv}\ h{\isadigit{2}}\ {\isacharhash}{\kern0pt}\ hs{\isadigit{4}}{\isachardoublequoteclose}\isanewline
\ \ \ \ \isacommand{define}\isamarkupfalse%
\ hs{\isacharprime}{\kern0pt}{\isacharunderscore}{\kern0pt}IH\ \isakeyword{where}\ {\isachardoublequoteopen}hs{\isacharprime}{\kern0pt}{\isacharunderscore}{\kern0pt}IH\ {\isasymequiv}\ {\isacharparenleft}{\kern0pt}fst\ h{\isadigit{2}}{\isacharcomma}{\kern0pt}\ A{\isadigit{2}}{\isacharprime}{\kern0pt}{\isacharparenright}{\kern0pt}\ {\isacharhash}{\kern0pt}\ hs{\isadigit{4}}{\isacharprime}{\kern0pt}{\isachardoublequoteclose}\isanewline
\ \ \ \ \isacommand{define}\isamarkupfalse%
\ htps{\isacharunderscore}{\kern0pt}IH\ \isakeyword{where}\ {\isachardoublequoteopen}htps{\isacharunderscore}{\kern0pt}IH\ {\isasymequiv}\ htps\ {\isacharminus}{\kern0pt}\ {\isacharbraceleft}{\kern0pt}t{\isacharbraceright}{\kern0pt}{\isachardoublequoteclose}\isanewline
\ \ \ \ \isacommand{define}\isamarkupfalse%
\ M{\isacharunderscore}{\kern0pt}IH\ \isakeyword{where}\ {\isachardoublequoteopen}M{\isacharunderscore}{\kern0pt}IH\ {\isasymequiv}\ {\isacharparenleft}{\kern0pt}apply{\isacharunderscore}{\kern0pt}happ\ {\isacharparenleft}{\kern0pt}t{\isacharcomma}{\kern0pt}\ A{\isacharprime}{\kern0pt}{\isacharparenright}{\kern0pt}\ M{\isacharparenright}{\kern0pt}{\isachardoublequoteclose}\isanewline
\ \ \ \ \isacommand{have}\isamarkupfalse%
\ {\isacharbrackleft}{\kern0pt}dest{\isacharbrackright}{\kern0pt}{\isacharcolon}{\kern0pt}{\isachardoublequoteopen}{\isasymAnd}h{\isacharprime}{\kern0pt}{\isachardot}{\kern0pt}\ h{\isacharprime}{\kern0pt}\ {\isasymin}\ set\ {\isacharparenleft}{\kern0pt}hs{\isacharunderscore}{\kern0pt}IH{\isacharparenright}{\kern0pt}\ {\isasymLongrightarrow}\ t\ {\isasymnoteq}\ fst\ h{\isacharprime}{\kern0pt}{\isachardoublequoteclose}\isanewline
\ \ \ \ \ \ \isacommand{using}\isamarkupfalse%
\ {\isacartoucheopen}strict{\isacharunderscore}{\kern0pt}sorted\ {\isacharparenleft}{\kern0pt}map\ fst\ hs{\isacharparenright}{\kern0pt}{\isacartoucheclose}\isanewline
\ \ \ \ \ \ \isacommand{by}\isamarkupfalse%
\ {\isacharparenleft}{\kern0pt}auto\ simp{\isacharcolon}{\kern0pt}\ hs{\isacharunderscore}{\kern0pt}IH{\isacharunderscore}{\kern0pt}def{\isacharparenright}{\kern0pt}\isanewline
\ \ \ \ \isacommand{have}\isamarkupfalse%
\ {\isachardoublequoteopen}{\isasymAnd}h{\isacharprime}{\kern0pt}{\isachardot}{\kern0pt}\ h{\isacharprime}{\kern0pt}\ {\isasymin}\ set\ {\isacharparenleft}{\kern0pt}hs{\isacharprime}{\kern0pt}{\isacharunderscore}{\kern0pt}IH{\isacharparenright}{\kern0pt}\ {\isasymLongrightarrow}\ t{\isacharprime}{\kern0pt}\ {\isasymnoteq}\ fst\ h{\isacharprime}{\kern0pt}{\isachardoublequoteclose}\isanewline
\ \ \ \ \ \ \isacommand{using}\isamarkupfalse%
\ {\isacartoucheopen}strict{\isacharunderscore}{\kern0pt}sorted\ {\isacharparenleft}{\kern0pt}map\ fst\ hs{\isacharprime}{\kern0pt}{\isacharparenright}{\kern0pt}{\isacartoucheclose}\isanewline
\ \ \ \ \ \ \isacommand{by}\isamarkupfalse%
\ {\isacharparenleft}{\kern0pt}auto\ simp{\isacharcolon}{\kern0pt}\ hs{\isacharprime}{\kern0pt}{\isacharunderscore}{\kern0pt}IH{\isacharunderscore}{\kern0pt}def{\isacharparenright}{\kern0pt}\isanewline
\isanewline
\ \ \ \ \isacommand{have}\isamarkupfalse%
\ {\isachardoublequoteopen}xs\ {\isacharequal}{\kern0pt}\ filter\ {\isacharparenleft}{\kern0pt}{\isasymlambda}h{\isachardot}{\kern0pt}\ fst\ h\ {\isasymin}\ htps{\isacharunderscore}{\kern0pt}IH{\isacharparenright}{\kern0pt}\ {\isacharparenleft}{\kern0pt}hs{\isacharunderscore}{\kern0pt}IH{\isacharparenright}{\kern0pt}{\isachardoublequoteclose}\ {\isacharparenleft}{\kern0pt}\isakeyword{is}\ {\isachardoublequoteopen}{\isacharquery}{\kern0pt}xs\ {\isacharequal}{\kern0pt}\ {\isacharquery}{\kern0pt}fil{\isacharunderscore}{\kern0pt}xs\ {\isacharparenleft}{\kern0pt}{\isasymlambda}h{\isachardot}{\kern0pt}\ fst\ h\ {\isasymin}\ htps{\isacharunderscore}{\kern0pt}IH{\isacharparenright}{\kern0pt}{\isachardoublequoteclose}{\isacharparenright}{\kern0pt}\isanewline
\ \ \ \ \isacommand{proof}\isamarkupfalse%
{\isacharminus}{\kern0pt}\isanewline
\ \ \ \ \ \ \isacommand{have}\isamarkupfalse%
\ {\isachardoublequoteopen}{\isacharquery}{\kern0pt}fil{\isacharunderscore}{\kern0pt}xs\ {\isacharparenleft}{\kern0pt}{\isasymlambda}h{\isachardot}{\kern0pt}\ fst\ h\ {\isasymin}\ htps{\isacharunderscore}{\kern0pt}IH{\isacharparenright}{\kern0pt}\ {\isacharequal}{\kern0pt}\ {\isacharquery}{\kern0pt}fil{\isacharunderscore}{\kern0pt}xs\ {\isacharparenleft}{\kern0pt}{\isasymlambda}h{\isachardot}{\kern0pt}\ fst\ h\ {\isasymin}\ htps{\isacharparenright}{\kern0pt}{\isachardoublequoteclose}\isanewline
\ \ \ \ \ \ \ \ \isacommand{by}\isamarkupfalse%
\ {\isacharparenleft}{\kern0pt}fastforce\ intro{\isacharbang}{\kern0pt}{\isacharcolon}{\kern0pt}\ filter{\isacharunderscore}{\kern0pt}cong\ simp{\isacharcolon}{\kern0pt}\ htps{\isacharunderscore}{\kern0pt}IH{\isacharunderscore}{\kern0pt}def{\isacharparenright}{\kern0pt}\isanewline
\ \ \ \ \ \ \isacommand{also}\isamarkupfalse%
\ \isacommand{have}\isamarkupfalse%
\ {\isachardoublequoteopen}{\isachardot}{\kern0pt}{\isachardot}{\kern0pt}{\isachardot}{\kern0pt}\ {\isacharequal}{\kern0pt}\ filter\ {\isacharparenleft}{\kern0pt}{\isasymlambda}h{\isachardot}{\kern0pt}\ fst\ h\ {\isasymin}\ htps{\isacharparenright}{\kern0pt}\ hs{\isadigit{2}}{\isachardoublequoteclose}\isanewline
\ \ \ \ \ \ \ \ \isacommand{using}\isamarkupfalse%
\ Cons\ {\isacartoucheopen}{\isacharparenleft}{\kern0pt}{\isasymforall}h{\isasymin}set\ hs{\isadigit{3}}{\isachardot}{\kern0pt}\ fst\ h\ {\isasymnotin}\ htps{\isacharparenright}{\kern0pt}{\isacartoucheclose}\isanewline
\ \ \ \ \ \ \ \ \isacommand{by}\isamarkupfalse%
\ {\isacharparenleft}{\kern0pt}auto\ simp{\isacharcolon}{\kern0pt}\ hs{\isacharunderscore}{\kern0pt}IH{\isacharunderscore}{\kern0pt}def{\isacharparenright}{\kern0pt}\isanewline
\ \ \ \ \ \ \isacommand{finally}\isamarkupfalse%
\ \isacommand{show}\isamarkupfalse%
\ {\isacharquery}{\kern0pt}thesis\isanewline
\ \ \ \ \ \ \ \ \isacommand{using}\isamarkupfalse%
\ Cons\isanewline
\ \ \ \ \ \ \ \ \isacommand{by}\isamarkupfalse%
\ {\isacharparenleft}{\kern0pt}simp\ add{\isacharcolon}{\kern0pt}\ {\isacartoucheopen}xs\ {\isacharequal}{\kern0pt}\ filter\ {\isacharparenleft}{\kern0pt}{\isasymlambda}h{\isachardot}{\kern0pt}\ fst\ h\ {\isasymin}\ htps{\isacharparenright}{\kern0pt}\ hs{\isadigit{2}}{\isacartoucheclose}{\isacharparenright}{\kern0pt}\isanewline
\ \ \ \ \isacommand{qed}\isamarkupfalse%
\isanewline
\ \ \ \ \isacommand{moreover}\isamarkupfalse%
\ \isacommand{have}\isamarkupfalse%
\ {\isachardoublequoteopen}ys\ {\isacharequal}{\kern0pt}\ filter\ {\isacharparenleft}{\kern0pt}{\isasymlambda}h{\isachardot}{\kern0pt}\ fst\ h\ {\isasymin}\ htps{\isacharunderscore}{\kern0pt}IH{\isacharparenright}{\kern0pt}\ {\isacharparenleft}{\kern0pt}hs{\isacharprime}{\kern0pt}{\isacharunderscore}{\kern0pt}IH{\isacharparenright}{\kern0pt}{\isachardoublequoteclose}\ {\isacharparenleft}{\kern0pt}\isakeyword{is}\ {\isachardoublequoteopen}{\isacharquery}{\kern0pt}ys\ {\isacharequal}{\kern0pt}\ {\isacharquery}{\kern0pt}fil{\isacharunderscore}{\kern0pt}ys\ {\isacharparenleft}{\kern0pt}{\isasymlambda}h{\isachardot}{\kern0pt}\ fst\ h\ {\isasymin}\ htps{\isacharunderscore}{\kern0pt}IH{\isacharparenright}{\kern0pt}{\isachardoublequoteclose}{\isacharparenright}{\kern0pt}\isanewline
\ \ \ \ \isacommand{proof}\isamarkupfalse%
{\isacharminus}{\kern0pt}\isanewline
\ \ \ \ \ \ \isacommand{have}\isamarkupfalse%
\ {\isachardoublequoteopen}{\isacharquery}{\kern0pt}fil{\isacharunderscore}{\kern0pt}ys\ {\isacharparenleft}{\kern0pt}{\isasymlambda}h{\isachardot}{\kern0pt}\ fst\ h\ {\isasymin}\ htps\ {\isacharminus}{\kern0pt}\ {\isacharbraceleft}{\kern0pt}t{\isacharprime}{\kern0pt}{\isacharbraceright}{\kern0pt}{\isacharparenright}{\kern0pt}\ {\isacharequal}{\kern0pt}\ {\isacharquery}{\kern0pt}fil{\isacharunderscore}{\kern0pt}ys\ {\isacharparenleft}{\kern0pt}{\isasymlambda}h{\isachardot}{\kern0pt}\ fst\ h\ {\isasymin}\ htps{\isacharparenright}{\kern0pt}{\isachardoublequoteclose}\isanewline
\ \ \ \ \ \ \ \ \isacommand{using}\isamarkupfalse%
\ {\isacartoucheopen}{\isasymAnd}h{\isacharprime}{\kern0pt}{\isachardot}{\kern0pt}\ h{\isacharprime}{\kern0pt}\ {\isasymin}\ set\ {\isacharparenleft}{\kern0pt}hs{\isacharprime}{\kern0pt}{\isacharunderscore}{\kern0pt}IH{\isacharparenright}{\kern0pt}\ {\isasymLongrightarrow}\ t{\isacharprime}{\kern0pt}\ {\isasymnoteq}\ fst\ h{\isacharprime}{\kern0pt}{\isacartoucheclose}\isanewline
\ \ \ \ \ \ \ \ \isacommand{by}\isamarkupfalse%
\ {\isacharparenleft}{\kern0pt}force\ intro{\isacharbang}{\kern0pt}{\isacharcolon}{\kern0pt}\ filter{\isacharunderscore}{\kern0pt}cong{\isacharparenright}{\kern0pt}\isanewline
\ \ \ \ \ \ \isacommand{also}\isamarkupfalse%
\ \isacommand{have}\isamarkupfalse%
\ {\isachardoublequoteopen}{\isachardot}{\kern0pt}{\isachardot}{\kern0pt}{\isachardot}{\kern0pt}\ {\isacharequal}{\kern0pt}\ filter\ {\isacharparenleft}{\kern0pt}{\isasymlambda}h{\isachardot}{\kern0pt}\ fst\ h\ {\isasymin}\ htps{\isacharparenright}{\kern0pt}\ hs{\isadigit{2}}{\isacharprime}{\kern0pt}{\isachardoublequoteclose}\isanewline
\ \ \ \ \ \ \ \ \isacommand{using}\isamarkupfalse%
\ Cons\ fst{\isacharunderscore}{\kern0pt}h{\isadigit{3}}{\isacharprime}{\kern0pt}{\isacharunderscore}{\kern0pt}n{\isacharunderscore}{\kern0pt}htp\isanewline
\ \ \ \ \ \ \ \ \isacommand{by}\isamarkupfalse%
\ {\isacharparenleft}{\kern0pt}auto\ simp{\isacharcolon}{\kern0pt}\ hs{\isacharprime}{\kern0pt}{\isacharunderscore}{\kern0pt}IH{\isacharunderscore}{\kern0pt}def{\isacharparenright}{\kern0pt}\isanewline
\ \ \ \ \ \ \isacommand{finally}\isamarkupfalse%
\ \isacommand{show}\isamarkupfalse%
\ {\isacharquery}{\kern0pt}thesis\isanewline
\ \ \ \ \ \ \ \ \isacommand{using}\isamarkupfalse%
\ Cons\isanewline
\ \ \ \ \ \ \ \ \isacommand{by}\isamarkupfalse%
\ {\isacharparenleft}{\kern0pt}simp\ add{\isacharcolon}{\kern0pt}\ {\isacartoucheopen}ys\ {\isacharequal}{\kern0pt}\ filter\ {\isacharparenleft}{\kern0pt}{\isasymlambda}h{\isachardot}{\kern0pt}\ fst\ h\ {\isasymin}\ htps{\isacharparenright}{\kern0pt}\ hs{\isadigit{2}}{\isacharprime}{\kern0pt}{\isacartoucheclose}\ htps{\isacharunderscore}{\kern0pt}IH{\isacharunderscore}{\kern0pt}def{\isacharparenright}{\kern0pt}\isanewline
\ \ \ \ \isacommand{qed}\isamarkupfalse%
\isanewline
\ \ \ \ \isacommand{moreover}\isamarkupfalse%
\ \isacommand{have}\isamarkupfalse%
\ {\isachardoublequoteopen}{\isasymAnd}t{\isacharprime}{\kern0pt}{\isachardot}{\kern0pt}\ t{\isacharprime}{\kern0pt}\ {\isasymin}\ htps{\isacharunderscore}{\kern0pt}IH\ {\isasymLongrightarrow}\ {\isasymexists}A{\isachardot}{\kern0pt}\ {\isacharparenleft}{\kern0pt}t{\isacharprime}{\kern0pt}{\isacharcomma}{\kern0pt}\ A{\isacharparenright}{\kern0pt}\ {\isasymin}\ set\ {\isacharparenleft}{\kern0pt}hs{\isacharunderscore}{\kern0pt}IH{\isacharparenright}{\kern0pt}{\isachardoublequoteclose}\isanewline
\ \ \ \ \ \ \isacommand{using}\isamarkupfalse%
\ {\isacartoucheopen}{\isacharparenleft}{\kern0pt}{\isasymforall}h{\isasymin}set\ hs{\isadigit{3}}{\isachardot}{\kern0pt}\ fst\ h\ {\isasymnotin}\ htps{\isacharparenright}{\kern0pt}{\isacartoucheclose}\isanewline
\ \ \ \ \ \ \ \ {\isacartoucheopen}{\isasymAnd}t{\isachardot}{\kern0pt}\ t\ {\isasymin}\ htps\ {\isasymLongrightarrow}\ {\isasymexists}A{\isachardot}{\kern0pt}\ {\isacharparenleft}{\kern0pt}t{\isacharcomma}{\kern0pt}\ A{\isacharparenright}{\kern0pt}\ {\isasymin}\ set\ {\isacharparenleft}{\kern0pt}h\ {\isacharhash}{\kern0pt}\ hs{\isadigit{2}}{\isacharparenright}{\kern0pt}{\isacartoucheclose}\isanewline
\ \ \ \ \ \ \ \ {\isacartoucheopen}{\isasymAnd}h{\isacharprime}{\kern0pt}{\isachardot}{\kern0pt}\ h{\isacharprime}{\kern0pt}\ {\isasymin}\ set\ {\isacharparenleft}{\kern0pt}hs{\isacharunderscore}{\kern0pt}IH{\isacharparenright}{\kern0pt}\ {\isasymLongrightarrow}\ t\ {\isasymnoteq}\ fst\ h{\isacharprime}{\kern0pt}{\isacartoucheclose}\isanewline
\ \ \ \ \ \ \isacommand{by}\isamarkupfalse%
\ {\isacharparenleft}{\kern0pt}fastforce\ simp{\isacharcolon}{\kern0pt}\ hs{\isacharunderscore}{\kern0pt}IH{\isacharunderscore}{\kern0pt}def\ htps{\isacharunderscore}{\kern0pt}IH{\isacharunderscore}{\kern0pt}def{\isacharparenright}{\kern0pt}\isanewline
\ \ \ \ \isacommand{moreover}\isamarkupfalse%
\ \isacommand{have}\isamarkupfalse%
\ {\isachardoublequoteopen}{\isasymAnd}t{\isacharprime}{\kern0pt}{\isacharprime}{\kern0pt}{\isachardot}{\kern0pt}\ t{\isacharprime}{\kern0pt}{\isacharprime}{\kern0pt}\ {\isasymin}\ htps{\isacharunderscore}{\kern0pt}IH\ {\isasymLongrightarrow}\ {\isasymexists}A{\isachardot}{\kern0pt}\ {\isacharparenleft}{\kern0pt}t{\isacharprime}{\kern0pt}{\isacharprime}{\kern0pt}{\isacharcomma}{\kern0pt}\ A{\isacharparenright}{\kern0pt}\ {\isasymin}\ set\ {\isacharparenleft}{\kern0pt}hs{\isacharprime}{\kern0pt}{\isacharunderscore}{\kern0pt}IH{\isacharparenright}{\kern0pt}{\isachardoublequoteclose}\isanewline
\ \ \ \ \ \ \isacommand{using}\isamarkupfalse%
\ fst{\isacharunderscore}{\kern0pt}h{\isadigit{3}}{\isacharprime}{\kern0pt}{\isacharunderscore}{\kern0pt}n{\isacharunderscore}{\kern0pt}htp\isanewline
\ \ \ \ \ \ \ \ {\isacartoucheopen}{\isasymAnd}t{\isachardot}{\kern0pt}\ t\ {\isasymin}\ htps\ {\isasymLongrightarrow}\ {\isasymexists}A{\isachardot}{\kern0pt}\ {\isacharparenleft}{\kern0pt}t{\isacharcomma}{\kern0pt}\ A{\isacharparenright}{\kern0pt}\ {\isasymin}\ set\ {\isacharparenleft}{\kern0pt}h{\isacharprime}{\kern0pt}\ {\isacharhash}{\kern0pt}\ hs{\isadigit{2}}{\isacharprime}{\kern0pt}{\isacharparenright}{\kern0pt}{\isacartoucheclose}\isanewline
\ \ \ \ \ \ \isacommand{by}\isamarkupfalse%
\ {\isacharparenleft}{\kern0pt}fastforce\ simp{\isacharcolon}{\kern0pt}\ hs{\isacharprime}{\kern0pt}{\isacharunderscore}{\kern0pt}IH{\isacharunderscore}{\kern0pt}def\ htps{\isacharunderscore}{\kern0pt}IH{\isacharunderscore}{\kern0pt}def{\isacharparenright}{\kern0pt}\isanewline
\ \ \ \ \isacommand{moreover}\isamarkupfalse%
\ \isacommand{have}\isamarkupfalse%
\ {\isachardoublequoteopen}strict{\isacharunderscore}{\kern0pt}sorted\ {\isacharparenleft}{\kern0pt}map\ fst\ {\isacharparenleft}{\kern0pt}hs{\isacharunderscore}{\kern0pt}IH{\isacharparenright}{\kern0pt}{\isacharparenright}{\kern0pt}{\isachardoublequoteclose}\isanewline
\ \ \ \ \ \ \isacommand{using}\isamarkupfalse%
\ {\isacartoucheopen}strict{\isacharunderscore}{\kern0pt}sorted\ {\isacharparenleft}{\kern0pt}map\ fst\ hs{\isacharparenright}{\kern0pt}{\isacartoucheclose}\isanewline
\ \ \ \ \ \ \isacommand{by}\isamarkupfalse%
\ {\isacharparenleft}{\kern0pt}auto\ simp{\isacharcolon}{\kern0pt}\ hs{\isacharunderscore}{\kern0pt}IH{\isacharunderscore}{\kern0pt}def\ htps{\isacharunderscore}{\kern0pt}IH{\isacharunderscore}{\kern0pt}def{\isacharparenright}{\kern0pt}\isanewline
\ \ \ \ \isacommand{moreover}\isamarkupfalse%
\ \isacommand{have}\isamarkupfalse%
\ {\isachardoublequoteopen}strict{\isacharunderscore}{\kern0pt}sorted\ {\isacharparenleft}{\kern0pt}map\ fst\ {\isacharparenleft}{\kern0pt}hs{\isacharprime}{\kern0pt}{\isacharunderscore}{\kern0pt}IH{\isacharparenright}{\kern0pt}{\isacharparenright}{\kern0pt}{\isachardoublequoteclose}\isanewline
\ \ \ \ \ \ \isacommand{using}\isamarkupfalse%
\ {\isacartoucheopen}strict{\isacharunderscore}{\kern0pt}sorted\ {\isacharparenleft}{\kern0pt}map\ fst\ hs{\isacharprime}{\kern0pt}{\isacharparenright}{\kern0pt}{\isacartoucheclose}\isanewline
\ \ \ \ \ \ \isacommand{by}\isamarkupfalse%
\ {\isacharparenleft}{\kern0pt}auto\ simp{\isacharcolon}{\kern0pt}\ hs{\isacharprime}{\kern0pt}{\isacharunderscore}{\kern0pt}IH{\isacharunderscore}{\kern0pt}def\ htps{\isacharunderscore}{\kern0pt}IH{\isacharunderscore}{\kern0pt}def{\isacharparenright}{\kern0pt}\isanewline
\ \ \ \ \isacommand{moreover}\isamarkupfalse%
\ \isacommand{have}\isamarkupfalse%
\isanewline
\ \ \ \ \ \ {\isachardoublequoteopen}{\isasymAnd}\ t{\isacharprime}{\kern0pt}{\isacharprime}{\kern0pt}\ A{\isacharprime}{\kern0pt}{\isacharprime}{\kern0pt}\ A{\isacharprime}{\kern0pt}{\isacharprime}{\kern0pt}{\isacharprime}{\kern0pt}{\isachardot}{\kern0pt}\ {\isasymlbrakk}t{\isacharprime}{\kern0pt}{\isacharprime}{\kern0pt}\ {\isasymin}\ htps{\isacharunderscore}{\kern0pt}IH{\isacharsemicolon}{\kern0pt}\ {\isacharparenleft}{\kern0pt}t{\isacharprime}{\kern0pt}{\isacharprime}{\kern0pt}{\isacharcomma}{\kern0pt}\ A{\isacharprime}{\kern0pt}{\isacharprime}{\kern0pt}{\isacharparenright}{\kern0pt}\ {\isasymin}\ set\ {\isacharparenleft}{\kern0pt}hs{\isacharunderscore}{\kern0pt}IH{\isacharparenright}{\kern0pt}{\isacharsemicolon}{\kern0pt}\ {\isacharparenleft}{\kern0pt}t{\isacharprime}{\kern0pt}{\isacharprime}{\kern0pt}{\isacharcomma}{\kern0pt}\ A{\isacharprime}{\kern0pt}{\isacharprime}{\kern0pt}{\isacharprime}{\kern0pt}{\isacharparenright}{\kern0pt}\ {\isasymin}\ set\ {\isacharparenleft}{\kern0pt}hs{\isacharprime}{\kern0pt}{\isacharunderscore}{\kern0pt}IH{\isacharparenright}{\kern0pt}{\isasymrbrakk}\ {\isasymLongrightarrow}\ set\ A{\isacharprime}{\kern0pt}{\isacharprime}{\kern0pt}\ {\isacharequal}{\kern0pt}\ set\ A{\isacharprime}{\kern0pt}{\isacharprime}{\kern0pt}{\isacharprime}{\kern0pt}{\isachardoublequoteclose}\isanewline
\ \ \ \ \ \ \isacommand{using}\isamarkupfalse%
\ {\isadigit{2}}{\isacharparenleft}{\kern0pt}{\isadigit{8}}{\isacharparenright}{\kern0pt}\isanewline
\ \ \ \ \ \ \isacommand{by}\isamarkupfalse%
\ {\isacharparenleft}{\kern0pt}auto\ simp{\isacharcolon}{\kern0pt}\ htps{\isacharunderscore}{\kern0pt}IH{\isacharunderscore}{\kern0pt}def\ hs{\isacharunderscore}{\kern0pt}IH{\isacharunderscore}{\kern0pt}def\ hs{\isacharprime}{\kern0pt}{\isacharunderscore}{\kern0pt}IH{\isacharunderscore}{\kern0pt}def{\isacharparenright}{\kern0pt}\isanewline
\ \ \ \ \isacommand{moreover}\isamarkupfalse%
\ \isacommand{have}\isamarkupfalse%
\ {\isachardoublequoteopen}{\isasymAnd}t{\isacharprime}{\kern0pt}{\isacharprime}{\kern0pt}\ A{\isacharprime}{\kern0pt}{\isacharprime}{\kern0pt}\ a{\isachardot}{\kern0pt}\ {\isasymlbrakk}t{\isacharprime}{\kern0pt}{\isacharprime}{\kern0pt}\ {\isasymnotin}\ htps{\isacharunderscore}{\kern0pt}IH{\isacharsemicolon}{\kern0pt}\ {\isacharparenleft}{\kern0pt}t{\isacharprime}{\kern0pt}{\isacharprime}{\kern0pt}{\isacharcomma}{\kern0pt}\ A{\isacharprime}{\kern0pt}{\isacharprime}{\kern0pt}{\isacharparenright}{\kern0pt}\ {\isasymin}\ set\ {\isacharparenleft}{\kern0pt}hs{\isacharunderscore}{\kern0pt}IH{\isacharparenright}{\kern0pt}{\isacharsemicolon}{\kern0pt}\ a\ {\isasymin}\ set\ A{\isacharprime}{\kern0pt}{\isacharprime}{\kern0pt}{\isasymrbrakk}\ {\isasymLongrightarrow}\ effect\ a\ {\isacharequal}{\kern0pt}\ Effect\ {\isacharbrackleft}{\kern0pt}{\isacharbrackright}{\kern0pt}\ {\isacharbrackleft}{\kern0pt}{\isacharbrackright}{\kern0pt}{\isachardoublequoteclose}\isanewline
\ \ \ \ \ \ \isacommand{using}\isamarkupfalse%
\ {\isadigit{2}}{\isacharparenleft}{\kern0pt}{\isadigit{9}}{\isacharparenright}{\kern0pt}\ {\isacartoucheopen}fst\ h{\isadigit{2}}\ {\isasymin}\ htps{\isacartoucheclose}\ {\isacartoucheopen}{\isasymAnd}h{\isacharprime}{\kern0pt}{\isachardot}{\kern0pt}\ h{\isacharprime}{\kern0pt}\ {\isasymin}\ set\ {\isacharparenleft}{\kern0pt}hs{\isacharunderscore}{\kern0pt}IH{\isacharparenright}{\kern0pt}\ {\isasymLongrightarrow}\ t\ {\isasymnoteq}\ fst\ h{\isacharprime}{\kern0pt}{\isacartoucheclose}\isanewline
\ \ \ \ \ \ \isacommand{by}\isamarkupfalse%
\ {\isacharparenleft}{\kern0pt}fastforce\ simp\ add{\isacharcolon}{\kern0pt}\ htps{\isacharunderscore}{\kern0pt}IH{\isacharunderscore}{\kern0pt}def\ hs{\isacharunderscore}{\kern0pt}IH{\isacharunderscore}{\kern0pt}def{\isacharparenright}{\kern0pt}\isanewline
\ \ \ \ \isacommand{moreover}\isamarkupfalse%
\ \isacommand{have}\isamarkupfalse%
\ {\isachardoublequoteopen}{\isasymAnd}t{\isacharprime}{\kern0pt}{\isacharprime}{\kern0pt}\ A{\isacharprime}{\kern0pt}{\isacharprime}{\kern0pt}\ a{\isachardot}{\kern0pt}\ {\isasymlbrakk}t{\isacharprime}{\kern0pt}{\isacharprime}{\kern0pt}\ {\isasymnotin}\ htps{\isacharunderscore}{\kern0pt}IH{\isacharsemicolon}{\kern0pt}\ {\isacharparenleft}{\kern0pt}t{\isacharprime}{\kern0pt}{\isacharprime}{\kern0pt}{\isacharcomma}{\kern0pt}\ A{\isacharprime}{\kern0pt}{\isacharprime}{\kern0pt}{\isacharparenright}{\kern0pt}\ {\isasymin}\ set\ {\isacharparenleft}{\kern0pt}hs{\isacharprime}{\kern0pt}{\isacharunderscore}{\kern0pt}IH{\isacharparenright}{\kern0pt}{\isacharsemicolon}{\kern0pt}\ a\ {\isasymin}\ set\ A{\isacharprime}{\kern0pt}{\isacharprime}{\kern0pt}{\isasymrbrakk}\ {\isasymLongrightarrow}\ effect\ a\ {\isacharequal}{\kern0pt}\ Effect\ {\isacharbrackleft}{\kern0pt}{\isacharbrackright}{\kern0pt}\ {\isacharbrackleft}{\kern0pt}{\isacharbrackright}{\kern0pt}{\isachardoublequoteclose}\isanewline
\ \ \ \ \ \ \isacommand{using}\isamarkupfalse%
\ {\isadigit{2}}{\isacharparenleft}{\kern0pt}{\isadigit{1}}{\isadigit{0}}{\isacharparenright}{\kern0pt}\ {\isacartoucheopen}fst\ h{\isadigit{2}}\ {\isasymin}\ htps{\isacartoucheclose}\ {\isacartoucheopen}{\isasymAnd}h{\isacharprime}{\kern0pt}{\isachardot}{\kern0pt}\ h{\isacharprime}{\kern0pt}\ {\isasymin}\ set\ {\isacharparenleft}{\kern0pt}hs{\isacharprime}{\kern0pt}{\isacharunderscore}{\kern0pt}IH{\isacharparenright}{\kern0pt}\ {\isasymLongrightarrow}\ t{\isacharprime}{\kern0pt}\ {\isasymnoteq}\ fst\ h{\isacharprime}{\kern0pt}{\isacartoucheclose}\isanewline
\ \ \ \ \ \ \isacommand{by}\isamarkupfalse%
\ {\isacharparenleft}{\kern0pt}force\ simp\ add{\isacharcolon}{\kern0pt}\ htps{\isacharunderscore}{\kern0pt}IH{\isacharunderscore}{\kern0pt}def\ hs{\isacharprime}{\kern0pt}{\isacharunderscore}{\kern0pt}IH{\isacharunderscore}{\kern0pt}def{\isacharparenright}{\kern0pt}\isanewline
\ \ \ \ \isacommand{moreover}\isamarkupfalse%
\ \isacommand{have}\isamarkupfalse%
\ {\isachardoublequoteopen}{\isasymexists}t{\isacharprime}{\kern0pt}{\isasymin}{\isacharbraceleft}{\kern0pt}t\isactrlsub {\isadigit{1}}{\isacharless}{\kern0pt}{\isachardot}{\kern0pt}{\isachardot}{\kern0pt}{\isacharless}{\kern0pt}\ t\isactrlsub {\isadigit{2}}{\isacharbraceright}{\kern0pt}{\isachardot}{\kern0pt}{\isasymexists}A{\isacharprime}{\kern0pt}{\isachardot}{\kern0pt}\ {\isacharparenleft}{\kern0pt}t{\isacharprime}{\kern0pt}{\isacharcomma}{\kern0pt}A{\isacharprime}{\kern0pt}{\isacharparenright}{\kern0pt}\ {\isasymin}\ set\ hs{\isacharprime}{\kern0pt}{\isacharunderscore}{\kern0pt}IH\ {\isasymand}\ a\ {\isasymin}\ set\ A{\isacharprime}{\kern0pt}{\isachardoublequoteclose}\isanewline
\ \ \ \ \ \ \isakeyword{if}\ {\isachardoublequoteopen}t{\isacharprime}{\kern0pt}{\isacharprime}{\kern0pt}\ {\isasymnotin}\ htps{\isacharunderscore}{\kern0pt}IH{\isachardoublequoteclose}\ {\isachardoublequoteopen}t\isactrlsub {\isadigit{1}}\ {\isasymin}\ htps{\isacharunderscore}{\kern0pt}IH{\isachardoublequoteclose}\ {\isachardoublequoteopen}t\isactrlsub {\isadigit{2}}\ {\isasymin}\ htps{\isacharunderscore}{\kern0pt}IH{\isachardoublequoteclose}\ {\isachardoublequoteopen}t{\isacharprime}{\kern0pt}{\isacharprime}{\kern0pt}\ {\isasymin}\ {\isacharbraceleft}{\kern0pt}t\isactrlsub {\isadigit{1}}{\isacharless}{\kern0pt}{\isachardot}{\kern0pt}{\isachardot}{\kern0pt}{\isacharless}{\kern0pt}\ t\isactrlsub {\isadigit{2}}{\isacharbraceright}{\kern0pt}{\isachardoublequoteclose}\ {\isachardoublequoteopen}{\isacharparenleft}{\kern0pt}t{\isacharprime}{\kern0pt}{\isacharprime}{\kern0pt}{\isacharcomma}{\kern0pt}\ A{\isacharprime}{\kern0pt}{\isacharprime}{\kern0pt}{\isacharparenright}{\kern0pt}\ {\isasymin}\ set\ hs{\isacharunderscore}{\kern0pt}IH{\isachardoublequoteclose}\isanewline
\ \ \ \ \ \ \ \ {\isachardoublequoteopen}a\ {\isasymin}\ set\ A{\isacharprime}{\kern0pt}{\isacharprime}{\kern0pt}{\isachardoublequoteclose}\isanewline
\ \ \ \ \ \ \isakeyword{for}\ t{\isacharprime}{\kern0pt}{\isacharprime}{\kern0pt}\ A{\isacharprime}{\kern0pt}{\isacharprime}{\kern0pt}\ t\isactrlsub {\isadigit{1}}\ t\isactrlsub {\isadigit{2}}\ a\isanewline
\ \ \ \ \isacommand{proof}\isamarkupfalse%
{\isacharminus}{\kern0pt}\isanewline
\ \ \ \ \ \ \isacommand{have}\isamarkupfalse%
\ {\isachardoublequoteopen}t{\isacharprime}{\kern0pt}{\isacharprime}{\kern0pt}\ {\isasymnotin}\ htps{\isachardoublequoteclose}\isanewline
\ \ \ \ \ \ \ \ \isacommand{using}\isamarkupfalse%
\ {\isacartoucheopen}{\isasymAnd}h{\isacharprime}{\kern0pt}{\isachardot}{\kern0pt}\ h{\isacharprime}{\kern0pt}\ {\isasymin}\ set\ {\isacharparenleft}{\kern0pt}hs{\isacharunderscore}{\kern0pt}IH{\isacharparenright}{\kern0pt}\ {\isasymLongrightarrow}\ t\ {\isasymnoteq}\ fst\ h{\isacharprime}{\kern0pt}{\isacartoucheclose}\ {\isacartoucheopen}t{\isacharprime}{\kern0pt}{\isacharprime}{\kern0pt}\ {\isasymnotin}\ htps{\isacharunderscore}{\kern0pt}IH{\isacartoucheclose}\ {\isacartoucheopen}{\isacharparenleft}{\kern0pt}t{\isacharprime}{\kern0pt}{\isacharprime}{\kern0pt}{\isacharcomma}{\kern0pt}\ A{\isacharprime}{\kern0pt}{\isacharprime}{\kern0pt}{\isacharparenright}{\kern0pt}\ {\isasymin}\ set\ hs{\isacharunderscore}{\kern0pt}IH{\isacartoucheclose}\isanewline
\ \ \ \ \ \ \ \ \isacommand{by}\isamarkupfalse%
\ {\isacharparenleft}{\kern0pt}fastforce\ simp{\isacharcolon}{\kern0pt}\ htps{\isacharunderscore}{\kern0pt}IH{\isacharunderscore}{\kern0pt}def{\isacharparenright}{\kern0pt}\isanewline
\ \ \ \ \ \ \isacommand{hence}\isamarkupfalse%
\ {\isachardoublequoteopen}{\isasymexists}t{\isacharprime}{\kern0pt}{\isasymin}{\isacharbraceleft}{\kern0pt}t\isactrlsub {\isadigit{1}}{\isacharless}{\kern0pt}{\isachardot}{\kern0pt}{\isachardot}{\kern0pt}{\isacharless}{\kern0pt}\ t\isactrlsub {\isadigit{2}}{\isacharbraceright}{\kern0pt}{\isachardot}{\kern0pt}{\isasymexists}A{\isacharprime}{\kern0pt}{\isachardot}{\kern0pt}\ {\isacharparenleft}{\kern0pt}t{\isacharprime}{\kern0pt}{\isacharcomma}{\kern0pt}A{\isacharprime}{\kern0pt}{\isacharparenright}{\kern0pt}\ {\isasymin}\ set\ hs{\isacharprime}{\kern0pt}\ {\isasymand}\ a\ {\isasymin}\ set\ A{\isacharprime}{\kern0pt}{\isachardoublequoteclose}\isanewline
\ \ \ \ \ \ \ \ \isacommand{using}\isamarkupfalse%
\ that\isanewline
\ \ \ \ \ \ \ \ \isacommand{by}\isamarkupfalse%
\ {\isacharparenleft}{\kern0pt}auto\ simp\ del{\isacharcolon}{\kern0pt}\ {\isacartoucheopen}hs{\isacharprime}{\kern0pt}\ {\isacharequal}{\kern0pt}\ hs{\isadigit{1}}{\isacharprime}{\kern0pt}\ {\isacharat}{\kern0pt}\ h{\isacharprime}{\kern0pt}\ {\isacharhash}{\kern0pt}\ hs{\isadigit{2}}{\isacharprime}{\kern0pt}{\isacartoucheclose}\ simp{\isacharcolon}{\kern0pt}\ htps{\isacharunderscore}{\kern0pt}IH{\isacharunderscore}{\kern0pt}def\ hs{\isacharunderscore}{\kern0pt}IH{\isacharunderscore}{\kern0pt}def\ intro{\isacharbang}{\kern0pt}{\isacharcolon}{\kern0pt}\ {\isadigit{2}}{\isacharparenleft}{\kern0pt}{\isadigit{1}}{\isadigit{1}}{\isacharparenright}{\kern0pt}{\isacharparenright}{\kern0pt}\isanewline
\ \ \ \ \ \ \isacommand{then}\isamarkupfalse%
\ \isacommand{obtain}\isamarkupfalse%
\ t{\isacharprime}{\kern0pt}{\isacharprime}{\kern0pt}{\isacharprime}{\kern0pt}\ A{\isacharprime}{\kern0pt}{\isacharprime}{\kern0pt}{\isacharprime}{\kern0pt}\ \isakeyword{where}\ {\isachardoublequoteopen}t{\isacharprime}{\kern0pt}{\isacharprime}{\kern0pt}{\isacharprime}{\kern0pt}\ {\isasymin}\ {\isacharbraceleft}{\kern0pt}t\isactrlsub {\isadigit{1}}{\isacharless}{\kern0pt}{\isachardot}{\kern0pt}{\isachardot}{\kern0pt}{\isacharless}{\kern0pt}\ t\isactrlsub {\isadigit{2}}{\isacharbraceright}{\kern0pt}{\isachardoublequoteclose}\ {\isachardoublequoteopen}{\isacharparenleft}{\kern0pt}t{\isacharprime}{\kern0pt}{\isacharprime}{\kern0pt}{\isacharprime}{\kern0pt}{\isacharcomma}{\kern0pt}A{\isacharprime}{\kern0pt}{\isacharprime}{\kern0pt}{\isacharprime}{\kern0pt}{\isacharparenright}{\kern0pt}\ {\isasymin}\ set\ hs{\isacharprime}{\kern0pt}{\isachardoublequoteclose}\ {\isachardoublequoteopen}a\ {\isasymin}\ set\ A{\isacharprime}{\kern0pt}{\isacharprime}{\kern0pt}{\isacharprime}{\kern0pt}{\isachardoublequoteclose}\isanewline
\ \ \ \ \ \ \ \ \isacommand{by}\isamarkupfalse%
\ {\isacharparenleft}{\kern0pt}auto\ simp\ del{\isacharcolon}{\kern0pt}\ {\isacartoucheopen}hs{\isacharprime}{\kern0pt}\ {\isacharequal}{\kern0pt}\ hs{\isadigit{1}}{\isacharprime}{\kern0pt}\ {\isacharat}{\kern0pt}\ h{\isacharprime}{\kern0pt}\ {\isacharhash}{\kern0pt}\ hs{\isadigit{2}}{\isacharprime}{\kern0pt}{\isacartoucheclose}{\isacharparenright}{\kern0pt}\isanewline
\ \ \ \ \ \ \isacommand{hence}\isamarkupfalse%
\ {\isachardoublequoteopen}t\isactrlsub {\isadigit{1}}\ {\isacharless}{\kern0pt}\ t{\isacharprime}{\kern0pt}{\isacharprime}{\kern0pt}{\isacharprime}{\kern0pt}{\isachardoublequoteclose}\isanewline
\ \ \ \ \ \ \ \ \isacommand{by}\isamarkupfalse%
\ auto\isanewline
\ \ \ \ \ \ \isacommand{moreover}\isamarkupfalse%
\ \isacommand{have}\isamarkupfalse%
\ {\isachardoublequoteopen}t{\isacharprime}{\kern0pt}\ {\isacharless}{\kern0pt}\ t\isactrlsub {\isadigit{1}}{\isachardoublequoteclose}\isanewline
\ \ \ \ \ \ \ \ \isacommand{using}\isamarkupfalse%
\ {\isacartoucheopen}t\isactrlsub {\isadigit{1}}\ {\isasymin}\ htps{\isacharunderscore}{\kern0pt}IH{\isacartoucheclose}\ {\isacartoucheopen}{\isasymAnd}t{\isachardot}{\kern0pt}\ t\ {\isasymin}\ htps\ {\isasymLongrightarrow}\ {\isasymexists}A{\isachardot}{\kern0pt}\ {\isacharparenleft}{\kern0pt}t{\isacharcomma}{\kern0pt}\ A{\isacharparenright}{\kern0pt}\ {\isasymin}\ set\ {\isacharparenleft}{\kern0pt}h\ {\isacharhash}{\kern0pt}\ hs{\isadigit{2}}{\isacharparenright}{\kern0pt}{\isacartoucheclose}\ {\isacartoucheopen}strict{\isacharunderscore}{\kern0pt}sorted\ {\isacharparenleft}{\kern0pt}map\ fst\ hs{\isacharparenright}{\kern0pt}{\isacartoucheclose}\isanewline
\ \ \ \ \ \ \ \ \isacommand{by}\isamarkupfalse%
{\isacharparenleft}{\kern0pt}fastforce\ simp{\isacharcolon}{\kern0pt}\ htps{\isacharunderscore}{\kern0pt}IH{\isacharunderscore}{\kern0pt}def{\isacharparenright}{\kern0pt}\isanewline
\ \ \ \ \ \ \isacommand{hence}\isamarkupfalse%
\ {\isachardoublequoteopen}t\isactrlsub {\isadigit{1}}\ {\isasymin}\ set\ {\isacharparenleft}{\kern0pt}map\ fst\ hs{\isacharprime}{\kern0pt}{\isacharunderscore}{\kern0pt}IH{\isacharparenright}{\kern0pt}{\isachardoublequoteclose}\isanewline
\ \ \ \ \ \ \ \ \isacommand{using}\isamarkupfalse%
\ fst{\isacharunderscore}{\kern0pt}h{\isadigit{3}}{\isacharprime}{\kern0pt}{\isacharunderscore}{\kern0pt}n{\isacharunderscore}{\kern0pt}htp\ {\isacartoucheopen}t\isactrlsub {\isadigit{1}}\ {\isasymin}\ htps{\isacharunderscore}{\kern0pt}IH{\isacartoucheclose}\ {\isacartoucheopen}{\isasymAnd}t{\isacharprime}{\kern0pt}{\isacharprime}{\kern0pt}{\isachardot}{\kern0pt}\ t{\isacharprime}{\kern0pt}{\isacharprime}{\kern0pt}\ {\isasymin}\ htps{\isacharunderscore}{\kern0pt}IH\ {\isasymLongrightarrow}\ {\isasymexists}A{\isachardot}{\kern0pt}\ {\isacharparenleft}{\kern0pt}t{\isacharprime}{\kern0pt}{\isacharprime}{\kern0pt}{\isacharcomma}{\kern0pt}\ A{\isacharparenright}{\kern0pt}\ {\isasymin}\ set\ hs{\isacharprime}{\kern0pt}{\isacharunderscore}{\kern0pt}IH{\isacartoucheclose}\isanewline
\ \ \ \ \ \ \ \ \isacommand{by}\isamarkupfalse%
\ {\isacharparenleft}{\kern0pt}fastforce\ simp{\isacharcolon}{\kern0pt}\ image{\isacharunderscore}{\kern0pt}def\ hs{\isacharprime}{\kern0pt}{\isacharunderscore}{\kern0pt}IH{\isacharunderscore}{\kern0pt}def\ htps{\isacharunderscore}{\kern0pt}IH{\isacharunderscore}{\kern0pt}def{\isacharparenright}{\kern0pt}{\isacharplus}{\kern0pt}\isanewline
\ \ \ \ \ \ \isacommand{hence}\isamarkupfalse%
\ {\isachardoublequoteopen}{\isasymAnd}h{\isadigit{3}}{\isacharprime}{\kern0pt}{\isachardot}{\kern0pt}\ h{\isadigit{3}}{\isacharprime}{\kern0pt}\ {\isasymin}\ set\ hs{\isadigit{3}}{\isacharprime}{\kern0pt}\ {\isasymLongrightarrow}\ fst\ h{\isadigit{3}}{\isacharprime}{\kern0pt}\ {\isacharless}{\kern0pt}\ t\isactrlsub {\isadigit{1}}{\isachardoublequoteclose}\isanewline
\ \ \ \ \ \ \ \ \isacommand{using}\isamarkupfalse%
\ {\isacartoucheopen}strict{\isacharunderscore}{\kern0pt}sorted\ {\isacharparenleft}{\kern0pt}map\ fst\ hs{\isacharprime}{\kern0pt}{\isacharparenright}{\kern0pt}{\isacartoucheclose}\isanewline
\ \ \ \ \ \ \ \ \isacommand{by}\isamarkupfalse%
\ {\isacharparenleft}{\kern0pt}auto\ simp{\isacharcolon}{\kern0pt}\ hs{\isacharprime}{\kern0pt}{\isacharunderscore}{\kern0pt}IH{\isacharunderscore}{\kern0pt}def{\isacharparenright}{\kern0pt}\isanewline
\ \ \ \ \ \ \isacommand{ultimately}\isamarkupfalse%
\ \isacommand{have}\isamarkupfalse%
\ {\isachardoublequoteopen}t{\isacharprime}{\kern0pt}{\isacharprime}{\kern0pt}{\isacharprime}{\kern0pt}\ {\isasymnotin}\ set\ {\isacharparenleft}{\kern0pt}map\ fst\ hs{\isadigit{3}}{\isacharprime}{\kern0pt}{\isacharparenright}{\kern0pt}{\isachardoublequoteclose}\isanewline
\ \ \ \ \ \ \ \ \isacommand{by}\isamarkupfalse%
\ fastforce\isanewline
\ \ \ \ \ \ \isacommand{thus}\isamarkupfalse%
\ {\isacharquery}{\kern0pt}thesis\ \isanewline
\ \ \ \ \ \ \ \ \isacommand{using}\isamarkupfalse%
\ {\isacartoucheopen}{\isacharparenleft}{\kern0pt}t{\isacharprime}{\kern0pt}{\isacharprime}{\kern0pt}{\isacharprime}{\kern0pt}{\isacharcomma}{\kern0pt}A{\isacharprime}{\kern0pt}{\isacharprime}{\kern0pt}{\isacharprime}{\kern0pt}{\isacharparenright}{\kern0pt}\ {\isasymin}\ set\ hs{\isacharprime}{\kern0pt}{\isacartoucheclose}\ {\isacartoucheopen}t{\isacharprime}{\kern0pt}\ {\isacharless}{\kern0pt}\ t\isactrlsub {\isadigit{1}}{\isacartoucheclose}\ {\isacartoucheopen}t\isactrlsub {\isadigit{1}}\ {\isacharless}{\kern0pt}\ t{\isacharprime}{\kern0pt}{\isacharprime}{\kern0pt}{\isacharprime}{\kern0pt}{\isacartoucheclose}\ {\isacartoucheopen}t{\isacharprime}{\kern0pt}{\isacharprime}{\kern0pt}{\isacharprime}{\kern0pt}\ {\isasymin}\ {\isacharbraceleft}{\kern0pt}t\isactrlsub {\isadigit{1}}{\isacharless}{\kern0pt}{\isachardot}{\kern0pt}{\isachardot}{\kern0pt}{\isacharless}{\kern0pt}\ t\isactrlsub {\isadigit{2}}{\isacharbraceright}{\kern0pt}{\isacartoucheclose}\ {\isacartoucheopen}a\ {\isasymin}\ set\ A{\isacharprime}{\kern0pt}{\isacharprime}{\kern0pt}{\isacharprime}{\kern0pt}{\isacartoucheclose}\isanewline
\ \ \ \ \ \ \ \ \isacommand{by}\isamarkupfalse%
\ {\isacharparenleft}{\kern0pt}auto\ simp{\isacharcolon}{\kern0pt}\ hs{\isacharprime}{\kern0pt}{\isacharunderscore}{\kern0pt}IH{\isacharunderscore}{\kern0pt}def{\isacharparenright}{\kern0pt}{\isacharplus}{\kern0pt}\isanewline
\ \ \ \ \isacommand{qed}\isamarkupfalse%
\isanewline
\ \ \ \ \ \ \isanewline
\ \ \ \ \isacommand{moreover}\isamarkupfalse%
\ \isacommand{have}\isamarkupfalse%
\ {\isachardoublequoteopen}{\isasymexists}t{\isacharprime}{\kern0pt}{\isasymin}{\isacharbraceleft}{\kern0pt}t\isactrlsub {\isadigit{1}}{\isacharless}{\kern0pt}{\isachardot}{\kern0pt}{\isachardot}{\kern0pt}{\isacharless}{\kern0pt}\ t\isactrlsub {\isadigit{2}}{\isacharbraceright}{\kern0pt}{\isachardot}{\kern0pt}{\isasymexists}A{\isacharprime}{\kern0pt}{\isachardot}{\kern0pt}\ {\isacharparenleft}{\kern0pt}t{\isacharprime}{\kern0pt}{\isacharcomma}{\kern0pt}A{\isacharprime}{\kern0pt}{\isacharparenright}{\kern0pt}\ {\isasymin}\ set\ hs{\isacharunderscore}{\kern0pt}IH\ {\isasymand}\ a\ {\isasymin}\ set\ A{\isacharprime}{\kern0pt}{\isachardoublequoteclose}\isanewline
\ \ \ \ \ \ \isakeyword{if}\ {\isachardoublequoteopen}t{\isacharprime}{\kern0pt}{\isacharprime}{\kern0pt}\ {\isasymnotin}\ htps{\isacharunderscore}{\kern0pt}IH{\isachardoublequoteclose}\ {\isachardoublequoteopen}t\isactrlsub {\isadigit{1}}\ {\isasymin}\ htps{\isacharunderscore}{\kern0pt}IH{\isachardoublequoteclose}\ {\isachardoublequoteopen}t\isactrlsub {\isadigit{2}}\ {\isasymin}\ htps{\isacharunderscore}{\kern0pt}IH{\isachardoublequoteclose}\ {\isachardoublequoteopen}t{\isacharprime}{\kern0pt}{\isacharprime}{\kern0pt}\ {\isasymin}\ {\isacharbraceleft}{\kern0pt}t\isactrlsub {\isadigit{1}}{\isacharless}{\kern0pt}{\isachardot}{\kern0pt}{\isachardot}{\kern0pt}{\isacharless}{\kern0pt}\ t\isactrlsub {\isadigit{2}}{\isacharbraceright}{\kern0pt}{\isachardoublequoteclose}\ {\isachardoublequoteopen}{\isacharparenleft}{\kern0pt}t{\isacharprime}{\kern0pt}{\isacharprime}{\kern0pt}{\isacharcomma}{\kern0pt}\ A{\isacharprime}{\kern0pt}{\isacharprime}{\kern0pt}{\isacharparenright}{\kern0pt}\ {\isasymin}\ set\ hs{\isacharprime}{\kern0pt}{\isacharunderscore}{\kern0pt}IH{\isachardoublequoteclose}\isanewline
\ \ \ \ \ \ \ \ {\isachardoublequoteopen}a\ {\isasymin}\ set\ A{\isacharprime}{\kern0pt}{\isacharprime}{\kern0pt}{\isachardoublequoteclose}\isanewline
\ \ \ \ \ \ \isakeyword{for}\ t{\isacharprime}{\kern0pt}{\isacharprime}{\kern0pt}\ A{\isacharprime}{\kern0pt}{\isacharprime}{\kern0pt}\ t\isactrlsub {\isadigit{1}}\ t\isactrlsub {\isadigit{2}}\ a\isanewline
\ \ \ \ \isacommand{proof}\isamarkupfalse%
{\isacharminus}{\kern0pt}\isanewline
\ \ \ \ \ \ \isacommand{have}\isamarkupfalse%
\ {\isachardoublequoteopen}t{\isacharprime}{\kern0pt}{\isacharprime}{\kern0pt}\ {\isasymnotin}\ htps{\isachardoublequoteclose}\isanewline
\ \ \ \ \ \ \ \ \isacommand{using}\isamarkupfalse%
\ {\isacartoucheopen}{\isasymAnd}h{\isacharprime}{\kern0pt}{\isachardot}{\kern0pt}\ h{\isacharprime}{\kern0pt}\ {\isasymin}\ set\ {\isacharparenleft}{\kern0pt}hs{\isacharprime}{\kern0pt}{\isacharunderscore}{\kern0pt}IH{\isacharparenright}{\kern0pt}\ {\isasymLongrightarrow}\ t{\isacharprime}{\kern0pt}\ {\isasymnoteq}\ fst\ h{\isacharprime}{\kern0pt}{\isacartoucheclose}\ {\isacartoucheopen}t{\isacharprime}{\kern0pt}{\isacharprime}{\kern0pt}\ {\isasymnotin}\ htps{\isacharunderscore}{\kern0pt}IH{\isacartoucheclose}\ {\isacartoucheopen}{\isacharparenleft}{\kern0pt}t{\isacharprime}{\kern0pt}{\isacharprime}{\kern0pt}{\isacharcomma}{\kern0pt}\ A{\isacharprime}{\kern0pt}{\isacharprime}{\kern0pt}{\isacharparenright}{\kern0pt}\ {\isasymin}\ set\ hs{\isacharprime}{\kern0pt}{\isacharunderscore}{\kern0pt}IH{\isacartoucheclose}\isanewline
\ \ \ \ \ \ \ \ \isacommand{by}\isamarkupfalse%
\ {\isacharparenleft}{\kern0pt}fastforce\ simp{\isacharcolon}{\kern0pt}\ htps{\isacharunderscore}{\kern0pt}IH{\isacharunderscore}{\kern0pt}def{\isacharparenright}{\kern0pt}\ \ \ \ \ \ \ \ \isanewline
\ \ \ \ \ \ \isacommand{hence}\isamarkupfalse%
\ {\isachardoublequoteopen}{\isasymexists}t{\isacharprime}{\kern0pt}{\isasymin}{\isacharbraceleft}{\kern0pt}t\isactrlsub {\isadigit{1}}{\isacharless}{\kern0pt}{\isachardot}{\kern0pt}{\isachardot}{\kern0pt}{\isacharless}{\kern0pt}\ t\isactrlsub {\isadigit{2}}{\isacharbraceright}{\kern0pt}{\isachardot}{\kern0pt}{\isasymexists}A{\isacharprime}{\kern0pt}{\isachardot}{\kern0pt}\ {\isacharparenleft}{\kern0pt}t{\isacharprime}{\kern0pt}{\isacharcomma}{\kern0pt}A{\isacharprime}{\kern0pt}{\isacharparenright}{\kern0pt}\ {\isasymin}\ set\ hs\ {\isasymand}\ a\ {\isasymin}\ set\ A{\isacharprime}{\kern0pt}{\isachardoublequoteclose}\isanewline
\ \ \ \ \ \ \ \ \isacommand{using}\isamarkupfalse%
\ that\isanewline
\ \ \ \ \ \ \ \ \isacommand{by}\isamarkupfalse%
\ {\isacharparenleft}{\kern0pt}auto\ simp\ del{\isacharcolon}{\kern0pt}\ {\isacartoucheopen}hs\ {\isacharequal}{\kern0pt}\ hs{\isadigit{1}}\ {\isacharat}{\kern0pt}\ h\ {\isacharhash}{\kern0pt}\ hs{\isadigit{2}}{\isacartoucheclose}\ simp{\isacharcolon}{\kern0pt}\ htps{\isacharunderscore}{\kern0pt}IH{\isacharunderscore}{\kern0pt}def\ hs{\isacharprime}{\kern0pt}{\isacharunderscore}{\kern0pt}IH{\isacharunderscore}{\kern0pt}def\ intro{\isacharbang}{\kern0pt}{\isacharcolon}{\kern0pt}\ {\isadigit{2}}{\isacharparenleft}{\kern0pt}{\isadigit{1}}{\isadigit{2}}{\isacharparenright}{\kern0pt}{\isacharparenright}{\kern0pt}\isanewline
\ \ \ \ \ \ \isacommand{then}\isamarkupfalse%
\ \isacommand{obtain}\isamarkupfalse%
\ t{\isacharprime}{\kern0pt}{\isacharprime}{\kern0pt}{\isacharprime}{\kern0pt}\ A{\isacharprime}{\kern0pt}{\isacharprime}{\kern0pt}{\isacharprime}{\kern0pt}\ \isakeyword{where}\ {\isachardoublequoteopen}t{\isacharprime}{\kern0pt}{\isacharprime}{\kern0pt}{\isacharprime}{\kern0pt}\ {\isasymin}\ {\isacharbraceleft}{\kern0pt}t\isactrlsub {\isadigit{1}}{\isacharless}{\kern0pt}{\isachardot}{\kern0pt}{\isachardot}{\kern0pt}{\isacharless}{\kern0pt}\ t\isactrlsub {\isadigit{2}}{\isacharbraceright}{\kern0pt}{\isachardoublequoteclose}\ {\isachardoublequoteopen}{\isacharparenleft}{\kern0pt}t{\isacharprime}{\kern0pt}{\isacharprime}{\kern0pt}{\isacharprime}{\kern0pt}{\isacharcomma}{\kern0pt}A{\isacharprime}{\kern0pt}{\isacharprime}{\kern0pt}{\isacharprime}{\kern0pt}{\isacharparenright}{\kern0pt}\ {\isasymin}\ set\ hs{\isachardoublequoteclose}\ {\isachardoublequoteopen}a\ {\isasymin}\ set\ A{\isacharprime}{\kern0pt}{\isacharprime}{\kern0pt}{\isacharprime}{\kern0pt}{\isachardoublequoteclose}\isanewline
\ \ \ \ \ \ \ \ \isacommand{by}\isamarkupfalse%
\ {\isacharparenleft}{\kern0pt}auto\ simp\ del{\isacharcolon}{\kern0pt}\ {\isacartoucheopen}hs\ {\isacharequal}{\kern0pt}\ hs{\isadigit{1}}\ {\isacharat}{\kern0pt}\ h\ {\isacharhash}{\kern0pt}\ hs{\isadigit{2}}{\isacartoucheclose}{\isacharparenright}{\kern0pt}\isanewline
\ \ \ \ \ \ \isacommand{hence}\isamarkupfalse%
\ {\isachardoublequoteopen}t\isactrlsub {\isadigit{1}}\ {\isacharless}{\kern0pt}\ t{\isacharprime}{\kern0pt}{\isacharprime}{\kern0pt}{\isacharprime}{\kern0pt}{\isachardoublequoteclose}\isanewline
\ \ \ \ \ \ \ \ \isacommand{by}\isamarkupfalse%
\ auto\isanewline
\ \ \ \ \ \ \isacommand{moreover}\isamarkupfalse%
\ \isacommand{have}\isamarkupfalse%
\ {\isachardoublequoteopen}t{\isacharprime}{\kern0pt}\ {\isacharless}{\kern0pt}\ t\isactrlsub {\isadigit{1}}{\isachardoublequoteclose}\isanewline
\ \ \ \ \ \ \ \ \isacommand{using}\isamarkupfalse%
\ {\isacartoucheopen}t\isactrlsub {\isadigit{1}}\ {\isasymin}\ htps{\isacharunderscore}{\kern0pt}IH{\isacartoucheclose}\ {\isacartoucheopen}{\isasymAnd}t{\isachardot}{\kern0pt}\ t\ {\isasymin}\ htps\ {\isasymLongrightarrow}\ {\isasymexists}A{\isachardot}{\kern0pt}\ {\isacharparenleft}{\kern0pt}t{\isacharcomma}{\kern0pt}\ A{\isacharparenright}{\kern0pt}\ {\isasymin}\ set\ {\isacharparenleft}{\kern0pt}h\ {\isacharhash}{\kern0pt}\ hs{\isadigit{2}}{\isacharparenright}{\kern0pt}{\isacartoucheclose}\ {\isacartoucheopen}strict{\isacharunderscore}{\kern0pt}sorted\ {\isacharparenleft}{\kern0pt}map\ fst\ hs{\isacharparenright}{\kern0pt}{\isacartoucheclose}\isanewline
\ \ \ \ \ \ \ \ \isacommand{by}\isamarkupfalse%
{\isacharparenleft}{\kern0pt}fastforce\ simp{\isacharcolon}{\kern0pt}\ htps{\isacharunderscore}{\kern0pt}IH{\isacharunderscore}{\kern0pt}def{\isacharparenright}{\kern0pt}\isanewline
\ \ \ \ \ \ \isacommand{hence}\isamarkupfalse%
\ {\isachardoublequoteopen}t\isactrlsub {\isadigit{1}}\ {\isasymin}\ set\ {\isacharparenleft}{\kern0pt}map\ fst\ hs{\isacharunderscore}{\kern0pt}IH{\isacharparenright}{\kern0pt}{\isachardoublequoteclose}\isanewline
\ \ \ \ \ \ \ \ \isacommand{using}\isamarkupfalse%
\ fst{\isacharunderscore}{\kern0pt}h{\isadigit{3}}{\isacharprime}{\kern0pt}{\isacharunderscore}{\kern0pt}n{\isacharunderscore}{\kern0pt}htp\ {\isacartoucheopen}t\isactrlsub {\isadigit{1}}\ {\isasymin}\ htps{\isacharunderscore}{\kern0pt}IH{\isacartoucheclose}\ {\isacartoucheopen}hs{\isadigit{1}}\ {\isacharequal}{\kern0pt}\ {\isacharbrackleft}{\kern0pt}{\isacharbrackright}{\kern0pt}{\isacartoucheclose}\ {\isacartoucheopen}{\isasymAnd}t{\isacharprime}{\kern0pt}{\isacharprime}{\kern0pt}{\isachardot}{\kern0pt}\ t{\isacharprime}{\kern0pt}{\isacharprime}{\kern0pt}\ {\isasymin}\ htps{\isacharunderscore}{\kern0pt}IH\ {\isasymLongrightarrow}\ {\isasymexists}A{\isachardot}{\kern0pt}\ {\isacharparenleft}{\kern0pt}t{\isacharprime}{\kern0pt}{\isacharprime}{\kern0pt}{\isacharcomma}{\kern0pt}\ A{\isacharparenright}{\kern0pt}\ {\isasymin}\ set\ hs{\isacharunderscore}{\kern0pt}IH{\isacartoucheclose}\isanewline
\ \ \ \ \ \ \ \ \isacommand{by}\isamarkupfalse%
{\isacharparenleft}{\kern0pt}fastforce\ simp{\isacharcolon}{\kern0pt}\ image{\isacharunderscore}{\kern0pt}def\ hs{\isacharprime}{\kern0pt}{\isacharunderscore}{\kern0pt}IH{\isacharunderscore}{\kern0pt}def\ htps{\isacharunderscore}{\kern0pt}IH{\isacharunderscore}{\kern0pt}def{\isacharparenright}{\kern0pt}{\isacharplus}{\kern0pt}\isanewline
\ \ \ \ \ \ \isacommand{hence}\isamarkupfalse%
\ {\isachardoublequoteopen}{\isasymAnd}h{\isadigit{3}}{\isacharprime}{\kern0pt}{\isachardot}{\kern0pt}\ h{\isadigit{3}}{\isacharprime}{\kern0pt}\ {\isasymin}\ set\ hs{\isadigit{3}}\ {\isasymLongrightarrow}\ fst\ h{\isadigit{3}}{\isacharprime}{\kern0pt}\ {\isacharless}{\kern0pt}\ t\isactrlsub {\isadigit{1}}{\isachardoublequoteclose}\isanewline
\ \ \ \ \ \ \ \ \isacommand{using}\isamarkupfalse%
\ {\isacartoucheopen}strict{\isacharunderscore}{\kern0pt}sorted\ {\isacharparenleft}{\kern0pt}map\ fst\ hs{\isacharparenright}{\kern0pt}{\isacartoucheclose}\isanewline
\ \ \ \ \ \ \ \ \isacommand{by}\isamarkupfalse%
\ {\isacharparenleft}{\kern0pt}auto\ simp{\isacharcolon}{\kern0pt}\ hs{\isacharunderscore}{\kern0pt}IH{\isacharunderscore}{\kern0pt}def{\isacharparenright}{\kern0pt}\isanewline
\ \ \ \ \ \ \isacommand{ultimately}\isamarkupfalse%
\ \isacommand{have}\isamarkupfalse%
\ {\isachardoublequoteopen}t{\isacharprime}{\kern0pt}{\isacharprime}{\kern0pt}{\isacharprime}{\kern0pt}\ {\isasymnotin}\ set\ {\isacharparenleft}{\kern0pt}map\ fst\ hs{\isadigit{3}}{\isacharparenright}{\kern0pt}{\isachardoublequoteclose}\isanewline
\ \ \ \ \ \ \ \ \isacommand{by}\isamarkupfalse%
\ fastforce\isanewline
\ \ \ \ \ \ \isacommand{thus}\isamarkupfalse%
\ {\isacharquery}{\kern0pt}thesis\isanewline
\ \ \ \ \ \ \ \ \isacommand{using}\isamarkupfalse%
\ {\isacartoucheopen}{\isacharparenleft}{\kern0pt}t{\isacharprime}{\kern0pt}{\isacharprime}{\kern0pt}{\isacharprime}{\kern0pt}{\isacharcomma}{\kern0pt}A{\isacharprime}{\kern0pt}{\isacharprime}{\kern0pt}{\isacharprime}{\kern0pt}{\isacharparenright}{\kern0pt}\ {\isasymin}\ set\ hs{\isacartoucheclose}\ {\isacartoucheopen}t{\isacharprime}{\kern0pt}{\isacharprime}{\kern0pt}{\isacharprime}{\kern0pt}\ {\isasymin}\ {\isacharbraceleft}{\kern0pt}t\isactrlsub {\isadigit{1}}{\isacharless}{\kern0pt}{\isachardot}{\kern0pt}{\isachardot}{\kern0pt}{\isacharless}{\kern0pt}\ t\isactrlsub {\isadigit{2}}{\isacharbraceright}{\kern0pt}{\isacartoucheclose}\ {\isacartoucheopen}a\ {\isasymin}\ set\ A{\isacharprime}{\kern0pt}{\isacharprime}{\kern0pt}{\isacharprime}{\kern0pt}{\isacartoucheclose}\ {\isacartoucheopen}t{\isacharprime}{\kern0pt}\ {\isacharless}{\kern0pt}\ t\isactrlsub {\isadigit{1}}{\isacartoucheclose}\ {\isacartoucheopen}t\isactrlsub {\isadigit{1}}\ {\isacharless}{\kern0pt}\ t{\isacharprime}{\kern0pt}{\isacharprime}{\kern0pt}{\isacharprime}{\kern0pt}{\isacartoucheclose}\isanewline
\ \ \ \ \ \ \ \ \isacommand{by}\isamarkupfalse%
\ {\isacharparenleft}{\kern0pt}auto\ simp{\isacharcolon}{\kern0pt}\ hs{\isacharunderscore}{\kern0pt}IH{\isacharunderscore}{\kern0pt}def{\isacharparenright}{\kern0pt}\isanewline
\ \ \ \ \isacommand{qed}\isamarkupfalse%
\isanewline
\ \ \ \ \isacommand{moreover}\isamarkupfalse%
\ \isacommand{have}\isamarkupfalse%
\ {\isachardoublequoteopen}valid{\isacharunderscore}{\kern0pt}happ{\isacharunderscore}{\kern0pt}seq\ M{\isacharunderscore}{\kern0pt}IH\ hs{\isacharunderscore}{\kern0pt}IH\ M{\isacharprime}{\kern0pt}{\isachardoublequoteclose}\isanewline
\ \ \ \ \ \ \isacommand{using}\isamarkupfalse%
\ {\isacartoucheopen}valid{\isacharunderscore}{\kern0pt}happ{\isacharunderscore}{\kern0pt}seq\ M\ hs\ M{\isacharprime}{\kern0pt}{\isacartoucheclose}\ {\isacartoucheopen}valid{\isacharunderscore}{\kern0pt}happ{\isacharunderscore}{\kern0pt}seq\ M\ {\isacharbrackleft}{\kern0pt}h{\isacharbrackright}{\kern0pt}\ {\isacharparenleft}{\kern0pt}apply{\isacharunderscore}{\kern0pt}happ\ {\isacharparenleft}{\kern0pt}t{\isacharcomma}{\kern0pt}\ A{\isacharprime}{\kern0pt}{\isacharparenright}{\kern0pt}\ M{\isacharparenright}{\kern0pt}{\isacartoucheclose}\isanewline
\ \ \ \ \ \ \ \ {\isacartoucheopen}{\isacharparenleft}{\kern0pt}{\isasymforall}h{\isasymin}set\ hs{\isadigit{3}}{\isachardot}{\kern0pt}\ fst\ h\ {\isasymnotin}\ htps{\isacharparenright}{\kern0pt}{\isacartoucheclose}\isanewline
\ \ \ \ \ \ \isacommand{by}\isamarkupfalse%
\ {\isacharparenleft}{\kern0pt}auto\ simp{\isacharcolon}{\kern0pt}\ M{\isacharunderscore}{\kern0pt}IH{\isacharunderscore}{\kern0pt}def\ hs{\isacharunderscore}{\kern0pt}IH{\isacharunderscore}{\kern0pt}def\ intro{\isacharcolon}{\kern0pt}\ {\isadigit{2}}{\isacharparenleft}{\kern0pt}{\isadigit{9}}{\isacharparenright}{\kern0pt}\ dest{\isacharbang}{\kern0pt}{\isacharcolon}{\kern0pt}\ valid{\isacharunderscore}{\kern0pt}happ{\isacharunderscore}{\kern0pt}seq{\isacharunderscore}{\kern0pt}append{\isacharunderscore}{\kern0pt}{\isadigit{2}}{\isacharparenright}{\kern0pt}\isanewline
\ \ \ \ \isacommand{moreover}\isamarkupfalse%
\ \isacommand{have}\isamarkupfalse%
\ {\isachardoublequoteopen}{\isacharparenleft}{\kern0pt}fst\ {\isacharparenleft}{\kern0pt}hd\ hs{\isacharunderscore}{\kern0pt}IH{\isacharparenright}{\kern0pt}{\isacharparenright}{\kern0pt}\ {\isasymin}\ htps{\isacharunderscore}{\kern0pt}IH{\isachardoublequoteclose}\isanewline
\ \ \ \ \ \ \isacommand{using}\isamarkupfalse%
\ {\isacartoucheopen}fst\ h{\isadigit{2}}\ {\isasymin}\ htps{\isacartoucheclose}\isanewline
\ \ \ \ \ \ \isacommand{by}\isamarkupfalse%
\ {\isacharparenleft}{\kern0pt}auto\ simp{\isacharcolon}{\kern0pt}\ hs{\isacharunderscore}{\kern0pt}IH{\isacharunderscore}{\kern0pt}def\ htps{\isacharunderscore}{\kern0pt}IH{\isacharunderscore}{\kern0pt}def\ {\isacartoucheopen}{\isasymAnd}h{\isacharprime}{\kern0pt}{\isachardot}{\kern0pt}\ h{\isacharprime}{\kern0pt}\ {\isasymin}\ set\ hs{\isacharunderscore}{\kern0pt}IH\ {\isasymLongrightarrow}\ t\ {\isasymnoteq}\ fst\ h{\isacharprime}{\kern0pt}{\isacartoucheclose}\ simp\ del{\isacharcolon}{\kern0pt}\ {\isacartoucheopen}t\ {\isacharequal}{\kern0pt}\ t{\isacharprime}{\kern0pt}{\isacartoucheclose}{\isacharparenright}{\kern0pt}\isanewline
\ \ \ \ \isacommand{moreover}\isamarkupfalse%
\ \isacommand{have}\isamarkupfalse%
\ {\isachardoublequoteopen}{\isacharparenleft}{\kern0pt}fst\ {\isacharparenleft}{\kern0pt}hd\ hs{\isacharprime}{\kern0pt}{\isacharunderscore}{\kern0pt}IH{\isacharparenright}{\kern0pt}{\isacharparenright}{\kern0pt}\ {\isasymin}\ htps{\isacharunderscore}{\kern0pt}IH{\isachardoublequoteclose}\isanewline
\ \ \ \ \ \ \isacommand{using}\isamarkupfalse%
\ {\isacartoucheopen}fst\ h{\isadigit{2}}\ {\isasymin}\ htps{\isacartoucheclose}\ {\isacartoucheopen}fst\ h{\isadigit{2}}\ {\isasymnoteq}\ fst\ h{\isacharprime}{\kern0pt}{\isacartoucheclose}\isanewline
\ \ \ \ \ \ \isacommand{by}\isamarkupfalse%
\ {\isacharparenleft}{\kern0pt}auto\ simp{\isacharcolon}{\kern0pt}\ hs{\isacharprime}{\kern0pt}{\isacharunderscore}{\kern0pt}IH{\isacharunderscore}{\kern0pt}def\ htps{\isacharunderscore}{\kern0pt}IH{\isacharunderscore}{\kern0pt}def{\isacharparenright}{\kern0pt}\isanewline
\ \ \ \ \isacommand{moreover}\isamarkupfalse%
\ \isacommand{have}\isamarkupfalse%
\ {\isachardoublequoteopen}{\isacharparenleft}{\kern0pt}fst\ {\isacharparenleft}{\kern0pt}last\ hs{\isacharunderscore}{\kern0pt}IH{\isacharparenright}{\kern0pt}{\isacharparenright}{\kern0pt}\ {\isasymin}\ htps{\isacharunderscore}{\kern0pt}IH{\isachardoublequoteclose}\isanewline
\ \ \ \ \ \ \isacommand{using}\isamarkupfalse%
\ {\isacartoucheopen}hs\ {\isasymnoteq}\ {\isacharbrackleft}{\kern0pt}{\isacharbrackright}{\kern0pt}\ {\isasymLongrightarrow}\ fst\ {\isacharparenleft}{\kern0pt}last\ hs{\isacharparenright}{\kern0pt}\ {\isasymin}\ htps{\isacartoucheclose}\isanewline
\ \ \ \ \ \ \isacommand{by}\isamarkupfalse%
\ {\isacharparenleft}{\kern0pt}auto\ simp{\isacharcolon}{\kern0pt}\ hs{\isacharunderscore}{\kern0pt}IH{\isacharunderscore}{\kern0pt}def\ htps{\isacharunderscore}{\kern0pt}IH{\isacharunderscore}{\kern0pt}def\ {\isacartoucheopen}{\isasymAnd}h{\isacharprime}{\kern0pt}{\isachardot}{\kern0pt}\ h{\isacharprime}{\kern0pt}\ {\isasymin}\ set\ hs{\isacharunderscore}{\kern0pt}IH\ {\isasymLongrightarrow}\ t\ {\isasymnoteq}\ fst\ h{\isacharprime}{\kern0pt}{\isacartoucheclose}\ simp\ del{\isacharcolon}{\kern0pt}\ {\isacartoucheopen}t\ {\isacharequal}{\kern0pt}\ t{\isacharprime}{\kern0pt}{\isacartoucheclose}{\isacharparenright}{\kern0pt}\isanewline
\ \ \ \ \isacommand{moreover}\isamarkupfalse%
\ \isacommand{have}\isamarkupfalse%
\ {\isachardoublequoteopen}{\isacharparenleft}{\kern0pt}fst\ {\isacharparenleft}{\kern0pt}last\ hs{\isacharprime}{\kern0pt}{\isacharunderscore}{\kern0pt}IH{\isacharparenright}{\kern0pt}{\isacharparenright}{\kern0pt}\ {\isasymin}\ htps{\isacharunderscore}{\kern0pt}IH{\isachardoublequoteclose}\isanewline
\ \ \ \ \ \ \isacommand{using}\isamarkupfalse%
\ {\isacartoucheopen}fst\ h{\isadigit{2}}\ {\isasymin}\ htps{\isacartoucheclose}\ \ {\isacartoucheopen}hs{\isacharprime}{\kern0pt}\ {\isasymnoteq}\ {\isacharbrackleft}{\kern0pt}{\isacharbrackright}{\kern0pt}\ {\isasymLongrightarrow}\ fst\ {\isacharparenleft}{\kern0pt}last\ hs{\isacharprime}{\kern0pt}{\isacharparenright}{\kern0pt}\ {\isasymin}\ htps{\isacartoucheclose}\ {\isacartoucheopen}{\isasymAnd}h{\isacharprime}{\kern0pt}{\isachardot}{\kern0pt}\ h{\isacharprime}{\kern0pt}\ {\isasymin}\ set\ hs{\isacharprime}{\kern0pt}{\isacharunderscore}{\kern0pt}IH\ {\isasymLongrightarrow}\ t{\isacharprime}{\kern0pt}\ {\isasymnoteq}\ fst\ h{\isacharprime}{\kern0pt}{\isacartoucheclose}\isanewline
\ \ \ \ \ \ \isacommand{by}\isamarkupfalse%
\ {\isacharparenleft}{\kern0pt}auto\ simp{\isacharcolon}{\kern0pt}\ htps{\isacharunderscore}{\kern0pt}IH{\isacharunderscore}{\kern0pt}def\ hs{\isacharprime}{\kern0pt}{\isacharunderscore}{\kern0pt}IH{\isacharunderscore}{\kern0pt}def\ split{\isacharcolon}{\kern0pt}\ if{\isacharunderscore}{\kern0pt}splits{\isacharparenright}{\kern0pt}{\isacharplus}{\kern0pt}\isanewline
\ \ \ \ \isacommand{ultimately}\isamarkupfalse%
\ \isacommand{have}\isamarkupfalse%
\ {\isachardoublequoteopen}valid{\isacharunderscore}{\kern0pt}happ{\isacharunderscore}{\kern0pt}seq\ M{\isacharunderscore}{\kern0pt}IH\ hs{\isacharprime}{\kern0pt}{\isacharunderscore}{\kern0pt}IH\ M{\isacharprime}{\kern0pt}{\isachardoublequoteclose}\isanewline
\ \ \ \ \ \ \isacommand{by}\isamarkupfalse%
\ {\isacharparenleft}{\kern0pt}rule\ {\isadigit{2}}{\isacharparenright}{\kern0pt}\isanewline
\ \ \ \ \isacommand{moreover}\isamarkupfalse%
\ \isacommand{have}\isamarkupfalse%
\ {\isachardoublequoteopen}{\isacharparenleft}{\kern0pt}apply{\isacharunderscore}{\kern0pt}happ\ {\isacharparenleft}{\kern0pt}t{\isacharcomma}{\kern0pt}\ A{\isacharprime}{\kern0pt}{\isacharparenright}{\kern0pt}\ M{\isacharparenright}{\kern0pt}\ \isactrlsup c{\isasymTTurnstile}\isactrlsub {\isacharequal}{\kern0pt}\ precondition\ a{\isachardoublequoteclose}\isanewline
\ \ \ \ \ \ \isakeyword{if}\ {\isachardoublequoteopen}a\ {\isasymin}\ set\ A{\isacharprime}{\kern0pt}{\isacharprime}{\kern0pt}{\isachardoublequoteclose}\ {\isachardoublequoteopen}{\isacharparenleft}{\kern0pt}t{\isacharprime}{\kern0pt}{\isacharprime}{\kern0pt}{\isacharcomma}{\kern0pt}A{\isacharprime}{\kern0pt}{\isacharprime}{\kern0pt}{\isacharparenright}{\kern0pt}\ {\isasymin}\ set\ hs{\isadigit{3}}{\isacharprime}{\kern0pt}{\isachardoublequoteclose}\isanewline
\ \ \ \ \ \ \isakeyword{for}\ t{\isacharprime}{\kern0pt}{\isacharprime}{\kern0pt}\ A{\isacharprime}{\kern0pt}{\isacharprime}{\kern0pt}\ a\isanewline
\ \ \ \ \isacommand{proof}\isamarkupfalse%
{\isacharminus}{\kern0pt}\isanewline
\ \ \ \ \ \ \isacommand{have}\isamarkupfalse%
\ {\isachardoublequoteopen}t{\isacharprime}{\kern0pt}{\isacharprime}{\kern0pt}\ {\isasymnotin}\ htps{\isachardoublequoteclose}\isanewline
\ \ \ \ \ \ \ \ \isacommand{using}\isamarkupfalse%
\ {\isacartoucheopen}{\isacharparenleft}{\kern0pt}t{\isacharprime}{\kern0pt}{\isacharprime}{\kern0pt}{\isacharcomma}{\kern0pt}A{\isacharprime}{\kern0pt}{\isacharprime}{\kern0pt}{\isacharparenright}{\kern0pt}\ {\isasymin}\ set\ hs{\isadigit{3}}{\isacharprime}{\kern0pt}{\isacartoucheclose}\ fst{\isacharunderscore}{\kern0pt}h{\isadigit{3}}{\isacharprime}{\kern0pt}{\isacharunderscore}{\kern0pt}n{\isacharunderscore}{\kern0pt}htp\isanewline
\ \ \ \ \ \ \ \ \isacommand{by}\isamarkupfalse%
\ auto\isanewline
\ \ \ \ \ \ \isacommand{moreover}\isamarkupfalse%
\ \isacommand{have}\isamarkupfalse%
\ {\isachardoublequoteopen}t\ {\isasymin}\ htps{\isachardoublequoteclose}\isanewline
\ \ \ \ \ \ \ \ \isacommand{using}\isamarkupfalse%
\ {\isacartoucheopen}t\ {\isasymin}\ htps{\isacartoucheclose}\ \isacommand{{\isachardot}{\kern0pt}}\isamarkupfalse%
\isanewline
\ \ \ \ \ \ \isacommand{moreover}\isamarkupfalse%
\ \isacommand{have}\isamarkupfalse%
\ {\isachardoublequoteopen}fst\ h{\isadigit{2}}\ {\isasymin}\ htps{\isachardoublequoteclose}\isanewline
\ \ \ \ \ \ \ \ \isacommand{using}\isamarkupfalse%
\ {\isacartoucheopen}fst\ h{\isadigit{2}}\ {\isasymin}\ htps{\isacartoucheclose}\ \isacommand{{\isachardot}{\kern0pt}}\isamarkupfalse%
\isanewline
\ \ \ \ \ \ \isacommand{moreover}\isamarkupfalse%
\ \isacommand{have}\isamarkupfalse%
\ {\isachardoublequoteopen}t{\isacharprime}{\kern0pt}{\isacharprime}{\kern0pt}\ {\isasymin}\ {\isacharbraceleft}{\kern0pt}t{\isacharless}{\kern0pt}{\isachardot}{\kern0pt}{\isachardot}{\kern0pt}{\isacharless}{\kern0pt}fst\ h{\isadigit{2}}{\isacharbraceright}{\kern0pt}{\isachardoublequoteclose}\isanewline
\ \ \ \ \ \ \ \ \isacommand{using}\isamarkupfalse%
\ {\isacartoucheopen}{\isacharparenleft}{\kern0pt}t{\isacharprime}{\kern0pt}{\isacharprime}{\kern0pt}{\isacharcomma}{\kern0pt}A{\isacharprime}{\kern0pt}{\isacharprime}{\kern0pt}{\isacharparenright}{\kern0pt}\ {\isasymin}\ set\ hs{\isadigit{3}}{\isacharprime}{\kern0pt}{\isacartoucheclose}\ {\isacartoucheopen}strict{\isacharunderscore}{\kern0pt}sorted\ {\isacharparenleft}{\kern0pt}map\ fst\ hs{\isacharprime}{\kern0pt}{\isacharparenright}{\kern0pt}{\isacartoucheclose}\isanewline
\ \ \ \ \ \ \ \ \isacommand{by}\isamarkupfalse%
\ auto\isanewline
\ \ \ \ \ \ \isacommand{moreover}\isamarkupfalse%
\ \isacommand{have}\isamarkupfalse%
\ {\isachardoublequoteopen}{\isacharparenleft}{\kern0pt}t{\isacharprime}{\kern0pt}{\isacharprime}{\kern0pt}{\isacharcomma}{\kern0pt}A{\isacharprime}{\kern0pt}{\isacharprime}{\kern0pt}{\isacharparenright}{\kern0pt}\ {\isasymin}\ set\ hs{\isacharprime}{\kern0pt}{\isachardoublequoteclose}\isanewline
\ \ \ \ \ \ \ \ \isacommand{using}\isamarkupfalse%
\ {\isacartoucheopen}{\isacharparenleft}{\kern0pt}t{\isacharprime}{\kern0pt}{\isacharprime}{\kern0pt}{\isacharcomma}{\kern0pt}A{\isacharprime}{\kern0pt}{\isacharprime}{\kern0pt}{\isacharparenright}{\kern0pt}\ {\isasymin}\ set\ hs{\isadigit{3}}{\isacharprime}{\kern0pt}{\isacartoucheclose}\ {\isacartoucheopen}strict{\isacharunderscore}{\kern0pt}sorted\ {\isacharparenleft}{\kern0pt}map\ fst\ hs{\isacharprime}{\kern0pt}{\isacharparenright}{\kern0pt}{\isacartoucheclose}\isanewline
\ \ \ \ \ \ \ \ \isacommand{by}\isamarkupfalse%
\ auto\isanewline
\ \ \ \ \ \ \isacommand{ultimately}\isamarkupfalse%
\ \isacommand{have}\isamarkupfalse%
\ {\isachardoublequoteopen}{\isasymexists}t{\isacharprime}{\kern0pt}{\isasymin}{\isacharbraceleft}{\kern0pt}t{\isacharless}{\kern0pt}{\isachardot}{\kern0pt}{\isachardot}{\kern0pt}{\isacharless}{\kern0pt}fst\ h{\isadigit{2}}{\isacharbraceright}{\kern0pt}{\isachardot}{\kern0pt}\ {\isasymexists}A{\isacharprime}{\kern0pt}{\isachardot}{\kern0pt}\ {\isacharparenleft}{\kern0pt}t{\isacharprime}{\kern0pt}{\isacharcomma}{\kern0pt}\ A{\isacharprime}{\kern0pt}{\isacharparenright}{\kern0pt}\ {\isasymin}\ set\ hs\ {\isasymand}\ a\ {\isasymin}\ set\ A{\isacharprime}{\kern0pt}{\isachardoublequoteclose}\isanewline
\ \ \ \ \ \ \ \ \isacommand{using}\isamarkupfalse%
\ {\isacartoucheopen}a\ {\isasymin}\ set\ A{\isacharprime}{\kern0pt}{\isacharprime}{\kern0pt}{\isacartoucheclose}\ \isanewline
\ \ \ \ \ \ \ \ \isacommand{by}\isamarkupfalse%
{\isacharparenleft}{\kern0pt}rule\ {\isadigit{2}}{\isacharparenright}{\kern0pt}\ \ \ \ \ \ \isanewline
\ \ \ \ \ \ \isacommand{then}\isamarkupfalse%
\ \isacommand{obtain}\isamarkupfalse%
\ t{\isacharprime}{\kern0pt}{\isacharprime}{\kern0pt}{\isacharprime}{\kern0pt}\ A{\isacharprime}{\kern0pt}{\isacharprime}{\kern0pt}{\isacharprime}{\kern0pt}\ \isakeyword{where}\ {\isachardoublequoteopen}t{\isacharprime}{\kern0pt}{\isacharprime}{\kern0pt}{\isacharprime}{\kern0pt}\ {\isasymin}\ {\isacharbraceleft}{\kern0pt}t{\isacharless}{\kern0pt}{\isachardot}{\kern0pt}{\isachardot}{\kern0pt}{\isacharless}{\kern0pt}fst\ h{\isadigit{2}}{\isacharbraceright}{\kern0pt}{\isachardoublequoteclose}\ {\isachardoublequoteopen}{\isacharparenleft}{\kern0pt}t{\isacharprime}{\kern0pt}{\isacharprime}{\kern0pt}{\isacharprime}{\kern0pt}{\isacharcomma}{\kern0pt}A{\isacharprime}{\kern0pt}{\isacharprime}{\kern0pt}{\isacharprime}{\kern0pt}{\isacharparenright}{\kern0pt}\ {\isasymin}\ set\ hs{\isachardoublequoteclose}\ {\isachardoublequoteopen}a\ {\isasymin}\ set\ A{\isacharprime}{\kern0pt}{\isacharprime}{\kern0pt}{\isacharprime}{\kern0pt}{\isachardoublequoteclose}\isanewline
\ \ \ \ \ \ \ \ \isacommand{by}\isamarkupfalse%
\ {\isacharparenleft}{\kern0pt}auto\ simp\ del{\isacharcolon}{\kern0pt}\ {\isacartoucheopen}hs\ {\isacharequal}{\kern0pt}\ hs{\isadigit{1}}\ {\isacharat}{\kern0pt}\ h\ {\isacharhash}{\kern0pt}\ hs{\isadigit{2}}{\isacartoucheclose}{\isacharparenright}{\kern0pt}\isanewline
\ \ \ \ \ \ \isacommand{hence}\isamarkupfalse%
\ {\isachardoublequoteopen}{\isacharparenleft}{\kern0pt}t{\isacharprime}{\kern0pt}{\isacharprime}{\kern0pt}{\isacharprime}{\kern0pt}{\isacharcomma}{\kern0pt}A{\isacharprime}{\kern0pt}{\isacharprime}{\kern0pt}{\isacharprime}{\kern0pt}{\isacharparenright}{\kern0pt}\ {\isasymin}\ set\ hs{\isadigit{3}}{\isachardoublequoteclose}\isanewline
\ \ \ \ \ \ \ \ \isacommand{using}\isamarkupfalse%
\ {\isacartoucheopen}strict{\isacharunderscore}{\kern0pt}sorted\ {\isacharparenleft}{\kern0pt}map\ fst\ hs{\isacharparenright}{\kern0pt}{\isacartoucheclose}\isanewline
\ \ \ \ \ \ \ \ \isacommand{by}\isamarkupfalse%
\ auto\isanewline
\ \ \ \ \ \ \isacommand{moreover}\isamarkupfalse%
\ \isacommand{have}\isamarkupfalse%
\ {\isachardoublequoteopen}valid{\isacharunderscore}{\kern0pt}happ{\isacharunderscore}{\kern0pt}seq\ {\isacharparenleft}{\kern0pt}apply{\isacharunderscore}{\kern0pt}happ\ {\isacharparenleft}{\kern0pt}t{\isacharcomma}{\kern0pt}\ A{\isacharprime}{\kern0pt}{\isacharparenright}{\kern0pt}\ M{\isacharparenright}{\kern0pt}\ {\isacharparenleft}{\kern0pt}hs{\isadigit{3}}\ {\isacharat}{\kern0pt}\ h{\isadigit{2}}\ {\isacharhash}{\kern0pt}\ hs{\isadigit{4}}{\isacharparenright}{\kern0pt}\ M{\isacharprime}{\kern0pt}{\isachardoublequoteclose}\isanewline
\ \ \ \ \ \ \ \ \isacommand{using}\isamarkupfalse%
\ {\isacartoucheopen}valid{\isacharunderscore}{\kern0pt}happ{\isacharunderscore}{\kern0pt}seq\ M\ hs\ M{\isacharprime}{\kern0pt}{\isacartoucheclose}\ {\isacartoucheopen}valid{\isacharunderscore}{\kern0pt}happ{\isacharunderscore}{\kern0pt}seq\ M\ {\isacharbrackleft}{\kern0pt}h{\isacharbrackright}{\kern0pt}\ {\isacharparenleft}{\kern0pt}apply{\isacharunderscore}{\kern0pt}happ\ {\isacharparenleft}{\kern0pt}t{\isacharcomma}{\kern0pt}\ A{\isacharprime}{\kern0pt}{\isacharparenright}{\kern0pt}\ M{\isacharparenright}{\kern0pt}{\isacartoucheclose}\isanewline
\ \ \ \ \ \ \ \ \isacommand{by}\isamarkupfalse%
\ auto\isanewline
\ \ \ \ \ \ \isacommand{ultimately}\isamarkupfalse%
\ \isacommand{show}\isamarkupfalse%
\ {\isacharquery}{\kern0pt}thesis\isanewline
\ \ \ \ \ \ \ \ \isacommand{using}\isamarkupfalse%
\ {\isadigit{2}}{\isacharparenleft}{\kern0pt}{\isadigit{9}}{\isacharparenright}{\kern0pt}\ {\isacartoucheopen}{\isacharparenleft}{\kern0pt}{\isasymforall}h{\isasymin}set\ hs{\isadigit{3}}{\isachardot}{\kern0pt}\ fst\ h\ {\isasymnotin}\ htps{\isacharparenright}{\kern0pt}{\isacartoucheclose}\ {\isacartoucheopen}a\ {\isasymin}\ set\ A{\isacharprime}{\kern0pt}{\isacharprime}{\kern0pt}{\isacharprime}{\kern0pt}{\isacartoucheclose}\isanewline
\ \ \ \ \ \ \ \ \isacommand{by}\isamarkupfalse%
\ {\isacharparenleft}{\kern0pt}force\ simp{\isacharcolon}{\kern0pt}\ intro{\isacharbang}{\kern0pt}{\isacharcolon}{\kern0pt}\ valid{\isacharunderscore}{\kern0pt}happ{\isacharunderscore}{\kern0pt}seq{\isacharunderscore}{\kern0pt}append{\isacharunderscore}{\kern0pt}{\isadigit{3}}{\isacharparenright}{\kern0pt}\isanewline
\ \ \ \ \isacommand{qed}\isamarkupfalse%
\isanewline
\ \ \ \ \isacommand{moreover}\isamarkupfalse%
\ \isacommand{have}\isamarkupfalse%
\ {\isachardoublequoteopen}{\isasymAnd}t{\isacharprime}{\kern0pt}{\isacharprime}{\kern0pt}\ A{\isacharprime}{\kern0pt}{\isacharprime}{\kern0pt}{\isachardot}{\kern0pt}\ {\isacharparenleft}{\kern0pt}t{\isacharprime}{\kern0pt}{\isacharprime}{\kern0pt}{\isacharcomma}{\kern0pt}A{\isacharprime}{\kern0pt}{\isacharprime}{\kern0pt}{\isacharparenright}{\kern0pt}\ {\isasymin}\ set\ hs{\isadigit{3}}{\isacharprime}{\kern0pt}\ {\isasymLongrightarrow}\ pairwise\ acts{\isacharunderscore}{\kern0pt}non{\isacharunderscore}{\kern0pt}intrf\ {\isacharparenleft}{\kern0pt}set\ A{\isacharprime}{\kern0pt}{\isacharprime}{\kern0pt}{\isacharparenright}{\kern0pt}{\isachardoublequoteclose}\isanewline
\ \ \ \ \ \ \isacommand{using}\isamarkupfalse%
\ fst{\isacharunderscore}{\kern0pt}h{\isadigit{3}}{\isacharprime}{\kern0pt}{\isacharunderscore}{\kern0pt}n{\isacharunderscore}{\kern0pt}htp\ {\isadigit{2}}{\isacharparenleft}{\kern0pt}{\isadigit{1}}{\isadigit{0}}{\isacharparenright}{\kern0pt}\isanewline
\ \ \ \ \ \ \isacommand{by}\isamarkupfalse%
\ {\isacharparenleft}{\kern0pt}force\ simp{\isacharcolon}{\kern0pt}\ intro{\isacharbang}{\kern0pt}{\isacharcolon}{\kern0pt}\ empty{\isacharunderscore}{\kern0pt}eff{\isacharunderscore}{\kern0pt}non{\isacharunderscore}{\kern0pt}intrf{\isacharparenright}{\kern0pt}\isanewline
\ \ \ \ \isacommand{ultimately}\isamarkupfalse%
\ \isacommand{show}\isamarkupfalse%
\ {\isacharquery}{\kern0pt}thesis\isanewline
\ \ \ \ \ \ \isacommand{using}\isamarkupfalse%
\ {\isacartoucheopen}valid{\isacharunderscore}{\kern0pt}happ{\isacharunderscore}{\kern0pt}seq\ M\ {\isacharbrackleft}{\kern0pt}h{\isacharprime}{\kern0pt}{\isacharbrackright}{\kern0pt}\ {\isacharparenleft}{\kern0pt}apply{\isacharunderscore}{\kern0pt}happ\ {\isacharparenleft}{\kern0pt}t{\isacharcomma}{\kern0pt}\ A{\isacharprime}{\kern0pt}{\isacharparenright}{\kern0pt}\ M{\isacharparenright}{\kern0pt}{\isacartoucheclose}\ fst{\isacharunderscore}{\kern0pt}h{\isadigit{3}}{\isacharprime}{\kern0pt}{\isacharunderscore}{\kern0pt}n{\isacharunderscore}{\kern0pt}htp\isanewline
\ \ \ \ \ \ \isacommand{by}\isamarkupfalse%
\ {\isacharparenleft}{\kern0pt}auto\ simp\ add{\isacharcolon}{\kern0pt}\ M{\isacharunderscore}{\kern0pt}IH{\isacharunderscore}{\kern0pt}def\ hs{\isacharprime}{\kern0pt}{\isacharunderscore}{\kern0pt}IH{\isacharunderscore}{\kern0pt}def\ intro{\isacharbang}{\kern0pt}{\isacharcolon}{\kern0pt}\ valid{\isacharunderscore}{\kern0pt}happ{\isacharunderscore}{\kern0pt}seq{\isacharunderscore}{\kern0pt}append{\isacharunderscore}{\kern0pt}{\isadigit{1}}\ \ {\isadigit{2}}{\isacharparenleft}{\kern0pt}{\isadigit{1}}{\isadigit{0}}{\isacharparenright}{\kern0pt}{\isacharparenright}{\kern0pt}\isanewline
\ \ \isacommand{qed}\isamarkupfalse%
\isanewline
\isacommand{qed}\isamarkupfalse%
\ {\isacharparenleft}{\kern0pt}fastforce\ simp{\isacharcolon}{\kern0pt}\ empty{\isacharunderscore}{\kern0pt}filter{\isacharunderscore}{\kern0pt}conv\ dest{\isacharbang}{\kern0pt}{\isacharcolon}{\kern0pt}\ filter{\isacharunderscore}{\kern0pt}eq{\isacharunderscore}{\kern0pt}consD{\isacharparenright}{\kern0pt}{\isacharplus}{\kern0pt}%
\endisatagproof
{\isafoldproof}%
%
\isadelimproof
\isanewline
%
\endisadelimproof
\isanewline
\isanewline
\isacommand{lemma}\isamarkupfalse%
\ not{\isacharunderscore}{\kern0pt}htp{\isacharunderscore}{\kern0pt}empty{\isacharunderscore}{\kern0pt}effect{\isacharcolon}{\kern0pt}\isanewline
\ \ \ \ \isakeyword{assumes}\ {\isachardoublequoteopen}wf{\isacharunderscore}{\kern0pt}plan\ {\isasympi}s{\isachardoublequoteclose}\isanewline
\ \ \ \ \ \ \isakeyword{and}\ {\isachardoublequoteopen}{\isasymnot}\ is{\isacharunderscore}{\kern0pt}htp\ {\isasympi}s\ t\isactrlsub i{\isachardoublequoteclose}\isanewline
\ \ \ \ \ \ \isakeyword{and}\ {\isachardoublequoteopen}ind{\isacharunderscore}{\kern0pt}happ{\isacharunderscore}{\kern0pt}seq\ {\isasympi}s\ hs{\isachardoublequoteclose}\isanewline
\ \ \ \ \ \ \isakeyword{and}\ {\isachardoublequoteopen}{\isacharparenleft}{\kern0pt}t\isactrlsub i{\isacharcomma}{\kern0pt}A\isactrlsub i{\isacharparenright}{\kern0pt}\ {\isasymin}\ set\ hs{\isachardoublequoteclose}\isanewline
\ \ \ \ \ \ \isakeyword{and}\ {\isachardoublequoteopen}a\ {\isasymin}\ set\ A\isactrlsub i{\isachardoublequoteclose}\isanewline
\ \ \ \ \isakeyword{shows}\ {\isachardoublequoteopen}{\isacharparenleft}{\kern0pt}effect\ a{\isacharparenright}{\kern0pt}\ {\isacharequal}{\kern0pt}\ Effect\ {\isacharbrackleft}{\kern0pt}{\isacharbrackright}{\kern0pt}\ {\isacharbrackleft}{\kern0pt}{\isacharbrackright}{\kern0pt}{\isachardoublequoteclose}\isanewline
%
\isadelimproof
\ \ %
\endisadelimproof
%
\isatagproof
\isacommand{proof}\isamarkupfalse%
\ {\isacharminus}{\kern0pt}\ \ \ \ \isanewline
\ \ \ \ \isacommand{have}\isamarkupfalse%
\ {\isacharasterisk}{\kern0pt}{\isacharcolon}{\kern0pt}\ {\isachardoublequoteopen}{\isacharparenleft}{\kern0pt}t\isactrlsub i{\isacharcomma}{\kern0pt}{\isasympi}{\isacharparenright}{\kern0pt}\ {\isasymnotin}\ set\ {\isasympi}s{\isachardoublequoteclose}\isanewline
\ \ \ \ \ \ \ \ \ \ \ \ {\isachardoublequoteopen}{\isasymAnd}t\isactrlsub {\isasympi}{\isachardot}{\kern0pt}\ {\isacharparenleft}{\kern0pt}t\isactrlsub {\isasympi}{\isacharcomma}{\kern0pt}{\isasympi}{\isacharparenright}{\kern0pt}\ {\isasymin}\ durative{\isacharunderscore}{\kern0pt}acts\ {\isasympi}s\ {\isasymLongrightarrow}\ t\isactrlsub i\ {\isasymnoteq}\ t\isactrlsub {\isasympi}\ {\isacharplus}{\kern0pt}\ {\isacharparenleft}{\kern0pt}duration\ {\isasympi}{\isacharparenright}{\kern0pt}{\isachardoublequoteclose}\isanewline
\ \ \ \ \ \ \ \ \ \isakeyword{for}\ {\isasympi}\isanewline
\ \ \ \ \ \ \isacommand{using}\isamarkupfalse%
\ {\isacartoucheopen}{\isasymnot}is{\isacharunderscore}{\kern0pt}htp\ {\isasympi}s\ t\isactrlsub i{\isacartoucheclose}\isanewline
\ \ \ \ \ \ \isacommand{by}\isamarkupfalse%
\ {\isacharparenleft}{\kern0pt}auto\ simp{\isacharcolon}{\kern0pt}\ is{\isacharunderscore}{\kern0pt}htp{\isacharunderscore}{\kern0pt}def{\isacharparenright}{\kern0pt}\isanewline
\ \ \ \ \isacommand{obtain}\isamarkupfalse%
\ t\isactrlsub {\isasympi}\ {\isasympi}\ \isakeyword{where}\ {\isachardoublequoteopen}{\isacharparenleft}{\kern0pt}t\isactrlsub {\isasympi}{\isacharcomma}{\kern0pt}\ {\isasympi}{\isacharparenright}{\kern0pt}{\isasymin}set\ {\isasympi}s{\isachardoublequoteclose}\ {\isachardoublequoteopen}inst{\isacharunderscore}{\kern0pt}of{\isacharunderscore}{\kern0pt}plan{\isacharunderscore}{\kern0pt}action\ {\isasympi}s\ {\isacharparenleft}{\kern0pt}t\isactrlsub {\isasympi}{\isacharcomma}{\kern0pt}\ {\isasympi}{\isacharparenright}{\kern0pt}\ {\isacharparenleft}{\kern0pt}t\isactrlsub i{\isacharcomma}{\kern0pt}\ a{\isacharparenright}{\kern0pt}{\isachardoublequoteclose}\isanewline
\ \ \ \ \ \ \isacommand{using}\isamarkupfalse%
\ ind{\isacharunderscore}{\kern0pt}happ{\isacharunderscore}{\kern0pt}seq{\isacharunderscore}{\kern0pt}props{\isacharparenleft}{\kern0pt}{\isadigit{7}}{\isacharparenright}{\kern0pt}{\isacharbrackleft}{\kern0pt}OF\ assms{\isacharparenleft}{\kern0pt}{\isadigit{3}}{\isacharminus}{\kern0pt}{\isadigit{5}}{\isacharparenright}{\kern0pt}{\isacharbrackright}{\kern0pt}\isanewline
\ \ \ \ \ \ \isacommand{by}\isamarkupfalse%
\ auto\isanewline
\ \ \ \ \isacommand{have}\isamarkupfalse%
\ {\isachardoublequoteopen}wf{\isacharunderscore}{\kern0pt}plan{\isacharunderscore}{\kern0pt}action\ {\isasympi}{\isachardoublequoteclose}\isanewline
\ \ \ \ \ \ \isacommand{using}\isamarkupfalse%
\ {\isacartoucheopen}wf{\isacharunderscore}{\kern0pt}plan\ {\isasympi}s{\isacartoucheclose}\ {\isacartoucheopen}{\isacharparenleft}{\kern0pt}t\isactrlsub {\isasympi}{\isacharcomma}{\kern0pt}\ {\isasympi}{\isacharparenright}{\kern0pt}\ {\isasymin}\ set\ {\isasympi}s{\isacartoucheclose}\isanewline
\ \ \ \ \ \ \isacommand{by}\isamarkupfalse%
\ {\isacharparenleft}{\kern0pt}auto\ simp{\isacharcolon}{\kern0pt}\ wf{\isacharunderscore}{\kern0pt}plan{\isacharunderscore}{\kern0pt}def{\isacharparenright}{\kern0pt}\isanewline
\ \ \ \ \isacommand{moreover}\isamarkupfalse%
\ \isacommand{then}\isamarkupfalse%
\ \isacommand{obtain}\isamarkupfalse%
\ a\isactrlsub s\isactrlsub c\isactrlsub h\isactrlsub e\isactrlsub m\isactrlsub a\ \isakeyword{where}\isanewline
\ \ \ \ \ \ {\isachardoublequoteopen}wf{\isacharunderscore}{\kern0pt}action{\isacharunderscore}{\kern0pt}schema\ a\isactrlsub s\isactrlsub c\isactrlsub h\isactrlsub e\isactrlsub m\isactrlsub a{\isachardoublequoteclose}\isanewline
\ \ \ \ \ \ {\isachardoublequoteopen}Some\ a\isactrlsub s\isactrlsub c\isactrlsub h\isactrlsub e\isactrlsub m\isactrlsub a\ {\isacharequal}{\kern0pt}\ resolve{\isacharunderscore}{\kern0pt}action{\isacharunderscore}{\kern0pt}schema\ {\isacharparenleft}{\kern0pt}plan{\isacharunderscore}{\kern0pt}action{\isachardot}{\kern0pt}name\ {\isasympi}{\isacharparenright}{\kern0pt}{\isachardoublequoteclose}\isanewline
\ \ \ \ \ \ \isacommand{by}\isamarkupfalse%
\ {\isacharparenleft}{\kern0pt}cases\ {\isasympi}{\isacharparenright}{\kern0pt}\ {\isacharparenleft}{\kern0pt}auto\ split{\isacharcolon}{\kern0pt}\ option{\isachardot}{\kern0pt}splits\ dest{\isacharcolon}{\kern0pt}\ resolve{\isacharunderscore}{\kern0pt}action{\isacharunderscore}{\kern0pt}wf{\isacharparenright}{\kern0pt}\isanewline
\ \ \ \ \isacommand{moreover}\isamarkupfalse%
\ \isacommand{have}\isamarkupfalse%
\ {\isachardoublequoteopen}Some\ a\ {\isacharequal}{\kern0pt}\ res{\isacharunderscore}{\kern0pt}inst{\isacharunderscore}{\kern0pt}snap{\isacharunderscore}{\kern0pt}action\ {\isasympi}\ Over{\isacharunderscore}{\kern0pt}All{\isachardoublequoteclose}\isanewline
\ \ \ \ \ \ \isacommand{using}\isamarkupfalse%
\ {\isacharasterisk}{\kern0pt}\ {\isacartoucheopen}{\isacharparenleft}{\kern0pt}t\isactrlsub {\isasympi}{\isacharcomma}{\kern0pt}\ {\isasympi}{\isacharparenright}{\kern0pt}{\isasymin}set\ {\isasympi}s{\isacartoucheclose}\ {\isacartoucheopen}inst{\isacharunderscore}{\kern0pt}of{\isacharunderscore}{\kern0pt}plan{\isacharunderscore}{\kern0pt}action\ {\isasympi}s\ {\isacharparenleft}{\kern0pt}t\isactrlsub {\isasympi}{\isacharcomma}{\kern0pt}\ {\isasympi}{\isacharparenright}{\kern0pt}\ {\isacharparenleft}{\kern0pt}t\isactrlsub i{\isacharcomma}{\kern0pt}\ a{\isacharparenright}{\kern0pt}{\isacartoucheclose}\isanewline
\ \ \ \ \ \ \isacommand{by}\isamarkupfalse%
\ {\isacharparenleft}{\kern0pt}cases\ {\isasympi}{\isacharparenright}{\kern0pt}\ {\isacharparenleft}{\kern0pt}fastforce\ simp{\isacharcolon}{\kern0pt}\ durative{\isacharunderscore}{\kern0pt}acts{\isacharunderscore}{\kern0pt}def\ is{\isacharunderscore}{\kern0pt}act{\isacharunderscore}{\kern0pt}simple{\isacharunderscore}{\kern0pt}def{\isacharparenright}{\kern0pt}{\isacharplus}{\kern0pt}\isanewline
\ \ \ \ \isacommand{ultimately}\isamarkupfalse%
\ \isacommand{show}\isamarkupfalse%
\ {\isacharquery}{\kern0pt}thesis\ \isanewline
\ \ \ \ \ \ \isacommand{by}\isamarkupfalse%
{\isacharparenleft}{\kern0pt}rule\ wf{\isacharunderscore}{\kern0pt}over{\isacharunderscore}{\kern0pt}all{\isacharunderscore}{\kern0pt}empty{\isacharunderscore}{\kern0pt}eff{\isacharparenright}{\kern0pt}\isanewline
\ \ \isacommand{qed}\isamarkupfalse%
%
\endisatagproof
{\isafoldproof}%
%
\isadelimproof
\isanewline
%
\endisadelimproof
\isanewline
\isacommand{lemma}\isamarkupfalse%
\ unique{\isacharunderscore}{\kern0pt}happ{\isacharcolon}{\kern0pt}\isanewline
\ \ \isakeyword{assumes}\ {\isachardoublequoteopen}ind{\isacharunderscore}{\kern0pt}happ{\isacharunderscore}{\kern0pt}seq\ {\isasympi}s\ hs{\isachardoublequoteclose}\isanewline
\ \ \isakeyword{shows}\ {\isachardoublequoteopen}{\isasymlbrakk}{\isacharparenleft}{\kern0pt}t{\isacharcomma}{\kern0pt}\ A{\isacharparenright}{\kern0pt}\ {\isasymin}\ set\ hs{\isacharsemicolon}{\kern0pt}\ {\isacharparenleft}{\kern0pt}t{\isacharcomma}{\kern0pt}\ A{\isacharprime}{\kern0pt}{\isacharparenright}{\kern0pt}\ {\isasymin}\ set\ hs{\isasymrbrakk}\ {\isasymLongrightarrow}\ A\ {\isacharequal}{\kern0pt}\ A{\isacharprime}{\kern0pt}{\isachardoublequoteclose}\isanewline
%
\isadelimproof
\ \ %
\endisadelimproof
%
\isatagproof
\isacommand{using}\isamarkupfalse%
\ ind{\isacharunderscore}{\kern0pt}happ{\isacharunderscore}{\kern0pt}seq{\isacharunderscore}{\kern0pt}props{\isacharparenleft}{\kern0pt}{\isadigit{1}}{\isacharparenright}{\kern0pt}{\isacharbrackleft}{\kern0pt}OF\ assms{\isacharbrackright}{\kern0pt}\isanewline
\ \ \isacommand{by}\isamarkupfalse%
\ {\isacharparenleft}{\kern0pt}induction\ hs{\isacharparenright}{\kern0pt}\ auto%
\endisatagproof
{\isafoldproof}%
%
\isadelimproof
\isanewline
%
\endisadelimproof
\isanewline
\isacommand{lemma}\isamarkupfalse%
\ inst{\isacharunderscore}{\kern0pt}of{\isacharunderscore}{\kern0pt}plan{\isacharunderscore}{\kern0pt}actionE{\isacharcolon}{\kern0pt}\isanewline
\ \ {\isachardoublequoteopen}{\isasymAnd}t\isactrlsub {\isasympi}\ {\isasympi}{\isachardot}{\kern0pt}\ {\isasymlbrakk}inst{\isacharunderscore}{\kern0pt}of{\isacharunderscore}{\kern0pt}plan{\isacharunderscore}{\kern0pt}action\ {\isasympi}s\ {\isacharparenleft}{\kern0pt}t\isactrlsub {\isasympi}{\isacharcomma}{\kern0pt}{\isasympi}{\isacharparenright}{\kern0pt}\ {\isacharparenleft}{\kern0pt}t\isactrlsub a{\isacharcomma}{\kern0pt}a{\isacharparenright}{\kern0pt}{\isacharsemicolon}{\kern0pt}\isanewline
\ \ \ \ \ \ \ \ \ \ \ {\isasymlbrakk}res{\isacharunderscore}{\kern0pt}inst\ {\isasympi}\ {\isacharequal}{\kern0pt}\ Some\ a{\isacharsemicolon}{\kern0pt}\ t\isactrlsub {\isasympi}\ {\isacharequal}{\kern0pt}\ t\isactrlsub a{\isasymrbrakk}\ {\isasymLongrightarrow}\ Q{\isacharsemicolon}{\kern0pt}\isanewline
\ \ \ \ \ \ \ \ \ \ \ {\isasymlbrakk}res{\isacharunderscore}{\kern0pt}inst{\isacharunderscore}{\kern0pt}snap{\isacharunderscore}{\kern0pt}action\ {\isasympi}\ At{\isacharunderscore}{\kern0pt}Start\ {\isacharequal}{\kern0pt}\ Some\ a{\isacharsemicolon}{\kern0pt}\ t\isactrlsub {\isasympi}\ {\isacharequal}{\kern0pt}\ t\isactrlsub a{\isasymrbrakk}\ {\isasymLongrightarrow}\ Q{\isacharsemicolon}{\kern0pt}\isanewline
\ \ \ \ \ \ \ \ \ \ \ {\isasymlbrakk}res{\isacharunderscore}{\kern0pt}inst{\isacharunderscore}{\kern0pt}snap{\isacharunderscore}{\kern0pt}action\ {\isasympi}\ At{\isacharunderscore}{\kern0pt}End\ {\isacharequal}{\kern0pt}\ Some\ a{\isacharsemicolon}{\kern0pt}\ t\isactrlsub {\isasympi}\ {\isacharplus}{\kern0pt}\ {\isacharparenleft}{\kern0pt}duration\ {\isasympi}{\isacharparenright}{\kern0pt}\ {\isacharequal}{\kern0pt}\ t\isactrlsub a{\isasymrbrakk}\ {\isasymLongrightarrow}\ Q{\isacharsemicolon}{\kern0pt}\isanewline
\ \ \ \ \ \ \ \ \ \ \ {\isasymAnd}t\isactrlsub i\ t\isactrlsub j{\isachardot}{\kern0pt}\ {\isasymlbrakk}res{\isacharunderscore}{\kern0pt}inst{\isacharunderscore}{\kern0pt}snap{\isacharunderscore}{\kern0pt}action\ {\isasympi}\ Over{\isacharunderscore}{\kern0pt}All\ {\isacharequal}{\kern0pt}\ Some\ a{\isacharsemicolon}{\kern0pt}\ consec{\isacharunderscore}{\kern0pt}htps\ {\isasympi}s\ t\isactrlsub i\ t\isactrlsub j{\isacharsemicolon}{\kern0pt}\ t\isactrlsub {\isasympi}\ {\isasymle}\ t\isactrlsub i{\isacharsemicolon}{\kern0pt}\ t\isactrlsub j\ {\isasymle}\ t\isactrlsub {\isasympi}\ {\isacharplus}{\kern0pt}\ {\isacharparenleft}{\kern0pt}duration\ {\isasympi}{\isacharparenright}{\kern0pt}{\isacharsemicolon}{\kern0pt}\ t\isactrlsub i\ {\isacharless}{\kern0pt}\ t\isactrlsub a{\isacharsemicolon}{\kern0pt}\ t\isactrlsub a\ {\isacharless}{\kern0pt}\ t\isactrlsub j{\isasymrbrakk}\ {\isasymLongrightarrow}\ Q{\isasymrbrakk}\isanewline
\ \ \ \ {\isasymLongrightarrow}\ Q{\isachardoublequoteclose}\isanewline
%
\isadelimproof
\ \ %
\endisadelimproof
%
\isatagproof
\isacommand{by}\isamarkupfalse%
\ auto%
\endisatagproof
{\isafoldproof}%
%
\isadelimproof
\isanewline
%
\endisadelimproof
\isanewline
\ \ \isacommand{lemma}\isamarkupfalse%
\ ind{\isacharunderscore}{\kern0pt}happ{\isacharunderscore}{\kern0pt}seq{\isacharunderscore}{\kern0pt}happ{\isacharunderscore}{\kern0pt}at{\isacharunderscore}{\kern0pt}htps{\isacharcolon}{\kern0pt}\isanewline
\ \ \ \ \isakeyword{assumes}\ {\isachardoublequoteopen}is{\isacharunderscore}{\kern0pt}htp\ {\isasympi}s\ t{\isachardoublequoteclose}\isanewline
\ \ \ \ \ \ \isakeyword{and}\ {\isachardoublequoteopen}ind{\isacharunderscore}{\kern0pt}happ{\isacharunderscore}{\kern0pt}seq\ {\isasympi}s\ hs{\isachardoublequoteclose}\isanewline
\ \ \ \ \ \ \isakeyword{and}\ {\isachardoublequoteopen}ind{\isacharunderscore}{\kern0pt}happ{\isacharunderscore}{\kern0pt}seq\ {\isasympi}s\ hs{\isacharprime}{\kern0pt}{\isachardoublequoteclose}\isanewline
\ \ \ \ \ \ \isakeyword{and}\ {\isachardoublequoteopen}{\isacharparenleft}{\kern0pt}t{\isacharcomma}{\kern0pt}A{\isacharparenright}{\kern0pt}\ {\isasymin}\ set\ hs{\isachardoublequoteclose}\isanewline
\ \ \ \ \ \ \isakeyword{and}\ {\isachardoublequoteopen}{\isacharparenleft}{\kern0pt}t{\isacharcomma}{\kern0pt}A{\isacharprime}{\kern0pt}{\isacharparenright}{\kern0pt}\ {\isasymin}\ set\ hs{\isacharprime}{\kern0pt}{\isachardoublequoteclose}\isanewline
\ \ \ \ \ \ \isakeyword{and}\ {\isachardoublequoteopen}a\ {\isasymin}\ set\ A{\isachardoublequoteclose}\isanewline
\ \ \ \ \isakeyword{shows}\ {\isachardoublequoteopen}a\ {\isasymin}\ set\ A{\isacharprime}{\kern0pt}{\isachardoublequoteclose}\isanewline
%
\isadelimproof
\ \ %
\endisadelimproof
%
\isatagproof
\isacommand{proof}\isamarkupfalse%
{\isacharminus}{\kern0pt}\isanewline
\ \ \ \ \isacommand{obtain}\isamarkupfalse%
\ t\isactrlsub {\isasympi}\ {\isasympi}\ \isakeyword{where}\ {\isachardoublequoteopen}{\isacharparenleft}{\kern0pt}t\isactrlsub {\isasympi}{\isacharcomma}{\kern0pt}\ {\isasympi}{\isacharparenright}{\kern0pt}{\isasymin}set\ {\isasympi}s{\isachardoublequoteclose}\isanewline
\ \ \ \ \ \ {\isachardoublequoteopen}{\isacharparenleft}{\kern0pt}res{\isacharunderscore}{\kern0pt}inst\ {\isasympi}\ {\isacharequal}{\kern0pt}\ Some\ a\ {\isasymand}\ t\isactrlsub {\isasympi}\ {\isacharequal}{\kern0pt}\ t{\isacharparenright}{\kern0pt}\isanewline
\ \ \ \ \ \ \ \ \ \ \ \ \ \ \ \ \ \ \ \ \ \ \ {\isasymor}\ {\isacharparenleft}{\kern0pt}res{\isacharunderscore}{\kern0pt}inst{\isacharunderscore}{\kern0pt}snap{\isacharunderscore}{\kern0pt}action\ {\isasympi}\ At{\isacharunderscore}{\kern0pt}Start\ {\isacharequal}{\kern0pt}\ Some\ a\ {\isasymand}\ t\isactrlsub {\isasympi}\ {\isacharequal}{\kern0pt}\ t{\isacharparenright}{\kern0pt}\isanewline
\ \ \ \ \ \ \ \ \ \ \ \ \ \ \ \ \ \ \ \ \ \ \ {\isasymor}\ {\isacharparenleft}{\kern0pt}res{\isacharunderscore}{\kern0pt}inst{\isacharunderscore}{\kern0pt}snap{\isacharunderscore}{\kern0pt}action\ {\isasympi}\ At{\isacharunderscore}{\kern0pt}End\ {\isacharequal}{\kern0pt}\ Some\ a\ {\isasymand}\ t\isactrlsub {\isasympi}\ {\isacharplus}{\kern0pt}\ {\isacharparenleft}{\kern0pt}duration\ {\isasympi}{\isacharparenright}{\kern0pt}\ {\isacharequal}{\kern0pt}\ t{\isacharparenright}{\kern0pt}{\isachardoublequoteclose}\isanewline
\ \ \ \ \ \ \isacommand{using}\isamarkupfalse%
\ ind{\isacharunderscore}{\kern0pt}happ{\isacharunderscore}{\kern0pt}seq{\isacharunderscore}{\kern0pt}props{\isacharparenleft}{\kern0pt}{\isadigit{7}}{\isacharparenright}{\kern0pt}{\isacharbrackleft}{\kern0pt}OF\ {\isacartoucheopen}ind{\isacharunderscore}{\kern0pt}happ{\isacharunderscore}{\kern0pt}seq\ {\isasympi}s\ hs{\isacartoucheclose}\ {\isacartoucheopen}{\isacharparenleft}{\kern0pt}t{\isacharcomma}{\kern0pt}A{\isacharparenright}{\kern0pt}\ {\isasymin}\ set\ hs{\isacartoucheclose}\ {\isacartoucheopen}a\ {\isasymin}\ set\ A{\isacartoucheclose}{\isacharbrackright}{\kern0pt}\isanewline
\ \ \ \ \ \ \ \ {\isacartoucheopen}is{\isacharunderscore}{\kern0pt}htp\ {\isasympi}s\ t{\isacartoucheclose}\isanewline
\ \ \ \ \ \ \isacommand{by}\isamarkupfalse%
\ {\isacharparenleft}{\kern0pt}auto\ simp{\isacharcolon}{\kern0pt}\ ast{\isacharunderscore}{\kern0pt}problem{\isachardot}{\kern0pt}consec{\isacharunderscore}{\kern0pt}htps{\isacharunderscore}{\kern0pt}def{\isacharparenright}{\kern0pt}\isanewline
\ \ \ \ \isacommand{then}\isamarkupfalse%
\ \isacommand{obtain}\isamarkupfalse%
\ A{\isacharprime}{\kern0pt}{\isacharprime}{\kern0pt}\ \isakeyword{where}\ {\isachardoublequoteopen}{\isacharparenleft}{\kern0pt}t{\isacharcomma}{\kern0pt}A{\isacharprime}{\kern0pt}{\isacharprime}{\kern0pt}{\isacharparenright}{\kern0pt}\ {\isasymin}\ set\ hs{\isacharprime}{\kern0pt}{\isachardoublequoteclose}\ {\isachardoublequoteopen}a\ {\isasymin}\ set\ A{\isacharprime}{\kern0pt}{\isacharprime}{\kern0pt}{\isachardoublequoteclose}\isanewline
\ \ \ \ \ \ \isacommand{using}\isamarkupfalse%
\ ind{\isacharunderscore}{\kern0pt}happ{\isacharunderscore}{\kern0pt}seq{\isacharunderscore}{\kern0pt}props{\isacharparenleft}{\kern0pt}{\isadigit{2}}{\isacharminus}{\kern0pt}{\isadigit{4}}{\isacharparenright}{\kern0pt}{\isacharbrackleft}{\kern0pt}OF\ {\isacartoucheopen}ind{\isacharunderscore}{\kern0pt}happ{\isacharunderscore}{\kern0pt}seq\ {\isasympi}s\ hs{\isacharprime}{\kern0pt}{\isacartoucheclose}{\isacharbrackright}{\kern0pt}\isanewline
\ \ \ \ \ \ \isacommand{by}\isamarkupfalse%
\ {\isacharparenleft}{\kern0pt}force\ simp{\isacharcolon}{\kern0pt}\ durative{\isacharunderscore}{\kern0pt}acts{\isacharunderscore}{\kern0pt}def\ simple{\isacharunderscore}{\kern0pt}acts{\isacharunderscore}{\kern0pt}def\ is{\isacharunderscore}{\kern0pt}act{\isacharunderscore}{\kern0pt}simple{\isacharunderscore}{\kern0pt}def\ split{\isacharcolon}{\kern0pt}\ plan{\isacharunderscore}{\kern0pt}action{\isachardot}{\kern0pt}splits{\isacharparenright}{\kern0pt}\isanewline
\ \ \ \ \isacommand{moreover}\isamarkupfalse%
\ \isacommand{hence}\isamarkupfalse%
\ {\isachardoublequoteopen}A{\isacharprime}{\kern0pt}{\isacharprime}{\kern0pt}\ {\isacharequal}{\kern0pt}\ A{\isacharprime}{\kern0pt}{\isachardoublequoteclose}\isanewline
\ \ \ \ \ \ \isacommand{using}\isamarkupfalse%
\ \ {\isacartoucheopen}ind{\isacharunderscore}{\kern0pt}happ{\isacharunderscore}{\kern0pt}seq\ {\isasympi}s\ hs{\isacharprime}{\kern0pt}{\isacartoucheclose}\ {\isacartoucheopen}{\isacharparenleft}{\kern0pt}t{\isacharcomma}{\kern0pt}\ A{\isacharprime}{\kern0pt}{\isacharparenright}{\kern0pt}\ {\isasymin}\ set\ hs{\isacharprime}{\kern0pt}{\isacartoucheclose}\isanewline
\ \ \ \ \ \ \isacommand{by}\isamarkupfalse%
{\isacharparenleft}{\kern0pt}auto\ intro{\isacharcolon}{\kern0pt}\ unique{\isacharunderscore}{\kern0pt}happ{\isacharparenright}{\kern0pt}\isanewline
\ \ \ \ \isacommand{ultimately}\isamarkupfalse%
\ \isacommand{show}\isamarkupfalse%
\ {\isacharquery}{\kern0pt}thesis\isanewline
\ \ \ \ \ \ \isacommand{by}\isamarkupfalse%
\ auto\isanewline
\ \ \isacommand{qed}\isamarkupfalse%
%
\endisatagproof
{\isafoldproof}%
%
\isadelimproof
\isanewline
%
\endisadelimproof
\isanewline
\ \ \isacommand{lemma}\isamarkupfalse%
\ ind{\isacharunderscore}{\kern0pt}happ{\isacharunderscore}{\kern0pt}seq{\isacharunderscore}{\kern0pt}happ{\isacharunderscore}{\kern0pt}at{\isacharunderscore}{\kern0pt}htps{\isacharprime}{\kern0pt}{\isacharcolon}{\kern0pt}\isanewline
\ \ \ \ \isakeyword{assumes}\ {\isachardoublequoteopen}is{\isacharunderscore}{\kern0pt}htp\ {\isasympi}s\ t{\isachardoublequoteclose}\isanewline
\ \ \ \ \ \ \isakeyword{and}\ {\isachardoublequoteopen}ind{\isacharunderscore}{\kern0pt}happ{\isacharunderscore}{\kern0pt}seq\ {\isasympi}s\ hs{\isachardoublequoteclose}\isanewline
\ \ \ \ \ \ \isakeyword{and}\ {\isachardoublequoteopen}ind{\isacharunderscore}{\kern0pt}happ{\isacharunderscore}{\kern0pt}seq\ {\isasympi}s\ hs{\isacharprime}{\kern0pt}{\isachardoublequoteclose}\isanewline
\ \ \ \ \ \ \isakeyword{and}\ {\isachardoublequoteopen}{\isacharparenleft}{\kern0pt}t{\isacharcomma}{\kern0pt}A{\isacharparenright}{\kern0pt}\ {\isasymin}\ set\ hs{\isachardoublequoteclose}\isanewline
\ \ \ \ \ \ \isakeyword{and}\ {\isachardoublequoteopen}{\isacharparenleft}{\kern0pt}t{\isacharcomma}{\kern0pt}A{\isacharprime}{\kern0pt}{\isacharparenright}{\kern0pt}\ {\isasymin}\ set\ hs{\isacharprime}{\kern0pt}{\isachardoublequoteclose}\isanewline
\ \ \ \ \isakeyword{shows}\ {\isachardoublequoteopen}set\ A\ {\isacharequal}{\kern0pt}\ set\ A{\isacharprime}{\kern0pt}{\isachardoublequoteclose}\isanewline
%
\isadelimproof
\ \ \ \ %
\endisadelimproof
%
\isatagproof
\isacommand{using}\isamarkupfalse%
\ assms\ ind{\isacharunderscore}{\kern0pt}happ{\isacharunderscore}{\kern0pt}seq{\isacharunderscore}{\kern0pt}happ{\isacharunderscore}{\kern0pt}at{\isacharunderscore}{\kern0pt}htps\isanewline
\ \ \ \ \isacommand{by}\isamarkupfalse%
\ blast%
\endisatagproof
{\isafoldproof}%
%
\isadelimproof
\isanewline
%
\endisadelimproof
\isanewline
\isacommand{lemma}\isamarkupfalse%
\ htp{\isacharunderscore}{\kern0pt}in{\isacharunderscore}{\kern0pt}ind{\isacharunderscore}{\kern0pt}happ{\isacharunderscore}{\kern0pt}seq{\isacharcolon}{\kern0pt}\isanewline
\ \ {\isachardoublequoteopen}{\isasymlbrakk}is{\isacharunderscore}{\kern0pt}htp\ {\isasympi}s\ t{\isacharsemicolon}{\kern0pt}\ ind{\isacharunderscore}{\kern0pt}happ{\isacharunderscore}{\kern0pt}seq\ {\isasympi}s\ hs{\isasymrbrakk}\ {\isasymLongrightarrow}\ {\isasymexists}A{\isachardot}{\kern0pt}\ {\isacharparenleft}{\kern0pt}t{\isacharcomma}{\kern0pt}A{\isacharparenright}{\kern0pt}\ {\isasymin}\ set\ hs{\isachardoublequoteclose}\isanewline
%
\isadelimproof
\ \ %
\endisadelimproof
%
\isatagproof
\isacommand{by}\isamarkupfalse%
{\isacharparenleft}{\kern0pt}fastforce\ simp{\isacharcolon}{\kern0pt}\ is{\isacharunderscore}{\kern0pt}htp{\isacharunderscore}{\kern0pt}def\ ind{\isacharunderscore}{\kern0pt}happ{\isacharunderscore}{\kern0pt}seq{\isacharunderscore}{\kern0pt}def\ durative{\isacharunderscore}{\kern0pt}acts{\isacharunderscore}{\kern0pt}def\ simple{\isacharunderscore}{\kern0pt}acts{\isacharunderscore}{\kern0pt}def\ split{\isacharcolon}{\kern0pt}\ prod{\isachardot}{\kern0pt}splits{\isacharparenright}{\kern0pt}%
\endisatagproof
{\isafoldproof}%
%
\isadelimproof
\isanewline
%
\endisadelimproof
\isanewline
\isacommand{lemma}\isamarkupfalse%
\ sorted{\isacharunderscore}{\kern0pt}takeWhile{\isacharcolon}{\kern0pt}\ {\isachardoublequoteopen}strict{\isacharunderscore}{\kern0pt}sorted\ xs\ {\isasymLongrightarrow}\ x\ {\isasymin}\ set\ {\isacharparenleft}{\kern0pt}takeWhile\ {\isacharparenleft}{\kern0pt}{\isacharparenleft}{\kern0pt}{\isasymlambda}x{\isachardot}{\kern0pt}\ x\ {\isacharless}{\kern0pt}\ c{\isacharparenright}{\kern0pt}{\isacharparenright}{\kern0pt}\ xs{\isacharparenright}{\kern0pt}\ {\isacharequal}{\kern0pt}\ {\isacharparenleft}{\kern0pt}x\ {\isasymin}\ set\ xs\ {\isasymand}\ {\isacharparenleft}{\kern0pt}x\ {\isacharless}{\kern0pt}\ c{\isacharparenright}{\kern0pt}{\isacharparenright}{\kern0pt}{\isachardoublequoteclose}\isanewline
%
\isadelimproof
\ \ %
\endisadelimproof
%
\isatagproof
\isacommand{by}\isamarkupfalse%
\ {\isacharparenleft}{\kern0pt}induction\ xs\ arbitrary{\isacharcolon}{\kern0pt}\ c\ x{\isacharparenright}{\kern0pt}\ auto%
\endisatagproof
{\isafoldproof}%
%
\isadelimproof
\isanewline
%
\endisadelimproof
\ \ \isanewline
\isacommand{lemma}\isamarkupfalse%
\ sorted{\isacharunderscore}{\kern0pt}dropWhile{\isacharcolon}{\kern0pt}\ {\isachardoublequoteopen}strict{\isacharunderscore}{\kern0pt}sorted\ xs\ {\isasymLongrightarrow}\ x\ {\isasymin}\ set\ {\isacharparenleft}{\kern0pt}dropWhile\ {\isacharparenleft}{\kern0pt}{\isacharparenleft}{\kern0pt}{\isasymlambda}x{\isachardot}{\kern0pt}\ x\ {\isasymle}\ c{\isacharparenright}{\kern0pt}{\isacharparenright}{\kern0pt}\ xs{\isacharparenright}{\kern0pt}\ {\isacharequal}{\kern0pt}\ {\isacharparenleft}{\kern0pt}x\ {\isasymin}\ set\ xs\ {\isasymand}\ {\isacharparenleft}{\kern0pt}c\ {\isacharless}{\kern0pt}\ x{\isacharparenright}{\kern0pt}{\isacharparenright}{\kern0pt}{\isachardoublequoteclose}\isanewline
%
\isadelimproof
\ \ %
\endisadelimproof
%
\isatagproof
\isacommand{by}\isamarkupfalse%
\ {\isacharparenleft}{\kern0pt}induction\ xs\ arbitrary{\isacharcolon}{\kern0pt}\ c\ x{\isacharparenright}{\kern0pt}\ auto%
\endisatagproof
{\isafoldproof}%
%
\isadelimproof
\isanewline
%
\endisadelimproof
\isanewline
\isacommand{lemma}\isamarkupfalse%
\ map{\isacharunderscore}{\kern0pt}fst{\isacharunderscore}{\kern0pt}dropWhile{\isacharcolon}{\kern0pt}\ {\isachardoublequoteopen}map\ fst\ {\isacharparenleft}{\kern0pt}dropWhile\ {\isacharparenleft}{\kern0pt}leq\ t{\isacharparenright}{\kern0pt}\ hs{\isacharparenright}{\kern0pt}\ {\isacharequal}{\kern0pt}\ {\isacharparenleft}{\kern0pt}dropWhile\ {\isacharparenleft}{\kern0pt}{\isacharparenleft}{\kern0pt}{\isasymlambda}x{\isachardot}{\kern0pt}\ x\ {\isasymle}\ t{\isacharparenright}{\kern0pt}{\isacharparenright}{\kern0pt}\ {\isacharparenleft}{\kern0pt}map\ fst\ hs{\isacharparenright}{\kern0pt}{\isacharparenright}{\kern0pt}{\isachardoublequoteclose}\isanewline
%
\isadelimproof
\ \ %
\endisadelimproof
%
\isatagproof
\isacommand{by}\isamarkupfalse%
\ {\isacharparenleft}{\kern0pt}induction\ hs{\isacharparenright}{\kern0pt}\ {\isacharparenleft}{\kern0pt}auto\ simp{\isacharcolon}{\kern0pt}\ leq{\isacharunderscore}{\kern0pt}def{\isacharparenright}{\kern0pt}%
\endisatagproof
{\isafoldproof}%
%
\isadelimproof
\isanewline
%
\endisadelimproof
\isanewline
\isacommand{lemma}\isamarkupfalse%
\ map{\isacharunderscore}{\kern0pt}fst{\isacharunderscore}{\kern0pt}takeWhile{\isacharcolon}{\kern0pt}\ {\isachardoublequoteopen}map\ fst\ {\isacharparenleft}{\kern0pt}takeWhile\ {\isacharparenleft}{\kern0pt}le\ t{\isacharparenright}{\kern0pt}\ hs{\isacharparenright}{\kern0pt}\ {\isacharequal}{\kern0pt}\ {\isacharparenleft}{\kern0pt}takeWhile\ {\isacharparenleft}{\kern0pt}{\isacharparenleft}{\kern0pt}{\isasymlambda}x{\isachardot}{\kern0pt}\ x\ {\isacharless}{\kern0pt}\ t{\isacharparenright}{\kern0pt}{\isacharparenright}{\kern0pt}\ {\isacharparenleft}{\kern0pt}map\ fst\ hs{\isacharparenright}{\kern0pt}{\isacharparenright}{\kern0pt}{\isachardoublequoteclose}\isanewline
%
\isadelimproof
\ \ %
\endisadelimproof
%
\isatagproof
\isacommand{by}\isamarkupfalse%
\ {\isacharparenleft}{\kern0pt}induction\ hs{\isacharparenright}{\kern0pt}\ {\isacharparenleft}{\kern0pt}auto\ simp{\isacharcolon}{\kern0pt}\ le{\isacharunderscore}{\kern0pt}def{\isacharparenright}{\kern0pt}%
\endisatagproof
{\isafoldproof}%
%
\isadelimproof
\isanewline
%
\endisadelimproof
\isanewline
\isacommand{lemma}\isamarkupfalse%
\ strict{\isacharunderscore}{\kern0pt}sorted{\isacharunderscore}{\kern0pt}fst{\isacharunderscore}{\kern0pt}unique{\isacharunderscore}{\kern0pt}snd{\isacharcolon}{\kern0pt}\isanewline
\ \ {\isachardoublequoteopen}{\isasymlbrakk}strict{\isacharunderscore}{\kern0pt}sorted\ {\isacharparenleft}{\kern0pt}map\ fst\ hs{\isacharparenright}{\kern0pt}{\isacharsemicolon}{\kern0pt}\ {\isacharparenleft}{\kern0pt}t{\isacharcomma}{\kern0pt}\ A{\isacharparenright}{\kern0pt}\ {\isasymin}\ set\ hs{\isacharsemicolon}{\kern0pt}\ {\isacharparenleft}{\kern0pt}t{\isacharcomma}{\kern0pt}\ A{\isacharprime}{\kern0pt}{\isacharparenright}{\kern0pt}\ {\isasymin}\ set\ hs{\isasymrbrakk}\ {\isasymLongrightarrow}\ A\ {\isacharequal}{\kern0pt}\ A{\isacharprime}{\kern0pt}{\isachardoublequoteclose}\isanewline
%
\isadelimproof
\ \ %
\endisadelimproof
%
\isatagproof
\isacommand{by}\isamarkupfalse%
\ {\isacharparenleft}{\kern0pt}induction\ hs{\isacharparenright}{\kern0pt}\ auto%
\endisatagproof
{\isafoldproof}%
%
\isadelimproof
\isanewline
%
\endisadelimproof
\isanewline
\isacommand{lemma}\isamarkupfalse%
\ strict{\isacharunderscore}{\kern0pt}sorted{\isacharunderscore}{\kern0pt}dropWhile{\isacharcolon}{\kern0pt}\isanewline
\ \ {\isachardoublequoteopen}strict{\isacharunderscore}{\kern0pt}sorted\ xs\ {\isasymLongrightarrow}\ strict{\isacharunderscore}{\kern0pt}sorted\ {\isacharparenleft}{\kern0pt}dropWhile\ Q\ xs{\isacharparenright}{\kern0pt}{\isachardoublequoteclose}\isanewline
%
\isadelimproof
\ \ %
\endisadelimproof
%
\isatagproof
\isacommand{by}\isamarkupfalse%
\ {\isacharparenleft}{\kern0pt}induction\ xs{\isacharparenright}{\kern0pt}\ auto%
\endisatagproof
{\isafoldproof}%
%
\isadelimproof
\isanewline
%
\endisadelimproof
\isanewline
\isacommand{lemma}\isamarkupfalse%
\ strict{\isacharunderscore}{\kern0pt}sorted{\isacharunderscore}{\kern0pt}takeWhile{\isacharcolon}{\kern0pt}\isanewline
\ \ {\isachardoublequoteopen}strict{\isacharunderscore}{\kern0pt}sorted\ xs\ {\isasymLongrightarrow}\ strict{\isacharunderscore}{\kern0pt}sorted\ {\isacharparenleft}{\kern0pt}takeWhile\ Q\ xs{\isacharparenright}{\kern0pt}{\isachardoublequoteclose}\isanewline
%
\isadelimproof
\ \ %
\endisadelimproof
%
\isatagproof
\isacommand{by}\isamarkupfalse%
\ {\isacharparenleft}{\kern0pt}induction\ xs{\isacharparenright}{\kern0pt}\ {\isacharparenleft}{\kern0pt}auto\ dest{\isacharbang}{\kern0pt}{\isacharcolon}{\kern0pt}\ set{\isacharunderscore}{\kern0pt}takeWhileD{\isacharparenright}{\kern0pt}%
\endisatagproof
{\isafoldproof}%
%
\isadelimproof
\isanewline
%
\endisadelimproof
\isanewline
\isacommand{lemma}\isamarkupfalse%
\ strict{\isacharunderscore}{\kern0pt}sorted{\isacharunderscore}{\kern0pt}drop{\isacharunderscore}{\kern0pt}leq{\isacharunderscore}{\kern0pt}take{\isacharunderscore}{\kern0pt}le{\isacharprime}{\kern0pt}{\isacharcolon}{\kern0pt}\isanewline
\ \ \ \ \isakeyword{assumes}\ {\isachardoublequoteopen}strict{\isacharunderscore}{\kern0pt}sorted\ {\isacharparenleft}{\kern0pt}map\ fst\ hs{\isacharparenright}{\kern0pt}{\isachardoublequoteclose}\isanewline
\ \ \ \ \isakeyword{shows}\ {\isachardoublequoteopen}{\isacharparenleft}{\kern0pt}t\isactrlsub a{\isacharcomma}{\kern0pt}A{\isacharparenright}{\kern0pt}\ {\isasymin}\ set\ {\isacharparenleft}{\kern0pt}dropWhile\ {\isacharparenleft}{\kern0pt}leq\ t\isactrlsub i{\isacharparenright}{\kern0pt}\ {\isacharparenleft}{\kern0pt}takeWhile\ {\isacharparenleft}{\kern0pt}le\ t\isactrlsub j{\isacharparenright}{\kern0pt}\ hs{\isacharparenright}{\kern0pt}{\isacharparenright}{\kern0pt}\ {\isasymlongleftrightarrow}\ \isanewline
\ \ \ \ \ \ \ \ \ \ \ \ \ \ {\isacharparenleft}{\kern0pt}t\isactrlsub i\ {\isacharless}{\kern0pt}\ t\isactrlsub a\ {\isasymand}\ t\isactrlsub a\ {\isacharless}{\kern0pt}\ t\isactrlsub j\ {\isasymand}\ {\isacharparenleft}{\kern0pt}t\isactrlsub a{\isacharcomma}{\kern0pt}A{\isacharparenright}{\kern0pt}\ {\isasymin}\ set\ hs{\isacharparenright}{\kern0pt}{\isachardoublequoteclose}\ {\isacharparenleft}{\kern0pt}\isakeyword{is}\ {\isachardoublequoteopen}{\isacharquery}{\kern0pt}l\ {\isasymlongleftrightarrow}\ {\isacharquery}{\kern0pt}r{\isachardoublequoteclose}{\isacharparenright}{\kern0pt}\isanewline
%
\isadelimproof
%
\endisadelimproof
%
\isatagproof
\isacommand{proof}\isamarkupfalse%
\isanewline
\ \ \isacommand{assume}\isamarkupfalse%
\ {\isacharquery}{\kern0pt}l\isanewline
\ \ \isacommand{hence}\isamarkupfalse%
\ {\isachardoublequoteopen}t\isactrlsub a\ {\isasymin}\ set\ {\isacharparenleft}{\kern0pt}map\ fst\ {\isacharparenleft}{\kern0pt}dropWhile\ {\isacharparenleft}{\kern0pt}leq\ t\isactrlsub i{\isacharparenright}{\kern0pt}\ {\isacharparenleft}{\kern0pt}takeWhile\ {\isacharparenleft}{\kern0pt}le\ t\isactrlsub j{\isacharparenright}{\kern0pt}\ hs{\isacharparenright}{\kern0pt}{\isacharparenright}{\kern0pt}{\isacharparenright}{\kern0pt}{\isachardoublequoteclose}\isanewline
\ \ \ \ \isacommand{by}\isamarkupfalse%
\ auto\isanewline
\ \ \isacommand{hence}\isamarkupfalse%
\ {\isachardoublequoteopen}{\isacharparenleft}{\kern0pt}t\isactrlsub i\ {\isacharless}{\kern0pt}\ t\isactrlsub a{\isacharparenright}{\kern0pt}{\isachardoublequoteclose}\ {\isachardoublequoteopen}t\isactrlsub a\ {\isacharless}{\kern0pt}\ t\isactrlsub j{\isachardoublequoteclose}\isanewline
\ \ \ \ \isacommand{using}\isamarkupfalse%
\ strict{\isacharunderscore}{\kern0pt}sorted{\isacharunderscore}{\kern0pt}takeWhile{\isacharbrackleft}{\kern0pt}OF\ {\isacartoucheopen}strict{\isacharunderscore}{\kern0pt}sorted\ {\isacharparenleft}{\kern0pt}map\ fst\ hs{\isacharparenright}{\kern0pt}{\isacartoucheclose}{\isacharbrackright}{\kern0pt}\ {\isacartoucheopen}strict{\isacharunderscore}{\kern0pt}sorted\ {\isacharparenleft}{\kern0pt}map\ fst\ hs{\isacharparenright}{\kern0pt}{\isacartoucheclose}\isanewline
\ \ \ \ \isacommand{by}\isamarkupfalse%
{\isacharparenleft}{\kern0pt}auto\ simp{\isacharcolon}{\kern0pt}\ sorted{\isacharunderscore}{\kern0pt}takeWhile\ sorted{\isacharunderscore}{\kern0pt}dropWhile\ map{\isacharunderscore}{\kern0pt}fst{\isacharunderscore}{\kern0pt}takeWhile\ map{\isacharunderscore}{\kern0pt}fst{\isacharunderscore}{\kern0pt}dropWhile{\isacharparenright}{\kern0pt}{\isacharplus}{\kern0pt}\isanewline
\ \ \isacommand{thus}\isamarkupfalse%
\ {\isacharquery}{\kern0pt}r\isanewline
\ \ \ \ \isacommand{using}\isamarkupfalse%
\ {\isacartoucheopen}{\isacharquery}{\kern0pt}l{\isacartoucheclose}\isanewline
\ \ \ \ \isacommand{by}\isamarkupfalse%
\ {\isacharparenleft}{\kern0pt}auto\ dest{\isacharcolon}{\kern0pt}\ set{\isacharunderscore}{\kern0pt}dropWhileD\ set{\isacharunderscore}{\kern0pt}takeWhileD{\isacharparenright}{\kern0pt}\isanewline
\isacommand{next}\isamarkupfalse%
\ \isanewline
\ \ \isacommand{assume}\isamarkupfalse%
\ {\isacharquery}{\kern0pt}r\isanewline
\ \ \isacommand{hence}\isamarkupfalse%
\ {\isachardoublequoteopen}t\isactrlsub a\ {\isasymin}\ set\ {\isacharparenleft}{\kern0pt}map\ fst\ hs{\isacharparenright}{\kern0pt}{\isachardoublequoteclose}\isanewline
\ \ \ \ \isacommand{by}\isamarkupfalse%
\ auto\isanewline
\ \ \isacommand{hence}\isamarkupfalse%
\ {\isachardoublequoteopen}t\isactrlsub a\ {\isasymin}\ set\ {\isacharparenleft}{\kern0pt}map\ fst\ {\isacharparenleft}{\kern0pt}{\isacharparenleft}{\kern0pt}dropWhile\ {\isacharparenleft}{\kern0pt}leq\ t\isactrlsub i{\isacharparenright}{\kern0pt}\ {\isacharparenleft}{\kern0pt}takeWhile\ {\isacharparenleft}{\kern0pt}le\ t\isactrlsub j{\isacharparenright}{\kern0pt}\ hs{\isacharparenright}{\kern0pt}{\isacharparenright}{\kern0pt}{\isacharparenright}{\kern0pt}{\isacharparenright}{\kern0pt}{\isachardoublequoteclose}\isanewline
\ \ \ \ \isacommand{using}\isamarkupfalse%
\ strict{\isacharunderscore}{\kern0pt}sorted{\isacharunderscore}{\kern0pt}takeWhile{\isacharbrackleft}{\kern0pt}OF\ {\isacartoucheopen}strict{\isacharunderscore}{\kern0pt}sorted\ {\isacharparenleft}{\kern0pt}map\ fst\ hs{\isacharparenright}{\kern0pt}{\isacartoucheclose}{\isacharbrackright}{\kern0pt}\ {\isacartoucheopen}strict{\isacharunderscore}{\kern0pt}sorted\ {\isacharparenleft}{\kern0pt}map\ fst\ hs{\isacharparenright}{\kern0pt}{\isacartoucheclose}\ {\isacartoucheopen}{\isacharquery}{\kern0pt}r{\isacartoucheclose}\isanewline
\ \ \ \ \isacommand{by}\isamarkupfalse%
{\isacharparenleft}{\kern0pt}auto\ simp{\isacharcolon}{\kern0pt}\ sorted{\isacharunderscore}{\kern0pt}takeWhile\ sorted{\isacharunderscore}{\kern0pt}dropWhile\ map{\isacharunderscore}{\kern0pt}fst{\isacharunderscore}{\kern0pt}takeWhile\ map{\isacharunderscore}{\kern0pt}fst{\isacharunderscore}{\kern0pt}dropWhile{\isacharparenright}{\kern0pt}\isanewline
\ \ \isacommand{then}\isamarkupfalse%
\ \isacommand{obtain}\isamarkupfalse%
\ A{\isacharprime}{\kern0pt}\ \isakeyword{where}\ {\isachardoublequoteopen}{\isacharparenleft}{\kern0pt}t\isactrlsub a{\isacharcomma}{\kern0pt}\ A{\isacharprime}{\kern0pt}{\isacharparenright}{\kern0pt}\ {\isasymin}\ set\ {\isacharparenleft}{\kern0pt}{\isacharparenleft}{\kern0pt}dropWhile\ {\isacharparenleft}{\kern0pt}leq\ t\isactrlsub i{\isacharparenright}{\kern0pt}\ {\isacharparenleft}{\kern0pt}takeWhile\ {\isacharparenleft}{\kern0pt}le\ t\isactrlsub j{\isacharparenright}{\kern0pt}\ hs{\isacharparenright}{\kern0pt}{\isacharparenright}{\kern0pt}{\isacharparenright}{\kern0pt}{\isachardoublequoteclose}\ {\isachardoublequoteopen}{\isacharparenleft}{\kern0pt}t\isactrlsub a{\isacharcomma}{\kern0pt}\ A{\isacharprime}{\kern0pt}{\isacharparenright}{\kern0pt}\ {\isasymin}\ set\ hs{\isachardoublequoteclose}\isanewline
\ \ \ \ \isacommand{by}\isamarkupfalse%
\ {\isacharparenleft}{\kern0pt}fastforce\ dest{\isacharcolon}{\kern0pt}\ set{\isacharunderscore}{\kern0pt}dropWhileD\ set{\isacharunderscore}{\kern0pt}takeWhileD{\isacharparenright}{\kern0pt}\isanewline
\ \ \isacommand{thus}\isamarkupfalse%
\ {\isacharquery}{\kern0pt}l\isanewline
\ \ \ \ \isacommand{using}\isamarkupfalse%
\ {\isacartoucheopen}{\isacharquery}{\kern0pt}r{\isacartoucheclose}\ strict{\isacharunderscore}{\kern0pt}sorted{\isacharunderscore}{\kern0pt}fst{\isacharunderscore}{\kern0pt}unique{\isacharunderscore}{\kern0pt}snd\ {\isacartoucheopen}strict{\isacharunderscore}{\kern0pt}sorted\ {\isacharparenleft}{\kern0pt}map\ fst\ hs{\isacharparenright}{\kern0pt}{\isacartoucheclose}\isanewline
\ \ \ \ \isacommand{by}\isamarkupfalse%
\ force\isanewline
\isacommand{qed}\isamarkupfalse%
%
\endisatagproof
{\isafoldproof}%
%
\isadelimproof
\isanewline
%
\endisadelimproof
\isanewline
\isacommand{lemma}\isamarkupfalse%
\ ind{\isacharunderscore}{\kern0pt}happ{\isacharunderscore}{\kern0pt}seqs{\isacharunderscore}{\kern0pt}same{\isacharunderscore}{\kern0pt}acts{\isacharunderscore}{\kern0pt}leq\isactrlsub i{\isacharunderscore}{\kern0pt}le\isactrlsub j{\isacharprime}{\kern0pt}{\isacharcolon}{\kern0pt}\isanewline
\ \ \isakeyword{assumes}\ {\isachardoublequoteopen}is{\isacharunderscore}{\kern0pt}htp\ {\isasympi}s\ t\isactrlsub i{\isachardoublequoteclose}\isanewline
\ \ \ \ \ \ \ \ \isakeyword{and}\ {\isachardoublequoteopen}is{\isacharunderscore}{\kern0pt}htp\ {\isasympi}s\ t\isactrlsub j{\isachardoublequoteclose}\isanewline
\ \ \ \ \ \ \ \ \isakeyword{and}\ {\isachardoublequoteopen}ind{\isacharunderscore}{\kern0pt}happ{\isacharunderscore}{\kern0pt}seq\ {\isasympi}s\ hs\isactrlsub {\isadigit{1}}{\isachardoublequoteclose}\ \isanewline
\ \ \ \ \ \ \ \ \isakeyword{and}\ {\isachardoublequoteopen}ind{\isacharunderscore}{\kern0pt}happ{\isacharunderscore}{\kern0pt}seq\ {\isasympi}s\ hs\isactrlsub {\isadigit{2}}{\isachardoublequoteclose}\isanewline
\ \ \ \ \ \ \ \ \isakeyword{and}\ {\isachardoublequoteopen}{\isacharparenleft}{\kern0pt}t\isactrlsub a{\isacharcomma}{\kern0pt}A{\isacharparenright}{\kern0pt}\ {\isasymin}\ set\ {\isacharparenleft}{\kern0pt}dropWhile\ {\isacharparenleft}{\kern0pt}leq\ t\isactrlsub i{\isacharparenright}{\kern0pt}\ {\isacharparenleft}{\kern0pt}takeWhile\ {\isacharparenleft}{\kern0pt}le\ t\isactrlsub j{\isacharparenright}{\kern0pt}\ hs\isactrlsub {\isadigit{1}}{\isacharparenright}{\kern0pt}{\isacharparenright}{\kern0pt}{\isachardoublequoteclose}\isanewline
\ \ \ \ \ \ \ \ \isakeyword{and}\ {\isachardoublequoteopen}a\ {\isasymin}\ set\ A{\isachardoublequoteclose}\isanewline
\ \ \ \ \ \ \isakeyword{shows}\ {\isachardoublequoteopen}{\isasymexists}t\isactrlsub a{\isacharprime}{\kern0pt}\ A{\isacharprime}{\kern0pt}{\isachardot}{\kern0pt}\ {\isacharparenleft}{\kern0pt}t\isactrlsub a{\isacharprime}{\kern0pt}{\isacharcomma}{\kern0pt}A{\isacharprime}{\kern0pt}{\isacharparenright}{\kern0pt}\ {\isasymin}\ set\ {\isacharparenleft}{\kern0pt}dropWhile\ {\isacharparenleft}{\kern0pt}leq\ t\isactrlsub i{\isacharparenright}{\kern0pt}\ {\isacharparenleft}{\kern0pt}takeWhile\ {\isacharparenleft}{\kern0pt}le\ t\isactrlsub j{\isacharparenright}{\kern0pt}\ hs\isactrlsub {\isadigit{2}}{\isacharparenright}{\kern0pt}{\isacharparenright}{\kern0pt}\ {\isasymand}\ a\ {\isasymin}\ set\ A{\isacharprime}{\kern0pt}{\isachardoublequoteclose}\isanewline
%
\isadelimproof
\ \ %
\endisadelimproof
%
\isatagproof
\isacommand{proof}\isamarkupfalse%
\ {\isacharminus}{\kern0pt}\isanewline
\ \ \ \ \isacommand{let}\isamarkupfalse%
\ {\isacharquery}{\kern0pt}hs\isactrlsub {\isadigit{1}}{\isacharprime}{\kern0pt}{\isacharequal}{\kern0pt}{\isachardoublequoteopen}dropWhile\ {\isacharparenleft}{\kern0pt}leq\ t\isactrlsub i{\isacharparenright}{\kern0pt}\ {\isacharparenleft}{\kern0pt}takeWhile\ {\isacharparenleft}{\kern0pt}le\ t\isactrlsub j{\isacharparenright}{\kern0pt}\ hs\isactrlsub {\isadigit{1}}{\isacharparenright}{\kern0pt}{\isachardoublequoteclose}\isanewline
\ \ \ \ \isacommand{let}\isamarkupfalse%
\ {\isacharquery}{\kern0pt}hs\isactrlsub {\isadigit{2}}{\isacharprime}{\kern0pt}{\isacharequal}{\kern0pt}{\isachardoublequoteopen}dropWhile\ {\isacharparenleft}{\kern0pt}leq\ t\isactrlsub i{\isacharparenright}{\kern0pt}\ {\isacharparenleft}{\kern0pt}takeWhile\ {\isacharparenleft}{\kern0pt}le\ t\isactrlsub j{\isacharparenright}{\kern0pt}\ hs\isactrlsub {\isadigit{2}}{\isacharparenright}{\kern0pt}{\isachardoublequoteclose}\isanewline
\ \ \ \ \isacommand{have}\isamarkupfalse%
\ {\isachardoublequoteopen}{\isacharparenleft}{\kern0pt}t\isactrlsub a{\isacharcomma}{\kern0pt}A{\isacharparenright}{\kern0pt}\ {\isasymin}\ set\ hs\isactrlsub {\isadigit{1}}{\isachardoublequoteclose}\isanewline
\ \ \ \ \ \ \isacommand{using}\isamarkupfalse%
\ {\isacartoucheopen}{\isacharparenleft}{\kern0pt}t\isactrlsub a{\isacharcomma}{\kern0pt}A{\isacharparenright}{\kern0pt}\ {\isasymin}\ set\ {\isacharquery}{\kern0pt}hs\isactrlsub {\isadigit{1}}{\isacharprime}{\kern0pt}{\isacartoucheclose}\ \isacommand{by}\isamarkupfalse%
\ {\isacharparenleft}{\kern0pt}meson\ set{\isacharunderscore}{\kern0pt}dropWhileD\ set{\isacharunderscore}{\kern0pt}takeWhileD{\isacharparenright}{\kern0pt}\ \isanewline
\ \ \ \ \isacommand{have}\isamarkupfalse%
\ {\isachardoublequoteopen}strict{\isacharunderscore}{\kern0pt}sorted\ {\isacharparenleft}{\kern0pt}map\ fst\ hs\isactrlsub {\isadigit{1}}{\isacharparenright}{\kern0pt}{\isachardoublequoteclose}\ \isakeyword{and}\ {\isachardoublequoteopen}strict{\isacharunderscore}{\kern0pt}sorted\ {\isacharparenleft}{\kern0pt}map\ fst\ hs\isactrlsub {\isadigit{2}}{\isacharparenright}{\kern0pt}{\isachardoublequoteclose}\isanewline
\ \ \ \ \ \ \isacommand{using}\isamarkupfalse%
\ {\isacartoucheopen}ind{\isacharunderscore}{\kern0pt}happ{\isacharunderscore}{\kern0pt}seq\ {\isasympi}s\ hs\isactrlsub {\isadigit{1}}{\isacartoucheclose}\ {\isacartoucheopen}ind{\isacharunderscore}{\kern0pt}happ{\isacharunderscore}{\kern0pt}seq\ {\isasympi}s\ hs\isactrlsub {\isadigit{2}}{\isacartoucheclose}\ \isacommand{unfolding}\isamarkupfalse%
\ ind{\isacharunderscore}{\kern0pt}happ{\isacharunderscore}{\kern0pt}seq{\isacharunderscore}{\kern0pt}def\ \isacommand{by}\isamarkupfalse%
\ auto\isanewline
\ \ \ \ \isacommand{moreover}\isamarkupfalse%
\ \isacommand{have}\isamarkupfalse%
\ {\isachardoublequoteopen}t\isactrlsub i\ {\isacharless}{\kern0pt}\ t\isactrlsub a\ {\isasymand}\ t\isactrlsub a\ {\isacharless}{\kern0pt}\ t\isactrlsub j{\isachardoublequoteclose}\isanewline
\ \ \ \ \ \ \isacommand{using}\isamarkupfalse%
\ {\isacartoucheopen}{\isacharparenleft}{\kern0pt}t\isactrlsub a{\isacharcomma}{\kern0pt}A{\isacharparenright}{\kern0pt}\ {\isasymin}\ set\ {\isacharquery}{\kern0pt}hs\isactrlsub {\isadigit{1}}{\isacharprime}{\kern0pt}{\isacartoucheclose}\ {\isacartoucheopen}strict{\isacharunderscore}{\kern0pt}sorted\ {\isacharparenleft}{\kern0pt}map\ fst\ hs\isactrlsub {\isadigit{1}}{\isacharparenright}{\kern0pt}{\isacartoucheclose}\ {\isacartoucheopen}{\isacharparenleft}{\kern0pt}t\isactrlsub a{\isacharcomma}{\kern0pt}A{\isacharparenright}{\kern0pt}\ {\isasymin}\ set\ hs\isactrlsub {\isadigit{1}}{\isacartoucheclose}\isanewline
\ \ \ \ \ \ \isacommand{by}\isamarkupfalse%
\ {\isacharparenleft}{\kern0pt}fastforce\ simp{\isacharcolon}{\kern0pt}\ strict{\isacharunderscore}{\kern0pt}sorted{\isacharunderscore}{\kern0pt}drop{\isacharunderscore}{\kern0pt}leq{\isacharunderscore}{\kern0pt}take{\isacharunderscore}{\kern0pt}le{\isacharprime}{\kern0pt}{\isacharparenright}{\kern0pt}\isanewline
\ \ \ \ \isacommand{moreover}\isamarkupfalse%
\ \isacommand{obtain}\isamarkupfalse%
\ t\isactrlsub {\isasympi}\ {\isasympi}\ \isakeyword{where}\ {\isachardoublequoteopen}{\isacharparenleft}{\kern0pt}t\isactrlsub {\isasympi}{\isacharcomma}{\kern0pt}{\isasympi}{\isacharparenright}{\kern0pt}\ {\isasymin}\ set\ {\isasympi}s{\isachardoublequoteclose}\ \isakeyword{and}\ {\isachardoublequoteopen}inst{\isacharunderscore}{\kern0pt}of{\isacharunderscore}{\kern0pt}plan{\isacharunderscore}{\kern0pt}action\ {\isasympi}s\ {\isacharparenleft}{\kern0pt}t\isactrlsub {\isasympi}{\isacharcomma}{\kern0pt}{\isasympi}{\isacharparenright}{\kern0pt}\ {\isacharparenleft}{\kern0pt}t\isactrlsub a{\isacharcomma}{\kern0pt}a{\isacharparenright}{\kern0pt}{\isachardoublequoteclose}\isanewline
\ \ \ \ \ \ \isacommand{using}\isamarkupfalse%
\ {\isacartoucheopen}ind{\isacharunderscore}{\kern0pt}happ{\isacharunderscore}{\kern0pt}seq\ {\isasympi}s\ hs\isactrlsub {\isadigit{1}}{\isacartoucheclose}\ {\isacartoucheopen}{\isacharparenleft}{\kern0pt}t\isactrlsub a{\isacharcomma}{\kern0pt}A{\isacharparenright}{\kern0pt}\ {\isasymin}\ set\ hs\isactrlsub {\isadigit{1}}{\isacartoucheclose}\ {\isacartoucheopen}a\ {\isasymin}\ set\ A{\isacartoucheclose}\isanewline
\ \ \ \ \ \ \isacommand{by}\isamarkupfalse%
\ {\isacharparenleft}{\kern0pt}fastforce\ simp{\isacharcolon}{\kern0pt}\ ind{\isacharunderscore}{\kern0pt}happ{\isacharunderscore}{\kern0pt}seq{\isacharunderscore}{\kern0pt}def{\isacharparenright}{\kern0pt}\ \ \ \ \ \ \ \ \ \ \ \ \ \ \ \ \ \ \ \ \ \ \ \ \ \ \ \ \ \ \isanewline
\ \ \ \ \isacommand{ultimately}\isamarkupfalse%
\ \isacommand{show}\isamarkupfalse%
\ {\isacharquery}{\kern0pt}thesis\isanewline
\ \ \ \ \isacommand{proof}\isamarkupfalse%
\ {\isacharparenleft}{\kern0pt}elim\ inst{\isacharunderscore}{\kern0pt}of{\isacharunderscore}{\kern0pt}plan{\isacharunderscore}{\kern0pt}actionE{\isacharcomma}{\kern0pt}\ goal{\isacharunderscore}{\kern0pt}cases{\isacharparenright}{\kern0pt}\isanewline
\ \ \ \ \ \ \isacommand{case}\isamarkupfalse%
\ {\isacharparenleft}{\kern0pt}{\isadigit{4}}\ t\isactrlsub i{\isacharprime}{\kern0pt}\ t\isactrlsub j{\isacharprime}{\kern0pt}{\isacharparenright}{\kern0pt}\isanewline
\ \ \ \ \ \ \isacommand{let}\isamarkupfalse%
\ {\isacharquery}{\kern0pt}a\isactrlsub i\isactrlsub n\isactrlsub v{\isacharequal}{\kern0pt}{\isachardoublequoteopen}the\ {\isacharparenleft}{\kern0pt}res{\isacharunderscore}{\kern0pt}inst{\isacharunderscore}{\kern0pt}snap{\isacharunderscore}{\kern0pt}action\ {\isasympi}\ Over{\isacharunderscore}{\kern0pt}All{\isacharparenright}{\kern0pt}{\isachardoublequoteclose}\isanewline
\ \ \ \ \ \ \isacommand{have}\isamarkupfalse%
\ {\isachardoublequoteopen}{\isacharparenleft}{\kern0pt}t\isactrlsub {\isasympi}{\isacharcomma}{\kern0pt}{\isasympi}{\isacharparenright}{\kern0pt}\ {\isasymin}\ durative{\isacharunderscore}{\kern0pt}acts\ {\isasympi}s{\isachardoublequoteclose}\isanewline
\ \ \ \ \ \ \ \ \isacommand{unfolding}\isamarkupfalse%
\ durative{\isacharunderscore}{\kern0pt}acts{\isacharunderscore}{\kern0pt}def\ is{\isacharunderscore}{\kern0pt}act{\isacharunderscore}{\kern0pt}simple{\isacharunderscore}{\kern0pt}def\isanewline
\ \ \ \ \ \ \ \ \isacommand{using}\isamarkupfalse%
\ {\isadigit{4}}\ {\isacartoucheopen}{\isacharparenleft}{\kern0pt}t\isactrlsub {\isasympi}{\isacharcomma}{\kern0pt}{\isasympi}{\isacharparenright}{\kern0pt}\ {\isasymin}\ set\ {\isasympi}s{\isacartoucheclose}\ \isacommand{by}\isamarkupfalse%
\ {\isacharparenleft}{\kern0pt}cases\ {\isasympi}{\isacharparenright}{\kern0pt}\ auto\isanewline
\ \ \ \ \ \ \isacommand{then}\isamarkupfalse%
\ \isacommand{have}\isamarkupfalse%
\ {\isachardoublequoteopen}{\isasymexists}{\isacharparenleft}{\kern0pt}t{\isacharprime}{\kern0pt}{\isacharcomma}{\kern0pt}as{\isacharparenright}{\kern0pt}\ {\isasymin}\ set\ hs\isactrlsub {\isadigit{2}}{\isachardot}{\kern0pt}\ t\isactrlsub i{\isacharprime}{\kern0pt}\ {\isacharless}{\kern0pt}\ t{\isacharprime}{\kern0pt}\ {\isasymand}\ t{\isacharprime}{\kern0pt}\ {\isacharless}{\kern0pt}\ t\isactrlsub j{\isacharprime}{\kern0pt}\ {\isasymand}\ {\isacharquery}{\kern0pt}a\isactrlsub i\isactrlsub n\isactrlsub v\ {\isasymin}\ set\ as{\isachardoublequoteclose}\isanewline
\ \ \ \ \ \ \ \ \isacommand{using}\isamarkupfalse%
\ {\isacartoucheopen}ind{\isacharunderscore}{\kern0pt}happ{\isacharunderscore}{\kern0pt}seq\ {\isasympi}s\ hs\isactrlsub {\isadigit{2}}{\isacartoucheclose}\ \isanewline
\ \ \ \ \ \ \ \ \isacommand{using}\isamarkupfalse%
\ {\isacartoucheopen}consec{\isacharunderscore}{\kern0pt}htps\ {\isasympi}s\ t\isactrlsub i{\isacharprime}{\kern0pt}\ t\isactrlsub j{\isacharprime}{\kern0pt}{\isacartoucheclose}\ {\isadigit{4}}\ \isacommand{by}\isamarkupfalse%
\ {\isacharparenleft}{\kern0pt}auto\ simp{\isacharcolon}{\kern0pt}\ ind{\isacharunderscore}{\kern0pt}happ{\isacharunderscore}{\kern0pt}seq{\isacharunderscore}{\kern0pt}def{\isacharparenright}{\kern0pt}\isanewline
\ \ \ \ \ \ \isacommand{then}\isamarkupfalse%
\ \isacommand{obtain}\isamarkupfalse%
\ t\isactrlsub a{\isacharprime}{\kern0pt}\ A{\isacharprime}{\kern0pt}\ \isakeyword{where}\ {\isachardoublequoteopen}{\isacharparenleft}{\kern0pt}t\isactrlsub a{\isacharprime}{\kern0pt}{\isacharcomma}{\kern0pt}A{\isacharprime}{\kern0pt}{\isacharparenright}{\kern0pt}\ {\isasymin}\ set\ hs\isactrlsub {\isadigit{2}}{\isachardoublequoteclose}\ \isakeyword{and}\ {\isachardoublequoteopen}t\isactrlsub i{\isacharprime}{\kern0pt}\ {\isacharless}{\kern0pt}\ t\isactrlsub a{\isacharprime}{\kern0pt}\ {\isasymand}\ t\isactrlsub a{\isacharprime}{\kern0pt}\ {\isacharless}{\kern0pt}\ t\isactrlsub j{\isacharprime}{\kern0pt}{\isachardoublequoteclose}\ \isakeyword{and}\ {\isachardoublequoteopen}{\isacharquery}{\kern0pt}a\isactrlsub i\isactrlsub n\isactrlsub v\ {\isasymin}\ set\ A{\isacharprime}{\kern0pt}{\isachardoublequoteclose}\isanewline
\ \ \ \ \ \ \ \ \isacommand{by}\isamarkupfalse%
\ auto\isanewline
\ \ \ \ \ \ \isacommand{then}\isamarkupfalse%
\ \isacommand{have}\isamarkupfalse%
\ {\isachardoublequoteopen}{\isacharparenleft}{\kern0pt}t\isactrlsub a{\isacharprime}{\kern0pt}{\isacharcomma}{\kern0pt}A{\isacharprime}{\kern0pt}{\isacharparenright}{\kern0pt}\ {\isasymin}\ set\ hs\isactrlsub {\isadigit{2}}{\isachardoublequoteclose}\isanewline
\ \ \ \ \ \ \ \ \isacommand{using}\isamarkupfalse%
\ {\isadigit{4}}\ \isacommand{by}\isamarkupfalse%
\ auto\isanewline
\ \ \ \ \ \ \isacommand{moreover}\isamarkupfalse%
\ \isacommand{have}\isamarkupfalse%
\ {\isachardoublequoteopen}t\isactrlsub i\ {\isacharless}{\kern0pt}\ t\isactrlsub a{\isacharprime}{\kern0pt}\ {\isasymand}\ t\isactrlsub a{\isacharprime}{\kern0pt}\ {\isacharless}{\kern0pt}\ t\isactrlsub j{\isachardoublequoteclose}\isanewline
\ \ \ \ \ \ \ \ \isacommand{using}\isamarkupfalse%
\ assms\ {\isacartoucheopen}consec{\isacharunderscore}{\kern0pt}htps\ {\isasympi}s\ t\isactrlsub i{\isacharprime}{\kern0pt}\ t\isactrlsub j{\isacharprime}{\kern0pt}{\isacartoucheclose}\ {\isacartoucheopen}t\isactrlsub i\ {\isacharless}{\kern0pt}\ t\isactrlsub a\ {\isasymand}\ t\isactrlsub a\ {\isacharless}{\kern0pt}\ t\isactrlsub j{\isacartoucheclose}\ {\isacartoucheopen}t\isactrlsub i{\isacharprime}{\kern0pt}\ {\isacharless}{\kern0pt}\ t\isactrlsub a{\isacartoucheclose}\ {\isacartoucheopen}t\isactrlsub a\ {\isacharless}{\kern0pt}\ t\isactrlsub j{\isacharprime}{\kern0pt}{\isacartoucheclose}\ \isanewline
\ \ \ \ \ \ \ \ \ \ \ \ \ \ {\isacartoucheopen}t\isactrlsub i{\isacharprime}{\kern0pt}\ {\isacharless}{\kern0pt}\ t\isactrlsub a{\isacharprime}{\kern0pt}\ {\isasymand}\ t\isactrlsub a{\isacharprime}{\kern0pt}\ {\isacharless}{\kern0pt}\ t\isactrlsub j{\isacharprime}{\kern0pt}{\isacartoucheclose}\isanewline
\ \ \ \ \ \ \ \ \isacommand{by}\isamarkupfalse%
\ {\isacharparenleft}{\kern0pt}auto\ simp{\isacharcolon}{\kern0pt}\ consec{\isacharunderscore}{\kern0pt}htps{\isacharunderscore}{\kern0pt}def{\isacharparenright}{\kern0pt}\isanewline
\ \ \ \ \ \ \isacommand{ultimately}\isamarkupfalse%
\ \isacommand{show}\isamarkupfalse%
\ {\isacharquery}{\kern0pt}thesis\isanewline
\ \ \ \ \ \ \ \ \isacommand{using}\isamarkupfalse%
\ {\isadigit{4}}\ {\isacartoucheopen}{\isacharquery}{\kern0pt}a\isactrlsub i\isactrlsub n\isactrlsub v\ {\isasymin}\ set\ A{\isacharprime}{\kern0pt}{\isacartoucheclose}\ {\isacartoucheopen}strict{\isacharunderscore}{\kern0pt}sorted\ {\isacharparenleft}{\kern0pt}map\ fst\ hs\isactrlsub {\isadigit{2}}{\isacharparenright}{\kern0pt}{\isacartoucheclose}\isanewline
\ \ \ \ \ \ \ \ \isacommand{by}\isamarkupfalse%
{\isacharparenleft}{\kern0pt}fastforce\ simp{\isacharcolon}{\kern0pt}\ strict{\isacharunderscore}{\kern0pt}sorted{\isacharunderscore}{\kern0pt}drop{\isacharunderscore}{\kern0pt}leq{\isacharunderscore}{\kern0pt}take{\isacharunderscore}{\kern0pt}le{\isacharprime}{\kern0pt}{\isacharparenright}{\kern0pt}\isanewline
\ \ \ \ \isacommand{qed}\isamarkupfalse%
\ {\isacharparenleft}{\kern0pt}insert\ {\isacartoucheopen}ind{\isacharunderscore}{\kern0pt}happ{\isacharunderscore}{\kern0pt}seq\ {\isasympi}s\ hs\isactrlsub {\isadigit{2}}{\isacartoucheclose}{\isacharcomma}{\kern0pt}\ cases\ {\isasympi}{\isacharsemicolon}{\kern0pt}\isanewline
\ \ \ \ \ \ \ \ \ fastforce\ simp{\isacharcolon}{\kern0pt}\ strict{\isacharunderscore}{\kern0pt}sorted{\isacharunderscore}{\kern0pt}drop{\isacharunderscore}{\kern0pt}leq{\isacharunderscore}{\kern0pt}take{\isacharunderscore}{\kern0pt}le{\isacharprime}{\kern0pt}\ ind{\isacharunderscore}{\kern0pt}happ{\isacharunderscore}{\kern0pt}seq{\isacharunderscore}{\kern0pt}def\ simple{\isacharunderscore}{\kern0pt}acts{\isacharunderscore}{\kern0pt}def\isanewline
\ \ \ \ \ \ \ \ \ \ \ \ \ \ \ \ \ \ \ \ \ \ \ \ \ durative{\isacharunderscore}{\kern0pt}acts{\isacharunderscore}{\kern0pt}def\ is{\isacharunderscore}{\kern0pt}act{\isacharunderscore}{\kern0pt}simple{\isacharunderscore}{\kern0pt}def{\isacharparenright}{\kern0pt}{\isacharplus}{\kern0pt}\isanewline
\ \ \isacommand{qed}\isamarkupfalse%
%
\endisatagproof
{\isafoldproof}%
%
\isadelimproof
\isanewline
%
\endisadelimproof
\isanewline
\isacommand{lemma}\isamarkupfalse%
\ happs{\isacharunderscore}{\kern0pt}same{\isacharunderscore}{\kern0pt}between{\isacharunderscore}{\kern0pt}htps{\isacharcolon}{\kern0pt}\isanewline
\ \ \isakeyword{assumes}\ {\isachardoublequoteopen}ind{\isacharunderscore}{\kern0pt}happ{\isacharunderscore}{\kern0pt}seq\ {\isasympi}s\ hs{\isachardoublequoteclose}\isanewline
\ \ \ \ \ \ \isakeyword{and}\ {\isachardoublequoteopen}ind{\isacharunderscore}{\kern0pt}happ{\isacharunderscore}{\kern0pt}seq\ {\isasympi}s\ hs{\isacharprime}{\kern0pt}{\isachardoublequoteclose}\isanewline
\ \ \ \ \ \ \isakeyword{and}\ {\isachardoublequoteopen}is{\isacharunderscore}{\kern0pt}htp\ {\isasympi}s\ t\isactrlsub {\isadigit{1}}{\isachardoublequoteclose}\ {\isachardoublequoteopen}is{\isacharunderscore}{\kern0pt}htp\ {\isasympi}s\ t\isactrlsub {\isadigit{2}}{\isachardoublequoteclose}\isanewline
\ \ \ \ \ \ \isakeyword{and}\ {\isachardoublequoteopen}{\isacharparenleft}{\kern0pt}t{\isacharcomma}{\kern0pt}\ A{\isacharparenright}{\kern0pt}\ {\isasymin}\ set\ hs{\isachardoublequoteclose}\isanewline
\ \ \ \ \ \ \isakeyword{and}\ {\isachardoublequoteopen}{\isasymnot}is{\isacharunderscore}{\kern0pt}htp\ {\isasympi}s\ t{\isachardoublequoteclose}\isanewline
\ \ \ \ \ \ \isakeyword{and}\ {\isachardoublequoteopen}t\isactrlsub {\isadigit{1}}\ {\isacharless}{\kern0pt}\ t{\isachardoublequoteclose}\isanewline
\ \ \ \ \ \ \isakeyword{and}\ {\isachardoublequoteopen}t\ {\isacharless}{\kern0pt}\ t\isactrlsub {\isadigit{2}}{\isachardoublequoteclose}\isanewline
\ \ \ \ \ \ \isakeyword{and}\ {\isachardoublequoteopen}a\ {\isasymin}\ set\ A{\isachardoublequoteclose}\isanewline
\ \ \isakeyword{shows}\ {\isachardoublequoteopen}{\isasymexists}t{\isacharprime}{\kern0pt}{\isasymin}{\isacharbraceleft}{\kern0pt}t\isactrlsub {\isadigit{1}}{\isacharless}{\kern0pt}{\isachardot}{\kern0pt}{\isachardot}{\kern0pt}{\isacharless}{\kern0pt}t\isactrlsub {\isadigit{2}}{\isacharbraceright}{\kern0pt}{\isachardot}{\kern0pt}\ {\isasymexists}A{\isacharprime}{\kern0pt}{\isachardot}{\kern0pt}\ {\isacharparenleft}{\kern0pt}t{\isacharprime}{\kern0pt}{\isacharcomma}{\kern0pt}A{\isacharprime}{\kern0pt}{\isacharparenright}{\kern0pt}{\isasymin}set\ hs{\isacharprime}{\kern0pt}{\isasymand}\ a\ {\isasymin}\ set\ A{\isacharprime}{\kern0pt}{\isachardoublequoteclose}\isanewline
%
\isadelimproof
%
\endisadelimproof
%
\isatagproof
\isacommand{proof}\isamarkupfalse%
{\isacharminus}{\kern0pt}\isanewline
\ \ \isacommand{have}\isamarkupfalse%
\ {\isachardoublequoteopen}{\isasymexists}t\isactrlsub a{\isacharprime}{\kern0pt}\ A{\isacharprime}{\kern0pt}{\isachardot}{\kern0pt}\ {\isacharparenleft}{\kern0pt}t\isactrlsub a{\isacharprime}{\kern0pt}{\isacharcomma}{\kern0pt}A{\isacharprime}{\kern0pt}{\isacharparenright}{\kern0pt}\ {\isasymin}\ set\ {\isacharparenleft}{\kern0pt}dropWhile\ {\isacharparenleft}{\kern0pt}leq\ t\isactrlsub {\isadigit{1}}{\isacharparenright}{\kern0pt}\ {\isacharparenleft}{\kern0pt}takeWhile\ {\isacharparenleft}{\kern0pt}le\ t\isactrlsub {\isadigit{2}}{\isacharparenright}{\kern0pt}\ hs{\isacharprime}{\kern0pt}{\isacharparenright}{\kern0pt}{\isacharparenright}{\kern0pt}\ {\isasymand}\ a\ {\isasymin}\ set\ A{\isacharprime}{\kern0pt}{\isachardoublequoteclose}\isanewline
\ \ \ \ \isacommand{using}\isamarkupfalse%
\ ind{\isacharunderscore}{\kern0pt}happ{\isacharunderscore}{\kern0pt}seq{\isacharunderscore}{\kern0pt}props{\isacharparenleft}{\kern0pt}{\isadigit{1}}{\isacharparenright}{\kern0pt}{\isacharbrackleft}{\kern0pt}OF\ {\isacartoucheopen}ind{\isacharunderscore}{\kern0pt}happ{\isacharunderscore}{\kern0pt}seq\ {\isasympi}s\ hs{\isacartoucheclose}{\isacharbrackright}{\kern0pt}\ assms\isanewline
\ \ \ \ \isacommand{by}\isamarkupfalse%
\ {\isacharparenleft}{\kern0pt}fastforce\ simp{\isacharcolon}{\kern0pt}\ strict{\isacharunderscore}{\kern0pt}sorted{\isacharunderscore}{\kern0pt}drop{\isacharunderscore}{\kern0pt}leq{\isacharunderscore}{\kern0pt}take{\isacharunderscore}{\kern0pt}le{\isacharprime}{\kern0pt}\ ind{\isacharunderscore}{\kern0pt}happ{\isacharunderscore}{\kern0pt}seqs{\isacharunderscore}{\kern0pt}same{\isacharunderscore}{\kern0pt}acts{\isacharunderscore}{\kern0pt}leq\isactrlsub i{\isacharunderscore}{\kern0pt}le\isactrlsub j{\isacharprime}{\kern0pt}{\isacharparenright}{\kern0pt}\isanewline
\ \ \isacommand{thus}\isamarkupfalse%
\ {\isacharquery}{\kern0pt}thesis\isanewline
\ \ \ \ \isacommand{using}\isamarkupfalse%
\ {\isacartoucheopen}{\isacharparenleft}{\kern0pt}t{\isacharcomma}{\kern0pt}\ A{\isacharparenright}{\kern0pt}\ {\isasymin}\ set\ hs{\isacartoucheclose}\ ind{\isacharunderscore}{\kern0pt}happ{\isacharunderscore}{\kern0pt}seq{\isacharunderscore}{\kern0pt}props{\isacharparenleft}{\kern0pt}{\isadigit{1}}{\isacharparenright}{\kern0pt}{\isacharbrackleft}{\kern0pt}OF\ {\isacartoucheopen}ind{\isacharunderscore}{\kern0pt}happ{\isacharunderscore}{\kern0pt}seq\ {\isasympi}s\ hs{\isacharprime}{\kern0pt}{\isacartoucheclose}{\isacharbrackright}{\kern0pt}\isanewline
\ \ \ \ \isacommand{by}\isamarkupfalse%
\ {\isacharparenleft}{\kern0pt}auto\ simp{\isacharcolon}{\kern0pt}\ strict{\isacharunderscore}{\kern0pt}sorted{\isacharunderscore}{\kern0pt}drop{\isacharunderscore}{\kern0pt}leq{\isacharunderscore}{\kern0pt}take{\isacharunderscore}{\kern0pt}le{\isacharprime}{\kern0pt}{\isacharparenright}{\kern0pt}\isanewline
\isacommand{qed}\isamarkupfalse%
%
\endisatagproof
{\isafoldproof}%
%
\isadelimproof
\isanewline
%
\endisadelimproof
\isanewline
\isacommand{lemma}\isamarkupfalse%
\ consec{\isacharunderscore}{\kern0pt}htps{\isacharunderscore}{\kern0pt}props{\isacharcolon}{\kern0pt}\isanewline
\ \ \isakeyword{assumes}\ {\isachardoublequoteopen}consec{\isacharunderscore}{\kern0pt}htps\ {\isasympi}s\ t\isactrlsub i\ t\isactrlsub j{\isachardoublequoteclose}\isanewline
\ \ \isakeyword{shows}\ {\isachardoublequoteopen}t\isactrlsub i\ {\isacharless}{\kern0pt}\ t\isactrlsub j{\isachardoublequoteclose}\isanewline
\ \ \ \ \ \ \ \ {\isachardoublequoteopen}is{\isacharunderscore}{\kern0pt}htp\ {\isasympi}s\ t\isactrlsub i{\isachardoublequoteclose}\isanewline
\ \ \ \ \ \ \ \ {\isachardoublequoteopen}is{\isacharunderscore}{\kern0pt}htp\ {\isasympi}s\ t\isactrlsub j{\isachardoublequoteclose}\isanewline
\ \ \ \ \ \ \ \ {\isachardoublequoteopen}is{\isacharunderscore}{\kern0pt}htp\ {\isasympi}s\ t\ {\isasymLongrightarrow}\ t\ {\isasymle}\ t\isactrlsub i\ {\isasymor}\ t\isactrlsub j\ {\isasymle}\ t{\isachardoublequoteclose}\isanewline
%
\isadelimproof
\ \ %
\endisadelimproof
%
\isatagproof
\isacommand{using}\isamarkupfalse%
\ assms\isanewline
\ \ \isacommand{by}\isamarkupfalse%
\ {\isacharparenleft}{\kern0pt}auto\ simp{\isacharcolon}{\kern0pt}\ consec{\isacharunderscore}{\kern0pt}htps{\isacharunderscore}{\kern0pt}def{\isacharparenright}{\kern0pt}%
\endisatagproof
{\isafoldproof}%
%
\isadelimproof
\isanewline
%
\endisadelimproof
\isanewline
\isacommand{lemma}\isamarkupfalse%
\ ind{\isacharunderscore}{\kern0pt}happ{\isacharunderscore}{\kern0pt}seq{\isacharunderscore}{\kern0pt}hd{\isacharunderscore}{\kern0pt}htp{\isacharcolon}{\kern0pt}\isanewline
\ \ \isakeyword{assumes}\ {\isachardoublequoteopen}ind{\isacharunderscore}{\kern0pt}happ{\isacharunderscore}{\kern0pt}seq\ {\isasympi}s\ hs{\isachardoublequoteclose}\ {\isachardoublequoteopen}hs\ {\isasymnoteq}\ {\isacharbrackleft}{\kern0pt}{\isacharbrackright}{\kern0pt}{\isachardoublequoteclose}\isanewline
\ \ \isakeyword{shows}\ {\isachardoublequoteopen}is{\isacharunderscore}{\kern0pt}htp\ {\isasympi}s\ {\isacharparenleft}{\kern0pt}fst\ {\isacharparenleft}{\kern0pt}hd\ hs{\isacharparenright}{\kern0pt}{\isacharparenright}{\kern0pt}{\isachardoublequoteclose}\isanewline
%
\isadelimproof
%
\endisadelimproof
%
\isatagproof
\isacommand{proof}\isamarkupfalse%
{\isacharminus}{\kern0pt}\isanewline
\ \ \isacommand{obtain}\isamarkupfalse%
\ h\ hs{\isacharprime}{\kern0pt}\ \isakeyword{where}\ hs{\isacharbrackleft}{\kern0pt}simp{\isacharbrackright}{\kern0pt}{\isacharcolon}{\kern0pt}\ {\isachardoublequoteopen}hs\ {\isacharequal}{\kern0pt}\ h{\isacharhash}{\kern0pt}hs{\isacharprime}{\kern0pt}{\isachardoublequoteclose}\isanewline
\ \ \ \ \isacommand{using}\isamarkupfalse%
\ {\isacartoucheopen}hs\ {\isasymnoteq}\ {\isacharbrackleft}{\kern0pt}{\isacharbrackright}{\kern0pt}{\isacartoucheclose}\isanewline
\ \ \ \ \isacommand{by}\isamarkupfalse%
\ {\isacharparenleft}{\kern0pt}auto\ simp{\isacharcolon}{\kern0pt}\ neq{\isacharunderscore}{\kern0pt}Nil{\isacharunderscore}{\kern0pt}conv{\isacharparenright}{\kern0pt}\isanewline
\ \ \isacommand{then}\isamarkupfalse%
\ \isacommand{obtain}\isamarkupfalse%
\ t\ A\ \isakeyword{where}\ h{\isacharbrackleft}{\kern0pt}simp{\isacharbrackright}{\kern0pt}{\isacharcolon}{\kern0pt}\ {\isachardoublequoteopen}h\ {\isacharequal}{\kern0pt}\ {\isacharparenleft}{\kern0pt}t{\isacharcomma}{\kern0pt}\ A{\isacharparenright}{\kern0pt}{\isachardoublequoteclose}\isanewline
\ \ \ \ \isacommand{by}\isamarkupfalse%
\ {\isacharparenleft}{\kern0pt}elim\ prod{\isachardot}{\kern0pt}exhaust{\isacharparenright}{\kern0pt}\isanewline
\ \ \isacommand{then}\isamarkupfalse%
\ \isacommand{obtain}\isamarkupfalse%
\ a\ \isakeyword{where}\ {\isachardoublequoteopen}a\ {\isasymin}\ set\ A{\isachardoublequoteclose}\isanewline
\ \ \ \ \isacommand{using}\isamarkupfalse%
\ ind{\isacharunderscore}{\kern0pt}happ{\isacharunderscore}{\kern0pt}seq{\isacharunderscore}{\kern0pt}props{\isacharparenleft}{\kern0pt}{\isadigit{6}}{\isacharparenright}{\kern0pt}{\isacharbrackleft}{\kern0pt}OF\ {\isacartoucheopen}ind{\isacharunderscore}{\kern0pt}happ{\isacharunderscore}{\kern0pt}seq\ {\isasympi}s\ hs{\isacartoucheclose}{\isacharbrackright}{\kern0pt}\isanewline
\ \ \ \ \isacommand{by}\isamarkupfalse%
\ auto\isanewline
\ \ \isacommand{then}\isamarkupfalse%
\ \isacommand{obtain}\isamarkupfalse%
\ t\isactrlsub {\isasympi}\ {\isasympi}\ \isakeyword{where}\ {\isachardoublequoteopen}{\isacharparenleft}{\kern0pt}t\isactrlsub {\isasympi}{\isacharcomma}{\kern0pt}{\isasympi}{\isacharparenright}{\kern0pt}\ {\isasymin}\ set\ {\isasympi}s{\isachardoublequoteclose}\ \isakeyword{and}\ {\isachardoublequoteopen}inst{\isacharunderscore}{\kern0pt}of{\isacharunderscore}{\kern0pt}plan{\isacharunderscore}{\kern0pt}action\ {\isasympi}s\ {\isacharparenleft}{\kern0pt}t\isactrlsub {\isasympi}{\isacharcomma}{\kern0pt}{\isasympi}{\isacharparenright}{\kern0pt}\ {\isacharparenleft}{\kern0pt}t{\isacharcomma}{\kern0pt}a{\isacharparenright}{\kern0pt}{\isachardoublequoteclose}\isanewline
\ \ \ \ \isacommand{using}\isamarkupfalse%
\ {\isacartoucheopen}ind{\isacharunderscore}{\kern0pt}happ{\isacharunderscore}{\kern0pt}seq\ {\isasympi}s\ hs{\isacartoucheclose}\ {\isacartoucheopen}a\ {\isasymin}\ set\ A{\isacartoucheclose}\isanewline
\ \ \ \ \isacommand{by}\isamarkupfalse%
\ {\isacharparenleft}{\kern0pt}fastforce\ simp{\isacharcolon}{\kern0pt}\ ind{\isacharunderscore}{\kern0pt}happ{\isacharunderscore}{\kern0pt}seq{\isacharunderscore}{\kern0pt}def{\isacharparenright}{\kern0pt}\isanewline
\ \ \isacommand{then}\isamarkupfalse%
\ \isacommand{show}\isamarkupfalse%
\ {\isacharquery}{\kern0pt}thesis\isanewline
\ \ \ \ \isacommand{using}\isamarkupfalse%
\ ind{\isacharunderscore}{\kern0pt}happ{\isacharunderscore}{\kern0pt}seq{\isacharunderscore}{\kern0pt}props{\isacharparenleft}{\kern0pt}{\isadigit{1}}{\isacharparenright}{\kern0pt}{\isacharbrackleft}{\kern0pt}OF\ {\isacartoucheopen}ind{\isacharunderscore}{\kern0pt}happ{\isacharunderscore}{\kern0pt}seq\ {\isasympi}s\ hs{\isacartoucheclose}{\isacharbrackright}{\kern0pt}\isanewline
\ \ \ \ \isacommand{by}\isamarkupfalse%
{\isacharparenleft}{\kern0pt}force\ simp{\isacharcolon}{\kern0pt}\ durative{\isacharunderscore}{\kern0pt}acts{\isacharunderscore}{\kern0pt}def\ is{\isacharunderscore}{\kern0pt}act{\isacharunderscore}{\kern0pt}simple{\isacharunderscore}{\kern0pt}def\ is{\isacharunderscore}{\kern0pt}htp{\isacharunderscore}{\kern0pt}def\isanewline
\ \ \ \ \ \ \ \ \ \ \ \ \ dest{\isacharbang}{\kern0pt}{\isacharcolon}{\kern0pt}\ consec{\isacharunderscore}{\kern0pt}htps{\isacharunderscore}{\kern0pt}props{\isacharparenleft}{\kern0pt}{\isadigit{2}}{\isacharparenright}{\kern0pt}\ htp{\isacharunderscore}{\kern0pt}in{\isacharunderscore}{\kern0pt}ind{\isacharunderscore}{\kern0pt}happ{\isacharunderscore}{\kern0pt}seq{\isacharbrackleft}{\kern0pt}OF\ {\isacharunderscore}{\kern0pt}\ {\isacartoucheopen}ind{\isacharunderscore}{\kern0pt}happ{\isacharunderscore}{\kern0pt}seq\ {\isasympi}s\ hs{\isacartoucheclose}{\isacharbrackright}{\kern0pt}\isanewline
\ \ \ \ \ \ \ \ \ \ \ \ \ elim{\isacharbang}{\kern0pt}{\isacharcolon}{\kern0pt}\ res{\isacharunderscore}{\kern0pt}inst{\isacharunderscore}{\kern0pt}snap{\isacharunderscore}{\kern0pt}action{\isachardot}{\kern0pt}elims{\isacharparenright}{\kern0pt}\isanewline
\isacommand{qed}\isamarkupfalse%
%
\endisatagproof
{\isafoldproof}%
%
\isadelimproof
\isanewline
%
\endisadelimproof
\isanewline
\isacommand{lemma}\isamarkupfalse%
\ sorted{\isacharunderscore}{\kern0pt}last{\isacharcolon}{\kern0pt}\ {\isachardoublequoteopen}{\isasymlbrakk}strict{\isacharunderscore}{\kern0pt}sorted\ xs{\isacharsemicolon}{\kern0pt}\ x\ {\isasymin}\ set\ xs{\isasymrbrakk}\ {\isasymLongrightarrow}\ x\ {\isasymle}\ last\ xs{\isachardoublequoteclose}\isanewline
%
\isadelimproof
\ \ %
\endisadelimproof
%
\isatagproof
\isacommand{by}\isamarkupfalse%
\ {\isacharparenleft}{\kern0pt}induction\ xs{\isacharparenright}{\kern0pt}\ {\isacharparenleft}{\kern0pt}auto\ simp{\isacharcolon}{\kern0pt}\ dest{\isacharbang}{\kern0pt}{\isacharcolon}{\kern0pt}\ last{\isacharunderscore}{\kern0pt}in{\isacharunderscore}{\kern0pt}set{\isacharparenright}{\kern0pt}%
\endisatagproof
{\isafoldproof}%
%
\isadelimproof
\isanewline
%
\endisadelimproof
\isanewline
\isacommand{lemma}\isamarkupfalse%
\ ind{\isacharunderscore}{\kern0pt}happ{\isacharunderscore}{\kern0pt}seq{\isacharunderscore}{\kern0pt}last{\isacharunderscore}{\kern0pt}htp{\isacharcolon}{\kern0pt}\isanewline
\ \ \isakeyword{assumes}\ {\isachardoublequoteopen}ind{\isacharunderscore}{\kern0pt}happ{\isacharunderscore}{\kern0pt}seq\ {\isasympi}s\ hs{\isachardoublequoteclose}\ {\isachardoublequoteopen}hs\ {\isasymnoteq}\ {\isacharbrackleft}{\kern0pt}{\isacharbrackright}{\kern0pt}{\isachardoublequoteclose}\isanewline
\ \ \isakeyword{shows}\ {\isachardoublequoteopen}is{\isacharunderscore}{\kern0pt}htp\ {\isasympi}s\ {\isacharparenleft}{\kern0pt}fst\ {\isacharparenleft}{\kern0pt}last\ hs{\isacharparenright}{\kern0pt}{\isacharparenright}{\kern0pt}{\isachardoublequoteclose}\isanewline
%
\isadelimproof
%
\endisadelimproof
%
\isatagproof
\isacommand{proof}\isamarkupfalse%
{\isacharminus}{\kern0pt}\isanewline
\ \ \isacommand{obtain}\isamarkupfalse%
\ h\ hs{\isacharprime}{\kern0pt}\ \isakeyword{where}\ hs{\isacharbrackleft}{\kern0pt}simp{\isacharbrackright}{\kern0pt}{\isacharcolon}{\kern0pt}\ {\isachardoublequoteopen}hs\ {\isacharequal}{\kern0pt}\ hs{\isacharprime}{\kern0pt}\ {\isacharat}{\kern0pt}\ {\isacharbrackleft}{\kern0pt}h{\isacharbrackright}{\kern0pt}{\isachardoublequoteclose}\isanewline
\ \ \ \ \isacommand{using}\isamarkupfalse%
\ {\isacartoucheopen}hs\ {\isasymnoteq}\ {\isacharbrackleft}{\kern0pt}{\isacharbrackright}{\kern0pt}{\isacartoucheclose}\isanewline
\ \ \ \ \isacommand{by}\isamarkupfalse%
\ {\isacharparenleft}{\kern0pt}auto\ simp{\isacharcolon}{\kern0pt}\ neq{\isacharunderscore}{\kern0pt}Nil{\isacharunderscore}{\kern0pt}rev{\isacharunderscore}{\kern0pt}conv{\isacharparenright}{\kern0pt}\isanewline
\ \ \isacommand{then}\isamarkupfalse%
\ \isacommand{obtain}\isamarkupfalse%
\ t\ A\ \isakeyword{where}\ h{\isacharbrackleft}{\kern0pt}simp{\isacharbrackright}{\kern0pt}{\isacharcolon}{\kern0pt}\ {\isachardoublequoteopen}h\ {\isacharequal}{\kern0pt}\ {\isacharparenleft}{\kern0pt}t{\isacharcomma}{\kern0pt}\ A{\isacharparenright}{\kern0pt}{\isachardoublequoteclose}\isanewline
\ \ \ \ \isacommand{by}\isamarkupfalse%
\ {\isacharparenleft}{\kern0pt}elim\ prod{\isachardot}{\kern0pt}exhaust{\isacharparenright}{\kern0pt}\isanewline
\ \ \isacommand{then}\isamarkupfalse%
\ \isacommand{obtain}\isamarkupfalse%
\ a\ \isakeyword{where}\ {\isachardoublequoteopen}a\ {\isasymin}\ set\ A{\isachardoublequoteclose}\isanewline
\ \ \ \ \isacommand{using}\isamarkupfalse%
\ ind{\isacharunderscore}{\kern0pt}happ{\isacharunderscore}{\kern0pt}seq{\isacharunderscore}{\kern0pt}props{\isacharparenleft}{\kern0pt}{\isadigit{6}}{\isacharparenright}{\kern0pt}{\isacharbrackleft}{\kern0pt}OF\ {\isacartoucheopen}ind{\isacharunderscore}{\kern0pt}happ{\isacharunderscore}{\kern0pt}seq\ {\isasympi}s\ hs{\isacartoucheclose}{\isacharbrackright}{\kern0pt}\isanewline
\ \ \ \ \isacommand{by}\isamarkupfalse%
\ auto\isanewline
\ \ \isacommand{then}\isamarkupfalse%
\ \isacommand{obtain}\isamarkupfalse%
\ t\isactrlsub {\isasympi}\ {\isasympi}\ \isakeyword{where}\ {\isachardoublequoteopen}{\isacharparenleft}{\kern0pt}t\isactrlsub {\isasympi}{\isacharcomma}{\kern0pt}{\isasympi}{\isacharparenright}{\kern0pt}\ {\isasymin}\ set\ {\isasympi}s{\isachardoublequoteclose}\ \isakeyword{and}\ {\isachardoublequoteopen}inst{\isacharunderscore}{\kern0pt}of{\isacharunderscore}{\kern0pt}plan{\isacharunderscore}{\kern0pt}action\ {\isasympi}s\ {\isacharparenleft}{\kern0pt}t\isactrlsub {\isasympi}{\isacharcomma}{\kern0pt}{\isasympi}{\isacharparenright}{\kern0pt}\ {\isacharparenleft}{\kern0pt}t{\isacharcomma}{\kern0pt}a{\isacharparenright}{\kern0pt}{\isachardoublequoteclose}\isanewline
\ \ \ \ \isacommand{using}\isamarkupfalse%
\ {\isacartoucheopen}ind{\isacharunderscore}{\kern0pt}happ{\isacharunderscore}{\kern0pt}seq\ {\isasympi}s\ hs{\isacartoucheclose}\ {\isacartoucheopen}a\ {\isasymin}\ set\ A{\isacartoucheclose}\isanewline
\ \ \ \ \isacommand{by}\isamarkupfalse%
\ {\isacharparenleft}{\kern0pt}fastforce\ simp{\isacharcolon}{\kern0pt}\ ind{\isacharunderscore}{\kern0pt}happ{\isacharunderscore}{\kern0pt}seq{\isacharunderscore}{\kern0pt}def{\isacharparenright}{\kern0pt}\isanewline
\ \ \isacommand{then}\isamarkupfalse%
\ \isacommand{show}\isamarkupfalse%
\ {\isacharquery}{\kern0pt}thesis\isanewline
\ \ \ \ \isacommand{using}\isamarkupfalse%
\ ind{\isacharunderscore}{\kern0pt}happ{\isacharunderscore}{\kern0pt}seq{\isacharunderscore}{\kern0pt}props{\isacharparenleft}{\kern0pt}{\isadigit{1}}{\isacharparenright}{\kern0pt}{\isacharbrackleft}{\kern0pt}OF\ {\isacartoucheopen}ind{\isacharunderscore}{\kern0pt}happ{\isacharunderscore}{\kern0pt}seq\ {\isasympi}s\ hs{\isacartoucheclose}{\isacharbrackright}{\kern0pt}\ sorted{\isacharunderscore}{\kern0pt}last\isanewline
\ \ \ \ \isacommand{by}\isamarkupfalse%
\ {\isacharparenleft}{\kern0pt}fastforce\ simp{\isacharcolon}{\kern0pt}\ durative{\isacharunderscore}{\kern0pt}acts{\isacharunderscore}{\kern0pt}def\ is{\isacharunderscore}{\kern0pt}act{\isacharunderscore}{\kern0pt}simple{\isacharunderscore}{\kern0pt}def\ is{\isacharunderscore}{\kern0pt}htp{\isacharunderscore}{\kern0pt}def\isanewline
\ \ \ \ \ \ \ \ dest{\isacharbang}{\kern0pt}{\isacharcolon}{\kern0pt}\ consec{\isacharunderscore}{\kern0pt}htps{\isacharunderscore}{\kern0pt}props{\isacharparenleft}{\kern0pt}{\isadigit{3}}{\isacharparenright}{\kern0pt}\ htp{\isacharunderscore}{\kern0pt}in{\isacharunderscore}{\kern0pt}ind{\isacharunderscore}{\kern0pt}happ{\isacharunderscore}{\kern0pt}seq{\isacharbrackleft}{\kern0pt}OF\ {\isacharunderscore}{\kern0pt}\ {\isacartoucheopen}ind{\isacharunderscore}{\kern0pt}happ{\isacharunderscore}{\kern0pt}seq\ {\isasympi}s\ hs{\isacartoucheclose}{\isacharbrackright}{\kern0pt}\isanewline
\ \ \ \ \ \ \ \ elim{\isacharbang}{\kern0pt}{\isacharcolon}{\kern0pt}\ res{\isacharunderscore}{\kern0pt}inst{\isacharunderscore}{\kern0pt}snap{\isacharunderscore}{\kern0pt}action{\isachardot}{\kern0pt}elims{\isacharparenright}{\kern0pt}\isanewline
\isacommand{qed}\isamarkupfalse%
%
\endisatagproof
{\isafoldproof}%
%
\isadelimproof
\isanewline
%
\endisadelimproof
\isanewline
\isacommand{definition}\isamarkupfalse%
\ {\isachardoublequoteopen}htps{\isacharprime}{\kern0pt}\ {\isasympi}s\ {\isasymequiv}\ {\isacharbraceleft}{\kern0pt}t\ {\isachardot}{\kern0pt}\ is{\isacharunderscore}{\kern0pt}htp\ {\isasympi}s\ t{\isacharbraceright}{\kern0pt}{\isachardoublequoteclose}\isanewline
\isanewline
\isacommand{lemma}\isamarkupfalse%
\ htps{\isacharunderscore}{\kern0pt}props{\isacharcolon}{\kern0pt}\ {\isachardoublequoteopen}t\ {\isasymin}\ htps{\isacharprime}{\kern0pt}\ {\isasympi}s\ {\isasymLongrightarrow}\ is{\isacharunderscore}{\kern0pt}htp\ {\isasympi}s\ t{\isachardoublequoteclose}\isanewline
%
\isadelimproof
\ \ %
\endisadelimproof
%
\isatagproof
\isacommand{by}\isamarkupfalse%
{\isacharparenleft}{\kern0pt}auto\ simp{\isacharcolon}{\kern0pt}\ htps{\isacharprime}{\kern0pt}{\isacharunderscore}{\kern0pt}def{\isacharparenright}{\kern0pt}%
\endisatagproof
{\isafoldproof}%
%
\isadelimproof
\isanewline
%
\endisadelimproof
\isanewline
\isacommand{lemma}\isamarkupfalse%
\ ind{\isacharunderscore}{\kern0pt}happ{\isacharunderscore}{\kern0pt}seq{\isacharunderscore}{\kern0pt}unique{\isacharunderscore}{\kern0pt}final{\isacharunderscore}{\kern0pt}state{\isacharprime}{\kern0pt}{\isacharcolon}{\kern0pt}\isanewline
\ \ \isakeyword{assumes}{\isachardoublequoteopen}ind{\isacharunderscore}{\kern0pt}happ{\isacharunderscore}{\kern0pt}seq\ {\isasympi}s\ hs{\isachardoublequoteclose}\ \isakeyword{and}\ {\isachardoublequoteopen}ind{\isacharunderscore}{\kern0pt}happ{\isacharunderscore}{\kern0pt}seq\ {\isasympi}s\ hs{\isacharprime}{\kern0pt}{\isachardoublequoteclose}\isanewline
\ \ \isakeyword{shows}\ {\isachardoublequoteopen}{\isasymlbrakk}wf{\isacharunderscore}{\kern0pt}plan\ {\isasympi}s{\isacharsemicolon}{\kern0pt}\ valid{\isacharunderscore}{\kern0pt}happ{\isacharunderscore}{\kern0pt}seq\ M\isactrlsub {\isadigit{0}}\ hs\ M\isactrlsub n{\isasymrbrakk}\ {\isasymLongrightarrow}\ valid{\isacharunderscore}{\kern0pt}happ{\isacharunderscore}{\kern0pt}seq\ M\isactrlsub {\isadigit{0}}\ hs{\isacharprime}{\kern0pt}\ M\isactrlsub n{\isachardoublequoteclose}\isanewline
%
\isadelimproof
\ \ %
\endisadelimproof
%
\isatagproof
\isacommand{using}\isamarkupfalse%
\ not{\isacharunderscore}{\kern0pt}htp{\isacharunderscore}{\kern0pt}empty{\isacharunderscore}{\kern0pt}effect\ assms\isanewline
\ \ \isacommand{by}\isamarkupfalse%
\ {\isacharparenleft}{\kern0pt}auto\ intro{\isacharbang}{\kern0pt}{\isacharcolon}{\kern0pt}\ ind{\isacharunderscore}{\kern0pt}happ{\isacharunderscore}{\kern0pt}seq{\isacharunderscore}{\kern0pt}last{\isacharunderscore}{\kern0pt}htp\ htp{\isacharunderscore}{\kern0pt}in{\isacharunderscore}{\kern0pt}ind{\isacharunderscore}{\kern0pt}happ{\isacharunderscore}{\kern0pt}seq\ ind{\isacharunderscore}{\kern0pt}happ{\isacharunderscore}{\kern0pt}seq{\isacharunderscore}{\kern0pt}hd{\isacharunderscore}{\kern0pt}htp\ ind{\isacharunderscore}{\kern0pt}happ{\isacharunderscore}{\kern0pt}seq{\isacharunderscore}{\kern0pt}props{\isacharparenleft}{\kern0pt}{\isadigit{1}}{\isacharparenright}{\kern0pt}\isanewline
\ \ \ \ \ \ \ \ \ \ \ \ \ \ \ \ \ \ \ valid{\isacharunderscore}{\kern0pt}hap{\isacharunderscore}{\kern0pt}seq{\isacharunderscore}{\kern0pt}construct{\isacharbrackleft}{\kern0pt}\isakeyword{where}\ htps\ {\isacharequal}{\kern0pt}\ {\isachardoublequoteopen}htps{\isacharprime}{\kern0pt}\ {\isasympi}s{\isachardoublequoteclose}\ \isakeyword{and}\ hs\ {\isacharequal}{\kern0pt}\ {\isachardoublequoteopen}hs{\isachardoublequoteclose}{\isacharbrackright}{\kern0pt}\isanewline
\ \ \ \ \ \ \ \ \ \ \ intro{\isacharcolon}{\kern0pt}\ happs{\isacharunderscore}{\kern0pt}same{\isacharunderscore}{\kern0pt}between{\isacharunderscore}{\kern0pt}htps\isanewline
\ \ \ \ \ \ \ \ \ \ \ simp{\isacharcolon}{\kern0pt}\ htps{\isacharprime}{\kern0pt}{\isacharunderscore}{\kern0pt}def\ ind{\isacharunderscore}{\kern0pt}happ{\isacharunderscore}{\kern0pt}seq{\isacharunderscore}{\kern0pt}happ{\isacharunderscore}{\kern0pt}at{\isacharunderscore}{\kern0pt}htps{\isacharprime}{\kern0pt}{\isacharbrackleft}{\kern0pt}OF\ htps{\isacharunderscore}{\kern0pt}props\ {\isacartoucheopen}ind{\isacharunderscore}{\kern0pt}happ{\isacharunderscore}{\kern0pt}seq\ {\isasympi}s\ hs{\isacartoucheclose}\ {\isacartoucheopen}ind{\isacharunderscore}{\kern0pt}happ{\isacharunderscore}{\kern0pt}seq\ {\isasympi}s\ hs{\isacharprime}{\kern0pt}{\isacartoucheclose}{\isacharbrackright}{\kern0pt}{\isacharparenright}{\kern0pt}%
\endisatagproof
{\isafoldproof}%
%
\isadelimproof
\isanewline
%
\endisadelimproof
\isanewline
\isacommand{end}\isamarkupfalse%
\isanewline
%
\isadelimtheory
\isanewline
%
\endisadelimtheory
%
\isatagtheory
\isacommand{end}\isamarkupfalse%
\ %
\isamarkupcmt{Theory%
}%
\endisatagtheory
{\isafoldtheory}%
%
\isadelimtheory
%
\endisadelimtheory
%
\end{isabellebody}%
\endinput
%:%file=TEMPORAL_PDDL_Semantics.tex%:%
%:%11=1%:%
%:%27=2%:%
%:%28=2%:%
%:%29=3%:%
%:%30=4%:%
%:%31=5%:%
%:%32=6%:%
%:%33=7%:%
%:%34=8%:%
%:%35=9%:%
%:%36=10%:%
%:%37=11%:%
%:%42=11%:%
%:%45=12%:%
%:%46=12%:%
%:%53=14%:%
%:%63=16%:%
%:%64=16%:%
%:%65=17%:%
%:%66=18%:%
%:%67=18%:%
%:%68=19%:%
%:%69=20%:%
%:%70=20%:%
%:%71=21%:%
%:%72=22%:%
%:%75=23%:%
%:%79=23%:%
%:%80=23%:%
%:%81=24%:%
%:%82=24%:%
%:%83=25%:%
%:%84=25%:%
%:%89=25%:%
%:%92=26%:%
%:%93=27%:%
%:%94=27%:%
%:%95=28%:%
%:%98=29%:%
%:%102=29%:%
%:%103=29%:%
%:%104=30%:%
%:%105=30%:%
%:%110=30%:%
%:%113=31%:%
%:%114=32%:%
%:%115=33%:%
%:%116=33%:%
%:%117=34%:%
%:%118=35%:%
%:%119=36%:%
%:%120=37%:%
%:%121=38%:%
%:%124=39%:%
%:%128=39%:%
%:%129=39%:%
%:%130=40%:%
%:%131=40%:%
%:%136=40%:%
%:%139=41%:%
%:%140=42%:%
%:%141=42%:%
%:%144=43%:%
%:%148=43%:%
%:%149=43%:%
%:%163=46%:%
%:%167=48%:%
%:%179=49%:%
%:%181=51%:%
%:%182=51%:%
%:%183=52%:%
%:%184=53%:%
%:%185=53%:%
%:%186=54%:%
%:%187=55%:%
%:%188=55%:%
%:%189=56%:%
%:%190=57%:%
%:%191=57%:%
%:%193=59%:%
%:%194=60%:%
%:%195=61%:%
%:%199=64%:%
%:%201=65%:%
%:%202=65%:%
%:%203=66%:%
%:%205=69%:%
%:%206=70%:%
%:%208=71%:%
%:%209=71%:%
%:%211=74%:%
%:%212=75%:%
%:%214=76%:%
%:%215=76%:%
%:%216=77%:%
%:%218=79%:%
%:%220=80%:%
%:%221=80%:%
%:%223=82%:%
%:%225=83%:%
%:%226=83%:%
%:%227=84%:%
%:%228=85%:%
%:%229=85%:%
%:%230=86%:%
%:%231=87%:%
%:%232=87%:%
%:%233=87%:%
%:%241=90%:%
%:%253=92%:%
%:%255=93%:%
%:%256=93%:%
%:%257=94%:%
%:%258=95%:%
%:%259=95%:%
%:%260=96%:%
%:%261=97%:%
%:%262=97%:%
%:%263=98%:%
%:%264=99%:%
%:%265=100%:%
%:%267=102%:%
%:%268=103%:%
%:%269=104%:%
%:%271=105%:%
%:%272=105%:%
%:%273=106%:%
%:%274=107%:%
%:%275=108%:%
%:%276=109%:%
%:%277=110%:%
%:%278=111%:%
%:%279=112%:%
%:%280=113%:%
%:%281=114%:%
%:%282=115%:%
%:%283=116%:%
%:%285=118%:%
%:%286=119%:%
%:%288=120%:%
%:%289=120%:%
%:%290=121%:%
%:%291=122%:%
%:%292=122%:%
%:%294=124%:%
%:%295=125%:%
%:%297=126%:%
%:%298=126%:%
%:%299=127%:%
%:%300=127%:%
%:%301=127%:%
%:%302=128%:%
%:%303=129%:%
%:%304=130%:%
%:%305=131%:%
%:%312=133%:%
%:%324=135%:%
%:%326=136%:%
%:%327=136%:%
%:%329=138%:%
%:%330=139%:%
%:%332=140%:%
%:%333=140%:%
%:%334=141%:%
%:%335=142%:%
%:%336=143%:%
%:%337=144%:%
%:%344=146%:%
%:%354=148%:%
%:%355=148%:%
%:%356=149%:%
%:%357=150%:%
%:%358=151%:%
%:%359=152%:%
%:%360=152%:%
%:%367=154%:%
%:%379=155%:%
%:%380=156%:%
%:%381=157%:%
%:%383=158%:%
%:%384=158%:%
%:%385=159%:%
%:%386=160%:%
%:%388=162%:%
%:%389=163%:%
%:%391=164%:%
%:%392=164%:%
%:%399=166%:%
%:%411=167%:%
%:%413=168%:%
%:%414=168%:%
%:%415=169%:%
%:%416=170%:%
%:%417=171%:%
%:%418=171%:%
%:%419=172%:%
%:%421=174%:%
%:%423=175%:%
%:%424=175%:%
%:%426=177%:%
%:%428=178%:%
%:%429=178%:%
%:%431=180%:%
%:%433=181%:%
%:%434=181%:%
%:%435=182%:%
%:%437=184%:%
%:%438=185%:%
%:%440=186%:%
%:%441=186%:%
%:%442=187%:%
%:%444=189%:%
%:%445=190%:%
%:%446=191%:%
%:%447=191%:%
%:%448=192%:%
%:%449=192%:%
%:%452=193%:%
%:%456=193%:%
%:%457=193%:%
%:%458=194%:%
%:%459=194%:%
%:%460=195%:%
%:%461=195%:%
%:%462=196%:%
%:%463=196%:%
%:%468=196%:%
%:%471=197%:%
%:%472=198%:%
%:%473=198%:%
%:%474=199%:%
%:%475=200%:%
%:%476=201%:%
%:%477=201%:%
%:%478=202%:%
%:%479=203%:%
%:%482=204%:%
%:%486=204%:%
%:%487=204%:%
%:%492=204%:%
%:%495=205%:%
%:%496=206%:%
%:%497=206%:%
%:%498=207%:%
%:%503=212%:%
%:%506=213%:%
%:%510=213%:%
%:%511=213%:%
%:%516=213%:%
%:%519=214%:%
%:%520=215%:%
%:%521=215%:%
%:%522=216%:%
%:%523=217%:%
%:%524=218%:%
%:%525=219%:%
%:%528=220%:%
%:%532=220%:%
%:%533=220%:%
%:%534=220%:%
%:%535=221%:%
%:%536=221%:%
%:%537=222%:%
%:%538=222%:%
%:%539=223%:%
%:%540=223%:%
%:%541=224%:%
%:%542=224%:%
%:%543=225%:%
%:%549=225%:%
%:%552=226%:%
%:%553=227%:%
%:%554=227%:%
%:%555=228%:%
%:%556=229%:%
%:%559=230%:%
%:%563=230%:%
%:%564=230%:%
%:%565=230%:%
%:%566=231%:%
%:%567=231%:%
%:%572=231%:%
%:%575=232%:%
%:%576=233%:%
%:%577=233%:%
%:%580=234%:%
%:%584=234%:%
%:%585=234%:%
%:%586=235%:%
%:%587=235%:%
%:%588=236%:%
%:%589=236%:%
%:%594=236%:%
%:%597=237%:%
%:%598=238%:%
%:%599=238%:%
%:%600=239%:%
%:%603=240%:%
%:%607=240%:%
%:%608=240%:%
%:%609=240%:%
%:%610=240%:%
%:%619=242%:%
%:%620=243%:%
%:%621=244%:%
%:%622=245%:%
%:%623=246%:%
%:%624=247%:%
%:%626=249%:%
%:%627=249%:%
%:%628=250%:%
%:%629=251%:%
%:%632=252%:%
%:%636=252%:%
%:%637=252%:%
%:%638=252%:%
%:%652=255%:%
%:%664=256%:%
%:%665=257%:%
%:%667=259%:%
%:%668=259%:%
%:%669=260%:%
%:%670=261%:%
%:%671=262%:%
%:%672=263%:%
%:%673=264%:%
%:%674=265%:%
%:%675=266%:%
%:%676=267%:%
%:%677=267%:%
%:%680=268%:%
%:%684=268%:%
%:%685=268%:%
%:%686=269%:%
%:%687=269%:%
%:%692=269%:%
%:%695=270%:%
%:%696=271%:%
%:%697=271%:%
%:%700=272%:%
%:%704=272%:%
%:%705=272%:%
%:%706=273%:%
%:%707=273%:%
%:%708=274%:%
%:%709=274%:%
%:%710=275%:%
%:%711=275%:%
%:%712=275%:%
%:%713=276%:%
%:%714=276%:%
%:%715=276%:%
%:%720=276%:%
%:%723=277%:%
%:%724=278%:%
%:%725=279%:%
%:%726=279%:%
%:%727=280%:%
%:%728=281%:%
%:%729=282%:%
%:%736=283%:%
%:%737=283%:%
%:%738=284%:%
%:%739=284%:%
%:%740=285%:%
%:%741=285%:%
%:%742=285%:%
%:%743=286%:%
%:%744=286%:%
%:%745=287%:%
%:%746=287%:%
%:%747=288%:%
%:%748=288%:%
%:%749=289%:%
%:%750=289%:%
%:%751=290%:%
%:%752=290%:%
%:%753=291%:%
%:%754=291%:%
%:%755=292%:%
%:%756=292%:%
%:%757=292%:%
%:%758=293%:%
%:%759=293%:%
%:%760=293%:%
%:%761=294%:%
%:%762=294%:%
%:%763=295%:%
%:%764=295%:%
%:%765=296%:%
%:%775=298%:%
%:%776=299%:%
%:%778=300%:%
%:%779=300%:%
%:%780=301%:%
%:%783=302%:%
%:%787=302%:%
%:%788=302%:%
%:%802=304%:%
%:%814=306%:%
%:%815=307%:%
%:%817=308%:%
%:%818=308%:%
%:%819=309%:%
%:%820=310%:%
%:%822=312%:%
%:%824=313%:%
%:%825=313%:%
%:%826=314%:%
%:%834=322%:%
%:%835=323%:%
%:%836=324%:%
%:%837=324%:%
%:%838=325%:%
%:%841=326%:%
%:%845=326%:%
%:%846=326%:%
%:%847=326%:%
%:%856=328%:%
%:%857=329%:%
%:%858=330%:%
%:%860=331%:%
%:%861=331%:%
%:%862=332%:%
%:%865=335%:%
%:%866=336%:%
%:%867=337%:%
%:%868=338%:%
%:%869=339%:%
%:%870=340%:%
%:%871=341%:%
%:%873=343%:%
%:%874=343%:%
%:%875=344%:%
%:%876=345%:%
%:%879=348%:%
%:%880=349%:%
%:%881=350%:%
%:%882=350%:%
%:%883=351%:%
%:%884=352%:%
%:%885=353%:%
%:%888=354%:%
%:%892=354%:%
%:%893=354%:%
%:%894=354%:%
%:%899=354%:%
%:%902=355%:%
%:%903=356%:%
%:%904=356%:%
%:%905=356%:%
%:%913=360%:%
%:%923=375%:%
%:%924=375%:%
%:%925=376%:%
%:%926=377%:%
%:%927=378%:%
%:%928=379%:%
%:%929=380%:%
%:%930=380%:%
%:%931=381%:%
%:%934=382%:%
%:%938=382%:%
%:%939=382%:%
%:%940=383%:%
%:%941=383%:%
%:%942=384%:%
%:%943=384%:%
%:%944=385%:%
%:%945=385%:%
%:%946=386%:%
%:%947=386%:%
%:%948=387%:%
%:%954=387%:%
%:%957=388%:%
%:%958=389%:%
%:%959=390%:%
%:%960=390%:%
%:%962=392%:%
%:%963=393%:%
%:%965=394%:%
%:%966=394%:%
%:%967=395%:%
%:%968=396%:%
%:%969=397%:%
%:%970=397%:%
%:%971=398%:%
%:%973=400%:%
%:%974=401%:%
%:%975=402%:%
%:%976=403%:%
%:%977=404%:%
%:%978=405%:%
%:%980=408%:%
%:%981=408%:%
%:%982=409%:%
%:%983=410%:%
%:%984=411%:%
%:%985=411%:%
%:%986=412%:%
%:%987=413%:%
%:%988=413%:%
%:%989=414%:%
%:%991=415%:%
%:%992=416%:%
%:%996=419%:%
%:%997=420%:%
%:%998=421%:%
%:%1000=422%:%
%:%1001=422%:%
%:%1002=423%:%
%:%1003=423%:%
%:%1004=423%:%
%:%1005=424%:%
%:%1007=426%:%
%:%1009=427%:%
%:%1010=427%:%
%:%1011=428%:%
%:%1014=431%:%
%:%1015=432%:%
%:%1016=433%:%
%:%1017=433%:%
%:%1018=434%:%
%:%1021=437%:%
%:%1022=438%:%
%:%1023=439%:%
%:%1024=439%:%
%:%1025=440%:%
%:%1030=445%:%
%:%1031=446%:%
%:%1032=447%:%
%:%1033=448%:%
%:%1034=449%:%
%:%1036=451%:%
%:%1037=451%:%
%:%1038=452%:%
%:%1039=453%:%
%:%1041=455%:%
%:%1042=456%:%
%:%1044=458%:%
%:%1045=458%:%
%:%1046=459%:%
%:%1047=460%:%
%:%1048=461%:%
%:%1049=462%:%
%:%1050=463%:%
%:%1051=464%:%
%:%1052=465%:%
%:%1053=466%:%
%:%1054=466%:%
%:%1057=467%:%
%:%1061=467%:%
%:%1062=467%:%
%:%1071=472%:%
%:%1073=473%:%
%:%1074=473%:%
%:%1075=474%:%
%:%1076=475%:%
%:%1077=476%:%
%:%1078=477%:%
%:%1079=477%:%
%:%1082=478%:%
%:%1086=478%:%
%:%1087=478%:%
%:%1096=480%:%
%:%1097=481%:%
%:%1099=482%:%
%:%1100=482%:%
%:%1101=483%:%
%:%1102=484%:%
%:%1103=485%:%
%:%1104=485%:%
%:%1105=486%:%
%:%1106=487%:%
%:%1107=488%:%
%:%1110=491%:%
%:%1111=492%:%
%:%1112=493%:%
%:%1113=493%:%
%:%1114=493%:%
%:%1115=493%:%
%:%1116=494%:%
%:%1117=495%:%
%:%1118=496%:%
%:%1119=496%:%
%:%1120=497%:%
%:%1122=499%:%
%:%1123=500%:%
%:%1125=502%:%
%:%1126=502%:%
%:%1127=503%:%
%:%1131=507%:%
%:%1132=508%:%
%:%1139=515%:%
%:%1140=516%:%
%:%1142=517%:%
%:%1143=517%:%
%:%1144=518%:%
%:%1146=520%:%
%:%1148=521%:%
%:%1149=521%:%
%:%1150=522%:%
%:%1152=524%:%
%:%1154=525%:%
%:%1155=525%:%
%:%1156=526%:%
%:%1158=528%:%
%:%1160=529%:%
%:%1161=529%:%
%:%1163=531%:%
%:%1167=532%:%
%:%1169=533%:%
%:%1171=534%:%
%:%1173=535%:%
%:%1175=536%:%
%:%1177=537%:%
%:%1180=539%:%
%:%1181=539%:%
%:%1182=540%:%
%:%1192=550%:%
%:%1193=551%:%
%:%1194=552%:%
%:%1195=552%:%
%:%1196=552%:%
%:%1199=554%:%
%:%1200=555%:%
%:%1202=556%:%
%:%1203=556%:%
%:%1204=557%:%
%:%1205=558%:%
%:%1207=559%:%
%:%1209=560%:%
%:%1210=560%:%
%:%1211=561%:%
%:%1212=562%:%
%:%1213=562%:%
%:%1214=563%:%
%:%1215=564%:%
%:%1216=565%:%
%:%1217=565%:%
%:%1220=566%:%
%:%1224=566%:%
%:%1225=566%:%
%:%1226=567%:%
%:%1227=567%:%
%:%1228=568%:%
%:%1234=568%:%
%:%1237=569%:%
%:%1238=570%:%
%:%1239=570%:%
%:%1240=571%:%
%:%1241=572%:%
%:%1242=573%:%
%:%1243=573%:%
%:%1244=574%:%
%:%1245=575%:%
%:%1246=576%:%
%:%1247=577%:%
%:%1248=577%:%
%:%1251=578%:%
%:%1255=578%:%
%:%1256=578%:%
%:%1257=579%:%
%:%1258=579%:%
%:%1259=580%:%
%:%1260=580%:%
%:%1261=581%:%
%:%1271=583%:%
%:%1272=584%:%
%:%1274=586%:%
%:%1275=586%:%
%:%1276=587%:%
%:%1277=588%:%
%:%1278=589%:%
%:%1279=590%:%
%:%1280=590%:%
%:%1281=591%:%
%:%1287=597%:%
%:%1288=598%:%
%:%1289=599%:%
%:%1290=599%:%
%:%1291=600%:%
%:%1292=601%:%
%:%1293=602%:%
%:%1294=602%:%
%:%1297=603%:%
%:%1301=603%:%
%:%1302=603%:%
%:%1307=603%:%
%:%1310=604%:%
%:%1311=605%:%
%:%1312=605%:%
%:%1313=605%:%
%:%1316=607%:%
%:%1318=608%:%
%:%1319=608%:%
%:%1320=609%:%
%:%1322=611%:%
%:%1324=612%:%
%:%1325=612%:%
%:%1326=613%:%
%:%1327=614%:%
%:%1328=615%:%
%:%1329=615%:%
%:%1332=616%:%
%:%1336=616%:%
%:%1337=616%:%
%:%1338=617%:%
%:%1339=617%:%
%:%1340=618%:%
%:%1341=618%:%
%:%1342=618%:%
%:%1347=618%:%
%:%1350=619%:%
%:%1351=620%:%
%:%1352=620%:%
%:%1353=620%:%
%:%1361=622%:%
%:%1371=626%:%
%:%1372=626%:%
%:%1373=627%:%
%:%1374=628%:%
%:%1375=628%:%
%:%1376=629%:%
%:%1377=630%:%
%:%1378=631%:%
%:%1379=631%:%
%:%1380=632%:%
%:%1381=633%:%
%:%1383=635%:%
%:%1384=636%:%
%:%1385=637%:%
%:%1386=638%:%
%:%1388=640%:%
%:%1389=640%:%
%:%1390=641%:%
%:%1399=650%:%
%:%1400=651%:%
%:%1402=652%:%
%:%1403=652%:%
%:%1404=653%:%
%:%1405=654%:%
%:%1406=655%:%
%:%1407=655%:%
%:%1410=656%:%
%:%1414=656%:%
%:%1415=656%:%
%:%1416=657%:%
%:%1417=657%:%
%:%1426=659%:%
%:%1427=660%:%
%:%1429=661%:%
%:%1430=661%:%
%:%1431=662%:%
%:%1433=664%:%
%:%1434=665%:%
%:%1435=666%:%
%:%1437=667%:%
%:%1438=667%:%
%:%1439=668%:%
%:%1446=675%:%
%:%1447=676%:%
%:%1448=677%:%
%:%1449=677%:%
%:%1450=677%:%
%:%1451=677%:%
%:%1452=678%:%
%:%1453=679%:%
%:%1454=679%:%
%:%1456=680%:%
%:%1458=683%:%
%:%1459=683%:%
%:%1460=684%:%
%:%1461=685%:%
%:%1464=686%:%
%:%1468=686%:%
%:%1469=686%:%
%:%1470=687%:%
%:%1471=687%:%
%:%1472=687%:%
%:%1473=688%:%
%:%1474=689%:%
%:%1475=690%:%
%:%1476=690%:%
%:%1477=690%:%
%:%1478=691%:%
%:%1479=692%:%
%:%1480=692%:%
%:%1481=693%:%
%:%1482=693%:%
%:%1483=693%:%
%:%1484=694%:%
%:%1485=694%:%
%:%1486=695%:%
%:%1492=695%:%
%:%1495=696%:%
%:%1496=697%:%
%:%1497=697%:%
%:%1498=697%:%
%:%1499=697%:%
%:%1500=698%:%
%:%1501=699%:%
%:%1502=699%:%
%:%1504=701%:%
%:%1506=702%:%
%:%1507=702%:%
%:%1508=703%:%
%:%1510=706%:%
%:%1512=707%:%
%:%1513=707%:%
%:%1514=708%:%
%:%1515=709%:%
%:%1516=710%:%
%:%1517=711%:%
%:%1518=712%:%
%:%1519=713%:%
%:%1520=713%:%
%:%1521=714%:%
%:%1522=715%:%
%:%1523=716%:%
%:%1525=718%:%
%:%1527=719%:%
%:%1528=719%:%
%:%1532=723%:%
%:%1534=724%:%
%:%1535=724%:%
%:%1538=725%:%
%:%1542=725%:%
%:%1543=725%:%
%:%1544=726%:%
%:%1545=726%:%
%:%1546=726%:%
%:%1555=729%:%
%:%1556=730%:%
%:%1557=731%:%
%:%1558=732%:%
%:%1559=733%:%
%:%1560=734%:%
%:%1561=735%:%
%:%1563=737%:%
%:%1564=737%:%
%:%1567=740%:%
%:%1568=741%:%
%:%1570=742%:%
%:%1571=742%:%
%:%1572=743%:%
%:%1573=744%:%
%:%1574=745%:%
%:%1575=745%:%
%:%1576=746%:%
%:%1577=747%:%
%:%1581=751%:%
%:%1582=752%:%
%:%1591=761%:%
%:%1592=762%:%
%:%1593=763%:%
%:%1595=766%:%
%:%1596=766%:%
%:%1597=767%:%
%:%1602=772%:%
%:%1603=773%:%
%:%1611=783%:%
%:%1612=784%:%
%:%1614=785%:%
%:%1615=785%:%
%:%1616=786%:%
%:%1617=787%:%
%:%1618=788%:%
%:%1619=788%:%
%:%1620=789%:%
%:%1621=790%:%
%:%1622=791%:%
%:%1623=791%:%
%:%1624=792%:%
%:%1625=793%:%
%:%1626=794%:%
%:%1627=794%:%
%:%1630=795%:%
%:%1634=795%:%
%:%1635=795%:%
%:%1636=795%:%
%:%1641=795%:%
%:%1644=796%:%
%:%1645=797%:%
%:%1646=797%:%
%:%1647=798%:%
%:%1648=799%:%
%:%1649=800%:%
%:%1650=800%:%
%:%1653=801%:%
%:%1657=801%:%
%:%1658=801%:%
%:%1659=801%:%
%:%1668=803%:%
%:%1669=804%:%
%:%1670=805%:%
%:%1672=806%:%
%:%1673=806%:%
%:%1674=807%:%
%:%1676=809%:%
%:%1678=810%:%
%:%1679=810%:%
%:%1680=811%:%
%:%1682=813%:%
%:%1683=814%:%
%:%1685=815%:%
%:%1686=815%:%
%:%1687=816%:%
%:%1690=829%:%
%:%1691=830%:%
%:%1693=831%:%
%:%1694=831%:%
%:%1695=832%:%
%:%1700=837%:%
%:%1701=838%:%
%:%1702=839%:%
%:%1703=840%:%
%:%1703=845%:%
%:%1704=846%:%
%:%1705=847%:%
%:%1706=847%:%
%:%1707=848%:%
%:%1709=850%:%
%:%1713=851%:%
%:%1715=852%:%
%:%1719=853%:%
%:%1720=854%:%
%:%1722=855%:%
%:%1723=856%:%
%:%1725=857%:%
%:%1726=858%:%
%:%1727=859%:%
%:%1730=860%:%
%:%1731=861%:%
%:%1734=863%:%
%:%1735=863%:%
%:%1736=864%:%
%:%1750=878%:%
%:%1751=879%:%
%:%1752=880%:%
%:%1753=880%:%
%:%1754=881%:%
%:%1755=882%:%
%:%1756=883%:%
%:%1757=884%:%
%:%1758=885%:%
%:%1759=886%:%
%:%1760=887%:%
%:%1761=888%:%
%:%1762=889%:%
%:%1765=890%:%
%:%1769=890%:%
%:%1770=890%:%
%:%1771=891%:%
%:%1772=891%:%
%:%1781=893%:%
%:%1782=894%:%
%:%1784=895%:%
%:%1785=895%:%
%:%1786=896%:%
%:%1788=898%:%
%:%1789=899%:%
%:%1791=900%:%
%:%1792=900%:%
%:%1793=901%:%
%:%1795=903%:%
%:%1796=904%:%
%:%1798=905%:%
%:%1799=905%:%
%:%1800=906%:%
%:%1802=908%:%
%:%1804=909%:%
%:%1805=909%:%
%:%1808=910%:%
%:%1812=910%:%
%:%1813=910%:%
%:%1814=910%:%
%:%1819=910%:%
%:%1822=911%:%
%:%1823=912%:%
%:%1824=912%:%
%:%1825=912%:%
%:%1833=915%:%
%:%1837=917%:%
%:%1849=918%:%
%:%1850=919%:%
%:%1852=922%:%
%:%1853=922%:%
%:%1855=924%:%
%:%1856=925%:%
%:%1857=926%:%
%:%1859=928%:%
%:%1860=928%:%
%:%1861=929%:%
%:%1863=931%:%
%:%1864=932%:%
%:%1865=933%:%
%:%1866=934%:%
%:%1868=937%:%
%:%1869=937%:%
%:%1872=938%:%
%:%1876=938%:%
%:%1877=938%:%
%:%1878=938%:%
%:%1883=938%:%
%:%1886=939%:%
%:%1887=940%:%
%:%1888=940%:%
%:%1889=941%:%
%:%1892=942%:%
%:%1896=942%:%
%:%1897=942%:%
%:%1898=943%:%
%:%1899=943%:%
%:%1904=943%:%
%:%1907=944%:%
%:%1908=945%:%
%:%1909=945%:%
%:%1910=946%:%
%:%1911=947%:%
%:%1914=948%:%
%:%1918=948%:%
%:%1919=948%:%
%:%1920=949%:%
%:%1921=949%:%
%:%1922=950%:%
%:%1923=950%:%
%:%1928=950%:%
%:%1931=951%:%
%:%1932=952%:%
%:%1933=952%:%
%:%1934=953%:%
%:%1935=954%:%
%:%1936=955%:%
%:%1939=956%:%
%:%1943=956%:%
%:%1944=956%:%
%:%1945=957%:%
%:%1946=957%:%
%:%1947=958%:%
%:%1948=958%:%
%:%1949=959%:%
%:%1950=959%:%
%:%1951=960%:%
%:%1952=960%:%
%:%1953=961%:%
%:%1954=961%:%
%:%1955=961%:%
%:%1956=962%:%
%:%1957=962%:%
%:%1958=963%:%
%:%1964=963%:%
%:%1967=964%:%
%:%1968=965%:%
%:%1969=965%:%
%:%1970=966%:%
%:%1971=967%:%
%:%1972=968%:%
%:%1975=969%:%
%:%1979=969%:%
%:%1980=969%:%
%:%1981=970%:%
%:%1982=970%:%
%:%1983=971%:%
%:%1984=971%:%
%:%1985=972%:%
%:%1986=972%:%
%:%1987=973%:%
%:%1988=973%:%
%:%1989=974%:%
%:%1990=974%:%
%:%1991=974%:%
%:%1992=975%:%
%:%1993=975%:%
%:%1994=976%:%
%:%2000=976%:%
%:%2003=977%:%
%:%2004=978%:%
%:%2005=979%:%
%:%2006=979%:%
%:%2009=980%:%
%:%2013=980%:%
%:%2014=980%:%
%:%2015=981%:%
%:%2016=981%:%
%:%2017=982%:%
%:%2018=982%:%
%:%2023=982%:%
%:%2026=983%:%
%:%2027=984%:%
%:%2028=984%:%
%:%2031=985%:%
%:%2035=985%:%
%:%2036=985%:%
%:%2037=986%:%
%:%2038=986%:%
%:%2039=987%:%
%:%2040=987%:%
%:%2045=987%:%
%:%2048=988%:%
%:%2049=989%:%
%:%2050=989%:%
%:%2053=990%:%
%:%2057=990%:%
%:%2058=990%:%
%:%2059=991%:%
%:%2060=991%:%
%:%2061=992%:%
%:%2062=992%:%
%:%2067=992%:%
%:%2070=993%:%
%:%2071=994%:%
%:%2072=994%:%
%:%2073=995%:%
%:%2074=996%:%
%:%2075=997%:%
%:%2076=998%:%
%:%2077=999%:%
%:%2078=999%:%
%:%2081=1000%:%
%:%2085=1000%:%
%:%2086=1000%:%
%:%2087=1000%:%
%:%2092=1000%:%
%:%2095=1001%:%
%:%2096=1002%:%
%:%2097=1002%:%
%:%2100=1003%:%
%:%2104=1003%:%
%:%2105=1003%:%
%:%2106=1004%:%
%:%2107=1004%:%
%:%2112=1004%:%
%:%2115=1005%:%
%:%2116=1006%:%
%:%2117=1006%:%
%:%2120=1007%:%
%:%2124=1007%:%
%:%2125=1007%:%
%:%2126=1007%:%
%:%2127=1007%:%
%:%2128=1008%:%
%:%2134=1008%:%
%:%2137=1009%:%
%:%2138=1010%:%
%:%2139=1011%:%
%:%2140=1011%:%
%:%2143=1012%:%
%:%2147=1012%:%
%:%2148=1012%:%
%:%2149=1012%:%
%:%2150=1013%:%
%:%2151=1013%:%
%:%2156=1013%:%
%:%2159=1014%:%
%:%2160=1015%:%
%:%2161=1015%:%
%:%2164=1016%:%
%:%2168=1016%:%
%:%2169=1016%:%
%:%2174=1016%:%
%:%2177=1017%:%
%:%2178=1018%:%
%:%2179=1019%:%
%:%2180=1019%:%
%:%2181=1020%:%
%:%2182=1021%:%
%:%2185=1022%:%
%:%2189=1022%:%
%:%2190=1022%:%
%:%2191=1023%:%
%:%2192=1023%:%
%:%2193=1024%:%
%:%2194=1025%:%
%:%2195=1025%:%
%:%2196=1026%:%
%:%2197=1026%:%
%:%2198=1027%:%
%:%2199=1027%:%
%:%2200=1028%:%
%:%2201=1028%:%
%:%2202=1029%:%
%:%2203=1029%:%
%:%2204=1029%:%
%:%2205=1030%:%
%:%2206=1030%:%
%:%2207=1031%:%
%:%2208=1031%:%
%:%2209=1032%:%
%:%2210=1033%:%
%:%2216=1033%:%
%:%2219=1034%:%
%:%2220=1035%:%
%:%2221=1035%:%
%:%2222=1036%:%
%:%2223=1037%:%
%:%2226=1038%:%
%:%2230=1038%:%
%:%2231=1038%:%
%:%2232=1039%:%
%:%2233=1039%:%
%:%2234=1040%:%
%:%2235=1040%:%
%:%2240=1040%:%
%:%2243=1041%:%
%:%2244=1042%:%
%:%2245=1042%:%
%:%2246=1043%:%
%:%2247=1044%:%
%:%2250=1045%:%
%:%2254=1045%:%
%:%2255=1045%:%
%:%2256=1046%:%
%:%2257=1046%:%
%:%2258=1047%:%
%:%2259=1047%:%
%:%2260=1048%:%
%:%2261=1048%:%
%:%2262=1048%:%
%:%2263=1048%:%
%:%2264=1048%:%
%:%2265=1049%:%
%:%2271=1049%:%
%:%2274=1050%:%
%:%2275=1051%:%
%:%2276=1051%:%
%:%2277=1052%:%
%:%2278=1053%:%
%:%2281=1054%:%
%:%2285=1054%:%
%:%2286=1054%:%
%:%2287=1055%:%
%:%2288=1055%:%
%:%2293=1055%:%
%:%2296=1056%:%
%:%2297=1057%:%
%:%2298=1057%:%
%:%2301=1058%:%
%:%2305=1058%:%
%:%2306=1058%:%
%:%2311=1058%:%
%:%2314=1059%:%
%:%2315=1060%:%
%:%2316=1060%:%
%:%2317=1061%:%
%:%2318=1062%:%
%:%2319=1063%:%
%:%2320=1064%:%
%:%2323=1065%:%
%:%2327=1065%:%
%:%2328=1065%:%
%:%2329=1066%:%
%:%2330=1066%:%
%:%2335=1066%:%
%:%2338=1067%:%
%:%2339=1068%:%
%:%2340=1068%:%
%:%2343=1069%:%
%:%2347=1069%:%
%:%2348=1069%:%
%:%2349=1070%:%
%:%2350=1070%:%
%:%2351=1071%:%
%:%2352=1071%:%
%:%2353=1072%:%
%:%2354=1072%:%
%:%2355=1073%:%
%:%2363=1073%:%
%:%2364=1074%:%
%:%2365=1075%:%
%:%2366=1075%:%
%:%2367=1076%:%
%:%2368=1077%:%
%:%2369=1078%:%
%:%2370=1079%:%
%:%2373=1080%:%
%:%2377=1080%:%
%:%2378=1080%:%
%:%2379=1081%:%
%:%2380=1081%:%
%:%2385=1081%:%
%:%2388=1082%:%
%:%2389=1083%:%
%:%2392=1085%:%
%:%2393=1086%:%
%:%2395=1087%:%
%:%2396=1087%:%
%:%2397=1088%:%
%:%2398=1089%:%
%:%2399=1090%:%
%:%2400=1091%:%
%:%2401=1092%:%
%:%2404=1093%:%
%:%2408=1093%:%
%:%2409=1093%:%
%:%2410=1094%:%
%:%2411=1094%:%
%:%2412=1095%:%
%:%2413=1096%:%
%:%2414=1097%:%
%:%2415=1097%:%
%:%2416=1098%:%
%:%2417=1098%:%
%:%2418=1099%:%
%:%2419=1099%:%
%:%2420=1100%:%
%:%2421=1100%:%
%:%2422=1101%:%
%:%2423=1101%:%
%:%2424=1102%:%
%:%2425=1102%:%
%:%2426=1103%:%
%:%2427=1103%:%
%:%2428=1104%:%
%:%2429=1105%:%
%:%2430=1105%:%
%:%2431=1106%:%
%:%2432=1106%:%
%:%2433=1107%:%
%:%2434=1108%:%
%:%2435=1109%:%
%:%2436=1110%:%
%:%2437=1110%:%
%:%2438=1111%:%
%:%2439=1111%:%
%:%2440=1112%:%
%:%2441=1112%:%
%:%2442=1112%:%
%:%2443=1113%:%
%:%2444=1113%:%
%:%2445=1114%:%
%:%2446=1114%:%
%:%2447=1115%:%
%:%2453=1115%:%
%:%2456=1116%:%
%:%2457=1117%:%
%:%2458=1117%:%
%:%2459=1118%:%
%:%2460=1119%:%
%:%2461=1120%:%
%:%2462=1121%:%
%:%2463=1122%:%
%:%2466=1123%:%
%:%2470=1123%:%
%:%2471=1123%:%
%:%2472=1124%:%
%:%2473=1124%:%
%:%2474=1125%:%
%:%2475=1126%:%
%:%2476=1127%:%
%:%2477=1127%:%
%:%2478=1128%:%
%:%2479=1128%:%
%:%2480=1129%:%
%:%2481=1129%:%
%:%2482=1130%:%
%:%2483=1130%:%
%:%2484=1131%:%
%:%2485=1131%:%
%:%2486=1132%:%
%:%2487=1132%:%
%:%2488=1133%:%
%:%2489=1133%:%
%:%2490=1134%:%
%:%2491=1135%:%
%:%2492=1135%:%
%:%2493=1136%:%
%:%2494=1136%:%
%:%2495=1137%:%
%:%2496=1138%:%
%:%2497=1139%:%
%:%2498=1140%:%
%:%2499=1140%:%
%:%2500=1141%:%
%:%2501=1141%:%
%:%2502=1142%:%
%:%2503=1142%:%
%:%2504=1142%:%
%:%2505=1143%:%
%:%2506=1143%:%
%:%2507=1144%:%
%:%2508=1145%:%
%:%2509=1146%:%
%:%2510=1147%:%
%:%2511=1148%:%
%:%2512=1149%:%
%:%2513=1150%:%
%:%2514=1150%:%
%:%2515=1151%:%
%:%2516=1151%:%
%:%2517=1152%:%
%:%2523=1152%:%
%:%2526=1153%:%
%:%2527=1154%:%
%:%2528=1154%:%
%:%2529=1155%:%
%:%2530=1156%:%
%:%2531=1157%:%
%:%2532=1158%:%
%:%2533=1159%:%
%:%2536=1160%:%
%:%2540=1160%:%
%:%2541=1160%:%
%:%2542=1161%:%
%:%2543=1161%:%
%:%2544=1162%:%
%:%2545=1162%:%
%:%2546=1163%:%
%:%2547=1163%:%
%:%2548=1163%:%
%:%2549=1164%:%
%:%2550=1164%:%
%:%2551=1165%:%
%:%2552=1165%:%
%:%2553=1166%:%
%:%2554=1166%:%
%:%2555=1167%:%
%:%2556=1167%:%
%:%2557=1168%:%
%:%2558=1168%:%
%:%2559=1168%:%
%:%2560=1169%:%
%:%2561=1169%:%
%:%2562=1169%:%
%:%2563=1170%:%
%:%2564=1170%:%
%:%2565=1170%:%
%:%2566=1171%:%
%:%2567=1171%:%
%:%2568=1171%:%
%:%2569=1172%:%
%:%2570=1172%:%
%:%2571=1172%:%
%:%2572=1173%:%
%:%2573=1173%:%
%:%2574=1174%:%
%:%2575=1174%:%
%:%2576=1175%:%
%:%2582=1175%:%
%:%2585=1176%:%
%:%2586=1177%:%
%:%2587=1177%:%
%:%2588=1177%:%
%:%2596=1181%:%
%:%2606=1183%:%
%:%2607=1183%:%
%:%2608=1184%:%
%:%2609=1185%:%
%:%2610=1185%:%
%:%2611=1186%:%
%:%2612=1187%:%
%:%2613=1188%:%
%:%2614=1188%:%
%:%2615=1189%:%
%:%2617=1191%:%
%:%2618=1192%:%
%:%2620=1193%:%
%:%2621=1193%:%
%:%2622=1194%:%
%:%2624=1196%:%
%:%2626=1197%:%
%:%2627=1197%:%
%:%2628=1198%:%
%:%2629=1199%:%
%:%2630=1200%:%
%:%2631=1200%:%
%:%2632=1200%:%
%:%2633=1200%:%
%:%2634=1201%:%
%:%2635=1202%:%
%:%2636=1202%:%
%:%2638=1203%:%
%:%2640=1204%:%
%:%2641=1204%:%
%:%2644=1205%:%
%:%2648=1205%:%
%:%2649=1205%:%
%:%2650=1206%:%
%:%2651=1206%:%
%:%2652=1207%:%
%:%2653=1207%:%
%:%2654=1208%:%
%:%2655=1208%:%
%:%2656=1209%:%
%:%2657=1209%:%
%:%2658=1210%:%
%:%2659=1210%:%
%:%2660=1211%:%
%:%2661=1211%:%
%:%2662=1212%:%
%:%2668=1212%:%
%:%2671=1213%:%
%:%2672=1214%:%
%:%2673=1215%:%
%:%2674=1215%:%
%:%2675=1216%:%
%:%2678=1217%:%
%:%2682=1217%:%
%:%2683=1217%:%
%:%2684=1218%:%
%:%2685=1218%:%
%:%2694=1220%:%
%:%2696=1221%:%
%:%2697=1221%:%
%:%2698=1222%:%
%:%2699=1223%:%
%:%2700=1224%:%
%:%2703=1225%:%
%:%2707=1225%:%
%:%2708=1225%:%
%:%2709=1225%:%
%:%2714=1225%:%
%:%2717=1226%:%
%:%2718=1227%:%
%:%2719=1227%:%
%:%2720=1228%:%
%:%2723=1229%:%
%:%2727=1229%:%
%:%2728=1229%:%
%:%2733=1229%:%
%:%2736=1230%:%
%:%2737=1231%:%
%:%2738=1231%:%
%:%2739=1232%:%
%:%2740=1233%:%
%:%2741=1234%:%
%:%2742=1235%:%
%:%2745=1236%:%
%:%2749=1236%:%
%:%2750=1236%:%
%:%2751=1236%:%
%:%2760=1238%:%
%:%2762=1239%:%
%:%2763=1239%:%
%:%2764=1240%:%
%:%2765=1241%:%
%:%2766=1242%:%
%:%2767=1243%:%
%:%2770=1244%:%
%:%2774=1244%:%
%:%2775=1244%:%
%:%2776=1244%:%
%:%2781=1244%:%
%:%2784=1245%:%
%:%2785=1246%:%
%:%2786=1246%:%
%:%2787=1247%:%
%:%2788=1248%:%
%:%2791=1249%:%
%:%2795=1249%:%
%:%2796=1249%:%
%:%2797=1250%:%
%:%2798=1250%:%
%:%2799=1251%:%
%:%2800=1251%:%
%:%2801=1251%:%
%:%2802=1252%:%
%:%2803=1252%:%
%:%2804=1252%:%
%:%2805=1253%:%
%:%2806=1253%:%
%:%2807=1254%:%
%:%2808=1254%:%
%:%2809=1255%:%
%:%2810=1255%:%
%:%2811=1255%:%
%:%2812=1256%:%
%:%2813=1256%:%
%:%2814=1256%:%
%:%2815=1257%:%
%:%2816=1257%:%
%:%2817=1257%:%
%:%2818=1258%:%
%:%2819=1258%:%
%:%2820=1259%:%
%:%2821=1259%:%
%:%2822=1260%:%
%:%2828=1260%:%
%:%2831=1261%:%
%:%2832=1262%:%
%:%2833=1262%:%
%:%2834=1263%:%
%:%2835=1264%:%
%:%2836=1265%:%
%:%2839=1266%:%
%:%2843=1266%:%
%:%2844=1266%:%
%:%2845=1267%:%
%:%2846=1267%:%
%:%2847=1268%:%
%:%2848=1268%:%
%:%2849=1268%:%
%:%2850=1269%:%
%:%2851=1269%:%
%:%2852=1269%:%
%:%2853=1270%:%
%:%2854=1270%:%
%:%2855=1271%:%
%:%2856=1271%:%
%:%2857=1272%:%
%:%2858=1272%:%
%:%2859=1272%:%
%:%2860=1273%:%
%:%2861=1273%:%
%:%2862=1273%:%
%:%2863=1274%:%
%:%2864=1274%:%
%:%2865=1274%:%
%:%2866=1275%:%
%:%2867=1275%:%
%:%2868=1276%:%
%:%2869=1276%:%
%:%2870=1277%:%
%:%2880=1279%:%
%:%2882=1280%:%
%:%2883=1280%:%
%:%2884=1281%:%
%:%2885=1282%:%
%:%2886=1283%:%
%:%2887=1284%:%
%:%2890=1285%:%
%:%2894=1285%:%
%:%2895=1285%:%
%:%2896=1286%:%
%:%2897=1286%:%
%:%2898=1287%:%
%:%2899=1287%:%
%:%2900=1288%:%
%:%2901=1288%:%
%:%2902=1288%:%
%:%2903=1289%:%
%:%2904=1289%:%
%:%2905=1290%:%
%:%2906=1290%:%
%:%2907=1290%:%
%:%2908=1291%:%
%:%2909=1291%:%
%:%2910=1291%:%
%:%2911=1292%:%
%:%2912=1292%:%
%:%2913=1293%:%
%:%2914=1293%:%
%:%2915=1294%:%
%:%2916=1294%:%
%:%2917=1294%:%
%:%2918=1295%:%
%:%2919=1295%:%
%:%2920=1296%:%
%:%2921=1296%:%
%:%2922=1297%:%
%:%2923=1297%:%
%:%2924=1297%:%
%:%2925=1298%:%
%:%2926=1299%:%
%:%2927=1299%:%
%:%2928=1299%:%
%:%2929=1300%:%
%:%2930=1300%:%
%:%2931=1300%:%
%:%2932=1301%:%
%:%2933=1301%:%
%:%2934=1301%:%
%:%2935=1301%:%
%:%2936=1302%:%
%:%2937=1302%:%
%:%2938=1302%:%
%:%2939=1303%:%
%:%2940=1303%:%
%:%2941=1304%:%
%:%2942=1304%:%
%:%2943=1305%:%
%:%2944=1305%:%
%:%2945=1306%:%
%:%2946=1306%:%
%:%2947=1307%:%
%:%2948=1307%:%
%:%2949=1307%:%
%:%2950=1307%:%
%:%2951=1307%:%
%:%2952=1308%:%
%:%2953=1308%:%
%:%2954=1309%:%
%:%2955=1309%:%
%:%2956=1310%:%
%:%2957=1310%:%
%:%2958=1311%:%
%:%2959=1311%:%
%:%2960=1311%:%
%:%2961=1312%:%
%:%2962=1312%:%
%:%2963=1313%:%
%:%2964=1313%:%
%:%2965=1314%:%
%:%2966=1314%:%
%:%2967=1314%:%
%:%2968=1314%:%
%:%2969=1314%:%
%:%2970=1315%:%
%:%2971=1315%:%
%:%2972=1316%:%
%:%2973=1316%:%
%:%2974=1317%:%
%:%2975=1317%:%
%:%2976=1317%:%
%:%2977=1318%:%
%:%2978=1319%:%
%:%2979=1319%:%
%:%2980=1319%:%
%:%2981=1320%:%
%:%2982=1320%:%
%:%2983=1320%:%
%:%2984=1321%:%
%:%2985=1321%:%
%:%2986=1321%:%
%:%2987=1321%:%
%:%2988=1322%:%
%:%2989=1322%:%
%:%2990=1322%:%
%:%2991=1323%:%
%:%2992=1323%:%
%:%2993=1324%:%
%:%2994=1324%:%
%:%2995=1325%:%
%:%2996=1325%:%
%:%2997=1326%:%
%:%2998=1326%:%
%:%2999=1327%:%
%:%3009=1329%:%
%:%3010=1330%:%
%:%3012=1331%:%
%:%3013=1331%:%
%:%3014=1332%:%
%:%3015=1333%:%
%:%3016=1334%:%
%:%3017=1335%:%
%:%3018=1336%:%
%:%3021=1337%:%
%:%3025=1337%:%
%:%3026=1337%:%
%:%3027=1338%:%
%:%3028=1338%:%
%:%3029=1339%:%
%:%3030=1339%:%
%:%3031=1340%:%
%:%3032=1340%:%
%:%3033=1340%:%
%:%3034=1341%:%
%:%3035=1341%:%
%:%3036=1341%:%
%:%3037=1342%:%
%:%3038=1342%:%
%:%3039=1343%:%
%:%3040=1343%:%
%:%3041=1343%:%
%:%3042=1344%:%
%:%3043=1344%:%
%:%3044=1345%:%
%:%3045=1345%:%
%:%3046=1346%:%
%:%3056=1348%:%
%:%3057=1349%:%
%:%3059=1350%:%
%:%3060=1350%:%
%:%3061=1351%:%
%:%3062=1352%:%
%:%3063=1353%:%
%:%3066=1354%:%
%:%3070=1354%:%
%:%3071=1354%:%
%:%3072=1355%:%
%:%3073=1355%:%
%:%3074=1355%:%
%:%3083=1357%:%
%:%3085=1358%:%
%:%3086=1358%:%
%:%3087=1359%:%
%:%3088=1360%:%
%:%3091=1361%:%
%:%3095=1361%:%
%:%3096=1361%:%
%:%3097=1362%:%
%:%3098=1362%:%
%:%3099=1362%:%
%:%3104=1362%:%
%:%3107=1363%:%
%:%3108=1364%:%
%:%3109=1364%:%
%:%3110=1364%:%
%:%3113=1367%:%
%:%3115=1368%:%
%:%3116=1368%:%
%:%3117=1369%:%
%:%3118=1370%:%
%:%3119=1370%:%
%:%3120=1371%:%
%:%3121=1372%:%
%:%3122=1372%:%
%:%3123=1373%:%
%:%3126=1374%:%
%:%3130=1374%:%
%:%3131=1374%:%
%:%3136=1374%:%
%:%3139=1375%:%
%:%3140=1376%:%
%:%3141=1376%:%
%:%3144=1377%:%
%:%3148=1377%:%
%:%3149=1377%:%
%:%3163=1381%:%
%:%3173=1384%:%
%:%3174=1384%:%
%:%3175=1385%:%
%:%3176=1386%:%
%:%3177=1387%:%
%:%3178=1387%:%
%:%3179=1388%:%
%:%3180=1389%:%
%:%3181=1390%:%
%:%3182=1391%:%
%:%3183=1391%:%
%:%3184=1392%:%
%:%3187=1393%:%
%:%3191=1393%:%
%:%3192=1393%:%
%:%3197=1393%:%
%:%3200=1394%:%
%:%3201=1395%:%
%:%3202=1395%:%
%:%3203=1396%:%
%:%3207=1400%:%
%:%3210=1401%:%
%:%3214=1401%:%
%:%3215=1401%:%
%:%3220=1401%:%
%:%3223=1402%:%
%:%3224=1403%:%
%:%3225=1403%:%
%:%3226=1404%:%
%:%3227=1405%:%
%:%3230=1406%:%
%:%3234=1406%:%
%:%3235=1406%:%
%:%3240=1406%:%
%:%3243=1407%:%
%:%3244=1408%:%
%:%3245=1408%:%
%:%3246=1409%:%
%:%3247=1410%:%
%:%3250=1411%:%
%:%3254=1411%:%
%:%3255=1411%:%
%:%3256=1412%:%
%:%3257=1412%:%
%:%3258=1413%:%
%:%3259=1413%:%
%:%3264=1413%:%
%:%3267=1414%:%
%:%3268=1415%:%
%:%3269=1415%:%
%:%3270=1416%:%
%:%3272=1418%:%
%:%3275=1419%:%
%:%3279=1419%:%
%:%3280=1419%:%
%:%3285=1419%:%
%:%3288=1420%:%
%:%3289=1421%:%
%:%3290=1421%:%
%:%3291=1422%:%
%:%3292=1423%:%
%:%3295=1424%:%
%:%3299=1424%:%
%:%3300=1424%:%
%:%3305=1424%:%
%:%3308=1425%:%
%:%3309=1426%:%
%:%3310=1426%:%
%:%3311=1427%:%
%:%3314=1430%:%
%:%3317=1431%:%
%:%3321=1431%:%
%:%3322=1431%:%
%:%3327=1431%:%
%:%3330=1432%:%
%:%3331=1433%:%
%:%3332=1433%:%
%:%3333=1434%:%
%:%3337=1438%:%
%:%3340=1439%:%
%:%3344=1439%:%
%:%3345=1439%:%
%:%3346=1440%:%
%:%3351=1440%:%
%:%3354=1441%:%
%:%3355=1442%:%
%:%3356=1442%:%
%:%3359=1443%:%
%:%3363=1443%:%
%:%3364=1443%:%
%:%3369=1443%:%
%:%3372=1444%:%
%:%3373=1445%:%
%:%3374=1445%:%
%:%3375=1446%:%
%:%3376=1447%:%
%:%3379=1448%:%
%:%3383=1448%:%
%:%3384=1448%:%
%:%3389=1448%:%
%:%3392=1449%:%
%:%3393=1450%:%
%:%3394=1450%:%
%:%3395=1451%:%
%:%3398=1452%:%
%:%3402=1452%:%
%:%3403=1452%:%
%:%3408=1452%:%
%:%3411=1453%:%
%:%3412=1454%:%
%:%3413=1454%:%
%:%3414=1455%:%
%:%3417=1456%:%
%:%3421=1456%:%
%:%3422=1456%:%
%:%3427=1456%:%
%:%3430=1457%:%
%:%3431=1458%:%
%:%3432=1458%:%
%:%3433=1459%:%
%:%3436=1460%:%
%:%3440=1460%:%
%:%3441=1460%:%
%:%3446=1460%:%
%:%3449=1461%:%
%:%3450=1462%:%
%:%3451=1462%:%
%:%3454=1463%:%
%:%3458=1463%:%
%:%3459=1463%:%
%:%3464=1463%:%
%:%3467=1464%:%
%:%3468=1465%:%
%:%3469=1465%:%
%:%3470=1466%:%
%:%3471=1467%:%
%:%3472=1468%:%
%:%3473=1469%:%
%:%3474=1470%:%
%:%3475=1471%:%
%:%3476=1472%:%
%:%3477=1473%:%
%:%3479=1475%:%
%:%3480=1476%:%
%:%3482=1478%:%
%:%3483=1479%:%
%:%3484=1480%:%
%:%3485=1481%:%
%:%3486=1482%:%
%:%3487=1483%:%
%:%3488=1484%:%
%:%3491=1485%:%
%:%3495=1485%:%
%:%3496=1485%:%
%:%3497=1486%:%
%:%3498=1486%:%
%:%3499=1487%:%
%:%3500=1488%:%
%:%3501=1488%:%
%:%3502=1489%:%
%:%3503=1489%:%
%:%3504=1489%:%
%:%3505=1490%:%
%:%3506=1491%:%
%:%3507=1492%:%
%:%3508=1493%:%
%:%3509=1494%:%
%:%3510=1495%:%
%:%3511=1495%:%
%:%3512=1496%:%
%:%3513=1496%:%
%:%3514=1497%:%
%:%3515=1497%:%
%:%3516=1498%:%
%:%3517=1498%:%
%:%3518=1499%:%
%:%3519=1499%:%
%:%3520=1500%:%
%:%3521=1501%:%
%:%3522=1501%:%
%:%3523=1502%:%
%:%3524=1503%:%
%:%3525=1503%:%
%:%3526=1504%:%
%:%3527=1504%:%
%:%3528=1505%:%
%:%3529=1505%:%
%:%3530=1506%:%
%:%3531=1506%:%
%:%3532=1507%:%
%:%3533=1507%:%
%:%3534=1508%:%
%:%3535=1508%:%
%:%3536=1508%:%
%:%3537=1509%:%
%:%3538=1510%:%
%:%3539=1510%:%
%:%3540=1511%:%
%:%3541=1511%:%
%:%3542=1512%:%
%:%3543=1512%:%
%:%3544=1512%:%
%:%3545=1513%:%
%:%3546=1513%:%
%:%3547=1514%:%
%:%3548=1514%:%
%:%3549=1515%:%
%:%3550=1515%:%
%:%3551=1515%:%
%:%3552=1516%:%
%:%3553=1516%:%
%:%3554=1517%:%
%:%3555=1517%:%
%:%3556=1518%:%
%:%3557=1518%:%
%:%3558=1519%:%
%:%3559=1519%:%
%:%3560=1520%:%
%:%3561=1520%:%
%:%3562=1520%:%
%:%3563=1521%:%
%:%3564=1521%:%
%:%3565=1522%:%
%:%3566=1522%:%
%:%3567=1523%:%
%:%3568=1523%:%
%:%3569=1523%:%
%:%3570=1524%:%
%:%3571=1524%:%
%:%3572=1525%:%
%:%3573=1525%:%
%:%3574=1526%:%
%:%3575=1526%:%
%:%3576=1527%:%
%:%3577=1527%:%
%:%3578=1528%:%
%:%3579=1528%:%
%:%3580=1529%:%
%:%3581=1530%:%
%:%3582=1530%:%
%:%3583=1531%:%
%:%3584=1531%:%
%:%3585=1532%:%
%:%3586=1532%:%
%:%3587=1533%:%
%:%3588=1533%:%
%:%3589=1534%:%
%:%3590=1534%:%
%:%3591=1535%:%
%:%3592=1535%:%
%:%3593=1536%:%
%:%3594=1536%:%
%:%3595=1537%:%
%:%3596=1537%:%
%:%3597=1538%:%
%:%3598=1539%:%
%:%3599=1539%:%
%:%3600=1540%:%
%:%3601=1540%:%
%:%3602=1540%:%
%:%3603=1540%:%
%:%3604=1541%:%
%:%3605=1541%:%
%:%3606=1542%:%
%:%3607=1543%:%
%:%3608=1543%:%
%:%3609=1544%:%
%:%3610=1544%:%
%:%3611=1544%:%
%:%3612=1545%:%
%:%3613=1545%:%
%:%3614=1546%:%
%:%3615=1546%:%
%:%3616=1547%:%
%:%3617=1547%:%
%:%3618=1548%:%
%:%3619=1548%:%
%:%3620=1549%:%
%:%3621=1549%:%
%:%3622=1549%:%
%:%3623=1550%:%
%:%3624=1551%:%
%:%3625=1552%:%
%:%3626=1552%:%
%:%3627=1553%:%
%:%3628=1553%:%
%:%3629=1554%:%
%:%3630=1554%:%
%:%3631=1554%:%
%:%3632=1555%:%
%:%3633=1555%:%
%:%3634=1556%:%
%:%3635=1556%:%
%:%3636=1557%:%
%:%3637=1557%:%
%:%3638=1558%:%
%:%3639=1558%:%
%:%3640=1559%:%
%:%3641=1559%:%
%:%3642=1560%:%
%:%3643=1560%:%
%:%3644=1560%:%
%:%3645=1561%:%
%:%3646=1561%:%
%:%3647=1562%:%
%:%3648=1562%:%
%:%3649=1563%:%
%:%3650=1563%:%
%:%3651=1563%:%
%:%3652=1564%:%
%:%3653=1564%:%
%:%3654=1565%:%
%:%3655=1566%:%
%:%3656=1566%:%
%:%3657=1567%:%
%:%3658=1568%:%
%:%3659=1568%:%
%:%3660=1569%:%
%:%3661=1569%:%
%:%3662=1570%:%
%:%3663=1570%:%
%:%3664=1571%:%
%:%3665=1571%:%
%:%3666=1572%:%
%:%3667=1572%:%
%:%3668=1573%:%
%:%3669=1573%:%
%:%3670=1574%:%
%:%3671=1574%:%
%:%3672=1574%:%
%:%3673=1575%:%
%:%3674=1575%:%
%:%3675=1576%:%
%:%3676=1576%:%
%:%3677=1577%:%
%:%3678=1577%:%
%:%3679=1577%:%
%:%3680=1578%:%
%:%3681=1578%:%
%:%3682=1579%:%
%:%3683=1579%:%
%:%3684=1580%:%
%:%3685=1581%:%
%:%3686=1581%:%
%:%3687=1582%:%
%:%3688=1582%:%
%:%3689=1583%:%
%:%3690=1583%:%
%:%3691=1584%:%
%:%3692=1584%:%
%:%3693=1585%:%
%:%3694=1586%:%
%:%3695=1586%:%
%:%3696=1587%:%
%:%3697=1587%:%
%:%3698=1588%:%
%:%3699=1588%:%
%:%3700=1589%:%
%:%3701=1589%:%
%:%3702=1590%:%
%:%3703=1590%:%
%:%3704=1591%:%
%:%3705=1591%:%
%:%3706=1592%:%
%:%3707=1592%:%
%:%3708=1593%:%
%:%3709=1593%:%
%:%3710=1594%:%
%:%3711=1594%:%
%:%3712=1595%:%
%:%3713=1595%:%
%:%3714=1596%:%
%:%3715=1597%:%
%:%3716=1597%:%
%:%3717=1598%:%
%:%3718=1598%:%
%:%3719=1599%:%
%:%3720=1599%:%
%:%3721=1600%:%
%:%3722=1600%:%
%:%3723=1601%:%
%:%3724=1601%:%
%:%3725=1601%:%
%:%3726=1602%:%
%:%3727=1602%:%
%:%3728=1603%:%
%:%3729=1603%:%
%:%3730=1604%:%
%:%3731=1604%:%
%:%3732=1604%:%
%:%3733=1605%:%
%:%3734=1605%:%
%:%3735=1606%:%
%:%3736=1606%:%
%:%3737=1607%:%
%:%3738=1607%:%
%:%3739=1608%:%
%:%3740=1608%:%
%:%3741=1608%:%
%:%3742=1609%:%
%:%3743=1609%:%
%:%3744=1610%:%
%:%3745=1610%:%
%:%3746=1611%:%
%:%3747=1611%:%
%:%3748=1612%:%
%:%3749=1612%:%
%:%3750=1613%:%
%:%3751=1613%:%
%:%3752=1613%:%
%:%3753=1614%:%
%:%3754=1614%:%
%:%3755=1615%:%
%:%3756=1615%:%
%:%3757=1616%:%
%:%3758=1616%:%
%:%3759=1616%:%
%:%3760=1617%:%
%:%3761=1617%:%
%:%3762=1618%:%
%:%3763=1618%:%
%:%3764=1619%:%
%:%3765=1619%:%
%:%3766=1620%:%
%:%3767=1620%:%
%:%3768=1620%:%
%:%3769=1621%:%
%:%3770=1621%:%
%:%3771=1622%:%
%:%3772=1623%:%
%:%3773=1624%:%
%:%3774=1624%:%
%:%3775=1625%:%
%:%3776=1625%:%
%:%3777=1625%:%
%:%3778=1626%:%
%:%3779=1626%:%
%:%3780=1627%:%
%:%3781=1628%:%
%:%3782=1628%:%
%:%3783=1629%:%
%:%3784=1629%:%
%:%3785=1629%:%
%:%3786=1630%:%
%:%3787=1630%:%
%:%3788=1631%:%
%:%3789=1631%:%
%:%3790=1632%:%
%:%3791=1632%:%
%:%3792=1632%:%
%:%3793=1633%:%
%:%3794=1633%:%
%:%3795=1634%:%
%:%3796=1634%:%
%:%3797=1635%:%
%:%3798=1635%:%
%:%3799=1635%:%
%:%3800=1636%:%
%:%3801=1637%:%
%:%3802=1637%:%
%:%3803=1638%:%
%:%3804=1638%:%
%:%3805=1639%:%
%:%3806=1639%:%
%:%3807=1639%:%
%:%3808=1640%:%
%:%3809=1640%:%
%:%3810=1641%:%
%:%3811=1641%:%
%:%3812=1642%:%
%:%3813=1642%:%
%:%3814=1642%:%
%:%3815=1643%:%
%:%3816=1643%:%
%:%3817=1644%:%
%:%3818=1644%:%
%:%3819=1645%:%
%:%3820=1645%:%
%:%3821=1645%:%
%:%3822=1646%:%
%:%3823=1647%:%
%:%3824=1648%:%
%:%3825=1649%:%
%:%3826=1649%:%
%:%3827=1650%:%
%:%3828=1650%:%
%:%3829=1651%:%
%:%3830=1651%:%
%:%3831=1652%:%
%:%3832=1652%:%
%:%3833=1653%:%
%:%3834=1653%:%
%:%3835=1654%:%
%:%3836=1654%:%
%:%3837=1655%:%
%:%3838=1655%:%
%:%3839=1656%:%
%:%3840=1656%:%
%:%3841=1656%:%
%:%3842=1657%:%
%:%3843=1657%:%
%:%3844=1658%:%
%:%3845=1658%:%
%:%3846=1659%:%
%:%3847=1659%:%
%:%3848=1660%:%
%:%3849=1660%:%
%:%3850=1660%:%
%:%3851=1661%:%
%:%3852=1661%:%
%:%3853=1662%:%
%:%3854=1662%:%
%:%3855=1663%:%
%:%3856=1663%:%
%:%3857=1664%:%
%:%3858=1664%:%
%:%3859=1665%:%
%:%3860=1665%:%
%:%3861=1666%:%
%:%3862=1666%:%
%:%3863=1667%:%
%:%3864=1667%:%
%:%3865=1668%:%
%:%3866=1668%:%
%:%3867=1669%:%
%:%3868=1669%:%
%:%3869=1669%:%
%:%3870=1670%:%
%:%3871=1670%:%
%:%3872=1671%:%
%:%3873=1671%:%
%:%3874=1672%:%
%:%3875=1672%:%
%:%3876=1673%:%
%:%3877=1673%:%
%:%3878=1674%:%
%:%3879=1674%:%
%:%3880=1675%:%
%:%3881=1676%:%
%:%3882=1676%:%
%:%3883=1676%:%
%:%3884=1677%:%
%:%3885=1678%:%
%:%3886=1679%:%
%:%3887=1680%:%
%:%3888=1680%:%
%:%3889=1681%:%
%:%3890=1681%:%
%:%3891=1682%:%
%:%3892=1682%:%
%:%3893=1683%:%
%:%3894=1683%:%
%:%3895=1684%:%
%:%3896=1684%:%
%:%3897=1685%:%
%:%3898=1685%:%
%:%3899=1686%:%
%:%3900=1686%:%
%:%3901=1687%:%
%:%3902=1687%:%
%:%3903=1687%:%
%:%3904=1688%:%
%:%3905=1688%:%
%:%3906=1689%:%
%:%3907=1689%:%
%:%3908=1690%:%
%:%3909=1690%:%
%:%3910=1691%:%
%:%3911=1691%:%
%:%3912=1691%:%
%:%3913=1692%:%
%:%3914=1692%:%
%:%3915=1693%:%
%:%3916=1693%:%
%:%3917=1694%:%
%:%3918=1694%:%
%:%3919=1695%:%
%:%3920=1695%:%
%:%3921=1696%:%
%:%3922=1696%:%
%:%3923=1697%:%
%:%3924=1697%:%
%:%3925=1698%:%
%:%3926=1698%:%
%:%3927=1699%:%
%:%3928=1699%:%
%:%3929=1700%:%
%:%3930=1700%:%
%:%3931=1700%:%
%:%3932=1701%:%
%:%3933=1701%:%
%:%3934=1702%:%
%:%3935=1702%:%
%:%3936=1703%:%
%:%3937=1703%:%
%:%3938=1704%:%
%:%3939=1704%:%
%:%3940=1705%:%
%:%3941=1705%:%
%:%3942=1706%:%
%:%3943=1706%:%
%:%3944=1706%:%
%:%3945=1707%:%
%:%3946=1707%:%
%:%3947=1708%:%
%:%3948=1709%:%
%:%3949=1709%:%
%:%3950=1710%:%
%:%3951=1710%:%
%:%3952=1710%:%
%:%3953=1711%:%
%:%3954=1711%:%
%:%3955=1712%:%
%:%3956=1712%:%
%:%3957=1713%:%
%:%3958=1713%:%
%:%3959=1713%:%
%:%3960=1714%:%
%:%3961=1714%:%
%:%3962=1715%:%
%:%3963=1715%:%
%:%3964=1716%:%
%:%3965=1716%:%
%:%3966=1716%:%
%:%3967=1717%:%
%:%3968=1717%:%
%:%3969=1718%:%
%:%3970=1718%:%
%:%3971=1719%:%
%:%3972=1719%:%
%:%3973=1719%:%
%:%3974=1720%:%
%:%3975=1720%:%
%:%3976=1721%:%
%:%3977=1721%:%
%:%3978=1722%:%
%:%3979=1722%:%
%:%3980=1722%:%
%:%3981=1723%:%
%:%3982=1723%:%
%:%3983=1724%:%
%:%3984=1724%:%
%:%3985=1724%:%
%:%3986=1725%:%
%:%3987=1726%:%
%:%3988=1727%:%
%:%3989=1727%:%
%:%3990=1728%:%
%:%3991=1728%:%
%:%3992=1729%:%
%:%3993=1729%:%
%:%3994=1730%:%
%:%3995=1730%:%
%:%3996=1731%:%
%:%3997=1731%:%
%:%3998=1731%:%
%:%3999=1732%:%
%:%4000=1732%:%
%:%4001=1732%:%
%:%4002=1733%:%
%:%4003=1733%:%
%:%4004=1733%:%
%:%4005=1734%:%
%:%4006=1734%:%
%:%4007=1734%:%
%:%4008=1735%:%
%:%4009=1735%:%
%:%4010=1735%:%
%:%4011=1736%:%
%:%4012=1736%:%
%:%4013=1737%:%
%:%4014=1737%:%
%:%4015=1738%:%
%:%4016=1738%:%
%:%4017=1738%:%
%:%4018=1739%:%
%:%4019=1739%:%
%:%4020=1740%:%
%:%4021=1740%:%
%:%4022=1741%:%
%:%4023=1741%:%
%:%4024=1741%:%
%:%4025=1742%:%
%:%4026=1742%:%
%:%4027=1743%:%
%:%4028=1743%:%
%:%4029=1744%:%
%:%4030=1744%:%
%:%4031=1744%:%
%:%4032=1745%:%
%:%4033=1745%:%
%:%4034=1746%:%
%:%4035=1746%:%
%:%4036=1747%:%
%:%4037=1747%:%
%:%4038=1748%:%
%:%4039=1748%:%
%:%4040=1749%:%
%:%4041=1749%:%
%:%4042=1749%:%
%:%4043=1750%:%
%:%4044=1750%:%
%:%4045=1751%:%
%:%4046=1751%:%
%:%4047=1752%:%
%:%4048=1752%:%
%:%4049=1752%:%
%:%4050=1753%:%
%:%4051=1753%:%
%:%4052=1754%:%
%:%4053=1754%:%
%:%4054=1755%:%
%:%4055=1755%:%
%:%4056=1756%:%
%:%4057=1756%:%
%:%4058=1756%:%
%:%4059=1757%:%
%:%4060=1757%:%
%:%4061=1758%:%
%:%4062=1758%:%
%:%4063=1759%:%
%:%4064=1759%:%
%:%4065=1759%:%
%:%4066=1760%:%
%:%4067=1760%:%
%:%4068=1761%:%
%:%4069=1761%:%
%:%4070=1762%:%
%:%4071=1762%:%
%:%4072=1763%:%
%:%4073=1763%:%
%:%4078=1763%:%
%:%4081=1764%:%
%:%4082=1765%:%
%:%4083=1766%:%
%:%4084=1766%:%
%:%4085=1767%:%
%:%4086=1768%:%
%:%4087=1769%:%
%:%4088=1770%:%
%:%4089=1771%:%
%:%4090=1772%:%
%:%4093=1773%:%
%:%4097=1773%:%
%:%4098=1773%:%
%:%4099=1774%:%
%:%4100=1774%:%
%:%4101=1775%:%
%:%4102=1776%:%
%:%4103=1777%:%
%:%4104=1777%:%
%:%4105=1778%:%
%:%4106=1778%:%
%:%4107=1779%:%
%:%4108=1779%:%
%:%4109=1780%:%
%:%4110=1780%:%
%:%4111=1781%:%
%:%4112=1781%:%
%:%4113=1782%:%
%:%4114=1782%:%
%:%4115=1783%:%
%:%4116=1783%:%
%:%4117=1784%:%
%:%4118=1784%:%
%:%4119=1785%:%
%:%4120=1785%:%
%:%4121=1785%:%
%:%4122=1785%:%
%:%4123=1786%:%
%:%4124=1787%:%
%:%4125=1788%:%
%:%4126=1788%:%
%:%4127=1789%:%
%:%4128=1789%:%
%:%4129=1789%:%
%:%4130=1790%:%
%:%4131=1790%:%
%:%4132=1791%:%
%:%4133=1791%:%
%:%4134=1792%:%
%:%4135=1792%:%
%:%4136=1792%:%
%:%4137=1793%:%
%:%4138=1793%:%
%:%4139=1794%:%
%:%4145=1794%:%
%:%4148=1795%:%
%:%4149=1796%:%
%:%4150=1796%:%
%:%4151=1797%:%
%:%4152=1798%:%
%:%4155=1799%:%
%:%4159=1799%:%
%:%4160=1799%:%
%:%4161=1800%:%
%:%4162=1800%:%
%:%4167=1800%:%
%:%4170=1801%:%
%:%4171=1802%:%
%:%4172=1802%:%
%:%4173=1803%:%
%:%4178=1808%:%
%:%4181=1809%:%
%:%4185=1809%:%
%:%4186=1809%:%
%:%4191=1809%:%
%:%4194=1810%:%
%:%4195=1811%:%
%:%4196=1811%:%
%:%4197=1812%:%
%:%4198=1813%:%
%:%4199=1814%:%
%:%4200=1815%:%
%:%4201=1816%:%
%:%4202=1817%:%
%:%4203=1818%:%
%:%4206=1819%:%
%:%4210=1819%:%
%:%4211=1819%:%
%:%4212=1820%:%
%:%4213=1820%:%
%:%4214=1821%:%
%:%4216=1823%:%
%:%4217=1824%:%
%:%4218=1824%:%
%:%4219=1825%:%
%:%4220=1826%:%
%:%4221=1826%:%
%:%4222=1827%:%
%:%4223=1827%:%
%:%4224=1827%:%
%:%4225=1828%:%
%:%4226=1828%:%
%:%4227=1829%:%
%:%4228=1829%:%
%:%4229=1830%:%
%:%4230=1830%:%
%:%4231=1830%:%
%:%4232=1831%:%
%:%4233=1831%:%
%:%4234=1832%:%
%:%4235=1832%:%
%:%4236=1833%:%
%:%4237=1833%:%
%:%4238=1833%:%
%:%4239=1834%:%
%:%4240=1834%:%
%:%4241=1835%:%
%:%4247=1835%:%
%:%4250=1836%:%
%:%4251=1837%:%
%:%4252=1837%:%
%:%4253=1838%:%
%:%4254=1839%:%
%:%4255=1840%:%
%:%4256=1841%:%
%:%4257=1842%:%
%:%4258=1843%:%
%:%4261=1844%:%
%:%4265=1844%:%
%:%4266=1844%:%
%:%4267=1845%:%
%:%4268=1845%:%
%:%4273=1845%:%
%:%4276=1846%:%
%:%4277=1847%:%
%:%4278=1847%:%
%:%4279=1848%:%
%:%4282=1849%:%
%:%4286=1849%:%
%:%4287=1849%:%
%:%4292=1849%:%
%:%4295=1850%:%
%:%4296=1851%:%
%:%4297=1851%:%
%:%4300=1852%:%
%:%4304=1852%:%
%:%4305=1852%:%
%:%4310=1852%:%
%:%4313=1853%:%
%:%4314=1854%:%
%:%4315=1854%:%
%:%4318=1855%:%
%:%4322=1855%:%
%:%4323=1855%:%
%:%4328=1855%:%
%:%4331=1856%:%
%:%4332=1857%:%
%:%4333=1857%:%
%:%4336=1858%:%
%:%4340=1858%:%
%:%4341=1858%:%
%:%4346=1858%:%
%:%4349=1859%:%
%:%4350=1860%:%
%:%4351=1860%:%
%:%4354=1861%:%
%:%4358=1861%:%
%:%4359=1861%:%
%:%4364=1861%:%
%:%4367=1862%:%
%:%4368=1863%:%
%:%4369=1863%:%
%:%4370=1864%:%
%:%4373=1865%:%
%:%4377=1865%:%
%:%4378=1865%:%
%:%4383=1865%:%
%:%4386=1866%:%
%:%4387=1867%:%
%:%4388=1867%:%
%:%4389=1868%:%
%:%4392=1869%:%
%:%4396=1869%:%
%:%4397=1869%:%
%:%4402=1869%:%
%:%4405=1870%:%
%:%4406=1871%:%
%:%4407=1871%:%
%:%4408=1872%:%
%:%4411=1873%:%
%:%4415=1873%:%
%:%4416=1873%:%
%:%4421=1873%:%
%:%4424=1874%:%
%:%4425=1875%:%
%:%4426=1875%:%
%:%4427=1876%:%
%:%4428=1877%:%
%:%4429=1878%:%
%:%4436=1879%:%
%:%4437=1879%:%
%:%4438=1880%:%
%:%4439=1880%:%
%:%4440=1881%:%
%:%4441=1881%:%
%:%4442=1882%:%
%:%4443=1882%:%
%:%4444=1883%:%
%:%4445=1883%:%
%:%4446=1884%:%
%:%4447=1884%:%
%:%4448=1885%:%
%:%4449=1885%:%
%:%4450=1886%:%
%:%4451=1886%:%
%:%4452=1887%:%
%:%4453=1887%:%
%:%4454=1888%:%
%:%4455=1888%:%
%:%4456=1889%:%
%:%4457=1889%:%
%:%4458=1890%:%
%:%4459=1890%:%
%:%4460=1891%:%
%:%4461=1891%:%
%:%4462=1892%:%
%:%4463=1892%:%
%:%4464=1893%:%
%:%4465=1893%:%
%:%4466=1894%:%
%:%4467=1894%:%
%:%4468=1895%:%
%:%4469=1895%:%
%:%4470=1896%:%
%:%4471=1896%:%
%:%4472=1896%:%
%:%4473=1897%:%
%:%4474=1897%:%
%:%4475=1898%:%
%:%4476=1898%:%
%:%4477=1899%:%
%:%4478=1899%:%
%:%4479=1900%:%
%:%4480=1900%:%
%:%4481=1901%:%
%:%4487=1901%:%
%:%4490=1902%:%
%:%4491=1903%:%
%:%4492=1903%:%
%:%4493=1904%:%
%:%4494=1905%:%
%:%4495=1906%:%
%:%4496=1907%:%
%:%4497=1908%:%
%:%4498=1909%:%
%:%4499=1910%:%
%:%4502=1911%:%
%:%4506=1911%:%
%:%4507=1911%:%
%:%4508=1912%:%
%:%4509=1912%:%
%:%4510=1913%:%
%:%4511=1913%:%
%:%4512=1914%:%
%:%4513=1914%:%
%:%4514=1915%:%
%:%4515=1915%:%
%:%4516=1915%:%
%:%4517=1916%:%
%:%4518=1916%:%
%:%4519=1917%:%
%:%4520=1917%:%
%:%4521=1917%:%
%:%4522=1917%:%
%:%4523=1918%:%
%:%4524=1918%:%
%:%4525=1918%:%
%:%4526=1919%:%
%:%4527=1919%:%
%:%4528=1920%:%
%:%4529=1920%:%
%:%4530=1921%:%
%:%4531=1921%:%
%:%4532=1921%:%
%:%4533=1922%:%
%:%4534=1922%:%
%:%4535=1923%:%
%:%4536=1923%:%
%:%4537=1924%:%
%:%4538=1924%:%
%:%4539=1924%:%
%:%4540=1925%:%
%:%4541=1925%:%
%:%4542=1926%:%
%:%4543=1926%:%
%:%4544=1927%:%
%:%4545=1927%:%
%:%4546=1928%:%
%:%4547=1928%:%
%:%4548=1929%:%
%:%4549=1929%:%
%:%4550=1930%:%
%:%4551=1930%:%
%:%4552=1930%:%
%:%4553=1931%:%
%:%4554=1931%:%
%:%4555=1931%:%
%:%4556=1932%:%
%:%4557=1932%:%
%:%4558=1933%:%
%:%4559=1933%:%
%:%4560=1933%:%
%:%4561=1934%:%
%:%4562=1934%:%
%:%4563=1934%:%
%:%4564=1935%:%
%:%4565=1935%:%
%:%4566=1936%:%
%:%4567=1936%:%
%:%4568=1936%:%
%:%4569=1937%:%
%:%4570=1937%:%
%:%4571=1937%:%
%:%4572=1938%:%
%:%4573=1938%:%
%:%4574=1938%:%
%:%4575=1939%:%
%:%4576=1939%:%
%:%4577=1940%:%
%:%4578=1941%:%
%:%4579=1941%:%
%:%4580=1942%:%
%:%4581=1942%:%
%:%4582=1942%:%
%:%4583=1943%:%
%:%4584=1943%:%
%:%4585=1944%:%
%:%4586=1944%:%
%:%4587=1945%:%
%:%4588=1945%:%
%:%4589=1946%:%
%:%4590=1947%:%
%:%4591=1948%:%
%:%4597=1948%:%
%:%4600=1949%:%
%:%4601=1950%:%
%:%4602=1950%:%
%:%4603=1951%:%
%:%4604=1952%:%
%:%4605=1953%:%
%:%4606=1954%:%
%:%4607=1955%:%
%:%4608=1956%:%
%:%4609=1957%:%
%:%4610=1958%:%
%:%4611=1959%:%
%:%4618=1960%:%
%:%4619=1960%:%
%:%4620=1961%:%
%:%4621=1961%:%
%:%4622=1962%:%
%:%4623=1962%:%
%:%4624=1963%:%
%:%4625=1963%:%
%:%4626=1964%:%
%:%4627=1964%:%
%:%4628=1965%:%
%:%4629=1965%:%
%:%4630=1966%:%
%:%4631=1966%:%
%:%4632=1967%:%
%:%4638=1967%:%
%:%4641=1968%:%
%:%4642=1969%:%
%:%4643=1969%:%
%:%4644=1970%:%
%:%4645=1971%:%
%:%4646=1972%:%
%:%4647=1973%:%
%:%4648=1974%:%
%:%4651=1975%:%
%:%4655=1975%:%
%:%4656=1975%:%
%:%4657=1976%:%
%:%4658=1976%:%
%:%4663=1976%:%
%:%4666=1977%:%
%:%4667=1978%:%
%:%4668=1978%:%
%:%4669=1979%:%
%:%4670=1980%:%
%:%4677=1981%:%
%:%4678=1981%:%
%:%4679=1982%:%
%:%4680=1982%:%
%:%4681=1983%:%
%:%4682=1983%:%
%:%4683=1984%:%
%:%4684=1984%:%
%:%4685=1985%:%
%:%4686=1985%:%
%:%4687=1985%:%
%:%4688=1986%:%
%:%4689=1986%:%
%:%4690=1987%:%
%:%4691=1987%:%
%:%4692=1987%:%
%:%4693=1988%:%
%:%4694=1988%:%
%:%4695=1989%:%
%:%4696=1989%:%
%:%4697=1990%:%
%:%4698=1990%:%
%:%4699=1990%:%
%:%4700=1991%:%
%:%4701=1991%:%
%:%4702=1992%:%
%:%4703=1992%:%
%:%4704=1993%:%
%:%4705=1993%:%
%:%4706=1993%:%
%:%4707=1994%:%
%:%4708=1994%:%
%:%4709=1995%:%
%:%4710=1995%:%
%:%4711=1996%:%
%:%4712=1997%:%
%:%4713=1998%:%
%:%4719=1998%:%
%:%4722=1999%:%
%:%4723=2000%:%
%:%4724=2000%:%
%:%4727=2001%:%
%:%4731=2001%:%
%:%4732=2001%:%
%:%4737=2001%:%
%:%4740=2002%:%
%:%4741=2003%:%
%:%4742=2003%:%
%:%4743=2004%:%
%:%4744=2005%:%
%:%4751=2006%:%
%:%4752=2006%:%
%:%4753=2007%:%
%:%4754=2007%:%
%:%4755=2008%:%
%:%4756=2008%:%
%:%4757=2009%:%
%:%4758=2009%:%
%:%4759=2010%:%
%:%4760=2010%:%
%:%4761=2010%:%
%:%4762=2011%:%
%:%4763=2011%:%
%:%4764=2012%:%
%:%4765=2012%:%
%:%4766=2012%:%
%:%4767=2013%:%
%:%4768=2013%:%
%:%4769=2014%:%
%:%4770=2014%:%
%:%4771=2015%:%
%:%4772=2015%:%
%:%4773=2015%:%
%:%4774=2016%:%
%:%4775=2016%:%
%:%4776=2017%:%
%:%4777=2017%:%
%:%4778=2018%:%
%:%4779=2018%:%
%:%4780=2018%:%
%:%4781=2019%:%
%:%4782=2019%:%
%:%4783=2020%:%
%:%4784=2020%:%
%:%4785=2021%:%
%:%4786=2022%:%
%:%4787=2023%:%
%:%4793=2023%:%
%:%4796=2024%:%
%:%4797=2025%:%
%:%4798=2025%:%
%:%4799=2026%:%
%:%4800=2027%:%
%:%4801=2027%:%
%:%4804=2028%:%
%:%4808=2028%:%
%:%4809=2028%:%
%:%4814=2028%:%
%:%4817=2029%:%
%:%4818=2030%:%
%:%4819=2030%:%
%:%4820=2031%:%
%:%4821=2032%:%
%:%4824=2033%:%
%:%4828=2033%:%
%:%4829=2033%:%
%:%4830=2034%:%
%:%4831=2034%:%
%:%4832=2035%:%
%:%4833=2036%:%
%:%4834=2037%:%
%:%4839=2037%:%
%:%4842=2038%:%
%:%4843=2039%:%
%:%4844=2039%:%
%:%4847=2040%:%
%:%4852=2041%:%
%:%4853=2041%:%
%:%4854=2041%:%

%
\begin{isabellebody}%
\setisabellecontext{Error{\isacharunderscore}{\kern0pt}Monad{\isacharunderscore}{\kern0pt}Add}%
%
\isadelimtheory
%
\endisadelimtheory
%
\isatagtheory
\isacommand{theory}\isamarkupfalse%
\ Error{\isacharunderscore}{\kern0pt}Monad{\isacharunderscore}{\kern0pt}Add\isanewline
\isakeyword{imports}\isanewline
\ \ {\isachardoublequoteopen}Certification{\isacharunderscore}{\kern0pt}Monads{\isachardot}{\kern0pt}Check{\isacharunderscore}{\kern0pt}Monad{\isachardoublequoteclose}\isanewline
\ \ {\isachardoublequoteopen}Show{\isachardot}{\kern0pt}Show{\isacharunderscore}{\kern0pt}Instances{\isachardoublequoteclose}\isanewline
\isakeyword{begin}%
\endisatagtheory
{\isafoldtheory}%
%
\isadelimtheory
\isanewline
%
\endisadelimtheory
\ \ \ \ \isanewline
\ \ \isacommand{abbreviation}\isamarkupfalse%
\ {\isachardoublequoteopen}assert{\isacharunderscore}{\kern0pt}opt\ {\isasymPhi}\ {\isasymequiv}\ if\ {\isasymPhi}\ then\ Some\ {\isacharparenleft}{\kern0pt}{\isacharparenright}{\kern0pt}\ else\ None{\isachardoublequoteclose}\ \ \isanewline
\isanewline
\ \ \isacommand{definition}\isamarkupfalse%
\ {\isachardoublequoteopen}lift{\isacharunderscore}{\kern0pt}opt\ m\ e\ {\isasymequiv}\ case\ m\ of\ Some\ x\ {\isasymRightarrow}\ Error{\isacharunderscore}{\kern0pt}Monad{\isachardot}{\kern0pt}return\ x\ {\isacharbar}{\kern0pt}\ None\ {\isasymRightarrow}\ Error{\isacharunderscore}{\kern0pt}Monad{\isachardot}{\kern0pt}error\ e{\isachardoublequoteclose}\isanewline
\ \ \ \ \isanewline
\ \ \isacommand{lemma}\isamarkupfalse%
\ lift{\isacharunderscore}{\kern0pt}opt{\isacharunderscore}{\kern0pt}simps{\isacharbrackleft}{\kern0pt}simp{\isacharbrackright}{\kern0pt}{\isacharcolon}{\kern0pt}\ \isanewline
\ \ \ \ {\isachardoublequoteopen}lift{\isacharunderscore}{\kern0pt}opt\ None\ e\ {\isacharequal}{\kern0pt}\ error\ e{\isachardoublequoteclose}\isanewline
\ \ \ \ {\isachardoublequoteopen}lift{\isacharunderscore}{\kern0pt}opt\ {\isacharparenleft}{\kern0pt}Some\ v{\isacharparenright}{\kern0pt}\ e\ {\isacharequal}{\kern0pt}\ return\ v{\isachardoublequoteclose}\isanewline
%
\isadelimproof
\ \ \ \ %
\endisadelimproof
%
\isatagproof
\isacommand{by}\isamarkupfalse%
\ {\isacharparenleft}{\kern0pt}auto\ simp{\isacharcolon}{\kern0pt}\ lift{\isacharunderscore}{\kern0pt}opt{\isacharunderscore}{\kern0pt}def{\isacharparenright}{\kern0pt}%
\endisatagproof
{\isafoldproof}%
%
\isadelimproof
\isanewline
%
\endisadelimproof
\ \ \isanewline
\ \ \ \ \isanewline
\ \ \isacommand{lemma}\isamarkupfalse%
\ reflcl{\isacharunderscore}{\kern0pt}image{\isacharunderscore}{\kern0pt}iff{\isacharbrackleft}{\kern0pt}simp{\isacharbrackright}{\kern0pt}{\isacharcolon}{\kern0pt}\ {\isachardoublequoteopen}R\isactrlsup {\isacharequal}{\kern0pt}{\isacharbackquote}{\kern0pt}{\isacharbackquote}{\kern0pt}S\ {\isacharequal}{\kern0pt}\ S{\isasymunion}R{\isacharbackquote}{\kern0pt}{\isacharbackquote}{\kern0pt}S{\isachardoublequoteclose}%
\isadelimproof
\ %
\endisadelimproof
%
\isatagproof
\isacommand{by}\isamarkupfalse%
\ blast%
\endisatagproof
{\isafoldproof}%
%
\isadelimproof
%
\endisadelimproof
\ \isanewline
\ \ \ \ \isanewline
\ \ \isacommand{named{\isacharunderscore}{\kern0pt}theorems}\isamarkupfalse%
\ return{\isacharunderscore}{\kern0pt}iff\isanewline
\ \ \ \ \ \ \isanewline
\ \ \isacommand{lemma}\isamarkupfalse%
\ bind{\isacharunderscore}{\kern0pt}return{\isacharunderscore}{\kern0pt}iff{\isacharbrackleft}{\kern0pt}return{\isacharunderscore}{\kern0pt}iff{\isacharbrackright}{\kern0pt}{\isacharcolon}{\kern0pt}\ {\isachardoublequoteopen}Error{\isacharunderscore}{\kern0pt}Monad{\isachardot}{\kern0pt}bind\ m\ f\ {\isacharequal}{\kern0pt}\ Inr\ y\ {\isasymlongleftrightarrow}\ {\isacharparenleft}{\kern0pt}{\isasymexists}x{\isachardot}{\kern0pt}\ m\ {\isacharequal}{\kern0pt}\ Inr\ x\ {\isasymand}\ f\ x\ {\isacharequal}{\kern0pt}\ Inr\ y{\isacharparenright}{\kern0pt}{\isachardoublequoteclose}\isanewline
%
\isadelimproof
\ \ \ \ %
\endisadelimproof
%
\isatagproof
\isacommand{by}\isamarkupfalse%
\ auto%
\endisatagproof
{\isafoldproof}%
%
\isadelimproof
\isanewline
%
\endisadelimproof
\ \ \ \ \isanewline
\ \ \isacommand{lemma}\isamarkupfalse%
\ lift{\isacharunderscore}{\kern0pt}opt{\isacharunderscore}{\kern0pt}return{\isacharunderscore}{\kern0pt}iff{\isacharbrackleft}{\kern0pt}return{\isacharunderscore}{\kern0pt}iff{\isacharbrackright}{\kern0pt}{\isacharcolon}{\kern0pt}\ {\isachardoublequoteopen}lift{\isacharunderscore}{\kern0pt}opt\ m\ e\ {\isacharequal}{\kern0pt}\ Inr\ x\ {\isasymlongleftrightarrow}\ m{\isacharequal}{\kern0pt}Some\ x{\isachardoublequoteclose}\isanewline
%
\isadelimproof
\ \ \ \ %
\endisadelimproof
%
\isatagproof
\isacommand{unfolding}\isamarkupfalse%
\ lift{\isacharunderscore}{\kern0pt}opt{\isacharunderscore}{\kern0pt}def\ \isacommand{by}\isamarkupfalse%
\ {\isacharparenleft}{\kern0pt}auto\ split{\isacharcolon}{\kern0pt}\ option{\isachardot}{\kern0pt}split{\isacharparenright}{\kern0pt}%
\endisatagproof
{\isafoldproof}%
%
\isadelimproof
\isanewline
%
\endisadelimproof
\ \ \ \ \ \ \isanewline
\ \ \isacommand{lemma}\isamarkupfalse%
\ check{\isacharunderscore}{\kern0pt}return{\isacharunderscore}{\kern0pt}iff{\isacharbrackleft}{\kern0pt}return{\isacharunderscore}{\kern0pt}iff{\isacharbrackright}{\kern0pt}{\isacharcolon}{\kern0pt}\ {\isachardoublequoteopen}check\ {\isasymPhi}\ e\ {\isacharequal}{\kern0pt}\ Inr\ uu\ {\isasymlongleftrightarrow}\ {\isasymPhi}{\isachardoublequoteclose}\ \ \ \ \isanewline
%
\isadelimproof
\ \ \ \ %
\endisadelimproof
%
\isatagproof
\isacommand{by}\isamarkupfalse%
\ {\isacharparenleft}{\kern0pt}auto\ simp{\isacharcolon}{\kern0pt}\ check{\isacharunderscore}{\kern0pt}def{\isacharparenright}{\kern0pt}%
\endisatagproof
{\isafoldproof}%
%
\isadelimproof
\isanewline
%
\endisadelimproof
\ \ \isanewline
\ \ \isanewline
\ \ \isacommand{lemma}\isamarkupfalse%
\ check{\isacharunderscore}{\kern0pt}simps{\isacharbrackleft}{\kern0pt}simp{\isacharbrackright}{\kern0pt}{\isacharcolon}{\kern0pt}\ \ \isanewline
\ \ \ \ {\isachardoublequoteopen}check\ True\ e\ {\isacharequal}{\kern0pt}\ succeed{\isachardoublequoteclose}\isanewline
\ \ \ \ {\isachardoublequoteopen}check\ False\ e\ {\isacharequal}{\kern0pt}\ error\ e{\isachardoublequoteclose}\isanewline
%
\isadelimproof
\ \ \ \ %
\endisadelimproof
%
\isatagproof
\isacommand{by}\isamarkupfalse%
\ {\isacharparenleft}{\kern0pt}auto\ simp{\isacharcolon}{\kern0pt}\ check{\isacharunderscore}{\kern0pt}def{\isacharparenright}{\kern0pt}%
\endisatagproof
{\isafoldproof}%
%
\isadelimproof
\isanewline
%
\endisadelimproof
\ \ \ \ \ \ \ \ \isanewline
\ \ \isacommand{lemma}\isamarkupfalse%
\ Let{\isacharunderscore}{\kern0pt}return{\isacharunderscore}{\kern0pt}iff{\isacharbrackleft}{\kern0pt}return{\isacharunderscore}{\kern0pt}iff{\isacharbrackright}{\kern0pt}{\isacharcolon}{\kern0pt}\ {\isachardoublequoteopen}{\isacharparenleft}{\kern0pt}let\ x{\isacharequal}{\kern0pt}v\ in\ f\ x{\isacharparenright}{\kern0pt}\ {\isacharequal}{\kern0pt}\ Inr\ w\ {\isasymlongleftrightarrow}\ f\ v\ {\isacharequal}{\kern0pt}\ Inr\ w{\isachardoublequoteclose}%
\isadelimproof
\ %
\endisadelimproof
%
\isatagproof
\isacommand{by}\isamarkupfalse%
\ simp%
\endisatagproof
{\isafoldproof}%
%
\isadelimproof
%
\endisadelimproof
\isanewline
\isanewline
\ \ \isanewline
\ \ \isacommand{abbreviation}\isamarkupfalse%
\ ERR\ {\isacharcolon}{\kern0pt}{\isacharcolon}{\kern0pt}\ {\isachardoublequoteopen}shows\ {\isasymRightarrow}\ {\isacharparenleft}{\kern0pt}unit\ {\isasymRightarrow}\ shows{\isacharparenright}{\kern0pt}{\isachardoublequoteclose}\ \isakeyword{where}\ {\isachardoublequoteopen}ERR\ s\ {\isasymequiv}\ {\isasymlambda}{\isacharunderscore}{\kern0pt}{\isachardot}{\kern0pt}\ s{\isachardoublequoteclose}\isanewline
\ \ \isacommand{abbreviation}\isamarkupfalse%
\ ERRS\ {\isacharcolon}{\kern0pt}{\isacharcolon}{\kern0pt}\ {\isachardoublequoteopen}string\ {\isasymRightarrow}\ {\isacharparenleft}{\kern0pt}unit\ {\isasymRightarrow}\ shows{\isacharparenright}{\kern0pt}{\isachardoublequoteclose}\ \isakeyword{where}\ {\isachardoublequoteopen}ERRS\ s\ {\isasymequiv}\ ERR\ {\isacharparenleft}{\kern0pt}shows\ s{\isacharparenright}{\kern0pt}{\isachardoublequoteclose}\isanewline
\ \ \isanewline
\ \ \isanewline
\ \ \isacommand{lemma}\isamarkupfalse%
\ error{\isacharunderscore}{\kern0pt}monad{\isacharunderscore}{\kern0pt}bind{\isacharunderscore}{\kern0pt}split{\isacharcolon}{\kern0pt}\ {\isachardoublequoteopen}P\ {\isacharparenleft}{\kern0pt}bind\ m\ f{\isacharparenright}{\kern0pt}\ {\isasymlongleftrightarrow}\ {\isacharparenleft}{\kern0pt}{\isasymforall}v{\isachardot}{\kern0pt}\ m\ {\isacharequal}{\kern0pt}\ Inl\ v\ {\isasymlongrightarrow}\ P\ {\isacharparenleft}{\kern0pt}Inl\ v{\isacharparenright}{\kern0pt}{\isacharparenright}{\kern0pt}\ {\isasymand}\ {\isacharparenleft}{\kern0pt}{\isasymforall}v{\isachardot}{\kern0pt}\ m\ {\isacharequal}{\kern0pt}\ Inr\ v\ {\isasymlongrightarrow}\ P\ {\isacharparenleft}{\kern0pt}f\ v{\isacharparenright}{\kern0pt}{\isacharparenright}{\kern0pt}{\isachardoublequoteclose}\isanewline
%
\isadelimproof
\ \ \ \ %
\endisadelimproof
%
\isatagproof
\isacommand{by}\isamarkupfalse%
\ {\isacharparenleft}{\kern0pt}cases\ m{\isacharparenright}{\kern0pt}\ auto%
\endisatagproof
{\isafoldproof}%
%
\isadelimproof
\isanewline
%
\endisadelimproof
\ \ \isanewline
\ \ \isacommand{lemma}\isamarkupfalse%
\ error{\isacharunderscore}{\kern0pt}monad{\isacharunderscore}{\kern0pt}bind{\isacharunderscore}{\kern0pt}split{\isacharunderscore}{\kern0pt}asm{\isacharcolon}{\kern0pt}\ {\isachardoublequoteopen}P\ {\isacharparenleft}{\kern0pt}bind\ m\ f{\isacharparenright}{\kern0pt}\ {\isasymlongleftrightarrow}\ {\isasymnot}\ {\isacharparenleft}{\kern0pt}{\isasymexists}x{\isachardot}{\kern0pt}\ m\ {\isacharequal}{\kern0pt}\ Inl\ x\ {\isasymand}\ {\isasymnot}\ P\ {\isacharparenleft}{\kern0pt}Inl\ x{\isacharparenright}{\kern0pt}\ {\isasymor}\ {\isacharparenleft}{\kern0pt}{\isasymexists}x{\isachardot}{\kern0pt}\ m\ {\isacharequal}{\kern0pt}\ Inr\ x\ {\isasymand}\ {\isasymnot}\ P\ {\isacharparenleft}{\kern0pt}f\ x{\isacharparenright}{\kern0pt}{\isacharparenright}{\kern0pt}{\isacharparenright}{\kern0pt}{\isachardoublequoteclose}\isanewline
%
\isadelimproof
\ \ \ \ %
\endisadelimproof
%
\isatagproof
\isacommand{by}\isamarkupfalse%
\ {\isacharparenleft}{\kern0pt}cases\ m{\isacharparenright}{\kern0pt}\ auto%
\endisatagproof
{\isafoldproof}%
%
\isadelimproof
\isanewline
%
\endisadelimproof
\ \ \isanewline
\ \ \isacommand{lemmas}\isamarkupfalse%
\ error{\isacharunderscore}{\kern0pt}monad{\isacharunderscore}{\kern0pt}bind{\isacharunderscore}{\kern0pt}splits\ {\isacharequal}{\kern0pt}error{\isacharunderscore}{\kern0pt}monad{\isacharunderscore}{\kern0pt}bind{\isacharunderscore}{\kern0pt}split\ error{\isacharunderscore}{\kern0pt}monad{\isacharunderscore}{\kern0pt}bind{\isacharunderscore}{\kern0pt}split{\isacharunderscore}{\kern0pt}asm\isanewline
\ \ \isanewline
%
\isadelimtheory
\ \ \isanewline
%
\endisadelimtheory
%
\isatagtheory
\isacommand{end}\isamarkupfalse%
%
\endisatagtheory
{\isafoldtheory}%
%
\isadelimtheory
%
\endisadelimtheory
%
\end{isabellebody}%
\endinput
%:%file=Error_Monad_Add.tex%:%
%:%10=1%:%
%:%11=1%:%
%:%12=2%:%
%:%13=3%:%
%:%14=4%:%
%:%15=5%:%
%:%20=5%:%
%:%23=6%:%
%:%24=7%:%
%:%25=7%:%
%:%26=8%:%
%:%27=9%:%
%:%28=9%:%
%:%29=10%:%
%:%30=11%:%
%:%31=11%:%
%:%32=12%:%
%:%33=13%:%
%:%36=14%:%
%:%40=14%:%
%:%41=14%:%
%:%46=14%:%
%:%49=15%:%
%:%50=16%:%
%:%51=17%:%
%:%52=17%:%
%:%54=17%:%
%:%58=17%:%
%:%59=17%:%
%:%66=17%:%
%:%67=18%:%
%:%68=19%:%
%:%69=19%:%
%:%70=20%:%
%:%71=21%:%
%:%72=21%:%
%:%75=22%:%
%:%79=22%:%
%:%80=22%:%
%:%85=22%:%
%:%88=23%:%
%:%89=24%:%
%:%90=24%:%
%:%93=25%:%
%:%97=25%:%
%:%98=25%:%
%:%99=25%:%
%:%104=25%:%
%:%107=26%:%
%:%108=27%:%
%:%109=27%:%
%:%112=28%:%
%:%116=28%:%
%:%117=28%:%
%:%122=28%:%
%:%125=29%:%
%:%126=30%:%
%:%127=31%:%
%:%128=31%:%
%:%129=32%:%
%:%130=33%:%
%:%133=34%:%
%:%137=34%:%
%:%138=34%:%
%:%143=34%:%
%:%146=35%:%
%:%147=36%:%
%:%148=36%:%
%:%150=36%:%
%:%154=36%:%
%:%155=36%:%
%:%162=36%:%
%:%163=37%:%
%:%164=38%:%
%:%165=39%:%
%:%166=39%:%
%:%167=40%:%
%:%168=40%:%
%:%169=41%:%
%:%170=42%:%
%:%171=43%:%
%:%172=43%:%
%:%175=44%:%
%:%179=44%:%
%:%180=44%:%
%:%185=44%:%
%:%188=45%:%
%:%189=46%:%
%:%190=46%:%
%:%193=47%:%
%:%197=47%:%
%:%198=47%:%
%:%203=47%:%
%:%206=48%:%
%:%207=49%:%
%:%208=49%:%
%:%209=50%:%
%:%212=51%:%
%:%217=52%:%

%
\begin{isabellebody}%
\setisabellecontext{TEMPORAL{\isacharunderscore}{\kern0pt}PDDL{\isacharunderscore}{\kern0pt}Checker}%
%
\isadelimdocument
%
\endisadelimdocument
%
\isatagdocument
%
\isamarkupsection{Executable Temporal PDDL Checker%
}
\isamarkuptrue%
%
\endisatagdocument
{\isafolddocument}%
%
\isadelimdocument
%
\endisadelimdocument
%
\isadelimtheory
%
\endisadelimtheory
%
\isatagtheory
\isacommand{theory}\isamarkupfalse%
\ TEMPORAL{\isacharunderscore}{\kern0pt}PDDL{\isacharunderscore}{\kern0pt}Checker\isanewline
\isakeyword{imports}\isanewline
\ \ TEMPORAL{\isacharunderscore}{\kern0pt}PDDL{\isacharunderscore}{\kern0pt}Semantics\isanewline
\ \ Error{\isacharunderscore}{\kern0pt}Monad{\isacharunderscore}{\kern0pt}Add\isanewline
\ \ {\isachardoublequoteopen}HOL{\isachardot}{\kern0pt}String{\isachardoublequoteclose}\isanewline
\ \ \isanewline
\ \ {\isachardoublequoteopen}HOL{\isacharminus}{\kern0pt}Library{\isachardot}{\kern0pt}Code{\isacharunderscore}{\kern0pt}Target{\isacharunderscore}{\kern0pt}Nat{\isachardoublequoteclose}\isanewline
\ \ {\isachardoublequoteopen}HOL{\isacharminus}{\kern0pt}Library{\isachardot}{\kern0pt}While{\isacharunderscore}{\kern0pt}Combinator{\isachardoublequoteclose}\isanewline
\ \ {\isachardoublequoteopen}Containers{\isachardot}{\kern0pt}Containers{\isachardoublequoteclose}\isanewline
\isanewline
\ \ {\isachardoublequoteopen}HOL{\isacharminus}{\kern0pt}Library{\isachardot}{\kern0pt}Tree{\isachardoublequoteclose}\isanewline
\isakeyword{begin}%
\endisatagtheory
{\isafoldtheory}%
%
\isadelimtheory
%
\endisadelimtheory
%
\isadelimdocument
%
\endisadelimdocument
%
\isatagdocument
%
\isamarkupsubsection{Generic DFS Reachability Checker%
}
\isamarkuptrue%
%
\endisatagdocument
{\isafolddocument}%
%
\isadelimdocument
%
\endisadelimdocument
%
\begin{isamarkuptext}%
Used for subtype checks%
\end{isamarkuptext}\isamarkuptrue%
\isacommand{definition}\isamarkupfalse%
\ {\isachardoublequoteopen}E{\isacharunderscore}{\kern0pt}of{\isacharunderscore}{\kern0pt}succ\ succ\ {\isasymequiv}\ {\isacharbraceleft}{\kern0pt}\ {\isacharparenleft}{\kern0pt}u{\isacharcomma}{\kern0pt}v{\isacharparenright}{\kern0pt}{\isachardot}{\kern0pt}\ v{\isasymin}set\ {\isacharparenleft}{\kern0pt}succ\ u{\isacharparenright}{\kern0pt}\ {\isacharbraceright}{\kern0pt}{\isachardoublequoteclose}\isanewline
\isacommand{lemma}\isamarkupfalse%
\ succ{\isacharunderscore}{\kern0pt}as{\isacharunderscore}{\kern0pt}E{\isacharcolon}{\kern0pt}\ {\isachardoublequoteopen}set\ {\isacharparenleft}{\kern0pt}succ\ x{\isacharparenright}{\kern0pt}\ {\isacharequal}{\kern0pt}\ E{\isacharunderscore}{\kern0pt}of{\isacharunderscore}{\kern0pt}succ\ succ\ {\isacharbackquote}{\kern0pt}{\isacharbackquote}{\kern0pt}\ {\isacharbraceleft}{\kern0pt}x{\isacharbraceright}{\kern0pt}{\isachardoublequoteclose}\isanewline
%
\isadelimproof
\ \ %
\endisadelimproof
%
\isatagproof
\isacommand{unfolding}\isamarkupfalse%
\ E{\isacharunderscore}{\kern0pt}of{\isacharunderscore}{\kern0pt}succ{\isacharunderscore}{\kern0pt}def\ \isacommand{by}\isamarkupfalse%
\ auto%
\endisatagproof
{\isafoldproof}%
%
\isadelimproof
\isanewline
%
\endisadelimproof
\isanewline
\isacommand{context}\isamarkupfalse%
\isanewline
\ \ \isakeyword{fixes}\ succ\ {\isacharcolon}{\kern0pt}{\isacharcolon}{\kern0pt}\ {\isachardoublequoteopen}{\isacharprime}{\kern0pt}a\ {\isasymRightarrow}\ {\isacharprime}{\kern0pt}a\ list{\isachardoublequoteclose}\isanewline
\isakeyword{begin}\isanewline
\isanewline
\ \ \isakeyword{private}\ \isacommand{abbreviation}\isamarkupfalse%
\ {\isacharparenleft}{\kern0pt}input{\isacharparenright}{\kern0pt}\ {\isachardoublequoteopen}E\ {\isasymequiv}\ E{\isacharunderscore}{\kern0pt}of{\isacharunderscore}{\kern0pt}succ\ succ{\isachardoublequoteclose}\isanewline
\isanewline
\isanewline
\isacommand{definition}\isamarkupfalse%
\ {\isachardoublequoteopen}dfs{\isacharunderscore}{\kern0pt}reachable\ D\ w\ {\isasymequiv}\isanewline
\ \ let\ {\isacharparenleft}{\kern0pt}V{\isacharcomma}{\kern0pt}w{\isacharcomma}{\kern0pt}brk{\isacharparenright}{\kern0pt}\ {\isacharequal}{\kern0pt}\ while\ {\isacharparenleft}{\kern0pt}{\isasymlambda}{\isacharparenleft}{\kern0pt}V{\isacharcomma}{\kern0pt}w{\isacharcomma}{\kern0pt}brk{\isacharparenright}{\kern0pt}{\isachardot}{\kern0pt}\ {\isasymnot}brk\ {\isasymand}\ w{\isasymnoteq}{\isacharbrackleft}{\kern0pt}{\isacharbrackright}{\kern0pt}{\isacharparenright}{\kern0pt}\ {\isacharparenleft}{\kern0pt}{\isasymlambda}{\isacharparenleft}{\kern0pt}V{\isacharcomma}{\kern0pt}w{\isacharcomma}{\kern0pt}{\isacharunderscore}{\kern0pt}{\isacharparenright}{\kern0pt}{\isachardot}{\kern0pt}\isanewline
\ \ \ \ case\ w\ of\ v{\isacharhash}{\kern0pt}w\ {\isasymRightarrow}\isanewline
\ \ \ \ if\ D\ v\ then\ {\isacharparenleft}{\kern0pt}V{\isacharcomma}{\kern0pt}v{\isacharhash}{\kern0pt}w{\isacharcomma}{\kern0pt}True{\isacharparenright}{\kern0pt}\isanewline
\ \ \ \ else\ if\ v{\isasymin}V\ then\ {\isacharparenleft}{\kern0pt}V{\isacharcomma}{\kern0pt}w{\isacharcomma}{\kern0pt}False{\isacharparenright}{\kern0pt}\isanewline
\ \ \ \ else\isanewline
\ \ \ \ \ \ let\ V\ {\isacharequal}{\kern0pt}\ insert\ v\ V\ in\isanewline
\ \ \ \ \ \ let\ w\ {\isacharequal}{\kern0pt}\ succ\ v\ {\isacharat}{\kern0pt}\ w\ in\isanewline
\ \ \ \ \ \ {\isacharparenleft}{\kern0pt}V{\isacharcomma}{\kern0pt}w{\isacharcomma}{\kern0pt}False{\isacharparenright}{\kern0pt}\isanewline
\ \ \ \ {\isacharparenright}{\kern0pt}\ {\isacharparenleft}{\kern0pt}{\isacharbraceleft}{\kern0pt}{\isacharbraceright}{\kern0pt}{\isacharcomma}{\kern0pt}w{\isacharcomma}{\kern0pt}False{\isacharparenright}{\kern0pt}\isanewline
\ \ in\ brk{\isachardoublequoteclose}\isanewline
\isanewline
\isanewline
\isacommand{context}\isamarkupfalse%
\isanewline
\ \ \isakeyword{fixes}\ w\isactrlsub {\isadigit{0}}\ {\isacharcolon}{\kern0pt}{\isacharcolon}{\kern0pt}\ {\isachardoublequoteopen}{\isacharprime}{\kern0pt}a\ list{\isachardoublequoteclose}\isanewline
\ \ \isakeyword{assumes}\ finite{\isacharunderscore}{\kern0pt}dfs{\isacharunderscore}{\kern0pt}reachable{\isacharbrackleft}{\kern0pt}simp{\isacharcomma}{\kern0pt}\ intro{\isacharbang}{\kern0pt}{\isacharbrackright}{\kern0pt}{\isacharcolon}{\kern0pt}\ {\isachardoublequoteopen}finite\ {\isacharparenleft}{\kern0pt}E\isactrlsup {\isacharasterisk}{\kern0pt}\ {\isacharbackquote}{\kern0pt}{\isacharbackquote}{\kern0pt}\ set\ w\isactrlsub {\isadigit{0}}{\isacharparenright}{\kern0pt}{\isachardoublequoteclose}\isanewline
\isakeyword{begin}\isanewline
\isanewline
\ \ \isakeyword{private}\ \isacommand{abbreviation}\isamarkupfalse%
\ {\isacharparenleft}{\kern0pt}input{\isacharparenright}{\kern0pt}\ {\isachardoublequoteopen}W\isactrlsub {\isadigit{0}}\ {\isasymequiv}\ set\ w\isactrlsub {\isadigit{0}}{\isachardoublequoteclose}\isanewline
\isanewline
\isacommand{definition}\isamarkupfalse%
\ {\isachardoublequoteopen}dfs{\isacharunderscore}{\kern0pt}reachable{\isacharunderscore}{\kern0pt}invar\ D\ V\ W\ brk\ {\isasymlongleftrightarrow}\isanewline
\ \ \ \ W\isactrlsub {\isadigit{0}}\ {\isasymsubseteq}\ W\ {\isasymunion}\ V\isanewline
\ \ {\isasymand}\ W\ {\isasymunion}\ V\ {\isasymsubseteq}\ E\isactrlsup {\isacharasterisk}{\kern0pt}\ {\isacharbackquote}{\kern0pt}{\isacharbackquote}{\kern0pt}\ W\isactrlsub {\isadigit{0}}\isanewline
\ \ {\isasymand}\ E{\isacharbackquote}{\kern0pt}{\isacharbackquote}{\kern0pt}V\ {\isasymsubseteq}\ W\ {\isasymunion}\ V\isanewline
\ \ {\isasymand}\ Collect\ D\ {\isasyminter}\ V\ {\isacharequal}{\kern0pt}\ {\isacharbraceleft}{\kern0pt}{\isacharbraceright}{\kern0pt}\isanewline
\ \ {\isasymand}\ {\isacharparenleft}{\kern0pt}brk\ {\isasymlongrightarrow}\ Collect\ D\ {\isasyminter}\ E\isactrlsup {\isacharasterisk}{\kern0pt}\ {\isacharbackquote}{\kern0pt}{\isacharbackquote}{\kern0pt}\ W\isactrlsub {\isadigit{0}}\ {\isasymnoteq}\ {\isacharbraceleft}{\kern0pt}{\isacharbraceright}{\kern0pt}{\isacharparenright}{\kern0pt}{\isachardoublequoteclose}\isanewline
\isanewline
\isacommand{lemma}\isamarkupfalse%
\ card{\isacharunderscore}{\kern0pt}decreases{\isacharcolon}{\kern0pt}\ {\isachardoublequoteopen}\isanewline
\ \ \ {\isasymlbrakk}finite\ V{\isacharsemicolon}{\kern0pt}\ y\ {\isasymnotin}\ V{\isacharsemicolon}{\kern0pt}\ dfs{\isacharunderscore}{\kern0pt}reachable{\isacharunderscore}{\kern0pt}invar\ D\ V\ {\isacharparenleft}{\kern0pt}Set{\isachardot}{\kern0pt}insert\ y\ W{\isacharparenright}{\kern0pt}\ brk\ {\isasymrbrakk}\isanewline
\ \ \ {\isasymLongrightarrow}\ card\ {\isacharparenleft}{\kern0pt}E\isactrlsup {\isacharasterisk}{\kern0pt}\ {\isacharbackquote}{\kern0pt}{\isacharbackquote}{\kern0pt}\ W\isactrlsub {\isadigit{0}}\ {\isacharminus}{\kern0pt}\ Set{\isachardot}{\kern0pt}insert\ y\ V{\isacharparenright}{\kern0pt}\ {\isacharless}{\kern0pt}\ card\ {\isacharparenleft}{\kern0pt}E\isactrlsup {\isacharasterisk}{\kern0pt}\ {\isacharbackquote}{\kern0pt}{\isacharbackquote}{\kern0pt}\ W\isactrlsub {\isadigit{0}}\ {\isacharminus}{\kern0pt}\ V{\isacharparenright}{\kern0pt}{\isachardoublequoteclose}\isanewline
%
\isadelimproof
\ \ %
\endisadelimproof
%
\isatagproof
\isacommand{apply}\isamarkupfalse%
\ {\isacharparenleft}{\kern0pt}rule\ psubset{\isacharunderscore}{\kern0pt}card{\isacharunderscore}{\kern0pt}mono{\isacharparenright}{\kern0pt}\isanewline
\ \ \isacommand{apply}\isamarkupfalse%
\ {\isacharparenleft}{\kern0pt}auto\ simp{\isacharcolon}{\kern0pt}\ dfs{\isacharunderscore}{\kern0pt}reachable{\isacharunderscore}{\kern0pt}invar{\isacharunderscore}{\kern0pt}def{\isacharparenright}{\kern0pt}\isanewline
\ \ \isacommand{done}\isamarkupfalse%
%
\endisatagproof
{\isafoldproof}%
%
\isadelimproof
\isanewline
%
\endisadelimproof
\isanewline
\isacommand{lemma}\isamarkupfalse%
\ all{\isacharunderscore}{\kern0pt}neq{\isacharunderscore}{\kern0pt}Cons{\isacharunderscore}{\kern0pt}is{\isacharunderscore}{\kern0pt}Nil{\isacharbrackleft}{\kern0pt}simp{\isacharbrackright}{\kern0pt}{\isacharcolon}{\kern0pt}\ \isanewline
\ \ {\isachardoublequoteopen}{\isacharparenleft}{\kern0pt}{\isasymforall}y\ ys{\isachardot}{\kern0pt}\ x{\isadigit{2}}\ {\isasymnoteq}\ y\ {\isacharhash}{\kern0pt}\ ys{\isacharparenright}{\kern0pt}\ {\isasymlongleftrightarrow}\ x{\isadigit{2}}\ {\isacharequal}{\kern0pt}\ {\isacharbrackleft}{\kern0pt}{\isacharbrackright}{\kern0pt}{\isachardoublequoteclose}%
\isadelimproof
\ %
\endisadelimproof
%
\isatagproof
\isacommand{by}\isamarkupfalse%
\ {\isacharparenleft}{\kern0pt}cases\ x{\isadigit{2}}{\isacharparenright}{\kern0pt}\ auto%
\endisatagproof
{\isafoldproof}%
%
\isadelimproof
%
\endisadelimproof
\isanewline
\isanewline
\isacommand{lemma}\isamarkupfalse%
\ dfs{\isacharunderscore}{\kern0pt}reachable{\isacharunderscore}{\kern0pt}correct{\isacharcolon}{\kern0pt}\ {\isachardoublequoteopen}dfs{\isacharunderscore}{\kern0pt}reachable\ D\ w\isactrlsub {\isadigit{0}}\ {\isasymlongleftrightarrow}\ Collect\ D\ {\isasyminter}\ E\isactrlsup {\isacharasterisk}{\kern0pt}\ {\isacharbackquote}{\kern0pt}{\isacharbackquote}{\kern0pt}\ set\ w\isactrlsub {\isadigit{0}}\ {\isasymnoteq}\ {\isacharbraceleft}{\kern0pt}{\isacharbraceright}{\kern0pt}{\isachardoublequoteclose}\isanewline
%
\isadelimproof
\ \ %
\endisadelimproof
%
\isatagproof
\isacommand{unfolding}\isamarkupfalse%
\ dfs{\isacharunderscore}{\kern0pt}reachable{\isacharunderscore}{\kern0pt}def\isanewline
\ \ \isacommand{apply}\isamarkupfalse%
\ {\isacharparenleft}{\kern0pt}rule\ while{\isacharunderscore}{\kern0pt}rule{\isacharbrackleft}{\kern0pt}\isakeyword{where}\isanewline
\ \ \ \ P{\isacharequal}{\kern0pt}{\isachardoublequoteopen}{\isasymlambda}{\isacharparenleft}{\kern0pt}V{\isacharcomma}{\kern0pt}w{\isacharcomma}{\kern0pt}brk{\isacharparenright}{\kern0pt}{\isachardot}{\kern0pt}\ dfs{\isacharunderscore}{\kern0pt}reachable{\isacharunderscore}{\kern0pt}invar\ D\ V\ {\isacharparenleft}{\kern0pt}set\ w{\isacharparenright}{\kern0pt}\ brk\ {\isasymand}\ finite\ V{\isachardoublequoteclose}\isanewline
\ \ \ \ \isakeyword{and}\ r{\isacharequal}{\kern0pt}{\isachardoublequoteopen}measure\ {\isacharparenleft}{\kern0pt}{\isasymlambda}V{\isachardot}{\kern0pt}\ card\ {\isacharparenleft}{\kern0pt}E\isactrlsup {\isacharasterisk}{\kern0pt}\ {\isacharbackquote}{\kern0pt}{\isacharbackquote}{\kern0pt}\ {\isacharparenleft}{\kern0pt}set\ w\isactrlsub {\isadigit{0}}{\isacharparenright}{\kern0pt}\ {\isacharminus}{\kern0pt}\ V{\isacharparenright}{\kern0pt}{\isacharparenright}{\kern0pt}\ {\isacharless}{\kern0pt}{\isacharasterisk}{\kern0pt}lex{\isacharasterisk}{\kern0pt}{\isachargreater}{\kern0pt}\ measure\ length\ {\isacharless}{\kern0pt}{\isacharasterisk}{\kern0pt}lex{\isacharasterisk}{\kern0pt}{\isachargreater}{\kern0pt}\ measure\ {\isacharparenleft}{\kern0pt}{\isasymlambda}True{\isasymRightarrow}{\isadigit{0}}\ {\isacharbar}{\kern0pt}\ False{\isasymRightarrow}{\isadigit{1}}{\isacharparenright}{\kern0pt}{\isachardoublequoteclose}\isanewline
\ \ \ \ {\isacharbrackright}{\kern0pt}{\isacharparenright}{\kern0pt}\isanewline
\ \ \isacommand{subgoal}\isamarkupfalse%
\ \isacommand{by}\isamarkupfalse%
\ {\isacharparenleft}{\kern0pt}auto\ simp{\isacharcolon}{\kern0pt}\ dfs{\isacharunderscore}{\kern0pt}reachable{\isacharunderscore}{\kern0pt}invar{\isacharunderscore}{\kern0pt}def{\isacharparenright}{\kern0pt}\isanewline
\ \ \isacommand{subgoal}\isamarkupfalse%
\isanewline
\ \ \ \ \isacommand{apply}\isamarkupfalse%
\ {\isacharparenleft}{\kern0pt}auto\ simp{\isacharcolon}{\kern0pt}\ neq{\isacharunderscore}{\kern0pt}Nil{\isacharunderscore}{\kern0pt}conv\ succ{\isacharunderscore}{\kern0pt}as{\isacharunderscore}{\kern0pt}E{\isacharbrackleft}{\kern0pt}of\ succ{\isacharbrackright}{\kern0pt}\ split{\isacharcolon}{\kern0pt}\ if{\isacharunderscore}{\kern0pt}splits{\isacharparenright}{\kern0pt}\isanewline
\ \ \ \ \isacommand{by}\isamarkupfalse%
\ {\isacharparenleft}{\kern0pt}auto\ simp{\isacharcolon}{\kern0pt}\ dfs{\isacharunderscore}{\kern0pt}reachable{\isacharunderscore}{\kern0pt}invar{\isacharunderscore}{\kern0pt}def\ Image{\isacharunderscore}{\kern0pt}iff\ intro{\isacharcolon}{\kern0pt}\ rtrancl{\isachardot}{\kern0pt}rtrancl{\isacharunderscore}{\kern0pt}into{\isacharunderscore}{\kern0pt}rtrancl{\isacharparenright}{\kern0pt}\isanewline
\ \ \isacommand{subgoal}\isamarkupfalse%
\ \isacommand{by}\isamarkupfalse%
\ {\isacharparenleft}{\kern0pt}fastforce\ simp{\isacharcolon}{\kern0pt}\ dfs{\isacharunderscore}{\kern0pt}reachable{\isacharunderscore}{\kern0pt}invar{\isacharunderscore}{\kern0pt}def\ dest{\isacharcolon}{\kern0pt}\ Image{\isacharunderscore}{\kern0pt}closed{\isacharunderscore}{\kern0pt}trancl{\isacharparenright}{\kern0pt}\isanewline
\ \ \isacommand{subgoal}\isamarkupfalse%
\ \isacommand{by}\isamarkupfalse%
\ blast\isanewline
\ \ \isacommand{subgoal}\isamarkupfalse%
\ \isacommand{by}\isamarkupfalse%
\ {\isacharparenleft}{\kern0pt}auto\ simp{\isacharcolon}{\kern0pt}\ neq{\isacharunderscore}{\kern0pt}Nil{\isacharunderscore}{\kern0pt}conv\ card{\isacharunderscore}{\kern0pt}decreases{\isacharparenright}{\kern0pt}\isanewline
\ \ \isacommand{done}\isamarkupfalse%
%
\endisatagproof
{\isafoldproof}%
%
\isadelimproof
\isanewline
%
\endisadelimproof
\isanewline
\isacommand{end}\isamarkupfalse%
\isanewline
\isanewline
\isacommand{definition}\isamarkupfalse%
\ {\isachardoublequoteopen}tab{\isacharunderscore}{\kern0pt}succ\ l\ {\isasymequiv}\ Mapping{\isachardot}{\kern0pt}lookup{\isacharunderscore}{\kern0pt}default\ {\isacharbrackleft}{\kern0pt}{\isacharbrackright}{\kern0pt}\ {\isacharparenleft}{\kern0pt}fold\ {\isacharparenleft}{\kern0pt}{\isasymlambda}{\isacharparenleft}{\kern0pt}u{\isacharcomma}{\kern0pt}v{\isacharparenright}{\kern0pt}{\isachardot}{\kern0pt}\ Mapping{\isachardot}{\kern0pt}map{\isacharunderscore}{\kern0pt}default\ u\ {\isacharbrackleft}{\kern0pt}{\isacharbrackright}{\kern0pt}\ {\isacharparenleft}{\kern0pt}Cons\ v{\isacharparenright}{\kern0pt}{\isacharparenright}{\kern0pt}\ l\ Mapping{\isachardot}{\kern0pt}empty{\isacharparenright}{\kern0pt}{\isachardoublequoteclose}\isanewline
\isanewline
\isanewline
\isacommand{lemma}\isamarkupfalse%
\ Some{\isacharunderscore}{\kern0pt}eq{\isacharunderscore}{\kern0pt}map{\isacharunderscore}{\kern0pt}option\ {\isacharbrackleft}{\kern0pt}iff{\isacharbrackright}{\kern0pt}{\isacharcolon}{\kern0pt}\ {\isachardoublequoteopen}{\isacharparenleft}{\kern0pt}Some\ y\ {\isacharequal}{\kern0pt}\ map{\isacharunderscore}{\kern0pt}option\ f\ xo{\isacharparenright}{\kern0pt}\ {\isacharequal}{\kern0pt}\ {\isacharparenleft}{\kern0pt}{\isasymexists}z{\isachardot}{\kern0pt}\ xo\ {\isacharequal}{\kern0pt}\ Some\ z\ {\isasymand}\ f\ z\ {\isacharequal}{\kern0pt}\ y{\isacharparenright}{\kern0pt}{\isachardoublequoteclose}\isanewline
%
\isadelimproof
\ \ %
\endisadelimproof
%
\isatagproof
\isacommand{by}\isamarkupfalse%
\ {\isacharparenleft}{\kern0pt}auto\ simp\ add{\isacharcolon}{\kern0pt}\ map{\isacharunderscore}{\kern0pt}option{\isacharunderscore}{\kern0pt}case\ split{\isacharcolon}{\kern0pt}\ option{\isachardot}{\kern0pt}split{\isacharparenright}{\kern0pt}%
\endisatagproof
{\isafoldproof}%
%
\isadelimproof
\isanewline
%
\endisadelimproof
\isanewline
\isanewline
\isacommand{lemma}\isamarkupfalse%
\ tab{\isacharunderscore}{\kern0pt}succ{\isacharunderscore}{\kern0pt}correct{\isacharcolon}{\kern0pt}\ {\isachardoublequoteopen}E{\isacharunderscore}{\kern0pt}of{\isacharunderscore}{\kern0pt}succ\ {\isacharparenleft}{\kern0pt}tab{\isacharunderscore}{\kern0pt}succ\ l{\isacharparenright}{\kern0pt}\ {\isacharequal}{\kern0pt}\ set\ l{\isachardoublequoteclose}\isanewline
%
\isadelimproof
%
\endisadelimproof
%
\isatagproof
\isacommand{proof}\isamarkupfalse%
\ {\isacharminus}{\kern0pt}\isanewline
\ \ \isacommand{have}\isamarkupfalse%
\ {\isachardoublequoteopen}set\ {\isacharparenleft}{\kern0pt}Mapping{\isachardot}{\kern0pt}lookup{\isacharunderscore}{\kern0pt}default\ {\isacharbrackleft}{\kern0pt}{\isacharbrackright}{\kern0pt}\ {\isacharparenleft}{\kern0pt}fold\ {\isacharparenleft}{\kern0pt}{\isasymlambda}{\isacharparenleft}{\kern0pt}u{\isacharcomma}{\kern0pt}v{\isacharparenright}{\kern0pt}{\isachardot}{\kern0pt}\ Mapping{\isachardot}{\kern0pt}map{\isacharunderscore}{\kern0pt}default\ u\ {\isacharbrackleft}{\kern0pt}{\isacharbrackright}{\kern0pt}\ {\isacharparenleft}{\kern0pt}Cons\ v{\isacharparenright}{\kern0pt}{\isacharparenright}{\kern0pt}\ l\ m{\isacharparenright}{\kern0pt}\ u{\isacharparenright}{\kern0pt}\ {\isacharequal}{\kern0pt}\ set\ l\ {\isacharbackquote}{\kern0pt}{\isacharbackquote}{\kern0pt}\ {\isacharbraceleft}{\kern0pt}u{\isacharbraceright}{\kern0pt}\ {\isasymunion}\ set\ {\isacharparenleft}{\kern0pt}Mapping{\isachardot}{\kern0pt}lookup{\isacharunderscore}{\kern0pt}default\ {\isacharbrackleft}{\kern0pt}{\isacharbrackright}{\kern0pt}\ m\ u{\isacharparenright}{\kern0pt}{\isachardoublequoteclose}\isanewline
\ \ \ \ \isakeyword{for}\ m\ u\isanewline
\ \ \ \ \isacommand{apply}\isamarkupfalse%
\ {\isacharparenleft}{\kern0pt}induction\ l\ arbitrary{\isacharcolon}{\kern0pt}\ m{\isacharparenright}{\kern0pt}\isanewline
\ \ \ \ \isacommand{by}\isamarkupfalse%
\ {\isacharparenleft}{\kern0pt}auto\isanewline
\ \ \ \ \ \ simp{\isacharcolon}{\kern0pt}\ Mapping{\isachardot}{\kern0pt}lookup{\isacharunderscore}{\kern0pt}default{\isacharunderscore}{\kern0pt}def\ Mapping{\isachardot}{\kern0pt}map{\isacharunderscore}{\kern0pt}default{\isacharunderscore}{\kern0pt}def\ Mapping{\isachardot}{\kern0pt}default{\isacharunderscore}{\kern0pt}def\isanewline
\ \ \ \ \ \ simp{\isacharcolon}{\kern0pt}\ lookup{\isacharunderscore}{\kern0pt}map{\isacharunderscore}{\kern0pt}entry{\isacharprime}{\kern0pt}\ lookup{\isacharunderscore}{\kern0pt}update{\isacharprime}{\kern0pt}\ keys{\isacharunderscore}{\kern0pt}is{\isacharunderscore}{\kern0pt}none{\isacharunderscore}{\kern0pt}rep\ Option{\isachardot}{\kern0pt}is{\isacharunderscore}{\kern0pt}none{\isacharunderscore}{\kern0pt}def\isanewline
\ \ \ \ \ \ split{\isacharcolon}{\kern0pt}\ if{\isacharunderscore}{\kern0pt}splits\isanewline
\ \ \ \ {\isacharparenright}{\kern0pt}\isanewline
\ \ \isacommand{from}\isamarkupfalse%
\ this{\isacharbrackleft}{\kern0pt}\isakeyword{where}\ m{\isacharequal}{\kern0pt}Mapping{\isachardot}{\kern0pt}empty{\isacharbrackright}{\kern0pt}\ \isacommand{show}\isamarkupfalse%
\ {\isacharquery}{\kern0pt}thesis\isanewline
\ \ \ \ \isacommand{by}\isamarkupfalse%
\ {\isacharparenleft}{\kern0pt}auto\ simp{\isacharcolon}{\kern0pt}\ E{\isacharunderscore}{\kern0pt}of{\isacharunderscore}{\kern0pt}succ{\isacharunderscore}{\kern0pt}def\ tab{\isacharunderscore}{\kern0pt}succ{\isacharunderscore}{\kern0pt}def\ lookup{\isacharunderscore}{\kern0pt}default{\isacharunderscore}{\kern0pt}empty{\isacharparenright}{\kern0pt}\isanewline
\isacommand{qed}\isamarkupfalse%
%
\endisatagproof
{\isafoldproof}%
%
\isadelimproof
\isanewline
%
\endisadelimproof
\isanewline
\isacommand{end}\isamarkupfalse%
\isanewline
\isanewline
\isacommand{lemma}\isamarkupfalse%
\ finite{\isacharunderscore}{\kern0pt}imp{\isacharunderscore}{\kern0pt}finite{\isacharunderscore}{\kern0pt}dfs{\isacharunderscore}{\kern0pt}reachable{\isacharcolon}{\kern0pt}\isanewline
\ \ {\isachardoublequoteopen}{\isasymlbrakk}finite\ E{\isacharsemicolon}{\kern0pt}\ finite\ S{\isasymrbrakk}\ {\isasymLongrightarrow}\ finite\ {\isacharparenleft}{\kern0pt}E\isactrlsup {\isacharasterisk}{\kern0pt}{\isacharbackquote}{\kern0pt}{\isacharbackquote}{\kern0pt}S{\isacharparenright}{\kern0pt}{\isachardoublequoteclose}\isanewline
%
\isadelimproof
\ \ %
\endisadelimproof
%
\isatagproof
\isacommand{apply}\isamarkupfalse%
\ {\isacharparenleft}{\kern0pt}rule\ finite{\isacharunderscore}{\kern0pt}subset{\isacharbrackleft}{\kern0pt}\isakeyword{where}\ B{\isacharequal}{\kern0pt}{\isachardoublequoteopen}S\ {\isasymunion}\ {\isacharparenleft}{\kern0pt}Relation{\isachardot}{\kern0pt}Domain\ E\ {\isasymunion}\ Relation{\isachardot}{\kern0pt}Range\ E{\isacharparenright}{\kern0pt}{\isachardoublequoteclose}{\isacharbrackright}{\kern0pt}{\isacharparenright}{\kern0pt}\isanewline
\ \ \isacommand{apply}\isamarkupfalse%
\ {\isacharparenleft}{\kern0pt}auto\ simp{\isacharcolon}{\kern0pt}\ intro{\isacharcolon}{\kern0pt}\ finite{\isacharunderscore}{\kern0pt}Domain\ finite{\isacharunderscore}{\kern0pt}Range\ elim{\isacharcolon}{\kern0pt}\ rtranclE{\isacharparenright}{\kern0pt}\isanewline
\ \ \isacommand{done}\isamarkupfalse%
%
\endisatagproof
{\isafoldproof}%
%
\isadelimproof
\isanewline
%
\endisadelimproof
\isanewline
\isacommand{lemma}\isamarkupfalse%
\ dfs{\isacharunderscore}{\kern0pt}reachable{\isacharunderscore}{\kern0pt}tab{\isacharunderscore}{\kern0pt}succ{\isacharunderscore}{\kern0pt}correct{\isacharcolon}{\kern0pt}\ {\isachardoublequoteopen}dfs{\isacharunderscore}{\kern0pt}reachable\ {\isacharparenleft}{\kern0pt}tab{\isacharunderscore}{\kern0pt}succ\ l{\isacharparenright}{\kern0pt}\ D\ vs\isactrlsub {\isadigit{0}}\ {\isasymlongleftrightarrow}\ Collect\ D\ {\isasyminter}\ {\isacharparenleft}{\kern0pt}set\ l{\isacharparenright}{\kern0pt}\isactrlsup {\isacharasterisk}{\kern0pt}{\isacharbackquote}{\kern0pt}{\isacharbackquote}{\kern0pt}set\ vs\isactrlsub {\isadigit{0}}\ {\isasymnoteq}\ {\isacharbraceleft}{\kern0pt}{\isacharbraceright}{\kern0pt}{\isachardoublequoteclose}\isanewline
%
\isadelimproof
\ \ %
\endisadelimproof
%
\isatagproof
\isacommand{apply}\isamarkupfalse%
\ {\isacharparenleft}{\kern0pt}subst\ dfs{\isacharunderscore}{\kern0pt}reachable{\isacharunderscore}{\kern0pt}correct{\isacharparenright}{\kern0pt}\isanewline
\ \ \isacommand{by}\isamarkupfalse%
\ {\isacharparenleft}{\kern0pt}simp{\isacharunderscore}{\kern0pt}all\ add{\isacharcolon}{\kern0pt}\ tab{\isacharunderscore}{\kern0pt}succ{\isacharunderscore}{\kern0pt}correct\ finite{\isacharunderscore}{\kern0pt}imp{\isacharunderscore}{\kern0pt}finite{\isacharunderscore}{\kern0pt}dfs{\isacharunderscore}{\kern0pt}reachable{\isacharparenright}{\kern0pt}%
\endisatagproof
{\isafoldproof}%
%
\isadelimproof
%
\endisadelimproof
%
\isadelimdocument
%
\endisadelimdocument
%
\isatagdocument
%
\isamarkupsubsection{Implementation Refinements%
}
\isamarkuptrue%
%
\isamarkupsubsubsection{Of-Type%
}
\isamarkuptrue%
%
\endisatagdocument
{\isafolddocument}%
%
\isadelimdocument
%
\endisadelimdocument
\isacommand{definition}\isamarkupfalse%
\ {\isachardoublequoteopen}of{\isacharunderscore}{\kern0pt}type{\isacharunderscore}{\kern0pt}impl\ G\ oT\ T\ {\isasymequiv}\ {\isacharparenleft}{\kern0pt}{\isasymforall}pt{\isasymin}set\ {\isacharparenleft}{\kern0pt}primitives\ oT{\isacharparenright}{\kern0pt}{\isachardot}{\kern0pt}\ dfs{\isacharunderscore}{\kern0pt}reachable\ G\ {\isacharparenleft}{\kern0pt}{\isacharparenleft}{\kern0pt}{\isacharequal}{\kern0pt}{\isacharparenright}{\kern0pt}\ pt{\isacharparenright}{\kern0pt}\ {\isacharparenleft}{\kern0pt}primitives\ T{\isacharparenright}{\kern0pt}{\isacharparenright}{\kern0pt}{\isachardoublequoteclose}\isanewline
\isanewline
\isanewline
\isacommand{fun}\isamarkupfalse%
\ ty{\isacharunderscore}{\kern0pt}term{\isacharprime}{\kern0pt}\ \isakeyword{where}\isanewline
\ \ {\isachardoublequoteopen}ty{\isacharunderscore}{\kern0pt}term{\isacharprime}{\kern0pt}\ varT\ objT\ {\isacharparenleft}{\kern0pt}term{\isachardot}{\kern0pt}VAR\ v{\isacharparenright}{\kern0pt}\ {\isacharequal}{\kern0pt}\ varT\ v{\isachardoublequoteclose}\isanewline
{\isacharbar}{\kern0pt}\ {\isachardoublequoteopen}ty{\isacharunderscore}{\kern0pt}term{\isacharprime}{\kern0pt}\ varT\ objT\ {\isacharparenleft}{\kern0pt}term{\isachardot}{\kern0pt}CONST\ c{\isacharparenright}{\kern0pt}\ {\isacharequal}{\kern0pt}\ Mapping{\isachardot}{\kern0pt}lookup\ objT\ c{\isachardoublequoteclose}\isanewline
\isanewline
\isacommand{lemma}\isamarkupfalse%
\ ty{\isacharunderscore}{\kern0pt}term{\isacharprime}{\kern0pt}{\isacharunderscore}{\kern0pt}correct{\isacharunderscore}{\kern0pt}aux{\isacharcolon}{\kern0pt}\ {\isachardoublequoteopen}ty{\isacharunderscore}{\kern0pt}term{\isacharprime}{\kern0pt}\ varT\ objT\ t\ {\isacharequal}{\kern0pt}\ ty{\isacharunderscore}{\kern0pt}term\ varT\ {\isacharparenleft}{\kern0pt}Mapping{\isachardot}{\kern0pt}lookup\ objT{\isacharparenright}{\kern0pt}\ t{\isachardoublequoteclose}\isanewline
%
\isadelimproof
\ \ %
\endisadelimproof
%
\isatagproof
\isacommand{by}\isamarkupfalse%
\ {\isacharparenleft}{\kern0pt}cases\ t{\isacharparenright}{\kern0pt}\ auto%
\endisatagproof
{\isafoldproof}%
%
\isadelimproof
\isanewline
%
\endisadelimproof
\isanewline
\isacommand{lemma}\isamarkupfalse%
\ ty{\isacharunderscore}{\kern0pt}term{\isacharprime}{\kern0pt}{\isacharunderscore}{\kern0pt}correct{\isacharbrackleft}{\kern0pt}simp{\isacharbrackright}{\kern0pt}{\isacharcolon}{\kern0pt}\ {\isachardoublequoteopen}ty{\isacharunderscore}{\kern0pt}term{\isacharprime}{\kern0pt}\ varT\ objT\ {\isacharequal}{\kern0pt}\ ty{\isacharunderscore}{\kern0pt}term\ varT\ {\isacharparenleft}{\kern0pt}Mapping{\isachardot}{\kern0pt}lookup\ objT{\isacharparenright}{\kern0pt}{\isachardoublequoteclose}\isanewline
%
\isadelimproof
\ \ %
\endisadelimproof
%
\isatagproof
\isacommand{using}\isamarkupfalse%
\ ty{\isacharunderscore}{\kern0pt}term{\isacharprime}{\kern0pt}{\isacharunderscore}{\kern0pt}correct{\isacharunderscore}{\kern0pt}aux\ \isacommand{by}\isamarkupfalse%
\ auto%
\endisatagproof
{\isafoldproof}%
%
\isadelimproof
\isanewline
%
\endisadelimproof
\isanewline
\isacommand{context}\isamarkupfalse%
\ ast{\isacharunderscore}{\kern0pt}domain\ \isakeyword{begin}\isanewline
\isanewline
\ \ \isacommand{definition}\isamarkupfalse%
\ {\isachardoublequoteopen}of{\isacharunderscore}{\kern0pt}type{\isadigit{1}}\ pt\ T\ {\isasymlongleftrightarrow}\ pt\ {\isasymin}\ subtype{\isacharunderscore}{\kern0pt}rel\isactrlsup {\isacharasterisk}{\kern0pt}\ {\isacharbackquote}{\kern0pt}{\isacharbackquote}{\kern0pt}\ set\ {\isacharparenleft}{\kern0pt}primitives\ T{\isacharparenright}{\kern0pt}{\isachardoublequoteclose}\isanewline
\isanewline
\ \ \isacommand{lemma}\isamarkupfalse%
\ of{\isacharunderscore}{\kern0pt}type{\isacharunderscore}{\kern0pt}refine{\isadigit{1}}{\isacharcolon}{\kern0pt}\ {\isachardoublequoteopen}of{\isacharunderscore}{\kern0pt}type\ oT\ T\ {\isasymlongleftrightarrow}\ {\isacharparenleft}{\kern0pt}{\isasymforall}pt{\isasymin}set\ {\isacharparenleft}{\kern0pt}primitives\ oT{\isacharparenright}{\kern0pt}{\isachardot}{\kern0pt}\ of{\isacharunderscore}{\kern0pt}type{\isadigit{1}}\ pt\ T{\isacharparenright}{\kern0pt}{\isachardoublequoteclose}\isanewline
%
\isadelimproof
\ \ \ \ %
\endisadelimproof
%
\isatagproof
\isacommand{unfolding}\isamarkupfalse%
\ of{\isacharunderscore}{\kern0pt}type{\isacharunderscore}{\kern0pt}def\ of{\isacharunderscore}{\kern0pt}type{\isadigit{1}}{\isacharunderscore}{\kern0pt}def\ \isacommand{by}\isamarkupfalse%
\ auto%
\endisatagproof
{\isafoldproof}%
%
\isadelimproof
\isanewline
%
\endisadelimproof
\isanewline
\ \ \isacommand{definition}\isamarkupfalse%
\ {\isachardoublequoteopen}STG\ {\isasymequiv}\ {\isacharparenleft}{\kern0pt}tab{\isacharunderscore}{\kern0pt}succ\ {\isacharparenleft}{\kern0pt}map\ subtype{\isacharunderscore}{\kern0pt}edge\ {\isacharparenleft}{\kern0pt}types\ D{\isacharparenright}{\kern0pt}{\isacharparenright}{\kern0pt}{\isacharparenright}{\kern0pt}{\isachardoublequoteclose}\isanewline
\isanewline
\ \ \isacommand{lemma}\isamarkupfalse%
\ subtype{\isacharunderscore}{\kern0pt}rel{\isacharunderscore}{\kern0pt}impl{\isacharcolon}{\kern0pt}\ {\isachardoublequoteopen}subtype{\isacharunderscore}{\kern0pt}rel\ {\isacharequal}{\kern0pt}\ E{\isacharunderscore}{\kern0pt}of{\isacharunderscore}{\kern0pt}succ\ {\isacharparenleft}{\kern0pt}tab{\isacharunderscore}{\kern0pt}succ\ {\isacharparenleft}{\kern0pt}map\ subtype{\isacharunderscore}{\kern0pt}edge\ {\isacharparenleft}{\kern0pt}types\ D{\isacharparenright}{\kern0pt}{\isacharparenright}{\kern0pt}{\isacharparenright}{\kern0pt}{\isachardoublequoteclose}\isanewline
%
\isadelimproof
\ \ \ \ %
\endisadelimproof
%
\isatagproof
\isacommand{by}\isamarkupfalse%
\ {\isacharparenleft}{\kern0pt}simp\ add{\isacharcolon}{\kern0pt}\ tab{\isacharunderscore}{\kern0pt}succ{\isacharunderscore}{\kern0pt}correct\ subtype{\isacharunderscore}{\kern0pt}rel{\isacharunderscore}{\kern0pt}def{\isacharparenright}{\kern0pt}%
\endisatagproof
{\isafoldproof}%
%
\isadelimproof
\isanewline
%
\endisadelimproof
\isanewline
\ \ \isacommand{lemma}\isamarkupfalse%
\ of{\isacharunderscore}{\kern0pt}type{\isadigit{1}}{\isacharunderscore}{\kern0pt}impl{\isacharcolon}{\kern0pt}\ {\isachardoublequoteopen}of{\isacharunderscore}{\kern0pt}type{\isadigit{1}}\ pt\ T\ {\isasymlongleftrightarrow}\ dfs{\isacharunderscore}{\kern0pt}reachable\ {\isacharparenleft}{\kern0pt}tab{\isacharunderscore}{\kern0pt}succ\ {\isacharparenleft}{\kern0pt}map\ subtype{\isacharunderscore}{\kern0pt}edge\ {\isacharparenleft}{\kern0pt}types\ D{\isacharparenright}{\kern0pt}{\isacharparenright}{\kern0pt}{\isacharparenright}{\kern0pt}\ {\isacharparenleft}{\kern0pt}{\isacharparenleft}{\kern0pt}{\isacharequal}{\kern0pt}{\isacharparenright}{\kern0pt}pt{\isacharparenright}{\kern0pt}\ {\isacharparenleft}{\kern0pt}primitives\ T{\isacharparenright}{\kern0pt}{\isachardoublequoteclose}\isanewline
%
\isadelimproof
\ \ \ \ %
\endisadelimproof
%
\isatagproof
\isacommand{by}\isamarkupfalse%
\ {\isacharparenleft}{\kern0pt}simp\ add{\isacharcolon}{\kern0pt}\ subtype{\isacharunderscore}{\kern0pt}rel{\isacharunderscore}{\kern0pt}impl\ of{\isacharunderscore}{\kern0pt}type{\isadigit{1}}{\isacharunderscore}{\kern0pt}def\ dfs{\isacharunderscore}{\kern0pt}reachable{\isacharunderscore}{\kern0pt}tab{\isacharunderscore}{\kern0pt}succ{\isacharunderscore}{\kern0pt}correct\ tab{\isacharunderscore}{\kern0pt}succ{\isacharunderscore}{\kern0pt}correct{\isacharparenright}{\kern0pt}%
\endisatagproof
{\isafoldproof}%
%
\isadelimproof
\isanewline
%
\endisadelimproof
\isanewline
\ \ \isacommand{lemma}\isamarkupfalse%
\ of{\isacharunderscore}{\kern0pt}type{\isacharunderscore}{\kern0pt}impl{\isacharunderscore}{\kern0pt}correct{\isacharcolon}{\kern0pt}\ {\isachardoublequoteopen}of{\isacharunderscore}{\kern0pt}type{\isacharunderscore}{\kern0pt}impl\ STG\ oT\ T\ {\isasymlongleftrightarrow}\ of{\isacharunderscore}{\kern0pt}type\ oT\ T{\isachardoublequoteclose}\isanewline
%
\isadelimproof
\ \ \ \ %
\endisadelimproof
%
\isatagproof
\isacommand{unfolding}\isamarkupfalse%
\ of{\isacharunderscore}{\kern0pt}type{\isadigit{1}}{\isacharunderscore}{\kern0pt}impl\ STG{\isacharunderscore}{\kern0pt}def\ of{\isacharunderscore}{\kern0pt}type{\isacharunderscore}{\kern0pt}impl{\isacharunderscore}{\kern0pt}def\ of{\isacharunderscore}{\kern0pt}type{\isacharunderscore}{\kern0pt}refine{\isadigit{1}}\ \isacommand{{\isachardot}{\kern0pt}{\isachardot}{\kern0pt}}\isamarkupfalse%
%
\endisatagproof
{\isafoldproof}%
%
\isadelimproof
\isanewline
%
\endisadelimproof
\isanewline
\ \ \isacommand{definition}\isamarkupfalse%
\ mp{\isacharunderscore}{\kern0pt}constT\ {\isacharcolon}{\kern0pt}{\isacharcolon}{\kern0pt}\ {\isachardoublequoteopen}{\isacharparenleft}{\kern0pt}object{\isacharcomma}{\kern0pt}\ type{\isacharparenright}{\kern0pt}\ mapping{\isachardoublequoteclose}\ \isakeyword{where}\isanewline
\ \ \ \ {\isachardoublequoteopen}mp{\isacharunderscore}{\kern0pt}constT\ {\isacharequal}{\kern0pt}\ Mapping{\isachardot}{\kern0pt}of{\isacharunderscore}{\kern0pt}alist\ {\isacharparenleft}{\kern0pt}consts\ D{\isacharparenright}{\kern0pt}{\isachardoublequoteclose}\isanewline
\isanewline
\ \ \isacommand{lemma}\isamarkupfalse%
\ mp{\isacharunderscore}{\kern0pt}objT{\isacharunderscore}{\kern0pt}correct{\isacharbrackleft}{\kern0pt}simp{\isacharbrackright}{\kern0pt}{\isacharcolon}{\kern0pt}\ {\isachardoublequoteopen}Mapping{\isachardot}{\kern0pt}lookup\ mp{\isacharunderscore}{\kern0pt}constT\ {\isacharequal}{\kern0pt}\ constT{\isachardoublequoteclose}\isanewline
%
\isadelimproof
\ \ \ \ %
\endisadelimproof
%
\isatagproof
\isacommand{unfolding}\isamarkupfalse%
\ mp{\isacharunderscore}{\kern0pt}constT{\isacharunderscore}{\kern0pt}def\ constT{\isacharunderscore}{\kern0pt}def\isanewline
\ \ \ \ \isacommand{by}\isamarkupfalse%
\ transfer\ {\isacharparenleft}{\kern0pt}simp\ add{\isacharcolon}{\kern0pt}\ Map{\isacharunderscore}{\kern0pt}To{\isacharunderscore}{\kern0pt}Mapping{\isachardot}{\kern0pt}map{\isacharunderscore}{\kern0pt}apply{\isacharunderscore}{\kern0pt}def{\isacharparenright}{\kern0pt}%
\endisatagproof
{\isafoldproof}%
%
\isadelimproof
%
\endisadelimproof
%
\begin{isamarkuptext}%
Lifting the subtype-graph through wf-checker%
\end{isamarkuptext}\isamarkuptrue%
\ \ \isacommand{context}\isamarkupfalse%
\isanewline
\ \ \ \ \isakeyword{fixes}\ ty{\isacharunderscore}{\kern0pt}ent\ {\isacharcolon}{\kern0pt}{\isacharcolon}{\kern0pt}\ {\isachardoublequoteopen}{\isacharprime}{\kern0pt}ent\ {\isasymrightharpoonup}\ type{\isachardoublequoteclose}\ \ %
\isamarkupcmt{Entity's type, None if invalid%
}\isanewline
\ \ \isakeyword{begin}\isanewline
\isanewline
\ \ \ \ \isacommand{definition}\isamarkupfalse%
\ {\isachardoublequoteopen}is{\isacharunderscore}{\kern0pt}of{\isacharunderscore}{\kern0pt}type{\isacharprime}{\kern0pt}\ stg\ v\ T\ {\isasymlongleftrightarrow}\ {\isacharparenleft}{\kern0pt}\isanewline
\ \ \ \ \ \ case\ ty{\isacharunderscore}{\kern0pt}ent\ v\ of\isanewline
\ \ \ \ \ \ \ \ Some\ vT\ {\isasymRightarrow}\ of{\isacharunderscore}{\kern0pt}type{\isacharunderscore}{\kern0pt}impl\ stg\ vT\ T\isanewline
\ \ \ \ \ \ {\isacharbar}{\kern0pt}\ None\ {\isasymRightarrow}\ False{\isacharparenright}{\kern0pt}{\isachardoublequoteclose}\isanewline
\isanewline
\ \ \ \ \isacommand{lemma}\isamarkupfalse%
\ is{\isacharunderscore}{\kern0pt}of{\isacharunderscore}{\kern0pt}type{\isacharprime}{\kern0pt}{\isacharunderscore}{\kern0pt}correct{\isacharcolon}{\kern0pt}\ {\isachardoublequoteopen}is{\isacharunderscore}{\kern0pt}of{\isacharunderscore}{\kern0pt}type{\isacharprime}{\kern0pt}\ STG\ v\ T\ {\isacharequal}{\kern0pt}\ is{\isacharunderscore}{\kern0pt}of{\isacharunderscore}{\kern0pt}type\ ty{\isacharunderscore}{\kern0pt}ent\ v\ T{\isachardoublequoteclose}\isanewline
%
\isadelimproof
\ \ \ \ \ \ %
\endisadelimproof
%
\isatagproof
\isacommand{unfolding}\isamarkupfalse%
\ is{\isacharunderscore}{\kern0pt}of{\isacharunderscore}{\kern0pt}type{\isacharprime}{\kern0pt}{\isacharunderscore}{\kern0pt}def\ is{\isacharunderscore}{\kern0pt}of{\isacharunderscore}{\kern0pt}type{\isacharunderscore}{\kern0pt}def\ of{\isacharunderscore}{\kern0pt}type{\isacharunderscore}{\kern0pt}impl{\isacharunderscore}{\kern0pt}correct\ \isacommand{{\isachardot}{\kern0pt}{\isachardot}{\kern0pt}}\isamarkupfalse%
%
\endisatagproof
{\isafoldproof}%
%
\isadelimproof
\isanewline
%
\endisadelimproof
\isanewline
\ \ \ \ \isacommand{fun}\isamarkupfalse%
\ wf{\isacharunderscore}{\kern0pt}pred{\isacharunderscore}{\kern0pt}atom{\isacharprime}{\kern0pt}\ \isakeyword{where}\ {\isachardoublequoteopen}wf{\isacharunderscore}{\kern0pt}pred{\isacharunderscore}{\kern0pt}atom{\isacharprime}{\kern0pt}\ stg\ {\isacharparenleft}{\kern0pt}p{\isacharcomma}{\kern0pt}vs{\isacharparenright}{\kern0pt}\ {\isasymlongleftrightarrow}\ \isanewline
\ \ \ \ \ \ {\isacharparenleft}{\kern0pt}case\ sig\ p\ of\isanewline
\ \ \ \ \ \ \ \ \ \ None\ {\isasymRightarrow}\ False\isanewline
\ \ \ \ \ \ \ \ {\isacharbar}{\kern0pt}\ Some\ Ts\ {\isasymRightarrow}\ list{\isacharunderscore}{\kern0pt}all{\isadigit{2}}\ {\isacharparenleft}{\kern0pt}is{\isacharunderscore}{\kern0pt}of{\isacharunderscore}{\kern0pt}type{\isacharprime}{\kern0pt}\ stg{\isacharparenright}{\kern0pt}\ vs\ Ts{\isacharparenright}{\kern0pt}{\isachardoublequoteclose}\isanewline
\isanewline
\ \ \ \ \isacommand{lemma}\isamarkupfalse%
\ wf{\isacharunderscore}{\kern0pt}pred{\isacharunderscore}{\kern0pt}atom{\isacharprime}{\kern0pt}{\isacharunderscore}{\kern0pt}correct{\isacharcolon}{\kern0pt}\ {\isachardoublequoteopen}wf{\isacharunderscore}{\kern0pt}pred{\isacharunderscore}{\kern0pt}atom{\isacharprime}{\kern0pt}\ STG\ pvs\ {\isacharequal}{\kern0pt}\ wf{\isacharunderscore}{\kern0pt}pred{\isacharunderscore}{\kern0pt}atom\ ty{\isacharunderscore}{\kern0pt}ent\ pvs{\isachardoublequoteclose}\isanewline
%
\isadelimproof
\ \ \ \ \ \ %
\endisadelimproof
%
\isatagproof
\isacommand{by}\isamarkupfalse%
\ {\isacharparenleft}{\kern0pt}cases\ pvs{\isacharparenright}{\kern0pt}\ {\isacharparenleft}{\kern0pt}auto\ simp{\isacharcolon}{\kern0pt}\ is{\isacharunderscore}{\kern0pt}of{\isacharunderscore}{\kern0pt}type{\isacharprime}{\kern0pt}{\isacharunderscore}{\kern0pt}correct{\isacharbrackleft}{\kern0pt}abs{\isacharunderscore}{\kern0pt}def{\isacharbrackright}{\kern0pt}\ split{\isacharcolon}{\kern0pt}option{\isachardot}{\kern0pt}split{\isacharparenright}{\kern0pt}%
\endisatagproof
{\isafoldproof}%
%
\isadelimproof
\isanewline
%
\endisadelimproof
\isanewline
\ \ \ \ \isacommand{fun}\isamarkupfalse%
\ wf{\isacharunderscore}{\kern0pt}func{\isacharunderscore}{\kern0pt}args{\isacharprime}{\kern0pt}\ \isakeyword{where}\ {\isachardoublequoteopen}wf{\isacharunderscore}{\kern0pt}func{\isacharunderscore}{\kern0pt}args{\isacharprime}{\kern0pt}\ stg\ {\isacharparenleft}{\kern0pt}f{\isacharcomma}{\kern0pt}args{\isacharparenright}{\kern0pt}\ {\isasymlongleftrightarrow}\ \isanewline
\ \ \ \ \ \ {\isacharparenleft}{\kern0pt}case\ func{\isacharunderscore}{\kern0pt}sig\ f\ of\isanewline
\ \ \ \ \ \ \ \ \ \ None\ {\isasymRightarrow}\ False\isanewline
\ \ \ \ \ \ \ \ {\isacharbar}{\kern0pt}\ Some\ Ts\ {\isasymRightarrow}\ list{\isacharunderscore}{\kern0pt}all{\isadigit{2}}\ {\isacharparenleft}{\kern0pt}is{\isacharunderscore}{\kern0pt}of{\isacharunderscore}{\kern0pt}type{\isacharprime}{\kern0pt}\ stg{\isacharparenright}{\kern0pt}\ args\ Ts{\isacharparenright}{\kern0pt}{\isachardoublequoteclose}\isanewline
\isanewline
\ \ \ \ \isacommand{lemma}\isamarkupfalse%
\ wf{\isacharunderscore}{\kern0pt}func{\isacharunderscore}{\kern0pt}args{\isacharprime}{\kern0pt}{\isacharunderscore}{\kern0pt}correct{\isacharcolon}{\kern0pt}\ {\isachardoublequoteopen}wf{\isacharunderscore}{\kern0pt}func{\isacharunderscore}{\kern0pt}args{\isacharprime}{\kern0pt}\ STG\ fargs\ {\isacharequal}{\kern0pt}\ wf{\isacharunderscore}{\kern0pt}func{\isacharunderscore}{\kern0pt}args\ ty{\isacharunderscore}{\kern0pt}ent\ fargs{\isachardoublequoteclose}\isanewline
%
\isadelimproof
\ \ \ \ \ \ %
\endisadelimproof
%
\isatagproof
\isacommand{by}\isamarkupfalse%
\ {\isacharparenleft}{\kern0pt}cases\ fargs{\isacharparenright}{\kern0pt}\ {\isacharparenleft}{\kern0pt}auto\ simp{\isacharcolon}{\kern0pt}\ is{\isacharunderscore}{\kern0pt}of{\isacharunderscore}{\kern0pt}type{\isacharprime}{\kern0pt}{\isacharunderscore}{\kern0pt}correct{\isacharbrackleft}{\kern0pt}abs{\isacharunderscore}{\kern0pt}def{\isacharbrackright}{\kern0pt}\ split{\isacharcolon}{\kern0pt}option{\isachardot}{\kern0pt}split{\isacharparenright}{\kern0pt}%
\endisatagproof
{\isafoldproof}%
%
\isadelimproof
\isanewline
%
\endisadelimproof
\isanewline
\ \ \ \ \isacommand{fun}\isamarkupfalse%
\ wf{\isacharunderscore}{\kern0pt}atom{\isacharprime}{\kern0pt}\ {\isacharcolon}{\kern0pt}{\isacharcolon}{\kern0pt}\ {\isachardoublequoteopen}{\isacharunderscore}{\kern0pt}\ {\isasymRightarrow}\ {\isacharprime}{\kern0pt}ent\ atom\ {\isasymRightarrow}\ bool{\isachardoublequoteclose}\ \isakeyword{where}\isanewline
\ \ \ \ \ \ {\isachardoublequoteopen}wf{\isacharunderscore}{\kern0pt}atom{\isacharprime}{\kern0pt}\ stg\ {\isacharparenleft}{\kern0pt}atom{\isachardot}{\kern0pt}predAtm\ p\ vs{\isacharparenright}{\kern0pt}\ {\isasymlongleftrightarrow}\ wf{\isacharunderscore}{\kern0pt}pred{\isacharunderscore}{\kern0pt}atom{\isacharprime}{\kern0pt}\ stg\ {\isacharparenleft}{\kern0pt}p{\isacharcomma}{\kern0pt}vs{\isacharparenright}{\kern0pt}{\isachardoublequoteclose}\isanewline
\ \ \ \ {\isacharbar}{\kern0pt}\ {\isachardoublequoteopen}wf{\isacharunderscore}{\kern0pt}atom{\isacharprime}{\kern0pt}\ stg\ {\isacharparenleft}{\kern0pt}atom{\isachardot}{\kern0pt}eqAtm\ a\ b{\isacharparenright}{\kern0pt}\ {\isacharequal}{\kern0pt}\ {\isacharparenleft}{\kern0pt}ty{\isacharunderscore}{\kern0pt}ent\ a\ {\isasymnoteq}\ None\ {\isasymand}\ ty{\isacharunderscore}{\kern0pt}ent\ b\ {\isasymnoteq}\ None{\isacharparenright}{\kern0pt}{\isachardoublequoteclose}\isanewline
\isanewline
\ \ \ \ \isacommand{lemma}\isamarkupfalse%
\ wf{\isacharunderscore}{\kern0pt}atom{\isacharprime}{\kern0pt}{\isacharunderscore}{\kern0pt}correct{\isacharcolon}{\kern0pt}\ {\isachardoublequoteopen}wf{\isacharunderscore}{\kern0pt}atom{\isacharprime}{\kern0pt}\ STG\ a\ {\isacharequal}{\kern0pt}\ wf{\isacharunderscore}{\kern0pt}atom\ ty{\isacharunderscore}{\kern0pt}ent\ a{\isachardoublequoteclose}\isanewline
%
\isadelimproof
\ \ \ \ \ \ %
\endisadelimproof
%
\isatagproof
\isacommand{by}\isamarkupfalse%
\ {\isacharparenleft}{\kern0pt}cases\ a{\isacharparenright}{\kern0pt}\ {\isacharparenleft}{\kern0pt}auto\ simp{\isacharcolon}{\kern0pt}\ wf{\isacharunderscore}{\kern0pt}pred{\isacharunderscore}{\kern0pt}atom{\isacharprime}{\kern0pt}{\isacharunderscore}{\kern0pt}correct\ is{\isacharunderscore}{\kern0pt}of{\isacharunderscore}{\kern0pt}type{\isacharprime}{\kern0pt}{\isacharunderscore}{\kern0pt}correct{\isacharbrackleft}{\kern0pt}abs{\isacharunderscore}{\kern0pt}def{\isacharbrackright}{\kern0pt}\ split{\isacharcolon}{\kern0pt}\ option{\isachardot}{\kern0pt}splits{\isacharparenright}{\kern0pt}%
\endisatagproof
{\isafoldproof}%
%
\isadelimproof
\isanewline
%
\endisadelimproof
\isanewline
\ \ \ \ \isacommand{fun}\isamarkupfalse%
\ wf{\isacharunderscore}{\kern0pt}fmla{\isacharprime}{\kern0pt}\ {\isacharcolon}{\kern0pt}{\isacharcolon}{\kern0pt}\ {\isachardoublequoteopen}{\isacharunderscore}{\kern0pt}\ {\isasymRightarrow}\ {\isacharparenleft}{\kern0pt}{\isacharprime}{\kern0pt}ent\ atom{\isacharparenright}{\kern0pt}\ formula\ {\isasymRightarrow}\ bool{\isachardoublequoteclose}\ \isakeyword{where}\isanewline
\ \ \ \ \ \ {\isachardoublequoteopen}wf{\isacharunderscore}{\kern0pt}fmla{\isacharprime}{\kern0pt}\ stg\ {\isacharparenleft}{\kern0pt}Atom\ a{\isacharparenright}{\kern0pt}\ {\isasymlongleftrightarrow}\ wf{\isacharunderscore}{\kern0pt}atom{\isacharprime}{\kern0pt}\ stg\ a{\isachardoublequoteclose}\isanewline
\ \ \ \ {\isacharbar}{\kern0pt}\ {\isachardoublequoteopen}wf{\isacharunderscore}{\kern0pt}fmla{\isacharprime}{\kern0pt}\ stg\ {\isasymbottom}\ {\isasymlongleftrightarrow}\ True{\isachardoublequoteclose}\isanewline
\ \ \ \ {\isacharbar}{\kern0pt}\ {\isachardoublequoteopen}wf{\isacharunderscore}{\kern0pt}fmla{\isacharprime}{\kern0pt}\ stg\ {\isacharparenleft}{\kern0pt}{\isasymphi}{\isadigit{1}}\ \isactrlbold {\isasymand}\ {\isasymphi}{\isadigit{2}}{\isacharparenright}{\kern0pt}\ {\isasymlongleftrightarrow}\ {\isacharparenleft}{\kern0pt}wf{\isacharunderscore}{\kern0pt}fmla{\isacharprime}{\kern0pt}\ stg\ {\isasymphi}{\isadigit{1}}\ {\isasymand}\ wf{\isacharunderscore}{\kern0pt}fmla{\isacharprime}{\kern0pt}\ stg\ {\isasymphi}{\isadigit{2}}{\isacharparenright}{\kern0pt}{\isachardoublequoteclose}\isanewline
\ \ \ \ {\isacharbar}{\kern0pt}\ {\isachardoublequoteopen}wf{\isacharunderscore}{\kern0pt}fmla{\isacharprime}{\kern0pt}\ stg\ {\isacharparenleft}{\kern0pt}{\isasymphi}{\isadigit{1}}\ \isactrlbold {\isasymor}\ {\isasymphi}{\isadigit{2}}{\isacharparenright}{\kern0pt}\ {\isasymlongleftrightarrow}\ {\isacharparenleft}{\kern0pt}wf{\isacharunderscore}{\kern0pt}fmla{\isacharprime}{\kern0pt}\ stg\ {\isasymphi}{\isadigit{1}}\ {\isasymand}\ wf{\isacharunderscore}{\kern0pt}fmla{\isacharprime}{\kern0pt}\ stg\ {\isasymphi}{\isadigit{2}}{\isacharparenright}{\kern0pt}{\isachardoublequoteclose}\isanewline
\ \ \ \ {\isacharbar}{\kern0pt}\ {\isachardoublequoteopen}wf{\isacharunderscore}{\kern0pt}fmla{\isacharprime}{\kern0pt}\ stg\ {\isacharparenleft}{\kern0pt}{\isasymphi}{\isadigit{1}}\ \isactrlbold {\isasymrightarrow}\ {\isasymphi}{\isadigit{2}}{\isacharparenright}{\kern0pt}\ {\isasymlongleftrightarrow}\ {\isacharparenleft}{\kern0pt}wf{\isacharunderscore}{\kern0pt}fmla{\isacharprime}{\kern0pt}\ stg\ {\isasymphi}{\isadigit{1}}\ {\isasymand}\ wf{\isacharunderscore}{\kern0pt}fmla{\isacharprime}{\kern0pt}\ stg\ {\isasymphi}{\isadigit{2}}{\isacharparenright}{\kern0pt}{\isachardoublequoteclose}\isanewline
\ \ \ \ {\isacharbar}{\kern0pt}\ {\isachardoublequoteopen}wf{\isacharunderscore}{\kern0pt}fmla{\isacharprime}{\kern0pt}\ stg\ {\isacharparenleft}{\kern0pt}\isactrlbold {\isasymnot}{\isasymphi}{\isacharparenright}{\kern0pt}\ {\isasymlongleftrightarrow}\ wf{\isacharunderscore}{\kern0pt}fmla{\isacharprime}{\kern0pt}\ stg\ {\isasymphi}{\isachardoublequoteclose}\isanewline
\isanewline
\ \ \ \ \isacommand{lemma}\isamarkupfalse%
\ wf{\isacharunderscore}{\kern0pt}fmla{\isacharprime}{\kern0pt}{\isacharunderscore}{\kern0pt}correct{\isacharcolon}{\kern0pt}\ {\isachardoublequoteopen}wf{\isacharunderscore}{\kern0pt}fmla{\isacharprime}{\kern0pt}\ STG\ {\isasymphi}\ {\isasymlongleftrightarrow}\ wf{\isacharunderscore}{\kern0pt}fmla\ ty{\isacharunderscore}{\kern0pt}ent\ {\isasymphi}{\isachardoublequoteclose}\isanewline
%
\isadelimproof
\ \ \ \ \ \ %
\endisadelimproof
%
\isatagproof
\isacommand{by}\isamarkupfalse%
\ {\isacharparenleft}{\kern0pt}induction\ {\isasymphi}\ rule{\isacharcolon}{\kern0pt}\ wf{\isacharunderscore}{\kern0pt}fmla{\isachardot}{\kern0pt}induct{\isacharparenright}{\kern0pt}\ {\isacharparenleft}{\kern0pt}auto\ simp{\isacharcolon}{\kern0pt}\ wf{\isacharunderscore}{\kern0pt}atom{\isacharprime}{\kern0pt}{\isacharunderscore}{\kern0pt}correct{\isacharparenright}{\kern0pt}%
\endisatagproof
{\isafoldproof}%
%
\isadelimproof
\isanewline
%
\endisadelimproof
\isanewline
\ \ \ \ \isacommand{fun}\isamarkupfalse%
\ wf{\isacharunderscore}{\kern0pt}fmla{\isacharunderscore}{\kern0pt}atom{\isadigit{1}}{\isacharprime}{\kern0pt}\ \isakeyword{where}\isanewline
\ \ \ \ \ \ {\isachardoublequoteopen}wf{\isacharunderscore}{\kern0pt}fmla{\isacharunderscore}{\kern0pt}atom{\isadigit{1}}{\isacharprime}{\kern0pt}\ stg\ {\isacharparenleft}{\kern0pt}Atom\ {\isacharparenleft}{\kern0pt}predAtm\ p\ vs{\isacharparenright}{\kern0pt}{\isacharparenright}{\kern0pt}\ {\isasymlongleftrightarrow}\ wf{\isacharunderscore}{\kern0pt}pred{\isacharunderscore}{\kern0pt}atom{\isacharprime}{\kern0pt}\ stg\ {\isacharparenleft}{\kern0pt}p{\isacharcomma}{\kern0pt}vs{\isacharparenright}{\kern0pt}{\isachardoublequoteclose}\isanewline
\ \ \ \ {\isacharbar}{\kern0pt}\ {\isachardoublequoteopen}wf{\isacharunderscore}{\kern0pt}fmla{\isacharunderscore}{\kern0pt}atom{\isadigit{1}}{\isacharprime}{\kern0pt}\ stg\ {\isacharunderscore}{\kern0pt}\ {\isasymlongleftrightarrow}\ False{\isachardoublequoteclose}\isanewline
\isanewline
\ \ \ \ \isacommand{lemma}\isamarkupfalse%
\ wf{\isacharunderscore}{\kern0pt}fmla{\isacharunderscore}{\kern0pt}atom{\isadigit{1}}{\isacharprime}{\kern0pt}{\isacharunderscore}{\kern0pt}correct{\isacharcolon}{\kern0pt}\ {\isachardoublequoteopen}wf{\isacharunderscore}{\kern0pt}fmla{\isacharunderscore}{\kern0pt}atom{\isadigit{1}}{\isacharprime}{\kern0pt}\ STG\ {\isasymphi}\ {\isacharequal}{\kern0pt}\ wf{\isacharunderscore}{\kern0pt}fmla{\isacharunderscore}{\kern0pt}atom\ ty{\isacharunderscore}{\kern0pt}ent\ {\isasymphi}{\isachardoublequoteclose}\isanewline
%
\isadelimproof
\ \ \ \ \ \ %
\endisadelimproof
%
\isatagproof
\isacommand{by}\isamarkupfalse%
\ {\isacharparenleft}{\kern0pt}cases\ {\isasymphi}\ rule{\isacharcolon}{\kern0pt}\ wf{\isacharunderscore}{\kern0pt}fmla{\isacharunderscore}{\kern0pt}atom{\isachardot}{\kern0pt}cases{\isacharparenright}{\kern0pt}\ {\isacharparenleft}{\kern0pt}auto\isanewline
\ \ \ \ \ \ \ \ simp{\isacharcolon}{\kern0pt}\ wf{\isacharunderscore}{\kern0pt}atom{\isacharprime}{\kern0pt}{\isacharunderscore}{\kern0pt}correct\ is{\isacharunderscore}{\kern0pt}of{\isacharunderscore}{\kern0pt}type{\isacharprime}{\kern0pt}{\isacharunderscore}{\kern0pt}correct{\isacharbrackleft}{\kern0pt}abs{\isacharunderscore}{\kern0pt}def{\isacharbrackright}{\kern0pt}\ split{\isacharcolon}{\kern0pt}\ option{\isachardot}{\kern0pt}splits{\isacharparenright}{\kern0pt}%
\endisatagproof
{\isafoldproof}%
%
\isadelimproof
\isanewline
%
\endisadelimproof
\isanewline
\ \ \ \ \isacommand{fun}\isamarkupfalse%
\ wf{\isacharunderscore}{\kern0pt}effect{\isacharprime}{\kern0pt}\ \isakeyword{where}\isanewline
\ \ \ \ \ \ {\isachardoublequoteopen}wf{\isacharunderscore}{\kern0pt}effect{\isacharprime}{\kern0pt}\ stg\ {\isacharparenleft}{\kern0pt}Effect\ a\ d{\isacharparenright}{\kern0pt}\ {\isasymlongleftrightarrow}\isanewline
\ \ \ \ \ \ \ \ \ \ {\isacharparenleft}{\kern0pt}{\isasymforall}ae{\isasymin}set\ a{\isachardot}{\kern0pt}\ wf{\isacharunderscore}{\kern0pt}fmla{\isacharunderscore}{\kern0pt}atom{\isadigit{1}}{\isacharprime}{\kern0pt}\ stg\ ae{\isacharparenright}{\kern0pt}\isanewline
\ \ \ \ \ \ \ \ {\isasymand}\ {\isacharparenleft}{\kern0pt}{\isasymforall}de{\isasymin}set\ d{\isachardot}{\kern0pt}\ \ wf{\isacharunderscore}{\kern0pt}fmla{\isacharunderscore}{\kern0pt}atom{\isadigit{1}}{\isacharprime}{\kern0pt}\ stg\ de{\isacharparenright}{\kern0pt}{\isachardoublequoteclose}\isanewline
\isanewline
\ \ \ \ \isacommand{lemma}\isamarkupfalse%
\ wf{\isacharunderscore}{\kern0pt}effect{\isacharprime}{\kern0pt}{\isacharunderscore}{\kern0pt}correct{\isacharcolon}{\kern0pt}\ {\isachardoublequoteopen}wf{\isacharunderscore}{\kern0pt}effect{\isacharprime}{\kern0pt}\ STG\ e\ {\isacharequal}{\kern0pt}\ wf{\isacharunderscore}{\kern0pt}effect\ ty{\isacharunderscore}{\kern0pt}ent\ e{\isachardoublequoteclose}\isanewline
%
\isadelimproof
\ \ \ \ \ \ %
\endisadelimproof
%
\isatagproof
\isacommand{by}\isamarkupfalse%
\ {\isacharparenleft}{\kern0pt}cases\ e{\isacharparenright}{\kern0pt}\ {\isacharparenleft}{\kern0pt}auto\ simp{\isacharcolon}{\kern0pt}\ wf{\isacharunderscore}{\kern0pt}fmla{\isacharunderscore}{\kern0pt}atom{\isadigit{1}}{\isacharprime}{\kern0pt}{\isacharunderscore}{\kern0pt}correct{\isacharparenright}{\kern0pt}%
\endisatagproof
{\isafoldproof}%
%
\isadelimproof
\isanewline
%
\endisadelimproof
\isanewline
\ \ \ \ \isacommand{fun}\isamarkupfalse%
\ wf{\isacharunderscore}{\kern0pt}duration{\isacharunderscore}{\kern0pt}const{\isacharprime}{\kern0pt}\ {\isacharcolon}{\kern0pt}{\isacharcolon}{\kern0pt}\ {\isachardoublequoteopen}{\isacharunderscore}{\kern0pt}\ {\isasymRightarrow}\ {\isacharprime}{\kern0pt}ent\ duration{\isacharunderscore}{\kern0pt}constraint\ {\isasymRightarrow}\ bool{\isachardoublequoteclose}\ \isakeyword{where}\isanewline
\ \ \ \ \ \ {\isachardoublequoteopen}wf{\isacharunderscore}{\kern0pt}duration{\isacharunderscore}{\kern0pt}const{\isacharprime}{\kern0pt}\ stg\ No{\isacharunderscore}{\kern0pt}Const\ {\isasymlongleftrightarrow}\ True{\isachardoublequoteclose}\isanewline
\ \ \ \ {\isacharbar}{\kern0pt}\ {\isachardoublequoteopen}wf{\isacharunderscore}{\kern0pt}duration{\isacharunderscore}{\kern0pt}const{\isacharprime}{\kern0pt}\ stg\ {\isacharparenleft}{\kern0pt}Time{\isacharunderscore}{\kern0pt}Const\ op\ d{\isacharparenright}{\kern0pt}\ {\isasymlongleftrightarrow}\ d\ {\isasymge}\ {\isadigit{0}}{\isachardoublequoteclose}\isanewline
\ \ \ \ {\isacharbar}{\kern0pt}\ {\isachardoublequoteopen}wf{\isacharunderscore}{\kern0pt}duration{\isacharunderscore}{\kern0pt}const{\isacharprime}{\kern0pt}\ stg\ {\isacharparenleft}{\kern0pt}Func{\isacharunderscore}{\kern0pt}Const\ op\ f\ vs{\isacharparenright}{\kern0pt}\ {\isasymlongleftrightarrow}\ \isanewline
\ \ \ \ \ \ \ \ {\isacharparenleft}{\kern0pt}case\ func{\isacharunderscore}{\kern0pt}sig\ f\ of\isanewline
\ \ \ \ \ \ \ \ \ \ \ \ None\ {\isasymRightarrow}\ False\isanewline
\ \ \ \ \ \ \ \ \ \ {\isacharbar}{\kern0pt}\ Some\ Ts\ {\isasymRightarrow}\ list{\isacharunderscore}{\kern0pt}all{\isadigit{2}}\ {\isacharparenleft}{\kern0pt}is{\isacharunderscore}{\kern0pt}of{\isacharunderscore}{\kern0pt}type{\isacharprime}{\kern0pt}\ stg{\isacharparenright}{\kern0pt}\ vs\ Ts{\isacharparenright}{\kern0pt}{\isachardoublequoteclose}\isanewline
\isanewline
\ \ \ \ \isacommand{lemma}\isamarkupfalse%
\ wf{\isacharunderscore}{\kern0pt}duration{\isacharunderscore}{\kern0pt}const{\isacharprime}{\kern0pt}{\isacharunderscore}{\kern0pt}correct{\isacharcolon}{\kern0pt}\ {\isachardoublequoteopen}wf{\isacharunderscore}{\kern0pt}duration{\isacharunderscore}{\kern0pt}const{\isacharprime}{\kern0pt}\ STG\ d\ {\isacharequal}{\kern0pt}\ wf{\isacharunderscore}{\kern0pt}duration{\isacharunderscore}{\kern0pt}const\ ty{\isacharunderscore}{\kern0pt}ent\ d{\isachardoublequoteclose}\isanewline
%
\isadelimproof
\ \ \ \ \ \ %
\endisadelimproof
%
\isatagproof
\isacommand{by}\isamarkupfalse%
\ {\isacharparenleft}{\kern0pt}cases\ d{\isacharparenright}{\kern0pt}\ {\isacharparenleft}{\kern0pt}auto\ simp{\isacharcolon}{\kern0pt}\ is{\isacharunderscore}{\kern0pt}of{\isacharunderscore}{\kern0pt}type{\isacharprime}{\kern0pt}{\isacharunderscore}{\kern0pt}correct{\isacharbrackleft}{\kern0pt}abs{\isacharunderscore}{\kern0pt}def{\isacharbrackright}{\kern0pt}\ split{\isacharcolon}{\kern0pt}option{\isachardot}{\kern0pt}split{\isacharparenright}{\kern0pt}%
\endisatagproof
{\isafoldproof}%
%
\isadelimproof
\isanewline
%
\endisadelimproof
\isanewline
\ \ \isacommand{end}\isamarkupfalse%
\ %
\isamarkupcmt{Context fixing \isa{ty{\isacharunderscore}{\kern0pt}ent}%
}\isanewline
\isanewline
\ \ \isacommand{fun}\isamarkupfalse%
\ wf{\isacharunderscore}{\kern0pt}action{\isacharunderscore}{\kern0pt}schema{\isacharprime}{\kern0pt}\ {\isacharcolon}{\kern0pt}{\isacharcolon}{\kern0pt}\ {\isachardoublequoteopen}{\isacharunderscore}{\kern0pt}\ {\isasymRightarrow}\ {\isacharunderscore}{\kern0pt}\ {\isasymRightarrow}\ ast{\isacharunderscore}{\kern0pt}action{\isacharunderscore}{\kern0pt}schema\ {\isasymRightarrow}\ bool{\isachardoublequoteclose}\ \isakeyword{where}\isanewline
\ \ \ \ {\isachardoublequoteopen}wf{\isacharunderscore}{\kern0pt}action{\isacharunderscore}{\kern0pt}schema{\isacharprime}{\kern0pt}\ stg\ conT\ {\isacharparenleft}{\kern0pt}Simple{\isacharunderscore}{\kern0pt}Action{\isacharunderscore}{\kern0pt}Schema\ n\ params\ pre\ eff{\isacharparenright}{\kern0pt}\ {\isasymlongleftrightarrow}\ {\isacharparenleft}{\kern0pt}\isanewline
\ \ \ \ \ \ let\isanewline
\ \ \ \ \ \ \ \ tyv\ {\isacharequal}{\kern0pt}\ ty{\isacharunderscore}{\kern0pt}term{\isacharprime}{\kern0pt}\ {\isacharparenleft}{\kern0pt}map{\isacharunderscore}{\kern0pt}of\ params{\isacharparenright}{\kern0pt}\ conT\isanewline
\ \ \ \ \ \ in\isanewline
\ \ \ \ \ \ \ \ distinct\ {\isacharparenleft}{\kern0pt}map\ fst\ params{\isacharparenright}{\kern0pt}\isanewline
\ \ \ \ \ \ {\isasymand}\ wf{\isacharunderscore}{\kern0pt}fmla{\isacharprime}{\kern0pt}\ tyv\ stg\ pre\isanewline
\ \ \ \ \ \ {\isasymand}\ wf{\isacharunderscore}{\kern0pt}effect{\isacharprime}{\kern0pt}\ tyv\ stg\ eff{\isacharparenright}{\kern0pt}{\isachardoublequoteclose}\isanewline
\ \ {\isacharbar}{\kern0pt}\ {\isachardoublequoteopen}wf{\isacharunderscore}{\kern0pt}action{\isacharunderscore}{\kern0pt}schema{\isacharprime}{\kern0pt}\ stg\ conT\ {\isacharparenleft}{\kern0pt}Durative{\isacharunderscore}{\kern0pt}Action{\isacharunderscore}{\kern0pt}Schema\ n\ params\ d\ cond\ eff{\isacharparenright}{\kern0pt}\ {\isasymlongleftrightarrow}\ {\isacharparenleft}{\kern0pt}\isanewline
\ \ \ \ \ \ let\isanewline
\ \ \ \ \ \ \ \ tyv\ {\isacharequal}{\kern0pt}\ ty{\isacharunderscore}{\kern0pt}term{\isacharprime}{\kern0pt}\ {\isacharparenleft}{\kern0pt}map{\isacharunderscore}{\kern0pt}of\ params{\isacharparenright}{\kern0pt}\ conT\isanewline
\ \ \ \ \ \ in\isanewline
\ \ \ \ \ \ \ \ distinct\ {\isacharparenleft}{\kern0pt}map\ fst\ params{\isacharparenright}{\kern0pt}\isanewline
\ \ \ \ \ \ {\isasymand}\ {\isacharparenleft}{\kern0pt}{\isasymforall}{\isacharparenleft}{\kern0pt}t{\isacharcomma}{\kern0pt}c{\isacharparenright}{\kern0pt}\ {\isasymin}\ set\ cond{\isachardot}{\kern0pt}\ wf{\isacharunderscore}{\kern0pt}fmla{\isacharprime}{\kern0pt}\ tyv\ stg\ c{\isacharparenright}{\kern0pt}\isanewline
\ \ \ \ \ \ {\isasymand}\ {\isacharparenleft}{\kern0pt}{\isasymforall}{\isacharparenleft}{\kern0pt}t{\isacharcomma}{\kern0pt}e{\isacharparenright}{\kern0pt}\ {\isasymin}\ set\ eff{\isachardot}{\kern0pt}\ wf{\isacharunderscore}{\kern0pt}effect{\isacharprime}{\kern0pt}\ tyv\ stg\ e\ {\isasymand}\ t\ {\isasymnoteq}\ Over{\isacharunderscore}{\kern0pt}All{\isacharparenright}{\kern0pt}\isanewline
\ \ \ \ \ \ {\isasymand}\ wf{\isacharunderscore}{\kern0pt}duration{\isacharunderscore}{\kern0pt}const{\isacharprime}{\kern0pt}\ tyv\ stg\ d{\isacharparenright}{\kern0pt}{\isachardoublequoteclose}\isanewline
\isanewline
\ \ \isacommand{lemma}\isamarkupfalse%
\ wf{\isacharunderscore}{\kern0pt}action{\isacharunderscore}{\kern0pt}schema{\isacharprime}{\kern0pt}{\isacharunderscore}{\kern0pt}correct{\isacharcolon}{\kern0pt}\ {\isachardoublequoteopen}wf{\isacharunderscore}{\kern0pt}action{\isacharunderscore}{\kern0pt}schema{\isacharprime}{\kern0pt}\ STG\ mp{\isacharunderscore}{\kern0pt}constT\ s\ {\isacharequal}{\kern0pt}\ wf{\isacharunderscore}{\kern0pt}action{\isacharunderscore}{\kern0pt}schema\ s{\isachardoublequoteclose}\isanewline
%
\isadelimproof
\ \ \ \ %
\endisadelimproof
%
\isatagproof
\isacommand{by}\isamarkupfalse%
\ {\isacharparenleft}{\kern0pt}cases\ s{\isacharparenright}{\kern0pt}\ \isanewline
\ \ \ \ \ \ \ {\isacharparenleft}{\kern0pt}auto\ simp{\isacharcolon}{\kern0pt}\ wf{\isacharunderscore}{\kern0pt}fmla{\isacharprime}{\kern0pt}{\isacharunderscore}{\kern0pt}correct\ wf{\isacharunderscore}{\kern0pt}effect{\isacharprime}{\kern0pt}{\isacharunderscore}{\kern0pt}correct\ wf{\isacharunderscore}{\kern0pt}duration{\isacharunderscore}{\kern0pt}const{\isacharprime}{\kern0pt}{\isacharunderscore}{\kern0pt}correct\ Let{\isacharunderscore}{\kern0pt}def\ \isanewline
\ \ \ \ \ \ \ \ \ \ \ \ \ split{\isacharcolon}{\kern0pt}\ option{\isachardot}{\kern0pt}splits{\isacharparenright}{\kern0pt}%
\endisatagproof
{\isafoldproof}%
%
\isadelimproof
\isanewline
%
\endisadelimproof
\isanewline
\ \ \isacommand{definition}\isamarkupfalse%
\ wf{\isacharunderscore}{\kern0pt}domain{\isacharprime}{\kern0pt}\ {\isacharcolon}{\kern0pt}{\isacharcolon}{\kern0pt}\ {\isachardoublequoteopen}{\isacharunderscore}{\kern0pt}\ {\isasymRightarrow}\ {\isacharunderscore}{\kern0pt}\ {\isasymRightarrow}\ bool{\isachardoublequoteclose}\ \isakeyword{where}\isanewline
\ \ \ \ {\isachardoublequoteopen}wf{\isacharunderscore}{\kern0pt}domain{\isacharprime}{\kern0pt}\ stg\ conT\ {\isasymequiv}\isanewline
\ \ \ \ \ \ wf{\isacharunderscore}{\kern0pt}types\isanewline
\ \ \ \ {\isasymand}\ distinct\ {\isacharparenleft}{\kern0pt}map\ {\isacharparenleft}{\kern0pt}predicate{\isacharunderscore}{\kern0pt}decl{\isachardot}{\kern0pt}pred{\isacharparenright}{\kern0pt}\ {\isacharparenleft}{\kern0pt}predicates\ D{\isacharparenright}{\kern0pt}{\isacharparenright}{\kern0pt}\isanewline
\ \ \ \ {\isasymand}\ {\isacharparenleft}{\kern0pt}{\isasymforall}p{\isasymin}set\ {\isacharparenleft}{\kern0pt}predicates\ D{\isacharparenright}{\kern0pt}{\isachardot}{\kern0pt}\ wf{\isacharunderscore}{\kern0pt}predicate{\isacharunderscore}{\kern0pt}decl\ p{\isacharparenright}{\kern0pt}\isanewline
\ \ \ \ {\isasymand}\ distinct\ {\isacharparenleft}{\kern0pt}map\ {\isacharparenleft}{\kern0pt}function{\isacharunderscore}{\kern0pt}decl{\isachardot}{\kern0pt}func{\isacharparenright}{\kern0pt}\ {\isacharparenleft}{\kern0pt}functions\ D{\isacharparenright}{\kern0pt}{\isacharparenright}{\kern0pt}\isanewline
\ \ \ \ {\isasymand}\ {\isacharparenleft}{\kern0pt}{\isasymforall}f{\isasymin}set\ {\isacharparenleft}{\kern0pt}functions\ D{\isacharparenright}{\kern0pt}{\isachardot}{\kern0pt}\ wf{\isacharunderscore}{\kern0pt}function{\isacharunderscore}{\kern0pt}decl\ f{\isacharparenright}{\kern0pt}\isanewline
\ \ \ \ {\isasymand}\ distinct\ {\isacharparenleft}{\kern0pt}map\ fst\ {\isacharparenleft}{\kern0pt}consts\ D{\isacharparenright}{\kern0pt}{\isacharparenright}{\kern0pt}\isanewline
\ \ \ \ {\isasymand}\ {\isacharparenleft}{\kern0pt}{\isasymforall}{\isacharparenleft}{\kern0pt}n{\isacharcomma}{\kern0pt}T{\isacharparenright}{\kern0pt}{\isasymin}set\ {\isacharparenleft}{\kern0pt}consts\ D{\isacharparenright}{\kern0pt}{\isachardot}{\kern0pt}\ wf{\isacharunderscore}{\kern0pt}type\ T{\isacharparenright}{\kern0pt}\isanewline
\ \ \ \ {\isasymand}\ distinct\ {\isacharparenleft}{\kern0pt}map\ ast{\isacharunderscore}{\kern0pt}action{\isacharunderscore}{\kern0pt}schema{\isachardot}{\kern0pt}name\ {\isacharparenleft}{\kern0pt}actions\ D{\isacharparenright}{\kern0pt}{\isacharparenright}{\kern0pt}\isanewline
\ \ \ \ {\isasymand}\ {\isacharparenleft}{\kern0pt}{\isasymforall}a{\isasymin}set\ {\isacharparenleft}{\kern0pt}actions\ D{\isacharparenright}{\kern0pt}{\isachardot}{\kern0pt}\ wf{\isacharunderscore}{\kern0pt}action{\isacharunderscore}{\kern0pt}schema{\isacharprime}{\kern0pt}\ stg\ conT\ a{\isacharparenright}{\kern0pt}\isanewline
\ \ \ \ {\isachardoublequoteclose}\isanewline
\isanewline
\ \ \isacommand{lemma}\isamarkupfalse%
\ wf{\isacharunderscore}{\kern0pt}domain{\isacharprime}{\kern0pt}{\isacharunderscore}{\kern0pt}correct{\isacharcolon}{\kern0pt}\ {\isachardoublequoteopen}wf{\isacharunderscore}{\kern0pt}domain{\isacharprime}{\kern0pt}\ STG\ mp{\isacharunderscore}{\kern0pt}constT\ {\isacharequal}{\kern0pt}\ wf{\isacharunderscore}{\kern0pt}domain{\isachardoublequoteclose}\isanewline
%
\isadelimproof
\ \ \ \ %
\endisadelimproof
%
\isatagproof
\isacommand{unfolding}\isamarkupfalse%
\ wf{\isacharunderscore}{\kern0pt}domain{\isacharunderscore}{\kern0pt}def\ wf{\isacharunderscore}{\kern0pt}domain{\isacharprime}{\kern0pt}{\isacharunderscore}{\kern0pt}def\isanewline
\ \ \ \ \isacommand{by}\isamarkupfalse%
\ {\isacharparenleft}{\kern0pt}auto\ simp{\isacharcolon}{\kern0pt}\ wf{\isacharunderscore}{\kern0pt}action{\isacharunderscore}{\kern0pt}schema{\isacharprime}{\kern0pt}{\isacharunderscore}{\kern0pt}correct{\isacharparenright}{\kern0pt}%
\endisatagproof
{\isafoldproof}%
%
\isadelimproof
\isanewline
%
\endisadelimproof
\isanewline
\isanewline
\isacommand{end}\isamarkupfalse%
\ %
\isamarkupcmt{Context of \isa{ast{\isacharunderscore}{\kern0pt}domain}%
}%
\isadelimdocument
%
\endisadelimdocument
%
\isatagdocument
%
\isamarkupsubsubsection{Ground Action Interference \& Application of Effects%
}
\isamarkuptrue%
%
\endisatagdocument
{\isafolddocument}%
%
\isadelimdocument
%
\endisadelimdocument
\isacommand{context}\isamarkupfalse%
\ ast{\isacharunderscore}{\kern0pt}problem\ \isakeyword{begin}%
\begin{isamarkuptext}%
Implementation of executable refinements for \isa{acts{\isacharunderscore}{\kern0pt}non{\isacharunderscore}{\kern0pt}intrf} 
  \& \isa{apply{\isacharunderscore}{\kern0pt}happ}%
\end{isamarkuptext}\isamarkuptrue%
%
\begin{isamarkuptext}%
a list of simulataneous \& non-interfering grounded actions are all applied at once%
\end{isamarkuptext}\isamarkuptrue%
\ \ \isacommand{fun}\isamarkupfalse%
\ apply{\isacharunderscore}{\kern0pt}happ{\isacharunderscore}{\kern0pt}exec\ {\isacharcolon}{\kern0pt}{\isacharcolon}{\kern0pt}\ {\isachardoublequoteopen}happening\ {\isasymRightarrow}\ world{\isacharunderscore}{\kern0pt}model\ {\isasymRightarrow}\ world{\isacharunderscore}{\kern0pt}model{\isachardoublequoteclose}\ \isakeyword{where}\isanewline
\ \ \ \ {\isachardoublequoteopen}apply{\isacharunderscore}{\kern0pt}happ{\isacharunderscore}{\kern0pt}exec\ {\isacharparenleft}{\kern0pt}t\isactrlsub i{\isacharcomma}{\kern0pt}A\isactrlsub i{\isacharparenright}{\kern0pt}\ s\ {\isacharequal}{\kern0pt}\ \isanewline
\ \ \ \ \ \ {\isacharparenleft}{\kern0pt}let\ \isanewline
\ \ \ \ \ \ \ \ adds\ {\isacharequal}{\kern0pt}\ {\isacharparenleft}{\kern0pt}concat\ o\ {\isacharparenleft}{\kern0pt}map\ {\isacharparenleft}{\kern0pt}adds\ o\ effect{\isacharparenright}{\kern0pt}{\isacharparenright}{\kern0pt}{\isacharparenright}{\kern0pt}\ A\isactrlsub i{\isacharsemicolon}{\kern0pt}\ \isanewline
\ \ \ \ \ \ \ \ dels\ {\isacharequal}{\kern0pt}\ {\isacharparenleft}{\kern0pt}concat\ o\ {\isacharparenleft}{\kern0pt}map\ {\isacharparenleft}{\kern0pt}dels\ o\ effect{\isacharparenright}{\kern0pt}{\isacharparenright}{\kern0pt}{\isacharparenright}{\kern0pt}\ A\isactrlsub i\ in\ \isanewline
\ \ \ \ \ \ \ \ fold\ Set{\isachardot}{\kern0pt}insert\ adds\ {\isacharparenleft}{\kern0pt}fold\ Set{\isachardot}{\kern0pt}remove\ dels\ s{\isacharparenright}{\kern0pt}{\isacharparenright}{\kern0pt}{\isachardoublequoteclose}\isanewline
\isanewline
\ \ \isacommand{lemma}\isamarkupfalse%
\ fold{\isacharunderscore}{\kern0pt}set{\isacharunderscore}{\kern0pt}remove{\isacharcolon}{\kern0pt}\ {\isachardoublequoteopen}fold\ Set{\isachardot}{\kern0pt}remove\ xs\ s\ {\isacharequal}{\kern0pt}\ s\ {\isacharminus}{\kern0pt}\ set\ xs{\isachardoublequoteclose}\isanewline
%
\isadelimproof
\ \ \ \ %
\endisadelimproof
%
\isatagproof
\isacommand{by}\isamarkupfalse%
\ {\isacharparenleft}{\kern0pt}induction\ xs\ arbitrary{\isacharcolon}{\kern0pt}\ s{\isacharparenright}{\kern0pt}\ auto%
\endisatagproof
{\isafoldproof}%
%
\isadelimproof
\isanewline
%
\endisadelimproof
\isanewline
\ \ \isacommand{lemma}\isamarkupfalse%
\ fold{\isacharunderscore}{\kern0pt}set{\isacharunderscore}{\kern0pt}insert{\isacharcolon}{\kern0pt}\ {\isachardoublequoteopen}fold\ Set{\isachardot}{\kern0pt}insert\ xs\ s\ {\isacharequal}{\kern0pt}\ s\ {\isasymunion}\ set\ xs{\isachardoublequoteclose}\isanewline
%
\isadelimproof
\ \ \ \ %
\endisadelimproof
%
\isatagproof
\isacommand{by}\isamarkupfalse%
\ {\isacharparenleft}{\kern0pt}induction\ xs\ arbitrary{\isacharcolon}{\kern0pt}\ s{\isacharparenright}{\kern0pt}\ auto%
\endisatagproof
{\isafoldproof}%
%
\isadelimproof
%
\endisadelimproof
%
\begin{isamarkuptext}%
Justification of refinement for \isa{apply{\isacharunderscore}{\kern0pt}happ}%
\end{isamarkuptext}\isamarkuptrue%
\ \ \isacommand{lemma}\isamarkupfalse%
\ apply{\isacharunderscore}{\kern0pt}happ{\isacharunderscore}{\kern0pt}exec{\isacharunderscore}{\kern0pt}refine{\isacharcolon}{\kern0pt}\ {\isachardoublequoteopen}apply{\isacharunderscore}{\kern0pt}happ{\isacharunderscore}{\kern0pt}exec\ h\ s\ {\isacharequal}{\kern0pt}\ apply{\isacharunderscore}{\kern0pt}happ\ h\ s{\isachardoublequoteclose}\isanewline
%
\isadelimproof
\ \ \ \ %
\endisadelimproof
%
\isatagproof
\isacommand{by}\isamarkupfalse%
\ {\isacharparenleft}{\kern0pt}cases\ h{\isacharparenright}{\kern0pt}\ {\isacharparenleft}{\kern0pt}auto\ simp{\isacharcolon}{\kern0pt}\ fold{\isacharunderscore}{\kern0pt}set{\isacharunderscore}{\kern0pt}remove\ fold{\isacharunderscore}{\kern0pt}set{\isacharunderscore}{\kern0pt}insert{\isacharparenright}{\kern0pt}%
\endisatagproof
{\isafoldproof}%
%
\isadelimproof
\isanewline
%
\endisadelimproof
\isanewline
\isacommand{end}\isamarkupfalse%
\ %
\isamarkupcmt{Context of \isa{ast{\isacharunderscore}{\kern0pt}problem}%
}%
\isadelimdocument
%
\endisadelimdocument
%
\isatagdocument
%
\isamarkupsubsubsection{Well-Formedness%
}
\isamarkuptrue%
%
\endisatagdocument
{\isafolddocument}%
%
\isadelimdocument
%
\endisadelimdocument
\isacommand{context}\isamarkupfalse%
\ ast{\isacharunderscore}{\kern0pt}problem\ \isakeyword{begin}%
\begin{isamarkuptext}%
We start by defining a mapping from objects to types. The container
    framework will generate efficient, red-black tree based code for that
    later.%
\end{isamarkuptext}\isamarkuptrue%
\ \ \isacommand{type{\isacharunderscore}{\kern0pt}synonym}\isamarkupfalse%
\ objT\ {\isacharequal}{\kern0pt}\ {\isachardoublequoteopen}{\isacharparenleft}{\kern0pt}object{\isacharcomma}{\kern0pt}\ type{\isacharparenright}{\kern0pt}\ mapping{\isachardoublequoteclose}\isanewline
\isanewline
\ \ \isacommand{definition}\isamarkupfalse%
\ mp{\isacharunderscore}{\kern0pt}objT\ {\isacharcolon}{\kern0pt}{\isacharcolon}{\kern0pt}\ {\isachardoublequoteopen}{\isacharparenleft}{\kern0pt}object{\isacharcomma}{\kern0pt}\ type{\isacharparenright}{\kern0pt}\ mapping{\isachardoublequoteclose}\ \isakeyword{where}\isanewline
\ \ \ \ {\isachardoublequoteopen}mp{\isacharunderscore}{\kern0pt}objT\ {\isacharequal}{\kern0pt}\ Mapping{\isachardot}{\kern0pt}of{\isacharunderscore}{\kern0pt}alist\ {\isacharparenleft}{\kern0pt}consts\ D\ {\isacharat}{\kern0pt}\ objects\ P{\isacharparenright}{\kern0pt}{\isachardoublequoteclose}\isanewline
\isanewline
\ \ \isacommand{lemma}\isamarkupfalse%
\ mp{\isacharunderscore}{\kern0pt}objT{\isacharunderscore}{\kern0pt}correct{\isacharbrackleft}{\kern0pt}simp{\isacharbrackright}{\kern0pt}{\isacharcolon}{\kern0pt}\ {\isachardoublequoteopen}Mapping{\isachardot}{\kern0pt}lookup\ mp{\isacharunderscore}{\kern0pt}objT\ {\isacharequal}{\kern0pt}\ objT{\isachardoublequoteclose}\isanewline
%
\isadelimproof
\ \ \ \ %
\endisadelimproof
%
\isatagproof
\isacommand{unfolding}\isamarkupfalse%
\ mp{\isacharunderscore}{\kern0pt}objT{\isacharunderscore}{\kern0pt}def\ objT{\isacharunderscore}{\kern0pt}alt\isanewline
\ \ \ \ \isacommand{by}\isamarkupfalse%
\ transfer\ {\isacharparenleft}{\kern0pt}simp\ add{\isacharcolon}{\kern0pt}\ Map{\isacharunderscore}{\kern0pt}To{\isacharunderscore}{\kern0pt}Mapping{\isachardot}{\kern0pt}map{\isacharunderscore}{\kern0pt}apply{\isacharunderscore}{\kern0pt}def{\isacharparenright}{\kern0pt}%
\endisatagproof
{\isafoldproof}%
%
\isadelimproof
%
\endisadelimproof
%
\begin{isamarkuptext}%
We refine the typecheck to use the mapping%
\end{isamarkuptext}\isamarkuptrue%
\ \ \isacommand{definition}\isamarkupfalse%
\ {\isachardoublequoteopen}is{\isacharunderscore}{\kern0pt}obj{\isacharunderscore}{\kern0pt}of{\isacharunderscore}{\kern0pt}type{\isacharunderscore}{\kern0pt}impl\ stg\ mp\ n\ T\ {\isacharequal}{\kern0pt}\ {\isacharparenleft}{\kern0pt}\isanewline
\ \ \ \ case\ Mapping{\isachardot}{\kern0pt}lookup\ mp\ n\ of\ None\ {\isasymRightarrow}\ False\ {\isacharbar}{\kern0pt}\ Some\ oT\ {\isasymRightarrow}\ of{\isacharunderscore}{\kern0pt}type{\isacharunderscore}{\kern0pt}impl\ stg\ oT\ T\isanewline
\ \ {\isacharparenright}{\kern0pt}{\isachardoublequoteclose}\isanewline
\isanewline
\ \ \isacommand{lemma}\isamarkupfalse%
\ is{\isacharunderscore}{\kern0pt}obj{\isacharunderscore}{\kern0pt}of{\isacharunderscore}{\kern0pt}type{\isacharunderscore}{\kern0pt}impl{\isacharunderscore}{\kern0pt}correct{\isacharbrackleft}{\kern0pt}simp{\isacharbrackright}{\kern0pt}{\isacharcolon}{\kern0pt}\isanewline
\ \ \ \ {\isachardoublequoteopen}is{\isacharunderscore}{\kern0pt}obj{\isacharunderscore}{\kern0pt}of{\isacharunderscore}{\kern0pt}type{\isacharunderscore}{\kern0pt}impl\ STG\ mp{\isacharunderscore}{\kern0pt}objT\ {\isacharequal}{\kern0pt}\ is{\isacharunderscore}{\kern0pt}obj{\isacharunderscore}{\kern0pt}of{\isacharunderscore}{\kern0pt}type{\isachardoublequoteclose}\isanewline
%
\isadelimproof
\ \ \ \ %
\endisadelimproof
%
\isatagproof
\isacommand{apply}\isamarkupfalse%
\ {\isacharparenleft}{\kern0pt}intro\ ext{\isacharparenright}{\kern0pt}\isanewline
\ \ \ \ \isacommand{apply}\isamarkupfalse%
\ {\isacharparenleft}{\kern0pt}auto\ simp{\isacharcolon}{\kern0pt}\ is{\isacharunderscore}{\kern0pt}obj{\isacharunderscore}{\kern0pt}of{\isacharunderscore}{\kern0pt}type{\isacharunderscore}{\kern0pt}impl{\isacharunderscore}{\kern0pt}def\ is{\isacharunderscore}{\kern0pt}obj{\isacharunderscore}{\kern0pt}of{\isacharunderscore}{\kern0pt}type{\isacharunderscore}{\kern0pt}def\ of{\isacharunderscore}{\kern0pt}type{\isacharunderscore}{\kern0pt}impl{\isacharunderscore}{\kern0pt}correct\ split{\isacharcolon}{\kern0pt}\ option{\isachardot}{\kern0pt}split{\isacharparenright}{\kern0pt}\isanewline
\ \ \ \ \isacommand{done}\isamarkupfalse%
%
\endisatagproof
{\isafoldproof}%
%
\isadelimproof
\isanewline
%
\endisadelimproof
\isanewline
\ \ \isacommand{fun}\isamarkupfalse%
\ wf{\isacharunderscore}{\kern0pt}func{\isacharunderscore}{\kern0pt}assign{\isacharprime}{\kern0pt}\ {\isacharcolon}{\kern0pt}{\isacharcolon}{\kern0pt}\ {\isachardoublequoteopen}objT\ {\isasymRightarrow}\ {\isacharunderscore}{\kern0pt}\ {\isasymRightarrow}\ object\ atom\ formula\ {\isasymRightarrow}\ bool{\isachardoublequoteclose}\ \isakeyword{where}\ \isanewline
\ \ \ \ {\isachardoublequoteopen}wf{\isacharunderscore}{\kern0pt}func{\isacharunderscore}{\kern0pt}assign{\isacharprime}{\kern0pt}\ ot\ stg\ {\isacharparenleft}{\kern0pt}Atom\ {\isacharparenleft}{\kern0pt}eqAtm\ {\isacharparenleft}{\kern0pt}FuncEnt\ f\ args{\isacharparenright}{\kern0pt}\ {\isacharparenleft}{\kern0pt}TimeEnt\ d{\isacharparenright}{\kern0pt}{\isacharparenright}{\kern0pt}{\isacharparenright}{\kern0pt}\ \isanewline
\ \ \ \ \ \ {\isasymlongleftrightarrow}\ wf{\isacharunderscore}{\kern0pt}func{\isacharunderscore}{\kern0pt}args{\isacharprime}{\kern0pt}\ {\isacharparenleft}{\kern0pt}Mapping{\isachardot}{\kern0pt}lookup\ ot{\isacharparenright}{\kern0pt}\ stg\ {\isacharparenleft}{\kern0pt}f{\isacharcomma}{\kern0pt}args{\isacharparenright}{\kern0pt}{\isachardoublequoteclose}\isanewline
\ \ {\isacharbar}{\kern0pt}\ {\isachardoublequoteopen}wf{\isacharunderscore}{\kern0pt}func{\isacharunderscore}{\kern0pt}assign{\isacharprime}{\kern0pt}\ ot\ stg\ {\isacharunderscore}{\kern0pt}\ {\isasymlongleftrightarrow}\ False{\isachardoublequoteclose}\isanewline
\isanewline
\ \ \isacommand{lemma}\isamarkupfalse%
\ wf{\isacharunderscore}{\kern0pt}func{\isacharunderscore}{\kern0pt}asign{\isacharprime}{\kern0pt}{\isacharunderscore}{\kern0pt}correct{\isacharbrackleft}{\kern0pt}simp{\isacharbrackright}{\kern0pt}{\isacharcolon}{\kern0pt}\ {\isachardoublequoteopen}wf{\isacharunderscore}{\kern0pt}func{\isacharunderscore}{\kern0pt}assign{\isacharprime}{\kern0pt}\ mp{\isacharunderscore}{\kern0pt}objT\ STG\ f\ {\isacharequal}{\kern0pt}\ wf{\isacharunderscore}{\kern0pt}func{\isacharunderscore}{\kern0pt}assign\ f{\isachardoublequoteclose}\isanewline
%
\isadelimproof
\ \ \ \ %
\endisadelimproof
%
\isatagproof
\isacommand{by}\isamarkupfalse%
\ {\isacharparenleft}{\kern0pt}cases\ f\ rule{\isacharcolon}{\kern0pt}\ wf{\isacharunderscore}{\kern0pt}func{\isacharunderscore}{\kern0pt}assign{\isachardot}{\kern0pt}cases{\isacharparenright}{\kern0pt}\ {\isacharparenleft}{\kern0pt}auto\ simp{\isacharcolon}{\kern0pt}\ wf{\isacharunderscore}{\kern0pt}func{\isacharunderscore}{\kern0pt}args{\isacharprime}{\kern0pt}{\isacharunderscore}{\kern0pt}correct{\isacharbrackleft}{\kern0pt}abs{\isacharunderscore}{\kern0pt}def{\isacharbrackright}{\kern0pt}{\isacharparenright}{\kern0pt}%
\endisatagproof
{\isafoldproof}%
%
\isadelimproof
%
\endisadelimproof
%
\begin{isamarkuptext}%
We refine the well-formedness checks to use the mapping%
\end{isamarkuptext}\isamarkuptrue%
\ \ \isacommand{definition}\isamarkupfalse%
\ wf{\isacharunderscore}{\kern0pt}fact{\isacharprime}{\kern0pt}\ {\isacharcolon}{\kern0pt}{\isacharcolon}{\kern0pt}\ {\isachardoublequoteopen}objT\ {\isasymRightarrow}\ {\isacharunderscore}{\kern0pt}\ {\isasymRightarrow}\ fact\ {\isasymRightarrow}\ bool{\isachardoublequoteclose}\ \isakeyword{where}\isanewline
\ \ \ \ {\isachardoublequoteopen}wf{\isacharunderscore}{\kern0pt}fact{\isacharprime}{\kern0pt}\ ot\ stg\ {\isasymequiv}\ wf{\isacharunderscore}{\kern0pt}pred{\isacharunderscore}{\kern0pt}atom{\isacharprime}{\kern0pt}\ {\isacharparenleft}{\kern0pt}Mapping{\isachardot}{\kern0pt}lookup\ ot{\isacharparenright}{\kern0pt}\ stg{\isachardoublequoteclose}\isanewline
\isanewline
\ \ \isacommand{lemma}\isamarkupfalse%
\ wf{\isacharunderscore}{\kern0pt}fact{\isacharprime}{\kern0pt}{\isacharunderscore}{\kern0pt}correct{\isacharbrackleft}{\kern0pt}simp{\isacharbrackright}{\kern0pt}{\isacharcolon}{\kern0pt}\ {\isachardoublequoteopen}wf{\isacharunderscore}{\kern0pt}fact{\isacharprime}{\kern0pt}\ mp{\isacharunderscore}{\kern0pt}objT\ STG\ {\isacharequal}{\kern0pt}\ wf{\isacharunderscore}{\kern0pt}fact{\isachardoublequoteclose}\isanewline
%
\isadelimproof
\ \ \ \ %
\endisadelimproof
%
\isatagproof
\isacommand{by}\isamarkupfalse%
\ {\isacharparenleft}{\kern0pt}auto\ simp{\isacharcolon}{\kern0pt}\ wf{\isacharunderscore}{\kern0pt}fact{\isacharprime}{\kern0pt}{\isacharunderscore}{\kern0pt}def\ wf{\isacharunderscore}{\kern0pt}fact{\isacharunderscore}{\kern0pt}def\ wf{\isacharunderscore}{\kern0pt}pred{\isacharunderscore}{\kern0pt}atom{\isacharprime}{\kern0pt}{\isacharunderscore}{\kern0pt}correct{\isacharbrackleft}{\kern0pt}abs{\isacharunderscore}{\kern0pt}def{\isacharbrackright}{\kern0pt}{\isacharparenright}{\kern0pt}%
\endisatagproof
{\isafoldproof}%
%
\isadelimproof
\isanewline
%
\endisadelimproof
\isanewline
\isanewline
\ \ \isacommand{definition}\isamarkupfalse%
\ {\isachardoublequoteopen}wf{\isacharunderscore}{\kern0pt}fmla{\isacharunderscore}{\kern0pt}atom{\isadigit{2}}{\isacharprime}{\kern0pt}\ mp\ stg\ f\isanewline
\ \ \ \ {\isacharequal}{\kern0pt}\ {\isacharparenleft}{\kern0pt}case\ f\ of\ formula{\isachardot}{\kern0pt}Atom\ {\isacharparenleft}{\kern0pt}predAtm\ p\ vs{\isacharparenright}{\kern0pt}\ {\isasymRightarrow}\ {\isacharparenleft}{\kern0pt}wf{\isacharunderscore}{\kern0pt}fact{\isacharprime}{\kern0pt}\ mp\ stg\ {\isacharparenleft}{\kern0pt}p{\isacharcomma}{\kern0pt}vs{\isacharparenright}{\kern0pt}{\isacharparenright}{\kern0pt}\ {\isacharbar}{\kern0pt}\ {\isacharunderscore}{\kern0pt}\ {\isasymRightarrow}\ False{\isacharparenright}{\kern0pt}{\isachardoublequoteclose}\isanewline
\isanewline
\ \ \isacommand{lemma}\isamarkupfalse%
\ wf{\isacharunderscore}{\kern0pt}fmla{\isacharunderscore}{\kern0pt}atom{\isadigit{2}}{\isacharprime}{\kern0pt}{\isacharunderscore}{\kern0pt}correct{\isacharbrackleft}{\kern0pt}simp{\isacharbrackright}{\kern0pt}{\isacharcolon}{\kern0pt}\isanewline
\ \ \ \ {\isachardoublequoteopen}wf{\isacharunderscore}{\kern0pt}fmla{\isacharunderscore}{\kern0pt}atom{\isadigit{2}}{\isacharprime}{\kern0pt}\ mp{\isacharunderscore}{\kern0pt}objT\ STG\ {\isasymphi}\ {\isacharequal}{\kern0pt}\ wf{\isacharunderscore}{\kern0pt}fmla{\isacharunderscore}{\kern0pt}atom\ objT\ {\isasymphi}{\isachardoublequoteclose}\isanewline
%
\isadelimproof
\ \ \ \ %
\endisadelimproof
%
\isatagproof
\isacommand{apply}\isamarkupfalse%
\ {\isacharparenleft}{\kern0pt}cases\ {\isasymphi}\ rule{\isacharcolon}{\kern0pt}\ wf{\isacharunderscore}{\kern0pt}fmla{\isacharunderscore}{\kern0pt}atom{\isachardot}{\kern0pt}cases{\isacharparenright}{\kern0pt}\isanewline
\ \ \ \ \isacommand{by}\isamarkupfalse%
\ {\isacharparenleft}{\kern0pt}auto\ simp{\isacharcolon}{\kern0pt}\ wf{\isacharunderscore}{\kern0pt}fmla{\isacharunderscore}{\kern0pt}atom{\isadigit{2}}{\isacharprime}{\kern0pt}{\isacharunderscore}{\kern0pt}def\ wf{\isacharunderscore}{\kern0pt}fact{\isacharunderscore}{\kern0pt}def\ split{\isacharcolon}{\kern0pt}\ option{\isachardot}{\kern0pt}splits{\isacharparenright}{\kern0pt}%
\endisatagproof
{\isafoldproof}%
%
\isadelimproof
\isanewline
%
\endisadelimproof
\isanewline
\ \ \isacommand{definition}\isamarkupfalse%
\ {\isachardoublequoteopen}wf{\isacharunderscore}{\kern0pt}problem{\isacharprime}{\kern0pt}\ stg\ conT\ mp\ {\isasymequiv}\isanewline
\ \ \ \ \ \ wf{\isacharunderscore}{\kern0pt}domain{\isacharprime}{\kern0pt}\ stg\ conT\isanewline
\ \ \ \ {\isasymand}\ distinct\ {\isacharparenleft}{\kern0pt}map\ fst\ {\isacharparenleft}{\kern0pt}objects\ P{\isacharparenright}{\kern0pt}\ {\isacharat}{\kern0pt}\ map\ fst\ {\isacharparenleft}{\kern0pt}consts\ D{\isacharparenright}{\kern0pt}{\isacharparenright}{\kern0pt}\isanewline
\ \ \ \ {\isasymand}\ {\isacharparenleft}{\kern0pt}{\isasymforall}{\isacharparenleft}{\kern0pt}n{\isacharcomma}{\kern0pt}T{\isacharparenright}{\kern0pt}{\isasymin}set\ {\isacharparenleft}{\kern0pt}objects\ P{\isacharparenright}{\kern0pt}{\isachardot}{\kern0pt}\ wf{\isacharunderscore}{\kern0pt}type\ T{\isacharparenright}{\kern0pt}\isanewline
\ \ \ \ {\isasymand}\ distinct\ {\isacharparenleft}{\kern0pt}init\ P{\isacharparenright}{\kern0pt}\isanewline
\ \ \ \ {\isasymand}\ {\isacharparenleft}{\kern0pt}{\isasymforall}f{\isasymin}set\ {\isacharparenleft}{\kern0pt}init\ P{\isacharparenright}{\kern0pt}{\isachardot}{\kern0pt}\ wf{\isacharunderscore}{\kern0pt}fmla{\isacharunderscore}{\kern0pt}atom{\isadigit{2}}{\isacharprime}{\kern0pt}\ mp\ stg\ f\ {\isasymor}\ wf{\isacharunderscore}{\kern0pt}func{\isacharunderscore}{\kern0pt}assign{\isacharprime}{\kern0pt}\ mp\ stg\ f{\isacharparenright}{\kern0pt}\isanewline
\ \ \ \ {\isasymand}\ wf{\isacharunderscore}{\kern0pt}fmla{\isacharprime}{\kern0pt}\ {\isacharparenleft}{\kern0pt}Mapping{\isachardot}{\kern0pt}lookup\ mp{\isacharparenright}{\kern0pt}\ stg\ {\isacharparenleft}{\kern0pt}goal\ P{\isacharparenright}{\kern0pt}{\isachardoublequoteclose}\isanewline
\isanewline
\ \ \isacommand{lemma}\isamarkupfalse%
\ wf{\isacharunderscore}{\kern0pt}problem{\isacharprime}{\kern0pt}{\isacharunderscore}{\kern0pt}correct{\isacharcolon}{\kern0pt}\isanewline
\ \ \ \ {\isachardoublequoteopen}wf{\isacharunderscore}{\kern0pt}problem{\isacharprime}{\kern0pt}\ STG\ mp{\isacharunderscore}{\kern0pt}constT\ mp{\isacharunderscore}{\kern0pt}objT\ {\isacharequal}{\kern0pt}\ wf{\isacharunderscore}{\kern0pt}problem{\isachardoublequoteclose}\isanewline
%
\isadelimproof
\ \ \ \ %
\endisadelimproof
%
\isatagproof
\isacommand{unfolding}\isamarkupfalse%
\ wf{\isacharunderscore}{\kern0pt}problem{\isacharunderscore}{\kern0pt}def\ wf{\isacharunderscore}{\kern0pt}problem{\isacharprime}{\kern0pt}{\isacharunderscore}{\kern0pt}def\ wf{\isacharunderscore}{\kern0pt}world{\isacharunderscore}{\kern0pt}model{\isacharunderscore}{\kern0pt}def\isanewline
\ \ \ \ \isacommand{by}\isamarkupfalse%
\ {\isacharparenleft}{\kern0pt}auto\ simp{\isacharcolon}{\kern0pt}\ wf{\isacharunderscore}{\kern0pt}domain{\isacharprime}{\kern0pt}{\isacharunderscore}{\kern0pt}correct\ wf{\isacharunderscore}{\kern0pt}fmla{\isacharprime}{\kern0pt}{\isacharunderscore}{\kern0pt}correct{\isacharparenright}{\kern0pt}%
\endisatagproof
{\isafoldproof}%
%
\isadelimproof
%
\endisadelimproof
%
\begin{isamarkuptext}%
Instantiating actions will yield well-founded effects.
    Corollary of \isa{{\isasymlbrakk}action{\isacharunderscore}{\kern0pt}params{\isacharunderscore}{\kern0pt}match\ {\isacharparenleft}{\kern0pt}Simple{\isacharunderscore}{\kern0pt}Action{\isacharunderscore}{\kern0pt}Schema\ {\isacharquery}{\kern0pt}n\ {\isacharquery}{\kern0pt}params\ {\isacharquery}{\kern0pt}pre\ {\isacharquery}{\kern0pt}eff{\isacharparenright}{\kern0pt}\ {\isacharquery}{\kern0pt}args{\isacharsemicolon}{\kern0pt}\ wf{\isacharunderscore}{\kern0pt}action{\isacharunderscore}{\kern0pt}schema\ {\isacharparenleft}{\kern0pt}Simple{\isacharunderscore}{\kern0pt}Action{\isacharunderscore}{\kern0pt}Schema\ {\isacharquery}{\kern0pt}n\ {\isacharquery}{\kern0pt}params\ {\isacharquery}{\kern0pt}pre\ {\isacharquery}{\kern0pt}eff{\isacharparenright}{\kern0pt}{\isasymrbrakk}\ {\isasymLongrightarrow}\ wf{\isacharunderscore}{\kern0pt}ground{\isacharunderscore}{\kern0pt}action\ {\isacharparenleft}{\kern0pt}instantiate{\isacharunderscore}{\kern0pt}action{\isacharunderscore}{\kern0pt}schema\ {\isacharparenleft}{\kern0pt}Simple{\isacharunderscore}{\kern0pt}Action{\isacharunderscore}{\kern0pt}Schema\ {\isacharquery}{\kern0pt}n\ {\isacharquery}{\kern0pt}params\ {\isacharquery}{\kern0pt}pre\ {\isacharquery}{\kern0pt}eff{\isacharparenright}{\kern0pt}\ {\isacharquery}{\kern0pt}args{\isacharparenright}{\kern0pt}} 
             and \isa{{\isasymlbrakk}action{\isacharunderscore}{\kern0pt}params{\isacharunderscore}{\kern0pt}match\ {\isacharparenleft}{\kern0pt}Durative{\isacharunderscore}{\kern0pt}Action{\isacharunderscore}{\kern0pt}Schema\ {\isacharquery}{\kern0pt}n\ {\isacharquery}{\kern0pt}params\ {\isacharquery}{\kern0pt}d\ {\isacharquery}{\kern0pt}cond\ {\isacharquery}{\kern0pt}eff{\isacharparenright}{\kern0pt}\ {\isacharquery}{\kern0pt}args{\isacharsemicolon}{\kern0pt}\ wf{\isacharunderscore}{\kern0pt}action{\isacharunderscore}{\kern0pt}schema\ {\isacharparenleft}{\kern0pt}Durative{\isacharunderscore}{\kern0pt}Action{\isacharunderscore}{\kern0pt}Schema\ {\isacharquery}{\kern0pt}n\ {\isacharquery}{\kern0pt}params\ {\isacharquery}{\kern0pt}d\ {\isacharquery}{\kern0pt}cond\ {\isacharquery}{\kern0pt}eff{\isacharparenright}{\kern0pt}{\isasymrbrakk}\ {\isasymLongrightarrow}\ wf{\isacharunderscore}{\kern0pt}ground{\isacharunderscore}{\kern0pt}action\ {\isacharparenleft}{\kern0pt}inst{\isacharunderscore}{\kern0pt}snap{\isacharunderscore}{\kern0pt}action\ {\isacharparenleft}{\kern0pt}Durative{\isacharunderscore}{\kern0pt}Action{\isacharunderscore}{\kern0pt}Schema\ {\isacharquery}{\kern0pt}n\ {\isacharquery}{\kern0pt}params\ {\isacharquery}{\kern0pt}d\ {\isacharquery}{\kern0pt}cond\ {\isacharquery}{\kern0pt}eff{\isacharparenright}{\kern0pt}\ {\isacharquery}{\kern0pt}args\ {\isacharquery}{\kern0pt}ta{\isacharparenright}{\kern0pt}}.%
\end{isamarkuptext}\isamarkuptrue%
\ \ \isacommand{lemma}\isamarkupfalse%
\ wf{\isacharunderscore}{\kern0pt}effect{\isacharunderscore}{\kern0pt}inst{\isacharunderscore}{\kern0pt}weak{\isacharcolon}{\kern0pt}\isanewline
\ \ \ \ \isakeyword{fixes}\ n\ params\ pre\ eff\isanewline
\ \ \ \ \isakeyword{defines}\ {\isachardoublequoteopen}a\isactrlsub s\isactrlsub c\isactrlsub h\isactrlsub e\isactrlsub m\isactrlsub a\ {\isasymequiv}\ Simple{\isacharunderscore}{\kern0pt}Action{\isacharunderscore}{\kern0pt}Schema\ n\ params\ pre\ eff{\isachardoublequoteclose}\ \ \isanewline
\ \ \ \ \isakeyword{assumes}\ {\isachardoublequoteopen}a\ {\isacharequal}{\kern0pt}\ instantiate{\isacharunderscore}{\kern0pt}action{\isacharunderscore}{\kern0pt}schema\ a\isactrlsub s\isactrlsub c\isactrlsub h\isactrlsub e\isactrlsub m\isactrlsub a\ args{\isachardoublequoteclose}\ \isanewline
\ \ \ \ \ \ \ \ \isakeyword{and}\ {\isachardoublequoteopen}action{\isacharunderscore}{\kern0pt}params{\isacharunderscore}{\kern0pt}match\ a\isactrlsub s\isactrlsub c\isactrlsub h\isactrlsub e\isactrlsub m\isactrlsub a\ args{\isachardoublequoteclose}\ \isanewline
\ \ \ \ \ \ \ \ \isakeyword{and}\ {\isachardoublequoteopen}wf{\isacharunderscore}{\kern0pt}action{\isacharunderscore}{\kern0pt}schema\ a\isactrlsub s\isactrlsub c\isactrlsub h\isactrlsub e\isactrlsub m\isactrlsub a{\isachardoublequoteclose}\isanewline
\ \ \ \ \isakeyword{shows}\ {\isachardoublequoteopen}wf{\isacharunderscore}{\kern0pt}effect{\isacharunderscore}{\kern0pt}inst\ {\isacharparenleft}{\kern0pt}effect\ a{\isacharparenright}{\kern0pt}{\isachardoublequoteclose}\isanewline
%
\isadelimproof
\ \ \ \ %
\endisadelimproof
%
\isatagproof
\isacommand{using}\isamarkupfalse%
\ assms\ wf{\isacharunderscore}{\kern0pt}inst{\isacharunderscore}{\kern0pt}action{\isacharunderscore}{\kern0pt}schema\isanewline
\ \ \ \ \isacommand{by}\isamarkupfalse%
\ {\isacharparenleft}{\kern0pt}cases\ a{\isacharparenright}{\kern0pt}\ {\isacharparenleft}{\kern0pt}auto\ simp{\isacharcolon}{\kern0pt}\ wf{\isacharunderscore}{\kern0pt}effect{\isacharunderscore}{\kern0pt}inst{\isacharunderscore}{\kern0pt}alt\ Let{\isacharunderscore}{\kern0pt}def{\isacharparenright}{\kern0pt}%
\endisatagproof
{\isafoldproof}%
%
\isadelimproof
\isanewline
%
\endisadelimproof
\isanewline
\ \ \isacommand{lemma}\isamarkupfalse%
\ wf{\isacharunderscore}{\kern0pt}effect{\isacharunderscore}{\kern0pt}durative{\isacharunderscore}{\kern0pt}inst{\isacharunderscore}{\kern0pt}weak{\isacharcolon}{\kern0pt}\isanewline
\ \ \ \ \isakeyword{fixes}\ n\ params\ d\ cond\ eff\isanewline
\ \ \ \ \isakeyword{defines}\ {\isachardoublequoteopen}a\isactrlsub s\isactrlsub c\isactrlsub h\isactrlsub e\isactrlsub m\isactrlsub a\ {\isasymequiv}\ Durative{\isacharunderscore}{\kern0pt}Action{\isacharunderscore}{\kern0pt}Schema\ n\ params\ d\ cond\ eff{\isachardoublequoteclose}\ \ \isanewline
\ \ \ \ \isakeyword{assumes}\ {\isachardoublequoteopen}a\ {\isacharequal}{\kern0pt}\ inst{\isacharunderscore}{\kern0pt}snap{\isacharunderscore}{\kern0pt}action\ a\isactrlsub s\isactrlsub c\isactrlsub h\isactrlsub e\isactrlsub m\isactrlsub a\ args\ ta{\isachardoublequoteclose}\ \isanewline
\ \ \ \ \ \ \ \ \isakeyword{and}\ {\isachardoublequoteopen}action{\isacharunderscore}{\kern0pt}params{\isacharunderscore}{\kern0pt}match\ a\isactrlsub s\isactrlsub c\isactrlsub h\isactrlsub e\isactrlsub m\isactrlsub a\ args{\isachardoublequoteclose}\ \isanewline
\ \ \ \ \ \ \ \ \isakeyword{and}\ {\isachardoublequoteopen}wf{\isacharunderscore}{\kern0pt}action{\isacharunderscore}{\kern0pt}schema\ a\isactrlsub s\isactrlsub c\isactrlsub h\isactrlsub e\isactrlsub m\isactrlsub a{\isachardoublequoteclose}\isanewline
\ \ \ \ \isakeyword{shows}\ {\isachardoublequoteopen}wf{\isacharunderscore}{\kern0pt}effect{\isacharunderscore}{\kern0pt}inst\ {\isacharparenleft}{\kern0pt}effect\ a{\isacharparenright}{\kern0pt}{\isachardoublequoteclose}\isanewline
%
\isadelimproof
\ \ \ \ %
\endisadelimproof
%
\isatagproof
\isacommand{using}\isamarkupfalse%
\ assms\ wf{\isacharunderscore}{\kern0pt}inst{\isacharunderscore}{\kern0pt}durative{\isacharunderscore}{\kern0pt}action{\isacharunderscore}{\kern0pt}schema\isanewline
\ \ \ \ \isacommand{by}\isamarkupfalse%
\ {\isacharparenleft}{\kern0pt}cases\ a{\isacharparenright}{\kern0pt}\ {\isacharparenleft}{\kern0pt}auto\ simp{\isacharcolon}{\kern0pt}\ wf{\isacharunderscore}{\kern0pt}effect{\isacharunderscore}{\kern0pt}inst{\isacharunderscore}{\kern0pt}alt\ Let{\isacharunderscore}{\kern0pt}def{\isacharparenright}{\kern0pt}%
\endisatagproof
{\isafoldproof}%
%
\isadelimproof
\isanewline
%
\endisadelimproof
\isanewline
\isacommand{end}\isamarkupfalse%
\ %
\isamarkupcmt{Context of \isa{ast{\isacharunderscore}{\kern0pt}problem}%
}%
\isadelimdocument
%
\endisadelimdocument
%
\isatagdocument
%
\isamarkupsubsubsection{Happening Execution%
}
\isamarkuptrue%
%
\endisatagdocument
{\isafolddocument}%
%
\isadelimdocument
%
\endisadelimdocument
\isacommand{context}\isamarkupfalse%
\ ast{\isacharunderscore}{\kern0pt}domain\ \isakeyword{begin}%
\begin{isamarkuptext}%
We first lift action schema lookup into the error monad.%
\end{isamarkuptext}\isamarkuptrue%
\ \ \isacommand{definition}\isamarkupfalse%
\ {\isachardoublequoteopen}resolve{\isacharunderscore}{\kern0pt}action{\isacharunderscore}{\kern0pt}schemaE\ n\ {\isasymequiv}\isanewline
\ \ \ \ lift{\isacharunderscore}{\kern0pt}opt\isanewline
\ \ \ \ \ \ {\isacharparenleft}{\kern0pt}resolve{\isacharunderscore}{\kern0pt}action{\isacharunderscore}{\kern0pt}schema\ n{\isacharparenright}{\kern0pt}\isanewline
\ \ \ \ \ \ {\isacharparenleft}{\kern0pt}ERR\ {\isacharparenleft}{\kern0pt}shows\ {\isacharprime}{\kern0pt}{\isacharprime}{\kern0pt}No\ such\ action\ schema\ {\isacharprime}{\kern0pt}{\isacharprime}{\kern0pt}\ o\ shows\ n{\isacharparenright}{\kern0pt}{\isacharparenright}{\kern0pt}{\isachardoublequoteclose}\isanewline
\isanewline
\ \ \isacommand{lemma}\isamarkupfalse%
\ resolve{\isacharunderscore}{\kern0pt}action{\isacharunderscore}{\kern0pt}schemaE{\isacharunderscore}{\kern0pt}return{\isacharunderscore}{\kern0pt}iff{\isacharbrackleft}{\kern0pt}return{\isacharunderscore}{\kern0pt}iff{\isacharbrackright}{\kern0pt}{\isacharcolon}{\kern0pt}\ \isanewline
\ \ \ \ {\isachardoublequoteopen}{\isacharparenleft}{\kern0pt}resolve{\isacharunderscore}{\kern0pt}action{\isacharunderscore}{\kern0pt}schemaE\ n\ {\isacharequal}{\kern0pt}\ Inr\ a{\isacharparenright}{\kern0pt}\ {\isasymlongleftrightarrow}\ {\isacharparenleft}{\kern0pt}resolve{\isacharunderscore}{\kern0pt}action{\isacharunderscore}{\kern0pt}schema\ n\ {\isacharequal}{\kern0pt}\ Some\ a{\isacharparenright}{\kern0pt}{\isachardoublequoteclose}\isanewline
%
\isadelimproof
\ \ \ \ %
\endisadelimproof
%
\isatagproof
\isacommand{using}\isamarkupfalse%
\ resolve{\isacharunderscore}{\kern0pt}action{\isacharunderscore}{\kern0pt}schemaE{\isacharunderscore}{\kern0pt}def\ \isacommand{by}\isamarkupfalse%
\ {\isacharparenleft}{\kern0pt}simp\ add{\isacharcolon}{\kern0pt}\ lift{\isacharunderscore}{\kern0pt}opt{\isacharunderscore}{\kern0pt}return{\isacharunderscore}{\kern0pt}iff{\isacharparenright}{\kern0pt}%
\endisatagproof
{\isafoldproof}%
%
\isadelimproof
\isanewline
%
\endisadelimproof
\isanewline
\isacommand{end}\isamarkupfalse%
\ %
\isamarkupcmt{Context of \isa{ast{\isacharunderscore}{\kern0pt}domain}%
}\isanewline
\isanewline
\isacommand{context}\isamarkupfalse%
\ ast{\isacharunderscore}{\kern0pt}problem\ \isakeyword{begin}%
\begin{isamarkuptext}%
We define a function to determine whether a formula holds in
    a world model%
\end{isamarkuptext}\isamarkuptrue%
\ \ \isacommand{definition}\isamarkupfalse%
\ {\isachardoublequoteopen}holds\ M\ F\ {\isasymequiv}\ {\isacharparenleft}{\kern0pt}valuation\ M{\isacharparenright}{\kern0pt}\ {\isasymTurnstile}\ F{\isachardoublequoteclose}%
\begin{isamarkuptext}%
Justification of this function%
\end{isamarkuptext}\isamarkuptrue%
\ \ \isacommand{lemma}\isamarkupfalse%
\ holds{\isacharunderscore}{\kern0pt}for{\isacharunderscore}{\kern0pt}wf{\isacharunderscore}{\kern0pt}fmlas{\isacharcolon}{\kern0pt}\isanewline
\ \ \ \ \isakeyword{assumes}\ {\isachardoublequoteopen}wm{\isacharunderscore}{\kern0pt}basic\ s{\isachardoublequoteclose}\isanewline
\ \ \ \ \isakeyword{shows}\ {\isachardoublequoteopen}holds\ s\ F\ {\isasymlongleftrightarrow}\ close{\isacharunderscore}{\kern0pt}world\ s\ {\isasymTTurnstile}\ F{\isachardoublequoteclose}\isanewline
%
\isadelimproof
\ \ \ \ %
\endisadelimproof
%
\isatagproof
\isacommand{unfolding}\isamarkupfalse%
\ holds{\isacharunderscore}{\kern0pt}def\ \isacommand{using}\isamarkupfalse%
\ assms\ valuation{\isacharunderscore}{\kern0pt}iff{\isacharunderscore}{\kern0pt}close{\isacharunderscore}{\kern0pt}world\ \isacommand{by}\isamarkupfalse%
\ blast%
\endisatagproof
{\isafoldproof}%
%
\isadelimproof
%
\endisadelimproof
%
\begin{isamarkuptext}%
The first refinement summarizes the enabledness check and the 
    application of a happening. Moreover, we implement the precondition 
    evaluation by our \isa{holds} function. This way, we can eliminate 
    redundant resolving and instantiation of the action.%
\end{isamarkuptext}\isamarkuptrue%
\ \ \isacommand{definition}\isamarkupfalse%
\ en{\isacharunderscore}{\kern0pt}exE\ {\isacharcolon}{\kern0pt}{\isacharcolon}{\kern0pt}\ {\isachardoublequoteopen}happening\ {\isasymRightarrow}\ world{\isacharunderscore}{\kern0pt}model\ {\isasymRightarrow}\ {\isacharunderscore}{\kern0pt}{\isacharplus}{\kern0pt}world{\isacharunderscore}{\kern0pt}model{\isachardoublequoteclose}\ \isakeyword{where}\isanewline
\ \ \ \ {\isachardoublequoteopen}en{\isacharunderscore}{\kern0pt}exE\ {\isasymequiv}\ {\isasymlambda}{\isacharparenleft}{\kern0pt}t\isactrlsub i{\isacharcomma}{\kern0pt}A\isactrlsub i{\isacharparenright}{\kern0pt}\ {\isasymRightarrow}\ {\isasymlambda}s{\isachardot}{\kern0pt}\ do\ {\isacharbraceleft}{\kern0pt}\isanewline
\ \ \ \ \ \ check{\isacharunderscore}{\kern0pt}allm\ {\isacharparenleft}{\kern0pt}{\isasymlambda}a{\isachardot}{\kern0pt}\ check\ {\isacharparenleft}{\kern0pt}holds\ s\ {\isacharparenleft}{\kern0pt}precondition\ a{\isacharparenright}{\kern0pt}{\isacharparenright}{\kern0pt}\ {\isacharparenleft}{\kern0pt}ERRS\ {\isacharprime}{\kern0pt}{\isacharprime}{\kern0pt}Precondition\ not\ satisfied{\isacharprime}{\kern0pt}{\isacharprime}{\kern0pt}{\isacharparenright}{\kern0pt}{\isacharparenright}{\kern0pt}\ A\isactrlsub i{\isacharsemicolon}{\kern0pt}\isanewline
\ \ \ \ \ \ check{\isacharunderscore}{\kern0pt}pairwise\ {\isacharparenleft}{\kern0pt}{\isasymlambda}a\ b{\isachardot}{\kern0pt}\ check\ {\isacharparenleft}{\kern0pt}a\ {\isasymnoteq}\ b\ {\isasymlongrightarrow}\ acts{\isacharunderscore}{\kern0pt}non{\isacharunderscore}{\kern0pt}intrf\ a\ b{\isacharparenright}{\kern0pt}\ {\isacharparenleft}{\kern0pt}ERRS\ {\isacharprime}{\kern0pt}{\isacharprime}{\kern0pt}Actions\ in\ happening\ interfering{\isacharprime}{\kern0pt}{\isacharprime}{\kern0pt}{\isacharparenright}{\kern0pt}{\isacharparenright}{\kern0pt}\ A\isactrlsub i{\isacharsemicolon}{\kern0pt}\isanewline
\ \ \ \ \ \ Error{\isacharunderscore}{\kern0pt}Monad{\isachardot}{\kern0pt}return\ {\isacharparenleft}{\kern0pt}apply{\isacharunderscore}{\kern0pt}happ{\isacharunderscore}{\kern0pt}exec\ {\isacharparenleft}{\kern0pt}t\isactrlsub i{\isacharcomma}{\kern0pt}A\isactrlsub i{\isacharparenright}{\kern0pt}\ s{\isacharparenright}{\kern0pt}\isanewline
\ \ \ \ {\isacharbraceright}{\kern0pt}{\isachardoublequoteclose}\isanewline
\isanewline
\ \ \isacommand{lemma}\isamarkupfalse%
\ symmetric{\isacharunderscore}{\kern0pt}pred{\isacharunderscore}{\kern0pt}check{\isacharunderscore}{\kern0pt}pairwise{\isacharcolon}{\kern0pt}\isanewline
\ \ \ \ \isakeyword{assumes}\ {\isachardoublequoteopen}{\isasymforall}x\isactrlsub {\isadigit{1}}\ x\isactrlsub {\isadigit{2}}{\isachardot}{\kern0pt}\ A\ x\isactrlsub {\isadigit{1}}\ x\isactrlsub {\isadigit{2}}\ {\isasymlongleftrightarrow}\ A\ x\isactrlsub {\isadigit{2}}\ x\isactrlsub {\isadigit{1}}{\isachardoublequoteclose}\ \isanewline
\ \ \ \ \ \ \ \ \isakeyword{and}\ {\isachardoublequoteopen}isOK\ {\isacharparenleft}{\kern0pt}check{\isacharunderscore}{\kern0pt}pairwise\ {\isacharparenleft}{\kern0pt}{\isasymlambda}x\isactrlsub {\isadigit{1}}\ x\isactrlsub {\isadigit{2}}{\isachardot}{\kern0pt}\ check\ {\isacharparenleft}{\kern0pt}x\isactrlsub {\isadigit{1}}\ {\isasymnoteq}\ x\isactrlsub {\isadigit{2}}\ {\isasymlongrightarrow}\ A\ x\isactrlsub {\isadigit{1}}\ x\isactrlsub {\isadigit{2}}{\isacharparenright}{\kern0pt}\ e{\isacharparenright}{\kern0pt}\ as{\isacharparenright}{\kern0pt}{\isachardoublequoteclose}\ \isanewline
\ \ \ \ \isakeyword{shows}\ {\isachardoublequoteopen}{\isasymforall}x\isactrlsub {\isadigit{1}}\ {\isasymin}\ set\ as{\isachardot}{\kern0pt}{\isasymforall}x\isactrlsub {\isadigit{2}}\ {\isasymin}\ set\ as{\isachardot}{\kern0pt}\ x\isactrlsub {\isadigit{1}}\ {\isasymnoteq}\ x\isactrlsub {\isadigit{2}}\ {\isasymlongrightarrow}\ A\ x\isactrlsub {\isadigit{1}}\ x\isactrlsub {\isadigit{2}}{\isachardoublequoteclose}\isanewline
%
\isadelimproof
\ \ \ \ %
\endisadelimproof
%
\isatagproof
\isacommand{using}\isamarkupfalse%
\ assms\ \isacommand{by}\isamarkupfalse%
\ {\isacharparenleft}{\kern0pt}auto\ simp{\isacharcolon}{\kern0pt}\ return{\isacharunderscore}{\kern0pt}iff{\isacharparenright}{\kern0pt}\ {\isacharparenleft}{\kern0pt}smt\ in{\isacharunderscore}{\kern0pt}set{\isacharunderscore}{\kern0pt}conv{\isacharunderscore}{\kern0pt}nth\ linorder{\isacharunderscore}{\kern0pt}neqE{\isacharunderscore}{\kern0pt}nat{\isacharparenright}{\kern0pt}%
\endisatagproof
{\isafoldproof}%
%
\isadelimproof
%
\endisadelimproof
%
\begin{isamarkuptext}%
Justification of implementation.%
\end{isamarkuptext}\isamarkuptrue%
\ \ \isacommand{lemma}\isamarkupfalse%
\ {\isacharparenleft}{\kern0pt}\isakeyword{in}\ wf{\isacharunderscore}{\kern0pt}ast{\isacharunderscore}{\kern0pt}problem{\isacharparenright}{\kern0pt}\ en{\isacharunderscore}{\kern0pt}exE{\isacharunderscore}{\kern0pt}return{\isacharunderscore}{\kern0pt}iff{\isacharcolon}{\kern0pt}\isanewline
\ \ \ \ \isakeyword{assumes}\ {\isachardoublequoteopen}wm{\isacharunderscore}{\kern0pt}basic\ s{\isachardoublequoteclose}\isanewline
\ \ \ \ \ \ \ \ \isakeyword{and}\ {\isachardoublequoteopen}{\isasymforall}a\ {\isasymin}\ set\ A\isactrlsub i{\isachardot}{\kern0pt}\ wf{\isacharunderscore}{\kern0pt}ground{\isacharunderscore}{\kern0pt}action\ a{\isachardoublequoteclose}\isanewline
\ \ \ \ \isakeyword{shows}\ {\isachardoublequoteopen}en{\isacharunderscore}{\kern0pt}exE\ {\isacharparenleft}{\kern0pt}t\isactrlsub i{\isacharcomma}{\kern0pt}A\isactrlsub i{\isacharparenright}{\kern0pt}\ s\ {\isacharequal}{\kern0pt}\ Inr\ s{\isacharprime}{\kern0pt}\ {\isasymlongleftrightarrow}\ happ{\isacharunderscore}{\kern0pt}enabled\ {\isacharparenleft}{\kern0pt}t\isactrlsub i{\isacharcomma}{\kern0pt}A\isactrlsub i{\isacharparenright}{\kern0pt}\ s\ {\isasymand}\ s{\isacharprime}{\kern0pt}\ {\isacharequal}{\kern0pt}\ apply{\isacharunderscore}{\kern0pt}happ\ {\isacharparenleft}{\kern0pt}t\isactrlsub i{\isacharcomma}{\kern0pt}A\isactrlsub i{\isacharparenright}{\kern0pt}\ s{\isachardoublequoteclose}\isanewline
%
\isadelimproof
\ \ \ \ %
\endisadelimproof
%
\isatagproof
\isacommand{unfolding}\isamarkupfalse%
\ en{\isacharunderscore}{\kern0pt}exE{\isacharunderscore}{\kern0pt}def\isanewline
\ \ \ \ \isacommand{using}\isamarkupfalse%
\ assms\ holds{\isacharunderscore}{\kern0pt}for{\isacharunderscore}{\kern0pt}wf{\isacharunderscore}{\kern0pt}fmlas{\isacharbrackleft}{\kern0pt}OF\ {\isacartoucheopen}wm{\isacharunderscore}{\kern0pt}basic\ s{\isacartoucheclose}{\isacharbrackright}{\kern0pt}\ \ \isanewline
\ \ \ \ \ \ \ \ \ \ symmetric{\isacharunderscore}{\kern0pt}pred{\isacharunderscore}{\kern0pt}check{\isacharunderscore}{\kern0pt}pairwise{\isacharbrackleft}{\kern0pt}OF\ acts{\isacharunderscore}{\kern0pt}non{\isacharunderscore}{\kern0pt}intrf{\isacharunderscore}{\kern0pt}symmetric{\isacharbrackright}{\kern0pt}\isanewline
\ \ \ \ \ \ \ \ \ \ apply{\isacharunderscore}{\kern0pt}happ{\isacharunderscore}{\kern0pt}exec{\isacharunderscore}{\kern0pt}refine\isanewline
\ \ \ \ \isacommand{by}\isamarkupfalse%
\ auto%
\endisatagproof
{\isafoldproof}%
%
\isadelimproof
%
\endisadelimproof
%
\begin{isamarkuptext}%
Next, we use the efficient implementation \isa{is{\isacharunderscore}{\kern0pt}obj{\isacharunderscore}{\kern0pt}of{\isacharunderscore}{\kern0pt}type{\isacharunderscore}{\kern0pt}impl}
    for the type check, and omit the well-formedness check, as effects obtained
    from instantiating well-formed action schemas are always well-formed
    (\isa{wf{\isacharunderscore}{\kern0pt}effect{\isacharunderscore}{\kern0pt}inst{\isacharunderscore}{\kern0pt}weak}).%
\end{isamarkuptext}\isamarkuptrue%
\ \ \isacommand{abbreviation}\isamarkupfalse%
\ {\isachardoublequoteopen}action{\isacharunderscore}{\kern0pt}params{\isacharunderscore}{\kern0pt}match{\isadigit{2}}\ stg\ mp\ a\ args\isanewline
\ \ \ \ {\isasymequiv}\ list{\isacharunderscore}{\kern0pt}all{\isadigit{2}}\ {\isacharparenleft}{\kern0pt}is{\isacharunderscore}{\kern0pt}obj{\isacharunderscore}{\kern0pt}of{\isacharunderscore}{\kern0pt}type{\isacharunderscore}{\kern0pt}impl\ stg\ mp{\isacharparenright}{\kern0pt}\isanewline
\ \ \ \ \ \ \ \ args\ {\isacharparenleft}{\kern0pt}map\ snd\ {\isacharparenleft}{\kern0pt}ast{\isacharunderscore}{\kern0pt}action{\isacharunderscore}{\kern0pt}schema{\isachardot}{\kern0pt}parameters\ a{\isacharparenright}{\kern0pt}{\isacharparenright}{\kern0pt}{\isachardoublequoteclose}\isanewline
\isanewline
\isacommand{end}\isamarkupfalse%
\ %
\isamarkupcmt{Context of \isa{ast{\isacharunderscore}{\kern0pt}problem}%
}%
\isadelimdocument
%
\endisadelimdocument
%
\isatagdocument
%
\isamarkupsubsubsection{Construction of Induced Happening Sequence%
}
\isamarkuptrue%
%
\endisatagdocument
{\isafolddocument}%
%
\isadelimdocument
%
\endisadelimdocument
\isacommand{context}\isamarkupfalse%
\ ast{\isacharunderscore}{\kern0pt}problem\ \isakeyword{begin}%
\begin{isamarkuptext}%
implemented own version of \isa{insort{\isacharunderscore}{\kern0pt}insert}, because \isa{insort{\isacharunderscore}{\kern0pt}insert} 
  requiers sort type; rat is not of sort type%
\end{isamarkuptext}\isamarkuptrue%
\ \ \isacommand{fun}\isamarkupfalse%
\ insort{\isacharunderscore}{\kern0pt}htp\ {\isacharcolon}{\kern0pt}{\isacharcolon}{\kern0pt}\ {\isachardoublequoteopen}time\ {\isasymRightarrow}\ time\ list\ {\isasymRightarrow}\ time\ list{\isachardoublequoteclose}\ \isakeyword{where}\isanewline
\ \ \ \ {\isachardoublequoteopen}insort{\isacharunderscore}{\kern0pt}htp\ x\ {\isacharbrackleft}{\kern0pt}{\isacharbrackright}{\kern0pt}\ {\isacharequal}{\kern0pt}\ {\isacharbrackleft}{\kern0pt}x{\isacharbrackright}{\kern0pt}{\isachardoublequoteclose}\isanewline
\ \ {\isacharbar}{\kern0pt}\ {\isachardoublequoteopen}insort{\isacharunderscore}{\kern0pt}htp\ x\ {\isacharparenleft}{\kern0pt}y{\isacharhash}{\kern0pt}ys{\isacharparenright}{\kern0pt}\ {\isacharequal}{\kern0pt}\ \isanewline
\ \ \ \ \ \ {\isacharparenleft}{\kern0pt}if\ x\ {\isacharless}{\kern0pt}\ y\ then\ x{\isacharhash}{\kern0pt}y{\isacharhash}{\kern0pt}ys\isanewline
\ \ \ \ \ \ else\ if\ x\ {\isacharequal}{\kern0pt}\ y\ then\ y{\isacharhash}{\kern0pt}ys\isanewline
\ \ \ \ \ \ else\ y\ {\isacharhash}{\kern0pt}\ {\isacharparenleft}{\kern0pt}insort{\isacharunderscore}{\kern0pt}htp\ x\ ys{\isacharparenright}{\kern0pt}{\isacharparenright}{\kern0pt}{\isachardoublequoteclose}%
\begin{isamarkuptext}%
Basic properties text of \isa{insort{\isacharunderscore}{\kern0pt}htp}: 
    preserves sorted-ness \& distinct-ness, inserts element%
\end{isamarkuptext}\isamarkuptrue%
\ \ \isacommand{lemma}\isamarkupfalse%
\ insort{\isacharunderscore}{\kern0pt}htps{\isacharunderscore}{\kern0pt}insort{\isacharunderscore}{\kern0pt}insert{\isacharcolon}{\kern0pt}\ {\isachardoublequoteopen}sorted\ xs\ {\isasymLongrightarrow}\ insort{\isacharunderscore}{\kern0pt}htp\ x\ xs\ {\isacharequal}{\kern0pt}\ insort{\isacharunderscore}{\kern0pt}insert\ x\ xs{\isachardoublequoteclose}\isanewline
%
\isadelimproof
\ \ \ \ %
\endisadelimproof
%
\isatagproof
\isacommand{by}\isamarkupfalse%
\ {\isacharparenleft}{\kern0pt}induction\ x\ xs\ rule{\isacharcolon}{\kern0pt}\ insort{\isacharunderscore}{\kern0pt}htp{\isachardot}{\kern0pt}induct{\isacharparenright}{\kern0pt}\ {\isacharparenleft}{\kern0pt}auto\ simp{\isacharcolon}{\kern0pt}\ insort{\isacharunderscore}{\kern0pt}insert{\isacharunderscore}{\kern0pt}key{\isacharunderscore}{\kern0pt}def{\isacharparenright}{\kern0pt}%
\endisatagproof
{\isafoldproof}%
%
\isadelimproof
\isanewline
%
\endisadelimproof
\isanewline
\ \ \isacommand{lemma}\isamarkupfalse%
\ sorted{\isacharunderscore}{\kern0pt}insort{\isacharunderscore}{\kern0pt}htps{\isacharcolon}{\kern0pt}\ {\isachardoublequoteopen}sorted\ xs\ {\isasymLongrightarrow}\ sorted\ {\isacharparenleft}{\kern0pt}insort{\isacharunderscore}{\kern0pt}htp\ x\ xs{\isacharparenright}{\kern0pt}{\isachardoublequoteclose}\isanewline
%
\isadelimproof
\ \ \ \ %
\endisadelimproof
%
\isatagproof
\isacommand{using}\isamarkupfalse%
\ insort{\isacharunderscore}{\kern0pt}htps{\isacharunderscore}{\kern0pt}insort{\isacharunderscore}{\kern0pt}insert\ sorted{\isacharunderscore}{\kern0pt}insort{\isacharunderscore}{\kern0pt}insert\ \isacommand{by}\isamarkupfalse%
\ auto%
\endisatagproof
{\isafoldproof}%
%
\isadelimproof
\isanewline
%
\endisadelimproof
\isanewline
\ \ \isacommand{lemma}\isamarkupfalse%
\ set{\isacharunderscore}{\kern0pt}insort{\isacharunderscore}{\kern0pt}htps{\isacharcolon}{\kern0pt}\ {\isachardoublequoteopen}set\ {\isacharparenleft}{\kern0pt}insort{\isacharunderscore}{\kern0pt}htp\ x\ xs{\isacharparenright}{\kern0pt}\ {\isacharequal}{\kern0pt}\ insert\ x\ {\isacharparenleft}{\kern0pt}set\ xs{\isacharparenright}{\kern0pt}{\isachardoublequoteclose}\isanewline
%
\isadelimproof
\ \ \ \ %
\endisadelimproof
%
\isatagproof
\isacommand{by}\isamarkupfalse%
\ {\isacharparenleft}{\kern0pt}induction\ x\ xs\ rule{\isacharcolon}{\kern0pt}\ insort{\isacharunderscore}{\kern0pt}htp{\isachardot}{\kern0pt}induct{\isacharparenright}{\kern0pt}\ auto%
\endisatagproof
{\isafoldproof}%
%
\isadelimproof
\isanewline
%
\endisadelimproof
\isanewline
\ \ \isacommand{lemma}\isamarkupfalse%
\ distinct{\isacharunderscore}{\kern0pt}insort{\isacharunderscore}{\kern0pt}htps{\isacharcolon}{\kern0pt}\ {\isachardoublequoteopen}sorted\ xs\ {\isasymLongrightarrow}\ distinct\ xs\ {\isasymLongrightarrow}\ distinct\ {\isacharparenleft}{\kern0pt}insort{\isacharunderscore}{\kern0pt}htp\ x\ xs{\isacharparenright}{\kern0pt}{\isachardoublequoteclose}\isanewline
%
\isadelimproof
\ \ \ \ %
\endisadelimproof
%
\isatagproof
\isacommand{using}\isamarkupfalse%
\ set{\isacharunderscore}{\kern0pt}insort{\isacharunderscore}{\kern0pt}htps\ \isacommand{by}\isamarkupfalse%
\ {\isacharparenleft}{\kern0pt}induction\ x\ xs\ rule{\isacharcolon}{\kern0pt}\ insort{\isacharunderscore}{\kern0pt}htp{\isachardot}{\kern0pt}induct{\isacharparenright}{\kern0pt}\ auto%
\endisatagproof
{\isafoldproof}%
%
\isadelimproof
\isanewline
%
\endisadelimproof
\isanewline
\ \ \isacommand{lemma}\isamarkupfalse%
\ insort{\isacharunderscore}{\kern0pt}htps{\isacharunderscore}{\kern0pt}non{\isacharunderscore}{\kern0pt}Nil{\isacharcolon}{\kern0pt}\ {\isachardoublequoteopen}insort{\isacharunderscore}{\kern0pt}htp\ x\ xs\ {\isasymnoteq}\ {\isacharbrackleft}{\kern0pt}{\isacharbrackright}{\kern0pt}{\isachardoublequoteclose}\isanewline
%
\isadelimproof
\ \ \ \ %
\endisadelimproof
%
\isatagproof
\isacommand{by}\isamarkupfalse%
\ {\isacharparenleft}{\kern0pt}induction\ x\ xs\ rule{\isacharcolon}{\kern0pt}\ insort{\isacharunderscore}{\kern0pt}htp{\isachardot}{\kern0pt}induct{\isacharparenright}{\kern0pt}\ auto%
\endisatagproof
{\isafoldproof}%
%
\isadelimproof
%
\endisadelimproof
%
\begin{isamarkuptext}%
Implementation of executable refinement for happening time points. Produces a 
  strictly sorted sequence of all \isa{is{\isacharunderscore}{\kern0pt}htp}%
\end{isamarkuptext}\isamarkuptrue%
\ \ \isacommand{fun}\isamarkupfalse%
\ htps{\isacharunderscore}{\kern0pt}exec\ {\isacharcolon}{\kern0pt}{\isacharcolon}{\kern0pt}\ {\isachardoublequoteopen}plan\ {\isasymRightarrow}\ time\ list{\isachardoublequoteclose}\ \isakeyword{where}\isanewline
\ \ \ \ {\isachardoublequoteopen}htps{\isacharunderscore}{\kern0pt}exec\ {\isacharbrackleft}{\kern0pt}{\isacharbrackright}{\kern0pt}\ {\isacharequal}{\kern0pt}\ {\isacharbrackleft}{\kern0pt}{\isacharbrackright}{\kern0pt}{\isachardoublequoteclose}\ \isanewline
\ \ {\isacharbar}{\kern0pt}\ {\isachardoublequoteopen}htps{\isacharunderscore}{\kern0pt}exec\ {\isacharparenleft}{\kern0pt}{\isacharparenleft}{\kern0pt}t\isactrlsub {\isasympi}{\isacharcomma}{\kern0pt}Simple{\isacharunderscore}{\kern0pt}Plan{\isacharunderscore}{\kern0pt}Action\ n\ args{\isacharparenright}{\kern0pt}{\isacharhash}{\kern0pt}{\isasympi}s{\isacharparenright}{\kern0pt}\ {\isacharequal}{\kern0pt}\ \isanewline
\ \ \ \ \ \ insort{\isacharunderscore}{\kern0pt}htp\ t\isactrlsub {\isasympi}\ {\isacharparenleft}{\kern0pt}htps{\isacharunderscore}{\kern0pt}exec\ {\isasympi}s{\isacharparenright}{\kern0pt}{\isachardoublequoteclose}\isanewline
\ \ {\isacharbar}{\kern0pt}\ {\isachardoublequoteopen}htps{\isacharunderscore}{\kern0pt}exec\ {\isacharparenleft}{\kern0pt}{\isacharparenleft}{\kern0pt}t\isactrlsub {\isasympi}{\isacharcomma}{\kern0pt}Durative{\isacharunderscore}{\kern0pt}Plan{\isacharunderscore}{\kern0pt}Action\ n\ args\ d{\isacharparenright}{\kern0pt}{\isacharhash}{\kern0pt}{\isasympi}s{\isacharparenright}{\kern0pt}\ {\isacharequal}{\kern0pt}\ \isanewline
\ \ \ \ \ \ insort{\isacharunderscore}{\kern0pt}htp\ t\isactrlsub {\isasympi}\ {\isacharparenleft}{\kern0pt}insort{\isacharunderscore}{\kern0pt}htp\ {\isacharparenleft}{\kern0pt}t\isactrlsub {\isasympi}{\isacharplus}{\kern0pt}d{\isacharparenright}{\kern0pt}\ {\isacharparenleft}{\kern0pt}htps{\isacharunderscore}{\kern0pt}exec\ {\isasympi}s{\isacharparenright}{\kern0pt}{\isacharparenright}{\kern0pt}{\isachardoublequoteclose}%
\begin{isamarkuptext}%
Basic properties of \isa{htps{\isacharunderscore}{\kern0pt}exec}%
\end{isamarkuptext}\isamarkuptrue%
\ \ \isacommand{lemma}\isamarkupfalse%
\ htps{\isacharunderscore}{\kern0pt}exec{\isacharunderscore}{\kern0pt}Nil{\isacharcolon}{\kern0pt}\ {\isachardoublequoteopen}htps{\isacharunderscore}{\kern0pt}exec\ {\isasympi}s\ {\isacharequal}{\kern0pt}\ {\isacharbrackleft}{\kern0pt}{\isacharbrackright}{\kern0pt}\ {\isasymlongleftrightarrow}\ {\isasympi}s\ {\isacharequal}{\kern0pt}\ {\isacharbrackleft}{\kern0pt}{\isacharbrackright}{\kern0pt}{\isachardoublequoteclose}\isanewline
%
\isadelimproof
\ \ \ \ %
\endisadelimproof
%
\isatagproof
\isacommand{by}\isamarkupfalse%
\ {\isacharparenleft}{\kern0pt}induction\ {\isasympi}s\ rule{\isacharcolon}{\kern0pt}\ htps{\isacharunderscore}{\kern0pt}exec{\isachardot}{\kern0pt}induct{\isacharparenright}{\kern0pt}\ {\isacharparenleft}{\kern0pt}auto\ simp{\isacharcolon}{\kern0pt}\ insort{\isacharunderscore}{\kern0pt}htps{\isacharunderscore}{\kern0pt}non{\isacharunderscore}{\kern0pt}Nil\ split{\isacharcolon}{\kern0pt}\ if{\isacharunderscore}{\kern0pt}splits{\isacharparenright}{\kern0pt}%
\endisatagproof
{\isafoldproof}%
%
\isadelimproof
\isanewline
%
\endisadelimproof
\isanewline
\ \ \isacommand{lemma}\isamarkupfalse%
\ htps{\isacharunderscore}{\kern0pt}exec{\isacharunderscore}{\kern0pt}sorted{\isacharcolon}{\kern0pt}\ {\isachardoublequoteopen}sorted\ {\isacharparenleft}{\kern0pt}htps{\isacharunderscore}{\kern0pt}exec\ {\isasympi}s{\isacharparenright}{\kern0pt}{\isachardoublequoteclose}\isanewline
%
\isadelimproof
\ \ \ \ %
\endisadelimproof
%
\isatagproof
\isacommand{using}\isamarkupfalse%
\ sorted{\isacharunderscore}{\kern0pt}insort{\isacharunderscore}{\kern0pt}htps\ \isacommand{by}\isamarkupfalse%
\ {\isacharparenleft}{\kern0pt}induction\ {\isasympi}s\ rule{\isacharcolon}{\kern0pt}\ htps{\isacharunderscore}{\kern0pt}exec{\isachardot}{\kern0pt}induct{\isacharparenright}{\kern0pt}\ auto%
\endisatagproof
{\isafoldproof}%
%
\isadelimproof
\isanewline
%
\endisadelimproof
\isanewline
\ \ \isacommand{lemma}\isamarkupfalse%
\ htps{\isacharunderscore}{\kern0pt}exec{\isacharunderscore}{\kern0pt}distinct{\isacharcolon}{\kern0pt}\ {\isachardoublequoteopen}distinct\ {\isacharparenleft}{\kern0pt}htps{\isacharunderscore}{\kern0pt}exec\ {\isasympi}s{\isacharparenright}{\kern0pt}{\isachardoublequoteclose}\isanewline
%
\isadelimproof
\ \ \ \ %
\endisadelimproof
%
\isatagproof
\isacommand{using}\isamarkupfalse%
\ htps{\isacharunderscore}{\kern0pt}exec{\isacharunderscore}{\kern0pt}sorted\ distinct{\isacharunderscore}{\kern0pt}insort{\isacharunderscore}{\kern0pt}htps\ sorted{\isacharunderscore}{\kern0pt}insort{\isacharunderscore}{\kern0pt}htps\isanewline
\ \ \ \ \isacommand{by}\isamarkupfalse%
\ {\isacharparenleft}{\kern0pt}induction\ {\isasympi}s\ rule{\isacharcolon}{\kern0pt}\ htps{\isacharunderscore}{\kern0pt}exec{\isachardot}{\kern0pt}induct{\isacharparenright}{\kern0pt}\ auto%
\endisatagproof
{\isafoldproof}%
%
\isadelimproof
%
\endisadelimproof
%
\begin{isamarkuptext}%
Correctness proof for \isa{htps{\isacharunderscore}{\kern0pt}exec}%
\end{isamarkuptext}\isamarkuptrue%
%
\begin{isamarkuptext}%
Every time point in \isa{htps{\isacharunderscore}{\kern0pt}exec} is a happening time point.%
\end{isamarkuptext}\isamarkuptrue%
\ \ \isacommand{lemma}\isamarkupfalse%
\ htps{\isacharunderscore}{\kern0pt}exec{\isacharunderscore}{\kern0pt}is{\isacharunderscore}{\kern0pt}htps{\isacharcolon}{\kern0pt}\ {\isachardoublequoteopen}t\isactrlsub i\ {\isasymin}\ set\ {\isacharparenleft}{\kern0pt}htps{\isacharunderscore}{\kern0pt}exec\ {\isasympi}s{\isacharparenright}{\kern0pt}\ {\isasymLongrightarrow}\ is{\isacharunderscore}{\kern0pt}htp\ {\isasympi}s\ t\isactrlsub i{\isachardoublequoteclose}\isanewline
%
\isadelimproof
\ \ %
\endisadelimproof
%
\isatagproof
\isacommand{proof}\isamarkupfalse%
\ {\isacharparenleft}{\kern0pt}induction\ {\isasympi}s{\isacharparenright}{\kern0pt}\isanewline
\ \ \ \ \isacommand{case}\isamarkupfalse%
\ Nil\isanewline
\ \ \ \ \isacommand{then}\isamarkupfalse%
\ \isacommand{show}\isamarkupfalse%
\ {\isacharquery}{\kern0pt}case\ \isacommand{by}\isamarkupfalse%
\ simp\isanewline
\ \ \isacommand{next}\isamarkupfalse%
\isanewline
\ \ \ \ \isacommand{case}\isamarkupfalse%
\ {\isacharparenleft}{\kern0pt}Cons\ {\isasympi}\ {\isasympi}s{\isacharparenright}{\kern0pt}\isanewline
\ \ \ \ \isacommand{obtain}\isamarkupfalse%
\ t\isactrlsub {\isasympi}\ {\isasympi}{\isacharprime}{\kern0pt}\ \isakeyword{where}\ {\isachardoublequoteopen}{\isasympi}\ {\isacharequal}{\kern0pt}\ {\isacharparenleft}{\kern0pt}t\isactrlsub {\isasympi}{\isacharcomma}{\kern0pt}{\isasympi}{\isacharprime}{\kern0pt}{\isacharparenright}{\kern0pt}{\isachardoublequoteclose}\isanewline
\ \ \ \ \ \ \isacommand{by}\isamarkupfalse%
\ {\isacharparenleft}{\kern0pt}meson\ surj{\isacharunderscore}{\kern0pt}pair{\isacharparenright}{\kern0pt}\isanewline
\ \ \ \ \isacommand{then}\isamarkupfalse%
\ \isacommand{show}\isamarkupfalse%
\ {\isacharquery}{\kern0pt}case\isanewline
\ \ \ \ \ \ \isacommand{using}\isamarkupfalse%
\ {\isacartoucheopen}{\isasympi}\ {\isacharequal}{\kern0pt}\ {\isacharparenleft}{\kern0pt}t\isactrlsub {\isasympi}{\isacharcomma}{\kern0pt}{\isasympi}{\isacharprime}{\kern0pt}{\isacharparenright}{\kern0pt}{\isacartoucheclose}\ set{\isacharunderscore}{\kern0pt}insort{\isacharunderscore}{\kern0pt}htps\ \ Cons{\isachardot}{\kern0pt}IH\ Cons{\isachardot}{\kern0pt}prems\isanewline
\ \ \ \ \ \ \isacommand{by}\isamarkupfalse%
\ {\isacharparenleft}{\kern0pt}cases\ {\isasympi}{\isacharprime}{\kern0pt}{\isacharparenright}{\kern0pt}\ {\isacharparenleft}{\kern0pt}auto\ simp{\isacharcolon}{\kern0pt}\ is{\isacharunderscore}{\kern0pt}htp{\isacharunderscore}{\kern0pt}def\ durative{\isacharunderscore}{\kern0pt}acts{\isacharunderscore}{\kern0pt}def\ is{\isacharunderscore}{\kern0pt}act{\isacharunderscore}{\kern0pt}simple{\isacharunderscore}{\kern0pt}def{\isacharparenright}{\kern0pt}\isanewline
\ \ \isacommand{qed}\isamarkupfalse%
%
\endisatagproof
{\isafoldproof}%
%
\isadelimproof
\isanewline
%
\endisadelimproof
\isanewline
\ \ \isacommand{lemma}\isamarkupfalse%
\ htps{\isacharunderscore}{\kern0pt}exec{\isacharunderscore}{\kern0pt}Cons{\isacharunderscore}{\kern0pt}subset{\isacharcolon}{\kern0pt}\ {\isachardoublequoteopen}set\ {\isacharparenleft}{\kern0pt}htps{\isacharunderscore}{\kern0pt}exec\ {\isasympi}s{\isacharparenright}{\kern0pt}\ {\isasymsubseteq}\ set\ {\isacharparenleft}{\kern0pt}htps{\isacharunderscore}{\kern0pt}exec\ {\isacharparenleft}{\kern0pt}{\isasympi}\ {\isacharhash}{\kern0pt}\ {\isasympi}s{\isacharparenright}{\kern0pt}{\isacharparenright}{\kern0pt}{\isachardoublequoteclose}\isanewline
%
\isadelimproof
\ \ \ \ %
\endisadelimproof
%
\isatagproof
\isacommand{by}\isamarkupfalse%
\ {\isacharparenleft}{\kern0pt}induction\ {\isachardoublequoteopen}{\isasympi}\ {\isacharhash}{\kern0pt}\ {\isasympi}s{\isachardoublequoteclose}\ rule{\isacharcolon}{\kern0pt}\ htps{\isacharunderscore}{\kern0pt}exec{\isachardot}{\kern0pt}induct{\isacharparenright}{\kern0pt}\ {\isacharparenleft}{\kern0pt}auto\ simp{\isacharcolon}{\kern0pt}\ set{\isacharunderscore}{\kern0pt}insort{\isacharunderscore}{\kern0pt}htps{\isacharparenright}{\kern0pt}%
\endisatagproof
{\isafoldproof}%
%
\isadelimproof
\isanewline
%
\endisadelimproof
\isanewline
\ \ \isacommand{lemma}\isamarkupfalse%
\ htps{\isacharunderscore}{\kern0pt}exec{\isacharunderscore}{\kern0pt}app{\isacharunderscore}{\kern0pt}subset{\isacharcolon}{\kern0pt}\ {\isachardoublequoteopen}set\ {\isacharparenleft}{\kern0pt}htps{\isacharunderscore}{\kern0pt}exec\ {\isasympi}s\isactrlsub {\isadigit{2}}{\isacharparenright}{\kern0pt}\ {\isasymsubseteq}\ set\ {\isacharparenleft}{\kern0pt}htps{\isacharunderscore}{\kern0pt}exec\ {\isacharparenleft}{\kern0pt}{\isasympi}s\isactrlsub {\isadigit{1}}\ {\isacharat}{\kern0pt}\ {\isasympi}s\isactrlsub {\isadigit{2}}{\isacharparenright}{\kern0pt}{\isacharparenright}{\kern0pt}{\isachardoublequoteclose}\isanewline
%
\isadelimproof
\ \ \ \ %
\endisadelimproof
%
\isatagproof
\isacommand{using}\isamarkupfalse%
\ htps{\isacharunderscore}{\kern0pt}exec{\isacharunderscore}{\kern0pt}Cons{\isacharunderscore}{\kern0pt}subset\isanewline
\ \ \ \ \isacommand{by}\isamarkupfalse%
\ {\isacharparenleft}{\kern0pt}induction\ {\isasympi}s\isactrlsub {\isadigit{1}}\ arbitrary{\isacharcolon}{\kern0pt}\ {\isasympi}s\isactrlsub {\isadigit{2}}\ rule{\isacharcolon}{\kern0pt}\ htps{\isacharunderscore}{\kern0pt}exec{\isachardot}{\kern0pt}induct{\isacharparenright}{\kern0pt}\ {\isacharparenleft}{\kern0pt}auto\ simp{\isacharcolon}{\kern0pt}\ set{\isacharunderscore}{\kern0pt}insort{\isacharunderscore}{\kern0pt}htps{\isacharparenright}{\kern0pt}%
\endisatagproof
{\isafoldproof}%
%
\isadelimproof
%
\endisadelimproof
%
\begin{isamarkuptext}%
Every happening time point is included in \isa{htps{\isacharunderscore}{\kern0pt}exec}.%
\end{isamarkuptext}\isamarkuptrue%
\ \ \isacommand{lemma}\isamarkupfalse%
\ is{\isacharunderscore}{\kern0pt}htps{\isacharunderscore}{\kern0pt}htps{\isacharunderscore}{\kern0pt}exec{\isacharcolon}{\kern0pt}\ {\isachardoublequoteopen}is{\isacharunderscore}{\kern0pt}htp\ {\isasympi}s\ t\isactrlsub i\ {\isasymLongrightarrow}\ t\isactrlsub i\ {\isasymin}\ set\ {\isacharparenleft}{\kern0pt}htps{\isacharunderscore}{\kern0pt}exec\ {\isasympi}s{\isacharparenright}{\kern0pt}{\isachardoublequoteclose}\isanewline
%
\isadelimproof
\ \ %
\endisadelimproof
%
\isatagproof
\isacommand{proof}\isamarkupfalse%
\ {\isacharparenleft}{\kern0pt}induction\ {\isasympi}s{\isacharparenright}{\kern0pt}\isanewline
\ \ \ \ \isacommand{case}\isamarkupfalse%
\ Nil\isanewline
\ \ \ \ \isacommand{then}\isamarkupfalse%
\ \isacommand{show}\isamarkupfalse%
\ {\isacharquery}{\kern0pt}case\ \isanewline
\ \ \ \ \ \ \isacommand{unfolding}\isamarkupfalse%
\ is{\isacharunderscore}{\kern0pt}htp{\isacharunderscore}{\kern0pt}def\ durative{\isacharunderscore}{\kern0pt}acts{\isacharunderscore}{\kern0pt}def\ is{\isacharunderscore}{\kern0pt}act{\isacharunderscore}{\kern0pt}simple{\isacharunderscore}{\kern0pt}def\ \isacommand{by}\isamarkupfalse%
\ simp\isanewline
\ \ \isacommand{next}\isamarkupfalse%
\isanewline
\ \ \ \ \isacommand{case}\isamarkupfalse%
\ {\isacharparenleft}{\kern0pt}Cons\ {\isasympi}\ {\isasympi}s{\isacharparenright}{\kern0pt}\isanewline
\ \ \ \ \isacommand{have}\isamarkupfalse%
\ {\isachardoublequoteopen}{\isacharparenleft}{\kern0pt}{\isasymexists}{\isasympi}{\isacharprime}{\kern0pt}{\isachardot}{\kern0pt}\ {\isacharparenleft}{\kern0pt}t\isactrlsub i{\isacharcomma}{\kern0pt}{\isasympi}{\isacharprime}{\kern0pt}{\isacharparenright}{\kern0pt}\ {\isasymin}\ set\ {\isacharparenleft}{\kern0pt}{\isasympi}{\isacharhash}{\kern0pt}{\isasympi}s{\isacharparenright}{\kern0pt}{\isacharparenright}{\kern0pt}\ {\isasymor}\ {\isacharparenleft}{\kern0pt}{\isasymexists}{\isacharparenleft}{\kern0pt}t\isactrlsub {\isasympi}{\isacharcomma}{\kern0pt}{\isasympi}{\isacharprime}{\kern0pt}{\isacharparenright}{\kern0pt}\ {\isasymin}\ durative{\isacharunderscore}{\kern0pt}acts\ {\isacharparenleft}{\kern0pt}{\isasympi}{\isacharhash}{\kern0pt}{\isasympi}s{\isacharparenright}{\kern0pt}{\isachardot}{\kern0pt}\ t\isactrlsub i\ {\isacharequal}{\kern0pt}\ t\isactrlsub {\isasympi}\ {\isacharplus}{\kern0pt}\ {\isacharparenleft}{\kern0pt}duration\ {\isasympi}{\isacharprime}{\kern0pt}{\isacharparenright}{\kern0pt}{\isacharparenright}{\kern0pt}{\isachardoublequoteclose}\isanewline
\ \ \ \ \ \ \isacommand{using}\isamarkupfalse%
\ Cons{\isachardot}{\kern0pt}prems\ \isacommand{unfolding}\isamarkupfalse%
\ is{\isacharunderscore}{\kern0pt}htp{\isacharunderscore}{\kern0pt}def\ \isacommand{by}\isamarkupfalse%
\ simp\isanewline
\ \ \ \ \isacommand{then}\isamarkupfalse%
\ \isacommand{show}\isamarkupfalse%
\ {\isacharquery}{\kern0pt}case\isanewline
\ \ \ \ \isacommand{proof}\isamarkupfalse%
\isanewline
\ \ \ \ \ \ \isacommand{assume}\isamarkupfalse%
\ {\isachardoublequoteopen}{\isasymexists}{\isasympi}{\isacharprime}{\kern0pt}{\isachardot}{\kern0pt}\ {\isacharparenleft}{\kern0pt}t\isactrlsub i{\isacharcomma}{\kern0pt}{\isasympi}{\isacharprime}{\kern0pt}{\isacharparenright}{\kern0pt}\ {\isasymin}\ set\ {\isacharparenleft}{\kern0pt}{\isasympi}{\isacharhash}{\kern0pt}{\isasympi}s{\isacharparenright}{\kern0pt}{\isachardoublequoteclose}\isanewline
\ \ \ \ \ \ \isacommand{then}\isamarkupfalse%
\ \isacommand{obtain}\isamarkupfalse%
\ {\isasympi}{\isacharprime}{\kern0pt}\ \isakeyword{where}\ {\isachardoublequoteopen}{\isacharparenleft}{\kern0pt}t\isactrlsub i{\isacharcomma}{\kern0pt}{\isasympi}{\isacharprime}{\kern0pt}{\isacharparenright}{\kern0pt}\ {\isasymin}\ set\ {\isacharparenleft}{\kern0pt}{\isasympi}{\isacharhash}{\kern0pt}{\isasympi}s{\isacharparenright}{\kern0pt}{\isachardoublequoteclose}\isanewline
\ \ \ \ \ \ \ \ \isacommand{by}\isamarkupfalse%
\ auto\isanewline
\ \ \ \ \ \ \isacommand{then}\isamarkupfalse%
\ \isacommand{obtain}\isamarkupfalse%
\ {\isasympi}s\isactrlsub {\isadigit{1}}\ {\isasympi}s\isactrlsub {\isadigit{2}}\ \isakeyword{where}\ {\isachardoublequoteopen}{\isacharparenleft}{\kern0pt}{\isasympi}{\isacharhash}{\kern0pt}{\isasympi}s{\isacharparenright}{\kern0pt}\ {\isacharequal}{\kern0pt}\ {\isasympi}s\isactrlsub {\isadigit{1}}\ {\isacharat}{\kern0pt}\ {\isacharparenleft}{\kern0pt}{\isacharparenleft}{\kern0pt}t\isactrlsub i{\isacharcomma}{\kern0pt}{\isasympi}{\isacharprime}{\kern0pt}{\isacharparenright}{\kern0pt}{\isacharhash}{\kern0pt}{\isasympi}s\isactrlsub {\isadigit{2}}{\isacharparenright}{\kern0pt}{\isachardoublequoteclose}\isanewline
\ \ \ \ \ \ \ \ \isacommand{by}\isamarkupfalse%
\ {\isacharparenleft}{\kern0pt}meson\ split{\isacharunderscore}{\kern0pt}list{\isacharparenright}{\kern0pt}\isanewline
\ \ \ \ \ \ \isacommand{then}\isamarkupfalse%
\ \isacommand{have}\isamarkupfalse%
\ {\isachardoublequoteopen}t\isactrlsub i\ {\isasymin}\ set\ {\isacharparenleft}{\kern0pt}htps{\isacharunderscore}{\kern0pt}exec\ {\isacharparenleft}{\kern0pt}{\isacharparenleft}{\kern0pt}t\isactrlsub i{\isacharcomma}{\kern0pt}{\isasympi}{\isacharprime}{\kern0pt}{\isacharparenright}{\kern0pt}{\isacharhash}{\kern0pt}{\isasympi}s\isactrlsub {\isadigit{2}}{\isacharparenright}{\kern0pt}{\isacharparenright}{\kern0pt}{\isachardoublequoteclose}\isanewline
\ \ \ \ \ \ \ \ \isacommand{using}\isamarkupfalse%
\ set{\isacharunderscore}{\kern0pt}insort{\isacharunderscore}{\kern0pt}htps\ \isacommand{by}\isamarkupfalse%
\ {\isacharparenleft}{\kern0pt}cases\ {\isasympi}{\isacharprime}{\kern0pt}{\isacharparenright}{\kern0pt}\ auto\isanewline
\ \ \ \ \ \ \isacommand{then}\isamarkupfalse%
\ \isacommand{show}\isamarkupfalse%
\ {\isacharquery}{\kern0pt}case\isanewline
\ \ \ \ \ \ \ \ \isacommand{using}\isamarkupfalse%
\ {\isacartoucheopen}{\isacharparenleft}{\kern0pt}{\isasympi}{\isacharhash}{\kern0pt}{\isasympi}s{\isacharparenright}{\kern0pt}\ {\isacharequal}{\kern0pt}\ {\isasympi}s\isactrlsub {\isadigit{1}}\ {\isacharat}{\kern0pt}\ {\isacharparenleft}{\kern0pt}{\isacharparenleft}{\kern0pt}t\isactrlsub i{\isacharcomma}{\kern0pt}{\isasympi}{\isacharprime}{\kern0pt}{\isacharparenright}{\kern0pt}{\isacharhash}{\kern0pt}{\isasympi}s\isactrlsub {\isadigit{2}}{\isacharparenright}{\kern0pt}{\isacartoucheclose}\ htps{\isacharunderscore}{\kern0pt}exec{\isacharunderscore}{\kern0pt}app{\isacharunderscore}{\kern0pt}subset\ \isacommand{by}\isamarkupfalse%
\ auto\isanewline
\ \ \ \ \isacommand{next}\isamarkupfalse%
\isanewline
\ \ \ \ \ \ \isacommand{assume}\isamarkupfalse%
\ {\isachardoublequoteopen}{\isasymexists}{\isacharparenleft}{\kern0pt}t\isactrlsub {\isasympi}{\isacharcomma}{\kern0pt}{\isasympi}{\isacharprime}{\kern0pt}{\isacharparenright}{\kern0pt}\ {\isasymin}\ durative{\isacharunderscore}{\kern0pt}acts\ {\isacharparenleft}{\kern0pt}{\isasympi}{\isacharhash}{\kern0pt}{\isasympi}s{\isacharparenright}{\kern0pt}{\isachardot}{\kern0pt}\ t\isactrlsub i\ {\isacharequal}{\kern0pt}\ t\isactrlsub {\isasympi}\ {\isacharplus}{\kern0pt}\ {\isacharparenleft}{\kern0pt}duration\ {\isasympi}{\isacharprime}{\kern0pt}{\isacharparenright}{\kern0pt}{\isachardoublequoteclose}\isanewline
\ \ \ \ \ \ \isacommand{then}\isamarkupfalse%
\ \isacommand{obtain}\isamarkupfalse%
\ t\isactrlsub {\isasympi}\ {\isasympi}{\isacharprime}{\kern0pt}\ \isakeyword{where}\ {\isachardoublequoteopen}{\isacharparenleft}{\kern0pt}t\isactrlsub {\isasympi}{\isacharcomma}{\kern0pt}{\isasympi}{\isacharprime}{\kern0pt}{\isacharparenright}{\kern0pt}\ {\isasymin}\ durative{\isacharunderscore}{\kern0pt}acts\ {\isacharparenleft}{\kern0pt}{\isasympi}{\isacharhash}{\kern0pt}{\isasympi}s{\isacharparenright}{\kern0pt}\ {\isasymand}\ t\isactrlsub i\ {\isacharequal}{\kern0pt}\ t\isactrlsub {\isasympi}\ {\isacharplus}{\kern0pt}\ {\isacharparenleft}{\kern0pt}duration\ {\isasympi}{\isacharprime}{\kern0pt}{\isacharparenright}{\kern0pt}{\isachardoublequoteclose}\isanewline
\ \ \ \ \ \ \ \ \isacommand{by}\isamarkupfalse%
\ auto\isanewline
\ \ \ \ \ \ \isacommand{then}\isamarkupfalse%
\ \isacommand{obtain}\isamarkupfalse%
\ {\isasympi}s\isactrlsub {\isadigit{1}}\ {\isasympi}s\isactrlsub {\isadigit{2}}\ \isakeyword{where}\ {\isachardoublequoteopen}{\isacharparenleft}{\kern0pt}{\isasympi}{\isacharhash}{\kern0pt}{\isasympi}s{\isacharparenright}{\kern0pt}\ {\isacharequal}{\kern0pt}\ {\isasympi}s\isactrlsub {\isadigit{1}}\ {\isacharat}{\kern0pt}\ {\isacharparenleft}{\kern0pt}{\isacharparenleft}{\kern0pt}t\isactrlsub {\isasympi}{\isacharcomma}{\kern0pt}{\isasympi}{\isacharprime}{\kern0pt}{\isacharparenright}{\kern0pt}{\isacharhash}{\kern0pt}{\isasympi}s\isactrlsub {\isadigit{2}}{\isacharparenright}{\kern0pt}{\isachardoublequoteclose}\isanewline
\ \ \ \ \ \ \ \ \isacommand{using}\isamarkupfalse%
\ dur{\isacharunderscore}{\kern0pt}acts{\isacharunderscore}{\kern0pt}subset\ \isacommand{by}\isamarkupfalse%
\ {\isacharparenleft}{\kern0pt}meson\ split{\isacharunderscore}{\kern0pt}list\ subsetD{\isacharparenright}{\kern0pt}\isanewline
\ \ \ \ \ \ \isacommand{then}\isamarkupfalse%
\ \isacommand{have}\isamarkupfalse%
\ {\isachardoublequoteopen}t\isactrlsub i\ {\isasymin}\ set\ {\isacharparenleft}{\kern0pt}htps{\isacharunderscore}{\kern0pt}exec\ {\isacharparenleft}{\kern0pt}{\isacharparenleft}{\kern0pt}t\isactrlsub {\isasympi}{\isacharcomma}{\kern0pt}{\isasympi}{\isacharprime}{\kern0pt}{\isacharparenright}{\kern0pt}{\isacharhash}{\kern0pt}{\isasympi}s\isactrlsub {\isadigit{2}}{\isacharparenright}{\kern0pt}{\isacharparenright}{\kern0pt}{\isachardoublequoteclose}\isanewline
\ \ \ \ \ \ \ \ \isacommand{using}\isamarkupfalse%
\ set{\isacharunderscore}{\kern0pt}insort{\isacharunderscore}{\kern0pt}htps\ {\isacartoucheopen}{\isacharparenleft}{\kern0pt}t\isactrlsub {\isasympi}{\isacharcomma}{\kern0pt}{\isasympi}{\isacharprime}{\kern0pt}{\isacharparenright}{\kern0pt}\ {\isasymin}\ durative{\isacharunderscore}{\kern0pt}acts\ {\isacharparenleft}{\kern0pt}{\isasympi}{\isacharhash}{\kern0pt}{\isasympi}s{\isacharparenright}{\kern0pt}\ {\isasymand}\ t\isactrlsub i\ {\isacharequal}{\kern0pt}\ t\isactrlsub {\isasympi}\ {\isacharplus}{\kern0pt}\ {\isacharparenleft}{\kern0pt}duration\ {\isasympi}{\isacharprime}{\kern0pt}{\isacharparenright}{\kern0pt}{\isacartoucheclose}\isanewline
\ \ \ \ \ \ \ \ \isacommand{unfolding}\isamarkupfalse%
\ durative{\isacharunderscore}{\kern0pt}acts{\isacharunderscore}{\kern0pt}def\ is{\isacharunderscore}{\kern0pt}act{\isacharunderscore}{\kern0pt}simple{\isacharunderscore}{\kern0pt}def\ \isacommand{by}\isamarkupfalse%
\ {\isacharparenleft}{\kern0pt}cases\ {\isasympi}{\isacharprime}{\kern0pt}{\isacharparenright}{\kern0pt}\ auto\isanewline
\ \ \ \ \ \ \isacommand{then}\isamarkupfalse%
\ \isacommand{show}\isamarkupfalse%
\ {\isacharquery}{\kern0pt}case\isanewline
\ \ \ \ \ \ \ \ \isacommand{using}\isamarkupfalse%
\ {\isacartoucheopen}{\isacharparenleft}{\kern0pt}{\isasympi}{\isacharhash}{\kern0pt}{\isasympi}s{\isacharparenright}{\kern0pt}\ {\isacharequal}{\kern0pt}\ {\isasympi}s\isactrlsub {\isadigit{1}}\ {\isacharat}{\kern0pt}\ {\isacharparenleft}{\kern0pt}{\isacharparenleft}{\kern0pt}t\isactrlsub {\isasympi}{\isacharcomma}{\kern0pt}{\isasympi}{\isacharprime}{\kern0pt}{\isacharparenright}{\kern0pt}{\isacharhash}{\kern0pt}{\isasympi}s\isactrlsub {\isadigit{2}}{\isacharparenright}{\kern0pt}{\isacartoucheclose}\ htps{\isacharunderscore}{\kern0pt}exec{\isacharunderscore}{\kern0pt}app{\isacharunderscore}{\kern0pt}subset\ \isacommand{by}\isamarkupfalse%
\ auto\isanewline
\ \ \ \ \isacommand{qed}\isamarkupfalse%
\isanewline
\ \ \isacommand{qed}\isamarkupfalse%
%
\endisatagproof
{\isafoldproof}%
%
\isadelimproof
\isanewline
%
\endisadelimproof
\isanewline
\ \ \isacommand{lemma}\isamarkupfalse%
\ is{\isacharunderscore}{\kern0pt}htp{\isacharunderscore}{\kern0pt}iff{\isacharunderscore}{\kern0pt}htps{\isacharunderscore}{\kern0pt}exec{\isacharcolon}{\kern0pt}\ {\isachardoublequoteopen}{\isasymforall}t{\isachardot}{\kern0pt}\ is{\isacharunderscore}{\kern0pt}htp\ {\isasympi}s\ t\ {\isasymlongleftrightarrow}\ t\ {\isasymin}\ set\ {\isacharparenleft}{\kern0pt}htps{\isacharunderscore}{\kern0pt}exec\ {\isasympi}s{\isacharparenright}{\kern0pt}{\isachardoublequoteclose}\isanewline
%
\isadelimproof
\ \ \ \ %
\endisadelimproof
%
\isatagproof
\isacommand{using}\isamarkupfalse%
\ htps{\isacharunderscore}{\kern0pt}exec{\isacharunderscore}{\kern0pt}is{\isacharunderscore}{\kern0pt}htps\ is{\isacharunderscore}{\kern0pt}htps{\isacharunderscore}{\kern0pt}htps{\isacharunderscore}{\kern0pt}exec\ \isacommand{by}\isamarkupfalse%
\ blast%
\endisatagproof
{\isafoldproof}%
%
\isadelimproof
\isanewline
%
\endisadelimproof
\isanewline
\ \ \isacommand{definition}\isamarkupfalse%
\ consec{\isacharunderscore}{\kern0pt}htps{\isacharunderscore}{\kern0pt}in{\isacharunderscore}{\kern0pt}interval\ {\isacharcolon}{\kern0pt}{\isacharcolon}{\kern0pt}\ {\isachardoublequoteopen}time\ list\ {\isasymRightarrow}\ time\ {\isasymRightarrow}\ time\ {\isasymRightarrow}\ {\isacharparenleft}{\kern0pt}time\ {\isasymtimes}\ time{\isacharparenright}{\kern0pt}\ list{\isachardoublequoteclose}\ \isakeyword{where}\ \isanewline
\ \ \ \ {\isachardoublequoteopen}consec{\isacharunderscore}{\kern0pt}htps{\isacharunderscore}{\kern0pt}in{\isacharunderscore}{\kern0pt}interval\ htps\ t\isactrlsub i\ t\isactrlsub j\ {\isacharequal}{\kern0pt}\ \isanewline
\ \ \ \ \ \ filter\ {\isacharparenleft}{\kern0pt}{\isasymlambda}{\isacharparenleft}{\kern0pt}t{\isacharcomma}{\kern0pt}t{\isacharprime}{\kern0pt}{\isacharparenright}{\kern0pt}{\isachardot}{\kern0pt}\ t\isactrlsub i\ {\isasymle}\ t\ {\isasymand}\ t{\isacharprime}{\kern0pt}\ {\isasymle}\ t\isactrlsub j{\isacharparenright}{\kern0pt}\ {\isacharparenleft}{\kern0pt}zip\ {\isacharparenleft}{\kern0pt}butlast\ htps{\isacharparenright}{\kern0pt}\ {\isacharparenleft}{\kern0pt}tl\ htps{\isacharparenright}{\kern0pt}{\isacharparenright}{\kern0pt}{\isachardoublequoteclose}%
\begin{isamarkuptext}%
Construct all ground actions for a given plan action%
\end{isamarkuptext}\isamarkuptrue%
\ \ \isacommand{fun}\isamarkupfalse%
\ simplify{\isacharunderscore}{\kern0pt}action\ {\isacharcolon}{\kern0pt}{\isacharcolon}{\kern0pt}\ {\isachardoublequoteopen}time\ list\ {\isasymRightarrow}\ {\isacharparenleft}{\kern0pt}time\ {\isasymtimes}\ plan{\isacharunderscore}{\kern0pt}action{\isacharparenright}{\kern0pt}\ {\isasymRightarrow}\ {\isacharparenleft}{\kern0pt}time\ {\isasymtimes}\ ground{\isacharunderscore}{\kern0pt}action{\isacharparenright}{\kern0pt}\ list{\isachardoublequoteclose}\ \isakeyword{where}\isanewline
\ \ \ \ {\isachardoublequoteopen}simplify{\isacharunderscore}{\kern0pt}action\ htps\ {\isacharparenleft}{\kern0pt}t{\isacharcomma}{\kern0pt}Simple{\isacharunderscore}{\kern0pt}Plan{\isacharunderscore}{\kern0pt}Action\ n\ args{\isacharparenright}{\kern0pt}\ {\isacharequal}{\kern0pt}\ {\isacharparenleft}{\kern0pt}\isanewline
\ \ \ \ \ \ let\ a\ {\isacharequal}{\kern0pt}\ the\ {\isacharparenleft}{\kern0pt}resolve{\isacharunderscore}{\kern0pt}action{\isacharunderscore}{\kern0pt}schema\ n{\isacharparenright}{\kern0pt}\ in\isanewline
\ \ \ \ \ \ \ \ {\isacharbrackleft}{\kern0pt}{\isacharparenleft}{\kern0pt}t{\isacharcomma}{\kern0pt}\ instantiate{\isacharunderscore}{\kern0pt}action{\isacharunderscore}{\kern0pt}schema\ a\ args{\isacharparenright}{\kern0pt}{\isacharbrackright}{\kern0pt}\isanewline
\ \ \ \ \ \ {\isacharparenright}{\kern0pt}{\isachardoublequoteclose}\isanewline
\ \ {\isacharbar}{\kern0pt}\ {\isachardoublequoteopen}simplify{\isacharunderscore}{\kern0pt}action\ htps\ {\isacharparenleft}{\kern0pt}t{\isacharcomma}{\kern0pt}Durative{\isacharunderscore}{\kern0pt}Plan{\isacharunderscore}{\kern0pt}Action\ n\ args\ d{\isacharparenright}{\kern0pt}\ {\isacharequal}{\kern0pt}\ {\isacharparenleft}{\kern0pt}\isanewline
\ \ \ \ \ \ let\ a\ {\isacharequal}{\kern0pt}\ the\ {\isacharparenleft}{\kern0pt}resolve{\isacharunderscore}{\kern0pt}action{\isacharunderscore}{\kern0pt}schema\ n{\isacharparenright}{\kern0pt}\ in\isanewline
\ \ \ \ \ \ \ \ {\isacharparenleft}{\kern0pt}t{\isacharcomma}{\kern0pt}inst{\isacharunderscore}{\kern0pt}snap{\isacharunderscore}{\kern0pt}action\ a\ args\ At{\isacharunderscore}{\kern0pt}Start{\isacharparenright}{\kern0pt}\ {\isacharhash}{\kern0pt}\ \isanewline
\ \ \ \ \ \ \ \ {\isacharparenleft}{\kern0pt}t{\isacharplus}{\kern0pt}d{\isacharcomma}{\kern0pt}inst{\isacharunderscore}{\kern0pt}snap{\isacharunderscore}{\kern0pt}action\ a\ args\ At{\isacharunderscore}{\kern0pt}End{\isacharparenright}{\kern0pt}\ {\isacharhash}{\kern0pt}\ \isanewline
\ \ \ \ \ \ \ \ {\isacharparenleft}{\kern0pt}map\ {\isacharparenleft}{\kern0pt}{\isasymlambda}{\isacharparenleft}{\kern0pt}t\isactrlsub i{\isacharcomma}{\kern0pt}t\isactrlsub j{\isacharparenright}{\kern0pt}{\isachardot}{\kern0pt}{\isacharparenleft}{\kern0pt}{\isacharparenleft}{\kern0pt}t\isactrlsub i{\isacharplus}{\kern0pt}t\isactrlsub j{\isacharparenright}{\kern0pt}\ {\isacharslash}{\kern0pt}\ {\isadigit{2}}{\isacharcomma}{\kern0pt}inst{\isacharunderscore}{\kern0pt}snap{\isacharunderscore}{\kern0pt}action\ a\ args\ Over{\isacharunderscore}{\kern0pt}All{\isacharparenright}{\kern0pt}{\isacharparenright}{\kern0pt}\ {\isacharparenleft}{\kern0pt}consec{\isacharunderscore}{\kern0pt}htps{\isacharunderscore}{\kern0pt}in{\isacharunderscore}{\kern0pt}interval\ htps\ t\ {\isacharparenleft}{\kern0pt}t{\isacharplus}{\kern0pt}d{\isacharparenright}{\kern0pt}{\isacharparenright}{\kern0pt}{\isacharparenright}{\kern0pt}\isanewline
\ \ \ \ \ \ {\isacharparenright}{\kern0pt}{\isachardoublequoteclose}%
\begin{isamarkuptext}%
Auxiliary lemmas to prove properties about 'zip (butlast xs) (tl xs)', 
  which are needed to ultimatly prove properties about \isa{consec{\isacharunderscore}{\kern0pt}htps}.%
\end{isamarkuptext}\isamarkuptrue%
\ \ \isacommand{lemma}\isamarkupfalse%
\ zip{\isacharunderscore}{\kern0pt}butlast{\isacharunderscore}{\kern0pt}tl{\isacharunderscore}{\kern0pt}hd{\isacharcolon}{\kern0pt}\ \isanewline
\ \ \ \ \isakeyword{assumes}\ {\isachardoublequoteopen}distinct\ {\isacharparenleft}{\kern0pt}x\isactrlsub {\isadigit{1}}\ {\isacharhash}{\kern0pt}\ xs{\isacharparenright}{\kern0pt}{\isachardoublequoteclose}\ \isanewline
\ \ \ \ \ \ \ \ \isakeyword{and}\ {\isachardoublequoteopen}{\isacharparenleft}{\kern0pt}x\isactrlsub {\isadigit{1}}{\isacharcomma}{\kern0pt}\ x\isactrlsub {\isadigit{2}}{\isacharparenright}{\kern0pt}\ {\isasymin}\ set\ {\isacharparenleft}{\kern0pt}zip\ {\isacharparenleft}{\kern0pt}butlast\ {\isacharparenleft}{\kern0pt}x\isactrlsub {\isadigit{1}}\ {\isacharhash}{\kern0pt}\ xs{\isacharparenright}{\kern0pt}{\isacharparenright}{\kern0pt}\ {\isacharparenleft}{\kern0pt}tl\ {\isacharparenleft}{\kern0pt}x\isactrlsub {\isadigit{1}}\ {\isacharhash}{\kern0pt}\ xs{\isacharparenright}{\kern0pt}{\isacharparenright}{\kern0pt}{\isacharparenright}{\kern0pt}{\isachardoublequoteclose}\ \isanewline
\ \ \ \ \isakeyword{shows}\ {\isachardoublequoteopen}xs\ {\isasymnoteq}\ {\isacharbrackleft}{\kern0pt}{\isacharbrackright}{\kern0pt}\ {\isasymand}\ x\isactrlsub {\isadigit{2}}\ {\isacharequal}{\kern0pt}\ hd\ xs{\isachardoublequoteclose}\isanewline
%
\isadelimproof
\ \ \ \ %
\endisadelimproof
%
\isatagproof
\isacommand{using}\isamarkupfalse%
\ assms\isanewline
\ \ \ \ \isacommand{apply}\isamarkupfalse%
\ {\isacharparenleft}{\kern0pt}induction\ xs{\isacharparenright}{\kern0pt}\isanewline
\ \ \ \ \isacommand{apply}\isamarkupfalse%
\ {\isacharparenleft}{\kern0pt}auto\ split{\isacharcolon}{\kern0pt}\ if{\isacharunderscore}{\kern0pt}splits{\isacharparenright}{\kern0pt}\isanewline
\ \ \ \ \isacommand{apply}\isamarkupfalse%
\ {\isacharparenleft}{\kern0pt}meson\ in{\isacharunderscore}{\kern0pt}set{\isacharunderscore}{\kern0pt}butlastD\ in{\isacharunderscore}{\kern0pt}set{\isacharunderscore}{\kern0pt}zipE\ set{\isacharunderscore}{\kern0pt}ConsD{\isacharparenright}{\kern0pt}\isanewline
\ \ \ \ \isacommand{done}\isamarkupfalse%
%
\endisatagproof
{\isafoldproof}%
%
\isadelimproof
%
\endisadelimproof
\ \isanewline
\isanewline
\ \ \isacommand{lemma}\isamarkupfalse%
\ zip{\isacharunderscore}{\kern0pt}butlast{\isacharunderscore}{\kern0pt}sublist{\isacharcolon}{\kern0pt}\isanewline
\ \ \ \ \isakeyword{assumes}\ {\isachardoublequoteopen}distinct\ {\isacharparenleft}{\kern0pt}xs\isactrlsub {\isadigit{1}}\ {\isacharat}{\kern0pt}\ x\isactrlsub {\isadigit{1}}\ {\isacharhash}{\kern0pt}\ xs\isactrlsub {\isadigit{2}}{\isacharparenright}{\kern0pt}{\isachardoublequoteclose}\isanewline
\ \ \ \ \ \ \ \ \isakeyword{and}\ {\isachardoublequoteopen}{\isacharparenleft}{\kern0pt}x\isactrlsub {\isadigit{1}}{\isacharcomma}{\kern0pt}\ x\isactrlsub {\isadigit{2}}{\isacharparenright}{\kern0pt}\ {\isasymin}\ set\ {\isacharparenleft}{\kern0pt}zip\ {\isacharparenleft}{\kern0pt}butlast\ {\isacharparenleft}{\kern0pt}xs\isactrlsub {\isadigit{1}}\ {\isacharat}{\kern0pt}\ x\isactrlsub {\isadigit{1}}\ {\isacharhash}{\kern0pt}\ xs\isactrlsub {\isadigit{2}}{\isacharparenright}{\kern0pt}{\isacharparenright}{\kern0pt}\ {\isacharparenleft}{\kern0pt}tl\ {\isacharparenleft}{\kern0pt}xs\isactrlsub {\isadigit{1}}\ {\isacharat}{\kern0pt}\ x\isactrlsub {\isadigit{1}}\ {\isacharhash}{\kern0pt}\ xs\isactrlsub {\isadigit{2}}{\isacharparenright}{\kern0pt}{\isacharparenright}{\kern0pt}{\isacharparenright}{\kern0pt}{\isachardoublequoteclose}\isanewline
\ \ \ \ \isakeyword{shows}\ {\isachardoublequoteopen}{\isacharparenleft}{\kern0pt}x\isactrlsub {\isadigit{1}}{\isacharcomma}{\kern0pt}\ x\isactrlsub {\isadigit{2}}{\isacharparenright}{\kern0pt}\ {\isasymin}\ set\ {\isacharparenleft}{\kern0pt}zip\ {\isacharparenleft}{\kern0pt}butlast\ {\isacharparenleft}{\kern0pt}x\isactrlsub {\isadigit{1}}\ {\isacharhash}{\kern0pt}\ xs\isactrlsub {\isadigit{2}}{\isacharparenright}{\kern0pt}{\isacharparenright}{\kern0pt}\ {\isacharparenleft}{\kern0pt}tl\ {\isacharparenleft}{\kern0pt}x\isactrlsub {\isadigit{1}}\ {\isacharhash}{\kern0pt}\ xs\isactrlsub {\isadigit{2}}{\isacharparenright}{\kern0pt}{\isacharparenright}{\kern0pt}{\isacharparenright}{\kern0pt}{\isachardoublequoteclose}\isanewline
%
\isadelimproof
\ \ \ \ %
\endisadelimproof
%
\isatagproof
\isacommand{using}\isamarkupfalse%
\ assms\isanewline
\ \ \ \ \isacommand{apply}\isamarkupfalse%
\ {\isacharparenleft}{\kern0pt}induction\ xs\isactrlsub {\isadigit{1}}{\isacharparenright}{\kern0pt}\isanewline
\ \ \ \ \isacommand{apply}\isamarkupfalse%
\ {\isacharparenleft}{\kern0pt}auto\ split{\isacharcolon}{\kern0pt}\ if{\isacharunderscore}{\kern0pt}splits{\isacharparenright}{\kern0pt}\isanewline
\ \ \ \ \isacommand{using}\isamarkupfalse%
\ set{\isacharunderscore}{\kern0pt}zip{\isacharunderscore}{\kern0pt}leftD\ \isacommand{apply}\isamarkupfalse%
\ fastforce\isanewline
\ \ \ \ \isacommand{apply}\isamarkupfalse%
\ {\isacharparenleft}{\kern0pt}smt\ append{\isacharunderscore}{\kern0pt}is{\isacharunderscore}{\kern0pt}Nil{\isacharunderscore}{\kern0pt}conv\ list{\isachardot}{\kern0pt}collapse\ list{\isachardot}{\kern0pt}discI\ old{\isachardot}{\kern0pt}prod{\isachardot}{\kern0pt}inject\ set{\isacharunderscore}{\kern0pt}ConsD\ set{\isacharunderscore}{\kern0pt}zip{\isacharunderscore}{\kern0pt}leftD\ zip{\isacharunderscore}{\kern0pt}Cons{\isacharunderscore}{\kern0pt}Cons{\isacharparenright}{\kern0pt}\isanewline
\ \ \ \ \isacommand{done}\isamarkupfalse%
%
\endisatagproof
{\isafoldproof}%
%
\isadelimproof
%
\endisadelimproof
\ \isanewline
\isanewline
\ \ \isacommand{lemma}\isamarkupfalse%
\ zip{\isacharunderscore}{\kern0pt}butlast{\isacharunderscore}{\kern0pt}tl{\isacharunderscore}{\kern0pt}app{\isadigit{1}}{\isacharcolon}{\kern0pt}\ \isanewline
\ \ \ \ \isakeyword{assumes}\ {\isachardoublequoteopen}distinct\ xs{\isachardoublequoteclose}\ \isanewline
\ \ \ \ \ \ \ \ \isakeyword{and}\ {\isachardoublequoteopen}{\isacharparenleft}{\kern0pt}x\isactrlsub {\isadigit{1}}{\isacharcomma}{\kern0pt}\ x\isactrlsub {\isadigit{2}}{\isacharparenright}{\kern0pt}\ {\isasymin}\ set\ {\isacharparenleft}{\kern0pt}zip\ {\isacharparenleft}{\kern0pt}butlast\ xs{\isacharparenright}{\kern0pt}\ {\isacharparenleft}{\kern0pt}tl\ xs{\isacharparenright}{\kern0pt}{\isacharparenright}{\kern0pt}{\isachardoublequoteclose}\ \isanewline
\ \ \ \ \isakeyword{shows}\ {\isachardoublequoteopen}{\isasymexists}xs\isactrlsub {\isadigit{1}}\ xs\isactrlsub {\isadigit{2}}{\isachardot}{\kern0pt}\ xs\ {\isacharequal}{\kern0pt}\ xs\isactrlsub {\isadigit{1}}\ {\isacharat}{\kern0pt}\ x\isactrlsub {\isadigit{1}}\ {\isacharhash}{\kern0pt}\ x\isactrlsub {\isadigit{2}}\ {\isacharhash}{\kern0pt}\ xs\isactrlsub {\isadigit{2}}{\isachardoublequoteclose}\isanewline
%
\isadelimproof
\ \ \ \ %
\endisadelimproof
%
\isatagproof
\isacommand{using}\isamarkupfalse%
\ assms\isanewline
\ \ \isacommand{proof}\isamarkupfalse%
\ {\isacharparenleft}{\kern0pt}induction\ xs{\isacharparenright}{\kern0pt}\isanewline
\ \ \ \ \isacommand{case}\isamarkupfalse%
\ Nil\isanewline
\ \ \ \ \isacommand{then}\isamarkupfalse%
\ \isacommand{show}\isamarkupfalse%
\ {\isacharquery}{\kern0pt}case\isanewline
\ \ \ \ \ \ \isacommand{by}\isamarkupfalse%
\ simp\isanewline
\ \ \isacommand{next}\isamarkupfalse%
\isanewline
\ \ \ \ \isacommand{case}\isamarkupfalse%
\ {\isacharparenleft}{\kern0pt}Cons\ x\ xs{\isacharparenright}{\kern0pt}\isanewline
\ \ \ \ \isacommand{have}\isamarkupfalse%
\ {\isachardoublequoteopen}x\isactrlsub {\isadigit{1}}\ {\isasymin}\ set\ {\isacharparenleft}{\kern0pt}x\ {\isacharhash}{\kern0pt}\ xs{\isacharparenright}{\kern0pt}{\isachardoublequoteclose}\isanewline
\ \ \ \ \ \ \isacommand{by}\isamarkupfalse%
\ {\isacharparenleft}{\kern0pt}meson\ Cons{\isachardot}{\kern0pt}prems\ in{\isacharunderscore}{\kern0pt}set{\isacharunderscore}{\kern0pt}butlastD\ set{\isacharunderscore}{\kern0pt}zip{\isacharunderscore}{\kern0pt}leftD{\isacharparenright}{\kern0pt}\isanewline
\ \ \ \ \isacommand{then}\isamarkupfalse%
\ \isacommand{obtain}\isamarkupfalse%
\ xs\isactrlsub {\isadigit{1}}\ xs\isactrlsub {\isadigit{2}}\ \isakeyword{where}\ {\isachardoublequoteopen}x\ {\isacharhash}{\kern0pt}\ xs\ {\isacharequal}{\kern0pt}\ xs\isactrlsub {\isadigit{1}}\ {\isacharat}{\kern0pt}\ x\isactrlsub {\isadigit{1}}\ {\isacharhash}{\kern0pt}\ xs\isactrlsub {\isadigit{2}}{\isachardoublequoteclose}\isanewline
\ \ \ \ \ \ \isacommand{by}\isamarkupfalse%
\ {\isacharparenleft}{\kern0pt}meson\ split{\isacharunderscore}{\kern0pt}list{\isacharparenright}{\kern0pt}\isanewline
\ \ \ \ \isacommand{then}\isamarkupfalse%
\ \isacommand{have}\isamarkupfalse%
\ {\isachardoublequoteopen}distinct\ {\isacharparenleft}{\kern0pt}xs\isactrlsub {\isadigit{1}}\ {\isacharat}{\kern0pt}\ x\isactrlsub {\isadigit{1}}\ {\isacharhash}{\kern0pt}\ xs\isactrlsub {\isadigit{2}}{\isacharparenright}{\kern0pt}{\isachardoublequoteclose}\isanewline
\ \ \ \ \ \ \isacommand{using}\isamarkupfalse%
\ Cons{\isachardot}{\kern0pt}prems\ \isacommand{by}\isamarkupfalse%
\ auto\isanewline
\ \ \ \ \isacommand{then}\isamarkupfalse%
\ \isacommand{have}\isamarkupfalse%
\ {\isachardoublequoteopen}distinct\ {\isacharparenleft}{\kern0pt}x\isactrlsub {\isadigit{1}}{\isacharhash}{\kern0pt}xs\isactrlsub {\isadigit{2}}{\isacharparenright}{\kern0pt}{\isachardoublequoteclose}\isanewline
\ \ \ \ \ \ \isacommand{using}\isamarkupfalse%
\ Cons{\isachardot}{\kern0pt}prems\ \isacommand{by}\isamarkupfalse%
\ auto\isanewline
\ \ \ \ \isacommand{have}\isamarkupfalse%
\ {\isachardoublequoteopen}{\isacharparenleft}{\kern0pt}x\isactrlsub {\isadigit{1}}{\isacharcomma}{\kern0pt}\ x\isactrlsub {\isadigit{2}}{\isacharparenright}{\kern0pt}\ {\isasymin}\ set\ {\isacharparenleft}{\kern0pt}zip\ {\isacharparenleft}{\kern0pt}butlast\ {\isacharparenleft}{\kern0pt}xs\isactrlsub {\isadigit{1}}\ {\isacharat}{\kern0pt}\ x\isactrlsub {\isadigit{1}}\ {\isacharhash}{\kern0pt}\ xs\isactrlsub {\isadigit{2}}{\isacharparenright}{\kern0pt}{\isacharparenright}{\kern0pt}\ {\isacharparenleft}{\kern0pt}tl\ {\isacharparenleft}{\kern0pt}xs\isactrlsub {\isadigit{1}}\ {\isacharat}{\kern0pt}\ x\isactrlsub {\isadigit{1}}\ {\isacharhash}{\kern0pt}\ xs\isactrlsub {\isadigit{2}}{\isacharparenright}{\kern0pt}{\isacharparenright}{\kern0pt}{\isacharparenright}{\kern0pt}{\isachardoublequoteclose}\isanewline
\ \ \ \ \ \ \isacommand{using}\isamarkupfalse%
\ {\isacartoucheopen}x\ {\isacharhash}{\kern0pt}\ xs\ {\isacharequal}{\kern0pt}\ xs\isactrlsub {\isadigit{1}}\ {\isacharat}{\kern0pt}\ x\isactrlsub {\isadigit{1}}\ {\isacharhash}{\kern0pt}\ xs\isactrlsub {\isadigit{2}}{\isacartoucheclose}\ Cons{\isachardot}{\kern0pt}prems\ \isacommand{by}\isamarkupfalse%
\ auto\isanewline
\ \ \ \ \isacommand{then}\isamarkupfalse%
\ \isacommand{have}\isamarkupfalse%
\ {\isachardoublequoteopen}{\isacharparenleft}{\kern0pt}x\isactrlsub {\isadigit{1}}{\isacharcomma}{\kern0pt}\ x\isactrlsub {\isadigit{2}}{\isacharparenright}{\kern0pt}\ {\isasymin}\ set\ {\isacharparenleft}{\kern0pt}zip\ {\isacharparenleft}{\kern0pt}butlast\ {\isacharparenleft}{\kern0pt}x\isactrlsub {\isadigit{1}}\ {\isacharhash}{\kern0pt}\ xs\isactrlsub {\isadigit{2}}{\isacharparenright}{\kern0pt}{\isacharparenright}{\kern0pt}\ {\isacharparenleft}{\kern0pt}tl\ {\isacharparenleft}{\kern0pt}x\isactrlsub {\isadigit{1}}\ {\isacharhash}{\kern0pt}\ xs\isactrlsub {\isadigit{2}}{\isacharparenright}{\kern0pt}{\isacharparenright}{\kern0pt}{\isacharparenright}{\kern0pt}{\isachardoublequoteclose}\isanewline
\ \ \ \ \ \ \isacommand{using}\isamarkupfalse%
\ zip{\isacharunderscore}{\kern0pt}butlast{\isacharunderscore}{\kern0pt}sublist{\isacharbrackleft}{\kern0pt}OF\ {\isacartoucheopen}distinct\ {\isacharparenleft}{\kern0pt}xs\isactrlsub {\isadigit{1}}\ {\isacharat}{\kern0pt}\ {\isacharparenleft}{\kern0pt}x\isactrlsub {\isadigit{1}}\ {\isacharhash}{\kern0pt}\ xs\isactrlsub {\isadigit{2}}{\isacharparenright}{\kern0pt}{\isacharparenright}{\kern0pt}{\isacartoucheclose}{\isacharbrackright}{\kern0pt}\ \isacommand{by}\isamarkupfalse%
\ auto\isanewline
\ \ \ \ \isacommand{obtain}\isamarkupfalse%
\ xs\isactrlsub {\isadigit{2}}{\isacharprime}{\kern0pt}\ \isakeyword{where}\ {\isachardoublequoteopen}xs\isactrlsub {\isadigit{2}}\ {\isacharequal}{\kern0pt}\ x\isactrlsub {\isadigit{2}}\ {\isacharhash}{\kern0pt}\ xs\isactrlsub {\isadigit{2}}{\isacharprime}{\kern0pt}{\isachardoublequoteclose}\isanewline
\ \ \ \ \ \ \isacommand{using}\isamarkupfalse%
\ zip{\isacharunderscore}{\kern0pt}butlast{\isacharunderscore}{\kern0pt}tl{\isacharunderscore}{\kern0pt}hd{\isacharbrackleft}{\kern0pt}OF\ {\isacartoucheopen}distinct\ {\isacharparenleft}{\kern0pt}x\isactrlsub {\isadigit{1}}\ {\isacharhash}{\kern0pt}\ xs\isactrlsub {\isadigit{2}}{\isacharparenright}{\kern0pt}{\isacartoucheclose}\ \isanewline
\ \ \ \ \ \ \ \ \ \ \ \ \ \ \ \ \ \ \ \ \ \ \ \ \ \ \ \ \ \ \ \ \ {\isacartoucheopen}{\isacharparenleft}{\kern0pt}x\isactrlsub {\isadigit{1}}{\isacharcomma}{\kern0pt}\ x\isactrlsub {\isadigit{2}}{\isacharparenright}{\kern0pt}\ {\isasymin}\ set\ {\isacharparenleft}{\kern0pt}zip\ {\isacharparenleft}{\kern0pt}butlast\ {\isacharparenleft}{\kern0pt}x\isactrlsub {\isadigit{1}}\ {\isacharhash}{\kern0pt}\ xs\isactrlsub {\isadigit{2}}{\isacharparenright}{\kern0pt}{\isacharparenright}{\kern0pt}\ {\isacharparenleft}{\kern0pt}tl\ {\isacharparenleft}{\kern0pt}x\isactrlsub {\isadigit{1}}\ {\isacharhash}{\kern0pt}\ xs\isactrlsub {\isadigit{2}}{\isacharparenright}{\kern0pt}{\isacharparenright}{\kern0pt}{\isacharparenright}{\kern0pt}{\isacartoucheclose}{\isacharbrackright}{\kern0pt}\isanewline
\ \ \ \ \ \ \ \ \ \ \ \ list{\isachardot}{\kern0pt}exhaust{\isacharunderscore}{\kern0pt}sel\ \isacommand{by}\isamarkupfalse%
\ blast\isanewline
\ \ \ \ \isacommand{then}\isamarkupfalse%
\ \isacommand{show}\isamarkupfalse%
\ {\isacharquery}{\kern0pt}case\isanewline
\ \ \ \ \ \ \isacommand{using}\isamarkupfalse%
\ {\isacartoucheopen}x\ {\isacharhash}{\kern0pt}\ xs\ {\isacharequal}{\kern0pt}\ xs\isactrlsub {\isadigit{1}}\ {\isacharat}{\kern0pt}\ x\isactrlsub {\isadigit{1}}\ {\isacharhash}{\kern0pt}\ xs\isactrlsub {\isadigit{2}}{\isacartoucheclose}\isanewline
\ \ \ \ \ \ \isacommand{by}\isamarkupfalse%
\ auto\isanewline
\ \ \isacommand{qed}\isamarkupfalse%
%
\endisatagproof
{\isafoldproof}%
%
\isadelimproof
\isanewline
%
\endisadelimproof
\isanewline
\ \ \isacommand{lemma}\isamarkupfalse%
\ zip{\isacharunderscore}{\kern0pt}butlast{\isacharunderscore}{\kern0pt}tl{\isacharunderscore}{\kern0pt}app{\isadigit{2}}{\isacharcolon}{\kern0pt}\ \isanewline
\ \ \ \ {\isachardoublequoteopen}{\isacharparenleft}{\kern0pt}x\isactrlsub {\isadigit{1}}{\isacharcomma}{\kern0pt}\ x\isactrlsub {\isadigit{2}}{\isacharparenright}{\kern0pt}\ {\isasymin}\ set\ {\isacharparenleft}{\kern0pt}zip\ {\isacharparenleft}{\kern0pt}butlast\ {\isacharparenleft}{\kern0pt}xs\isactrlsub {\isadigit{1}}\ {\isacharat}{\kern0pt}\ x\isactrlsub {\isadigit{1}}\ {\isacharhash}{\kern0pt}\ x\isactrlsub {\isadigit{2}}\ {\isacharhash}{\kern0pt}\ xs\isactrlsub {\isadigit{2}}{\isacharparenright}{\kern0pt}{\isacharparenright}{\kern0pt}\ {\isacharparenleft}{\kern0pt}tl\ {\isacharparenleft}{\kern0pt}xs\isactrlsub {\isadigit{1}}\ {\isacharat}{\kern0pt}\ x\isactrlsub {\isadigit{1}}\ {\isacharhash}{\kern0pt}\ x\isactrlsub {\isadigit{2}}\ {\isacharhash}{\kern0pt}\ xs\isactrlsub {\isadigit{2}}{\isacharparenright}{\kern0pt}{\isacharparenright}{\kern0pt}{\isacharparenright}{\kern0pt}{\isachardoublequoteclose}\ \isanewline
%
\isadelimproof
\ \ %
\endisadelimproof
%
\isatagproof
\isacommand{proof}\isamarkupfalse%
\ {\isacharminus}{\kern0pt}\isanewline
\ \ \ \ \isacommand{let}\isamarkupfalse%
\ {\isacharquery}{\kern0pt}xs{\isacharequal}{\kern0pt}{\isachardoublequoteopen}xs\isactrlsub {\isadigit{1}}\ {\isacharat}{\kern0pt}\ x\isactrlsub {\isadigit{1}}\ {\isacharhash}{\kern0pt}\ x\isactrlsub {\isadigit{2}}\ {\isacharhash}{\kern0pt}\ xs\isactrlsub {\isadigit{2}}{\isachardoublequoteclose}\isanewline
\ \ \ \ \isacommand{let}\isamarkupfalse%
\ {\isacharquery}{\kern0pt}butlast{\isacharequal}{\kern0pt}{\isachardoublequoteopen}butlast\ {\isacharquery}{\kern0pt}xs{\isachardoublequoteclose}\isanewline
\ \ \ \ \isacommand{let}\isamarkupfalse%
\ {\isacharquery}{\kern0pt}tl{\isacharequal}{\kern0pt}{\isachardoublequoteopen}tl\ {\isacharquery}{\kern0pt}xs{\isachardoublequoteclose}\isanewline
\ \ \ \ \isacommand{have}\isamarkupfalse%
\ {\isachardoublequoteopen}x\isactrlsub {\isadigit{1}}\ {\isacharequal}{\kern0pt}\ {\isacharquery}{\kern0pt}xs\ {\isacharbang}{\kern0pt}\ {\isacharparenleft}{\kern0pt}length\ xs\isactrlsub {\isadigit{1}}{\isacharparenright}{\kern0pt}{\isachardoublequoteclose}\ \isakeyword{and}\ {\isachardoublequoteopen}{\isacharparenleft}{\kern0pt}length\ xs\isactrlsub {\isadigit{1}}{\isacharparenright}{\kern0pt}\ {\isacharless}{\kern0pt}\ length\ {\isacharquery}{\kern0pt}butlast{\isachardoublequoteclose}\isanewline
\ \ \ \ \ \ \isacommand{using}\isamarkupfalse%
\ nth{\isacharunderscore}{\kern0pt}append{\isacharunderscore}{\kern0pt}length\ \isacommand{by}\isamarkupfalse%
\ auto\isanewline
\ \ \ \ \isacommand{then}\isamarkupfalse%
\ \isacommand{have}\isamarkupfalse%
\ {\isachardoublequoteopen}x\isactrlsub {\isadigit{1}}\ {\isacharequal}{\kern0pt}\ {\isacharquery}{\kern0pt}butlast\ {\isacharbang}{\kern0pt}\ {\isacharparenleft}{\kern0pt}length\ xs\isactrlsub {\isadigit{1}}{\isacharparenright}{\kern0pt}{\isachardoublequoteclose}\isanewline
\ \ \ \ \ \ \isacommand{using}\isamarkupfalse%
\ nth{\isacharunderscore}{\kern0pt}butlast\ \isacommand{by}\isamarkupfalse%
\ fastforce\isanewline
\ \ \ \ \isacommand{have}\isamarkupfalse%
\ {\isachardoublequoteopen}x\isactrlsub {\isadigit{2}}\ {\isacharequal}{\kern0pt}\ {\isacharquery}{\kern0pt}xs\ {\isacharbang}{\kern0pt}\ {\isacharparenleft}{\kern0pt}length\ xs\isactrlsub {\isadigit{1}}\ {\isacharplus}{\kern0pt}\ {\isadigit{1}}{\isacharparenright}{\kern0pt}{\isachardoublequoteclose}\ \isakeyword{and}\ {\isachardoublequoteopen}{\isacharquery}{\kern0pt}xs\ {\isasymnoteq}\ {\isacharbrackleft}{\kern0pt}{\isacharbrackright}{\kern0pt}{\isachardoublequoteclose}\isanewline
\ \ \ \ \ \ \isacommand{using}\isamarkupfalse%
\ nth{\isacharunderscore}{\kern0pt}append{\isacharunderscore}{\kern0pt}length{\isacharbrackleft}{\kern0pt}\isakeyword{where}\ xs{\isacharequal}{\kern0pt}{\isachardoublequoteopen}xs\isactrlsub {\isadigit{1}}\ {\isacharat}{\kern0pt}\ {\isacharbrackleft}{\kern0pt}x\isactrlsub {\isadigit{1}}{\isacharbrackright}{\kern0pt}{\isachardoublequoteclose}{\isacharbrackright}{\kern0pt}\ \isacommand{by}\isamarkupfalse%
\ auto\isanewline
\ \ \ \ \isacommand{then}\isamarkupfalse%
\ \isacommand{have}\isamarkupfalse%
\ {\isachardoublequoteopen}x\isactrlsub {\isadigit{2}}\ {\isacharequal}{\kern0pt}\ {\isacharquery}{\kern0pt}tl\ {\isacharbang}{\kern0pt}\ {\isacharparenleft}{\kern0pt}length\ xs\isactrlsub {\isadigit{1}}{\isacharparenright}{\kern0pt}{\isachardoublequoteclose}\isanewline
\ \ \ \ \ \ \isacommand{using}\isamarkupfalse%
\ nth{\isacharunderscore}{\kern0pt}tl{\isacharbrackleft}{\kern0pt}OF\ {\isacartoucheopen}{\isacharquery}{\kern0pt}xs\ {\isasymnoteq}\ {\isacharbrackleft}{\kern0pt}{\isacharbrackright}{\kern0pt}{\isacartoucheclose}{\isacharbrackright}{\kern0pt}\ \isacommand{by}\isamarkupfalse%
\ auto\isanewline
\ \ \ \ \isacommand{have}\isamarkupfalse%
\ {\isachardoublequoteopen}{\isacharparenleft}{\kern0pt}length\ xs\isactrlsub {\isadigit{1}}{\isacharparenright}{\kern0pt}\ {\isacharless}{\kern0pt}\ length\ {\isacharquery}{\kern0pt}butlast{\isachardoublequoteclose}\ \isakeyword{and}\ {\isachardoublequoteopen}{\isacharparenleft}{\kern0pt}length\ xs\isactrlsub {\isadigit{1}}{\isacharparenright}{\kern0pt}\ {\isacharless}{\kern0pt}\ length\ {\isacharquery}{\kern0pt}tl{\isachardoublequoteclose}\isanewline
\ \ \ \ \ \ \isacommand{by}\isamarkupfalse%
\ auto\isanewline
\ \ \ \ \isacommand{then}\isamarkupfalse%
\ \isacommand{show}\isamarkupfalse%
\ {\isacharquery}{\kern0pt}thesis\isanewline
\ \ \ \ \ \ \isacommand{using}\isamarkupfalse%
\ {\isacartoucheopen}x\isactrlsub {\isadigit{1}}\ {\isacharequal}{\kern0pt}\ {\isacharquery}{\kern0pt}butlast\ {\isacharbang}{\kern0pt}\ {\isacharparenleft}{\kern0pt}length\ xs\isactrlsub {\isadigit{1}}{\isacharparenright}{\kern0pt}{\isacartoucheclose}\ {\isacartoucheopen}x\isactrlsub {\isadigit{2}}\ {\isacharequal}{\kern0pt}\ {\isacharquery}{\kern0pt}tl\ {\isacharbang}{\kern0pt}\ {\isacharparenleft}{\kern0pt}length\ xs\isactrlsub {\isadigit{1}}{\isacharparenright}{\kern0pt}{\isacartoucheclose}\ in{\isacharunderscore}{\kern0pt}set{\isacharunderscore}{\kern0pt}zip\isanewline
\ \ \ \ \ \ \isacommand{by}\isamarkupfalse%
\ {\isacharparenleft}{\kern0pt}metis\ fst{\isacharunderscore}{\kern0pt}conv\ snd{\isacharunderscore}{\kern0pt}conv{\isacharparenright}{\kern0pt}\isanewline
\ \ \isacommand{qed}\isamarkupfalse%
%
\endisatagproof
{\isafoldproof}%
%
\isadelimproof
\isanewline
%
\endisadelimproof
\isanewline
\ \ \isacommand{lemma}\isamarkupfalse%
\ zip{\isacharunderscore}{\kern0pt}butlast{\isacharunderscore}{\kern0pt}tl{\isacharunderscore}{\kern0pt}app{\isacharcolon}{\kern0pt}\ \isanewline
\ \ \ \ \isakeyword{assumes}\ {\isachardoublequoteopen}distinct\ xs{\isachardoublequoteclose}\ \isanewline
\ \ \ \ \isakeyword{shows}\ {\isachardoublequoteopen}{\isacharparenleft}{\kern0pt}x\isactrlsub {\isadigit{1}}{\isacharcomma}{\kern0pt}\ x\isactrlsub {\isadigit{2}}{\isacharparenright}{\kern0pt}\ {\isasymin}\ set\ {\isacharparenleft}{\kern0pt}zip\ {\isacharparenleft}{\kern0pt}butlast\ xs{\isacharparenright}{\kern0pt}\ {\isacharparenleft}{\kern0pt}tl\ xs{\isacharparenright}{\kern0pt}{\isacharparenright}{\kern0pt}\ {\isasymlongleftrightarrow}\ {\isacharparenleft}{\kern0pt}{\isasymexists}xs\isactrlsub {\isadigit{1}}\ xs\isactrlsub {\isadigit{2}}{\isachardot}{\kern0pt}\ xs\ {\isacharequal}{\kern0pt}\ xs\isactrlsub {\isadigit{1}}\ {\isacharat}{\kern0pt}\ x\isactrlsub {\isadigit{1}}\ {\isacharhash}{\kern0pt}\ x\isactrlsub {\isadigit{2}}\ {\isacharhash}{\kern0pt}\ xs\isactrlsub {\isadigit{2}}{\isacharparenright}{\kern0pt}{\isachardoublequoteclose}\isanewline
%
\isadelimproof
\ \ \ \ %
\endisadelimproof
%
\isatagproof
\isacommand{using}\isamarkupfalse%
\ zip{\isacharunderscore}{\kern0pt}butlast{\isacharunderscore}{\kern0pt}tl{\isacharunderscore}{\kern0pt}app{\isadigit{1}}{\isacharbrackleft}{\kern0pt}OF\ {\isacartoucheopen}distinct\ xs{\isacartoucheclose}{\isacharbrackright}{\kern0pt}\ zip{\isacharunderscore}{\kern0pt}butlast{\isacharunderscore}{\kern0pt}tl{\isacharunderscore}{\kern0pt}app{\isadigit{2}}\ \isacommand{by}\isamarkupfalse%
\ fastforce%
\endisatagproof
{\isafoldproof}%
%
\isadelimproof
\isanewline
%
\endisadelimproof
\isanewline
\ \ \isacommand{lemma}\isamarkupfalse%
\ sorted{\isacharunderscore}{\kern0pt}distinct{\isacharunderscore}{\kern0pt}zip{\isacharunderscore}{\kern0pt}le{\isacharcolon}{\kern0pt}\ \isanewline
\ \ \ \ \isakeyword{assumes}\ {\isachardoublequoteopen}strict{\isacharunderscore}{\kern0pt}sorted\ xs{\isachardoublequoteclose}\ \isanewline
\ \ \ \ \ \ \ \ \isakeyword{and}\ {\isachardoublequoteopen}{\isacharparenleft}{\kern0pt}x\isactrlsub {\isadigit{1}}{\isacharcomma}{\kern0pt}\ x\isactrlsub {\isadigit{2}}{\isacharparenright}{\kern0pt}\ {\isasymin}\ set\ {\isacharparenleft}{\kern0pt}zip\ {\isacharparenleft}{\kern0pt}butlast\ xs{\isacharparenright}{\kern0pt}\ {\isacharparenleft}{\kern0pt}tl\ xs{\isacharparenright}{\kern0pt}{\isacharparenright}{\kern0pt}{\isachardoublequoteclose}\isanewline
\ \ \ \ \isakeyword{shows}\ {\isachardoublequoteopen}x\isactrlsub {\isadigit{1}}\ {\isacharless}{\kern0pt}\ x\isactrlsub {\isadigit{2}}{\isachardoublequoteclose}\isanewline
%
\isadelimproof
\ \ \ \ %
\endisadelimproof
%
\isatagproof
\isacommand{using}\isamarkupfalse%
\ assms\isanewline
\ \ \isacommand{proof}\isamarkupfalse%
\ {\isacharparenleft}{\kern0pt}induction\ xs{\isacharparenright}{\kern0pt}\isanewline
\ \ \ \ \isacommand{case}\isamarkupfalse%
\ Nil\isanewline
\ \ \ \ \isacommand{then}\isamarkupfalse%
\ \isacommand{show}\isamarkupfalse%
\ {\isacharquery}{\kern0pt}case\ \isanewline
\ \ \ \ \ \ \isacommand{by}\isamarkupfalse%
\ simp\isanewline
\ \ \isacommand{next}\isamarkupfalse%
\isanewline
\ \ \ \ \isacommand{case}\isamarkupfalse%
\ {\isacharparenleft}{\kern0pt}Cons\ x\ xs{\isacharparenright}{\kern0pt}\isanewline
\ \ \ \ \isacommand{obtain}\isamarkupfalse%
\ xs\isactrlsub {\isadigit{1}}\ xs\isactrlsub {\isadigit{2}}\ \isakeyword{where}\ {\isachardoublequoteopen}x\ {\isacharhash}{\kern0pt}\ xs\ {\isacharequal}{\kern0pt}\ xs\isactrlsub {\isadigit{1}}\ {\isacharat}{\kern0pt}\ {\isacharparenleft}{\kern0pt}x\isactrlsub {\isadigit{1}}\ {\isacharhash}{\kern0pt}\ x\isactrlsub {\isadigit{2}}\ {\isacharhash}{\kern0pt}\ xs\isactrlsub {\isadigit{2}}{\isacharparenright}{\kern0pt}{\isachardoublequoteclose}\isanewline
\ \ \ \ \ \ \isacommand{by}\isamarkupfalse%
\ {\isacharparenleft}{\kern0pt}meson\ Cons{\isachardot}{\kern0pt}prems\ zip{\isacharunderscore}{\kern0pt}butlast{\isacharunderscore}{\kern0pt}tl{\isacharunderscore}{\kern0pt}app\ strict{\isacharunderscore}{\kern0pt}sorted{\isacharunderscore}{\kern0pt}iff{\isacharparenright}{\kern0pt}\isanewline
\ \ \ \ \isacommand{then}\isamarkupfalse%
\ \isacommand{show}\isamarkupfalse%
\ {\isacharquery}{\kern0pt}case\isanewline
\ \ \ \ \ \ \isacommand{using}\isamarkupfalse%
\ Cons{\isachardot}{\kern0pt}prems\ \isacommand{by}\isamarkupfalse%
\ {\isacharparenleft}{\kern0pt}auto\ simp{\isacharcolon}{\kern0pt}\ sorted{\isacharunderscore}{\kern0pt}append\ strict{\isacharunderscore}{\kern0pt}sorted{\isacharunderscore}{\kern0pt}iff{\isacharparenright}{\kern0pt}\isanewline
\ \ \isacommand{qed}\isamarkupfalse%
%
\endisatagproof
{\isafoldproof}%
%
\isadelimproof
\isanewline
%
\endisadelimproof
\isanewline
\ \ \isacommand{lemma}\isamarkupfalse%
\ strict{\isacharunderscore}{\kern0pt}sorted{\isacharunderscore}{\kern0pt}zip{\isacharunderscore}{\kern0pt}butlast{\isacharunderscore}{\kern0pt}tl{\isacharunderscore}{\kern0pt}leq{\isacharunderscore}{\kern0pt}or{\isacharunderscore}{\kern0pt}geq{\isacharcolon}{\kern0pt}\isanewline
\ \ \ \ \isakeyword{assumes}\ {\isachardoublequoteopen}strict{\isacharunderscore}{\kern0pt}sorted\ xs{\isachardoublequoteclose}\isanewline
\ \ \ \ \ \ \ \ \isakeyword{and}\ {\isachardoublequoteopen}{\isacharparenleft}{\kern0pt}x\isactrlsub {\isadigit{1}}{\isacharcomma}{\kern0pt}x\isactrlsub {\isadigit{2}}{\isacharparenright}{\kern0pt}\ {\isasymin}\ set\ {\isacharparenleft}{\kern0pt}zip\ {\isacharparenleft}{\kern0pt}butlast\ xs{\isacharparenright}{\kern0pt}\ {\isacharparenleft}{\kern0pt}tl\ xs{\isacharparenright}{\kern0pt}{\isacharparenright}{\kern0pt}{\isachardoublequoteclose}\ \isanewline
\ \ \ \ \ \ \ \ \isakeyword{and}\ {\isachardoublequoteopen}x\ {\isasymin}\ set\ xs{\isachardoublequoteclose}\isanewline
\ \ \ \ \isakeyword{shows}\ {\isachardoublequoteopen}x\ {\isasymle}\ x\isactrlsub {\isadigit{1}}\ {\isasymor}\ x\isactrlsub {\isadigit{2}}\ {\isasymle}\ x{\isachardoublequoteclose}\isanewline
%
\isadelimproof
\ \ %
\endisadelimproof
%
\isatagproof
\isacommand{proof}\isamarkupfalse%
\ {\isacharminus}{\kern0pt}\isanewline
\ \ \ \ \isacommand{have}\isamarkupfalse%
\ {\isachardoublequoteopen}distinct\ xs{\isachardoublequoteclose}\isanewline
\ \ \ \ \ \ \isacommand{using}\isamarkupfalse%
\ {\isacartoucheopen}strict{\isacharunderscore}{\kern0pt}sorted\ xs{\isacartoucheclose}\ \isacommand{by}\isamarkupfalse%
\ {\isacharparenleft}{\kern0pt}auto\ simp{\isacharcolon}{\kern0pt}\ strict{\isacharunderscore}{\kern0pt}sorted{\isacharunderscore}{\kern0pt}iff{\isacharparenright}{\kern0pt}\isanewline
\ \ \ \ \isacommand{then}\isamarkupfalse%
\ \isacommand{obtain}\isamarkupfalse%
\ xs\isactrlsub {\isadigit{1}}\ xs\isactrlsub {\isadigit{2}}\ \isakeyword{where}\ {\isachardoublequoteopen}xs\ {\isacharequal}{\kern0pt}\ xs\isactrlsub {\isadigit{1}}\ {\isacharat}{\kern0pt}\ x\isactrlsub {\isadigit{1}}\ {\isacharhash}{\kern0pt}\ x\isactrlsub {\isadigit{2}}\ {\isacharhash}{\kern0pt}\ xs\isactrlsub {\isadigit{2}}{\isachardoublequoteclose}\isanewline
\ \ \ \ \ \ \isacommand{using}\isamarkupfalse%
\ assms\ zip{\isacharunderscore}{\kern0pt}butlast{\isacharunderscore}{\kern0pt}tl{\isacharunderscore}{\kern0pt}app{\isacharbrackleft}{\kern0pt}OF\ {\isacartoucheopen}distinct\ xs{\isacartoucheclose}{\isacharbrackright}{\kern0pt}\ \isacommand{by}\isamarkupfalse%
\ auto\isanewline
\ \ \ \ \isacommand{then}\isamarkupfalse%
\ \isacommand{have}\isamarkupfalse%
\ {\isachardoublequoteopen}x\ {\isasymin}\ set\ xs\isactrlsub {\isadigit{1}}\ {\isasymor}\ x\ {\isacharequal}{\kern0pt}\ x\isactrlsub {\isadigit{1}}\ {\isasymor}\ x\ {\isacharequal}{\kern0pt}\ x\isactrlsub {\isadigit{2}}\ {\isasymor}\ x\ {\isasymin}\ set\ xs\isactrlsub {\isadigit{2}}{\isachardoublequoteclose}\isanewline
\ \ \ \ \ \ \isacommand{using}\isamarkupfalse%
\ assms\ \isacommand{by}\isamarkupfalse%
\ auto\isanewline
\ \ \ \ \isacommand{then}\isamarkupfalse%
\ \isacommand{show}\isamarkupfalse%
\ {\isacharquery}{\kern0pt}thesis\isanewline
\ \ \ \ \isacommand{proof}\isamarkupfalse%
\ {\isacharparenleft}{\kern0pt}elim\ disjE{\isacharparenright}{\kern0pt}\isanewline
\ \ \ \ \ \ \isacommand{assume}\isamarkupfalse%
\ {\isachardoublequoteopen}x\ {\isasymin}\ set\ xs\isactrlsub {\isadigit{1}}{\isachardoublequoteclose}\isanewline
\ \ \ \ \ \ \isacommand{then}\isamarkupfalse%
\ \isacommand{show}\isamarkupfalse%
\ {\isacharquery}{\kern0pt}thesis\isanewline
\ \ \ \ \ \ \ \ \isacommand{using}\isamarkupfalse%
\ assms\ {\isacartoucheopen}xs\ {\isacharequal}{\kern0pt}\ xs\isactrlsub {\isadigit{1}}\ {\isacharat}{\kern0pt}\ x\isactrlsub {\isadigit{1}}\ {\isacharhash}{\kern0pt}\ x\isactrlsub {\isadigit{2}}\ {\isacharhash}{\kern0pt}\ xs\isactrlsub {\isadigit{2}}{\isacartoucheclose}\isanewline
\ \ \ \ \ \ \ \ \isacommand{by}\isamarkupfalse%
\ {\isacharparenleft}{\kern0pt}auto\ simp{\isacharcolon}{\kern0pt}\ sorted{\isacharunderscore}{\kern0pt}wrt{\isacharunderscore}{\kern0pt}append{\isacharparenright}{\kern0pt}\isanewline
\ \ \ \ \isacommand{next}\isamarkupfalse%
\isanewline
\ \ \ \ \ \ \isacommand{assume}\isamarkupfalse%
\ {\isachardoublequoteopen}x\ {\isacharequal}{\kern0pt}\ x\isactrlsub {\isadigit{1}}{\isachardoublequoteclose}\isanewline
\ \ \ \ \ \ \isacommand{then}\isamarkupfalse%
\ \isacommand{show}\isamarkupfalse%
\ {\isacharquery}{\kern0pt}thesis\isanewline
\ \ \ \ \ \ \ \ \isacommand{by}\isamarkupfalse%
\ auto\isanewline
\ \ \ \ \isacommand{next}\isamarkupfalse%
\isanewline
\ \ \ \ \ \ \isacommand{assume}\isamarkupfalse%
\ {\isachardoublequoteopen}\ x\ {\isacharequal}{\kern0pt}\ x\isactrlsub {\isadigit{2}}{\isachardoublequoteclose}\isanewline
\ \ \ \ \ \ \isacommand{then}\isamarkupfalse%
\ \isacommand{show}\isamarkupfalse%
\ {\isacharquery}{\kern0pt}thesis\isanewline
\ \ \ \ \ \ \ \ \isacommand{by}\isamarkupfalse%
\ auto\isanewline
\ \ \ \ \isacommand{next}\isamarkupfalse%
\isanewline
\ \ \ \ \ \ \isacommand{assume}\isamarkupfalse%
\ {\isachardoublequoteopen}x\ {\isasymin}\ set\ xs\isactrlsub {\isadigit{2}}{\isachardoublequoteclose}\isanewline
\ \ \ \ \ \ \isacommand{then}\isamarkupfalse%
\ \isacommand{show}\isamarkupfalse%
\ {\isacharquery}{\kern0pt}thesis\isanewline
\ \ \ \ \ \ \ \ \isacommand{using}\isamarkupfalse%
\ assms\ {\isacartoucheopen}xs\ {\isacharequal}{\kern0pt}\ xs\isactrlsub {\isadigit{1}}\ {\isacharat}{\kern0pt}\ x\isactrlsub {\isadigit{1}}\ {\isacharhash}{\kern0pt}\ x\isactrlsub {\isadigit{2}}\ {\isacharhash}{\kern0pt}\ xs\isactrlsub {\isadigit{2}}{\isacartoucheclose}\isanewline
\ \ \ \ \ \ \ \ \isacommand{by}\isamarkupfalse%
\ {\isacharparenleft}{\kern0pt}auto\ simp{\isacharcolon}{\kern0pt}\ sorted{\isacharunderscore}{\kern0pt}wrt{\isacharunderscore}{\kern0pt}append{\isacharparenright}{\kern0pt}\isanewline
\ \ \ \ \isacommand{qed}\isamarkupfalse%
\isanewline
\ \ \isacommand{qed}\isamarkupfalse%
%
\endisatagproof
{\isafoldproof}%
%
\isadelimproof
\isanewline
%
\endisadelimproof
\isanewline
\ \ \isacommand{lemma}\isamarkupfalse%
\ consec{\isacharunderscore}{\kern0pt}htps{\isacharunderscore}{\kern0pt}in{\isacharunderscore}{\kern0pt}zip{\isacharunderscore}{\kern0pt}butlast{\isacharunderscore}{\kern0pt}tl{\isacharcolon}{\kern0pt}\isanewline
\ \ \ \ \isakeyword{assumes}\ {\isachardoublequoteopen}htps\ {\isacharequal}{\kern0pt}\ htps{\isacharunderscore}{\kern0pt}exec\ {\isasympi}s{\isachardoublequoteclose}\isanewline
\ \ \ \ \ \ \ \ \isakeyword{and}\ {\isachardoublequoteopen}consec{\isacharunderscore}{\kern0pt}htps\ {\isasympi}s\ t\isactrlsub i\ t\isactrlsub j{\isachardoublequoteclose}\isanewline
\ \ \ \ \ \ \isakeyword{shows}\ {\isachardoublequoteopen}{\isacharparenleft}{\kern0pt}t\isactrlsub i{\isacharcomma}{\kern0pt}t\isactrlsub j{\isacharparenright}{\kern0pt}\ {\isasymin}\ set\ {\isacharparenleft}{\kern0pt}zip\ {\isacharparenleft}{\kern0pt}butlast\ htps{\isacharparenright}{\kern0pt}\ {\isacharparenleft}{\kern0pt}tl\ htps{\isacharparenright}{\kern0pt}{\isacharparenright}{\kern0pt}{\isachardoublequoteclose}\isanewline
%
\isadelimproof
\ \ %
\endisadelimproof
%
\isatagproof
\isacommand{proof}\isamarkupfalse%
\ {\isacharminus}{\kern0pt}\isanewline
\ \ \ \ \isacommand{have}\isamarkupfalse%
\ {\isachardoublequoteopen}strict{\isacharunderscore}{\kern0pt}sorted\ htps{\isachardoublequoteclose}\isanewline
\ \ \ \ \ \ \isacommand{using}\isamarkupfalse%
\ assms\ htps{\isacharunderscore}{\kern0pt}exec{\isacharunderscore}{\kern0pt}sorted\ htps{\isacharunderscore}{\kern0pt}exec{\isacharunderscore}{\kern0pt}distinct\ \isacommand{by}\isamarkupfalse%
\ {\isacharparenleft}{\kern0pt}simp\ add{\isacharcolon}{\kern0pt}\ strict{\isacharunderscore}{\kern0pt}sorted{\isacharunderscore}{\kern0pt}iff{\isacharparenright}{\kern0pt}\isanewline
\ \ \ \ \isacommand{have}\isamarkupfalse%
\ {\isachardoublequoteopen}{\isasymforall}t\ {\isasymin}\ set\ htps{\isachardot}{\kern0pt}\ t\ {\isasymle}\ t\isactrlsub i\ {\isasymor}\ t\isactrlsub j\ {\isasymle}\ t{\isachardoublequoteclose}\ \isakeyword{and}\ {\isachardoublequoteopen}t\isactrlsub i\ {\isasymin}\ set\ htps{\isachardoublequoteclose}\ \isakeyword{and}\ {\isachardoublequoteopen}t\isactrlsub j\ {\isasymin}\ set\ htps{\isachardoublequoteclose}\ \isakeyword{and}\ {\isachardoublequoteopen}t\isactrlsub i\ {\isacharless}{\kern0pt}\ t\isactrlsub j{\isachardoublequoteclose}\isanewline
\ \ \ \ \ \ \isacommand{using}\isamarkupfalse%
\ assms\ is{\isacharunderscore}{\kern0pt}htp{\isacharunderscore}{\kern0pt}iff{\isacharunderscore}{\kern0pt}htps{\isacharunderscore}{\kern0pt}exec\ \isacommand{unfolding}\isamarkupfalse%
\ consec{\isacharunderscore}{\kern0pt}htps{\isacharunderscore}{\kern0pt}def\ \isacommand{by}\isamarkupfalse%
\ auto\isanewline
\ \ \ \ \isacommand{then}\isamarkupfalse%
\ \isacommand{obtain}\isamarkupfalse%
\ htps\isactrlsub {\isadigit{1}}\ htps\isactrlsub {\isadigit{2}}{\isacharprime}{\kern0pt}\ \isakeyword{where}\ {\isachardoublequoteopen}htps\ {\isacharequal}{\kern0pt}\ htps\isactrlsub {\isadigit{1}}\ {\isacharat}{\kern0pt}\ t\isactrlsub i\ {\isacharhash}{\kern0pt}\ htps\isactrlsub {\isadigit{2}}{\isacharprime}{\kern0pt}{\isachardoublequoteclose}\isanewline
\ \ \ \ \ \ \isacommand{by}\isamarkupfalse%
\ {\isacharparenleft}{\kern0pt}meson\ split{\isacharunderscore}{\kern0pt}list{\isacharparenright}{\kern0pt}\isanewline
\ \ \ \ \isacommand{then}\isamarkupfalse%
\ \isacommand{have}\isamarkupfalse%
\ {\isachardoublequoteopen}t\isactrlsub j\ {\isasymin}\ set\ htps\isactrlsub {\isadigit{2}}{\isacharprime}{\kern0pt}{\isachardoublequoteclose}\isanewline
\ \ \ \ \ \ \isacommand{using}\isamarkupfalse%
\ {\isacartoucheopen}strict{\isacharunderscore}{\kern0pt}sorted\ htps{\isacartoucheclose}\ {\isacartoucheopen}t\isactrlsub j\ {\isasymin}\ set\ htps{\isacartoucheclose}\ {\isacartoucheopen}t\isactrlsub i\ {\isacharless}{\kern0pt}\ t\isactrlsub j{\isacartoucheclose}\isanewline
\ \ \ \ \ \ \isacommand{by}\isamarkupfalse%
\ {\isacharparenleft}{\kern0pt}auto\ simp{\isacharcolon}{\kern0pt}\ sorted{\isacharunderscore}{\kern0pt}wrt{\isacharunderscore}{\kern0pt}append{\isacharparenright}{\kern0pt}\isanewline
\ \ \ \ \isacommand{then}\isamarkupfalse%
\ \isacommand{obtain}\isamarkupfalse%
\ htps\isactrlsub {\isadigit{2}}\ htps\isactrlsub {\isadigit{3}}\ \isakeyword{where}\ {\isachardoublequoteopen}htps\isactrlsub {\isadigit{2}}{\isacharprime}{\kern0pt}\ {\isacharequal}{\kern0pt}\ htps\isactrlsub {\isadigit{2}}\ {\isacharat}{\kern0pt}\ t\isactrlsub j\ {\isacharhash}{\kern0pt}\ htps\isactrlsub {\isadigit{3}}{\isachardoublequoteclose}\isanewline
\ \ \ \ \ \ \isacommand{by}\isamarkupfalse%
\ {\isacharparenleft}{\kern0pt}meson\ split{\isacharunderscore}{\kern0pt}list{\isacharparenright}{\kern0pt}\isanewline
\ \ \ \ \isacommand{then}\isamarkupfalse%
\ \isacommand{have}\isamarkupfalse%
\ {\isachardoublequoteopen}{\isasymforall}t\ {\isasymin}\ set\ htps\isactrlsub {\isadigit{2}}{\isachardot}{\kern0pt}\ t\ {\isasymle}\ t\isactrlsub i\ {\isasymor}\ t\isactrlsub j\ {\isasymle}\ t{\isachardoublequoteclose}\isanewline
\ \ \ \ \ \ \isacommand{using}\isamarkupfalse%
\ {\isacartoucheopen}{\isasymforall}t\ {\isasymin}\ set\ htps{\isachardot}{\kern0pt}\ t\ {\isasymle}\ t\isactrlsub i\ {\isasymor}\ t\isactrlsub j\ {\isasymle}\ t{\isacartoucheclose}\ {\isacartoucheopen}htps\ {\isacharequal}{\kern0pt}\ htps\isactrlsub {\isadigit{1}}\ {\isacharat}{\kern0pt}\ t\isactrlsub i\ {\isacharhash}{\kern0pt}\ htps\isactrlsub {\isadigit{2}}{\isacharprime}{\kern0pt}{\isacartoucheclose}\ \isacommand{by}\isamarkupfalse%
\ auto\isanewline
\ \ \ \ \isacommand{then}\isamarkupfalse%
\ \isacommand{have}\isamarkupfalse%
\ {\isachardoublequoteopen}htps\isactrlsub {\isadigit{2}}\ {\isacharequal}{\kern0pt}\ {\isacharbrackleft}{\kern0pt}{\isacharbrackright}{\kern0pt}{\isachardoublequoteclose}\isanewline
\ \ \ \ \ \ \isacommand{using}\isamarkupfalse%
\ {\isacartoucheopen}strict{\isacharunderscore}{\kern0pt}sorted\ htps{\isacartoucheclose}\ {\isacartoucheopen}htps\ {\isacharequal}{\kern0pt}\ htps\isactrlsub {\isadigit{1}}\ {\isacharat}{\kern0pt}\ t\isactrlsub i\ {\isacharhash}{\kern0pt}\ htps\isactrlsub {\isadigit{2}}{\isacharprime}{\kern0pt}{\isacartoucheclose}\ {\isacartoucheopen}htps\isactrlsub {\isadigit{2}}{\isacharprime}{\kern0pt}\ {\isacharequal}{\kern0pt}\ htps\isactrlsub {\isadigit{2}}\ {\isacharat}{\kern0pt}\ t\isactrlsub j\ {\isacharhash}{\kern0pt}\ htps\isactrlsub {\isadigit{3}}{\isacartoucheclose}\isanewline
\ \ \ \ \ \ \isacommand{by}\isamarkupfalse%
\ {\isacharparenleft}{\kern0pt}induction\ htps\isactrlsub {\isadigit{2}}{\isacharparenright}{\kern0pt}\ {\isacharparenleft}{\kern0pt}auto\ simp{\isacharcolon}{\kern0pt}\ sorted{\isacharunderscore}{\kern0pt}wrt{\isacharunderscore}{\kern0pt}append{\isacharparenright}{\kern0pt}\isanewline
\ \ \ \ \isacommand{then}\isamarkupfalse%
\ \isacommand{show}\isamarkupfalse%
\ {\isacharquery}{\kern0pt}thesis\isanewline
\ \ \ \ \ \ \isacommand{using}\isamarkupfalse%
\ {\isacartoucheopen}htps\ {\isacharequal}{\kern0pt}\ htps\isactrlsub {\isadigit{1}}\ {\isacharat}{\kern0pt}\ t\isactrlsub i\ {\isacharhash}{\kern0pt}\ htps\isactrlsub {\isadigit{2}}{\isacharprime}{\kern0pt}{\isacartoucheclose}\ {\isacartoucheopen}htps\isactrlsub {\isadigit{2}}{\isacharprime}{\kern0pt}\ {\isacharequal}{\kern0pt}\ htps\isactrlsub {\isadigit{2}}\ {\isacharat}{\kern0pt}\ t\isactrlsub j\ {\isacharhash}{\kern0pt}\ htps\isactrlsub {\isadigit{3}}{\isacartoucheclose}\isanewline
\ \ \ \ \ \ \isacommand{by}\isamarkupfalse%
\ {\isacharparenleft}{\kern0pt}auto\ simp{\isacharcolon}{\kern0pt}\ zip{\isacharunderscore}{\kern0pt}butlast{\isacharunderscore}{\kern0pt}tl{\isacharunderscore}{\kern0pt}app{\isadigit{2}}{\isacharparenright}{\kern0pt}\isanewline
\ \ \isacommand{qed}\isamarkupfalse%
%
\endisatagproof
{\isafoldproof}%
%
\isadelimproof
%
\endisadelimproof
%
\begin{isamarkuptext}%
Executable refinement for \isa{consec{\isacharunderscore}{\kern0pt}htps}.%
\end{isamarkuptext}\isamarkuptrue%
\ \ \isacommand{lemma}\isamarkupfalse%
\ happ{\isacharunderscore}{\kern0pt}time{\isacharunderscore}{\kern0pt}pts{\isacharunderscore}{\kern0pt}exec{\isacharunderscore}{\kern0pt}iff{\isacharunderscore}{\kern0pt}consec{\isacharunderscore}{\kern0pt}htps{\isacharcolon}{\kern0pt}\isanewline
\ \ \ \ {\isachardoublequoteopen}{\isacharparenleft}{\kern0pt}t\isactrlsub i{\isacharcomma}{\kern0pt}t\isactrlsub j{\isacharparenright}{\kern0pt}\ {\isasymin}\ set\ {\isacharparenleft}{\kern0pt}zip\ {\isacharparenleft}{\kern0pt}butlast\ {\isacharparenleft}{\kern0pt}htps{\isacharunderscore}{\kern0pt}exec\ {\isasympi}s{\isacharparenright}{\kern0pt}{\isacharparenright}{\kern0pt}\ {\isacharparenleft}{\kern0pt}tl\ {\isacharparenleft}{\kern0pt}htps{\isacharunderscore}{\kern0pt}exec\ {\isasympi}s{\isacharparenright}{\kern0pt}{\isacharparenright}{\kern0pt}{\isacharparenright}{\kern0pt}\ {\isasymlongleftrightarrow}\ consec{\isacharunderscore}{\kern0pt}htps\ {\isasympi}s\ t\isactrlsub i\ t\isactrlsub j{\isachardoublequoteclose}\isanewline
%
\isadelimproof
\ \ %
\endisadelimproof
%
\isatagproof
\isacommand{proof}\isamarkupfalse%
\ {\isacharminus}{\kern0pt}\isanewline
\ \ \ \ \isacommand{have}\isamarkupfalse%
\ {\isachardoublequoteopen}strict{\isacharunderscore}{\kern0pt}sorted\ {\isacharparenleft}{\kern0pt}htps{\isacharunderscore}{\kern0pt}exec\ {\isasympi}s{\isacharparenright}{\kern0pt}{\isachardoublequoteclose}\isanewline
\ \ \ \ \ \ \isacommand{using}\isamarkupfalse%
\ htps{\isacharunderscore}{\kern0pt}exec{\isacharunderscore}{\kern0pt}sorted\ htps{\isacharunderscore}{\kern0pt}exec{\isacharunderscore}{\kern0pt}distinct\ \isacommand{by}\isamarkupfalse%
\ {\isacharparenleft}{\kern0pt}simp\ add{\isacharcolon}{\kern0pt}\ strict{\isacharunderscore}{\kern0pt}sorted{\isacharunderscore}{\kern0pt}iff{\isacharparenright}{\kern0pt}\isanewline
\ \ \ \ \isacommand{then}\isamarkupfalse%
\ \isacommand{show}\isamarkupfalse%
\ {\isacharquery}{\kern0pt}thesis\isanewline
\ \ \ \ \ \ \isacommand{using}\isamarkupfalse%
\ is{\isacharunderscore}{\kern0pt}htp{\isacharunderscore}{\kern0pt}iff{\isacharunderscore}{\kern0pt}htps{\isacharunderscore}{\kern0pt}exec\ sorted{\isacharunderscore}{\kern0pt}distinct{\isacharunderscore}{\kern0pt}zip{\isacharunderscore}{\kern0pt}le\isanewline
\ \ \ \ \ \ \ \ \ \ \ \ strict{\isacharunderscore}{\kern0pt}sorted{\isacharunderscore}{\kern0pt}zip{\isacharunderscore}{\kern0pt}butlast{\isacharunderscore}{\kern0pt}tl{\isacharunderscore}{\kern0pt}leq{\isacharunderscore}{\kern0pt}or{\isacharunderscore}{\kern0pt}geq\ consec{\isacharunderscore}{\kern0pt}htps{\isacharunderscore}{\kern0pt}in{\isacharunderscore}{\kern0pt}zip{\isacharunderscore}{\kern0pt}butlast{\isacharunderscore}{\kern0pt}tl\isanewline
\ \ \ \ \ \ \isacommand{by}\isamarkupfalse%
\ {\isacharparenleft}{\kern0pt}metis\ consec{\isacharunderscore}{\kern0pt}htps{\isacharunderscore}{\kern0pt}def\ in{\isacharunderscore}{\kern0pt}set{\isacharunderscore}{\kern0pt}butlastD\ in{\isacharunderscore}{\kern0pt}set{\isacharunderscore}{\kern0pt}tlD\ in{\isacharunderscore}{\kern0pt}set{\isacharunderscore}{\kern0pt}zipE{\isacharparenright}{\kern0pt}\isanewline
\ \ \isacommand{qed}\isamarkupfalse%
%
\endisatagproof
{\isafoldproof}%
%
\isadelimproof
%
\endisadelimproof
%
\begin{isamarkuptext}%
Justification for \isa{consec{\isacharunderscore}{\kern0pt}htps{\isacharunderscore}{\kern0pt}in{\isacharunderscore}{\kern0pt}interval}.%
\end{isamarkuptext}\isamarkuptrue%
\ \ \isacommand{lemma}\isamarkupfalse%
\ consec{\isacharunderscore}{\kern0pt}htps{\isacharunderscore}{\kern0pt}in{\isacharunderscore}{\kern0pt}interval{\isacharunderscore}{\kern0pt}iff{\isacharcolon}{\kern0pt}\isanewline
\ \ \ \ \isakeyword{assumes}\ {\isachardoublequoteopen}t\isactrlsub i\ {\isasymle}\ t\isactrlsub j{\isachardoublequoteclose}\isanewline
\ \ \ \ \isakeyword{shows}\ {\isachardoublequoteopen}{\isasymforall}t\ t{\isacharprime}{\kern0pt}{\isachardot}{\kern0pt}\ {\isacharparenleft}{\kern0pt}consec{\isacharunderscore}{\kern0pt}htps\ {\isasympi}s\ t\ t{\isacharprime}{\kern0pt}\ {\isasymand}\ t\isactrlsub i\ {\isasymle}\ t\ {\isasymand}\ t{\isacharprime}{\kern0pt}\ {\isasymle}\ t\isactrlsub j{\isacharparenright}{\kern0pt}\ \isanewline
\ \ \ \ \ \ \ \ \ \ \ \ {\isasymlongleftrightarrow}\ {\isacharparenleft}{\kern0pt}t{\isacharcomma}{\kern0pt}t{\isacharprime}{\kern0pt}{\isacharparenright}{\kern0pt}\ {\isasymin}\ set\ {\isacharparenleft}{\kern0pt}consec{\isacharunderscore}{\kern0pt}htps{\isacharunderscore}{\kern0pt}in{\isacharunderscore}{\kern0pt}interval\ {\isacharparenleft}{\kern0pt}htps{\isacharunderscore}{\kern0pt}exec\ {\isasympi}s{\isacharparenright}{\kern0pt}\ t\isactrlsub i\ t\isactrlsub j{\isacharparenright}{\kern0pt}{\isachardoublequoteclose}\isanewline
%
\isadelimproof
\ \ \ \ %
\endisadelimproof
%
\isatagproof
\isacommand{unfolding}\isamarkupfalse%
\ consec{\isacharunderscore}{\kern0pt}htps{\isacharunderscore}{\kern0pt}in{\isacharunderscore}{\kern0pt}interval{\isacharunderscore}{\kern0pt}def\isanewline
\ \ \ \ \isacommand{using}\isamarkupfalse%
\ assms\ happ{\isacharunderscore}{\kern0pt}time{\isacharunderscore}{\kern0pt}pts{\isacharunderscore}{\kern0pt}exec{\isacharunderscore}{\kern0pt}iff{\isacharunderscore}{\kern0pt}consec{\isacharunderscore}{\kern0pt}htps\ is{\isacharunderscore}{\kern0pt}htps{\isacharunderscore}{\kern0pt}htps{\isacharunderscore}{\kern0pt}exec\isanewline
\ \ \ \ \isacommand{by}\isamarkupfalse%
\ {\isacharparenleft}{\kern0pt}auto\ simp{\isacharcolon}{\kern0pt}\ consec{\isacharunderscore}{\kern0pt}htps{\isacharunderscore}{\kern0pt}def{\isacharparenright}{\kern0pt}%
\endisatagproof
{\isafoldproof}%
%
\isadelimproof
%
\endisadelimproof
%
\begin{isamarkuptext}%
All returned actions from \isa{simplify{\isacharunderscore}{\kern0pt}action} are correct instantiations of a \isa{plan{\isacharunderscore}{\kern0pt}action} 
  for an induced simple plan.%
\end{isamarkuptext}\isamarkuptrue%
\ \ \isacommand{lemma}\isamarkupfalse%
\ {\isacharparenleft}{\kern0pt}\isakeyword{in}\ wf{\isacharunderscore}{\kern0pt}ast{\isacharunderscore}{\kern0pt}problem{\isacharparenright}{\kern0pt}\ simp{\isacharunderscore}{\kern0pt}act{\isacharunderscore}{\kern0pt}inst{\isacharunderscore}{\kern0pt}of{\isacharunderscore}{\kern0pt}plan{\isacharunderscore}{\kern0pt}act{\isacharcolon}{\kern0pt}\ \isanewline
\ \ \ \ \isakeyword{assumes}\ {\isachardoublequoteopen}wf{\isacharunderscore}{\kern0pt}plan\ {\isasympi}s{\isachardoublequoteclose}\ \isanewline
\ \ \ \ \ \ \ \ \isakeyword{and}\ {\isachardoublequoteopen}{\isacharparenleft}{\kern0pt}t\isactrlsub {\isasympi}{\isacharcomma}{\kern0pt}{\isasympi}{\isacharparenright}{\kern0pt}\ {\isasymin}\ set\ {\isasympi}s{\isachardoublequoteclose}\ \isanewline
\ \ \ \ \ \ \ \ \isakeyword{and}\ {\isachardoublequoteopen}{\isacharparenleft}{\kern0pt}t\isactrlsub a{\isacharcomma}{\kern0pt}a{\isacharparenright}{\kern0pt}\ {\isasymin}\ set\ {\isacharparenleft}{\kern0pt}simplify{\isacharunderscore}{\kern0pt}action\ {\isacharparenleft}{\kern0pt}htps{\isacharunderscore}{\kern0pt}exec\ {\isasympi}s{\isacharparenright}{\kern0pt}\ {\isacharparenleft}{\kern0pt}t\isactrlsub {\isasympi}{\isacharcomma}{\kern0pt}{\isasympi}{\isacharparenright}{\kern0pt}{\isacharparenright}{\kern0pt}{\isachardoublequoteclose}\ \isanewline
\ \ \ \ \isakeyword{shows}\ {\isachardoublequoteopen}inst{\isacharunderscore}{\kern0pt}of{\isacharunderscore}{\kern0pt}plan{\isacharunderscore}{\kern0pt}action\ {\isasympi}s\ {\isacharparenleft}{\kern0pt}t\isactrlsub {\isasympi}{\isacharcomma}{\kern0pt}{\isasympi}{\isacharparenright}{\kern0pt}\ {\isacharparenleft}{\kern0pt}t\isactrlsub a{\isacharcomma}{\kern0pt}a{\isacharparenright}{\kern0pt}{\isachardoublequoteclose}\isanewline
%
\isadelimproof
\ \ \ \ %
\endisadelimproof
%
\isatagproof
\isacommand{using}\isamarkupfalse%
\ assms\ \isanewline
\ \ \isacommand{proof}\isamarkupfalse%
\ {\isacharparenleft}{\kern0pt}cases\ {\isasympi}{\isacharparenright}{\kern0pt}\isanewline
\ \ \ \ \isacommand{case}\isamarkupfalse%
\ {\isacharparenleft}{\kern0pt}Simple{\isacharunderscore}{\kern0pt}Plan{\isacharunderscore}{\kern0pt}Action\ n\ args{\isacharparenright}{\kern0pt}\isanewline
\ \ \ \ \isacommand{then}\isamarkupfalse%
\ \isacommand{show}\isamarkupfalse%
\ {\isacharquery}{\kern0pt}thesis\ \isanewline
\ \ \ \ \ \ \isacommand{using}\isamarkupfalse%
\ {\isacartoucheopen}{\isacharparenleft}{\kern0pt}t\isactrlsub a{\isacharcomma}{\kern0pt}a{\isacharparenright}{\kern0pt}\ {\isasymin}\ set\ {\isacharparenleft}{\kern0pt}simplify{\isacharunderscore}{\kern0pt}action\ {\isacharparenleft}{\kern0pt}htps{\isacharunderscore}{\kern0pt}exec\ {\isasympi}s{\isacharparenright}{\kern0pt}\ {\isacharparenleft}{\kern0pt}t\isactrlsub {\isasympi}{\isacharcomma}{\kern0pt}{\isasympi}{\isacharparenright}{\kern0pt}{\isacharparenright}{\kern0pt}{\isacartoucheclose}\ \isacommand{by}\isamarkupfalse%
\ auto\isanewline
\ \ \isacommand{next}\isamarkupfalse%
\isanewline
\ \ \ \ \isacommand{case}\isamarkupfalse%
\ {\isacharparenleft}{\kern0pt}Durative{\isacharunderscore}{\kern0pt}Plan{\isacharunderscore}{\kern0pt}Action\ n\ args\ d{\isacharparenright}{\kern0pt}\isanewline
\ \ \ \ \isacommand{let}\isamarkupfalse%
\ {\isacharquery}{\kern0pt}htps{\isacharequal}{\kern0pt}{\isachardoublequoteopen}htps{\isacharunderscore}{\kern0pt}exec\ {\isasympi}s{\isachardoublequoteclose}\isanewline
\ \ \ \ \isacommand{let}\isamarkupfalse%
\ {\isacharquery}{\kern0pt}a\isactrlsub s\isactrlsub c\isactrlsub h\isactrlsub e\isactrlsub m\isactrlsub a{\isacharequal}{\kern0pt}{\isachardoublequoteopen}the\ {\isacharparenleft}{\kern0pt}resolve{\isacharunderscore}{\kern0pt}action{\isacharunderscore}{\kern0pt}schema\ n{\isacharparenright}{\kern0pt}{\isachardoublequoteclose}\isanewline
\ \ \ \ \isacommand{let}\isamarkupfalse%
\ {\isacharquery}{\kern0pt}a\isactrlsub s\isactrlsub t\isactrlsub a\isactrlsub r\isactrlsub t{\isacharequal}{\kern0pt}{\isachardoublequoteopen}inst{\isacharunderscore}{\kern0pt}snap{\isacharunderscore}{\kern0pt}action\ {\isacharquery}{\kern0pt}a\isactrlsub s\isactrlsub c\isactrlsub h\isactrlsub e\isactrlsub m\isactrlsub a\ args\ At{\isacharunderscore}{\kern0pt}Start{\isachardoublequoteclose}\isanewline
\ \ \ \ \isacommand{let}\isamarkupfalse%
\ {\isacharquery}{\kern0pt}a\isactrlsub e\isactrlsub n\isactrlsub d{\isacharequal}{\kern0pt}{\isachardoublequoteopen}inst{\isacharunderscore}{\kern0pt}snap{\isacharunderscore}{\kern0pt}action\ {\isacharquery}{\kern0pt}a\isactrlsub s\isactrlsub c\isactrlsub h\isactrlsub e\isactrlsub m\isactrlsub a\ args\ At{\isacharunderscore}{\kern0pt}End{\isachardoublequoteclose}\isanewline
\ \ \ \ \isacommand{let}\isamarkupfalse%
\ {\isacharquery}{\kern0pt}a\isactrlsub i\isactrlsub n\isactrlsub v{\isacharequal}{\kern0pt}{\isachardoublequoteopen}inst{\isacharunderscore}{\kern0pt}snap{\isacharunderscore}{\kern0pt}action\ {\isacharquery}{\kern0pt}a\isactrlsub s\isactrlsub c\isactrlsub h\isactrlsub e\isactrlsub m\isactrlsub a\ args\ Over{\isacharunderscore}{\kern0pt}All{\isachardoublequoteclose}\isanewline
\ \ \ \ \isacommand{let}\isamarkupfalse%
\ {\isacharquery}{\kern0pt}t\isactrlsub i\isactrlsub n\isactrlsub v{\isacharequal}{\kern0pt}{\isachardoublequoteopen}{\isasymexists}t\isactrlsub i\ t\isactrlsub j{\isachardot}{\kern0pt}\ {\isacharparenleft}{\kern0pt}t\isactrlsub i{\isacharcomma}{\kern0pt}t\isactrlsub j{\isacharparenright}{\kern0pt}\ {\isasymin}\ set\ {\isacharparenleft}{\kern0pt}consec{\isacharunderscore}{\kern0pt}htps{\isacharunderscore}{\kern0pt}in{\isacharunderscore}{\kern0pt}interval\ {\isacharquery}{\kern0pt}htps\ t\isactrlsub {\isasympi}\ {\isacharparenleft}{\kern0pt}t\isactrlsub {\isasympi}\ {\isacharplus}{\kern0pt}\ d{\isacharparenright}{\kern0pt}{\isacharparenright}{\kern0pt}\ {\isasymand}\ t\isactrlsub a\ {\isacharequal}{\kern0pt}\ {\isacharparenleft}{\kern0pt}t\isactrlsub i{\isacharplus}{\kern0pt}t\isactrlsub j{\isacharparenright}{\kern0pt}\ {\isacharslash}{\kern0pt}\ {\isadigit{2}}{\isachardoublequoteclose}\isanewline
\ \ \ \ \isacommand{have}\isamarkupfalse%
\ {\isachardoublequoteopen}{\isacharparenleft}{\kern0pt}a\ {\isacharequal}{\kern0pt}\ {\isacharquery}{\kern0pt}a\isactrlsub s\isactrlsub t\isactrlsub a\isactrlsub r\isactrlsub t\ {\isasymand}\ t\isactrlsub {\isasympi}\ {\isacharequal}{\kern0pt}\ t\isactrlsub a{\isacharparenright}{\kern0pt}\ {\isasymor}\ {\isacharparenleft}{\kern0pt}a\ {\isacharequal}{\kern0pt}\ {\isacharquery}{\kern0pt}a\isactrlsub e\isactrlsub n\isactrlsub d\ {\isasymand}\ t\isactrlsub a\ {\isacharequal}{\kern0pt}\ t\isactrlsub {\isasympi}\ {\isacharplus}{\kern0pt}\ d{\isacharparenright}{\kern0pt}\ {\isasymor}\ {\isacharparenleft}{\kern0pt}a\ {\isacharequal}{\kern0pt}\ {\isacharquery}{\kern0pt}a\isactrlsub i\isactrlsub n\isactrlsub v\ {\isasymand}\ {\isacharquery}{\kern0pt}t\isactrlsub i\isactrlsub n\isactrlsub v{\isacharparenright}{\kern0pt}{\isachardoublequoteclose}\isanewline
\ \ \ \ \ \ \isacommand{using}\isamarkupfalse%
\ {\isacartoucheopen}{\isacharparenleft}{\kern0pt}t\isactrlsub a{\isacharcomma}{\kern0pt}a{\isacharparenright}{\kern0pt}\ {\isasymin}\ set\ {\isacharparenleft}{\kern0pt}simplify{\isacharunderscore}{\kern0pt}action\ {\isacharquery}{\kern0pt}htps\ {\isacharparenleft}{\kern0pt}t\isactrlsub {\isasympi}{\isacharcomma}{\kern0pt}{\isasympi}{\isacharparenright}{\kern0pt}{\isacharparenright}{\kern0pt}{\isacartoucheclose}\ Durative{\isacharunderscore}{\kern0pt}Plan{\isacharunderscore}{\kern0pt}Action\isanewline
\ \ \ \ \ \ \isacommand{by}\isamarkupfalse%
\ {\isacharparenleft}{\kern0pt}auto\ simp{\isacharcolon}{\kern0pt}\ Let{\isacharunderscore}{\kern0pt}def{\isacharparenright}{\kern0pt}\ blast{\isacharplus}{\kern0pt}\isanewline
\ \ \ \ \isacommand{then}\isamarkupfalse%
\ \isacommand{show}\isamarkupfalse%
\ {\isacharquery}{\kern0pt}thesis\isanewline
\ \ \ \ \isacommand{proof}\isamarkupfalse%
\ {\isacharparenleft}{\kern0pt}elim\ disjE{\isacharparenright}{\kern0pt}\isanewline
\ \ \ \ \ \ \isacommand{assume}\isamarkupfalse%
\ {\isachardoublequoteopen}a\ {\isacharequal}{\kern0pt}\ {\isacharquery}{\kern0pt}a\isactrlsub s\isactrlsub t\isactrlsub a\isactrlsub r\isactrlsub t\ {\isasymand}\ t\isactrlsub {\isasympi}\ {\isacharequal}{\kern0pt}\ t\isactrlsub a{\isachardoublequoteclose}\isanewline
\ \ \ \ \ \ \isacommand{then}\isamarkupfalse%
\ \isacommand{show}\isamarkupfalse%
\ {\isacharquery}{\kern0pt}thesis\isanewline
\ \ \ \ \ \ \ \ \isacommand{using}\isamarkupfalse%
\ Durative{\isacharunderscore}{\kern0pt}Plan{\isacharunderscore}{\kern0pt}Action\ \isacommand{by}\isamarkupfalse%
\ auto\isanewline
\ \ \ \ \isacommand{next}\isamarkupfalse%
\isanewline
\ \ \ \ \ \ \isacommand{assume}\isamarkupfalse%
\ {\isachardoublequoteopen}a\ {\isacharequal}{\kern0pt}\ {\isacharquery}{\kern0pt}a\isactrlsub e\isactrlsub n\isactrlsub d\ {\isasymand}\ t\isactrlsub a\ {\isacharequal}{\kern0pt}\ t\isactrlsub {\isasympi}{\isacharplus}{\kern0pt}d{\isachardoublequoteclose}\isanewline
\ \ \ \ \ \ \isacommand{then}\isamarkupfalse%
\ \isacommand{show}\isamarkupfalse%
\ {\isacharquery}{\kern0pt}thesis\isanewline
\ \ \ \ \ \ \ \ \isacommand{using}\isamarkupfalse%
\ Durative{\isacharunderscore}{\kern0pt}Plan{\isacharunderscore}{\kern0pt}Action\ \isacommand{by}\isamarkupfalse%
\ auto\isanewline
\ \ \ \ \isacommand{next}\isamarkupfalse%
\isanewline
\ \ \ \ \ \ \isacommand{assume}\isamarkupfalse%
\ {\isachardoublequoteopen}a\ {\isacharequal}{\kern0pt}\ {\isacharquery}{\kern0pt}a\isactrlsub i\isactrlsub n\isactrlsub v\ {\isasymand}\ {\isacharquery}{\kern0pt}t\isactrlsub i\isactrlsub n\isactrlsub v{\isachardoublequoteclose}\isanewline
\ \ \ \ \ \ \isacommand{then}\isamarkupfalse%
\ \isacommand{have}\isamarkupfalse%
\ {\isachardoublequoteopen}res{\isacharunderscore}{\kern0pt}inst{\isacharunderscore}{\kern0pt}snap{\isacharunderscore}{\kern0pt}action\ {\isasympi}\ Over{\isacharunderscore}{\kern0pt}All\ {\isacharequal}{\kern0pt}\ Some\ a{\isachardoublequoteclose}\isanewline
\ \ \ \ \ \ \ \ \isacommand{using}\isamarkupfalse%
\ Durative{\isacharunderscore}{\kern0pt}Plan{\isacharunderscore}{\kern0pt}Action\ \isacommand{by}\isamarkupfalse%
\ auto\isanewline
\ \ \ \ \ \ \isacommand{obtain}\isamarkupfalse%
\ t\isactrlsub i\ t\isactrlsub j\ \isakeyword{where}\ {\isachardoublequoteopen}{\isacharparenleft}{\kern0pt}t\isactrlsub i{\isacharcomma}{\kern0pt}t\isactrlsub j{\isacharparenright}{\kern0pt}\ {\isasymin}\ set\ {\isacharparenleft}{\kern0pt}consec{\isacharunderscore}{\kern0pt}htps{\isacharunderscore}{\kern0pt}in{\isacharunderscore}{\kern0pt}interval\ {\isacharquery}{\kern0pt}htps\ t\isactrlsub {\isasympi}\ {\isacharparenleft}{\kern0pt}t\isactrlsub {\isasympi}{\isacharplus}{\kern0pt}d{\isacharparenright}{\kern0pt}{\isacharparenright}{\kern0pt}{\isachardoublequoteclose}\ \isanewline
\ \ \ \ \ \ \ \ \ \ \ \ \ \ \ \ \ \ \ \ \isakeyword{and}\ {\isachardoublequoteopen}t\isactrlsub a\ {\isacharequal}{\kern0pt}\ {\isacharparenleft}{\kern0pt}t\isactrlsub i{\isacharplus}{\kern0pt}t\isactrlsub j{\isacharparenright}{\kern0pt}\ {\isacharslash}{\kern0pt}\ {\isadigit{2}}{\isachardoublequoteclose}\isanewline
\ \ \ \ \ \ \ \ \isacommand{using}\isamarkupfalse%
\ {\isacartoucheopen}a\ {\isacharequal}{\kern0pt}\ {\isacharquery}{\kern0pt}a\isactrlsub i\isactrlsub n\isactrlsub v\ {\isasymand}\ {\isacharquery}{\kern0pt}t\isactrlsub i\isactrlsub n\isactrlsub v{\isacartoucheclose}\ \isacommand{by}\isamarkupfalse%
\ auto\isanewline
\ \ \ \ \ \ \isacommand{have}\isamarkupfalse%
\ {\isachardoublequoteopen}t\isactrlsub {\isasympi}\ {\isasymle}\ t\isactrlsub {\isasympi}\ {\isacharplus}{\kern0pt}\ duration\ {\isasympi}{\isachardoublequoteclose}\isanewline
\ \ \ \ \ \ \ \ \isacommand{using}\isamarkupfalse%
\ assms\ Durative{\isacharunderscore}{\kern0pt}Plan{\isacharunderscore}{\kern0pt}Action\ \isacommand{unfolding}\isamarkupfalse%
\ wf{\isacharunderscore}{\kern0pt}plan{\isacharunderscore}{\kern0pt}def\isanewline
\ \ \ \ \ \ \ \ \isacommand{by}\isamarkupfalse%
\ {\isacharparenleft}{\kern0pt}auto\ simp{\isacharcolon}{\kern0pt}\ Let{\isacharunderscore}{\kern0pt}def\ split{\isacharcolon}{\kern0pt}\ option{\isachardot}{\kern0pt}splits\ ast{\isacharunderscore}{\kern0pt}action{\isacharunderscore}{\kern0pt}schema{\isachardot}{\kern0pt}splits{\isacharparenright}{\kern0pt}\isanewline
\ \ \ \ \ \ \isacommand{then}\isamarkupfalse%
\ \isacommand{have}\isamarkupfalse%
\ {\isachardoublequoteopen}consec{\isacharunderscore}{\kern0pt}htps\ {\isasympi}s\ t\isactrlsub i\ t\isactrlsub j\ {\isasymand}\ t\isactrlsub {\isasympi}\ {\isasymle}\ t\isactrlsub i\ {\isasymand}\ t\isactrlsub j\ {\isasymle}\ t\isactrlsub {\isasympi}\ {\isacharplus}{\kern0pt}\ d{\isachardoublequoteclose}\isanewline
\ \ \ \ \ \ \ \ \isacommand{using}\isamarkupfalse%
\ Durative{\isacharunderscore}{\kern0pt}Plan{\isacharunderscore}{\kern0pt}Action\ {\isacartoucheopen}{\isacharparenleft}{\kern0pt}t\isactrlsub i{\isacharcomma}{\kern0pt}t\isactrlsub j{\isacharparenright}{\kern0pt}\ {\isasymin}\ set\ {\isacharparenleft}{\kern0pt}consec{\isacharunderscore}{\kern0pt}htps{\isacharunderscore}{\kern0pt}in{\isacharunderscore}{\kern0pt}interval\ {\isacharquery}{\kern0pt}htps\ t\isactrlsub {\isasympi}\ {\isacharparenleft}{\kern0pt}t\isactrlsub {\isasympi}\ {\isacharplus}{\kern0pt}\ d{\isacharparenright}{\kern0pt}{\isacharparenright}{\kern0pt}{\isacartoucheclose}\isanewline
\ \ \ \ \ \ \ \ \ \ \ \ \ \ consec{\isacharunderscore}{\kern0pt}htps{\isacharunderscore}{\kern0pt}in{\isacharunderscore}{\kern0pt}interval{\isacharunderscore}{\kern0pt}iff\ \isanewline
\ \ \ \ \ \ \ \ \isacommand{by}\isamarkupfalse%
\ {\isacharparenleft}{\kern0pt}auto\ simp{\isacharcolon}{\kern0pt}\ consec{\isacharunderscore}{\kern0pt}htps{\isacharunderscore}{\kern0pt}def{\isacharparenright}{\kern0pt}\isanewline
\ \ \ \ \ \ \isacommand{then}\isamarkupfalse%
\ \isacommand{have}\isamarkupfalse%
\ {\isachardoublequoteopen}t\isactrlsub i\ {\isacharless}{\kern0pt}\ t\isactrlsub a\ {\isasymand}\ t\isactrlsub a\ {\isacharless}{\kern0pt}\ t\isactrlsub j{\isachardoublequoteclose}\isanewline
\ \ \ \ \ \ \ \ \isacommand{using}\isamarkupfalse%
\ {\isacartoucheopen}t\isactrlsub a\ {\isacharequal}{\kern0pt}\ {\isacharparenleft}{\kern0pt}t\isactrlsub i{\isacharplus}{\kern0pt}t\isactrlsub j{\isacharparenright}{\kern0pt}\ {\isacharslash}{\kern0pt}\ {\isadigit{2}}{\isacartoucheclose}\isanewline
\ \ \ \ \ \ \ \ \isacommand{by}\isamarkupfalse%
\ {\isacharparenleft}{\kern0pt}auto\ simp{\isacharcolon}{\kern0pt}\ consec{\isacharunderscore}{\kern0pt}htps{\isacharunderscore}{\kern0pt}def{\isacharparenright}{\kern0pt}\isanewline
\ \ \ \ \ \ \isacommand{then}\isamarkupfalse%
\ \isacommand{show}\isamarkupfalse%
\ {\isacharquery}{\kern0pt}thesis\isanewline
\ \ \ \ \ \ \ \ \isacommand{using}\isamarkupfalse%
\ Durative{\isacharunderscore}{\kern0pt}Plan{\isacharunderscore}{\kern0pt}Action\ {\isacartoucheopen}res{\isacharunderscore}{\kern0pt}inst{\isacharunderscore}{\kern0pt}snap{\isacharunderscore}{\kern0pt}action\ {\isasympi}\ Over{\isacharunderscore}{\kern0pt}All\ {\isacharequal}{\kern0pt}\ Some\ a{\isacartoucheclose}\ \isanewline
\ \ \ \ \ \ \ \ \ \ \ \ \ \ {\isacartoucheopen}consec{\isacharunderscore}{\kern0pt}htps\ {\isasympi}s\ t\isactrlsub i\ t\isactrlsub j\ {\isasymand}\ t\isactrlsub {\isasympi}\ {\isasymle}\ t\isactrlsub i\ {\isasymand}\ t\isactrlsub j\ {\isasymle}\ t\isactrlsub {\isasympi}\ {\isacharplus}{\kern0pt}\ d{\isacartoucheclose}\isanewline
\ \ \ \ \ \ \ \ \isacommand{by}\isamarkupfalse%
\ auto\isanewline
\ \ \ \ \isacommand{qed}\isamarkupfalse%
\isanewline
\ \ \isacommand{qed}\isamarkupfalse%
%
\endisatagproof
{\isafoldproof}%
%
\isadelimproof
%
\endisadelimproof
%
\begin{isamarkuptext}%
Function to insort a happening into a happening sequence.%
\end{isamarkuptext}\isamarkuptrue%
\ \ \isacommand{fun}\isamarkupfalse%
\ insort{\isacharunderscore}{\kern0pt}happ\ {\isacharcolon}{\kern0pt}{\isacharcolon}{\kern0pt}\ {\isachardoublequoteopen}happening\ {\isasymRightarrow}\ happening\ list\ {\isasymRightarrow}\ happening\ list{\isachardoublequoteclose}\ \isakeyword{where}\isanewline
\ \ \ \ {\isachardoublequoteopen}insort{\isacharunderscore}{\kern0pt}happ\ {\isacharparenleft}{\kern0pt}t\isactrlsub i{\isacharcomma}{\kern0pt}A\isactrlsub i{\isacharparenright}{\kern0pt}\ {\isacharbrackleft}{\kern0pt}{\isacharbrackright}{\kern0pt}\ {\isacharequal}{\kern0pt}\ {\isacharbrackleft}{\kern0pt}{\isacharparenleft}{\kern0pt}t\isactrlsub i{\isacharcomma}{\kern0pt}A\isactrlsub i{\isacharparenright}{\kern0pt}{\isacharbrackright}{\kern0pt}{\isachardoublequoteclose}\isanewline
\ \ {\isacharbar}{\kern0pt}\ {\isachardoublequoteopen}insort{\isacharunderscore}{\kern0pt}happ\ {\isacharparenleft}{\kern0pt}t\isactrlsub i{\isacharcomma}{\kern0pt}A\isactrlsub i{\isacharparenright}{\kern0pt}\ {\isacharparenleft}{\kern0pt}{\isacharparenleft}{\kern0pt}t\isactrlsub j{\isacharcomma}{\kern0pt}A\isactrlsub j{\isacharparenright}{\kern0pt}{\isacharhash}{\kern0pt}hs{\isacharparenright}{\kern0pt}\ {\isacharequal}{\kern0pt}\ \isanewline
\ \ \ \ \ \ {\isacharparenleft}{\kern0pt}if\ t\isactrlsub i\ {\isacharless}{\kern0pt}\ t\isactrlsub j\ then\ {\isacharparenleft}{\kern0pt}t\isactrlsub i{\isacharcomma}{\kern0pt}A\isactrlsub i{\isacharparenright}{\kern0pt}{\isacharhash}{\kern0pt}{\isacharparenleft}{\kern0pt}t\isactrlsub j{\isacharcomma}{\kern0pt}A\isactrlsub j{\isacharparenright}{\kern0pt}{\isacharhash}{\kern0pt}hs\isanewline
\ \ \ \ \ \ else\ if\ t\isactrlsub i\ {\isacharequal}{\kern0pt}\ t\isactrlsub j\ then\ {\isacharparenleft}{\kern0pt}t\isactrlsub j{\isacharcomma}{\kern0pt}A\isactrlsub i\ {\isacharat}{\kern0pt}\ A\isactrlsub j{\isacharparenright}{\kern0pt}{\isacharhash}{\kern0pt}hs\isanewline
\ \ \ \ \ \ else\ {\isacharparenleft}{\kern0pt}t\isactrlsub j{\isacharcomma}{\kern0pt}A\isactrlsub j{\isacharparenright}{\kern0pt}{\isacharhash}{\kern0pt}{\isacharparenleft}{\kern0pt}insort{\isacharunderscore}{\kern0pt}happ\ {\isacharparenleft}{\kern0pt}t\isactrlsub i{\isacharcomma}{\kern0pt}A\isactrlsub i{\isacharparenright}{\kern0pt}\ hs{\isacharparenright}{\kern0pt}{\isacharparenright}{\kern0pt}{\isachardoublequoteclose}%
\begin{isamarkuptext}%
Properties of \isa{insort{\isacharunderscore}{\kern0pt}happ}%
\end{isamarkuptext}\isamarkuptrue%
\ \ \isacommand{lemma}\isamarkupfalse%
\ insort{\isacharunderscore}{\kern0pt}happ{\isacharunderscore}{\kern0pt}insort{\isacharunderscore}{\kern0pt}insert{\isacharcolon}{\kern0pt}\ \isanewline
\ \ \ \ {\isachardoublequoteopen}sorted\ {\isacharparenleft}{\kern0pt}map\ fst\ hs{\isacharparenright}{\kern0pt}\ {\isasymLongrightarrow}\ insort{\isacharunderscore}{\kern0pt}insert\ {\isacharparenleft}{\kern0pt}fst\ h{\isacharparenright}{\kern0pt}\ {\isacharparenleft}{\kern0pt}map\ fst\ hs{\isacharparenright}{\kern0pt}\ {\isacharequal}{\kern0pt}\ map\ fst\ {\isacharparenleft}{\kern0pt}insort{\isacharunderscore}{\kern0pt}happ\ h\ hs{\isacharparenright}{\kern0pt}{\isachardoublequoteclose}\isanewline
%
\isadelimproof
\ \ \ \ %
\endisadelimproof
%
\isatagproof
\isacommand{by}\isamarkupfalse%
\ {\isacharparenleft}{\kern0pt}induction\ h\ hs\ rule{\isacharcolon}{\kern0pt}\ insort{\isacharunderscore}{\kern0pt}happ{\isachardot}{\kern0pt}induct{\isacharparenright}{\kern0pt}\ {\isacharparenleft}{\kern0pt}auto\ simp{\isacharcolon}{\kern0pt}\ insort{\isacharunderscore}{\kern0pt}insert{\isacharunderscore}{\kern0pt}key{\isacharunderscore}{\kern0pt}def\ split{\isacharcolon}{\kern0pt}\ if{\isacharunderscore}{\kern0pt}splits{\isacharparenright}{\kern0pt}%
\endisatagproof
{\isafoldproof}%
%
\isadelimproof
%
\endisadelimproof
%
\begin{isamarkuptext}%
\isa{insort{\isacharunderscore}{\kern0pt}happ} preserves sorted-ness.%
\end{isamarkuptext}\isamarkuptrue%
\ \ \isacommand{lemma}\isamarkupfalse%
\ insort{\isacharunderscore}{\kern0pt}happ{\isacharunderscore}{\kern0pt}distinct{\isacharcolon}{\kern0pt}\ {\isachardoublequoteopen}strict{\isacharunderscore}{\kern0pt}sorted\ {\isacharparenleft}{\kern0pt}map\ fst\ hs{\isacharparenright}{\kern0pt}\ {\isasymLongrightarrow}\ distinct\ {\isacharparenleft}{\kern0pt}map\ fst\ {\isacharparenleft}{\kern0pt}insort{\isacharunderscore}{\kern0pt}happ\ h\ hs{\isacharparenright}{\kern0pt}{\isacharparenright}{\kern0pt}{\isachardoublequoteclose}\isanewline
%
\isadelimproof
\ \ \ \ %
\endisadelimproof
%
\isatagproof
\isacommand{using}\isamarkupfalse%
\ strict{\isacharunderscore}{\kern0pt}sorted{\isacharunderscore}{\kern0pt}iff\isanewline
\ \ \ \ \isacommand{by}\isamarkupfalse%
\ {\isacharparenleft}{\kern0pt}metis\ insort{\isacharunderscore}{\kern0pt}happ{\isacharunderscore}{\kern0pt}insort{\isacharunderscore}{\kern0pt}insert\ distinct{\isacharunderscore}{\kern0pt}insort{\isacharunderscore}{\kern0pt}insert{\isacharparenright}{\kern0pt}%
\endisatagproof
{\isafoldproof}%
%
\isadelimproof
\isanewline
%
\endisadelimproof
\isanewline
\ \ \isacommand{lemma}\isamarkupfalse%
\ insort{\isacharunderscore}{\kern0pt}happ{\isacharunderscore}{\kern0pt}not{\isacharunderscore}{\kern0pt}Nil{\isacharcolon}{\kern0pt}\isanewline
\ \ \ \ \isakeyword{assumes}\ {\isachardoublequoteopen}A\isactrlsub i\ {\isasymnoteq}\ {\isacharbrackleft}{\kern0pt}{\isacharbrackright}{\kern0pt}{\isachardoublequoteclose}\ \isakeyword{and}\ {\isachardoublequoteopen}{\isasymforall}{\isacharparenleft}{\kern0pt}t\isactrlsub j{\isacharcomma}{\kern0pt}A\isactrlsub j{\isacharparenright}{\kern0pt}\ {\isasymin}\ set\ hs{\isachardot}{\kern0pt}\ A\isactrlsub j\ {\isasymnoteq}\ {\isacharbrackleft}{\kern0pt}{\isacharbrackright}{\kern0pt}{\isachardoublequoteclose}\ \isanewline
\ \ \ \ \isakeyword{shows}\ {\isachardoublequoteopen}{\isasymforall}{\isacharparenleft}{\kern0pt}t\isactrlsub j{\isacharcomma}{\kern0pt}A\isactrlsub j{\isacharparenright}{\kern0pt}\ {\isasymin}\ set\ {\isacharparenleft}{\kern0pt}insort{\isacharunderscore}{\kern0pt}happ\ {\isacharparenleft}{\kern0pt}t\isactrlsub i{\isacharcomma}{\kern0pt}A\isactrlsub i{\isacharparenright}{\kern0pt}\ hs{\isacharparenright}{\kern0pt}{\isachardot}{\kern0pt}\ A\isactrlsub j\ {\isasymnoteq}\ {\isacharbrackleft}{\kern0pt}{\isacharbrackright}{\kern0pt}{\isachardoublequoteclose}\isanewline
%
\isadelimproof
\ \ \ \ %
\endisadelimproof
%
\isatagproof
\isacommand{using}\isamarkupfalse%
\ assms\ \isacommand{by}\isamarkupfalse%
\ {\isacharparenleft}{\kern0pt}induction\ {\isachardoublequoteopen}{\isacharparenleft}{\kern0pt}t\isactrlsub i{\isacharcomma}{\kern0pt}A\isactrlsub i{\isacharparenright}{\kern0pt}{\isachardoublequoteclose}\ hs\ rule{\isacharcolon}{\kern0pt}\ insort{\isacharunderscore}{\kern0pt}happ{\isachardot}{\kern0pt}induct{\isacharparenright}{\kern0pt}\ auto%
\endisatagproof
{\isafoldproof}%
%
\isadelimproof
\isanewline
%
\endisadelimproof
\isanewline
\ \ \isacommand{lemma}\isamarkupfalse%
\ insort{\isacharunderscore}{\kern0pt}happ{\isacharunderscore}{\kern0pt}action{\isacharunderscore}{\kern0pt}subset{\isacharcolon}{\kern0pt}\ \isanewline
\ \ \ \ {\isachardoublequoteopen}{\isasymforall}{\isacharparenleft}{\kern0pt}t\isactrlsub a{\isacharcomma}{\kern0pt}A{\isacharparenright}{\kern0pt}\ {\isasymin}\ set\ hs\ {\isasymunion}\ {\isacharbraceleft}{\kern0pt}h{\isacharbraceright}{\kern0pt}{\isachardot}{\kern0pt}\ {\isasymexists}A{\isacharprime}{\kern0pt}{\isachardot}{\kern0pt}\ {\isacharparenleft}{\kern0pt}t\isactrlsub a{\isacharcomma}{\kern0pt}A{\isacharprime}{\kern0pt}{\isacharparenright}{\kern0pt}\ {\isasymin}\ set\ {\isacharparenleft}{\kern0pt}insort{\isacharunderscore}{\kern0pt}happ\ h\ hs{\isacharparenright}{\kern0pt}\ {\isasymand}\ set\ A\ {\isasymsubseteq}\ set\ A{\isacharprime}{\kern0pt}{\isachardoublequoteclose}\isanewline
%
\isadelimproof
\ \ \ \ %
\endisadelimproof
%
\isatagproof
\isacommand{by}\isamarkupfalse%
\ {\isacharparenleft}{\kern0pt}induction\ h\ hs\ rule{\isacharcolon}{\kern0pt}\ insort{\isacharunderscore}{\kern0pt}happ{\isachardot}{\kern0pt}induct{\isacharparenright}{\kern0pt}\ auto{\isacharplus}{\kern0pt}%
\endisatagproof
{\isafoldproof}%
%
\isadelimproof
\isanewline
%
\endisadelimproof
\isanewline
\ \ \isacommand{lemma}\isamarkupfalse%
\ insort{\isacharunderscore}{\kern0pt}happ{\isadigit{2}}{\isacharunderscore}{\kern0pt}action{\isacharunderscore}{\kern0pt}subset{\isacharcolon}{\kern0pt}\ \isanewline
\ \ \ \ {\isachardoublequoteopen}{\isasymexists}A{\isacharprime}{\kern0pt}{\isachardot}{\kern0pt}\ {\isacharparenleft}{\kern0pt}t\isactrlsub {\isadigit{2}}{\isacharcomma}{\kern0pt}A{\isacharprime}{\kern0pt}{\isacharparenright}{\kern0pt}\ {\isasymin}\ set\ {\isacharparenleft}{\kern0pt}insort{\isacharunderscore}{\kern0pt}happ\ {\isacharparenleft}{\kern0pt}t\isactrlsub {\isadigit{1}}{\isacharcomma}{\kern0pt}A\isactrlsub {\isadigit{1}}{\isacharparenright}{\kern0pt}\ {\isacharparenleft}{\kern0pt}insort{\isacharunderscore}{\kern0pt}happ\ {\isacharparenleft}{\kern0pt}t\isactrlsub {\isadigit{2}}{\isacharcomma}{\kern0pt}A\isactrlsub {\isadigit{2}}{\isacharparenright}{\kern0pt}\ hs{\isacharparenright}{\kern0pt}{\isacharparenright}{\kern0pt}\ {\isasymand}\ set\ A\isactrlsub {\isadigit{2}}\ {\isasymsubseteq}\ set\ A{\isacharprime}{\kern0pt}{\isachardoublequoteclose}\isanewline
%
\isadelimproof
\ \ \ \ %
\endisadelimproof
%
\isatagproof
\isacommand{using}\isamarkupfalse%
\ insort{\isacharunderscore}{\kern0pt}happ{\isacharunderscore}{\kern0pt}action{\isacharunderscore}{\kern0pt}subset\isanewline
\ \ \ \ \isacommand{by}\isamarkupfalse%
\ {\isacharparenleft}{\kern0pt}induction\ {\isachardoublequoteopen}{\isacharparenleft}{\kern0pt}t\isactrlsub {\isadigit{2}}{\isacharcomma}{\kern0pt}A\isactrlsub {\isadigit{2}}{\isacharparenright}{\kern0pt}{\isachardoublequoteclose}\ hs\ rule{\isacharcolon}{\kern0pt}\ insort{\isacharunderscore}{\kern0pt}happ{\isachardot}{\kern0pt}induct{\isacharparenright}{\kern0pt}\ auto%
\endisatagproof
{\isafoldproof}%
%
\isadelimproof
\isanewline
%
\endisadelimproof
\ \ \isanewline
\isanewline
\ \ \isacommand{lemma}\isamarkupfalse%
\ insort{\isacharunderscore}{\kern0pt}happ{\isacharunderscore}{\kern0pt}sorted{\isacharcolon}{\kern0pt}\ {\isachardoublequoteopen}sorted\ {\isacharparenleft}{\kern0pt}map\ fst\ hs{\isacharparenright}{\kern0pt}\ {\isasymLongrightarrow}\ sorted\ {\isacharparenleft}{\kern0pt}map\ fst\ {\isacharparenleft}{\kern0pt}insort{\isacharunderscore}{\kern0pt}happ\ h\ hs{\isacharparenright}{\kern0pt}{\isacharparenright}{\kern0pt}{\isachardoublequoteclose}\isanewline
%
\isadelimproof
\ \ \ \ %
\endisadelimproof
%
\isatagproof
\isacommand{by}\isamarkupfalse%
\ {\isacharparenleft}{\kern0pt}metis\ insort{\isacharunderscore}{\kern0pt}happ{\isacharunderscore}{\kern0pt}insort{\isacharunderscore}{\kern0pt}insert\ sorted{\isacharunderscore}{\kern0pt}insort{\isacharunderscore}{\kern0pt}insert{\isacharparenright}{\kern0pt}%
\endisatagproof
{\isafoldproof}%
%
\isadelimproof
\isanewline
%
\endisadelimproof
\isanewline
\ \ \isacommand{lemma}\isamarkupfalse%
\ insort{\isacharunderscore}{\kern0pt}happ{\isacharunderscore}{\kern0pt}strict{\isacharunderscore}{\kern0pt}sorted{\isacharcolon}{\kern0pt}\ \isanewline
\ \ \ \ {\isachardoublequoteopen}strict{\isacharunderscore}{\kern0pt}sorted\ {\isacharparenleft}{\kern0pt}map\ fst\ hs{\isacharparenright}{\kern0pt}\ {\isasymLongrightarrow}\ strict{\isacharunderscore}{\kern0pt}sorted\ {\isacharparenleft}{\kern0pt}map\ fst\ {\isacharparenleft}{\kern0pt}insort{\isacharunderscore}{\kern0pt}happ\ h\ hs{\isacharparenright}{\kern0pt}{\isacharparenright}{\kern0pt}{\isachardoublequoteclose}\isanewline
%
\isadelimproof
\ \ \ \ %
\endisadelimproof
%
\isatagproof
\isacommand{using}\isamarkupfalse%
\ insort{\isacharunderscore}{\kern0pt}happ{\isacharunderscore}{\kern0pt}sorted\ insort{\isacharunderscore}{\kern0pt}happ{\isacharunderscore}{\kern0pt}distinct\isanewline
\ \ \ \ \isacommand{by}\isamarkupfalse%
\ {\isacharparenleft}{\kern0pt}simp\ add{\isacharcolon}{\kern0pt}\ strict{\isacharunderscore}{\kern0pt}sorted{\isacharunderscore}{\kern0pt}iff{\isacharparenright}{\kern0pt}%
\endisatagproof
{\isafoldproof}%
%
\isadelimproof
%
\endisadelimproof
%
\begin{isamarkuptext}%
Function to insort multiple happening into a happening sequence.%
\end{isamarkuptext}\isamarkuptrue%
\ \ \isacommand{fun}\isamarkupfalse%
\ insort{\isacharunderscore}{\kern0pt}mult{\isacharunderscore}{\kern0pt}happs\ {\isacharcolon}{\kern0pt}{\isacharcolon}{\kern0pt}\ {\isachardoublequoteopen}happening\ list\ {\isasymRightarrow}\ happening\ list\ {\isasymRightarrow}\ happening\ list{\isachardoublequoteclose}\ \isakeyword{where}\isanewline
\ \ \ \ {\isachardoublequoteopen}insort{\isacharunderscore}{\kern0pt}mult{\isacharunderscore}{\kern0pt}happs\ {\isacharbrackleft}{\kern0pt}{\isacharbrackright}{\kern0pt}\ hs\isactrlsub {\isadigit{2}}\ {\isacharequal}{\kern0pt}\ hs\isactrlsub {\isadigit{2}}{\isachardoublequoteclose}\isanewline
\ \ {\isacharbar}{\kern0pt}\ {\isachardoublequoteopen}insort{\isacharunderscore}{\kern0pt}mult{\isacharunderscore}{\kern0pt}happs\ {\isacharparenleft}{\kern0pt}h{\isacharhash}{\kern0pt}hs\isactrlsub {\isadigit{1}}{\isacharparenright}{\kern0pt}\ hs\isactrlsub {\isadigit{2}}\ {\isacharequal}{\kern0pt}\ insort{\isacharunderscore}{\kern0pt}happ\ h\ {\isacharparenleft}{\kern0pt}insort{\isacharunderscore}{\kern0pt}mult{\isacharunderscore}{\kern0pt}happs\ hs\isactrlsub {\isadigit{1}}\ hs\isactrlsub {\isadigit{2}}{\isacharparenright}{\kern0pt}{\isachardoublequoteclose}\isanewline
\isanewline
\ \ \isacommand{fun}\isamarkupfalse%
\ insort{\isacharunderscore}{\kern0pt}timed{\isacharunderscore}{\kern0pt}ground{\isacharunderscore}{\kern0pt}acts\ {\isacharcolon}{\kern0pt}{\isacharcolon}{\kern0pt}\ {\isachardoublequoteopen}{\isacharparenleft}{\kern0pt}time\ {\isasymtimes}\ ground{\isacharunderscore}{\kern0pt}action{\isacharparenright}{\kern0pt}\ list\ {\isasymRightarrow}\ happening\ list\ {\isasymRightarrow}\ happening\ list{\isachardoublequoteclose}\ \isakeyword{where}\isanewline
\ \ \ \ {\isachardoublequoteopen}insort{\isacharunderscore}{\kern0pt}timed{\isacharunderscore}{\kern0pt}ground{\isacharunderscore}{\kern0pt}acts\ {\isacharbrackleft}{\kern0pt}{\isacharbrackright}{\kern0pt}\ hs\ {\isacharequal}{\kern0pt}\ hs{\isachardoublequoteclose}\isanewline
\ \ {\isacharbar}{\kern0pt}\ {\isachardoublequoteopen}insort{\isacharunderscore}{\kern0pt}timed{\isacharunderscore}{\kern0pt}ground{\isacharunderscore}{\kern0pt}acts\ {\isacharparenleft}{\kern0pt}{\isacharparenleft}{\kern0pt}t\isactrlsub a{\isacharcomma}{\kern0pt}a{\isacharparenright}{\kern0pt}{\isacharhash}{\kern0pt}as{\isacharparenright}{\kern0pt}\ hs\ {\isacharequal}{\kern0pt}\ \isanewline
\ \ \ \ insort{\isacharunderscore}{\kern0pt}happ\ {\isacharparenleft}{\kern0pt}t\isactrlsub a{\isacharcomma}{\kern0pt}{\isacharbrackleft}{\kern0pt}a{\isacharbrackright}{\kern0pt}{\isacharparenright}{\kern0pt}\ {\isacharparenleft}{\kern0pt}insort{\isacharunderscore}{\kern0pt}timed{\isacharunderscore}{\kern0pt}ground{\isacharunderscore}{\kern0pt}acts\ as\ hs{\isacharparenright}{\kern0pt}{\isachardoublequoteclose}\isanewline
\isanewline
\ \ \isacommand{lemma}\isamarkupfalse%
\ insort{\isacharunderscore}{\kern0pt}timed{\isacharunderscore}{\kern0pt}ground{\isacharunderscore}{\kern0pt}acts{\isacharunderscore}{\kern0pt}correct{\isacharcolon}{\kern0pt}\ \isanewline
\ \ \ \ {\isachardoublequoteopen}insort{\isacharunderscore}{\kern0pt}mult{\isacharunderscore}{\kern0pt}happs\ {\isacharparenleft}{\kern0pt}map\ {\isacharparenleft}{\kern0pt}{\isasymlambda}{\isacharparenleft}{\kern0pt}t{\isacharcomma}{\kern0pt}a{\isacharparenright}{\kern0pt}{\isachardot}{\kern0pt}\ {\isacharparenleft}{\kern0pt}t{\isacharcomma}{\kern0pt}{\isacharbrackleft}{\kern0pt}a{\isacharbrackright}{\kern0pt}{\isacharparenright}{\kern0pt}{\isacharparenright}{\kern0pt}\ as{\isacharparenright}{\kern0pt}\ hs\ {\isacharequal}{\kern0pt}\ insort{\isacharunderscore}{\kern0pt}timed{\isacharunderscore}{\kern0pt}ground{\isacharunderscore}{\kern0pt}acts\ as\ hs{\isachardoublequoteclose}\isanewline
%
\isadelimproof
\ \ \ \ %
\endisadelimproof
%
\isatagproof
\isacommand{by}\isamarkupfalse%
\ {\isacharparenleft}{\kern0pt}induction\ as\ hs\ rule{\isacharcolon}{\kern0pt}\ insort{\isacharunderscore}{\kern0pt}timed{\isacharunderscore}{\kern0pt}ground{\isacharunderscore}{\kern0pt}acts{\isachardot}{\kern0pt}induct{\isacharparenright}{\kern0pt}\ auto%
\endisatagproof
{\isafoldproof}%
%
\isadelimproof
%
\endisadelimproof
%
\begin{isamarkuptext}%
Properties of \isa{insort{\isacharunderscore}{\kern0pt}mult{\isacharunderscore}{\kern0pt}happs}%
\end{isamarkuptext}\isamarkuptrue%
\ \ \isacommand{lemma}\isamarkupfalse%
\ insort{\isacharunderscore}{\kern0pt}mult{\isacharunderscore}{\kern0pt}happs{\isacharunderscore}{\kern0pt}strict{\isacharunderscore}{\kern0pt}sorted{\isacharcolon}{\kern0pt}\ \isanewline
\ \ \ \ {\isachardoublequoteopen}strict{\isacharunderscore}{\kern0pt}sorted\ {\isacharparenleft}{\kern0pt}map\ fst\ hs\isactrlsub {\isadigit{2}}{\isacharparenright}{\kern0pt}\ {\isasymLongrightarrow}\ strict{\isacharunderscore}{\kern0pt}sorted\ {\isacharparenleft}{\kern0pt}map\ fst\ {\isacharparenleft}{\kern0pt}insort{\isacharunderscore}{\kern0pt}mult{\isacharunderscore}{\kern0pt}happs\ hs\isactrlsub {\isadigit{1}}\ hs\isactrlsub {\isadigit{2}}{\isacharparenright}{\kern0pt}{\isacharparenright}{\kern0pt}{\isachardoublequoteclose}\isanewline
%
\isadelimproof
\ \ \ \ %
\endisadelimproof
%
\isatagproof
\isacommand{by}\isamarkupfalse%
\ {\isacharparenleft}{\kern0pt}induction\ hs\isactrlsub {\isadigit{1}}\ arbitrary{\isacharcolon}{\kern0pt}\ hs\isactrlsub {\isadigit{2}}{\isacharparenright}{\kern0pt}\ {\isacharparenleft}{\kern0pt}auto\ simp{\isacharcolon}{\kern0pt}\ insort{\isacharunderscore}{\kern0pt}happ{\isacharunderscore}{\kern0pt}strict{\isacharunderscore}{\kern0pt}sorted{\isacharparenright}{\kern0pt}%
\endisatagproof
{\isafoldproof}%
%
\isadelimproof
\isanewline
%
\endisadelimproof
\isanewline
\ \ \isacommand{lemma}\isamarkupfalse%
\ insort{\isacharunderscore}{\kern0pt}mult{\isacharunderscore}{\kern0pt}happs{\isacharunderscore}{\kern0pt}not{\isacharunderscore}{\kern0pt}Nil{\isacharcolon}{\kern0pt}\ \isanewline
\ \ \ \ \isakeyword{assumes}\ {\isachardoublequoteopen}{\isasymforall}{\isacharparenleft}{\kern0pt}t\isactrlsub a{\isacharcomma}{\kern0pt}A{\isacharparenright}{\kern0pt}\ {\isasymin}\ set\ xs{\isachardot}{\kern0pt}\ A\ {\isasymnoteq}\ {\isacharbrackleft}{\kern0pt}{\isacharbrackright}{\kern0pt}{\isachardoublequoteclose}\ \isakeyword{and}\ {\isachardoublequoteopen}{\isasymforall}{\isacharparenleft}{\kern0pt}t\isactrlsub a{\isacharcomma}{\kern0pt}A{\isacharparenright}{\kern0pt}\ {\isasymin}\ set\ ys{\isachardot}{\kern0pt}\ A\ {\isasymnoteq}\ {\isacharbrackleft}{\kern0pt}{\isacharbrackright}{\kern0pt}{\isachardoublequoteclose}\ \isanewline
\ \ \ \ \isakeyword{shows}\ {\isachardoublequoteopen}{\isasymforall}{\isacharparenleft}{\kern0pt}t\isactrlsub a{\isacharcomma}{\kern0pt}A{\isacharparenright}{\kern0pt}\ {\isasymin}\ set\ {\isacharparenleft}{\kern0pt}insort{\isacharunderscore}{\kern0pt}mult{\isacharunderscore}{\kern0pt}happs\ xs\ ys{\isacharparenright}{\kern0pt}{\isachardot}{\kern0pt}\ A\ {\isasymnoteq}\ {\isacharbrackleft}{\kern0pt}{\isacharbrackright}{\kern0pt}{\isachardoublequoteclose}\isanewline
%
\isadelimproof
\ \ \ \ %
\endisadelimproof
%
\isatagproof
\isacommand{using}\isamarkupfalse%
\ assms\isanewline
\ \ \isacommand{proof}\isamarkupfalse%
\ {\isacharparenleft}{\kern0pt}induction\ xs\ ys\ rule{\isacharcolon}{\kern0pt}\ insort{\isacharunderscore}{\kern0pt}mult{\isacharunderscore}{\kern0pt}happs{\isachardot}{\kern0pt}induct{\isacharparenright}{\kern0pt}\isanewline
\ \ \ \ \isacommand{case}\isamarkupfalse%
\ {\isacharparenleft}{\kern0pt}{\isadigit{1}}\ ys{\isacharparenright}{\kern0pt}\isanewline
\ \ \ \ \isacommand{then}\isamarkupfalse%
\ \isacommand{show}\isamarkupfalse%
\ {\isacharquery}{\kern0pt}case\ \isacommand{by}\isamarkupfalse%
\ simp\isanewline
\ \ \isacommand{next}\isamarkupfalse%
\isanewline
\ \ \ \ \isacommand{case}\isamarkupfalse%
\ {\isacharparenleft}{\kern0pt}{\isadigit{2}}\ x\ xs\ ys{\isacharparenright}{\kern0pt}\isanewline
\ \ \ \ \isacommand{then}\isamarkupfalse%
\ \isacommand{show}\isamarkupfalse%
\ {\isacharquery}{\kern0pt}case\isanewline
\ \ \ \ \isacommand{proof}\isamarkupfalse%
\ {\isacharparenleft}{\kern0pt}cases\ x{\isacharparenright}{\kern0pt}\isanewline
\ \ \ \ \ \ \isacommand{case}\isamarkupfalse%
\ {\isacharparenleft}{\kern0pt}Pair\ t\isactrlsub a\ A{\isacharparenright}{\kern0pt}\isanewline
\ \ \ \ \ \ \isacommand{then}\isamarkupfalse%
\ \isacommand{have}\isamarkupfalse%
\ {\isachardoublequoteopen}{\isasymforall}{\isacharparenleft}{\kern0pt}t\isactrlsub a{\isacharcomma}{\kern0pt}A{\isacharparenright}{\kern0pt}\ {\isasymin}\ set\ xs{\isachardot}{\kern0pt}\ A\ {\isasymnoteq}\ {\isacharbrackleft}{\kern0pt}{\isacharbrackright}{\kern0pt}{\isachardoublequoteclose}\ \isakeyword{and}\ {\isachardoublequoteopen}A\ {\isasymnoteq}\ {\isacharbrackleft}{\kern0pt}{\isacharbrackright}{\kern0pt}{\isachardoublequoteclose}\isanewline
\ \ \ \ \ \ \ \ \isacommand{using}\isamarkupfalse%
\ {\isacartoucheopen}{\isasymforall}{\isacharparenleft}{\kern0pt}t\isactrlsub a{\isacharcomma}{\kern0pt}A{\isacharparenright}{\kern0pt}\ {\isasymin}\ set\ {\isacharparenleft}{\kern0pt}x{\isacharhash}{\kern0pt}xs{\isacharparenright}{\kern0pt}{\isachardot}{\kern0pt}\ A\ {\isasymnoteq}\ {\isacharbrackleft}{\kern0pt}{\isacharbrackright}{\kern0pt}{\isacartoucheclose}\isanewline
\ \ \ \ \ \ \ \ \isacommand{by}\isamarkupfalse%
\ auto\isanewline
\ \ \ \ \ \ \isacommand{then}\isamarkupfalse%
\ \isacommand{show}\isamarkupfalse%
\ {\isacharquery}{\kern0pt}thesis\ \isanewline
\ \ \ \ \ \ \ \ \isacommand{using}\isamarkupfalse%
\ Pair\ {\isachardoublequoteopen}{\isadigit{2}}{\isachardot}{\kern0pt}IH{\isachardoublequoteclose}{\isacharbrackleft}{\kern0pt}OF\ {\isacartoucheopen}{\isasymforall}{\isacharparenleft}{\kern0pt}t\isactrlsub a{\isacharcomma}{\kern0pt}A{\isacharparenright}{\kern0pt}\ {\isasymin}\ set\ xs{\isachardot}{\kern0pt}\ A\ {\isasymnoteq}\ {\isacharbrackleft}{\kern0pt}{\isacharbrackright}{\kern0pt}{\isacartoucheclose}\ {\isacartoucheopen}{\isasymforall}{\isacharparenleft}{\kern0pt}t\isactrlsub a{\isacharcomma}{\kern0pt}A{\isacharparenright}{\kern0pt}\ {\isasymin}\ set\ ys{\isachardot}{\kern0pt}\ A\ {\isasymnoteq}\ {\isacharbrackleft}{\kern0pt}{\isacharbrackright}{\kern0pt}{\isacartoucheclose}{\isacharbrackright}{\kern0pt}\isanewline
\ \ \ \ \ \ \ \ \ \ \ \ \ \ insort{\isacharunderscore}{\kern0pt}happ{\isacharunderscore}{\kern0pt}not{\isacharunderscore}{\kern0pt}Nil{\isacharbrackleft}{\kern0pt}OF\ {\isacartoucheopen}A\ {\isasymnoteq}\ {\isacharbrackleft}{\kern0pt}{\isacharbrackright}{\kern0pt}{\isacartoucheclose}{\isacharbrackright}{\kern0pt}\ \isacommand{by}\isamarkupfalse%
\ auto\isanewline
\ \ \ \ \isacommand{qed}\isamarkupfalse%
\isanewline
\ \ \isacommand{qed}\isamarkupfalse%
%
\endisatagproof
{\isafoldproof}%
%
\isadelimproof
\isanewline
%
\endisadelimproof
\isanewline
\ \ \isacommand{lemma}\isamarkupfalse%
\ insort{\isacharunderscore}{\kern0pt}mult{\isacharunderscore}{\kern0pt}happs{\isacharunderscore}{\kern0pt}action{\isacharunderscore}{\kern0pt}subset{\isadigit{1}}{\isacharcolon}{\kern0pt}\ \isanewline
\ \ \ \ {\isachardoublequoteopen}{\isacharparenleft}{\kern0pt}t\isactrlsub a{\isacharcomma}{\kern0pt}A{\isacharparenright}{\kern0pt}\ {\isasymin}\ set\ hs\isactrlsub {\isadigit{1}}\ {\isasymLongrightarrow}\ {\isasymexists}as{\isacharprime}{\kern0pt}{\isachardot}{\kern0pt}\ {\isacharparenleft}{\kern0pt}t\isactrlsub a{\isacharcomma}{\kern0pt}as{\isacharprime}{\kern0pt}{\isacharparenright}{\kern0pt}\ {\isasymin}\ set\ {\isacharparenleft}{\kern0pt}insort{\isacharunderscore}{\kern0pt}mult{\isacharunderscore}{\kern0pt}happs\ hs\isactrlsub {\isadigit{1}}\ hs\isactrlsub {\isadigit{2}}{\isacharparenright}{\kern0pt}\ {\isasymand}\ set\ A\ {\isasymsubseteq}\ set\ as{\isacharprime}{\kern0pt}{\isachardoublequoteclose}\isanewline
%
\isadelimproof
\ \ %
\endisadelimproof
%
\isatagproof
\isacommand{proof}\isamarkupfalse%
\ {\isacharparenleft}{\kern0pt}induction\ hs\isactrlsub {\isadigit{1}}{\isacharparenright}{\kern0pt}\isanewline
\ \ \ \ \isacommand{case}\isamarkupfalse%
\ Nil\isanewline
\ \ \ \ \isacommand{then}\isamarkupfalse%
\ \isacommand{show}\isamarkupfalse%
\ {\isacharquery}{\kern0pt}case\isanewline
\ \ \ \ \ \ \isacommand{by}\isamarkupfalse%
\ auto\isanewline
\ \ \isacommand{next}\isamarkupfalse%
\isanewline
\ \ \ \ \isacommand{case}\isamarkupfalse%
\ {\isacharparenleft}{\kern0pt}Cons\ h\ hs\isactrlsub {\isadigit{1}}{\isacharparenright}{\kern0pt}\isanewline
\ \ \ \ \isacommand{have}\isamarkupfalse%
\ {\isachardoublequoteopen}{\isacharparenleft}{\kern0pt}t\isactrlsub a{\isacharcomma}{\kern0pt}A{\isacharparenright}{\kern0pt}\ {\isacharequal}{\kern0pt}\ h\ {\isasymor}\ {\isacharparenleft}{\kern0pt}t\isactrlsub a{\isacharcomma}{\kern0pt}A{\isacharparenright}{\kern0pt}\ {\isasymnoteq}\ h{\isachardoublequoteclose}\isanewline
\ \ \ \ \ \ \isacommand{by}\isamarkupfalse%
\ simp\isanewline
\ \ \ \ \isacommand{then}\isamarkupfalse%
\ \isacommand{show}\isamarkupfalse%
\ {\isacharquery}{\kern0pt}case\isanewline
\ \ \ \ \isacommand{proof}\isamarkupfalse%
\isanewline
\ \ \ \ \ \ \isacommand{assume}\isamarkupfalse%
\ {\isachardoublequoteopen}{\isacharparenleft}{\kern0pt}t\isactrlsub a{\isacharcomma}{\kern0pt}A{\isacharparenright}{\kern0pt}\ {\isacharequal}{\kern0pt}\ h{\isachardoublequoteclose}\isanewline
\ \ \ \ \ \ \isacommand{then}\isamarkupfalse%
\ \isacommand{show}\isamarkupfalse%
\ {\isacharquery}{\kern0pt}case\isanewline
\ \ \ \ \ \ \ \ \isacommand{using}\isamarkupfalse%
\ insort{\isacharunderscore}{\kern0pt}happ{\isacharunderscore}{\kern0pt}action{\isacharunderscore}{\kern0pt}subset\ \isacommand{by}\isamarkupfalse%
\ auto\isanewline
\ \ \ \ \isacommand{next}\isamarkupfalse%
\isanewline
\ \ \ \ \ \ \isacommand{assume}\isamarkupfalse%
\ {\isachardoublequoteopen}{\isacharparenleft}{\kern0pt}t\isactrlsub a{\isacharcomma}{\kern0pt}A{\isacharparenright}{\kern0pt}\ {\isasymnoteq}\ h{\isachardoublequoteclose}\isanewline
\ \ \ \ \ \ \isacommand{then}\isamarkupfalse%
\ \isacommand{have}\isamarkupfalse%
\ {\isachardoublequoteopen}{\isacharparenleft}{\kern0pt}t\isactrlsub a{\isacharcomma}{\kern0pt}A{\isacharparenright}{\kern0pt}\ {\isasymin}\ set\ hs\isactrlsub {\isadigit{1}}{\isachardoublequoteclose}\isanewline
\ \ \ \ \ \ \ \ \isacommand{using}\isamarkupfalse%
\ Cons{\isachardot}{\kern0pt}prems\ \isacommand{by}\isamarkupfalse%
\ auto\isanewline
\ \ \ \ \ \ \isacommand{then}\isamarkupfalse%
\ \isacommand{obtain}\isamarkupfalse%
\ A{\isacharprime}{\kern0pt}\ \isakeyword{where}\ {\isachardoublequoteopen}{\isacharparenleft}{\kern0pt}t\isactrlsub a{\isacharcomma}{\kern0pt}A{\isacharprime}{\kern0pt}{\isacharparenright}{\kern0pt}\ {\isasymin}\ set\ {\isacharparenleft}{\kern0pt}insort{\isacharunderscore}{\kern0pt}mult{\isacharunderscore}{\kern0pt}happs\ hs\isactrlsub {\isadigit{1}}\ hs\isactrlsub {\isadigit{2}}{\isacharparenright}{\kern0pt}{\isachardoublequoteclose}\ \isakeyword{and}\ {\isachardoublequoteopen}set\ A\ {\isasymsubseteq}\ set\ A{\isacharprime}{\kern0pt}{\isachardoublequoteclose}\isanewline
\ \ \ \ \ \ \ \ \isacommand{using}\isamarkupfalse%
\ Cons{\isachardot}{\kern0pt}IH\ \isacommand{by}\isamarkupfalse%
\ auto\isanewline
\ \ \ \ \ \ \isacommand{then}\isamarkupfalse%
\ \isacommand{obtain}\isamarkupfalse%
\ A{\isacharprime}{\kern0pt}{\isacharprime}{\kern0pt}\ \isakeyword{where}\ {\isachardoublequoteopen}{\isacharparenleft}{\kern0pt}t\isactrlsub a{\isacharcomma}{\kern0pt}A{\isacharprime}{\kern0pt}{\isacharprime}{\kern0pt}{\isacharparenright}{\kern0pt}\ {\isasymin}\ set\ {\isacharparenleft}{\kern0pt}insort{\isacharunderscore}{\kern0pt}happ\ h\ {\isacharparenleft}{\kern0pt}insort{\isacharunderscore}{\kern0pt}mult{\isacharunderscore}{\kern0pt}happs\ hs\isactrlsub {\isadigit{1}}\ hs\isactrlsub {\isadigit{2}}{\isacharparenright}{\kern0pt}{\isacharparenright}{\kern0pt}{\isachardoublequoteclose}\ \isanewline
\ \ \ \ \ \ \ \ \ \ \ \ \ \ \ \ \ \ \ \ \ \ \ \isakeyword{and}\ {\isachardoublequoteopen}set\ A{\isacharprime}{\kern0pt}\ {\isasymsubseteq}\ set\ A{\isacharprime}{\kern0pt}{\isacharprime}{\kern0pt}{\isachardoublequoteclose}\isanewline
\ \ \ \ \ \ \ \ \isacommand{using}\isamarkupfalse%
\ insort{\isacharunderscore}{\kern0pt}happ{\isacharunderscore}{\kern0pt}action{\isacharunderscore}{\kern0pt}subset{\isacharbrackleft}{\kern0pt}\isakeyword{where}\ h{\isacharequal}{\kern0pt}{\isachardoublequoteopen}h{\isachardoublequoteclose}\ \isakeyword{and}\ hs{\isacharequal}{\kern0pt}{\isachardoublequoteopen}insort{\isacharunderscore}{\kern0pt}mult{\isacharunderscore}{\kern0pt}happs\ hs\isactrlsub {\isadigit{1}}\ hs\isactrlsub {\isadigit{2}}{\isachardoublequoteclose}{\isacharbrackright}{\kern0pt}\ \isacommand{by}\isamarkupfalse%
\ auto\isanewline
\ \ \ \ \ \ \isacommand{then}\isamarkupfalse%
\ \isacommand{show}\isamarkupfalse%
\ {\isacharquery}{\kern0pt}case\isanewline
\ \ \ \ \ \ \ \ \isacommand{using}\isamarkupfalse%
\ {\isacartoucheopen}set\ A\ {\isasymsubseteq}\ set\ A{\isacharprime}{\kern0pt}{\isacartoucheclose}\ \isacommand{by}\isamarkupfalse%
\ auto\isanewline
\ \ \ \ \isacommand{qed}\isamarkupfalse%
\isanewline
\ \ \isacommand{qed}\isamarkupfalse%
%
\endisatagproof
{\isafoldproof}%
%
\isadelimproof
\isanewline
%
\endisadelimproof
\isanewline
\ \ \isacommand{lemma}\isamarkupfalse%
\ insort{\isacharunderscore}{\kern0pt}mult{\isacharunderscore}{\kern0pt}happs{\isacharunderscore}{\kern0pt}action{\isacharunderscore}{\kern0pt}subset{\isadigit{2}}{\isacharcolon}{\kern0pt}\ \isanewline
\ \ \ \ \isakeyword{assumes}\ {\isachardoublequoteopen}{\isacharparenleft}{\kern0pt}t\isactrlsub a{\isacharcomma}{\kern0pt}A{\isacharparenright}{\kern0pt}\ {\isasymin}\ set\ hs\isactrlsub {\isadigit{2}}{\isachardoublequoteclose}\ \isanewline
\ \ \ \ \isakeyword{shows}\ {\isachardoublequoteopen}{\isasymexists}A{\isacharprime}{\kern0pt}{\isachardot}{\kern0pt}\ {\isacharparenleft}{\kern0pt}t\isactrlsub a{\isacharcomma}{\kern0pt}A{\isacharprime}{\kern0pt}{\isacharparenright}{\kern0pt}\ {\isasymin}\ set\ {\isacharparenleft}{\kern0pt}insort{\isacharunderscore}{\kern0pt}mult{\isacharunderscore}{\kern0pt}happs\ hs\isactrlsub {\isadigit{1}}\ hs\isactrlsub {\isadigit{2}}{\isacharparenright}{\kern0pt}\ {\isasymand}\ set\ A\ {\isasymsubseteq}\ set\ A{\isacharprime}{\kern0pt}{\isachardoublequoteclose}\isanewline
%
\isadelimproof
\ \ %
\endisadelimproof
%
\isatagproof
\isacommand{proof}\isamarkupfalse%
\ {\isacharparenleft}{\kern0pt}induction\ hs\isactrlsub {\isadigit{1}}{\isacharparenright}{\kern0pt}\isanewline
\ \ \ \ \isacommand{case}\isamarkupfalse%
\ Nil\isanewline
\ \ \ \ \isacommand{then}\isamarkupfalse%
\ \isacommand{show}\isamarkupfalse%
\ {\isacharquery}{\kern0pt}case\ \isanewline
\ \ \ \ \ \ \isacommand{using}\isamarkupfalse%
\ assms\ \isacommand{by}\isamarkupfalse%
\ auto\isanewline
\ \ \isacommand{next}\isamarkupfalse%
\isanewline
\ \ \ \ \isacommand{case}\isamarkupfalse%
\ {\isacharparenleft}{\kern0pt}Cons\ h\ hs\isactrlsub {\isadigit{1}}{\isacharparenright}{\kern0pt}\isanewline
\ \ \ \ \isacommand{obtain}\isamarkupfalse%
\ A{\isacharprime}{\kern0pt}\ \isakeyword{where}\ {\isachardoublequoteopen}{\isacharparenleft}{\kern0pt}t\isactrlsub a{\isacharcomma}{\kern0pt}A{\isacharprime}{\kern0pt}{\isacharparenright}{\kern0pt}\ {\isasymin}\ set\ {\isacharparenleft}{\kern0pt}insort{\isacharunderscore}{\kern0pt}mult{\isacharunderscore}{\kern0pt}happs\ hs\isactrlsub {\isadigit{1}}\ hs\isactrlsub {\isadigit{2}}{\isacharparenright}{\kern0pt}{\isachardoublequoteclose}\ \isakeyword{and}\ {\isachardoublequoteopen}set\ A\ {\isasymsubseteq}\ set\ A{\isacharprime}{\kern0pt}{\isachardoublequoteclose}\isanewline
\ \ \ \ \ \ \isacommand{using}\isamarkupfalse%
\ assms\ Cons{\isachardot}{\kern0pt}IH\ \isacommand{by}\isamarkupfalse%
\ auto\isanewline
\ \ \ \ \isacommand{then}\isamarkupfalse%
\ \isacommand{obtain}\isamarkupfalse%
\ A{\isacharprime}{\kern0pt}{\isacharprime}{\kern0pt}\ \isakeyword{where}\ {\isachardoublequoteopen}{\isacharparenleft}{\kern0pt}t\isactrlsub a{\isacharcomma}{\kern0pt}A{\isacharprime}{\kern0pt}{\isacharprime}{\kern0pt}{\isacharparenright}{\kern0pt}\ {\isasymin}\ set\ {\isacharparenleft}{\kern0pt}insort{\isacharunderscore}{\kern0pt}happ\ h\ {\isacharparenleft}{\kern0pt}insort{\isacharunderscore}{\kern0pt}mult{\isacharunderscore}{\kern0pt}happs\ hs\isactrlsub {\isadigit{1}}\ hs\isactrlsub {\isadigit{2}}{\isacharparenright}{\kern0pt}{\isacharparenright}{\kern0pt}{\isachardoublequoteclose}\ \isanewline
\ \ \ \ \ \ \ \ \ \ \ \ \ \ \ \ \ \ \ \ \ \ \ \isakeyword{and}\ {\isachardoublequoteopen}set\ A{\isacharprime}{\kern0pt}\ {\isasymsubseteq}\ set\ A{\isacharprime}{\kern0pt}{\isacharprime}{\kern0pt}{\isachardoublequoteclose}\isanewline
\ \ \ \ \ \ \isacommand{using}\isamarkupfalse%
\ insort{\isacharunderscore}{\kern0pt}happ{\isacharunderscore}{\kern0pt}action{\isacharunderscore}{\kern0pt}subset{\isacharbrackleft}{\kern0pt}\isakeyword{where}\ h{\isacharequal}{\kern0pt}{\isachardoublequoteopen}h{\isachardoublequoteclose}\ \isakeyword{and}\ hs{\isacharequal}{\kern0pt}{\isachardoublequoteopen}insort{\isacharunderscore}{\kern0pt}mult{\isacharunderscore}{\kern0pt}happs\ hs\isactrlsub {\isadigit{1}}\ hs\isactrlsub {\isadigit{2}}{\isachardoublequoteclose}{\isacharbrackright}{\kern0pt}\ \isacommand{by}\isamarkupfalse%
\ auto\isanewline
\ \ \ \ \isacommand{then}\isamarkupfalse%
\ \isacommand{show}\isamarkupfalse%
\ {\isacharquery}{\kern0pt}case\isanewline
\ \ \ \ \ \ \isacommand{using}\isamarkupfalse%
\ {\isacartoucheopen}set\ A\ {\isasymsubseteq}\ set\ A{\isacharprime}{\kern0pt}{\isacartoucheclose}\ \isacommand{by}\isamarkupfalse%
\ auto\isanewline
\ \ \isacommand{qed}\isamarkupfalse%
%
\endisatagproof
{\isafoldproof}%
%
\isadelimproof
\isanewline
%
\endisadelimproof
\isanewline
\ \ \isacommand{lemma}\isamarkupfalse%
\ insort{\isacharunderscore}{\kern0pt}mult{\isacharunderscore}{\kern0pt}happs{\isacharunderscore}{\kern0pt}action{\isacharunderscore}{\kern0pt}subset{\isacharcolon}{\kern0pt}\ \isanewline
\ \ \ \ {\isachardoublequoteopen}{\isasymforall}{\isacharparenleft}{\kern0pt}t\isactrlsub a{\isacharcomma}{\kern0pt}A{\isacharparenright}{\kern0pt}\ {\isasymin}\ set\ hs\isactrlsub {\isadigit{1}}\ {\isasymunion}\ set\ hs\isactrlsub {\isadigit{2}}{\isachardot}{\kern0pt}\ {\isasymexists}A{\isacharprime}{\kern0pt}{\isachardot}{\kern0pt}\ {\isacharparenleft}{\kern0pt}t\isactrlsub a{\isacharcomma}{\kern0pt}A{\isacharprime}{\kern0pt}{\isacharparenright}{\kern0pt}\ {\isasymin}\ set\ {\isacharparenleft}{\kern0pt}insort{\isacharunderscore}{\kern0pt}mult{\isacharunderscore}{\kern0pt}happs\ hs\isactrlsub {\isadigit{1}}\ hs\isactrlsub {\isadigit{2}}{\isacharparenright}{\kern0pt}\ {\isasymand}\ set\ A\ {\isasymsubseteq}\ set\ A{\isacharprime}{\kern0pt}{\isachardoublequoteclose}\isanewline
%
\isadelimproof
\ \ \ \ %
\endisadelimproof
%
\isatagproof
\isacommand{using}\isamarkupfalse%
\ insort{\isacharunderscore}{\kern0pt}mult{\isacharunderscore}{\kern0pt}happs{\isacharunderscore}{\kern0pt}action{\isacharunderscore}{\kern0pt}subset{\isadigit{1}}\ insort{\isacharunderscore}{\kern0pt}mult{\isacharunderscore}{\kern0pt}happs{\isacharunderscore}{\kern0pt}action{\isacharunderscore}{\kern0pt}subset{\isadigit{2}}\ \isacommand{by}\isamarkupfalse%
\ fastforce%
\endisatagproof
{\isafoldproof}%
%
\isadelimproof
\isanewline
%
\endisadelimproof
\isanewline
\ \ \isacommand{lemma}\isamarkupfalse%
\ insort{\isacharunderscore}{\kern0pt}happ{\isadigit{2}}{\isacharunderscore}{\kern0pt}insort{\isacharunderscore}{\kern0pt}mult{\isacharunderscore}{\kern0pt}happs{\isacharunderscore}{\kern0pt}action{\isacharunderscore}{\kern0pt}subset{\isacharunderscore}{\kern0pt}aux{\isacharcolon}{\kern0pt}\isanewline
\ \ \ \ \isakeyword{assumes}\ {\isachardoublequoteopen}{\isacharparenleft}{\kern0pt}t\isactrlsub a{\isacharcomma}{\kern0pt}A{\isacharparenright}{\kern0pt}\ {\isasymin}\ set\ hs\isactrlsub {\isadigit{1}}{\isachardoublequoteclose}\ \isanewline
\ \ \ \ \isakeyword{shows}\ {\isachardoublequoteopen}{\isasymexists}A{\isacharprime}{\kern0pt}{\isachardot}{\kern0pt}\ {\isacharparenleft}{\kern0pt}t\isactrlsub a{\isacharcomma}{\kern0pt}A{\isacharprime}{\kern0pt}{\isacharparenright}{\kern0pt}\ {\isasymin}\ set\ {\isacharparenleft}{\kern0pt}insort{\isacharunderscore}{\kern0pt}happ\ {\isacharparenleft}{\kern0pt}t\isactrlsub {\isadigit{1}}{\isacharcomma}{\kern0pt}A\isactrlsub {\isadigit{1}}{\isacharparenright}{\kern0pt}\ {\isacharparenleft}{\kern0pt}insort{\isacharunderscore}{\kern0pt}happ\ {\isacharparenleft}{\kern0pt}t\isactrlsub {\isadigit{2}}{\isacharcomma}{\kern0pt}A\isactrlsub {\isadigit{2}}{\isacharparenright}{\kern0pt}\ {\isacharparenleft}{\kern0pt}insort{\isacharunderscore}{\kern0pt}mult{\isacharunderscore}{\kern0pt}happs\ hs\isactrlsub {\isadigit{1}}\ hs\isactrlsub {\isadigit{2}}{\isacharparenright}{\kern0pt}{\isacharparenright}{\kern0pt}{\isacharparenright}{\kern0pt}\ {\isasymand}\ set\ A\ {\isasymsubseteq}\ set\ A{\isacharprime}{\kern0pt}{\isachardoublequoteclose}\isanewline
%
\isadelimproof
\ \ %
\endisadelimproof
%
\isatagproof
\isacommand{proof}\isamarkupfalse%
\ {\isacharminus}{\kern0pt}\isanewline
\ \ \ \ \isacommand{have}\isamarkupfalse%
\ {\isachardoublequoteopen}{\isasymexists}A{\isacharprime}{\kern0pt}{\isachardot}{\kern0pt}\ {\isacharparenleft}{\kern0pt}t\isactrlsub a{\isacharcomma}{\kern0pt}A{\isacharprime}{\kern0pt}{\isacharparenright}{\kern0pt}\ {\isasymin}\ set\ {\isacharparenleft}{\kern0pt}insort{\isacharunderscore}{\kern0pt}mult{\isacharunderscore}{\kern0pt}happs\ hs\isactrlsub {\isadigit{1}}\ hs\isactrlsub {\isadigit{2}}{\isacharparenright}{\kern0pt}\ {\isasymand}\ set\ A\ {\isasymsubseteq}\ set\ A{\isacharprime}{\kern0pt}{\isachardoublequoteclose}\isanewline
\ \ \ \ \ \ \isacommand{using}\isamarkupfalse%
\ assms\ insort{\isacharunderscore}{\kern0pt}mult{\isacharunderscore}{\kern0pt}happs{\isacharunderscore}{\kern0pt}action{\isacharunderscore}{\kern0pt}subset{\isacharbrackleft}{\kern0pt}\isakeyword{where}\ hs\isactrlsub {\isadigit{1}}{\isacharequal}{\kern0pt}{\isachardoublequoteopen}hs\isactrlsub {\isadigit{1}}{\isachardoublequoteclose}\ \isakeyword{and}\ hs\isactrlsub {\isadigit{2}}{\isacharequal}{\kern0pt}{\isachardoublequoteopen}hs\isactrlsub {\isadigit{2}}{\isachardoublequoteclose}{\isacharbrackright}{\kern0pt}\isanewline
\ \ \ \ \ \ \isacommand{by}\isamarkupfalse%
\ {\isacharparenleft}{\kern0pt}auto\ split{\isacharcolon}{\kern0pt}\ prod{\isachardot}{\kern0pt}splits{\isacharparenright}{\kern0pt}\isanewline
\ \ \ \ \isacommand{then}\isamarkupfalse%
\ \isacommand{obtain}\isamarkupfalse%
\ A{\isacharprime}{\kern0pt}\ \isakeyword{where}\ {\isachardoublequoteopen}{\isacharparenleft}{\kern0pt}t\isactrlsub a{\isacharcomma}{\kern0pt}A{\isacharprime}{\kern0pt}{\isacharparenright}{\kern0pt}\ {\isasymin}\ set\ {\isacharparenleft}{\kern0pt}insort{\isacharunderscore}{\kern0pt}mult{\isacharunderscore}{\kern0pt}happs\ hs\isactrlsub {\isadigit{1}}\ hs\isactrlsub {\isadigit{2}}{\isacharparenright}{\kern0pt}{\isachardoublequoteclose}\ \isakeyword{and}\ {\isachardoublequoteopen}set\ A\ {\isasymsubseteq}\ set\ A{\isacharprime}{\kern0pt}{\isachardoublequoteclose}\isanewline
\ \ \ \ \ \ \isacommand{by}\isamarkupfalse%
\ auto\isanewline
\ \ \ \ \isacommand{then}\isamarkupfalse%
\ \isacommand{obtain}\isamarkupfalse%
\ A{\isacharprime}{\kern0pt}{\isacharprime}{\kern0pt}\ \isakeyword{where}\ {\isachardoublequoteopen}{\isacharparenleft}{\kern0pt}t\isactrlsub a{\isacharcomma}{\kern0pt}A{\isacharprime}{\kern0pt}{\isacharprime}{\kern0pt}{\isacharparenright}{\kern0pt}\ {\isasymin}\ set\ {\isacharparenleft}{\kern0pt}insort{\isacharunderscore}{\kern0pt}happ\ {\isacharparenleft}{\kern0pt}t\isactrlsub {\isadigit{2}}{\isacharcomma}{\kern0pt}A\isactrlsub {\isadigit{2}}{\isacharparenright}{\kern0pt}\ {\isacharparenleft}{\kern0pt}insort{\isacharunderscore}{\kern0pt}mult{\isacharunderscore}{\kern0pt}happs\ hs\isactrlsub {\isadigit{1}}\ hs\isactrlsub {\isadigit{2}}{\isacharparenright}{\kern0pt}{\isacharparenright}{\kern0pt}{\isachardoublequoteclose}\ \isanewline
\ \ \ \ \ \ \ \ \ \ \ \ \ \ \ \ \ \ \ \ \ \ \ \isakeyword{and}\ {\isachardoublequoteopen}set\ A{\isacharprime}{\kern0pt}\ {\isasymsubseteq}\ set\ A{\isacharprime}{\kern0pt}{\isacharprime}{\kern0pt}{\isachardoublequoteclose}\isanewline
\ \ \ \ \ \ \isacommand{using}\isamarkupfalse%
\ insort{\isacharunderscore}{\kern0pt}happ{\isacharunderscore}{\kern0pt}action{\isacharunderscore}{\kern0pt}subset{\isacharbrackleft}{\kern0pt}\isakeyword{where}\ h{\isacharequal}{\kern0pt}{\isachardoublequoteopen}{\isacharparenleft}{\kern0pt}t\isactrlsub {\isadigit{2}}{\isacharcomma}{\kern0pt}A\isactrlsub {\isadigit{2}}{\isacharparenright}{\kern0pt}{\isachardoublequoteclose}\ \isakeyword{and}\ hs{\isacharequal}{\kern0pt}{\isachardoublequoteopen}insort{\isacharunderscore}{\kern0pt}mult{\isacharunderscore}{\kern0pt}happs\ hs\isactrlsub {\isadigit{1}}\ hs\isactrlsub {\isadigit{2}}{\isachardoublequoteclose}{\isacharbrackright}{\kern0pt}\ \isacommand{by}\isamarkupfalse%
\ auto\isanewline
\ \ \ \ \isacommand{then}\isamarkupfalse%
\ \isacommand{obtain}\isamarkupfalse%
\ A{\isacharprime}{\kern0pt}{\isacharprime}{\kern0pt}{\isacharprime}{\kern0pt}\ \isakeyword{where}\ {\isachardoublequoteopen}{\isacharparenleft}{\kern0pt}t\isactrlsub a{\isacharcomma}{\kern0pt}A{\isacharprime}{\kern0pt}{\isacharprime}{\kern0pt}{\isacharprime}{\kern0pt}{\isacharparenright}{\kern0pt}\ {\isasymin}\ set\ {\isacharparenleft}{\kern0pt}insort{\isacharunderscore}{\kern0pt}happ\ {\isacharparenleft}{\kern0pt}t\isactrlsub {\isadigit{1}}{\isacharcomma}{\kern0pt}A\isactrlsub {\isadigit{1}}{\isacharparenright}{\kern0pt}\ {\isacharparenleft}{\kern0pt}insort{\isacharunderscore}{\kern0pt}happ\ {\isacharparenleft}{\kern0pt}t\isactrlsub {\isadigit{2}}{\isacharcomma}{\kern0pt}A\isactrlsub {\isadigit{2}}{\isacharparenright}{\kern0pt}\ {\isacharparenleft}{\kern0pt}insort{\isacharunderscore}{\kern0pt}mult{\isacharunderscore}{\kern0pt}happs\ hs\isactrlsub {\isadigit{1}}\ hs\isactrlsub {\isadigit{2}}{\isacharparenright}{\kern0pt}{\isacharparenright}{\kern0pt}{\isacharparenright}{\kern0pt}{\isachardoublequoteclose}\ \isanewline
\ \ \ \ \ \ \ \ \ \ \ \ \ \ \ \ \ \ \ \ \ \ \ \ \isakeyword{and}\ {\isachardoublequoteopen}set\ A{\isacharprime}{\kern0pt}{\isacharprime}{\kern0pt}\ {\isasymsubseteq}\ set\ A{\isacharprime}{\kern0pt}{\isacharprime}{\kern0pt}{\isacharprime}{\kern0pt}{\isachardoublequoteclose}\isanewline
\ \ \ \ \ \ \isacommand{using}\isamarkupfalse%
\ insort{\isacharunderscore}{\kern0pt}happ{\isacharunderscore}{\kern0pt}action{\isacharunderscore}{\kern0pt}subset{\isacharbrackleft}{\kern0pt}\isakeyword{where}\ h{\isacharequal}{\kern0pt}{\isachardoublequoteopen}{\isacharparenleft}{\kern0pt}t\isactrlsub {\isadigit{1}}{\isacharcomma}{\kern0pt}A\isactrlsub {\isadigit{1}}{\isacharparenright}{\kern0pt}{\isachardoublequoteclose}\ \isanewline
\ \ \ \ \ \ \ \ \ \ \ \ \ \ \ \ \ \ \ \ \ \ \ \ \ \ \ \ \ \ \ \ \ \ \ \ \ \ \ \isakeyword{and}\ hs{\isacharequal}{\kern0pt}{\isachardoublequoteopen}insort{\isacharunderscore}{\kern0pt}happ\ {\isacharparenleft}{\kern0pt}t\isactrlsub {\isadigit{2}}{\isacharcomma}{\kern0pt}A\isactrlsub {\isadigit{2}}{\isacharparenright}{\kern0pt}\ {\isacharparenleft}{\kern0pt}insort{\isacharunderscore}{\kern0pt}mult{\isacharunderscore}{\kern0pt}happs\ hs\isactrlsub {\isadigit{1}}\ hs\isactrlsub {\isadigit{2}}{\isacharparenright}{\kern0pt}{\isachardoublequoteclose}{\isacharbrackright}{\kern0pt}\isanewline
\ \ \ \ \ \ \isacommand{by}\isamarkupfalse%
\ auto\isanewline
\ \ \ \ \isacommand{show}\isamarkupfalse%
\ {\isacharquery}{\kern0pt}thesis\isanewline
\ \ \ \ \ \ \isacommand{using}\isamarkupfalse%
\ {\isacartoucheopen}{\isacharparenleft}{\kern0pt}t\isactrlsub a{\isacharcomma}{\kern0pt}A{\isacharparenright}{\kern0pt}\ {\isasymin}\ set\ hs\isactrlsub {\isadigit{1}}{\isacartoucheclose}\ \isanewline
\ \ \ \ \ \ \ \ \ \ \ \ {\isacartoucheopen}{\isacharparenleft}{\kern0pt}t\isactrlsub a{\isacharcomma}{\kern0pt}A{\isacharprime}{\kern0pt}{\isacharprime}{\kern0pt}{\isacharprime}{\kern0pt}{\isacharparenright}{\kern0pt}\ {\isasymin}\ set\ {\isacharparenleft}{\kern0pt}insort{\isacharunderscore}{\kern0pt}happ\ {\isacharparenleft}{\kern0pt}t\isactrlsub {\isadigit{1}}{\isacharcomma}{\kern0pt}A\isactrlsub {\isadigit{1}}{\isacharparenright}{\kern0pt}\ {\isacharparenleft}{\kern0pt}insort{\isacharunderscore}{\kern0pt}happ\ {\isacharparenleft}{\kern0pt}t\isactrlsub {\isadigit{2}}{\isacharcomma}{\kern0pt}A\isactrlsub {\isadigit{2}}{\isacharparenright}{\kern0pt}\ {\isacharparenleft}{\kern0pt}insort{\isacharunderscore}{\kern0pt}mult{\isacharunderscore}{\kern0pt}happs\ hs\isactrlsub {\isadigit{1}}\ hs\isactrlsub {\isadigit{2}}{\isacharparenright}{\kern0pt}{\isacharparenright}{\kern0pt}{\isacharparenright}{\kern0pt}{\isacartoucheclose}\ \isanewline
\ \ \ \ \ \ \ \ \ \ \ \ {\isacartoucheopen}set\ A\ {\isasymsubseteq}\ set\ A{\isacharprime}{\kern0pt}{\isacartoucheclose}\ {\isacartoucheopen}set\ A{\isacharprime}{\kern0pt}\ {\isasymsubseteq}\ set\ A{\isacharprime}{\kern0pt}{\isacharprime}{\kern0pt}{\isacartoucheclose}\ {\isacartoucheopen}set\ A{\isacharprime}{\kern0pt}{\isacharprime}{\kern0pt}\ {\isasymsubseteq}\ set\ A{\isacharprime}{\kern0pt}{\isacharprime}{\kern0pt}{\isacharprime}{\kern0pt}{\isacartoucheclose}\ \isacommand{by}\isamarkupfalse%
\ auto\isanewline
\ \ \isacommand{qed}\isamarkupfalse%
%
\endisatagproof
{\isafoldproof}%
%
\isadelimproof
\ \ \ \ \ \ \ \ \ \ \ \ \ \ \ \ \ \ \ \ \ \ \ \ \ \ \ \ \ \ \ \ \ \ \ \ \ \ \ \ \ \ \ \ \ \ \ \ \ \ \ \ \ \ \ \ \ \ \ \ \ \ \ \ \ \ \ \ \ \ \ \ \ \ \ \ \ \ \ \ \ \ \ \ \ \ \ \ \ \ \ \ \ \ \ \ \ \ \ \ \ \ \ \ \ \ \ \ \ \ \ \ \ \ \ \ \ \ \ \isanewline
%
\endisadelimproof
\ \ \isanewline
\isanewline
\ \ \isacommand{lemma}\isamarkupfalse%
\ insort{\isacharunderscore}{\kern0pt}happ{\isadigit{2}}{\isacharunderscore}{\kern0pt}insort{\isacharunderscore}{\kern0pt}mult{\isacharunderscore}{\kern0pt}happs{\isacharunderscore}{\kern0pt}action{\isacharunderscore}{\kern0pt}subset{\isacharcolon}{\kern0pt}\isanewline
\ \ \ \ {\isachardoublequoteopen}{\isasymforall}{\isacharparenleft}{\kern0pt}t\isactrlsub a{\isacharcomma}{\kern0pt}A{\isacharparenright}{\kern0pt}\ {\isasymin}\ set\ hs\isactrlsub {\isadigit{1}}{\isachardot}{\kern0pt}\ {\isasymexists}A{\isacharprime}{\kern0pt}{\isachardot}{\kern0pt}\ {\isacharparenleft}{\kern0pt}t\isactrlsub a{\isacharcomma}{\kern0pt}A{\isacharprime}{\kern0pt}{\isacharparenright}{\kern0pt}\ {\isasymin}\ set\ {\isacharparenleft}{\kern0pt}insort{\isacharunderscore}{\kern0pt}happ\ {\isacharparenleft}{\kern0pt}t\isactrlsub {\isadigit{1}}{\isacharcomma}{\kern0pt}A\isactrlsub {\isadigit{1}}{\isacharparenright}{\kern0pt}\ {\isacharparenleft}{\kern0pt}insort{\isacharunderscore}{\kern0pt}happ\ {\isacharparenleft}{\kern0pt}t\isactrlsub {\isadigit{2}}{\isacharcomma}{\kern0pt}A\isactrlsub {\isadigit{2}}{\isacharparenright}{\kern0pt}\ {\isacharparenleft}{\kern0pt}insort{\isacharunderscore}{\kern0pt}mult{\isacharunderscore}{\kern0pt}happs\ hs\isactrlsub {\isadigit{1}}\ hs\isactrlsub {\isadigit{2}}{\isacharparenright}{\kern0pt}{\isacharparenright}{\kern0pt}{\isacharparenright}{\kern0pt}\ {\isasymand}\ set\ A\ {\isasymsubseteq}\ set\ A{\isacharprime}{\kern0pt}{\isachardoublequoteclose}\isanewline
%
\isadelimproof
\ \ \ \ %
\endisadelimproof
%
\isatagproof
\isacommand{using}\isamarkupfalse%
\ insort{\isacharunderscore}{\kern0pt}happ{\isadigit{2}}{\isacharunderscore}{\kern0pt}insort{\isacharunderscore}{\kern0pt}mult{\isacharunderscore}{\kern0pt}happs{\isacharunderscore}{\kern0pt}action{\isacharunderscore}{\kern0pt}subset{\isacharunderscore}{\kern0pt}aux\ \isacommand{by}\isamarkupfalse%
\ auto%
\endisatagproof
{\isafoldproof}%
%
\isadelimproof
%
\endisadelimproof
%
\begin{isamarkuptext}%
We can define a function to construct a simplified plan for a given temporal plan.%
\end{isamarkuptext}\isamarkuptrue%
\ \ \isacommand{fun}\isamarkupfalse%
\ simplify{\isacharunderscore}{\kern0pt}plan\ {\isacharcolon}{\kern0pt}{\isacharcolon}{\kern0pt}\ {\isachardoublequoteopen}time\ list\ {\isasymRightarrow}\ plan\ {\isasymRightarrow}\ happening\ list{\isachardoublequoteclose}\ \isakeyword{where}\isanewline
\ \ \ \ {\isachardoublequoteopen}simplify{\isacharunderscore}{\kern0pt}plan\ htps\ {\isacharbrackleft}{\kern0pt}{\isacharbrackright}{\kern0pt}\ {\isacharequal}{\kern0pt}\ {\isacharbrackleft}{\kern0pt}{\isacharbrackright}{\kern0pt}{\isachardoublequoteclose}\isanewline
\ \ {\isacharbar}{\kern0pt}\ {\isachardoublequoteopen}simplify{\isacharunderscore}{\kern0pt}plan\ htps\ {\isacharparenleft}{\kern0pt}{\isasympi}{\isacharhash}{\kern0pt}{\isasympi}s{\isacharparenright}{\kern0pt}\ {\isacharequal}{\kern0pt}\ \isanewline
\ \ \ \ \ \ insort{\isacharunderscore}{\kern0pt}mult{\isacharunderscore}{\kern0pt}happs\ {\isacharparenleft}{\kern0pt}map\ {\isacharparenleft}{\kern0pt}{\isasymlambda}{\isacharparenleft}{\kern0pt}t\isactrlsub a{\isacharcomma}{\kern0pt}a{\isacharparenright}{\kern0pt}{\isachardot}{\kern0pt}\ {\isacharparenleft}{\kern0pt}t\isactrlsub a{\isacharcomma}{\kern0pt}{\isacharbrackleft}{\kern0pt}a{\isacharbrackright}{\kern0pt}{\isacharparenright}{\kern0pt}{\isacharparenright}{\kern0pt}\ {\isacharparenleft}{\kern0pt}simplify{\isacharunderscore}{\kern0pt}action\ htps\ {\isasympi}{\isacharparenright}{\kern0pt}{\isacharparenright}{\kern0pt}\ {\isacharparenleft}{\kern0pt}simplify{\isacharunderscore}{\kern0pt}plan\ htps\ {\isasympi}s{\isacharparenright}{\kern0pt}{\isachardoublequoteclose}%
\begin{isamarkuptext}%
Definition of subset on happening sequences. This is used to simplify some proofs.%
\end{isamarkuptext}\isamarkuptrue%
\ \ \isacommand{definition}\isamarkupfalse%
\ happ{\isacharunderscore}{\kern0pt}seq{\isacharunderscore}{\kern0pt}member\ {\isacharcolon}{\kern0pt}{\isacharcolon}{\kern0pt}\ {\isachardoublequoteopen}happening\ {\isasymRightarrow}\ happening\ list\ {\isasymRightarrow}\ bool{\isachardoublequoteclose}\ {\isacharparenleft}{\kern0pt}\isakeyword{infix}\ {\isachardoublequoteopen}{\isasymin}\isactrlsub h\isactrlsub s{\isachardoublequoteclose}\ {\isadigit{5}}{\isadigit{3}}{\isacharparenright}{\kern0pt}\ \isakeyword{where}\isanewline
\ \ \ \ {\isachardoublequoteopen}happ{\isacharunderscore}{\kern0pt}seq{\isacharunderscore}{\kern0pt}member\ h\ hs\ {\isasymlongleftrightarrow}\ {\isacharparenleft}{\kern0pt}case\ h\ of\ {\isacharparenleft}{\kern0pt}t\isactrlsub a{\isacharcomma}{\kern0pt}A{\isacharparenright}{\kern0pt}\ {\isasymRightarrow}\ {\isasymexists}A\isactrlsub {\isadigit{2}}{\isachardot}{\kern0pt}\ set\ A\ {\isasymsubseteq}\ set\ A\isactrlsub {\isadigit{2}}\ {\isasymand}\ {\isacharparenleft}{\kern0pt}t\isactrlsub a{\isacharcomma}{\kern0pt}A\isactrlsub {\isadigit{2}}{\isacharparenright}{\kern0pt}\ {\isasymin}\ set\ hs{\isacharparenright}{\kern0pt}{\isachardoublequoteclose}\isanewline
\isanewline
\ \ \isacommand{definition}\isamarkupfalse%
\ happ{\isacharunderscore}{\kern0pt}seq{\isacharunderscore}{\kern0pt}subset\ {\isacharcolon}{\kern0pt}{\isacharcolon}{\kern0pt}\ {\isachardoublequoteopen}happening\ list\ {\isasymRightarrow}\ happening\ list\ {\isasymRightarrow}\ bool{\isachardoublequoteclose}\ {\isacharparenleft}{\kern0pt}\isakeyword{infix}\ {\isachardoublequoteopen}{\isasymsubseteq}\isactrlsub h\isactrlsub s{\isachardoublequoteclose}\ {\isadigit{5}}{\isadigit{3}}{\isacharparenright}{\kern0pt}\ \isakeyword{where}\isanewline
\ \ \ \ {\isachardoublequoteopen}happ{\isacharunderscore}{\kern0pt}seq{\isacharunderscore}{\kern0pt}subset\ hs\isactrlsub {\isadigit{1}}\ hs\isactrlsub {\isadigit{2}}\ {\isasymlongleftrightarrow}\ {\isacharparenleft}{\kern0pt}{\isasymforall}h\isactrlsub {\isadigit{1}}\ {\isasymin}\ set\ hs\isactrlsub {\isadigit{1}}{\isachardot}{\kern0pt}\ h\isactrlsub {\isadigit{1}}\ {\isasymin}\isactrlsub h\isactrlsub s\ hs\isactrlsub {\isadigit{2}}{\isacharparenright}{\kern0pt}{\isachardoublequoteclose}%
\begin{isamarkuptext}%
Properties of \isa{{\isacharparenleft}{\kern0pt}{\isasymsubseteq}\isactrlsub h\isactrlsub s{\isacharparenright}{\kern0pt}}%
\end{isamarkuptext}\isamarkuptrue%
\ \ \isacommand{lemma}\isamarkupfalse%
\ happ{\isacharunderscore}{\kern0pt}seq{\isacharunderscore}{\kern0pt}subset{\isacharunderscore}{\kern0pt}Nil{\isacharcolon}{\kern0pt}\ {\isachardoublequoteopen}{\isacharbrackleft}{\kern0pt}{\isacharbrackright}{\kern0pt}\ {\isasymsubseteq}\isactrlsub h\isactrlsub s\ hs{\isachardoublequoteclose}\isanewline
%
\isadelimproof
\ \ \ \ %
\endisadelimproof
%
\isatagproof
\isacommand{unfolding}\isamarkupfalse%
\ happ{\isacharunderscore}{\kern0pt}seq{\isacharunderscore}{\kern0pt}subset{\isacharunderscore}{\kern0pt}def\ \isacommand{by}\isamarkupfalse%
\ auto%
\endisatagproof
{\isafoldproof}%
%
\isadelimproof
\isanewline
%
\endisadelimproof
\isanewline
\ \ \isacommand{lemma}\isamarkupfalse%
\ happ{\isacharunderscore}{\kern0pt}seq{\isacharunderscore}{\kern0pt}subset{\isacharunderscore}{\kern0pt}refl{\isacharcolon}{\kern0pt}\ {\isachardoublequoteopen}hs\ {\isasymsubseteq}\isactrlsub h\isactrlsub s\ hs{\isachardoublequoteclose}\isanewline
%
\isadelimproof
\ \ \ \ %
\endisadelimproof
%
\isatagproof
\isacommand{unfolding}\isamarkupfalse%
\ happ{\isacharunderscore}{\kern0pt}seq{\isacharunderscore}{\kern0pt}subset{\isacharunderscore}{\kern0pt}def\ happ{\isacharunderscore}{\kern0pt}seq{\isacharunderscore}{\kern0pt}member{\isacharunderscore}{\kern0pt}def\ \isacommand{by}\isamarkupfalse%
\ auto%
\endisatagproof
{\isafoldproof}%
%
\isadelimproof
\isanewline
%
\endisadelimproof
\isanewline
\ \ \isacommand{lemma}\isamarkupfalse%
\ happ{\isacharunderscore}{\kern0pt}seq{\isacharunderscore}{\kern0pt}subset{\isacharunderscore}{\kern0pt}trans{\isacharcolon}{\kern0pt}\ {\isachardoublequoteopen}hs\isactrlsub {\isadigit{1}}\ {\isasymsubseteq}\isactrlsub h\isactrlsub s\ hs\isactrlsub {\isadigit{2}}\ {\isasymand}\ hs\isactrlsub {\isadigit{2}}\ {\isasymsubseteq}\isactrlsub h\isactrlsub s\ hs\isactrlsub {\isadigit{3}}\ {\isasymLongrightarrow}\ hs\isactrlsub {\isadigit{1}}\ {\isasymsubseteq}\isactrlsub h\isactrlsub s\ hs\isactrlsub {\isadigit{3}}{\isachardoublequoteclose}\isanewline
%
\isadelimproof
\ \ \ \ %
\endisadelimproof
%
\isatagproof
\isacommand{unfolding}\isamarkupfalse%
\ happ{\isacharunderscore}{\kern0pt}seq{\isacharunderscore}{\kern0pt}subset{\isacharunderscore}{\kern0pt}def\ happ{\isacharunderscore}{\kern0pt}seq{\isacharunderscore}{\kern0pt}member{\isacharunderscore}{\kern0pt}def\isanewline
\ \ \ \ \isacommand{by}\isamarkupfalse%
\ {\isacharparenleft}{\kern0pt}auto\ split{\isacharcolon}{\kern0pt}\ prod{\isachardot}{\kern0pt}splits{\isacharparenright}{\kern0pt}\ {\isacharparenleft}{\kern0pt}metis\ {\isacharparenleft}{\kern0pt}full{\isacharunderscore}{\kern0pt}types{\isacharparenright}{\kern0pt}\ dual{\isacharunderscore}{\kern0pt}order{\isachardot}{\kern0pt}trans{\isacharparenright}{\kern0pt}%
\endisatagproof
{\isafoldproof}%
%
\isadelimproof
\isanewline
%
\endisadelimproof
\isanewline
\ \ \isacommand{lemma}\isamarkupfalse%
\ insort{\isacharunderscore}{\kern0pt}happ{\isacharunderscore}{\kern0pt}subset{\isacharcolon}{\kern0pt}\ {\isachardoublequoteopen}hs\ {\isasymsubseteq}\isactrlsub h\isactrlsub s\ insort{\isacharunderscore}{\kern0pt}happ\ h\ hs{\isachardoublequoteclose}\isanewline
%
\isadelimproof
\ \ %
\endisadelimproof
%
\isatagproof
\isacommand{proof}\isamarkupfalse%
\ {\isacharminus}{\kern0pt}\isanewline
\ \ \ \ \isacommand{have}\isamarkupfalse%
\ {\isachardoublequoteopen}hs\ {\isacharequal}{\kern0pt}\ {\isacharbrackleft}{\kern0pt}{\isacharbrackright}{\kern0pt}\ {\isasymor}\ hs\ {\isasymnoteq}\ {\isacharbrackleft}{\kern0pt}{\isacharbrackright}{\kern0pt}{\isachardoublequoteclose}\isanewline
\ \ \ \ \ \ \isacommand{by}\isamarkupfalse%
\ simp\isanewline
\ \ \ \ \isacommand{then}\isamarkupfalse%
\ \isacommand{show}\isamarkupfalse%
\ {\isacharquery}{\kern0pt}thesis\isanewline
\ \ \ \ \isacommand{proof}\isamarkupfalse%
\isanewline
\ \ \ \ \ \ \isacommand{assume}\isamarkupfalse%
\ {\isachardoublequoteopen}hs\ {\isacharequal}{\kern0pt}\ {\isacharbrackleft}{\kern0pt}{\isacharbrackright}{\kern0pt}{\isachardoublequoteclose}\isanewline
\ \ \ \ \ \ \isacommand{then}\isamarkupfalse%
\ \isacommand{show}\isamarkupfalse%
\ {\isacharquery}{\kern0pt}thesis\isanewline
\ \ \ \ \ \ \ \ \isacommand{unfolding}\isamarkupfalse%
\ happ{\isacharunderscore}{\kern0pt}seq{\isacharunderscore}{\kern0pt}subset{\isacharunderscore}{\kern0pt}def\ \isacommand{by}\isamarkupfalse%
\ auto\isanewline
\ \ \ \ \isacommand{next}\isamarkupfalse%
\isanewline
\ \ \ \ \ \ \isacommand{assume}\isamarkupfalse%
\ {\isachardoublequoteopen}hs\ {\isasymnoteq}\ {\isacharbrackleft}{\kern0pt}{\isacharbrackright}{\kern0pt}{\isachardoublequoteclose}\isanewline
\ \ \ \ \ \ \isacommand{then}\isamarkupfalse%
\ \isacommand{show}\isamarkupfalse%
\ {\isacharquery}{\kern0pt}thesis\isanewline
\ \ \ \ \ \ \ \ \isacommand{unfolding}\isamarkupfalse%
\ happ{\isacharunderscore}{\kern0pt}seq{\isacharunderscore}{\kern0pt}subset{\isacharunderscore}{\kern0pt}def\ happ{\isacharunderscore}{\kern0pt}seq{\isacharunderscore}{\kern0pt}member{\isacharunderscore}{\kern0pt}def\isanewline
\ \ \ \ \ \ \ \ \isacommand{using}\isamarkupfalse%
\ insort{\isacharunderscore}{\kern0pt}happ{\isacharunderscore}{\kern0pt}action{\isacharunderscore}{\kern0pt}subset\ \isacommand{by}\isamarkupfalse%
\ blast\isanewline
\ \ \ \ \isacommand{qed}\isamarkupfalse%
\isanewline
\ \ \isacommand{qed}\isamarkupfalse%
%
\endisatagproof
{\isafoldproof}%
%
\isadelimproof
\isanewline
%
\endisadelimproof
\isanewline
\ \ \isacommand{lemma}\isamarkupfalse%
\ insort{\isacharunderscore}{\kern0pt}mult{\isacharunderscore}{\kern0pt}happs{\isacharunderscore}{\kern0pt}subset{\isacharcolon}{\kern0pt}\ {\isachardoublequoteopen}hs\isactrlsub {\isadigit{1}}\ {\isasymsubseteq}\isactrlsub h\isactrlsub s\ insort{\isacharunderscore}{\kern0pt}mult{\isacharunderscore}{\kern0pt}happs\ hs\isactrlsub {\isadigit{2}}\ hs\isactrlsub {\isadigit{1}}{\isachardoublequoteclose}\isanewline
%
\isadelimproof
\ \ %
\endisadelimproof
%
\isatagproof
\isacommand{proof}\isamarkupfalse%
\ {\isacharparenleft}{\kern0pt}induction\ hs\isactrlsub {\isadigit{2}}\ hs\isactrlsub {\isadigit{1}}\ rule{\isacharcolon}{\kern0pt}\ insort{\isacharunderscore}{\kern0pt}mult{\isacharunderscore}{\kern0pt}happs{\isachardot}{\kern0pt}induct{\isacharparenright}{\kern0pt}\isanewline
\ \ \ \ \isacommand{case}\isamarkupfalse%
\ {\isacharparenleft}{\kern0pt}{\isadigit{1}}\ hs{\isacharparenright}{\kern0pt}\isanewline
\ \ \ \ \isacommand{then}\isamarkupfalse%
\ \isacommand{show}\isamarkupfalse%
\ {\isacharquery}{\kern0pt}case\isanewline
\ \ \ \ \ \ \isacommand{using}\isamarkupfalse%
\ happ{\isacharunderscore}{\kern0pt}seq{\isacharunderscore}{\kern0pt}subset{\isacharunderscore}{\kern0pt}refl\ \isacommand{by}\isamarkupfalse%
\ simp\isanewline
\ \ \isacommand{next}\isamarkupfalse%
\isanewline
\ \ \ \ \isacommand{case}\isamarkupfalse%
\ {\isacharparenleft}{\kern0pt}{\isadigit{2}}\ a\ A\ hs{\isacharparenright}{\kern0pt}\isanewline
\ \ \ \ \isacommand{have}\isamarkupfalse%
\ {\isachardoublequoteopen}insort{\isacharunderscore}{\kern0pt}mult{\isacharunderscore}{\kern0pt}happs\ {\isacharparenleft}{\kern0pt}a\ {\isacharhash}{\kern0pt}\ A{\isacharparenright}{\kern0pt}\ hs\ {\isacharequal}{\kern0pt}\ insort{\isacharunderscore}{\kern0pt}happ\ a\ {\isacharparenleft}{\kern0pt}insort{\isacharunderscore}{\kern0pt}mult{\isacharunderscore}{\kern0pt}happs\ A\ hs{\isacharparenright}{\kern0pt}{\isachardoublequoteclose}\isanewline
\ \ \ \ \ \ \isacommand{by}\isamarkupfalse%
\ simp\isanewline
\ \ \ \ \isacommand{then}\isamarkupfalse%
\ \isacommand{have}\isamarkupfalse%
\ {\isachardoublequoteopen}insort{\isacharunderscore}{\kern0pt}mult{\isacharunderscore}{\kern0pt}happs\ A\ hs\ {\isasymsubseteq}\isactrlsub h\isactrlsub s\ insort{\isacharunderscore}{\kern0pt}mult{\isacharunderscore}{\kern0pt}happs\ {\isacharparenleft}{\kern0pt}a\ {\isacharhash}{\kern0pt}\ A{\isacharparenright}{\kern0pt}\ hs{\isachardoublequoteclose}\isanewline
\ \ \ \ \ \ \isacommand{using}\isamarkupfalse%
\ insort{\isacharunderscore}{\kern0pt}happ{\isacharunderscore}{\kern0pt}subset\ \isacommand{by}\isamarkupfalse%
\ auto\isanewline
\ \ \ \ \isacommand{then}\isamarkupfalse%
\ \isacommand{show}\isamarkupfalse%
\ {\isacharquery}{\kern0pt}case\isanewline
\ \ \ \ \ \ \isacommand{using}\isamarkupfalse%
\ {\isachardoublequoteopen}{\isadigit{2}}{\isachardot}{\kern0pt}IH{\isachardoublequoteclose}\ happ{\isacharunderscore}{\kern0pt}seq{\isacharunderscore}{\kern0pt}subset{\isacharunderscore}{\kern0pt}trans\ \isacommand{by}\isamarkupfalse%
\ blast\isanewline
\ \ \isacommand{qed}\isamarkupfalse%
%
\endisatagproof
{\isafoldproof}%
%
\isadelimproof
%
\endisadelimproof
%
\begin{isamarkuptext}%
Happening sequence subset for simplify plan.%
\end{isamarkuptext}\isamarkuptrue%
\ \ \isacommand{lemma}\isamarkupfalse%
\ simplify{\isacharunderscore}{\kern0pt}plan{\isacharunderscore}{\kern0pt}happs{\isacharunderscore}{\kern0pt}subset{\isacharunderscore}{\kern0pt}Cons{\isacharcolon}{\kern0pt}\ {\isachardoublequoteopen}simplify{\isacharunderscore}{\kern0pt}plan\ htps\ {\isasympi}s\ {\isasymsubseteq}\isactrlsub h\isactrlsub s\ simplify{\isacharunderscore}{\kern0pt}plan\ htps\ {\isacharparenleft}{\kern0pt}{\isasympi}{\isacharhash}{\kern0pt}{\isasympi}s{\isacharparenright}{\kern0pt}{\isachardoublequoteclose}\isanewline
%
\isadelimproof
\ \ \ \ %
\endisadelimproof
%
\isatagproof
\isacommand{using}\isamarkupfalse%
\ insort{\isacharunderscore}{\kern0pt}mult{\isacharunderscore}{\kern0pt}happs{\isacharunderscore}{\kern0pt}subset\ happ{\isacharunderscore}{\kern0pt}seq{\isacharunderscore}{\kern0pt}subset{\isacharunderscore}{\kern0pt}refl\isanewline
\ \ \ \ \isacommand{by}\isamarkupfalse%
\ {\isacharparenleft}{\kern0pt}induction\ {\isasympi}s\ arbitrary{\isacharcolon}{\kern0pt}\ {\isasympi}{\isacharparenright}{\kern0pt}\ auto%
\endisatagproof
{\isafoldproof}%
%
\isadelimproof
\isanewline
%
\endisadelimproof
\isanewline
\ \ \isacommand{lemma}\isamarkupfalse%
\ simplify{\isacharunderscore}{\kern0pt}plan{\isacharunderscore}{\kern0pt}happs{\isacharunderscore}{\kern0pt}subset{\isacharunderscore}{\kern0pt}app{\isacharcolon}{\kern0pt}\ {\isachardoublequoteopen}simplify{\isacharunderscore}{\kern0pt}plan\ htps\ {\isasympi}s\isactrlsub {\isadigit{2}}\ {\isasymsubseteq}\isactrlsub h\isactrlsub s\ simplify{\isacharunderscore}{\kern0pt}plan\ htps\ {\isacharparenleft}{\kern0pt}{\isasympi}s\isactrlsub {\isadigit{1}}\ {\isacharat}{\kern0pt}\ {\isasympi}s\isactrlsub {\isadigit{2}}{\isacharparenright}{\kern0pt}{\isachardoublequoteclose}\isanewline
%
\isadelimproof
\ \ \ \ %
\endisadelimproof
%
\isatagproof
\isacommand{using}\isamarkupfalse%
\ simplify{\isacharunderscore}{\kern0pt}plan{\isacharunderscore}{\kern0pt}happs{\isacharunderscore}{\kern0pt}subset{\isacharunderscore}{\kern0pt}Cons\isanewline
\ \ \ \ \isacommand{apply}\isamarkupfalse%
\ {\isacharparenleft}{\kern0pt}induction\ {\isasympi}s\isactrlsub {\isadigit{1}}\ arbitrary{\isacharcolon}{\kern0pt}\ {\isasympi}s\isactrlsub {\isadigit{2}}{\isacharparenright}{\kern0pt}\isanewline
\ \ \ \ \isacommand{apply}\isamarkupfalse%
\ {\isacharparenleft}{\kern0pt}auto\ simp{\isacharcolon}{\kern0pt}\ happ{\isacharunderscore}{\kern0pt}seq{\isacharunderscore}{\kern0pt}subset{\isacharunderscore}{\kern0pt}refl{\isacharparenright}{\kern0pt}\isanewline
\ \ \ \ \isacommand{using}\isamarkupfalse%
\ happ{\isacharunderscore}{\kern0pt}seq{\isacharunderscore}{\kern0pt}subset{\isacharunderscore}{\kern0pt}trans\ \isacommand{apply}\isamarkupfalse%
\ blast\isanewline
\ \ \ \ \isacommand{done}\isamarkupfalse%
%
\endisatagproof
{\isafoldproof}%
%
\isadelimproof
%
\endisadelimproof
%
\begin{isamarkuptext}%
Properties of \isa{simplify{\isacharunderscore}{\kern0pt}plan}%
\end{isamarkuptext}\isamarkuptrue%
%
\begin{isamarkuptext}%
All happenings in a simplified plan are always non-Nil.%
\end{isamarkuptext}\isamarkuptrue%
\ \ \isacommand{lemma}\isamarkupfalse%
\ simp{\isacharunderscore}{\kern0pt}plan{\isacharunderscore}{\kern0pt}happ{\isacharunderscore}{\kern0pt}not{\isacharunderscore}{\kern0pt}Nil{\isacharcolon}{\kern0pt}\isanewline
\ \ \ \ \isakeyword{fixes}\ {\isasympi}s\isanewline
\ \ \ \ \isakeyword{defines}\ {\isachardoublequoteopen}htps\ {\isasymequiv}\ htps{\isacharunderscore}{\kern0pt}exec\ {\isasympi}s{\isachardoublequoteclose}\isanewline
\ \ \ \ \isakeyword{shows}\ {\isachardoublequoteopen}{\isasymforall}{\isacharparenleft}{\kern0pt}t\isactrlsub a{\isacharcomma}{\kern0pt}A{\isacharparenright}{\kern0pt}\ {\isasymin}\ set\ {\isacharparenleft}{\kern0pt}simplify{\isacharunderscore}{\kern0pt}plan\ htps\ {\isasympi}s{\isacharparenright}{\kern0pt}{\isachardot}{\kern0pt}\ A\ {\isasymnoteq}\ {\isacharbrackleft}{\kern0pt}{\isacharbrackright}{\kern0pt}{\isachardoublequoteclose}\isanewline
%
\isadelimproof
\ \ %
\endisadelimproof
%
\isatagproof
\isacommand{proof}\isamarkupfalse%
\ {\isacharparenleft}{\kern0pt}induction\ {\isasympi}s{\isacharparenright}{\kern0pt}\isanewline
\ \ \ \ \isacommand{case}\isamarkupfalse%
\ Nil\isanewline
\ \ \ \ \isacommand{then}\isamarkupfalse%
\ \isacommand{show}\isamarkupfalse%
\ {\isacharquery}{\kern0pt}case\ \isacommand{by}\isamarkupfalse%
\ simp\isanewline
\ \ \isacommand{next}\isamarkupfalse%
\isanewline
\ \ \ \ \isacommand{case}\isamarkupfalse%
\ {\isacharparenleft}{\kern0pt}Cons\ {\isasympi}\ {\isasympi}s{\isacharparenright}{\kern0pt}\ \isanewline
\ \ \ \ \isacommand{have}\isamarkupfalse%
\ {\isachardoublequoteopen}{\isasymforall}{\isacharparenleft}{\kern0pt}t\isactrlsub a{\isacharcomma}{\kern0pt}A{\isacharparenright}{\kern0pt}\ {\isasymin}\ set\ {\isacharparenleft}{\kern0pt}map\ {\isacharparenleft}{\kern0pt}{\isasymlambda}{\isacharparenleft}{\kern0pt}t\isactrlsub a{\isacharcomma}{\kern0pt}a{\isacharparenright}{\kern0pt}{\isachardot}{\kern0pt}\ {\isacharparenleft}{\kern0pt}t\isactrlsub a{\isacharcomma}{\kern0pt}{\isacharbrackleft}{\kern0pt}a{\isacharbrackright}{\kern0pt}{\isacharparenright}{\kern0pt}{\isacharparenright}{\kern0pt}\ {\isacharparenleft}{\kern0pt}simplify{\isacharunderscore}{\kern0pt}action\ htps\ {\isasympi}{\isacharparenright}{\kern0pt}{\isacharparenright}{\kern0pt}{\isachardot}{\kern0pt}\ A\ {\isasymnoteq}\ {\isacharbrackleft}{\kern0pt}{\isacharbrackright}{\kern0pt}{\isachardoublequoteclose}\isanewline
\ \ \ \ \ \ \isacommand{by}\isamarkupfalse%
\ auto\isanewline
\ \ \ \ \isacommand{then}\isamarkupfalse%
\ \isacommand{show}\isamarkupfalse%
\ {\isacharquery}{\kern0pt}case\ \isanewline
\ \ \ \ \ \ \isacommand{using}\isamarkupfalse%
\ {\isachardoublequoteopen}Cons{\isachardot}{\kern0pt}IH{\isachardoublequoteclose}\ insort{\isacharunderscore}{\kern0pt}mult{\isacharunderscore}{\kern0pt}happs{\isacharunderscore}{\kern0pt}not{\isacharunderscore}{\kern0pt}Nil\ simplify{\isacharunderscore}{\kern0pt}plan{\isachardot}{\kern0pt}simps{\isacharparenleft}{\kern0pt}{\isadigit{2}}{\isacharparenright}{\kern0pt}\ \isanewline
\ \ \ \ \ \ \isacommand{by}\isamarkupfalse%
\ presburger\isanewline
\ \ \isacommand{qed}\isamarkupfalse%
%
\endisatagproof
{\isafoldproof}%
%
\isadelimproof
%
\endisadelimproof
%
\begin{isamarkuptext}%
A simplified plan produced by \isa{simplify{\isacharunderscore}{\kern0pt}plan} is strictly sorted.%
\end{isamarkuptext}\isamarkuptrue%
\ \ \isacommand{lemma}\isamarkupfalse%
\ simp{\isacharunderscore}{\kern0pt}plan{\isacharunderscore}{\kern0pt}strict{\isacharunderscore}{\kern0pt}sorted{\isacharcolon}{\kern0pt}\ {\isachardoublequoteopen}strict{\isacharunderscore}{\kern0pt}sorted\ {\isacharparenleft}{\kern0pt}map\ fst\ {\isacharparenleft}{\kern0pt}simplify{\isacharunderscore}{\kern0pt}plan\ htps\ {\isasympi}s{\isacharparenright}{\kern0pt}{\isacharparenright}{\kern0pt}{\isachardoublequoteclose}\isanewline
%
\isadelimproof
\ \ \ \ %
\endisadelimproof
%
\isatagproof
\isacommand{by}\isamarkupfalse%
\ {\isacharparenleft}{\kern0pt}induction\ {\isasympi}s{\isacharparenright}{\kern0pt}\ {\isacharparenleft}{\kern0pt}auto\ simp{\isacharcolon}{\kern0pt}\ insort{\isacharunderscore}{\kern0pt}mult{\isacharunderscore}{\kern0pt}happs{\isacharunderscore}{\kern0pt}strict{\isacharunderscore}{\kern0pt}sorted{\isacharparenright}{\kern0pt}%
\endisatagproof
{\isafoldproof}%
%
\isadelimproof
%
\endisadelimproof
%
\begin{isamarkuptext}%
A simplified plan contains all grounded actions for every \isa{Simple{\isacharunderscore}{\kern0pt}Plan{\isacharunderscore}{\kern0pt}Action}.%
\end{isamarkuptext}\isamarkuptrue%
\ \ \isacommand{lemma}\isamarkupfalse%
\ simp{\isacharunderscore}{\kern0pt}plan{\isacharunderscore}{\kern0pt}corr{\isacharunderscore}{\kern0pt}simp{\isacharunderscore}{\kern0pt}acts{\isacharcolon}{\kern0pt}\ \isanewline
\ \ \ \ \isakeyword{assumes}\ {\isachardoublequoteopen}hs\ {\isacharequal}{\kern0pt}\ simplify{\isacharunderscore}{\kern0pt}plan\ htps\ {\isasympi}s{\isachardoublequoteclose}\ \isanewline
\ \ \ \ \ \ \ \ \isakeyword{and}\ {\isachardoublequoteopen}{\isacharparenleft}{\kern0pt}t\isactrlsub {\isasympi}{\isacharcomma}{\kern0pt}{\isasympi}{\isacharparenright}{\kern0pt}\ {\isasymin}\ simple{\isacharunderscore}{\kern0pt}acts\ {\isasympi}s{\isachardoublequoteclose}\isanewline
\ \ \ \ \isakeyword{shows}\ {\isachardoublequoteopen}{\isasymexists}A{\isachardot}{\kern0pt}\ {\isacharparenleft}{\kern0pt}t\isactrlsub {\isasympi}{\isacharcomma}{\kern0pt}A{\isacharparenright}{\kern0pt}\ {\isasymin}\ set\ hs\ {\isasymand}\ the\ {\isacharparenleft}{\kern0pt}res{\isacharunderscore}{\kern0pt}inst\ {\isasympi}{\isacharparenright}{\kern0pt}\ {\isasymin}\ set\ A{\isachardoublequoteclose}\isanewline
%
\isadelimproof
\ \ \ \ %
\endisadelimproof
%
\isatagproof
\isacommand{using}\isamarkupfalse%
\ assms\isanewline
\ \ \isacommand{proof}\isamarkupfalse%
\ {\isacharminus}{\kern0pt}\isanewline
\ \ \ \ \isacommand{obtain}\isamarkupfalse%
\ {\isasympi}s\isactrlsub {\isadigit{1}}\ {\isasympi}s\isactrlsub {\isadigit{2}}\ \isakeyword{where}\ {\isachardoublequoteopen}{\isasympi}s\ {\isacharequal}{\kern0pt}\ {\isasympi}s\isactrlsub {\isadigit{1}}\ {\isacharat}{\kern0pt}\ {\isacharparenleft}{\kern0pt}{\isacharparenleft}{\kern0pt}t\isactrlsub {\isasympi}{\isacharcomma}{\kern0pt}{\isasympi}{\isacharparenright}{\kern0pt}{\isacharhash}{\kern0pt}{\isasympi}s\isactrlsub {\isadigit{2}}{\isacharparenright}{\kern0pt}{\isachardoublequoteclose}\isanewline
\ \ \ \ \ \ \isacommand{using}\isamarkupfalse%
\ assms\ simple{\isacharunderscore}{\kern0pt}acts{\isacharunderscore}{\kern0pt}subset\ \isacommand{by}\isamarkupfalse%
\ {\isacharparenleft}{\kern0pt}meson\ split{\isacharunderscore}{\kern0pt}list\ subsetD{\isacharparenright}{\kern0pt}\isanewline
\ \ \ \ \isacommand{then}\isamarkupfalse%
\ \isacommand{obtain}\isamarkupfalse%
\ hs{\isacharprime}{\kern0pt}\ \isakeyword{where}\ {\isachardoublequoteopen}hs{\isacharprime}{\kern0pt}\ {\isacharequal}{\kern0pt}\ simplify{\isacharunderscore}{\kern0pt}plan\ htps\ {\isacharparenleft}{\kern0pt}{\isacharparenleft}{\kern0pt}t\isactrlsub {\isasympi}{\isacharcomma}{\kern0pt}{\isasympi}{\isacharparenright}{\kern0pt}{\isacharhash}{\kern0pt}{\isasympi}s\isactrlsub {\isadigit{2}}{\isacharparenright}{\kern0pt}{\isachardoublequoteclose}\isanewline
\ \ \ \ \ \ \isacommand{by}\isamarkupfalse%
\ simp\isanewline
\ \ \ \ \isacommand{have}\isamarkupfalse%
\ {\isachardoublequoteopen}{\isacharparenleft}{\kern0pt}t\isactrlsub {\isasympi}{\isacharcomma}{\kern0pt}{\isasympi}{\isacharparenright}{\kern0pt}\ {\isacharequal}{\kern0pt}\ hd\ {\isacharparenleft}{\kern0pt}{\isacharparenleft}{\kern0pt}t\isactrlsub {\isasympi}{\isacharcomma}{\kern0pt}{\isasympi}{\isacharparenright}{\kern0pt}{\isacharhash}{\kern0pt}{\isasympi}s\isactrlsub {\isadigit{2}}{\isacharparenright}{\kern0pt}{\isachardoublequoteclose}\isanewline
\ \ \ \ \ \ \isacommand{by}\isamarkupfalse%
\ auto\isanewline
\ \ \ \ \isacommand{have}\isamarkupfalse%
\ {\isachardoublequoteopen}{\isacharparenleft}{\kern0pt}t\isactrlsub {\isasympi}{\isacharcomma}{\kern0pt}\ {\isasympi}{\isacharparenright}{\kern0pt}\ {\isasymin}\ simple{\isacharunderscore}{\kern0pt}acts\ {\isacharparenleft}{\kern0pt}{\isacharparenleft}{\kern0pt}t\isactrlsub {\isasympi}{\isacharcomma}{\kern0pt}\ {\isasympi}{\isacharparenright}{\kern0pt}\ {\isacharhash}{\kern0pt}\ {\isasympi}s\isactrlsub {\isadigit{2}}{\isacharparenright}{\kern0pt}{\isachardoublequoteclose}\isanewline
\ \ \ \ \ \ \isacommand{using}\isamarkupfalse%
\ assms\ {\isacartoucheopen}{\isasympi}s\ {\isacharequal}{\kern0pt}\ {\isasympi}s\isactrlsub {\isadigit{1}}\ {\isacharat}{\kern0pt}\ {\isacharparenleft}{\kern0pt}{\isacharparenleft}{\kern0pt}t\isactrlsub {\isasympi}{\isacharcomma}{\kern0pt}{\isasympi}{\isacharparenright}{\kern0pt}{\isacharhash}{\kern0pt}{\isasympi}s\isactrlsub {\isadigit{2}}{\isacharparenright}{\kern0pt}{\isacartoucheclose}\isanewline
\ \ \ \ \ \ \isacommand{unfolding}\isamarkupfalse%
\ simple{\isacharunderscore}{\kern0pt}acts{\isacharunderscore}{\kern0pt}def\ is{\isacharunderscore}{\kern0pt}act{\isacharunderscore}{\kern0pt}simple{\isacharunderscore}{\kern0pt}def\ \isacommand{by}\isamarkupfalse%
\ {\isacharparenleft}{\kern0pt}auto\ split{\isacharcolon}{\kern0pt}\ plan{\isacharunderscore}{\kern0pt}action{\isachardot}{\kern0pt}splits{\isacharparenright}{\kern0pt}\isanewline
\ \ \ \ \isacommand{have}\isamarkupfalse%
\ {\isachardoublequoteopen}{\isasymexists}as{\isachardot}{\kern0pt}\ {\isacharparenleft}{\kern0pt}t\isactrlsub {\isasympi}{\isacharcomma}{\kern0pt}as{\isacharparenright}{\kern0pt}\ {\isasymin}\ set\ hs{\isacharprime}{\kern0pt}\ {\isasymand}\ the\ {\isacharparenleft}{\kern0pt}res{\isacharunderscore}{\kern0pt}inst\ {\isasympi}{\isacharparenright}{\kern0pt}\ {\isasymin}\ set\ as{\isachardoublequoteclose}\isanewline
\ \ \ \ \ \ \isacommand{using}\isamarkupfalse%
\ {\isacartoucheopen}hs{\isacharprime}{\kern0pt}\ {\isacharequal}{\kern0pt}\ simplify{\isacharunderscore}{\kern0pt}plan\ htps\ {\isacharparenleft}{\kern0pt}{\isacharparenleft}{\kern0pt}t\isactrlsub {\isasympi}{\isacharcomma}{\kern0pt}{\isasympi}{\isacharparenright}{\kern0pt}{\isacharhash}{\kern0pt}{\isasympi}s\isactrlsub {\isadigit{2}}{\isacharparenright}{\kern0pt}{\isacartoucheclose}\ \isanewline
\ \ \ \ \ \ \ \ \ \ \ \ {\isacartoucheopen}{\isacharparenleft}{\kern0pt}t\isactrlsub {\isasympi}{\isacharcomma}{\kern0pt}{\isasympi}{\isacharparenright}{\kern0pt}\ {\isacharequal}{\kern0pt}\ hd\ {\isacharparenleft}{\kern0pt}{\isacharparenleft}{\kern0pt}t\isactrlsub {\isasympi}{\isacharcomma}{\kern0pt}{\isasympi}{\isacharparenright}{\kern0pt}{\isacharhash}{\kern0pt}{\isasympi}s\isactrlsub {\isadigit{2}}{\isacharparenright}{\kern0pt}{\isacartoucheclose}\ \isanewline
\ \ \ \ \ \ \ \ \ \ \ \ {\isacartoucheopen}{\isacharparenleft}{\kern0pt}t\isactrlsub {\isasympi}{\isacharcomma}{\kern0pt}\ {\isasympi}{\isacharparenright}{\kern0pt}\ {\isasymin}\ simple{\isacharunderscore}{\kern0pt}acts\ {\isacharparenleft}{\kern0pt}{\isacharparenleft}{\kern0pt}t\isactrlsub {\isasympi}{\isacharcomma}{\kern0pt}\ {\isasympi}{\isacharparenright}{\kern0pt}\ {\isacharhash}{\kern0pt}\ {\isasympi}s\isactrlsub {\isadigit{2}}{\isacharparenright}{\kern0pt}{\isacartoucheclose}\isanewline
\ \ \ \ \ \ \ \ \ \ \ \ insort{\isacharunderscore}{\kern0pt}happ{\isacharunderscore}{\kern0pt}action{\isacharunderscore}{\kern0pt}subset\isanewline
\ \ \ \ \ \ \isacommand{unfolding}\isamarkupfalse%
\ simple{\isacharunderscore}{\kern0pt}acts{\isacharunderscore}{\kern0pt}def\ is{\isacharunderscore}{\kern0pt}act{\isacharunderscore}{\kern0pt}simple{\isacharunderscore}{\kern0pt}def\isanewline
\ \ \ \ \ \ \isacommand{apply}\isamarkupfalse%
\ {\isacharparenleft}{\kern0pt}induction\ {\isasympi}s{\isacharparenright}{\kern0pt}\isanewline
\ \ \ \ \ \ \isacommand{apply}\isamarkupfalse%
\ {\isacharparenleft}{\kern0pt}auto\ split{\isacharcolon}{\kern0pt}\ plan{\isacharunderscore}{\kern0pt}action{\isachardot}{\kern0pt}splits{\isacharparenright}{\kern0pt}\isanewline
\ \ \ \ \ \ \isacommand{apply}\isamarkupfalse%
\ {\isacharparenleft}{\kern0pt}meson\ list{\isachardot}{\kern0pt}set{\isacharunderscore}{\kern0pt}intros{\isacharparenleft}{\kern0pt}{\isadigit{1}}{\isacharparenright}{\kern0pt}\ subsetD{\isacharparenright}{\kern0pt}\isanewline
\ \ \ \ \ \ \isacommand{done}\isamarkupfalse%
\isanewline
\ \ \ \ \isacommand{have}\isamarkupfalse%
\ {\isachardoublequoteopen}hs{\isacharprime}{\kern0pt}\ {\isasymsubseteq}\isactrlsub h\isactrlsub s\ hs{\isachardoublequoteclose}\isanewline
\ \ \ \ \ \ \isacommand{using}\isamarkupfalse%
\ assms\ simplify{\isacharunderscore}{\kern0pt}plan{\isacharunderscore}{\kern0pt}happs{\isacharunderscore}{\kern0pt}subset{\isacharunderscore}{\kern0pt}app\ \isanewline
\ \ \ \ \ \ \ \ \ \ \ \ {\isacartoucheopen}hs{\isacharprime}{\kern0pt}\ {\isacharequal}{\kern0pt}\ simplify{\isacharunderscore}{\kern0pt}plan\ htps\ {\isacharparenleft}{\kern0pt}{\isacharparenleft}{\kern0pt}t\isactrlsub {\isasympi}{\isacharcomma}{\kern0pt}{\isasympi}{\isacharparenright}{\kern0pt}{\isacharhash}{\kern0pt}{\isasympi}s\isactrlsub {\isadigit{2}}{\isacharparenright}{\kern0pt}{\isacartoucheclose}\ \isanewline
\ \ \ \ \ \ \ \ \ \ \ \ {\isacartoucheopen}{\isasympi}s\ {\isacharequal}{\kern0pt}\ {\isasympi}s\isactrlsub {\isadigit{1}}\ {\isacharat}{\kern0pt}\ {\isacharparenleft}{\kern0pt}{\isacharparenleft}{\kern0pt}t\isactrlsub {\isasympi}{\isacharcomma}{\kern0pt}{\isasympi}{\isacharparenright}{\kern0pt}{\isacharhash}{\kern0pt}{\isasympi}s\isactrlsub {\isadigit{2}}{\isacharparenright}{\kern0pt}{\isacartoucheclose}\isanewline
\ \ \ \ \ \ \isacommand{by}\isamarkupfalse%
\ blast\isanewline
\ \ \ \ \isacommand{show}\isamarkupfalse%
\ {\isacharquery}{\kern0pt}thesis\isanewline
\ \ \ \ \ \ \isacommand{using}\isamarkupfalse%
\ {\isacartoucheopen}{\isasymexists}as{\isachardot}{\kern0pt}\ {\isacharparenleft}{\kern0pt}t\isactrlsub {\isasympi}{\isacharcomma}{\kern0pt}as{\isacharparenright}{\kern0pt}\ {\isasymin}\ set\ hs{\isacharprime}{\kern0pt}\ {\isasymand}\ the\ {\isacharparenleft}{\kern0pt}res{\isacharunderscore}{\kern0pt}inst\ {\isasympi}{\isacharparenright}{\kern0pt}\ {\isasymin}\ set\ as{\isacartoucheclose}\isanewline
\ \ \ \ \ \ \ \ \ \ \ \ {\isacartoucheopen}hs{\isacharprime}{\kern0pt}\ {\isasymsubseteq}\isactrlsub h\isactrlsub s\ hs{\isacartoucheclose}\isanewline
\ \ \ \ \ \ \isacommand{unfolding}\isamarkupfalse%
\ happ{\isacharunderscore}{\kern0pt}seq{\isacharunderscore}{\kern0pt}subset{\isacharunderscore}{\kern0pt}def\ happ{\isacharunderscore}{\kern0pt}seq{\isacharunderscore}{\kern0pt}member{\isacharunderscore}{\kern0pt}def\ \isacommand{by}\isamarkupfalse%
\ blast\isanewline
\ \ \isacommand{qed}\isamarkupfalse%
%
\endisatagproof
{\isafoldproof}%
%
\isadelimproof
\isanewline
%
\endisadelimproof
\isanewline
\ \ \isacommand{lemma}\isamarkupfalse%
\ simp{\isacharunderscore}{\kern0pt}plan{\isacharunderscore}{\kern0pt}corr{\isacharunderscore}{\kern0pt}dur{\isacharunderscore}{\kern0pt}acts{\isacharunderscore}{\kern0pt}start{\isacharunderscore}{\kern0pt}aux{\isacharcolon}{\kern0pt}\ \isanewline
\ \ \ \ \isakeyword{assumes}\ {\isachardoublequoteopen}hs\ {\isacharequal}{\kern0pt}\ simplify{\isacharunderscore}{\kern0pt}plan\ htps\ {\isasympi}s{\isachardoublequoteclose}\ \isanewline
\ \ \ \ \ \ \ \ \isakeyword{and}\ {\isachardoublequoteopen}{\isacharparenleft}{\kern0pt}t\isactrlsub {\isasympi}{\isacharcomma}{\kern0pt}{\isasympi}{\isacharparenright}{\kern0pt}\ {\isacharequal}{\kern0pt}\ hd\ {\isasympi}s{\isachardoublequoteclose}\ \isanewline
\ \ \ \ \ \ \ \ \isakeyword{and}\ {\isachardoublequoteopen}{\isacharparenleft}{\kern0pt}t\isactrlsub {\isasympi}{\isacharcomma}{\kern0pt}{\isasympi}{\isacharparenright}{\kern0pt}\ {\isasymin}\ durative{\isacharunderscore}{\kern0pt}acts\ {\isasympi}s{\isachardoublequoteclose}\isanewline
\ \ \ \ \isakeyword{shows}\ {\isachardoublequoteopen}{\isacharparenleft}{\kern0pt}{\isasymexists}A{\isachardot}{\kern0pt}\ {\isacharparenleft}{\kern0pt}t\isactrlsub {\isasympi}{\isacharcomma}{\kern0pt}A{\isacharparenright}{\kern0pt}\ {\isasymin}\ set\ hs\ {\isasymand}\ {\isacharparenleft}{\kern0pt}the\ {\isacharparenleft}{\kern0pt}res{\isacharunderscore}{\kern0pt}inst{\isacharunderscore}{\kern0pt}snap{\isacharunderscore}{\kern0pt}action\ {\isasympi}\ At{\isacharunderscore}{\kern0pt}Start{\isacharparenright}{\kern0pt}{\isacharparenright}{\kern0pt}\ {\isasymin}\ set\ A{\isacharparenright}{\kern0pt}{\isachardoublequoteclose}\isanewline
%
\isadelimproof
\ \ \ \ %
\endisadelimproof
%
\isatagproof
\isacommand{using}\isamarkupfalse%
\ assms\ insort{\isacharunderscore}{\kern0pt}happ{\isacharunderscore}{\kern0pt}action{\isacharunderscore}{\kern0pt}subset\isanewline
\ \ \ \ \isacommand{unfolding}\isamarkupfalse%
\ durative{\isacharunderscore}{\kern0pt}acts{\isacharunderscore}{\kern0pt}def\ is{\isacharunderscore}{\kern0pt}act{\isacharunderscore}{\kern0pt}simple{\isacharunderscore}{\kern0pt}def\isanewline
\ \ \ \ \isacommand{apply}\isamarkupfalse%
\ {\isacharparenleft}{\kern0pt}induction\ {\isasympi}s{\isacharparenright}{\kern0pt}\isanewline
\ \ \ \ \isacommand{apply}\isamarkupfalse%
\ {\isacharparenleft}{\kern0pt}auto\ simp{\isacharcolon}{\kern0pt}\ Let{\isacharunderscore}{\kern0pt}def\ split{\isacharcolon}{\kern0pt}\ plan{\isacharunderscore}{\kern0pt}action{\isachardot}{\kern0pt}splits{\isacharparenright}{\kern0pt}\isanewline
\ \ \ \ \isacommand{apply}\isamarkupfalse%
\ {\isacharparenleft}{\kern0pt}meson\ list{\isachardot}{\kern0pt}set{\isacharunderscore}{\kern0pt}intros{\isacharparenleft}{\kern0pt}{\isadigit{1}}{\isacharparenright}{\kern0pt}\ subsetD{\isacharparenright}{\kern0pt}\isanewline
\ \ \ \ \isacommand{done}\isamarkupfalse%
%
\endisatagproof
{\isafoldproof}%
%
\isadelimproof
%
\endisadelimproof
%
\begin{isamarkuptext}%
A simplified plan contains all \isa{At{\isacharunderscore}{\kern0pt}Start}-snap-actions.%
\end{isamarkuptext}\isamarkuptrue%
\ \ \isacommand{lemma}\isamarkupfalse%
\ simp{\isacharunderscore}{\kern0pt}plan{\isacharunderscore}{\kern0pt}corr{\isacharunderscore}{\kern0pt}dur{\isacharunderscore}{\kern0pt}acts{\isacharunderscore}{\kern0pt}start{\isacharcolon}{\kern0pt}\ \isanewline
\ \ \ \ \isakeyword{fixes}\ {\isasympi}\isanewline
\ \ \ \ \isakeyword{defines}\ {\isachardoublequoteopen}a\isactrlsub s\isactrlsub t\isactrlsub a\isactrlsub r\isactrlsub t\ {\isasymequiv}\ the\ {\isacharparenleft}{\kern0pt}res{\isacharunderscore}{\kern0pt}inst{\isacharunderscore}{\kern0pt}snap{\isacharunderscore}{\kern0pt}action\ {\isasympi}\ At{\isacharunderscore}{\kern0pt}Start{\isacharparenright}{\kern0pt}{\isachardoublequoteclose}\isanewline
\ \ \ \ \isakeyword{assumes}\ {\isachardoublequoteopen}hs\ {\isacharequal}{\kern0pt}\ simplify{\isacharunderscore}{\kern0pt}plan\ htps\ {\isasympi}s{\isachardoublequoteclose}\ \isanewline
\ \ \ \ \ \ \ \ \isakeyword{and}\ {\isachardoublequoteopen}{\isacharparenleft}{\kern0pt}t\isactrlsub {\isasympi}{\isacharcomma}{\kern0pt}{\isasympi}{\isacharparenright}{\kern0pt}\ {\isasymin}\ durative{\isacharunderscore}{\kern0pt}acts\ {\isasympi}s{\isachardoublequoteclose}\isanewline
\ \ \ \ \isakeyword{shows}\ {\isachardoublequoteopen}{\isasymexists}A{\isachardot}{\kern0pt}\ {\isacharparenleft}{\kern0pt}t\isactrlsub {\isasympi}{\isacharcomma}{\kern0pt}A{\isacharparenright}{\kern0pt}\ {\isasymin}\ set\ hs\ {\isasymand}\ a\isactrlsub s\isactrlsub t\isactrlsub a\isactrlsub r\isactrlsub t\ {\isasymin}\ set\ A{\isachardoublequoteclose}\isanewline
%
\isadelimproof
\ \ \ \ %
\endisadelimproof
%
\isatagproof
\isacommand{using}\isamarkupfalse%
\ assms\isanewline
\ \ \isacommand{proof}\isamarkupfalse%
\ {\isacharminus}{\kern0pt}\isanewline
\ \ \ \ \isacommand{obtain}\isamarkupfalse%
\ {\isasympi}s\isactrlsub {\isadigit{1}}\ {\isasympi}s\isactrlsub {\isadigit{2}}\ \isakeyword{where}\ {\isachardoublequoteopen}{\isasympi}s\ {\isacharequal}{\kern0pt}\ {\isasympi}s\isactrlsub {\isadigit{1}}\ {\isacharat}{\kern0pt}\ {\isacharparenleft}{\kern0pt}{\isacharparenleft}{\kern0pt}t\isactrlsub {\isasympi}{\isacharcomma}{\kern0pt}{\isasympi}{\isacharparenright}{\kern0pt}{\isacharhash}{\kern0pt}{\isasympi}s\isactrlsub {\isadigit{2}}{\isacharparenright}{\kern0pt}{\isachardoublequoteclose}\isanewline
\ \ \ \ \ \ \isacommand{using}\isamarkupfalse%
\ assms\ dur{\isacharunderscore}{\kern0pt}acts{\isacharunderscore}{\kern0pt}subset\ \isacommand{by}\isamarkupfalse%
\ {\isacharparenleft}{\kern0pt}meson\ split{\isacharunderscore}{\kern0pt}list\ subsetD{\isacharparenright}{\kern0pt}\isanewline
\ \ \ \ \isacommand{then}\isamarkupfalse%
\ \isacommand{obtain}\isamarkupfalse%
\ hs{\isacharprime}{\kern0pt}\ \isakeyword{where}\ {\isachardoublequoteopen}hs{\isacharprime}{\kern0pt}\ {\isacharequal}{\kern0pt}\ simplify{\isacharunderscore}{\kern0pt}plan\ htps\ {\isacharparenleft}{\kern0pt}{\isacharparenleft}{\kern0pt}t\isactrlsub {\isasympi}{\isacharcomma}{\kern0pt}{\isasympi}{\isacharparenright}{\kern0pt}{\isacharhash}{\kern0pt}{\isasympi}s\isactrlsub {\isadigit{2}}{\isacharparenright}{\kern0pt}{\isachardoublequoteclose}\isanewline
\ \ \ \ \ \ \isacommand{by}\isamarkupfalse%
\ simp\isanewline
\ \ \ \ \isacommand{have}\isamarkupfalse%
\ {\isachardoublequoteopen}{\isacharparenleft}{\kern0pt}t\isactrlsub {\isasympi}{\isacharcomma}{\kern0pt}{\isasympi}{\isacharparenright}{\kern0pt}\ {\isacharequal}{\kern0pt}\ hd\ {\isacharparenleft}{\kern0pt}{\isacharparenleft}{\kern0pt}t\isactrlsub {\isasympi}{\isacharcomma}{\kern0pt}{\isasympi}{\isacharparenright}{\kern0pt}{\isacharhash}{\kern0pt}{\isasympi}s\isactrlsub {\isadigit{2}}{\isacharparenright}{\kern0pt}{\isachardoublequoteclose}\isanewline
\ \ \ \ \ \ \isacommand{by}\isamarkupfalse%
\ auto\isanewline
\ \ \ \ \isacommand{have}\isamarkupfalse%
\ {\isachardoublequoteopen}{\isacharparenleft}{\kern0pt}t\isactrlsub {\isasympi}{\isacharcomma}{\kern0pt}\ {\isasympi}{\isacharparenright}{\kern0pt}\ {\isasymin}\ durative{\isacharunderscore}{\kern0pt}acts\ {\isacharparenleft}{\kern0pt}{\isacharparenleft}{\kern0pt}t\isactrlsub {\isasympi}{\isacharcomma}{\kern0pt}\ {\isasympi}{\isacharparenright}{\kern0pt}\ {\isacharhash}{\kern0pt}\ {\isasympi}s\isactrlsub {\isadigit{2}}{\isacharparenright}{\kern0pt}{\isachardoublequoteclose}\isanewline
\ \ \ \ \ \ \isacommand{using}\isamarkupfalse%
\ assms\ {\isacartoucheopen}{\isasympi}s\ {\isacharequal}{\kern0pt}\ {\isasympi}s\isactrlsub {\isadigit{1}}\ {\isacharat}{\kern0pt}\ {\isacharparenleft}{\kern0pt}{\isacharparenleft}{\kern0pt}t\isactrlsub {\isasympi}{\isacharcomma}{\kern0pt}{\isasympi}{\isacharparenright}{\kern0pt}{\isacharhash}{\kern0pt}{\isasympi}s\isactrlsub {\isadigit{2}}{\isacharparenright}{\kern0pt}{\isacartoucheclose}\isanewline
\ \ \ \ \ \ \isacommand{unfolding}\isamarkupfalse%
\ durative{\isacharunderscore}{\kern0pt}acts{\isacharunderscore}{\kern0pt}def\ is{\isacharunderscore}{\kern0pt}act{\isacharunderscore}{\kern0pt}simple{\isacharunderscore}{\kern0pt}def\ \isacommand{by}\isamarkupfalse%
\ {\isacharparenleft}{\kern0pt}auto\ split{\isacharcolon}{\kern0pt}\ plan{\isacharunderscore}{\kern0pt}action{\isachardot}{\kern0pt}splits{\isacharparenright}{\kern0pt}\isanewline
\ \ \ \ \isacommand{have}\isamarkupfalse%
\ {\isachardoublequoteopen}{\isasymexists}as{\isachardot}{\kern0pt}\ {\isacharparenleft}{\kern0pt}t\isactrlsub {\isasympi}{\isacharcomma}{\kern0pt}as{\isacharparenright}{\kern0pt}\ {\isasymin}\ set\ hs{\isacharprime}{\kern0pt}\ {\isasymand}\ a\isactrlsub s\isactrlsub t\isactrlsub a\isactrlsub r\isactrlsub t\ {\isasymin}\ set\ as{\isachardoublequoteclose}\isanewline
\ \ \ \ \ \ \isacommand{using}\isamarkupfalse%
\ assms\ simp{\isacharunderscore}{\kern0pt}plan{\isacharunderscore}{\kern0pt}corr{\isacharunderscore}{\kern0pt}dur{\isacharunderscore}{\kern0pt}acts{\isacharunderscore}{\kern0pt}start{\isacharunderscore}{\kern0pt}aux{\isacharbrackleft}{\kern0pt}OF\ {\isacartoucheopen}hs{\isacharprime}{\kern0pt}\ {\isacharequal}{\kern0pt}\ simplify{\isacharunderscore}{\kern0pt}plan\ htps\ {\isacharparenleft}{\kern0pt}{\isacharparenleft}{\kern0pt}t\isactrlsub {\isasympi}{\isacharcomma}{\kern0pt}{\isasympi}{\isacharparenright}{\kern0pt}{\isacharhash}{\kern0pt}{\isasympi}s\isactrlsub {\isadigit{2}}{\isacharparenright}{\kern0pt}{\isacartoucheclose}\ \isanewline
\ \ \ \ \ \ \ \ \ \ \ \ \ \ \ \ \ \ \ \ \ \ \ \ \ \ \ \ \ \ \ \ \ \ \ \ \ \ \ \ \ \ \ \ \ \ \ \ \ \ \ \ \ \ \ {\isacartoucheopen}{\isacharparenleft}{\kern0pt}t\isactrlsub {\isasympi}{\isacharcomma}{\kern0pt}{\isasympi}{\isacharparenright}{\kern0pt}\ {\isacharequal}{\kern0pt}\ hd\ {\isacharparenleft}{\kern0pt}{\isacharparenleft}{\kern0pt}t\isactrlsub {\isasympi}{\isacharcomma}{\kern0pt}{\isasympi}{\isacharparenright}{\kern0pt}{\isacharhash}{\kern0pt}{\isasympi}s\isactrlsub {\isadigit{2}}{\isacharparenright}{\kern0pt}{\isacartoucheclose}\ \isanewline
\ \ \ \ \ \ \ \ \ \ \ \ \ \ \ \ \ \ \ \ \ \ \ \ \ \ \ \ \ \ \ \ \ \ \ \ \ \ \ \ \ \ \ \ \ \ \ \ \ \ \ \ \ \ \ {\isacartoucheopen}{\isacharparenleft}{\kern0pt}t\isactrlsub {\isasympi}{\isacharcomma}{\kern0pt}{\isasympi}{\isacharparenright}{\kern0pt}\ {\isasymin}\ durative{\isacharunderscore}{\kern0pt}acts\ {\isacharparenleft}{\kern0pt}{\isacharparenleft}{\kern0pt}t\isactrlsub {\isasympi}{\isacharcomma}{\kern0pt}{\isasympi}{\isacharparenright}{\kern0pt}{\isacharhash}{\kern0pt}{\isasympi}s\isactrlsub {\isadigit{2}}{\isacharparenright}{\kern0pt}{\isacartoucheclose}{\isacharbrackright}{\kern0pt}\ \isanewline
\ \ \ \ \ \ \isacommand{by}\isamarkupfalse%
\ simp\isanewline
\ \ \ \ \isacommand{have}\isamarkupfalse%
\ {\isachardoublequoteopen}hs{\isacharprime}{\kern0pt}\ {\isasymsubseteq}\isactrlsub h\isactrlsub s\ hs{\isachardoublequoteclose}\isanewline
\ \ \ \ \ \ \isacommand{using}\isamarkupfalse%
\ assms\ simplify{\isacharunderscore}{\kern0pt}plan{\isacharunderscore}{\kern0pt}happs{\isacharunderscore}{\kern0pt}subset{\isacharunderscore}{\kern0pt}app\ \isanewline
\ \ \ \ \ \ \ \ \ \ \ \ {\isacartoucheopen}hs{\isacharprime}{\kern0pt}\ {\isacharequal}{\kern0pt}\ simplify{\isacharunderscore}{\kern0pt}plan\ htps\ {\isacharparenleft}{\kern0pt}{\isacharparenleft}{\kern0pt}t\isactrlsub {\isasympi}{\isacharcomma}{\kern0pt}{\isasympi}{\isacharparenright}{\kern0pt}{\isacharhash}{\kern0pt}{\isasympi}s\isactrlsub {\isadigit{2}}{\isacharparenright}{\kern0pt}{\isacartoucheclose}\ \isanewline
\ \ \ \ \ \ \ \ \ \ \ \ {\isacartoucheopen}{\isasympi}s\ {\isacharequal}{\kern0pt}\ {\isasympi}s\isactrlsub {\isadigit{1}}\ {\isacharat}{\kern0pt}\ {\isacharparenleft}{\kern0pt}{\isacharparenleft}{\kern0pt}t\isactrlsub {\isasympi}{\isacharcomma}{\kern0pt}{\isasympi}{\isacharparenright}{\kern0pt}{\isacharhash}{\kern0pt}{\isasympi}s\isactrlsub {\isadigit{2}}{\isacharparenright}{\kern0pt}{\isacartoucheclose}\isanewline
\ \ \ \ \ \ \isacommand{by}\isamarkupfalse%
\ blast\isanewline
\ \ \ \ \isacommand{show}\isamarkupfalse%
\ {\isacharquery}{\kern0pt}thesis\isanewline
\ \ \ \ \ \ \isacommand{using}\isamarkupfalse%
\ {\isacartoucheopen}{\isasymexists}as{\isachardot}{\kern0pt}\ {\isacharparenleft}{\kern0pt}t\isactrlsub {\isasympi}{\isacharcomma}{\kern0pt}as{\isacharparenright}{\kern0pt}\ {\isasymin}\ set\ hs{\isacharprime}{\kern0pt}\ {\isasymand}\ a\isactrlsub s\isactrlsub t\isactrlsub a\isactrlsub r\isactrlsub t\ {\isasymin}\ set\ as{\isacartoucheclose}\isanewline
\ \ \ \ \ \ \ \ \ \ \ \ {\isacartoucheopen}hs{\isacharprime}{\kern0pt}\ {\isasymsubseteq}\isactrlsub h\isactrlsub s\ hs{\isacartoucheclose}\isanewline
\ \ \ \ \ \ \isacommand{unfolding}\isamarkupfalse%
\ happ{\isacharunderscore}{\kern0pt}seq{\isacharunderscore}{\kern0pt}subset{\isacharunderscore}{\kern0pt}def\ happ{\isacharunderscore}{\kern0pt}seq{\isacharunderscore}{\kern0pt}member{\isacharunderscore}{\kern0pt}def\ \isacommand{by}\isamarkupfalse%
\ blast\isanewline
\ \ \isacommand{qed}\isamarkupfalse%
%
\endisatagproof
{\isafoldproof}%
%
\isadelimproof
\isanewline
%
\endisadelimproof
\isanewline
\ \ \isacommand{lemma}\isamarkupfalse%
\ simp{\isacharunderscore}{\kern0pt}plan{\isacharunderscore}{\kern0pt}corr{\isacharunderscore}{\kern0pt}dur{\isacharunderscore}{\kern0pt}acts{\isacharunderscore}{\kern0pt}end{\isacharunderscore}{\kern0pt}aux{\isacharcolon}{\kern0pt}\ \isanewline
\ \ \ \ \isakeyword{assumes}\ {\isachardoublequoteopen}hs\ {\isacharequal}{\kern0pt}\ simplify{\isacharunderscore}{\kern0pt}plan\ htps\ {\isasympi}s{\isachardoublequoteclose}\ \isanewline
\ \ \ \ \ \ \ \ \isakeyword{and}\ {\isachardoublequoteopen}{\isacharparenleft}{\kern0pt}t\isactrlsub {\isasympi}{\isacharcomma}{\kern0pt}{\isasympi}{\isacharparenright}{\kern0pt}\ {\isacharequal}{\kern0pt}\ hd\ {\isasympi}s{\isachardoublequoteclose}\ \isanewline
\ \ \ \ \ \ \ \ \isakeyword{and}\ {\isachardoublequoteopen}{\isacharparenleft}{\kern0pt}t\isactrlsub {\isasympi}{\isacharcomma}{\kern0pt}{\isasympi}{\isacharparenright}{\kern0pt}\ {\isasymin}\ durative{\isacharunderscore}{\kern0pt}acts\ {\isasympi}s{\isachardoublequoteclose}\isanewline
\ \ \ \ \isakeyword{shows}\ {\isachardoublequoteopen}{\isacharparenleft}{\kern0pt}{\isasymexists}A{\isachardot}{\kern0pt}\ {\isacharparenleft}{\kern0pt}t\isactrlsub {\isasympi}\ {\isacharplus}{\kern0pt}\ {\isacharparenleft}{\kern0pt}duration\ {\isasympi}{\isacharparenright}{\kern0pt}{\isacharcomma}{\kern0pt}A{\isacharparenright}{\kern0pt}\ {\isasymin}\ set\ hs\ {\isasymand}\ {\isacharparenleft}{\kern0pt}the\ {\isacharparenleft}{\kern0pt}res{\isacharunderscore}{\kern0pt}inst{\isacharunderscore}{\kern0pt}snap{\isacharunderscore}{\kern0pt}action\ {\isasympi}\ At{\isacharunderscore}{\kern0pt}End{\isacharparenright}{\kern0pt}{\isacharparenright}{\kern0pt}\ {\isasymin}\ set\ A{\isacharparenright}{\kern0pt}{\isachardoublequoteclose}\isanewline
%
\isadelimproof
\ \ \ \ %
\endisadelimproof
%
\isatagproof
\isacommand{using}\isamarkupfalse%
\ assms\ insort{\isacharunderscore}{\kern0pt}happ{\isadigit{2}}{\isacharunderscore}{\kern0pt}action{\isacharunderscore}{\kern0pt}subset\isanewline
\ \ \ \ \isacommand{unfolding}\isamarkupfalse%
\ durative{\isacharunderscore}{\kern0pt}acts{\isacharunderscore}{\kern0pt}def\ is{\isacharunderscore}{\kern0pt}act{\isacharunderscore}{\kern0pt}simple{\isacharunderscore}{\kern0pt}def\isanewline
\ \ \ \ \isacommand{apply}\isamarkupfalse%
\ {\isacharparenleft}{\kern0pt}induction\ {\isasympi}s{\isacharparenright}{\kern0pt}\isanewline
\ \ \ \ \isacommand{apply}\isamarkupfalse%
\ {\isacharparenleft}{\kern0pt}auto\ simp{\isacharcolon}{\kern0pt}\ Let{\isacharunderscore}{\kern0pt}def\ split{\isacharcolon}{\kern0pt}\ plan{\isacharunderscore}{\kern0pt}action{\isachardot}{\kern0pt}splits{\isacharparenright}{\kern0pt}\isanewline
\ \ \ \ \isacommand{apply}\isamarkupfalse%
\ {\isacharparenleft}{\kern0pt}meson\ list{\isachardot}{\kern0pt}set{\isacharunderscore}{\kern0pt}intros{\isacharparenleft}{\kern0pt}{\isadigit{1}}{\isacharparenright}{\kern0pt}\ subsetD{\isacharparenright}{\kern0pt}\isanewline
\ \ \ \ \isacommand{done}\isamarkupfalse%
%
\endisatagproof
{\isafoldproof}%
%
\isadelimproof
%
\endisadelimproof
%
\begin{isamarkuptext}%
A simplified plan contains all \isa{At{\isacharunderscore}{\kern0pt}End}-snap-actions.%
\end{isamarkuptext}\isamarkuptrue%
\ \ \isacommand{lemma}\isamarkupfalse%
\ simp{\isacharunderscore}{\kern0pt}plan{\isacharunderscore}{\kern0pt}corr{\isacharunderscore}{\kern0pt}dur{\isacharunderscore}{\kern0pt}acts{\isacharunderscore}{\kern0pt}end{\isacharcolon}{\kern0pt}\ \isanewline
\ \ \ \ \isakeyword{fixes}\ {\isasympi}\isanewline
\ \ \ \ \isakeyword{defines}\ {\isachardoublequoteopen}a\isactrlsub e\isactrlsub n\isactrlsub d\ {\isasymequiv}\ the\ {\isacharparenleft}{\kern0pt}res{\isacharunderscore}{\kern0pt}inst{\isacharunderscore}{\kern0pt}snap{\isacharunderscore}{\kern0pt}action\ {\isasympi}\ At{\isacharunderscore}{\kern0pt}End{\isacharparenright}{\kern0pt}{\isachardoublequoteclose}\isanewline
\ \ \ \ \isakeyword{assumes}\ {\isachardoublequoteopen}hs\ {\isacharequal}{\kern0pt}\ simplify{\isacharunderscore}{\kern0pt}plan\ htps\ {\isasympi}s{\isachardoublequoteclose}\ \isanewline
\ \ \ \ \ \ \ \ \isakeyword{and}\ {\isachardoublequoteopen}{\isacharparenleft}{\kern0pt}t\isactrlsub {\isasympi}{\isacharcomma}{\kern0pt}{\isasympi}{\isacharparenright}{\kern0pt}\ {\isasymin}\ durative{\isacharunderscore}{\kern0pt}acts\ {\isasympi}s{\isachardoublequoteclose}\isanewline
\ \ \ \ \isakeyword{shows}\ {\isachardoublequoteopen}{\isasymexists}A{\isachardot}{\kern0pt}\ {\isacharparenleft}{\kern0pt}t\isactrlsub {\isasympi}\ {\isacharplus}{\kern0pt}\ {\isacharparenleft}{\kern0pt}duration\ {\isasympi}{\isacharparenright}{\kern0pt}{\isacharcomma}{\kern0pt}A{\isacharparenright}{\kern0pt}\ {\isasymin}\ set\ hs\ {\isasymand}\ a\isactrlsub e\isactrlsub n\isactrlsub d\ {\isasymin}\ set\ A{\isachardoublequoteclose}\isanewline
%
\isadelimproof
\ \ \ \ %
\endisadelimproof
%
\isatagproof
\isacommand{using}\isamarkupfalse%
\ assms\isanewline
\ \ \isacommand{proof}\isamarkupfalse%
\ {\isacharminus}{\kern0pt}\isanewline
\ \ \ \ \isacommand{obtain}\isamarkupfalse%
\ {\isasympi}s\isactrlsub {\isadigit{1}}\ {\isasympi}s\isactrlsub {\isadigit{2}}\ \isakeyword{where}\ {\isachardoublequoteopen}{\isasympi}s\ {\isacharequal}{\kern0pt}\ {\isasympi}s\isactrlsub {\isadigit{1}}\ {\isacharat}{\kern0pt}\ {\isacharparenleft}{\kern0pt}{\isacharparenleft}{\kern0pt}t\isactrlsub {\isasympi}{\isacharcomma}{\kern0pt}{\isasympi}{\isacharparenright}{\kern0pt}{\isacharhash}{\kern0pt}{\isasympi}s\isactrlsub {\isadigit{2}}{\isacharparenright}{\kern0pt}{\isachardoublequoteclose}\isanewline
\ \ \ \ \ \ \isacommand{using}\isamarkupfalse%
\ assms\ dur{\isacharunderscore}{\kern0pt}acts{\isacharunderscore}{\kern0pt}subset\ \isacommand{by}\isamarkupfalse%
\ {\isacharparenleft}{\kern0pt}meson\ split{\isacharunderscore}{\kern0pt}list\ subsetD{\isacharparenright}{\kern0pt}\ \isanewline
\ \ \ \ \isacommand{then}\isamarkupfalse%
\ \isacommand{obtain}\isamarkupfalse%
\ hs{\isacharprime}{\kern0pt}\ \isakeyword{where}\ {\isachardoublequoteopen}hs{\isacharprime}{\kern0pt}\ {\isacharequal}{\kern0pt}\ simplify{\isacharunderscore}{\kern0pt}plan\ htps\ {\isacharparenleft}{\kern0pt}{\isacharparenleft}{\kern0pt}t\isactrlsub {\isasympi}{\isacharcomma}{\kern0pt}{\isasympi}{\isacharparenright}{\kern0pt}{\isacharhash}{\kern0pt}{\isasympi}s\isactrlsub {\isadigit{2}}{\isacharparenright}{\kern0pt}{\isachardoublequoteclose}\isanewline
\ \ \ \ \ \ \isacommand{by}\isamarkupfalse%
\ simp\isanewline
\ \ \ \ \isacommand{have}\isamarkupfalse%
\ {\isachardoublequoteopen}{\isacharparenleft}{\kern0pt}t\isactrlsub {\isasympi}{\isacharcomma}{\kern0pt}{\isasympi}{\isacharparenright}{\kern0pt}\ {\isacharequal}{\kern0pt}\ hd\ {\isacharparenleft}{\kern0pt}{\isacharparenleft}{\kern0pt}t\isactrlsub {\isasympi}{\isacharcomma}{\kern0pt}{\isasympi}{\isacharparenright}{\kern0pt}{\isacharhash}{\kern0pt}{\isasympi}s\isactrlsub {\isadigit{2}}{\isacharparenright}{\kern0pt}{\isachardoublequoteclose}\isanewline
\ \ \ \ \ \ \isacommand{by}\isamarkupfalse%
\ auto\isanewline
\ \ \ \ \isacommand{have}\isamarkupfalse%
\ {\isachardoublequoteopen}{\isacharparenleft}{\kern0pt}t\isactrlsub {\isasympi}{\isacharcomma}{\kern0pt}\ {\isasympi}{\isacharparenright}{\kern0pt}\ {\isasymin}\ durative{\isacharunderscore}{\kern0pt}acts\ {\isacharparenleft}{\kern0pt}{\isacharparenleft}{\kern0pt}t\isactrlsub {\isasympi}{\isacharcomma}{\kern0pt}\ {\isasympi}{\isacharparenright}{\kern0pt}\ {\isacharhash}{\kern0pt}\ {\isasympi}s\isactrlsub {\isadigit{2}}{\isacharparenright}{\kern0pt}{\isachardoublequoteclose}\isanewline
\ \ \ \ \ \ \isacommand{using}\isamarkupfalse%
\ assms\ {\isacartoucheopen}{\isasympi}s\ {\isacharequal}{\kern0pt}\ {\isasympi}s\isactrlsub {\isadigit{1}}\ {\isacharat}{\kern0pt}\ {\isacharparenleft}{\kern0pt}{\isacharparenleft}{\kern0pt}t\isactrlsub {\isasympi}{\isacharcomma}{\kern0pt}{\isasympi}{\isacharparenright}{\kern0pt}{\isacharhash}{\kern0pt}{\isasympi}s\isactrlsub {\isadigit{2}}{\isacharparenright}{\kern0pt}{\isacartoucheclose}\isanewline
\ \ \ \ \ \ \isacommand{unfolding}\isamarkupfalse%
\ durative{\isacharunderscore}{\kern0pt}acts{\isacharunderscore}{\kern0pt}def\ is{\isacharunderscore}{\kern0pt}act{\isacharunderscore}{\kern0pt}simple{\isacharunderscore}{\kern0pt}def\ \isacommand{by}\isamarkupfalse%
\ {\isacharparenleft}{\kern0pt}auto\ split{\isacharcolon}{\kern0pt}\ plan{\isacharunderscore}{\kern0pt}action{\isachardot}{\kern0pt}splits{\isacharparenright}{\kern0pt}\isanewline
\ \ \ \ \isacommand{have}\isamarkupfalse%
\ {\isachardoublequoteopen}{\isasymexists}as{\isachardot}{\kern0pt}\ {\isacharparenleft}{\kern0pt}t\isactrlsub {\isasympi}\ {\isacharplus}{\kern0pt}\ {\isacharparenleft}{\kern0pt}duration\ {\isasympi}{\isacharparenright}{\kern0pt}{\isacharcomma}{\kern0pt}as{\isacharparenright}{\kern0pt}\ {\isasymin}\ set\ hs{\isacharprime}{\kern0pt}\ {\isasymand}\ a\isactrlsub e\isactrlsub n\isactrlsub d\ {\isasymin}\ set\ as{\isachardoublequoteclose}\isanewline
\ \ \ \ \ \ \isacommand{using}\isamarkupfalse%
\ assms\ simp{\isacharunderscore}{\kern0pt}plan{\isacharunderscore}{\kern0pt}corr{\isacharunderscore}{\kern0pt}dur{\isacharunderscore}{\kern0pt}acts{\isacharunderscore}{\kern0pt}end{\isacharunderscore}{\kern0pt}aux{\isacharbrackleft}{\kern0pt}OF\ {\isacartoucheopen}hs{\isacharprime}{\kern0pt}\ {\isacharequal}{\kern0pt}\ simplify{\isacharunderscore}{\kern0pt}plan\ htps\ {\isacharparenleft}{\kern0pt}{\isacharparenleft}{\kern0pt}t\isactrlsub {\isasympi}{\isacharcomma}{\kern0pt}{\isasympi}{\isacharparenright}{\kern0pt}{\isacharhash}{\kern0pt}{\isasympi}s\isactrlsub {\isadigit{2}}{\isacharparenright}{\kern0pt}{\isacartoucheclose}\ \isanewline
\ \ \ \ \ \ \ \ \ \ \ \ \ \ \ \ \ \ \ \ \ \ \ \ \ \ \ \ \ \ \ \ \ \ \ \ \ \ \ \ \ \ \ \ \ \ \ \ \ \ \ \ \ {\isacartoucheopen}{\isacharparenleft}{\kern0pt}t\isactrlsub {\isasympi}{\isacharcomma}{\kern0pt}{\isasympi}{\isacharparenright}{\kern0pt}\ {\isacharequal}{\kern0pt}\ hd\ {\isacharparenleft}{\kern0pt}{\isacharparenleft}{\kern0pt}t\isactrlsub {\isasympi}{\isacharcomma}{\kern0pt}{\isasympi}{\isacharparenright}{\kern0pt}{\isacharhash}{\kern0pt}{\isasympi}s\isactrlsub {\isadigit{2}}{\isacharparenright}{\kern0pt}{\isacartoucheclose}\ \isanewline
\ \ \ \ \ \ \ \ \ \ \ \ \ \ \ \ \ \ \ \ \ \ \ \ \ \ \ \ \ \ \ \ \ \ \ \ \ \ \ \ \ \ \ \ \ \ \ \ \ \ \ \ \ {\isacartoucheopen}{\isacharparenleft}{\kern0pt}t\isactrlsub {\isasympi}{\isacharcomma}{\kern0pt}{\isasympi}{\isacharparenright}{\kern0pt}\ {\isasymin}\ durative{\isacharunderscore}{\kern0pt}acts\ {\isacharparenleft}{\kern0pt}{\isacharparenleft}{\kern0pt}t\isactrlsub {\isasympi}{\isacharcomma}{\kern0pt}{\isasympi}{\isacharparenright}{\kern0pt}{\isacharhash}{\kern0pt}{\isasympi}s\isactrlsub {\isadigit{2}}{\isacharparenright}{\kern0pt}{\isacartoucheclose}{\isacharbrackright}{\kern0pt}\ \isanewline
\ \ \ \ \ \ \isacommand{by}\isamarkupfalse%
\ simp\isanewline
\ \ \ \ \isacommand{have}\isamarkupfalse%
\ {\isachardoublequoteopen}hs{\isacharprime}{\kern0pt}\ {\isasymsubseteq}\isactrlsub h\isactrlsub s\ hs{\isachardoublequoteclose}\isanewline
\ \ \ \ \ \ \isacommand{using}\isamarkupfalse%
\ assms\ simplify{\isacharunderscore}{\kern0pt}plan{\isacharunderscore}{\kern0pt}happs{\isacharunderscore}{\kern0pt}subset{\isacharunderscore}{\kern0pt}app\ \isanewline
\ \ \ \ \ \ \ \ \ \ \ \ {\isacartoucheopen}hs{\isacharprime}{\kern0pt}\ {\isacharequal}{\kern0pt}\ simplify{\isacharunderscore}{\kern0pt}plan\ htps\ {\isacharparenleft}{\kern0pt}{\isacharparenleft}{\kern0pt}t\isactrlsub {\isasympi}{\isacharcomma}{\kern0pt}{\isasympi}{\isacharparenright}{\kern0pt}{\isacharhash}{\kern0pt}{\isasympi}s\isactrlsub {\isadigit{2}}{\isacharparenright}{\kern0pt}{\isacartoucheclose}\ \isanewline
\ \ \ \ \ \ \ \ \ \ \ \ {\isacartoucheopen}{\isasympi}s\ {\isacharequal}{\kern0pt}\ {\isasympi}s\isactrlsub {\isadigit{1}}\ {\isacharat}{\kern0pt}\ {\isacharparenleft}{\kern0pt}{\isacharparenleft}{\kern0pt}t\isactrlsub {\isasympi}{\isacharcomma}{\kern0pt}{\isasympi}{\isacharparenright}{\kern0pt}{\isacharhash}{\kern0pt}{\isasympi}s\isactrlsub {\isadigit{2}}{\isacharparenright}{\kern0pt}{\isacartoucheclose}\isanewline
\ \ \ \ \ \ \isacommand{by}\isamarkupfalse%
\ blast\isanewline
\ \ \ \ \isacommand{show}\isamarkupfalse%
\ {\isacharquery}{\kern0pt}thesis\isanewline
\ \ \ \ \ \ \isacommand{using}\isamarkupfalse%
\ {\isacartoucheopen}{\isasymexists}as{\isachardot}{\kern0pt}\ {\isacharparenleft}{\kern0pt}t\isactrlsub {\isasympi}\ {\isacharplus}{\kern0pt}\ {\isacharparenleft}{\kern0pt}duration\ {\isasympi}{\isacharparenright}{\kern0pt}{\isacharcomma}{\kern0pt}as{\isacharparenright}{\kern0pt}\ {\isasymin}\ set\ hs{\isacharprime}{\kern0pt}\ {\isasymand}\ a\isactrlsub e\isactrlsub n\isactrlsub d\ {\isasymin}\ set\ as{\isacartoucheclose}\isanewline
\ \ \ \ \ \ \ \ \ \ \ \ {\isacartoucheopen}hs{\isacharprime}{\kern0pt}\ {\isasymsubseteq}\isactrlsub h\isactrlsub s\ hs{\isacartoucheclose}\isanewline
\ \ \ \ \ \ \isacommand{unfolding}\isamarkupfalse%
\ happ{\isacharunderscore}{\kern0pt}seq{\isacharunderscore}{\kern0pt}subset{\isacharunderscore}{\kern0pt}def\ happ{\isacharunderscore}{\kern0pt}seq{\isacharunderscore}{\kern0pt}member{\isacharunderscore}{\kern0pt}def\ \isacommand{by}\isamarkupfalse%
\ blast\isanewline
\ \ \isacommand{qed}\isamarkupfalse%
%
\endisatagproof
{\isafoldproof}%
%
\isadelimproof
\isanewline
%
\endisadelimproof
\isanewline
\ \ \isanewline
\ \ \isacommand{lemma}\isamarkupfalse%
\ member{\isacharunderscore}{\kern0pt}map{\isacharcolon}{\kern0pt}\ {\isachardoublequoteopen}a\ {\isasymin}\ set\ A\ {\isasymLongrightarrow}\ f\ a\ {\isasymin}\ set\ {\isacharparenleft}{\kern0pt}map\ f\ A{\isacharparenright}{\kern0pt}{\isachardoublequoteclose}\isanewline
%
\isadelimproof
\ \ \ \ %
\endisadelimproof
%
\isatagproof
\isacommand{by}\isamarkupfalse%
\ simp%
\endisatagproof
{\isafoldproof}%
%
\isadelimproof
\isanewline
%
\endisadelimproof
\isanewline
\ \ \isacommand{lemma}\isamarkupfalse%
\ simp{\isacharunderscore}{\kern0pt}plan{\isacharunderscore}{\kern0pt}corr{\isacharunderscore}{\kern0pt}dur{\isacharunderscore}{\kern0pt}acts{\isacharunderscore}{\kern0pt}overall{\isacharunderscore}{\kern0pt}aux{\isacharcolon}{\kern0pt}\ \isanewline
\ \ \ \ \isakeyword{assumes}\ {\isachardoublequoteopen}strict{\isacharunderscore}{\kern0pt}sorted\ htps{\isachardoublequoteclose}\isanewline
\ \ \ \ \ \ \ \ \isakeyword{and}\ {\isachardoublequoteopen}hs\ {\isacharequal}{\kern0pt}\ simplify{\isacharunderscore}{\kern0pt}plan\ htps\ {\isasympi}s{\isachardoublequoteclose}\isanewline
\ \ \ \ \ \ \ \ \isakeyword{and}\ {\isachardoublequoteopen}{\isacharparenleft}{\kern0pt}t\isactrlsub {\isasympi}{\isacharcomma}{\kern0pt}{\isasympi}{\isacharparenright}{\kern0pt}\ {\isacharequal}{\kern0pt}\ hd\ {\isasympi}s{\isachardoublequoteclose}\ \isanewline
\ \ \ \ \ \ \ \ \isakeyword{and}\ {\isachardoublequoteopen}{\isacharparenleft}{\kern0pt}t\isactrlsub {\isasympi}{\isacharcomma}{\kern0pt}{\isasympi}{\isacharparenright}{\kern0pt}\ {\isasymin}\ durative{\isacharunderscore}{\kern0pt}acts\ {\isasympi}s{\isachardoublequoteclose}\isanewline
\ \ \ \ \isakeyword{shows}\ {\isachardoublequoteopen}{\isasymforall}t\isactrlsub i\ t\isactrlsub j{\isachardot}{\kern0pt}\ {\isacharparenleft}{\kern0pt}t\isactrlsub i{\isacharcomma}{\kern0pt}t\isactrlsub j{\isacharparenright}{\kern0pt}\ {\isasymin}\ set\ {\isacharparenleft}{\kern0pt}consec{\isacharunderscore}{\kern0pt}htps{\isacharunderscore}{\kern0pt}in{\isacharunderscore}{\kern0pt}interval\ htps\ t\isactrlsub {\isasympi}\ {\isacharparenleft}{\kern0pt}t\isactrlsub {\isasympi}\ {\isacharplus}{\kern0pt}\ duration\ {\isasympi}{\isacharparenright}{\kern0pt}{\isacharparenright}{\kern0pt}\ {\isasymlongrightarrow}\isanewline
\ \ \ \ \ \ \ \ \ \ \ \ {\isacharparenleft}{\kern0pt}{\isasymexists}{\isacharparenleft}{\kern0pt}t{\isacharprime}{\kern0pt}{\isacharcomma}{\kern0pt}A{\isacharparenright}{\kern0pt}\ {\isasymin}\ set\ hs{\isachardot}{\kern0pt}\ t\isactrlsub i\ {\isacharless}{\kern0pt}\ t{\isacharprime}{\kern0pt}\ {\isasymand}\ t{\isacharprime}{\kern0pt}\ {\isacharless}{\kern0pt}\ t\isactrlsub j\ {\isasymand}\ {\isacharparenleft}{\kern0pt}the\ {\isacharparenleft}{\kern0pt}res{\isacharunderscore}{\kern0pt}inst{\isacharunderscore}{\kern0pt}snap{\isacharunderscore}{\kern0pt}action\ {\isasympi}\ Over{\isacharunderscore}{\kern0pt}All{\isacharparenright}{\kern0pt}{\isacharparenright}{\kern0pt}\ {\isasymin}\ set\ A{\isacharparenright}{\kern0pt}{\isachardoublequoteclose}\isanewline
%
\isadelimproof
\ \ \ \ %
\endisadelimproof
%
\isatagproof
\isacommand{using}\isamarkupfalse%
\ assms\isanewline
\ \ \ \ \isacommand{unfolding}\isamarkupfalse%
\ durative{\isacharunderscore}{\kern0pt}acts{\isacharunderscore}{\kern0pt}def\ is{\isacharunderscore}{\kern0pt}act{\isacharunderscore}{\kern0pt}simple{\isacharunderscore}{\kern0pt}def\isanewline
\ \ \isacommand{proof}\isamarkupfalse%
\ {\isacharparenleft}{\kern0pt}induction\ {\isasympi}s{\isacharparenright}{\kern0pt}\isanewline
\ \ \ \ \isacommand{case}\isamarkupfalse%
\ Nil\isanewline
\ \ \ \ \isacommand{then}\isamarkupfalse%
\ \isacommand{show}\isamarkupfalse%
\ {\isacharquery}{\kern0pt}case\ \isanewline
\ \ \ \ \ \ \isacommand{by}\isamarkupfalse%
\ simp\isanewline
\ \ \isacommand{next}\isamarkupfalse%
\isanewline
\ \ \ \ \isacommand{case}\isamarkupfalse%
\ {\isacharparenleft}{\kern0pt}Cons\ {\isasympi}\ {\isasympi}s{\isacharparenright}{\kern0pt}\isanewline
\ \ \ \ \isacommand{fix}\isamarkupfalse%
\ t\isactrlsub i\ t\isactrlsub j\isanewline
\ \ \ \ \isacommand{obtain}\isamarkupfalse%
\ t\isactrlsub {\isasympi}\ {\isasympi}{\isacharprime}{\kern0pt}\ \isakeyword{where}\ {\isachardoublequoteopen}{\isasympi}\ {\isacharequal}{\kern0pt}\ {\isacharparenleft}{\kern0pt}t\isactrlsub {\isasympi}{\isacharcomma}{\kern0pt}{\isasympi}{\isacharprime}{\kern0pt}{\isacharparenright}{\kern0pt}{\isachardoublequoteclose}\isanewline
\ \ \ \ \ \ \isacommand{by}\isamarkupfalse%
\ {\isacharparenleft}{\kern0pt}meson\ surj{\isacharunderscore}{\kern0pt}pair{\isacharparenright}{\kern0pt}\isanewline
\ \ \ \ \isacommand{then}\isamarkupfalse%
\ \isacommand{show}\isamarkupfalse%
\ {\isacharquery}{\kern0pt}case\isanewline
\ \ \ \ \isacommand{proof}\isamarkupfalse%
\ {\isacharparenleft}{\kern0pt}cases\ {\isasympi}{\isacharprime}{\kern0pt}{\isacharparenright}{\kern0pt}\isanewline
\ \ \ \ \ \ \isacommand{case}\isamarkupfalse%
\ {\isacharparenleft}{\kern0pt}Simple{\isacharunderscore}{\kern0pt}Plan{\isacharunderscore}{\kern0pt}Action\ n\ args{\isacharparenright}{\kern0pt}\isanewline
\ \ \ \ \ \ \isacommand{then}\isamarkupfalse%
\ \isacommand{show}\isamarkupfalse%
\ {\isacharquery}{\kern0pt}thesis\ \isanewline
\ \ \ \ \ \ \ \ \isacommand{using}\isamarkupfalse%
\ Cons{\isachardot}{\kern0pt}prems\ {\isacartoucheopen}{\isasympi}\ {\isacharequal}{\kern0pt}\ {\isacharparenleft}{\kern0pt}t\isactrlsub {\isasympi}{\isacharcomma}{\kern0pt}{\isasympi}{\isacharprime}{\kern0pt}{\isacharparenright}{\kern0pt}{\isacartoucheclose}\ \isacommand{by}\isamarkupfalse%
\ simp\isanewline
\ \ \ \ \ \ \isacommand{next}\isamarkupfalse%
\isanewline
\ \ \ \ \ \ \ \ \isacommand{case}\isamarkupfalse%
\ {\isacharparenleft}{\kern0pt}Durative{\isacharunderscore}{\kern0pt}Plan{\isacharunderscore}{\kern0pt}Action\ n\ args\ d{\isacharparenright}{\kern0pt}\isanewline
\ \ \ \ \ \ \ \ \isacommand{let}\isamarkupfalse%
\ {\isacharquery}{\kern0pt}to{\isacharunderscore}{\kern0pt}happ\ {\isacharequal}{\kern0pt}\ {\isachardoublequoteopen}{\isasymlambda}{\isacharparenleft}{\kern0pt}t\isactrlsub a{\isacharcomma}{\kern0pt}a{\isacharparenright}{\kern0pt}{\isachardot}{\kern0pt}\ {\isacharparenleft}{\kern0pt}t\isactrlsub a{\isacharcomma}{\kern0pt}{\isacharbrackleft}{\kern0pt}a{\isacharbrackright}{\kern0pt}{\isacharparenright}{\kern0pt}{\isachardoublequoteclose}\isanewline
\ \ \ \ \ \ \ \ \isacommand{let}\isamarkupfalse%
\ {\isacharquery}{\kern0pt}snap{\isacharunderscore}{\kern0pt}act\ {\isacharequal}{\kern0pt}\ {\isachardoublequoteopen}inst{\isacharunderscore}{\kern0pt}snap{\isacharunderscore}{\kern0pt}action\ {\isacharparenleft}{\kern0pt}the\ {\isacharparenleft}{\kern0pt}resolve{\isacharunderscore}{\kern0pt}action{\isacharunderscore}{\kern0pt}schema\ n{\isacharparenright}{\kern0pt}{\isacharparenright}{\kern0pt}\ args{\isachardoublequoteclose}\isanewline
\ \ \ \ \ \ \ \ \isacommand{let}\isamarkupfalse%
\ {\isacharquery}{\kern0pt}chtps{\isacharunderscore}{\kern0pt}in{\isacharunderscore}{\kern0pt}intv\ {\isacharequal}{\kern0pt}\ {\isachardoublequoteopen}consec{\isacharunderscore}{\kern0pt}htps{\isacharunderscore}{\kern0pt}in{\isacharunderscore}{\kern0pt}interval\ htps\ t\isactrlsub {\isasympi}\ {\isacharparenleft}{\kern0pt}t\isactrlsub {\isasympi}\ {\isacharplus}{\kern0pt}\ duration\ {\isasympi}{\isacharprime}{\kern0pt}{\isacharparenright}{\kern0pt}{\isachardoublequoteclose}\isanewline
\ \ \ \ \ \ \ \ \isacommand{let}\isamarkupfalse%
\ {\isacharquery}{\kern0pt}hs\isactrlsub {\isadigit{1}}\ {\isacharequal}{\kern0pt}\ {\isachardoublequoteopen}map\ {\isacharparenleft}{\kern0pt}{\isasymlambda}{\isacharparenleft}{\kern0pt}t\isactrlsub i{\isacharcomma}{\kern0pt}t\isactrlsub j{\isacharparenright}{\kern0pt}{\isachardot}{\kern0pt}{\isacharparenleft}{\kern0pt}{\isacharparenleft}{\kern0pt}t\isactrlsub i{\isacharplus}{\kern0pt}t\isactrlsub j{\isacharparenright}{\kern0pt}\ {\isacharslash}{\kern0pt}\ {\isadigit{2}}{\isacharcomma}{\kern0pt}{\isacharquery}{\kern0pt}snap{\isacharunderscore}{\kern0pt}act\ Over{\isacharunderscore}{\kern0pt}All{\isacharparenright}{\kern0pt}{\isacharparenright}{\kern0pt}\ {\isacharquery}{\kern0pt}chtps{\isacharunderscore}{\kern0pt}in{\isacharunderscore}{\kern0pt}intv{\isachardoublequoteclose}\isanewline
\ \ \ \ \ \ \ \ \isacommand{have}\isamarkupfalse%
\ {\isadigit{1}}{\isacharcolon}{\kern0pt}\ {\isachardoublequoteopen}{\isasymforall}{\isacharparenleft}{\kern0pt}t\isactrlsub a{\isacharcomma}{\kern0pt}A{\isacharparenright}{\kern0pt}\ {\isasymin}\ set\ {\isacharparenleft}{\kern0pt}map\ {\isacharquery}{\kern0pt}to{\isacharunderscore}{\kern0pt}happ\ {\isacharquery}{\kern0pt}hs\isactrlsub {\isadigit{1}}{\isacharparenright}{\kern0pt}{\isachardot}{\kern0pt}\ {\isasymexists}A{\isacharprime}{\kern0pt}{\isachardot}{\kern0pt}\ {\isacharparenleft}{\kern0pt}t\isactrlsub a{\isacharcomma}{\kern0pt}A{\isacharprime}{\kern0pt}{\isacharparenright}{\kern0pt}\ {\isasymin}\ set\ {\isacharparenleft}{\kern0pt}simplify{\isacharunderscore}{\kern0pt}plan\ htps\ {\isacharparenleft}{\kern0pt}{\isasympi}{\isacharhash}{\kern0pt}{\isasympi}s{\isacharparenright}{\kern0pt}{\isacharparenright}{\kern0pt}\ {\isasymand}\ set\ A\ {\isasymsubseteq}\ set\ A{\isacharprime}{\kern0pt}{\isachardoublequoteclose}\isanewline
\ \ \ \ \ \ \ \ \ \ \isacommand{using}\isamarkupfalse%
\ {\isacartoucheopen}{\isasympi}\ {\isacharequal}{\kern0pt}\ {\isacharparenleft}{\kern0pt}t\isactrlsub {\isasympi}{\isacharcomma}{\kern0pt}{\isasympi}{\isacharprime}{\kern0pt}{\isacharparenright}{\kern0pt}{\isacartoucheclose}\ Durative{\isacharunderscore}{\kern0pt}Plan{\isacharunderscore}{\kern0pt}Action\ insort{\isacharunderscore}{\kern0pt}happ{\isadigit{2}}{\isacharunderscore}{\kern0pt}insort{\isacharunderscore}{\kern0pt}mult{\isacharunderscore}{\kern0pt}happs{\isacharunderscore}{\kern0pt}action{\isacharunderscore}{\kern0pt}subset\isanewline
\ \ \ \ \ \ \ \ \ \ \isacommand{by}\isamarkupfalse%
\ {\isacharparenleft}{\kern0pt}auto\ simp{\isacharcolon}{\kern0pt}\ Let{\isacharunderscore}{\kern0pt}def{\isacharparenright}{\kern0pt}\isanewline
\ \ \ \ \ \ \ \ \isacommand{have}\isamarkupfalse%
\ {\isachardoublequoteopen}{\isasymforall}{\isacharparenleft}{\kern0pt}t\isactrlsub i{\isacharcomma}{\kern0pt}\ t\isactrlsub j{\isacharparenright}{\kern0pt}\ {\isasymin}\ set\ {\isacharquery}{\kern0pt}chtps{\isacharunderscore}{\kern0pt}in{\isacharunderscore}{\kern0pt}intv{\isachardot}{\kern0pt}\ {\isacharparenleft}{\kern0pt}{\isacharparenleft}{\kern0pt}t\isactrlsub i{\isacharplus}{\kern0pt}t\isactrlsub j{\isacharparenright}{\kern0pt}\ {\isacharslash}{\kern0pt}\ {\isadigit{2}}{\isacharcomma}{\kern0pt}{\isacharquery}{\kern0pt}snap{\isacharunderscore}{\kern0pt}act\ Over{\isacharunderscore}{\kern0pt}All{\isacharparenright}{\kern0pt}\ {\isasymin}\ set\ {\isacharquery}{\kern0pt}hs\isactrlsub {\isadigit{1}}{\isachardoublequoteclose}\isanewline
\ \ \ \ \ \ \ \ \ \ \isacommand{by}\isamarkupfalse%
\ auto\isanewline
\ \ \ \ \ \ \ \ \isacommand{then}\isamarkupfalse%
\ \isacommand{have}\isamarkupfalse%
\ {\isachardoublequoteopen}{\isasymforall}{\isacharparenleft}{\kern0pt}t\isactrlsub i{\isacharcomma}{\kern0pt}\ t\isactrlsub j{\isacharparenright}{\kern0pt}\ {\isasymin}\ set\ {\isacharquery}{\kern0pt}chtps{\isacharunderscore}{\kern0pt}in{\isacharunderscore}{\kern0pt}intv{\isachardot}{\kern0pt}\ \isanewline
\ \ \ \ \ \ \ \ \ \ \ \ \ \ \ \ \ \ \ \ {\isacharquery}{\kern0pt}to{\isacharunderscore}{\kern0pt}happ\ {\isacharparenleft}{\kern0pt}{\isacharparenleft}{\kern0pt}t\isactrlsub i{\isacharplus}{\kern0pt}t\isactrlsub j{\isacharparenright}{\kern0pt}\ {\isacharslash}{\kern0pt}\ {\isadigit{2}}{\isacharcomma}{\kern0pt}{\isacharquery}{\kern0pt}snap{\isacharunderscore}{\kern0pt}act\ Over{\isacharunderscore}{\kern0pt}All{\isacharparenright}{\kern0pt}\ {\isasymin}\ set\ {\isacharparenleft}{\kern0pt}map\ {\isacharquery}{\kern0pt}to{\isacharunderscore}{\kern0pt}happ\ {\isacharquery}{\kern0pt}hs\isactrlsub {\isadigit{1}}{\isacharparenright}{\kern0pt}{\isachardoublequoteclose}\isanewline
\ \ \ \ \ \ \ \ \ \ \isacommand{by}\isamarkupfalse%
\ {\isacharparenleft}{\kern0pt}smt\ member{\isacharunderscore}{\kern0pt}map\ case{\isacharunderscore}{\kern0pt}prodI{\isadigit{2}}\ case{\isacharunderscore}{\kern0pt}prod{\isacharunderscore}{\kern0pt}conv{\isacharparenright}{\kern0pt}\ \isanewline
\ \ \ \ \ \ \ \ \isacommand{then}\isamarkupfalse%
\ \isacommand{have}\isamarkupfalse%
\ {\isachardoublequoteopen}{\isasymforall}{\isacharparenleft}{\kern0pt}t\isactrlsub i{\isacharcomma}{\kern0pt}\ t\isactrlsub j{\isacharparenright}{\kern0pt}\ {\isasymin}\ set\ {\isacharquery}{\kern0pt}chtps{\isacharunderscore}{\kern0pt}in{\isacharunderscore}{\kern0pt}intv{\isachardot}{\kern0pt}\ \isanewline
\ \ \ \ \ \ \ \ \ \ \ \ \ \ \ \ \ \ \ \ {\isacharparenleft}{\kern0pt}{\isasymexists}A{\isachardot}{\kern0pt}\ {\isacharparenleft}{\kern0pt}{\isacharparenleft}{\kern0pt}t\isactrlsub i{\isacharplus}{\kern0pt}t\isactrlsub j{\isacharparenright}{\kern0pt}\ {\isacharslash}{\kern0pt}\ {\isadigit{2}}{\isacharcomma}{\kern0pt}A{\isacharparenright}{\kern0pt}\ {\isasymin}\ set\ hs\ {\isasymand}\ \isanewline
\ \ \ \ \ \ \ \ \ \ \ \ \ \ \ \ \ \ \ \ \ \ t\isactrlsub i\ {\isacharless}{\kern0pt}\ {\isacharparenleft}{\kern0pt}t\isactrlsub i{\isacharplus}{\kern0pt}t\isactrlsub j{\isacharparenright}{\kern0pt}\ {\isacharslash}{\kern0pt}\ {\isadigit{2}}\ {\isasymand}\ {\isacharparenleft}{\kern0pt}t\isactrlsub i{\isacharplus}{\kern0pt}t\isactrlsub j{\isacharparenright}{\kern0pt}\ {\isacharslash}{\kern0pt}\ {\isadigit{2}}\ {\isacharless}{\kern0pt}\ t\isactrlsub j\ {\isasymand}\ \isanewline
\ \ \ \ \ \ \ \ \ \ \ \ \ \ \ \ \ \ \ \ \ \ \ \ {\isacharparenleft}{\kern0pt}the\ {\isacharparenleft}{\kern0pt}res{\isacharunderscore}{\kern0pt}inst{\isacharunderscore}{\kern0pt}snap{\isacharunderscore}{\kern0pt}action\ {\isasympi}{\isacharprime}{\kern0pt}\ Over{\isacharunderscore}{\kern0pt}All{\isacharparenright}{\kern0pt}{\isacharparenright}{\kern0pt}\ {\isasymin}\ set\ A{\isacharparenright}{\kern0pt}{\isachardoublequoteclose}\isanewline
\ \ \ \ \ \ \ \ \ \ \isacommand{using}\isamarkupfalse%
\ {\isadigit{1}}\ Durative{\isacharunderscore}{\kern0pt}Plan{\isacharunderscore}{\kern0pt}Action\ {\isacartoucheopen}hs\ {\isacharequal}{\kern0pt}\ simplify{\isacharunderscore}{\kern0pt}plan\ htps\ {\isacharparenleft}{\kern0pt}{\isasympi}{\isacharhash}{\kern0pt}{\isasympi}s{\isacharparenright}{\kern0pt}{\isacartoucheclose}\ \isanewline
\ \ \ \ \ \ \ \ \ \ \ \ \ \ \ \ sorted{\isacharunderscore}{\kern0pt}distinct{\isacharunderscore}{\kern0pt}zip{\isacharunderscore}{\kern0pt}le{\isacharbrackleft}{\kern0pt}OF\ {\isacartoucheopen}strict{\isacharunderscore}{\kern0pt}sorted\ htps{\isacartoucheclose}{\isacharbrackright}{\kern0pt}\isanewline
\ \ \ \ \ \ \ \ \ \ \isacommand{by}\isamarkupfalse%
\ {\isacharparenleft}{\kern0pt}auto\ simp{\isacharcolon}{\kern0pt}\ consec{\isacharunderscore}{\kern0pt}htps{\isacharunderscore}{\kern0pt}in{\isacharunderscore}{\kern0pt}interval{\isacharunderscore}{\kern0pt}def{\isacharparenright}{\kern0pt}\isanewline
\ \ \ \ \ \ \ \ \isacommand{then}\isamarkupfalse%
\ \isacommand{have}\isamarkupfalse%
\ {\isachardoublequoteopen}{\isasymforall}t\isactrlsub i\ t\isactrlsub j{\isachardot}{\kern0pt}\ {\isacharparenleft}{\kern0pt}t\isactrlsub i{\isacharcomma}{\kern0pt}\ t\isactrlsub j{\isacharparenright}{\kern0pt}\ {\isasymin}\ set\ {\isacharquery}{\kern0pt}chtps{\isacharunderscore}{\kern0pt}in{\isacharunderscore}{\kern0pt}intv\ {\isasymlongrightarrow}\ \isanewline
\ \ \ \ \ \ \ \ \ \ \ \ \ \ \ \ \ \ \ \ {\isacharparenleft}{\kern0pt}{\isasymexists}{\isacharparenleft}{\kern0pt}t{\isacharprime}{\kern0pt}{\isacharcomma}{\kern0pt}A{\isacharparenright}{\kern0pt}\ {\isasymin}\ set\ hs{\isachardot}{\kern0pt}\ t\isactrlsub i\ {\isacharless}{\kern0pt}\ t{\isacharprime}{\kern0pt}\ {\isasymand}\ t{\isacharprime}{\kern0pt}\ {\isacharless}{\kern0pt}\ t\isactrlsub j\ {\isasymand}\ \isanewline
\ \ \ \ \ \ \ \ \ \ \ \ \ \ \ \ \ \ \ \ \ \ {\isacharparenleft}{\kern0pt}the\ {\isacharparenleft}{\kern0pt}res{\isacharunderscore}{\kern0pt}inst{\isacharunderscore}{\kern0pt}snap{\isacharunderscore}{\kern0pt}action\ {\isasympi}{\isacharprime}{\kern0pt}\ Over{\isacharunderscore}{\kern0pt}All{\isacharparenright}{\kern0pt}{\isacharparenright}{\kern0pt}\ {\isasymin}\ set\ A{\isacharparenright}{\kern0pt}{\isachardoublequoteclose}\isanewline
\ \ \ \ \ \ \ \ \ \ \isacommand{by}\isamarkupfalse%
\ fastforce\isanewline
\ \ \ \ \ \ \ \ \isacommand{then}\isamarkupfalse%
\ \isacommand{show}\isamarkupfalse%
\ {\isacharquery}{\kern0pt}thesis\isanewline
\ \ \ \ \ \ \ \ \ \ \isacommand{using}\isamarkupfalse%
\ Cons{\isachardot}{\kern0pt}prems\ {\isacartoucheopen}{\isasympi}\ {\isacharequal}{\kern0pt}\ {\isacharparenleft}{\kern0pt}t\isactrlsub {\isasympi}{\isacharcomma}{\kern0pt}{\isasympi}{\isacharprime}{\kern0pt}{\isacharparenright}{\kern0pt}{\isacartoucheclose}\ Durative{\isacharunderscore}{\kern0pt}Plan{\isacharunderscore}{\kern0pt}Action\ \isacommand{by}\isamarkupfalse%
\ simp\isanewline
\ \ \ \ \isacommand{qed}\isamarkupfalse%
\isanewline
\ \ \isacommand{qed}\isamarkupfalse%
%
\endisatagproof
{\isafoldproof}%
%
\isadelimproof
%
\endisadelimproof
%
\begin{isamarkuptext}%
A simplified plan contains all \isa{Over{\isacharunderscore}{\kern0pt}All}-snap-actions.%
\end{isamarkuptext}\isamarkuptrue%
\ \ \isacommand{lemma}\isamarkupfalse%
\ {\isacharparenleft}{\kern0pt}\isakeyword{in}\ wf{\isacharunderscore}{\kern0pt}ast{\isacharunderscore}{\kern0pt}problem{\isacharparenright}{\kern0pt}\ simp{\isacharunderscore}{\kern0pt}plan{\isacharunderscore}{\kern0pt}corr{\isacharunderscore}{\kern0pt}dur{\isacharunderscore}{\kern0pt}acts{\isacharunderscore}{\kern0pt}overall{\isacharcolon}{\kern0pt}\isanewline
\ \ \ \ \isakeyword{fixes}\ {\isasympi}\ {\isasympi}s\isanewline
\ \ \ \ \isakeyword{defines}\ {\isachardoublequoteopen}htps\ {\isasymequiv}\ htps{\isacharunderscore}{\kern0pt}exec\ {\isasympi}s{\isachardoublequoteclose}\ \isanewline
\ \ \ \ \ \ \ \ \isakeyword{and}\ {\isachardoublequoteopen}a\isactrlsub i\isactrlsub n\isactrlsub v\ {\isasymequiv}\ the\ {\isacharparenleft}{\kern0pt}res{\isacharunderscore}{\kern0pt}inst{\isacharunderscore}{\kern0pt}snap{\isacharunderscore}{\kern0pt}action\ {\isasympi}\ Over{\isacharunderscore}{\kern0pt}All{\isacharparenright}{\kern0pt}{\isachardoublequoteclose}\isanewline
\ \ \ \ \isakeyword{assumes}\ {\isachardoublequoteopen}wf{\isacharunderscore}{\kern0pt}plan\ {\isasympi}s{\isachardoublequoteclose}\ \isanewline
\ \ \ \ \ \ \ \ \isakeyword{and}\ {\isachardoublequoteopen}hs\ {\isacharequal}{\kern0pt}\ simplify{\isacharunderscore}{\kern0pt}plan\ htps\ {\isasympi}s{\isachardoublequoteclose}\ \isanewline
\ \ \ \ \ \ \ \ \isakeyword{and}\ {\isachardoublequoteopen}{\isacharparenleft}{\kern0pt}t\isactrlsub {\isasympi}{\isacharcomma}{\kern0pt}{\isasympi}{\isacharparenright}{\kern0pt}\ {\isasymin}\ durative{\isacharunderscore}{\kern0pt}acts\ {\isasympi}s{\isachardoublequoteclose}\isanewline
\ \ \ \ \isakeyword{shows}\ {\isachardoublequoteopen}{\isasymforall}t\isactrlsub i\ t\isactrlsub j{\isachardot}{\kern0pt}\ {\isacharparenleft}{\kern0pt}consec{\isacharunderscore}{\kern0pt}htps\ {\isasympi}s\ t\isactrlsub i\ t\isactrlsub j\ {\isasymand}\ t\isactrlsub {\isasympi}\ {\isasymle}\ t\isactrlsub i\ {\isasymand}\ t\isactrlsub j\ {\isasymle}\ t\isactrlsub {\isasympi}\ {\isacharplus}{\kern0pt}\ duration\ {\isasympi}{\isacharparenright}{\kern0pt}\ {\isasymlongrightarrow}\ \isanewline
\ \ \ \ \ \ \ \ \ \ \ \ {\isacharparenleft}{\kern0pt}{\isasymexists}{\isacharparenleft}{\kern0pt}t{\isacharprime}{\kern0pt}{\isacharcomma}{\kern0pt}A{\isacharparenright}{\kern0pt}\ {\isasymin}\ set\ hs{\isachardot}{\kern0pt}\ t\isactrlsub i\ {\isacharless}{\kern0pt}\ t{\isacharprime}{\kern0pt}\ {\isasymand}\ t{\isacharprime}{\kern0pt}\ {\isacharless}{\kern0pt}\ t\isactrlsub j\ {\isasymand}\ a\isactrlsub i\isactrlsub n\isactrlsub v\ {\isasymin}\ set\ A{\isacharparenright}{\kern0pt}{\isachardoublequoteclose}\isanewline
\ \ \ \ \isanewline
%
\isadelimproof
\ \ \ \ %
\endisadelimproof
%
\isatagproof
\isacommand{using}\isamarkupfalse%
\ assms\isanewline
\ \ \isacommand{proof}\isamarkupfalse%
\ {\isacharminus}{\kern0pt}\isanewline
\ \ \ \ \isacommand{have}\isamarkupfalse%
\ {\isachardoublequoteopen}strict{\isacharunderscore}{\kern0pt}sorted\ htps{\isachardoublequoteclose}\isanewline
\ \ \ \ \ \ \isacommand{using}\isamarkupfalse%
\ assms\ htps{\isacharunderscore}{\kern0pt}exec{\isacharunderscore}{\kern0pt}sorted\ htps{\isacharunderscore}{\kern0pt}exec{\isacharunderscore}{\kern0pt}distinct\ \isacommand{by}\isamarkupfalse%
\ {\isacharparenleft}{\kern0pt}simp\ add{\isacharcolon}{\kern0pt}\ strict{\isacharunderscore}{\kern0pt}sorted{\isacharunderscore}{\kern0pt}iff{\isacharparenright}{\kern0pt}\isanewline
\ \ \ \ \isacommand{have}\isamarkupfalse%
\ {\isachardoublequoteopen}t\isactrlsub {\isasympi}\ {\isasymle}\ t\isactrlsub {\isasympi}\ {\isacharplus}{\kern0pt}\ duration\ {\isasympi}{\isachardoublequoteclose}\isanewline
\ \ \ \ \ \ \isacommand{using}\isamarkupfalse%
\ assms\ \isacommand{unfolding}\isamarkupfalse%
\ wf{\isacharunderscore}{\kern0pt}plan{\isacharunderscore}{\kern0pt}def\ durative{\isacharunderscore}{\kern0pt}acts{\isacharunderscore}{\kern0pt}def\ is{\isacharunderscore}{\kern0pt}act{\isacharunderscore}{\kern0pt}simple{\isacharunderscore}{\kern0pt}def\isanewline
\ \ \ \ \ \ \isacommand{by}\isamarkupfalse%
\ {\isacharparenleft}{\kern0pt}auto\ simp{\isacharcolon}{\kern0pt}\ Let{\isacharunderscore}{\kern0pt}def\ split{\isacharcolon}{\kern0pt}\ option{\isachardot}{\kern0pt}splits\ plan{\isacharunderscore}{\kern0pt}action{\isachardot}{\kern0pt}splits\ ast{\isacharunderscore}{\kern0pt}action{\isacharunderscore}{\kern0pt}schema{\isachardot}{\kern0pt}splits{\isacharparenright}{\kern0pt}\isanewline
\ \ \ \ \isacommand{obtain}\isamarkupfalse%
\ {\isasympi}s\isactrlsub {\isadigit{1}}\ {\isasympi}s\isactrlsub {\isadigit{2}}\ \isakeyword{where}\ {\isachardoublequoteopen}{\isasympi}s\ {\isacharequal}{\kern0pt}\ {\isasympi}s\isactrlsub {\isadigit{1}}\ {\isacharat}{\kern0pt}\ {\isacharparenleft}{\kern0pt}{\isacharparenleft}{\kern0pt}t\isactrlsub {\isasympi}{\isacharcomma}{\kern0pt}{\isasympi}{\isacharparenright}{\kern0pt}{\isacharhash}{\kern0pt}{\isasympi}s\isactrlsub {\isadigit{2}}{\isacharparenright}{\kern0pt}{\isachardoublequoteclose}\isanewline
\ \ \ \ \ \ \isacommand{using}\isamarkupfalse%
\ assms\ dur{\isacharunderscore}{\kern0pt}acts{\isacharunderscore}{\kern0pt}subset\ \isacommand{by}\isamarkupfalse%
\ {\isacharparenleft}{\kern0pt}meson\ split{\isacharunderscore}{\kern0pt}list\ subsetD{\isacharparenright}{\kern0pt}\ \isanewline
\ \ \ \ \isacommand{then}\isamarkupfalse%
\ \isacommand{have}\isamarkupfalse%
\ {\isachardoublequoteopen}{\isacharparenleft}{\kern0pt}t\isactrlsub {\isasympi}{\isacharcomma}{\kern0pt}{\isasympi}{\isacharparenright}{\kern0pt}\ {\isacharequal}{\kern0pt}\ hd\ {\isacharparenleft}{\kern0pt}{\isacharparenleft}{\kern0pt}t\isactrlsub {\isasympi}{\isacharcomma}{\kern0pt}{\isasympi}{\isacharparenright}{\kern0pt}{\isacharhash}{\kern0pt}{\isasympi}s\isactrlsub {\isadigit{2}}{\isacharparenright}{\kern0pt}{\isachardoublequoteclose}\ \isakeyword{and}\ {\isachardoublequoteopen}{\isacharparenleft}{\kern0pt}t\isactrlsub {\isasympi}{\isacharcomma}{\kern0pt}\ {\isasympi}{\isacharparenright}{\kern0pt}\ {\isasymin}\ durative{\isacharunderscore}{\kern0pt}acts\ {\isacharparenleft}{\kern0pt}{\isacharparenleft}{\kern0pt}t\isactrlsub {\isasympi}{\isacharcomma}{\kern0pt}\ {\isasympi}{\isacharparenright}{\kern0pt}\ {\isacharhash}{\kern0pt}\ {\isasympi}s\isactrlsub {\isadigit{2}}{\isacharparenright}{\kern0pt}{\isachardoublequoteclose}\isanewline
\ \ \ \ \ \ \isacommand{using}\isamarkupfalse%
\ assms\ \isacommand{unfolding}\isamarkupfalse%
\ durative{\isacharunderscore}{\kern0pt}acts{\isacharunderscore}{\kern0pt}def\ is{\isacharunderscore}{\kern0pt}act{\isacharunderscore}{\kern0pt}simple{\isacharunderscore}{\kern0pt}def\ \isacommand{by}\isamarkupfalse%
\ {\isacharparenleft}{\kern0pt}auto\ split{\isacharcolon}{\kern0pt}\ plan{\isacharunderscore}{\kern0pt}action{\isachardot}{\kern0pt}splits{\isacharparenright}{\kern0pt}\isanewline
\ \ \ \ \isacommand{obtain}\isamarkupfalse%
\ hs{\isacharprime}{\kern0pt}\ \isakeyword{where}\ {\isachardoublequoteopen}hs{\isacharprime}{\kern0pt}\ {\isacharequal}{\kern0pt}\ simplify{\isacharunderscore}{\kern0pt}plan\ htps\ {\isacharparenleft}{\kern0pt}{\isacharparenleft}{\kern0pt}t\isactrlsub {\isasympi}{\isacharcomma}{\kern0pt}{\isasympi}{\isacharparenright}{\kern0pt}{\isacharhash}{\kern0pt}{\isasympi}s\isactrlsub {\isadigit{2}}{\isacharparenright}{\kern0pt}{\isachardoublequoteclose}\isanewline
\ \ \ \ \ \ \isacommand{by}\isamarkupfalse%
\ simp\isanewline
\ \ \ \ \isacommand{then}\isamarkupfalse%
\ \isacommand{have}\isamarkupfalse%
\ {\isachardoublequoteopen}hs{\isacharprime}{\kern0pt}\ {\isasymsubseteq}\isactrlsub h\isactrlsub s\ hs{\isachardoublequoteclose}\isanewline
\ \ \ \ \ \ \isacommand{using}\isamarkupfalse%
\ assms\ simplify{\isacharunderscore}{\kern0pt}plan{\isacharunderscore}{\kern0pt}happs{\isacharunderscore}{\kern0pt}subset{\isacharunderscore}{\kern0pt}app\ {\isacartoucheopen}{\isasympi}s\ {\isacharequal}{\kern0pt}\ {\isasympi}s\isactrlsub {\isadigit{1}}\ {\isacharat}{\kern0pt}\ {\isacharparenleft}{\kern0pt}{\isacharparenleft}{\kern0pt}t\isactrlsub {\isasympi}{\isacharcomma}{\kern0pt}{\isasympi}{\isacharparenright}{\kern0pt}{\isacharhash}{\kern0pt}{\isasympi}s\isactrlsub {\isadigit{2}}{\isacharparenright}{\kern0pt}{\isacartoucheclose}\ \isacommand{by}\isamarkupfalse%
\ blast\isanewline
\ \ \ \ \isacommand{have}\isamarkupfalse%
\ {\isachardoublequoteopen}{\isasymforall}t\isactrlsub i\ t\isactrlsub j{\isachardot}{\kern0pt}\ {\isacharparenleft}{\kern0pt}t\isactrlsub i{\isacharcomma}{\kern0pt}t\isactrlsub j{\isacharparenright}{\kern0pt}\ {\isasymin}\ set\ {\isacharparenleft}{\kern0pt}consec{\isacharunderscore}{\kern0pt}htps{\isacharunderscore}{\kern0pt}in{\isacharunderscore}{\kern0pt}interval\ htps\ t\isactrlsub {\isasympi}\ {\isacharparenleft}{\kern0pt}t\isactrlsub {\isasympi}\ {\isacharplus}{\kern0pt}\ {\isacharparenleft}{\kern0pt}duration\ {\isasympi}{\isacharparenright}{\kern0pt}{\isacharparenright}{\kern0pt}{\isacharparenright}{\kern0pt}\ \isanewline
\ \ \ \ \ \ \ \ \ \ {\isasymlongrightarrow}\ {\isacharparenleft}{\kern0pt}{\isasymexists}{\isacharparenleft}{\kern0pt}t{\isacharprime}{\kern0pt}{\isacharcomma}{\kern0pt}A{\isacharparenright}{\kern0pt}\ {\isasymin}\ set\ hs{\isacharprime}{\kern0pt}{\isachardot}{\kern0pt}\ t\isactrlsub i\ {\isacharless}{\kern0pt}\ t{\isacharprime}{\kern0pt}\ {\isasymand}\ t{\isacharprime}{\kern0pt}\ {\isacharless}{\kern0pt}\ t\isactrlsub j\ {\isasymand}\ a\isactrlsub i\isactrlsub n\isactrlsub v\ {\isasymin}\ set\ A{\isacharparenright}{\kern0pt}{\isachardoublequoteclose}\isanewline
\ \ \ \ \ \ \isacommand{using}\isamarkupfalse%
\ assms\ simp{\isacharunderscore}{\kern0pt}plan{\isacharunderscore}{\kern0pt}corr{\isacharunderscore}{\kern0pt}dur{\isacharunderscore}{\kern0pt}acts{\isacharunderscore}{\kern0pt}overall{\isacharunderscore}{\kern0pt}aux{\isacharbrackleft}{\kern0pt}OF\ {\isacartoucheopen}strict{\isacharunderscore}{\kern0pt}sorted\ htps{\isacartoucheclose}\isanewline
\ \ \ \ \ \ \ \ \ \ \ \ \ \ \ \ \ \ \ \ \ \ \ \ \ \ \ \ \ \ \ \ \ \ \ \ \ \ \ \ \ \ \ \ \ \ \ \ \ \ \ \ \ \ \ \ \ {\isacartoucheopen}hs{\isacharprime}{\kern0pt}\ {\isacharequal}{\kern0pt}\ simplify{\isacharunderscore}{\kern0pt}plan\ htps\ {\isacharparenleft}{\kern0pt}{\isacharparenleft}{\kern0pt}t\isactrlsub {\isasympi}{\isacharcomma}{\kern0pt}{\isasympi}{\isacharparenright}{\kern0pt}{\isacharhash}{\kern0pt}{\isasympi}s\isactrlsub {\isadigit{2}}{\isacharparenright}{\kern0pt}{\isacartoucheclose}\ \isanewline
\ \ \ \ \ \ \ \ \ \ \ \ \ \ \ \ \ \ \ \ \ \ \ \ \ \ \ \ \ \ \ \ \ \ \ \ \ \ \ \ \ \ \ \ \ \ \ \ \ \ \ \ \ \ \ \ \ {\isacartoucheopen}{\isacharparenleft}{\kern0pt}t\isactrlsub {\isasympi}{\isacharcomma}{\kern0pt}{\isasympi}{\isacharparenright}{\kern0pt}\ {\isacharequal}{\kern0pt}\ hd\ {\isacharparenleft}{\kern0pt}{\isacharparenleft}{\kern0pt}t\isactrlsub {\isasympi}{\isacharcomma}{\kern0pt}{\isasympi}{\isacharparenright}{\kern0pt}{\isacharhash}{\kern0pt}{\isasympi}s\isactrlsub {\isadigit{2}}{\isacharparenright}{\kern0pt}{\isacartoucheclose}\ \isanewline
\ \ \ \ \ \ \ \ \ \ \ \ \ \ \ \ \ \ \ \ \ \ \ \ \ \ \ \ \ \ \ \ \ \ \ \ \ \ \ \ \ \ \ \ \ \ \ \ \ \ \ \ \ \ \ \ \ {\isacartoucheopen}{\isacharparenleft}{\kern0pt}t\isactrlsub {\isasympi}{\isacharcomma}{\kern0pt}{\isasympi}{\isacharparenright}{\kern0pt}\ {\isasymin}\ durative{\isacharunderscore}{\kern0pt}acts\ {\isacharparenleft}{\kern0pt}{\isacharparenleft}{\kern0pt}t\isactrlsub {\isasympi}{\isacharcomma}{\kern0pt}{\isasympi}{\isacharparenright}{\kern0pt}{\isacharhash}{\kern0pt}{\isasympi}s\isactrlsub {\isadigit{2}}{\isacharparenright}{\kern0pt}{\isacartoucheclose}{\isacharbrackright}{\kern0pt}\isanewline
\ \ \ \ \ \ \isacommand{by}\isamarkupfalse%
\ auto\isanewline
\ \ \ \ \isacommand{then}\isamarkupfalse%
\ \isacommand{have}\isamarkupfalse%
\ {\isachardoublequoteopen}{\isasymforall}t\isactrlsub i\ t\isactrlsub j{\isachardot}{\kern0pt}\ {\isacharparenleft}{\kern0pt}consec{\isacharunderscore}{\kern0pt}htps\ {\isasympi}s\ t\isactrlsub i\ t\isactrlsub j\ {\isasymand}\ t\isactrlsub {\isasympi}\ {\isasymle}\ t\isactrlsub i\ {\isasymand}\ t\isactrlsub j\ {\isasymle}\ t\isactrlsub {\isasympi}\ {\isacharplus}{\kern0pt}\ {\isacharparenleft}{\kern0pt}duration\ {\isasympi}{\isacharparenright}{\kern0pt}{\isacharparenright}{\kern0pt}\ \isanewline
\ \ \ \ \ \ \ \ \ \ \ \ \ \ {\isasymlongrightarrow}\ {\isacharparenleft}{\kern0pt}{\isasymexists}{\isacharparenleft}{\kern0pt}t{\isacharprime}{\kern0pt}{\isacharcomma}{\kern0pt}as{\isacharparenright}{\kern0pt}\ {\isasymin}\ set\ hs{\isacharprime}{\kern0pt}{\isachardot}{\kern0pt}\ t\isactrlsub i\ {\isacharless}{\kern0pt}\ t{\isacharprime}{\kern0pt}\ {\isasymand}\ t{\isacharprime}{\kern0pt}\ {\isacharless}{\kern0pt}\ t\isactrlsub j\ {\isasymand}\ a\isactrlsub i\isactrlsub n\isactrlsub v\ {\isasymin}\ set\ as{\isacharparenright}{\kern0pt}{\isachardoublequoteclose}\isanewline
\ \ \ \ \ \ \isacommand{using}\isamarkupfalse%
\ htps{\isacharunderscore}{\kern0pt}def\ consec{\isacharunderscore}{\kern0pt}htps{\isacharunderscore}{\kern0pt}in{\isacharunderscore}{\kern0pt}interval{\isacharunderscore}{\kern0pt}iff{\isacharbrackleft}{\kern0pt}OF\ {\isacartoucheopen}t\isactrlsub {\isasympi}\ {\isasymle}\ t\isactrlsub {\isasympi}\ {\isacharplus}{\kern0pt}\ {\isacharparenleft}{\kern0pt}duration\ {\isasympi}{\isacharparenright}{\kern0pt}{\isacartoucheclose}{\isacharbrackright}{\kern0pt}\ \isacommand{by}\isamarkupfalse%
\ auto\isanewline
\ \ \ \ \isacommand{then}\isamarkupfalse%
\ \isacommand{show}\isamarkupfalse%
\ {\isacharquery}{\kern0pt}thesis\isanewline
\ \ \ \ \ \ \isacommand{using}\isamarkupfalse%
\ {\isacartoucheopen}hs{\isacharprime}{\kern0pt}\ {\isasymsubseteq}\isactrlsub h\isactrlsub s\ hs{\isacartoucheclose}\ \isacommand{unfolding}\isamarkupfalse%
\ happ{\isacharunderscore}{\kern0pt}seq{\isacharunderscore}{\kern0pt}subset{\isacharunderscore}{\kern0pt}def\ happ{\isacharunderscore}{\kern0pt}seq{\isacharunderscore}{\kern0pt}member{\isacharunderscore}{\kern0pt}def\ \isacommand{by}\isamarkupfalse%
\ blast\isanewline
\ \ \isacommand{qed}\isamarkupfalse%
%
\endisatagproof
{\isafoldproof}%
%
\isadelimproof
%
\endisadelimproof
%
\begin{isamarkuptext}%
Flatten a happening into individual timed grounded actions. This is used to simplify proofs.%
\end{isamarkuptext}\isamarkuptrue%
\ \ \isacommand{fun}\isamarkupfalse%
\ set{\isacharunderscore}{\kern0pt}happ{\isacharunderscore}{\kern0pt}acts\ {\isacharcolon}{\kern0pt}{\isacharcolon}{\kern0pt}\ {\isachardoublequoteopen}happening\ list\ {\isasymRightarrow}\ {\isacharparenleft}{\kern0pt}time\ {\isasymtimes}\ ground{\isacharunderscore}{\kern0pt}action{\isacharparenright}{\kern0pt}\ set{\isachardoublequoteclose}\ \isakeyword{where}\isanewline
\ \ \ \ {\isachardoublequoteopen}set{\isacharunderscore}{\kern0pt}happ{\isacharunderscore}{\kern0pt}acts\ {\isacharbrackleft}{\kern0pt}{\isacharbrackright}{\kern0pt}\ {\isacharequal}{\kern0pt}\ {\isacharbraceleft}{\kern0pt}{\isacharbraceright}{\kern0pt}{\isachardoublequoteclose}\isanewline
\ \ {\isacharbar}{\kern0pt}\ {\isachardoublequoteopen}set{\isacharunderscore}{\kern0pt}happ{\isacharunderscore}{\kern0pt}acts\ {\isacharparenleft}{\kern0pt}{\isacharparenleft}{\kern0pt}t\isactrlsub a{\isacharcomma}{\kern0pt}A{\isacharparenright}{\kern0pt}{\isacharhash}{\kern0pt}hs{\isacharparenright}{\kern0pt}\ {\isacharequal}{\kern0pt}\ {\isacharbraceleft}{\kern0pt}{\isacharparenleft}{\kern0pt}t\isactrlsub a{\isacharprime}{\kern0pt}{\isacharcomma}{\kern0pt}a{\isacharprime}{\kern0pt}{\isacharparenright}{\kern0pt}{\isachardot}{\kern0pt}\ t\isactrlsub a\ {\isacharequal}{\kern0pt}\ t\isactrlsub a{\isacharprime}{\kern0pt}\ {\isasymand}\ a{\isacharprime}{\kern0pt}\ {\isasymin}\ set\ A{\isacharbraceright}{\kern0pt}\ {\isasymunion}\ set{\isacharunderscore}{\kern0pt}happ{\isacharunderscore}{\kern0pt}acts\ hs{\isachardoublequoteclose}\isanewline
\isanewline
\ \ \isacommand{lemma}\isamarkupfalse%
\ set{\isacharunderscore}{\kern0pt}happ{\isacharunderscore}{\kern0pt}acts{\isacharunderscore}{\kern0pt}refine{\isacharcolon}{\kern0pt}\isanewline
\ \ \ \ {\isachardoublequoteopen}{\isacharparenleft}{\kern0pt}{\isasymexists}A{\isachardot}{\kern0pt}\ {\isacharparenleft}{\kern0pt}t\isactrlsub a{\isacharcomma}{\kern0pt}A{\isacharparenright}{\kern0pt}\ {\isasymin}\ set\ hs\ {\isasymand}\ a\ {\isasymin}\ set\ A{\isacharparenright}{\kern0pt}\ {\isasymlongleftrightarrow}\ {\isacharparenleft}{\kern0pt}{\isacharparenleft}{\kern0pt}t\isactrlsub a{\isacharcomma}{\kern0pt}a{\isacharparenright}{\kern0pt}\ {\isasymin}\ set{\isacharunderscore}{\kern0pt}happ{\isacharunderscore}{\kern0pt}acts\ hs{\isacharparenright}{\kern0pt}{\isachardoublequoteclose}\isanewline
%
\isadelimproof
\ \ \ \ %
\endisadelimproof
%
\isatagproof
\isacommand{by}\isamarkupfalse%
\ {\isacharparenleft}{\kern0pt}induction\ hs{\isacharparenright}{\kern0pt}\ auto%
\endisatagproof
{\isafoldproof}%
%
\isadelimproof
\isanewline
%
\endisadelimproof
\isanewline
\ \ \isacommand{lemma}\isamarkupfalse%
\ set{\isacharunderscore}{\kern0pt}happ{\isacharunderscore}{\kern0pt}acts{\isacharunderscore}{\kern0pt}insort{\isacharunderscore}{\kern0pt}happ{\isacharcolon}{\kern0pt}\ \isanewline
\ \ \ \ {\isachardoublequoteopen}set{\isacharunderscore}{\kern0pt}happ{\isacharunderscore}{\kern0pt}acts\ {\isacharparenleft}{\kern0pt}insort{\isacharunderscore}{\kern0pt}happ\ h\ hs{\isacharparenright}{\kern0pt}\ {\isacharequal}{\kern0pt}\ set{\isacharunderscore}{\kern0pt}happ{\isacharunderscore}{\kern0pt}acts\ {\isacharbrackleft}{\kern0pt}h{\isacharbrackright}{\kern0pt}\ {\isasymunion}\ set{\isacharunderscore}{\kern0pt}happ{\isacharunderscore}{\kern0pt}acts\ hs{\isachardoublequoteclose}\isanewline
%
\isadelimproof
\ \ \ \ %
\endisadelimproof
%
\isatagproof
\isacommand{by}\isamarkupfalse%
\ {\isacharparenleft}{\kern0pt}induction\ h\ hs\ rule{\isacharcolon}{\kern0pt}\ insort{\isacharunderscore}{\kern0pt}happ{\isachardot}{\kern0pt}induct{\isacharparenright}{\kern0pt}\ auto%
\endisatagproof
{\isafoldproof}%
%
\isadelimproof
\isanewline
%
\endisadelimproof
\isanewline
\ \ \isacommand{lemma}\isamarkupfalse%
\ set{\isacharunderscore}{\kern0pt}happ{\isacharunderscore}{\kern0pt}acts{\isacharunderscore}{\kern0pt}insort{\isacharunderscore}{\kern0pt}mult{\isacharunderscore}{\kern0pt}happs{\isacharcolon}{\kern0pt}\ \isanewline
\ \ \ \ {\isachardoublequoteopen}set{\isacharunderscore}{\kern0pt}happ{\isacharunderscore}{\kern0pt}acts\ {\isacharparenleft}{\kern0pt}insort{\isacharunderscore}{\kern0pt}mult{\isacharunderscore}{\kern0pt}happs\ hs\isactrlsub {\isadigit{1}}\ hs\isactrlsub {\isadigit{2}}{\isacharparenright}{\kern0pt}\ {\isacharequal}{\kern0pt}\ set{\isacharunderscore}{\kern0pt}happ{\isacharunderscore}{\kern0pt}acts\ hs\isactrlsub {\isadigit{1}}\ {\isasymunion}\ set{\isacharunderscore}{\kern0pt}happ{\isacharunderscore}{\kern0pt}acts\ hs\isactrlsub {\isadigit{2}}{\isachardoublequoteclose}\isanewline
%
\isadelimproof
\ \ \ \ %
\endisadelimproof
%
\isatagproof
\isacommand{using}\isamarkupfalse%
\ set{\isacharunderscore}{\kern0pt}happ{\isacharunderscore}{\kern0pt}acts{\isacharunderscore}{\kern0pt}insort{\isacharunderscore}{\kern0pt}happ\isanewline
\ \ \ \ \isacommand{by}\isamarkupfalse%
\ {\isacharparenleft}{\kern0pt}induction\ hs\isactrlsub {\isadigit{1}}{\isacharparenright}{\kern0pt}\ auto%
\endisatagproof
{\isafoldproof}%
%
\isadelimproof
\isanewline
%
\endisadelimproof
\isanewline
\ \ \isacommand{lemma}\isamarkupfalse%
\ set{\isacharunderscore}{\kern0pt}happ{\isacharunderscore}{\kern0pt}acts{\isacharunderscore}{\kern0pt}to{\isacharunderscore}{\kern0pt}happ{\isacharcolon}{\kern0pt}\ {\isachardoublequoteopen}set{\isacharunderscore}{\kern0pt}happ{\isacharunderscore}{\kern0pt}acts\ {\isacharparenleft}{\kern0pt}map\ {\isacharparenleft}{\kern0pt}{\isasymlambda}{\isacharparenleft}{\kern0pt}t\isactrlsub a{\isacharcomma}{\kern0pt}a{\isacharparenright}{\kern0pt}{\isachardot}{\kern0pt}\ {\isacharparenleft}{\kern0pt}t\isactrlsub a{\isacharcomma}{\kern0pt}{\isacharbrackleft}{\kern0pt}a{\isacharbrackright}{\kern0pt}{\isacharparenright}{\kern0pt}{\isacharparenright}{\kern0pt}\ xs{\isacharparenright}{\kern0pt}\ {\isacharequal}{\kern0pt}\ set\ xs{\isachardoublequoteclose}\isanewline
%
\isadelimproof
\ \ %
\endisadelimproof
%
\isatagproof
\isacommand{proof}\isamarkupfalse%
\ {\isacharparenleft}{\kern0pt}induction\ xs{\isacharparenright}{\kern0pt}\isanewline
\ \ \ \ \isacommand{case}\isamarkupfalse%
\ Nil\isanewline
\ \ \ \ \isacommand{then}\isamarkupfalse%
\ \isacommand{show}\isamarkupfalse%
\ {\isacharquery}{\kern0pt}case\ \isanewline
\ \ \ \ \ \ \isacommand{by}\isamarkupfalse%
\ auto\isanewline
\ \ \isacommand{next}\isamarkupfalse%
\isanewline
\ \ \ \ \isacommand{case}\isamarkupfalse%
\ {\isacharparenleft}{\kern0pt}Cons\ x\ xs{\isacharparenright}{\kern0pt}\isanewline
\ \ \ \ \isacommand{then}\isamarkupfalse%
\ \isacommand{show}\isamarkupfalse%
\ {\isacharquery}{\kern0pt}case\ \isanewline
\ \ \ \ \ \ \isacommand{by}\isamarkupfalse%
\ {\isacharparenleft}{\kern0pt}cases\ x{\isacharparenright}{\kern0pt}\ {\isacharparenleft}{\kern0pt}auto\ split{\isacharcolon}{\kern0pt}\ prod{\isachardot}{\kern0pt}splits{\isacharparenright}{\kern0pt}\isanewline
\ \ \isacommand{qed}\isamarkupfalse%
%
\endisatagproof
{\isafoldproof}%
%
\isadelimproof
\isanewline
%
\endisadelimproof
\isanewline
\ \ \isacommand{lemma}\isamarkupfalse%
\ simp{\isacharunderscore}{\kern0pt}plan{\isacharunderscore}{\kern0pt}corr{\isacharunderscore}{\kern0pt}inst{\isacharunderscore}{\kern0pt}plan{\isacharunderscore}{\kern0pt}act{\isacharunderscore}{\kern0pt}aux{\isacharcolon}{\kern0pt}\isanewline
\ \ \ \ \isakeyword{assumes}\ {\isachardoublequoteopen}strict{\isacharunderscore}{\kern0pt}sorted\ htps{\isachardoublequoteclose}\isanewline
\ \ \ \ \ \ \ \ \isakeyword{and}\ {\isachardoublequoteopen}hs\ {\isacharequal}{\kern0pt}\ simplify{\isacharunderscore}{\kern0pt}plan\ htps\ {\isasympi}s{\isachardoublequoteclose}\ \isanewline
\ \ \ \ \isakeyword{shows}\ {\isachardoublequoteopen}{\isasymforall}{\isacharparenleft}{\kern0pt}t\isactrlsub a{\isacharcomma}{\kern0pt}a{\isacharparenright}{\kern0pt}\ {\isasymin}\ set{\isacharunderscore}{\kern0pt}happ{\isacharunderscore}{\kern0pt}acts\ hs{\isachardot}{\kern0pt}\ {\isasymexists}{\isacharparenleft}{\kern0pt}t\isactrlsub {\isasympi}{\isacharcomma}{\kern0pt}{\isasympi}{\isacharparenright}{\kern0pt}\ {\isasymin}\ set\ {\isasympi}s{\isachardot}{\kern0pt}\ {\isacharparenleft}{\kern0pt}t\isactrlsub a{\isacharcomma}{\kern0pt}a{\isacharparenright}{\kern0pt}\ {\isasymin}\ set\ {\isacharparenleft}{\kern0pt}simplify{\isacharunderscore}{\kern0pt}action\ htps\ {\isacharparenleft}{\kern0pt}t\isactrlsub {\isasympi}{\isacharcomma}{\kern0pt}{\isasympi}{\isacharparenright}{\kern0pt}{\isacharparenright}{\kern0pt}{\isachardoublequoteclose}\isanewline
%
\isadelimproof
\ \ \ \ %
\endisadelimproof
%
\isatagproof
\isacommand{using}\isamarkupfalse%
\ assms\ set{\isacharunderscore}{\kern0pt}happ{\isacharunderscore}{\kern0pt}acts{\isacharunderscore}{\kern0pt}insort{\isacharunderscore}{\kern0pt}mult{\isacharunderscore}{\kern0pt}happs\ set{\isacharunderscore}{\kern0pt}happ{\isacharunderscore}{\kern0pt}acts{\isacharunderscore}{\kern0pt}to{\isacharunderscore}{\kern0pt}happ\isanewline
\ \ \ \ \isacommand{by}\isamarkupfalse%
\ {\isacharparenleft}{\kern0pt}induction\ {\isasympi}s\ arbitrary{\isacharcolon}{\kern0pt}\ hs{\isacharparenright}{\kern0pt}\ auto%
\endisatagproof
{\isafoldproof}%
%
\isadelimproof
%
\endisadelimproof
%
\begin{isamarkuptext}%
Every ground action in a simplified plan is the instantiation of a plan action.%
\end{isamarkuptext}\isamarkuptrue%
\ \ \isacommand{lemma}\isamarkupfalse%
\ {\isacharparenleft}{\kern0pt}\isakeyword{in}\ wf{\isacharunderscore}{\kern0pt}ast{\isacharunderscore}{\kern0pt}problem{\isacharparenright}{\kern0pt}\ simp{\isacharunderscore}{\kern0pt}plan{\isacharunderscore}{\kern0pt}corr{\isacharunderscore}{\kern0pt}inst{\isacharunderscore}{\kern0pt}plan{\isacharunderscore}{\kern0pt}act{\isacharcolon}{\kern0pt}\ \isanewline
\ \ \ \ \isakeyword{fixes}\ {\isasympi}s\isanewline
\ \ \ \ \isakeyword{defines}\ {\isachardoublequoteopen}htps\ {\isasymequiv}\ htps{\isacharunderscore}{\kern0pt}exec\ {\isasympi}s{\isachardoublequoteclose}\isanewline
\ \ \ \ \isakeyword{assumes}\ {\isachardoublequoteopen}wf{\isacharunderscore}{\kern0pt}plan\ {\isasympi}s{\isachardoublequoteclose}\ \isanewline
\ \ \ \ \ \ \ \ \isakeyword{and}\ {\isachardoublequoteopen}hs\ {\isacharequal}{\kern0pt}\ simplify{\isacharunderscore}{\kern0pt}plan\ htps\ {\isasympi}s{\isachardoublequoteclose}\ \isanewline
\ \ \ \ \ \ \ \ \isakeyword{and}\ {\isachardoublequoteopen}{\isacharparenleft}{\kern0pt}t\isactrlsub a{\isacharcomma}{\kern0pt}A{\isacharparenright}{\kern0pt}\ {\isasymin}\ set\ hs{\isachardoublequoteclose}\ \isanewline
\ \ \ \ \ \ \ \ \isakeyword{and}\ {\isachardoublequoteopen}a\ {\isasymin}\ set\ A{\isachardoublequoteclose}\isanewline
\ \ \ \ \isakeyword{shows}\ {\isachardoublequoteopen}{\isasymexists}{\isacharparenleft}{\kern0pt}t\isactrlsub {\isasympi}{\isacharcomma}{\kern0pt}{\isasympi}{\isacharparenright}{\kern0pt}\ {\isasymin}\ set\ {\isasympi}s{\isachardot}{\kern0pt}\ inst{\isacharunderscore}{\kern0pt}of{\isacharunderscore}{\kern0pt}plan{\isacharunderscore}{\kern0pt}action\ {\isasympi}s\ {\isacharparenleft}{\kern0pt}t\isactrlsub {\isasympi}{\isacharcomma}{\kern0pt}{\isasympi}{\isacharparenright}{\kern0pt}\ {\isacharparenleft}{\kern0pt}t\isactrlsub a{\isacharcomma}{\kern0pt}a{\isacharparenright}{\kern0pt}{\isachardoublequoteclose}\isanewline
%
\isadelimproof
\ \ %
\endisadelimproof
%
\isatagproof
\isacommand{proof}\isamarkupfalse%
\ {\isacharminus}{\kern0pt}\isanewline
\ \ \ \ \isacommand{have}\isamarkupfalse%
\ {\isachardoublequoteopen}strict{\isacharunderscore}{\kern0pt}sorted\ htps{\isachardoublequoteclose}\isanewline
\ \ \ \ \ \ \isacommand{using}\isamarkupfalse%
\ assms\ htps{\isacharunderscore}{\kern0pt}exec{\isacharunderscore}{\kern0pt}sorted\ htps{\isacharunderscore}{\kern0pt}exec{\isacharunderscore}{\kern0pt}distinct\ \isacommand{by}\isamarkupfalse%
\ {\isacharparenleft}{\kern0pt}simp\ add{\isacharcolon}{\kern0pt}\ strict{\isacharunderscore}{\kern0pt}sorted{\isacharunderscore}{\kern0pt}iff{\isacharparenright}{\kern0pt}\isanewline
\ \ \ \ \isacommand{have}\isamarkupfalse%
\ {\isachardoublequoteopen}{\isasymexists}{\isacharparenleft}{\kern0pt}t\isactrlsub {\isasympi}{\isacharcomma}{\kern0pt}{\isasympi}{\isacharparenright}{\kern0pt}\ {\isasymin}\ set\ {\isasympi}s{\isachardot}{\kern0pt}\ {\isacharparenleft}{\kern0pt}t\isactrlsub a{\isacharcomma}{\kern0pt}a{\isacharparenright}{\kern0pt}\ {\isasymin}\ set\ {\isacharparenleft}{\kern0pt}simplify{\isacharunderscore}{\kern0pt}action\ htps\ {\isacharparenleft}{\kern0pt}t\isactrlsub {\isasympi}{\isacharcomma}{\kern0pt}{\isasympi}{\isacharparenright}{\kern0pt}{\isacharparenright}{\kern0pt}{\isachardoublequoteclose}\isanewline
\ \ \ \ \ \ \isacommand{using}\isamarkupfalse%
\ assms\isanewline
\ \ \ \ \ \ \ \ \ \ \ \ simp{\isacharunderscore}{\kern0pt}plan{\isacharunderscore}{\kern0pt}corr{\isacharunderscore}{\kern0pt}inst{\isacharunderscore}{\kern0pt}plan{\isacharunderscore}{\kern0pt}act{\isacharunderscore}{\kern0pt}aux{\isacharbrackleft}{\kern0pt}OF\ {\isacartoucheopen}strict{\isacharunderscore}{\kern0pt}sorted\ htps{\isacartoucheclose}\ {\isacartoucheopen}hs\ {\isacharequal}{\kern0pt}\ simplify{\isacharunderscore}{\kern0pt}plan\ htps\ {\isasympi}s{\isacartoucheclose}{\isacharbrackright}{\kern0pt}\isanewline
\ \ \ \ \ \ \ \ \ \ \ \ set{\isacharunderscore}{\kern0pt}happ{\isacharunderscore}{\kern0pt}acts{\isacharunderscore}{\kern0pt}refine\ \isacommand{by}\isamarkupfalse%
\ blast\isanewline
\ \ \ \ \isacommand{then}\isamarkupfalse%
\ \isacommand{obtain}\isamarkupfalse%
\ t\isactrlsub {\isasympi}\ {\isasympi}\ \isakeyword{where}\ {\isachardoublequoteopen}{\isacharparenleft}{\kern0pt}t\isactrlsub {\isasympi}{\isacharcomma}{\kern0pt}{\isasympi}{\isacharparenright}{\kern0pt}\ {\isasymin}\ set\ {\isasympi}s{\isachardoublequoteclose}\ \isakeyword{and}\ {\isachardoublequoteopen}{\isacharparenleft}{\kern0pt}t\isactrlsub a{\isacharcomma}{\kern0pt}a{\isacharparenright}{\kern0pt}\ {\isasymin}\ set\ {\isacharparenleft}{\kern0pt}simplify{\isacharunderscore}{\kern0pt}action\ htps\ {\isacharparenleft}{\kern0pt}t\isactrlsub {\isasympi}{\isacharcomma}{\kern0pt}{\isasympi}{\isacharparenright}{\kern0pt}{\isacharparenright}{\kern0pt}{\isachardoublequoteclose}\isanewline
\ \ \ \ \ \ \isacommand{by}\isamarkupfalse%
\ auto\isanewline
\ \ \ \ \isacommand{show}\isamarkupfalse%
\ {\isacharquery}{\kern0pt}thesis\isanewline
\ \ \ \ \ \ \isacommand{using}\isamarkupfalse%
\ assms\ {\isacartoucheopen}{\isacharparenleft}{\kern0pt}t\isactrlsub {\isasympi}{\isacharcomma}{\kern0pt}{\isasympi}{\isacharparenright}{\kern0pt}\ {\isasymin}\ set\ {\isasympi}s{\isacartoucheclose}\ {\isacartoucheopen}{\isacharparenleft}{\kern0pt}t\isactrlsub a{\isacharcomma}{\kern0pt}a{\isacharparenright}{\kern0pt}\ {\isasymin}\ set\ {\isacharparenleft}{\kern0pt}simplify{\isacharunderscore}{\kern0pt}action\ htps\ {\isacharparenleft}{\kern0pt}t\isactrlsub {\isasympi}{\isacharcomma}{\kern0pt}{\isasympi}{\isacharparenright}{\kern0pt}{\isacharparenright}{\kern0pt}{\isacartoucheclose}\isanewline
\ \ \ \ \ \ \ \ \ \ \ \ simp{\isacharunderscore}{\kern0pt}act{\isacharunderscore}{\kern0pt}inst{\isacharunderscore}{\kern0pt}of{\isacharunderscore}{\kern0pt}plan{\isacharunderscore}{\kern0pt}act\ \isacommand{by}\isamarkupfalse%
\ blast\isanewline
\ \ \isacommand{qed}\isamarkupfalse%
%
\endisatagproof
{\isafoldproof}%
%
\isadelimproof
%
\endisadelimproof
%
\begin{isamarkuptext}%
Finally the correctness proof for \isa{simplify{\isacharunderscore}{\kern0pt}plan}%
\end{isamarkuptext}\isamarkuptrue%
\ \ \isacommand{lemma}\isamarkupfalse%
\ {\isacharparenleft}{\kern0pt}\isakeyword{in}\ wf{\isacharunderscore}{\kern0pt}ast{\isacharunderscore}{\kern0pt}problem{\isacharparenright}{\kern0pt}\ simplify{\isacharunderscore}{\kern0pt}plan{\isacharunderscore}{\kern0pt}correct{\isacharcolon}{\kern0pt}\ \isanewline
\ \ \ \ \isakeyword{fixes}\ {\isasympi}s\isanewline
\ \ \ \ \isakeyword{defines}\ {\isachardoublequoteopen}htps\ {\isasymequiv}\ htps{\isacharunderscore}{\kern0pt}exec\ {\isasympi}s{\isachardoublequoteclose}\isanewline
\ \ \ \ \isakeyword{assumes}\ {\isachardoublequoteopen}wf{\isacharunderscore}{\kern0pt}plan\ {\isasympi}s{\isachardoublequoteclose}\isanewline
\ \ \ \ \isakeyword{shows}\ {\isachardoublequoteopen}ind{\isacharunderscore}{\kern0pt}happ{\isacharunderscore}{\kern0pt}seq\ {\isasympi}s\ {\isacharparenleft}{\kern0pt}simplify{\isacharunderscore}{\kern0pt}plan\ htps\ {\isasympi}s{\isacharparenright}{\kern0pt}{\isachardoublequoteclose}\isanewline
%
\isadelimproof
\ \ \ \ %
\endisadelimproof
%
\isatagproof
\isacommand{unfolding}\isamarkupfalse%
\ ind{\isacharunderscore}{\kern0pt}happ{\isacharunderscore}{\kern0pt}seq{\isacharunderscore}{\kern0pt}def\ \ \ \ \ \ \ \ \ \ \ \ \isanewline
\ \ \ \ \isacommand{using}\isamarkupfalse%
\ assms\isanewline
\ \ \ \ \ \ \ \ \ \ simp{\isacharunderscore}{\kern0pt}plan{\isacharunderscore}{\kern0pt}happ{\isacharunderscore}{\kern0pt}not{\isacharunderscore}{\kern0pt}Nil\isanewline
\ \ \ \ \ \ \ \ \ \ simp{\isacharunderscore}{\kern0pt}plan{\isacharunderscore}{\kern0pt}strict{\isacharunderscore}{\kern0pt}sorted\ \isanewline
\ \ \ \ \ \ \ \ \ \ simp{\isacharunderscore}{\kern0pt}plan{\isacharunderscore}{\kern0pt}corr{\isacharunderscore}{\kern0pt}simp{\isacharunderscore}{\kern0pt}acts\ \isanewline
\ \ \ \ \ \ \ \ \ \ simp{\isacharunderscore}{\kern0pt}plan{\isacharunderscore}{\kern0pt}corr{\isacharunderscore}{\kern0pt}dur{\isacharunderscore}{\kern0pt}acts{\isacharunderscore}{\kern0pt}start\isanewline
\ \ \ \ \ \ \ \ \ \ simp{\isacharunderscore}{\kern0pt}plan{\isacharunderscore}{\kern0pt}corr{\isacharunderscore}{\kern0pt}dur{\isacharunderscore}{\kern0pt}acts{\isacharunderscore}{\kern0pt}end\isanewline
\ \ \ \ \ \ \ \ \ \ simp{\isacharunderscore}{\kern0pt}plan{\isacharunderscore}{\kern0pt}corr{\isacharunderscore}{\kern0pt}dur{\isacharunderscore}{\kern0pt}acts{\isacharunderscore}{\kern0pt}overall\isanewline
\ \ \ \ \ \ \ \ \ \ simp{\isacharunderscore}{\kern0pt}plan{\isacharunderscore}{\kern0pt}corr{\isacharunderscore}{\kern0pt}inst{\isacharunderscore}{\kern0pt}plan{\isacharunderscore}{\kern0pt}act\isanewline
\ \ \ \ \isacommand{by}\isamarkupfalse%
\ auto%
\endisatagproof
{\isafoldproof}%
%
\isadelimproof
\isanewline
%
\endisadelimproof
\isanewline
\ \ \isanewline
\ \ \isacommand{lemma}\isamarkupfalse%
\ {\isacharparenleft}{\kern0pt}\isakeyword{in}\ wf{\isacharunderscore}{\kern0pt}ast{\isacharunderscore}{\kern0pt}problem{\isacharparenright}{\kern0pt}\isanewline
\ \ \ \ \isakeyword{assumes}\ {\isachardoublequoteopen}wf{\isacharunderscore}{\kern0pt}plan\ {\isasympi}s{\isachardoublequoteclose}\isanewline
\ \ \ \ \isakeyword{shows}\ {\isachardoublequoteopen}ind{\isacharunderscore}{\kern0pt}happ{\isacharunderscore}{\kern0pt}seq\ {\isasympi}s\ {\isacharparenleft}{\kern0pt}simplify{\isacharunderscore}{\kern0pt}plan\ {\isacharparenleft}{\kern0pt}htps{\isacharunderscore}{\kern0pt}exec\ {\isasympi}s{\isacharparenright}{\kern0pt}\ {\isasympi}s{\isacharparenright}{\kern0pt}{\isachardoublequoteclose}\isanewline
%
\isadelimproof
\ \ \ \ %
\endisadelimproof
%
\isatagproof
\isacommand{using}\isamarkupfalse%
\ assms\ \isacommand{by}\isamarkupfalse%
\ {\isacharparenleft}{\kern0pt}rule\ simplify{\isacharunderscore}{\kern0pt}plan{\isacharunderscore}{\kern0pt}correct{\isacharparenright}{\kern0pt}%
\endisatagproof
{\isafoldproof}%
%
\isadelimproof
%
\endisadelimproof
%
\begin{isamarkuptext}%
With the correctness proof of \isa{simplify{\isacharunderscore}{\kern0pt}plan}. 
  We can easily proof that for every plan there exists an induced happening sequence.%
\end{isamarkuptext}\isamarkuptrue%
\ \ \isacommand{lemma}\isamarkupfalse%
\ {\isacharparenleft}{\kern0pt}\isakeyword{in}\ wf{\isacharunderscore}{\kern0pt}ast{\isacharunderscore}{\kern0pt}problem{\isacharparenright}{\kern0pt}\ ind{\isacharunderscore}{\kern0pt}happ{\isacharunderscore}{\kern0pt}seq{\isacharunderscore}{\kern0pt}exists{\isacharcolon}{\kern0pt}\ \isanewline
\ \ \ \ \isakeyword{assumes}\ {\isachardoublequoteopen}wf{\isacharunderscore}{\kern0pt}plan\ {\isasympi}s{\isachardoublequoteclose}\isanewline
\ \ \ \ \isakeyword{shows}\ {\isachardoublequoteopen}{\isasymexists}hs{\isachardot}{\kern0pt}\ ind{\isacharunderscore}{\kern0pt}happ{\isacharunderscore}{\kern0pt}seq\ {\isasympi}s\ hs{\isachardoublequoteclose}\isanewline
%
\isadelimproof
\ \ \ \ %
\endisadelimproof
%
\isatagproof
\isacommand{using}\isamarkupfalse%
\ assms\ simplify{\isacharunderscore}{\kern0pt}plan{\isacharunderscore}{\kern0pt}correct\ \isacommand{by}\isamarkupfalse%
\ auto%
\endisatagproof
{\isafoldproof}%
%
\isadelimproof
\isanewline
%
\endisadelimproof
\isanewline
\ \ \isacommand{lemma}\isamarkupfalse%
\ {\isacharparenleft}{\kern0pt}\isakeyword{in}\ wf{\isacharunderscore}{\kern0pt}ast{\isacharunderscore}{\kern0pt}problem{\isacharparenright}{\kern0pt}\ simplify{\isacharunderscore}{\kern0pt}plan{\isacharunderscore}{\kern0pt}wf{\isacharunderscore}{\kern0pt}ground{\isacharunderscore}{\kern0pt}action{\isacharunderscore}{\kern0pt}aux{\isacharcolon}{\kern0pt}\isanewline
\ \ \ \ \isakeyword{fixes}\ {\isasympi}s\isanewline
\ \ \ \ \isakeyword{defines}\ {\isachardoublequoteopen}htps\ {\isasymequiv}\ htps{\isacharunderscore}{\kern0pt}exec\ {\isasympi}s{\isachardoublequoteclose}\isanewline
\ \ \ \ \isakeyword{assumes}\ {\isachardoublequoteopen}wf{\isacharunderscore}{\kern0pt}plan\ {\isasympi}s{\isachardoublequoteclose}\isanewline
\ \ \ \ \ \ \ \ \isakeyword{and}\ {\isachardoublequoteopen}hs\ {\isacharequal}{\kern0pt}\ simplify{\isacharunderscore}{\kern0pt}plan\ htps\ {\isasympi}s{\isachardoublequoteclose}\isanewline
\ \ \ \ \ \ \ \ \isakeyword{and}\ {\isachardoublequoteopen}{\isacharparenleft}{\kern0pt}t\isactrlsub a{\isacharcomma}{\kern0pt}A{\isacharparenright}{\kern0pt}\ {\isasymin}\ set\ hs{\isachardoublequoteclose}\isanewline
\ \ \ \ \ \ \ \ \isakeyword{and}\ {\isachardoublequoteopen}a\ {\isasymin}\ set\ A{\isachardoublequoteclose}\isanewline
\ \ \ \ \isakeyword{shows}\ {\isachardoublequoteopen}wf{\isacharunderscore}{\kern0pt}ground{\isacharunderscore}{\kern0pt}action\ a{\isachardoublequoteclose}\isanewline
%
\isadelimproof
\ \ %
\endisadelimproof
%
\isatagproof
\isacommand{proof}\isamarkupfalse%
\ {\isacharminus}{\kern0pt}\isanewline
\ \ \ \ \isacommand{obtain}\isamarkupfalse%
\ t\isactrlsub {\isasympi}\ {\isasympi}\ \isakeyword{where}\ {\isachardoublequoteopen}{\isacharparenleft}{\kern0pt}t\isactrlsub {\isasympi}{\isacharcomma}{\kern0pt}{\isasympi}{\isacharparenright}{\kern0pt}\ {\isasymin}\ set\ {\isasympi}s{\isachardoublequoteclose}\ \isakeyword{and}\ {\isachardoublequoteopen}inst{\isacharunderscore}{\kern0pt}of{\isacharunderscore}{\kern0pt}plan{\isacharunderscore}{\kern0pt}action\ {\isasympi}s\ {\isacharparenleft}{\kern0pt}t\isactrlsub {\isasympi}{\isacharcomma}{\kern0pt}\ {\isasympi}{\isacharparenright}{\kern0pt}\ {\isacharparenleft}{\kern0pt}t\isactrlsub a{\isacharcomma}{\kern0pt}\ a{\isacharparenright}{\kern0pt}{\isachardoublequoteclose}\isanewline
\ \ \ \ \ \ \isacommand{using}\isamarkupfalse%
\ assms\ simp{\isacharunderscore}{\kern0pt}plan{\isacharunderscore}{\kern0pt}corr{\isacharunderscore}{\kern0pt}inst{\isacharunderscore}{\kern0pt}plan{\isacharunderscore}{\kern0pt}act\ \isacommand{by}\isamarkupfalse%
\ blast\isanewline
\ \ \ \ \isacommand{then}\isamarkupfalse%
\ \isacommand{show}\isamarkupfalse%
\ {\isacharquery}{\kern0pt}thesis\isanewline
\ \ \ \ \ \ \isacommand{using}\isamarkupfalse%
\ assms\ wf{\isacharunderscore}{\kern0pt}inst{\isacharunderscore}{\kern0pt}of{\isacharunderscore}{\kern0pt}plan{\isacharunderscore}{\kern0pt}action\ \isacommand{by}\isamarkupfalse%
\ simp\isanewline
\ \ \isacommand{qed}\isamarkupfalse%
%
\endisatagproof
{\isafoldproof}%
%
\isadelimproof
\isanewline
%
\endisadelimproof
\isanewline
\ \ \isacommand{lemma}\isamarkupfalse%
\ {\isacharparenleft}{\kern0pt}\isakeyword{in}\ wf{\isacharunderscore}{\kern0pt}ast{\isacharunderscore}{\kern0pt}problem{\isacharparenright}{\kern0pt}\ simplify{\isacharunderscore}{\kern0pt}plan{\isacharunderscore}{\kern0pt}wf{\isacharunderscore}{\kern0pt}ground{\isacharunderscore}{\kern0pt}action{\isacharcolon}{\kern0pt}\isanewline
\ \ \ \ \isakeyword{fixes}\ {\isasympi}s\isanewline
\ \ \ \ \isakeyword{defines}\ {\isachardoublequoteopen}htps\ {\isasymequiv}\ htps{\isacharunderscore}{\kern0pt}exec\ {\isasympi}s{\isachardoublequoteclose}\isanewline
\ \ \ \ \isakeyword{assumes}\ {\isachardoublequoteopen}wf{\isacharunderscore}{\kern0pt}plan\ {\isasympi}s{\isachardoublequoteclose}\isanewline
\ \ \ \ \ \ \ \ \isakeyword{and}\ {\isachardoublequoteopen}hs\ {\isacharequal}{\kern0pt}\ simplify{\isacharunderscore}{\kern0pt}plan\ htps\ {\isasympi}s{\isachardoublequoteclose}\isanewline
\ \ \ \ \isakeyword{shows}\ {\isachardoublequoteopen}{\isasymforall}{\isacharparenleft}{\kern0pt}t\isactrlsub a{\isacharcomma}{\kern0pt}A{\isacharparenright}{\kern0pt}\ {\isasymin}\ set\ hs{\isachardot}{\kern0pt}\ {\isasymforall}a\ {\isasymin}\ set\ A{\isachardot}{\kern0pt}\ wf{\isacharunderscore}{\kern0pt}ground{\isacharunderscore}{\kern0pt}action\ a{\isachardoublequoteclose}\isanewline
%
\isadelimproof
\ \ \ \ %
\endisadelimproof
%
\isatagproof
\isacommand{using}\isamarkupfalse%
\ assms\ simplify{\isacharunderscore}{\kern0pt}plan{\isacharunderscore}{\kern0pt}wf{\isacharunderscore}{\kern0pt}ground{\isacharunderscore}{\kern0pt}action{\isacharunderscore}{\kern0pt}aux\ \isacommand{by}\isamarkupfalse%
\ {\isacharparenleft}{\kern0pt}induction\ hs{\isacharparenright}{\kern0pt}\ auto%
\endisatagproof
{\isafoldproof}%
%
\isadelimproof
%
\endisadelimproof
%
\begin{isamarkuptext}%
Recursive refinement for \isa{valid{\isacharunderscore}{\kern0pt}happ{\isacharunderscore}{\kern0pt}seq}.%
\end{isamarkuptext}\isamarkuptrue%
\ \ \isacommand{fun}\isamarkupfalse%
\ valid{\isacharunderscore}{\kern0pt}happ{\isacharunderscore}{\kern0pt}seq{\isacharunderscore}{\kern0pt}from\ {\isacharcolon}{\kern0pt}{\isacharcolon}{\kern0pt}\ {\isachardoublequoteopen}world{\isacharunderscore}{\kern0pt}model\ {\isasymRightarrow}\ happening\ list\ {\isasymRightarrow}\ bool{\isachardoublequoteclose}\ \isakeyword{where}\isanewline
\ \ \ \ {\isachardoublequoteopen}valid{\isacharunderscore}{\kern0pt}happ{\isacharunderscore}{\kern0pt}seq{\isacharunderscore}{\kern0pt}from\ s\ {\isacharbrackleft}{\kern0pt}{\isacharbrackright}{\kern0pt}\ {\isasymlongleftrightarrow}\ close{\isacharunderscore}{\kern0pt}world\ s\ {\isasymTTurnstile}\ {\isacharparenleft}{\kern0pt}goal\ P{\isacharparenright}{\kern0pt}{\isachardoublequoteclose}\isanewline
\ \ {\isacharbar}{\kern0pt}\ {\isachardoublequoteopen}valid{\isacharunderscore}{\kern0pt}happ{\isacharunderscore}{\kern0pt}seq{\isacharunderscore}{\kern0pt}from\ s\ {\isacharparenleft}{\kern0pt}h{\isacharhash}{\kern0pt}hs{\isacharparenright}{\kern0pt}\ {\isasymlongleftrightarrow}\ \isanewline
\ \ \ \ {\isacharparenleft}{\kern0pt}happ{\isacharunderscore}{\kern0pt}enabled\ h\ s\ {\isasymand}\ valid{\isacharunderscore}{\kern0pt}happ{\isacharunderscore}{\kern0pt}seq{\isacharunderscore}{\kern0pt}from\ {\isacharparenleft}{\kern0pt}apply{\isacharunderscore}{\kern0pt}happ{\isacharunderscore}{\kern0pt}exec\ h\ s{\isacharparenright}{\kern0pt}\ hs{\isacharparenright}{\kern0pt}{\isachardoublequoteclose}%
\begin{isamarkuptext}%
Justification of the refinement for \isa{valid{\isacharunderscore}{\kern0pt}happ{\isacharunderscore}{\kern0pt}seq}.%
\end{isamarkuptext}\isamarkuptrue%
\ \ \isacommand{lemma}\isamarkupfalse%
\ valid{\isacharunderscore}{\kern0pt}happ{\isacharunderscore}{\kern0pt}seq{\isacharunderscore}{\kern0pt}from{\isacharunderscore}{\kern0pt}refine{\isacharcolon}{\kern0pt}\ \isanewline
\ \ \ \ \isakeyword{assumes}\ {\isachardoublequoteopen}wf{\isacharunderscore}{\kern0pt}happ{\isacharunderscore}{\kern0pt}seq\ hs{\isachardoublequoteclose}\isanewline
\ \ \ \ \isakeyword{shows}\ {\isachardoublequoteopen}valid{\isacharunderscore}{\kern0pt}happ{\isacharunderscore}{\kern0pt}seq{\isacharunderscore}{\kern0pt}from\ M\ hs\ {\isasymlongleftrightarrow}\ {\isacharparenleft}{\kern0pt}{\isasymexists}M{\isacharprime}{\kern0pt}{\isachardot}{\kern0pt}\ valid{\isacharunderscore}{\kern0pt}happ{\isacharunderscore}{\kern0pt}seq\ M\ hs\ M{\isacharprime}{\kern0pt}\ {\isasymand}\ M{\isacharprime}{\kern0pt}\ \isactrlsup c{\isasymTTurnstile}\isactrlsub {\isacharequal}{\kern0pt}\ {\isacharparenleft}{\kern0pt}goal\ P{\isacharparenright}{\kern0pt}{\isacharparenright}{\kern0pt}{\isachardoublequoteclose}\isanewline
%
\isadelimproof
\ \ \ \ %
\endisadelimproof
%
\isatagproof
\isacommand{using}\isamarkupfalse%
\ assms\ apply{\isacharunderscore}{\kern0pt}happ{\isacharunderscore}{\kern0pt}exec{\isacharunderscore}{\kern0pt}refine\isanewline
\ \ \ \ \isacommand{by}\isamarkupfalse%
\ {\isacharparenleft}{\kern0pt}induction\ hs\ arbitrary{\isacharcolon}{\kern0pt}\ M{\isacharparenright}{\kern0pt}\ {\isacharparenleft}{\kern0pt}auto\ split{\isacharcolon}{\kern0pt}\ prod{\isachardot}{\kern0pt}splits\ simp{\isacharcolon}{\kern0pt}\ pairwise{\isacharunderscore}{\kern0pt}def{\isacharparenright}{\kern0pt}%
\endisatagproof
{\isafoldproof}%
%
\isadelimproof
\isanewline
%
\endisadelimproof
\isanewline
\ \ \isacommand{definition}\isamarkupfalse%
\ mp{\isacharunderscore}{\kern0pt}Fevl\ {\isacharcolon}{\kern0pt}{\isacharcolon}{\kern0pt}\ {\isachardoublequoteopen}{\isacharparenleft}{\kern0pt}object{\isacharcomma}{\kern0pt}\ rat{\isacharparenright}{\kern0pt}\ mapping{\isachardoublequoteclose}\ \isakeyword{where}\isanewline
\ \ \ \ {\isachardoublequoteopen}mp{\isacharunderscore}{\kern0pt}Fevl\ {\isacharequal}{\kern0pt}\ Mapping{\isachardot}{\kern0pt}of{\isacharunderscore}{\kern0pt}alist\ {\isacharparenleft}{\kern0pt}fold\ {\isacharparenleft}{\kern0pt}{\isasymlambda}a\ fev{\isachardot}{\kern0pt}\ case\ a\ of\ Atom\ {\isacharparenleft}{\kern0pt}eqAtm\ f\ {\isacharparenleft}{\kern0pt}TimeEnt\ t{\isacharparenright}{\kern0pt}{\isacharparenright}{\kern0pt}\ {\isasymRightarrow}\ {\isacharparenleft}{\kern0pt}f{\isacharcomma}{\kern0pt}t{\isacharparenright}{\kern0pt}{\isacharhash}{\kern0pt}fev\ {\isacharbar}{\kern0pt}\ {\isacharunderscore}{\kern0pt}\ {\isasymRightarrow}\ fev{\isacharparenright}{\kern0pt}\ {\isacharparenleft}{\kern0pt}init\ P{\isacharparenright}{\kern0pt}\ {\isacharbrackleft}{\kern0pt}{\isacharbrackright}{\kern0pt}{\isacharparenright}{\kern0pt}{\isachardoublequoteclose}\isanewline
\isanewline
\ \ \isacommand{lemma}\isamarkupfalse%
\ mp{\isacharunderscore}{\kern0pt}func{\isacharunderscore}{\kern0pt}eval{\isacharunderscore}{\kern0pt}correct{\isacharbrackleft}{\kern0pt}simp{\isacharbrackright}{\kern0pt}{\isacharcolon}{\kern0pt}\ {\isachardoublequoteopen}Mapping{\isachardot}{\kern0pt}lookup\ mp{\isacharunderscore}{\kern0pt}Fevl\ {\isacharequal}{\kern0pt}\ func{\isacharunderscore}{\kern0pt}eval{\isachardoublequoteclose}\isanewline
%
\isadelimproof
\ \ \ \ %
\endisadelimproof
%
\isatagproof
\isacommand{unfolding}\isamarkupfalse%
\ mp{\isacharunderscore}{\kern0pt}Fevl{\isacharunderscore}{\kern0pt}def\ func{\isacharunderscore}{\kern0pt}eval{\isacharunderscore}{\kern0pt}def\ \isacommand{by}\isamarkupfalse%
\ transfer\ {\isacharparenleft}{\kern0pt}simp\ add{\isacharcolon}{\kern0pt}\ Map{\isacharunderscore}{\kern0pt}To{\isacharunderscore}{\kern0pt}Mapping{\isachardot}{\kern0pt}map{\isacharunderscore}{\kern0pt}apply{\isacharunderscore}{\kern0pt}def{\isacharparenright}{\kern0pt}%
\endisatagproof
{\isafoldproof}%
%
\isadelimproof
\isanewline
%
\endisadelimproof
\isanewline
\ \ \isacommand{fun}\isamarkupfalse%
\ duration{\isacharunderscore}{\kern0pt}matches{\isacharprime}{\kern0pt}\ {\isacharcolon}{\kern0pt}{\isacharcolon}{\kern0pt}\ {\isachardoublequoteopen}{\isacharparenleft}{\kern0pt}object{\isacharcomma}{\kern0pt}\ rat{\isacharparenright}{\kern0pt}\ mapping\ {\isasymRightarrow}\ time\ {\isasymRightarrow}\ term\ duration{\isacharunderscore}{\kern0pt}constraint\ {\isasymRightarrow}\ {\isacharparenleft}{\kern0pt}variable\ {\isasymtimes}\ type{\isacharparenright}{\kern0pt}\ list\ {\isasymRightarrow}\ object\ list\ {\isasymRightarrow}\ bool{\isachardoublequoteclose}\ \isakeyword{where}\isanewline
\ \ \ \ {\isachardoublequoteopen}duration{\isacharunderscore}{\kern0pt}matches{\isacharprime}{\kern0pt}\ mp{\isacharunderscore}{\kern0pt}fe\ d\ No{\isacharunderscore}{\kern0pt}Const\ params\ args\ {\isasymlongleftrightarrow}\ d\ {\isasymge}\ {\isadigit{0}}{\isachardoublequoteclose}\isanewline
\ \ {\isacharbar}{\kern0pt}\ {\isachardoublequoteopen}duration{\isacharunderscore}{\kern0pt}matches{\isacharprime}{\kern0pt}\ mp{\isacharunderscore}{\kern0pt}fe\ d\ {\isacharparenleft}{\kern0pt}Time{\isacharunderscore}{\kern0pt}Const\ op\ d{\isacharprime}{\kern0pt}{\isacharparenright}{\kern0pt}\ params\ args\ {\isasymlongleftrightarrow}\ \isanewline
\ \ \ \ \ \ {\isacharparenleft}{\kern0pt}case\ op\ of\isanewline
\ \ \ \ \ \ \ \ LEQ\ {\isasymRightarrow}\ d\ {\isasymle}\ d{\isacharprime}{\kern0pt}\isanewline
\ \ \ \ \ \ {\isacharbar}{\kern0pt}\ EQ\ {\isasymRightarrow}\ d\ {\isacharequal}{\kern0pt}\ d{\isacharprime}{\kern0pt}\isanewline
\ \ \ \ \ \ {\isacharbar}{\kern0pt}\ GEQ\ {\isasymRightarrow}\ d\ {\isasymge}\ d{\isacharprime}{\kern0pt}{\isacharparenright}{\kern0pt}{\isachardoublequoteclose}\isanewline
\ \ {\isacharbar}{\kern0pt}\ {\isachardoublequoteopen}duration{\isacharunderscore}{\kern0pt}matches{\isacharprime}{\kern0pt}\ mp{\isacharunderscore}{\kern0pt}fe\ d\ {\isacharparenleft}{\kern0pt}Func{\isacharunderscore}{\kern0pt}Const\ op\ f\ vs{\isacharparenright}{\kern0pt}\ params\ args\ {\isasymlongleftrightarrow}\ \isanewline
\ \ \ \ \ \ {\isacharparenleft}{\kern0pt}let\ tsubst\ {\isacharequal}{\kern0pt}\ subst{\isacharunderscore}{\kern0pt}term\ {\isacharparenleft}{\kern0pt}the\ o\ {\isacharparenleft}{\kern0pt}map{\isacharunderscore}{\kern0pt}of\ {\isacharparenleft}{\kern0pt}zip\ {\isacharparenleft}{\kern0pt}map\ fst\ params{\isacharparenright}{\kern0pt}\ args{\isacharparenright}{\kern0pt}{\isacharparenright}{\kern0pt}{\isacharparenright}{\kern0pt}\ in\isanewline
\ \ \ \ \ \ \ \ case\ Mapping{\isachardot}{\kern0pt}lookup\ mp{\isacharunderscore}{\kern0pt}fe\ {\isacharparenleft}{\kern0pt}FuncEnt\ f\ {\isacharparenleft}{\kern0pt}map\ tsubst\ vs{\isacharparenright}{\kern0pt}{\isacharparenright}{\kern0pt}\ of\ \isanewline
\ \ \ \ \ \ \ \ \ \ None\ {\isasymRightarrow}\ False\ \isanewline
\ \ \ \ \ \ \ \ {\isacharbar}{\kern0pt}\ Some\ d{\isacharprime}{\kern0pt}\ {\isasymRightarrow}\ {\isacharparenleft}{\kern0pt}case\ op\ of\isanewline
\ \ \ \ \ \ \ \ \ \ \ \ LEQ\ {\isasymRightarrow}\ d\ {\isasymle}\ d{\isacharprime}{\kern0pt}\isanewline
\ \ \ \ \ \ \ \ \ \ {\isacharbar}{\kern0pt}\ EQ\ {\isasymRightarrow}\ d\ {\isacharequal}{\kern0pt}\ d{\isacharprime}{\kern0pt}\isanewline
\ \ \ \ \ \ \ \ \ \ {\isacharbar}{\kern0pt}\ GEQ\ {\isasymRightarrow}\ d\ {\isasymge}\ d{\isacharprime}{\kern0pt}{\isacharparenright}{\kern0pt}{\isacharparenright}{\kern0pt}{\isachardoublequoteclose}\isanewline
\isanewline
\ \ \isacommand{lemma}\isamarkupfalse%
\ {\isacharparenleft}{\kern0pt}\isakeyword{in}\ wf{\isacharunderscore}{\kern0pt}ast{\isacharunderscore}{\kern0pt}problem{\isacharparenright}{\kern0pt}\ duration{\isacharunderscore}{\kern0pt}matches{\isacharprime}{\kern0pt}{\isacharunderscore}{\kern0pt}correct{\isacharcolon}{\kern0pt}\ \isanewline
\ \ \ \ {\isachardoublequoteopen}duration{\isacharunderscore}{\kern0pt}matches{\isacharprime}{\kern0pt}\ mp{\isacharunderscore}{\kern0pt}Fevl\ d\ dc\ params\ args\ {\isasymlongleftrightarrow}\ duration{\isacharunderscore}{\kern0pt}matches\ d\ dc\ params\ args{\isachardoublequoteclose}\isanewline
%
\isadelimproof
\ \ \ \ %
\endisadelimproof
%
\isatagproof
\isacommand{by}\isamarkupfalse%
\ {\isacharparenleft}{\kern0pt}cases\ dc{\isacharparenright}{\kern0pt}\ {\isacharparenleft}{\kern0pt}auto\ split{\isacharcolon}{\kern0pt}\ option{\isachardot}{\kern0pt}splits{\isacharparenright}{\kern0pt}%
\endisatagproof
{\isafoldproof}%
%
\isadelimproof
\isanewline
%
\endisadelimproof
\isanewline
\ \ \isacommand{fun}\isamarkupfalse%
\ consec{\isacharunderscore}{\kern0pt}htps{\isacharunderscore}{\kern0pt}in{\isacharunderscore}{\kern0pt}interval{\isacharunderscore}{\kern0pt}exec\ {\isacharcolon}{\kern0pt}{\isacharcolon}{\kern0pt}\ {\isachardoublequoteopen}time\ list\ {\isasymRightarrow}\ time\ {\isasymRightarrow}\ time\ {\isasymRightarrow}\ {\isacharparenleft}{\kern0pt}time\ {\isasymtimes}\ time{\isacharparenright}{\kern0pt}\ list{\isachardoublequoteclose}\ \isakeyword{where}\isanewline
\ \ \ \ {\isachardoublequoteopen}consec{\isacharunderscore}{\kern0pt}htps{\isacharunderscore}{\kern0pt}in{\isacharunderscore}{\kern0pt}interval{\isacharunderscore}{\kern0pt}exec\ {\isacharparenleft}{\kern0pt}t{\isacharhash}{\kern0pt}t{\isacharprime}{\kern0pt}{\isacharhash}{\kern0pt}ts{\isacharparenright}{\kern0pt}\ t\isactrlsub i\ t\isactrlsub j\ {\isacharequal}{\kern0pt}\ \isanewline
\ \ \ \ {\isacharparenleft}{\kern0pt}if\ t\ {\isacharless}{\kern0pt}\ t\isactrlsub i\ then\ consec{\isacharunderscore}{\kern0pt}htps{\isacharunderscore}{\kern0pt}in{\isacharunderscore}{\kern0pt}interval{\isacharunderscore}{\kern0pt}exec\ {\isacharparenleft}{\kern0pt}t{\isacharprime}{\kern0pt}{\isacharhash}{\kern0pt}ts{\isacharparenright}{\kern0pt}\ t\isactrlsub i\ t\isactrlsub j\isanewline
\ \ \ \ else\ if\ t\isactrlsub i\ {\isasymle}\ t\ {\isasymand}\ t{\isacharprime}{\kern0pt}\ {\isasymle}\ t\isactrlsub j\ then\ {\isacharparenleft}{\kern0pt}t{\isacharcomma}{\kern0pt}t{\isacharprime}{\kern0pt}{\isacharparenright}{\kern0pt}{\isacharhash}{\kern0pt}{\isacharparenleft}{\kern0pt}consec{\isacharunderscore}{\kern0pt}htps{\isacharunderscore}{\kern0pt}in{\isacharunderscore}{\kern0pt}interval{\isacharunderscore}{\kern0pt}exec\ {\isacharparenleft}{\kern0pt}t{\isacharprime}{\kern0pt}{\isacharhash}{\kern0pt}ts{\isacharparenright}{\kern0pt}\ t\isactrlsub i\ t\isactrlsub j{\isacharparenright}{\kern0pt}\isanewline
\ \ \ \ else\ {\isacharbrackleft}{\kern0pt}{\isacharbrackright}{\kern0pt}{\isacharparenright}{\kern0pt}{\isachardoublequoteclose}\isanewline
\ \ {\isacharbar}{\kern0pt}\ {\isachardoublequoteopen}consec{\isacharunderscore}{\kern0pt}htps{\isacharunderscore}{\kern0pt}in{\isacharunderscore}{\kern0pt}interval{\isacharunderscore}{\kern0pt}exec\ ts\ t\isactrlsub i\ t\isactrlsub j\ {\isacharequal}{\kern0pt}\ {\isacharbrackleft}{\kern0pt}{\isacharbrackright}{\kern0pt}{\isachardoublequoteclose}\isanewline
\isanewline
\ \ \isacommand{lemma}\isamarkupfalse%
\ consec{\isacharunderscore}{\kern0pt}htps{\isacharunderscore}{\kern0pt}in{\isacharunderscore}{\kern0pt}interval{\isacharunderscore}{\kern0pt}exec{\isacharunderscore}{\kern0pt}correct{\isacharcolon}{\kern0pt}\ \isanewline
\ \ \ \ \isakeyword{assumes}\ {\isachardoublequoteopen}strict{\isacharunderscore}{\kern0pt}sorted\ ts{\isachardoublequoteclose}\isanewline
\ \ \ \ \isakeyword{shows}\ {\isachardoublequoteopen}consec{\isacharunderscore}{\kern0pt}htps{\isacharunderscore}{\kern0pt}in{\isacharunderscore}{\kern0pt}interval\ ts\ t\isactrlsub i\ t\isactrlsub j\ {\isacharequal}{\kern0pt}\ consec{\isacharunderscore}{\kern0pt}htps{\isacharunderscore}{\kern0pt}in{\isacharunderscore}{\kern0pt}interval{\isacharunderscore}{\kern0pt}exec\ ts\ t\isactrlsub i\ t\isactrlsub j{\isachardoublequoteclose}\isanewline
%
\isadelimproof
\ \ \ \ %
\endisadelimproof
%
\isatagproof
\isacommand{unfolding}\isamarkupfalse%
\ consec{\isacharunderscore}{\kern0pt}htps{\isacharunderscore}{\kern0pt}in{\isacharunderscore}{\kern0pt}interval{\isacharunderscore}{\kern0pt}def\isanewline
\ \ \ \ \isacommand{using}\isamarkupfalse%
\ assms\isanewline
\ \ \isacommand{proof}\isamarkupfalse%
\ {\isacharparenleft}{\kern0pt}induction\ ts\ t\isactrlsub i\ t\isactrlsub j\ rule{\isacharcolon}{\kern0pt}\ consec{\isacharunderscore}{\kern0pt}htps{\isacharunderscore}{\kern0pt}in{\isacharunderscore}{\kern0pt}interval{\isacharunderscore}{\kern0pt}exec{\isachardot}{\kern0pt}induct{\isacharparenright}{\kern0pt}\isanewline
\ \ \ \ \isacommand{case}\isamarkupfalse%
\ {\isacharparenleft}{\kern0pt}{\isadigit{1}}\ t\ t{\isacharprime}{\kern0pt}\ ts\ t\isactrlsub i\ t\isactrlsub j{\isacharparenright}{\kern0pt}\isanewline
\ \ \ \ \isacommand{have}\isamarkupfalse%
\ {\isachardoublequoteopen}t\ {\isacharless}{\kern0pt}\ t\isactrlsub i\ {\isasymor}\ {\isacharparenleft}{\kern0pt}t\isactrlsub i\ {\isasymle}\ t\ {\isasymand}\ t{\isacharprime}{\kern0pt}\ {\isasymle}\ t\isactrlsub j{\isacharparenright}{\kern0pt}\ {\isasymor}\ {\isacharparenleft}{\kern0pt}t\isactrlsub i\ {\isasymle}\ t\ {\isasymand}\ t\isactrlsub j\ {\isacharless}{\kern0pt}\ t{\isacharprime}{\kern0pt}{\isacharparenright}{\kern0pt}{\isachardoublequoteclose}\isanewline
\ \ \ \ \ \ \isacommand{by}\isamarkupfalse%
\ auto\isanewline
\ \ \ \ \isacommand{then}\isamarkupfalse%
\ \isacommand{show}\isamarkupfalse%
\ {\isacharquery}{\kern0pt}case\ \isanewline
\ \ \ \ \isacommand{proof}\isamarkupfalse%
\ {\isacharparenleft}{\kern0pt}elim\ disjE{\isacharparenright}{\kern0pt}\isanewline
\ \ \ \ \ \ \isacommand{assume}\isamarkupfalse%
\ {\isachardoublequoteopen}t\ {\isacharless}{\kern0pt}\ t\isactrlsub i{\isachardoublequoteclose}\isanewline
\ \ \ \ \ \ \isacommand{then}\isamarkupfalse%
\ \isacommand{show}\isamarkupfalse%
\ {\isacharquery}{\kern0pt}case\isanewline
\ \ \ \ \ \ \ \ \isacommand{using}\isamarkupfalse%
\ {\isachardoublequoteopen}{\isadigit{1}}{\isachardot}{\kern0pt}prems{\isachardoublequoteclose}\ {\isachardoublequoteopen}{\isadigit{1}}{\isachardot}{\kern0pt}IH{\isachardoublequoteclose}\ \isacommand{by}\isamarkupfalse%
\ auto\isanewline
\ \ \ \ \isacommand{next}\isamarkupfalse%
\isanewline
\ \ \ \ \ \ \isacommand{assume}\isamarkupfalse%
\ {\isachardoublequoteopen}t\isactrlsub i\ {\isasymle}\ t\ {\isasymand}\ t{\isacharprime}{\kern0pt}\ {\isasymle}\ t\isactrlsub j{\isachardoublequoteclose}\isanewline
\ \ \ \ \ \ \isacommand{then}\isamarkupfalse%
\ \isacommand{show}\isamarkupfalse%
\ {\isacharquery}{\kern0pt}case\isanewline
\ \ \ \ \ \ \ \ \isacommand{using}\isamarkupfalse%
\ {\isachardoublequoteopen}{\isadigit{1}}{\isachardot}{\kern0pt}prems{\isachardoublequoteclose}\ {\isachardoublequoteopen}{\isadigit{1}}{\isachardot}{\kern0pt}IH{\isachardoublequoteclose}\ \isacommand{by}\isamarkupfalse%
\ auto\isanewline
\ \ \ \ \isacommand{next}\isamarkupfalse%
\isanewline
\ \ \ \ \ \ \isacommand{assume}\isamarkupfalse%
\ {\isachardoublequoteopen}t\isactrlsub i\ {\isasymle}\ t\ {\isasymand}\ t\isactrlsub j\ {\isacharless}{\kern0pt}\ t{\isacharprime}{\kern0pt}{\isachardoublequoteclose}\isanewline
\ \ \ \ \ \ \isacommand{then}\isamarkupfalse%
\ \isacommand{have}\isamarkupfalse%
\ {\isachardoublequoteopen}{\isasymforall}{\isacharparenleft}{\kern0pt}t{\isacharcomma}{\kern0pt}\ t{\isacharprime}{\kern0pt}{\isacharparenright}{\kern0pt}\ {\isasymin}\ set\ {\isacharparenleft}{\kern0pt}zip\ {\isacharparenleft}{\kern0pt}butlast\ {\isacharparenleft}{\kern0pt}t{\isacharhash}{\kern0pt}t{\isacharprime}{\kern0pt}{\isacharhash}{\kern0pt}ts{\isacharparenright}{\kern0pt}{\isacharparenright}{\kern0pt}\ {\isacharparenleft}{\kern0pt}tl\ {\isacharparenleft}{\kern0pt}t{\isacharhash}{\kern0pt}t{\isacharprime}{\kern0pt}{\isacharhash}{\kern0pt}ts{\isacharparenright}{\kern0pt}{\isacharparenright}{\kern0pt}{\isacharparenright}{\kern0pt}{\isachardot}{\kern0pt}\ t\isactrlsub j\ {\isacharless}{\kern0pt}\ t{\isacharprime}{\kern0pt}{\isachardoublequoteclose}\isanewline
\ \ \ \ \ \ \ \ \isacommand{using}\isamarkupfalse%
\ {\isacartoucheopen}strict{\isacharunderscore}{\kern0pt}sorted\ {\isacharparenleft}{\kern0pt}t\ {\isacharhash}{\kern0pt}\ t{\isacharprime}{\kern0pt}\ {\isacharhash}{\kern0pt}\ ts{\isacharparenright}{\kern0pt}{\isacartoucheclose}\isanewline
\ \ \ \ \ \ \ \ \isacommand{apply}\isamarkupfalse%
\ {\isacharparenleft}{\kern0pt}induction\ ts{\isacharparenright}{\kern0pt}\isanewline
\ \ \ \ \ \ \ \ \isacommand{apply}\isamarkupfalse%
\ auto\isanewline
\ \ \ \ \ \ \ \ \isacommand{apply}\isamarkupfalse%
\ {\isacharparenleft}{\kern0pt}metis\ dual{\isacharunderscore}{\kern0pt}order{\isachardot}{\kern0pt}antisym\ less{\isacharunderscore}{\kern0pt}eq{\isacharunderscore}{\kern0pt}rat{\isacharunderscore}{\kern0pt}def\ order{\isacharunderscore}{\kern0pt}trans\ set{\isacharunderscore}{\kern0pt}zip{\isacharunderscore}{\kern0pt}rightD{\isacharparenright}{\kern0pt}\isanewline
\ \ \ \ \ \ \ \ \isacommand{done}\isamarkupfalse%
\isanewline
\ \ \ \ \ \ \isacommand{then}\isamarkupfalse%
\ \isacommand{have}\isamarkupfalse%
\ {\isachardoublequoteopen}{\isasymnot}{\isacharparenleft}{\kern0pt}{\isasymexists}{\isacharparenleft}{\kern0pt}t{\isacharcomma}{\kern0pt}\ t{\isacharprime}{\kern0pt}{\isacharparenright}{\kern0pt}\ {\isasymin}\ set\ {\isacharparenleft}{\kern0pt}zip\ {\isacharparenleft}{\kern0pt}butlast\ {\isacharparenleft}{\kern0pt}t{\isacharhash}{\kern0pt}t{\isacharprime}{\kern0pt}{\isacharhash}{\kern0pt}ts{\isacharparenright}{\kern0pt}{\isacharparenright}{\kern0pt}\ {\isacharparenleft}{\kern0pt}tl\ {\isacharparenleft}{\kern0pt}t{\isacharhash}{\kern0pt}t{\isacharprime}{\kern0pt}{\isacharhash}{\kern0pt}ts{\isacharparenright}{\kern0pt}{\isacharparenright}{\kern0pt}{\isacharparenright}{\kern0pt}{\isachardot}{\kern0pt}\ t\isactrlsub i\ {\isasymle}\ t\ {\isasymand}\ t{\isacharprime}{\kern0pt}\ {\isasymle}\ t\isactrlsub j{\isacharparenright}{\kern0pt}{\isachardoublequoteclose}\isanewline
\ \ \ \ \ \ \ \ \isacommand{by}\isamarkupfalse%
\ fastforce\isanewline
\ \ \ \ \ \ \isacommand{then}\isamarkupfalse%
\ \isacommand{have}\isamarkupfalse%
\ {\isachardoublequoteopen}filter\ {\isacharparenleft}{\kern0pt}{\isasymlambda}{\isacharparenleft}{\kern0pt}t{\isacharcomma}{\kern0pt}\ t{\isacharprime}{\kern0pt}{\isacharparenright}{\kern0pt}{\isachardot}{\kern0pt}\ t\isactrlsub i\ {\isasymle}\ t\ {\isasymand}\ t{\isacharprime}{\kern0pt}\ {\isasymle}\ t\isactrlsub j{\isacharparenright}{\kern0pt}\ {\isacharparenleft}{\kern0pt}zip\ {\isacharparenleft}{\kern0pt}t{\isacharprime}{\kern0pt}\ {\isacharhash}{\kern0pt}\ butlast\ ts{\isacharparenright}{\kern0pt}\ ts{\isacharparenright}{\kern0pt}\ {\isacharequal}{\kern0pt}\ {\isacharbrackleft}{\kern0pt}{\isacharbrackright}{\kern0pt}{\isachardoublequoteclose}\isanewline
\ \ \ \ \ \ \ \ \isacommand{by}\isamarkupfalse%
\ {\isacharparenleft}{\kern0pt}induction\ ts{\isacharparenright}{\kern0pt}\ auto\isanewline
\ \ \ \ \ \ \isacommand{then}\isamarkupfalse%
\ \isacommand{show}\isamarkupfalse%
\ {\isacharquery}{\kern0pt}case\isanewline
\ \ \ \ \ \ \ \ \isacommand{using}\isamarkupfalse%
\ {\isachardoublequoteopen}{\isadigit{1}}{\isachardot}{\kern0pt}prems{\isachardoublequoteclose}\ {\isachardoublequoteopen}{\isadigit{1}}{\isachardot}{\kern0pt}IH{\isachardoublequoteclose}\ \isanewline
\ \ \ \ \ \ \ \ \isacommand{by}\isamarkupfalse%
\ auto\isanewline
\ \ \ \ \isacommand{qed}\isamarkupfalse%
\isanewline
\ \ \isacommand{next}\isamarkupfalse%
\isanewline
\ \ \ \ \isacommand{case}\isamarkupfalse%
\ {\isacharparenleft}{\kern0pt}{\isachardoublequoteopen}{\isadigit{2}}{\isacharunderscore}{\kern0pt}{\isadigit{1}}{\isachardoublequoteclose}\ t\isactrlsub i\ t\isactrlsub j{\isacharparenright}{\kern0pt}\isanewline
\ \ \ \ \isacommand{then}\isamarkupfalse%
\ \isacommand{show}\isamarkupfalse%
\ {\isacharquery}{\kern0pt}case\ \isacommand{by}\isamarkupfalse%
\ auto\isanewline
\ \ \isacommand{next}\isamarkupfalse%
\isanewline
\ \ \ \ \isacommand{case}\isamarkupfalse%
\ {\isacharparenleft}{\kern0pt}{\isachardoublequoteopen}{\isadigit{2}}{\isacharunderscore}{\kern0pt}{\isadigit{2}}{\isachardoublequoteclose}\ v\ t\isactrlsub i\ t\isactrlsub j{\isacharparenright}{\kern0pt}\isanewline
\ \ \ \ \isacommand{then}\isamarkupfalse%
\ \isacommand{show}\isamarkupfalse%
\ {\isacharquery}{\kern0pt}case\ \isacommand{by}\isamarkupfalse%
\ auto\isanewline
\ \ \isacommand{qed}\isamarkupfalse%
%
\endisatagproof
{\isafoldproof}%
%
\isadelimproof
\isanewline
%
\endisadelimproof
\isanewline
\ \ \isacommand{fun}\isamarkupfalse%
\ place{\isacharunderscore}{\kern0pt}inv{\isacharunderscore}{\kern0pt}snap{\isacharunderscore}{\kern0pt}acts\ {\isacharcolon}{\kern0pt}{\isacharcolon}{\kern0pt}\ {\isachardoublequoteopen}time\ list\ {\isasymRightarrow}\ time\ {\isasymRightarrow}\ time\ {\isasymRightarrow}\ ground{\isacharunderscore}{\kern0pt}action\ {\isasymRightarrow}\ {\isacharparenleft}{\kern0pt}time\ {\isasymtimes}\ ground{\isacharunderscore}{\kern0pt}action{\isacharparenright}{\kern0pt}\ list{\isachardoublequoteclose}\ \isakeyword{where}\isanewline
\ \ \ \ {\isachardoublequoteopen}place{\isacharunderscore}{\kern0pt}inv{\isacharunderscore}{\kern0pt}snap{\isacharunderscore}{\kern0pt}acts\ {\isacharparenleft}{\kern0pt}t{\isacharhash}{\kern0pt}t{\isacharprime}{\kern0pt}{\isacharhash}{\kern0pt}ts{\isacharparenright}{\kern0pt}\ t\isactrlsub i\ t\isactrlsub j\ a\isactrlsub i\isactrlsub n\isactrlsub v\ {\isacharequal}{\kern0pt}\ \isanewline
\ \ \ \ {\isacharparenleft}{\kern0pt}if\ t\ {\isacharless}{\kern0pt}\ t\isactrlsub i\ then\ place{\isacharunderscore}{\kern0pt}inv{\isacharunderscore}{\kern0pt}snap{\isacharunderscore}{\kern0pt}acts\ {\isacharparenleft}{\kern0pt}t{\isacharprime}{\kern0pt}{\isacharhash}{\kern0pt}ts{\isacharparenright}{\kern0pt}\ t\isactrlsub i\ t\isactrlsub j\ a\isactrlsub i\isactrlsub n\isactrlsub v\isanewline
\ \ \ \ else\ if\ t\isactrlsub i\ {\isasymle}\ t\ {\isasymand}\ t{\isacharprime}{\kern0pt}\ {\isasymle}\ t\isactrlsub j\ then\ {\isacharparenleft}{\kern0pt}{\isacharparenleft}{\kern0pt}t{\isacharplus}{\kern0pt}t{\isacharprime}{\kern0pt}{\isacharparenright}{\kern0pt}\ div\ {\isadigit{2}}{\isacharcomma}{\kern0pt}a\isactrlsub i\isactrlsub n\isactrlsub v{\isacharparenright}{\kern0pt}{\isacharhash}{\kern0pt}{\isacharparenleft}{\kern0pt}place{\isacharunderscore}{\kern0pt}inv{\isacharunderscore}{\kern0pt}snap{\isacharunderscore}{\kern0pt}acts\ {\isacharparenleft}{\kern0pt}t{\isacharprime}{\kern0pt}{\isacharhash}{\kern0pt}ts{\isacharparenright}{\kern0pt}\ t\isactrlsub i\ t\isactrlsub j\ a\isactrlsub i\isactrlsub n\isactrlsub v{\isacharparenright}{\kern0pt}\isanewline
\ \ \ \ else\ {\isacharbrackleft}{\kern0pt}{\isacharbrackright}{\kern0pt}{\isacharparenright}{\kern0pt}{\isachardoublequoteclose}\isanewline
\ \ {\isacharbar}{\kern0pt}\ {\isachardoublequoteopen}place{\isacharunderscore}{\kern0pt}inv{\isacharunderscore}{\kern0pt}snap{\isacharunderscore}{\kern0pt}acts\ ts\ t\isactrlsub i\ t\isactrlsub j\ a\isactrlsub i\isactrlsub n\isactrlsub v\ {\isacharequal}{\kern0pt}\ {\isacharbrackleft}{\kern0pt}{\isacharbrackright}{\kern0pt}{\isachardoublequoteclose}\isanewline
\ \ \isanewline
\isanewline
\ \ \isacommand{lemma}\isamarkupfalse%
\ place{\isacharunderscore}{\kern0pt}inv{\isacharunderscore}{\kern0pt}snap{\isacharunderscore}{\kern0pt}acts{\isacharunderscore}{\kern0pt}correct{\isacharcolon}{\kern0pt}\isanewline
\ \ \ \ \isakeyword{assumes}\ {\isachardoublequoteopen}strict{\isacharunderscore}{\kern0pt}sorted\ ts{\isachardoublequoteclose}\isanewline
\ \ \ \ \isakeyword{shows}\ {\isachardoublequoteopen}place{\isacharunderscore}{\kern0pt}inv{\isacharunderscore}{\kern0pt}snap{\isacharunderscore}{\kern0pt}acts\ ts\ t\isactrlsub {\isadigit{1}}\ t\isactrlsub {\isadigit{2}}\ a\isactrlsub i\isactrlsub n\isactrlsub v\ {\isacharequal}{\kern0pt}\ map\ {\isacharparenleft}{\kern0pt}{\isasymlambda}{\isacharparenleft}{\kern0pt}t\isactrlsub i{\isacharcomma}{\kern0pt}t\isactrlsub j{\isacharparenright}{\kern0pt}{\isachardot}{\kern0pt}{\isacharparenleft}{\kern0pt}{\isacharparenleft}{\kern0pt}t\isactrlsub i{\isacharplus}{\kern0pt}t\isactrlsub j{\isacharparenright}{\kern0pt}\ div\ {\isadigit{2}}{\isacharcomma}{\kern0pt}a\isactrlsub i\isactrlsub n\isactrlsub v{\isacharparenright}{\kern0pt}{\isacharparenright}{\kern0pt}\ {\isacharparenleft}{\kern0pt}consec{\isacharunderscore}{\kern0pt}htps{\isacharunderscore}{\kern0pt}in{\isacharunderscore}{\kern0pt}interval\ ts\ t\isactrlsub {\isadigit{1}}\ t\isactrlsub {\isadigit{2}}{\isacharparenright}{\kern0pt}{\isachardoublequoteclose}\isanewline
%
\isadelimproof
\ \ \ \ %
\endisadelimproof
%
\isatagproof
\isacommand{using}\isamarkupfalse%
\ assms\ consec{\isacharunderscore}{\kern0pt}htps{\isacharunderscore}{\kern0pt}in{\isacharunderscore}{\kern0pt}interval{\isacharunderscore}{\kern0pt}exec{\isacharunderscore}{\kern0pt}correct\isanewline
\ \ \ \ \isacommand{by}\isamarkupfalse%
\ {\isacharparenleft}{\kern0pt}induction\ ts\ t\isactrlsub {\isadigit{1}}\ t\isactrlsub {\isadigit{2}}\ a\isactrlsub i\isactrlsub n\isactrlsub v\ rule{\isacharcolon}{\kern0pt}\ place{\isacharunderscore}{\kern0pt}inv{\isacharunderscore}{\kern0pt}snap{\isacharunderscore}{\kern0pt}acts{\isachardot}{\kern0pt}induct{\isacharparenright}{\kern0pt}\ auto%
\endisatagproof
{\isafoldproof}%
%
\isadelimproof
%
\endisadelimproof
%
\begin{isamarkuptext}%
Lift \isa{simplify{\isacharunderscore}{\kern0pt}action} into an Error Monad.%
\end{isamarkuptext}\isamarkuptrue%
\ \ \isacommand{fun}\isamarkupfalse%
\ simplify{\isacharunderscore}{\kern0pt}actionE\ {\isacharcolon}{\kern0pt}{\isacharcolon}{\kern0pt}\ {\isachardoublequoteopen}{\isacharunderscore}{\kern0pt}\ {\isasymRightarrow}\ {\isacharparenleft}{\kern0pt}object{\isacharcomma}{\kern0pt}\ type{\isacharparenright}{\kern0pt}\ mapping\ {\isasymRightarrow}\ {\isacharparenleft}{\kern0pt}object{\isacharcomma}{\kern0pt}\ rat{\isacharparenright}{\kern0pt}\ mapping\ {\isasymRightarrow}\ time\ list\ {\isasymRightarrow}\ {\isacharparenleft}{\kern0pt}time\ {\isasymtimes}\ plan{\isacharunderscore}{\kern0pt}action{\isacharparenright}{\kern0pt}\ {\isasymRightarrow}\ {\isacharunderscore}{\kern0pt}{\isacharplus}{\kern0pt}{\isacharparenleft}{\kern0pt}time\ {\isasymtimes}\ ground{\isacharunderscore}{\kern0pt}action{\isacharparenright}{\kern0pt}\ list{\isachardoublequoteclose}\ \isakeyword{where}\isanewline
\ \ \ \ {\isachardoublequoteopen}simplify{\isacharunderscore}{\kern0pt}actionE\ stg\ mp\ mp{\isacharunderscore}{\kern0pt}fe\ htps\ {\isacharparenleft}{\kern0pt}t\isactrlsub {\isasympi}{\isacharcomma}{\kern0pt}Simple{\isacharunderscore}{\kern0pt}Plan{\isacharunderscore}{\kern0pt}Action\ n\ args{\isacharparenright}{\kern0pt}\ {\isacharequal}{\kern0pt}\ do\ {\isacharbraceleft}{\kern0pt}\isanewline
\ \ \ \ \ \ check\ {\isacharparenleft}{\kern0pt}t\isactrlsub {\isasympi}\ {\isasymge}\ {\isadigit{0}}{\isacharparenright}{\kern0pt}\ {\isacharparenleft}{\kern0pt}ERRS\ {\isacharprime}{\kern0pt}{\isacharprime}{\kern0pt}Invalid\ starting\ time\ point{\isacharprime}{\kern0pt}{\isacharprime}{\kern0pt}{\isacharparenright}{\kern0pt}{\isacharsemicolon}{\kern0pt}\isanewline
\ \ \ \ \ \ a\ {\isasymleftarrow}\ resolve{\isacharunderscore}{\kern0pt}action{\isacharunderscore}{\kern0pt}schemaE\ n{\isacharsemicolon}{\kern0pt}\isanewline
\ \ \ \ \ \ check\ {\isacharparenleft}{\kern0pt}case\ a\ of\ Simple{\isacharunderscore}{\kern0pt}Action{\isacharunderscore}{\kern0pt}Schema\ {\isacharunderscore}{\kern0pt}\ {\isacharunderscore}{\kern0pt}\ {\isacharunderscore}{\kern0pt}\ {\isacharunderscore}{\kern0pt}\ {\isasymRightarrow}\ True\ {\isacharbar}{\kern0pt}\ {\isacharunderscore}{\kern0pt}\ {\isasymRightarrow}\ False{\isacharparenright}{\kern0pt}\ {\isacharparenleft}{\kern0pt}ERRS\ {\isacharprime}{\kern0pt}{\isacharprime}{\kern0pt}Invalid\ Plan\ action\ type{\isacharprime}{\kern0pt}{\isacharprime}{\kern0pt}{\isacharparenright}{\kern0pt}{\isacharsemicolon}{\kern0pt}\isanewline
\ \ \ \ \ \ check\ {\isacharparenleft}{\kern0pt}action{\isacharunderscore}{\kern0pt}params{\isacharunderscore}{\kern0pt}match{\isadigit{2}}\ stg\ mp\ a\ args{\isacharparenright}{\kern0pt}\ {\isacharparenleft}{\kern0pt}ERRS\ {\isacharprime}{\kern0pt}{\isacharprime}{\kern0pt}Parameter\ mismatch{\isacharprime}{\kern0pt}{\isacharprime}{\kern0pt}{\isacharparenright}{\kern0pt}{\isacharsemicolon}{\kern0pt}\isanewline
\ \ \ \ \ \ Error{\isacharunderscore}{\kern0pt}Monad{\isachardot}{\kern0pt}return\ {\isacharbrackleft}{\kern0pt}{\isacharparenleft}{\kern0pt}t\isactrlsub {\isasympi}{\isacharcomma}{\kern0pt}instantiate{\isacharunderscore}{\kern0pt}action{\isacharunderscore}{\kern0pt}schema\ a\ args{\isacharparenright}{\kern0pt}{\isacharbrackright}{\kern0pt}\isanewline
\ \ \ \ {\isacharbraceright}{\kern0pt}{\isachardoublequoteclose}\isanewline
\ \ {\isacharbar}{\kern0pt}\ {\isachardoublequoteopen}simplify{\isacharunderscore}{\kern0pt}actionE\ stg\ mp\ mp{\isacharunderscore}{\kern0pt}fe\ htps\ {\isacharparenleft}{\kern0pt}t\isactrlsub {\isasympi}{\isacharcomma}{\kern0pt}Durative{\isacharunderscore}{\kern0pt}Plan{\isacharunderscore}{\kern0pt}Action\ n\ args\ d{\isacharparenright}{\kern0pt}\ {\isacharequal}{\kern0pt}\ do\ {\isacharbraceleft}{\kern0pt}\isanewline
\ \ \ \ \ \ check\ {\isacharparenleft}{\kern0pt}t\isactrlsub {\isasympi}\ {\isasymge}\ {\isadigit{0}}{\isacharparenright}{\kern0pt}\ {\isacharparenleft}{\kern0pt}ERRS\ {\isacharprime}{\kern0pt}{\isacharprime}{\kern0pt}Invalid\ starting\ time\ point{\isacharprime}{\kern0pt}{\isacharprime}{\kern0pt}{\isacharparenright}{\kern0pt}{\isacharsemicolon}{\kern0pt}\isanewline
\ \ \ \ \ \ a\ {\isasymleftarrow}\ resolve{\isacharunderscore}{\kern0pt}action{\isacharunderscore}{\kern0pt}schemaE\ n{\isacharsemicolon}{\kern0pt}\isanewline
\ \ \ \ \ \ check\ {\isacharparenleft}{\kern0pt}case\ a\ of\ Durative{\isacharunderscore}{\kern0pt}Action{\isacharunderscore}{\kern0pt}Schema\ {\isacharunderscore}{\kern0pt}\ {\isacharunderscore}{\kern0pt}\ {\isacharunderscore}{\kern0pt}\ {\isacharunderscore}{\kern0pt}\ {\isacharunderscore}{\kern0pt}\ {\isasymRightarrow}\ True\ {\isacharbar}{\kern0pt}\ {\isacharunderscore}{\kern0pt}\ {\isasymRightarrow}\ False{\isacharparenright}{\kern0pt}\ {\isacharparenleft}{\kern0pt}ERRS\ {\isacharprime}{\kern0pt}{\isacharprime}{\kern0pt}Invalid\ Plan\ action\ type{\isacharprime}{\kern0pt}{\isacharprime}{\kern0pt}{\isacharparenright}{\kern0pt}{\isacharsemicolon}{\kern0pt}\isanewline
\ \ \ \ \ \ check\ {\isacharparenleft}{\kern0pt}action{\isacharunderscore}{\kern0pt}params{\isacharunderscore}{\kern0pt}match{\isadigit{2}}\ stg\ mp\ a\ args{\isacharparenright}{\kern0pt}\ {\isacharparenleft}{\kern0pt}ERRS\ {\isacharprime}{\kern0pt}{\isacharprime}{\kern0pt}Parameter\ mismatch{\isacharprime}{\kern0pt}{\isacharprime}{\kern0pt}{\isacharparenright}{\kern0pt}{\isacharsemicolon}{\kern0pt}\isanewline
\ \ \ \ \ \ check\ {\isacharparenleft}{\kern0pt}d\ {\isasymge}\ {\isadigit{0}}{\isacharparenright}{\kern0pt}\ {\isacharparenleft}{\kern0pt}ERRS\ {\isacharprime}{\kern0pt}{\isacharprime}{\kern0pt}Invalid\ duration\ {\isacharparenleft}{\kern0pt}negative{\isacharparenright}{\kern0pt}{\isacharprime}{\kern0pt}{\isacharprime}{\kern0pt}{\isacharparenright}{\kern0pt}{\isacharsemicolon}{\kern0pt}\isanewline
\ \ \ \ \ \ check\ {\isacharparenleft}{\kern0pt}duration{\isacharunderscore}{\kern0pt}matches{\isacharprime}{\kern0pt}\ mp{\isacharunderscore}{\kern0pt}fe\ d\ {\isacharparenleft}{\kern0pt}duration{\isacharunderscore}{\kern0pt}constraint\ a{\isacharparenright}{\kern0pt}\ {\isacharparenleft}{\kern0pt}parameters\ a{\isacharparenright}{\kern0pt}\ args{\isacharparenright}{\kern0pt}\ {\isacharparenleft}{\kern0pt}ERRS\ {\isacharprime}{\kern0pt}{\isacharprime}{\kern0pt}Duration\ constraint\ not\ satisfied{\isacharprime}{\kern0pt}{\isacharprime}{\kern0pt}{\isacharparenright}{\kern0pt}{\isacharsemicolon}{\kern0pt}\isanewline
\ \ \ \ \ \ let\ a\isactrlsub i\isactrlsub n\isactrlsub v\ {\isacharequal}{\kern0pt}\ inst{\isacharunderscore}{\kern0pt}snap{\isacharunderscore}{\kern0pt}action\ a\ args\ Over{\isacharunderscore}{\kern0pt}All\ in\isanewline
\ \ \ \ \ \ \ \ Error{\isacharunderscore}{\kern0pt}Monad{\isachardot}{\kern0pt}return\ {\isacharparenleft}{\kern0pt}\isanewline
\ \ \ \ \ \ \ \ \ \ {\isacharparenleft}{\kern0pt}t\isactrlsub {\isasympi}{\isacharcomma}{\kern0pt}inst{\isacharunderscore}{\kern0pt}snap{\isacharunderscore}{\kern0pt}action\ a\ args\ At{\isacharunderscore}{\kern0pt}Start{\isacharparenright}{\kern0pt}\ {\isacharhash}{\kern0pt}\ \isanewline
\ \ \ \ \ \ \ \ \ \ {\isacharparenleft}{\kern0pt}t\isactrlsub {\isasympi}{\isacharplus}{\kern0pt}d{\isacharcomma}{\kern0pt}inst{\isacharunderscore}{\kern0pt}snap{\isacharunderscore}{\kern0pt}action\ a\ args\ At{\isacharunderscore}{\kern0pt}End{\isacharparenright}{\kern0pt}\ {\isacharhash}{\kern0pt}\ \isanewline
\ \ \ \ \ \ \ \ \ \ {\isacharparenleft}{\kern0pt}place{\isacharunderscore}{\kern0pt}inv{\isacharunderscore}{\kern0pt}snap{\isacharunderscore}{\kern0pt}acts\ htps\ t\isactrlsub {\isasympi}\ {\isacharparenleft}{\kern0pt}t\isactrlsub {\isasympi}{\isacharplus}{\kern0pt}d{\isacharparenright}{\kern0pt}\ a\isactrlsub i\isactrlsub n\isactrlsub v{\isacharparenright}{\kern0pt}\isanewline
\ \ \ \ \ \ \ \ {\isacharparenright}{\kern0pt}\isanewline
\ \ \ \ {\isacharbraceright}{\kern0pt}{\isachardoublequoteclose}%
\begin{isamarkuptext}%
Justification of refinement for \isa{simplify{\isacharunderscore}{\kern0pt}action}.%
\end{isamarkuptext}\isamarkuptrue%
\ \ \isacommand{lemma}\isamarkupfalse%
\ {\isacharparenleft}{\kern0pt}\isakeyword{in}\ wf{\isacharunderscore}{\kern0pt}ast{\isacharunderscore}{\kern0pt}problem{\isacharparenright}{\kern0pt}\ simplify{\isacharunderscore}{\kern0pt}actionE{\isacharunderscore}{\kern0pt}return{\isacharunderscore}{\kern0pt}iff{\isacharcolon}{\kern0pt}\isanewline
\ \ \ \ \isakeyword{assumes}\ {\isachardoublequoteopen}wf{\isacharunderscore}{\kern0pt}domain{\isachardoublequoteclose}\ \isakeyword{and}\ {\isachardoublequoteopen}strict{\isacharunderscore}{\kern0pt}sorted\ htps{\isachardoublequoteclose}\isanewline
\ \ \ \ \isakeyword{shows}\ {\isachardoublequoteopen}simplify{\isacharunderscore}{\kern0pt}actionE\ STG\ mp{\isacharunderscore}{\kern0pt}objT\ mp{\isacharunderscore}{\kern0pt}Fevl\ htps\ {\isacharparenleft}{\kern0pt}t\isactrlsub {\isasympi}{\isacharcomma}{\kern0pt}{\isasympi}{\isacharparenright}{\kern0pt}\ {\isacharequal}{\kern0pt}\ Inr\ as\ \isanewline
\ \ \ \ \ \ {\isasymlongleftrightarrow}\ {\isacharparenleft}{\kern0pt}wf{\isacharunderscore}{\kern0pt}plan{\isacharunderscore}{\kern0pt}action\ {\isasympi}\ {\isasymand}\ t\isactrlsub {\isasympi}\ {\isasymge}\ {\isadigit{0}}\ {\isasymand}\ simplify{\isacharunderscore}{\kern0pt}action\ htps\ {\isacharparenleft}{\kern0pt}t\isactrlsub {\isasympi}{\isacharcomma}{\kern0pt}{\isasympi}{\isacharparenright}{\kern0pt}\ {\isacharequal}{\kern0pt}\ as{\isacharparenright}{\kern0pt}{\isachardoublequoteclose}\isanewline
%
\isadelimproof
\ \ %
\endisadelimproof
%
\isatagproof
\isacommand{proof}\isamarkupfalse%
\isanewline
\ \ \ \ \isacommand{assume}\isamarkupfalse%
\ assm{\isadigit{1}}{\isacharcolon}{\kern0pt}\ {\isachardoublequoteopen}simplify{\isacharunderscore}{\kern0pt}actionE\ STG\ mp{\isacharunderscore}{\kern0pt}objT\ mp{\isacharunderscore}{\kern0pt}Fevl\ htps\ {\isacharparenleft}{\kern0pt}t\isactrlsub {\isasympi}{\isacharcomma}{\kern0pt}{\isasympi}{\isacharparenright}{\kern0pt}\ {\isacharequal}{\kern0pt}\ Inr\ as{\isachardoublequoteclose}\isanewline
\ \ \ \ \isacommand{obtain}\isamarkupfalse%
\ a\isactrlsub s\isactrlsub c\isactrlsub h\isactrlsub e\isactrlsub m\isactrlsub a\ \isakeyword{where}\ a\isactrlsub s\isactrlsub c\isactrlsub h\isactrlsub e\isactrlsub m\isactrlsub a{\isacharunderscore}{\kern0pt}obt{\isacharcolon}{\kern0pt}\ {\isachardoublequoteopen}Some\ a\isactrlsub s\isactrlsub c\isactrlsub h\isactrlsub e\isactrlsub m\isactrlsub a\ {\isacharequal}{\kern0pt}\ resolve{\isacharunderscore}{\kern0pt}action{\isacharunderscore}{\kern0pt}schema\ {\isacharparenleft}{\kern0pt}name\ {\isasympi}{\isacharparenright}{\kern0pt}{\isachardoublequoteclose}\isanewline
\ \ \ \ \ \ \isacommand{using}\isamarkupfalse%
\ assm{\isadigit{1}}\ resolve{\isacharunderscore}{\kern0pt}action{\isacharunderscore}{\kern0pt}schema{\isacharunderscore}{\kern0pt}def\isanewline
\ \ \ \ \ \ \isacommand{unfolding}\isamarkupfalse%
\ wf{\isacharunderscore}{\kern0pt}domain{\isacharunderscore}{\kern0pt}def\ \isacommand{by}\isamarkupfalse%
\ {\isacharparenleft}{\kern0pt}cases\ {\isasympi}{\isacharparenright}{\kern0pt}\ {\isacharparenleft}{\kern0pt}auto\ simp{\isacharcolon}{\kern0pt}\ return{\isacharunderscore}{\kern0pt}iff{\isacharparenright}{\kern0pt}\isanewline
\ \ \ \ \isacommand{then}\isamarkupfalse%
\ \isacommand{have}\isamarkupfalse%
\ wf{\isacharunderscore}{\kern0pt}act{\isacharunderscore}{\kern0pt}schema{\isacharunderscore}{\kern0pt}{\isasympi}{\isacharcolon}{\kern0pt}\ {\isachardoublequoteopen}wf{\isacharunderscore}{\kern0pt}action{\isacharunderscore}{\kern0pt}schema\ {\isacharparenleft}{\kern0pt}the\ {\isacharparenleft}{\kern0pt}resolve{\isacharunderscore}{\kern0pt}action{\isacharunderscore}{\kern0pt}schema\ {\isacharparenleft}{\kern0pt}name\ {\isasympi}{\isacharparenright}{\kern0pt}{\isacharparenright}{\kern0pt}{\isacharparenright}{\kern0pt}{\isachardoublequoteclose}\isanewline
\ \ \ \ \ \ \isacommand{using}\isamarkupfalse%
\ assms\ a\isactrlsub s\isactrlsub c\isactrlsub h\isactrlsub e\isactrlsub m\isactrlsub a{\isacharunderscore}{\kern0pt}obt\ \isanewline
\ \ \ \ \ \ \ \ \ \ \ \ resolve{\isacharunderscore}{\kern0pt}action{\isacharunderscore}{\kern0pt}schema{\isacharunderscore}{\kern0pt}def\ \isanewline
\ \ \ \ \ \ \ \ \ \ \ \ index{\isacharunderscore}{\kern0pt}by{\isacharunderscore}{\kern0pt}eq{\isacharunderscore}{\kern0pt}SomeD{\isacharbrackleft}{\kern0pt}\isakeyword{where}\ n{\isacharequal}{\kern0pt}{\isachardoublequoteopen}name\ {\isasympi}{\isachardoublequoteclose}\ \isakeyword{and}\ x{\isacharequal}{\kern0pt}{\isachardoublequoteopen}a\isactrlsub s\isactrlsub c\isactrlsub h\isactrlsub e\isactrlsub m\isactrlsub a{\isachardoublequoteclose}{\isacharbrackright}{\kern0pt}\isanewline
\ \ \ \ \ \ \isacommand{unfolding}\isamarkupfalse%
\ wf{\isacharunderscore}{\kern0pt}domain{\isacharunderscore}{\kern0pt}def\ \isacommand{by}\isamarkupfalse%
\ {\isacharparenleft}{\kern0pt}metis\ option{\isachardot}{\kern0pt}sel{\isacharparenright}{\kern0pt}\isanewline
\ \ \ \ \isacommand{then}\isamarkupfalse%
\ \isacommand{show}\isamarkupfalse%
\ {\isachardoublequoteopen}wf{\isacharunderscore}{\kern0pt}plan{\isacharunderscore}{\kern0pt}action\ {\isasympi}\ {\isasymand}\ t\isactrlsub {\isasympi}\ {\isasymge}\ {\isadigit{0}}\ {\isasymand}\ simplify{\isacharunderscore}{\kern0pt}action\ htps\ {\isacharparenleft}{\kern0pt}t\isactrlsub {\isasympi}{\isacharcomma}{\kern0pt}{\isasympi}{\isacharparenright}{\kern0pt}\ {\isacharequal}{\kern0pt}\ as{\isachardoublequoteclose}\isanewline
\ \ \ \ \isacommand{proof}\isamarkupfalse%
\ {\isacharparenleft}{\kern0pt}cases\ {\isasympi}{\isacharparenright}{\kern0pt}\isanewline
\ \ \ \ \ \ \isacommand{case}\isamarkupfalse%
\ case{\isacharunderscore}{\kern0pt}{\isasympi}{\isacharcolon}{\kern0pt}\ {\isacharparenleft}{\kern0pt}Simple{\isacharunderscore}{\kern0pt}Plan{\isacharunderscore}{\kern0pt}Action\ n\ args{\isacharparenright}{\kern0pt}\isanewline
\ \ \ \ \ \ \isacommand{obtain}\isamarkupfalse%
\ a\ \isakeyword{where}\ a{\isacharunderscore}{\kern0pt}obt{\isacharcolon}{\kern0pt}\ {\isachardoublequoteopen}a\ {\isacharequal}{\kern0pt}\ instantiate{\isacharunderscore}{\kern0pt}action{\isacharunderscore}{\kern0pt}schema\ a\isactrlsub s\isactrlsub c\isactrlsub h\isactrlsub e\isactrlsub m\isactrlsub a\ args\ {\isasymand}\ as\ {\isacharequal}{\kern0pt}\ {\isacharbrackleft}{\kern0pt}{\isacharparenleft}{\kern0pt}t\isactrlsub {\isasympi}{\isacharcomma}{\kern0pt}a{\isacharparenright}{\kern0pt}{\isacharbrackright}{\kern0pt}{\isachardoublequoteclose}\isanewline
\ \ \ \ \ \ \ \ \isacommand{using}\isamarkupfalse%
\ assm{\isadigit{1}}\ a\isactrlsub s\isactrlsub c\isactrlsub h\isactrlsub e\isactrlsub m\isactrlsub a{\isacharunderscore}{\kern0pt}obt\ case{\isacharunderscore}{\kern0pt}{\isasympi}\ \isacommand{by}\isamarkupfalse%
\ {\isacharparenleft}{\kern0pt}auto\ simp{\isacharcolon}{\kern0pt}\ return{\isacharunderscore}{\kern0pt}iff{\isacharparenright}{\kern0pt}\isanewline
\ \ \ \ \ \ \isacommand{have}\isamarkupfalse%
\ wf{\isacharunderscore}{\kern0pt}p{\isacharunderscore}{\kern0pt}act{\isacharcolon}{\kern0pt}\ {\isachardoublequoteopen}wf{\isacharunderscore}{\kern0pt}plan{\isacharunderscore}{\kern0pt}action\ {\isasympi}\ {\isasymand}\ t\isactrlsub {\isasympi}\ {\isasymge}\ {\isadigit{0}}{\isachardoublequoteclose}\isanewline
\ \ \ \ \ \ \ \ \isacommand{using}\isamarkupfalse%
\ assm{\isadigit{1}}\ wf{\isacharunderscore}{\kern0pt}act{\isacharunderscore}{\kern0pt}schema{\isacharunderscore}{\kern0pt}{\isasympi}\ case{\isacharunderscore}{\kern0pt}{\isasympi}\ action{\isacharunderscore}{\kern0pt}params{\isacharunderscore}{\kern0pt}match{\isacharunderscore}{\kern0pt}def\ wf{\isacharunderscore}{\kern0pt}effect{\isacharunderscore}{\kern0pt}inst{\isacharunderscore}{\kern0pt}weak\isanewline
\ \ \ \ \ \ \ \ \isacommand{apply}\isamarkupfalse%
\ {\isacharparenleft}{\kern0pt}auto\ simp{\isacharcolon}{\kern0pt}\ return{\isacharunderscore}{\kern0pt}iff\ split{\isacharcolon}{\kern0pt}\ option{\isachardot}{\kern0pt}splits\ ast{\isacharunderscore}{\kern0pt}action{\isacharunderscore}{\kern0pt}schema{\isachardot}{\kern0pt}splits{\isacharparenright}{\kern0pt}\isanewline
\ \ \ \ \ \ \ \ \isacommand{apply}\isamarkupfalse%
\ {\isacharparenleft}{\kern0pt}meson\ ast{\isacharunderscore}{\kern0pt}action{\isacharunderscore}{\kern0pt}schema{\isachardot}{\kern0pt}exhaust{\isacharparenright}{\kern0pt}{\isacharplus}{\kern0pt}\isanewline
\ \ \ \ \ \ \ \ \isacommand{done}\isamarkupfalse%
\isanewline
\ \ \ \ \ \ \isacommand{show}\isamarkupfalse%
\ {\isacharquery}{\kern0pt}thesis\ \isanewline
\ \ \ \ \ \ \ \ \isacommand{using}\isamarkupfalse%
\ a{\isacharunderscore}{\kern0pt}obt\ a\isactrlsub s\isactrlsub c\isactrlsub h\isactrlsub e\isactrlsub m\isactrlsub a{\isacharunderscore}{\kern0pt}obt\ case{\isacharunderscore}{\kern0pt}{\isasympi}\ wf{\isacharunderscore}{\kern0pt}p{\isacharunderscore}{\kern0pt}act\isanewline
\ \ \ \ \ \ \ \ \isacommand{by}\isamarkupfalse%
\ {\isacharparenleft}{\kern0pt}auto\ split{\isacharcolon}{\kern0pt}\ option{\isachardot}{\kern0pt}splits\ ast{\isacharunderscore}{\kern0pt}action{\isacharunderscore}{\kern0pt}schema{\isachardot}{\kern0pt}splits{\isacharparenright}{\kern0pt}\isanewline
\ \ \ \ \isacommand{next}\isamarkupfalse%
\isanewline
\ \ \ \ \ \ \isacommand{case}\isamarkupfalse%
\ case{\isacharunderscore}{\kern0pt}{\isasympi}{\isacharcolon}{\kern0pt}\ {\isacharparenleft}{\kern0pt}Durative{\isacharunderscore}{\kern0pt}Plan{\isacharunderscore}{\kern0pt}Action\ n\ args\ d{\isacharparenright}{\kern0pt}\isanewline
\ \ \ \ \ \ \isacommand{show}\isamarkupfalse%
\ {\isacharquery}{\kern0pt}thesis\ \isanewline
\ \ \ \ \ \ \ \ \isacommand{using}\isamarkupfalse%
\ assm{\isadigit{1}}\ wf{\isacharunderscore}{\kern0pt}act{\isacharunderscore}{\kern0pt}schema{\isacharunderscore}{\kern0pt}{\isasympi}\ case{\isacharunderscore}{\kern0pt}{\isasympi}\ action{\isacharunderscore}{\kern0pt}params{\isacharunderscore}{\kern0pt}match{\isacharunderscore}{\kern0pt}def\ \isanewline
\ \ \ \ \ \ \ \ \ \ \ \ \ \ wf{\isacharunderscore}{\kern0pt}effect{\isacharunderscore}{\kern0pt}durative{\isacharunderscore}{\kern0pt}inst{\isacharunderscore}{\kern0pt}weak\ \isanewline
\ \ \ \ \ \ \ \ \ \ \ \ \ \ duration{\isacharunderscore}{\kern0pt}matches{\isacharprime}{\kern0pt}{\isacharunderscore}{\kern0pt}correct\ \isanewline
\ \ \ \ \ \ \ \ \ \ \ \ \ \ place{\isacharunderscore}{\kern0pt}inv{\isacharunderscore}{\kern0pt}snap{\isacharunderscore}{\kern0pt}acts{\isacharunderscore}{\kern0pt}correct{\isacharbrackleft}{\kern0pt}OF\ {\isacartoucheopen}strict{\isacharunderscore}{\kern0pt}sorted\ htps{\isacartoucheclose}{\isacharbrackright}{\kern0pt}\isanewline
\ \ \ \ \ \ \ \ \isacommand{apply}\isamarkupfalse%
\ {\isacharparenleft}{\kern0pt}auto\ simp{\isacharcolon}{\kern0pt}\ return{\isacharunderscore}{\kern0pt}iff\ split{\isacharcolon}{\kern0pt}\ ast{\isacharunderscore}{\kern0pt}action{\isacharunderscore}{\kern0pt}schema{\isachardot}{\kern0pt}splits{\isacharparenright}{\kern0pt}\isanewline
\ \ \ \ \ \ \ \ \isacommand{apply}\isamarkupfalse%
\ {\isacharparenleft}{\kern0pt}auto\ simp{\isacharcolon}{\kern0pt}\ return{\isacharunderscore}{\kern0pt}iff\ split{\isacharcolon}{\kern0pt}\ option{\isachardot}{\kern0pt}splits{\isacharparenright}{\kern0pt}\isanewline
\ \ \ \ \ \ \ \ \isacommand{apply}\isamarkupfalse%
\ {\isacharparenleft}{\kern0pt}meson\ ast{\isacharunderscore}{\kern0pt}action{\isacharunderscore}{\kern0pt}schema{\isachardot}{\kern0pt}exhaust{\isacharparenright}{\kern0pt}{\isacharplus}{\kern0pt}\isanewline
\ \ \ \ \ \ \ \ \isacommand{done}\isamarkupfalse%
\isanewline
\ \ \ \ \isacommand{qed}\isamarkupfalse%
\isanewline
\ \ \isacommand{next}\isamarkupfalse%
\isanewline
\ \ \ \ \isacommand{assume}\isamarkupfalse%
\ assm{\isadigit{2}}{\isacharcolon}{\kern0pt}\ {\isachardoublequoteopen}wf{\isacharunderscore}{\kern0pt}plan{\isacharunderscore}{\kern0pt}action\ {\isasympi}\ {\isasymand}\ t\isactrlsub {\isasympi}\ {\isasymge}\ {\isadigit{0}}\ {\isasymand}\ simplify{\isacharunderscore}{\kern0pt}action\ htps\ {\isacharparenleft}{\kern0pt}t\isactrlsub {\isasympi}{\isacharcomma}{\kern0pt}{\isasympi}{\isacharparenright}{\kern0pt}\ {\isacharequal}{\kern0pt}\ as{\isachardoublequoteclose}\isanewline
\ \ \ \ \isacommand{obtain}\isamarkupfalse%
\ a\isactrlsub s\isactrlsub c\isactrlsub h\isactrlsub e\isactrlsub m\isactrlsub a\ \isakeyword{where}\ a\isactrlsub s\isactrlsub c\isactrlsub h\isactrlsub e\isactrlsub m\isactrlsub a{\isacharunderscore}{\kern0pt}obt{\isacharcolon}{\kern0pt}\ {\isachardoublequoteopen}Some\ a\isactrlsub s\isactrlsub c\isactrlsub h\isactrlsub e\isactrlsub m\isactrlsub a\ {\isacharequal}{\kern0pt}\ resolve{\isacharunderscore}{\kern0pt}action{\isacharunderscore}{\kern0pt}schema\ {\isacharparenleft}{\kern0pt}name\ {\isasympi}{\isacharparenright}{\kern0pt}{\isachardoublequoteclose}\isanewline
\ \ \ \ \ \ \isacommand{using}\isamarkupfalse%
\ assm{\isadigit{2}}\ \isacommand{by}\isamarkupfalse%
\ {\isacharparenleft}{\kern0pt}cases\ {\isasympi}{\isacharparenright}{\kern0pt}\ {\isacharparenleft}{\kern0pt}auto\ split{\isacharcolon}{\kern0pt}\ option{\isachardot}{\kern0pt}splits{\isacharparenright}{\kern0pt}\isanewline
\ \ \ \ \isacommand{then}\isamarkupfalse%
\ \isacommand{show}\isamarkupfalse%
\ {\isachardoublequoteopen}simplify{\isacharunderscore}{\kern0pt}actionE\ STG\ mp{\isacharunderscore}{\kern0pt}objT\ mp{\isacharunderscore}{\kern0pt}Fevl\ htps\ {\isacharparenleft}{\kern0pt}t\isactrlsub {\isasympi}{\isacharcomma}{\kern0pt}{\isasympi}{\isacharparenright}{\kern0pt}\ {\isacharequal}{\kern0pt}\ Inr\ as{\isachardoublequoteclose}\isanewline
\ \ \ \ \ \ \isacommand{using}\isamarkupfalse%
\ assm{\isadigit{2}}\ action{\isacharunderscore}{\kern0pt}params{\isacharunderscore}{\kern0pt}match{\isacharunderscore}{\kern0pt}def\ \isanewline
\ \ \ \ \ \ \ \ \ \ \ \ duration{\isacharunderscore}{\kern0pt}matches{\isacharprime}{\kern0pt}{\isacharunderscore}{\kern0pt}correct\ \isanewline
\ \ \ \ \ \ \ \ \ \ \ \ place{\isacharunderscore}{\kern0pt}inv{\isacharunderscore}{\kern0pt}snap{\isacharunderscore}{\kern0pt}acts{\isacharunderscore}{\kern0pt}correct{\isacharbrackleft}{\kern0pt}OF\ {\isacartoucheopen}strict{\isacharunderscore}{\kern0pt}sorted\ htps{\isacartoucheclose}{\isacharbrackright}{\kern0pt}\isanewline
\ \ \ \ \ \ \isacommand{by}\isamarkupfalse%
\ {\isacharparenleft}{\kern0pt}cases\ {\isasympi}{\isacharparenright}{\kern0pt}\ {\isacharparenleft}{\kern0pt}auto\ simp{\isacharcolon}{\kern0pt}\ return{\isacharunderscore}{\kern0pt}iff\ split{\isacharcolon}{\kern0pt}\ option{\isachardot}{\kern0pt}splits\ ast{\isacharunderscore}{\kern0pt}action{\isacharunderscore}{\kern0pt}schema{\isachardot}{\kern0pt}splits{\isacharparenright}{\kern0pt}\isanewline
\ \ \isacommand{qed}\isamarkupfalse%
%
\endisatagproof
{\isafoldproof}%
%
\isadelimproof
%
\endisadelimproof
%
\begin{isamarkuptext}%
Lift \isa{simplify{\isacharunderscore}{\kern0pt}plan} into an Error Monad.%
\end{isamarkuptext}\isamarkuptrue%
\ \ \isacommand{fun}\isamarkupfalse%
\ simplify{\isacharunderscore}{\kern0pt}planE\ {\isacharcolon}{\kern0pt}{\isacharcolon}{\kern0pt}\ {\isachardoublequoteopen}{\isacharunderscore}{\kern0pt}\ {\isasymRightarrow}\ {\isacharparenleft}{\kern0pt}object{\isacharcomma}{\kern0pt}\ type{\isacharparenright}{\kern0pt}\ mapping\ {\isasymRightarrow}\ {\isacharparenleft}{\kern0pt}object{\isacharcomma}{\kern0pt}\ rat{\isacharparenright}{\kern0pt}\ mapping\ {\isasymRightarrow}\ time\ list\ {\isasymRightarrow}\ plan\ {\isasymRightarrow}\ {\isacharunderscore}{\kern0pt}{\isacharplus}{\kern0pt}happening\ list{\isachardoublequoteclose}\ \isakeyword{where}\isanewline
\ \ \ \ {\isachardoublequoteopen}simplify{\isacharunderscore}{\kern0pt}planE\ stg\ mp\ mp{\isacharunderscore}{\kern0pt}fe\ htps\ {\isacharbrackleft}{\kern0pt}{\isacharbrackright}{\kern0pt}\ {\isacharequal}{\kern0pt}\ do\ {\isacharbraceleft}{\kern0pt}\ \isanewline
\ \ \ \ \ \ Error{\isacharunderscore}{\kern0pt}Monad{\isachardot}{\kern0pt}return\ {\isacharbrackleft}{\kern0pt}{\isacharbrackright}{\kern0pt}\ \isanewline
\ \ \ \ {\isacharbraceright}{\kern0pt}{\isachardoublequoteclose}\ \isanewline
\ \ {\isacharbar}{\kern0pt}\ {\isachardoublequoteopen}simplify{\isacharunderscore}{\kern0pt}planE\ stg\ mp\ mp{\isacharunderscore}{\kern0pt}fe\ htps\ {\isacharparenleft}{\kern0pt}{\isasympi}{\isacharhash}{\kern0pt}{\isasympi}s{\isacharparenright}{\kern0pt}\ {\isacharequal}{\kern0pt}\ do\ {\isacharbraceleft}{\kern0pt}\isanewline
\ \ \ \ \ \ as\ {\isasymleftarrow}\ simplify{\isacharunderscore}{\kern0pt}actionE\ stg\ mp\ mp{\isacharunderscore}{\kern0pt}fe\ htps\ {\isasympi}{\isacharsemicolon}{\kern0pt}\isanewline
\ \ \ \ \ \ hs\ {\isasymleftarrow}\ simplify{\isacharunderscore}{\kern0pt}planE\ stg\ mp\ mp{\isacharunderscore}{\kern0pt}fe\ htps\ {\isasympi}s{\isacharsemicolon}{\kern0pt}\isanewline
\ \ \ \ \ \ Error{\isacharunderscore}{\kern0pt}Monad{\isachardot}{\kern0pt}return\ {\isacharparenleft}{\kern0pt}insort{\isacharunderscore}{\kern0pt}timed{\isacharunderscore}{\kern0pt}ground{\isacharunderscore}{\kern0pt}acts\ as\ hs{\isacharparenright}{\kern0pt}\isanewline
\ \ \ \ {\isacharbraceright}{\kern0pt}{\isachardoublequoteclose}%
\begin{isamarkuptext}%
Justification of refinement for \isa{simplify{\isacharunderscore}{\kern0pt}plan}%
\end{isamarkuptext}\isamarkuptrue%
\ \ \isacommand{lemma}\isamarkupfalse%
\ {\isacharparenleft}{\kern0pt}\isakeyword{in}\ wf{\isacharunderscore}{\kern0pt}ast{\isacharunderscore}{\kern0pt}problem{\isacharparenright}{\kern0pt}\ simplify{\isacharunderscore}{\kern0pt}planE{\isacharunderscore}{\kern0pt}return{\isacharunderscore}{\kern0pt}iff{\isacharbrackleft}{\kern0pt}return{\isacharunderscore}{\kern0pt}iff{\isacharbrackright}{\kern0pt}{\isacharcolon}{\kern0pt}\ \isanewline
\ \ \ \ \isakeyword{assumes}\ {\isachardoublequoteopen}wf{\isacharunderscore}{\kern0pt}domain{\isachardoublequoteclose}\ \isakeyword{and}\ {\isachardoublequoteopen}strict{\isacharunderscore}{\kern0pt}sorted\ htps{\isachardoublequoteclose}\isanewline
\ \ \ \ \isakeyword{shows}\ {\isachardoublequoteopen}simplify{\isacharunderscore}{\kern0pt}planE\ STG\ mp{\isacharunderscore}{\kern0pt}objT\ mp{\isacharunderscore}{\kern0pt}Fevl\ htps\ {\isasympi}s\ {\isacharequal}{\kern0pt}\ Inr\ hs\ \isanewline
\ \ \ \ \ \ \ \ \ \ \ \ {\isasymlongleftrightarrow}\ {\isacharparenleft}{\kern0pt}simplify{\isacharunderscore}{\kern0pt}plan\ htps\ {\isasympi}s\ {\isacharequal}{\kern0pt}\ hs\ {\isasymand}\ wf{\isacharunderscore}{\kern0pt}plan\ {\isasympi}s{\isacharparenright}{\kern0pt}{\isachardoublequoteclose}\isanewline
%
\isadelimproof
\ \ \ \ %
\endisadelimproof
%
\isatagproof
\isacommand{using}\isamarkupfalse%
\ assms\ simplify{\isacharunderscore}{\kern0pt}actionE{\isacharunderscore}{\kern0pt}return{\isacharunderscore}{\kern0pt}iff\ wf{\isacharunderscore}{\kern0pt}plan{\isacharunderscore}{\kern0pt}def\ insort{\isacharunderscore}{\kern0pt}timed{\isacharunderscore}{\kern0pt}ground{\isacharunderscore}{\kern0pt}acts{\isacharunderscore}{\kern0pt}correct\isanewline
\ \ \ \ \isacommand{by}\isamarkupfalse%
\ {\isacharparenleft}{\kern0pt}induction\ {\isasympi}s\ arbitrary{\isacharcolon}{\kern0pt}\ hs{\isacharparenright}{\kern0pt}\ {\isacharparenleft}{\kern0pt}auto\ simp{\isacharcolon}{\kern0pt}\ return{\isacharunderscore}{\kern0pt}iff{\isacharparenright}{\kern0pt}%
\endisatagproof
{\isafoldproof}%
%
\isadelimproof
\isanewline
%
\endisadelimproof
\isanewline
\ \ \isanewline
\ \ \isacommand{type{\isacharunderscore}{\kern0pt}synonym}\isamarkupfalse%
\ {\isacharprime}{\kern0pt}a\ tree{\isacharunderscore}{\kern0pt}ht\ {\isacharequal}{\kern0pt}\ {\isachardoublequoteopen}{\isacharparenleft}{\kern0pt}{\isacharprime}{\kern0pt}a\ {\isasymtimes}\ nat{\isacharparenright}{\kern0pt}\ tree{\isachardoublequoteclose}\isanewline
\isanewline
\ \ \isacommand{fun}\isamarkupfalse%
\ avl{\isacharunderscore}{\kern0pt}height\ {\isacharcolon}{\kern0pt}{\isacharcolon}{\kern0pt}\ {\isachardoublequoteopen}happening\ tree{\isacharunderscore}{\kern0pt}ht\ {\isasymRightarrow}\ nat{\isachardoublequoteclose}\ \isakeyword{where}\isanewline
\ \ \ \ {\isachardoublequoteopen}avl{\isacharunderscore}{\kern0pt}height\ {\isasymlangle}{\isasymrangle}\ {\isacharequal}{\kern0pt}\ {\isadigit{0}}{\isachardoublequoteclose}\ \isanewline
\ \ {\isacharbar}{\kern0pt}\ {\isachardoublequoteopen}avl{\isacharunderscore}{\kern0pt}height\ {\isasymlangle}l{\isacharcomma}{\kern0pt}{\isacharparenleft}{\kern0pt}{\isacharparenleft}{\kern0pt}t{\isacharcomma}{\kern0pt}as{\isacharparenright}{\kern0pt}{\isacharcomma}{\kern0pt}ht{\isacharparenright}{\kern0pt}{\isacharcomma}{\kern0pt}r{\isasymrangle}\ {\isacharequal}{\kern0pt}\ ht{\isachardoublequoteclose}\isanewline
\isanewline
\ \ \isacommand{definition}\isamarkupfalse%
\ avl{\isacharunderscore}{\kern0pt}node\ {\isacharcolon}{\kern0pt}{\isacharcolon}{\kern0pt}\ {\isachardoublequoteopen}happening\ tree{\isacharunderscore}{\kern0pt}ht\ {\isasymRightarrow}\ happening\ {\isasymRightarrow}\ happening\ tree{\isacharunderscore}{\kern0pt}ht\ {\isasymRightarrow}\ happening\ tree{\isacharunderscore}{\kern0pt}ht{\isachardoublequoteclose}\ \isakeyword{where}\isanewline
\ \ \ \ {\isachardoublequoteopen}avl{\isacharunderscore}{\kern0pt}node\ l\ h\ r\ {\isacharequal}{\kern0pt}\ {\isasymlangle}l{\isacharcomma}{\kern0pt}{\isacharparenleft}{\kern0pt}h{\isacharcomma}{\kern0pt}{\isacharparenleft}{\kern0pt}max\ {\isacharparenleft}{\kern0pt}avl{\isacharunderscore}{\kern0pt}height\ l{\isacharparenright}{\kern0pt}\ {\isacharparenleft}{\kern0pt}avl{\isacharunderscore}{\kern0pt}height\ r{\isacharparenright}{\kern0pt}\ {\isacharplus}{\kern0pt}\ {\isadigit{1}}{\isacharparenright}{\kern0pt}{\isacharparenright}{\kern0pt}{\isacharcomma}{\kern0pt}r{\isasymrangle}{\isachardoublequoteclose}\isanewline
\isanewline
\ \ \isacommand{definition}\isamarkupfalse%
\ avl{\isacharunderscore}{\kern0pt}balL\ {\isacharcolon}{\kern0pt}{\isacharcolon}{\kern0pt}\ {\isachardoublequoteopen}happening\ tree{\isacharunderscore}{\kern0pt}ht\ {\isasymRightarrow}\ happening\ {\isasymRightarrow}\ happening\ tree{\isacharunderscore}{\kern0pt}ht\ {\isasymRightarrow}\ happening\ tree{\isacharunderscore}{\kern0pt}ht{\isachardoublequoteclose}\ \isakeyword{where}\isanewline
\ \ \ \ {\isachardoublequoteopen}avl{\isacharunderscore}{\kern0pt}balL\ l\ a\ r\ {\isacharequal}{\kern0pt}\isanewline
\ \ \ \ \ \ {\isacharparenleft}{\kern0pt}if\ avl{\isacharunderscore}{\kern0pt}height\ l\ {\isacharequal}{\kern0pt}\ avl{\isacharunderscore}{\kern0pt}height\ r\ {\isacharplus}{\kern0pt}\ {\isadigit{2}}\ then\isanewline
\ \ \ \ \ \ \ \ case\ l\ of\ {\isasymlangle}bl{\isacharcomma}{\kern0pt}{\isacharparenleft}{\kern0pt}b{\isacharcomma}{\kern0pt}{\isacharunderscore}{\kern0pt}{\isacharparenright}{\kern0pt}{\isacharcomma}{\kern0pt}br{\isasymrangle}\ {\isasymRightarrow}\isanewline
\ \ \ \ \ \ \ \ \ \ \ \ if\ avl{\isacharunderscore}{\kern0pt}height\ bl\ {\isacharless}{\kern0pt}\ avl{\isacharunderscore}{\kern0pt}height\ br\ then\isanewline
\ \ \ \ \ \ \ \ \ \ \ \ \ \ case\ br\ of\ {\isasymlangle}cl{\isacharcomma}{\kern0pt}{\isacharparenleft}{\kern0pt}c{\isacharcomma}{\kern0pt}{\isacharunderscore}{\kern0pt}{\isacharparenright}{\kern0pt}{\isacharcomma}{\kern0pt}cr{\isasymrangle}\ {\isasymRightarrow}\ avl{\isacharunderscore}{\kern0pt}node\ {\isacharparenleft}{\kern0pt}avl{\isacharunderscore}{\kern0pt}node\ bl\ b\ cl{\isacharparenright}{\kern0pt}\ c\ {\isacharparenleft}{\kern0pt}avl{\isacharunderscore}{\kern0pt}node\ cr\ a\ r\ {\isacharparenright}{\kern0pt}\isanewline
\ \ \ \ \ \ \ \ \ \ \ \ else\ avl{\isacharunderscore}{\kern0pt}node\ bl\ b\ {\isacharparenleft}{\kern0pt}avl{\isacharunderscore}{\kern0pt}node\ br\ a\ r{\isacharparenright}{\kern0pt}\isanewline
\ \ \ \ \ \ \ else\ avl{\isacharunderscore}{\kern0pt}node\ l\ a\ r{\isacharparenright}{\kern0pt}{\isachardoublequoteclose}\isanewline
\isanewline
\ \ \isacommand{definition}\isamarkupfalse%
\ avl{\isacharunderscore}{\kern0pt}balR\ {\isacharcolon}{\kern0pt}{\isacharcolon}{\kern0pt}\ {\isachardoublequoteopen}happening\ tree{\isacharunderscore}{\kern0pt}ht\ {\isasymRightarrow}\ happening\ {\isasymRightarrow}\ happening\ tree{\isacharunderscore}{\kern0pt}ht\ {\isasymRightarrow}\ happening\ tree{\isacharunderscore}{\kern0pt}ht{\isachardoublequoteclose}\ \isakeyword{where}\isanewline
\ \ \ \ {\isachardoublequoteopen}avl{\isacharunderscore}{\kern0pt}balR\ l\ a\ r\ {\isacharequal}{\kern0pt}\isanewline
\ \ \ \ \ \ {\isacharparenleft}{\kern0pt}if\ avl{\isacharunderscore}{\kern0pt}height\ r\ {\isacharequal}{\kern0pt}\ avl{\isacharunderscore}{\kern0pt}height\ l\ {\isacharplus}{\kern0pt}\ {\isadigit{2}}\ then\isanewline
\ \ \ \ \ \ \ \ case\ r\ of\ {\isasymlangle}bl{\isacharcomma}{\kern0pt}{\isacharparenleft}{\kern0pt}b{\isacharcomma}{\kern0pt}{\isacharunderscore}{\kern0pt}{\isacharparenright}{\kern0pt}{\isacharcomma}{\kern0pt}br{\isasymrangle}\ {\isasymRightarrow}\isanewline
\ \ \ \ \ \ \ \ \ \ if\ avl{\isacharunderscore}{\kern0pt}height\ bl\ {\isachargreater}{\kern0pt}\ avl{\isacharunderscore}{\kern0pt}height\ br\ then\isanewline
\ \ \ \ \ \ \ \ \ \ \ \ case\ bl\ of\ {\isasymlangle}cl{\isacharcomma}{\kern0pt}{\isacharparenleft}{\kern0pt}c{\isacharcomma}{\kern0pt}{\isacharunderscore}{\kern0pt}{\isacharparenright}{\kern0pt}{\isacharcomma}{\kern0pt}cr{\isasymrangle}\ {\isasymRightarrow}\ avl{\isacharunderscore}{\kern0pt}node\ {\isacharparenleft}{\kern0pt}avl{\isacharunderscore}{\kern0pt}node\ l\ a\ cl{\isacharparenright}{\kern0pt}\ c\ {\isacharparenleft}{\kern0pt}avl{\isacharunderscore}{\kern0pt}node\ cr\ b\ br{\isacharparenright}{\kern0pt}\isanewline
\ \ \ \ \ \ \ \ \ \ else\ avl{\isacharunderscore}{\kern0pt}node\ {\isacharparenleft}{\kern0pt}avl{\isacharunderscore}{\kern0pt}node\ l\ a\ bl{\isacharparenright}{\kern0pt}\ b\ br\isanewline
\ \ \ \ \ \ else\ avl{\isacharunderscore}{\kern0pt}node\ l\ a\ r{\isacharparenright}{\kern0pt}{\isachardoublequoteclose}\isanewline
\isanewline
\ \ \isacommand{fun}\isamarkupfalse%
\ insert{\isacharunderscore}{\kern0pt}happ{\isacharunderscore}{\kern0pt}to{\isacharunderscore}{\kern0pt}tree\ {\isacharcolon}{\kern0pt}{\isacharcolon}{\kern0pt}\ {\isachardoublequoteopen}happening\ {\isasymRightarrow}\ happening\ tree{\isacharunderscore}{\kern0pt}ht\ {\isasymRightarrow}\ happening\ tree{\isacharunderscore}{\kern0pt}ht{\isachardoublequoteclose}\ \isakeyword{where}\isanewline
\ \ \ \ {\isachardoublequoteopen}insert{\isacharunderscore}{\kern0pt}happ{\isacharunderscore}{\kern0pt}to{\isacharunderscore}{\kern0pt}tree\ {\isacharparenleft}{\kern0pt}t\isactrlsub i{\isacharcomma}{\kern0pt}A\isactrlsub i{\isacharparenright}{\kern0pt}\ {\isasymlangle}{\isasymrangle}\ {\isacharequal}{\kern0pt}\ avl{\isacharunderscore}{\kern0pt}node\ {\isasymlangle}{\isasymrangle}\ {\isacharparenleft}{\kern0pt}t\isactrlsub i{\isacharcomma}{\kern0pt}A\isactrlsub i{\isacharparenright}{\kern0pt}\ {\isasymlangle}{\isasymrangle}{\isachardoublequoteclose}\isanewline
\ \ {\isacharbar}{\kern0pt}\ {\isachardoublequoteopen}insert{\isacharunderscore}{\kern0pt}happ{\isacharunderscore}{\kern0pt}to{\isacharunderscore}{\kern0pt}tree\ {\isacharparenleft}{\kern0pt}t\isactrlsub i{\isacharcomma}{\kern0pt}A\isactrlsub i{\isacharparenright}{\kern0pt}\ {\isasymlangle}l{\isacharcomma}{\kern0pt}{\isacharparenleft}{\kern0pt}{\isacharparenleft}{\kern0pt}t\isactrlsub j{\isacharcomma}{\kern0pt}A\isactrlsub j{\isacharparenright}{\kern0pt}{\isacharcomma}{\kern0pt}h{\isacharparenright}{\kern0pt}{\isacharcomma}{\kern0pt}r{\isasymrangle}\ {\isacharequal}{\kern0pt}\ {\isacharparenleft}{\kern0pt}\isanewline
\ \ \ \ if\ t\isactrlsub i\ {\isacharless}{\kern0pt}\ t\isactrlsub j\ then\ avl{\isacharunderscore}{\kern0pt}balL\ {\isacharparenleft}{\kern0pt}insert{\isacharunderscore}{\kern0pt}happ{\isacharunderscore}{\kern0pt}to{\isacharunderscore}{\kern0pt}tree\ {\isacharparenleft}{\kern0pt}t\isactrlsub i{\isacharcomma}{\kern0pt}A\isactrlsub i{\isacharparenright}{\kern0pt}\ l{\isacharparenright}{\kern0pt}\ {\isacharparenleft}{\kern0pt}t\isactrlsub j{\isacharcomma}{\kern0pt}A\isactrlsub j{\isacharparenright}{\kern0pt}\ r\isanewline
\ \ \ \ else\ if\ t\isactrlsub i\ {\isacharequal}{\kern0pt}\ t\isactrlsub j\ then\ avl{\isacharunderscore}{\kern0pt}node\ l\ {\isacharparenleft}{\kern0pt}t\isactrlsub j{\isacharcomma}{\kern0pt}A\isactrlsub i\ {\isacharat}{\kern0pt}\ A\isactrlsub j{\isacharparenright}{\kern0pt}\ r\isanewline
\ \ \ \ else\ avl{\isacharunderscore}{\kern0pt}balR\ l\ {\isacharparenleft}{\kern0pt}t\isactrlsub j{\isacharcomma}{\kern0pt}A\isactrlsub j{\isacharparenright}{\kern0pt}\ {\isacharparenleft}{\kern0pt}insert{\isacharunderscore}{\kern0pt}happ{\isacharunderscore}{\kern0pt}to{\isacharunderscore}{\kern0pt}tree\ {\isacharparenleft}{\kern0pt}t\isactrlsub i{\isacharcomma}{\kern0pt}A\isactrlsub i{\isacharparenright}{\kern0pt}\ r{\isacharparenright}{\kern0pt}{\isacharparenright}{\kern0pt}{\isachardoublequoteclose}\isanewline
\isanewline
\ \ \isacommand{fun}\isamarkupfalse%
\ insert{\isacharunderscore}{\kern0pt}timed{\isacharunderscore}{\kern0pt}ground{\isacharunderscore}{\kern0pt}acts{\isacharunderscore}{\kern0pt}to{\isacharunderscore}{\kern0pt}tree\ {\isacharcolon}{\kern0pt}{\isacharcolon}{\kern0pt}\ {\isachardoublequoteopen}{\isacharparenleft}{\kern0pt}time\ {\isasymtimes}\ ground{\isacharunderscore}{\kern0pt}action{\isacharparenright}{\kern0pt}\ list\ {\isasymRightarrow}\ happening\ tree{\isacharunderscore}{\kern0pt}ht\ {\isasymRightarrow}\ happening\ tree{\isacharunderscore}{\kern0pt}ht{\isachardoublequoteclose}\ \isakeyword{where}\isanewline
\ \ \ \ {\isachardoublequoteopen}insert{\isacharunderscore}{\kern0pt}timed{\isacharunderscore}{\kern0pt}ground{\isacharunderscore}{\kern0pt}acts{\isacharunderscore}{\kern0pt}to{\isacharunderscore}{\kern0pt}tree\ {\isacharbrackleft}{\kern0pt}{\isacharbrackright}{\kern0pt}\ htree\ {\isacharequal}{\kern0pt}\ htree{\isachardoublequoteclose}\isanewline
\ \ {\isacharbar}{\kern0pt}\ {\isachardoublequoteopen}insert{\isacharunderscore}{\kern0pt}timed{\isacharunderscore}{\kern0pt}ground{\isacharunderscore}{\kern0pt}acts{\isacharunderscore}{\kern0pt}to{\isacharunderscore}{\kern0pt}tree\ {\isacharparenleft}{\kern0pt}{\isacharparenleft}{\kern0pt}t\isactrlsub a{\isacharcomma}{\kern0pt}a{\isacharparenright}{\kern0pt}{\isacharhash}{\kern0pt}as{\isacharparenright}{\kern0pt}\ htree\ {\isacharequal}{\kern0pt}\ \isanewline
\ \ \ \ insert{\isacharunderscore}{\kern0pt}happ{\isacharunderscore}{\kern0pt}to{\isacharunderscore}{\kern0pt}tree\ {\isacharparenleft}{\kern0pt}t\isactrlsub a{\isacharcomma}{\kern0pt}{\isacharbrackleft}{\kern0pt}a{\isacharbrackright}{\kern0pt}{\isacharparenright}{\kern0pt}\ {\isacharparenleft}{\kern0pt}insert{\isacharunderscore}{\kern0pt}timed{\isacharunderscore}{\kern0pt}ground{\isacharunderscore}{\kern0pt}acts{\isacharunderscore}{\kern0pt}to{\isacharunderscore}{\kern0pt}tree\ as\ htree{\isacharparenright}{\kern0pt}{\isachardoublequoteclose}\isanewline
\isanewline
\ \ \isacommand{definition}\isamarkupfalse%
\ avl{\isacharunderscore}{\kern0pt}inorder\ {\isacharcolon}{\kern0pt}{\isacharcolon}{\kern0pt}\ {\isachardoublequoteopen}happening\ tree{\isacharunderscore}{\kern0pt}ht\ {\isasymRightarrow}\ happening\ list{\isachardoublequoteclose}\ \isakeyword{where}\ \isanewline
\ \ \ \ {\isachardoublequoteopen}avl{\isacharunderscore}{\kern0pt}inorder\ htree\ {\isacharequal}{\kern0pt}\ map\ fst\ {\isacharparenleft}{\kern0pt}inorder\ htree{\isacharparenright}{\kern0pt}{\isachardoublequoteclose}\isanewline
\isanewline
\ \ \isacommand{fun}\isamarkupfalse%
\ simplify{\isacharunderscore}{\kern0pt}planE{\isacharunderscore}{\kern0pt}avl{\isacharprime}{\kern0pt}\ {\isacharcolon}{\kern0pt}{\isacharcolon}{\kern0pt}\ {\isachardoublequoteopen}{\isacharunderscore}{\kern0pt}\ {\isasymRightarrow}\ {\isacharparenleft}{\kern0pt}object{\isacharcomma}{\kern0pt}\ type{\isacharparenright}{\kern0pt}\ mapping\ {\isasymRightarrow}\ {\isacharparenleft}{\kern0pt}object{\isacharcomma}{\kern0pt}\ rat{\isacharparenright}{\kern0pt}\ mapping\ {\isasymRightarrow}\ time\ list\ {\isasymRightarrow}\ plan\ {\isasymRightarrow}\ {\isacharunderscore}{\kern0pt}{\isacharplus}{\kern0pt}happening\ tree{\isacharunderscore}{\kern0pt}ht{\isachardoublequoteclose}\ \isakeyword{where}\isanewline
\ \ \ \ {\isachardoublequoteopen}simplify{\isacharunderscore}{\kern0pt}planE{\isacharunderscore}{\kern0pt}avl{\isacharprime}{\kern0pt}\ G\ mp\ mp{\isacharunderscore}{\kern0pt}fe\ htps\ {\isacharbrackleft}{\kern0pt}{\isacharbrackright}{\kern0pt}\ {\isacharequal}{\kern0pt}\ do\ {\isacharbraceleft}{\kern0pt}\ \isanewline
\ \ \ \ \ \ Error{\isacharunderscore}{\kern0pt}Monad{\isachardot}{\kern0pt}return\ {\isasymlangle}{\isasymrangle}\isanewline
\ \ \ \ {\isacharbraceright}{\kern0pt}{\isachardoublequoteclose}\ \isanewline
\ \ {\isacharbar}{\kern0pt}\ {\isachardoublequoteopen}simplify{\isacharunderscore}{\kern0pt}planE{\isacharunderscore}{\kern0pt}avl{\isacharprime}{\kern0pt}\ G\ mp\ mp{\isacharunderscore}{\kern0pt}fe\ htps\ {\isacharparenleft}{\kern0pt}{\isasympi}{\isacharhash}{\kern0pt}{\isasympi}s{\isacharparenright}{\kern0pt}\ {\isacharequal}{\kern0pt}\ do\ {\isacharbraceleft}{\kern0pt}\isanewline
\ \ \ \ \ \ as\ {\isasymleftarrow}\ simplify{\isacharunderscore}{\kern0pt}actionE\ G\ mp\ mp{\isacharunderscore}{\kern0pt}fe\ htps\ {\isasympi}{\isacharsemicolon}{\kern0pt}\isanewline
\ \ \ \ \ \ htree\ {\isasymleftarrow}\ simplify{\isacharunderscore}{\kern0pt}planE{\isacharunderscore}{\kern0pt}avl{\isacharprime}{\kern0pt}\ G\ mp\ mp{\isacharunderscore}{\kern0pt}fe\ htps\ {\isasympi}s{\isacharsemicolon}{\kern0pt}\isanewline
\ \ \ \ \ \ Error{\isacharunderscore}{\kern0pt}Monad{\isachardot}{\kern0pt}return\ {\isacharparenleft}{\kern0pt}insert{\isacharunderscore}{\kern0pt}timed{\isacharunderscore}{\kern0pt}ground{\isacharunderscore}{\kern0pt}acts{\isacharunderscore}{\kern0pt}to{\isacharunderscore}{\kern0pt}tree\ as\ htree{\isacharparenright}{\kern0pt}\isanewline
\ \ \ \ {\isacharbraceright}{\kern0pt}{\isachardoublequoteclose}\isanewline
\isanewline
\ \ \isacommand{fun}\isamarkupfalse%
\ simplify{\isacharunderscore}{\kern0pt}planE{\isacharunderscore}{\kern0pt}avl\ {\isacharcolon}{\kern0pt}{\isacharcolon}{\kern0pt}\ {\isachardoublequoteopen}{\isacharunderscore}{\kern0pt}\ {\isasymRightarrow}\ {\isacharparenleft}{\kern0pt}object{\isacharcomma}{\kern0pt}\ type{\isacharparenright}{\kern0pt}\ mapping\ {\isasymRightarrow}\ {\isacharparenleft}{\kern0pt}object{\isacharcomma}{\kern0pt}\ rat{\isacharparenright}{\kern0pt}\ mapping\ {\isasymRightarrow}\ time\ list\ {\isasymRightarrow}\ plan\ {\isasymRightarrow}\ {\isacharunderscore}{\kern0pt}{\isacharplus}{\kern0pt}happening\ list{\isachardoublequoteclose}\ \isakeyword{where}\isanewline
\ \ \ \ {\isachardoublequoteopen}simplify{\isacharunderscore}{\kern0pt}planE{\isacharunderscore}{\kern0pt}avl\ G\ mp\ mp{\isacharunderscore}{\kern0pt}fe\ htps\ {\isasympi}s\ {\isacharequal}{\kern0pt}\ do\ {\isacharbraceleft}{\kern0pt}\ \isanewline
\ \ \ \ \ \ htree\ {\isasymleftarrow}\ simplify{\isacharunderscore}{\kern0pt}planE{\isacharunderscore}{\kern0pt}avl{\isacharprime}{\kern0pt}\ G\ mp\ mp{\isacharunderscore}{\kern0pt}fe\ htps\ {\isasympi}s{\isacharsemicolon}{\kern0pt}\isanewline
\ \ \ \ \ \ Error{\isacharunderscore}{\kern0pt}Monad{\isachardot}{\kern0pt}return\ {\isacharparenleft}{\kern0pt}avl{\isacharunderscore}{\kern0pt}inorder\ htree{\isacharparenright}{\kern0pt}\isanewline
\ \ \ \ {\isacharbraceright}{\kern0pt}{\isachardoublequoteclose}\isanewline
\isanewline
\ \ \isacommand{lemma}\isamarkupfalse%
\ insort{\isacharunderscore}{\kern0pt}happ{\isacharunderscore}{\kern0pt}consL{\isacharcolon}{\kern0pt}\isanewline
\ \ \ \ \isakeyword{assumes}\ {\isachardoublequoteopen}{\isasymforall}{\isacharparenleft}{\kern0pt}t\isactrlsub j{\isacharcomma}{\kern0pt}A\isactrlsub j{\isacharparenright}{\kern0pt}\ {\isasymin}\ set\ hs\isactrlsub {\isadigit{1}}{\isachardot}{\kern0pt}\ t\isactrlsub j\ {\isacharless}{\kern0pt}\ t\isactrlsub i{\isachardoublequoteclose}\isanewline
\ \ \ \ \isakeyword{shows}\ {\isachardoublequoteopen}hs\isactrlsub {\isadigit{1}}\ {\isacharat}{\kern0pt}\ {\isacharparenleft}{\kern0pt}insort{\isacharunderscore}{\kern0pt}happ\ {\isacharparenleft}{\kern0pt}t\isactrlsub i{\isacharcomma}{\kern0pt}A\isactrlsub i{\isacharparenright}{\kern0pt}\ hs\isactrlsub {\isadigit{2}}{\isacharparenright}{\kern0pt}\ {\isacharequal}{\kern0pt}\ insort{\isacharunderscore}{\kern0pt}happ\ {\isacharparenleft}{\kern0pt}t\isactrlsub i{\isacharcomma}{\kern0pt}A\isactrlsub i{\isacharparenright}{\kern0pt}\ {\isacharparenleft}{\kern0pt}hs\isactrlsub {\isadigit{1}}\ {\isacharat}{\kern0pt}\ hs\isactrlsub {\isadigit{2}}{\isacharparenright}{\kern0pt}{\isachardoublequoteclose}\isanewline
%
\isadelimproof
\ \ \ \ %
\endisadelimproof
%
\isatagproof
\isacommand{using}\isamarkupfalse%
\ assms\ \isacommand{by}\isamarkupfalse%
\ {\isacharparenleft}{\kern0pt}induction\ hs\isactrlsub {\isadigit{1}}\ arbitrary{\isacharcolon}{\kern0pt}\ t\isactrlsub i\ A\isactrlsub i\ hs\isactrlsub {\isadigit{2}}{\isacharparenright}{\kern0pt}\ auto%
\endisatagproof
{\isafoldproof}%
%
\isadelimproof
\isanewline
%
\endisadelimproof
\isanewline
\ \ \isacommand{lemma}\isamarkupfalse%
\ insort{\isacharunderscore}{\kern0pt}happ{\isacharunderscore}{\kern0pt}consR{\isacharcolon}{\kern0pt}\isanewline
\ \ \ \ \isakeyword{assumes}\ {\isachardoublequoteopen}{\isasymforall}{\isacharparenleft}{\kern0pt}t\isactrlsub j{\isacharcomma}{\kern0pt}A\isactrlsub j{\isacharparenright}{\kern0pt}\ {\isasymin}\ set\ hs\isactrlsub {\isadigit{2}}{\isachardot}{\kern0pt}\ t\isactrlsub i\ {\isacharless}{\kern0pt}\ t\isactrlsub j{\isachardoublequoteclose}\isanewline
\ \ \ \ \isakeyword{shows}\ {\isachardoublequoteopen}{\isacharparenleft}{\kern0pt}insort{\isacharunderscore}{\kern0pt}happ\ {\isacharparenleft}{\kern0pt}t\isactrlsub i{\isacharcomma}{\kern0pt}A\isactrlsub i{\isacharparenright}{\kern0pt}\ hs\isactrlsub {\isadigit{1}}{\isacharparenright}{\kern0pt}\ {\isacharat}{\kern0pt}\ hs\isactrlsub {\isadigit{2}}\ {\isacharequal}{\kern0pt}\ insort{\isacharunderscore}{\kern0pt}happ\ {\isacharparenleft}{\kern0pt}t\isactrlsub i{\isacharcomma}{\kern0pt}A\isactrlsub i{\isacharparenright}{\kern0pt}\ {\isacharparenleft}{\kern0pt}hs\isactrlsub {\isadigit{1}}\ {\isacharat}{\kern0pt}\ hs\isactrlsub {\isadigit{2}}{\isacharparenright}{\kern0pt}{\isachardoublequoteclose}\isanewline
%
\isadelimproof
\ \ \ \ %
\endisadelimproof
%
\isatagproof
\isacommand{using}\isamarkupfalse%
\ assms\isanewline
\ \ \isacommand{proof}\isamarkupfalse%
\ {\isacharparenleft}{\kern0pt}induction\ hs\isactrlsub {\isadigit{2}}\ arbitrary{\isacharcolon}{\kern0pt}\ t\isactrlsub i\ A\isactrlsub i\ hs\isactrlsub {\isadigit{1}}{\isacharparenright}{\kern0pt}\isanewline
\ \ \ \ \isacommand{case}\isamarkupfalse%
\ Nil\isanewline
\ \ \ \ \isacommand{then}\isamarkupfalse%
\ \isacommand{show}\isamarkupfalse%
\ {\isacharquery}{\kern0pt}case\ \isacommand{by}\isamarkupfalse%
\ auto\isanewline
\ \ \isacommand{next}\isamarkupfalse%
\isanewline
\ \ \ \ \isacommand{case}\isamarkupfalse%
\ {\isacharparenleft}{\kern0pt}Cons\ h\ hs\isactrlsub {\isadigit{2}}{\isacharparenright}{\kern0pt}\isanewline
\ \ \ \ \isacommand{then}\isamarkupfalse%
\ \isacommand{show}\isamarkupfalse%
\ {\isacharquery}{\kern0pt}case\ \isacommand{by}\isamarkupfalse%
\ {\isacharparenleft}{\kern0pt}induction\ hs\isactrlsub {\isadigit{1}}{\isacharparenright}{\kern0pt}\ auto\isanewline
\ \ \isacommand{qed}\isamarkupfalse%
%
\endisatagproof
{\isafoldproof}%
%
\isadelimproof
\isanewline
%
\endisadelimproof
\isanewline
\ \ \isacommand{lemma}\isamarkupfalse%
\ inorder{\isacharunderscore}{\kern0pt}avl{\isacharunderscore}{\kern0pt}balL{\isacharcolon}{\kern0pt}\ {\isachardoublequoteopen}avl{\isacharunderscore}{\kern0pt}inorder\ {\isacharparenleft}{\kern0pt}avl{\isacharunderscore}{\kern0pt}balL\ l\ a\ r{\isacharparenright}{\kern0pt}\ {\isacharequal}{\kern0pt}\ avl{\isacharunderscore}{\kern0pt}inorder\ l\ {\isacharat}{\kern0pt}\ a\ {\isacharhash}{\kern0pt}\ avl{\isacharunderscore}{\kern0pt}inorder\ r{\isachardoublequoteclose}\isanewline
%
\isadelimproof
\ \ \ \ %
\endisadelimproof
%
\isatagproof
\isacommand{unfolding}\isamarkupfalse%
\ avl{\isacharunderscore}{\kern0pt}inorder{\isacharunderscore}{\kern0pt}def\ avl{\isacharunderscore}{\kern0pt}node{\isacharunderscore}{\kern0pt}def\ avl{\isacharunderscore}{\kern0pt}balL{\isacharunderscore}{\kern0pt}def\ \isacommand{by}\isamarkupfalse%
\ {\isacharparenleft}{\kern0pt}auto\ split{\isacharcolon}{\kern0pt}\ tree{\isachardot}{\kern0pt}splits{\isacharparenright}{\kern0pt}%
\endisatagproof
{\isafoldproof}%
%
\isadelimproof
\isanewline
%
\endisadelimproof
\isanewline
\ \ \isacommand{lemma}\isamarkupfalse%
\ inorder{\isacharunderscore}{\kern0pt}avl{\isacharunderscore}{\kern0pt}balR{\isacharcolon}{\kern0pt}\ {\isachardoublequoteopen}avl{\isacharunderscore}{\kern0pt}inorder\ {\isacharparenleft}{\kern0pt}avl{\isacharunderscore}{\kern0pt}balR\ l\ a\ r{\isacharparenright}{\kern0pt}\ {\isacharequal}{\kern0pt}\ avl{\isacharunderscore}{\kern0pt}inorder\ l\ {\isacharat}{\kern0pt}\ a\ {\isacharhash}{\kern0pt}\ avl{\isacharunderscore}{\kern0pt}inorder\ r{\isachardoublequoteclose}\isanewline
%
\isadelimproof
\ \ \ \ %
\endisadelimproof
%
\isatagproof
\isacommand{unfolding}\isamarkupfalse%
\ avl{\isacharunderscore}{\kern0pt}inorder{\isacharunderscore}{\kern0pt}def\ avl{\isacharunderscore}{\kern0pt}node{\isacharunderscore}{\kern0pt}def\ avl{\isacharunderscore}{\kern0pt}balR{\isacharunderscore}{\kern0pt}def\ \isacommand{by}\isamarkupfalse%
\ {\isacharparenleft}{\kern0pt}auto\ split{\isacharcolon}{\kern0pt}\ tree{\isachardot}{\kern0pt}splits{\isacharparenright}{\kern0pt}%
\endisatagproof
{\isafoldproof}%
%
\isadelimproof
\isanewline
%
\endisadelimproof
\isanewline
\ \ \isanewline
\isanewline
\ \ \isacommand{lemma}\isamarkupfalse%
\ strict{\isacharunderscore}{\kern0pt}sorted{\isacharunderscore}{\kern0pt}appl{\isacharcolon}{\kern0pt}\ {\isachardoublequoteopen}strict{\isacharunderscore}{\kern0pt}sorted\ {\isacharparenleft}{\kern0pt}a\ {\isacharat}{\kern0pt}\ {\isacharbrackleft}{\kern0pt}y{\isacharbrackright}{\kern0pt}{\isacharparenright}{\kern0pt}\ {\isasymLongrightarrow}\ {\isasymforall}x\ {\isasymin}\ set\ a{\isachardot}{\kern0pt}\ x\ {\isacharless}{\kern0pt}\ y{\isachardoublequoteclose}\isanewline
%
\isadelimproof
\ \ \ \ %
\endisadelimproof
%
\isatagproof
\isacommand{by}\isamarkupfalse%
\ {\isacharparenleft}{\kern0pt}induction\ a{\isacharparenright}{\kern0pt}\ auto%
\endisatagproof
{\isafoldproof}%
%
\isadelimproof
\isanewline
%
\endisadelimproof
\isanewline
\ \ \isacommand{lemma}\isamarkupfalse%
\ strict{\isacharunderscore}{\kern0pt}sorted{\isacharunderscore}{\kern0pt}consg{\isacharcolon}{\kern0pt}\ {\isachardoublequoteopen}strict{\isacharunderscore}{\kern0pt}sorted\ {\isacharparenleft}{\kern0pt}y\ {\isacharhash}{\kern0pt}\ a{\isacharparenright}{\kern0pt}\ {\isasymLongrightarrow}\ {\isasymforall}x\ {\isasymin}\ set\ a{\isachardot}{\kern0pt}\ y\ {\isacharless}{\kern0pt}\ x{\isachardoublequoteclose}\isanewline
%
\isadelimproof
\ \ \ \ %
\endisadelimproof
%
\isatagproof
\isacommand{by}\isamarkupfalse%
\ {\isacharparenleft}{\kern0pt}induction\ a{\isacharparenright}{\kern0pt}\ auto%
\endisatagproof
{\isafoldproof}%
%
\isadelimproof
\isanewline
%
\endisadelimproof
\isanewline
\ \ \isacommand{lemma}\isamarkupfalse%
\ strict{\isacharunderscore}{\kern0pt}sorted{\isacharunderscore}{\kern0pt}avl{\isacharunderscore}{\kern0pt}inorder{\isacharunderscore}{\kern0pt}l{\isacharcolon}{\kern0pt}\isanewline
\ \ \ \ \isakeyword{assumes}\ {\isachardoublequoteopen}strict{\isacharunderscore}{\kern0pt}sorted\ {\isacharparenleft}{\kern0pt}map\ fst\ {\isacharparenleft}{\kern0pt}avl{\isacharunderscore}{\kern0pt}inorder\ {\isasymlangle}l{\isacharcomma}{\kern0pt}\ {\isacharparenleft}{\kern0pt}{\isacharparenleft}{\kern0pt}t\isactrlsub i{\isacharcomma}{\kern0pt}\ A\isactrlsub i{\isacharparenright}{\kern0pt}{\isacharcomma}{\kern0pt}\ h{\isacharparenright}{\kern0pt}{\isacharcomma}{\kern0pt}\ r{\isasymrangle}{\isacharparenright}{\kern0pt}{\isacharparenright}{\kern0pt}{\isachardoublequoteclose}\isanewline
\ \ \ \ \isakeyword{shows}\ {\isachardoublequoteopen}{\isasymforall}{\isacharparenleft}{\kern0pt}t\isactrlsub j{\isacharcomma}{\kern0pt}\ A\isactrlsub j{\isacharparenright}{\kern0pt}\ {\isasymin}\ set\ {\isacharparenleft}{\kern0pt}avl{\isacharunderscore}{\kern0pt}inorder\ l{\isacharparenright}{\kern0pt}{\isachardot}{\kern0pt}\ t\isactrlsub j\ {\isacharless}{\kern0pt}\ t\isactrlsub i{\isachardoublequoteclose}\isanewline
%
\isadelimproof
\ \ %
\endisadelimproof
%
\isatagproof
\isacommand{proof}\isamarkupfalse%
\ {\isacharminus}{\kern0pt}\isanewline
\ \ \ \ \isacommand{have}\isamarkupfalse%
\ {\isachardoublequoteopen}strict{\isacharunderscore}{\kern0pt}sorted\ {\isacharparenleft}{\kern0pt}map\ fst\ {\isacharparenleft}{\kern0pt}avl{\isacharunderscore}{\kern0pt}inorder\ l{\isacharparenright}{\kern0pt}\ {\isacharat}{\kern0pt}\ {\isacharbrackleft}{\kern0pt}t\isactrlsub i{\isacharbrackright}{\kern0pt}{\isacharparenright}{\kern0pt}{\isachardoublequoteclose}\isanewline
\ \ \ \ \ \ \isacommand{using}\isamarkupfalse%
\ assms\ strict{\isacharunderscore}{\kern0pt}sorted{\isacharunderscore}{\kern0pt}app{\isacharunderscore}{\kern0pt}iff{\isacharbrackleft}{\kern0pt}\isakeyword{where}\ a{\isacharequal}{\kern0pt}{\isachardoublequoteopen}map\ fst\ {\isacharparenleft}{\kern0pt}avl{\isacharunderscore}{\kern0pt}inorder\ l{\isacharparenright}{\kern0pt}\ {\isacharat}{\kern0pt}\ {\isacharbrackleft}{\kern0pt}t\isactrlsub i{\isacharbrackright}{\kern0pt}{\isachardoublequoteclose}{\isacharbrackright}{\kern0pt}\ \isanewline
\ \ \ \ \ \ \isacommand{by}\isamarkupfalse%
\ {\isacharparenleft}{\kern0pt}auto\ simp{\isacharcolon}{\kern0pt}\ avl{\isacharunderscore}{\kern0pt}inorder{\isacharunderscore}{\kern0pt}def{\isacharparenright}{\kern0pt}\isanewline
\ \ \ \ \isacommand{then}\isamarkupfalse%
\ \isacommand{show}\isamarkupfalse%
\ {\isacharquery}{\kern0pt}thesis\isanewline
\ \ \ \ \ \ \isacommand{using}\isamarkupfalse%
\ strict{\isacharunderscore}{\kern0pt}sorted{\isacharunderscore}{\kern0pt}appl{\isacharbrackleft}{\kern0pt}OF\ {\isacartoucheopen}strict{\isacharunderscore}{\kern0pt}sorted\ {\isacharparenleft}{\kern0pt}map\ fst\ {\isacharparenleft}{\kern0pt}avl{\isacharunderscore}{\kern0pt}inorder\ l{\isacharparenright}{\kern0pt}\ {\isacharat}{\kern0pt}\ {\isacharbrackleft}{\kern0pt}t\isactrlsub i{\isacharbrackright}{\kern0pt}{\isacharparenright}{\kern0pt}{\isacartoucheclose}{\isacharbrackright}{\kern0pt}\ \isacommand{by}\isamarkupfalse%
\ auto\isanewline
\ \ \isacommand{qed}\isamarkupfalse%
%
\endisatagproof
{\isafoldproof}%
%
\isadelimproof
\isanewline
%
\endisadelimproof
\isanewline
\ \ \isacommand{lemma}\isamarkupfalse%
\ strict{\isacharunderscore}{\kern0pt}sorted{\isacharunderscore}{\kern0pt}avl{\isacharunderscore}{\kern0pt}inorder{\isacharunderscore}{\kern0pt}r{\isacharcolon}{\kern0pt}\isanewline
\ \ \ \ \isakeyword{assumes}\ {\isachardoublequoteopen}strict{\isacharunderscore}{\kern0pt}sorted\ {\isacharparenleft}{\kern0pt}map\ fst\ {\isacharparenleft}{\kern0pt}avl{\isacharunderscore}{\kern0pt}inorder\ {\isasymlangle}l{\isacharcomma}{\kern0pt}\ {\isacharparenleft}{\kern0pt}{\isacharparenleft}{\kern0pt}t\isactrlsub i{\isacharcomma}{\kern0pt}\ A\isactrlsub i{\isacharparenright}{\kern0pt}{\isacharcomma}{\kern0pt}\ h{\isacharparenright}{\kern0pt}{\isacharcomma}{\kern0pt}\ r{\isasymrangle}{\isacharparenright}{\kern0pt}{\isacharparenright}{\kern0pt}{\isachardoublequoteclose}\isanewline
\ \ \ \ \isakeyword{shows}\ {\isachardoublequoteopen}{\isasymforall}{\isacharparenleft}{\kern0pt}t\isactrlsub j{\isacharcomma}{\kern0pt}\ A\isactrlsub j{\isacharparenright}{\kern0pt}\ {\isasymin}\ set\ {\isacharparenleft}{\kern0pt}avl{\isacharunderscore}{\kern0pt}inorder\ r{\isacharparenright}{\kern0pt}{\isachardot}{\kern0pt}\ t\isactrlsub i\ {\isacharless}{\kern0pt}\ t\isactrlsub j{\isachardoublequoteclose}\isanewline
%
\isadelimproof
\ \ %
\endisadelimproof
%
\isatagproof
\isacommand{proof}\isamarkupfalse%
\ {\isacharminus}{\kern0pt}\isanewline
\ \ \ \ \isacommand{have}\isamarkupfalse%
\ {\isachardoublequoteopen}strict{\isacharunderscore}{\kern0pt}sorted\ {\isacharparenleft}{\kern0pt}t\isactrlsub i\ {\isacharhash}{\kern0pt}\ map\ fst\ {\isacharparenleft}{\kern0pt}avl{\isacharunderscore}{\kern0pt}inorder\ r{\isacharparenright}{\kern0pt}{\isacharparenright}{\kern0pt}{\isachardoublequoteclose}\isanewline
\ \ \ \ \ \ \isacommand{using}\isamarkupfalse%
\ assms\ strict{\isacharunderscore}{\kern0pt}sorted{\isacharunderscore}{\kern0pt}app{\isacharunderscore}{\kern0pt}iff{\isacharbrackleft}{\kern0pt}\isakeyword{where}\ b{\isacharequal}{\kern0pt}{\isachardoublequoteopen}t\isactrlsub i\ {\isacharhash}{\kern0pt}\ map\ fst\ {\isacharparenleft}{\kern0pt}avl{\isacharunderscore}{\kern0pt}inorder\ r{\isacharparenright}{\kern0pt}{\isachardoublequoteclose}{\isacharbrackright}{\kern0pt}\ \isanewline
\ \ \ \ \ \ \isacommand{by}\isamarkupfalse%
\ {\isacharparenleft}{\kern0pt}auto\ simp{\isacharcolon}{\kern0pt}\ avl{\isacharunderscore}{\kern0pt}inorder{\isacharunderscore}{\kern0pt}def{\isacharparenright}{\kern0pt}\isanewline
\ \ \ \ \isacommand{then}\isamarkupfalse%
\ \isacommand{show}\isamarkupfalse%
\ {\isacharquery}{\kern0pt}thesis\isanewline
\ \ \ \ \ \ \isacommand{using}\isamarkupfalse%
\ strict{\isacharunderscore}{\kern0pt}sorted{\isacharunderscore}{\kern0pt}consg{\isacharbrackleft}{\kern0pt}OF\ {\isacartoucheopen}strict{\isacharunderscore}{\kern0pt}sorted\ {\isacharparenleft}{\kern0pt}t\isactrlsub i\ {\isacharhash}{\kern0pt}\ map\ fst\ {\isacharparenleft}{\kern0pt}avl{\isacharunderscore}{\kern0pt}inorder\ r{\isacharparenright}{\kern0pt}{\isacharparenright}{\kern0pt}{\isacartoucheclose}{\isacharbrackright}{\kern0pt}\ \isacommand{by}\isamarkupfalse%
\ auto\isanewline
\ \ \isacommand{qed}\isamarkupfalse%
%
\endisatagproof
{\isafoldproof}%
%
\isadelimproof
\isanewline
%
\endisadelimproof
\isanewline
\ \ \isacommand{lemma}\isamarkupfalse%
\ insert{\isacharunderscore}{\kern0pt}happ{\isacharunderscore}{\kern0pt}to{\isacharunderscore}{\kern0pt}tree{\isacharunderscore}{\kern0pt}set{\isacharcolon}{\kern0pt}\isanewline
\ \ \ \ \isakeyword{assumes}\ {\isachardoublequoteopen}strict{\isacharunderscore}{\kern0pt}sorted\ {\isacharparenleft}{\kern0pt}map\ fst\ {\isacharparenleft}{\kern0pt}avl{\isacharunderscore}{\kern0pt}inorder\ htree{\isacharparenright}{\kern0pt}{\isacharparenright}{\kern0pt}{\isachardoublequoteclose}\isanewline
\ \ \ \ \isakeyword{shows}\ {\isachardoublequoteopen}set\ {\isacharparenleft}{\kern0pt}map\ fst\ {\isacharparenleft}{\kern0pt}avl{\isacharunderscore}{\kern0pt}inorder\ {\isacharparenleft}{\kern0pt}insert{\isacharunderscore}{\kern0pt}happ{\isacharunderscore}{\kern0pt}to{\isacharunderscore}{\kern0pt}tree\ {\isacharparenleft}{\kern0pt}t\isactrlsub i{\isacharcomma}{\kern0pt}A\isactrlsub i{\isacharparenright}{\kern0pt}\ htree{\isacharparenright}{\kern0pt}{\isacharparenright}{\kern0pt}{\isacharparenright}{\kern0pt}\ {\isacharequal}{\kern0pt}\ set\ {\isacharparenleft}{\kern0pt}map\ fst\ {\isacharparenleft}{\kern0pt}avl{\isacharunderscore}{\kern0pt}inorder\ htree{\isacharparenright}{\kern0pt}{\isacharparenright}{\kern0pt}\ {\isasymunion}\ {\isacharbraceleft}{\kern0pt}t\isactrlsub i{\isacharbraceright}{\kern0pt}{\isachardoublequoteclose}\isanewline
%
\isadelimproof
\ \ \ \ %
\endisadelimproof
%
\isatagproof
\isacommand{using}\isamarkupfalse%
\ assms\isanewline
\ \ \isacommand{proof}\isamarkupfalse%
\ {\isacharparenleft}{\kern0pt}induction\ {\isachardoublequoteopen}{\isacharparenleft}{\kern0pt}t\isactrlsub i{\isacharcomma}{\kern0pt}A\isactrlsub i{\isacharparenright}{\kern0pt}{\isachardoublequoteclose}\ htree\ rule{\isacharcolon}{\kern0pt}\ insert{\isacharunderscore}{\kern0pt}happ{\isacharunderscore}{\kern0pt}to{\isacharunderscore}{\kern0pt}tree{\isachardot}{\kern0pt}induct{\isacharparenright}{\kern0pt}\isanewline
\ \ \ \ \isacommand{case}\isamarkupfalse%
\ {\isadigit{1}}\isanewline
\ \ \ \ \isacommand{then}\isamarkupfalse%
\ \isacommand{show}\isamarkupfalse%
\ {\isacharquery}{\kern0pt}case\ \isacommand{by}\isamarkupfalse%
\ {\isacharparenleft}{\kern0pt}auto\ simp{\isacharcolon}{\kern0pt}\ avl{\isacharunderscore}{\kern0pt}inorder{\isacharunderscore}{\kern0pt}def\ avl{\isacharunderscore}{\kern0pt}node{\isacharunderscore}{\kern0pt}def{\isacharparenright}{\kern0pt}\isanewline
\ \ \isacommand{next}\isamarkupfalse%
\isanewline
\ \ \ \ \isacommand{case}\isamarkupfalse%
\ {\isacharparenleft}{\kern0pt}{\isadigit{2}}\ l\ t\isactrlsub j\ A\isactrlsub j\ h\ r{\isacharparenright}{\kern0pt}\isanewline
\ \ \ \ \isacommand{then}\isamarkupfalse%
\ \isacommand{have}\isamarkupfalse%
\ {\isachardoublequoteopen}strict{\isacharunderscore}{\kern0pt}sorted\ {\isacharparenleft}{\kern0pt}map\ fst\ {\isacharparenleft}{\kern0pt}avl{\isacharunderscore}{\kern0pt}inorder\ {\isasymlangle}l{\isacharcomma}{\kern0pt}\ {\isacharparenleft}{\kern0pt}{\isacharparenleft}{\kern0pt}t\isactrlsub j{\isacharcomma}{\kern0pt}\ A\isactrlsub j{\isacharparenright}{\kern0pt}{\isacharcomma}{\kern0pt}\ h{\isacharparenright}{\kern0pt}{\isacharcomma}{\kern0pt}\ r{\isasymrangle}{\isacharparenright}{\kern0pt}{\isacharparenright}{\kern0pt}{\isachardoublequoteclose}\isanewline
\ \ \ \ \ \ \isacommand{by}\isamarkupfalse%
\ auto\isanewline
\ \ \ \ \isacommand{then}\isamarkupfalse%
\ \isacommand{have}\isamarkupfalse%
\ {\isachardoublequoteopen}strict{\isacharunderscore}{\kern0pt}sorted\ {\isacharparenleft}{\kern0pt}map\ fst\ {\isacharparenleft}{\kern0pt}avl{\isacharunderscore}{\kern0pt}inorder\ l\ {\isacharat}{\kern0pt}\ {\isacharparenleft}{\kern0pt}t\isactrlsub j{\isacharcomma}{\kern0pt}\ A\isactrlsub j{\isacharparenright}{\kern0pt}\ {\isacharhash}{\kern0pt}\ \ avl{\isacharunderscore}{\kern0pt}inorder\ r{\isacharparenright}{\kern0pt}{\isacharparenright}{\kern0pt}{\isachardoublequoteclose}\isanewline
\ \ \ \ \ \ \isacommand{unfolding}\isamarkupfalse%
\ avl{\isacharunderscore}{\kern0pt}inorder{\isacharunderscore}{\kern0pt}def\ \isacommand{by}\isamarkupfalse%
\ auto\isanewline
\ \ \ \ \isacommand{then}\isamarkupfalse%
\ \isacommand{have}\isamarkupfalse%
\ {\isachardoublequoteopen}strict{\isacharunderscore}{\kern0pt}sorted\ {\isacharparenleft}{\kern0pt}map\ fst\ {\isacharparenleft}{\kern0pt}avl{\isacharunderscore}{\kern0pt}inorder\ l{\isacharparenright}{\kern0pt}\ {\isacharat}{\kern0pt}\ t\isactrlsub j\ {\isacharhash}{\kern0pt}\ map\ fst\ {\isacharparenleft}{\kern0pt}avl{\isacharunderscore}{\kern0pt}inorder\ r{\isacharparenright}{\kern0pt}{\isacharparenright}{\kern0pt}{\isachardoublequoteclose}\isanewline
\ \ \ \ \ \ \ \ \ \ \isakeyword{and}\ {\isachardoublequoteopen}strict{\isacharunderscore}{\kern0pt}sorted\ {\isacharparenleft}{\kern0pt}map\ fst\ {\isacharparenleft}{\kern0pt}avl{\isacharunderscore}{\kern0pt}inorder\ l{\isacharparenright}{\kern0pt}\ {\isacharat}{\kern0pt}\ {\isacharbrackleft}{\kern0pt}t\isactrlsub j{\isacharbrackright}{\kern0pt}\ {\isacharat}{\kern0pt}\ map\ fst\ {\isacharparenleft}{\kern0pt}avl{\isacharunderscore}{\kern0pt}inorder\ r{\isacharparenright}{\kern0pt}{\isacharparenright}{\kern0pt}{\isachardoublequoteclose}\isanewline
\ \ \ \ \ \ \isacommand{by}\isamarkupfalse%
\ auto\isanewline
\ \ \ \ \isacommand{then}\isamarkupfalse%
\ \isacommand{have}\isamarkupfalse%
\ {\isachardoublequoteopen}strict{\isacharunderscore}{\kern0pt}sorted\ {\isacharparenleft}{\kern0pt}map\ fst\ {\isacharparenleft}{\kern0pt}avl{\isacharunderscore}{\kern0pt}inorder\ l{\isacharparenright}{\kern0pt}{\isacharparenright}{\kern0pt}{\isachardoublequoteclose}\ \isanewline
\ \ \ \ \ \ \ \ \ \ \isakeyword{and}\ {\isachardoublequoteopen}strict{\isacharunderscore}{\kern0pt}sorted\ {\isacharparenleft}{\kern0pt}map\ fst\ {\isacharparenleft}{\kern0pt}avl{\isacharunderscore}{\kern0pt}inorder\ r{\isacharparenright}{\kern0pt}{\isacharparenright}{\kern0pt}{\isachardoublequoteclose}\isanewline
\ \ \ \ \ \ \isacommand{using}\isamarkupfalse%
\ strict{\isacharunderscore}{\kern0pt}sorted{\isacharunderscore}{\kern0pt}app{\isacharunderscore}{\kern0pt}iff{\isacharbrackleft}{\kern0pt}\isakeyword{where}\ a{\isacharequal}{\kern0pt}{\isachardoublequoteopen}map\ fst\ {\isacharparenleft}{\kern0pt}avl{\isacharunderscore}{\kern0pt}inorder\ l{\isacharparenright}{\kern0pt}{\isachardoublequoteclose}\ \isakeyword{and}\ b{\isacharequal}{\kern0pt}{\isachardoublequoteopen}t\isactrlsub j\ {\isacharhash}{\kern0pt}\ map\ fst\ {\isacharparenleft}{\kern0pt}avl{\isacharunderscore}{\kern0pt}inorder\ r{\isacharparenright}{\kern0pt}{\isachardoublequoteclose}{\isacharbrackright}{\kern0pt}\ \ \ \ \ \ \ \ \ \ \ \ \isanewline
\ \ \ \ \ \ \ \ \ \ \ \ strict{\isacharunderscore}{\kern0pt}sorted{\isacharunderscore}{\kern0pt}app{\isacharunderscore}{\kern0pt}iff{\isacharbrackleft}{\kern0pt}\isakeyword{where}\ a{\isacharequal}{\kern0pt}{\isachardoublequoteopen}map\ fst\ {\isacharparenleft}{\kern0pt}avl{\isacharunderscore}{\kern0pt}inorder\ l{\isacharparenright}{\kern0pt}\ {\isacharat}{\kern0pt}\ {\isacharbrackleft}{\kern0pt}t\isactrlsub j{\isacharbrackright}{\kern0pt}{\isachardoublequoteclose}\ \isakeyword{and}\ b{\isacharequal}{\kern0pt}{\isachardoublequoteopen}map\ fst\ {\isacharparenleft}{\kern0pt}avl{\isacharunderscore}{\kern0pt}inorder\ r{\isacharparenright}{\kern0pt}{\isachardoublequoteclose}{\isacharbrackright}{\kern0pt}\isanewline
\ \ \ \ \ \ \isacommand{by}\isamarkupfalse%
\ auto\isanewline
\ \ \ \ \isacommand{have}\isamarkupfalse%
\ {\isachardoublequoteopen}t\isactrlsub i\ {\isacharless}{\kern0pt}\ t\isactrlsub j\ {\isasymor}\ t\isactrlsub i\ {\isacharequal}{\kern0pt}\ t\isactrlsub j\ {\isasymor}\ t\isactrlsub j\ {\isacharless}{\kern0pt}\ t\isactrlsub i{\isachardoublequoteclose}\isanewline
\ \ \ \ \ \ \isacommand{by}\isamarkupfalse%
\ auto\isanewline
\ \ \ \ \isacommand{then}\isamarkupfalse%
\ \isacommand{show}\isamarkupfalse%
\ {\isacharquery}{\kern0pt}case\isanewline
\ \ \ \ \isacommand{proof}\isamarkupfalse%
\ {\isacharparenleft}{\kern0pt}elim\ disjE{\isacharparenright}{\kern0pt}\isanewline
\ \ \ \ \ \ \isacommand{assume}\isamarkupfalse%
\ {\isachardoublequoteopen}t\isactrlsub i\ {\isacharless}{\kern0pt}\ t\isactrlsub j{\isachardoublequoteclose}\isanewline
\ \ \ \ \ \ \isacommand{then}\isamarkupfalse%
\ \isacommand{show}\isamarkupfalse%
\ {\isacharquery}{\kern0pt}case\isanewline
\ \ \ \ \ \ \ \ \isacommand{using}\isamarkupfalse%
\ inorder{\isacharunderscore}{\kern0pt}avl{\isacharunderscore}{\kern0pt}balL\ {\isacartoucheopen}strict{\isacharunderscore}{\kern0pt}sorted\ {\isacharparenleft}{\kern0pt}map\ fst\ {\isacharparenleft}{\kern0pt}avl{\isacharunderscore}{\kern0pt}inorder\ l{\isacharparenright}{\kern0pt}{\isacharparenright}{\kern0pt}{\isacartoucheclose}\ {\isachardoublequoteopen}{\isadigit{2}}{\isachardot}{\kern0pt}hyps{\isachardoublequoteclose}\ \isacommand{by}\isamarkupfalse%
\ {\isacharparenleft}{\kern0pt}auto\ simp{\isacharcolon}{\kern0pt}\ avl{\isacharunderscore}{\kern0pt}inorder{\isacharunderscore}{\kern0pt}def{\isacharparenright}{\kern0pt}\isanewline
\ \ \ \ \isacommand{next}\isamarkupfalse%
\isanewline
\ \ \ \ \ \ \isacommand{assume}\isamarkupfalse%
\ {\isachardoublequoteopen}t\isactrlsub i\ {\isacharequal}{\kern0pt}\ t\isactrlsub j{\isachardoublequoteclose}\isanewline
\ \ \ \ \ \ \isacommand{then}\isamarkupfalse%
\ \isacommand{show}\isamarkupfalse%
\ {\isacharquery}{\kern0pt}case\isanewline
\ \ \ \ \ \ \ \ \isacommand{by}\isamarkupfalse%
\ {\isacharparenleft}{\kern0pt}auto\ simp{\isacharcolon}{\kern0pt}\ avl{\isacharunderscore}{\kern0pt}inorder{\isacharunderscore}{\kern0pt}def\ avl{\isacharunderscore}{\kern0pt}node{\isacharunderscore}{\kern0pt}def{\isacharparenright}{\kern0pt}\isanewline
\ \ \ \ \isacommand{next}\isamarkupfalse%
\isanewline
\ \ \ \ \ \ \isacommand{assume}\isamarkupfalse%
\ {\isachardoublequoteopen}t\isactrlsub j\ {\isacharless}{\kern0pt}\ t\isactrlsub i{\isachardoublequoteclose}\ \ \ \ \ \ \ \ \ \ \ \ \ \ \ \ \ \ \ \ \ \ \ \ \ \ \ \ \ \ \ \ \isanewline
\ \ \ \ \ \ \isacommand{then}\isamarkupfalse%
\ \isacommand{show}\isamarkupfalse%
\ {\isacharquery}{\kern0pt}case\isanewline
\ \ \ \ \ \ \ \ \isacommand{using}\isamarkupfalse%
\ inorder{\isacharunderscore}{\kern0pt}avl{\isacharunderscore}{\kern0pt}balR\ {\isacartoucheopen}strict{\isacharunderscore}{\kern0pt}sorted\ {\isacharparenleft}{\kern0pt}map\ fst\ {\isacharparenleft}{\kern0pt}avl{\isacharunderscore}{\kern0pt}inorder\ r{\isacharparenright}{\kern0pt}{\isacharparenright}{\kern0pt}{\isacartoucheclose}\ {\isachardoublequoteopen}{\isadigit{2}}{\isachardot}{\kern0pt}hyps{\isachardoublequoteclose}\ \isacommand{by}\isamarkupfalse%
\ {\isacharparenleft}{\kern0pt}auto\ simp{\isacharcolon}{\kern0pt}\ avl{\isacharunderscore}{\kern0pt}inorder{\isacharunderscore}{\kern0pt}def{\isacharparenright}{\kern0pt}\isanewline
\ \ \ \ \isacommand{qed}\isamarkupfalse%
\isanewline
\ \ \isacommand{qed}\isamarkupfalse%
%
\endisatagproof
{\isafoldproof}%
%
\isadelimproof
\isanewline
%
\endisadelimproof
\isanewline
\ \ \isacommand{lemma}\isamarkupfalse%
\ insert{\isacharunderscore}{\kern0pt}timed{\isacharunderscore}{\kern0pt}ground{\isacharunderscore}{\kern0pt}act{\isacharunderscore}{\kern0pt}to{\isacharunderscore}{\kern0pt}tree{\isacharunderscore}{\kern0pt}strict{\isacharunderscore}{\kern0pt}sorted{\isacharcolon}{\kern0pt}\ \isanewline
\ \ \ \ \isakeyword{assumes}\ {\isachardoublequoteopen}strict{\isacharunderscore}{\kern0pt}sorted\ {\isacharparenleft}{\kern0pt}map\ fst\ {\isacharparenleft}{\kern0pt}avl{\isacharunderscore}{\kern0pt}inorder\ htree{\isacharparenright}{\kern0pt}{\isacharparenright}{\kern0pt}{\isachardoublequoteclose}\isanewline
\ \ \ \ \isakeyword{shows}\ {\isachardoublequoteopen}strict{\isacharunderscore}{\kern0pt}sorted\ {\isacharparenleft}{\kern0pt}map\ fst\ {\isacharparenleft}{\kern0pt}avl{\isacharunderscore}{\kern0pt}inorder\ {\isacharparenleft}{\kern0pt}insert{\isacharunderscore}{\kern0pt}happ{\isacharunderscore}{\kern0pt}to{\isacharunderscore}{\kern0pt}tree\ h\ htree{\isacharparenright}{\kern0pt}{\isacharparenright}{\kern0pt}{\isacharparenright}{\kern0pt}{\isachardoublequoteclose}\isanewline
%
\isadelimproof
\ \ \ \ %
\endisadelimproof
%
\isatagproof
\isacommand{using}\isamarkupfalse%
\ assms\isanewline
\ \ \isacommand{proof}\isamarkupfalse%
\ {\isacharparenleft}{\kern0pt}induction\ h\ htree\ rule{\isacharcolon}{\kern0pt}\ insert{\isacharunderscore}{\kern0pt}happ{\isacharunderscore}{\kern0pt}to{\isacharunderscore}{\kern0pt}tree{\isachardot}{\kern0pt}induct{\isacharparenright}{\kern0pt}\isanewline
\ \ \ \ \isacommand{case}\isamarkupfalse%
\ {\isacharparenleft}{\kern0pt}{\isadigit{1}}\ t\isactrlsub i\ A\isactrlsub i{\isacharparenright}{\kern0pt}\isanewline
\ \ \ \ \isacommand{then}\isamarkupfalse%
\ \isacommand{show}\isamarkupfalse%
\ {\isacharquery}{\kern0pt}case\ \isanewline
\ \ \ \ \ \ \isacommand{by}\isamarkupfalse%
\ {\isacharparenleft}{\kern0pt}auto\ simp{\isacharcolon}{\kern0pt}\ avl{\isacharunderscore}{\kern0pt}node{\isacharunderscore}{\kern0pt}def\ avl{\isacharunderscore}{\kern0pt}inorder{\isacharunderscore}{\kern0pt}def{\isacharparenright}{\kern0pt}\ \ \ \ \ \ \ \ \ \ \ \ \ \ \ \ \ \ \ \ \ \ \ \ \ \ \ \ \ \ \ \ \ \ \ \ \ \ \ \ \ \ \ \ \ \ \ \ \ \ \ \ \ \ \ \ \ \ \ \ \ \ \ \ \ \ \ \ \ \ \ \ \ \ \ \ \ \ \ \ \ \ \ \ \ \ \ \ \ \isanewline
\ \ \isacommand{next}\isamarkupfalse%
\isanewline
\ \ \ \ \isacommand{case}\isamarkupfalse%
\ {\isacharparenleft}{\kern0pt}{\isadigit{2}}\ t\isactrlsub i\ A\isactrlsub i\ l\ t\isactrlsub j\ A\isactrlsub j\ h\ r{\isacharparenright}{\kern0pt}\isanewline
\ \ \ \ \isacommand{then}\isamarkupfalse%
\ \isacommand{have}\isamarkupfalse%
\ {\isachardoublequoteopen}strict{\isacharunderscore}{\kern0pt}sorted\ {\isacharparenleft}{\kern0pt}map\ fst\ {\isacharparenleft}{\kern0pt}avl{\isacharunderscore}{\kern0pt}inorder\ {\isasymlangle}l{\isacharcomma}{\kern0pt}\ {\isacharparenleft}{\kern0pt}{\isacharparenleft}{\kern0pt}t\isactrlsub j{\isacharcomma}{\kern0pt}\ A\isactrlsub j{\isacharparenright}{\kern0pt}{\isacharcomma}{\kern0pt}\ h{\isacharparenright}{\kern0pt}{\isacharcomma}{\kern0pt}\ r{\isasymrangle}{\isacharparenright}{\kern0pt}{\isacharparenright}{\kern0pt}{\isachardoublequoteclose}\isanewline
\ \ \ \ \ \ \isacommand{by}\isamarkupfalse%
\ auto\isanewline
\ \ \ \ \isacommand{then}\isamarkupfalse%
\ \isacommand{have}\isamarkupfalse%
\ {\isachardoublequoteopen}strict{\isacharunderscore}{\kern0pt}sorted\ {\isacharparenleft}{\kern0pt}map\ fst\ {\isacharparenleft}{\kern0pt}avl{\isacharunderscore}{\kern0pt}inorder\ l\ {\isacharat}{\kern0pt}\ {\isacharparenleft}{\kern0pt}t\isactrlsub j{\isacharcomma}{\kern0pt}\ A\isactrlsub j{\isacharparenright}{\kern0pt}\ {\isacharhash}{\kern0pt}\ avl{\isacharunderscore}{\kern0pt}inorder\ r{\isacharparenright}{\kern0pt}{\isacharparenright}{\kern0pt}{\isachardoublequoteclose}\isanewline
\ \ \ \ \ \ \isacommand{unfolding}\isamarkupfalse%
\ avl{\isacharunderscore}{\kern0pt}inorder{\isacharunderscore}{\kern0pt}def\ \isacommand{by}\isamarkupfalse%
\ auto\isanewline
\ \ \ \ \isacommand{then}\isamarkupfalse%
\ \isacommand{have}\isamarkupfalse%
\ {\isachardoublequoteopen}strict{\isacharunderscore}{\kern0pt}sorted\ {\isacharparenleft}{\kern0pt}map\ fst\ {\isacharparenleft}{\kern0pt}avl{\isacharunderscore}{\kern0pt}inorder\ l{\isacharparenright}{\kern0pt}\ {\isacharat}{\kern0pt}\ t\isactrlsub j\ {\isacharhash}{\kern0pt}\ map\ fst\ {\isacharparenleft}{\kern0pt}avl{\isacharunderscore}{\kern0pt}inorder\ r{\isacharparenright}{\kern0pt}{\isacharparenright}{\kern0pt}{\isachardoublequoteclose}\isanewline
\ \ \ \ \ \ \ \ \ \ \isakeyword{and}\ {\isachardoublequoteopen}strict{\isacharunderscore}{\kern0pt}sorted\ {\isacharparenleft}{\kern0pt}map\ fst\ {\isacharparenleft}{\kern0pt}avl{\isacharunderscore}{\kern0pt}inorder\ l{\isacharparenright}{\kern0pt}\ {\isacharat}{\kern0pt}\ {\isacharbrackleft}{\kern0pt}t\isactrlsub j{\isacharbrackright}{\kern0pt}\ {\isacharat}{\kern0pt}\ map\ fst\ {\isacharparenleft}{\kern0pt}avl{\isacharunderscore}{\kern0pt}inorder\ r{\isacharparenright}{\kern0pt}{\isacharparenright}{\kern0pt}{\isachardoublequoteclose}\isanewline
\ \ \ \ \ \ \isacommand{by}\isamarkupfalse%
\ auto\isanewline
\ \ \ \ \isacommand{then}\isamarkupfalse%
\ \isacommand{have}\isamarkupfalse%
\ {\isachardoublequoteopen}strict{\isacharunderscore}{\kern0pt}sorted\ {\isacharparenleft}{\kern0pt}map\ fst\ {\isacharparenleft}{\kern0pt}avl{\isacharunderscore}{\kern0pt}inorder\ l{\isacharparenright}{\kern0pt}{\isacharparenright}{\kern0pt}{\isachardoublequoteclose}\ \isanewline
\ \ \ \ \ \ \ \ \ \ \isakeyword{and}\ {\isachardoublequoteopen}strict{\isacharunderscore}{\kern0pt}sorted\ {\isacharparenleft}{\kern0pt}map\ fst\ {\isacharparenleft}{\kern0pt}avl{\isacharunderscore}{\kern0pt}inorder\ l\ {\isacharat}{\kern0pt}\ {\isacharbrackleft}{\kern0pt}{\isacharparenleft}{\kern0pt}t\isactrlsub j{\isacharcomma}{\kern0pt}A\isactrlsub j{\isacharparenright}{\kern0pt}{\isacharbrackright}{\kern0pt}{\isacharparenright}{\kern0pt}{\isacharparenright}{\kern0pt}{\isachardoublequoteclose}\isanewline
\ \ \ \ \ \ \ \ \ \ \isakeyword{and}\ {\isachardoublequoteopen}strict{\isacharunderscore}{\kern0pt}sorted\ {\isacharparenleft}{\kern0pt}map\ fst\ {\isacharparenleft}{\kern0pt}avl{\isacharunderscore}{\kern0pt}inorder\ r{\isacharparenright}{\kern0pt}{\isacharparenright}{\kern0pt}{\isachardoublequoteclose}\isanewline
\ \ \ \ \ \ \ \ \ \ \isakeyword{and}\ {\isachardoublequoteopen}strict{\isacharunderscore}{\kern0pt}sorted\ {\isacharparenleft}{\kern0pt}map\ fst\ {\isacharparenleft}{\kern0pt}{\isacharparenleft}{\kern0pt}t\isactrlsub j{\isacharcomma}{\kern0pt}A\isactrlsub j{\isacharparenright}{\kern0pt}\ {\isacharhash}{\kern0pt}\ avl{\isacharunderscore}{\kern0pt}inorder\ r{\isacharparenright}{\kern0pt}{\isacharparenright}{\kern0pt}{\isachardoublequoteclose}\isanewline
\ \ \ \ \ \ \isacommand{using}\isamarkupfalse%
\ strict{\isacharunderscore}{\kern0pt}sorted{\isacharunderscore}{\kern0pt}app{\isacharunderscore}{\kern0pt}iff{\isacharbrackleft}{\kern0pt}\isakeyword{where}\ a{\isacharequal}{\kern0pt}{\isachardoublequoteopen}map\ fst\ {\isacharparenleft}{\kern0pt}avl{\isacharunderscore}{\kern0pt}inorder\ l{\isacharparenright}{\kern0pt}{\isachardoublequoteclose}\ \isakeyword{and}\ b{\isacharequal}{\kern0pt}{\isachardoublequoteopen}t\isactrlsub j\ {\isacharhash}{\kern0pt}\ map\ fst\ {\isacharparenleft}{\kern0pt}avl{\isacharunderscore}{\kern0pt}inorder\ r{\isacharparenright}{\kern0pt}{\isachardoublequoteclose}{\isacharbrackright}{\kern0pt}\ \ \ \ \ \ \ \ \ \ \ \ \isanewline
\ \ \ \ \ \ \ \ \ \ \ \ strict{\isacharunderscore}{\kern0pt}sorted{\isacharunderscore}{\kern0pt}app{\isacharunderscore}{\kern0pt}iff{\isacharbrackleft}{\kern0pt}\isakeyword{where}\ a{\isacharequal}{\kern0pt}{\isachardoublequoteopen}map\ fst\ {\isacharparenleft}{\kern0pt}avl{\isacharunderscore}{\kern0pt}inorder\ l{\isacharparenright}{\kern0pt}\ {\isacharat}{\kern0pt}\ {\isacharbrackleft}{\kern0pt}t\isactrlsub j{\isacharbrackright}{\kern0pt}{\isachardoublequoteclose}\ \isakeyword{and}\ b{\isacharequal}{\kern0pt}{\isachardoublequoteopen}map\ fst\ {\isacharparenleft}{\kern0pt}avl{\isacharunderscore}{\kern0pt}inorder\ r{\isacharparenright}{\kern0pt}{\isachardoublequoteclose}{\isacharbrackright}{\kern0pt}\isanewline
\ \ \ \ \ \ \isacommand{by}\isamarkupfalse%
\ auto\isanewline
\ \ \ \ \isacommand{have}\isamarkupfalse%
\ {\isachardoublequoteopen}t\isactrlsub i\ {\isacharless}{\kern0pt}\ t\isactrlsub j\ {\isasymor}\ t\isactrlsub i\ {\isacharequal}{\kern0pt}\ t\isactrlsub j\ {\isasymor}\ t\isactrlsub j\ {\isacharless}{\kern0pt}\ t\isactrlsub i{\isachardoublequoteclose}\isanewline
\ \ \ \ \ \ \isacommand{by}\isamarkupfalse%
\ auto\isanewline
\ \ \ \ \isacommand{then}\isamarkupfalse%
\ \isacommand{show}\isamarkupfalse%
\ {\isacharquery}{\kern0pt}case\isanewline
\ \ \ \ \isacommand{proof}\isamarkupfalse%
\ {\isacharparenleft}{\kern0pt}elim\ disjE{\isacharparenright}{\kern0pt}\isanewline
\ \ \ \ \ \ \isacommand{assume}\isamarkupfalse%
\ {\isachardoublequoteopen}t\isactrlsub i\ {\isacharless}{\kern0pt}\ t\isactrlsub j{\isachardoublequoteclose}\isanewline
\ \ \ \ \ \ \isacommand{then}\isamarkupfalse%
\ \isacommand{have}\isamarkupfalse%
\ {\isachardoublequoteopen}strict{\isacharunderscore}{\kern0pt}sorted\ {\isacharparenleft}{\kern0pt}map\ fst\ {\isacharparenleft}{\kern0pt}avl{\isacharunderscore}{\kern0pt}inorder\ {\isacharparenleft}{\kern0pt}insert{\isacharunderscore}{\kern0pt}happ{\isacharunderscore}{\kern0pt}to{\isacharunderscore}{\kern0pt}tree\ {\isacharparenleft}{\kern0pt}t\isactrlsub i{\isacharcomma}{\kern0pt}A\isactrlsub i{\isacharparenright}{\kern0pt}\ l{\isacharparenright}{\kern0pt}{\isacharparenright}{\kern0pt}{\isacharparenright}{\kern0pt}{\isachardoublequoteclose}\isanewline
\ \ \ \ \ \ \ \ \isacommand{using}\isamarkupfalse%
\ {\isacartoucheopen}strict{\isacharunderscore}{\kern0pt}sorted\ {\isacharparenleft}{\kern0pt}map\ fst\ {\isacharparenleft}{\kern0pt}avl{\isacharunderscore}{\kern0pt}inorder\ l{\isacharparenright}{\kern0pt}{\isacharparenright}{\kern0pt}{\isacartoucheclose}\ {\isachardoublequoteopen}{\isadigit{2}}{\isachardot}{\kern0pt}IH{\isachardoublequoteclose}\ \isacommand{by}\isamarkupfalse%
\ auto\isanewline
\isanewline
\ \ \ \ \ \ \isacommand{have}\isamarkupfalse%
\ {\isachardoublequoteopen}strict{\isacharunderscore}{\kern0pt}sorted\ {\isacharparenleft}{\kern0pt}map\ fst\ {\isacharparenleft}{\kern0pt}avl{\isacharunderscore}{\kern0pt}inorder\ l{\isacharparenright}{\kern0pt}\ {\isacharat}{\kern0pt}\ map\ fst\ {\isacharparenleft}{\kern0pt}{\isacharparenleft}{\kern0pt}t\isactrlsub j{\isacharcomma}{\kern0pt}\ A\isactrlsub j{\isacharparenright}{\kern0pt}\ {\isacharhash}{\kern0pt}\ avl{\isacharunderscore}{\kern0pt}inorder\ r{\isacharparenright}{\kern0pt}{\isacharparenright}{\kern0pt}{\isachardoublequoteclose}\isanewline
\ \ \ \ \ \ \ \ \isacommand{using}\isamarkupfalse%
\ {\isacartoucheopen}strict{\isacharunderscore}{\kern0pt}sorted\ {\isacharparenleft}{\kern0pt}map\ fst\ {\isacharparenleft}{\kern0pt}avl{\isacharunderscore}{\kern0pt}inorder\ l\ {\isacharat}{\kern0pt}\ {\isacharparenleft}{\kern0pt}t\isactrlsub j{\isacharcomma}{\kern0pt}\ A\isactrlsub j{\isacharparenright}{\kern0pt}\ {\isacharhash}{\kern0pt}\ avl{\isacharunderscore}{\kern0pt}inorder\ r{\isacharparenright}{\kern0pt}{\isacharparenright}{\kern0pt}{\isacartoucheclose}\ \isacommand{by}\isamarkupfalse%
\ auto\isanewline
\ \ \ \ \ \ \isacommand{then}\isamarkupfalse%
\ \isacommand{have}\isamarkupfalse%
\ {\isachardoublequoteopen}{\isasymforall}x\ {\isasymin}\ set\ {\isacharparenleft}{\kern0pt}map\ fst\ {\isacharparenleft}{\kern0pt}avl{\isacharunderscore}{\kern0pt}inorder\ l{\isacharparenright}{\kern0pt}{\isacharparenright}{\kern0pt}{\isachardot}{\kern0pt}\ {\isasymforall}y\ {\isasymin}\ set\ {\isacharparenleft}{\kern0pt}map\ fst\ {\isacharparenleft}{\kern0pt}{\isacharparenleft}{\kern0pt}t\isactrlsub j{\isacharcomma}{\kern0pt}\ A\isactrlsub j{\isacharparenright}{\kern0pt}\ {\isacharhash}{\kern0pt}\ avl{\isacharunderscore}{\kern0pt}inorder\ r{\isacharparenright}{\kern0pt}{\isacharparenright}{\kern0pt}{\isachardot}{\kern0pt}\ x\ {\isacharless}{\kern0pt}\ y{\isachardoublequoteclose}\isanewline
\ \ \ \ \ \ \ \ \isacommand{using}\isamarkupfalse%
\ strict{\isacharunderscore}{\kern0pt}sorted{\isacharunderscore}{\kern0pt}app{\isacharunderscore}{\kern0pt}iff{\isacharbrackleft}{\kern0pt}\isakeyword{where}\ a{\isacharequal}{\kern0pt}{\isachardoublequoteopen}map\ fst\ {\isacharparenleft}{\kern0pt}avl{\isacharunderscore}{\kern0pt}inorder\ l{\isacharparenright}{\kern0pt}{\isachardoublequoteclose}\ \isakeyword{and}\ b{\isacharequal}{\kern0pt}{\isachardoublequoteopen}map\ fst\ {\isacharparenleft}{\kern0pt}{\isacharparenleft}{\kern0pt}t\isactrlsub j{\isacharcomma}{\kern0pt}\ A\isactrlsub j{\isacharparenright}{\kern0pt}\ {\isacharhash}{\kern0pt}\ avl{\isacharunderscore}{\kern0pt}inorder\ r{\isacharparenright}{\kern0pt}{\isachardoublequoteclose}{\isacharbrackright}{\kern0pt}\ \isanewline
\ \ \ \ \ \ \ \ \isacommand{by}\isamarkupfalse%
\ auto\isanewline
\ \ \ \ \ \ \isacommand{then}\isamarkupfalse%
\ \isacommand{have}\isamarkupfalse%
\ {\isachardoublequoteopen}{\isasymforall}x\ {\isasymin}\ set\ {\isacharparenleft}{\kern0pt}map\ fst\ {\isacharparenleft}{\kern0pt}avl{\isacharunderscore}{\kern0pt}inorder\ l{\isacharparenright}{\kern0pt}{\isacharparenright}{\kern0pt}\ {\isasymunion}\ {\isacharbraceleft}{\kern0pt}t\isactrlsub i{\isacharbraceright}{\kern0pt}{\isachardot}{\kern0pt}\ {\isasymforall}y\ {\isasymin}\ set\ {\isacharparenleft}{\kern0pt}map\ fst\ {\isacharparenleft}{\kern0pt}{\isacharparenleft}{\kern0pt}t\isactrlsub j{\isacharcomma}{\kern0pt}\ A\isactrlsub j{\isacharparenright}{\kern0pt}\ {\isacharhash}{\kern0pt}\ avl{\isacharunderscore}{\kern0pt}inorder\ r{\isacharparenright}{\kern0pt}{\isacharparenright}{\kern0pt}{\isachardot}{\kern0pt}\ x\ {\isacharless}{\kern0pt}\ y{\isachardoublequoteclose}\isanewline
\ \ \ \ \ \ \ \ \isacommand{using}\isamarkupfalse%
\ {\isacartoucheopen}t\isactrlsub i\ {\isacharless}{\kern0pt}\ t\isactrlsub j{\isacartoucheclose}\ strict{\isacharunderscore}{\kern0pt}sorted{\isacharunderscore}{\kern0pt}avl{\isacharunderscore}{\kern0pt}inorder{\isacharunderscore}{\kern0pt}r{\isacharbrackleft}{\kern0pt}OF\ {\isacartoucheopen}strict{\isacharunderscore}{\kern0pt}sorted\ {\isacharparenleft}{\kern0pt}map\ fst\ {\isacharparenleft}{\kern0pt}avl{\isacharunderscore}{\kern0pt}inorder\ {\isasymlangle}l{\isacharcomma}{\kern0pt}\ {\isacharparenleft}{\kern0pt}{\isacharparenleft}{\kern0pt}t\isactrlsub j{\isacharcomma}{\kern0pt}\ A\isactrlsub j{\isacharparenright}{\kern0pt}{\isacharcomma}{\kern0pt}\ h{\isacharparenright}{\kern0pt}{\isacharcomma}{\kern0pt}\ r{\isasymrangle}{\isacharparenright}{\kern0pt}{\isacharparenright}{\kern0pt}{\isacartoucheclose}{\isacharbrackright}{\kern0pt}\ \isacommand{by}\isamarkupfalse%
\ auto\isanewline
\ \ \ \ \ \ \isacommand{then}\isamarkupfalse%
\ \isacommand{have}\isamarkupfalse%
\ {\isachardoublequoteopen}{\isasymforall}x\ {\isasymin}\ set\ {\isacharparenleft}{\kern0pt}map\ fst\ {\isacharparenleft}{\kern0pt}avl{\isacharunderscore}{\kern0pt}inorder\ {\isacharparenleft}{\kern0pt}insert{\isacharunderscore}{\kern0pt}happ{\isacharunderscore}{\kern0pt}to{\isacharunderscore}{\kern0pt}tree\ {\isacharparenleft}{\kern0pt}t\isactrlsub i{\isacharcomma}{\kern0pt}A\isactrlsub i{\isacharparenright}{\kern0pt}\ l{\isacharparenright}{\kern0pt}{\isacharparenright}{\kern0pt}{\isacharparenright}{\kern0pt}{\isachardot}{\kern0pt}\ {\isasymforall}y\ {\isasymin}\ set\ {\isacharparenleft}{\kern0pt}map\ fst\ {\isacharparenleft}{\kern0pt}{\isacharparenleft}{\kern0pt}t\isactrlsub j{\isacharcomma}{\kern0pt}\ A\isactrlsub j{\isacharparenright}{\kern0pt}\ {\isacharhash}{\kern0pt}\ avl{\isacharunderscore}{\kern0pt}inorder\ r{\isacharparenright}{\kern0pt}{\isacharparenright}{\kern0pt}{\isachardot}{\kern0pt}\ x\ {\isacharless}{\kern0pt}\ y{\isachardoublequoteclose}\isanewline
\ \ \ \ \ \ \ \ \isacommand{using}\isamarkupfalse%
\ insert{\isacharunderscore}{\kern0pt}happ{\isacharunderscore}{\kern0pt}to{\isacharunderscore}{\kern0pt}tree{\isacharunderscore}{\kern0pt}set{\isacharbrackleft}{\kern0pt}OF\ {\isacartoucheopen}strict{\isacharunderscore}{\kern0pt}sorted\ {\isacharparenleft}{\kern0pt}map\ fst\ {\isacharparenleft}{\kern0pt}avl{\isacharunderscore}{\kern0pt}inorder\ l{\isacharparenright}{\kern0pt}{\isacharparenright}{\kern0pt}{\isacartoucheclose}{\isacharbrackright}{\kern0pt}\ \isacommand{by}\isamarkupfalse%
\ auto\isanewline
\ \ \ \ \ \ \isacommand{then}\isamarkupfalse%
\ \isacommand{have}\isamarkupfalse%
\ {\isachardoublequoteopen}strict{\isacharunderscore}{\kern0pt}sorted\ {\isacharparenleft}{\kern0pt}map\ fst\ {\isacharparenleft}{\kern0pt}avl{\isacharunderscore}{\kern0pt}inorder\ {\isacharparenleft}{\kern0pt}insert{\isacharunderscore}{\kern0pt}happ{\isacharunderscore}{\kern0pt}to{\isacharunderscore}{\kern0pt}tree\ {\isacharparenleft}{\kern0pt}t\isactrlsub i{\isacharcomma}{\kern0pt}A\isactrlsub i{\isacharparenright}{\kern0pt}\ l{\isacharparenright}{\kern0pt}\ {\isacharat}{\kern0pt}\ {\isacharparenleft}{\kern0pt}t\isactrlsub j{\isacharcomma}{\kern0pt}A\isactrlsub j{\isacharparenright}{\kern0pt}\ {\isacharhash}{\kern0pt}\ avl{\isacharunderscore}{\kern0pt}inorder\ r{\isacharparenright}{\kern0pt}{\isacharparenright}{\kern0pt}{\isachardoublequoteclose}\isanewline
\ \ \ \ \ \ \ \ \isacommand{using}\isamarkupfalse%
\ {\isacartoucheopen}strict{\isacharunderscore}{\kern0pt}sorted\ {\isacharparenleft}{\kern0pt}map\ fst\ {\isacharparenleft}{\kern0pt}avl{\isacharunderscore}{\kern0pt}inorder\ {\isacharparenleft}{\kern0pt}insert{\isacharunderscore}{\kern0pt}happ{\isacharunderscore}{\kern0pt}to{\isacharunderscore}{\kern0pt}tree\ {\isacharparenleft}{\kern0pt}t\isactrlsub i{\isacharcomma}{\kern0pt}A\isactrlsub i{\isacharparenright}{\kern0pt}\ l{\isacharparenright}{\kern0pt}{\isacharparenright}{\kern0pt}{\isacharparenright}{\kern0pt}{\isacartoucheclose}\isanewline
\ \ \ \ \ \ \ \ \ \ \ \ \ \ {\isacartoucheopen}strict{\isacharunderscore}{\kern0pt}sorted\ {\isacharparenleft}{\kern0pt}map\ fst\ {\isacharparenleft}{\kern0pt}{\isacharparenleft}{\kern0pt}t\isactrlsub j{\isacharcomma}{\kern0pt}A\isactrlsub j{\isacharparenright}{\kern0pt}\ {\isacharhash}{\kern0pt}\ avl{\isacharunderscore}{\kern0pt}inorder\ r{\isacharparenright}{\kern0pt}{\isacharparenright}{\kern0pt}{\isacartoucheclose}\isanewline
\ \ \ \ \ \ \ \ \ \ \ \ \ \ strict{\isacharunderscore}{\kern0pt}sorted{\isacharunderscore}{\kern0pt}app{\isacharunderscore}{\kern0pt}iff{\isacharbrackleft}{\kern0pt}\isakeyword{where}\ a{\isacharequal}{\kern0pt}{\isachardoublequoteopen}map\ fst\ {\isacharparenleft}{\kern0pt}avl{\isacharunderscore}{\kern0pt}inorder\ {\isacharparenleft}{\kern0pt}insert{\isacharunderscore}{\kern0pt}happ{\isacharunderscore}{\kern0pt}to{\isacharunderscore}{\kern0pt}tree\ {\isacharparenleft}{\kern0pt}t\isactrlsub i{\isacharcomma}{\kern0pt}A\isactrlsub i{\isacharparenright}{\kern0pt}\ l{\isacharparenright}{\kern0pt}{\isacharparenright}{\kern0pt}{\isachardoublequoteclose}\ \isakeyword{and}\ b{\isacharequal}{\kern0pt}{\isachardoublequoteopen}map\ fst\ {\isacharparenleft}{\kern0pt}{\isacharparenleft}{\kern0pt}t\isactrlsub j{\isacharcomma}{\kern0pt}\ A\isactrlsub j{\isacharparenright}{\kern0pt}\ {\isacharhash}{\kern0pt}\ avl{\isacharunderscore}{\kern0pt}inorder\ r{\isacharparenright}{\kern0pt}{\isachardoublequoteclose}{\isacharbrackright}{\kern0pt}\ \isanewline
\ \ \ \ \ \ \ \ \isacommand{by}\isamarkupfalse%
\ auto\isanewline
\ \ \ \ \ \ \isacommand{then}\isamarkupfalse%
\ \isacommand{have}\isamarkupfalse%
\ {\isachardoublequoteopen}strict{\isacharunderscore}{\kern0pt}sorted\ {\isacharparenleft}{\kern0pt}map\ fst\ {\isacharparenleft}{\kern0pt}avl{\isacharunderscore}{\kern0pt}inorder\ {\isacharparenleft}{\kern0pt}avl{\isacharunderscore}{\kern0pt}balL\ {\isacharparenleft}{\kern0pt}insert{\isacharunderscore}{\kern0pt}happ{\isacharunderscore}{\kern0pt}to{\isacharunderscore}{\kern0pt}tree\ {\isacharparenleft}{\kern0pt}t\isactrlsub i{\isacharcomma}{\kern0pt}A\isactrlsub i{\isacharparenright}{\kern0pt}\ l{\isacharparenright}{\kern0pt}\ {\isacharparenleft}{\kern0pt}t\isactrlsub j{\isacharcomma}{\kern0pt}A\isactrlsub j{\isacharparenright}{\kern0pt}\ r{\isacharparenright}{\kern0pt}{\isacharparenright}{\kern0pt}{\isacharparenright}{\kern0pt}{\isachardoublequoteclose}\isanewline
\ \ \ \ \ \ \ \ \isacommand{using}\isamarkupfalse%
\ inorder{\isacharunderscore}{\kern0pt}avl{\isacharunderscore}{\kern0pt}balL\ \isacommand{by}\isamarkupfalse%
\ auto\isanewline
\ \ \ \ \ \ \isacommand{then}\isamarkupfalse%
\ \isacommand{have}\isamarkupfalse%
\ {\isachardoublequoteopen}strict{\isacharunderscore}{\kern0pt}sorted\ {\isacharparenleft}{\kern0pt}map\ fst\ {\isacharparenleft}{\kern0pt}avl{\isacharunderscore}{\kern0pt}inorder\ {\isacharparenleft}{\kern0pt}insert{\isacharunderscore}{\kern0pt}happ{\isacharunderscore}{\kern0pt}to{\isacharunderscore}{\kern0pt}tree\ {\isacharparenleft}{\kern0pt}t\isactrlsub i{\isacharcomma}{\kern0pt}\ A\isactrlsub i{\isacharparenright}{\kern0pt}\ {\isasymlangle}l{\isacharcomma}{\kern0pt}\ {\isacharparenleft}{\kern0pt}{\isacharparenleft}{\kern0pt}t\isactrlsub j{\isacharcomma}{\kern0pt}\ A\isactrlsub j{\isacharparenright}{\kern0pt}{\isacharcomma}{\kern0pt}\ h{\isacharparenright}{\kern0pt}{\isacharcomma}{\kern0pt}\ r{\isasymrangle}{\isacharparenright}{\kern0pt}{\isacharparenright}{\kern0pt}{\isacharparenright}{\kern0pt}{\isachardoublequoteclose}\isanewline
\ \ \ \ \ \ \ \ \isacommand{using}\isamarkupfalse%
\ {\isacartoucheopen}t\isactrlsub i\ {\isacharless}{\kern0pt}\ t\isactrlsub j{\isacartoucheclose}\ \isacommand{by}\isamarkupfalse%
\ auto\isanewline
\ \ \ \ \ \ \isacommand{then}\isamarkupfalse%
\ \isacommand{show}\isamarkupfalse%
\ {\isacharquery}{\kern0pt}case\isanewline
\ \ \ \ \ \ \ \ \isacommand{by}\isamarkupfalse%
\ auto\isanewline
\ \ \ \ \isacommand{next}\isamarkupfalse%
\isanewline
\ \ \ \ \ \ \isacommand{assume}\isamarkupfalse%
\ {\isachardoublequoteopen}t\isactrlsub i\ {\isacharequal}{\kern0pt}\ t\isactrlsub j{\isachardoublequoteclose}\isanewline
\ \ \ \ \ \ \isacommand{then}\isamarkupfalse%
\ \isacommand{show}\isamarkupfalse%
\ {\isacharquery}{\kern0pt}case\isanewline
\ \ \ \ \ \ \ \ \isacommand{using}\isamarkupfalse%
\ {\isachardoublequoteopen}{\isadigit{2}}{\isachardot}{\kern0pt}prems{\isachardoublequoteclose}\ \isacommand{by}\isamarkupfalse%
\ {\isacharparenleft}{\kern0pt}auto\ simp{\isacharcolon}{\kern0pt}\ avl{\isacharunderscore}{\kern0pt}inorder{\isacharunderscore}{\kern0pt}def\ avl{\isacharunderscore}{\kern0pt}node{\isacharunderscore}{\kern0pt}def{\isacharparenright}{\kern0pt}\isanewline
\ \ \ \ \isacommand{next}\isamarkupfalse%
\isanewline
\ \ \ \ \ \ \isacommand{assume}\isamarkupfalse%
\ {\isachardoublequoteopen}t\isactrlsub j\ {\isacharless}{\kern0pt}\ t\isactrlsub i{\isachardoublequoteclose}\isanewline
\ \ \ \ \ \ \isacommand{then}\isamarkupfalse%
\ \isacommand{have}\isamarkupfalse%
\ {\isachardoublequoteopen}strict{\isacharunderscore}{\kern0pt}sorted\ {\isacharparenleft}{\kern0pt}map\ fst\ {\isacharparenleft}{\kern0pt}avl{\isacharunderscore}{\kern0pt}inorder\ {\isacharparenleft}{\kern0pt}insert{\isacharunderscore}{\kern0pt}happ{\isacharunderscore}{\kern0pt}to{\isacharunderscore}{\kern0pt}tree\ {\isacharparenleft}{\kern0pt}t\isactrlsub i{\isacharcomma}{\kern0pt}A\isactrlsub i{\isacharparenright}{\kern0pt}\ r{\isacharparenright}{\kern0pt}{\isacharparenright}{\kern0pt}{\isacharparenright}{\kern0pt}{\isachardoublequoteclose}\isanewline
\ \ \ \ \ \ \ \ \isacommand{using}\isamarkupfalse%
\ {\isacartoucheopen}strict{\isacharunderscore}{\kern0pt}sorted\ {\isacharparenleft}{\kern0pt}map\ fst\ {\isacharparenleft}{\kern0pt}avl{\isacharunderscore}{\kern0pt}inorder\ r{\isacharparenright}{\kern0pt}{\isacharparenright}{\kern0pt}{\isacartoucheclose}\ {\isachardoublequoteopen}{\isadigit{2}}{\isachardot}{\kern0pt}IH{\isachardoublequoteclose}\ \isacommand{by}\isamarkupfalse%
\ auto\isanewline
\isanewline
\ \ \ \ \ \ \isacommand{have}\isamarkupfalse%
\ {\isachardoublequoteopen}strict{\isacharunderscore}{\kern0pt}sorted\ {\isacharparenleft}{\kern0pt}map\ fst\ {\isacharparenleft}{\kern0pt}avl{\isacharunderscore}{\kern0pt}inorder\ l\ {\isacharat}{\kern0pt}\ {\isacharbrackleft}{\kern0pt}{\isacharparenleft}{\kern0pt}t\isactrlsub j{\isacharcomma}{\kern0pt}\ A\isactrlsub j{\isacharparenright}{\kern0pt}{\isacharbrackright}{\kern0pt}{\isacharparenright}{\kern0pt}\ {\isacharat}{\kern0pt}\ map\ fst\ {\isacharparenleft}{\kern0pt}avl{\isacharunderscore}{\kern0pt}inorder\ r{\isacharparenright}{\kern0pt}{\isacharparenright}{\kern0pt}{\isachardoublequoteclose}\isanewline
\ \ \ \ \ \ \ \ \isacommand{using}\isamarkupfalse%
\ {\isacartoucheopen}strict{\isacharunderscore}{\kern0pt}sorted\ {\isacharparenleft}{\kern0pt}map\ fst\ {\isacharparenleft}{\kern0pt}avl{\isacharunderscore}{\kern0pt}inorder\ l\ {\isacharat}{\kern0pt}\ {\isacharparenleft}{\kern0pt}t\isactrlsub j{\isacharcomma}{\kern0pt}\ A\isactrlsub j{\isacharparenright}{\kern0pt}\ {\isacharhash}{\kern0pt}\ avl{\isacharunderscore}{\kern0pt}inorder\ r{\isacharparenright}{\kern0pt}{\isacharparenright}{\kern0pt}{\isacartoucheclose}\ \isacommand{by}\isamarkupfalse%
\ auto\isanewline
\ \ \ \ \ \ \isacommand{then}\isamarkupfalse%
\ \isacommand{have}\isamarkupfalse%
\ {\isachardoublequoteopen}{\isasymforall}x\ {\isasymin}\ set\ {\isacharparenleft}{\kern0pt}map\ fst\ {\isacharparenleft}{\kern0pt}avl{\isacharunderscore}{\kern0pt}inorder\ l\ {\isacharat}{\kern0pt}\ {\isacharbrackleft}{\kern0pt}{\isacharparenleft}{\kern0pt}t\isactrlsub j{\isacharcomma}{\kern0pt}\ A\isactrlsub j{\isacharparenright}{\kern0pt}{\isacharbrackright}{\kern0pt}{\isacharparenright}{\kern0pt}{\isacharparenright}{\kern0pt}{\isachardot}{\kern0pt}\ {\isasymforall}y\ {\isasymin}\ set\ {\isacharparenleft}{\kern0pt}map\ fst\ {\isacharparenleft}{\kern0pt}avl{\isacharunderscore}{\kern0pt}inorder\ r{\isacharparenright}{\kern0pt}{\isacharparenright}{\kern0pt}{\isachardot}{\kern0pt}\ x\ {\isacharless}{\kern0pt}\ y{\isachardoublequoteclose}\isanewline
\ \ \ \ \ \ \ \ \isacommand{using}\isamarkupfalse%
\ strict{\isacharunderscore}{\kern0pt}sorted{\isacharunderscore}{\kern0pt}app{\isacharunderscore}{\kern0pt}iff{\isacharbrackleft}{\kern0pt}\isakeyword{where}\ a{\isacharequal}{\kern0pt}{\isachardoublequoteopen}map\ fst\ {\isacharparenleft}{\kern0pt}avl{\isacharunderscore}{\kern0pt}inorder\ l\ {\isacharat}{\kern0pt}\ {\isacharbrackleft}{\kern0pt}{\isacharparenleft}{\kern0pt}t\isactrlsub j{\isacharcomma}{\kern0pt}\ A\isactrlsub j{\isacharparenright}{\kern0pt}{\isacharbrackright}{\kern0pt}{\isacharparenright}{\kern0pt}{\isachardoublequoteclose}\ \isakeyword{and}\ b{\isacharequal}{\kern0pt}{\isachardoublequoteopen}map\ fst\ {\isacharparenleft}{\kern0pt}avl{\isacharunderscore}{\kern0pt}inorder\ r{\isacharparenright}{\kern0pt}{\isachardoublequoteclose}{\isacharbrackright}{\kern0pt}\ \isanewline
\ \ \ \ \ \ \ \ \isacommand{by}\isamarkupfalse%
\ auto\isanewline
\ \ \ \ \ \ \isacommand{then}\isamarkupfalse%
\ \isacommand{have}\isamarkupfalse%
\ {\isachardoublequoteopen}{\isasymforall}x\ {\isasymin}\ set\ {\isacharparenleft}{\kern0pt}map\ fst\ {\isacharparenleft}{\kern0pt}avl{\isacharunderscore}{\kern0pt}inorder\ l\ {\isacharat}{\kern0pt}\ {\isacharbrackleft}{\kern0pt}{\isacharparenleft}{\kern0pt}t\isactrlsub j{\isacharcomma}{\kern0pt}\ A\isactrlsub j{\isacharparenright}{\kern0pt}{\isacharbrackright}{\kern0pt}{\isacharparenright}{\kern0pt}{\isacharparenright}{\kern0pt}{\isachardot}{\kern0pt}\ {\isasymforall}y\ {\isasymin}\ set\ {\isacharparenleft}{\kern0pt}map\ fst\ {\isacharparenleft}{\kern0pt}avl{\isacharunderscore}{\kern0pt}inorder\ r{\isacharparenright}{\kern0pt}{\isacharparenright}{\kern0pt}\ {\isasymunion}\ {\isacharbraceleft}{\kern0pt}t\isactrlsub i{\isacharbraceright}{\kern0pt}{\isachardot}{\kern0pt}\ x\ {\isacharless}{\kern0pt}\ y{\isachardoublequoteclose}\isanewline
\ \ \ \ \ \ \ \ \isacommand{using}\isamarkupfalse%
\ {\isacartoucheopen}t\isactrlsub j\ {\isacharless}{\kern0pt}\ t\isactrlsub i{\isacartoucheclose}\ strict{\isacharunderscore}{\kern0pt}sorted{\isacharunderscore}{\kern0pt}avl{\isacharunderscore}{\kern0pt}inorder{\isacharunderscore}{\kern0pt}l{\isacharbrackleft}{\kern0pt}OF\ {\isacartoucheopen}strict{\isacharunderscore}{\kern0pt}sorted\ {\isacharparenleft}{\kern0pt}map\ fst\ {\isacharparenleft}{\kern0pt}avl{\isacharunderscore}{\kern0pt}inorder\ {\isasymlangle}l{\isacharcomma}{\kern0pt}\ {\isacharparenleft}{\kern0pt}{\isacharparenleft}{\kern0pt}t\isactrlsub j{\isacharcomma}{\kern0pt}\ A\isactrlsub j{\isacharparenright}{\kern0pt}{\isacharcomma}{\kern0pt}\ h{\isacharparenright}{\kern0pt}{\isacharcomma}{\kern0pt}\ r{\isasymrangle}{\isacharparenright}{\kern0pt}{\isacharparenright}{\kern0pt}{\isacartoucheclose}{\isacharbrackright}{\kern0pt}\ \isacommand{by}\isamarkupfalse%
\ auto\isanewline
\ \ \ \ \ \ \isacommand{then}\isamarkupfalse%
\ \isacommand{have}\isamarkupfalse%
\ {\isachardoublequoteopen}{\isasymforall}x\ {\isasymin}\ set\ {\isacharparenleft}{\kern0pt}map\ fst\ {\isacharparenleft}{\kern0pt}avl{\isacharunderscore}{\kern0pt}inorder\ l\ {\isacharat}{\kern0pt}\ {\isacharbrackleft}{\kern0pt}{\isacharparenleft}{\kern0pt}t\isactrlsub j{\isacharcomma}{\kern0pt}\ A\isactrlsub j{\isacharparenright}{\kern0pt}{\isacharbrackright}{\kern0pt}{\isacharparenright}{\kern0pt}{\isacharparenright}{\kern0pt}{\isachardot}{\kern0pt}\ {\isasymforall}y\ {\isasymin}\ set\ {\isacharparenleft}{\kern0pt}map\ fst\ {\isacharparenleft}{\kern0pt}avl{\isacharunderscore}{\kern0pt}inorder\ {\isacharparenleft}{\kern0pt}insert{\isacharunderscore}{\kern0pt}happ{\isacharunderscore}{\kern0pt}to{\isacharunderscore}{\kern0pt}tree\ {\isacharparenleft}{\kern0pt}t\isactrlsub i{\isacharcomma}{\kern0pt}A\isactrlsub i{\isacharparenright}{\kern0pt}\ r{\isacharparenright}{\kern0pt}{\isacharparenright}{\kern0pt}{\isacharparenright}{\kern0pt}{\isachardot}{\kern0pt}\ x\ {\isacharless}{\kern0pt}\ y{\isachardoublequoteclose}\isanewline
\ \ \ \ \ \ \ \ \isacommand{using}\isamarkupfalse%
\ insert{\isacharunderscore}{\kern0pt}happ{\isacharunderscore}{\kern0pt}to{\isacharunderscore}{\kern0pt}tree{\isacharunderscore}{\kern0pt}set{\isacharbrackleft}{\kern0pt}OF\ {\isacartoucheopen}strict{\isacharunderscore}{\kern0pt}sorted\ {\isacharparenleft}{\kern0pt}map\ fst\ {\isacharparenleft}{\kern0pt}avl{\isacharunderscore}{\kern0pt}inorder\ r{\isacharparenright}{\kern0pt}{\isacharparenright}{\kern0pt}{\isacartoucheclose}{\isacharbrackright}{\kern0pt}\ \isacommand{by}\isamarkupfalse%
\ auto\isanewline
\ \ \ \ \ \ \isacommand{then}\isamarkupfalse%
\ \isacommand{have}\isamarkupfalse%
\ {\isachardoublequoteopen}strict{\isacharunderscore}{\kern0pt}sorted\ {\isacharparenleft}{\kern0pt}map\ fst\ {\isacharparenleft}{\kern0pt}avl{\isacharunderscore}{\kern0pt}inorder\ l\ {\isacharat}{\kern0pt}\ {\isacharparenleft}{\kern0pt}t\isactrlsub j{\isacharcomma}{\kern0pt}A\isactrlsub j{\isacharparenright}{\kern0pt}\ {\isacharhash}{\kern0pt}\ avl{\isacharunderscore}{\kern0pt}inorder\ {\isacharparenleft}{\kern0pt}insert{\isacharunderscore}{\kern0pt}happ{\isacharunderscore}{\kern0pt}to{\isacharunderscore}{\kern0pt}tree\ {\isacharparenleft}{\kern0pt}t\isactrlsub i{\isacharcomma}{\kern0pt}A\isactrlsub i{\isacharparenright}{\kern0pt}\ r{\isacharparenright}{\kern0pt}{\isacharparenright}{\kern0pt}{\isacharparenright}{\kern0pt}{\isachardoublequoteclose}\isanewline
\ \ \ \ \ \ \ \ \isacommand{using}\isamarkupfalse%
\ {\isacartoucheopen}strict{\isacharunderscore}{\kern0pt}sorted\ {\isacharparenleft}{\kern0pt}map\ fst\ {\isacharparenleft}{\kern0pt}avl{\isacharunderscore}{\kern0pt}inorder\ l\ {\isacharat}{\kern0pt}\ {\isacharbrackleft}{\kern0pt}{\isacharparenleft}{\kern0pt}t\isactrlsub j{\isacharcomma}{\kern0pt}A\isactrlsub j{\isacharparenright}{\kern0pt}{\isacharbrackright}{\kern0pt}{\isacharparenright}{\kern0pt}{\isacharparenright}{\kern0pt}{\isacartoucheclose}\isanewline
\ \ \ \ \ \ \ \ \ \ \ \ \ \ {\isacartoucheopen}strict{\isacharunderscore}{\kern0pt}sorted\ {\isacharparenleft}{\kern0pt}map\ fst\ {\isacharparenleft}{\kern0pt}avl{\isacharunderscore}{\kern0pt}inorder\ {\isacharparenleft}{\kern0pt}insert{\isacharunderscore}{\kern0pt}happ{\isacharunderscore}{\kern0pt}to{\isacharunderscore}{\kern0pt}tree\ {\isacharparenleft}{\kern0pt}t\isactrlsub i{\isacharcomma}{\kern0pt}A\isactrlsub i{\isacharparenright}{\kern0pt}\ r{\isacharparenright}{\kern0pt}{\isacharparenright}{\kern0pt}{\isacharparenright}{\kern0pt}{\isacartoucheclose}\isanewline
\ \ \ \ \ \ \ \ \ \ \ \ \ \ strict{\isacharunderscore}{\kern0pt}sorted{\isacharunderscore}{\kern0pt}app{\isacharunderscore}{\kern0pt}iff{\isacharbrackleft}{\kern0pt}\isakeyword{where}\ a{\isacharequal}{\kern0pt}{\isachardoublequoteopen}map\ fst\ {\isacharparenleft}{\kern0pt}avl{\isacharunderscore}{\kern0pt}inorder\ l\ {\isacharat}{\kern0pt}\ {\isacharbrackleft}{\kern0pt}{\isacharparenleft}{\kern0pt}t\isactrlsub j{\isacharcomma}{\kern0pt}A\isactrlsub j{\isacharparenright}{\kern0pt}{\isacharbrackright}{\kern0pt}{\isacharparenright}{\kern0pt}{\isachardoublequoteclose}\ \isakeyword{and}\ b{\isacharequal}{\kern0pt}{\isachardoublequoteopen}map\ fst\ {\isacharparenleft}{\kern0pt}avl{\isacharunderscore}{\kern0pt}inorder\ {\isacharparenleft}{\kern0pt}insert{\isacharunderscore}{\kern0pt}happ{\isacharunderscore}{\kern0pt}to{\isacharunderscore}{\kern0pt}tree\ {\isacharparenleft}{\kern0pt}t\isactrlsub i{\isacharcomma}{\kern0pt}A\isactrlsub i{\isacharparenright}{\kern0pt}\ r{\isacharparenright}{\kern0pt}{\isacharparenright}{\kern0pt}{\isachardoublequoteclose}{\isacharbrackright}{\kern0pt}\ \isanewline
\ \ \ \ \ \ \ \ \isacommand{by}\isamarkupfalse%
\ auto\isanewline
\ \ \ \ \ \ \isacommand{then}\isamarkupfalse%
\ \isacommand{have}\isamarkupfalse%
\ {\isachardoublequoteopen}strict{\isacharunderscore}{\kern0pt}sorted\ {\isacharparenleft}{\kern0pt}map\ fst\ {\isacharparenleft}{\kern0pt}avl{\isacharunderscore}{\kern0pt}inorder\ {\isacharparenleft}{\kern0pt}avl{\isacharunderscore}{\kern0pt}balR\ l\ {\isacharparenleft}{\kern0pt}t\isactrlsub j{\isacharcomma}{\kern0pt}A\isactrlsub j{\isacharparenright}{\kern0pt}\ {\isacharparenleft}{\kern0pt}insert{\isacharunderscore}{\kern0pt}happ{\isacharunderscore}{\kern0pt}to{\isacharunderscore}{\kern0pt}tree\ {\isacharparenleft}{\kern0pt}t\isactrlsub i{\isacharcomma}{\kern0pt}A\isactrlsub i{\isacharparenright}{\kern0pt}\ r{\isacharparenright}{\kern0pt}{\isacharparenright}{\kern0pt}{\isacharparenright}{\kern0pt}{\isacharparenright}{\kern0pt}{\isachardoublequoteclose}\isanewline
\ \ \ \ \ \ \ \ \isacommand{using}\isamarkupfalse%
\ inorder{\isacharunderscore}{\kern0pt}avl{\isacharunderscore}{\kern0pt}balR\ \isacommand{by}\isamarkupfalse%
\ auto\isanewline
\ \ \ \ \ \ \isacommand{then}\isamarkupfalse%
\ \isacommand{have}\isamarkupfalse%
\ {\isachardoublequoteopen}strict{\isacharunderscore}{\kern0pt}sorted\ {\isacharparenleft}{\kern0pt}map\ fst\ {\isacharparenleft}{\kern0pt}avl{\isacharunderscore}{\kern0pt}inorder\ {\isacharparenleft}{\kern0pt}insert{\isacharunderscore}{\kern0pt}happ{\isacharunderscore}{\kern0pt}to{\isacharunderscore}{\kern0pt}tree\ {\isacharparenleft}{\kern0pt}t\isactrlsub i{\isacharcomma}{\kern0pt}\ A\isactrlsub i{\isacharparenright}{\kern0pt}\ {\isasymlangle}l{\isacharcomma}{\kern0pt}\ {\isacharparenleft}{\kern0pt}{\isacharparenleft}{\kern0pt}t\isactrlsub j{\isacharcomma}{\kern0pt}\ A\isactrlsub j{\isacharparenright}{\kern0pt}{\isacharcomma}{\kern0pt}\ h{\isacharparenright}{\kern0pt}{\isacharcomma}{\kern0pt}\ r{\isasymrangle}{\isacharparenright}{\kern0pt}{\isacharparenright}{\kern0pt}{\isacharparenright}{\kern0pt}{\isachardoublequoteclose}\isanewline
\ \ \ \ \ \ \ \ \isacommand{using}\isamarkupfalse%
\ {\isacartoucheopen}t\isactrlsub j\ {\isacharless}{\kern0pt}\ t\isactrlsub i{\isacartoucheclose}\ \isacommand{by}\isamarkupfalse%
\ auto\isanewline
\ \ \ \ \ \ \isacommand{then}\isamarkupfalse%
\ \isacommand{show}\isamarkupfalse%
\ {\isacharquery}{\kern0pt}case\isanewline
\ \ \ \ \ \ \ \ \isacommand{by}\isamarkupfalse%
\ auto\isanewline
\ \ \ \ \isacommand{qed}\isamarkupfalse%
\isanewline
\ \ \isacommand{qed}\isamarkupfalse%
%
\endisatagproof
{\isafoldproof}%
%
\isadelimproof
\isanewline
%
\endisadelimproof
\isanewline
\ \ \isacommand{lemma}\isamarkupfalse%
\ insert{\isacharunderscore}{\kern0pt}timed{\isacharunderscore}{\kern0pt}ground{\isacharunderscore}{\kern0pt}acts{\isacharunderscore}{\kern0pt}to{\isacharunderscore}{\kern0pt}tree{\isacharunderscore}{\kern0pt}strict{\isacharunderscore}{\kern0pt}sorted{\isacharcolon}{\kern0pt}\ \isanewline
\ \ \ \ \isakeyword{assumes}\ {\isachardoublequoteopen}strict{\isacharunderscore}{\kern0pt}sorted\ {\isacharparenleft}{\kern0pt}map\ fst\ {\isacharparenleft}{\kern0pt}avl{\isacharunderscore}{\kern0pt}inorder\ htree{\isacharparenright}{\kern0pt}{\isacharparenright}{\kern0pt}{\isachardoublequoteclose}\isanewline
\ \ \ \ \isakeyword{shows}\ {\isachardoublequoteopen}strict{\isacharunderscore}{\kern0pt}sorted\ {\isacharparenleft}{\kern0pt}map\ fst\ {\isacharparenleft}{\kern0pt}avl{\isacharunderscore}{\kern0pt}inorder\ {\isacharparenleft}{\kern0pt}insert{\isacharunderscore}{\kern0pt}timed{\isacharunderscore}{\kern0pt}ground{\isacharunderscore}{\kern0pt}acts{\isacharunderscore}{\kern0pt}to{\isacharunderscore}{\kern0pt}tree\ as\ htree{\isacharparenright}{\kern0pt}{\isacharparenright}{\kern0pt}{\isacharparenright}{\kern0pt}{\isachardoublequoteclose}\isanewline
%
\isadelimproof
\ \ \ \ %
\endisadelimproof
%
\isatagproof
\isacommand{using}\isamarkupfalse%
\ assms\ insert{\isacharunderscore}{\kern0pt}timed{\isacharunderscore}{\kern0pt}ground{\isacharunderscore}{\kern0pt}act{\isacharunderscore}{\kern0pt}to{\isacharunderscore}{\kern0pt}tree{\isacharunderscore}{\kern0pt}strict{\isacharunderscore}{\kern0pt}sorted\isanewline
\ \ \ \ \isacommand{by}\isamarkupfalse%
\ {\isacharparenleft}{\kern0pt}induction\ as{\isacharparenright}{\kern0pt}\ auto%
\endisatagproof
{\isafoldproof}%
%
\isadelimproof
\isanewline
%
\endisadelimproof
\isanewline
\ \ \isanewline
\ \ \isacommand{lemma}\isamarkupfalse%
\ insert{\isacharunderscore}{\kern0pt}happ{\isacharunderscore}{\kern0pt}to{\isacharunderscore}{\kern0pt}tree{\isacharunderscore}{\kern0pt}equiv{\isacharcolon}{\kern0pt}\ \isanewline
\ \ \ \ \isakeyword{assumes}\ {\isachardoublequoteopen}strict{\isacharunderscore}{\kern0pt}sorted\ {\isacharparenleft}{\kern0pt}map\ fst\ {\isacharparenleft}{\kern0pt}avl{\isacharunderscore}{\kern0pt}inorder\ htree{\isacharparenright}{\kern0pt}{\isacharparenright}{\kern0pt}{\isachardoublequoteclose}\isanewline
\ \ \ \ \isakeyword{shows}\ {\isachardoublequoteopen}avl{\isacharunderscore}{\kern0pt}inorder\ {\isacharparenleft}{\kern0pt}insert{\isacharunderscore}{\kern0pt}happ{\isacharunderscore}{\kern0pt}to{\isacharunderscore}{\kern0pt}tree\ {\isacharparenleft}{\kern0pt}t\isactrlsub i{\isacharcomma}{\kern0pt}A\isactrlsub i{\isacharparenright}{\kern0pt}\ htree{\isacharparenright}{\kern0pt}\ {\isacharequal}{\kern0pt}\ insort{\isacharunderscore}{\kern0pt}happ\ {\isacharparenleft}{\kern0pt}t\isactrlsub i{\isacharcomma}{\kern0pt}A\isactrlsub i{\isacharparenright}{\kern0pt}\ {\isacharparenleft}{\kern0pt}avl{\isacharunderscore}{\kern0pt}inorder\ htree{\isacharparenright}{\kern0pt}{\isachardoublequoteclose}\isanewline
%
\isadelimproof
\ \ \ \ %
\endisadelimproof
%
\isatagproof
\isacommand{using}\isamarkupfalse%
\ assms\isanewline
\ \ \isacommand{proof}\isamarkupfalse%
\ {\isacharparenleft}{\kern0pt}induction\ htree{\isacharparenright}{\kern0pt}\isanewline
\ \ \ \ \isacommand{case}\isamarkupfalse%
\ Leaf\isanewline
\ \ \ \ \isacommand{then}\isamarkupfalse%
\ \isacommand{show}\isamarkupfalse%
\ {\isacharquery}{\kern0pt}case\isanewline
\ \ \ \ \ \ \isacommand{by}\isamarkupfalse%
\ {\isacharparenleft}{\kern0pt}auto\ simp{\isacharcolon}{\kern0pt}\ avl{\isacharunderscore}{\kern0pt}node{\isacharunderscore}{\kern0pt}def\ avl{\isacharunderscore}{\kern0pt}inorder{\isacharunderscore}{\kern0pt}def{\isacharparenright}{\kern0pt}\isanewline
\ \ \isacommand{next}\isamarkupfalse%
\isanewline
\ \ \ \ \isacommand{case}\isamarkupfalse%
\ {\isacharparenleft}{\kern0pt}Node\ l\ a\ r{\isacharparenright}{\kern0pt}\isanewline
\ \ \ \ \isacommand{show}\isamarkupfalse%
\ {\isacharquery}{\kern0pt}case\ \isanewline
\ \ \ \ \isacommand{proof}\isamarkupfalse%
\ {\isacharparenleft}{\kern0pt}cases\ a{\isacharparenright}{\kern0pt}\isanewline
\ \ \ \ \ \ \isacommand{case}\isamarkupfalse%
\ {\isacharparenleft}{\kern0pt}Pair\ a{\isacharprime}{\kern0pt}\ ht{\isacharparenright}{\kern0pt}\isanewline
\ \ \ \ \ \ \ \ \isacommand{then}\isamarkupfalse%
\ \isacommand{have}\isamarkupfalse%
\ {\isachardoublequoteopen}{\isasymlangle}l{\isacharcomma}{\kern0pt}a{\isacharcomma}{\kern0pt}r{\isasymrangle}\ {\isacharequal}{\kern0pt}\ {\isasymlangle}l{\isacharcomma}{\kern0pt}{\isacharparenleft}{\kern0pt}a{\isacharprime}{\kern0pt}{\isacharcomma}{\kern0pt}ht{\isacharparenright}{\kern0pt}{\isacharcomma}{\kern0pt}r{\isasymrangle}{\isachardoublequoteclose}\isanewline
\ \ \ \ \ \ \ \ \ \ \isacommand{by}\isamarkupfalse%
\ auto\isanewline
\ \ \ \ \ \ \ \ \isacommand{show}\isamarkupfalse%
\ {\isacharquery}{\kern0pt}thesis\isanewline
\ \ \ \ \ \ \ \ \ \ \isacommand{proof}\isamarkupfalse%
\ {\isacharparenleft}{\kern0pt}cases\ a{\isacharprime}{\kern0pt}{\isacharparenright}{\kern0pt}\isanewline
\ \ \ \ \ \ \ \ \ \ \ \ \isacommand{case}\isamarkupfalse%
\ {\isacharparenleft}{\kern0pt}Pair\ t\isactrlsub j\ A\isactrlsub j{\isacharparenright}{\kern0pt}\isanewline
\ \ \ \ \ \ \ \ \ \ \ \ \isacommand{then}\isamarkupfalse%
\ \isacommand{have}\isamarkupfalse%
\ {\isachardoublequoteopen}{\isasymlangle}l{\isacharcomma}{\kern0pt}a{\isacharcomma}{\kern0pt}r{\isasymrangle}\ {\isacharequal}{\kern0pt}\ {\isasymlangle}l{\isacharcomma}{\kern0pt}{\isacharparenleft}{\kern0pt}{\isacharparenleft}{\kern0pt}t\isactrlsub j{\isacharcomma}{\kern0pt}A\isactrlsub j{\isacharparenright}{\kern0pt}{\isacharcomma}{\kern0pt}ht{\isacharparenright}{\kern0pt}{\isacharcomma}{\kern0pt}r{\isasymrangle}{\isachardoublequoteclose}\isanewline
\ \ \ \ \ \ \ \ \ \ \ \ \ \ \isacommand{using}\isamarkupfalse%
\ {\isacartoucheopen}{\isasymlangle}l{\isacharcomma}{\kern0pt}a{\isacharcomma}{\kern0pt}r{\isasymrangle}\ {\isacharequal}{\kern0pt}\ {\isasymlangle}l{\isacharcomma}{\kern0pt}{\isacharparenleft}{\kern0pt}a{\isacharprime}{\kern0pt}{\isacharcomma}{\kern0pt}ht{\isacharparenright}{\kern0pt}{\isacharcomma}{\kern0pt}r{\isasymrangle}{\isacartoucheclose}\ \isacommand{by}\isamarkupfalse%
\ auto\isanewline
\ \ \ \ \ \ \ \ \ \ \ \ \isacommand{then}\isamarkupfalse%
\ \isacommand{have}\isamarkupfalse%
\ {\isachardoublequoteopen}avl{\isacharunderscore}{\kern0pt}inorder\ {\isasymlangle}l{\isacharcomma}{\kern0pt}a{\isacharcomma}{\kern0pt}r{\isasymrangle}\ {\isacharequal}{\kern0pt}\ avl{\isacharunderscore}{\kern0pt}inorder\ l\ {\isacharat}{\kern0pt}\ {\isacharparenleft}{\kern0pt}t\isactrlsub j{\isacharcomma}{\kern0pt}A\isactrlsub j{\isacharparenright}{\kern0pt}\ {\isacharhash}{\kern0pt}\ avl{\isacharunderscore}{\kern0pt}inorder\ r{\isachardoublequoteclose}\isanewline
\ \ \ \ \ \ \ \ \ \ \ \ \ \ \isacommand{unfolding}\isamarkupfalse%
\ avl{\isacharunderscore}{\kern0pt}inorder{\isacharunderscore}{\kern0pt}def\ \isacommand{by}\isamarkupfalse%
\ auto\isanewline
\ \ \ \ \ \ \ \ \ \ \ \ \isacommand{then}\isamarkupfalse%
\ \isacommand{have}\isamarkupfalse%
\ {\isachardoublequoteopen}strict{\isacharunderscore}{\kern0pt}sorted\ {\isacharparenleft}{\kern0pt}map\ fst\ {\isacharparenleft}{\kern0pt}avl{\isacharunderscore}{\kern0pt}inorder\ l\ {\isacharat}{\kern0pt}\ {\isacharparenleft}{\kern0pt}t\isactrlsub j{\isacharcomma}{\kern0pt}A\isactrlsub j{\isacharparenright}{\kern0pt}\ {\isacharhash}{\kern0pt}\ avl{\isacharunderscore}{\kern0pt}inorder\ r{\isacharparenright}{\kern0pt}{\isacharparenright}{\kern0pt}{\isachardoublequoteclose}\isanewline
\ \ \ \ \ \ \ \ \ \ \ \ \ \ \isacommand{using}\isamarkupfalse%
\ Node{\isachardot}{\kern0pt}prems\ \isacommand{by}\isamarkupfalse%
\ auto\isanewline
\ \ \ \ \ \ \ \ \ \ \ \ \isacommand{then}\isamarkupfalse%
\ \isacommand{have}\isamarkupfalse%
\ {\isachardoublequoteopen}strict{\isacharunderscore}{\kern0pt}sorted\ {\isacharparenleft}{\kern0pt}map\ fst\ {\isacharparenleft}{\kern0pt}avl{\isacharunderscore}{\kern0pt}inorder\ l\ {\isacharat}{\kern0pt}\ {\isacharbrackleft}{\kern0pt}{\isacharparenleft}{\kern0pt}t\isactrlsub j{\isacharcomma}{\kern0pt}A\isactrlsub j{\isacharparenright}{\kern0pt}{\isacharbrackright}{\kern0pt}{\isacharparenright}{\kern0pt}\ {\isacharat}{\kern0pt}\ map\ fst\ {\isacharparenleft}{\kern0pt}avl{\isacharunderscore}{\kern0pt}inorder\ r{\isacharparenright}{\kern0pt}{\isacharparenright}{\kern0pt}{\isachardoublequoteclose}\isanewline
\ \ \ \ \ \ \ \ \ \ \ \ \ \ \ \ \ \ \isakeyword{and}\ {\isachardoublequoteopen}strict{\isacharunderscore}{\kern0pt}sorted\ {\isacharparenleft}{\kern0pt}map\ fst\ {\isacharparenleft}{\kern0pt}avl{\isacharunderscore}{\kern0pt}inorder\ l{\isacharparenright}{\kern0pt}\ {\isacharat}{\kern0pt}\ map\ fst\ {\isacharparenleft}{\kern0pt}{\isacharparenleft}{\kern0pt}t\isactrlsub j{\isacharcomma}{\kern0pt}A\isactrlsub j{\isacharparenright}{\kern0pt}\ {\isacharhash}{\kern0pt}\ avl{\isacharunderscore}{\kern0pt}inorder\ r{\isacharparenright}{\kern0pt}{\isacharparenright}{\kern0pt}{\isachardoublequoteclose}\isanewline
\ \ \ \ \ \ \ \ \ \ \ \ \ \ \isacommand{by}\isamarkupfalse%
\ auto\isanewline
\ \ \ \ \ \ \ \ \ \ \ \ \isacommand{then}\isamarkupfalse%
\ \isacommand{have}\isamarkupfalse%
\ {\isachardoublequoteopen}strict{\isacharunderscore}{\kern0pt}sorted\ {\isacharparenleft}{\kern0pt}map\ fst\ {\isacharparenleft}{\kern0pt}avl{\isacharunderscore}{\kern0pt}inorder\ l\ {\isacharat}{\kern0pt}\ {\isacharbrackleft}{\kern0pt}{\isacharparenleft}{\kern0pt}t\isactrlsub j{\isacharcomma}{\kern0pt}A\isactrlsub j{\isacharparenright}{\kern0pt}{\isacharbrackright}{\kern0pt}{\isacharparenright}{\kern0pt}{\isacharparenright}{\kern0pt}{\isachardoublequoteclose}\ \isanewline
\ \ \ \ \ \ \ \ \ \ \ \ \ \ \ \ \ \ \isakeyword{and}\ {\isachardoublequoteopen}strict{\isacharunderscore}{\kern0pt}sorted\ {\isacharparenleft}{\kern0pt}map\ fst\ {\isacharparenleft}{\kern0pt}avl{\isacharunderscore}{\kern0pt}inorder\ l{\isacharparenright}{\kern0pt}{\isacharparenright}{\kern0pt}{\isachardoublequoteclose}\isanewline
\ \ \ \ \ \ \ \ \ \ \ \ \ \ \ \ \ \ \isakeyword{and}\ {\isachardoublequoteopen}strict{\isacharunderscore}{\kern0pt}sorted\ {\isacharparenleft}{\kern0pt}map\ fst\ {\isacharparenleft}{\kern0pt}{\isacharparenleft}{\kern0pt}t\isactrlsub j{\isacharcomma}{\kern0pt}A\isactrlsub j{\isacharparenright}{\kern0pt}\ {\isacharhash}{\kern0pt}\ avl{\isacharunderscore}{\kern0pt}inorder\ r{\isacharparenright}{\kern0pt}{\isacharparenright}{\kern0pt}{\isachardoublequoteclose}\isanewline
\ \ \ \ \ \ \ \ \ \ \ \ \ \ \ \ \ \ \isakeyword{and}\ {\isachardoublequoteopen}strict{\isacharunderscore}{\kern0pt}sorted\ {\isacharparenleft}{\kern0pt}map\ fst\ {\isacharparenleft}{\kern0pt}avl{\isacharunderscore}{\kern0pt}inorder\ r{\isacharparenright}{\kern0pt}{\isacharparenright}{\kern0pt}{\isachardoublequoteclose}\isanewline
\ \ \ \ \ \ \ \ \ \ \ \ \ \ \isacommand{using}\isamarkupfalse%
\ strict{\isacharunderscore}{\kern0pt}sorted{\isacharunderscore}{\kern0pt}app{\isacharunderscore}{\kern0pt}iff\ \isacommand{by}\isamarkupfalse%
\ blast{\isacharplus}{\kern0pt}\isanewline
\ \ \ \ \ \ \ \ \ \ \ \ \isacommand{then}\isamarkupfalse%
\ \isacommand{have}\isamarkupfalse%
\ {\isachardoublequoteopen}strict{\isacharunderscore}{\kern0pt}sorted\ {\isacharparenleft}{\kern0pt}map\ fst\ {\isacharparenleft}{\kern0pt}avl{\isacharunderscore}{\kern0pt}inorder\ l{\isacharparenright}{\kern0pt}\ {\isacharat}{\kern0pt}\ {\isacharbrackleft}{\kern0pt}t\isactrlsub j{\isacharbrackright}{\kern0pt}{\isacharparenright}{\kern0pt}{\isachardoublequoteclose}\ \isanewline
\ \ \ \ \ \ \ \ \ \ \ \ \ \ \ \ \ \ \isakeyword{and}\ {\isachardoublequoteopen}strict{\isacharunderscore}{\kern0pt}sorted\ {\isacharparenleft}{\kern0pt}t\isactrlsub j\ {\isacharhash}{\kern0pt}\ map\ fst\ {\isacharparenleft}{\kern0pt}avl{\isacharunderscore}{\kern0pt}inorder\ r{\isacharparenright}{\kern0pt}{\isacharparenright}{\kern0pt}{\isachardoublequoteclose}\isanewline
\ \ \ \ \ \ \ \ \ \ \ \ \ \ \isacommand{by}\isamarkupfalse%
\ auto\isanewline
\ \ \ \ \ \ \ \ \ \ \ \ \isacommand{then}\isamarkupfalse%
\ \isacommand{have}\isamarkupfalse%
\ {\isachardoublequoteopen}{\isasymforall}{\isacharparenleft}{\kern0pt}t\isactrlsub k{\isacharcomma}{\kern0pt}A\isactrlsub k{\isacharparenright}{\kern0pt}\ {\isasymin}\ set\ {\isacharparenleft}{\kern0pt}avl{\isacharunderscore}{\kern0pt}inorder\ l{\isacharparenright}{\kern0pt}{\isachardot}{\kern0pt}\ t\isactrlsub k\ {\isacharless}{\kern0pt}\ t\isactrlsub j{\isachardoublequoteclose}\ \isanewline
\ \ \ \ \ \ \ \ \ \ \ \ \ \ \ \ \ \ \isakeyword{and}\ {\isachardoublequoteopen}{\isasymforall}{\isacharparenleft}{\kern0pt}t\isactrlsub k{\isacharcomma}{\kern0pt}A\isactrlsub k{\isacharparenright}{\kern0pt}\ {\isasymin}\ set\ {\isacharparenleft}{\kern0pt}avl{\isacharunderscore}{\kern0pt}inorder\ r{\isacharparenright}{\kern0pt}{\isachardot}{\kern0pt}\ t\isactrlsub j\ {\isacharless}{\kern0pt}\ t\isactrlsub k{\isachardoublequoteclose}\isanewline
\ \ \ \ \ \ \ \ \ \ \ \ \ \ \isacommand{using}\isamarkupfalse%
\ strict{\isacharunderscore}{\kern0pt}sorted{\isacharunderscore}{\kern0pt}appl{\isacharbrackleft}{\kern0pt}OF\ {\isacartoucheopen}strict{\isacharunderscore}{\kern0pt}sorted\ {\isacharparenleft}{\kern0pt}map\ fst\ {\isacharparenleft}{\kern0pt}avl{\isacharunderscore}{\kern0pt}inorder\ l{\isacharparenright}{\kern0pt}\ {\isacharat}{\kern0pt}\ {\isacharbrackleft}{\kern0pt}t\isactrlsub j{\isacharbrackright}{\kern0pt}{\isacharparenright}{\kern0pt}{\isacartoucheclose}{\isacharbrackright}{\kern0pt}\ \isanewline
\ \ \ \ \ \ \ \ \ \ \ \ \ \ \ \ \ \ \ \ strict{\isacharunderscore}{\kern0pt}sorted{\isacharunderscore}{\kern0pt}consg{\isacharbrackleft}{\kern0pt}OF\ {\isacartoucheopen}strict{\isacharunderscore}{\kern0pt}sorted\ {\isacharparenleft}{\kern0pt}t\isactrlsub j\ {\isacharhash}{\kern0pt}\ map\ fst\ {\isacharparenleft}{\kern0pt}avl{\isacharunderscore}{\kern0pt}inorder\ r{\isacharparenright}{\kern0pt}{\isacharparenright}{\kern0pt}{\isacartoucheclose}{\isacharbrackright}{\kern0pt}\isanewline
\ \ \ \ \ \ \ \ \ \ \ \ \ \ \isacommand{by}\isamarkupfalse%
\ auto\isanewline
\isanewline
\ \ \ \ \ \ \ \ \ \ \ \ \isacommand{have}\isamarkupfalse%
\ {\isachardoublequoteopen}t\isactrlsub i\ {\isacharless}{\kern0pt}\ t\isactrlsub j\ {\isasymor}\ t\isactrlsub i\ {\isacharequal}{\kern0pt}\ t\isactrlsub j\ {\isasymor}\ t\isactrlsub j\ {\isacharless}{\kern0pt}\ t\isactrlsub i{\isachardoublequoteclose}\isanewline
\ \ \ \ \ \ \ \ \ \ \ \ \ \ \isacommand{by}\isamarkupfalse%
\ auto\isanewline
\ \ \ \ \ \ \ \ \ \ \ \ \isacommand{then}\isamarkupfalse%
\ \isacommand{show}\isamarkupfalse%
\ {\isacharquery}{\kern0pt}thesis\isanewline
\ \ \ \ \ \ \ \ \ \ \ \ \isacommand{proof}\isamarkupfalse%
\ {\isacharparenleft}{\kern0pt}elim\ disjE{\isacharparenright}{\kern0pt}\isanewline
\ \ \ \ \ \ \ \ \ \ \ \ \ \ \isacommand{assume}\isamarkupfalse%
\ {\isachardoublequoteopen}t\isactrlsub i\ {\isacharless}{\kern0pt}\ t\isactrlsub j{\isachardoublequoteclose}\isanewline
\ \ \ \ \ \ \ \ \ \ \ \ \ \ \isacommand{then}\isamarkupfalse%
\ \isacommand{have}\isamarkupfalse%
\ {\isachardoublequoteopen}{\isasymforall}{\isacharparenleft}{\kern0pt}t\isactrlsub k{\isacharcomma}{\kern0pt}A\isactrlsub k{\isacharparenright}{\kern0pt}\ {\isasymin}\ set\ {\isacharparenleft}{\kern0pt}{\isacharparenleft}{\kern0pt}t\isactrlsub j{\isacharcomma}{\kern0pt}A\isactrlsub j{\isacharparenright}{\kern0pt}\ {\isacharhash}{\kern0pt}\ avl{\isacharunderscore}{\kern0pt}inorder\ r{\isacharparenright}{\kern0pt}{\isachardot}{\kern0pt}\ t\isactrlsub i\ {\isacharless}{\kern0pt}\ t\isactrlsub k{\isachardoublequoteclose}\isanewline
\ \ \ \ \ \ \ \ \ \ \ \ \ \ \ \ \isacommand{using}\isamarkupfalse%
\ {\isacartoucheopen}{\isasymforall}{\isacharparenleft}{\kern0pt}t\isactrlsub k{\isacharcomma}{\kern0pt}A\isactrlsub k{\isacharparenright}{\kern0pt}\ {\isasymin}\ set\ {\isacharparenleft}{\kern0pt}avl{\isacharunderscore}{\kern0pt}inorder\ r{\isacharparenright}{\kern0pt}{\isachardot}{\kern0pt}\ t\isactrlsub j\ {\isacharless}{\kern0pt}\ t\isactrlsub k{\isacartoucheclose}\ \isacommand{by}\isamarkupfalse%
\ auto\isanewline
\ \ \ \ \ \ \ \ \ \ \ \ \ \ \isacommand{then}\isamarkupfalse%
\ \isacommand{show}\isamarkupfalse%
\ {\isacharquery}{\kern0pt}thesis\isanewline
\ \ \ \ \ \ \ \ \ \ \ \ \ \ \ \ \isacommand{using}\isamarkupfalse%
\ {\isacartoucheopen}{\isasymlangle}l{\isacharcomma}{\kern0pt}a{\isacharcomma}{\kern0pt}r{\isasymrangle}\ {\isacharequal}{\kern0pt}\ {\isasymlangle}l{\isacharcomma}{\kern0pt}{\isacharparenleft}{\kern0pt}{\isacharparenleft}{\kern0pt}t\isactrlsub j{\isacharcomma}{\kern0pt}A\isactrlsub j{\isacharparenright}{\kern0pt}{\isacharcomma}{\kern0pt}ht{\isacharparenright}{\kern0pt}{\isacharcomma}{\kern0pt}r{\isasymrangle}{\isacartoucheclose}\ {\isacartoucheopen}t\isactrlsub i\ {\isacharless}{\kern0pt}\ t\isactrlsub j{\isacartoucheclose}\ {\isacartoucheopen}strict{\isacharunderscore}{\kern0pt}sorted\ {\isacharparenleft}{\kern0pt}map\ fst\ {\isacharparenleft}{\kern0pt}avl{\isacharunderscore}{\kern0pt}inorder\ l{\isacharparenright}{\kern0pt}{\isacharparenright}{\kern0pt}{\isacartoucheclose}\ \isanewline
\ \ \ \ \ \ \ \ \ \ \ \ \ \ \ \ \ \ \ \ \ \ Node{\isachardot}{\kern0pt}IH\ inorder{\isacharunderscore}{\kern0pt}avl{\isacharunderscore}{\kern0pt}balL\ insort{\isacharunderscore}{\kern0pt}happ{\isacharunderscore}{\kern0pt}consR\isanewline
\ \ \ \ \ \ \ \ \ \ \ \ \ \ \ \ \isacommand{by}\isamarkupfalse%
\ {\isacharparenleft}{\kern0pt}auto\ simp{\isacharcolon}{\kern0pt}\ avl{\isacharunderscore}{\kern0pt}inorder{\isacharunderscore}{\kern0pt}def{\isacharparenright}{\kern0pt}\ \ \isanewline
\ \ \ \ \ \ \ \ \ \ \ \ \isacommand{next}\isamarkupfalse%
\isanewline
\ \ \ \ \ \ \ \ \ \ \ \ \ \ \isacommand{assume}\isamarkupfalse%
\ {\isachardoublequoteopen}t\isactrlsub i\ {\isacharequal}{\kern0pt}\ t\isactrlsub j{\isachardoublequoteclose}\isanewline
\ \ \ \ \ \ \ \ \ \ \ \ \ \ \isacommand{have}\isamarkupfalse%
\ {\isachardoublequoteopen}avl{\isacharunderscore}{\kern0pt}inorder\ {\isacharparenleft}{\kern0pt}insert{\isacharunderscore}{\kern0pt}happ{\isacharunderscore}{\kern0pt}to{\isacharunderscore}{\kern0pt}tree\ {\isacharparenleft}{\kern0pt}t\isactrlsub i{\isacharcomma}{\kern0pt}A\isactrlsub i{\isacharparenright}{\kern0pt}\ {\isasymlangle}l{\isacharcomma}{\kern0pt}a{\isacharcomma}{\kern0pt}r{\isasymrangle}{\isacharparenright}{\kern0pt}\ {\isacharequal}{\kern0pt}\ avl{\isacharunderscore}{\kern0pt}inorder\ {\isacharparenleft}{\kern0pt}avl{\isacharunderscore}{\kern0pt}node\ l\ {\isacharparenleft}{\kern0pt}t\isactrlsub j{\isacharcomma}{\kern0pt}A\isactrlsub i\ {\isacharat}{\kern0pt}\ A\isactrlsub j{\isacharparenright}{\kern0pt}\ r{\isacharparenright}{\kern0pt}{\isachardoublequoteclose}\isanewline
\ \ \ \ \ \ \ \ \ \ \ \ \ \ \ \ \isacommand{using}\isamarkupfalse%
\ {\isacartoucheopen}{\isasymlangle}l{\isacharcomma}{\kern0pt}a{\isacharcomma}{\kern0pt}r{\isasymrangle}\ {\isacharequal}{\kern0pt}\ {\isasymlangle}l{\isacharcomma}{\kern0pt}{\isacharparenleft}{\kern0pt}{\isacharparenleft}{\kern0pt}t\isactrlsub j{\isacharcomma}{\kern0pt}A\isactrlsub j{\isacharparenright}{\kern0pt}{\isacharcomma}{\kern0pt}ht{\isacharparenright}{\kern0pt}{\isacharcomma}{\kern0pt}r{\isasymrangle}{\isacartoucheclose}\ {\isacartoucheopen}t\isactrlsub i\ {\isacharequal}{\kern0pt}\ t\isactrlsub j{\isacartoucheclose}\ \isacommand{by}\isamarkupfalse%
\ auto\isanewline
\ \ \ \ \ \ \ \ \ \ \ \ \ \ \isacommand{also}\isamarkupfalse%
\ \isacommand{have}\isamarkupfalse%
\ {\isachardoublequoteopen}{\isachardot}{\kern0pt}{\isachardot}{\kern0pt}{\isachardot}{\kern0pt}\ {\isacharequal}{\kern0pt}\ avl{\isacharunderscore}{\kern0pt}inorder\ l\ {\isacharat}{\kern0pt}\ {\isacharparenleft}{\kern0pt}t\isactrlsub j{\isacharcomma}{\kern0pt}A\isactrlsub i\ {\isacharat}{\kern0pt}\ A\isactrlsub j{\isacharparenright}{\kern0pt}\ {\isacharhash}{\kern0pt}\ avl{\isacharunderscore}{\kern0pt}inorder\ r{\isachardoublequoteclose}\isanewline
\ \ \ \ \ \ \ \ \ \ \ \ \ \ \ \ \isacommand{unfolding}\isamarkupfalse%
\ {\isacartoucheopen}{\isasymlangle}l{\isacharcomma}{\kern0pt}a{\isacharcomma}{\kern0pt}r{\isasymrangle}\ {\isacharequal}{\kern0pt}\ {\isasymlangle}l{\isacharcomma}{\kern0pt}{\isacharparenleft}{\kern0pt}{\isacharparenleft}{\kern0pt}t\isactrlsub j{\isacharcomma}{\kern0pt}A\isactrlsub j{\isacharparenright}{\kern0pt}{\isacharcomma}{\kern0pt}ht{\isacharparenright}{\kern0pt}{\isacharcomma}{\kern0pt}r{\isasymrangle}{\isacartoucheclose}\ {\isacartoucheopen}t\isactrlsub i\ {\isacharequal}{\kern0pt}\ t\isactrlsub j{\isacartoucheclose}\ avl{\isacharunderscore}{\kern0pt}node{\isacharunderscore}{\kern0pt}def\ \isacommand{by}\isamarkupfalse%
\ {\isacharparenleft}{\kern0pt}auto\ simp{\isacharcolon}{\kern0pt}\ avl{\isacharunderscore}{\kern0pt}inorder{\isacharunderscore}{\kern0pt}def{\isacharparenright}{\kern0pt}\ \ \ \isanewline
\isanewline
\ \ \ \ \ \ \ \ \ \ \ \ \ \ \isacommand{also}\isamarkupfalse%
\ \isacommand{have}\isamarkupfalse%
\ {\isachardoublequoteopen}{\isachardot}{\kern0pt}{\isachardot}{\kern0pt}{\isachardot}{\kern0pt}\ {\isacharequal}{\kern0pt}\ insort{\isacharunderscore}{\kern0pt}happ\ {\isacharparenleft}{\kern0pt}t\isactrlsub i{\isacharcomma}{\kern0pt}A\isactrlsub i{\isacharparenright}{\kern0pt}\ {\isacharparenleft}{\kern0pt}avl{\isacharunderscore}{\kern0pt}inorder\ l\ {\isacharat}{\kern0pt}\ {\isacharbrackleft}{\kern0pt}{\isacharparenleft}{\kern0pt}t\isactrlsub j{\isacharcomma}{\kern0pt}A\isactrlsub j{\isacharparenright}{\kern0pt}{\isacharbrackright}{\kern0pt}{\isacharparenright}{\kern0pt}\ {\isacharat}{\kern0pt}\ avl{\isacharunderscore}{\kern0pt}inorder\ r{\isachardoublequoteclose}\isanewline
\ \ \ \ \ \ \ \ \ \ \ \ \ \ \ \ \isacommand{using}\isamarkupfalse%
\ {\isacartoucheopen}t\isactrlsub i\ {\isacharequal}{\kern0pt}\ t\isactrlsub j{\isacartoucheclose}\ insort{\isacharunderscore}{\kern0pt}happ{\isacharunderscore}{\kern0pt}consL{\isacharbrackleft}{\kern0pt}OF\ {\isacartoucheopen}{\isasymforall}{\isacharparenleft}{\kern0pt}t\isactrlsub k{\isacharcomma}{\kern0pt}A\isactrlsub k{\isacharparenright}{\kern0pt}\ {\isasymin}\ set\ {\isacharparenleft}{\kern0pt}avl{\isacharunderscore}{\kern0pt}inorder\ l{\isacharparenright}{\kern0pt}{\isachardot}{\kern0pt}\ t\isactrlsub k\ {\isacharless}{\kern0pt}\ t\isactrlsub j{\isacartoucheclose}{\isacharcomma}{\kern0pt}\ \isakeyword{where}\ hs\isactrlsub {\isadigit{2}}{\isacharequal}{\kern0pt}{\isachardoublequoteopen}{\isacharbrackleft}{\kern0pt}{\isacharparenleft}{\kern0pt}t\isactrlsub j{\isacharcomma}{\kern0pt}A\isactrlsub j{\isacharparenright}{\kern0pt}{\isacharbrackright}{\kern0pt}{\isachardoublequoteclose}{\isacharbrackright}{\kern0pt}\ \isanewline
\ \ \ \ \ \ \ \ \ \ \ \ \ \ \ \ \isacommand{by}\isamarkupfalse%
\ auto\isanewline
\ \ \ \ \ \ \ \ \ \ \ \ \ \ \isacommand{also}\isamarkupfalse%
\ \isacommand{have}\isamarkupfalse%
\ {\isachardoublequoteopen}{\isachardot}{\kern0pt}{\isachardot}{\kern0pt}{\isachardot}{\kern0pt}\ {\isacharequal}{\kern0pt}\ insort{\isacharunderscore}{\kern0pt}happ\ {\isacharparenleft}{\kern0pt}t\isactrlsub i{\isacharcomma}{\kern0pt}A\isactrlsub i{\isacharparenright}{\kern0pt}\ {\isacharparenleft}{\kern0pt}avl{\isacharunderscore}{\kern0pt}inorder\ {\isasymlangle}l{\isacharcomma}{\kern0pt}a{\isacharcomma}{\kern0pt}r{\isasymrangle}{\isacharparenright}{\kern0pt}{\isachardoublequoteclose}\isanewline
\ \ \ \ \ \ \ \ \ \ \ \ \ \ \ \ \isacommand{using}\isamarkupfalse%
\ {\isacartoucheopen}t\isactrlsub i\ {\isacharequal}{\kern0pt}\ t\isactrlsub j{\isacartoucheclose}\ {\isacartoucheopen}{\isasymlangle}l{\isacharcomma}{\kern0pt}a{\isacharcomma}{\kern0pt}r{\isasymrangle}\ {\isacharequal}{\kern0pt}\ {\isasymlangle}l{\isacharcomma}{\kern0pt}{\isacharparenleft}{\kern0pt}{\isacharparenleft}{\kern0pt}t\isactrlsub j{\isacharcomma}{\kern0pt}A\isactrlsub j{\isacharparenright}{\kern0pt}{\isacharcomma}{\kern0pt}ht{\isacharparenright}{\kern0pt}{\isacharcomma}{\kern0pt}r{\isasymrangle}{\isacartoucheclose}\ \isanewline
\ \ \ \ \ \ \ \ \ \ \ \ \ \ \ \ \ \ \ \ \ \ insort{\isacharunderscore}{\kern0pt}happ{\isacharunderscore}{\kern0pt}consR{\isacharbrackleft}{\kern0pt}OF\ {\isacartoucheopen}{\isasymforall}{\isacharparenleft}{\kern0pt}t\isactrlsub k{\isacharcomma}{\kern0pt}A\isactrlsub k{\isacharparenright}{\kern0pt}\ {\isasymin}\ set\ {\isacharparenleft}{\kern0pt}avl{\isacharunderscore}{\kern0pt}inorder\ r{\isacharparenright}{\kern0pt}{\isachardot}{\kern0pt}\ t\isactrlsub j\ {\isacharless}{\kern0pt}\ t\isactrlsub k{\isacartoucheclose}{\isacharbrackright}{\kern0pt}\ \isanewline
\ \ \ \ \ \ \ \ \ \ \ \ \ \ \ \ \isacommand{by}\isamarkupfalse%
\ {\isacharparenleft}{\kern0pt}auto\ simp{\isacharcolon}{\kern0pt}\ avl{\isacharunderscore}{\kern0pt}inorder{\isacharunderscore}{\kern0pt}def{\isacharparenright}{\kern0pt}\isanewline
\ \ \ \ \ \ \ \ \ \ \ \ \ \ \isacommand{finally}\isamarkupfalse%
\ \isacommand{show}\isamarkupfalse%
\ {\isacharquery}{\kern0pt}thesis\isanewline
\ \ \ \ \ \ \ \ \ \ \ \ \ \ \ \ \isacommand{by}\isamarkupfalse%
\ auto\ \ \isanewline
\ \ \ \ \ \ \ \ \ \ \ \ \isacommand{next}\isamarkupfalse%
\isanewline
\ \ \ \ \ \ \ \ \ \ \ \ \ \ \isacommand{assume}\isamarkupfalse%
\ {\isachardoublequoteopen}t\isactrlsub j\ {\isacharless}{\kern0pt}\ t\isactrlsub i{\isachardoublequoteclose}\isanewline
\ \ \ \ \ \ \ \ \ \ \ \ \ \ \isacommand{then}\isamarkupfalse%
\ \isacommand{have}\isamarkupfalse%
\ {\isachardoublequoteopen}{\isasymforall}{\isacharparenleft}{\kern0pt}t\isactrlsub k{\isacharcomma}{\kern0pt}A\isactrlsub k{\isacharparenright}{\kern0pt}\ {\isasymin}\ set\ {\isacharparenleft}{\kern0pt}avl{\isacharunderscore}{\kern0pt}inorder\ l\ {\isacharat}{\kern0pt}\ {\isacharbrackleft}{\kern0pt}{\isacharparenleft}{\kern0pt}t\isactrlsub j{\isacharcomma}{\kern0pt}A\isactrlsub j{\isacharparenright}{\kern0pt}{\isacharbrackright}{\kern0pt}{\isacharparenright}{\kern0pt}{\isachardot}{\kern0pt}\ t\isactrlsub k\ {\isacharless}{\kern0pt}\ t\isactrlsub i{\isachardoublequoteclose}\isanewline
\ \ \ \ \ \ \ \ \ \ \ \ \ \ \ \ \isacommand{using}\isamarkupfalse%
\ {\isacartoucheopen}{\isasymforall}{\isacharparenleft}{\kern0pt}t\isactrlsub k{\isacharcomma}{\kern0pt}A\isactrlsub k{\isacharparenright}{\kern0pt}\ {\isasymin}\ set\ {\isacharparenleft}{\kern0pt}avl{\isacharunderscore}{\kern0pt}inorder\ l{\isacharparenright}{\kern0pt}{\isachardot}{\kern0pt}\ t\isactrlsub k\ {\isacharless}{\kern0pt}\ t\isactrlsub j{\isacartoucheclose}\ \isacommand{by}\isamarkupfalse%
\ auto\isanewline
\ \ \ \ \ \ \ \ \ \ \ \ \ \ \isacommand{then}\isamarkupfalse%
\ \isacommand{show}\isamarkupfalse%
\ {\isacharquery}{\kern0pt}thesis\isanewline
\ \ \ \ \ \ \ \ \ \ \ \ \ \ \ \ \isacommand{using}\isamarkupfalse%
\ {\isacartoucheopen}{\isasymlangle}l{\isacharcomma}{\kern0pt}a{\isacharcomma}{\kern0pt}r{\isasymrangle}\ {\isacharequal}{\kern0pt}\ {\isasymlangle}l{\isacharcomma}{\kern0pt}{\isacharparenleft}{\kern0pt}{\isacharparenleft}{\kern0pt}t\isactrlsub j{\isacharcomma}{\kern0pt}A\isactrlsub j{\isacharparenright}{\kern0pt}{\isacharcomma}{\kern0pt}ht{\isacharparenright}{\kern0pt}{\isacharcomma}{\kern0pt}r{\isasymrangle}{\isacartoucheclose}\ {\isacartoucheopen}t\isactrlsub j\ {\isacharless}{\kern0pt}\ t\isactrlsub i{\isacartoucheclose}\ {\isacartoucheopen}strict{\isacharunderscore}{\kern0pt}sorted\ {\isacharparenleft}{\kern0pt}map\ fst\ {\isacharparenleft}{\kern0pt}avl{\isacharunderscore}{\kern0pt}inorder\ r{\isacharparenright}{\kern0pt}{\isacharparenright}{\kern0pt}{\isacartoucheclose}\ \isanewline
\ \ \ \ \ \ \ \ \ \ \ \ \ \ \ \ \ \ \ \ \ \ Node{\isachardot}{\kern0pt}IH\ inorder{\isacharunderscore}{\kern0pt}avl{\isacharunderscore}{\kern0pt}balR\ \isanewline
\ \ \ \ \ \ \ \ \ \ \ \ \ \ \ \ \ \ \ \ \ \ insort{\isacharunderscore}{\kern0pt}happ{\isacharunderscore}{\kern0pt}consL{\isacharbrackleft}{\kern0pt}OF\ {\isacartoucheopen}{\isasymforall}{\isacharparenleft}{\kern0pt}t\isactrlsub k{\isacharcomma}{\kern0pt}A\isactrlsub k{\isacharparenright}{\kern0pt}\ {\isasymin}\ set\ {\isacharparenleft}{\kern0pt}avl{\isacharunderscore}{\kern0pt}inorder\ l\ {\isacharat}{\kern0pt}\ {\isacharbrackleft}{\kern0pt}{\isacharparenleft}{\kern0pt}t\isactrlsub j{\isacharcomma}{\kern0pt}A\isactrlsub j{\isacharparenright}{\kern0pt}{\isacharbrackright}{\kern0pt}{\isacharparenright}{\kern0pt}{\isachardot}{\kern0pt}\ t\isactrlsub k\ {\isacharless}{\kern0pt}\ t\isactrlsub i{\isacartoucheclose}{\isacharbrackright}{\kern0pt}\ \isanewline
\ \ \ \ \ \ \ \ \ \ \ \ \ \ \ \ \isacommand{by}\isamarkupfalse%
\ {\isacharparenleft}{\kern0pt}auto\ simp{\isacharcolon}{\kern0pt}\ avl{\isacharunderscore}{\kern0pt}inorder{\isacharunderscore}{\kern0pt}def{\isacharparenright}{\kern0pt}\isanewline
\ \ \ \ \ \ \ \ \ \ \ \ \isacommand{qed}\isamarkupfalse%
\isanewline
\ \ \ \ \ \ \ \ \ \ \isacommand{qed}\isamarkupfalse%
\isanewline
\ \ \ \ \ \isacommand{qed}\isamarkupfalse%
\isanewline
\ \ \isacommand{qed}\isamarkupfalse%
%
\endisatagproof
{\isafoldproof}%
%
\isadelimproof
\isanewline
%
\endisadelimproof
\isanewline
\ \ \isacommand{lemma}\isamarkupfalse%
\ insert{\isacharunderscore}{\kern0pt}timed{\isacharunderscore}{\kern0pt}ground{\isacharunderscore}{\kern0pt}acts{\isacharunderscore}{\kern0pt}to{\isacharunderscore}{\kern0pt}tree{\isacharunderscore}{\kern0pt}equiv{\isacharcolon}{\kern0pt}\isanewline
\ \ \ \ \isakeyword{assumes}\ {\isachardoublequoteopen}strict{\isacharunderscore}{\kern0pt}sorted\ {\isacharparenleft}{\kern0pt}map\ fst\ {\isacharparenleft}{\kern0pt}avl{\isacharunderscore}{\kern0pt}inorder\ htree{\isacharparenright}{\kern0pt}{\isacharparenright}{\kern0pt}{\isachardoublequoteclose}\isanewline
\ \ \ \ \isakeyword{shows}\ {\isachardoublequoteopen}insort{\isacharunderscore}{\kern0pt}timed{\isacharunderscore}{\kern0pt}ground{\isacharunderscore}{\kern0pt}acts\ as\ {\isacharparenleft}{\kern0pt}avl{\isacharunderscore}{\kern0pt}inorder\ htree{\isacharparenright}{\kern0pt}\ \isanewline
\ \ \ \ \ \ \ \ \ \ \ \ {\isacharequal}{\kern0pt}\ avl{\isacharunderscore}{\kern0pt}inorder\ {\isacharparenleft}{\kern0pt}insert{\isacharunderscore}{\kern0pt}timed{\isacharunderscore}{\kern0pt}ground{\isacharunderscore}{\kern0pt}acts{\isacharunderscore}{\kern0pt}to{\isacharunderscore}{\kern0pt}tree\ as\ htree{\isacharparenright}{\kern0pt}{\isachardoublequoteclose}\isanewline
%
\isadelimproof
\ \ \ \ %
\endisadelimproof
%
\isatagproof
\isacommand{using}\isamarkupfalse%
\ assms\isanewline
\ \ \isacommand{proof}\isamarkupfalse%
\ {\isacharparenleft}{\kern0pt}induction\ as\ arbitrary{\isacharcolon}{\kern0pt}\ htree{\isacharparenright}{\kern0pt}\isanewline
\ \ \ \ \isacommand{case}\isamarkupfalse%
\ Nil\isanewline
\ \ \ \ \isacommand{then}\isamarkupfalse%
\ \isacommand{show}\isamarkupfalse%
\ {\isacharquery}{\kern0pt}case\ \isacommand{by}\isamarkupfalse%
\ auto\isanewline
\ \ \isacommand{next}\isamarkupfalse%
\isanewline
\ \ \ \ \isacommand{case}\isamarkupfalse%
\ {\isacharparenleft}{\kern0pt}Cons\ a{\isacharprime}{\kern0pt}\ as{\isacharparenright}{\kern0pt}\isanewline
\ \ \ \ \isacommand{then}\isamarkupfalse%
\ \isacommand{show}\isamarkupfalse%
\ {\isacharquery}{\kern0pt}case\ \isanewline
\ \ \ \ \ \ \isacommand{using}\isamarkupfalse%
\ insert{\isacharunderscore}{\kern0pt}happ{\isacharunderscore}{\kern0pt}to{\isacharunderscore}{\kern0pt}tree{\isacharunderscore}{\kern0pt}equiv{\isacharbrackleft}{\kern0pt}OF\ insert{\isacharunderscore}{\kern0pt}timed{\isacharunderscore}{\kern0pt}ground{\isacharunderscore}{\kern0pt}acts{\isacharunderscore}{\kern0pt}to{\isacharunderscore}{\kern0pt}tree{\isacharunderscore}{\kern0pt}strict{\isacharunderscore}{\kern0pt}sorted{\isacharbrackleft}{\kern0pt}OF\ Cons{\isachardot}{\kern0pt}prems{\isacharbrackright}{\kern0pt}{\isacharbrackright}{\kern0pt}\isanewline
\ \ \ \ \ \ \isacommand{by}\isamarkupfalse%
\ {\isacharparenleft}{\kern0pt}cases\ a{\isacharprime}{\kern0pt}{\isacharparenright}{\kern0pt}\ auto\isanewline
\ \ \isacommand{qed}\isamarkupfalse%
%
\endisatagproof
{\isafoldproof}%
%
\isadelimproof
\isanewline
%
\endisadelimproof
\isanewline
\ \ \isacommand{lemma}\isamarkupfalse%
\ simplify{\isacharunderscore}{\kern0pt}planE{\isacharunderscore}{\kern0pt}avl{\isacharprime}{\kern0pt}{\isacharunderscore}{\kern0pt}strict{\isacharunderscore}{\kern0pt}sorted{\isacharcolon}{\kern0pt}\isanewline
\ \ \ \ \isakeyword{assumes}\ {\isachardoublequoteopen}simplify{\isacharunderscore}{\kern0pt}planE{\isacharunderscore}{\kern0pt}avl{\isacharprime}{\kern0pt}\ STG\ mp{\isacharunderscore}{\kern0pt}objT\ mp{\isacharunderscore}{\kern0pt}Fevl\ htps\ {\isasympi}s\ {\isacharequal}{\kern0pt}\ Inr\ htree{\isachardoublequoteclose}\ \isanewline
\ \ \ \ \isakeyword{shows}\ {\isachardoublequoteopen}strict{\isacharunderscore}{\kern0pt}sorted\ {\isacharparenleft}{\kern0pt}map\ fst\ {\isacharparenleft}{\kern0pt}avl{\isacharunderscore}{\kern0pt}inorder\ htree{\isacharparenright}{\kern0pt}{\isacharparenright}{\kern0pt}{\isachardoublequoteclose}\isanewline
%
\isadelimproof
\ \ \ \ %
\endisadelimproof
%
\isatagproof
\isacommand{using}\isamarkupfalse%
\ assms\isanewline
\ \ \isacommand{proof}\isamarkupfalse%
\ {\isacharparenleft}{\kern0pt}induction\ {\isasympi}s\ arbitrary{\isacharcolon}{\kern0pt}\ htree{\isacharparenright}{\kern0pt}\isanewline
\ \ \ \ \isacommand{case}\isamarkupfalse%
\ Nil\isanewline
\ \ \ \ \isacommand{then}\isamarkupfalse%
\ \isacommand{show}\isamarkupfalse%
\ {\isacharquery}{\kern0pt}case\ \isanewline
\ \ \ \ \ \ \isacommand{by}\isamarkupfalse%
\ {\isacharparenleft}{\kern0pt}auto\ simp{\isacharcolon}{\kern0pt}\ avl{\isacharunderscore}{\kern0pt}inorder{\isacharunderscore}{\kern0pt}def{\isacharparenright}{\kern0pt}\isanewline
\ \ \isacommand{next}\isamarkupfalse%
\isanewline
\ \ \ \ \isacommand{case}\isamarkupfalse%
\ {\isacharparenleft}{\kern0pt}Cons\ {\isasympi}\ {\isasympi}s{\isacharparenright}{\kern0pt}\isanewline
\ \ \ \ \isacommand{then}\isamarkupfalse%
\ \isacommand{obtain}\isamarkupfalse%
\ htree{\isacharprime}{\kern0pt}\ \isakeyword{where}\ {\isachardoublequoteopen}simplify{\isacharunderscore}{\kern0pt}planE{\isacharunderscore}{\kern0pt}avl{\isacharprime}{\kern0pt}\ STG\ mp{\isacharunderscore}{\kern0pt}objT\ mp{\isacharunderscore}{\kern0pt}Fevl\ htps\ {\isasympi}s\ {\isacharequal}{\kern0pt}\ Inr\ htree{\isacharprime}{\kern0pt}{\isachardoublequoteclose}\isanewline
\ \ \ \ \ \ \isacommand{by}\isamarkupfalse%
\ {\isacharparenleft}{\kern0pt}auto\ simp{\isacharcolon}{\kern0pt}\ return{\isacharunderscore}{\kern0pt}iff{\isacharparenright}{\kern0pt}\isanewline
\ \ \ \ \isacommand{then}\isamarkupfalse%
\ \isacommand{show}\isamarkupfalse%
\ {\isacharquery}{\kern0pt}case\ \isanewline
\ \ \ \ \ \ \isacommand{using}\isamarkupfalse%
\ Cons{\isachardot}{\kern0pt}prems\ Cons{\isachardot}{\kern0pt}IH\ insert{\isacharunderscore}{\kern0pt}timed{\isacharunderscore}{\kern0pt}ground{\isacharunderscore}{\kern0pt}acts{\isacharunderscore}{\kern0pt}to{\isacharunderscore}{\kern0pt}tree{\isacharunderscore}{\kern0pt}strict{\isacharunderscore}{\kern0pt}sorted\ \isacommand{by}\isamarkupfalse%
\ auto\isanewline
\ \ \isacommand{qed}\isamarkupfalse%
%
\endisatagproof
{\isafoldproof}%
%
\isadelimproof
\isanewline
%
\endisadelimproof
\isanewline
\ \ \isacommand{lemma}\isamarkupfalse%
\ simplify{\isacharunderscore}{\kern0pt}planE{\isacharunderscore}{\kern0pt}avl{\isacharprime}{\kern0pt}{\isacharunderscore}{\kern0pt}equiv{\isacharcolon}{\kern0pt}\ \isanewline
\ \ \ \ \isakeyword{assumes}\ {\isachardoublequoteopen}simplify{\isacharunderscore}{\kern0pt}planE\ STG\ mp{\isacharunderscore}{\kern0pt}objT\ mp{\isacharunderscore}{\kern0pt}Fevl\ htps\ {\isasympi}s\ {\isacharequal}{\kern0pt}\ Inr\ hs{\isachardoublequoteclose}\ \isanewline
\ \ \ \ \ \ \ \ \isakeyword{and}\ {\isachardoublequoteopen}simplify{\isacharunderscore}{\kern0pt}planE{\isacharunderscore}{\kern0pt}avl{\isacharprime}{\kern0pt}\ STG\ mp{\isacharunderscore}{\kern0pt}objT\ mp{\isacharunderscore}{\kern0pt}Fevl\ htps\ {\isasympi}s\ {\isacharequal}{\kern0pt}\ Inr\ htree{\isachardoublequoteclose}\ \isanewline
\ \ \ \ \isakeyword{shows}\ {\isachardoublequoteopen}hs\ {\isacharequal}{\kern0pt}\ avl{\isacharunderscore}{\kern0pt}inorder\ htree{\isachardoublequoteclose}\isanewline
%
\isadelimproof
\ \ \ \ %
\endisadelimproof
%
\isatagproof
\isacommand{using}\isamarkupfalse%
\ assms\ \isacommand{unfolding}\isamarkupfalse%
\ avl{\isacharunderscore}{\kern0pt}inorder{\isacharunderscore}{\kern0pt}def\isanewline
\ \ \isacommand{proof}\isamarkupfalse%
\ {\isacharparenleft}{\kern0pt}induction\ {\isasympi}s\ arbitrary{\isacharcolon}{\kern0pt}\ hs\ htree{\isacharparenright}{\kern0pt}\isanewline
\ \ \ \ \isacommand{case}\isamarkupfalse%
\ Nil\isanewline
\ \ \ \ \isacommand{then}\isamarkupfalse%
\ \isacommand{show}\isamarkupfalse%
\ {\isacharquery}{\kern0pt}case\ \isacommand{by}\isamarkupfalse%
\ auto\isanewline
\ \ \isacommand{next}\isamarkupfalse%
\isanewline
\ \ \ \ \isacommand{case}\isamarkupfalse%
\ {\isacharparenleft}{\kern0pt}Cons\ {\isasympi}\ {\isasympi}s{\isacharparenright}{\kern0pt}\isanewline
\ \ \ \ \isacommand{then}\isamarkupfalse%
\ \isacommand{obtain}\isamarkupfalse%
\ hs{\isacharprime}{\kern0pt}\ htree{\isacharprime}{\kern0pt}\ \isakeyword{where}\ {\isachardoublequoteopen}simplify{\isacharunderscore}{\kern0pt}planE\ STG\ mp{\isacharunderscore}{\kern0pt}objT\ mp{\isacharunderscore}{\kern0pt}Fevl\ htps\ {\isasympi}s\ {\isacharequal}{\kern0pt}\ Inr\ hs{\isacharprime}{\kern0pt}{\isachardoublequoteclose}\isanewline
\ \ \ \ \ \ \ \ \ \ \ \ \ \ \ \ \ \ \ \ \ \ \ \ \ \ \ \ \ \isakeyword{and}\ {\isachardoublequoteopen}simplify{\isacharunderscore}{\kern0pt}planE{\isacharunderscore}{\kern0pt}avl{\isacharprime}{\kern0pt}\ STG\ mp{\isacharunderscore}{\kern0pt}objT\ mp{\isacharunderscore}{\kern0pt}Fevl\ htps\ {\isasympi}s\ {\isacharequal}{\kern0pt}\ Inr\ htree{\isacharprime}{\kern0pt}{\isachardoublequoteclose}\isanewline
\ \ \ \ \ \ \isacommand{by}\isamarkupfalse%
\ {\isacharparenleft}{\kern0pt}auto\ simp{\isacharcolon}{\kern0pt}\ return{\isacharunderscore}{\kern0pt}iff{\isacharparenright}{\kern0pt}\isanewline
\ \ \ \ \isacommand{then}\isamarkupfalse%
\ \isacommand{have}\isamarkupfalse%
\ {\isachardoublequoteopen}hs{\isacharprime}{\kern0pt}\ {\isacharequal}{\kern0pt}\ map\ fst\ {\isacharparenleft}{\kern0pt}inorder\ htree{\isacharprime}{\kern0pt}{\isacharparenright}{\kern0pt}{\isachardoublequoteclose}\isanewline
\ \ \ \ \ \ \isacommand{using}\isamarkupfalse%
\ Cons{\isachardot}{\kern0pt}prems\ Cons{\isachardot}{\kern0pt}IH\ \isacommand{by}\isamarkupfalse%
\ auto\isanewline
\ \ \ \ \isacommand{then}\isamarkupfalse%
\ \isacommand{show}\isamarkupfalse%
\ {\isacharquery}{\kern0pt}case\ \isanewline
\ \ \ \ \ \ \isacommand{using}\isamarkupfalse%
\ Cons{\isachardot}{\kern0pt}prems\ \isanewline
\ \ \ \ \ \ \ \ \ \ \ \ insert{\isacharunderscore}{\kern0pt}timed{\isacharunderscore}{\kern0pt}ground{\isacharunderscore}{\kern0pt}acts{\isacharunderscore}{\kern0pt}to{\isacharunderscore}{\kern0pt}tree{\isacharunderscore}{\kern0pt}equiv{\isacharbrackleft}{\kern0pt}OF\ simplify{\isacharunderscore}{\kern0pt}planE{\isacharunderscore}{\kern0pt}avl{\isacharprime}{\kern0pt}{\isacharunderscore}{\kern0pt}strict{\isacharunderscore}{\kern0pt}sorted{\isacharbrackright}{\kern0pt}\isanewline
\ \ \ \ \ \ \ \ \ \ \ \ {\isacartoucheopen}simplify{\isacharunderscore}{\kern0pt}planE\ STG\ mp{\isacharunderscore}{\kern0pt}objT\ mp{\isacharunderscore}{\kern0pt}Fevl\ htps\ {\isasympi}s\ {\isacharequal}{\kern0pt}\ Inr\ hs{\isacharprime}{\kern0pt}{\isacartoucheclose}\isanewline
\ \ \ \ \ \ \ \ \ \ \ \ {\isacartoucheopen}simplify{\isacharunderscore}{\kern0pt}planE{\isacharunderscore}{\kern0pt}avl{\isacharprime}{\kern0pt}\ STG\ mp{\isacharunderscore}{\kern0pt}objT\ mp{\isacharunderscore}{\kern0pt}Fevl\ htps\ {\isasympi}s\ {\isacharequal}{\kern0pt}\ Inr\ htree{\isacharprime}{\kern0pt}{\isacartoucheclose}\isanewline
\ \ \ \ \ \ \isacommand{by}\isamarkupfalse%
\ {\isacharparenleft}{\kern0pt}auto\ simp{\isacharcolon}{\kern0pt}\ avl{\isacharunderscore}{\kern0pt}inorder{\isacharunderscore}{\kern0pt}def{\isacharparenright}{\kern0pt}\isanewline
\ \ \isacommand{qed}\isamarkupfalse%
%
\endisatagproof
{\isafoldproof}%
%
\isadelimproof
\ \isanewline
%
\endisadelimproof
\isanewline
\ \ \isacommand{lemma}\isamarkupfalse%
\ simplify{\isacharunderscore}{\kern0pt}planE{\isacharunderscore}{\kern0pt}avl{\isacharunderscore}{\kern0pt}equiv{\isacharunderscore}{\kern0pt}Inr{\isacharcolon}{\kern0pt}\ \isanewline
\ \ \ \ \isakeyword{shows}\ {\isachardoublequoteopen}{\isacharparenleft}{\kern0pt}{\isasymexists}hs{\isachardot}{\kern0pt}\ simplify{\isacharunderscore}{\kern0pt}planE\ STG\ mp{\isacharunderscore}{\kern0pt}objT\ mp{\isacharunderscore}{\kern0pt}Fevl\ htps\ {\isasympi}s\ {\isacharequal}{\kern0pt}\ Inr\ hs{\isacharparenright}{\kern0pt}\ \isanewline
\ \ \ \ \ \ \ \ {\isasymlongleftrightarrow}\ {\isacharparenleft}{\kern0pt}{\isasymexists}htree{\isachardot}{\kern0pt}\ simplify{\isacharunderscore}{\kern0pt}planE{\isacharunderscore}{\kern0pt}avl{\isacharprime}{\kern0pt}\ STG\ mp{\isacharunderscore}{\kern0pt}objT\ mp{\isacharunderscore}{\kern0pt}Fevl\ htps\ {\isasympi}s\ {\isacharequal}{\kern0pt}\ Inr\ htree{\isacharparenright}{\kern0pt}{\isachardoublequoteclose}\isanewline
%
\isadelimproof
\ \ \ \ %
\endisadelimproof
%
\isatagproof
\isacommand{by}\isamarkupfalse%
\ {\isacharparenleft}{\kern0pt}induction\ {\isasympi}s{\isacharparenright}{\kern0pt}\ {\isacharparenleft}{\kern0pt}auto\ simp{\isacharcolon}{\kern0pt}\ return{\isacharunderscore}{\kern0pt}iff{\isacharparenright}{\kern0pt}%
\endisatagproof
{\isafoldproof}%
%
\isadelimproof
\isanewline
%
\endisadelimproof
\isanewline
\ \ \isanewline
\ \ \isacommand{lemma}\isamarkupfalse%
\ simplify{\isacharunderscore}{\kern0pt}planE{\isacharunderscore}{\kern0pt}avl{\isacharunderscore}{\kern0pt}equiv{\isacharcolon}{\kern0pt}\ \isanewline
\ \ \ \ \isakeyword{shows}\ {\isachardoublequoteopen}simplify{\isacharunderscore}{\kern0pt}planE\ STG\ mp{\isacharunderscore}{\kern0pt}objT\ mp{\isacharunderscore}{\kern0pt}Fevl\ htps\ {\isasympi}s\ {\isacharequal}{\kern0pt}\ Inr\ hs\ \isanewline
\ \ \ \ \ \ \ \ \ {\isasymlongleftrightarrow}\ simplify{\isacharunderscore}{\kern0pt}planE{\isacharunderscore}{\kern0pt}avl\ STG\ mp{\isacharunderscore}{\kern0pt}objT\ mp{\isacharunderscore}{\kern0pt}Fevl\ htps\ {\isasympi}s\ {\isacharequal}{\kern0pt}\ Inr\ hs{\isachardoublequoteclose}\isanewline
%
\isadelimproof
\ \ \ \ %
\endisadelimproof
%
\isatagproof
\isacommand{using}\isamarkupfalse%
\ simplify{\isacharunderscore}{\kern0pt}planE{\isacharunderscore}{\kern0pt}avl{\isacharprime}{\kern0pt}{\isacharunderscore}{\kern0pt}equiv\ simplify{\isacharunderscore}{\kern0pt}planE{\isacharunderscore}{\kern0pt}avl{\isacharunderscore}{\kern0pt}equiv{\isacharunderscore}{\kern0pt}Inr\ \isacommand{by}\isamarkupfalse%
\ fastforce%
\endisatagproof
{\isafoldproof}%
%
\isadelimproof
\isanewline
%
\endisadelimproof
\isanewline
\ \ \isanewline
\ \ \isacommand{lemma}\isamarkupfalse%
\ {\isacharparenleft}{\kern0pt}\isakeyword{in}\ wf{\isacharunderscore}{\kern0pt}ast{\isacharunderscore}{\kern0pt}problem{\isacharparenright}{\kern0pt}\isanewline
\ \ \ \ \isakeyword{assumes}\ {\isachardoublequoteopen}wf{\isacharunderscore}{\kern0pt}domain{\isachardoublequoteclose}\isanewline
\ \ \ \ \ \ \ \ \isakeyword{and}\ {\isachardoublequoteopen}simplify{\isacharunderscore}{\kern0pt}planE{\isacharunderscore}{\kern0pt}avl\ STG\ mp{\isacharunderscore}{\kern0pt}objT\ mp{\isacharunderscore}{\kern0pt}Fevl\ {\isacharparenleft}{\kern0pt}htps{\isacharunderscore}{\kern0pt}exec\ {\isasympi}s{\isacharparenright}{\kern0pt}\ {\isasympi}s\ {\isacharequal}{\kern0pt}\ Inr\ hs\ {\isachardoublequoteclose}\isanewline
\ \ \ \ \isakeyword{shows}\ {\isachardoublequoteopen}ind{\isacharunderscore}{\kern0pt}happ{\isacharunderscore}{\kern0pt}seq\ {\isasympi}s\ hs{\isachardoublequoteclose}\ \isakeyword{and}\ {\isachardoublequoteopen}wf{\isacharunderscore}{\kern0pt}plan\ {\isasympi}s{\isachardoublequoteclose}\isanewline
%
\isadelimproof
\ \ \ \ %
\endisadelimproof
%
\isatagproof
\isacommand{using}\isamarkupfalse%
\ assms\ simplify{\isacharunderscore}{\kern0pt}planE{\isacharunderscore}{\kern0pt}avl{\isacharunderscore}{\kern0pt}equiv\ simplify{\isacharunderscore}{\kern0pt}planE{\isacharunderscore}{\kern0pt}return{\isacharunderscore}{\kern0pt}iff\ \isanewline
\ \ \ \ \ \ \ \ \ \ htps{\isacharunderscore}{\kern0pt}exec{\isacharunderscore}{\kern0pt}sorted\ htps{\isacharunderscore}{\kern0pt}exec{\isacharunderscore}{\kern0pt}distinct\ simplify{\isacharunderscore}{\kern0pt}plan{\isacharunderscore}{\kern0pt}correct\isanewline
\ \ \ \ \isacommand{by}\isamarkupfalse%
\ {\isacharparenleft}{\kern0pt}auto\ simp{\isacharcolon}{\kern0pt}\ strict{\isacharunderscore}{\kern0pt}sorted{\isacharunderscore}{\kern0pt}iff{\isacharparenright}{\kern0pt}%
\endisatagproof
{\isafoldproof}%
%
\isadelimproof
%
\endisadelimproof
%
\begin{isamarkuptext}%
Lift \isa{valid{\isacharunderscore}{\kern0pt}happ{\isacharunderscore}{\kern0pt}seq{\isacharunderscore}{\kern0pt}from} into an Error Monad, and use efficient 
  happening execution \isa{en{\isacharunderscore}{\kern0pt}exE}%
\end{isamarkuptext}\isamarkuptrue%
\ \ \isacommand{fun}\isamarkupfalse%
\ valid{\isacharunderscore}{\kern0pt}happ{\isacharunderscore}{\kern0pt}seq{\isacharunderscore}{\kern0pt}fromE\ {\isacharcolon}{\kern0pt}{\isacharcolon}{\kern0pt}\ {\isachardoublequoteopen}nat\ {\isasymRightarrow}\ world{\isacharunderscore}{\kern0pt}model\ {\isasymRightarrow}\ happening\ list\ {\isasymRightarrow}\ {\isacharunderscore}{\kern0pt}{\isacharplus}{\kern0pt}unit{\isachardoublequoteclose}\ \isakeyword{where}\isanewline
\ \ \ \ {\isachardoublequoteopen}valid{\isacharunderscore}{\kern0pt}happ{\isacharunderscore}{\kern0pt}seq{\isacharunderscore}{\kern0pt}fromE\ si\ s\ {\isacharbrackleft}{\kern0pt}{\isacharbrackright}{\kern0pt}\ {\isacharequal}{\kern0pt}\ \isanewline
\ \ \ \ \ \ check\ {\isacharparenleft}{\kern0pt}holds\ s\ {\isacharparenleft}{\kern0pt}goal\ P{\isacharparenright}{\kern0pt}{\isacharparenright}{\kern0pt}\ {\isacharparenleft}{\kern0pt}ERRS\ {\isacharprime}{\kern0pt}{\isacharprime}{\kern0pt}Postcondition\ does\ not\ hold{\isacharprime}{\kern0pt}{\isacharprime}{\kern0pt}{\isacharparenright}{\kern0pt}{\isachardoublequoteclose}\isanewline
\ \ {\isacharbar}{\kern0pt}\ {\isachardoublequoteopen}valid{\isacharunderscore}{\kern0pt}happ{\isacharunderscore}{\kern0pt}seq{\isacharunderscore}{\kern0pt}fromE\ si\ s\ {\isacharparenleft}{\kern0pt}h{\isacharhash}{\kern0pt}hs{\isacharparenright}{\kern0pt}\ {\isacharequal}{\kern0pt}\ do\ {\isacharbraceleft}{\kern0pt}\isanewline
\ \ \ \ \ \ s\ {\isasymleftarrow}\ en{\isacharunderscore}{\kern0pt}exE\ h\ s\ {\isacharless}{\kern0pt}{\isacharplus}{\kern0pt}{\isacharquery}{\kern0pt}\ {\isacharparenleft}{\kern0pt}{\isasymlambda}e\ {\isacharunderscore}{\kern0pt}{\isachardot}{\kern0pt}\ shows\ {\isacharprime}{\kern0pt}{\isacharprime}{\kern0pt}at\ step\ {\isacharprime}{\kern0pt}{\isacharprime}{\kern0pt}\ o\ shows\ si\ o\ shows\ {\isacharprime}{\kern0pt}{\isacharprime}{\kern0pt}{\isacharcolon}{\kern0pt}\ {\isacharprime}{\kern0pt}{\isacharprime}{\kern0pt}\ o\ e\ {\isacharparenleft}{\kern0pt}{\isacharparenright}{\kern0pt}{\isacharparenright}{\kern0pt}{\isacharsemicolon}{\kern0pt}\isanewline
\ \ \ \ \ \ valid{\isacharunderscore}{\kern0pt}happ{\isacharunderscore}{\kern0pt}seq{\isacharunderscore}{\kern0pt}fromE\ {\isacharparenleft}{\kern0pt}si{\isacharplus}{\kern0pt}{\isadigit{1}}{\isacharparenright}{\kern0pt}\ s\ hs\isanewline
\ \ \ \ {\isacharbraceright}{\kern0pt}{\isachardoublequoteclose}%
\begin{isamarkuptext}%
For the refinement, we need to show that the world models only
    contain atoms, i.e., containing only atoms is an invariant under execution
    of well-formed happenings.%
\end{isamarkuptext}\isamarkuptrue%
\ \ \isacommand{lemma}\isamarkupfalse%
\ {\isacharparenleft}{\kern0pt}\isakeyword{in}\ wf{\isacharunderscore}{\kern0pt}ast{\isacharunderscore}{\kern0pt}problem{\isacharparenright}{\kern0pt}\ wf{\isacharunderscore}{\kern0pt}happ{\isacharunderscore}{\kern0pt}only{\isacharunderscore}{\kern0pt}add{\isacharunderscore}{\kern0pt}atoms{\isacharcolon}{\kern0pt}\isanewline
\ \ \ \ \isakeyword{assumes}\ {\isachardoublequoteopen}wm{\isacharunderscore}{\kern0pt}basic\ s{\isachardoublequoteclose}\ \isanewline
\ \ \ \ \ \ \ \ \isakeyword{and}\ {\isachardoublequoteopen}{\isasymforall}a\ {\isasymin}\ set\ as{\isachardot}{\kern0pt}\ wf{\isacharunderscore}{\kern0pt}ground{\isacharunderscore}{\kern0pt}action\ a{\isachardoublequoteclose}\isanewline
\ \ \ \ \isakeyword{shows}\ {\isachardoublequoteopen}wm{\isacharunderscore}{\kern0pt}basic\ {\isacharparenleft}{\kern0pt}apply{\isacharunderscore}{\kern0pt}happ{\isacharunderscore}{\kern0pt}exec\ {\isacharparenleft}{\kern0pt}t{\isacharcomma}{\kern0pt}as{\isacharparenright}{\kern0pt}\ s{\isacharparenright}{\kern0pt}{\isachardoublequoteclose}\isanewline
%
\isadelimproof
\ \ \ \ %
\endisadelimproof
%
\isatagproof
\isacommand{using}\isamarkupfalse%
\ wf{\isacharunderscore}{\kern0pt}problem\ wf{\isacharunderscore}{\kern0pt}domain\isanewline
\ \ \ \ \isacommand{unfolding}\isamarkupfalse%
\ wf{\isacharunderscore}{\kern0pt}problem{\isacharunderscore}{\kern0pt}def\ wf{\isacharunderscore}{\kern0pt}domain{\isacharunderscore}{\kern0pt}def\isanewline
\ \ \ \ \isacommand{using}\isamarkupfalse%
\ assms\ \isanewline
\ \ \ \ \ \ \ \ \ \ apply{\isacharunderscore}{\kern0pt}happ{\isacharunderscore}{\kern0pt}exec{\isacharunderscore}{\kern0pt}refine\ \isanewline
\ \ \ \ \ \ \ \ \ \ wf{\isacharunderscore}{\kern0pt}fmla{\isacharunderscore}{\kern0pt}atom{\isacharunderscore}{\kern0pt}alt\ \isanewline
\ \ \ \ \ \ \ \ \ \ wf{\isacharunderscore}{\kern0pt}ground{\isacharunderscore}{\kern0pt}action{\isacharunderscore}{\kern0pt}wf{\isacharunderscore}{\kern0pt}fmla{\isacharunderscore}{\kern0pt}atom\isanewline
\ \ \ \ \isacommand{by}\isamarkupfalse%
\ {\isacharparenleft}{\kern0pt}auto\ simp{\isacharcolon}{\kern0pt}\ wm{\isacharunderscore}{\kern0pt}basic{\isacharunderscore}{\kern0pt}def{\isacharparenright}{\kern0pt}\ fastforce%
\endisatagproof
{\isafoldproof}%
%
\isadelimproof
\isanewline
%
\endisadelimproof
\isanewline
\ \ \isacommand{lemma}\isamarkupfalse%
\ {\isacharparenleft}{\kern0pt}\isakeyword{in}\ wf{\isacharunderscore}{\kern0pt}ast{\isacharunderscore}{\kern0pt}problem{\isacharparenright}{\kern0pt}\ wf{\isacharunderscore}{\kern0pt}happ{\isacharunderscore}{\kern0pt}only{\isacharunderscore}{\kern0pt}add{\isacharunderscore}{\kern0pt}atoms{\isacharprime}{\kern0pt}{\isacharcolon}{\kern0pt}\isanewline
\ \ \ \ \isakeyword{assumes}\ {\isachardoublequoteopen}wf{\isacharunderscore}{\kern0pt}world{\isacharunderscore}{\kern0pt}model\ s{\isachardoublequoteclose}\ \isanewline
\ \ \ \ \ \ \ \ \isakeyword{and}\ {\isachardoublequoteopen}{\isasymforall}a\ {\isasymin}\ set\ as{\isachardot}{\kern0pt}\ wf{\isacharunderscore}{\kern0pt}ground{\isacharunderscore}{\kern0pt}action\ a{\isachardoublequoteclose}\isanewline
\ \ \ \ \isakeyword{shows}\ {\isachardoublequoteopen}wf{\isacharunderscore}{\kern0pt}world{\isacharunderscore}{\kern0pt}model\ {\isacharparenleft}{\kern0pt}apply{\isacharunderscore}{\kern0pt}happ{\isacharunderscore}{\kern0pt}exec\ {\isacharparenleft}{\kern0pt}t{\isacharcomma}{\kern0pt}as{\isacharparenright}{\kern0pt}\ s{\isacharparenright}{\kern0pt}{\isachardoublequoteclose}\isanewline
%
\isadelimproof
\ \ \ \ %
\endisadelimproof
%
\isatagproof
\isacommand{using}\isamarkupfalse%
\ wf{\isacharunderscore}{\kern0pt}problem\ wf{\isacharunderscore}{\kern0pt}domain\isanewline
\ \ \ \ \isacommand{unfolding}\isamarkupfalse%
\ wf{\isacharunderscore}{\kern0pt}problem{\isacharunderscore}{\kern0pt}def\ wf{\isacharunderscore}{\kern0pt}domain{\isacharunderscore}{\kern0pt}def\isanewline
\ \ \ \ \isacommand{using}\isamarkupfalse%
\ assms\ \isanewline
\ \ \ \ \ \ \ \ \ \ apply{\isacharunderscore}{\kern0pt}happ{\isacharunderscore}{\kern0pt}exec{\isacharunderscore}{\kern0pt}refine\ \isanewline
\ \ \ \ \ \ \ \ \ \ wf{\isacharunderscore}{\kern0pt}fmla{\isacharunderscore}{\kern0pt}atom{\isacharunderscore}{\kern0pt}alt\ \isanewline
\ \ \ \ \ \ \ \ \ \ wf{\isacharunderscore}{\kern0pt}ground{\isacharunderscore}{\kern0pt}action{\isacharunderscore}{\kern0pt}wf{\isacharunderscore}{\kern0pt}fmla{\isacharunderscore}{\kern0pt}atom\isanewline
\ \ \ \ \isacommand{by}\isamarkupfalse%
\ {\isacharparenleft}{\kern0pt}auto\ simp{\isacharcolon}{\kern0pt}\ wf{\isacharunderscore}{\kern0pt}world{\isacharunderscore}{\kern0pt}model{\isacharunderscore}{\kern0pt}def{\isacharparenright}{\kern0pt}%
\endisatagproof
{\isafoldproof}%
%
\isadelimproof
%
\endisadelimproof
%
\begin{isamarkuptext}%
Justification for refinement of \isa{valid{\isacharunderscore}{\kern0pt}happ{\isacharunderscore}{\kern0pt}seq{\isacharunderscore}{\kern0pt}from}.%
\end{isamarkuptext}\isamarkuptrue%
\ \ \isacommand{lemma}\isamarkupfalse%
\ {\isacharparenleft}{\kern0pt}\isakeyword{in}\ wf{\isacharunderscore}{\kern0pt}ast{\isacharunderscore}{\kern0pt}problem{\isacharparenright}{\kern0pt}\ valid{\isacharunderscore}{\kern0pt}happ{\isacharunderscore}{\kern0pt}seq{\isacharunderscore}{\kern0pt}fromE{\isacharunderscore}{\kern0pt}return{\isacharunderscore}{\kern0pt}iff{\isacharbrackleft}{\kern0pt}return{\isacharunderscore}{\kern0pt}iff{\isacharbrackright}{\kern0pt}{\isacharcolon}{\kern0pt}\ \isanewline
\ \ \ \ \isakeyword{assumes}\ {\isachardoublequoteopen}wm{\isacharunderscore}{\kern0pt}basic\ s{\isachardoublequoteclose}\ \isanewline
\ \ \ \ \ \ \ \ \isakeyword{and}\ {\isachardoublequoteopen}{\isasymforall}{\isacharparenleft}{\kern0pt}t{\isacharcomma}{\kern0pt}as{\isacharparenright}{\kern0pt}\ {\isasymin}\ set\ hs{\isachardot}{\kern0pt}\ {\isasymforall}a\ {\isasymin}\ set\ as{\isachardot}{\kern0pt}\ wf{\isacharunderscore}{\kern0pt}ground{\isacharunderscore}{\kern0pt}action\ a{\isachardoublequoteclose}\isanewline
\ \ \ \ \isakeyword{shows}\ {\isachardoublequoteopen}valid{\isacharunderscore}{\kern0pt}happ{\isacharunderscore}{\kern0pt}seq{\isacharunderscore}{\kern0pt}from\ s\ hs\ {\isasymlongleftrightarrow}\ valid{\isacharunderscore}{\kern0pt}happ{\isacharunderscore}{\kern0pt}seq{\isacharunderscore}{\kern0pt}fromE\ k\ s\ hs\ {\isacharequal}{\kern0pt}\ Inr\ {\isacharparenleft}{\kern0pt}{\isacharparenright}{\kern0pt}{\isachardoublequoteclose}\isanewline
%
\isadelimproof
\ \ \ \ %
\endisadelimproof
%
\isatagproof
\isacommand{using}\isamarkupfalse%
\ assms\ \isanewline
\ \ \ \ \ \ \ \ \ \ holds{\isacharunderscore}{\kern0pt}for{\isacharunderscore}{\kern0pt}wf{\isacharunderscore}{\kern0pt}fmlas\ \isanewline
\ \ \ \ \ \ \ \ \ \ wf{\isacharunderscore}{\kern0pt}happ{\isacharunderscore}{\kern0pt}only{\isacharunderscore}{\kern0pt}add{\isacharunderscore}{\kern0pt}atoms\ \isanewline
\ \ \ \ \ \ \ \ \ \ apply{\isacharunderscore}{\kern0pt}happ{\isacharunderscore}{\kern0pt}exec{\isacharunderscore}{\kern0pt}refine\ \ \isanewline
\ \ \ \ \isacommand{by}\isamarkupfalse%
\ {\isacharparenleft}{\kern0pt}induction\ hs\ arbitrary{\isacharcolon}{\kern0pt}\ s\ k{\isacharparenright}{\kern0pt}\ {\isacharparenleft}{\kern0pt}auto\ simp{\isacharcolon}{\kern0pt}\ return{\isacharunderscore}{\kern0pt}iff\ en{\isacharunderscore}{\kern0pt}exE{\isacharunderscore}{\kern0pt}return{\isacharunderscore}{\kern0pt}iff{\isacharparenright}{\kern0pt}%
\endisatagproof
{\isafoldproof}%
%
\isadelimproof
%
\endisadelimproof
%
\begin{isamarkuptext}%
Now we can define a refinement for \isa{valid{\isacharunderscore}{\kern0pt}plan{\isacharunderscore}{\kern0pt}from{\isadigit{2}}}.%
\end{isamarkuptext}\isamarkuptrue%
\ \ \isacommand{fun}\isamarkupfalse%
\ valid{\isacharunderscore}{\kern0pt}plan{\isacharunderscore}{\kern0pt}from{\isadigit{2}}{\isacharunderscore}{\kern0pt}exec\ {\isacharcolon}{\kern0pt}{\isacharcolon}{\kern0pt}\ {\isachardoublequoteopen}world{\isacharunderscore}{\kern0pt}model\ {\isasymRightarrow}\ plan\ {\isasymRightarrow}\ bool{\isachardoublequoteclose}\ \isakeyword{where}\isanewline
\ \ \ \ {\isachardoublequoteopen}valid{\isacharunderscore}{\kern0pt}plan{\isacharunderscore}{\kern0pt}from{\isadigit{2}}{\isacharunderscore}{\kern0pt}exec\ s\ {\isasympi}s\ {\isasymlongleftrightarrow}\ \isanewline
\ \ \ \ \ \ {\isacharparenleft}{\kern0pt}wf{\isacharunderscore}{\kern0pt}plan\ {\isasympi}s\ {\isasymand}\ valid{\isacharunderscore}{\kern0pt}happ{\isacharunderscore}{\kern0pt}seq{\isacharunderscore}{\kern0pt}from\ s\ {\isacharparenleft}{\kern0pt}simplify{\isacharunderscore}{\kern0pt}plan\ {\isacharparenleft}{\kern0pt}htps{\isacharunderscore}{\kern0pt}exec\ {\isasympi}s{\isacharparenright}{\kern0pt}\ {\isasympi}s{\isacharparenright}{\kern0pt}{\isacharparenright}{\kern0pt}{\isachardoublequoteclose}%
\begin{isamarkuptext}%
With \isa{{\isasymlbrakk}wf{\isacharunderscore}{\kern0pt}ast{\isacharunderscore}{\kern0pt}problem\ {\isacharquery}{\kern0pt}P{\isacharsemicolon}{\kern0pt}\ ast{\isacharunderscore}{\kern0pt}problem{\isachardot}{\kern0pt}ind{\isacharunderscore}{\kern0pt}happ{\isacharunderscore}{\kern0pt}seq\ {\isacharquery}{\kern0pt}P\ {\isacharquery}{\kern0pt}{\isasympi}s\ {\isacharquery}{\kern0pt}hs{\isacharsemicolon}{\kern0pt}\ ast{\isacharunderscore}{\kern0pt}problem{\isachardot}{\kern0pt}ind{\isacharunderscore}{\kern0pt}happ{\isacharunderscore}{\kern0pt}seq\ {\isacharquery}{\kern0pt}P\ {\isacharquery}{\kern0pt}{\isasympi}s\ {\isacharquery}{\kern0pt}hs{\isacharprime}{\kern0pt}{\isacharsemicolon}{\kern0pt}\ ast{\isacharunderscore}{\kern0pt}problem{\isachardot}{\kern0pt}wf{\isacharunderscore}{\kern0pt}plan\ {\isacharquery}{\kern0pt}P\ {\isacharquery}{\kern0pt}{\isasympi}s{\isacharsemicolon}{\kern0pt}\ valid{\isacharunderscore}{\kern0pt}happ{\isacharunderscore}{\kern0pt}seq\ {\isacharquery}{\kern0pt}M\isactrlsub {\isadigit{0}}\ {\isacharquery}{\kern0pt}hs\ {\isacharquery}{\kern0pt}M\isactrlsub n{\isasymrbrakk}\ {\isasymLongrightarrow}\ valid{\isacharunderscore}{\kern0pt}happ{\isacharunderscore}{\kern0pt}seq\ {\isacharquery}{\kern0pt}M\isactrlsub {\isadigit{0}}\ {\isacharquery}{\kern0pt}hs{\isacharprime}{\kern0pt}\ {\isacharquery}{\kern0pt}M\isactrlsub n} we can prove, that 
  if there exists a valid induced happening sequence, then simplify plan is also valid.%
\end{isamarkuptext}\isamarkuptrue%
\ \ \isacommand{lemma}\isamarkupfalse%
\ {\isacharparenleft}{\kern0pt}\isakeyword{in}\ wf{\isacharunderscore}{\kern0pt}ast{\isacharunderscore}{\kern0pt}problem{\isacharparenright}{\kern0pt}\ simplify{\isacharunderscore}{\kern0pt}plan{\isacharunderscore}{\kern0pt}valid{\isacharunderscore}{\kern0pt}exists{\isadigit{1}}{\isacharcolon}{\kern0pt}\isanewline
\ \ \ \ \isakeyword{fixes}\ {\isasympi}s\isanewline
\ \ \ \ \isakeyword{defines}\ {\isachardoublequoteopen}htps\ {\isasymequiv}\ htps{\isacharunderscore}{\kern0pt}exec\ {\isasympi}s{\isachardoublequoteclose}\isanewline
\ \ \ \ \isakeyword{assumes}\ {\isachardoublequoteopen}wf{\isacharunderscore}{\kern0pt}plan\ {\isasympi}s{\isachardoublequoteclose}\isanewline
\ \ \ \ \ \ \ \ \isakeyword{and}\ {\isachardoublequoteopen}{\isasymexists}hs{\isachardot}{\kern0pt}\ ind{\isacharunderscore}{\kern0pt}happ{\isacharunderscore}{\kern0pt}seq\ {\isasympi}s\ hs\ {\isasymand}\ valid{\isacharunderscore}{\kern0pt}happ{\isacharunderscore}{\kern0pt}seq\ M\ hs\ M{\isacharprime}{\kern0pt}\ {\isasymand}\ M{\isacharprime}{\kern0pt}\ \isactrlsup c{\isasymTTurnstile}\isactrlsub {\isacharequal}{\kern0pt}\ {\isacharparenleft}{\kern0pt}goal\ P{\isacharparenright}{\kern0pt}{\isachardoublequoteclose}\isanewline
\ \ \ \ \isakeyword{shows}\ {\isachardoublequoteopen}valid{\isacharunderscore}{\kern0pt}happ{\isacharunderscore}{\kern0pt}seq\ M\ {\isacharparenleft}{\kern0pt}simplify{\isacharunderscore}{\kern0pt}plan\ htps\ {\isasympi}s{\isacharparenright}{\kern0pt}\ M{\isacharprime}{\kern0pt}\ {\isasymand}\ M{\isacharprime}{\kern0pt}\ \isactrlsup c{\isasymTTurnstile}\isactrlsub {\isacharequal}{\kern0pt}\ {\isacharparenleft}{\kern0pt}goal\ P{\isacharparenright}{\kern0pt}{\isachardoublequoteclose}\isanewline
%
\isadelimproof
\ \ \ \ %
\endisadelimproof
%
\isatagproof
\isacommand{using}\isamarkupfalse%
\ assms\ simplify{\isacharunderscore}{\kern0pt}plan{\isacharunderscore}{\kern0pt}correct\ ind{\isacharunderscore}{\kern0pt}happ{\isacharunderscore}{\kern0pt}seq{\isacharunderscore}{\kern0pt}unique{\isacharunderscore}{\kern0pt}final{\isacharunderscore}{\kern0pt}state{\isacharprime}{\kern0pt}\ \isacommand{by}\isamarkupfalse%
\ blast%
\endisatagproof
{\isafoldproof}%
%
\isadelimproof
%
\endisadelimproof
%
\begin{isamarkuptext}%
Other way round is easy, since \isa{simplify{\isacharunderscore}{\kern0pt}plan} is an induced 
  happening sequence.%
\end{isamarkuptext}\isamarkuptrue%
\ \ \isacommand{lemma}\isamarkupfalse%
\ {\isacharparenleft}{\kern0pt}\isakeyword{in}\ wf{\isacharunderscore}{\kern0pt}ast{\isacharunderscore}{\kern0pt}problem{\isacharparenright}{\kern0pt}\ simplify{\isacharunderscore}{\kern0pt}plan{\isacharunderscore}{\kern0pt}valid{\isacharunderscore}{\kern0pt}exists{\isadigit{2}}{\isacharcolon}{\kern0pt}\ \isanewline
\ \ \ \ \isakeyword{fixes}\ {\isasympi}s\isanewline
\ \ \ \ \isakeyword{defines}\ {\isachardoublequoteopen}htps\ {\isasymequiv}\ htps{\isacharunderscore}{\kern0pt}exec\ {\isasympi}s{\isachardoublequoteclose}\isanewline
\ \ \ \ \isakeyword{assumes}\ {\isachardoublequoteopen}wf{\isacharunderscore}{\kern0pt}plan\ {\isasympi}s{\isachardoublequoteclose}\ \isanewline
\ \ \ \ \ \ \ \ \isakeyword{and}\ {\isachardoublequoteopen}valid{\isacharunderscore}{\kern0pt}happ{\isacharunderscore}{\kern0pt}seq\ M\ {\isacharparenleft}{\kern0pt}simplify{\isacharunderscore}{\kern0pt}plan\ htps\ {\isasympi}s{\isacharparenright}{\kern0pt}\ M{\isacharprime}{\kern0pt}\ {\isasymand}\ M{\isacharprime}{\kern0pt}\ \isactrlsup c{\isasymTTurnstile}\isactrlsub {\isacharequal}{\kern0pt}\ {\isacharparenleft}{\kern0pt}goal\ P{\isacharparenright}{\kern0pt}{\isachardoublequoteclose}\isanewline
\ \ \ \ \isakeyword{shows}\ {\isachardoublequoteopen}{\isasymexists}hs{\isachardot}{\kern0pt}\ ind{\isacharunderscore}{\kern0pt}happ{\isacharunderscore}{\kern0pt}seq\ {\isasympi}s\ hs\ {\isasymand}\ valid{\isacharunderscore}{\kern0pt}happ{\isacharunderscore}{\kern0pt}seq\ M\ hs\ M{\isacharprime}{\kern0pt}\ {\isasymand}\ M{\isacharprime}{\kern0pt}\ \isactrlsup c{\isasymTTurnstile}\isactrlsub {\isacharequal}{\kern0pt}\ {\isacharparenleft}{\kern0pt}goal\ P{\isacharparenright}{\kern0pt}{\isachardoublequoteclose}\isanewline
%
\isadelimproof
\ \ \ \ %
\endisadelimproof
%
\isatagproof
\isacommand{using}\isamarkupfalse%
\ assms\ simplify{\isacharunderscore}{\kern0pt}plan{\isacharunderscore}{\kern0pt}correct\ \isacommand{by}\isamarkupfalse%
\ blast%
\endisatagproof
{\isafoldproof}%
%
\isadelimproof
%
\endisadelimproof
%
\begin{isamarkuptext}%
Justification for the refinement of \isa{valid{\isacharunderscore}{\kern0pt}plan{\isacharunderscore}{\kern0pt}from{\isadigit{2}}}.%
\end{isamarkuptext}\isamarkuptrue%
\ \ \isacommand{lemma}\isamarkupfalse%
\ {\isacharparenleft}{\kern0pt}\isakeyword{in}\ wf{\isacharunderscore}{\kern0pt}ast{\isacharunderscore}{\kern0pt}problem{\isacharparenright}{\kern0pt}\ valid{\isacharunderscore}{\kern0pt}plan{\isacharunderscore}{\kern0pt}from{\isadigit{2}}{\isacharunderscore}{\kern0pt}exec{\isacharunderscore}{\kern0pt}refine{\isacharcolon}{\kern0pt}\ \isanewline
\ \ \ \ {\isachardoublequoteopen}valid{\isacharunderscore}{\kern0pt}plan{\isacharunderscore}{\kern0pt}from{\isadigit{2}}\ s\ {\isasympi}s\ {\isasymlongleftrightarrow}\ valid{\isacharunderscore}{\kern0pt}plan{\isacharunderscore}{\kern0pt}from{\isadigit{2}}{\isacharunderscore}{\kern0pt}exec\ s\ {\isasympi}s{\isachardoublequoteclose}\isanewline
%
\isadelimproof
\ \ \ \ %
\endisadelimproof
%
\isatagproof
\isacommand{using}\isamarkupfalse%
\ simplify{\isacharunderscore}{\kern0pt}plan{\isacharunderscore}{\kern0pt}valid{\isacharunderscore}{\kern0pt}exists{\isadigit{1}}\ \isanewline
\ \ \ \ \ \ \ \ \ \ simplify{\isacharunderscore}{\kern0pt}plan{\isacharunderscore}{\kern0pt}valid{\isacharunderscore}{\kern0pt}exists{\isadigit{2}}\isanewline
\ \ \ \ \ \ \ \ \ \ valid{\isacharunderscore}{\kern0pt}happ{\isacharunderscore}{\kern0pt}seq{\isacharunderscore}{\kern0pt}from{\isacharunderscore}{\kern0pt}refine{\isacharbrackleft}{\kern0pt}OF\ inst{\isacharunderscore}{\kern0pt}wf{\isacharunderscore}{\kern0pt}plan{\isacharunderscore}{\kern0pt}wf{\isacharunderscore}{\kern0pt}happ{\isacharunderscore}{\kern0pt}seq{\isacharbrackright}{\kern0pt}\isanewline
\ \ \ \ \ \ \ \ \ \ simplify{\isacharunderscore}{\kern0pt}plan{\isacharunderscore}{\kern0pt}correct\isanewline
\ \ \ \ \isacommand{unfolding}\isamarkupfalse%
\ valid{\isacharunderscore}{\kern0pt}plan{\isacharunderscore}{\kern0pt}from{\isadigit{2}}{\isacharunderscore}{\kern0pt}def\ plan{\isacharunderscore}{\kern0pt}happ{\isacharunderscore}{\kern0pt}path{\isacharunderscore}{\kern0pt}def\isanewline
\ \ \ \ \isacommand{by}\isamarkupfalse%
\ {\isacharparenleft}{\kern0pt}meson\ valid{\isacharunderscore}{\kern0pt}plan{\isacharunderscore}{\kern0pt}from{\isadigit{2}}{\isacharunderscore}{\kern0pt}exec{\isachardot}{\kern0pt}simps{\isacharparenright}{\kern0pt}%
\endisatagproof
{\isafoldproof}%
%
\isadelimproof
%
\endisadelimproof
%
\begin{isamarkuptext}%
Now we can also lift \isa{valid{\isacharunderscore}{\kern0pt}plan{\isacharunderscore}{\kern0pt}from{\isadigit{2}}{\isacharunderscore}{\kern0pt}exec} into an Error Monad, 
  using \isa{simplify{\isacharunderscore}{\kern0pt}planE} and \isa{valid{\isacharunderscore}{\kern0pt}happ{\isacharunderscore}{\kern0pt}seq{\isacharunderscore}{\kern0pt}fromE}%
\end{isamarkuptext}\isamarkuptrue%
\ \ \isacommand{fun}\isamarkupfalse%
\ valid{\isacharunderscore}{\kern0pt}plan{\isacharunderscore}{\kern0pt}from{\isadigit{2}}E\ {\isacharcolon}{\kern0pt}{\isacharcolon}{\kern0pt}\ {\isachardoublequoteopen}{\isacharunderscore}{\kern0pt}\ {\isasymRightarrow}\ {\isacharparenleft}{\kern0pt}object{\isacharcomma}{\kern0pt}\ type{\isacharparenright}{\kern0pt}\ mapping\ {\isasymRightarrow}\ {\isacharparenleft}{\kern0pt}object{\isacharcomma}{\kern0pt}\ rat{\isacharparenright}{\kern0pt}\ mapping\ {\isasymRightarrow}\ world{\isacharunderscore}{\kern0pt}model\ {\isasymRightarrow}\ plan\ {\isasymRightarrow}\ {\isacharunderscore}{\kern0pt}{\isacharplus}{\kern0pt}unit{\isachardoublequoteclose}\ \isakeyword{where}\isanewline
\ \ \ \ {\isachardoublequoteopen}valid{\isacharunderscore}{\kern0pt}plan{\isacharunderscore}{\kern0pt}from{\isadigit{2}}E\ stg\ mp\ mp{\isacharunderscore}{\kern0pt}fe\ s\ {\isasympi}s\ {\isacharequal}{\kern0pt}\ do\ {\isacharbraceleft}{\kern0pt}\isanewline
\ \ \ \ \ \ hs\ {\isasymleftarrow}\ simplify{\isacharunderscore}{\kern0pt}planE{\isacharunderscore}{\kern0pt}avl\ stg\ mp\ mp{\isacharunderscore}{\kern0pt}fe\ {\isacharparenleft}{\kern0pt}htps{\isacharunderscore}{\kern0pt}exec\ {\isasympi}s{\isacharparenright}{\kern0pt}\ {\isasympi}s{\isacharsemicolon}{\kern0pt}\ \ \isanewline
\ \ \ \ \ \ valid{\isacharunderscore}{\kern0pt}happ{\isacharunderscore}{\kern0pt}seq{\isacharunderscore}{\kern0pt}fromE\ {\isadigit{1}}\ s\ hs\isanewline
\ \ \ \ {\isacharbraceright}{\kern0pt}{\isachardoublequoteclose}\isanewline
\isanewline
\ \ \isanewline
\ \ \isanewline
\isanewline
\ \ \isacommand{lemma}\isamarkupfalse%
\ {\isacharparenleft}{\kern0pt}\isakeyword{in}\ wf{\isacharunderscore}{\kern0pt}ast{\isacharunderscore}{\kern0pt}problem{\isacharparenright}{\kern0pt}\ valid{\isacharunderscore}{\kern0pt}plan{\isacharunderscore}{\kern0pt}fromE{\isacharunderscore}{\kern0pt}wf{\isacharunderscore}{\kern0pt}plan{\isacharcolon}{\kern0pt}\isanewline
\ \ \ \ \isakeyword{assumes}\ {\isachardoublequoteopen}valid{\isacharunderscore}{\kern0pt}plan{\isacharunderscore}{\kern0pt}from{\isadigit{2}}E\ STG\ mp{\isacharunderscore}{\kern0pt}objT\ mp{\isacharunderscore}{\kern0pt}Fevl\ s\ {\isasympi}s\ {\isacharequal}{\kern0pt}\ Inr\ {\isacharparenleft}{\kern0pt}{\isacharparenright}{\kern0pt}{\isachardoublequoteclose}\ \isanewline
\ \ \ \ \isakeyword{shows}\ {\isachardoublequoteopen}wf{\isacharunderscore}{\kern0pt}plan\ {\isasympi}s{\isachardoublequoteclose}\isanewline
%
\isadelimproof
\ \ %
\endisadelimproof
%
\isatagproof
\isacommand{proof}\isamarkupfalse%
\ {\isacharminus}{\kern0pt}\isanewline
\ \ \ \ \isacommand{have}\isamarkupfalse%
\ {\isachardoublequoteopen}strict{\isacharunderscore}{\kern0pt}sorted\ {\isacharparenleft}{\kern0pt}htps{\isacharunderscore}{\kern0pt}exec\ {\isasympi}s{\isacharparenright}{\kern0pt}{\isachardoublequoteclose}\isanewline
\ \ \ \ \ \ \isacommand{using}\isamarkupfalse%
\ htps{\isacharunderscore}{\kern0pt}exec{\isacharunderscore}{\kern0pt}distinct\ htps{\isacharunderscore}{\kern0pt}exec{\isacharunderscore}{\kern0pt}sorted\ \isacommand{by}\isamarkupfalse%
\ {\isacharparenleft}{\kern0pt}simp\ add{\isacharcolon}{\kern0pt}\ strict{\isacharunderscore}{\kern0pt}sorted{\isacharunderscore}{\kern0pt}iff{\isacharparenright}{\kern0pt}\isanewline
\ \ \ \ \isacommand{then}\isamarkupfalse%
\ \isacommand{show}\isamarkupfalse%
\ {\isacharquery}{\kern0pt}thesis\isanewline
\ \ \ \ \ \ \isacommand{using}\isamarkupfalse%
\ assms\ wf{\isacharunderscore}{\kern0pt}problem\ simplify{\isacharunderscore}{\kern0pt}planE{\isacharunderscore}{\kern0pt}return{\isacharunderscore}{\kern0pt}iff\ simplify{\isacharunderscore}{\kern0pt}planE{\isacharunderscore}{\kern0pt}avl{\isacharunderscore}{\kern0pt}equiv\isanewline
\ \ \ \ \ \ \isacommand{by}\isamarkupfalse%
\ {\isacharparenleft}{\kern0pt}auto\ simp{\isacharcolon}{\kern0pt}\ return{\isacharunderscore}{\kern0pt}iff\ wf{\isacharunderscore}{\kern0pt}problem{\isacharunderscore}{\kern0pt}def{\isacharparenright}{\kern0pt}\isanewline
\ \ \isacommand{qed}\isamarkupfalse%
%
\endisatagproof
{\isafoldproof}%
%
\isadelimproof
%
\endisadelimproof
%
\begin{isamarkuptext}%
Refinement lemma for our plan checking algorithm%
\end{isamarkuptext}\isamarkuptrue%
\ \ \isacommand{lemma}\isamarkupfalse%
\ {\isacharparenleft}{\kern0pt}\isakeyword{in}\ wf{\isacharunderscore}{\kern0pt}ast{\isacharunderscore}{\kern0pt}problem{\isacharparenright}{\kern0pt}\ valid{\isacharunderscore}{\kern0pt}plan{\isacharunderscore}{\kern0pt}from{\isadigit{2}}E{\isacharunderscore}{\kern0pt}return{\isacharunderscore}{\kern0pt}iff{\isacharcolon}{\kern0pt}\isanewline
\ \ \ \ \isakeyword{assumes}\ {\isachardoublequoteopen}wf{\isacharunderscore}{\kern0pt}world{\isacharunderscore}{\kern0pt}model\ s{\isachardoublequoteclose}\isanewline
\ \ \ \ \isakeyword{shows}\ {\isachardoublequoteopen}valid{\isacharunderscore}{\kern0pt}plan{\isacharunderscore}{\kern0pt}from{\isadigit{2}}E\ STG\ mp{\isacharunderscore}{\kern0pt}objT\ mp{\isacharunderscore}{\kern0pt}Fevl\ s\ {\isasympi}s\ {\isacharequal}{\kern0pt}\ Inr\ {\isacharparenleft}{\kern0pt}{\isacharparenright}{\kern0pt}\ {\isasymlongleftrightarrow}\ valid{\isacharunderscore}{\kern0pt}plan{\isacharunderscore}{\kern0pt}from{\isadigit{2}}\ s\ {\isasympi}s{\isachardoublequoteclose}\isanewline
%
\isadelimproof
\ \ %
\endisadelimproof
%
\isatagproof
\isacommand{proof}\isamarkupfalse%
\ {\isacharminus}{\kern0pt}\isanewline
\ \ \ \ \isacommand{have}\isamarkupfalse%
\ {\isachardoublequoteopen}wf{\isacharunderscore}{\kern0pt}domain{\isachardoublequoteclose}\isanewline
\ \ \ \ \ \ \isacommand{using}\isamarkupfalse%
\ wf{\isacharunderscore}{\kern0pt}problem\ \isacommand{unfolding}\isamarkupfalse%
\ wf{\isacharunderscore}{\kern0pt}problem{\isacharunderscore}{\kern0pt}def\ \isacommand{by}\isamarkupfalse%
\ simp\isanewline
\ \ \ \ \isacommand{have}\isamarkupfalse%
\ {\isachardoublequoteopen}wm{\isacharunderscore}{\kern0pt}basic\ s{\isachardoublequoteclose}\isanewline
\ \ \ \ \ \ \isacommand{using}\isamarkupfalse%
\ {\isacartoucheopen}wf{\isacharunderscore}{\kern0pt}world{\isacharunderscore}{\kern0pt}model\ s{\isacartoucheclose}\ wf{\isacharunderscore}{\kern0pt}fmla{\isacharunderscore}{\kern0pt}atom{\isacharunderscore}{\kern0pt}alt\ wf{\isacharunderscore}{\kern0pt}func{\isacharunderscore}{\kern0pt}assign{\isacharunderscore}{\kern0pt}is{\isacharunderscore}{\kern0pt}eqAtom\ \isanewline
\ \ \ \ \ \ \isacommand{unfolding}\isamarkupfalse%
\ wf{\isacharunderscore}{\kern0pt}world{\isacharunderscore}{\kern0pt}model{\isacharunderscore}{\kern0pt}def\ wm{\isacharunderscore}{\kern0pt}basic{\isacharunderscore}{\kern0pt}def\ \isacommand{by}\isamarkupfalse%
\ auto\isanewline
\ \ \ \ \isacommand{have}\isamarkupfalse%
\ {\isachardoublequoteopen}strict{\isacharunderscore}{\kern0pt}sorted\ {\isacharparenleft}{\kern0pt}htps{\isacharunderscore}{\kern0pt}exec\ {\isasympi}s{\isacharparenright}{\kern0pt}{\isachardoublequoteclose}\isanewline
\ \ \ \ \ \ \isacommand{using}\isamarkupfalse%
\ htps{\isacharunderscore}{\kern0pt}exec{\isacharunderscore}{\kern0pt}distinct\ htps{\isacharunderscore}{\kern0pt}exec{\isacharunderscore}{\kern0pt}sorted\ \isacommand{by}\isamarkupfalse%
\ {\isacharparenleft}{\kern0pt}simp\ add{\isacharcolon}{\kern0pt}\ strict{\isacharunderscore}{\kern0pt}sorted{\isacharunderscore}{\kern0pt}iff{\isacharparenright}{\kern0pt}\isanewline
\ \ \ \ \isacommand{show}\isamarkupfalse%
\ {\isacharquery}{\kern0pt}thesis\isanewline
\ \ \ \ \ \ \isacommand{using}\isamarkupfalse%
\ {\isacartoucheopen}wf{\isacharunderscore}{\kern0pt}domain{\isacartoucheclose}\ {\isacartoucheopen}strict{\isacharunderscore}{\kern0pt}sorted\ {\isacharparenleft}{\kern0pt}htps{\isacharunderscore}{\kern0pt}exec\ {\isasympi}s{\isacharparenright}{\kern0pt}{\isacartoucheclose}\ valid{\isacharunderscore}{\kern0pt}plan{\isacharunderscore}{\kern0pt}from{\isadigit{2}}{\isacharunderscore}{\kern0pt}exec{\isacharunderscore}{\kern0pt}refine\isanewline
\ \ \ \ \ \ \ \ \ \ \ \ valid{\isacharunderscore}{\kern0pt}happ{\isacharunderscore}{\kern0pt}seq{\isacharunderscore}{\kern0pt}fromE{\isacharunderscore}{\kern0pt}return{\isacharunderscore}{\kern0pt}iff{\isacharbrackleft}{\kern0pt}OF\ {\isacartoucheopen}wm{\isacharunderscore}{\kern0pt}basic\ s{\isacartoucheclose}\ simplify{\isacharunderscore}{\kern0pt}plan{\isacharunderscore}{\kern0pt}wf{\isacharunderscore}{\kern0pt}ground{\isacharunderscore}{\kern0pt}action{\isacharbrackright}{\kern0pt}\ \isanewline
\ \ \ \ \ \ \ \ \ \ \ \ simplify{\isacharunderscore}{\kern0pt}planE{\isacharunderscore}{\kern0pt}return{\isacharunderscore}{\kern0pt}iff{\isacharbrackleft}{\kern0pt}OF\ {\isacartoucheopen}wf{\isacharunderscore}{\kern0pt}domain{\isacartoucheclose}\ {\isacartoucheopen}strict{\isacharunderscore}{\kern0pt}sorted\ {\isacharparenleft}{\kern0pt}htps{\isacharunderscore}{\kern0pt}exec\ {\isasympi}s{\isacharparenright}{\kern0pt}{\isacartoucheclose}{\isacharbrackright}{\kern0pt}\isanewline
\ \ \ \ \ \ \ \ \ \ \ \ simplify{\isacharunderscore}{\kern0pt}planE{\isacharunderscore}{\kern0pt}avl{\isacharunderscore}{\kern0pt}equiv\isanewline
\ \ \ \ \ \ \isacommand{apply}\isamarkupfalse%
\ {\isacharparenleft}{\kern0pt}simp\ add{\isacharcolon}{\kern0pt}\ valid{\isacharunderscore}{\kern0pt}plan{\isacharunderscore}{\kern0pt}fromE{\isacharunderscore}{\kern0pt}wf{\isacharunderscore}{\kern0pt}plan{\isacharparenright}{\kern0pt}\ \isanewline
\ \ \ \ \ \ \isacommand{apply}\isamarkupfalse%
\ {\isacharparenleft}{\kern0pt}auto\ simp{\isacharcolon}{\kern0pt}\ return{\isacharunderscore}{\kern0pt}iff{\isacharparenright}{\kern0pt}\isanewline
\ \ \ \ \ \ \isacommand{apply}\isamarkupfalse%
\ metis{\isacharplus}{\kern0pt}\isanewline
\ \ \ \ \ \ \isacommand{done}\isamarkupfalse%
\isanewline
\ \ \isacommand{qed}\isamarkupfalse%
%
\endisatagproof
{\isafoldproof}%
%
\isadelimproof
\isanewline
%
\endisadelimproof
\isanewline
\ \ \isanewline
\ \ \isacommand{lemma}\isamarkupfalse%
\ {\isacharparenleft}{\kern0pt}\isakeyword{in}\ wf{\isacharunderscore}{\kern0pt}ast{\isacharunderscore}{\kern0pt}problem{\isacharparenright}{\kern0pt}\isanewline
\ \ \ \ \isakeyword{assumes}\ {\isachardoublequoteopen}wf{\isacharunderscore}{\kern0pt}world{\isacharunderscore}{\kern0pt}model\ s{\isachardoublequoteclose}\isanewline
\ \ \ \ \isakeyword{shows}\ {\isachardoublequoteopen}valid{\isacharunderscore}{\kern0pt}plan{\isacharunderscore}{\kern0pt}from{\isadigit{2}}E\ STG\ mp{\isacharunderscore}{\kern0pt}objT\ mp{\isacharunderscore}{\kern0pt}Fevl\ s\ {\isasympi}s\ {\isacharequal}{\kern0pt}\ Inr\ {\isacharparenleft}{\kern0pt}{\isacharparenright}{\kern0pt}\ {\isasymlongleftrightarrow}\ valid{\isacharunderscore}{\kern0pt}plan{\isacharunderscore}{\kern0pt}from{\isadigit{2}}\ s\ {\isasympi}s{\isachardoublequoteclose}\isanewline
%
\isadelimproof
\ \ \ \ %
\endisadelimproof
%
\isatagproof
\isacommand{using}\isamarkupfalse%
\ assms\ \isacommand{by}\isamarkupfalse%
\ {\isacharparenleft}{\kern0pt}rule\ valid{\isacharunderscore}{\kern0pt}plan{\isacharunderscore}{\kern0pt}from{\isadigit{2}}E{\isacharunderscore}{\kern0pt}return{\isacharunderscore}{\kern0pt}iff{\isacharparenright}{\kern0pt}%
\endisatagproof
{\isafoldproof}%
%
\isadelimproof
\isanewline
%
\endisadelimproof
\isanewline
\ \ \isacommand{lemmas}\isamarkupfalse%
\ valid{\isacharunderscore}{\kern0pt}plan{\isacharunderscore}{\kern0pt}from{\isadigit{2}}E{\isacharunderscore}{\kern0pt}return{\isacharunderscore}{\kern0pt}iff{\isacharprime}{\kern0pt}{\isacharbrackleft}{\kern0pt}return{\isacharunderscore}{\kern0pt}iff{\isacharbrackright}{\kern0pt}\isanewline
\ \ \ \ {\isacharequal}{\kern0pt}\ wf{\isacharunderscore}{\kern0pt}ast{\isacharunderscore}{\kern0pt}problem{\isachardot}{\kern0pt}valid{\isacharunderscore}{\kern0pt}plan{\isacharunderscore}{\kern0pt}from{\isadigit{2}}E{\isacharunderscore}{\kern0pt}return{\isacharunderscore}{\kern0pt}iff{\isacharbrackleft}{\kern0pt}of\ P{\isacharcomma}{\kern0pt}\ OF\ wf{\isacharunderscore}{\kern0pt}ast{\isacharunderscore}{\kern0pt}problem{\isachardot}{\kern0pt}intro{\isacharbrackright}{\kern0pt}\isanewline
\isanewline
\ \ \isanewline
\isanewline
\ \ \isanewline
\ \ \isanewline
\isanewline
\isacommand{end}\isamarkupfalse%
\ %
\isamarkupcmt{Context of \isa{ast{\isacharunderscore}{\kern0pt}problem}%
}%
\isadelimdocument
%
\endisadelimdocument
%
\isatagdocument
%
\isamarkupsubsection{Executable Plan Checker%
}
\isamarkuptrue%
%
\endisatagdocument
{\isafolddocument}%
%
\isadelimdocument
%
\endisadelimdocument
%
\begin{isamarkuptext}%
We obtain the main plan checker by combining the well-formedness check
  and executability check.%
\end{isamarkuptext}\isamarkuptrue%
\isacommand{definition}\isamarkupfalse%
\ {\isachardoublequoteopen}check{\isacharunderscore}{\kern0pt}all{\isacharunderscore}{\kern0pt}list\ P\ l\ msg\ msgf\ {\isasymequiv}\isanewline
\ \ forallM\ {\isacharparenleft}{\kern0pt}{\isasymlambda}x{\isachardot}{\kern0pt}\ check\ {\isacharparenleft}{\kern0pt}P\ x{\isacharparenright}{\kern0pt}\ {\isacharparenleft}{\kern0pt}{\isasymlambda}{\isacharunderscore}{\kern0pt}{\isacharcolon}{\kern0pt}{\isacharcolon}{\kern0pt}unit{\isachardot}{\kern0pt}\ shows\ msg\ o\ shows\ {\isacharprime}{\kern0pt}{\isacharprime}{\kern0pt}{\isacharcolon}{\kern0pt}\ {\isacharprime}{\kern0pt}{\isacharprime}{\kern0pt}\ o\ msgf\ x{\isacharparenright}{\kern0pt}\ {\isacharparenright}{\kern0pt}\ l\ {\isacharless}{\kern0pt}{\isacharplus}{\kern0pt}{\isacharquery}{\kern0pt}\ snd{\isachardoublequoteclose}\isanewline
\isanewline
\isacommand{lemma}\isamarkupfalse%
\ check{\isacharunderscore}{\kern0pt}all{\isacharunderscore}{\kern0pt}list{\isacharunderscore}{\kern0pt}return{\isacharunderscore}{\kern0pt}iff{\isacharbrackleft}{\kern0pt}return{\isacharunderscore}{\kern0pt}iff{\isacharbrackright}{\kern0pt}{\isacharcolon}{\kern0pt}\ {\isachardoublequoteopen}check{\isacharunderscore}{\kern0pt}all{\isacharunderscore}{\kern0pt}list\ P\ l\ msg\ msgf\ {\isacharequal}{\kern0pt}\ Inr\ {\isacharparenleft}{\kern0pt}{\isacharparenright}{\kern0pt}\ {\isasymlongleftrightarrow}\ {\isacharparenleft}{\kern0pt}{\isasymforall}x{\isasymin}set\ l{\isachardot}{\kern0pt}\ P\ x{\isacharparenright}{\kern0pt}{\isachardoublequoteclose}\isanewline
%
\isadelimproof
\ \ %
\endisadelimproof
%
\isatagproof
\isacommand{unfolding}\isamarkupfalse%
\ check{\isacharunderscore}{\kern0pt}all{\isacharunderscore}{\kern0pt}list{\isacharunderscore}{\kern0pt}def\ \isacommand{by}\isamarkupfalse%
\ {\isacharparenleft}{\kern0pt}induction\ l{\isacharparenright}{\kern0pt}\ {\isacharparenleft}{\kern0pt}auto{\isacharparenright}{\kern0pt}%
\endisatagproof
{\isafoldproof}%
%
\isadelimproof
\isanewline
%
\endisadelimproof
\isanewline
\isacommand{definition}\isamarkupfalse%
\ {\isachardoublequoteopen}check{\isacharunderscore}{\kern0pt}wf{\isacharunderscore}{\kern0pt}types\ D\ {\isasymequiv}\ do\ {\isacharbraceleft}{\kern0pt}\isanewline
\ \ check{\isacharunderscore}{\kern0pt}all{\isacharunderscore}{\kern0pt}list\ {\isacharparenleft}{\kern0pt}{\isasymlambda}{\isacharparenleft}{\kern0pt}{\isacharunderscore}{\kern0pt}{\isacharcomma}{\kern0pt}t{\isacharparenright}{\kern0pt}{\isachardot}{\kern0pt}\ t{\isacharequal}{\kern0pt}{\isacharprime}{\kern0pt}{\isacharprime}{\kern0pt}object{\isacharprime}{\kern0pt}{\isacharprime}{\kern0pt}\ {\isasymor}\ t{\isasymin}fst{\isacharbackquote}{\kern0pt}set\ {\isacharparenleft}{\kern0pt}types\ D{\isacharparenright}{\kern0pt}{\isacharparenright}{\kern0pt}\ {\isacharparenleft}{\kern0pt}types\ D{\isacharparenright}{\kern0pt}\ {\isacharprime}{\kern0pt}{\isacharprime}{\kern0pt}Undeclared\ supertype{\isacharprime}{\kern0pt}{\isacharprime}{\kern0pt}\ {\isacharparenleft}{\kern0pt}shows\ o\ snd{\isacharparenright}{\kern0pt}\isanewline
{\isacharbraceright}{\kern0pt}{\isachardoublequoteclose}\isanewline
\isanewline
\isacommand{lemma}\isamarkupfalse%
\ check{\isacharunderscore}{\kern0pt}wf{\isacharunderscore}{\kern0pt}types{\isacharunderscore}{\kern0pt}return{\isacharunderscore}{\kern0pt}iff{\isacharbrackleft}{\kern0pt}return{\isacharunderscore}{\kern0pt}iff{\isacharbrackright}{\kern0pt}{\isacharcolon}{\kern0pt}\ {\isachardoublequoteopen}check{\isacharunderscore}{\kern0pt}wf{\isacharunderscore}{\kern0pt}types\ D\ {\isacharequal}{\kern0pt}\ Inr\ {\isacharparenleft}{\kern0pt}{\isacharparenright}{\kern0pt}\ {\isasymlongleftrightarrow}\ ast{\isacharunderscore}{\kern0pt}domain{\isachardot}{\kern0pt}wf{\isacharunderscore}{\kern0pt}types\ D{\isachardoublequoteclose}\isanewline
%
\isadelimproof
\ \ %
\endisadelimproof
%
\isatagproof
\isacommand{unfolding}\isamarkupfalse%
\ ast{\isacharunderscore}{\kern0pt}domain{\isachardot}{\kern0pt}wf{\isacharunderscore}{\kern0pt}types{\isacharunderscore}{\kern0pt}def\ check{\isacharunderscore}{\kern0pt}wf{\isacharunderscore}{\kern0pt}types{\isacharunderscore}{\kern0pt}def\ \isacommand{by}\isamarkupfalse%
\ {\isacharparenleft}{\kern0pt}force\ simp{\isacharcolon}{\kern0pt}\ return{\isacharunderscore}{\kern0pt}iff{\isacharparenright}{\kern0pt}%
\endisatagproof
{\isafoldproof}%
%
\isadelimproof
\isanewline
%
\endisadelimproof
\isanewline
\isacommand{definition}\isamarkupfalse%
\ {\isachardoublequoteopen}check{\isacharunderscore}{\kern0pt}wf{\isacharunderscore}{\kern0pt}domain\ D\ stg\ conT\ {\isasymequiv}\ do\ {\isacharbraceleft}{\kern0pt}\isanewline
\ \ check{\isacharunderscore}{\kern0pt}wf{\isacharunderscore}{\kern0pt}types\ D{\isacharsemicolon}{\kern0pt}\isanewline
\ \ check\ {\isacharparenleft}{\kern0pt}distinct\ {\isacharparenleft}{\kern0pt}map\ {\isacharparenleft}{\kern0pt}predicate{\isacharunderscore}{\kern0pt}decl{\isachardot}{\kern0pt}pred{\isacharparenright}{\kern0pt}\ {\isacharparenleft}{\kern0pt}predicates\ D{\isacharparenright}{\kern0pt}{\isacharparenright}{\kern0pt}{\isacharparenright}{\kern0pt}\ {\isacharparenleft}{\kern0pt}ERRS\ {\isacharprime}{\kern0pt}{\isacharprime}{\kern0pt}Duplicate\ predicate\ declaration{\isacharprime}{\kern0pt}{\isacharprime}{\kern0pt}{\isacharparenright}{\kern0pt}{\isacharsemicolon}{\kern0pt}\isanewline
\ \ check{\isacharunderscore}{\kern0pt}all{\isacharunderscore}{\kern0pt}list\ {\isacharparenleft}{\kern0pt}ast{\isacharunderscore}{\kern0pt}domain{\isachardot}{\kern0pt}wf{\isacharunderscore}{\kern0pt}predicate{\isacharunderscore}{\kern0pt}decl\ D{\isacharparenright}{\kern0pt}\ {\isacharparenleft}{\kern0pt}predicates\ D{\isacharparenright}{\kern0pt}\ {\isacharprime}{\kern0pt}{\isacharprime}{\kern0pt}Malformed\ predicate\ declaration{\isacharprime}{\kern0pt}{\isacharprime}{\kern0pt}\ {\isacharparenleft}{\kern0pt}shows\ o\ predicate{\isachardot}{\kern0pt}name\ o\ predicate{\isacharunderscore}{\kern0pt}decl{\isachardot}{\kern0pt}pred{\isacharparenright}{\kern0pt}{\isacharsemicolon}{\kern0pt}\isanewline
\ \ check\ {\isacharparenleft}{\kern0pt}distinct\ {\isacharparenleft}{\kern0pt}map\ {\isacharparenleft}{\kern0pt}function{\isacharunderscore}{\kern0pt}decl{\isachardot}{\kern0pt}func{\isacharparenright}{\kern0pt}\ {\isacharparenleft}{\kern0pt}functions\ D{\isacharparenright}{\kern0pt}{\isacharparenright}{\kern0pt}{\isacharparenright}{\kern0pt}\ {\isacharparenleft}{\kern0pt}ERRS\ {\isacharprime}{\kern0pt}{\isacharprime}{\kern0pt}Duplicate\ function\ declaration{\isacharprime}{\kern0pt}{\isacharprime}{\kern0pt}{\isacharparenright}{\kern0pt}{\isacharsemicolon}{\kern0pt}\isanewline
\ \ check{\isacharunderscore}{\kern0pt}all{\isacharunderscore}{\kern0pt}list\ {\isacharparenleft}{\kern0pt}ast{\isacharunderscore}{\kern0pt}domain{\isachardot}{\kern0pt}wf{\isacharunderscore}{\kern0pt}function{\isacharunderscore}{\kern0pt}decl\ D{\isacharparenright}{\kern0pt}\ {\isacharparenleft}{\kern0pt}functions\ D{\isacharparenright}{\kern0pt}\ {\isacharprime}{\kern0pt}{\isacharprime}{\kern0pt}Malformed\ function\ declaration{\isacharprime}{\kern0pt}{\isacharprime}{\kern0pt}\ {\isacharparenleft}{\kern0pt}shows\ o\ func{\isachardot}{\kern0pt}name\ o\ function{\isacharunderscore}{\kern0pt}decl{\isachardot}{\kern0pt}func{\isacharparenright}{\kern0pt}{\isacharsemicolon}{\kern0pt}\isanewline
\ \ check\ {\isacharparenleft}{\kern0pt}distinct\ {\isacharparenleft}{\kern0pt}map\ fst\ {\isacharparenleft}{\kern0pt}consts\ D{\isacharparenright}{\kern0pt}{\isacharparenright}{\kern0pt}{\isacharparenright}{\kern0pt}\ {\isacharparenleft}{\kern0pt}ERRS\ {\isacharprime}{\kern0pt}{\isacharprime}{\kern0pt}Duplicate\ constant\ declaration{\isacharprime}{\kern0pt}{\isacharprime}{\kern0pt}{\isacharparenright}{\kern0pt}{\isacharsemicolon}{\kern0pt}\isanewline
\ \ check\ {\isacharparenleft}{\kern0pt}{\isasymforall}{\isacharparenleft}{\kern0pt}n{\isacharcomma}{\kern0pt}T{\isacharparenright}{\kern0pt}{\isasymin}set\ {\isacharparenleft}{\kern0pt}consts\ D{\isacharparenright}{\kern0pt}{\isachardot}{\kern0pt}\ ast{\isacharunderscore}{\kern0pt}domain{\isachardot}{\kern0pt}wf{\isacharunderscore}{\kern0pt}type\ D\ T{\isacharparenright}{\kern0pt}\ {\isacharparenleft}{\kern0pt}ERRS\ {\isacharprime}{\kern0pt}{\isacharprime}{\kern0pt}Malformed\ type{\isacharprime}{\kern0pt}{\isacharprime}{\kern0pt}{\isacharparenright}{\kern0pt}{\isacharsemicolon}{\kern0pt}\isanewline
\ \ check\ {\isacharparenleft}{\kern0pt}distinct\ {\isacharparenleft}{\kern0pt}map\ ast{\isacharunderscore}{\kern0pt}action{\isacharunderscore}{\kern0pt}schema{\isachardot}{\kern0pt}name\ {\isacharparenleft}{\kern0pt}actions\ D{\isacharparenright}{\kern0pt}{\isacharparenright}{\kern0pt}\ \ {\isacharparenright}{\kern0pt}\ {\isacharparenleft}{\kern0pt}ERRS\ {\isacharprime}{\kern0pt}{\isacharprime}{\kern0pt}Duplicate\ action\ name{\isacharprime}{\kern0pt}{\isacharprime}{\kern0pt}{\isacharparenright}{\kern0pt}{\isacharsemicolon}{\kern0pt}\isanewline
\ \ check{\isacharunderscore}{\kern0pt}all{\isacharunderscore}{\kern0pt}list\ {\isacharparenleft}{\kern0pt}ast{\isacharunderscore}{\kern0pt}domain{\isachardot}{\kern0pt}wf{\isacharunderscore}{\kern0pt}action{\isacharunderscore}{\kern0pt}schema{\isacharprime}{\kern0pt}\ D\ stg\ conT{\isacharparenright}{\kern0pt}\ {\isacharparenleft}{\kern0pt}actions\ D{\isacharparenright}{\kern0pt}\ {\isacharprime}{\kern0pt}{\isacharprime}{\kern0pt}Malformed\ action{\isacharprime}{\kern0pt}{\isacharprime}{\kern0pt}\ {\isacharparenleft}{\kern0pt}shows\ o\ ast{\isacharunderscore}{\kern0pt}action{\isacharunderscore}{\kern0pt}schema{\isachardot}{\kern0pt}name{\isacharparenright}{\kern0pt}\isanewline
{\isacharbraceright}{\kern0pt}{\isachardoublequoteclose}\isanewline
\isanewline
\isacommand{lemma}\isamarkupfalse%
\ check{\isacharunderscore}{\kern0pt}wf{\isacharunderscore}{\kern0pt}domain{\isacharunderscore}{\kern0pt}return{\isacharunderscore}{\kern0pt}iff{\isacharbrackleft}{\kern0pt}return{\isacharunderscore}{\kern0pt}iff{\isacharbrackright}{\kern0pt}{\isacharcolon}{\kern0pt}\isanewline
\ \ {\isachardoublequoteopen}check{\isacharunderscore}{\kern0pt}wf{\isacharunderscore}{\kern0pt}domain\ D\ stg\ conT\ {\isacharequal}{\kern0pt}\ Inr\ {\isacharparenleft}{\kern0pt}{\isacharparenright}{\kern0pt}\ {\isasymlongleftrightarrow}\ ast{\isacharunderscore}{\kern0pt}domain{\isachardot}{\kern0pt}wf{\isacharunderscore}{\kern0pt}domain{\isacharprime}{\kern0pt}\ D\ stg\ conT{\isachardoublequoteclose}\isanewline
%
\isadelimproof
%
\endisadelimproof
%
\isatagproof
\isacommand{proof}\isamarkupfalse%
\ {\isacharminus}{\kern0pt}\isanewline
\ \ \isacommand{interpret}\isamarkupfalse%
\ ast{\isacharunderscore}{\kern0pt}domain\ D\ \isacommand{{\isachardot}{\kern0pt}}\isamarkupfalse%
\isanewline
\ \ \isacommand{show}\isamarkupfalse%
\ {\isacharquery}{\kern0pt}thesis\isanewline
\ \ \ \ \isacommand{unfolding}\isamarkupfalse%
\ check{\isacharunderscore}{\kern0pt}wf{\isacharunderscore}{\kern0pt}domain{\isacharunderscore}{\kern0pt}def\ wf{\isacharunderscore}{\kern0pt}domain{\isacharprime}{\kern0pt}{\isacharunderscore}{\kern0pt}def\ \isacommand{by}\isamarkupfalse%
\ {\isacharparenleft}{\kern0pt}auto\ simp{\isacharcolon}{\kern0pt}\ return{\isacharunderscore}{\kern0pt}iff{\isacharparenright}{\kern0pt}\isanewline
\isacommand{qed}\isamarkupfalse%
%
\endisatagproof
{\isafoldproof}%
%
\isadelimproof
\isanewline
%
\endisadelimproof
\isanewline
\isacommand{definition}\isamarkupfalse%
\ {\isachardoublequoteopen}prepend{\isacharunderscore}{\kern0pt}err{\isacharunderscore}{\kern0pt}msg\ msg\ e\ {\isasymequiv}\ {\isasymlambda}{\isacharunderscore}{\kern0pt}{\isacharcolon}{\kern0pt}{\isacharcolon}{\kern0pt}unit{\isachardot}{\kern0pt}\ shows\ msg\ o\ shows\ {\isacharprime}{\kern0pt}{\isacharprime}{\kern0pt}{\isacharcolon}{\kern0pt}\ {\isacharprime}{\kern0pt}{\isacharprime}{\kern0pt}\ o\ e\ {\isacharparenleft}{\kern0pt}{\isacharparenright}{\kern0pt}{\isachardoublequoteclose}\isanewline
\isanewline
\isacommand{definition}\isamarkupfalse%
\ {\isachardoublequoteopen}check{\isacharunderscore}{\kern0pt}wf{\isacharunderscore}{\kern0pt}problem\ P\ stg\ conT\ mp\ {\isasymequiv}\ do\ {\isacharbraceleft}{\kern0pt}\isanewline
\ \ let\ D\ {\isacharequal}{\kern0pt}\ ast{\isacharunderscore}{\kern0pt}problem{\isachardot}{\kern0pt}domain\ P{\isacharsemicolon}{\kern0pt}\isanewline
\ \ check{\isacharunderscore}{\kern0pt}wf{\isacharunderscore}{\kern0pt}domain\ D\ stg\ conT\ {\isacharless}{\kern0pt}{\isacharplus}{\kern0pt}{\isacharquery}{\kern0pt}\ prepend{\isacharunderscore}{\kern0pt}err{\isacharunderscore}{\kern0pt}msg\ {\isacharprime}{\kern0pt}{\isacharprime}{\kern0pt}Domain\ not\ well{\isacharminus}{\kern0pt}formed{\isacharprime}{\kern0pt}{\isacharprime}{\kern0pt}{\isacharsemicolon}{\kern0pt}\isanewline
\ \ check\ {\isacharparenleft}{\kern0pt}distinct\ {\isacharparenleft}{\kern0pt}map\ fst\ {\isacharparenleft}{\kern0pt}objects\ P{\isacharparenright}{\kern0pt}\ {\isacharat}{\kern0pt}\ map\ fst\ {\isacharparenleft}{\kern0pt}consts\ D{\isacharparenright}{\kern0pt}{\isacharparenright}{\kern0pt}{\isacharparenright}{\kern0pt}\ {\isacharparenleft}{\kern0pt}ERRS\ {\isacharprime}{\kern0pt}{\isacharprime}{\kern0pt}Duplicate\ object\ declaration{\isacharprime}{\kern0pt}{\isacharprime}{\kern0pt}{\isacharparenright}{\kern0pt}{\isacharsemicolon}{\kern0pt}\isanewline
\ \ check\ {\isacharparenleft}{\kern0pt}{\isacharparenleft}{\kern0pt}{\isasymforall}{\isacharparenleft}{\kern0pt}n{\isacharcomma}{\kern0pt}T{\isacharparenright}{\kern0pt}{\isasymin}set\ {\isacharparenleft}{\kern0pt}objects\ P{\isacharparenright}{\kern0pt}{\isachardot}{\kern0pt}\ ast{\isacharunderscore}{\kern0pt}domain{\isachardot}{\kern0pt}wf{\isacharunderscore}{\kern0pt}type\ D\ T{\isacharparenright}{\kern0pt}{\isacharparenright}{\kern0pt}\ {\isacharparenleft}{\kern0pt}ERRS\ {\isacharprime}{\kern0pt}{\isacharprime}{\kern0pt}Malformed\ type{\isacharprime}{\kern0pt}{\isacharprime}{\kern0pt}{\isacharparenright}{\kern0pt}{\isacharsemicolon}{\kern0pt}\isanewline
\ \ check\ {\isacharparenleft}{\kern0pt}distinct\ {\isacharparenleft}{\kern0pt}init\ P{\isacharparenright}{\kern0pt}{\isacharparenright}{\kern0pt}\ {\isacharparenleft}{\kern0pt}ERRS\ {\isacharprime}{\kern0pt}{\isacharprime}{\kern0pt}Duplicate\ fact\ in\ initial\ state{\isacharprime}{\kern0pt}{\isacharprime}{\kern0pt}{\isacharparenright}{\kern0pt}{\isacharsemicolon}{\kern0pt}\isanewline
\ \ check\ {\isacharparenleft}{\kern0pt}{\isasymforall}f{\isasymin}set\ {\isacharparenleft}{\kern0pt}init\ P{\isacharparenright}{\kern0pt}{\isachardot}{\kern0pt}\ ast{\isacharunderscore}{\kern0pt}problem{\isachardot}{\kern0pt}wf{\isacharunderscore}{\kern0pt}fmla{\isacharunderscore}{\kern0pt}atom{\isadigit{2}}{\isacharprime}{\kern0pt}\ P\ mp\ stg\ f\ {\isasymor}\ ast{\isacharunderscore}{\kern0pt}problem{\isachardot}{\kern0pt}wf{\isacharunderscore}{\kern0pt}func{\isacharunderscore}{\kern0pt}assign{\isacharprime}{\kern0pt}\ P\ mp\ stg\ f{\isacharparenright}{\kern0pt}\ {\isacharparenleft}{\kern0pt}ERRS\ {\isacharprime}{\kern0pt}{\isacharprime}{\kern0pt}Malformed\ formula\ in\ initial\ state{\isacharprime}{\kern0pt}{\isacharprime}{\kern0pt}{\isacharparenright}{\kern0pt}{\isacharsemicolon}{\kern0pt}\isanewline
\ \ check\ {\isacharparenleft}{\kern0pt}ast{\isacharunderscore}{\kern0pt}domain{\isachardot}{\kern0pt}wf{\isacharunderscore}{\kern0pt}fmla{\isacharprime}{\kern0pt}\ D\ {\isacharparenleft}{\kern0pt}Mapping{\isachardot}{\kern0pt}lookup\ mp{\isacharparenright}{\kern0pt}\ stg\ {\isacharparenleft}{\kern0pt}goal\ P{\isacharparenright}{\kern0pt}{\isacharparenright}{\kern0pt}\ {\isacharparenleft}{\kern0pt}ERRS\ {\isacharprime}{\kern0pt}{\isacharprime}{\kern0pt}Malformed\ goal\ formula{\isacharprime}{\kern0pt}{\isacharprime}{\kern0pt}{\isacharparenright}{\kern0pt}\isanewline
{\isacharbraceright}{\kern0pt}{\isachardoublequoteclose}\isanewline
\isanewline
\isacommand{lemma}\isamarkupfalse%
\ check{\isacharunderscore}{\kern0pt}wf{\isacharunderscore}{\kern0pt}problem{\isacharunderscore}{\kern0pt}return{\isacharunderscore}{\kern0pt}iff{\isacharbrackleft}{\kern0pt}return{\isacharunderscore}{\kern0pt}iff{\isacharbrackright}{\kern0pt}{\isacharcolon}{\kern0pt}\isanewline
\ \ {\isachardoublequoteopen}check{\isacharunderscore}{\kern0pt}wf{\isacharunderscore}{\kern0pt}problem\ P\ stg\ conT\ mp\ {\isacharequal}{\kern0pt}\ Inr\ {\isacharparenleft}{\kern0pt}{\isacharparenright}{\kern0pt}\ {\isasymlongleftrightarrow}\ ast{\isacharunderscore}{\kern0pt}problem{\isachardot}{\kern0pt}wf{\isacharunderscore}{\kern0pt}problem{\isacharprime}{\kern0pt}\ P\ stg\ conT\ mp{\isachardoublequoteclose}\isanewline
%
\isadelimproof
%
\endisadelimproof
%
\isatagproof
\isacommand{proof}\isamarkupfalse%
\ {\isacharminus}{\kern0pt}\isanewline
\ \ \isacommand{interpret}\isamarkupfalse%
\ ast{\isacharunderscore}{\kern0pt}problem\ P\ \isacommand{{\isachardot}{\kern0pt}}\isamarkupfalse%
\isanewline
\ \ \isacommand{show}\isamarkupfalse%
\ {\isacharquery}{\kern0pt}thesis\isanewline
\ \ \ \ \isacommand{unfolding}\isamarkupfalse%
\ check{\isacharunderscore}{\kern0pt}wf{\isacharunderscore}{\kern0pt}problem{\isacharunderscore}{\kern0pt}def\ wf{\isacharunderscore}{\kern0pt}problem{\isacharprime}{\kern0pt}{\isacharunderscore}{\kern0pt}def\ \isacommand{by}\isamarkupfalse%
\ {\isacharparenleft}{\kern0pt}auto\ simp{\isacharcolon}{\kern0pt}\ return{\isacharunderscore}{\kern0pt}iff{\isacharparenright}{\kern0pt}\isanewline
\isacommand{qed}\isamarkupfalse%
%
\endisatagproof
{\isafoldproof}%
%
\isadelimproof
\isanewline
%
\endisadelimproof
\isanewline
\isacommand{definition}\isamarkupfalse%
\ {\isachardoublequoteopen}check{\isacharunderscore}{\kern0pt}plan\ P\ {\isasympi}s\ {\isasymequiv}\ do\ {\isacharbraceleft}{\kern0pt}\isanewline
\ \ let\ stg\ {\isacharequal}{\kern0pt}\ ast{\isacharunderscore}{\kern0pt}domain{\isachardot}{\kern0pt}STG\ {\isacharparenleft}{\kern0pt}ast{\isacharunderscore}{\kern0pt}problem{\isachardot}{\kern0pt}domain\ P{\isacharparenright}{\kern0pt}{\isacharsemicolon}{\kern0pt}\isanewline
\ \ let\ conT\ {\isacharequal}{\kern0pt}\ ast{\isacharunderscore}{\kern0pt}domain{\isachardot}{\kern0pt}mp{\isacharunderscore}{\kern0pt}constT\ {\isacharparenleft}{\kern0pt}ast{\isacharunderscore}{\kern0pt}problem{\isachardot}{\kern0pt}domain\ P{\isacharparenright}{\kern0pt}{\isacharsemicolon}{\kern0pt}\isanewline
\ \ let\ mp\ {\isacharequal}{\kern0pt}\ ast{\isacharunderscore}{\kern0pt}problem{\isachardot}{\kern0pt}mp{\isacharunderscore}{\kern0pt}objT\ P{\isacharsemicolon}{\kern0pt}\isanewline
\ \ let\ mp{\isacharunderscore}{\kern0pt}fe\ {\isacharequal}{\kern0pt}\ ast{\isacharunderscore}{\kern0pt}problem{\isachardot}{\kern0pt}mp{\isacharunderscore}{\kern0pt}Fevl\ P{\isacharsemicolon}{\kern0pt}\isanewline
\ \ check{\isacharunderscore}{\kern0pt}wf{\isacharunderscore}{\kern0pt}problem\ P\ stg\ conT\ mp{\isacharsemicolon}{\kern0pt}\isanewline
\ \ ast{\isacharunderscore}{\kern0pt}problem{\isachardot}{\kern0pt}valid{\isacharunderscore}{\kern0pt}plan{\isacharunderscore}{\kern0pt}from{\isadigit{2}}E\ P\ stg\ mp\ mp{\isacharunderscore}{\kern0pt}fe\ {\isacharparenleft}{\kern0pt}ast{\isacharunderscore}{\kern0pt}problem{\isachardot}{\kern0pt}I\ P{\isacharparenright}{\kern0pt}\ {\isasympi}s\isanewline
{\isacharbraceright}{\kern0pt}\ {\isacharless}{\kern0pt}{\isacharplus}{\kern0pt}{\isacharquery}{\kern0pt}\ {\isacharparenleft}{\kern0pt}{\isasymlambda}e{\isachardot}{\kern0pt}\ String{\isachardot}{\kern0pt}implode\ {\isacharparenleft}{\kern0pt}e\ {\isacharparenleft}{\kern0pt}{\isacharparenright}{\kern0pt}\ {\isacharprime}{\kern0pt}{\isacharprime}{\kern0pt}{\isacharprime}{\kern0pt}{\isacharprime}{\kern0pt}{\isacharparenright}{\kern0pt}{\isacharparenright}{\kern0pt}{\isachardoublequoteclose}%
\begin{isamarkuptext}%
Correctness theorem of the plan checker: It returns \isa{Inr\ {\isacharparenleft}{\kern0pt}{\isacharparenright}{\kern0pt}}
  if and only if the problem is well-formed and the plan is valid.%
\end{isamarkuptext}\isamarkuptrue%
\isacommand{theorem}\isamarkupfalse%
\ check{\isacharunderscore}{\kern0pt}plan{\isacharunderscore}{\kern0pt}return{\isacharunderscore}{\kern0pt}iff{\isacharbrackleft}{\kern0pt}return{\isacharunderscore}{\kern0pt}iff{\isacharbrackright}{\kern0pt}{\isacharcolon}{\kern0pt}\ \isanewline
\ \ {\isachardoublequoteopen}check{\isacharunderscore}{\kern0pt}plan\ P\ {\isasympi}s\ {\isacharequal}{\kern0pt}\ Inr\ {\isacharparenleft}{\kern0pt}{\isacharparenright}{\kern0pt}\ {\isasymlongleftrightarrow}\ ast{\isacharunderscore}{\kern0pt}problem{\isachardot}{\kern0pt}wf{\isacharunderscore}{\kern0pt}problem\ P\ {\isasymand}\ ast{\isacharunderscore}{\kern0pt}problem{\isachardot}{\kern0pt}valid{\isacharunderscore}{\kern0pt}plan{\isadigit{2}}\ P\ {\isasympi}s{\isachardoublequoteclose}\isanewline
%
\isadelimproof
%
\endisadelimproof
%
\isatagproof
\isacommand{proof}\isamarkupfalse%
\ {\isacharminus}{\kern0pt}\isanewline
\ \ \isacommand{interpret}\isamarkupfalse%
\ ast{\isacharunderscore}{\kern0pt}problem\ P\ \isacommand{{\isachardot}{\kern0pt}}\isamarkupfalse%
\isanewline
\ \ \isacommand{show}\isamarkupfalse%
\ {\isacharquery}{\kern0pt}thesis\isanewline
\ \ \ \ \isacommand{unfolding}\isamarkupfalse%
\ check{\isacharunderscore}{\kern0pt}plan{\isacharunderscore}{\kern0pt}def\ ast{\isacharunderscore}{\kern0pt}problem{\isachardot}{\kern0pt}valid{\isacharunderscore}{\kern0pt}plan{\isadigit{2}}{\isacharunderscore}{\kern0pt}def\isanewline
\ \ \ \ \isacommand{using}\isamarkupfalse%
\ valid{\isacharunderscore}{\kern0pt}plan{\isacharunderscore}{\kern0pt}from{\isadigit{2}}E{\isacharunderscore}{\kern0pt}return{\isacharunderscore}{\kern0pt}iff{\isacharprime}{\kern0pt}\ wf{\isacharunderscore}{\kern0pt}ast{\isacharunderscore}{\kern0pt}problem{\isacharunderscore}{\kern0pt}def\ wf{\isacharunderscore}{\kern0pt}ast{\isacharunderscore}{\kern0pt}problem{\isachardot}{\kern0pt}wf{\isacharunderscore}{\kern0pt}I\isanewline
\ \ \ \ \isacommand{by}\isamarkupfalse%
\ {\isacharparenleft}{\kern0pt}fastforce\ simp{\isacharcolon}{\kern0pt}\ wf{\isacharunderscore}{\kern0pt}problem{\isacharprime}{\kern0pt}{\isacharunderscore}{\kern0pt}correct\ return{\isacharunderscore}{\kern0pt}iff{\isacharparenright}{\kern0pt}\isanewline
\isacommand{qed}\isamarkupfalse%
%
\endisatagproof
{\isafoldproof}%
%
\isadelimproof
\isanewline
%
\endisadelimproof
\isanewline
\isacommand{find{\isacharunderscore}{\kern0pt}theorems}\isamarkupfalse%
\ ast{\isacharunderscore}{\kern0pt}problem{\isachardot}{\kern0pt}valid{\isacharunderscore}{\kern0pt}plan{\isacharunderscore}{\kern0pt}from{\isadigit{2}}%
\isadelimdocument
%
\endisadelimdocument
%
\isatagdocument
%
\isamarkupsubsection{Code Setup%
}
\isamarkuptrue%
%
\endisatagdocument
{\isafolddocument}%
%
\isadelimdocument
%
\endisadelimdocument
%
\begin{isamarkuptext}%
In this section, we set up the code generator to generate verified
  code for our plan checker.%
\end{isamarkuptext}\isamarkuptrue%
%
\isadelimdocument
%
\endisadelimdocument
%
\isatagdocument
%
\isamarkupsubsubsection{Code Equations%
}
\isamarkuptrue%
%
\endisatagdocument
{\isafolddocument}%
%
\isadelimdocument
%
\endisadelimdocument
%
\begin{isamarkuptext}%
We first register the code equations for the functions of the checker.
  Note that we not necessarily register the original code equations, but also
  optimized ones.%
\end{isamarkuptext}\isamarkuptrue%
\isacommand{lemmas}\isamarkupfalse%
\ wf{\isacharunderscore}{\kern0pt}domain{\isacharunderscore}{\kern0pt}code\ {\isacharequal}{\kern0pt}\isanewline
\ \ ast{\isacharunderscore}{\kern0pt}domain{\isachardot}{\kern0pt}sig{\isacharunderscore}{\kern0pt}def\isanewline
\ \ ast{\isacharunderscore}{\kern0pt}domain{\isachardot}{\kern0pt}wf{\isacharunderscore}{\kern0pt}types{\isacharunderscore}{\kern0pt}def\isanewline
\ \ ast{\isacharunderscore}{\kern0pt}domain{\isachardot}{\kern0pt}wf{\isacharunderscore}{\kern0pt}type{\isachardot}{\kern0pt}simps\isanewline
\ \ ast{\isacharunderscore}{\kern0pt}domain{\isachardot}{\kern0pt}wf{\isacharunderscore}{\kern0pt}predicate{\isacharunderscore}{\kern0pt}decl{\isachardot}{\kern0pt}simps\isanewline
\ \ ast{\isacharunderscore}{\kern0pt}domain{\isachardot}{\kern0pt}STG{\isacharunderscore}{\kern0pt}def\isanewline
\ \ ast{\isacharunderscore}{\kern0pt}domain{\isachardot}{\kern0pt}is{\isacharunderscore}{\kern0pt}of{\isacharunderscore}{\kern0pt}type{\isacharprime}{\kern0pt}{\isacharunderscore}{\kern0pt}def\isanewline
\ \ ast{\isacharunderscore}{\kern0pt}domain{\isachardot}{\kern0pt}wf{\isacharunderscore}{\kern0pt}atom{\isacharprime}{\kern0pt}{\isachardot}{\kern0pt}simps\isanewline
\ \ ast{\isacharunderscore}{\kern0pt}domain{\isachardot}{\kern0pt}wf{\isacharunderscore}{\kern0pt}pred{\isacharunderscore}{\kern0pt}atom{\isacharprime}{\kern0pt}{\isachardot}{\kern0pt}simps\isanewline
\ \ ast{\isacharunderscore}{\kern0pt}domain{\isachardot}{\kern0pt}wf{\isacharunderscore}{\kern0pt}fmla{\isacharprime}{\kern0pt}{\isachardot}{\kern0pt}simps\isanewline
\ \ ast{\isacharunderscore}{\kern0pt}domain{\isachardot}{\kern0pt}wf{\isacharunderscore}{\kern0pt}fmla{\isacharunderscore}{\kern0pt}atom{\isadigit{1}}{\isacharprime}{\kern0pt}{\isachardot}{\kern0pt}simps\isanewline
\ \ ast{\isacharunderscore}{\kern0pt}domain{\isachardot}{\kern0pt}wf{\isacharunderscore}{\kern0pt}effect{\isacharprime}{\kern0pt}{\isachardot}{\kern0pt}simps\isanewline
\ \ ast{\isacharunderscore}{\kern0pt}domain{\isachardot}{\kern0pt}wf{\isacharunderscore}{\kern0pt}action{\isacharunderscore}{\kern0pt}schema{\isacharprime}{\kern0pt}{\isachardot}{\kern0pt}simps\isanewline
\ \ ast{\isacharunderscore}{\kern0pt}domain{\isachardot}{\kern0pt}wf{\isacharunderscore}{\kern0pt}domain{\isacharprime}{\kern0pt}{\isacharunderscore}{\kern0pt}def\isanewline
\ \ ast{\isacharunderscore}{\kern0pt}domain{\isachardot}{\kern0pt}subst{\isacharunderscore}{\kern0pt}term{\isachardot}{\kern0pt}simps\isanewline
\ \ ast{\isacharunderscore}{\kern0pt}domain{\isachardot}{\kern0pt}mp{\isacharunderscore}{\kern0pt}constT{\isacharunderscore}{\kern0pt}def\isanewline
\ \ \isanewline
\ \ ast{\isacharunderscore}{\kern0pt}domain{\isachardot}{\kern0pt}wf{\isacharunderscore}{\kern0pt}function{\isacharunderscore}{\kern0pt}decl{\isachardot}{\kern0pt}simps\isanewline
\ \ ast{\isacharunderscore}{\kern0pt}domain{\isachardot}{\kern0pt}func{\isacharunderscore}{\kern0pt}sig{\isacharunderscore}{\kern0pt}def\isanewline
\ \ ast{\isacharunderscore}{\kern0pt}domain{\isachardot}{\kern0pt}wf{\isacharunderscore}{\kern0pt}func{\isacharunderscore}{\kern0pt}args{\isacharprime}{\kern0pt}{\isachardot}{\kern0pt}simps\isanewline
\ \ ast{\isacharunderscore}{\kern0pt}domain{\isachardot}{\kern0pt}wf{\isacharunderscore}{\kern0pt}duration{\isacharunderscore}{\kern0pt}const{\isacharprime}{\kern0pt}{\isachardot}{\kern0pt}simps\isanewline
\ \ ast{\isacharunderscore}{\kern0pt}domain{\isachardot}{\kern0pt}acts{\isacharunderscore}{\kern0pt}non{\isacharunderscore}{\kern0pt}intrf{\isacharunderscore}{\kern0pt}def\isanewline
\ \ \isanewline
\isanewline
\isacommand{declare}\isamarkupfalse%
\ wf{\isacharunderscore}{\kern0pt}domain{\isacharunderscore}{\kern0pt}code{\isacharbrackleft}{\kern0pt}code{\isacharbrackright}{\kern0pt}\isanewline
\isanewline
\isacommand{lemmas}\isamarkupfalse%
\ wf{\isacharunderscore}{\kern0pt}problem{\isacharunderscore}{\kern0pt}code\ {\isacharequal}{\kern0pt}\isanewline
\ \ ast{\isacharunderscore}{\kern0pt}problem{\isachardot}{\kern0pt}wf{\isacharunderscore}{\kern0pt}problem{\isacharprime}{\kern0pt}{\isacharunderscore}{\kern0pt}def\isanewline
\ \ ast{\isacharunderscore}{\kern0pt}problem{\isachardot}{\kern0pt}wf{\isacharunderscore}{\kern0pt}fact{\isacharprime}{\kern0pt}{\isacharunderscore}{\kern0pt}def\isanewline
\ \ ast{\isacharunderscore}{\kern0pt}problem{\isachardot}{\kern0pt}is{\isacharunderscore}{\kern0pt}obj{\isacharunderscore}{\kern0pt}of{\isacharunderscore}{\kern0pt}type{\isacharunderscore}{\kern0pt}alt\isanewline
\ \ ast{\isacharunderscore}{\kern0pt}problem{\isachardot}{\kern0pt}wf{\isacharunderscore}{\kern0pt}fact{\isacharunderscore}{\kern0pt}def\isanewline
\ \ ast{\isacharunderscore}{\kern0pt}problem{\isachardot}{\kern0pt}wf{\isacharunderscore}{\kern0pt}plan{\isacharunderscore}{\kern0pt}action{\isachardot}{\kern0pt}simps\isanewline
\ \ ast{\isacharunderscore}{\kern0pt}domain{\isachardot}{\kern0pt}subtype{\isacharunderscore}{\kern0pt}edge{\isachardot}{\kern0pt}simps\isanewline
\isanewline
\isacommand{declare}\isamarkupfalse%
\ wf{\isacharunderscore}{\kern0pt}problem{\isacharunderscore}{\kern0pt}code{\isacharbrackleft}{\kern0pt}code{\isacharbrackright}{\kern0pt}\isanewline
\isanewline
\isacommand{lemmas}\isamarkupfalse%
\ check{\isacharunderscore}{\kern0pt}code\ {\isacharequal}{\kern0pt}\isanewline
\ \ ast{\isacharunderscore}{\kern0pt}problem{\isachardot}{\kern0pt}valid{\isacharunderscore}{\kern0pt}plan{\isadigit{2}}{\isacharunderscore}{\kern0pt}def\isanewline
\ \ ast{\isacharunderscore}{\kern0pt}problem{\isachardot}{\kern0pt}valid{\isacharunderscore}{\kern0pt}plan{\isacharunderscore}{\kern0pt}from{\isadigit{2}}E{\isachardot}{\kern0pt}simps\isanewline
\ \ ast{\isacharunderscore}{\kern0pt}problem{\isachardot}{\kern0pt}res{\isacharunderscore}{\kern0pt}inst{\isachardot}{\kern0pt}simps\isanewline
\ \ ast{\isacharunderscore}{\kern0pt}domain{\isachardot}{\kern0pt}resolve{\isacharunderscore}{\kern0pt}action{\isacharunderscore}{\kern0pt}schema{\isacharunderscore}{\kern0pt}def\isanewline
\ \ ast{\isacharunderscore}{\kern0pt}domain{\isachardot}{\kern0pt}resolve{\isacharunderscore}{\kern0pt}action{\isacharunderscore}{\kern0pt}schemaE{\isacharunderscore}{\kern0pt}def\isanewline
\ \ ast{\isacharunderscore}{\kern0pt}problem{\isachardot}{\kern0pt}I{\isacharunderscore}{\kern0pt}def\isanewline
\ \ ast{\isacharunderscore}{\kern0pt}domain{\isachardot}{\kern0pt}instantiate{\isacharunderscore}{\kern0pt}action{\isacharunderscore}{\kern0pt}schema{\isachardot}{\kern0pt}simps\isanewline
\ \ ast{\isacharunderscore}{\kern0pt}problem{\isachardot}{\kern0pt}holds{\isacharunderscore}{\kern0pt}def\isanewline
\ \ ast{\isacharunderscore}{\kern0pt}problem{\isachardot}{\kern0pt}mp{\isacharunderscore}{\kern0pt}objT{\isacharunderscore}{\kern0pt}def\isanewline
\ \ ast{\isacharunderscore}{\kern0pt}problem{\isachardot}{\kern0pt}is{\isacharunderscore}{\kern0pt}obj{\isacharunderscore}{\kern0pt}of{\isacharunderscore}{\kern0pt}type{\isacharunderscore}{\kern0pt}impl{\isacharunderscore}{\kern0pt}def\isanewline
\ \ ast{\isacharunderscore}{\kern0pt}problem{\isachardot}{\kern0pt}wf{\isacharunderscore}{\kern0pt}fmla{\isacharunderscore}{\kern0pt}atom{\isadigit{2}}{\isacharprime}{\kern0pt}{\isacharunderscore}{\kern0pt}def\isanewline
\ \ valuation{\isacharunderscore}{\kern0pt}def\isanewline
\ \ \isanewline
\ \ ast{\isacharunderscore}{\kern0pt}problem{\isachardot}{\kern0pt}valid{\isacharunderscore}{\kern0pt}happ{\isacharunderscore}{\kern0pt}seq{\isacharunderscore}{\kern0pt}fromE{\isachardot}{\kern0pt}simps\ \isanewline
\ \ ast{\isacharunderscore}{\kern0pt}problem{\isachardot}{\kern0pt}htps{\isacharunderscore}{\kern0pt}exec{\isachardot}{\kern0pt}simps\isanewline
\ \ ast{\isacharunderscore}{\kern0pt}problem{\isachardot}{\kern0pt}insort{\isacharunderscore}{\kern0pt}htp{\isachardot}{\kern0pt}simps\isanewline
\ \ ast{\isacharunderscore}{\kern0pt}problem{\isachardot}{\kern0pt}simplify{\isacharunderscore}{\kern0pt}planE{\isachardot}{\kern0pt}simps\isanewline
\ \ ast{\isacharunderscore}{\kern0pt}problem{\isachardot}{\kern0pt}simplify{\isacharunderscore}{\kern0pt}actionE{\isachardot}{\kern0pt}simps\isanewline
\ \ ast{\isacharunderscore}{\kern0pt}problem{\isachardot}{\kern0pt}insort{\isacharunderscore}{\kern0pt}happ{\isachardot}{\kern0pt}simps\isanewline
\ \ ast{\isacharunderscore}{\kern0pt}domain{\isachardot}{\kern0pt}inst{\isacharunderscore}{\kern0pt}snap{\isacharunderscore}{\kern0pt}action{\isachardot}{\kern0pt}simps\isanewline
\ \ ast{\isacharunderscore}{\kern0pt}domain{\isachardot}{\kern0pt}filter{\isacharunderscore}{\kern0pt}time{\isacharunderscore}{\kern0pt}spec{\isacharunderscore}{\kern0pt}def\isanewline
\ \ ast{\isacharunderscore}{\kern0pt}domain{\isachardot}{\kern0pt}conjunct{\isacharunderscore}{\kern0pt}effects{\isacharunderscore}{\kern0pt}def\isanewline
\ \ \isanewline
\ \ ast{\isacharunderscore}{\kern0pt}problem{\isachardot}{\kern0pt}apply{\isacharunderscore}{\kern0pt}happ{\isacharunderscore}{\kern0pt}exec{\isachardot}{\kern0pt}simps\isanewline
\ \ ast{\isacharunderscore}{\kern0pt}problem{\isachardot}{\kern0pt}en{\isacharunderscore}{\kern0pt}exE{\isacharunderscore}{\kern0pt}def\isanewline
\ \ ast{\isacharunderscore}{\kern0pt}problem{\isachardot}{\kern0pt}insort{\isacharunderscore}{\kern0pt}mult{\isacharunderscore}{\kern0pt}happs{\isachardot}{\kern0pt}simps\isanewline
\ \ ast{\isacharunderscore}{\kern0pt}problem{\isachardot}{\kern0pt}duration{\isacharunderscore}{\kern0pt}matches{\isacharprime}{\kern0pt}{\isachardot}{\kern0pt}simps\isanewline
\ \ ast{\isacharunderscore}{\kern0pt}problem{\isachardot}{\kern0pt}wf{\isacharunderscore}{\kern0pt}func{\isacharunderscore}{\kern0pt}assign{\isacharprime}{\kern0pt}{\isachardot}{\kern0pt}simps\isanewline
\ \ \isanewline
\ \ ast{\isacharunderscore}{\kern0pt}problem{\isachardot}{\kern0pt}mp{\isacharunderscore}{\kern0pt}Fevl{\isacharunderscore}{\kern0pt}def\isanewline
\ \ ast{\isacharunderscore}{\kern0pt}problem{\isachardot}{\kern0pt}place{\isacharunderscore}{\kern0pt}inv{\isacharunderscore}{\kern0pt}snap{\isacharunderscore}{\kern0pt}acts{\isachardot}{\kern0pt}simps\isanewline
\ \ \isanewline
\ \ ast{\isacharunderscore}{\kern0pt}problem{\isachardot}{\kern0pt}insort{\isacharunderscore}{\kern0pt}timed{\isacharunderscore}{\kern0pt}ground{\isacharunderscore}{\kern0pt}acts{\isachardot}{\kern0pt}simps\isanewline
\ \ ast{\isacharunderscore}{\kern0pt}problem{\isachardot}{\kern0pt}insert{\isacharunderscore}{\kern0pt}happ{\isacharunderscore}{\kern0pt}to{\isacharunderscore}{\kern0pt}tree{\isachardot}{\kern0pt}simps\isanewline
\ \ ast{\isacharunderscore}{\kern0pt}problem{\isachardot}{\kern0pt}insert{\isacharunderscore}{\kern0pt}timed{\isacharunderscore}{\kern0pt}ground{\isacharunderscore}{\kern0pt}acts{\isacharunderscore}{\kern0pt}to{\isacharunderscore}{\kern0pt}tree{\isachardot}{\kern0pt}simps\isanewline
\ \ ast{\isacharunderscore}{\kern0pt}problem{\isachardot}{\kern0pt}avl{\isacharunderscore}{\kern0pt}inorder{\isacharunderscore}{\kern0pt}def\isanewline
\ \ ast{\isacharunderscore}{\kern0pt}problem{\isachardot}{\kern0pt}simplify{\isacharunderscore}{\kern0pt}planE{\isacharunderscore}{\kern0pt}avl{\isachardot}{\kern0pt}simps\isanewline
\ \ ast{\isacharunderscore}{\kern0pt}problem{\isachardot}{\kern0pt}simplify{\isacharunderscore}{\kern0pt}planE{\isacharunderscore}{\kern0pt}avl{\isacharprime}{\kern0pt}{\isachardot}{\kern0pt}simps\isanewline
\ \ ast{\isacharunderscore}{\kern0pt}problem{\isachardot}{\kern0pt}avl{\isacharunderscore}{\kern0pt}height{\isachardot}{\kern0pt}simps\isanewline
\ \ ast{\isacharunderscore}{\kern0pt}problem{\isachardot}{\kern0pt}avl{\isacharunderscore}{\kern0pt}node{\isacharunderscore}{\kern0pt}def\isanewline
\ \ ast{\isacharunderscore}{\kern0pt}problem{\isachardot}{\kern0pt}avl{\isacharunderscore}{\kern0pt}balR{\isacharunderscore}{\kern0pt}def\isanewline
\ \ ast{\isacharunderscore}{\kern0pt}problem{\isachardot}{\kern0pt}avl{\isacharunderscore}{\kern0pt}balL{\isacharunderscore}{\kern0pt}def\isanewline
\ \ \isanewline
\isanewline
\isacommand{declare}\isamarkupfalse%
\ check{\isacharunderscore}{\kern0pt}code{\isacharbrackleft}{\kern0pt}code{\isacharbrackright}{\kern0pt}%
\isadelimdocument
%
\endisadelimdocument
%
\isatagdocument
%
\isamarkupsubsubsection{Setup for Containers Framework%
}
\isamarkuptrue%
%
\endisatagdocument
{\isafolddocument}%
%
\isadelimdocument
%
\endisadelimdocument
\isacommand{derive}\isamarkupfalse%
\ {\isacharparenleft}{\kern0pt}linorder{\isacharparenright}{\kern0pt}\ compare\ rat\isanewline
\isanewline
\isacommand{derive}\isamarkupfalse%
\ {\isacharparenleft}{\kern0pt}eq{\isacharparenright}{\kern0pt}\ ceq\ predicate\ func\ atom\ object\ formula\ \isanewline
\isacommand{derive}\isamarkupfalse%
\ ccompare\ predicate\ func\ atom\ object\ formula\isanewline
\isacommand{derive}\isamarkupfalse%
\ {\isacharparenleft}{\kern0pt}no{\isacharparenright}{\kern0pt}\ cenum\ atom\ object\ formula\isanewline
\isacommand{derive}\isamarkupfalse%
\ {\isacharparenleft}{\kern0pt}rbt{\isacharparenright}{\kern0pt}\ set{\isacharunderscore}{\kern0pt}impl\ func\ atom\ object\ formula\isanewline
\isanewline
\isacommand{derive}\isamarkupfalse%
\ {\isacharparenleft}{\kern0pt}rbt{\isacharparenright}{\kern0pt}\ mapping{\isacharunderscore}{\kern0pt}impl\ object\isanewline
\isanewline
\isacommand{derive}\isamarkupfalse%
\ linorder\ predicate\ func\ object\ atom\ {\isachardoublequoteopen}object\ atom\ formula{\isachardoublequoteclose}%
\isadelimdocument
%
\endisadelimdocument
%
\isatagdocument
%
\isamarkupsubsubsection{More Efficient Distinctness Check for Linorders%
}
\isamarkuptrue%
%
\endisatagdocument
{\isafolddocument}%
%
\isadelimdocument
%
\endisadelimdocument
\isacommand{fun}\isamarkupfalse%
\ no{\isacharunderscore}{\kern0pt}stutter\ {\isacharcolon}{\kern0pt}{\isacharcolon}{\kern0pt}\ {\isachardoublequoteopen}{\isacharprime}{\kern0pt}a\ list\ {\isasymRightarrow}\ bool{\isachardoublequoteclose}\ \isakeyword{where}\isanewline
\ \ {\isachardoublequoteopen}no{\isacharunderscore}{\kern0pt}stutter\ {\isacharbrackleft}{\kern0pt}{\isacharbrackright}{\kern0pt}\ {\isacharequal}{\kern0pt}\ True{\isachardoublequoteclose}\isanewline
{\isacharbar}{\kern0pt}\ {\isachardoublequoteopen}no{\isacharunderscore}{\kern0pt}stutter\ {\isacharbrackleft}{\kern0pt}{\isacharunderscore}{\kern0pt}{\isacharbrackright}{\kern0pt}\ {\isacharequal}{\kern0pt}\ True{\isachardoublequoteclose}\isanewline
{\isacharbar}{\kern0pt}\ {\isachardoublequoteopen}no{\isacharunderscore}{\kern0pt}stutter\ {\isacharparenleft}{\kern0pt}a{\isacharhash}{\kern0pt}b{\isacharhash}{\kern0pt}l{\isacharparenright}{\kern0pt}\ {\isacharequal}{\kern0pt}\ {\isacharparenleft}{\kern0pt}a{\isasymnoteq}b\ {\isasymand}\ no{\isacharunderscore}{\kern0pt}stutter\ {\isacharparenleft}{\kern0pt}b{\isacharhash}{\kern0pt}l{\isacharparenright}{\kern0pt}{\isacharparenright}{\kern0pt}{\isachardoublequoteclose}\isanewline
\isanewline
\isacommand{lemma}\isamarkupfalse%
\ sorted{\isacharunderscore}{\kern0pt}no{\isacharunderscore}{\kern0pt}stutter{\isacharunderscore}{\kern0pt}eq{\isacharunderscore}{\kern0pt}distinct{\isacharcolon}{\kern0pt}\ {\isachardoublequoteopen}sorted\ l\ {\isasymLongrightarrow}\ no{\isacharunderscore}{\kern0pt}stutter\ l\ {\isasymlongleftrightarrow}\ distinct\ l{\isachardoublequoteclose}\isanewline
%
\isadelimproof
\ \ %
\endisadelimproof
%
\isatagproof
\isacommand{apply}\isamarkupfalse%
\ {\isacharparenleft}{\kern0pt}induction\ l\ rule{\isacharcolon}{\kern0pt}\ no{\isacharunderscore}{\kern0pt}stutter{\isachardot}{\kern0pt}induct{\isacharparenright}{\kern0pt}\isanewline
\ \ \isacommand{apply}\isamarkupfalse%
\ {\isacharparenleft}{\kern0pt}auto\ simp{\isacharcolon}{\kern0pt}\ {\isacharparenright}{\kern0pt}\isanewline
\ \ \isacommand{done}\isamarkupfalse%
%
\endisatagproof
{\isafoldproof}%
%
\isadelimproof
\isanewline
%
\endisadelimproof
\isanewline
\isacommand{definition}\isamarkupfalse%
\ distinct{\isacharunderscore}{\kern0pt}ds\ {\isacharcolon}{\kern0pt}{\isacharcolon}{\kern0pt}\ {\isachardoublequoteopen}{\isacharprime}{\kern0pt}a{\isacharcolon}{\kern0pt}{\isacharcolon}{\kern0pt}linorder\ list\ {\isasymRightarrow}\ bool{\isachardoublequoteclose}\isanewline
\ \ \isakeyword{where}\ {\isachardoublequoteopen}distinct{\isacharunderscore}{\kern0pt}ds\ l\ {\isasymequiv}\ no{\isacharunderscore}{\kern0pt}stutter\ {\isacharparenleft}{\kern0pt}quicksort\ l{\isacharparenright}{\kern0pt}{\isachardoublequoteclose}\isanewline
\isanewline
\isacommand{lemma}\isamarkupfalse%
\ {\isacharbrackleft}{\kern0pt}code{\isacharunderscore}{\kern0pt}unfold{\isacharbrackright}{\kern0pt}{\isacharcolon}{\kern0pt}\ {\isachardoublequoteopen}distinct\ {\isacharequal}{\kern0pt}\ distinct{\isacharunderscore}{\kern0pt}ds{\isachardoublequoteclose}\isanewline
%
\isadelimproof
\ \ %
\endisadelimproof
%
\isatagproof
\isacommand{apply}\isamarkupfalse%
\ {\isacharparenleft}{\kern0pt}intro\ ext{\isacharparenright}{\kern0pt}\isanewline
\ \ \isacommand{unfolding}\isamarkupfalse%
\ distinct{\isacharunderscore}{\kern0pt}ds{\isacharunderscore}{\kern0pt}def\isanewline
\ \ \isacommand{apply}\isamarkupfalse%
\ {\isacharparenleft}{\kern0pt}auto\ simp{\isacharcolon}{\kern0pt}\ sorted{\isacharunderscore}{\kern0pt}no{\isacharunderscore}{\kern0pt}stutter{\isacharunderscore}{\kern0pt}eq{\isacharunderscore}{\kern0pt}distinct{\isacharparenright}{\kern0pt}\isanewline
\ \ \isacommand{done}\isamarkupfalse%
%
\endisatagproof
{\isafoldproof}%
%
\isadelimproof
%
\endisadelimproof
%
\isadelimdocument
%
\endisadelimdocument
%
\isatagdocument
%
\isamarkupsubsubsection{Parsing Rational Numbers%
}
\isamarkuptrue%
%
\endisatagdocument
{\isafolddocument}%
%
\isadelimdocument
%
\endisadelimdocument
\isacommand{type{\isacharunderscore}{\kern0pt}synonym}\isamarkupfalse%
\ digit\ {\isacharequal}{\kern0pt}\ nat%
\begin{isamarkuptext}%
Well-formedness conditions for a digit and a sequence of digits.%
\end{isamarkuptext}\isamarkuptrue%
\isacommand{definition}\isamarkupfalse%
\ {\isachardoublequoteopen}wf{\isacharunderscore}{\kern0pt}digit\ d\ {\isasymlongleftrightarrow}\ d\ {\isacharless}{\kern0pt}\ {\isadigit{1}}{\isadigit{0}}{\isachardoublequoteclose}%
\begin{isamarkuptext}%
A sequence of digits is well-formed iff it is normalized and only contains well-formed digits.%
\end{isamarkuptext}\isamarkuptrue%
\isacommand{definition}\isamarkupfalse%
\ {\isachardoublequoteopen}wf{\isacharunderscore}{\kern0pt}digits\ ds\ {\isasymlongleftrightarrow}\ {\isacharparenleft}{\kern0pt}{\isasymforall}d\ {\isasymin}\ set\ ds{\isachardot}{\kern0pt}\ wf{\isacharunderscore}{\kern0pt}digit\ d{\isacharparenright}{\kern0pt}{\isachardoublequoteclose}%
\begin{isamarkuptext}%
Functions to trim leading and trailing zeros.%
\end{isamarkuptext}\isamarkuptrue%
\isacommand{fun}\isamarkupfalse%
\ trim{\isacharunderscore}{\kern0pt}ld{\isacharunderscore}{\kern0pt}zs\ {\isacharcolon}{\kern0pt}{\isacharcolon}{\kern0pt}\ {\isachardoublequoteopen}digit\ list\ {\isasymRightarrow}\ digit\ list{\isachardoublequoteclose}\ \isakeyword{where}\isanewline
\ \ {\isachardoublequoteopen}trim{\isacharunderscore}{\kern0pt}ld{\isacharunderscore}{\kern0pt}zs\ {\isacharbrackleft}{\kern0pt}{\isacharbrackright}{\kern0pt}\ {\isacharequal}{\kern0pt}\ {\isacharbrackleft}{\kern0pt}{\isacharbrackright}{\kern0pt}{\isachardoublequoteclose}\isanewline
{\isacharbar}{\kern0pt}\ {\isachardoublequoteopen}trim{\isacharunderscore}{\kern0pt}ld{\isacharunderscore}{\kern0pt}zs\ {\isacharparenleft}{\kern0pt}d{\isacharhash}{\kern0pt}ds{\isacharparenright}{\kern0pt}\ {\isacharequal}{\kern0pt}\ \isanewline
\ \ {\isacharparenleft}{\kern0pt}if\ d\ {\isacharequal}{\kern0pt}\ {\isadigit{0}}\ then\ trim{\isacharunderscore}{\kern0pt}ld{\isacharunderscore}{\kern0pt}zs\ ds\ else\ d{\isacharhash}{\kern0pt}ds{\isacharparenright}{\kern0pt}{\isachardoublequoteclose}\isanewline
\isanewline
\isacommand{lemma}\isamarkupfalse%
\ tlz{\isacharunderscore}{\kern0pt}hd{\isacharunderscore}{\kern0pt}neq{\isacharunderscore}{\kern0pt}z{\isacharcolon}{\kern0pt}\ {\isachardoublequoteopen}trim{\isacharunderscore}{\kern0pt}ld{\isacharunderscore}{\kern0pt}zs\ ds\ {\isacharequal}{\kern0pt}\ ds{\isacharprime}{\kern0pt}\ {\isasymLongrightarrow}\ hd\ ds{\isacharprime}{\kern0pt}\ {\isasymnoteq}\ {\isadigit{0}}\ {\isasymor}\ ds{\isacharprime}{\kern0pt}\ {\isacharequal}{\kern0pt}\ {\isacharbrackleft}{\kern0pt}{\isacharbrackright}{\kern0pt}{\isachardoublequoteclose}\isanewline
%
\isadelimproof
\ \ %
\endisadelimproof
%
\isatagproof
\isacommand{by}\isamarkupfalse%
\ {\isacharparenleft}{\kern0pt}induction\ ds{\isacharparenright}{\kern0pt}\ auto%
\endisatagproof
{\isafoldproof}%
%
\isadelimproof
\isanewline
%
\endisadelimproof
\isanewline
\isacommand{lemma}\isamarkupfalse%
\ trim{\isacharunderscore}{\kern0pt}ld{\isacharunderscore}{\kern0pt}zs{\isacharunderscore}{\kern0pt}idem{\isacharcolon}{\kern0pt}\ {\isachardoublequoteopen}trim{\isacharunderscore}{\kern0pt}ld{\isacharunderscore}{\kern0pt}zs\ {\isacharparenleft}{\kern0pt}trim{\isacharunderscore}{\kern0pt}ld{\isacharunderscore}{\kern0pt}zs\ ds{\isacharparenright}{\kern0pt}\ {\isacharequal}{\kern0pt}\ trim{\isacharunderscore}{\kern0pt}ld{\isacharunderscore}{\kern0pt}zs\ ds{\isachardoublequoteclose}\isanewline
%
\isadelimproof
\ \ %
\endisadelimproof
%
\isatagproof
\isacommand{by}\isamarkupfalse%
\ {\isacharparenleft}{\kern0pt}induction\ ds{\isacharparenright}{\kern0pt}\ auto%
\endisatagproof
{\isafoldproof}%
%
\isadelimproof
\isanewline
%
\endisadelimproof
\isanewline
\isacommand{lemma}\isamarkupfalse%
\ trim{\isacharunderscore}{\kern0pt}ld{\isacharunderscore}{\kern0pt}zs{\isacharunderscore}{\kern0pt}wf{\isacharcolon}{\kern0pt}\ {\isachardoublequoteopen}wf{\isacharunderscore}{\kern0pt}digits\ ds\ {\isasymLongrightarrow}\ wf{\isacharunderscore}{\kern0pt}digits\ {\isacharparenleft}{\kern0pt}trim{\isacharunderscore}{\kern0pt}ld{\isacharunderscore}{\kern0pt}zs\ ds{\isacharparenright}{\kern0pt}{\isachardoublequoteclose}\isanewline
%
\isadelimproof
\ \ %
\endisadelimproof
%
\isatagproof
\isacommand{unfolding}\isamarkupfalse%
\ wf{\isacharunderscore}{\kern0pt}digits{\isacharunderscore}{\kern0pt}def\ \isacommand{by}\isamarkupfalse%
\ {\isacharparenleft}{\kern0pt}induction\ ds{\isacharparenright}{\kern0pt}\ auto%
\endisatagproof
{\isafoldproof}%
%
\isadelimproof
\isanewline
%
\endisadelimproof
\isanewline
\isacommand{lemma}\isamarkupfalse%
\ trim{\isacharunderscore}{\kern0pt}ld{\isacharunderscore}{\kern0pt}zs{\isacharunderscore}{\kern0pt}app{\isacharcolon}{\kern0pt}\ {\isachardoublequoteopen}trim{\isacharunderscore}{\kern0pt}ld{\isacharunderscore}{\kern0pt}zs\ {\isacharparenleft}{\kern0pt}ds{\isadigit{1}}\ {\isacharat}{\kern0pt}\ ds{\isadigit{2}}{\isacharparenright}{\kern0pt}\ {\isacharequal}{\kern0pt}\ {\isacharparenleft}{\kern0pt}trim{\isacharunderscore}{\kern0pt}ld{\isacharunderscore}{\kern0pt}zs\ ds{\isadigit{1}}{\isacharparenright}{\kern0pt}\ {\isacharat}{\kern0pt}\ ds{\isadigit{2}}\ {\isasymor}\ {\isacharparenleft}{\kern0pt}{\isasymforall}d\ {\isasymin}\ set\ ds{\isadigit{1}}{\isachardot}{\kern0pt}\ d\ {\isacharequal}{\kern0pt}\ {\isadigit{0}}{\isacharparenright}{\kern0pt}{\isachardoublequoteclose}\isanewline
%
\isadelimproof
\ \ %
\endisadelimproof
%
\isatagproof
\isacommand{by}\isamarkupfalse%
\ {\isacharparenleft}{\kern0pt}induction\ ds{\isadigit{1}}{\isacharparenright}{\kern0pt}\ auto%
\endisatagproof
{\isafoldproof}%
%
\isadelimproof
\isanewline
%
\endisadelimproof
\isanewline
\isacommand{fun}\isamarkupfalse%
\ trim{\isacharunderscore}{\kern0pt}tr{\isacharunderscore}{\kern0pt}zs{\isacharprime}{\kern0pt}\ {\isacharcolon}{\kern0pt}{\isacharcolon}{\kern0pt}\ {\isachardoublequoteopen}digit\ list\ {\isasymRightarrow}\ digit\ list\ {\isasymRightarrow}\ digit\ list{\isachardoublequoteclose}\ \isakeyword{where}\isanewline
\ \ {\isachardoublequoteopen}trim{\isacharunderscore}{\kern0pt}tr{\isacharunderscore}{\kern0pt}zs{\isacharprime}{\kern0pt}\ {\isacharbrackleft}{\kern0pt}{\isacharbrackright}{\kern0pt}\ zs\ {\isacharequal}{\kern0pt}\ {\isacharbrackleft}{\kern0pt}{\isacharbrackright}{\kern0pt}{\isachardoublequoteclose}\isanewline
{\isacharbar}{\kern0pt}\ {\isachardoublequoteopen}trim{\isacharunderscore}{\kern0pt}tr{\isacharunderscore}{\kern0pt}zs{\isacharprime}{\kern0pt}\ {\isacharparenleft}{\kern0pt}d{\isacharhash}{\kern0pt}ds{\isacharparenright}{\kern0pt}\ zs\ {\isacharequal}{\kern0pt}\ \isanewline
\ \ {\isacharparenleft}{\kern0pt}if\ d\ {\isacharequal}{\kern0pt}\ {\isadigit{0}}\ then\ trim{\isacharunderscore}{\kern0pt}tr{\isacharunderscore}{\kern0pt}zs{\isacharprime}{\kern0pt}\ ds\ {\isacharparenleft}{\kern0pt}{\isadigit{0}}{\isacharhash}{\kern0pt}zs{\isacharparenright}{\kern0pt}\isanewline
\ \ else\ zs\ {\isacharat}{\kern0pt}\ d\ {\isacharhash}{\kern0pt}\ {\isacharparenleft}{\kern0pt}trim{\isacharunderscore}{\kern0pt}tr{\isacharunderscore}{\kern0pt}zs{\isacharprime}{\kern0pt}\ ds\ {\isacharbrackleft}{\kern0pt}{\isacharbrackright}{\kern0pt}{\isacharparenright}{\kern0pt}{\isacharparenright}{\kern0pt}{\isachardoublequoteclose}\isanewline
\isanewline
\isacommand{fun}\isamarkupfalse%
\ trim{\isacharunderscore}{\kern0pt}tr{\isacharunderscore}{\kern0pt}zs\ {\isacharcolon}{\kern0pt}{\isacharcolon}{\kern0pt}\ {\isachardoublequoteopen}digit\ list\ {\isasymRightarrow}\ digit\ list{\isachardoublequoteclose}\ \isakeyword{where}\isanewline
\ \ {\isachardoublequoteopen}trim{\isacharunderscore}{\kern0pt}tr{\isacharunderscore}{\kern0pt}zs\ ds\ {\isacharequal}{\kern0pt}\ trim{\isacharunderscore}{\kern0pt}tr{\isacharunderscore}{\kern0pt}zs{\isacharprime}{\kern0pt}\ ds\ {\isacharbrackleft}{\kern0pt}{\isacharbrackright}{\kern0pt}{\isachardoublequoteclose}\isanewline
\isanewline
\isacommand{lemma}\isamarkupfalse%
\ ttz{\isacharunderscore}{\kern0pt}last{\isacharunderscore}{\kern0pt}neq{\isacharunderscore}{\kern0pt}z{\isacharunderscore}{\kern0pt}aux{\isacharcolon}{\kern0pt}\ \isanewline
\ \ \isakeyword{assumes}\ {\isachardoublequoteopen}trim{\isacharunderscore}{\kern0pt}tr{\isacharunderscore}{\kern0pt}zs{\isacharprime}{\kern0pt}\ ds\ zs\ {\isacharequal}{\kern0pt}\ ds{\isacharprime}{\kern0pt}{\isachardoublequoteclose}\ \isanewline
\ \ \isakeyword{shows}\ {\isachardoublequoteopen}last\ ds{\isacharprime}{\kern0pt}\ {\isasymnoteq}\ {\isadigit{0}}\ {\isasymor}\ ds{\isacharprime}{\kern0pt}\ {\isacharequal}{\kern0pt}\ {\isacharbrackleft}{\kern0pt}{\isacharbrackright}{\kern0pt}{\isachardoublequoteclose}\isanewline
%
\isadelimproof
\ \ %
\endisadelimproof
%
\isatagproof
\isacommand{using}\isamarkupfalse%
\ assms\isanewline
\isacommand{proof}\isamarkupfalse%
\ {\isacharparenleft}{\kern0pt}induction\ ds\ arbitrary{\isacharcolon}{\kern0pt}\ ds{\isacharprime}{\kern0pt}\ zs{\isacharparenright}{\kern0pt}\isanewline
\ \ \isacommand{case}\isamarkupfalse%
\ Nil\isanewline
\ \ \isacommand{then}\isamarkupfalse%
\ \isacommand{show}\isamarkupfalse%
\ {\isacharquery}{\kern0pt}case\ \isacommand{by}\isamarkupfalse%
\ auto\isanewline
\isacommand{next}\isamarkupfalse%
\isanewline
\ \ \isacommand{case}\isamarkupfalse%
\ {\isacharparenleft}{\kern0pt}Cons\ d\ ds{\isacharparenright}{\kern0pt}\isanewline
\ \ \isacommand{then}\isamarkupfalse%
\ \isacommand{show}\isamarkupfalse%
\ {\isacharquery}{\kern0pt}case\isanewline
\ \ \ \ \isacommand{by}\isamarkupfalse%
\ {\isacharparenleft}{\kern0pt}cases\ {\isachardoublequoteopen}d\ {\isacharequal}{\kern0pt}\ {\isadigit{0}}{\isachardoublequoteclose}{\isacharparenright}{\kern0pt}\ force{\isacharplus}{\kern0pt}\isanewline
\isacommand{qed}\isamarkupfalse%
%
\endisatagproof
{\isafoldproof}%
%
\isadelimproof
\isanewline
%
\endisadelimproof
\isanewline
\isacommand{lemma}\isamarkupfalse%
\ ttz{\isacharunderscore}{\kern0pt}last{\isacharunderscore}{\kern0pt}neq{\isacharunderscore}{\kern0pt}z{\isacharcolon}{\kern0pt}\ {\isachardoublequoteopen}trim{\isacharunderscore}{\kern0pt}tr{\isacharunderscore}{\kern0pt}zs\ ds\ {\isacharequal}{\kern0pt}\ ds{\isacharprime}{\kern0pt}\ {\isasymLongrightarrow}\ last\ ds{\isacharprime}{\kern0pt}\ {\isasymnoteq}\ {\isadigit{0}}\ {\isasymor}\ ds{\isacharprime}{\kern0pt}\ {\isacharequal}{\kern0pt}\ {\isacharbrackleft}{\kern0pt}{\isacharbrackright}{\kern0pt}{\isachardoublequoteclose}\isanewline
%
\isadelimproof
\ \ %
\endisadelimproof
%
\isatagproof
\isacommand{using}\isamarkupfalse%
\ ttz{\isacharunderscore}{\kern0pt}last{\isacharunderscore}{\kern0pt}neq{\isacharunderscore}{\kern0pt}z{\isacharunderscore}{\kern0pt}aux\ \isacommand{by}\isamarkupfalse%
\ {\isacharparenleft}{\kern0pt}auto\ simp{\isacharcolon}{\kern0pt}\ Let{\isacharunderscore}{\kern0pt}def\ split{\isacharcolon}{\kern0pt}\ if{\isacharunderscore}{\kern0pt}splits{\isacharparenright}{\kern0pt}%
\endisatagproof
{\isafoldproof}%
%
\isadelimproof
\isanewline
%
\endisadelimproof
\isanewline
\isacommand{lemma}\isamarkupfalse%
\ trim{\isacharunderscore}{\kern0pt}tr{\isacharunderscore}{\kern0pt}zs{\isacharprime}{\kern0pt}{\isacharunderscore}{\kern0pt}wf{\isacharcolon}{\kern0pt}\ {\isachardoublequoteopen}wf{\isacharunderscore}{\kern0pt}digits\ ds\ {\isasymand}\ wf{\isacharunderscore}{\kern0pt}digits\ zs\ {\isasymLongrightarrow}\ wf{\isacharunderscore}{\kern0pt}digits\ {\isacharparenleft}{\kern0pt}trim{\isacharunderscore}{\kern0pt}tr{\isacharunderscore}{\kern0pt}zs{\isacharprime}{\kern0pt}\ ds\ zs{\isacharparenright}{\kern0pt}{\isachardoublequoteclose}\isanewline
%
\isadelimproof
\ \ %
\endisadelimproof
%
\isatagproof
\isacommand{unfolding}\isamarkupfalse%
\ wf{\isacharunderscore}{\kern0pt}digits{\isacharunderscore}{\kern0pt}def\ \isacommand{by}\isamarkupfalse%
\ {\isacharparenleft}{\kern0pt}induction\ ds\ zs\ rule{\isacharcolon}{\kern0pt}\ trim{\isacharunderscore}{\kern0pt}tr{\isacharunderscore}{\kern0pt}zs{\isacharprime}{\kern0pt}{\isachardot}{\kern0pt}induct{\isacharparenright}{\kern0pt}\ auto%
\endisatagproof
{\isafoldproof}%
%
\isadelimproof
\isanewline
%
\endisadelimproof
\isanewline
\isacommand{lemma}\isamarkupfalse%
\ trim{\isacharunderscore}{\kern0pt}tr{\isacharunderscore}{\kern0pt}zs{\isacharunderscore}{\kern0pt}wf{\isacharcolon}{\kern0pt}\ {\isachardoublequoteopen}wf{\isacharunderscore}{\kern0pt}digits\ ds\ {\isasymLongrightarrow}\ wf{\isacharunderscore}{\kern0pt}digits\ {\isacharparenleft}{\kern0pt}trim{\isacharunderscore}{\kern0pt}tr{\isacharunderscore}{\kern0pt}zs\ ds{\isacharparenright}{\kern0pt}{\isachardoublequoteclose}\isanewline
%
\isadelimproof
\ \ %
\endisadelimproof
%
\isatagproof
\isacommand{using}\isamarkupfalse%
\ trim{\isacharunderscore}{\kern0pt}tr{\isacharunderscore}{\kern0pt}zs{\isacharprime}{\kern0pt}{\isacharunderscore}{\kern0pt}wf{\isacharbrackleft}{\kern0pt}\isakeyword{where}\ zs{\isacharequal}{\kern0pt}{\isachardoublequoteopen}{\isacharbrackleft}{\kern0pt}{\isacharbrackright}{\kern0pt}{\isachardoublequoteclose}{\isacharbrackright}{\kern0pt}\ \isacommand{by}\isamarkupfalse%
\ {\isacharparenleft}{\kern0pt}auto\ simp{\isacharcolon}{\kern0pt}\ wf{\isacharunderscore}{\kern0pt}digits{\isacharunderscore}{\kern0pt}def{\isacharparenright}{\kern0pt}%
\endisatagproof
{\isafoldproof}%
%
\isadelimproof
%
\endisadelimproof
%
\begin{isamarkuptext}%
Functions to convert between integers and a sequence of digits.%
\end{isamarkuptext}\isamarkuptrue%
\isacommand{fun}\isamarkupfalse%
\ int{\isacharunderscore}{\kern0pt}of{\isacharunderscore}{\kern0pt}digits\ {\isacharcolon}{\kern0pt}{\isacharcolon}{\kern0pt}\ {\isachardoublequoteopen}digit\ list\ {\isasymRightarrow}\ int\ {\isasymRightarrow}\ int{\isachardoublequoteclose}\ \isakeyword{where}\isanewline
\ \ {\isachardoublequoteopen}int{\isacharunderscore}{\kern0pt}of{\isacharunderscore}{\kern0pt}digits\ {\isacharbrackleft}{\kern0pt}{\isacharbrackright}{\kern0pt}\ acc\ {\isacharequal}{\kern0pt}\ acc{\isachardoublequoteclose}\isanewline
{\isacharbar}{\kern0pt}\ {\isachardoublequoteopen}int{\isacharunderscore}{\kern0pt}of{\isacharunderscore}{\kern0pt}digits\ {\isacharparenleft}{\kern0pt}d{\isacharhash}{\kern0pt}ds{\isacharparenright}{\kern0pt}\ acc\ {\isacharequal}{\kern0pt}\ int{\isacharunderscore}{\kern0pt}of{\isacharunderscore}{\kern0pt}digits\ ds\ {\isacharparenleft}{\kern0pt}acc\ {\isacharasterisk}{\kern0pt}\ {\isadigit{1}}{\isadigit{0}}\ {\isacharplus}{\kern0pt}\ d{\isacharparenright}{\kern0pt}{\isachardoublequoteclose}\isanewline
\isanewline
\isacommand{fun}\isamarkupfalse%
\ digits{\isacharunderscore}{\kern0pt}of{\isacharunderscore}{\kern0pt}int\ {\isacharcolon}{\kern0pt}{\isacharcolon}{\kern0pt}\ {\isachardoublequoteopen}int\ {\isasymRightarrow}\ digit\ list{\isachardoublequoteclose}\ \isakeyword{where}\isanewline
\ \ {\isachardoublequoteopen}digits{\isacharunderscore}{\kern0pt}of{\isacharunderscore}{\kern0pt}int\ i\ {\isacharequal}{\kern0pt}\ \isanewline
\ \ \ \ {\isacharparenleft}{\kern0pt}if\ i\ {\isasymle}\ {\isadigit{0}}\ then\ {\isacharbrackleft}{\kern0pt}{\isacharbrackright}{\kern0pt}\isanewline
\ \ \ \ else\ if\ i\ {\isacharless}{\kern0pt}\ {\isadigit{1}}{\isadigit{0}}\ then\ {\isacharbrackleft}{\kern0pt}nat\ i{\isacharbrackright}{\kern0pt}\isanewline
\ \ \ \ else\ digits{\isacharunderscore}{\kern0pt}of{\isacharunderscore}{\kern0pt}int\ {\isacharparenleft}{\kern0pt}i\ div\ {\isadigit{1}}{\isadigit{0}}{\isacharparenright}{\kern0pt}\ {\isacharat}{\kern0pt}\ {\isacharbrackleft}{\kern0pt}nat\ {\isacharparenleft}{\kern0pt}i\ mod\ {\isadigit{1}}{\isadigit{0}}{\isacharparenright}{\kern0pt}{\isacharbrackright}{\kern0pt}{\isacharparenright}{\kern0pt}{\isachardoublequoteclose}\isanewline
\isanewline
\isacommand{value}\isamarkupfalse%
\ {\isachardoublequoteopen}digits{\isacharunderscore}{\kern0pt}of{\isacharunderscore}{\kern0pt}int\ {\isacharparenleft}{\kern0pt}int{\isacharunderscore}{\kern0pt}of{\isacharunderscore}{\kern0pt}digits\ {\isacharbrackleft}{\kern0pt}{\isadigit{1}}{\isacharcomma}{\kern0pt}{\isadigit{2}}{\isacharcomma}{\kern0pt}{\isadigit{3}}{\isacharcomma}{\kern0pt}{\isadigit{4}}{\isacharbrackright}{\kern0pt}\ {\isadigit{0}}{\isacharparenright}{\kern0pt}{\isachardoublequoteclose}\isanewline
\isacommand{value}\isamarkupfalse%
\ {\isachardoublequoteopen}int{\isacharunderscore}{\kern0pt}of{\isacharunderscore}{\kern0pt}digits\ {\isacharparenleft}{\kern0pt}digits{\isacharunderscore}{\kern0pt}of{\isacharunderscore}{\kern0pt}int\ {\isadigit{1}}{\isadigit{2}}{\isadigit{3}}{\isadigit{4}}{\isacharparenright}{\kern0pt}\ {\isadigit{0}}{\isachardoublequoteclose}\isanewline
\isanewline
\isacommand{value}\isamarkupfalse%
\ {\isachardoublequoteopen}digits{\isacharunderscore}{\kern0pt}of{\isacharunderscore}{\kern0pt}int\ {\isacharparenleft}{\kern0pt}int{\isacharunderscore}{\kern0pt}of{\isacharunderscore}{\kern0pt}digits\ {\isacharbrackleft}{\kern0pt}{\isadigit{1}}{\isacharbrackright}{\kern0pt}\ {\isadigit{1}}{\isadigit{2}}{\isadigit{3}}{\isacharparenright}{\kern0pt}{\isachardoublequoteclose}\isanewline
\isanewline
\isacommand{lemma}\isamarkupfalse%
\ digits{\isacharunderscore}{\kern0pt}of{\isacharunderscore}{\kern0pt}int{\isacharunderscore}{\kern0pt}int{\isacharunderscore}{\kern0pt}of{\isacharunderscore}{\kern0pt}digits{\isacharunderscore}{\kern0pt}aux{\isacharcolon}{\kern0pt}\isanewline
\ \ \isakeyword{assumes}\ {\isachardoublequoteopen}acc\ {\isachargreater}{\kern0pt}\ {\isadigit{0}}{\isachardoublequoteclose}\ \isakeyword{and}\ {\isachardoublequoteopen}{\isasymforall}d\ {\isasymin}\ set\ ds{\isachardot}{\kern0pt}\ wf{\isacharunderscore}{\kern0pt}digit\ d{\isachardoublequoteclose}\isanewline
\ \ \isakeyword{shows}\ {\isachardoublequoteopen}digits{\isacharunderscore}{\kern0pt}of{\isacharunderscore}{\kern0pt}int\ {\isacharparenleft}{\kern0pt}int{\isacharunderscore}{\kern0pt}of{\isacharunderscore}{\kern0pt}digits\ ds\ acc{\isacharparenright}{\kern0pt}\ {\isacharequal}{\kern0pt}\ digits{\isacharunderscore}{\kern0pt}of{\isacharunderscore}{\kern0pt}int\ acc\ {\isacharat}{\kern0pt}\ ds{\isachardoublequoteclose}\isanewline
%
\isadelimproof
\ \ %
\endisadelimproof
%
\isatagproof
\isacommand{using}\isamarkupfalse%
\ assms\isanewline
\ \ \isacommand{by}\isamarkupfalse%
\ {\isacharparenleft}{\kern0pt}induction\ ds\ acc\ arbitrary{\isacharcolon}{\kern0pt}\ acc\ rule{\isacharcolon}{\kern0pt}\ int{\isacharunderscore}{\kern0pt}of{\isacharunderscore}{\kern0pt}digits{\isachardot}{\kern0pt}induct{\isacharparenright}{\kern0pt}\ {\isacharparenleft}{\kern0pt}auto\ simp{\isacharcolon}{\kern0pt}\ wf{\isacharunderscore}{\kern0pt}digit{\isacharunderscore}{\kern0pt}def{\isacharparenright}{\kern0pt}%
\endisatagproof
{\isafoldproof}%
%
\isadelimproof
%
\endisadelimproof
%
\begin{isamarkuptext}%
Proof for correctness of function \isa{int{\isacharunderscore}{\kern0pt}of{\isacharunderscore}{\kern0pt}digits} and \isa{digits{\isacharunderscore}{\kern0pt}of{\isacharunderscore}{\kern0pt}int}.%
\end{isamarkuptext}\isamarkuptrue%
\isacommand{lemma}\isamarkupfalse%
\ digits{\isacharunderscore}{\kern0pt}of{\isacharunderscore}{\kern0pt}int{\isacharunderscore}{\kern0pt}int{\isacharunderscore}{\kern0pt}of{\isacharunderscore}{\kern0pt}digits{\isacharcolon}{\kern0pt}\isanewline
\ \ \isakeyword{assumes}\ {\isachardoublequoteopen}wf{\isacharunderscore}{\kern0pt}digits\ ds{\isachardoublequoteclose}\isanewline
\ \ \isakeyword{shows}\ {\isachardoublequoteopen}digits{\isacharunderscore}{\kern0pt}of{\isacharunderscore}{\kern0pt}int\ {\isacharparenleft}{\kern0pt}int{\isacharunderscore}{\kern0pt}of{\isacharunderscore}{\kern0pt}digits\ ds\ {\isadigit{0}}{\isacharparenright}{\kern0pt}\ {\isacharequal}{\kern0pt}\ trim{\isacharunderscore}{\kern0pt}ld{\isacharunderscore}{\kern0pt}zs\ ds{\isachardoublequoteclose}\isanewline
%
\isadelimproof
\ \ %
\endisadelimproof
%
\isatagproof
\isacommand{using}\isamarkupfalse%
\ assms\ digits{\isacharunderscore}{\kern0pt}of{\isacharunderscore}{\kern0pt}int{\isacharunderscore}{\kern0pt}int{\isacharunderscore}{\kern0pt}of{\isacharunderscore}{\kern0pt}digits{\isacharunderscore}{\kern0pt}aux\isanewline
\isacommand{proof}\isamarkupfalse%
\ {\isacharparenleft}{\kern0pt}induction\ ds{\isacharparenright}{\kern0pt}\isanewline
\ \ \isacommand{case}\isamarkupfalse%
\ Nil\isanewline
\ \ \isacommand{then}\isamarkupfalse%
\ \isacommand{show}\isamarkupfalse%
\ {\isacharquery}{\kern0pt}case\ \isacommand{by}\isamarkupfalse%
\ auto\isanewline
\isacommand{next}\isamarkupfalse%
\isanewline
\ \ \isacommand{case}\isamarkupfalse%
\ {\isacharparenleft}{\kern0pt}Cons\ d\ ds{\isacharparenright}{\kern0pt}\isanewline
\ \ \isacommand{then}\isamarkupfalse%
\ \isacommand{show}\isamarkupfalse%
\ {\isacharquery}{\kern0pt}case\isanewline
\ \ \ \ \isacommand{by}\isamarkupfalse%
\ {\isacharparenleft}{\kern0pt}cases\ {\isachardoublequoteopen}d\ {\isacharequal}{\kern0pt}\ {\isadigit{0}}{\isachardoublequoteclose}{\isacharparenright}{\kern0pt}\ {\isacharparenleft}{\kern0pt}auto\ simp{\isacharcolon}{\kern0pt}\ wf{\isacharunderscore}{\kern0pt}digits{\isacharunderscore}{\kern0pt}def\ wf{\isacharunderscore}{\kern0pt}digit{\isacharunderscore}{\kern0pt}def{\isacharparenright}{\kern0pt}\isanewline
\isacommand{qed}\isamarkupfalse%
%
\endisatagproof
{\isafoldproof}%
%
\isadelimproof
%
\endisadelimproof
%
\begin{isamarkuptext}%
Function to convert from a sequences of digits to a rational number.%
\end{isamarkuptext}\isamarkuptrue%
\isacommand{primrec}\isamarkupfalse%
\ rat{\isacharunderscore}{\kern0pt}of{\isacharunderscore}{\kern0pt}digits{\isacharunderscore}{\kern0pt}pair\ {\isacharcolon}{\kern0pt}{\isacharcolon}{\kern0pt}\ {\isachardoublequoteopen}digit\ list\ {\isasymtimes}\ digit\ list\ {\isasymRightarrow}\ rat{\isachardoublequoteclose}\ \isakeyword{where}\isanewline
\ \ {\isachardoublequoteopen}rat{\isacharunderscore}{\kern0pt}of{\isacharunderscore}{\kern0pt}digits{\isacharunderscore}{\kern0pt}pair\ {\isacharparenleft}{\kern0pt}ds\isactrlsub {\isadigit{1}}{\isacharcomma}{\kern0pt}ds\isactrlsub {\isadigit{2}}{\isacharparenright}{\kern0pt}\ {\isacharequal}{\kern0pt}\ {\isacharparenleft}{\kern0pt}\isanewline
\ \ \ \ let\ ds\isactrlsub {\isadigit{1}}{\isacharprime}{\kern0pt}\ {\isacharequal}{\kern0pt}\ trim{\isacharunderscore}{\kern0pt}ld{\isacharunderscore}{\kern0pt}zs\ ds\isactrlsub {\isadigit{1}}{\isacharsemicolon}{\kern0pt}\ ds\isactrlsub {\isadigit{2}}{\isacharprime}{\kern0pt}\ {\isacharequal}{\kern0pt}\ trim{\isacharunderscore}{\kern0pt}tr{\isacharunderscore}{\kern0pt}zs\ ds\isactrlsub {\isadigit{2}}\ in\isanewline
\ \ \ \ \ \ Rat{\isachardot}{\kern0pt}Fract\ {\isacharparenleft}{\kern0pt}int{\isacharunderscore}{\kern0pt}of{\isacharunderscore}{\kern0pt}digits\ {\isacharparenleft}{\kern0pt}ds\isactrlsub {\isadigit{1}}{\isacharprime}{\kern0pt}\ {\isacharat}{\kern0pt}\ ds\isactrlsub {\isadigit{2}}{\isacharprime}{\kern0pt}{\isacharparenright}{\kern0pt}\ {\isadigit{0}}{\isacharparenright}{\kern0pt}\ {\isacharparenleft}{\kern0pt}{\isadigit{1}}{\isadigit{0}}\ {\isacharcircum}{\kern0pt}\ length\ ds\isactrlsub {\isadigit{2}}{\isacharprime}{\kern0pt}{\isacharparenright}{\kern0pt}\isanewline
\ \ {\isacharparenright}{\kern0pt}{\isachardoublequoteclose}%
\begin{isamarkuptext}%
For proofs about the function \isa{rat{\isacharunderscore}{\kern0pt}of{\isacharunderscore}{\kern0pt}digits{\isacharunderscore}{\kern0pt}pair} see \texttt{Rat\_Parsing.thy}.%
\end{isamarkuptext}\isamarkuptrue%
%
\isadelimdocument
%
\endisadelimdocument
%
\isatagdocument
%
\isamarkupsubsubsection{Code Generation%
}
\isamarkuptrue%
%
\endisatagdocument
{\isafolddocument}%
%
\isadelimdocument
%
\endisadelimdocument
\isacommand{export{\isacharunderscore}{\kern0pt}code}\isamarkupfalse%
\isanewline
\ \ check{\isacharunderscore}{\kern0pt}plan\isanewline
\ \ nat{\isacharunderscore}{\kern0pt}of{\isacharunderscore}{\kern0pt}integer\ integer{\isacharunderscore}{\kern0pt}of{\isacharunderscore}{\kern0pt}nat\ int{\isacharunderscore}{\kern0pt}of{\isacharunderscore}{\kern0pt}integer\ integer{\isacharunderscore}{\kern0pt}of{\isacharunderscore}{\kern0pt}int\ Inl\ Inr\ Rat{\isachardot}{\kern0pt}Fract\ Rat{\isachardot}{\kern0pt}of{\isacharunderscore}{\kern0pt}int\ rat{\isacharunderscore}{\kern0pt}of{\isacharunderscore}{\kern0pt}digits{\isacharunderscore}{\kern0pt}pair\isanewline
\ \ predAtm\ eqAtm\ predicate\ Pred\ Func\ Either\ Var\ Obj\ PredDecl\ FuncDecl\ BigAnd\ BigOr\isanewline
\ \ formula{\isachardot}{\kern0pt}Not\ formula{\isachardot}{\kern0pt}Bot\ Effect\ No{\isacharunderscore}{\kern0pt}Const\ Time{\isacharunderscore}{\kern0pt}Const\ Func{\isacharunderscore}{\kern0pt}Const\ LEQ\ EQ\ GEQ\isanewline
\ \ Simple{\isacharunderscore}{\kern0pt}Action{\isacharunderscore}{\kern0pt}Schema\ Durative{\isacharunderscore}{\kern0pt}Action{\isacharunderscore}{\kern0pt}Schema\ At{\isacharunderscore}{\kern0pt}Start\ At{\isacharunderscore}{\kern0pt}End\ Over{\isacharunderscore}{\kern0pt}All\isanewline
\ \ map{\isacharunderscore}{\kern0pt}atom\ Domain\ Problem\ Simple{\isacharunderscore}{\kern0pt}Plan{\isacharunderscore}{\kern0pt}Action\ Durative{\isacharunderscore}{\kern0pt}Plan{\isacharunderscore}{\kern0pt}Action\isanewline
\ \ term{\isachardot}{\kern0pt}CONST\ term{\isachardot}{\kern0pt}VAR\ \isanewline
\ \ String{\isachardot}{\kern0pt}explode\ String{\isachardot}{\kern0pt}implode\isanewline
\ \ \isakeyword{in}\ SML\isanewline
\ \ \isakeyword{module{\isacharunderscore}{\kern0pt}name}\ TEMPORAL{\isacharunderscore}{\kern0pt}PDDL{\isacharunderscore}{\kern0pt}Checker{\isacharunderscore}{\kern0pt}Exported\isanewline
\ \ \isakeyword{file}\ {\isachardoublequoteopen}code{\isacharslash}{\kern0pt}TEMPORAL{\isacharunderscore}{\kern0pt}PDDL{\isacharunderscore}{\kern0pt}Checker{\isacharunderscore}{\kern0pt}Exported{\isachardot}{\kern0pt}sml{\isachardoublequoteclose}\isanewline
\ \ \ \ \isanewline
\isanewline
\isanewline
\isanewline
\isanewline
%
\isadelimtheory
\isanewline
%
\endisadelimtheory
%
\isatagtheory
\isacommand{end}\isamarkupfalse%
\ %
\isamarkupcmt{Theory%
}%
\endisatagtheory
{\isafoldtheory}%
%
\isadelimtheory
%
\endisadelimtheory
%
\end{isabellebody}%
\endinput
%:%file=TEMPORAL_PDDL_Checker.tex%:%
%:%11=2%:%
%:%27=3%:%
%:%28=3%:%
%:%29=4%:%
%:%30=5%:%
%:%31=6%:%
%:%32=7%:%
%:%33=8%:%
%:%34=9%:%
%:%35=10%:%
%:%36=11%:%
%:%37=12%:%
%:%38=13%:%
%:%39=14%:%
%:%53=16%:%
%:%65=17%:%
%:%67=19%:%
%:%68=19%:%
%:%69=20%:%
%:%70=20%:%
%:%73=21%:%
%:%77=21%:%
%:%78=21%:%
%:%79=21%:%
%:%84=21%:%
%:%87=22%:%
%:%88=23%:%
%:%89=23%:%
%:%90=24%:%
%:%91=25%:%
%:%92=26%:%
%:%93=27%:%
%:%94=27%:%
%:%95=28%:%
%:%96=29%:%
%:%97=30%:%
%:%98=30%:%
%:%108=40%:%
%:%109=41%:%
%:%110=42%:%
%:%111=43%:%
%:%112=43%:%
%:%113=44%:%
%:%114=45%:%
%:%115=46%:%
%:%116=47%:%
%:%117=48%:%
%:%118=48%:%
%:%119=49%:%
%:%120=50%:%
%:%121=50%:%
%:%126=55%:%
%:%127=56%:%
%:%128=57%:%
%:%129=57%:%
%:%131=59%:%
%:%134=60%:%
%:%138=60%:%
%:%139=60%:%
%:%140=61%:%
%:%141=61%:%
%:%142=62%:%
%:%148=62%:%
%:%151=63%:%
%:%152=64%:%
%:%153=64%:%
%:%154=65%:%
%:%156=65%:%
%:%160=65%:%
%:%161=65%:%
%:%168=65%:%
%:%169=66%:%
%:%170=67%:%
%:%171=67%:%
%:%174=68%:%
%:%178=68%:%
%:%179=68%:%
%:%180=69%:%
%:%181=69%:%
%:%182=70%:%
%:%183=71%:%
%:%184=72%:%
%:%185=73%:%
%:%186=73%:%
%:%187=73%:%
%:%188=74%:%
%:%189=74%:%
%:%190=75%:%
%:%191=75%:%
%:%192=76%:%
%:%193=76%:%
%:%194=77%:%
%:%195=77%:%
%:%196=77%:%
%:%197=78%:%
%:%198=78%:%
%:%199=78%:%
%:%200=79%:%
%:%201=79%:%
%:%202=79%:%
%:%203=80%:%
%:%209=80%:%
%:%212=81%:%
%:%213=82%:%
%:%214=82%:%
%:%215=83%:%
%:%216=84%:%
%:%217=84%:%
%:%218=85%:%
%:%219=86%:%
%:%220=87%:%
%:%221=87%:%
%:%224=88%:%
%:%228=88%:%
%:%229=88%:%
%:%234=88%:%
%:%237=89%:%
%:%238=90%:%
%:%239=91%:%
%:%240=91%:%
%:%247=92%:%
%:%248=92%:%
%:%249=93%:%
%:%250=93%:%
%:%251=94%:%
%:%252=95%:%
%:%253=95%:%
%:%254=96%:%
%:%255=96%:%
%:%256=97%:%
%:%257=98%:%
%:%258=99%:%
%:%259=100%:%
%:%260=101%:%
%:%261=101%:%
%:%262=101%:%
%:%263=102%:%
%:%264=102%:%
%:%265=103%:%
%:%271=103%:%
%:%274=104%:%
%:%275=105%:%
%:%276=105%:%
%:%277=106%:%
%:%278=107%:%
%:%279=107%:%
%:%280=108%:%
%:%283=109%:%
%:%287=109%:%
%:%288=109%:%
%:%289=110%:%
%:%290=110%:%
%:%291=111%:%
%:%297=111%:%
%:%300=112%:%
%:%301=113%:%
%:%302=113%:%
%:%305=114%:%
%:%309=114%:%
%:%310=114%:%
%:%311=115%:%
%:%312=115%:%
%:%326=119%:%
%:%330=121%:%
%:%340=123%:%
%:%341=123%:%
%:%342=124%:%
%:%343=125%:%
%:%344=126%:%
%:%345=126%:%
%:%346=127%:%
%:%347=128%:%
%:%348=129%:%
%:%349=130%:%
%:%350=130%:%
%:%353=131%:%
%:%357=131%:%
%:%358=131%:%
%:%363=131%:%
%:%366=132%:%
%:%367=133%:%
%:%368=133%:%
%:%371=134%:%
%:%375=134%:%
%:%376=134%:%
%:%377=134%:%
%:%382=134%:%
%:%385=135%:%
%:%386=136%:%
%:%387=136%:%
%:%388=137%:%
%:%389=138%:%
%:%390=138%:%
%:%391=139%:%
%:%392=140%:%
%:%393=140%:%
%:%396=141%:%
%:%400=141%:%
%:%401=141%:%
%:%402=141%:%
%:%407=141%:%
%:%410=142%:%
%:%411=143%:%
%:%412=143%:%
%:%413=144%:%
%:%414=145%:%
%:%415=145%:%
%:%418=146%:%
%:%422=146%:%
%:%423=146%:%
%:%428=146%:%
%:%431=147%:%
%:%432=148%:%
%:%433=148%:%
%:%436=149%:%
%:%440=149%:%
%:%441=149%:%
%:%446=149%:%
%:%449=150%:%
%:%450=151%:%
%:%451=151%:%
%:%454=152%:%
%:%458=152%:%
%:%459=152%:%
%:%465=152%:%
%:%468=153%:%
%:%469=154%:%
%:%470=154%:%
%:%471=155%:%
%:%472=156%:%
%:%473=157%:%
%:%474=157%:%
%:%477=158%:%
%:%481=158%:%
%:%482=158%:%
%:%483=159%:%
%:%484=159%:%
%:%493=162%:%
%:%495=163%:%
%:%496=163%:%
%:%497=164%:%
%:%498=164%:%
%:%499=164%:%
%:%500=165%:%
%:%501=166%:%
%:%502=167%:%
%:%503=167%:%
%:%506=170%:%
%:%507=171%:%
%:%508=172%:%
%:%509=172%:%
%:%512=173%:%
%:%516=173%:%
%:%517=173%:%
%:%523=173%:%
%:%526=174%:%
%:%527=175%:%
%:%528=175%:%
%:%531=178%:%
%:%532=179%:%
%:%533=180%:%
%:%534=180%:%
%:%537=181%:%
%:%541=181%:%
%:%542=181%:%
%:%547=181%:%
%:%550=182%:%
%:%551=183%:%
%:%552=183%:%
%:%555=186%:%
%:%556=187%:%
%:%557=188%:%
%:%558=188%:%
%:%561=189%:%
%:%565=189%:%
%:%566=189%:%
%:%571=189%:%
%:%574=190%:%
%:%575=191%:%
%:%576=191%:%
%:%577=192%:%
%:%578=193%:%
%:%579=194%:%
%:%580=195%:%
%:%581=195%:%
%:%584=196%:%
%:%588=196%:%
%:%589=196%:%
%:%594=196%:%
%:%597=197%:%
%:%598=198%:%
%:%599=198%:%
%:%600=199%:%
%:%601=200%:%
%:%602=201%:%
%:%603=202%:%
%:%604=203%:%
%:%605=204%:%
%:%606=205%:%
%:%607=206%:%
%:%608=206%:%
%:%611=207%:%
%:%615=207%:%
%:%616=207%:%
%:%621=207%:%
%:%624=208%:%
%:%625=209%:%
%:%626=209%:%
%:%627=210%:%
%:%628=211%:%
%:%629=212%:%
%:%630=213%:%
%:%631=213%:%
%:%634=214%:%
%:%638=214%:%
%:%639=214%:%
%:%640=215%:%
%:%645=215%:%
%:%648=216%:%
%:%649=217%:%
%:%650=217%:%
%:%651=218%:%
%:%653=220%:%
%:%654=221%:%
%:%655=222%:%
%:%656=222%:%
%:%659=223%:%
%:%663=223%:%
%:%664=223%:%
%:%669=223%:%
%:%672=224%:%
%:%673=225%:%
%:%674=225%:%
%:%675=226%:%
%:%676=227%:%
%:%677=228%:%
%:%680=231%:%
%:%681=232%:%
%:%682=233%:%
%:%683=233%:%
%:%686=234%:%
%:%690=234%:%
%:%691=234%:%
%:%696=234%:%
%:%699=235%:%
%:%700=236%:%
%:%701=236%:%
%:%702=236%:%
%:%703=236%:%
%:%704=237%:%
%:%705=238%:%
%:%706=238%:%
%:%707=239%:%
%:%713=245%:%
%:%714=246%:%
%:%721=253%:%
%:%722=254%:%
%:%723=255%:%
%:%724=255%:%
%:%727=256%:%
%:%731=256%:%
%:%732=256%:%
%:%733=257%:%
%:%734=258%:%
%:%739=258%:%
%:%742=259%:%
%:%743=260%:%
%:%744=260%:%
%:%745=261%:%
%:%755=271%:%
%:%756=272%:%
%:%757=273%:%
%:%758=273%:%
%:%761=274%:%
%:%765=274%:%
%:%766=274%:%
%:%767=275%:%
%:%768=275%:%
%:%773=275%:%
%:%776=276%:%
%:%777=277%:%
%:%778=278%:%
%:%779=278%:%
%:%780=278%:%
%:%788=280%:%
%:%798=282%:%
%:%799=282%:%
%:%801=284%:%
%:%802=285%:%
%:%806=309%:%
%:%808=310%:%
%:%809=310%:%
%:%810=311%:%
%:%814=315%:%
%:%815=316%:%
%:%816=317%:%
%:%817=317%:%
%:%820=318%:%
%:%824=318%:%
%:%825=318%:%
%:%830=318%:%
%:%833=319%:%
%:%834=320%:%
%:%835=320%:%
%:%838=321%:%
%:%842=321%:%
%:%843=321%:%
%:%852=323%:%
%:%854=324%:%
%:%855=324%:%
%:%858=325%:%
%:%862=325%:%
%:%863=325%:%
%:%868=325%:%
%:%871=326%:%
%:%872=327%:%
%:%873=327%:%
%:%874=327%:%
%:%882=329%:%
%:%892=331%:%
%:%893=331%:%
%:%895=333%:%
%:%896=334%:%
%:%897=335%:%
%:%899=337%:%
%:%900=337%:%
%:%901=338%:%
%:%902=339%:%
%:%903=339%:%
%:%904=340%:%
%:%905=341%:%
%:%906=342%:%
%:%907=342%:%
%:%910=343%:%
%:%914=343%:%
%:%915=343%:%
%:%916=344%:%
%:%917=344%:%
%:%926=346%:%
%:%928=348%:%
%:%929=348%:%
%:%931=350%:%
%:%932=351%:%
%:%933=352%:%
%:%934=352%:%
%:%935=353%:%
%:%938=354%:%
%:%942=354%:%
%:%943=354%:%
%:%944=355%:%
%:%945=355%:%
%:%946=356%:%
%:%952=356%:%
%:%955=357%:%
%:%956=358%:%
%:%957=358%:%
%:%958=359%:%
%:%959=360%:%
%:%960=361%:%
%:%961=362%:%
%:%962=363%:%
%:%963=363%:%
%:%966=364%:%
%:%970=364%:%
%:%971=364%:%
%:%980=366%:%
%:%982=368%:%
%:%983=368%:%
%:%984=369%:%
%:%985=370%:%
%:%986=371%:%
%:%987=371%:%
%:%990=372%:%
%:%994=372%:%
%:%995=372%:%
%:%1000=372%:%
%:%1003=373%:%
%:%1004=374%:%
%:%1005=375%:%
%:%1006=375%:%
%:%1007=376%:%
%:%1008=377%:%
%:%1009=378%:%
%:%1010=378%:%
%:%1011=379%:%
%:%1014=380%:%
%:%1018=380%:%
%:%1019=380%:%
%:%1020=381%:%
%:%1021=381%:%
%:%1026=381%:%
%:%1029=382%:%
%:%1030=383%:%
%:%1031=383%:%
%:%1037=389%:%
%:%1038=390%:%
%:%1039=391%:%
%:%1040=391%:%
%:%1041=392%:%
%:%1044=393%:%
%:%1048=393%:%
%:%1049=393%:%
%:%1050=394%:%
%:%1051=394%:%
%:%1060=396%:%
%:%1061=397%:%
%:%1062=398%:%
%:%1064=400%:%
%:%1065=400%:%
%:%1066=401%:%
%:%1067=402%:%
%:%1068=403%:%
%:%1069=404%:%
%:%1070=405%:%
%:%1071=406%:%
%:%1074=407%:%
%:%1078=407%:%
%:%1079=407%:%
%:%1080=408%:%
%:%1081=408%:%
%:%1086=408%:%
%:%1089=409%:%
%:%1090=410%:%
%:%1091=410%:%
%:%1092=411%:%
%:%1093=412%:%
%:%1094=413%:%
%:%1095=414%:%
%:%1096=415%:%
%:%1097=416%:%
%:%1100=417%:%
%:%1104=417%:%
%:%1105=417%:%
%:%1106=418%:%
%:%1107=418%:%
%:%1112=418%:%
%:%1115=419%:%
%:%1116=420%:%
%:%1117=420%:%
%:%1118=420%:%
%:%1126=423%:%
%:%1136=425%:%
%:%1137=425%:%
%:%1139=427%:%
%:%1141=428%:%
%:%1142=428%:%
%:%1145=431%:%
%:%1146=432%:%
%:%1147=433%:%
%:%1148=433%:%
%:%1149=434%:%
%:%1152=435%:%
%:%1156=435%:%
%:%1157=435%:%
%:%1158=435%:%
%:%1163=435%:%
%:%1166=436%:%
%:%1167=437%:%
%:%1168=437%:%
%:%1169=437%:%
%:%1170=437%:%
%:%1171=438%:%
%:%1172=439%:%
%:%1173=439%:%
%:%1175=441%:%
%:%1176=442%:%
%:%1178=443%:%
%:%1179=443%:%
%:%1181=445%:%
%:%1183=447%:%
%:%1184=447%:%
%:%1185=448%:%
%:%1186=449%:%
%:%1189=450%:%
%:%1193=450%:%
%:%1194=450%:%
%:%1195=450%:%
%:%1196=450%:%
%:%1205=452%:%
%:%1206=453%:%
%:%1207=454%:%
%:%1208=455%:%
%:%1210=457%:%
%:%1211=457%:%
%:%1212=458%:%
%:%1216=462%:%
%:%1217=463%:%
%:%1218=464%:%
%:%1219=464%:%
%:%1220=465%:%
%:%1221=466%:%
%:%1222=467%:%
%:%1225=468%:%
%:%1229=468%:%
%:%1230=468%:%
%:%1231=468%:%
%:%1240=470%:%
%:%1242=471%:%
%:%1243=471%:%
%:%1244=472%:%
%:%1245=473%:%
%:%1246=474%:%
%:%1249=475%:%
%:%1253=475%:%
%:%1254=475%:%
%:%1255=476%:%
%:%1256=476%:%
%:%1257=477%:%
%:%1258=478%:%
%:%1259=479%:%
%:%1260=479%:%
%:%1269=481%:%
%:%1270=482%:%
%:%1271=483%:%
%:%1272=484%:%
%:%1274=485%:%
%:%1275=485%:%
%:%1277=487%:%
%:%1278=488%:%
%:%1279=489%:%
%:%1280=489%:%
%:%1281=489%:%
%:%1289=491%:%
%:%1299=493%:%
%:%1300=493%:%
%:%1302=495%:%
%:%1303=496%:%
%:%1305=497%:%
%:%1306=497%:%
%:%1307=498%:%
%:%1308=499%:%
%:%1313=504%:%
%:%1314=505%:%
%:%1316=507%:%
%:%1317=507%:%
%:%1320=508%:%
%:%1324=508%:%
%:%1325=508%:%
%:%1330=508%:%
%:%1333=509%:%
%:%1334=510%:%
%:%1335=510%:%
%:%1338=511%:%
%:%1342=511%:%
%:%1343=511%:%
%:%1344=511%:%
%:%1349=511%:%
%:%1352=512%:%
%:%1353=513%:%
%:%1354=513%:%
%:%1357=514%:%
%:%1361=514%:%
%:%1362=514%:%
%:%1367=514%:%
%:%1370=515%:%
%:%1371=516%:%
%:%1372=516%:%
%:%1375=517%:%
%:%1379=517%:%
%:%1380=517%:%
%:%1381=517%:%
%:%1386=517%:%
%:%1389=518%:%
%:%1390=519%:%
%:%1391=519%:%
%:%1394=520%:%
%:%1398=520%:%
%:%1399=520%:%
%:%1408=522%:%
%:%1409=523%:%
%:%1411=524%:%
%:%1412=524%:%
%:%1413=525%:%
%:%1414=526%:%
%:%1415=527%:%
%:%1416=528%:%
%:%1419=531%:%
%:%1421=533%:%
%:%1422=533%:%
%:%1425=534%:%
%:%1429=534%:%
%:%1430=534%:%
%:%1435=534%:%
%:%1438=535%:%
%:%1439=536%:%
%:%1440=536%:%
%:%1443=537%:%
%:%1447=537%:%
%:%1448=537%:%
%:%1449=537%:%
%:%1454=537%:%
%:%1457=538%:%
%:%1458=539%:%
%:%1459=539%:%
%:%1462=540%:%
%:%1466=540%:%
%:%1467=540%:%
%:%1468=541%:%
%:%1469=541%:%
%:%1478=543%:%
%:%1482=545%:%
%:%1484=546%:%
%:%1485=546%:%
%:%1488=547%:%
%:%1492=547%:%
%:%1493=547%:%
%:%1494=548%:%
%:%1495=548%:%
%:%1496=549%:%
%:%1497=549%:%
%:%1498=549%:%
%:%1499=549%:%
%:%1500=550%:%
%:%1501=550%:%
%:%1502=551%:%
%:%1503=551%:%
%:%1504=552%:%
%:%1505=552%:%
%:%1506=553%:%
%:%1507=553%:%
%:%1508=554%:%
%:%1509=554%:%
%:%1510=554%:%
%:%1511=555%:%
%:%1512=555%:%
%:%1513=556%:%
%:%1514=556%:%
%:%1515=557%:%
%:%1521=557%:%
%:%1524=558%:%
%:%1525=559%:%
%:%1526=559%:%
%:%1529=560%:%
%:%1533=560%:%
%:%1534=560%:%
%:%1539=560%:%
%:%1542=561%:%
%:%1543=562%:%
%:%1544=562%:%
%:%1547=563%:%
%:%1551=563%:%
%:%1552=563%:%
%:%1553=564%:%
%:%1554=564%:%
%:%1563=566%:%
%:%1565=567%:%
%:%1566=567%:%
%:%1569=568%:%
%:%1573=568%:%
%:%1574=568%:%
%:%1575=569%:%
%:%1576=569%:%
%:%1577=570%:%
%:%1578=570%:%
%:%1579=570%:%
%:%1580=571%:%
%:%1581=571%:%
%:%1582=571%:%
%:%1583=572%:%
%:%1584=572%:%
%:%1585=573%:%
%:%1586=573%:%
%:%1587=574%:%
%:%1588=574%:%
%:%1589=575%:%
%:%1590=575%:%
%:%1591=575%:%
%:%1592=575%:%
%:%1593=576%:%
%:%1594=576%:%
%:%1595=576%:%
%:%1596=577%:%
%:%1597=577%:%
%:%1598=578%:%
%:%1599=578%:%
%:%1600=579%:%
%:%1601=579%:%
%:%1602=579%:%
%:%1603=580%:%
%:%1604=580%:%
%:%1605=581%:%
%:%1606=581%:%
%:%1607=581%:%
%:%1608=582%:%
%:%1609=582%:%
%:%1610=583%:%
%:%1611=583%:%
%:%1612=583%:%
%:%1613=584%:%
%:%1614=584%:%
%:%1615=584%:%
%:%1616=585%:%
%:%1617=585%:%
%:%1618=585%:%
%:%1619=586%:%
%:%1620=586%:%
%:%1621=586%:%
%:%1622=587%:%
%:%1623=587%:%
%:%1624=588%:%
%:%1625=588%:%
%:%1626=589%:%
%:%1627=589%:%
%:%1628=589%:%
%:%1629=590%:%
%:%1630=590%:%
%:%1631=591%:%
%:%1632=591%:%
%:%1633=591%:%
%:%1634=592%:%
%:%1635=592%:%
%:%1636=592%:%
%:%1637=593%:%
%:%1638=593%:%
%:%1639=593%:%
%:%1640=594%:%
%:%1641=594%:%
%:%1642=595%:%
%:%1643=595%:%
%:%1644=595%:%
%:%1645=596%:%
%:%1646=596%:%
%:%1647=596%:%
%:%1648=597%:%
%:%1649=597%:%
%:%1650=597%:%
%:%1651=598%:%
%:%1652=598%:%
%:%1653=599%:%
%:%1659=599%:%
%:%1662=600%:%
%:%1663=601%:%
%:%1664=601%:%
%:%1667=602%:%
%:%1671=602%:%
%:%1672=602%:%
%:%1673=602%:%
%:%1678=602%:%
%:%1681=603%:%
%:%1682=604%:%
%:%1683=604%:%
%:%1684=605%:%
%:%1687=608%:%
%:%1689=609%:%
%:%1690=609%:%
%:%1691=610%:%
%:%1694=613%:%
%:%1695=614%:%
%:%1702=621%:%
%:%1703=622%:%
%:%1705=624%:%
%:%1706=624%:%
%:%1707=625%:%
%:%1708=626%:%
%:%1709=627%:%
%:%1712=628%:%
%:%1716=628%:%
%:%1717=628%:%
%:%1718=629%:%
%:%1719=629%:%
%:%1720=630%:%
%:%1721=630%:%
%:%1722=631%:%
%:%1723=631%:%
%:%1724=632%:%
%:%1732=632%:%
%:%1733=633%:%
%:%1734=634%:%
%:%1735=634%:%
%:%1736=635%:%
%:%1737=636%:%
%:%1738=637%:%
%:%1741=638%:%
%:%1745=638%:%
%:%1746=638%:%
%:%1747=639%:%
%:%1748=639%:%
%:%1749=640%:%
%:%1750=640%:%
%:%1751=641%:%
%:%1752=641%:%
%:%1753=641%:%
%:%1754=642%:%
%:%1755=642%:%
%:%1756=643%:%
%:%1764=643%:%
%:%1765=644%:%
%:%1766=645%:%
%:%1767=645%:%
%:%1768=646%:%
%:%1769=647%:%
%:%1770=648%:%
%:%1773=649%:%
%:%1777=649%:%
%:%1778=649%:%
%:%1779=650%:%
%:%1780=650%:%
%:%1781=651%:%
%:%1782=651%:%
%:%1783=652%:%
%:%1784=652%:%
%:%1785=652%:%
%:%1786=653%:%
%:%1787=653%:%
%:%1788=654%:%
%:%1789=654%:%
%:%1790=655%:%
%:%1791=655%:%
%:%1792=656%:%
%:%1793=656%:%
%:%1794=657%:%
%:%1795=657%:%
%:%1796=658%:%
%:%1797=658%:%
%:%1798=658%:%
%:%1799=659%:%
%:%1800=659%:%
%:%1801=660%:%
%:%1802=660%:%
%:%1803=660%:%
%:%1804=661%:%
%:%1805=661%:%
%:%1806=661%:%
%:%1807=662%:%
%:%1808=662%:%
%:%1809=662%:%
%:%1810=663%:%
%:%1811=663%:%
%:%1812=663%:%
%:%1813=664%:%
%:%1814=664%:%
%:%1815=665%:%
%:%1816=665%:%
%:%1817=665%:%
%:%1818=666%:%
%:%1819=666%:%
%:%1820=666%:%
%:%1821=667%:%
%:%1822=667%:%
%:%1823=667%:%
%:%1824=668%:%
%:%1825=668%:%
%:%1826=669%:%
%:%1827=669%:%
%:%1828=670%:%
%:%1829=671%:%
%:%1830=671%:%
%:%1831=672%:%
%:%1832=672%:%
%:%1833=672%:%
%:%1834=673%:%
%:%1835=673%:%
%:%1836=674%:%
%:%1837=674%:%
%:%1838=675%:%
%:%1844=675%:%
%:%1847=676%:%
%:%1848=677%:%
%:%1849=677%:%
%:%1850=678%:%
%:%1853=679%:%
%:%1857=679%:%
%:%1858=679%:%
%:%1859=680%:%
%:%1860=680%:%
%:%1861=681%:%
%:%1862=681%:%
%:%1863=682%:%
%:%1864=682%:%
%:%1865=683%:%
%:%1866=683%:%
%:%1867=684%:%
%:%1868=684%:%
%:%1869=684%:%
%:%1870=685%:%
%:%1871=685%:%
%:%1872=685%:%
%:%1873=686%:%
%:%1874=686%:%
%:%1875=686%:%
%:%1876=687%:%
%:%1877=687%:%
%:%1878=688%:%
%:%1879=688%:%
%:%1880=688%:%
%:%1881=689%:%
%:%1882=689%:%
%:%1883=689%:%
%:%1884=690%:%
%:%1885=690%:%
%:%1886=690%:%
%:%1887=691%:%
%:%1888=691%:%
%:%1889=692%:%
%:%1890=692%:%
%:%1891=693%:%
%:%1892=693%:%
%:%1893=693%:%
%:%1894=694%:%
%:%1895=694%:%
%:%1896=695%:%
%:%1897=695%:%
%:%1898=696%:%
%:%1904=696%:%
%:%1907=697%:%
%:%1908=698%:%
%:%1909=698%:%
%:%1910=699%:%
%:%1911=700%:%
%:%1914=701%:%
%:%1918=701%:%
%:%1919=701%:%
%:%1920=701%:%
%:%1925=701%:%
%:%1928=702%:%
%:%1929=703%:%
%:%1930=703%:%
%:%1931=704%:%
%:%1932=705%:%
%:%1933=706%:%
%:%1936=707%:%
%:%1940=707%:%
%:%1941=707%:%
%:%1942=708%:%
%:%1943=708%:%
%:%1944=709%:%
%:%1945=709%:%
%:%1946=710%:%
%:%1947=710%:%
%:%1948=710%:%
%:%1949=711%:%
%:%1950=711%:%
%:%1951=712%:%
%:%1952=712%:%
%:%1953=713%:%
%:%1954=713%:%
%:%1955=714%:%
%:%1956=714%:%
%:%1957=715%:%
%:%1958=715%:%
%:%1959=716%:%
%:%1960=716%:%
%:%1961=716%:%
%:%1962=717%:%
%:%1963=717%:%
%:%1964=717%:%
%:%1965=718%:%
%:%1971=718%:%
%:%1974=719%:%
%:%1975=720%:%
%:%1976=720%:%
%:%1977=721%:%
%:%1978=722%:%
%:%1979=723%:%
%:%1980=724%:%
%:%1983=725%:%
%:%1987=725%:%
%:%1988=725%:%
%:%1989=726%:%
%:%1990=726%:%
%:%1991=727%:%
%:%1992=727%:%
%:%1993=727%:%
%:%1994=728%:%
%:%1995=728%:%
%:%1996=728%:%
%:%1997=729%:%
%:%1998=729%:%
%:%1999=729%:%
%:%2000=730%:%
%:%2001=730%:%
%:%2002=730%:%
%:%2003=731%:%
%:%2004=731%:%
%:%2005=731%:%
%:%2006=732%:%
%:%2007=732%:%
%:%2008=732%:%
%:%2009=733%:%
%:%2010=733%:%
%:%2011=734%:%
%:%2012=734%:%
%:%2013=735%:%
%:%2014=735%:%
%:%2015=735%:%
%:%2016=736%:%
%:%2017=736%:%
%:%2018=737%:%
%:%2019=737%:%
%:%2020=738%:%
%:%2021=738%:%
%:%2022=739%:%
%:%2023=739%:%
%:%2024=740%:%
%:%2025=740%:%
%:%2026=740%:%
%:%2027=741%:%
%:%2028=741%:%
%:%2029=742%:%
%:%2030=742%:%
%:%2031=743%:%
%:%2032=743%:%
%:%2033=744%:%
%:%2034=744%:%
%:%2035=744%:%
%:%2036=745%:%
%:%2037=745%:%
%:%2038=746%:%
%:%2039=746%:%
%:%2040=747%:%
%:%2041=747%:%
%:%2042=748%:%
%:%2043=748%:%
%:%2044=748%:%
%:%2045=749%:%
%:%2046=749%:%
%:%2047=750%:%
%:%2048=750%:%
%:%2049=751%:%
%:%2050=751%:%
%:%2051=752%:%
%:%2057=752%:%
%:%2060=753%:%
%:%2061=754%:%
%:%2062=754%:%
%:%2063=755%:%
%:%2064=756%:%
%:%2065=757%:%
%:%2068=758%:%
%:%2072=758%:%
%:%2073=758%:%
%:%2074=759%:%
%:%2075=759%:%
%:%2076=760%:%
%:%2077=760%:%
%:%2078=760%:%
%:%2079=761%:%
%:%2080=761%:%
%:%2081=762%:%
%:%2082=762%:%
%:%2083=762%:%
%:%2084=762%:%
%:%2085=763%:%
%:%2086=763%:%
%:%2087=763%:%
%:%2088=764%:%
%:%2089=764%:%
%:%2090=765%:%
%:%2091=765%:%
%:%2092=765%:%
%:%2093=766%:%
%:%2094=766%:%
%:%2095=767%:%
%:%2096=767%:%
%:%2097=768%:%
%:%2098=768%:%
%:%2099=768%:%
%:%2100=769%:%
%:%2101=769%:%
%:%2102=770%:%
%:%2103=770%:%
%:%2104=770%:%
%:%2105=771%:%
%:%2106=771%:%
%:%2107=771%:%
%:%2108=772%:%
%:%2109=772%:%
%:%2110=772%:%
%:%2111=773%:%
%:%2112=773%:%
%:%2113=774%:%
%:%2114=774%:%
%:%2115=775%:%
%:%2116=775%:%
%:%2117=775%:%
%:%2118=776%:%
%:%2119=776%:%
%:%2120=777%:%
%:%2121=777%:%
%:%2122=778%:%
%:%2132=780%:%
%:%2134=781%:%
%:%2135=781%:%
%:%2136=782%:%
%:%2139=783%:%
%:%2143=783%:%
%:%2144=783%:%
%:%2145=784%:%
%:%2146=784%:%
%:%2147=785%:%
%:%2148=785%:%
%:%2149=785%:%
%:%2150=786%:%
%:%2151=786%:%
%:%2152=786%:%
%:%2153=787%:%
%:%2154=787%:%
%:%2155=788%:%
%:%2156=789%:%
%:%2157=789%:%
%:%2158=790%:%
%:%2168=792%:%
%:%2170=793%:%
%:%2171=793%:%
%:%2172=794%:%
%:%2173=795%:%
%:%2174=796%:%
%:%2177=797%:%
%:%2181=797%:%
%:%2182=797%:%
%:%2183=798%:%
%:%2184=798%:%
%:%2185=799%:%
%:%2186=799%:%
%:%2195=801%:%
%:%2196=802%:%
%:%2198=803%:%
%:%2199=803%:%
%:%2200=804%:%
%:%2201=805%:%
%:%2202=806%:%
%:%2203=807%:%
%:%2206=808%:%
%:%2210=808%:%
%:%2211=808%:%
%:%2212=809%:%
%:%2213=809%:%
%:%2214=810%:%
%:%2215=810%:%
%:%2216=811%:%
%:%2217=811%:%
%:%2218=811%:%
%:%2219=812%:%
%:%2220=812%:%
%:%2221=812%:%
%:%2222=813%:%
%:%2223=813%:%
%:%2224=814%:%
%:%2225=814%:%
%:%2226=815%:%
%:%2227=815%:%
%:%2228=816%:%
%:%2229=816%:%
%:%2230=817%:%
%:%2231=817%:%
%:%2232=818%:%
%:%2233=818%:%
%:%2234=819%:%
%:%2235=819%:%
%:%2236=820%:%
%:%2237=820%:%
%:%2238=821%:%
%:%2239=821%:%
%:%2240=822%:%
%:%2241=822%:%
%:%2242=823%:%
%:%2243=823%:%
%:%2244=824%:%
%:%2245=824%:%
%:%2246=824%:%
%:%2247=825%:%
%:%2248=825%:%
%:%2249=826%:%
%:%2250=826%:%
%:%2251=827%:%
%:%2252=827%:%
%:%2253=827%:%
%:%2254=828%:%
%:%2255=828%:%
%:%2256=828%:%
%:%2257=829%:%
%:%2258=829%:%
%:%2259=830%:%
%:%2260=830%:%
%:%2261=831%:%
%:%2262=831%:%
%:%2263=831%:%
%:%2264=832%:%
%:%2265=832%:%
%:%2266=832%:%
%:%2267=833%:%
%:%2268=833%:%
%:%2269=834%:%
%:%2270=834%:%
%:%2271=835%:%
%:%2272=835%:%
%:%2273=835%:%
%:%2274=836%:%
%:%2275=836%:%
%:%2276=836%:%
%:%2277=837%:%
%:%2278=837%:%
%:%2279=838%:%
%:%2280=839%:%
%:%2281=839%:%
%:%2282=839%:%
%:%2283=840%:%
%:%2284=840%:%
%:%2285=841%:%
%:%2286=841%:%
%:%2287=841%:%
%:%2288=842%:%
%:%2289=842%:%
%:%2290=843%:%
%:%2291=843%:%
%:%2292=843%:%
%:%2293=844%:%
%:%2294=844%:%
%:%2295=845%:%
%:%2296=846%:%
%:%2297=846%:%
%:%2298=847%:%
%:%2299=847%:%
%:%2300=847%:%
%:%2301=848%:%
%:%2302=848%:%
%:%2303=849%:%
%:%2304=849%:%
%:%2305=850%:%
%:%2306=850%:%
%:%2307=850%:%
%:%2308=851%:%
%:%2309=851%:%
%:%2310=852%:%
%:%2311=853%:%
%:%2312=853%:%
%:%2313=854%:%
%:%2314=854%:%
%:%2315=855%:%
%:%2325=857%:%
%:%2327=858%:%
%:%2328=858%:%
%:%2329=859%:%
%:%2330=860%:%
%:%2335=865%:%
%:%2337=867%:%
%:%2338=867%:%
%:%2339=868%:%
%:%2342=869%:%
%:%2346=869%:%
%:%2347=869%:%
%:%2356=871%:%
%:%2358=872%:%
%:%2359=872%:%
%:%2362=873%:%
%:%2366=873%:%
%:%2367=873%:%
%:%2368=874%:%
%:%2369=874%:%
%:%2374=874%:%
%:%2377=875%:%
%:%2378=876%:%
%:%2379=876%:%
%:%2380=877%:%
%:%2381=878%:%
%:%2384=879%:%
%:%2388=879%:%
%:%2389=879%:%
%:%2390=879%:%
%:%2395=879%:%
%:%2398=880%:%
%:%2399=881%:%
%:%2400=881%:%
%:%2401=882%:%
%:%2404=883%:%
%:%2408=883%:%
%:%2409=883%:%
%:%2414=883%:%
%:%2417=884%:%
%:%2418=885%:%
%:%2419=885%:%
%:%2420=886%:%
%:%2423=887%:%
%:%2427=887%:%
%:%2428=887%:%
%:%2429=888%:%
%:%2430=888%:%
%:%2435=888%:%
%:%2438=889%:%
%:%2439=890%:%
%:%2440=891%:%
%:%2441=891%:%
%:%2444=892%:%
%:%2448=892%:%
%:%2449=892%:%
%:%2454=892%:%
%:%2457=893%:%
%:%2458=894%:%
%:%2459=894%:%
%:%2460=895%:%
%:%2463=896%:%
%:%2467=896%:%
%:%2468=896%:%
%:%2469=897%:%
%:%2470=897%:%
%:%2479=899%:%
%:%2481=900%:%
%:%2482=900%:%
%:%2483=901%:%
%:%2484=902%:%
%:%2485=903%:%
%:%2486=904%:%
%:%2487=904%:%
%:%2488=905%:%
%:%2489=906%:%
%:%2490=907%:%
%:%2491=908%:%
%:%2492=909%:%
%:%2493=909%:%
%:%2494=910%:%
%:%2497=911%:%
%:%2501=911%:%
%:%2502=911%:%
%:%2511=913%:%
%:%2513=915%:%
%:%2514=915%:%
%:%2515=916%:%
%:%2518=917%:%
%:%2522=917%:%
%:%2523=917%:%
%:%2528=917%:%
%:%2531=918%:%
%:%2532=919%:%
%:%2533=919%:%
%:%2534=920%:%
%:%2535=921%:%
%:%2538=922%:%
%:%2542=922%:%
%:%2543=922%:%
%:%2544=923%:%
%:%2545=923%:%
%:%2546=924%:%
%:%2547=924%:%
%:%2548=925%:%
%:%2549=925%:%
%:%2550=925%:%
%:%2551=925%:%
%:%2552=926%:%
%:%2553=926%:%
%:%2554=927%:%
%:%2555=927%:%
%:%2556=928%:%
%:%2557=928%:%
%:%2558=928%:%
%:%2559=929%:%
%:%2560=929%:%
%:%2561=930%:%
%:%2562=930%:%
%:%2563=931%:%
%:%2564=931%:%
%:%2565=931%:%
%:%2566=932%:%
%:%2567=932%:%
%:%2568=933%:%
%:%2569=933%:%
%:%2570=934%:%
%:%2571=934%:%
%:%2572=934%:%
%:%2573=935%:%
%:%2574=935%:%
%:%2575=936%:%
%:%2576=936%:%
%:%2577=937%:%
%:%2578=937%:%
%:%2579=938%:%
%:%2585=938%:%
%:%2588=939%:%
%:%2589=940%:%
%:%2590=940%:%
%:%2591=941%:%
%:%2594=942%:%
%:%2598=942%:%
%:%2599=942%:%
%:%2600=943%:%
%:%2601=943%:%
%:%2602=944%:%
%:%2603=944%:%
%:%2604=944%:%
%:%2605=945%:%
%:%2606=945%:%
%:%2607=946%:%
%:%2608=946%:%
%:%2609=947%:%
%:%2610=947%:%
%:%2611=948%:%
%:%2612=948%:%
%:%2613=949%:%
%:%2614=949%:%
%:%2615=950%:%
%:%2616=950%:%
%:%2617=950%:%
%:%2618=951%:%
%:%2619=951%:%
%:%2620=952%:%
%:%2621=952%:%
%:%2622=953%:%
%:%2623=953%:%
%:%2624=953%:%
%:%2625=954%:%
%:%2626=954%:%
%:%2627=954%:%
%:%2628=955%:%
%:%2629=955%:%
%:%2630=956%:%
%:%2631=956%:%
%:%2632=957%:%
%:%2633=957%:%
%:%2634=957%:%
%:%2635=958%:%
%:%2636=958%:%
%:%2637=958%:%
%:%2638=959%:%
%:%2639=959%:%
%:%2640=959%:%
%:%2641=960%:%
%:%2642=960%:%
%:%2643=960%:%
%:%2644=961%:%
%:%2645=961%:%
%:%2646=961%:%
%:%2647=962%:%
%:%2648=963%:%
%:%2649=963%:%
%:%2650=963%:%
%:%2651=964%:%
%:%2652=964%:%
%:%2653=964%:%
%:%2654=965%:%
%:%2655=965%:%
%:%2656=965%:%
%:%2657=966%:%
%:%2658=966%:%
%:%2659=967%:%
%:%2665=967%:%
%:%2668=968%:%
%:%2669=969%:%
%:%2670=969%:%
%:%2671=970%:%
%:%2672=971%:%
%:%2675=972%:%
%:%2679=972%:%
%:%2680=972%:%
%:%2681=973%:%
%:%2682=973%:%
%:%2683=974%:%
%:%2684=974%:%
%:%2685=974%:%
%:%2686=975%:%
%:%2687=975%:%
%:%2688=975%:%
%:%2689=976%:%
%:%2690=976%:%
%:%2691=977%:%
%:%2692=977%:%
%:%2693=978%:%
%:%2694=978%:%
%:%2695=979%:%
%:%2696=979%:%
%:%2697=979%:%
%:%2698=980%:%
%:%2699=980%:%
%:%2700=980%:%
%:%2701=981%:%
%:%2702=982%:%
%:%2703=982%:%
%:%2704=982%:%
%:%2705=983%:%
%:%2706=983%:%
%:%2707=983%:%
%:%2708=984%:%
%:%2709=984%:%
%:%2710=984%:%
%:%2711=985%:%
%:%2717=985%:%
%:%2720=986%:%
%:%2721=987%:%
%:%2722=987%:%
%:%2723=988%:%
%:%2726=989%:%
%:%2730=989%:%
%:%2731=989%:%
%:%2732=989%:%
%:%2737=989%:%
%:%2740=990%:%
%:%2741=991%:%
%:%2742=991%:%
%:%2743=992%:%
%:%2744=993%:%
%:%2747=994%:%
%:%2751=994%:%
%:%2752=994%:%
%:%2753=995%:%
%:%2754=995%:%
%:%2755=996%:%
%:%2756=996%:%
%:%2757=997%:%
%:%2758=997%:%
%:%2759=998%:%
%:%2760=998%:%
%:%2761=998%:%
%:%2762=999%:%
%:%2763=999%:%
%:%2764=1000%:%
%:%2765=1000%:%
%:%2766=1000%:%
%:%2767=1001%:%
%:%2768=1002%:%
%:%2769=1002%:%
%:%2770=1002%:%
%:%2771=1003%:%
%:%2772=1003%:%
%:%2773=1003%:%
%:%2774=1004%:%
%:%2775=1005%:%
%:%2776=1005%:%
%:%2777=1006%:%
%:%2778=1007%:%
%:%2779=1007%:%
%:%2780=1008%:%
%:%2781=1008%:%
%:%2782=1009%:%
%:%2783=1009%:%
%:%2784=1010%:%
%:%2785=1011%:%
%:%2786=1011%:%
%:%2787=1012%:%
%:%2793=1012%:%
%:%2796=1013%:%
%:%2797=1014%:%
%:%2798=1015%:%
%:%2799=1015%:%
%:%2800=1016%:%
%:%2803=1017%:%
%:%2807=1017%:%
%:%2808=1017%:%
%:%2809=1017%:%
%:%2818=1020%:%
%:%2820=1021%:%
%:%2821=1021%:%
%:%2822=1022%:%
%:%2823=1023%:%
%:%2826=1026%:%
%:%2828=1028%:%
%:%2829=1028%:%
%:%2830=1029%:%
%:%2831=1030%:%
%:%2832=1031%:%
%:%2833=1031%:%
%:%2834=1032%:%
%:%2836=1034%:%
%:%2838=1036%:%
%:%2839=1036%:%
%:%2842=1037%:%
%:%2846=1037%:%
%:%2847=1037%:%
%:%2848=1037%:%
%:%2853=1037%:%
%:%2856=1038%:%
%:%2857=1039%:%
%:%2858=1039%:%
%:%2861=1040%:%
%:%2865=1040%:%
%:%2866=1040%:%
%:%2867=1040%:%
%:%2872=1040%:%
%:%2875=1041%:%
%:%2876=1042%:%
%:%2877=1042%:%
%:%2880=1043%:%
%:%2884=1043%:%
%:%2885=1043%:%
%:%2886=1044%:%
%:%2887=1044%:%
%:%2892=1044%:%
%:%2895=1045%:%
%:%2896=1046%:%
%:%2897=1046%:%
%:%2900=1047%:%
%:%2904=1047%:%
%:%2905=1047%:%
%:%2906=1048%:%
%:%2907=1048%:%
%:%2908=1049%:%
%:%2909=1049%:%
%:%2910=1050%:%
%:%2911=1050%:%
%:%2912=1050%:%
%:%2913=1051%:%
%:%2914=1051%:%
%:%2915=1052%:%
%:%2916=1052%:%
%:%2917=1053%:%
%:%2918=1053%:%
%:%2919=1053%:%
%:%2920=1054%:%
%:%2921=1054%:%
%:%2922=1054%:%
%:%2923=1055%:%
%:%2924=1055%:%
%:%2925=1056%:%
%:%2926=1056%:%
%:%2927=1057%:%
%:%2928=1057%:%
%:%2929=1057%:%
%:%2930=1058%:%
%:%2931=1058%:%
%:%2932=1059%:%
%:%2933=1059%:%
%:%2934=1059%:%
%:%2935=1060%:%
%:%2936=1060%:%
%:%2937=1061%:%
%:%2943=1061%:%
%:%2946=1062%:%
%:%2947=1063%:%
%:%2948=1063%:%
%:%2951=1064%:%
%:%2955=1064%:%
%:%2956=1064%:%
%:%2957=1065%:%
%:%2958=1065%:%
%:%2959=1066%:%
%:%2960=1066%:%
%:%2961=1066%:%
%:%2962=1067%:%
%:%2963=1067%:%
%:%2964=1067%:%
%:%2965=1068%:%
%:%2966=1068%:%
%:%2967=1069%:%
%:%2968=1069%:%
%:%2969=1070%:%
%:%2970=1070%:%
%:%2971=1071%:%
%:%2972=1071%:%
%:%2973=1072%:%
%:%2974=1072%:%
%:%2975=1072%:%
%:%2976=1073%:%
%:%2977=1073%:%
%:%2978=1073%:%
%:%2979=1074%:%
%:%2980=1074%:%
%:%2981=1074%:%
%:%2982=1075%:%
%:%2983=1075%:%
%:%2984=1075%:%
%:%2985=1076%:%
%:%2995=1078%:%
%:%2997=1079%:%
%:%2998=1079%:%
%:%3001=1080%:%
%:%3005=1080%:%
%:%3006=1080%:%
%:%3007=1081%:%
%:%3008=1081%:%
%:%3013=1081%:%
%:%3016=1082%:%
%:%3017=1083%:%
%:%3018=1083%:%
%:%3021=1084%:%
%:%3025=1084%:%
%:%3026=1084%:%
%:%3027=1085%:%
%:%3028=1085%:%
%:%3029=1086%:%
%:%3030=1086%:%
%:%3031=1087%:%
%:%3032=1087%:%
%:%3033=1087%:%
%:%3034=1088%:%
%:%3044=1090%:%
%:%3048=1092%:%
%:%3050=1093%:%
%:%3051=1093%:%
%:%3052=1094%:%
%:%3053=1095%:%
%:%3054=1096%:%
%:%3057=1097%:%
%:%3061=1097%:%
%:%3062=1097%:%
%:%3063=1098%:%
%:%3064=1098%:%
%:%3065=1099%:%
%:%3066=1099%:%
%:%3067=1099%:%
%:%3068=1099%:%
%:%3069=1100%:%
%:%3070=1100%:%
%:%3071=1101%:%
%:%3072=1101%:%
%:%3073=1102%:%
%:%3074=1102%:%
%:%3075=1103%:%
%:%3076=1103%:%
%:%3077=1104%:%
%:%3078=1104%:%
%:%3079=1104%:%
%:%3080=1105%:%
%:%3081=1105%:%
%:%3082=1106%:%
%:%3083=1106%:%
%:%3084=1107%:%
%:%3094=1109%:%
%:%3096=1110%:%
%:%3097=1110%:%
%:%3100=1111%:%
%:%3104=1111%:%
%:%3105=1111%:%
%:%3114=1113%:%
%:%3116=1114%:%
%:%3117=1114%:%
%:%3118=1115%:%
%:%3119=1116%:%
%:%3120=1117%:%
%:%3123=1118%:%
%:%3127=1118%:%
%:%3128=1118%:%
%:%3129=1119%:%
%:%3130=1119%:%
%:%3131=1120%:%
%:%3132=1120%:%
%:%3133=1121%:%
%:%3134=1121%:%
%:%3135=1121%:%
%:%3136=1122%:%
%:%3137=1122%:%
%:%3138=1122%:%
%:%3139=1123%:%
%:%3140=1123%:%
%:%3141=1124%:%
%:%3142=1124%:%
%:%3143=1125%:%
%:%3144=1125%:%
%:%3145=1126%:%
%:%3146=1126%:%
%:%3147=1127%:%
%:%3148=1127%:%
%:%3149=1128%:%
%:%3150=1128%:%
%:%3151=1128%:%
%:%3152=1129%:%
%:%3153=1129%:%
%:%3154=1130%:%
%:%3155=1130%:%
%:%3156=1131%:%
%:%3157=1132%:%
%:%3158=1133%:%
%:%3159=1134%:%
%:%3160=1134%:%
%:%3161=1135%:%
%:%3162=1135%:%
%:%3163=1136%:%
%:%3164=1136%:%
%:%3165=1137%:%
%:%3166=1137%:%
%:%3167=1138%:%
%:%3168=1138%:%
%:%3169=1139%:%
%:%3170=1139%:%
%:%3171=1140%:%
%:%3172=1140%:%
%:%3173=1141%:%
%:%3174=1142%:%
%:%3175=1143%:%
%:%3176=1143%:%
%:%3177=1144%:%
%:%3178=1144%:%
%:%3179=1145%:%
%:%3180=1145%:%
%:%3181=1146%:%
%:%3182=1147%:%
%:%3183=1147%:%
%:%3184=1147%:%
%:%3185=1148%:%
%:%3191=1148%:%
%:%3194=1149%:%
%:%3195=1150%:%
%:%3196=1150%:%
%:%3197=1151%:%
%:%3198=1152%:%
%:%3199=1153%:%
%:%3200=1154%:%
%:%3203=1155%:%
%:%3207=1155%:%
%:%3208=1155%:%
%:%3209=1156%:%
%:%3210=1156%:%
%:%3211=1157%:%
%:%3212=1157%:%
%:%3213=1158%:%
%:%3214=1158%:%
%:%3215=1159%:%
%:%3216=1159%:%
%:%3217=1160%:%
%:%3227=1162%:%
%:%3229=1163%:%
%:%3230=1163%:%
%:%3231=1164%:%
%:%3232=1165%:%
%:%3233=1166%:%
%:%3234=1167%:%
%:%3235=1168%:%
%:%3238=1169%:%
%:%3242=1169%:%
%:%3243=1169%:%
%:%3244=1170%:%
%:%3245=1170%:%
%:%3246=1171%:%
%:%3247=1171%:%
%:%3248=1172%:%
%:%3249=1172%:%
%:%3250=1172%:%
%:%3251=1173%:%
%:%3252=1173%:%
%:%3253=1173%:%
%:%3254=1174%:%
%:%3255=1174%:%
%:%3256=1175%:%
%:%3257=1175%:%
%:%3258=1176%:%
%:%3259=1176%:%
%:%3260=1177%:%
%:%3261=1177%:%
%:%3262=1178%:%
%:%3263=1178%:%
%:%3264=1179%:%
%:%3265=1179%:%
%:%3266=1179%:%
%:%3267=1180%:%
%:%3268=1180%:%
%:%3269=1181%:%
%:%3270=1181%:%
%:%3271=1182%:%
%:%3272=1183%:%
%:%3273=1184%:%
%:%3274=1184%:%
%:%3275=1185%:%
%:%3276=1185%:%
%:%3277=1186%:%
%:%3278=1186%:%
%:%3279=1187%:%
%:%3280=1188%:%
%:%3281=1189%:%
%:%3282=1189%:%
%:%3283=1190%:%
%:%3284=1190%:%
%:%3285=1191%:%
%:%3286=1191%:%
%:%3287=1192%:%
%:%3288=1193%:%
%:%3289=1193%:%
%:%3290=1193%:%
%:%3291=1194%:%
%:%3297=1194%:%
%:%3300=1195%:%
%:%3301=1196%:%
%:%3302=1196%:%
%:%3303=1197%:%
%:%3304=1198%:%
%:%3305=1199%:%
%:%3306=1200%:%
%:%3309=1201%:%
%:%3313=1201%:%
%:%3314=1201%:%
%:%3315=1202%:%
%:%3316=1202%:%
%:%3317=1203%:%
%:%3318=1203%:%
%:%3319=1204%:%
%:%3320=1204%:%
%:%3321=1205%:%
%:%3322=1205%:%
%:%3323=1206%:%
%:%3333=1208%:%
%:%3335=1209%:%
%:%3336=1209%:%
%:%3337=1210%:%
%:%3338=1211%:%
%:%3339=1212%:%
%:%3340=1213%:%
%:%3341=1214%:%
%:%3344=1215%:%
%:%3348=1215%:%
%:%3349=1215%:%
%:%3350=1216%:%
%:%3351=1216%:%
%:%3352=1217%:%
%:%3353=1217%:%
%:%3354=1218%:%
%:%3355=1218%:%
%:%3356=1218%:%
%:%3357=1219%:%
%:%3358=1219%:%
%:%3359=1219%:%
%:%3360=1220%:%
%:%3361=1220%:%
%:%3362=1221%:%
%:%3363=1221%:%
%:%3364=1222%:%
%:%3365=1222%:%
%:%3366=1223%:%
%:%3367=1223%:%
%:%3368=1224%:%
%:%3369=1224%:%
%:%3370=1225%:%
%:%3371=1225%:%
%:%3372=1225%:%
%:%3373=1226%:%
%:%3374=1226%:%
%:%3375=1227%:%
%:%3376=1227%:%
%:%3377=1228%:%
%:%3378=1229%:%
%:%3379=1230%:%
%:%3380=1230%:%
%:%3381=1231%:%
%:%3382=1231%:%
%:%3383=1232%:%
%:%3384=1232%:%
%:%3385=1233%:%
%:%3386=1234%:%
%:%3387=1235%:%
%:%3388=1235%:%
%:%3389=1236%:%
%:%3390=1236%:%
%:%3391=1237%:%
%:%3392=1237%:%
%:%3393=1238%:%
%:%3394=1239%:%
%:%3395=1239%:%
%:%3396=1239%:%
%:%3397=1240%:%
%:%3403=1240%:%
%:%3406=1241%:%
%:%3407=1242%:%
%:%3408=1243%:%
%:%3409=1243%:%
%:%3412=1244%:%
%:%3416=1244%:%
%:%3417=1244%:%
%:%3422=1244%:%
%:%3425=1245%:%
%:%3426=1246%:%
%:%3427=1246%:%
%:%3428=1247%:%
%:%3429=1248%:%
%:%3430=1249%:%
%:%3431=1250%:%
%:%3432=1251%:%
%:%3433=1252%:%
%:%3436=1253%:%
%:%3440=1253%:%
%:%3441=1253%:%
%:%3442=1254%:%
%:%3443=1254%:%
%:%3444=1255%:%
%:%3445=1255%:%
%:%3446=1256%:%
%:%3447=1256%:%
%:%3448=1257%:%
%:%3449=1257%:%
%:%3450=1257%:%
%:%3451=1258%:%
%:%3452=1258%:%
%:%3453=1259%:%
%:%3454=1259%:%
%:%3455=1260%:%
%:%3456=1260%:%
%:%3457=1261%:%
%:%3458=1261%:%
%:%3459=1262%:%
%:%3460=1262%:%
%:%3461=1263%:%
%:%3462=1263%:%
%:%3463=1264%:%
%:%3464=1264%:%
%:%3465=1264%:%
%:%3466=1265%:%
%:%3467=1265%:%
%:%3468=1266%:%
%:%3469=1266%:%
%:%3470=1267%:%
%:%3471=1267%:%
%:%3472=1267%:%
%:%3473=1268%:%
%:%3474=1268%:%
%:%3475=1268%:%
%:%3476=1269%:%
%:%3477=1269%:%
%:%3478=1270%:%
%:%3479=1270%:%
%:%3480=1271%:%
%:%3481=1271%:%
%:%3482=1272%:%
%:%3483=1272%:%
%:%3484=1273%:%
%:%3485=1273%:%
%:%3486=1274%:%
%:%3487=1274%:%
%:%3488=1275%:%
%:%3489=1275%:%
%:%3490=1276%:%
%:%3491=1276%:%
%:%3492=1277%:%
%:%3493=1277%:%
%:%3494=1278%:%
%:%3495=1278%:%
%:%3496=1279%:%
%:%3497=1279%:%
%:%3498=1280%:%
%:%3499=1280%:%
%:%3500=1280%:%
%:%3501=1281%:%
%:%3502=1282%:%
%:%3503=1282%:%
%:%3504=1283%:%
%:%3505=1283%:%
%:%3506=1283%:%
%:%3509=1286%:%
%:%3510=1287%:%
%:%3511=1287%:%
%:%3512=1288%:%
%:%3513=1289%:%
%:%3514=1289%:%
%:%3515=1290%:%
%:%3516=1290%:%
%:%3517=1290%:%
%:%3519=1292%:%
%:%3520=1293%:%
%:%3521=1293%:%
%:%3522=1294%:%
%:%3523=1294%:%
%:%3524=1294%:%
%:%3525=1295%:%
%:%3526=1295%:%
%:%3527=1295%:%
%:%3528=1296%:%
%:%3529=1296%:%
%:%3530=1297%:%
%:%3540=1299%:%
%:%3542=1300%:%
%:%3543=1300%:%
%:%3544=1301%:%
%:%3545=1302%:%
%:%3546=1303%:%
%:%3547=1304%:%
%:%3548=1305%:%
%:%3549=1306%:%
%:%3550=1307%:%
%:%3551=1308%:%
%:%3552=1309%:%
%:%3552=1310%:%
%:%3555=1311%:%
%:%3559=1311%:%
%:%3560=1311%:%
%:%3561=1312%:%
%:%3562=1312%:%
%:%3563=1313%:%
%:%3564=1313%:%
%:%3565=1314%:%
%:%3566=1314%:%
%:%3567=1314%:%
%:%3568=1315%:%
%:%3569=1315%:%
%:%3570=1316%:%
%:%3571=1316%:%
%:%3572=1316%:%
%:%3573=1317%:%
%:%3574=1317%:%
%:%3575=1318%:%
%:%3576=1318%:%
%:%3577=1319%:%
%:%3578=1319%:%
%:%3579=1319%:%
%:%3580=1320%:%
%:%3581=1320%:%
%:%3582=1320%:%
%:%3583=1321%:%
%:%3584=1321%:%
%:%3585=1321%:%
%:%3586=1321%:%
%:%3587=1322%:%
%:%3588=1322%:%
%:%3589=1323%:%
%:%3590=1323%:%
%:%3591=1324%:%
%:%3592=1324%:%
%:%3593=1324%:%
%:%3594=1325%:%
%:%3595=1325%:%
%:%3596=1325%:%
%:%3597=1326%:%
%:%3598=1326%:%
%:%3599=1327%:%
%:%3600=1328%:%
%:%3601=1328%:%
%:%3602=1329%:%
%:%3603=1330%:%
%:%3604=1331%:%
%:%3605=1332%:%
%:%3606=1332%:%
%:%3607=1333%:%
%:%3608=1333%:%
%:%3609=1333%:%
%:%3610=1334%:%
%:%3611=1335%:%
%:%3612=1335%:%
%:%3613=1335%:%
%:%3614=1336%:%
%:%3615=1336%:%
%:%3616=1336%:%
%:%3617=1337%:%
%:%3618=1337%:%
%:%3619=1337%:%
%:%3620=1337%:%
%:%3621=1338%:%
%:%3631=1340%:%
%:%3633=1341%:%
%:%3634=1341%:%
%:%3635=1342%:%
%:%3636=1343%:%
%:%3637=1344%:%
%:%3638=1345%:%
%:%3639=1345%:%
%:%3640=1346%:%
%:%3643=1347%:%
%:%3647=1347%:%
%:%3648=1347%:%
%:%3653=1347%:%
%:%3656=1348%:%
%:%3657=1349%:%
%:%3658=1349%:%
%:%3659=1350%:%
%:%3662=1351%:%
%:%3666=1351%:%
%:%3667=1351%:%
%:%3672=1351%:%
%:%3675=1352%:%
%:%3676=1353%:%
%:%3677=1353%:%
%:%3678=1354%:%
%:%3681=1355%:%
%:%3685=1355%:%
%:%3686=1355%:%
%:%3687=1356%:%
%:%3688=1356%:%
%:%3693=1356%:%
%:%3696=1357%:%
%:%3697=1358%:%
%:%3698=1358%:%
%:%3701=1359%:%
%:%3705=1359%:%
%:%3706=1359%:%
%:%3707=1360%:%
%:%3708=1360%:%
%:%3709=1361%:%
%:%3710=1361%:%
%:%3711=1361%:%
%:%3712=1362%:%
%:%3713=1362%:%
%:%3714=1363%:%
%:%3715=1363%:%
%:%3716=1364%:%
%:%3717=1364%:%
%:%3718=1365%:%
%:%3719=1365%:%
%:%3720=1365%:%
%:%3721=1366%:%
%:%3722=1366%:%
%:%3723=1367%:%
%:%3729=1367%:%
%:%3732=1368%:%
%:%3733=1369%:%
%:%3734=1369%:%
%:%3735=1370%:%
%:%3736=1371%:%
%:%3737=1372%:%
%:%3740=1373%:%
%:%3744=1373%:%
%:%3745=1373%:%
%:%3746=1374%:%
%:%3747=1374%:%
%:%3756=1376%:%
%:%3758=1377%:%
%:%3759=1377%:%
%:%3760=1378%:%
%:%3761=1379%:%
%:%3762=1380%:%
%:%3763=1381%:%
%:%3764=1382%:%
%:%3765=1383%:%
%:%3766=1384%:%
%:%3769=1385%:%
%:%3773=1385%:%
%:%3774=1385%:%
%:%3775=1386%:%
%:%3776=1386%:%
%:%3777=1387%:%
%:%3778=1387%:%
%:%3779=1387%:%
%:%3780=1388%:%
%:%3781=1388%:%
%:%3782=1389%:%
%:%3783=1389%:%
%:%3784=1390%:%
%:%3785=1391%:%
%:%3786=1391%:%
%:%3787=1392%:%
%:%3788=1392%:%
%:%3789=1392%:%
%:%3790=1393%:%
%:%3791=1393%:%
%:%3792=1394%:%
%:%3793=1394%:%
%:%3794=1395%:%
%:%3795=1395%:%
%:%3796=1396%:%
%:%3797=1396%:%
%:%3798=1397%:%
%:%3808=1399%:%
%:%3810=1400%:%
%:%3811=1400%:%
%:%3812=1401%:%
%:%3813=1402%:%
%:%3814=1403%:%
%:%3815=1404%:%
%:%3818=1405%:%
%:%3822=1405%:%
%:%3823=1405%:%
%:%3824=1406%:%
%:%3825=1406%:%
%:%3826=1407%:%
%:%3827=1408%:%
%:%3828=1409%:%
%:%3829=1410%:%
%:%3830=1411%:%
%:%3831=1412%:%
%:%3832=1413%:%
%:%3833=1414%:%
%:%3834=1414%:%
%:%3839=1414%:%
%:%3842=1415%:%
%:%3843=1416%:%
%:%3844=1417%:%
%:%3845=1417%:%
%:%3846=1418%:%
%:%3847=1419%:%
%:%3850=1420%:%
%:%3854=1420%:%
%:%3855=1420%:%
%:%3856=1420%:%
%:%3865=1433%:%
%:%3866=1434%:%
%:%3868=1435%:%
%:%3869=1435%:%
%:%3870=1436%:%
%:%3871=1437%:%
%:%3874=1438%:%
%:%3878=1438%:%
%:%3879=1438%:%
%:%3880=1438%:%
%:%3885=1438%:%
%:%3888=1439%:%
%:%3889=1440%:%
%:%3890=1440%:%
%:%3891=1441%:%
%:%3892=1442%:%
%:%3893=1443%:%
%:%3894=1444%:%
%:%3895=1445%:%
%:%3896=1446%:%
%:%3897=1447%:%
%:%3900=1448%:%
%:%3904=1448%:%
%:%3905=1448%:%
%:%3906=1449%:%
%:%3907=1449%:%
%:%3908=1450%:%
%:%3909=1450%:%
%:%3910=1450%:%
%:%3911=1451%:%
%:%3912=1451%:%
%:%3913=1451%:%
%:%3914=1452%:%
%:%3915=1452%:%
%:%3916=1452%:%
%:%3917=1453%:%
%:%3923=1453%:%
%:%3926=1454%:%
%:%3927=1455%:%
%:%3928=1455%:%
%:%3929=1456%:%
%:%3930=1457%:%
%:%3931=1458%:%
%:%3932=1459%:%
%:%3933=1460%:%
%:%3936=1461%:%
%:%3940=1461%:%
%:%3941=1461%:%
%:%3942=1461%:%
%:%3951=1463%:%
%:%3953=1464%:%
%:%3954=1464%:%
%:%3955=1465%:%
%:%3956=1466%:%
%:%3959=1469%:%
%:%3961=1470%:%
%:%3962=1470%:%
%:%3963=1471%:%
%:%3964=1472%:%
%:%3967=1473%:%
%:%3971=1473%:%
%:%3972=1473%:%
%:%3973=1474%:%
%:%3974=1474%:%
%:%3979=1474%:%
%:%3982=1475%:%
%:%3983=1476%:%
%:%3984=1476%:%
%:%3985=1477%:%
%:%3986=1478%:%
%:%3987=1479%:%
%:%3988=1479%:%
%:%3991=1480%:%
%:%3995=1480%:%
%:%3996=1480%:%
%:%3997=1480%:%
%:%4002=1480%:%
%:%4005=1481%:%
%:%4006=1482%:%
%:%4007=1482%:%
%:%4008=1483%:%
%:%4009=1484%:%
%:%4013=1488%:%
%:%4014=1489%:%
%:%4021=1496%:%
%:%4022=1497%:%
%:%4023=1498%:%
%:%4024=1498%:%
%:%4025=1499%:%
%:%4028=1500%:%
%:%4032=1500%:%
%:%4033=1500%:%
%:%4038=1500%:%
%:%4041=1501%:%
%:%4042=1502%:%
%:%4043=1502%:%
%:%4044=1503%:%
%:%4047=1506%:%
%:%4048=1507%:%
%:%4049=1508%:%
%:%4050=1509%:%
%:%4051=1509%:%
%:%4052=1510%:%
%:%4053=1511%:%
%:%4056=1512%:%
%:%4060=1512%:%
%:%4061=1512%:%
%:%4062=1513%:%
%:%4063=1513%:%
%:%4064=1514%:%
%:%4065=1514%:%
%:%4066=1515%:%
%:%4067=1515%:%
%:%4068=1516%:%
%:%4069=1516%:%
%:%4070=1517%:%
%:%4071=1517%:%
%:%4072=1518%:%
%:%4073=1518%:%
%:%4074=1518%:%
%:%4075=1519%:%
%:%4076=1519%:%
%:%4077=1520%:%
%:%4078=1520%:%
%:%4079=1521%:%
%:%4080=1521%:%
%:%4081=1521%:%
%:%4082=1522%:%
%:%4083=1522%:%
%:%4084=1522%:%
%:%4085=1523%:%
%:%4086=1523%:%
%:%4087=1524%:%
%:%4088=1524%:%
%:%4089=1525%:%
%:%4090=1525%:%
%:%4091=1525%:%
%:%4092=1526%:%
%:%4093=1526%:%
%:%4094=1526%:%
%:%4095=1527%:%
%:%4096=1527%:%
%:%4097=1528%:%
%:%4098=1528%:%
%:%4099=1529%:%
%:%4100=1529%:%
%:%4101=1529%:%
%:%4102=1530%:%
%:%4103=1530%:%
%:%4104=1531%:%
%:%4105=1531%:%
%:%4106=1532%:%
%:%4107=1532%:%
%:%4108=1533%:%
%:%4109=1533%:%
%:%4110=1534%:%
%:%4111=1534%:%
%:%4112=1535%:%
%:%4113=1535%:%
%:%4114=1535%:%
%:%4115=1536%:%
%:%4116=1536%:%
%:%4117=1537%:%
%:%4118=1537%:%
%:%4119=1537%:%
%:%4120=1538%:%
%:%4121=1538%:%
%:%4122=1539%:%
%:%4123=1539%:%
%:%4124=1539%:%
%:%4125=1540%:%
%:%4126=1540%:%
%:%4127=1541%:%
%:%4128=1541%:%
%:%4129=1542%:%
%:%4130=1542%:%
%:%4131=1543%:%
%:%4132=1543%:%
%:%4133=1544%:%
%:%4134=1544%:%
%:%4135=1545%:%
%:%4136=1545%:%
%:%4137=1545%:%
%:%4138=1545%:%
%:%4139=1546%:%
%:%4140=1546%:%
%:%4141=1547%:%
%:%4142=1547%:%
%:%4143=1548%:%
%:%4144=1548%:%
%:%4145=1548%:%
%:%4146=1548%:%
%:%4147=1549%:%
%:%4153=1549%:%
%:%4156=1550%:%
%:%4157=1551%:%
%:%4158=1551%:%
%:%4159=1552%:%
%:%4162=1555%:%
%:%4163=1556%:%
%:%4164=1557%:%
%:%4165=1558%:%
%:%4166=1559%:%
%:%4167=1559%:%
%:%4168=1560%:%
%:%4169=1561%:%
%:%4172=1562%:%
%:%4176=1562%:%
%:%4177=1562%:%
%:%4178=1563%:%
%:%4179=1563%:%
%:%4188=1579%:%
%:%4190=1580%:%
%:%4191=1580%:%
%:%4192=1581%:%
%:%4198=1587%:%
%:%4199=1588%:%
%:%4214=1603%:%
%:%4216=1604%:%
%:%4217=1604%:%
%:%4218=1605%:%
%:%4219=1606%:%
%:%4220=1607%:%
%:%4223=1608%:%
%:%4227=1608%:%
%:%4228=1608%:%
%:%4229=1609%:%
%:%4230=1609%:%
%:%4231=1610%:%
%:%4232=1610%:%
%:%4233=1611%:%
%:%4234=1611%:%
%:%4235=1612%:%
%:%4236=1612%:%
%:%4237=1612%:%
%:%4238=1613%:%
%:%4239=1613%:%
%:%4240=1613%:%
%:%4241=1614%:%
%:%4242=1614%:%
%:%4243=1615%:%
%:%4244=1616%:%
%:%4245=1617%:%
%:%4246=1617%:%
%:%4247=1617%:%
%:%4248=1618%:%
%:%4249=1618%:%
%:%4250=1618%:%
%:%4251=1619%:%
%:%4252=1619%:%
%:%4253=1620%:%
%:%4254=1620%:%
%:%4255=1621%:%
%:%4256=1621%:%
%:%4257=1622%:%
%:%4258=1622%:%
%:%4259=1622%:%
%:%4260=1623%:%
%:%4261=1623%:%
%:%4262=1624%:%
%:%4263=1624%:%
%:%4264=1625%:%
%:%4265=1625%:%
%:%4266=1626%:%
%:%4267=1626%:%
%:%4268=1627%:%
%:%4269=1627%:%
%:%4270=1628%:%
%:%4271=1628%:%
%:%4272=1629%:%
%:%4273=1629%:%
%:%4274=1630%:%
%:%4275=1630%:%
%:%4276=1631%:%
%:%4277=1631%:%
%:%4278=1632%:%
%:%4279=1632%:%
%:%4280=1633%:%
%:%4281=1633%:%
%:%4282=1634%:%
%:%4283=1634%:%
%:%4284=1635%:%
%:%4285=1636%:%
%:%4286=1637%:%
%:%4287=1638%:%
%:%4288=1638%:%
%:%4289=1639%:%
%:%4290=1639%:%
%:%4291=1640%:%
%:%4292=1640%:%
%:%4293=1641%:%
%:%4294=1641%:%
%:%4295=1642%:%
%:%4296=1642%:%
%:%4297=1643%:%
%:%4298=1643%:%
%:%4299=1644%:%
%:%4300=1644%:%
%:%4301=1645%:%
%:%4302=1645%:%
%:%4303=1646%:%
%:%4304=1646%:%
%:%4305=1646%:%
%:%4306=1647%:%
%:%4307=1647%:%
%:%4308=1647%:%
%:%4309=1648%:%
%:%4310=1648%:%
%:%4311=1649%:%
%:%4312=1650%:%
%:%4313=1651%:%
%:%4314=1651%:%
%:%4315=1652%:%
%:%4325=1654%:%
%:%4327=1655%:%
%:%4328=1655%:%
%:%4329=1656%:%
%:%4331=1658%:%
%:%4332=1659%:%
%:%4338=1665%:%
%:%4340=1666%:%
%:%4341=1666%:%
%:%4342=1667%:%
%:%4343=1668%:%
%:%4344=1669%:%
%:%4347=1670%:%
%:%4351=1670%:%
%:%4352=1670%:%
%:%4353=1671%:%
%:%4354=1671%:%
%:%4359=1671%:%
%:%4362=1672%:%
%:%4363=1673%:%
%:%4364=1674%:%
%:%4365=1674%:%
%:%4366=1675%:%
%:%4367=1676%:%
%:%4368=1676%:%
%:%4369=1677%:%
%:%4370=1678%:%
%:%4371=1679%:%
%:%4372=1680%:%
%:%4373=1680%:%
%:%4374=1681%:%
%:%4375=1682%:%
%:%4376=1683%:%
%:%4377=1683%:%
%:%4378=1684%:%
%:%4384=1690%:%
%:%4385=1691%:%
%:%4386=1692%:%
%:%4387=1692%:%
%:%4388=1693%:%
%:%4394=1699%:%
%:%4395=1700%:%
%:%4396=1701%:%
%:%4397=1701%:%
%:%4398=1702%:%
%:%4399=1703%:%
%:%4402=1706%:%
%:%4403=1707%:%
%:%4404=1708%:%
%:%4405=1708%:%
%:%4406=1709%:%
%:%4407=1710%:%
%:%4408=1711%:%
%:%4409=1712%:%
%:%4410=1713%:%
%:%4411=1713%:%
%:%4412=1714%:%
%:%4413=1715%:%
%:%4414=1716%:%
%:%4415=1716%:%
%:%4416=1717%:%
%:%4418=1719%:%
%:%4419=1720%:%
%:%4423=1724%:%
%:%4424=1725%:%
%:%4425=1726%:%
%:%4426=1726%:%
%:%4427=1727%:%
%:%4430=1730%:%
%:%4431=1731%:%
%:%4432=1732%:%
%:%4433=1732%:%
%:%4434=1733%:%
%:%4435=1734%:%
%:%4438=1735%:%
%:%4442=1735%:%
%:%4443=1735%:%
%:%4444=1735%:%
%:%4449=1735%:%
%:%4452=1736%:%
%:%4453=1737%:%
%:%4454=1737%:%
%:%4455=1738%:%
%:%4456=1739%:%
%:%4459=1740%:%
%:%4463=1740%:%
%:%4464=1740%:%
%:%4465=1741%:%
%:%4466=1741%:%
%:%4467=1742%:%
%:%4468=1742%:%
%:%4469=1743%:%
%:%4470=1743%:%
%:%4471=1743%:%
%:%4472=1743%:%
%:%4473=1744%:%
%:%4474=1744%:%
%:%4475=1745%:%
%:%4476=1745%:%
%:%4477=1746%:%
%:%4478=1746%:%
%:%4479=1746%:%
%:%4480=1746%:%
%:%4481=1747%:%
%:%4487=1747%:%
%:%4490=1748%:%
%:%4491=1749%:%
%:%4492=1749%:%
%:%4495=1750%:%
%:%4499=1750%:%
%:%4500=1750%:%
%:%4501=1750%:%
%:%4506=1750%:%
%:%4509=1751%:%
%:%4510=1752%:%
%:%4511=1752%:%
%:%4514=1753%:%
%:%4518=1753%:%
%:%4519=1753%:%
%:%4520=1753%:%
%:%4525=1753%:%
%:%4528=1754%:%
%:%4529=1755%:%
%:%4529=1756%:%
%:%4530=1757%:%
%:%4531=1758%:%
%:%4532=1758%:%
%:%4535=1759%:%
%:%4539=1759%:%
%:%4540=1759%:%
%:%4545=1759%:%
%:%4548=1760%:%
%:%4549=1761%:%
%:%4550=1761%:%
%:%4553=1762%:%
%:%4557=1762%:%
%:%4558=1762%:%
%:%4563=1762%:%
%:%4566=1763%:%
%:%4567=1764%:%
%:%4568=1764%:%
%:%4569=1765%:%
%:%4570=1766%:%
%:%4573=1767%:%
%:%4577=1767%:%
%:%4578=1767%:%
%:%4579=1768%:%
%:%4580=1768%:%
%:%4581=1769%:%
%:%4582=1769%:%
%:%4583=1770%:%
%:%4584=1770%:%
%:%4585=1771%:%
%:%4586=1771%:%
%:%4587=1771%:%
%:%4588=1772%:%
%:%4589=1772%:%
%:%4590=1772%:%
%:%4591=1773%:%
%:%4597=1773%:%
%:%4600=1774%:%
%:%4601=1775%:%
%:%4602=1775%:%
%:%4603=1776%:%
%:%4604=1777%:%
%:%4607=1778%:%
%:%4611=1778%:%
%:%4612=1778%:%
%:%4613=1779%:%
%:%4614=1779%:%
%:%4615=1780%:%
%:%4616=1780%:%
%:%4617=1781%:%
%:%4618=1781%:%
%:%4619=1782%:%
%:%4620=1782%:%
%:%4621=1782%:%
%:%4622=1783%:%
%:%4623=1783%:%
%:%4624=1783%:%
%:%4625=1784%:%
%:%4631=1784%:%
%:%4634=1785%:%
%:%4635=1786%:%
%:%4636=1786%:%
%:%4637=1787%:%
%:%4638=1788%:%
%:%4641=1789%:%
%:%4645=1789%:%
%:%4646=1789%:%
%:%4647=1790%:%
%:%4648=1790%:%
%:%4649=1791%:%
%:%4650=1791%:%
%:%4651=1792%:%
%:%4652=1792%:%
%:%4653=1792%:%
%:%4654=1792%:%
%:%4655=1793%:%
%:%4656=1793%:%
%:%4657=1794%:%
%:%4658=1794%:%
%:%4659=1795%:%
%:%4660=1795%:%
%:%4661=1795%:%
%:%4662=1796%:%
%:%4663=1796%:%
%:%4664=1797%:%
%:%4665=1797%:%
%:%4666=1797%:%
%:%4667=1798%:%
%:%4668=1798%:%
%:%4669=1798%:%
%:%4670=1799%:%
%:%4671=1799%:%
%:%4672=1799%:%
%:%4673=1800%:%
%:%4674=1801%:%
%:%4675=1801%:%
%:%4676=1802%:%
%:%4677=1802%:%
%:%4678=1802%:%
%:%4679=1803%:%
%:%4680=1804%:%
%:%4681=1804%:%
%:%4682=1805%:%
%:%4683=1806%:%
%:%4684=1806%:%
%:%4685=1807%:%
%:%4686=1807%:%
%:%4687=1808%:%
%:%4688=1808%:%
%:%4689=1809%:%
%:%4690=1809%:%
%:%4691=1809%:%
%:%4692=1810%:%
%:%4693=1810%:%
%:%4694=1811%:%
%:%4695=1811%:%
%:%4696=1812%:%
%:%4697=1812%:%
%:%4698=1812%:%
%:%4699=1813%:%
%:%4700=1813%:%
%:%4701=1813%:%
%:%4702=1814%:%
%:%4703=1814%:%
%:%4704=1815%:%
%:%4705=1815%:%
%:%4706=1816%:%
%:%4707=1816%:%
%:%4708=1816%:%
%:%4709=1817%:%
%:%4710=1817%:%
%:%4711=1818%:%
%:%4712=1818%:%
%:%4713=1819%:%
%:%4714=1819%:%
%:%4715=1820%:%
%:%4716=1820%:%
%:%4717=1820%:%
%:%4718=1821%:%
%:%4719=1821%:%
%:%4720=1821%:%
%:%4721=1822%:%
%:%4722=1822%:%
%:%4723=1823%:%
%:%4729=1823%:%
%:%4732=1824%:%
%:%4733=1825%:%
%:%4734=1825%:%
%:%4735=1826%:%
%:%4736=1827%:%
%:%4739=1828%:%
%:%4743=1828%:%
%:%4744=1828%:%
%:%4745=1829%:%
%:%4746=1829%:%
%:%4747=1830%:%
%:%4748=1830%:%
%:%4749=1831%:%
%:%4750=1831%:%
%:%4751=1831%:%
%:%4752=1832%:%
%:%4753=1832%:%
%:%4754=1833%:%
%:%4755=1833%:%
%:%4756=1834%:%
%:%4757=1834%:%
%:%4758=1835%:%
%:%4759=1835%:%
%:%4760=1835%:%
%:%4761=1836%:%
%:%4762=1836%:%
%:%4763=1837%:%
%:%4764=1837%:%
%:%4765=1837%:%
%:%4766=1838%:%
%:%4767=1838%:%
%:%4768=1838%:%
%:%4769=1839%:%
%:%4770=1839%:%
%:%4771=1839%:%
%:%4772=1840%:%
%:%4773=1841%:%
%:%4774=1841%:%
%:%4775=1842%:%
%:%4776=1842%:%
%:%4777=1842%:%
%:%4778=1843%:%
%:%4779=1844%:%
%:%4780=1845%:%
%:%4781=1846%:%
%:%4782=1846%:%
%:%4783=1847%:%
%:%4784=1848%:%
%:%4785=1848%:%
%:%4786=1849%:%
%:%4787=1849%:%
%:%4788=1850%:%
%:%4789=1850%:%
%:%4790=1851%:%
%:%4791=1851%:%
%:%4792=1851%:%
%:%4793=1852%:%
%:%4794=1852%:%
%:%4795=1853%:%
%:%4796=1853%:%
%:%4797=1854%:%
%:%4798=1854%:%
%:%4799=1854%:%
%:%4800=1855%:%
%:%4801=1855%:%
%:%4802=1855%:%
%:%4803=1856%:%
%:%4804=1857%:%
%:%4805=1857%:%
%:%4806=1858%:%
%:%4807=1858%:%
%:%4808=1858%:%
%:%4809=1859%:%
%:%4810=1859%:%
%:%4811=1859%:%
%:%4812=1860%:%
%:%4813=1860%:%
%:%4814=1861%:%
%:%4815=1861%:%
%:%4816=1862%:%
%:%4817=1862%:%
%:%4818=1862%:%
%:%4819=1863%:%
%:%4820=1863%:%
%:%4821=1863%:%
%:%4822=1864%:%
%:%4823=1864%:%
%:%4824=1864%:%
%:%4825=1865%:%
%:%4826=1865%:%
%:%4827=1865%:%
%:%4828=1866%:%
%:%4829=1866%:%
%:%4830=1866%:%
%:%4831=1867%:%
%:%4832=1867%:%
%:%4833=1868%:%
%:%4834=1869%:%
%:%4835=1870%:%
%:%4836=1870%:%
%:%4837=1871%:%
%:%4838=1871%:%
%:%4839=1871%:%
%:%4840=1872%:%
%:%4841=1872%:%
%:%4842=1872%:%
%:%4843=1873%:%
%:%4844=1873%:%
%:%4845=1873%:%
%:%4846=1874%:%
%:%4847=1874%:%
%:%4848=1874%:%
%:%4849=1875%:%
%:%4850=1875%:%
%:%4851=1875%:%
%:%4852=1876%:%
%:%4853=1876%:%
%:%4854=1877%:%
%:%4855=1877%:%
%:%4856=1878%:%
%:%4857=1878%:%
%:%4858=1879%:%
%:%4859=1879%:%
%:%4860=1879%:%
%:%4861=1880%:%
%:%4862=1880%:%
%:%4863=1880%:%
%:%4864=1881%:%
%:%4865=1881%:%
%:%4866=1882%:%
%:%4867=1882%:%
%:%4868=1883%:%
%:%4869=1883%:%
%:%4870=1883%:%
%:%4871=1884%:%
%:%4872=1884%:%
%:%4873=1884%:%
%:%4874=1885%:%
%:%4875=1886%:%
%:%4876=1886%:%
%:%4877=1887%:%
%:%4878=1887%:%
%:%4879=1887%:%
%:%4880=1888%:%
%:%4881=1888%:%
%:%4882=1888%:%
%:%4883=1889%:%
%:%4884=1889%:%
%:%4885=1890%:%
%:%4886=1890%:%
%:%4887=1891%:%
%:%4888=1891%:%
%:%4889=1891%:%
%:%4890=1892%:%
%:%4891=1892%:%
%:%4892=1892%:%
%:%4893=1893%:%
%:%4894=1893%:%
%:%4895=1893%:%
%:%4896=1894%:%
%:%4897=1894%:%
%:%4898=1894%:%
%:%4899=1895%:%
%:%4900=1895%:%
%:%4901=1895%:%
%:%4902=1896%:%
%:%4903=1896%:%
%:%4904=1897%:%
%:%4905=1898%:%
%:%4906=1899%:%
%:%4907=1899%:%
%:%4908=1900%:%
%:%4909=1900%:%
%:%4910=1900%:%
%:%4911=1901%:%
%:%4912=1901%:%
%:%4913=1901%:%
%:%4914=1902%:%
%:%4915=1902%:%
%:%4916=1902%:%
%:%4917=1903%:%
%:%4918=1903%:%
%:%4919=1903%:%
%:%4920=1904%:%
%:%4921=1904%:%
%:%4922=1904%:%
%:%4923=1905%:%
%:%4924=1905%:%
%:%4925=1906%:%
%:%4926=1906%:%
%:%4927=1907%:%
%:%4933=1907%:%
%:%4936=1908%:%
%:%4937=1909%:%
%:%4938=1909%:%
%:%4939=1910%:%
%:%4940=1911%:%
%:%4943=1912%:%
%:%4947=1912%:%
%:%4948=1912%:%
%:%4949=1913%:%
%:%4950=1913%:%
%:%4955=1913%:%
%:%4958=1914%:%
%:%4959=1915%:%
%:%4960=1916%:%
%:%4961=1916%:%
%:%4962=1917%:%
%:%4963=1918%:%
%:%4966=1919%:%
%:%4970=1919%:%
%:%4971=1919%:%
%:%4972=1920%:%
%:%4973=1920%:%
%:%4974=1921%:%
%:%4975=1921%:%
%:%4976=1922%:%
%:%4977=1922%:%
%:%4978=1922%:%
%:%4979=1923%:%
%:%4980=1923%:%
%:%4981=1924%:%
%:%4982=1924%:%
%:%4983=1925%:%
%:%4984=1925%:%
%:%4985=1926%:%
%:%4986=1926%:%
%:%4987=1927%:%
%:%4988=1927%:%
%:%4989=1928%:%
%:%4990=1928%:%
%:%4991=1929%:%
%:%4992=1929%:%
%:%4993=1929%:%
%:%4994=1930%:%
%:%4995=1930%:%
%:%4996=1931%:%
%:%4997=1931%:%
%:%4998=1932%:%
%:%4999=1932%:%
%:%5000=1933%:%
%:%5001=1933%:%
%:%5002=1934%:%
%:%5003=1934%:%
%:%5004=1934%:%
%:%5005=1935%:%
%:%5006=1935%:%
%:%5007=1935%:%
%:%5008=1936%:%
%:%5009=1936%:%
%:%5010=1936%:%
%:%5011=1937%:%
%:%5012=1937%:%
%:%5013=1937%:%
%:%5014=1938%:%
%:%5015=1938%:%
%:%5016=1938%:%
%:%5017=1939%:%
%:%5018=1939%:%
%:%5019=1939%:%
%:%5020=1940%:%
%:%5021=1940%:%
%:%5022=1940%:%
%:%5023=1941%:%
%:%5024=1942%:%
%:%5025=1942%:%
%:%5026=1943%:%
%:%5027=1943%:%
%:%5028=1943%:%
%:%5029=1944%:%
%:%5030=1945%:%
%:%5031=1946%:%
%:%5032=1947%:%
%:%5033=1947%:%
%:%5034=1947%:%
%:%5035=1948%:%
%:%5036=1948%:%
%:%5037=1948%:%
%:%5038=1949%:%
%:%5039=1950%:%
%:%5040=1950%:%
%:%5041=1951%:%
%:%5042=1951%:%
%:%5043=1951%:%
%:%5044=1952%:%
%:%5045=1953%:%
%:%5046=1953%:%
%:%5047=1954%:%
%:%5048=1955%:%
%:%5049=1955%:%
%:%5050=1956%:%
%:%5051=1957%:%
%:%5052=1957%:%
%:%5053=1958%:%
%:%5054=1958%:%
%:%5055=1959%:%
%:%5056=1959%:%
%:%5057=1959%:%
%:%5058=1960%:%
%:%5059=1960%:%
%:%5060=1961%:%
%:%5061=1961%:%
%:%5062=1962%:%
%:%5063=1962%:%
%:%5064=1962%:%
%:%5065=1963%:%
%:%5066=1963%:%
%:%5067=1963%:%
%:%5068=1964%:%
%:%5069=1964%:%
%:%5070=1964%:%
%:%5071=1965%:%
%:%5072=1965%:%
%:%5073=1966%:%
%:%5074=1967%:%
%:%5075=1967%:%
%:%5076=1968%:%
%:%5077=1968%:%
%:%5078=1969%:%
%:%5079=1969%:%
%:%5080=1970%:%
%:%5081=1970%:%
%:%5082=1971%:%
%:%5083=1971%:%
%:%5084=1971%:%
%:%5085=1972%:%
%:%5086=1972%:%
%:%5087=1972%:%
%:%5088=1973%:%
%:%5089=1973%:%
%:%5090=1973%:%
%:%5091=1974%:%
%:%5092=1975%:%
%:%5093=1975%:%
%:%5094=1975%:%
%:%5095=1976%:%
%:%5096=1976%:%
%:%5097=1977%:%
%:%5098=1977%:%
%:%5099=1978%:%
%:%5100=1978%:%
%:%5101=1978%:%
%:%5102=1979%:%
%:%5103=1979%:%
%:%5104=1980%:%
%:%5105=1981%:%
%:%5106=1981%:%
%:%5107=1982%:%
%:%5108=1982%:%
%:%5109=1982%:%
%:%5110=1983%:%
%:%5111=1983%:%
%:%5112=1984%:%
%:%5113=1984%:%
%:%5114=1985%:%
%:%5115=1985%:%
%:%5116=1986%:%
%:%5117=1986%:%
%:%5118=1986%:%
%:%5119=1987%:%
%:%5120=1987%:%
%:%5121=1987%:%
%:%5122=1988%:%
%:%5123=1988%:%
%:%5124=1988%:%
%:%5125=1989%:%
%:%5126=1989%:%
%:%5127=1990%:%
%:%5128=1991%:%
%:%5129=1992%:%
%:%5130=1992%:%
%:%5131=1993%:%
%:%5132=1993%:%
%:%5133=1994%:%
%:%5134=1994%:%
%:%5135=1995%:%
%:%5136=1995%:%
%:%5137=1996%:%
%:%5143=1996%:%
%:%5146=1997%:%
%:%5147=1998%:%
%:%5148=1998%:%
%:%5149=1999%:%
%:%5150=2000%:%
%:%5151=2001%:%
%:%5154=2002%:%
%:%5158=2002%:%
%:%5159=2002%:%
%:%5160=2003%:%
%:%5161=2003%:%
%:%5162=2004%:%
%:%5163=2004%:%
%:%5164=2005%:%
%:%5165=2005%:%
%:%5166=2005%:%
%:%5167=2005%:%
%:%5168=2006%:%
%:%5169=2006%:%
%:%5170=2007%:%
%:%5171=2007%:%
%:%5172=2008%:%
%:%5173=2008%:%
%:%5174=2008%:%
%:%5175=2009%:%
%:%5176=2009%:%
%:%5177=2010%:%
%:%5178=2010%:%
%:%5179=2011%:%
%:%5185=2011%:%
%:%5188=2012%:%
%:%5189=2013%:%
%:%5190=2013%:%
%:%5191=2014%:%
%:%5192=2015%:%
%:%5195=2016%:%
%:%5199=2016%:%
%:%5200=2016%:%
%:%5201=2017%:%
%:%5202=2017%:%
%:%5203=2018%:%
%:%5204=2018%:%
%:%5205=2019%:%
%:%5206=2019%:%
%:%5207=2019%:%
%:%5208=2020%:%
%:%5209=2020%:%
%:%5210=2021%:%
%:%5211=2021%:%
%:%5212=2022%:%
%:%5213=2022%:%
%:%5214=2023%:%
%:%5215=2023%:%
%:%5216=2023%:%
%:%5217=2024%:%
%:%5218=2024%:%
%:%5219=2025%:%
%:%5220=2025%:%
%:%5221=2025%:%
%:%5222=2026%:%
%:%5223=2026%:%
%:%5224=2026%:%
%:%5225=2027%:%
%:%5231=2027%:%
%:%5234=2028%:%
%:%5235=2029%:%
%:%5236=2029%:%
%:%5237=2030%:%
%:%5238=2031%:%
%:%5239=2032%:%
%:%5242=2033%:%
%:%5246=2033%:%
%:%5247=2033%:%
%:%5248=2033%:%
%:%5249=2034%:%
%:%5250=2034%:%
%:%5251=2035%:%
%:%5252=2035%:%
%:%5253=2036%:%
%:%5254=2036%:%
%:%5255=2036%:%
%:%5256=2036%:%
%:%5257=2037%:%
%:%5258=2037%:%
%:%5259=2038%:%
%:%5260=2038%:%
%:%5261=2039%:%
%:%5262=2039%:%
%:%5263=2039%:%
%:%5264=2040%:%
%:%5265=2041%:%
%:%5266=2041%:%
%:%5267=2042%:%
%:%5268=2042%:%
%:%5269=2042%:%
%:%5270=2043%:%
%:%5271=2043%:%
%:%5272=2043%:%
%:%5273=2044%:%
%:%5274=2044%:%
%:%5275=2044%:%
%:%5276=2045%:%
%:%5277=2045%:%
%:%5278=2046%:%
%:%5279=2047%:%
%:%5280=2048%:%
%:%5281=2049%:%
%:%5282=2049%:%
%:%5283=2050%:%
%:%5289=2050%:%
%:%5292=2051%:%
%:%5293=2052%:%
%:%5294=2052%:%
%:%5295=2053%:%
%:%5296=2054%:%
%:%5299=2055%:%
%:%5303=2055%:%
%:%5304=2055%:%
%:%5309=2055%:%
%:%5312=2056%:%
%:%5313=2057%:%
%:%5314=2058%:%
%:%5315=2058%:%
%:%5316=2059%:%
%:%5317=2060%:%
%:%5320=2061%:%
%:%5324=2061%:%
%:%5325=2061%:%
%:%5326=2061%:%
%:%5331=2061%:%
%:%5334=2062%:%
%:%5335=2063%:%
%:%5336=2064%:%
%:%5337=2064%:%
%:%5338=2065%:%
%:%5339=2066%:%
%:%5340=2067%:%
%:%5343=2068%:%
%:%5347=2068%:%
%:%5348=2068%:%
%:%5349=2069%:%
%:%5350=2070%:%
%:%5351=2070%:%
%:%5360=2072%:%
%:%5361=2073%:%
%:%5363=2074%:%
%:%5364=2074%:%
%:%5365=2075%:%
%:%5366=2076%:%
%:%5367=2077%:%
%:%5372=2096%:%
%:%5373=2097%:%
%:%5374=2098%:%
%:%5376=2099%:%
%:%5377=2099%:%
%:%5378=2100%:%
%:%5379=2101%:%
%:%5380=2102%:%
%:%5383=2103%:%
%:%5387=2103%:%
%:%5388=2103%:%
%:%5389=2104%:%
%:%5390=2104%:%
%:%5391=2105%:%
%:%5392=2105%:%
%:%5393=2106%:%
%:%5394=2107%:%
%:%5395=2108%:%
%:%5396=2109%:%
%:%5397=2109%:%
%:%5402=2109%:%
%:%5405=2110%:%
%:%5406=2111%:%
%:%5407=2111%:%
%:%5408=2112%:%
%:%5409=2113%:%
%:%5410=2114%:%
%:%5413=2115%:%
%:%5417=2115%:%
%:%5418=2115%:%
%:%5419=2116%:%
%:%5420=2116%:%
%:%5421=2117%:%
%:%5422=2117%:%
%:%5423=2118%:%
%:%5424=2119%:%
%:%5425=2120%:%
%:%5426=2121%:%
%:%5427=2121%:%
%:%5436=2123%:%
%:%5438=2124%:%
%:%5439=2124%:%
%:%5440=2125%:%
%:%5441=2126%:%
%:%5442=2127%:%
%:%5445=2128%:%
%:%5449=2128%:%
%:%5450=2128%:%
%:%5451=2129%:%
%:%5452=2130%:%
%:%5453=2131%:%
%:%5454=2132%:%
%:%5455=2132%:%
%:%5464=2134%:%
%:%5466=2135%:%
%:%5467=2135%:%
%:%5468=2136%:%
%:%5471=2139%:%
%:%5472=2140%:%
%:%5474=2141%:%
%:%5475=2141%:%
%:%5476=2142%:%
%:%5477=2143%:%
%:%5478=2144%:%
%:%5479=2145%:%
%:%5480=2146%:%
%:%5483=2147%:%
%:%5487=2147%:%
%:%5488=2147%:%
%:%5489=2147%:%
%:%5498=2149%:%
%:%5499=2150%:%
%:%5501=2151%:%
%:%5502=2151%:%
%:%5503=2152%:%
%:%5504=2153%:%
%:%5505=2154%:%
%:%5506=2155%:%
%:%5507=2156%:%
%:%5510=2157%:%
%:%5514=2157%:%
%:%5515=2157%:%
%:%5516=2157%:%
%:%5525=2159%:%
%:%5527=2160%:%
%:%5528=2160%:%
%:%5529=2161%:%
%:%5532=2162%:%
%:%5536=2162%:%
%:%5537=2162%:%
%:%5538=2163%:%
%:%5539=2164%:%
%:%5540=2165%:%
%:%5541=2166%:%
%:%5542=2166%:%
%:%5543=2167%:%
%:%5544=2167%:%
%:%5553=2169%:%
%:%5554=2170%:%
%:%5556=2171%:%
%:%5557=2171%:%
%:%5558=2172%:%
%:%5561=2175%:%
%:%5562=2176%:%
%:%5563=2177%:%
%:%5564=2178%:%
%:%5564=2182%:%
%:%5565=2183%:%
%:%5566=2184%:%
%:%5567=2184%:%
%:%5568=2185%:%
%:%5569=2186%:%
%:%5572=2187%:%
%:%5576=2187%:%
%:%5577=2187%:%
%:%5578=2188%:%
%:%5579=2188%:%
%:%5580=2189%:%
%:%5581=2189%:%
%:%5582=2189%:%
%:%5583=2190%:%
%:%5584=2190%:%
%:%5585=2190%:%
%:%5586=2191%:%
%:%5587=2191%:%
%:%5588=2192%:%
%:%5589=2192%:%
%:%5590=2193%:%
%:%5600=2195%:%
%:%5602=2196%:%
%:%5603=2196%:%
%:%5604=2197%:%
%:%5605=2198%:%
%:%5608=2199%:%
%:%5612=2199%:%
%:%5613=2199%:%
%:%5614=2200%:%
%:%5615=2200%:%
%:%5616=2201%:%
%:%5617=2201%:%
%:%5618=2201%:%
%:%5619=2201%:%
%:%5620=2202%:%
%:%5621=2202%:%
%:%5622=2203%:%
%:%5623=2203%:%
%:%5624=2204%:%
%:%5625=2204%:%
%:%5626=2204%:%
%:%5627=2205%:%
%:%5628=2205%:%
%:%5629=2206%:%
%:%5630=2206%:%
%:%5631=2206%:%
%:%5632=2207%:%
%:%5633=2207%:%
%:%5634=2208%:%
%:%5635=2208%:%
%:%5636=2209%:%
%:%5637=2210%:%
%:%5638=2211%:%
%:%5639=2212%:%
%:%5640=2212%:%
%:%5641=2213%:%
%:%5642=2213%:%
%:%5643=2214%:%
%:%5644=2214%:%
%:%5645=2215%:%
%:%5646=2215%:%
%:%5647=2216%:%
%:%5653=2216%:%
%:%5656=2217%:%
%:%5657=2218%:%
%:%5658=2219%:%
%:%5659=2219%:%
%:%5660=2220%:%
%:%5661=2221%:%
%:%5664=2222%:%
%:%5668=2222%:%
%:%5669=2222%:%
%:%5670=2222%:%
%:%5675=2222%:%
%:%5678=2223%:%
%:%5679=2224%:%
%:%5680=2224%:%
%:%5681=2225%:%
%:%5682=2226%:%
%:%5683=2227%:%
%:%5684=2228%:%
%:%5685=2229%:%
%:%5685=2230%:%
%:%5686=2231%:%
%:%5686=2291%:%
%:%5687=2292%:%
%:%5688=2293%:%
%:%5689=2293%:%
%:%5690=2293%:%
%:%5698=2295%:%
%:%5710=2296%:%
%:%5711=2297%:%
%:%5713=2300%:%
%:%5714=2300%:%
%:%5715=2301%:%
%:%5716=2302%:%
%:%5717=2303%:%
%:%5718=2303%:%
%:%5721=2304%:%
%:%5725=2304%:%
%:%5726=2304%:%
%:%5727=2304%:%
%:%5732=2304%:%
%:%5735=2305%:%
%:%5736=2306%:%
%:%5737=2306%:%
%:%5739=2308%:%
%:%5740=2309%:%
%:%5741=2310%:%
%:%5742=2310%:%
%:%5745=2311%:%
%:%5749=2311%:%
%:%5750=2311%:%
%:%5751=2311%:%
%:%5756=2311%:%
%:%5759=2312%:%
%:%5760=2313%:%
%:%5761=2313%:%
%:%5771=2323%:%
%:%5772=2324%:%
%:%5773=2325%:%
%:%5774=2325%:%
%:%5775=2326%:%
%:%5782=2327%:%
%:%5783=2327%:%
%:%5784=2328%:%
%:%5785=2328%:%
%:%5786=2328%:%
%:%5787=2329%:%
%:%5788=2329%:%
%:%5789=2330%:%
%:%5790=2330%:%
%:%5791=2330%:%
%:%5792=2331%:%
%:%5798=2331%:%
%:%5801=2332%:%
%:%5802=2333%:%
%:%5803=2333%:%
%:%5804=2334%:%
%:%5805=2335%:%
%:%5806=2335%:%
%:%5814=2343%:%
%:%5815=2344%:%
%:%5816=2345%:%
%:%5817=2345%:%
%:%5818=2346%:%
%:%5825=2347%:%
%:%5826=2347%:%
%:%5827=2348%:%
%:%5828=2348%:%
%:%5829=2348%:%
%:%5830=2349%:%
%:%5831=2349%:%
%:%5832=2350%:%
%:%5833=2350%:%
%:%5834=2350%:%
%:%5835=2351%:%
%:%5841=2351%:%
%:%5844=2352%:%
%:%5845=2353%:%
%:%5846=2353%:%
%:%5855=2362%:%
%:%5856=2363%:%
%:%5858=2365%:%
%:%5859=2365%:%
%:%5860=2366%:%
%:%5867=2367%:%
%:%5868=2367%:%
%:%5869=2368%:%
%:%5870=2368%:%
%:%5871=2368%:%
%:%5872=2369%:%
%:%5873=2369%:%
%:%5874=2370%:%
%:%5875=2370%:%
%:%5876=2371%:%
%:%5877=2371%:%
%:%5878=2372%:%
%:%5879=2372%:%
%:%5880=2373%:%
%:%5886=2373%:%
%:%5889=2374%:%
%:%5890=2375%:%
%:%5891=2375%:%
%:%5898=2377%:%
%:%5910=2379%:%
%:%5911=2380%:%
%:%5920=2382%:%
%:%5932=2383%:%
%:%5933=2384%:%
%:%5934=2385%:%
%:%5936=2388%:%
%:%5937=2388%:%
%:%5938=2389%:%
%:%5939=2390%:%
%:%5940=2391%:%
%:%5941=2392%:%
%:%5942=2393%:%
%:%5943=2394%:%
%:%5944=2395%:%
%:%5945=2396%:%
%:%5946=2397%:%
%:%5947=2398%:%
%:%5948=2399%:%
%:%5949=2400%:%
%:%5950=2401%:%
%:%5951=2402%:%
%:%5952=2403%:%
%:%5953=2404%:%
%:%5954=2405%:%
%:%5955=2406%:%
%:%5956=2407%:%
%:%5957=2408%:%
%:%5958=2409%:%
%:%5959=2410%:%
%:%5960=2411%:%
%:%5961=2412%:%
%:%5962=2412%:%
%:%5963=2413%:%
%:%5964=2414%:%
%:%5965=2414%:%
%:%5966=2415%:%
%:%5967=2416%:%
%:%5968=2417%:%
%:%5969=2418%:%
%:%5970=2419%:%
%:%5971=2420%:%
%:%5972=2421%:%
%:%5973=2422%:%
%:%5974=2422%:%
%:%5975=2423%:%
%:%5976=2424%:%
%:%5977=2424%:%
%:%5978=2425%:%
%:%5979=2426%:%
%:%5980=2427%:%
%:%5981=2428%:%
%:%5982=2429%:%
%:%5983=2430%:%
%:%5984=2431%:%
%:%5985=2432%:%
%:%5986=2433%:%
%:%5987=2434%:%
%:%5988=2435%:%
%:%5989=2436%:%
%:%5990=2437%:%
%:%5991=2438%:%
%:%5992=2439%:%
%:%5993=2440%:%
%:%5994=2441%:%
%:%5995=2442%:%
%:%5996=2443%:%
%:%5997=2444%:%
%:%5998=2445%:%
%:%5999=2446%:%
%:%6000=2447%:%
%:%6001=2448%:%
%:%6002=2449%:%
%:%6003=2450%:%
%:%6004=2451%:%
%:%6005=2452%:%
%:%6006=2453%:%
%:%6007=2454%:%
%:%6008=2455%:%
%:%6009=2456%:%
%:%6010=2457%:%
%:%6011=2458%:%
%:%6012=2459%:%
%:%6013=2460%:%
%:%6014=2461%:%
%:%6015=2462%:%
%:%6016=2463%:%
%:%6017=2464%:%
%:%6018=2465%:%
%:%6019=2466%:%
%:%6020=2467%:%
%:%6020=2469%:%
%:%6021=2470%:%
%:%6022=2471%:%
%:%6023=2471%:%
%:%6030=2473%:%
%:%6040=2475%:%
%:%6041=2475%:%
%:%6042=2476%:%
%:%6043=2477%:%
%:%6044=2477%:%
%:%6045=2478%:%
%:%6046=2478%:%
%:%6047=2479%:%
%:%6048=2479%:%
%:%6049=2480%:%
%:%6050=2480%:%
%:%6051=2481%:%
%:%6052=2482%:%
%:%6053=2482%:%
%:%6054=2483%:%
%:%6055=2484%:%
%:%6056=2484%:%
%:%6063=2486%:%
%:%6073=2488%:%
%:%6074=2488%:%
%:%6075=2489%:%
%:%6076=2490%:%
%:%6077=2491%:%
%:%6078=2492%:%
%:%6079=2493%:%
%:%6080=2493%:%
%:%6083=2494%:%
%:%6087=2494%:%
%:%6088=2494%:%
%:%6089=2495%:%
%:%6090=2495%:%
%:%6091=2496%:%
%:%6097=2496%:%
%:%6100=2497%:%
%:%6101=2498%:%
%:%6102=2498%:%
%:%6103=2499%:%
%:%6104=2500%:%
%:%6105=2501%:%
%:%6106=2501%:%
%:%6109=2502%:%
%:%6113=2502%:%
%:%6114=2502%:%
%:%6115=2503%:%
%:%6116=2503%:%
%:%6117=2504%:%
%:%6118=2504%:%
%:%6119=2505%:%
%:%6134=2507%:%
%:%6144=2509%:%
%:%6145=2509%:%
%:%6147=2511%:%
%:%6149=2512%:%
%:%6150=2512%:%
%:%6152=2513%:%
%:%6154=2514%:%
%:%6155=2514%:%
%:%6157=2516%:%
%:%6159=2518%:%
%:%6160=2518%:%
%:%6161=2519%:%
%:%6162=2520%:%
%:%6163=2521%:%
%:%6164=2522%:%
%:%6165=2523%:%
%:%6166=2523%:%
%:%6169=2524%:%
%:%6173=2524%:%
%:%6174=2524%:%
%:%6179=2524%:%
%:%6182=2525%:%
%:%6183=2526%:%
%:%6184=2526%:%
%:%6187=2527%:%
%:%6191=2527%:%
%:%6192=2527%:%
%:%6197=2527%:%
%:%6200=2528%:%
%:%6201=2529%:%
%:%6202=2529%:%
%:%6205=2530%:%
%:%6209=2530%:%
%:%6210=2530%:%
%:%6211=2530%:%
%:%6216=2530%:%
%:%6219=2531%:%
%:%6220=2532%:%
%:%6221=2532%:%
%:%6224=2533%:%
%:%6228=2533%:%
%:%6229=2533%:%
%:%6234=2533%:%
%:%6237=2534%:%
%:%6238=2535%:%
%:%6239=2535%:%
%:%6240=2536%:%
%:%6241=2537%:%
%:%6243=2539%:%
%:%6244=2540%:%
%:%6245=2541%:%
%:%6246=2541%:%
%:%6247=2542%:%
%:%6248=2543%:%
%:%6249=2544%:%
%:%6250=2544%:%
%:%6251=2545%:%
%:%6252=2546%:%
%:%6255=2547%:%
%:%6259=2547%:%
%:%6260=2547%:%
%:%6261=2548%:%
%:%6262=2548%:%
%:%6263=2549%:%
%:%6264=2549%:%
%:%6265=2550%:%
%:%6266=2550%:%
%:%6267=2550%:%
%:%6268=2550%:%
%:%6269=2551%:%
%:%6270=2551%:%
%:%6271=2552%:%
%:%6272=2552%:%
%:%6273=2553%:%
%:%6274=2553%:%
%:%6275=2553%:%
%:%6276=2554%:%
%:%6277=2554%:%
%:%6278=2555%:%
%:%6284=2555%:%
%:%6287=2556%:%
%:%6288=2557%:%
%:%6289=2557%:%
%:%6292=2558%:%
%:%6296=2558%:%
%:%6297=2558%:%
%:%6298=2558%:%
%:%6303=2558%:%
%:%6306=2559%:%
%:%6307=2560%:%
%:%6308=2560%:%
%:%6311=2561%:%
%:%6315=2561%:%
%:%6316=2561%:%
%:%6317=2561%:%
%:%6322=2561%:%
%:%6325=2562%:%
%:%6326=2563%:%
%:%6327=2563%:%
%:%6330=2564%:%
%:%6334=2564%:%
%:%6335=2564%:%
%:%6336=2564%:%
%:%6345=2566%:%
%:%6347=2568%:%
%:%6348=2568%:%
%:%6349=2569%:%
%:%6350=2570%:%
%:%6351=2571%:%
%:%6352=2572%:%
%:%6353=2572%:%
%:%6354=2573%:%
%:%6357=2576%:%
%:%6358=2577%:%
%:%6359=2578%:%
%:%6360=2578%:%
%:%6361=2579%:%
%:%6362=2579%:%
%:%6363=2580%:%
%:%6364=2581%:%
%:%6365=2581%:%
%:%6366=2582%:%
%:%6367=2583%:%
%:%6368=2583%:%
%:%6369=2584%:%
%:%6370=2585%:%
%:%6373=2586%:%
%:%6377=2586%:%
%:%6378=2586%:%
%:%6379=2587%:%
%:%6380=2587%:%
%:%6389=2589%:%
%:%6391=2590%:%
%:%6392=2590%:%
%:%6393=2591%:%
%:%6394=2592%:%
%:%6397=2593%:%
%:%6401=2593%:%
%:%6402=2593%:%
%:%6403=2594%:%
%:%6404=2594%:%
%:%6405=2595%:%
%:%6406=2595%:%
%:%6407=2596%:%
%:%6408=2596%:%
%:%6409=2596%:%
%:%6410=2596%:%
%:%6411=2597%:%
%:%6412=2597%:%
%:%6413=2598%:%
%:%6414=2598%:%
%:%6415=2599%:%
%:%6416=2599%:%
%:%6417=2599%:%
%:%6418=2600%:%
%:%6419=2600%:%
%:%6420=2601%:%
%:%6430=2603%:%
%:%6432=2604%:%
%:%6433=2604%:%
%:%6434=2605%:%
%:%6439=2610%:%
%:%6448=2612%:%
%:%6458=2617%:%
%:%6459=2617%:%
%:%6460=2618%:%
%:%6461=2619%:%
%:%6462=2620%:%
%:%6463=2621%:%
%:%6464=2622%:%
%:%6465=2623%:%
%:%6466=2624%:%
%:%6467=2625%:%
%:%6468=2626%:%
%:%6469=2627%:%
%:%6470=2628%:%
%:%6471=2629%:%
%:%6472=2630%:%
%:%6473=2631%:%
%:%6474=2632%:%
%:%6475=2649%:%
%:%6478=2650%:%
%:%6483=2651%:%
%:%6484=2651%:%
%:%6485=2651%:%

%
\begin{isabellebody}%
\setisabellecontext{TEMPORAL{\isacharunderscore}{\kern0pt}PDDL{\isacharunderscore}{\kern0pt}Semantics{\isacharunderscore}{\kern0pt}Alt}%
%
\isadelimdocument
%
\endisadelimdocument
%
\isatagdocument
%
\isamarkupsection{Connecting the Semantics to State-Transition Based Semantics%
}
\isamarkuptrue%
%
\endisatagdocument
{\isafolddocument}%
%
\isadelimdocument
%
\endisadelimdocument
%
\isadelimtheory
%
\endisadelimtheory
%
\isatagtheory
\isakeywordONE{theory}\isamarkupfalse%
\ TEMPORAL{\isacharunderscore}{\kern0pt}PDDL{\isacharunderscore}{\kern0pt}Semantics{\isacharunderscore}{\kern0pt}Alt\isanewline
\isakeywordTWO{imports}\isanewline
\ \ TEMPORAL{\isacharunderscore}{\kern0pt}PDDL{\isacharunderscore}{\kern0pt}Semantics\isanewline
\ \ TEMPORAL{\isacharunderscore}{\kern0pt}PDDL{\isacharunderscore}{\kern0pt}Checker\ \isanewline
\isakeywordTWO{begin}%
\endisatagtheory
{\isafoldtheory}%
%
\isadelimtheory
%
\endisadelimtheory
%
\begin{isamarkuptext}%
Auxiliary lemma with about \isa{\isaconst{dropWhile}}, \isa{\isaconst{takeWhile}} 
  with \isa{\isaconst{le}} and \isa{\isaconst{leq}}%
\end{isamarkuptext}\isamarkuptrue%
\ \ \isakeywordONE{lemma}\isamarkupfalse%
\ strict{\isacharunderscore}{\kern0pt}sorted{\isacharunderscore}{\kern0pt}drop{\isacharunderscore}{\kern0pt}le{\isacharunderscore}{\kern0pt}take{\isacharunderscore}{\kern0pt}leq{\isacharunderscore}{\kern0pt}comm{\isacharcolon}{\kern0pt}\isanewline
\ \ \ \ \isakeywordTWO{assumes}\ {\isachardoublequoteopen}t\isactrlsub i\ {\isasymle}\ t\isactrlsub j{\isachardoublequoteclose}\ \isanewline
\ \ \ \ \ \ \ \ \isakeywordTWO{and}\ {\isachardoublequoteopen}strict{\isacharunderscore}{\kern0pt}sorted\ {\isacharparenleft}{\kern0pt}map\ fst\ hs{\isacharparenright}{\kern0pt}{\isachardoublequoteclose}\isanewline
\ \ \ \ \isakeywordTWO{shows}\ {\isachardoublequoteopen}dropWhile\ {\isacharparenleft}{\kern0pt}le\ t\isactrlsub i{\isacharparenright}{\kern0pt}\ {\isacharparenleft}{\kern0pt}takeWhile\ {\isacharparenleft}{\kern0pt}leq\ t\isactrlsub j{\isacharparenright}{\kern0pt}\ hs{\isacharparenright}{\kern0pt}\ {\isacharequal}{\kern0pt}\ takeWhile\ {\isacharparenleft}{\kern0pt}leq\ t\isactrlsub j{\isacharparenright}{\kern0pt}\ {\isacharparenleft}{\kern0pt}dropWhile\ {\isacharparenleft}{\kern0pt}le\ t\isactrlsub i{\isacharparenright}{\kern0pt}\ hs{\isacharparenright}{\kern0pt}{\isachardoublequoteclose}\isanewline
%
\isadelimproof
\ \ \ \ %
\endisadelimproof
%
\isatagproof
\isakeywordONE{using}\isamarkupfalse%
\ assms\ \isanewline
\ \ \ \ \isakeywordONE{by}\isamarkupfalse%
\ {\isacharparenleft}{\kern0pt}induction\ hs{\isacharparenright}{\kern0pt}\ {\isacharparenleft}{\kern0pt}auto\ simp{\isacharcolon}{\kern0pt}\ le{\isacharunderscore}{\kern0pt}def\ leq{\isacharunderscore}{\kern0pt}def{\isacharparenright}{\kern0pt}%
\endisatagproof
{\isafoldproof}%
%
\isadelimproof
\isanewline
%
\endisadelimproof
\isanewline
\ \ \isakeywordONE{lemma}\isamarkupfalse%
\ strict{\isacharunderscore}{\kern0pt}sorted{\isacharunderscore}{\kern0pt}drop{\isacharunderscore}{\kern0pt}leq{\isacharunderscore}{\kern0pt}take{\isacharunderscore}{\kern0pt}le{\isacharunderscore}{\kern0pt}comm{\isacharcolon}{\kern0pt}\isanewline
\ \ \ \ \isakeywordTWO{assumes}\ {\isachardoublequoteopen}t\isactrlsub i\ {\isacharless}{\kern0pt}\ t\isactrlsub j{\isachardoublequoteclose}\ \isanewline
\ \ \ \ \ \ \ \ \isakeywordTWO{and}\ {\isachardoublequoteopen}strict{\isacharunderscore}{\kern0pt}sorted\ {\isacharparenleft}{\kern0pt}map\ fst\ hs{\isacharparenright}{\kern0pt}{\isachardoublequoteclose}\isanewline
\ \ \ \ \isakeywordTWO{shows}\ {\isachardoublequoteopen}dropWhile\ {\isacharparenleft}{\kern0pt}leq\ t\isactrlsub i{\isacharparenright}{\kern0pt}\ {\isacharparenleft}{\kern0pt}takeWhile\ {\isacharparenleft}{\kern0pt}le\ t\isactrlsub j{\isacharparenright}{\kern0pt}\ hs{\isacharparenright}{\kern0pt}\ {\isacharequal}{\kern0pt}\ takeWhile\ {\isacharparenleft}{\kern0pt}le\ t\isactrlsub j{\isacharparenright}{\kern0pt}\ {\isacharparenleft}{\kern0pt}dropWhile\ {\isacharparenleft}{\kern0pt}leq\ t\isactrlsub i{\isacharparenright}{\kern0pt}\ hs{\isacharparenright}{\kern0pt}{\isachardoublequoteclose}\isanewline
%
\isadelimproof
\ \ \ \ %
\endisadelimproof
%
\isatagproof
\isakeywordONE{using}\isamarkupfalse%
\ assms\ \isanewline
\ \ \ \ \isakeywordONE{by}\isamarkupfalse%
\ {\isacharparenleft}{\kern0pt}induction\ hs{\isacharparenright}{\kern0pt}\ {\isacharparenleft}{\kern0pt}auto\ simp{\isacharcolon}{\kern0pt}\ le{\isacharunderscore}{\kern0pt}def\ leq{\isacharunderscore}{\kern0pt}def{\isacharparenright}{\kern0pt}%
\endisatagproof
{\isafoldproof}%
%
\isadelimproof
\isanewline
%
\endisadelimproof
\isanewline
\ \ \isakeywordONE{lemma}\isamarkupfalse%
\ strict{\isacharunderscore}{\kern0pt}sorted{\isacharunderscore}{\kern0pt}drop{\isacharunderscore}{\kern0pt}leq{\isacharunderscore}{\kern0pt}take{\isacharunderscore}{\kern0pt}leq{\isacharunderscore}{\kern0pt}comm{\isacharcolon}{\kern0pt}\isanewline
\ \ \ \ \isakeywordTWO{assumes}\ {\isachardoublequoteopen}t\isactrlsub i\ {\isasymle}\ t\isactrlsub j{\isachardoublequoteclose}\ \isanewline
\ \ \ \ \ \ \ \ \isakeywordTWO{and}\ {\isachardoublequoteopen}strict{\isacharunderscore}{\kern0pt}sorted\ {\isacharparenleft}{\kern0pt}map\ fst\ hs{\isacharparenright}{\kern0pt}{\isachardoublequoteclose}\isanewline
\ \ \ \ \isakeywordTWO{shows}\ {\isachardoublequoteopen}dropWhile\ {\isacharparenleft}{\kern0pt}leq\ t\isactrlsub i{\isacharparenright}{\kern0pt}\ {\isacharparenleft}{\kern0pt}takeWhile\ {\isacharparenleft}{\kern0pt}leq\ t\isactrlsub j{\isacharparenright}{\kern0pt}\ hs{\isacharparenright}{\kern0pt}\ {\isacharequal}{\kern0pt}\ takeWhile\ {\isacharparenleft}{\kern0pt}leq\ t\isactrlsub j{\isacharparenright}{\kern0pt}\ {\isacharparenleft}{\kern0pt}dropWhile\ {\isacharparenleft}{\kern0pt}leq\ t\isactrlsub i{\isacharparenright}{\kern0pt}\ hs{\isacharparenright}{\kern0pt}{\isachardoublequoteclose}\isanewline
%
\isadelimproof
\ \ \ \ %
\endisadelimproof
%
\isatagproof
\isakeywordONE{using}\isamarkupfalse%
\ assms\ \isakeywordONE{by}\isamarkupfalse%
\ {\isacharparenleft}{\kern0pt}induction\ hs{\isacharparenright}{\kern0pt}\ {\isacharparenleft}{\kern0pt}auto\ simp{\isacharcolon}{\kern0pt}\ leq{\isacharunderscore}{\kern0pt}def{\isacharparenright}{\kern0pt}%
\endisatagproof
{\isafoldproof}%
%
\isadelimproof
\isanewline
%
\endisadelimproof
\isanewline
\ \ \isakeywordONE{lemma}\isamarkupfalse%
\ strict{\isacharunderscore}{\kern0pt}sorted{\isacharunderscore}{\kern0pt}drop{\isacharunderscore}{\kern0pt}le{\isacharunderscore}{\kern0pt}take{\isacharunderscore}{\kern0pt}leq{\isacharcolon}{\kern0pt}\isanewline
\ \ \ \ \isakeywordTWO{assumes}\ {\isachardoublequoteopen}strict{\isacharunderscore}{\kern0pt}sorted\ {\isacharparenleft}{\kern0pt}map\ fst\ hs{\isacharparenright}{\kern0pt}{\isachardoublequoteclose}\isanewline
\ \ \ \ \ \ \ \ \isakeywordTWO{and}\ {\isachardoublequoteopen}{\isacharparenleft}{\kern0pt}t\isactrlsub a{\isacharcomma}{\kern0pt}A{\isacharparenright}{\kern0pt}\ {\isasymin}\ set\ hs{\isachardoublequoteclose}\isanewline
\ \ \ \ \isakeywordTWO{shows}\ {\isachardoublequoteopen}{\isacharparenleft}{\kern0pt}t\isactrlsub a{\isacharcomma}{\kern0pt}A{\isacharparenright}{\kern0pt}\ {\isasymin}\ set\ {\isacharparenleft}{\kern0pt}dropWhile\ {\isacharparenleft}{\kern0pt}le\ t\isactrlsub i{\isacharparenright}{\kern0pt}\ {\isacharparenleft}{\kern0pt}takeWhile\ {\isacharparenleft}{\kern0pt}leq\ t\isactrlsub j{\isacharparenright}{\kern0pt}\ hs{\isacharparenright}{\kern0pt}{\isacharparenright}{\kern0pt}\ {\isasymlongleftrightarrow}\ t\isactrlsub i\ {\isasymle}\ t\isactrlsub a\ {\isasymand}\ t\isactrlsub a\ {\isasymle}\ t\isactrlsub j{\isachardoublequoteclose}\isanewline
%
\isadelimproof
\ \ \ \ %
\endisadelimproof
%
\isatagproof
\isakeywordONE{unfolding}\isamarkupfalse%
\ leq{\isacharunderscore}{\kern0pt}def\ le{\isacharunderscore}{\kern0pt}def\isanewline
\ \ \ \ \isakeywordONE{using}\isamarkupfalse%
\ assms\isanewline
\ \ \ \ \isakeywordONE{apply}\isamarkupfalse%
\ {\isacharparenleft}{\kern0pt}induction\ hs{\isacharparenright}{\kern0pt}\isanewline
\ \ \ \ \isakeywordONE{apply}\isamarkupfalse%
\ {\isacharparenleft}{\kern0pt}auto\ simp{\isacharcolon}{\kern0pt}\ leq{\isacharunderscore}{\kern0pt}def\ le{\isacharunderscore}{\kern0pt}def{\isacharparenright}{\kern0pt}\isanewline
\ \ \ \ \isakeywordONE{apply}\isamarkupfalse%
\ {\isacharparenleft}{\kern0pt}meson\ set{\isacharunderscore}{\kern0pt}dropWhileD\ set{\isacharunderscore}{\kern0pt}takeWhileD{\isacharparenright}{\kern0pt}\isanewline
\ \ \ \ \isakeywordONE{apply}\isamarkupfalse%
\ auto{\isacharbrackleft}{\kern0pt}{\isadigit{1}}{\isacharbrackright}{\kern0pt}\isanewline
\ \ \ \ \isakeywordONE{using}\isamarkupfalse%
\ set{\isacharunderscore}{\kern0pt}takeWhileD\ \isakeywordONE{apply}\isamarkupfalse%
\ fastforce\isanewline
\ \ \ \ \isakeywordONE{apply}\isamarkupfalse%
\ {\isacharparenleft}{\kern0pt}meson\ set{\isacharunderscore}{\kern0pt}dropWhileD{\isacharparenright}{\kern0pt}{\isacharplus}{\kern0pt}\isanewline
\ \ \ \ \isakeywordONE{done}\isamarkupfalse%
%
\endisatagproof
{\isafoldproof}%
%
\isadelimproof
%
\endisadelimproof
\ \isanewline
\ \ \isanewline
\ \ \isakeywordONE{lemma}\isamarkupfalse%
\ strict{\isacharunderscore}{\kern0pt}sorted{\isacharunderscore}{\kern0pt}drop{\isacharunderscore}{\kern0pt}leq{\isacharunderscore}{\kern0pt}take{\isacharunderscore}{\kern0pt}leq{\isacharcolon}{\kern0pt}\isanewline
\ \ \ \ \isakeywordTWO{assumes}\ {\isachardoublequoteopen}strict{\isacharunderscore}{\kern0pt}sorted\ {\isacharparenleft}{\kern0pt}map\ fst\ hs{\isacharparenright}{\kern0pt}{\isachardoublequoteclose}\isanewline
\ \ \ \ \ \ \ \ \isakeywordTWO{and}\ {\isachardoublequoteopen}{\isacharparenleft}{\kern0pt}t\isactrlsub a{\isacharcomma}{\kern0pt}A{\isacharparenright}{\kern0pt}\ {\isasymin}\ set\ hs{\isachardoublequoteclose}\isanewline
\ \ \ \ \isakeywordTWO{shows}\ {\isachardoublequoteopen}{\isacharparenleft}{\kern0pt}t\isactrlsub a{\isacharcomma}{\kern0pt}A{\isacharparenright}{\kern0pt}\ {\isasymin}\ set\ {\isacharparenleft}{\kern0pt}dropWhile\ {\isacharparenleft}{\kern0pt}leq\ t\isactrlsub i{\isacharparenright}{\kern0pt}\ {\isacharparenleft}{\kern0pt}takeWhile\ {\isacharparenleft}{\kern0pt}leq\ t\isactrlsub j{\isacharparenright}{\kern0pt}\ hs{\isacharparenright}{\kern0pt}{\isacharparenright}{\kern0pt}\ {\isasymlongleftrightarrow}\ t\isactrlsub i\ {\isacharless}{\kern0pt}\ t\isactrlsub a\ {\isasymand}\ t\isactrlsub a\ {\isasymle}\ t\isactrlsub j{\isachardoublequoteclose}\isanewline
%
\isadelimproof
\ \ \ \ %
\endisadelimproof
%
\isatagproof
\isakeywordONE{unfolding}\isamarkupfalse%
\ leq{\isacharunderscore}{\kern0pt}def\isanewline
\ \ \ \ \isakeywordONE{using}\isamarkupfalse%
\ assms\isanewline
\ \ \ \ \isakeywordONE{apply}\isamarkupfalse%
\ {\isacharparenleft}{\kern0pt}induction\ hs{\isacharparenright}{\kern0pt}\isanewline
\ \ \ \ \isakeywordONE{apply}\isamarkupfalse%
\ auto\isanewline
\ \ \ \ \isakeywordONE{apply}\isamarkupfalse%
\ {\isacharparenleft}{\kern0pt}meson\ set{\isacharunderscore}{\kern0pt}dropWhileD\ set{\isacharunderscore}{\kern0pt}takeWhileD{\isacharparenright}{\kern0pt}\isanewline
\ \ \ \ \isakeywordONE{apply}\isamarkupfalse%
\ auto{\isacharbrackleft}{\kern0pt}{\isadigit{1}}{\isacharbrackright}{\kern0pt}\isanewline
\ \ \ \ \isakeywordONE{using}\isamarkupfalse%
\ set{\isacharunderscore}{\kern0pt}takeWhileD\ \isakeywordONE{apply}\isamarkupfalse%
\ fastforce\isanewline
\ \ \ \ \isakeywordONE{apply}\isamarkupfalse%
\ {\isacharparenleft}{\kern0pt}meson\ set{\isacharunderscore}{\kern0pt}dropWhileD{\isacharparenright}{\kern0pt}{\isacharplus}{\kern0pt}\isanewline
\ \ \ \ \isakeywordONE{done}\isamarkupfalse%
%
\endisatagproof
{\isafoldproof}%
%
\isadelimproof
%
\endisadelimproof
\ \isanewline
\isanewline
\ \ \isakeywordONE{lemma}\isamarkupfalse%
\ strict{\isacharunderscore}{\kern0pt}sorted{\isacharunderscore}{\kern0pt}drop{\isacharunderscore}{\kern0pt}leq{\isacharunderscore}{\kern0pt}take{\isacharunderscore}{\kern0pt}le{\isacharcolon}{\kern0pt}\isanewline
\ \ \ \ \isakeywordTWO{assumes}\ {\isachardoublequoteopen}strict{\isacharunderscore}{\kern0pt}sorted\ {\isacharparenleft}{\kern0pt}map\ fst\ hs{\isacharparenright}{\kern0pt}{\isachardoublequoteclose}\isanewline
\ \ \ \ \ \ \ \ \isakeywordTWO{and}\ {\isachardoublequoteopen}{\isacharparenleft}{\kern0pt}t\isactrlsub a{\isacharcomma}{\kern0pt}A{\isacharparenright}{\kern0pt}\ {\isasymin}\ set\ hs{\isachardoublequoteclose}\isanewline
\ \ \ \ \isakeywordTWO{shows}\ {\isachardoublequoteopen}{\isacharparenleft}{\kern0pt}t\isactrlsub a{\isacharcomma}{\kern0pt}A{\isacharparenright}{\kern0pt}\ {\isasymin}\ set\ {\isacharparenleft}{\kern0pt}dropWhile\ {\isacharparenleft}{\kern0pt}leq\ t\isactrlsub i{\isacharparenright}{\kern0pt}\ {\isacharparenleft}{\kern0pt}takeWhile\ {\isacharparenleft}{\kern0pt}le\ t\isactrlsub j{\isacharparenright}{\kern0pt}\ hs{\isacharparenright}{\kern0pt}{\isacharparenright}{\kern0pt}\ {\isasymlongleftrightarrow}\ t\isactrlsub i\ {\isacharless}{\kern0pt}\ t\isactrlsub a\ {\isasymand}\ t\isactrlsub a\ {\isacharless}{\kern0pt}\ t\isactrlsub j{\isachardoublequoteclose}\isanewline
%
\isadelimproof
\ \ \ \ %
\endisadelimproof
%
\isatagproof
\isakeywordONE{unfolding}\isamarkupfalse%
\ leq{\isacharunderscore}{\kern0pt}def\ le{\isacharunderscore}{\kern0pt}def\isanewline
\ \ \ \ \isakeywordONE{using}\isamarkupfalse%
\ assms\isanewline
\ \ \ \ \isakeywordONE{apply}\isamarkupfalse%
\ {\isacharparenleft}{\kern0pt}induction\ hs{\isacharparenright}{\kern0pt}\isanewline
\ \ \ \ \isakeywordONE{apply}\isamarkupfalse%
\ auto\isanewline
\ \ \ \ \isakeywordONE{apply}\isamarkupfalse%
\ {\isacharparenleft}{\kern0pt}meson\ set{\isacharunderscore}{\kern0pt}dropWhileD\ set{\isacharunderscore}{\kern0pt}takeWhileD{\isacharparenright}{\kern0pt}\isanewline
\ \ \ \ \isakeywordONE{apply}\isamarkupfalse%
\ auto{\isacharbrackleft}{\kern0pt}{\isadigit{1}}{\isacharbrackright}{\kern0pt}\isanewline
\ \ \ \ \isakeywordONE{using}\isamarkupfalse%
\ set{\isacharunderscore}{\kern0pt}takeWhileD\ \isakeywordONE{apply}\isamarkupfalse%
\ fastforce\isanewline
\ \ \ \ \isakeywordONE{apply}\isamarkupfalse%
\ {\isacharparenleft}{\kern0pt}meson\ set{\isacharunderscore}{\kern0pt}dropWhileD{\isacharparenright}{\kern0pt}{\isacharplus}{\kern0pt}\isanewline
\ \ \ \ \isakeywordONE{done}\isamarkupfalse%
%
\endisatagproof
{\isafoldproof}%
%
\isadelimproof
%
\endisadelimproof
\ \isanewline
\isanewline
\isanewline
\ \ \isakeywordONE{lemma}\isamarkupfalse%
\ split{\isacharunderscore}{\kern0pt}list{\isadigit{2}}{\isacharunderscore}{\kern0pt}le{\isacharcolon}{\kern0pt}\isanewline
\ \ \ \ \isakeywordTWO{assumes}\ {\isachardoublequoteopen}strict{\isacharunderscore}{\kern0pt}sorted\ xs{\isachardoublequoteclose}\isanewline
\ \ \ \ \ \ \ \ \isakeywordTWO{and}\ {\isachardoublequoteopen}x\isactrlsub {\isadigit{1}}\ {\isasymin}\ set\ xs{\isachardoublequoteclose}\isanewline
\ \ \ \ \ \ \ \ \isakeywordTWO{and}\ {\isachardoublequoteopen}x\isactrlsub {\isadigit{2}}\ {\isasymin}\ set\ xs{\isachardoublequoteclose}\isanewline
\ \ \ \ \isakeywordTWO{shows}\ {\isachardoublequoteopen}x\isactrlsub {\isadigit{1}}\ {\isacharless}{\kern0pt}\ x\isactrlsub {\isadigit{2}}\ {\isasymlongleftrightarrow}\ {\isacharparenleft}{\kern0pt}{\isasymexists}xs\isactrlsub {\isadigit{1}}\ xs\isactrlsub {\isadigit{2}}\ xs\isactrlsub {\isadigit{3}}{\isachardot}{\kern0pt}\ xs\ {\isacharequal}{\kern0pt}\ xs\isactrlsub {\isadigit{1}}\ {\isacharat}{\kern0pt}\ x\isactrlsub {\isadigit{1}}\ {\isacharhash}{\kern0pt}\ xs\isactrlsub {\isadigit{2}}\ {\isacharat}{\kern0pt}\ x\isactrlsub {\isadigit{2}}\ {\isacharhash}{\kern0pt}\ xs\isactrlsub {\isadigit{3}}{\isacharparenright}{\kern0pt}{\isachardoublequoteclose}\isanewline
%
\isadelimproof
\ \ %
\endisadelimproof
%
\isatagproof
\isakeywordONE{proof}\isamarkupfalse%
\isanewline
\ \ \ \ \isakeywordTHREE{assume}\isamarkupfalse%
\ {\isachardoublequoteopen}x\isactrlsub {\isadigit{1}}\ {\isacharless}{\kern0pt}\ x\isactrlsub {\isadigit{2}}{\isachardoublequoteclose}\isanewline
\ \ \ \ \isakeywordTHREE{obtain}\isamarkupfalse%
\ xs\isactrlsub {\isadigit{2}}{\isacharprime}{\kern0pt}\ xs\isactrlsub {\isadigit{3}}\ \isakeywordTWO{where}\ {\isachardoublequoteopen}xs\ {\isacharequal}{\kern0pt}\ xs\isactrlsub {\isadigit{2}}{\isacharprime}{\kern0pt}\ {\isacharat}{\kern0pt}\ x\isactrlsub {\isadigit{2}}\ {\isacharhash}{\kern0pt}\ xs\isactrlsub {\isadigit{3}}{\isachardoublequoteclose}\isanewline
\ \ \ \ \ \ \isakeywordONE{using}\isamarkupfalse%
\ {\isacartoucheopen}x\isactrlsub {\isadigit{2}}\ {\isasymin}\ set\ xs{\isacartoucheclose}\isanewline
\ \ \ \ \ \ \isakeywordONE{by}\isamarkupfalse%
\ {\isacharparenleft}{\kern0pt}meson\ split{\isacharunderscore}{\kern0pt}list{\isacharparenright}{\kern0pt}\isanewline
\ \ \ \ \isakeywordONE{then}\isamarkupfalse%
\ \isakeywordONE{have}\isamarkupfalse%
\ {\isachardoublequoteopen}{\isasymforall}x\ {\isasymin}\ set\ xs\isactrlsub {\isadigit{3}}{\isachardot}{\kern0pt}\ x\isactrlsub {\isadigit{2}}\ {\isacharless}{\kern0pt}\ x{\isachardoublequoteclose}\isanewline
\ \ \ \ \ \ \isakeywordONE{using}\isamarkupfalse%
\ {\isacartoucheopen}strict{\isacharunderscore}{\kern0pt}sorted\ xs{\isacartoucheclose}\ sorted{\isacharunderscore}{\kern0pt}wrt{\isacharunderscore}{\kern0pt}append{\isacharbrackleft}{\kern0pt}\isakeywordTWO{where}\ xs{\isacharequal}{\kern0pt}{\isachardoublequoteopen}xs\isactrlsub {\isadigit{2}}{\isacharprime}{\kern0pt}{\isachardoublequoteclose}\ \isakeywordTWO{and}\ ys{\isacharequal}{\kern0pt}{\isachardoublequoteopen}x\isactrlsub {\isadigit{2}}\ {\isacharhash}{\kern0pt}\ xs\isactrlsub {\isadigit{3}}{\isachardoublequoteclose}{\isacharbrackright}{\kern0pt}\isanewline
\ \ \ \ \ \ \isakeywordONE{by}\isamarkupfalse%
\ {\isacharparenleft}{\kern0pt}simp\ add{\isacharcolon}{\kern0pt}\ {\isacartoucheopen}xs\ {\isacharequal}{\kern0pt}\ xs\isactrlsub {\isadigit{2}}{\isacharprime}{\kern0pt}\ {\isacharat}{\kern0pt}\ x\isactrlsub {\isadigit{2}}\ {\isacharhash}{\kern0pt}\ xs\isactrlsub {\isadigit{3}}{\isacartoucheclose}{\isacharparenright}{\kern0pt}\isanewline
\ \ \ \ \isakeywordONE{then}\isamarkupfalse%
\ \isakeywordONE{have}\isamarkupfalse%
\ {\isachardoublequoteopen}x\isactrlsub {\isadigit{1}}\ {\isasymin}\ set\ xs\isactrlsub {\isadigit{2}}{\isacharprime}{\kern0pt}{\isachardoublequoteclose}\isanewline
\ \ \ \ \ \ \isakeywordONE{using}\isamarkupfalse%
\ {\isacartoucheopen}x\isactrlsub {\isadigit{1}}\ {\isacharless}{\kern0pt}\ x\isactrlsub {\isadigit{2}}{\isacartoucheclose}\ {\isacartoucheopen}x\isactrlsub {\isadigit{1}}\ {\isasymin}\ set\ xs{\isacartoucheclose}\ {\isacartoucheopen}xs\ {\isacharequal}{\kern0pt}\ xs\isactrlsub {\isadigit{2}}{\isacharprime}{\kern0pt}\ {\isacharat}{\kern0pt}\ x\isactrlsub {\isadigit{2}}\ {\isacharhash}{\kern0pt}\ xs\isactrlsub {\isadigit{3}}{\isacartoucheclose}\ \isakeywordONE{by}\isamarkupfalse%
\ auto\isanewline
\ \ \ \ \isakeywordONE{then}\isamarkupfalse%
\ \isakeywordTHREE{obtain}\isamarkupfalse%
\ xs\isactrlsub {\isadigit{1}}\ xs\isactrlsub {\isadigit{2}}\ \isakeywordTWO{where}\ {\isachardoublequoteopen}xs\isactrlsub {\isadigit{2}}{\isacharprime}{\kern0pt}\ {\isacharequal}{\kern0pt}\ xs\isactrlsub {\isadigit{1}}\ {\isacharat}{\kern0pt}\ x\isactrlsub {\isadigit{1}}\ {\isacharhash}{\kern0pt}\ xs\isactrlsub {\isadigit{2}}{\isachardoublequoteclose}\isanewline
\ \ \ \ \ \ \isakeywordONE{by}\isamarkupfalse%
\ {\isacharparenleft}{\kern0pt}meson\ split{\isacharunderscore}{\kern0pt}list{\isacharparenright}{\kern0pt}\isanewline
\ \ \ \ \isakeywordONE{then}\isamarkupfalse%
\ \isakeywordTHREE{show}\isamarkupfalse%
\ {\isachardoublequoteopen}{\isasymexists}xs\isactrlsub {\isadigit{1}}\ xs\isactrlsub {\isadigit{2}}\ xs\isactrlsub {\isadigit{3}}{\isachardot}{\kern0pt}\ xs\ {\isacharequal}{\kern0pt}\ xs\isactrlsub {\isadigit{1}}\ {\isacharat}{\kern0pt}\ x\isactrlsub {\isadigit{1}}\ {\isacharhash}{\kern0pt}\ xs\isactrlsub {\isadigit{2}}\ {\isacharat}{\kern0pt}\ x\isactrlsub {\isadigit{2}}\ {\isacharhash}{\kern0pt}\ xs\isactrlsub {\isadigit{3}}{\isachardoublequoteclose}\isanewline
\ \ \ \ \ \ \isakeywordONE{using}\isamarkupfalse%
\ {\isacartoucheopen}xs\ {\isacharequal}{\kern0pt}\ xs\isactrlsub {\isadigit{2}}{\isacharprime}{\kern0pt}\ {\isacharat}{\kern0pt}\ x\isactrlsub {\isadigit{2}}\ {\isacharhash}{\kern0pt}\ xs\isactrlsub {\isadigit{3}}{\isacartoucheclose}\ {\isacartoucheopen}xs\isactrlsub {\isadigit{2}}{\isacharprime}{\kern0pt}\ {\isacharequal}{\kern0pt}\ xs\isactrlsub {\isadigit{1}}\ {\isacharat}{\kern0pt}\ x\isactrlsub {\isadigit{1}}\ {\isacharhash}{\kern0pt}\ xs\isactrlsub {\isadigit{2}}{\isacartoucheclose}\ \isakeywordONE{by}\isamarkupfalse%
\ auto\isanewline
\ \ \isakeywordONE{next}\isamarkupfalse%
\isanewline
\ \ \ \ \isakeywordTHREE{assume}\isamarkupfalse%
\ {\isachardoublequoteopen}{\isasymexists}xs\isactrlsub {\isadigit{1}}\ xs\isactrlsub {\isadigit{2}}\ xs\isactrlsub {\isadigit{3}}{\isachardot}{\kern0pt}\ xs\ {\isacharequal}{\kern0pt}\ xs\isactrlsub {\isadigit{1}}\ {\isacharat}{\kern0pt}\ x\isactrlsub {\isadigit{1}}\ {\isacharhash}{\kern0pt}\ xs\isactrlsub {\isadigit{2}}\ {\isacharat}{\kern0pt}\ x\isactrlsub {\isadigit{2}}\ {\isacharhash}{\kern0pt}\ xs\isactrlsub {\isadigit{3}}{\isachardoublequoteclose}\isanewline
\ \ \ \ \isakeywordONE{then}\isamarkupfalse%
\ \isakeywordTHREE{obtain}\isamarkupfalse%
\ xs\isactrlsub {\isadigit{1}}\ xs\isactrlsub {\isadigit{2}}\ xs\isactrlsub {\isadigit{3}}\ \isakeywordTWO{where}\ {\isachardoublequoteopen}xs\ {\isacharequal}{\kern0pt}\ xs\isactrlsub {\isadigit{1}}\ {\isacharat}{\kern0pt}\ x\isactrlsub {\isadigit{1}}\ {\isacharhash}{\kern0pt}\ xs\isactrlsub {\isadigit{2}}\ {\isacharat}{\kern0pt}\ x\isactrlsub {\isadigit{2}}\ {\isacharhash}{\kern0pt}\ xs\isactrlsub {\isadigit{3}}{\isachardoublequoteclose}\isanewline
\ \ \ \ \ \ \isakeywordONE{by}\isamarkupfalse%
\ auto\isanewline
\ \ \ \ \isakeywordONE{then}\isamarkupfalse%
\ \isakeywordTHREE{show}\isamarkupfalse%
\ {\isachardoublequoteopen}x\isactrlsub {\isadigit{1}}\ {\isacharless}{\kern0pt}\ x\isactrlsub {\isadigit{2}}{\isachardoublequoteclose}\isanewline
\ \ \ \ \ \ \isakeywordONE{using}\isamarkupfalse%
\ {\isacartoucheopen}strict{\isacharunderscore}{\kern0pt}sorted\ xs{\isacartoucheclose}\isanewline
\ \ \ \ \ \ \ \ \ \ \ \ sorted{\isacharunderscore}{\kern0pt}wrt{\isacharunderscore}{\kern0pt}append{\isacharbrackleft}{\kern0pt}\isakeywordTWO{where}\ xs{\isacharequal}{\kern0pt}{\isachardoublequoteopen}xs\isactrlsub {\isadigit{1}}\ {\isacharat}{\kern0pt}\ x\isactrlsub {\isadigit{1}}\ {\isacharhash}{\kern0pt}\ xs\isactrlsub {\isadigit{2}}{\isachardoublequoteclose}\ \isakeywordTWO{and}\ ys{\isacharequal}{\kern0pt}{\isachardoublequoteopen}x\isactrlsub {\isadigit{2}}\ {\isacharhash}{\kern0pt}\ xs\isactrlsub {\isadigit{3}}{\isachardoublequoteclose}{\isacharbrackright}{\kern0pt}\isanewline
\ \ \ \ \ \ \isakeywordONE{by}\isamarkupfalse%
\ {\isacharparenleft}{\kern0pt}simp\ add{\isacharcolon}{\kern0pt}\ {\isacartoucheopen}xs\ {\isacharequal}{\kern0pt}\ xs\isactrlsub {\isadigit{1}}\ {\isacharat}{\kern0pt}\ x\isactrlsub {\isadigit{1}}\ {\isacharhash}{\kern0pt}\ xs\isactrlsub {\isadigit{2}}\ {\isacharat}{\kern0pt}\ x\isactrlsub {\isadigit{2}}\ {\isacharhash}{\kern0pt}\ xs\isactrlsub {\isadigit{3}}{\isacartoucheclose}{\isacharparenright}{\kern0pt}\isanewline
\ \ \isakeywordONE{qed}\isamarkupfalse%
%
\endisatagproof
{\isafoldproof}%
%
\isadelimproof
\isanewline
%
\endisadelimproof
\isanewline
\ \ \isakeywordONE{lemma}\isamarkupfalse%
\ drop{\isacharunderscore}{\kern0pt}take{\isacharunderscore}{\kern0pt}leq\isactrlsub i{\isacharunderscore}{\kern0pt}leq{\isacharunderscore}{\kern0pt}leq\isactrlsub j{\isacharunderscore}{\kern0pt}aux{\isacharcolon}{\kern0pt}\isanewline
\ \ \ \ \isakeywordTWO{assumes}\ {\isachardoublequoteopen}t\isactrlsub k\ {\isasymle}\ t\isactrlsub j{\isachardoublequoteclose}\isanewline
\ \ \ \ \ \ \ \ \isakeywordTWO{and}\ {\isachardoublequoteopen}strict{\isacharunderscore}{\kern0pt}sorted\ {\isacharparenleft}{\kern0pt}map\ fst\ hs{\isacharparenright}{\kern0pt}{\isachardoublequoteclose}\isanewline
\ \ \ \ \isakeywordTWO{shows}\ {\isachardoublequoteopen}takeWhile\ {\isacharparenleft}{\kern0pt}leq\ t\isactrlsub k{\isacharparenright}{\kern0pt}\ hs\ {\isacharat}{\kern0pt}\ takeWhile\ {\isacharparenleft}{\kern0pt}leq\ t\isactrlsub j{\isacharparenright}{\kern0pt}\ {\isacharparenleft}{\kern0pt}dropWhile\ {\isacharparenleft}{\kern0pt}leq\ t\isactrlsub k{\isacharparenright}{\kern0pt}\ hs{\isacharparenright}{\kern0pt}\ {\isacharequal}{\kern0pt}\ takeWhile\ {\isacharparenleft}{\kern0pt}leq\ t\isactrlsub j{\isacharparenright}{\kern0pt}\ hs{\isachardoublequoteclose}\isanewline
%
\isadelimproof
\ \ \ \ %
\endisadelimproof
%
\isatagproof
\isakeywordONE{using}\isamarkupfalse%
\ assms\ \isakeywordONE{by}\isamarkupfalse%
\ {\isacharparenleft}{\kern0pt}induction\ hs{\isacharparenright}{\kern0pt}\ {\isacharparenleft}{\kern0pt}auto\ simp{\isacharcolon}{\kern0pt}\ leq{\isacharunderscore}{\kern0pt}def{\isacharparenright}{\kern0pt}%
\endisatagproof
{\isafoldproof}%
%
\isadelimproof
\isanewline
%
\endisadelimproof
\isanewline
\ \ \isakeywordONE{lemma}\isamarkupfalse%
\ drop{\isacharunderscore}{\kern0pt}take{\isacharunderscore}{\kern0pt}leq\isactrlsub i{\isacharunderscore}{\kern0pt}leq{\isacharunderscore}{\kern0pt}leq\isactrlsub j{\isacharunderscore}{\kern0pt}simp{\isacharcolon}{\kern0pt}\isanewline
\ \ \ \ \isakeywordTWO{assumes}\ {\isachardoublequoteopen}t\isactrlsub i\ {\isasymle}\ t\isactrlsub k{\isachardoublequoteclose}\ \isakeywordTWO{and}\ {\isachardoublequoteopen}t\isactrlsub k\ {\isasymle}\ t\isactrlsub j{\isachardoublequoteclose}\isanewline
\ \ \ \ \ \ \ \ \isakeywordTWO{and}\ {\isachardoublequoteopen}strict{\isacharunderscore}{\kern0pt}sorted\ {\isacharparenleft}{\kern0pt}map\ fst\ hs{\isacharparenright}{\kern0pt}{\isachardoublequoteclose}\isanewline
\ \ \ \ \ \ \isakeywordTWO{shows}\ {\isachardoublequoteopen}dropWhile\ {\isacharparenleft}{\kern0pt}leq\ t\isactrlsub i{\isacharparenright}{\kern0pt}\ {\isacharparenleft}{\kern0pt}takeWhile\ {\isacharparenleft}{\kern0pt}leq\ t\isactrlsub k{\isacharparenright}{\kern0pt}\ hs{\isacharparenright}{\kern0pt}\ {\isacharat}{\kern0pt}\ dropWhile\ {\isacharparenleft}{\kern0pt}leq\ t\isactrlsub k{\isacharparenright}{\kern0pt}\ {\isacharparenleft}{\kern0pt}takeWhile\ {\isacharparenleft}{\kern0pt}leq\ t\isactrlsub j{\isacharparenright}{\kern0pt}\ hs{\isacharparenright}{\kern0pt}\ \isanewline
\ \ \ \ \ \ \ \ \ \ \ \ \ \ {\isacharequal}{\kern0pt}\ dropWhile\ {\isacharparenleft}{\kern0pt}leq\ t\isactrlsub i{\isacharparenright}{\kern0pt}\ {\isacharparenleft}{\kern0pt}takeWhile\ {\isacharparenleft}{\kern0pt}leq\ t\isactrlsub j{\isacharparenright}{\kern0pt}\ hs{\isacharparenright}{\kern0pt}{\isachardoublequoteclose}\isanewline
%
\isadelimproof
\ \ \ \ %
\endisadelimproof
%
\isatagproof
\isakeywordONE{using}\isamarkupfalse%
\ {\isacartoucheopen}t\isactrlsub i\ {\isasymle}\ t\isactrlsub k{\isacartoucheclose}\ drop{\isacharunderscore}{\kern0pt}take{\isacharunderscore}{\kern0pt}leq\isactrlsub i{\isacharunderscore}{\kern0pt}leq{\isacharunderscore}{\kern0pt}leq\isactrlsub j{\isacharunderscore}{\kern0pt}aux{\isacharbrackleft}{\kern0pt}OF\ {\isacartoucheopen}t\isactrlsub k\ {\isasymle}\ t\isactrlsub j{\isacartoucheclose}\ {\isacartoucheopen}strict{\isacharunderscore}{\kern0pt}sorted\ {\isacharparenleft}{\kern0pt}map\ fst\ hs{\isacharparenright}{\kern0pt}{\isacartoucheclose}{\isacharbrackright}{\kern0pt}\isanewline
\ \ \ \ \ \ \ \ \ \ strict{\isacharunderscore}{\kern0pt}sorted{\isacharunderscore}{\kern0pt}drop{\isacharunderscore}{\kern0pt}leq{\isacharunderscore}{\kern0pt}take{\isacharunderscore}{\kern0pt}leq{\isacharunderscore}{\kern0pt}comm{\isacharbrackleft}{\kern0pt}OF\ {\isacartoucheopen}t\isactrlsub k\ {\isasymle}\ t\isactrlsub j{\isacartoucheclose}\ {\isacartoucheopen}strict{\isacharunderscore}{\kern0pt}sorted\ {\isacharparenleft}{\kern0pt}map\ fst\ hs{\isacharparenright}{\kern0pt}{\isacartoucheclose}{\isacharbrackright}{\kern0pt}\isanewline
\ \ \ \ \isakeywordONE{by}\isamarkupfalse%
\ {\isacharparenleft}{\kern0pt}induction\ hs{\isacharparenright}{\kern0pt}\ {\isacharparenleft}{\kern0pt}auto\ simp{\isacharcolon}{\kern0pt}\ leq{\isacharunderscore}{\kern0pt}def{\isacharparenright}{\kern0pt}%
\endisatagproof
{\isafoldproof}%
%
\isadelimproof
%
\endisadelimproof
%
\isadelimdocument
%
\endisadelimdocument
%
\isatagdocument
%
\isamarkupsubsection{State-Transition Based Semantics%
}
\isamarkuptrue%
%
\endisatagdocument
{\isafolddocument}%
%
\isadelimdocument
%
\endisadelimdocument
\isakeywordONE{context}\isamarkupfalse%
\ ast{\isacharunderscore}{\kern0pt}problem\ \isakeywordTWO{begin}\isanewline
\isanewline
\ \ \isakeywordONE{fun}\isamarkupfalse%
\ inst{\isacharunderscore}{\kern0pt}cond\ {\isacharcolon}{\kern0pt}{\isacharcolon}{\kern0pt}\ {\isachardoublequoteopen}ast{\isacharunderscore}{\kern0pt}action{\isacharunderscore}{\kern0pt}schema\ {\isasymRightarrow}\ object\ list\ {\isasymRightarrow}\ temporal{\isacharunderscore}{\kern0pt}annotation\ {\isasymRightarrow}\ {\isacharparenleft}{\kern0pt}object\ atom{\isacharparenright}{\kern0pt}\ formula{\isachardoublequoteclose}\ \isakeywordTWO{where}\isanewline
\ \ \ \ {\isachardoublequoteopen}inst{\isacharunderscore}{\kern0pt}cond\ {\isacharparenleft}{\kern0pt}Durative{\isacharunderscore}{\kern0pt}Action{\isacharunderscore}{\kern0pt}Schema\ n\ params\ d\ cond\ eff{\isacharparenright}{\kern0pt}\ args\ ts\ {\isacharequal}{\kern0pt}\ \isanewline
\ \ \ \ \ \ {\isacharparenleft}{\kern0pt}let\ tsubst\ {\isacharequal}{\kern0pt}\ subst{\isacharunderscore}{\kern0pt}term\ {\isacharparenleft}{\kern0pt}the\ o\ {\isacharparenleft}{\kern0pt}map{\isacharunderscore}{\kern0pt}of\ {\isacharparenleft}{\kern0pt}zip\ {\isacharparenleft}{\kern0pt}map\ fst\ params{\isacharparenright}{\kern0pt}\ args{\isacharparenright}{\kern0pt}{\isacharparenright}{\kern0pt}{\isacharparenright}{\kern0pt}\ in\ \isanewline
\ \ \ \ \ \ \ \ \ \ {\isacharparenleft}{\kern0pt}map{\isacharunderscore}{\kern0pt}formula\ o\ map{\isacharunderscore}{\kern0pt}atom{\isacharparenright}{\kern0pt}\ tsubst\ {\isacharparenleft}{\kern0pt}BigAnd\ {\isacharparenleft}{\kern0pt}filter{\isacharunderscore}{\kern0pt}time{\isacharunderscore}{\kern0pt}spec\ ts\ cond{\isacharparenright}{\kern0pt}{\isacharparenright}{\kern0pt}{\isacharparenright}{\kern0pt}{\isachardoublequoteclose}\isanewline
\isanewline
\ \ \isakeywordONE{fun}\isamarkupfalse%
\ res{\isacharunderscore}{\kern0pt}inst{\isacharunderscore}{\kern0pt}inv\ {\isacharcolon}{\kern0pt}{\isacharcolon}{\kern0pt}\ {\isachardoublequoteopen}plan{\isacharunderscore}{\kern0pt}action\ {\isasymRightarrow}\ {\isacharparenleft}{\kern0pt}object\ atom{\isacharparenright}{\kern0pt}\ formula\ option{\isachardoublequoteclose}\ \isakeywordTWO{where}\ \isanewline
\ \ \ \ {\isachardoublequoteopen}res{\isacharunderscore}{\kern0pt}inst{\isacharunderscore}{\kern0pt}inv\ {\isacharparenleft}{\kern0pt}Simple{\isacharunderscore}{\kern0pt}Plan{\isacharunderscore}{\kern0pt}Action\ n\ args{\isacharparenright}{\kern0pt}\ {\isacharequal}{\kern0pt}\ None{\isachardoublequoteclose}\isanewline
\ \ {\isacharbar}{\kern0pt}\ {\isachardoublequoteopen}res{\isacharunderscore}{\kern0pt}inst{\isacharunderscore}{\kern0pt}inv\ {\isacharparenleft}{\kern0pt}Durative{\isacharunderscore}{\kern0pt}Plan{\isacharunderscore}{\kern0pt}Action\ n\ args\ d{\isacharparenright}{\kern0pt}\ {\isacharequal}{\kern0pt}\ \isanewline
\ \ \ \ \ \ Some\ {\isacharparenleft}{\kern0pt}inst{\isacharunderscore}{\kern0pt}cond\ {\isacharparenleft}{\kern0pt}the\ {\isacharparenleft}{\kern0pt}resolve{\isacharunderscore}{\kern0pt}action{\isacharunderscore}{\kern0pt}schema\ n{\isacharparenright}{\kern0pt}{\isacharparenright}{\kern0pt}\ args\ Over{\isacharunderscore}{\kern0pt}All{\isacharparenright}{\kern0pt}{\isachardoublequoteclose}\isanewline
\isanewline
\ \ \isakeywordONE{lemma}\isamarkupfalse%
\ res{\isacharunderscore}{\kern0pt}inst{\isacharunderscore}{\kern0pt}inv{\isacharunderscore}{\kern0pt}refine{\isacharcolon}{\kern0pt}\isanewline
\ \ \ \ \isakeywordTWO{assumes}\ {\isachardoublequoteopen}wf{\isacharunderscore}{\kern0pt}plan{\isacharunderscore}{\kern0pt}action\ {\isasympi}{\isachardoublequoteclose}\isanewline
\ \ \ \ \ \ \ \ \isakeywordTWO{and}\ {\isachardoublequoteopen}Some\ a\ {\isacharequal}{\kern0pt}\ res{\isacharunderscore}{\kern0pt}inst{\isacharunderscore}{\kern0pt}snap{\isacharunderscore}{\kern0pt}action\ {\isasympi}\ Over{\isacharunderscore}{\kern0pt}All{\isachardoublequoteclose}\isanewline
\ \ \ \ \isakeywordTWO{shows}\ {\isachardoublequoteopen}Some\ {\isacharparenleft}{\kern0pt}precondition\ a{\isacharparenright}{\kern0pt}\ {\isacharequal}{\kern0pt}\ res{\isacharunderscore}{\kern0pt}inst{\isacharunderscore}{\kern0pt}inv\ {\isasympi}{\isachardoublequoteclose}\isanewline
%
\isadelimproof
\ \ \ \ %
\endisadelimproof
%
\isatagproof
\isakeywordONE{using}\isamarkupfalse%
\ assms\isanewline
\ \ \ \ \isakeywordONE{by}\isamarkupfalse%
\ {\isacharparenleft}{\kern0pt}cases\ {\isasympi}{\isacharparenright}{\kern0pt}\ {\isacharparenleft}{\kern0pt}auto\ simp{\isacharcolon}{\kern0pt}\ Let{\isacharunderscore}{\kern0pt}def\ split{\isacharcolon}{\kern0pt}\ option{\isachardot}{\kern0pt}splits\ ast{\isacharunderscore}{\kern0pt}action{\isacharunderscore}{\kern0pt}schema{\isachardot}{\kern0pt}splits{\isacharparenright}{\kern0pt}%
\endisatagproof
{\isafoldproof}%
%
\isadelimproof
%
\endisadelimproof
%
\begin{isamarkuptext}%
Set of all ground actions that, are executed at a given time point in the 
  induced happening sequence.%
\end{isamarkuptext}\isamarkuptrue%
\ \ \isakeywordONE{definition}\isamarkupfalse%
\ acts{\isacharunderscore}{\kern0pt}of{\isacharunderscore}{\kern0pt}plan{\isacharunderscore}{\kern0pt}at\ {\isacharcolon}{\kern0pt}{\isacharcolon}{\kern0pt}\ {\isachardoublequoteopen}time\ {\isasymRightarrow}\ plan\ {\isasymRightarrow}\ ground{\isacharunderscore}{\kern0pt}action\ set{\isachardoublequoteclose}\ \isakeywordTWO{where}\isanewline
\ \ \ \ {\isachardoublequoteopen}acts{\isacharunderscore}{\kern0pt}of{\isacharunderscore}{\kern0pt}plan{\isacharunderscore}{\kern0pt}at\ t\isactrlsub i\ {\isasympi}s\ {\isacharequal}{\kern0pt}\ \isanewline
\ \ \ \ \ \ {\isacharbraceleft}{\kern0pt}a\isactrlsub {\isasympi}{\isachardot}{\kern0pt}\ {\isasymexists}{\isasympi}{\isachardot}{\kern0pt}\ {\isacharparenleft}{\kern0pt}t\isactrlsub i{\isacharcomma}{\kern0pt}{\isasympi}{\isacharparenright}{\kern0pt}\ {\isasymin}\ simple{\isacharunderscore}{\kern0pt}acts\ {\isasympi}s\ {\isasymand}\ Some\ a\isactrlsub {\isasympi}\ {\isacharequal}{\kern0pt}\ res{\isacharunderscore}{\kern0pt}inst\ {\isasympi}\ At{\isacharunderscore}{\kern0pt}Start{\isacharbraceright}{\kern0pt}\ \isanewline
\ \ \ \ {\isasymunion}\ {\isacharbraceleft}{\kern0pt}a\isactrlsub s\isactrlsub t\isactrlsub a\isactrlsub r\isactrlsub t{\isachardot}{\kern0pt}\ {\isasymexists}{\isasympi}{\isachardot}{\kern0pt}\ {\isacharparenleft}{\kern0pt}t\isactrlsub i{\isacharcomma}{\kern0pt}{\isasympi}{\isacharparenright}{\kern0pt}\ {\isasymin}\ durative{\isacharunderscore}{\kern0pt}acts\ {\isasympi}s\ {\isasymand}\ Some\ a\isactrlsub s\isactrlsub t\isactrlsub a\isactrlsub r\isactrlsub t\ {\isacharequal}{\kern0pt}\ res{\isacharunderscore}{\kern0pt}inst{\isacharunderscore}{\kern0pt}snap{\isacharunderscore}{\kern0pt}action\ {\isasympi}\ At{\isacharunderscore}{\kern0pt}Start{\isacharbraceright}{\kern0pt}\ \isanewline
\ \ \ \ {\isasymunion}\ {\isacharbraceleft}{\kern0pt}a\isactrlsub e\isactrlsub n\isactrlsub d{\isachardot}{\kern0pt}\ {\isasymexists}t{\isacharprime}{\kern0pt}\ {\isasympi}{\isachardot}{\kern0pt}\ {\isacharparenleft}{\kern0pt}t{\isacharprime}{\kern0pt}{\isacharcomma}{\kern0pt}{\isasympi}{\isacharparenright}{\kern0pt}\ {\isasymin}\ durative{\isacharunderscore}{\kern0pt}acts\ {\isasympi}s\ {\isasymand}\ t\isactrlsub i\ {\isacharequal}{\kern0pt}\ t{\isacharprime}{\kern0pt}\ {\isacharplus}{\kern0pt}\ duration\ {\isasympi}\ {\isasymand}\ Some\ a\isactrlsub e\isactrlsub n\isactrlsub d\ {\isacharequal}{\kern0pt}\ res{\isacharunderscore}{\kern0pt}inst{\isacharunderscore}{\kern0pt}snap{\isacharunderscore}{\kern0pt}action\ {\isasympi}\ At{\isacharunderscore}{\kern0pt}End{\isacharbraceright}{\kern0pt}{\isachardoublequoteclose}%
\begin{isamarkuptext}%
Set of all invariant that have to hold in between consecutive happening time points \isa{t{\isacharprime}{\kern0pt}\ {\isacharless}{\kern0pt}\ t}.%
\end{isamarkuptext}\isamarkuptrue%
\ \ \isakeywordONE{definition}\isamarkupfalse%
\ invs{\isacharunderscore}{\kern0pt}of{\isacharunderscore}{\kern0pt}plan{\isacharunderscore}{\kern0pt}at\ {\isacharcolon}{\kern0pt}{\isacharcolon}{\kern0pt}\ {\isachardoublequoteopen}time\ {\isasymRightarrow}\ plan\ {\isasymRightarrow}\ {\isacharparenleft}{\kern0pt}object\ atom{\isacharparenright}{\kern0pt}\ formula\ set{\isachardoublequoteclose}\ \isakeywordTWO{where}\isanewline
\ \ \ \ {\isachardoublequoteopen}invs{\isacharunderscore}{\kern0pt}of{\isacharunderscore}{\kern0pt}plan{\isacharunderscore}{\kern0pt}at\ t\isactrlsub i\ {\isasympi}s\ {\isacharequal}{\kern0pt}\ \isanewline
\ \ \ \ \ \ {\isacharbraceleft}{\kern0pt}inv{\isachardot}{\kern0pt}\ {\isasymexists}t\isactrlsub {\isasympi}\ {\isasympi}{\isachardot}{\kern0pt}\ {\isacharparenleft}{\kern0pt}t\isactrlsub {\isasympi}{\isacharcomma}{\kern0pt}{\isasympi}{\isacharparenright}{\kern0pt}\ {\isasymin}\ durative{\isacharunderscore}{\kern0pt}acts\ {\isasympi}s\ {\isasymand}\ t\isactrlsub {\isasympi}\ {\isacharless}{\kern0pt}\ t\isactrlsub i\ {\isasymand}\ t\isactrlsub i\ {\isasymle}\ t\isactrlsub {\isasympi}\ {\isacharplus}{\kern0pt}\ {\isacharparenleft}{\kern0pt}duration\ {\isasympi}{\isacharparenright}{\kern0pt}\ {\isasymand}\ Some\ inv\ {\isacharequal}{\kern0pt}\ res{\isacharunderscore}{\kern0pt}inst{\isacharunderscore}{\kern0pt}inv\ {\isasympi}{\isacharbraceright}{\kern0pt}{\isachardoublequoteclose}%
\begin{isamarkuptext}%
Predicate for sequence of happening time points.%
\end{isamarkuptext}\isamarkuptrue%
\ \ \isakeywordONE{fun}\isamarkupfalse%
\ apply{\isacharunderscore}{\kern0pt}eff\ {\isacharcolon}{\kern0pt}{\isacharcolon}{\kern0pt}\ {\isachardoublequoteopen}ground{\isacharunderscore}{\kern0pt}action\ set\ {\isasymRightarrow}\ world{\isacharunderscore}{\kern0pt}model\ {\isasymRightarrow}\ world{\isacharunderscore}{\kern0pt}model{\isachardoublequoteclose}\ \isakeywordTWO{where}\isanewline
\ \ \ \ {\isachardoublequoteopen}apply{\isacharunderscore}{\kern0pt}eff\ A\isactrlsub i\ M\ {\isacharequal}{\kern0pt}\ {\isacharparenleft}{\kern0pt}M\ {\isacharminus}{\kern0pt}\ {\isasymUnion}\ {\isacharparenleft}{\kern0pt}set\ {\isacharbackquote}{\kern0pt}\ dels\ {\isacharbackquote}{\kern0pt}\ effect\ {\isacharbackquote}{\kern0pt}\ A\isactrlsub i{\isacharparenright}{\kern0pt}{\isacharparenright}{\kern0pt}\ {\isasymunion}\ {\isasymUnion}\ {\isacharparenleft}{\kern0pt}set\ {\isacharbackquote}{\kern0pt}\ adds\ {\isacharbackquote}{\kern0pt}\ effect\ {\isacharbackquote}{\kern0pt}\ A\isactrlsub i{\isacharparenright}{\kern0pt}{\isachardoublequoteclose}%
\begin{isamarkuptext}%
Instead of explicitly defining an induced happening sequence. Here the 
  happening sequence is only implicitly constructed. Only happening time points are considered 
  for the valid execution, invariants are treated like preconditions of an entire happening.%
\end{isamarkuptext}\isamarkuptrue%
\ \ \isakeywordONE{fun}\isamarkupfalse%
\ valid{\isacharunderscore}{\kern0pt}state{\isacharunderscore}{\kern0pt}seq\ {\isacharcolon}{\kern0pt}{\isacharcolon}{\kern0pt}\ {\isachardoublequoteopen}world{\isacharunderscore}{\kern0pt}model\ {\isasymRightarrow}\ time\ list\ {\isasymRightarrow}\ plan\ {\isasymRightarrow}\ world{\isacharunderscore}{\kern0pt}model\ {\isasymRightarrow}\ bool{\isachardoublequoteclose}\ \isakeywordTWO{where}\isanewline
\ \ \ \ {\isachardoublequoteopen}valid{\isacharunderscore}{\kern0pt}state{\isacharunderscore}{\kern0pt}seq\ M\ {\isacharbrackleft}{\kern0pt}{\isacharbrackright}{\kern0pt}\ {\isasympi}s\ M{\isacharprime}{\kern0pt}\ {\isasymlongleftrightarrow}\ {\isacharparenleft}{\kern0pt}M\ {\isacharequal}{\kern0pt}\ M{\isacharprime}{\kern0pt}{\isacharparenright}{\kern0pt}{\isachardoublequoteclose}\isanewline
\ \ {\isacharbar}{\kern0pt}\ {\isachardoublequoteopen}valid{\isacharunderscore}{\kern0pt}state{\isacharunderscore}{\kern0pt}seq\ M\ {\isacharparenleft}{\kern0pt}t\isactrlsub i{\isacharhash}{\kern0pt}ts{\isacharparenright}{\kern0pt}\ {\isasympi}s\ M{\isacharprime}{\kern0pt}\ {\isasymlongleftrightarrow}\ \isanewline
\ \ \ \ {\isacharparenleft}{\kern0pt}let\ A\isactrlsub i\ {\isacharequal}{\kern0pt}\ acts{\isacharunderscore}{\kern0pt}of{\isacharunderscore}{\kern0pt}plan{\isacharunderscore}{\kern0pt}at\ t\isactrlsub i\ {\isasympi}s\ in\isanewline
\ \ \ \ \ \ \ \ {\isacharparenleft}{\kern0pt}{\isasymforall}i\ {\isasymin}\ invs{\isacharunderscore}{\kern0pt}of{\isacharunderscore}{\kern0pt}plan{\isacharunderscore}{\kern0pt}at\ t\isactrlsub i\ {\isasympi}s{\isachardot}{\kern0pt}\ M\ \isactrlsup c{\isasymTTurnstile}\isactrlsub {\isacharequal}{\kern0pt}\ i{\isacharparenright}{\kern0pt}\isanewline
\ \ \ \ \ \ {\isasymand}\ {\isacharparenleft}{\kern0pt}{\isasymforall}a\ {\isasymin}\ A\isactrlsub i{\isachardot}{\kern0pt}\ M\ \isactrlsup c{\isasymTTurnstile}\isactrlsub {\isacharequal}{\kern0pt}\ precondition\ a{\isacharparenright}{\kern0pt}\isanewline
\ \ \ \ \ \ {\isasymand}\ {\isacharparenleft}{\kern0pt}{\isasymforall}a\ {\isasymin}\ A\isactrlsub i{\isachardot}{\kern0pt}\ {\isasymforall}b\ {\isasymin}\ A\isactrlsub i{\isachardot}{\kern0pt}\ a\ {\isasymnoteq}\ b\ {\isasymlongrightarrow}\ acts{\isacharunderscore}{\kern0pt}non{\isacharunderscore}{\kern0pt}intrf\ a\ b{\isacharparenright}{\kern0pt}\isanewline
\ \ \ \ \ \ {\isasymand}\ valid{\isacharunderscore}{\kern0pt}state{\isacharunderscore}{\kern0pt}seq\ {\isacharparenleft}{\kern0pt}apply{\isacharunderscore}{\kern0pt}eff\ A\isactrlsub i\ M{\isacharparenright}{\kern0pt}\ ts\ {\isasympi}s\ M{\isacharprime}{\kern0pt}{\isacharparenright}{\kern0pt}{\isachardoublequoteclose}\isanewline
\isanewline
\ \ \isakeywordONE{definition}\isamarkupfalse%
\ valid{\isacharunderscore}{\kern0pt}plan{\isacharunderscore}{\kern0pt}from\ {\isacharcolon}{\kern0pt}{\isacharcolon}{\kern0pt}\ {\isachardoublequoteopen}world{\isacharunderscore}{\kern0pt}model\ {\isasymRightarrow}\ plan\ {\isasymRightarrow}\ bool{\isachardoublequoteclose}\ \isakeywordTWO{where}\isanewline
\ \ \ \ {\isachardoublequoteopen}valid{\isacharunderscore}{\kern0pt}plan{\isacharunderscore}{\kern0pt}from\ M\ {\isasympi}s\ {\isasymlongleftrightarrow}\ {\isacharparenleft}{\kern0pt}wf{\isacharunderscore}{\kern0pt}plan\ {\isasympi}s\ {\isasymand}\ {\isacharparenleft}{\kern0pt}{\isasymexists}htps\ M{\isacharprime}{\kern0pt}{\isachardot}{\kern0pt}\ htps{\isacharunderscore}{\kern0pt}seq\ {\isasympi}s\ htps\ {\isasymand}\ valid{\isacharunderscore}{\kern0pt}state{\isacharunderscore}{\kern0pt}seq\ M\ htps\ {\isasympi}s\ M{\isacharprime}{\kern0pt}\ {\isasymand}\ M{\isacharprime}{\kern0pt}\ \isactrlsup c{\isasymTTurnstile}\isactrlsub {\isacharequal}{\kern0pt}\ {\isacharparenleft}{\kern0pt}goal\ P{\isacharparenright}{\kern0pt}{\isacharparenright}{\kern0pt}{\isacharparenright}{\kern0pt}{\isachardoublequoteclose}\isanewline
\isanewline
\ \ \isakeywordONE{definition}\isamarkupfalse%
\ valid{\isacharunderscore}{\kern0pt}plan\ {\isacharcolon}{\kern0pt}{\isacharcolon}{\kern0pt}\ {\isachardoublequoteopen}plan\ {\isasymRightarrow}\ bool{\isachardoublequoteclose}\isanewline
\ \ \ \ \isakeywordTWO{where}\ {\isachardoublequoteopen}valid{\isacharunderscore}{\kern0pt}plan\ {\isasymequiv}\ valid{\isacharunderscore}{\kern0pt}plan{\isacharunderscore}{\kern0pt}from\ I{\isachardoublequoteclose}%
\begin{isamarkuptext}%
Concise definition used in paper:%
\end{isamarkuptext}\isamarkuptrue%
\ \ \isakeywordONE{lemma}\isamarkupfalse%
\ {\isachardoublequoteopen}valid{\isacharunderscore}{\kern0pt}plan\ {\isasympi}s\ {\isasymequiv}\ wf{\isacharunderscore}{\kern0pt}plan\ {\isasympi}s\ {\isasymand}\ {\isacharparenleft}{\kern0pt}{\isasymexists}htps\ M{\isacharprime}{\kern0pt}{\isachardot}{\kern0pt}\ htps{\isacharunderscore}{\kern0pt}seq\ {\isasympi}s\ htps\ {\isasymand}\ valid{\isacharunderscore}{\kern0pt}state{\isacharunderscore}{\kern0pt}seq\ I\ htps\ {\isasympi}s\ M{\isacharprime}{\kern0pt}\ {\isasymand}\ M{\isacharprime}{\kern0pt}\ \isactrlsup c{\isasymTTurnstile}\isactrlsub {\isacharequal}{\kern0pt}\ {\isacharparenleft}{\kern0pt}goal\ P{\isacharparenright}{\kern0pt}{\isacharparenright}{\kern0pt}{\isachardoublequoteclose}\isanewline
%
\isadelimproof
\ \ \ \ %
\endisadelimproof
%
\isatagproof
\isakeywordONE{unfolding}\isamarkupfalse%
\ valid{\isacharunderscore}{\kern0pt}plan{\isacharunderscore}{\kern0pt}def\ valid{\isacharunderscore}{\kern0pt}plan{\isacharunderscore}{\kern0pt}from{\isacharunderscore}{\kern0pt}def\ \isakeywordONE{by}\isamarkupfalse%
\ auto%
\endisatagproof
{\isafoldproof}%
%
\isadelimproof
%
\endisadelimproof
%
\isadelimdocument
%
\endisadelimdocument
%
\isatagdocument
%
\isamarkupsubsection{Connecting the State-Transition Based Semantics to Happening Sequence Based Semantics%
}
\isamarkuptrue%
%
\endisatagdocument
{\isafolddocument}%
%
\isadelimdocument
%
\endisadelimdocument
%
\begin{isamarkuptext}%
Justification for Semantic refinement%
\end{isamarkuptext}\isamarkuptrue%
\ \ \isakeywordONE{lemma}\isamarkupfalse%
\ {\isachardoublequoteopen}apply{\isacharunderscore}{\kern0pt}happ\ {\isacharparenleft}{\kern0pt}t{\isacharcomma}{\kern0pt}A{\isacharparenright}{\kern0pt}\ M\ {\isacharequal}{\kern0pt}\ apply{\isacharunderscore}{\kern0pt}eff\ {\isacharparenleft}{\kern0pt}set\ A{\isacharparenright}{\kern0pt}\ M{\isachardoublequoteclose}\isanewline
%
\isadelimproof
\ \ \ \ %
\endisadelimproof
%
\isatagproof
\isakeywordONE{by}\isamarkupfalse%
\ auto%
\endisatagproof
{\isafoldproof}%
%
\isadelimproof
\isanewline
%
\endisadelimproof
\ \ \ \ \isanewline
\ \ \isakeywordONE{lemma}\isamarkupfalse%
\ valid{\isacharunderscore}{\kern0pt}state{\isacharunderscore}{\kern0pt}seq{\isacharunderscore}{\kern0pt}app{\isacharunderscore}{\kern0pt}iff{\isacharcolon}{\kern0pt}\isanewline
\ \ \ \ {\isachardoublequoteopen}valid{\isacharunderscore}{\kern0pt}state{\isacharunderscore}{\kern0pt}seq\ M\isactrlsub i\ {\isacharparenleft}{\kern0pt}ts\isactrlsub {\isadigit{1}}{\isacharat}{\kern0pt}ts\isactrlsub {\isadigit{2}}{\isacharparenright}{\kern0pt}\ {\isasympi}s\ M\isactrlsub k\ \isanewline
\ \ \ \ \ \ {\isasymlongleftrightarrow}\ {\isacharparenleft}{\kern0pt}{\isasymexists}M\isactrlsub j{\isachardot}{\kern0pt}\ valid{\isacharunderscore}{\kern0pt}state{\isacharunderscore}{\kern0pt}seq\ M\isactrlsub i\ ts\isactrlsub {\isadigit{1}}\ {\isasympi}s\ M\isactrlsub j\ {\isasymand}\ valid{\isacharunderscore}{\kern0pt}state{\isacharunderscore}{\kern0pt}seq\ M\isactrlsub j\ ts\isactrlsub {\isadigit{2}}\ {\isasympi}s\ M\isactrlsub k{\isacharparenright}{\kern0pt}{\isachardoublequoteclose}\isanewline
%
\isadelimproof
\ \ \ \ \isanewline
\ \ \ \ %
\endisadelimproof
%
\isatagproof
\isakeywordONE{by}\isamarkupfalse%
\ {\isacharparenleft}{\kern0pt}induction\ ts\isactrlsub {\isadigit{1}}\ arbitrary{\isacharcolon}{\kern0pt}\ ts\isactrlsub {\isadigit{2}}\ M\isactrlsub i\ M\isactrlsub k{\isacharparenright}{\kern0pt}\ {\isacharparenleft}{\kern0pt}auto\ simp{\isacharcolon}{\kern0pt}\ Let{\isacharunderscore}{\kern0pt}def{\isacharparenright}{\kern0pt}%
\endisatagproof
{\isafoldproof}%
%
\isadelimproof
\isanewline
%
\endisadelimproof
\isanewline
\ \ \isakeywordONE{lemma}\isamarkupfalse%
\ valid{\isacharunderscore}{\kern0pt}state{\isacharunderscore}{\kern0pt}seq{\isacharunderscore}{\kern0pt}state{\isacharunderscore}{\kern0pt}unique{\isacharcolon}{\kern0pt}\isanewline
\ \ \ \ \isakeywordTWO{assumes}\ {\isachardoublequoteopen}valid{\isacharunderscore}{\kern0pt}state{\isacharunderscore}{\kern0pt}seq\ M\isactrlsub {\isadigit{0}}\ ts\ {\isasympi}s\ M\isactrlsub i{\isachardoublequoteclose}\ \isakeywordTWO{and}\ {\isachardoublequoteopen}valid{\isacharunderscore}{\kern0pt}state{\isacharunderscore}{\kern0pt}seq\ M\isactrlsub {\isadigit{0}}\ ts\ {\isasympi}s\ M\isactrlsub i{\isacharprime}{\kern0pt}{\isachardoublequoteclose}\isanewline
\ \ \ \ \isakeywordTWO{shows}\ {\isachardoublequoteopen}M\isactrlsub i\ {\isacharequal}{\kern0pt}\ M\isactrlsub i{\isacharprime}{\kern0pt}{\isachardoublequoteclose}\isanewline
%
\isadelimproof
\ \ \ \ %
\endisadelimproof
%
\isatagproof
\isakeywordONE{using}\isamarkupfalse%
\ assms\ \isakeywordONE{by}\isamarkupfalse%
\ {\isacharparenleft}{\kern0pt}induction\ M\isactrlsub {\isadigit{0}}\ ts\ {\isasympi}s\ M\isactrlsub i\ rule{\isacharcolon}{\kern0pt}\ valid{\isacharunderscore}{\kern0pt}state{\isacharunderscore}{\kern0pt}seq{\isachardot}{\kern0pt}induct{\isacharparenright}{\kern0pt}\ {\isacharparenleft}{\kern0pt}auto\ simp{\isacharcolon}{\kern0pt}\ Let{\isacharunderscore}{\kern0pt}def{\isacharparenright}{\kern0pt}%
\endisatagproof
{\isafoldproof}%
%
\isadelimproof
%
\endisadelimproof
%
\begin{isamarkuptext}%
\isa{\isaconst{acts{\isacharunderscore}{\kern0pt}of{\isacharunderscore}{\kern0pt}plan{\isacharunderscore}{\kern0pt}at}} returns exactly the set of actions that would be executed 
  for a happening time point in the induced happening sequence%
\end{isamarkuptext}\isamarkuptrue%
\ \ \isakeywordONE{lemma}\isamarkupfalse%
\ ind{\isacharunderscore}{\kern0pt}happ{\isacharunderscore}{\kern0pt}seq{\isacharunderscore}{\kern0pt}htp{\isacharunderscore}{\kern0pt}acts{\isacharunderscore}{\kern0pt}of{\isacharunderscore}{\kern0pt}plan{\isacharunderscore}{\kern0pt}at{\isacharcolon}{\kern0pt}\isanewline
\ \ \ \ \isakeywordTWO{assumes}\ {\isachardoublequoteopen}ind{\isacharunderscore}{\kern0pt}happ{\isacharunderscore}{\kern0pt}seq\ {\isasympi}s\ hs{\isachardoublequoteclose}\isanewline
\ \ \ \ \ \ \ \ \isakeywordTWO{and}\ {\isachardoublequoteopen}is{\isacharunderscore}{\kern0pt}htp\ {\isasympi}s\ t{\isachardoublequoteclose}\isanewline
\ \ \ \ \ \ \ \ \isakeywordTWO{and}\ {\isachardoublequoteopen}{\isacharparenleft}{\kern0pt}t{\isacharcomma}{\kern0pt}A{\isacharparenright}{\kern0pt}\ {\isasymin}\ set\ hs{\isachardoublequoteclose}\isanewline
\ \ \ \ \isakeywordTWO{shows}\ {\isachardoublequoteopen}set\ A\ {\isacharequal}{\kern0pt}\ acts{\isacharunderscore}{\kern0pt}of{\isacharunderscore}{\kern0pt}plan{\isacharunderscore}{\kern0pt}at\ t\ {\isasympi}s{\isachardoublequoteclose}\isanewline
%
\isadelimproof
\ \ %
\endisadelimproof
%
\isatagproof
\isakeywordONE{proof}\isamarkupfalse%
\isanewline
\ \ \ \ \isakeywordTHREE{show}\isamarkupfalse%
\ {\isachardoublequoteopen}set\ A\ {\isasymsubseteq}\ acts{\isacharunderscore}{\kern0pt}of{\isacharunderscore}{\kern0pt}plan{\isacharunderscore}{\kern0pt}at\ t\ {\isasympi}s{\isachardoublequoteclose}\isanewline
\ \ \ \ \isakeywordONE{proof}\isamarkupfalse%
\isanewline
\ \ \ \ \ \ \isakeywordTHREE{fix}\isamarkupfalse%
\ a\isanewline
\ \ \ \ \ \ \isakeywordTHREE{assume}\isamarkupfalse%
\ {\isachardoublequoteopen}a\ {\isasymin}\ set\ A{\isachardoublequoteclose}\isanewline
\ \ \ \ \ \ \isakeywordONE{then}\isamarkupfalse%
\ \isakeywordONE{have}\isamarkupfalse%
\ {\isachardoublequoteopen}{\isasymexists}{\isacharparenleft}{\kern0pt}t\isactrlsub {\isasympi}{\isacharcomma}{\kern0pt}{\isasympi}{\isacharparenright}{\kern0pt}\ {\isasymin}\ set\ {\isasympi}s{\isachardot}{\kern0pt}\ inst{\isacharunderscore}{\kern0pt}of{\isacharunderscore}{\kern0pt}plan{\isacharunderscore}{\kern0pt}action\ {\isasympi}s\ {\isacharparenleft}{\kern0pt}t\isactrlsub {\isasympi}{\isacharcomma}{\kern0pt}{\isasympi}{\isacharparenright}{\kern0pt}\ {\isacharparenleft}{\kern0pt}t{\isacharcomma}{\kern0pt}a{\isacharparenright}{\kern0pt}{\isachardoublequoteclose}\isanewline
\ \ \ \ \ \ \ \ \isakeywordONE{using}\isamarkupfalse%
\ assms\ \isakeywordONE{by}\isamarkupfalse%
\ {\isacharparenleft}{\kern0pt}auto\ simp{\isacharcolon}{\kern0pt}\ ind{\isacharunderscore}{\kern0pt}happ{\isacharunderscore}{\kern0pt}seq{\isacharunderscore}{\kern0pt}def{\isacharparenright}{\kern0pt}\isanewline
\ \ \ \ \ \ \isakeywordONE{then}\isamarkupfalse%
\ \isakeywordTHREE{obtain}\isamarkupfalse%
\ t\isactrlsub {\isasympi}\ {\isasympi}\ \isakeywordTWO{where}\ {\isachardoublequoteopen}{\isacharparenleft}{\kern0pt}t\isactrlsub {\isasympi}{\isacharcomma}{\kern0pt}{\isasympi}{\isacharparenright}{\kern0pt}\ {\isasymin}\ set\ {\isasympi}s{\isachardoublequoteclose}\ \isakeywordTWO{and}\ {\isachardoublequoteopen}inst{\isacharunderscore}{\kern0pt}of{\isacharunderscore}{\kern0pt}plan{\isacharunderscore}{\kern0pt}action\ {\isasympi}s\ {\isacharparenleft}{\kern0pt}t\isactrlsub {\isasympi}{\isacharcomma}{\kern0pt}{\isasympi}{\isacharparenright}{\kern0pt}\ {\isacharparenleft}{\kern0pt}t{\isacharcomma}{\kern0pt}a{\isacharparenright}{\kern0pt}{\isachardoublequoteclose}\isanewline
\ \ \ \ \ \ \ \ \isakeywordONE{by}\isamarkupfalse%
\ auto\isanewline
\ \ \ \ \ \ \isakeywordONE{let}\isamarkupfalse%
\ {\isacharquery}{\kern0pt}case{\isadigit{1}}{\isacharequal}{\kern0pt}{\isachardoublequoteopen}res{\isacharunderscore}{\kern0pt}inst\ {\isasympi}\ At{\isacharunderscore}{\kern0pt}Start\ {\isacharequal}{\kern0pt}\ Some\ a\ {\isasymand}\ t\isactrlsub {\isasympi}\ {\isacharequal}{\kern0pt}\ t{\isachardoublequoteclose}\isanewline
\ \ \ \ \ \ \isakeywordONE{let}\isamarkupfalse%
\ {\isacharquery}{\kern0pt}case{\isadigit{2}}{\isacharequal}{\kern0pt}{\isachardoublequoteopen}res{\isacharunderscore}{\kern0pt}inst{\isacharunderscore}{\kern0pt}snap{\isacharunderscore}{\kern0pt}action\ {\isasympi}\ At{\isacharunderscore}{\kern0pt}Start\ {\isacharequal}{\kern0pt}\ Some\ a\ {\isasymand}\ t\isactrlsub {\isasympi}\ {\isacharequal}{\kern0pt}\ t{\isachardoublequoteclose}\isanewline
\ \ \ \ \ \ \isakeywordONE{let}\isamarkupfalse%
\ {\isacharquery}{\kern0pt}case{\isadigit{3}}{\isacharequal}{\kern0pt}{\isachardoublequoteopen}res{\isacharunderscore}{\kern0pt}inst{\isacharunderscore}{\kern0pt}snap{\isacharunderscore}{\kern0pt}action\ {\isasympi}\ At{\isacharunderscore}{\kern0pt}End\ {\isacharequal}{\kern0pt}\ Some\ a\ {\isasymand}\ t\isactrlsub {\isasympi}\ {\isacharplus}{\kern0pt}\ duration\ {\isasympi}\ {\isacharequal}{\kern0pt}\ t{\isachardoublequoteclose}\isanewline
\ \ \ \ \ \ \isakeywordONE{let}\isamarkupfalse%
\ {\isacharquery}{\kern0pt}case{\isadigit{4}}{\isacharequal}{\kern0pt}{\isachardoublequoteopen}res{\isacharunderscore}{\kern0pt}inst{\isacharunderscore}{\kern0pt}snap{\isacharunderscore}{\kern0pt}action\ {\isasympi}\ Over{\isacharunderscore}{\kern0pt}All\ {\isacharequal}{\kern0pt}\ Some\ a\ {\isasymand}\ \isanewline
\ \ \ \ \ \ \ \ \ \ \ \ {\isacharparenleft}{\kern0pt}{\isasymexists}t\isactrlsub i\ t\isactrlsub j{\isachardot}{\kern0pt}\ consec{\isacharunderscore}{\kern0pt}htps\ {\isasympi}s\ t\isactrlsub i\ t\isactrlsub j\ {\isasymand}\ t\isactrlsub {\isasympi}\ {\isasymle}\ t\isactrlsub i\ {\isasymand}\ t\isactrlsub j\ {\isasymle}\ t\isactrlsub {\isasympi}\ {\isacharplus}{\kern0pt}\ {\isacharparenleft}{\kern0pt}duration\ {\isasympi}{\isacharparenright}{\kern0pt}\ {\isasymand}\ t\isactrlsub i\ {\isacharless}{\kern0pt}\ t\ {\isasymand}\ t\ {\isacharless}{\kern0pt}\ t\isactrlsub j{\isacharparenright}{\kern0pt}{\isachardoublequoteclose}\isanewline
\ \ \ \ \ \ \isakeywordONE{have}\isamarkupfalse%
\ {\isachardoublequoteopen}{\isacharquery}{\kern0pt}case{\isadigit{1}}\ {\isasymor}\ {\isacharquery}{\kern0pt}case{\isadigit{2}}\ {\isasymor}\ {\isacharquery}{\kern0pt}case{\isadigit{3}}\ {\isasymor}\ {\isacharquery}{\kern0pt}case{\isadigit{4}}{\isachardoublequoteclose}\isanewline
\ \ \ \ \ \ \ \ \isakeywordONE{using}\isamarkupfalse%
\ {\isacartoucheopen}inst{\isacharunderscore}{\kern0pt}of{\isacharunderscore}{\kern0pt}plan{\isacharunderscore}{\kern0pt}action\ {\isasympi}s\ {\isacharparenleft}{\kern0pt}t\isactrlsub {\isasympi}{\isacharcomma}{\kern0pt}{\isasympi}{\isacharparenright}{\kern0pt}\ {\isacharparenleft}{\kern0pt}t{\isacharcomma}{\kern0pt}a{\isacharparenright}{\kern0pt}{\isacartoucheclose}\ \isakeywordONE{by}\isamarkupfalse%
\ simp\isanewline
\ \ \ \ \ \ \isakeywordONE{then}\isamarkupfalse%
\ \isakeywordTHREE{show}\isamarkupfalse%
\ {\isachardoublequoteopen}a\ {\isasymin}\ acts{\isacharunderscore}{\kern0pt}of{\isacharunderscore}{\kern0pt}plan{\isacharunderscore}{\kern0pt}at\ t\ {\isasympi}s{\isachardoublequoteclose}\isanewline
\ \ \ \ \ \ \ \ \isakeywordONE{proof}\isamarkupfalse%
\ {\isacharparenleft}{\kern0pt}elim\ disjE{\isacharparenright}{\kern0pt}\isanewline
\ \ \ \ \ \ \ \ \ \ \isakeywordTHREE{assume}\isamarkupfalse%
\ {\isacharquery}{\kern0pt}case{\isadigit{1}}\isanewline
\ \ \ \ \ \ \ \ \ \ \isakeywordONE{then}\isamarkupfalse%
\ \isakeywordTHREE{show}\isamarkupfalse%
\ {\isacharquery}{\kern0pt}thesis\isanewline
\ \ \ \ \ \ \ \ \ \ \ \ \isakeywordONE{unfolding}\isamarkupfalse%
\ acts{\isacharunderscore}{\kern0pt}of{\isacharunderscore}{\kern0pt}plan{\isacharunderscore}{\kern0pt}at{\isacharunderscore}{\kern0pt}def\ simple{\isacharunderscore}{\kern0pt}acts{\isacharunderscore}{\kern0pt}def\ is{\isacharunderscore}{\kern0pt}act{\isacharunderscore}{\kern0pt}simple{\isacharunderscore}{\kern0pt}def\isanewline
\ \ \ \ \ \ \ \ \ \ \ \ \isakeywordONE{using}\isamarkupfalse%
\ {\isacartoucheopen}{\isacharparenleft}{\kern0pt}t\isactrlsub {\isasympi}{\isacharcomma}{\kern0pt}{\isasympi}{\isacharparenright}{\kern0pt}\ {\isasymin}\ set\ {\isasympi}s{\isacartoucheclose}\ \isakeywordONE{by}\isamarkupfalse%
\ {\isacharparenleft}{\kern0pt}cases\ {\isasympi}{\isacharparenright}{\kern0pt}\ auto\isanewline
\ \ \ \ \ \ \ \ \isakeywordONE{next}\isamarkupfalse%
\isanewline
\ \ \ \ \ \ \ \ \ \ \isakeywordTHREE{assume}\isamarkupfalse%
\ {\isacharquery}{\kern0pt}case{\isadigit{2}}\isanewline
\ \ \ \ \ \ \ \ \ \ \isakeywordONE{then}\isamarkupfalse%
\ \isakeywordTHREE{show}\isamarkupfalse%
\ {\isacharquery}{\kern0pt}thesis\isanewline
\ \ \ \ \ \ \ \ \ \ \ \ \isakeywordONE{unfolding}\isamarkupfalse%
\ acts{\isacharunderscore}{\kern0pt}of{\isacharunderscore}{\kern0pt}plan{\isacharunderscore}{\kern0pt}at{\isacharunderscore}{\kern0pt}def\ durative{\isacharunderscore}{\kern0pt}acts{\isacharunderscore}{\kern0pt}def\ is{\isacharunderscore}{\kern0pt}act{\isacharunderscore}{\kern0pt}simple{\isacharunderscore}{\kern0pt}def\isanewline
\ \ \ \ \ \ \ \ \ \ \ \ \isakeywordONE{using}\isamarkupfalse%
\ {\isacartoucheopen}{\isacharparenleft}{\kern0pt}t\isactrlsub {\isasympi}{\isacharcomma}{\kern0pt}{\isasympi}{\isacharparenright}{\kern0pt}\ {\isasymin}\ set\ {\isasympi}s{\isacartoucheclose}\ \isakeywordONE{by}\isamarkupfalse%
\ {\isacharparenleft}{\kern0pt}cases\ {\isasympi}{\isacharparenright}{\kern0pt}\ auto\isanewline
\ \ \ \ \ \ \ \ \isakeywordONE{next}\isamarkupfalse%
\isanewline
\ \ \ \ \ \ \ \ \ \ \isakeywordTHREE{assume}\isamarkupfalse%
\ {\isacharquery}{\kern0pt}case{\isadigit{3}}\isanewline
\ \ \ \ \ \ \ \ \ \ \isakeywordONE{then}\isamarkupfalse%
\ \isakeywordTHREE{show}\isamarkupfalse%
\ {\isacharquery}{\kern0pt}thesis\isanewline
\ \ \ \ \ \ \ \ \ \ \ \ \isakeywordONE{unfolding}\isamarkupfalse%
\ acts{\isacharunderscore}{\kern0pt}of{\isacharunderscore}{\kern0pt}plan{\isacharunderscore}{\kern0pt}at{\isacharunderscore}{\kern0pt}def\ durative{\isacharunderscore}{\kern0pt}acts{\isacharunderscore}{\kern0pt}def\ is{\isacharunderscore}{\kern0pt}act{\isacharunderscore}{\kern0pt}simple{\isacharunderscore}{\kern0pt}def\isanewline
\ \ \ \ \ \ \ \ \ \ \ \ \isakeywordONE{using}\isamarkupfalse%
\ {\isacartoucheopen}{\isacharparenleft}{\kern0pt}t\isactrlsub {\isasympi}{\isacharcomma}{\kern0pt}{\isasympi}{\isacharparenright}{\kern0pt}\ {\isasymin}\ set\ {\isasympi}s{\isacartoucheclose}\ \isakeywordONE{by}\isamarkupfalse%
\ {\isacharparenleft}{\kern0pt}cases\ {\isasympi}{\isacharparenright}{\kern0pt}\ fastforce{\isacharplus}{\kern0pt}\isanewline
\ \ \ \ \ \ \ \ \isakeywordONE{next}\isamarkupfalse%
\isanewline
\ \ \ \ \ \ \ \ \ \ \isakeywordTHREE{assume}\isamarkupfalse%
\ {\isacharquery}{\kern0pt}case{\isadigit{4}}\isanewline
\ \ \ \ \ \ \ \ \ \ \isakeywordONE{then}\isamarkupfalse%
\ \isakeywordTHREE{show}\isamarkupfalse%
\ {\isacharquery}{\kern0pt}thesis\isanewline
\ \ \ \ \ \ \ \ \ \ \ \ \isakeywordONE{using}\isamarkupfalse%
\ assms\ \isakeywordONE{by}\isamarkupfalse%
\ {\isacharparenleft}{\kern0pt}auto\ simp{\isacharcolon}{\kern0pt}\ consec{\isacharunderscore}{\kern0pt}htps{\isacharunderscore}{\kern0pt}def{\isacharparenright}{\kern0pt}\isanewline
\ \ \ \ \ \ \ \ \isakeywordONE{qed}\isamarkupfalse%
\isanewline
\ \ \ \ \ \ \isakeywordONE{qed}\isamarkupfalse%
\isanewline
\ \ \isakeywordONE{next}\isamarkupfalse%
\isanewline
\ \ \ \ \isakeywordTHREE{show}\isamarkupfalse%
\ {\isachardoublequoteopen}set\ A\ {\isasymsupseteq}\ acts{\isacharunderscore}{\kern0pt}of{\isacharunderscore}{\kern0pt}plan{\isacharunderscore}{\kern0pt}at\ t\ {\isasympi}s{\isachardoublequoteclose}\isanewline
\ \ \ \ \isakeywordONE{proof}\isamarkupfalse%
\isanewline
\ \ \ \ \ \ \isakeywordTHREE{fix}\isamarkupfalse%
\ a\isanewline
\ \ \ \ \ \ \isakeywordONE{let}\isamarkupfalse%
\ {\isacharquery}{\kern0pt}case{\isadigit{1}}{\isacharequal}{\kern0pt}{\isachardoublequoteopen}a\ {\isasymin}\ {\isacharbraceleft}{\kern0pt}a\isactrlsub {\isasympi}{\isachardot}{\kern0pt}\ {\isasymexists}{\isasympi}{\isachardot}{\kern0pt}\ {\isacharparenleft}{\kern0pt}t{\isacharcomma}{\kern0pt}{\isasympi}{\isacharparenright}{\kern0pt}\ {\isasymin}\ simple{\isacharunderscore}{\kern0pt}acts\ {\isasympi}s\ {\isasymand}\ Some\ a\isactrlsub {\isasympi}\ {\isacharequal}{\kern0pt}\ res{\isacharunderscore}{\kern0pt}inst\ {\isasympi}\ At{\isacharunderscore}{\kern0pt}Start{\isacharbraceright}{\kern0pt}{\isachardoublequoteclose}\isanewline
\ \ \ \ \ \ \isakeywordONE{let}\isamarkupfalse%
\ {\isacharquery}{\kern0pt}case{\isadigit{2}}{\isacharequal}{\kern0pt}{\isachardoublequoteopen}a\ {\isasymin}\ {\isacharbraceleft}{\kern0pt}a\isactrlsub s\isactrlsub t\isactrlsub a\isactrlsub r\isactrlsub t{\isachardot}{\kern0pt}\ {\isasymexists}{\isasympi}{\isachardot}{\kern0pt}\ {\isacharparenleft}{\kern0pt}t{\isacharcomma}{\kern0pt}{\isasympi}{\isacharparenright}{\kern0pt}\ {\isasymin}\ durative{\isacharunderscore}{\kern0pt}acts\ {\isasympi}s\ {\isasymand}\ Some\ a\isactrlsub s\isactrlsub t\isactrlsub a\isactrlsub r\isactrlsub t\ {\isacharequal}{\kern0pt}\ res{\isacharunderscore}{\kern0pt}inst{\isacharunderscore}{\kern0pt}snap{\isacharunderscore}{\kern0pt}action\ {\isasympi}\ At{\isacharunderscore}{\kern0pt}Start{\isacharbraceright}{\kern0pt}{\isachardoublequoteclose}\isanewline
\ \ \ \ \ \ \isakeywordONE{let}\isamarkupfalse%
\ {\isacharquery}{\kern0pt}case{\isadigit{3}}{\isacharequal}{\kern0pt}{\isachardoublequoteopen}a\ {\isasymin}\ {\isacharbraceleft}{\kern0pt}a\isactrlsub e\isactrlsub n\isactrlsub d{\isachardot}{\kern0pt}\ {\isasymexists}t{\isacharprime}{\kern0pt}\ {\isasympi}{\isachardot}{\kern0pt}\ {\isacharparenleft}{\kern0pt}t{\isacharprime}{\kern0pt}{\isacharcomma}{\kern0pt}{\isasympi}{\isacharparenright}{\kern0pt}\ {\isasymin}\ durative{\isacharunderscore}{\kern0pt}acts\ {\isasympi}s\ {\isasymand}\ t\ {\isacharequal}{\kern0pt}\ t{\isacharprime}{\kern0pt}\ {\isacharplus}{\kern0pt}\ duration\ {\isasympi}\ {\isasymand}\ Some\ a\isactrlsub e\isactrlsub n\isactrlsub d\ {\isacharequal}{\kern0pt}\ res{\isacharunderscore}{\kern0pt}inst{\isacharunderscore}{\kern0pt}snap{\isacharunderscore}{\kern0pt}action\ {\isasympi}\ At{\isacharunderscore}{\kern0pt}End{\isacharbraceright}{\kern0pt}{\isachardoublequoteclose}\isanewline
\ \ \ \ \ \ \isakeywordTHREE{assume}\isamarkupfalse%
\ {\isachardoublequoteopen}a\ {\isasymin}\ acts{\isacharunderscore}{\kern0pt}of{\isacharunderscore}{\kern0pt}plan{\isacharunderscore}{\kern0pt}at\ t\ {\isasympi}s{\isachardoublequoteclose}\isanewline
\ \ \ \ \ \ \isakeywordONE{then}\isamarkupfalse%
\ \isakeywordONE{have}\isamarkupfalse%
\ {\isachardoublequoteopen}{\isacharquery}{\kern0pt}case{\isadigit{1}}\ {\isasymor}\ {\isacharquery}{\kern0pt}case{\isadigit{2}}\ {\isasymor}\ {\isacharquery}{\kern0pt}case{\isadigit{3}}{\isachardoublequoteclose}\isanewline
\ \ \ \ \ \ \ \ \isakeywordONE{unfolding}\isamarkupfalse%
\ acts{\isacharunderscore}{\kern0pt}of{\isacharunderscore}{\kern0pt}plan{\isacharunderscore}{\kern0pt}at{\isacharunderscore}{\kern0pt}def\ \isakeywordONE{by}\isamarkupfalse%
\ simp\isanewline
\ \ \ \ \ \ \isakeywordONE{then}\isamarkupfalse%
\ \isakeywordTHREE{show}\isamarkupfalse%
\ {\isachardoublequoteopen}a\ {\isasymin}\ set\ A{\isachardoublequoteclose}\isanewline
\ \ \ \ \ \ \ \ \isakeywordONE{proof}\isamarkupfalse%
\ {\isacharparenleft}{\kern0pt}elim\ disjE{\isacharparenright}{\kern0pt}\isanewline
\ \ \ \ \ \ \ \ \ \ \isakeywordTHREE{assume}\isamarkupfalse%
\ {\isacharquery}{\kern0pt}case{\isadigit{1}}\isanewline
\ \ \ \ \ \ \ \ \ \ \isakeywordONE{then}\isamarkupfalse%
\ \isakeywordTHREE{obtain}\isamarkupfalse%
\ {\isasympi}\ \isakeywordTWO{where}\ {\isachardoublequoteopen}{\isacharparenleft}{\kern0pt}t{\isacharcomma}{\kern0pt}{\isasympi}{\isacharparenright}{\kern0pt}\ {\isasymin}\ simple{\isacharunderscore}{\kern0pt}acts\ {\isasympi}s{\isachardoublequoteclose}\ \isakeywordTWO{and}\ {\isachardoublequoteopen}the\ {\isacharparenleft}{\kern0pt}res{\isacharunderscore}{\kern0pt}inst\ {\isasympi}\ At{\isacharunderscore}{\kern0pt}Start{\isacharparenright}{\kern0pt}\ {\isacharequal}{\kern0pt}\ a{\isachardoublequoteclose}\isanewline
\ \ \ \ \ \ \ \ \ \ \ \ \isakeywordONE{by}\isamarkupfalse%
\ auto\ {\isacharparenleft}{\kern0pt}metis\ option{\isachardot}{\kern0pt}sel{\isacharparenright}{\kern0pt}\isanewline
\ \ \ \ \ \ \ \ \ \ \isakeywordONE{then}\isamarkupfalse%
\ \isakeywordONE{have}\isamarkupfalse%
\ {\isachardoublequoteopen}{\isasymexists}as{\isachardot}{\kern0pt}\ {\isacharparenleft}{\kern0pt}t{\isacharcomma}{\kern0pt}\ as{\isacharparenright}{\kern0pt}\ {\isasymin}\ set\ hs\ {\isasymand}\ a\ {\isasymin}\ set\ as{\isachardoublequoteclose}\isanewline
\ \ \ \ \ \ \ \ \ \ \ \ \isakeywordONE{using}\isamarkupfalse%
\ assms\ \isakeywordONE{unfolding}\isamarkupfalse%
\ ind{\isacharunderscore}{\kern0pt}happ{\isacharunderscore}{\kern0pt}seq{\isacharunderscore}{\kern0pt}def\ \isakeywordONE{by}\isamarkupfalse%
\ auto\isanewline
\ \ \ \ \ \ \ \ \ \ \isakeywordONE{then}\isamarkupfalse%
\ \isakeywordTHREE{obtain}\isamarkupfalse%
\ A{\isacharprime}{\kern0pt}\ \isakeywordTWO{where}\ {\isachardoublequoteopen}{\isacharparenleft}{\kern0pt}t{\isacharcomma}{\kern0pt}\ A{\isacharprime}{\kern0pt}{\isacharparenright}{\kern0pt}\ {\isasymin}\ set\ hs{\isachardoublequoteclose}\ \isakeywordTWO{and}\ {\isachardoublequoteopen}a\ {\isasymin}\ set\ A{\isacharprime}{\kern0pt}{\isachardoublequoteclose}\isanewline
\ \ \ \ \ \ \ \ \ \ \ \ \isakeywordONE{by}\isamarkupfalse%
\ auto\isanewline
\ \ \ \ \ \ \ \ \ \ \isakeywordONE{have}\isamarkupfalse%
\ {\isachardoublequoteopen}A\ {\isacharequal}{\kern0pt}\ A{\isacharprime}{\kern0pt}{\isachardoublequoteclose}\isanewline
\ \ \ \ \ \ \ \ \ \ \ \ \isakeywordONE{using}\isamarkupfalse%
\ assms\ {\isacartoucheopen}{\isacharparenleft}{\kern0pt}t{\isacharcomma}{\kern0pt}\ A{\isacharprime}{\kern0pt}{\isacharparenright}{\kern0pt}\ {\isasymin}\ set\ hs{\isacartoucheclose}\ \isakeywordONE{unfolding}\isamarkupfalse%
\ ind{\isacharunderscore}{\kern0pt}happ{\isacharunderscore}{\kern0pt}seq{\isacharunderscore}{\kern0pt}def\isanewline
\ \ \ \ \ \ \ \ \ \ \ \ \isakeywordONE{by}\isamarkupfalse%
\ {\isacharparenleft}{\kern0pt}auto\ simp{\isacharcolon}{\kern0pt}\ strict{\isacharunderscore}{\kern0pt}sorted{\isacharunderscore}{\kern0pt}iff\ distinct{\isacharunderscore}{\kern0pt}map{\isacharunderscore}{\kern0pt}fstD{\isacharparenright}{\kern0pt}\isanewline
\ \ \ \ \ \ \ \ \ \ \isakeywordONE{then}\isamarkupfalse%
\ \isakeywordTHREE{show}\isamarkupfalse%
\ {\isacharquery}{\kern0pt}thesis\isanewline
\ \ \ \ \ \ \ \ \ \ \ \ \isakeywordONE{using}\isamarkupfalse%
\ {\isacartoucheopen}a\ {\isasymin}\ set\ A{\isacharprime}{\kern0pt}{\isacartoucheclose}\ \isakeywordONE{by}\isamarkupfalse%
\ auto\isanewline
\ \ \ \ \ \ \ \ \isakeywordONE{next}\isamarkupfalse%
\isanewline
\ \ \ \ \ \ \ \ \ \ \isakeywordTHREE{assume}\isamarkupfalse%
\ {\isacharquery}{\kern0pt}case{\isadigit{2}}\isanewline
\ \ \ \ \ \ \ \ \ \ \isakeywordONE{then}\isamarkupfalse%
\ \isakeywordTHREE{obtain}\isamarkupfalse%
\ {\isasympi}\ \isakeywordTWO{where}\ {\isachardoublequoteopen}{\isacharparenleft}{\kern0pt}t{\isacharcomma}{\kern0pt}{\isasympi}{\isacharparenright}{\kern0pt}\ {\isasymin}\ durative{\isacharunderscore}{\kern0pt}acts\ {\isasympi}s{\isachardoublequoteclose}\ \isanewline
\ \ \ \ \ \ \ \ \ \ \ \ \ \ \ \ \ \ \ \ \ \ \ \ \ \ \isakeywordTWO{and}\ {\isachardoublequoteopen}the\ {\isacharparenleft}{\kern0pt}res{\isacharunderscore}{\kern0pt}inst{\isacharunderscore}{\kern0pt}snap{\isacharunderscore}{\kern0pt}action\ {\isasympi}\ At{\isacharunderscore}{\kern0pt}Start{\isacharparenright}{\kern0pt}\ {\isacharequal}{\kern0pt}\ a{\isachardoublequoteclose}\isanewline
\ \ \ \ \ \ \ \ \ \ \ \ \isakeywordONE{by}\isamarkupfalse%
\ auto\ {\isacharparenleft}{\kern0pt}metis\ option{\isachardot}{\kern0pt}sel{\isacharparenright}{\kern0pt}\isanewline
\ \ \ \ \ \ \ \ \ \ \isakeywordONE{then}\isamarkupfalse%
\ \isakeywordONE{have}\isamarkupfalse%
\ {\isachardoublequoteopen}{\isasymexists}as{\isachardot}{\kern0pt}\ {\isacharparenleft}{\kern0pt}t{\isacharcomma}{\kern0pt}\ as{\isacharparenright}{\kern0pt}\ {\isasymin}\ set\ hs\ {\isasymand}\ a\ {\isasymin}\ set\ as{\isachardoublequoteclose}\isanewline
\ \ \ \ \ \ \ \ \ \ \ \ \isakeywordONE{using}\isamarkupfalse%
\ assms\ \isakeywordONE{unfolding}\isamarkupfalse%
\ ind{\isacharunderscore}{\kern0pt}happ{\isacharunderscore}{\kern0pt}seq{\isacharunderscore}{\kern0pt}def\ \isakeywordONE{by}\isamarkupfalse%
\ auto\isanewline
\ \ \ \ \ \ \ \ \ \ \isakeywordONE{then}\isamarkupfalse%
\ \isakeywordTHREE{obtain}\isamarkupfalse%
\ A{\isacharprime}{\kern0pt}\ \isakeywordTWO{where}\ {\isachardoublequoteopen}{\isacharparenleft}{\kern0pt}t{\isacharcomma}{\kern0pt}\ A{\isacharprime}{\kern0pt}{\isacharparenright}{\kern0pt}\ {\isasymin}\ set\ hs{\isachardoublequoteclose}\ \isakeywordTWO{and}\ {\isachardoublequoteopen}a\ {\isasymin}\ set\ A{\isacharprime}{\kern0pt}{\isachardoublequoteclose}\isanewline
\ \ \ \ \ \ \ \ \ \ \ \ \isakeywordONE{by}\isamarkupfalse%
\ auto\isanewline
\ \ \ \ \ \ \ \ \ \ \isakeywordONE{have}\isamarkupfalse%
\ {\isachardoublequoteopen}A\ {\isacharequal}{\kern0pt}\ A{\isacharprime}{\kern0pt}{\isachardoublequoteclose}\isanewline
\ \ \ \ \ \ \ \ \ \ \ \ \isakeywordONE{using}\isamarkupfalse%
\ assms\ {\isacartoucheopen}{\isacharparenleft}{\kern0pt}t{\isacharcomma}{\kern0pt}\ A{\isacharprime}{\kern0pt}{\isacharparenright}{\kern0pt}\ {\isasymin}\ set\ hs{\isacartoucheclose}\ \isakeywordONE{unfolding}\isamarkupfalse%
\ ind{\isacharunderscore}{\kern0pt}happ{\isacharunderscore}{\kern0pt}seq{\isacharunderscore}{\kern0pt}def\isanewline
\ \ \ \ \ \ \ \ \ \ \ \ \isakeywordONE{by}\isamarkupfalse%
\ {\isacharparenleft}{\kern0pt}auto\ simp{\isacharcolon}{\kern0pt}\ strict{\isacharunderscore}{\kern0pt}sorted{\isacharunderscore}{\kern0pt}iff\ distinct{\isacharunderscore}{\kern0pt}map{\isacharunderscore}{\kern0pt}fstD{\isacharparenright}{\kern0pt}\isanewline
\ \ \ \ \ \ \ \ \ \ \isakeywordONE{then}\isamarkupfalse%
\ \isakeywordTHREE{show}\isamarkupfalse%
\ {\isacharquery}{\kern0pt}thesis\isanewline
\ \ \ \ \ \ \ \ \ \ \ \ \isakeywordONE{using}\isamarkupfalse%
\ {\isacartoucheopen}a\ {\isasymin}\ set\ A{\isacharprime}{\kern0pt}{\isacartoucheclose}\ \isakeywordONE{by}\isamarkupfalse%
\ auto\isanewline
\ \ \ \ \ \ \ \ \isakeywordONE{next}\isamarkupfalse%
\isanewline
\ \ \ \ \ \ \ \ \ \ \isakeywordTHREE{assume}\isamarkupfalse%
\ {\isacharquery}{\kern0pt}case{\isadigit{3}}\isanewline
\ \ \ \ \ \ \ \ \ \ \isakeywordONE{then}\isamarkupfalse%
\ \isakeywordTHREE{obtain}\isamarkupfalse%
\ t{\isacharprime}{\kern0pt}\ {\isasympi}\ \isakeywordTWO{where}\ {\isachardoublequoteopen}{\isacharparenleft}{\kern0pt}t{\isacharprime}{\kern0pt}{\isacharcomma}{\kern0pt}{\isasympi}{\isacharparenright}{\kern0pt}\ {\isasymin}\ durative{\isacharunderscore}{\kern0pt}acts\ {\isasympi}s{\isachardoublequoteclose}\ \isakeywordTWO{and}\ {\isachardoublequoteopen}t\ {\isacharequal}{\kern0pt}\ t{\isacharprime}{\kern0pt}\ {\isacharplus}{\kern0pt}\ duration\ {\isasympi}{\isachardoublequoteclose}\ \isanewline
\ \ \ \ \ \ \ \ \ \ \ \ \ \ \ \ \ \ \ \ \ \ \ \ \ \ \ \ \ \isakeywordTWO{and}\ {\isachardoublequoteopen}the\ {\isacharparenleft}{\kern0pt}res{\isacharunderscore}{\kern0pt}inst{\isacharunderscore}{\kern0pt}snap{\isacharunderscore}{\kern0pt}action\ {\isasympi}\ At{\isacharunderscore}{\kern0pt}End{\isacharparenright}{\kern0pt}\ {\isacharequal}{\kern0pt}\ a{\isachardoublequoteclose}\isanewline
\ \ \ \ \ \ \ \ \ \ \ \ \isakeywordONE{by}\isamarkupfalse%
\ auto\ {\isacharparenleft}{\kern0pt}metis\ option{\isachardot}{\kern0pt}sel{\isacharparenright}{\kern0pt}\isanewline
\ \ \ \ \ \ \ \ \ \ \isakeywordONE{then}\isamarkupfalse%
\ \isakeywordONE{have}\isamarkupfalse%
\ {\isachardoublequoteopen}{\isasymexists}as{\isachardot}{\kern0pt}\ {\isacharparenleft}{\kern0pt}t{\isacharcomma}{\kern0pt}\ as{\isacharparenright}{\kern0pt}\ {\isasymin}\ set\ hs\ {\isasymand}\ a\ {\isasymin}\ set\ as{\isachardoublequoteclose}\isanewline
\ \ \ \ \ \ \ \ \ \ \ \ \isakeywordONE{using}\isamarkupfalse%
\ assms\ \isakeywordONE{unfolding}\isamarkupfalse%
\ ind{\isacharunderscore}{\kern0pt}happ{\isacharunderscore}{\kern0pt}seq{\isacharunderscore}{\kern0pt}def\ \isakeywordONE{by}\isamarkupfalse%
\ auto\isanewline
\ \ \ \ \ \ \ \ \ \ \isakeywordONE{then}\isamarkupfalse%
\ \isakeywordTHREE{obtain}\isamarkupfalse%
\ A{\isacharprime}{\kern0pt}\ \isakeywordTWO{where}\ {\isachardoublequoteopen}{\isacharparenleft}{\kern0pt}t{\isacharcomma}{\kern0pt}\ A{\isacharprime}{\kern0pt}{\isacharparenright}{\kern0pt}\ {\isasymin}\ set\ hs{\isachardoublequoteclose}\ \isakeywordTWO{and}\ {\isachardoublequoteopen}a\ {\isasymin}\ set\ A{\isacharprime}{\kern0pt}{\isachardoublequoteclose}\isanewline
\ \ \ \ \ \ \ \ \ \ \ \ \isakeywordONE{by}\isamarkupfalse%
\ auto\isanewline
\ \ \ \ \ \ \ \ \ \ \isakeywordONE{have}\isamarkupfalse%
\ {\isachardoublequoteopen}A\ {\isacharequal}{\kern0pt}\ A{\isacharprime}{\kern0pt}{\isachardoublequoteclose}\isanewline
\ \ \ \ \ \ \ \ \ \ \ \ \isakeywordONE{using}\isamarkupfalse%
\ assms\ {\isacartoucheopen}{\isacharparenleft}{\kern0pt}t{\isacharcomma}{\kern0pt}\ A{\isacharprime}{\kern0pt}{\isacharparenright}{\kern0pt}\ {\isasymin}\ set\ hs{\isacartoucheclose}\ \isakeywordONE{unfolding}\isamarkupfalse%
\ ind{\isacharunderscore}{\kern0pt}happ{\isacharunderscore}{\kern0pt}seq{\isacharunderscore}{\kern0pt}def\isanewline
\ \ \ \ \ \ \ \ \ \ \ \ \isakeywordONE{by}\isamarkupfalse%
\ {\isacharparenleft}{\kern0pt}auto\ simp{\isacharcolon}{\kern0pt}\ strict{\isacharunderscore}{\kern0pt}sorted{\isacharunderscore}{\kern0pt}iff\ distinct{\isacharunderscore}{\kern0pt}map{\isacharunderscore}{\kern0pt}fstD{\isacharparenright}{\kern0pt}\isanewline
\ \ \ \ \ \ \ \ \ \ \isakeywordONE{then}\isamarkupfalse%
\ \isakeywordTHREE{show}\isamarkupfalse%
\ {\isacharquery}{\kern0pt}thesis\isanewline
\ \ \ \ \ \ \ \ \ \ \ \ \isakeywordONE{using}\isamarkupfalse%
\ {\isacartoucheopen}a\ {\isasymin}\ set\ A{\isacharprime}{\kern0pt}{\isacartoucheclose}\ \isakeywordONE{by}\isamarkupfalse%
\ auto\isanewline
\ \ \ \ \ \ \isakeywordONE{qed}\isamarkupfalse%
\isanewline
\ \ \ \ \isakeywordONE{qed}\isamarkupfalse%
\isanewline
\ \ \isakeywordONE{qed}\isamarkupfalse%
%
\endisatagproof
{\isafoldproof}%
%
\isadelimproof
\isanewline
%
\endisadelimproof
\isakeywordTWO{end}\isamarkupfalse%
\isanewline
\isanewline
\isakeywordONE{context}\isamarkupfalse%
\ wf{\isacharunderscore}{\kern0pt}ast{\isacharunderscore}{\kern0pt}problem\isanewline
\isakeywordTWO{begin}%
\begin{isamarkuptext}%
The precondition of every invariant ground action in between two consecutive 
  happening time points is included in \isa{\isaconst{invs{\isacharunderscore}{\kern0pt}of{\isacharunderscore}{\kern0pt}plan{\isacharunderscore}{\kern0pt}at}}.%
\end{isamarkuptext}\isamarkuptrue%
\ \ \isakeywordONE{lemma}\isamarkupfalse%
\ ind{\isacharunderscore}{\kern0pt}happ{\isacharunderscore}{\kern0pt}seq{\isacharunderscore}{\kern0pt}invs{\isacharunderscore}{\kern0pt}of{\isacharunderscore}{\kern0pt}plan{\isacharunderscore}{\kern0pt}at{\isacharcolon}{\kern0pt}\isanewline
\ \ \ \ \isakeywordTWO{assumes}\ {\isachardoublequoteopen}wf{\isacharunderscore}{\kern0pt}plan\ {\isasympi}s{\isachardoublequoteclose}\isanewline
\ \ \ \ \ \ \ \ \isakeywordTWO{and}\ {\isachardoublequoteopen}consec{\isacharunderscore}{\kern0pt}htps\ {\isasympi}s\ t\isactrlsub i\ t\isactrlsub j{\isachardoublequoteclose}\isanewline
\ \ \ \ \ \ \ \ \isakeywordTWO{and}\ {\isachardoublequoteopen}ind{\isacharunderscore}{\kern0pt}happ{\isacharunderscore}{\kern0pt}seq\ {\isasympi}s\ hs{\isachardoublequoteclose}\isanewline
\ \ \ \ \ \ \ \ \isakeywordTWO{and}\ {\isachardoublequoteopen}{\isacharparenleft}{\kern0pt}t{\isacharcomma}{\kern0pt}A{\isacharparenright}{\kern0pt}\ {\isasymin}\ set\ {\isacharparenleft}{\kern0pt}dropWhile\ {\isacharparenleft}{\kern0pt}leq\ t\isactrlsub i{\isacharparenright}{\kern0pt}\ {\isacharparenleft}{\kern0pt}takeWhile\ {\isacharparenleft}{\kern0pt}le\ t\isactrlsub j{\isacharparenright}{\kern0pt}\ hs{\isacharparenright}{\kern0pt}{\isacharparenright}{\kern0pt}{\isachardoublequoteclose}\isanewline
\ \ \ \ \ \ \ \ \isakeywordTWO{and}\ {\isachardoublequoteopen}a\ {\isasymin}\ set\ A{\isachardoublequoteclose}\isanewline
\ \ \ \ \isakeywordTWO{shows}\ {\isachardoublequoteopen}{\isacharparenleft}{\kern0pt}precondition\ a{\isacharparenright}{\kern0pt}\ {\isasymin}\ invs{\isacharunderscore}{\kern0pt}of{\isacharunderscore}{\kern0pt}plan{\isacharunderscore}{\kern0pt}at\ t\isactrlsub j\ {\isasympi}s{\isachardoublequoteclose}\isanewline
%
\isadelimproof
\ \ %
\endisadelimproof
%
\isatagproof
\isakeywordONE{proof}\isamarkupfalse%
\ {\isacharminus}{\kern0pt}\isanewline
\ \ \ \ \isakeywordONE{let}\isamarkupfalse%
\ {\isacharquery}{\kern0pt}hs{\isacharprime}{\kern0pt}{\isacharequal}{\kern0pt}{\isachardoublequoteopen}dropWhile\ {\isacharparenleft}{\kern0pt}leq\ t\isactrlsub i{\isacharparenright}{\kern0pt}\ {\isacharparenleft}{\kern0pt}takeWhile\ {\isacharparenleft}{\kern0pt}le\ t\isactrlsub j{\isacharparenright}{\kern0pt}\ hs{\isacharparenright}{\kern0pt}{\isachardoublequoteclose}\isanewline
\ \ \ \ \isakeywordONE{have}\isamarkupfalse%
\ {\isachardoublequoteopen}strict{\isacharunderscore}{\kern0pt}sorted\ {\isacharparenleft}{\kern0pt}map\ fst\ hs{\isacharparenright}{\kern0pt}{\isachardoublequoteclose}\isanewline
\ \ \ \ \ \ \isakeywordONE{using}\isamarkupfalse%
\ {\isacartoucheopen}ind{\isacharunderscore}{\kern0pt}happ{\isacharunderscore}{\kern0pt}seq\ {\isasympi}s\ hs{\isacartoucheclose}\ \isakeywordONE{unfolding}\isamarkupfalse%
\ ind{\isacharunderscore}{\kern0pt}happ{\isacharunderscore}{\kern0pt}seq{\isacharunderscore}{\kern0pt}def\ \isakeywordONE{by}\isamarkupfalse%
\ auto\isanewline
\ \ \ \ \isakeywordONE{have}\isamarkupfalse%
\ {\isachardoublequoteopen}{\isacharparenleft}{\kern0pt}t{\isacharcomma}{\kern0pt}A{\isacharparenright}{\kern0pt}\ {\isasymin}\ set\ hs{\isachardoublequoteclose}\isanewline
\ \ \ \ \ \ \isakeywordONE{using}\isamarkupfalse%
\ assms\ \isakeywordONE{by}\isamarkupfalse%
\ {\isacharparenleft}{\kern0pt}meson\ set{\isacharunderscore}{\kern0pt}dropWhileD\ set{\isacharunderscore}{\kern0pt}takeWhileD{\isacharparenright}{\kern0pt}\isanewline
\ \ \ \ \isakeywordONE{then}\isamarkupfalse%
\ \isakeywordONE{have}\isamarkupfalse%
\ {\isachardoublequoteopen}{\isasymexists}{\isacharparenleft}{\kern0pt}t\isactrlsub {\isasympi}{\isacharcomma}{\kern0pt}{\isasympi}{\isacharparenright}{\kern0pt}\ {\isasymin}\ set\ {\isasympi}s{\isachardot}{\kern0pt}\ inst{\isacharunderscore}{\kern0pt}of{\isacharunderscore}{\kern0pt}plan{\isacharunderscore}{\kern0pt}action\ {\isasympi}s\ {\isacharparenleft}{\kern0pt}t\isactrlsub {\isasympi}{\isacharcomma}{\kern0pt}{\isasympi}{\isacharparenright}{\kern0pt}\ {\isacharparenleft}{\kern0pt}t{\isacharcomma}{\kern0pt}a{\isacharparenright}{\kern0pt}{\isachardoublequoteclose}\isanewline
\ \ \ \ \ \ \isakeywordONE{using}\isamarkupfalse%
\ assms\ \isakeywordONE{by}\isamarkupfalse%
\ {\isacharparenleft}{\kern0pt}auto\ simp{\isacharcolon}{\kern0pt}\ ind{\isacharunderscore}{\kern0pt}happ{\isacharunderscore}{\kern0pt}seq{\isacharunderscore}{\kern0pt}def{\isacharparenright}{\kern0pt}\isanewline
\ \ \ \ \isakeywordONE{then}\isamarkupfalse%
\ \isakeywordTHREE{obtain}\isamarkupfalse%
\ t\isactrlsub {\isasympi}\ {\isasympi}\ \isakeywordTWO{where}\ {\isachardoublequoteopen}{\isacharparenleft}{\kern0pt}t\isactrlsub {\isasympi}{\isacharcomma}{\kern0pt}{\isasympi}{\isacharparenright}{\kern0pt}\ {\isasymin}\ set\ {\isasympi}s{\isachardoublequoteclose}\ \isakeywordTWO{and}\ {\isachardoublequoteopen}inst{\isacharunderscore}{\kern0pt}of{\isacharunderscore}{\kern0pt}plan{\isacharunderscore}{\kern0pt}action\ {\isasympi}s\ {\isacharparenleft}{\kern0pt}t\isactrlsub {\isasympi}{\isacharcomma}{\kern0pt}{\isasympi}{\isacharparenright}{\kern0pt}\ {\isacharparenleft}{\kern0pt}t{\isacharcomma}{\kern0pt}a{\isacharparenright}{\kern0pt}{\isachardoublequoteclose}\isanewline
\ \ \ \ \ \ \isakeywordONE{by}\isamarkupfalse%
\ auto\isanewline
\ \ \ \ \isakeywordONE{then}\isamarkupfalse%
\ \isakeywordONE{have}\isamarkupfalse%
\ {\isachardoublequoteopen}is{\isacharunderscore}{\kern0pt}htp\ {\isasympi}s\ t\isactrlsub {\isasympi}{\isachardoublequoteclose}\isanewline
\ \ \ \ \ \ \isakeywordONE{unfolding}\isamarkupfalse%
\ is{\isacharunderscore}{\kern0pt}htp{\isacharunderscore}{\kern0pt}def\ \isakeywordONE{by}\isamarkupfalse%
\ auto\isanewline
\ \ \ \ \isakeywordONE{have}\isamarkupfalse%
\ {\isachardoublequoteopen}{\isasymnot}is{\isacharunderscore}{\kern0pt}htp\ {\isasympi}s\ t{\isachardoublequoteclose}\isanewline
\ \ \ \ \ \ \isakeywordONE{using}\isamarkupfalse%
\ assms\ strict{\isacharunderscore}{\kern0pt}sorted{\isacharunderscore}{\kern0pt}drop{\isacharunderscore}{\kern0pt}leq{\isacharunderscore}{\kern0pt}take{\isacharunderscore}{\kern0pt}le{\isacharprime}{\kern0pt}{\isacharbrackleft}{\kern0pt}OF\ {\isacartoucheopen}strict{\isacharunderscore}{\kern0pt}sorted\ {\isacharparenleft}{\kern0pt}map\ fst\ hs{\isacharparenright}{\kern0pt}{\isacartoucheclose}{\isacharbrackright}{\kern0pt}\ \isanewline
\ \ \ \ \ \ \isakeywordONE{unfolding}\isamarkupfalse%
\ consec{\isacharunderscore}{\kern0pt}htps{\isacharunderscore}{\kern0pt}def\ \isakeywordONE{by}\isamarkupfalse%
\ auto\isanewline
\ \ \ \ \isakeywordONE{let}\isamarkupfalse%
\ {\isacharquery}{\kern0pt}case{\isadigit{1}}{\isacharequal}{\kern0pt}{\isachardoublequoteopen}res{\isacharunderscore}{\kern0pt}inst\ {\isasympi}\ At{\isacharunderscore}{\kern0pt}Start{\isacharequal}{\kern0pt}\ Some\ a\ {\isasymand}\ t\isactrlsub {\isasympi}\ {\isacharequal}{\kern0pt}\ t{\isachardoublequoteclose}\isanewline
\ \ \ \ \isakeywordONE{let}\isamarkupfalse%
\ {\isacharquery}{\kern0pt}case{\isadigit{2}}{\isacharequal}{\kern0pt}{\isachardoublequoteopen}res{\isacharunderscore}{\kern0pt}inst{\isacharunderscore}{\kern0pt}snap{\isacharunderscore}{\kern0pt}action\ {\isasympi}\ At{\isacharunderscore}{\kern0pt}Start\ {\isacharequal}{\kern0pt}\ Some\ a\ {\isasymand}\ t\isactrlsub {\isasympi}\ {\isacharequal}{\kern0pt}\ t{\isachardoublequoteclose}\isanewline
\ \ \ \ \isakeywordONE{let}\isamarkupfalse%
\ {\isacharquery}{\kern0pt}case{\isadigit{3}}{\isacharequal}{\kern0pt}{\isachardoublequoteopen}res{\isacharunderscore}{\kern0pt}inst{\isacharunderscore}{\kern0pt}snap{\isacharunderscore}{\kern0pt}action\ {\isasympi}\ At{\isacharunderscore}{\kern0pt}End\ {\isacharequal}{\kern0pt}\ Some\ a\ {\isasymand}\ t\isactrlsub {\isasympi}\ {\isacharplus}{\kern0pt}\ duration\ {\isasympi}\ {\isacharequal}{\kern0pt}\ t{\isachardoublequoteclose}\isanewline
\ \ \ \ \isakeywordONE{let}\isamarkupfalse%
\ {\isacharquery}{\kern0pt}case{\isadigit{4}}{\isacharequal}{\kern0pt}{\isachardoublequoteopen}res{\isacharunderscore}{\kern0pt}inst{\isacharunderscore}{\kern0pt}snap{\isacharunderscore}{\kern0pt}action\ {\isasympi}\ Over{\isacharunderscore}{\kern0pt}All\ {\isacharequal}{\kern0pt}\ Some\ a\ {\isasymand}\ \isanewline
\ \ \ \ \ \ \ \ \ \ \ \ {\isacharparenleft}{\kern0pt}{\isasymexists}t\isactrlsub i\ t\isactrlsub j{\isachardot}{\kern0pt}\ consec{\isacharunderscore}{\kern0pt}htps\ {\isasympi}s\ t\isactrlsub i\ t\isactrlsub j\ {\isasymand}\ t\isactrlsub {\isasympi}\ {\isasymle}\ t\isactrlsub i\ {\isasymand}\ t\isactrlsub j\ {\isasymle}\ t\isactrlsub {\isasympi}\ {\isacharplus}{\kern0pt}\ {\isacharparenleft}{\kern0pt}duration\ {\isasympi}{\isacharparenright}{\kern0pt}\ {\isasymand}\ t\isactrlsub i\ {\isacharless}{\kern0pt}\ t\ {\isasymand}\ t\ {\isacharless}{\kern0pt}\ t\isactrlsub j{\isacharparenright}{\kern0pt}{\isachardoublequoteclose}\isanewline
\ \ \ \ \isakeywordONE{have}\isamarkupfalse%
\ {\isachardoublequoteopen}{\isacharquery}{\kern0pt}case{\isadigit{1}}\ {\isasymor}\ {\isacharquery}{\kern0pt}case{\isadigit{2}}\ {\isasymor}\ {\isacharquery}{\kern0pt}case{\isadigit{3}}\ {\isasymor}\ {\isacharquery}{\kern0pt}case{\isadigit{4}}{\isachardoublequoteclose}\isanewline
\ \ \ \ \ \ \isakeywordONE{using}\isamarkupfalse%
\ {\isacartoucheopen}inst{\isacharunderscore}{\kern0pt}of{\isacharunderscore}{\kern0pt}plan{\isacharunderscore}{\kern0pt}action\ {\isasympi}s\ {\isacharparenleft}{\kern0pt}t\isactrlsub {\isasympi}{\isacharcomma}{\kern0pt}{\isasympi}{\isacharparenright}{\kern0pt}\ {\isacharparenleft}{\kern0pt}t{\isacharcomma}{\kern0pt}a{\isacharparenright}{\kern0pt}{\isacartoucheclose}\ \isakeywordONE{by}\isamarkupfalse%
\ simp\isanewline
\ \ \ \ \isakeywordONE{then}\isamarkupfalse%
\ \isakeywordTHREE{show}\isamarkupfalse%
\ {\isacharquery}{\kern0pt}thesis\isanewline
\ \ \ \ \ \ \isakeywordONE{proof}\isamarkupfalse%
\ {\isacharparenleft}{\kern0pt}elim\ disjE{\isacharparenright}{\kern0pt}\isanewline
\ \ \ \ \ \ \ \ \isakeywordTHREE{assume}\isamarkupfalse%
\ {\isacharquery}{\kern0pt}case{\isadigit{1}}\isanewline
\ \ \ \ \ \ \ \ \isakeywordONE{then}\isamarkupfalse%
\ \isakeywordTHREE{show}\isamarkupfalse%
\ {\isacharquery}{\kern0pt}thesis\isanewline
\ \ \ \ \ \ \ \ \ \ \isakeywordONE{using}\isamarkupfalse%
\ assms\ {\isacartoucheopen}{\isasymnot}is{\isacharunderscore}{\kern0pt}htp\ {\isasympi}s\ t{\isacartoucheclose}\ {\isacartoucheopen}is{\isacharunderscore}{\kern0pt}htp\ {\isasympi}s\ t\isactrlsub {\isasympi}{\isacartoucheclose}\ \isakeywordONE{by}\isamarkupfalse%
\ auto\isanewline
\ \ \ \ \ \ \isakeywordONE{next}\isamarkupfalse%
\isanewline
\ \ \ \ \ \ \ \ \isakeywordTHREE{assume}\isamarkupfalse%
\ {\isacharquery}{\kern0pt}case{\isadigit{2}}\isanewline
\ \ \ \ \ \ \ \ \isakeywordONE{then}\isamarkupfalse%
\ \isakeywordTHREE{show}\isamarkupfalse%
\ {\isacharquery}{\kern0pt}thesis\isanewline
\ \ \ \ \ \ \ \ \ \ \isakeywordONE{using}\isamarkupfalse%
\ assms\ {\isacartoucheopen}{\isasymnot}is{\isacharunderscore}{\kern0pt}htp\ {\isasympi}s\ t{\isacartoucheclose}\ {\isacartoucheopen}is{\isacharunderscore}{\kern0pt}htp\ {\isasympi}s\ t\isactrlsub {\isasympi}{\isacartoucheclose}\ \isakeywordONE{by}\isamarkupfalse%
\ auto\isanewline
\ \ \ \ \ \ \isakeywordONE{next}\isamarkupfalse%
\isanewline
\ \ \ \ \ \ \ \ \isakeywordTHREE{assume}\isamarkupfalse%
\ {\isacharquery}{\kern0pt}case{\isadigit{3}}\isanewline
\ \ \ \ \ \ \ \ \isakeywordONE{then}\isamarkupfalse%
\ \isakeywordONE{have}\isamarkupfalse%
\ {\isachardoublequoteopen}{\isacharparenleft}{\kern0pt}t\isactrlsub {\isasympi}{\isacharcomma}{\kern0pt}{\isasympi}{\isacharparenright}{\kern0pt}\ {\isasymin}\ durative{\isacharunderscore}{\kern0pt}acts\ {\isasympi}s{\isachardoublequoteclose}\isanewline
\ \ \ \ \ \ \ \ \ \ \isakeywordONE{unfolding}\isamarkupfalse%
\ durative{\isacharunderscore}{\kern0pt}acts{\isacharunderscore}{\kern0pt}def\ is{\isacharunderscore}{\kern0pt}act{\isacharunderscore}{\kern0pt}simple{\isacharunderscore}{\kern0pt}def\isanewline
\ \ \ \ \ \ \ \ \ \ \isakeywordONE{using}\isamarkupfalse%
\ {\isacartoucheopen}{\isacharparenleft}{\kern0pt}t\isactrlsub {\isasympi}{\isacharcomma}{\kern0pt}{\isasympi}{\isacharparenright}{\kern0pt}\ {\isasymin}\ set\ {\isasympi}s{\isacartoucheclose}\ \isakeywordONE{by}\isamarkupfalse%
\ {\isacharparenleft}{\kern0pt}cases\ {\isasympi}{\isacharparenright}{\kern0pt}\ auto\isanewline
\ \ \ \ \ \ \ \ \isakeywordONE{then}\isamarkupfalse%
\ \isakeywordONE{have}\isamarkupfalse%
\ {\isachardoublequoteopen}is{\isacharunderscore}{\kern0pt}htp\ {\isasympi}s\ {\isacharparenleft}{\kern0pt}t\isactrlsub {\isasympi}\ {\isacharplus}{\kern0pt}\ duration\ {\isasympi}{\isacharparenright}{\kern0pt}{\isachardoublequoteclose}\isanewline
\ \ \ \ \ \ \ \ \ \ \isakeywordONE{unfolding}\isamarkupfalse%
\ is{\isacharunderscore}{\kern0pt}htp{\isacharunderscore}{\kern0pt}def\ \isakeywordONE{by}\isamarkupfalse%
\ auto\isanewline
\ \ \ \ \ \ \ \ \isakeywordONE{then}\isamarkupfalse%
\ \isakeywordTHREE{show}\isamarkupfalse%
\ {\isacharquery}{\kern0pt}thesis\isanewline
\ \ \ \ \ \ \ \ \ \ \isakeywordONE{using}\isamarkupfalse%
\ {\isacartoucheopen}{\isacharquery}{\kern0pt}case{\isadigit{3}}{\isacartoucheclose}\ {\isacartoucheopen}{\isasymnot}is{\isacharunderscore}{\kern0pt}htp\ {\isasympi}s\ t{\isacartoucheclose}\ \isakeywordONE{by}\isamarkupfalse%
\ auto\isanewline
\ \ \ \ \ \ \isakeywordONE{next}\isamarkupfalse%
\isanewline
\ \ \ \ \ \ \ \ \isakeywordTHREE{assume}\isamarkupfalse%
\ {\isacharquery}{\kern0pt}case{\isadigit{4}}\isanewline
\ \ \ \ \ \ \ \ \isakeywordTHREE{obtain}\isamarkupfalse%
\ t\isactrlsub i{\isacharprime}{\kern0pt}\ t\isactrlsub j{\isacharprime}{\kern0pt}\ \isakeywordTWO{where}\ {\isachardoublequoteopen}consec{\isacharunderscore}{\kern0pt}htps\ {\isasympi}s\ t\isactrlsub i{\isacharprime}{\kern0pt}\ t\isactrlsub j{\isacharprime}{\kern0pt}{\isachardoublequoteclose}\ \isanewline
\ \ \ \ \ \ \ \ \ \ \ \ \ \ \ \ \ \ \ \ \ \ \isakeywordTWO{and}\ {\isachardoublequoteopen}t\isactrlsub {\isasympi}\ {\isasymle}\ t\isactrlsub i{\isacharprime}{\kern0pt}\ {\isasymand}\ t\isactrlsub i{\isacharprime}{\kern0pt}\ {\isacharless}{\kern0pt}\ t\isactrlsub {\isasympi}\ {\isacharplus}{\kern0pt}\ duration\ {\isasympi}{\isachardoublequoteclose}\ \isanewline
\ \ \ \ \ \ \ \ \ \ \ \ \ \ \ \ \ \ \ \ \ \ \isakeywordTWO{and}\ {\isachardoublequoteopen}t\isactrlsub i{\isacharprime}{\kern0pt}\ {\isacharless}{\kern0pt}\ t\ {\isasymand}\ t\ {\isacharless}{\kern0pt}\ t\isactrlsub j{\isacharprime}{\kern0pt}{\isachardoublequoteclose}\isanewline
\ \ \ \ \ \ \ \ \isakeywordONE{using}\isamarkupfalse%
\ {\isacartoucheopen}{\isacharquery}{\kern0pt}case{\isadigit{4}}{\isacartoucheclose}\ \isakeywordONE{by}\isamarkupfalse%
\ auto\isanewline
\ \ \ \ \ \ \ \ \isakeywordONE{have}\isamarkupfalse%
\ {\isachardoublequoteopen}{\isacharparenleft}{\kern0pt}t\isactrlsub {\isasympi}{\isacharcomma}{\kern0pt}{\isasympi}{\isacharparenright}{\kern0pt}\ {\isasymin}\ durative{\isacharunderscore}{\kern0pt}acts\ {\isasympi}s{\isachardoublequoteclose}\isanewline
\ \ \ \ \ \ \ \ \ \ \isakeywordONE{unfolding}\isamarkupfalse%
\ durative{\isacharunderscore}{\kern0pt}acts{\isacharunderscore}{\kern0pt}def\ is{\isacharunderscore}{\kern0pt}act{\isacharunderscore}{\kern0pt}simple{\isacharunderscore}{\kern0pt}def\isanewline
\ \ \ \ \ \ \ \ \ \ \isakeywordONE{using}\isamarkupfalse%
\ {\isacartoucheopen}{\isacharquery}{\kern0pt}case{\isadigit{4}}{\isacartoucheclose}\ {\isacartoucheopen}{\isacharparenleft}{\kern0pt}t\isactrlsub {\isasympi}{\isacharcomma}{\kern0pt}{\isasympi}{\isacharparenright}{\kern0pt}\ {\isasymin}\ set\ {\isasympi}s{\isacartoucheclose}\ \isakeywordONE{by}\isamarkupfalse%
\ {\isacharparenleft}{\kern0pt}cases\ {\isasympi}{\isacharparenright}{\kern0pt}\ auto\isanewline
\ \ \ \ \ \ \ \ \isakeywordONE{then}\isamarkupfalse%
\ \isakeywordONE{have}\isamarkupfalse%
\ {\isachardoublequoteopen}is{\isacharunderscore}{\kern0pt}htp\ {\isasympi}s\ {\isacharparenleft}{\kern0pt}t\isactrlsub {\isasympi}\ {\isacharplus}{\kern0pt}\ duration\ {\isasympi}{\isacharparenright}{\kern0pt}{\isachardoublequoteclose}\isanewline
\ \ \ \ \ \ \ \ \ \ \isakeywordONE{unfolding}\isamarkupfalse%
\ is{\isacharunderscore}{\kern0pt}htp{\isacharunderscore}{\kern0pt}def\ \isakeywordONE{by}\isamarkupfalse%
\ auto\isanewline
\ \ \ \ \ \ \ \ \isakeywordONE{have}\isamarkupfalse%
\ {\isachardoublequoteopen}t\isactrlsub i\ {\isacharless}{\kern0pt}\ t\ {\isasymand}\ t\ {\isacharless}{\kern0pt}\ t\isactrlsub j{\isachardoublequoteclose}\isanewline
\ \ \ \ \ \ \ \ \ \ \isakeywordONE{using}\isamarkupfalse%
\ assms\ \isanewline
\ \ \ \ \ \ \ \ \ \ \ \ \ \ \ \ strict{\isacharunderscore}{\kern0pt}sorted{\isacharunderscore}{\kern0pt}drop{\isacharunderscore}{\kern0pt}leq{\isacharunderscore}{\kern0pt}take{\isacharunderscore}{\kern0pt}le{\isacharprime}{\kern0pt}{\isacharbrackleft}{\kern0pt}OF\ {\isacartoucheopen}strict{\isacharunderscore}{\kern0pt}sorted\ {\isacharparenleft}{\kern0pt}map\ fst\ hs{\isacharparenright}{\kern0pt}{\isacartoucheclose}{\isacharbrackright}{\kern0pt}\ \isanewline
\ \ \ \ \ \ \ \ \ \ \isakeywordONE{by}\isamarkupfalse%
\ simp\isanewline
\ \ \ \ \ \ \ \ \isakeywordONE{then}\isamarkupfalse%
\ \isakeywordONE{have}\isamarkupfalse%
\ {\isachardoublequoteopen}{\isacharparenleft}{\kern0pt}t\isactrlsub i\ {\isasymle}\ t\isactrlsub i{\isacharprime}{\kern0pt}\ {\isasymand}\ t\isactrlsub i{\isacharprime}{\kern0pt}\ {\isasymle}\ t\isactrlsub j{\isacharparenright}{\kern0pt}\ {\isasymor}\ {\isacharparenleft}{\kern0pt}t\isactrlsub i{\isacharprime}{\kern0pt}\ {\isasymle}\ t\isactrlsub i\ {\isasymand}\ t\isactrlsub i\ {\isasymle}\ t\isactrlsub j{\isacharprime}{\kern0pt}{\isacharparenright}{\kern0pt}{\isachardoublequoteclose}\isanewline
\ \ \ \ \ \ \ \ \ \ \isakeywordONE{using}\isamarkupfalse%
\ {\isacartoucheopen}t\isactrlsub i{\isacharprime}{\kern0pt}\ {\isacharless}{\kern0pt}\ t\ {\isasymand}\ t\ {\isacharless}{\kern0pt}\ t\isactrlsub j{\isacharprime}{\kern0pt}{\isacartoucheclose}\ \isakeywordONE{by}\isamarkupfalse%
\ auto\isanewline
\ \ \ \ \ \ \ \ \isakeywordONE{then}\isamarkupfalse%
\ \isakeywordONE{have}\isamarkupfalse%
\ {\isachardoublequoteopen}t\isactrlsub i{\isacharprime}{\kern0pt}\ {\isacharequal}{\kern0pt}\ t\isactrlsub i{\isachardoublequoteclose}\ \isakeywordTWO{and}\ {\isachardoublequoteopen}t\isactrlsub j{\isacharprime}{\kern0pt}\ {\isacharequal}{\kern0pt}\ t\isactrlsub j{\isachardoublequoteclose}\isanewline
\ \ \ \ \ \ \ \ \ \ \isakeywordONE{using}\isamarkupfalse%
\ {\isacartoucheopen}consec{\isacharunderscore}{\kern0pt}htps\ {\isasympi}s\ t\isactrlsub i\ t\isactrlsub j{\isacartoucheclose}\ {\isacartoucheopen}consec{\isacharunderscore}{\kern0pt}htps\ {\isasympi}s\ t\isactrlsub i{\isacharprime}{\kern0pt}\ t\isactrlsub j{\isacharprime}{\kern0pt}{\isacartoucheclose}\ {\isacartoucheopen}t\isactrlsub i\ {\isacharless}{\kern0pt}\ t\ {\isasymand}\ t\ {\isacharless}{\kern0pt}\ t\isactrlsub j{\isacartoucheclose}\ {\isacartoucheopen}t\isactrlsub i{\isacharprime}{\kern0pt}\ {\isacharless}{\kern0pt}\ t\ {\isasymand}\ t\ {\isacharless}{\kern0pt}\ t\isactrlsub j{\isacharprime}{\kern0pt}{\isacartoucheclose}\ \isanewline
\ \ \ \ \ \ \ \ \ \ \isakeywordONE{unfolding}\isamarkupfalse%
\ consec{\isacharunderscore}{\kern0pt}htps{\isacharunderscore}{\kern0pt}def\ \isakeywordONE{by}\isamarkupfalse%
\ auto\ {\isacharparenleft}{\kern0pt}metis\ dual{\isacharunderscore}{\kern0pt}order{\isachardot}{\kern0pt}strict{\isacharunderscore}{\kern0pt}trans\ eq{\isacharunderscore}{\kern0pt}iff\ leD{\isacharparenright}{\kern0pt}{\isacharplus}{\kern0pt}\isanewline
\ \ \ \ \ \ \ \ \isakeywordONE{then}\isamarkupfalse%
\ \isakeywordONE{have}\isamarkupfalse%
\ {\isachardoublequoteopen}t\isactrlsub {\isasympi}\ {\isasymle}\ t\isactrlsub i\ {\isasymand}\ t\isactrlsub i\ {\isacharless}{\kern0pt}\ t\isactrlsub {\isasympi}\ {\isacharplus}{\kern0pt}\ duration\ {\isasympi}{\isachardoublequoteclose}\isanewline
\ \ \ \ \ \ \ \ \ \ \isakeywordONE{using}\isamarkupfalse%
\ {\isacartoucheopen}t\isactrlsub {\isasympi}\ {\isasymle}\ t\isactrlsub i{\isacharprime}{\kern0pt}\ {\isasymand}\ t\isactrlsub i{\isacharprime}{\kern0pt}\ {\isacharless}{\kern0pt}\ t\isactrlsub {\isasympi}\ {\isacharplus}{\kern0pt}\ duration\ {\isasympi}{\isacartoucheclose}\ \isakeywordONE{by}\isamarkupfalse%
\ simp\isanewline
\ \ \ \ \ \ \ \ \isakeywordONE{then}\isamarkupfalse%
\ \isakeywordONE{have}\isamarkupfalse%
\ {\isachardoublequoteopen}t\isactrlsub {\isasympi}\ {\isacharless}{\kern0pt}\ t\isactrlsub j\ {\isasymand}\ t\isactrlsub j\ {\isasymle}\ t\isactrlsub {\isasympi}\ {\isacharplus}{\kern0pt}\ {\isacharparenleft}{\kern0pt}duration\ {\isasympi}{\isacharparenright}{\kern0pt}{\isachardoublequoteclose}\isanewline
\ \ \ \ \ \ \ \ \ \ \isakeywordONE{using}\isamarkupfalse%
\ assms\ {\isacartoucheopen}is{\isacharunderscore}{\kern0pt}htp\ {\isasympi}s\ {\isacharparenleft}{\kern0pt}t\isactrlsub {\isasympi}\ {\isacharplus}{\kern0pt}\ duration\ {\isasympi}{\isacharparenright}{\kern0pt}{\isacartoucheclose}\ \isakeywordONE{unfolding}\isamarkupfalse%
\ consec{\isacharunderscore}{\kern0pt}htps{\isacharunderscore}{\kern0pt}def\ \isakeywordONE{by}\isamarkupfalse%
\ auto\isanewline
\ \ \ \ \ \ \ \ \isanewline
\ \ \ \ \ \ \ \ \isakeywordONE{have}\isamarkupfalse%
\ {\isachardoublequoteopen}Some\ {\isacharparenleft}{\kern0pt}precondition\ a{\isacharparenright}{\kern0pt}\ {\isacharequal}{\kern0pt}\ res{\isacharunderscore}{\kern0pt}inst{\isacharunderscore}{\kern0pt}inv\ {\isasympi}{\isachardoublequoteclose}\isanewline
\ \ \ \ \ \ \ \ \ \ \isakeywordONE{using}\isamarkupfalse%
\ assms\ {\isacartoucheopen}{\isacharparenleft}{\kern0pt}t\isactrlsub {\isasympi}{\isacharcomma}{\kern0pt}{\isasympi}{\isacharparenright}{\kern0pt}\ {\isasymin}\ set\ {\isasympi}s{\isacartoucheclose}\ {\isacartoucheopen}{\isacharquery}{\kern0pt}case{\isadigit{4}}{\isacartoucheclose}\ res{\isacharunderscore}{\kern0pt}inst{\isacharunderscore}{\kern0pt}inv{\isacharunderscore}{\kern0pt}refine\ \isakeywordONE{unfolding}\isamarkupfalse%
\ wf{\isacharunderscore}{\kern0pt}plan{\isacharunderscore}{\kern0pt}def\ \isakeywordONE{by}\isamarkupfalse%
\ auto\isanewline
\isanewline
\ \ \ \ \ \ \ \ \isakeywordONE{then}\isamarkupfalse%
\ \isakeywordTHREE{show}\isamarkupfalse%
\ {\isacharquery}{\kern0pt}thesis\ \isanewline
\ \ \ \ \ \ \ \ \ \ \isakeywordONE{unfolding}\isamarkupfalse%
\ invs{\isacharunderscore}{\kern0pt}of{\isacharunderscore}{\kern0pt}plan{\isacharunderscore}{\kern0pt}at{\isacharunderscore}{\kern0pt}def\isanewline
\ \ \ \ \ \ \ \ \ \ \isakeywordONE{using}\isamarkupfalse%
\ {\isacartoucheopen}{\isacharparenleft}{\kern0pt}t\isactrlsub {\isasympi}{\isacharcomma}{\kern0pt}{\isasympi}{\isacharparenright}{\kern0pt}\ {\isasymin}\ durative{\isacharunderscore}{\kern0pt}acts\ {\isasympi}s{\isacartoucheclose}\ \isanewline
\ \ \ \ \ \ \ \ \ \ \ \ \ \ \ \ {\isacartoucheopen}t\isactrlsub {\isasympi}\ {\isacharless}{\kern0pt}\ t\isactrlsub j\ {\isasymand}\ t\isactrlsub j\ {\isasymle}\ t\isactrlsub {\isasympi}\ {\isacharplus}{\kern0pt}\ {\isacharparenleft}{\kern0pt}duration\ {\isasympi}{\isacharparenright}{\kern0pt}{\isacartoucheclose}\ \isanewline
\ \ \ \ \ \ \ \ \ \ \ \ \ \ \ \ {\isacartoucheopen}Some\ {\isacharparenleft}{\kern0pt}precondition\ a{\isacharparenright}{\kern0pt}\ {\isacharequal}{\kern0pt}\ res{\isacharunderscore}{\kern0pt}inst{\isacharunderscore}{\kern0pt}inv\ {\isasympi}{\isacartoucheclose}\isanewline
\ \ \ \ \ \ \ \ \ \ \isakeywordONE{by}\isamarkupfalse%
\ auto\isanewline
\ \ \ \ \ \isakeywordONE{qed}\isamarkupfalse%
\isanewline
\ \ \isakeywordONE{qed}\isamarkupfalse%
%
\endisatagproof
{\isafoldproof}%
%
\isadelimproof
%
\endisadelimproof
%
\begin{isamarkuptext}%
Every invariant from \isa{\isaconst{invs{\isacharunderscore}{\kern0pt}of{\isacharunderscore}{\kern0pt}plan{\isacharunderscore}{\kern0pt}at}} is placed as an invariant ground 
  action in between two consecutive happening time points in the induced happening 
  sequence.%
\end{isamarkuptext}\isamarkuptrue%
\ \ \isakeywordONE{lemma}\isamarkupfalse%
\ invs{\isacharunderscore}{\kern0pt}of{\isacharunderscore}{\kern0pt}plan{\isacharunderscore}{\kern0pt}at{\isacharunderscore}{\kern0pt}ind{\isacharunderscore}{\kern0pt}happ{\isacharunderscore}{\kern0pt}seq{\isacharcolon}{\kern0pt}\isanewline
\ \ \ \ \isakeywordTWO{assumes}\ {\isachardoublequoteopen}wf{\isacharunderscore}{\kern0pt}plan\ {\isasympi}s{\isachardoublequoteclose}\isanewline
\ \ \ \ \ \ \ \ \isakeywordTWO{and}\ {\isachardoublequoteopen}consec{\isacharunderscore}{\kern0pt}htps\ {\isasympi}s\ t\isactrlsub i\ t\isactrlsub j{\isachardoublequoteclose}\isanewline
\ \ \ \ \ \ \ \ \isakeywordTWO{and}\ {\isachardoublequoteopen}ind{\isacharunderscore}{\kern0pt}happ{\isacharunderscore}{\kern0pt}seq\ {\isasympi}s\ hs{\isachardoublequoteclose}\isanewline
\ \ \ \ \ \ \ \ \isakeywordTWO{and}\ {\isachardoublequoteopen}i\ {\isasymin}\ invs{\isacharunderscore}{\kern0pt}of{\isacharunderscore}{\kern0pt}plan{\isacharunderscore}{\kern0pt}at\ t\isactrlsub j\ {\isasympi}s{\isachardoublequoteclose}\isanewline
\ \ \ \ \isakeywordTWO{shows}\ {\isachardoublequoteopen}{\isasymexists}{\isacharparenleft}{\kern0pt}t{\isacharcomma}{\kern0pt}A{\isacharparenright}{\kern0pt}\ {\isasymin}\ set\ {\isacharparenleft}{\kern0pt}dropWhile\ {\isacharparenleft}{\kern0pt}leq\ t\isactrlsub i{\isacharparenright}{\kern0pt}\ {\isacharparenleft}{\kern0pt}takeWhile\ {\isacharparenleft}{\kern0pt}le\ t\isactrlsub j{\isacharparenright}{\kern0pt}\ hs{\isacharparenright}{\kern0pt}{\isacharparenright}{\kern0pt}{\isachardot}{\kern0pt}\ {\isasymexists}a\ {\isasymin}\ set\ A{\isachardot}{\kern0pt}\ i\ {\isacharequal}{\kern0pt}\ precondition\ a{\isachardoublequoteclose}\isanewline
%
\isadelimproof
\ \ %
\endisadelimproof
%
\isatagproof
\isakeywordONE{proof}\isamarkupfalse%
\ {\isacharminus}{\kern0pt}\isanewline
\ \ \ \ \isakeywordONE{let}\isamarkupfalse%
\ {\isacharquery}{\kern0pt}hs{\isacharprime}{\kern0pt}{\isacharequal}{\kern0pt}{\isachardoublequoteopen}dropWhile\ {\isacharparenleft}{\kern0pt}leq\ t\isactrlsub i{\isacharparenright}{\kern0pt}\ {\isacharparenleft}{\kern0pt}takeWhile\ {\isacharparenleft}{\kern0pt}le\ t\isactrlsub j{\isacharparenright}{\kern0pt}\ hs{\isacharparenright}{\kern0pt}{\isachardoublequoteclose}\isanewline
\ \ \ \ \isakeywordONE{have}\isamarkupfalse%
\ {\isachardoublequoteopen}strict{\isacharunderscore}{\kern0pt}sorted\ {\isacharparenleft}{\kern0pt}map\ fst\ hs{\isacharparenright}{\kern0pt}{\isachardoublequoteclose}\isanewline
\ \ \ \ \ \ \isakeywordONE{using}\isamarkupfalse%
\ {\isacartoucheopen}ind{\isacharunderscore}{\kern0pt}happ{\isacharunderscore}{\kern0pt}seq\ {\isasympi}s\ hs{\isacartoucheclose}\ \isakeywordONE{unfolding}\isamarkupfalse%
\ ind{\isacharunderscore}{\kern0pt}happ{\isacharunderscore}{\kern0pt}seq{\isacharunderscore}{\kern0pt}def\ \isakeywordONE{by}\isamarkupfalse%
\ auto\isanewline
\ \ \ \ \isakeywordTHREE{obtain}\isamarkupfalse%
\ t\isactrlsub {\isasympi}\ {\isasympi}\ \isakeywordTWO{where}\ {\isachardoublequoteopen}{\isacharparenleft}{\kern0pt}t\isactrlsub {\isasympi}{\isacharcomma}{\kern0pt}{\isasympi}{\isacharparenright}{\kern0pt}\ {\isasymin}\ durative{\isacharunderscore}{\kern0pt}acts\ {\isasympi}s{\isachardoublequoteclose}\ \isakeywordTWO{and}\ {\isachardoublequoteopen}t\isactrlsub {\isasympi}\ {\isacharless}{\kern0pt}\ t\isactrlsub j\ {\isasymand}\ t\isactrlsub j\ {\isasymle}\ t\isactrlsub {\isasympi}\ {\isacharplus}{\kern0pt}\ {\isacharparenleft}{\kern0pt}duration\ {\isasympi}{\isacharparenright}{\kern0pt}{\isachardoublequoteclose}\isanewline
\ \ \ \ \ \ \ \ \ \ \ \ \ \ \ \ \ \ \ \ \ \ \ \isakeywordTWO{and}\ {\isachardoublequoteopen}Some\ i\ {\isacharequal}{\kern0pt}\ res{\isacharunderscore}{\kern0pt}inst{\isacharunderscore}{\kern0pt}inv\ {\isasympi}{\isachardoublequoteclose}\isanewline
\ \ \ \ \ \ \isakeywordONE{using}\isamarkupfalse%
\ assms\ \isakeywordONE{unfolding}\isamarkupfalse%
\ invs{\isacharunderscore}{\kern0pt}of{\isacharunderscore}{\kern0pt}plan{\isacharunderscore}{\kern0pt}at{\isacharunderscore}{\kern0pt}def\ \isakeywordONE{by}\isamarkupfalse%
\ auto\isanewline
\ \ \ \ \isakeywordONE{then}\isamarkupfalse%
\ \isakeywordONE{have}\isamarkupfalse%
\ {\isachardoublequoteopen}is{\isacharunderscore}{\kern0pt}htp\ {\isasympi}s\ t\isactrlsub {\isasympi}{\isachardoublequoteclose}\isanewline
\ \ \ \ \ \ \isakeywordONE{unfolding}\isamarkupfalse%
\ is{\isacharunderscore}{\kern0pt}htp{\isacharunderscore}{\kern0pt}def\ \isakeywordONE{by}\isamarkupfalse%
\ {\isacharparenleft}{\kern0pt}auto\ simp{\isacharcolon}{\kern0pt}\ durative{\isacharunderscore}{\kern0pt}acts{\isacharunderscore}{\kern0pt}def{\isacharparenright}{\kern0pt}\isanewline
\ \ \ \ \isakeywordONE{then}\isamarkupfalse%
\ \isakeywordONE{have}\isamarkupfalse%
\ {\isachardoublequoteopen}t\isactrlsub {\isasympi}\ {\isasymle}\ t\isactrlsub i\ {\isasymand}\ t\isactrlsub j\ {\isasymle}\ t\isactrlsub {\isasympi}\ {\isacharplus}{\kern0pt}\ {\isacharparenleft}{\kern0pt}duration\ {\isasympi}{\isacharparenright}{\kern0pt}{\isachardoublequoteclose}\isanewline
\ \ \ \ \ \ \isakeywordONE{using}\isamarkupfalse%
\ assms\ {\isacartoucheopen}t\isactrlsub {\isasympi}\ {\isacharless}{\kern0pt}\ t\isactrlsub j\ {\isasymand}\ t\isactrlsub j\ {\isasymle}\ t\isactrlsub {\isasympi}\ {\isacharplus}{\kern0pt}\ {\isacharparenleft}{\kern0pt}duration\ {\isasympi}{\isacharparenright}{\kern0pt}{\isacartoucheclose}\ \isakeywordONE{unfolding}\isamarkupfalse%
\ consec{\isacharunderscore}{\kern0pt}htps{\isacharunderscore}{\kern0pt}def\ \isakeywordONE{by}\isamarkupfalse%
\ auto\isanewline
\ \ \ \ \isakeywordONE{then}\isamarkupfalse%
\ \isakeywordONE{have}\isamarkupfalse%
\ {\isachardoublequoteopen}{\isasymexists}{\isacharparenleft}{\kern0pt}t{\isacharprime}{\kern0pt}{\isacharcomma}{\kern0pt}as{\isacharparenright}{\kern0pt}\ {\isasymin}\ set\ hs{\isachardot}{\kern0pt}\ t\isactrlsub i\ {\isacharless}{\kern0pt}\ t{\isacharprime}{\kern0pt}\ {\isasymand}\ t{\isacharprime}{\kern0pt}\ {\isacharless}{\kern0pt}\ t\isactrlsub j\ {\isasymand}\ the\ {\isacharparenleft}{\kern0pt}res{\isacharunderscore}{\kern0pt}inst{\isacharunderscore}{\kern0pt}snap{\isacharunderscore}{\kern0pt}action\ {\isasympi}\ Over{\isacharunderscore}{\kern0pt}All{\isacharparenright}{\kern0pt}\ {\isasymin}\ set\ as{\isachardoublequoteclose}\isanewline
\ \ \ \ \ \ \isakeywordONE{using}\isamarkupfalse%
\ assms\ {\isacartoucheopen}{\isacharparenleft}{\kern0pt}t\isactrlsub {\isasympi}{\isacharcomma}{\kern0pt}{\isasympi}{\isacharparenright}{\kern0pt}\ {\isasymin}\ durative{\isacharunderscore}{\kern0pt}acts\ {\isasympi}s{\isacartoucheclose}\ \isakeywordONE{unfolding}\isamarkupfalse%
\ ind{\isacharunderscore}{\kern0pt}happ{\isacharunderscore}{\kern0pt}seq{\isacharunderscore}{\kern0pt}def\ \isakeywordONE{by}\isamarkupfalse%
\ auto\isanewline
\ \ \ \ \isanewline
\ \ \ \ \isakeywordONE{then}\isamarkupfalse%
\ \isakeywordTHREE{obtain}\isamarkupfalse%
\ t\ A\ \isakeywordTWO{where}\ {\isachardoublequoteopen}{\isacharparenleft}{\kern0pt}t{\isacharcomma}{\kern0pt}A{\isacharparenright}{\kern0pt}\ {\isasymin}\ set\ hs{\isachardoublequoteclose}\ \isakeywordTWO{and}\ {\isachardoublequoteopen}t\isactrlsub i\ {\isacharless}{\kern0pt}\ t\ {\isasymand}\ t\ {\isacharless}{\kern0pt}\ t\isactrlsub j{\isachardoublequoteclose}\ \isanewline
\ \ \ \ \ \ \ \ \ \ \ \ \ \ \ \ \ \ \ \ \ \ \ \isakeywordTWO{and}\ {\isachardoublequoteopen}the\ {\isacharparenleft}{\kern0pt}res{\isacharunderscore}{\kern0pt}inst{\isacharunderscore}{\kern0pt}snap{\isacharunderscore}{\kern0pt}action\ {\isasympi}\ Over{\isacharunderscore}{\kern0pt}All{\isacharparenright}{\kern0pt}\ {\isasymin}\ set\ A{\isachardoublequoteclose}\isanewline
\ \ \ \ \ \ \isakeywordONE{by}\isamarkupfalse%
\ auto\isanewline
\ \ \ \ \isakeywordTHREE{obtain}\isamarkupfalse%
\ a\ \isakeywordTWO{where}\ {\isachardoublequoteopen}Some\ a\ {\isacharequal}{\kern0pt}\ res{\isacharunderscore}{\kern0pt}inst{\isacharunderscore}{\kern0pt}snap{\isacharunderscore}{\kern0pt}action\ {\isasympi}\ Over{\isacharunderscore}{\kern0pt}All{\isachardoublequoteclose}\isanewline
\ \ \ \ \ \ \isakeywordONE{using}\isamarkupfalse%
\ {\isacartoucheopen}Some\ i\ {\isacharequal}{\kern0pt}\ res{\isacharunderscore}{\kern0pt}inst{\isacharunderscore}{\kern0pt}inv\ {\isasympi}{\isacartoucheclose}\isanewline
\ \ \ \ \ \ \isakeywordONE{by}\isamarkupfalse%
\ {\isacharparenleft}{\kern0pt}metis\ option{\isachardot}{\kern0pt}discI\ res{\isacharunderscore}{\kern0pt}inst{\isacharunderscore}{\kern0pt}inv{\isachardot}{\kern0pt}simps{\isacharparenleft}{\kern0pt}{\isadigit{1}}{\isacharparenright}{\kern0pt}\ res{\isacharunderscore}{\kern0pt}inst{\isacharunderscore}{\kern0pt}snap{\isacharunderscore}{\kern0pt}action{\isachardot}{\kern0pt}elims{\isacharparenright}{\kern0pt}\isanewline
\ \ \ \ \isakeywordONE{then}\isamarkupfalse%
\ \isakeywordONE{have}\isamarkupfalse%
\ {\isachardoublequoteopen}a\ {\isasymin}\ set\ A{\isachardoublequoteclose}\isanewline
\ \ \ \ \ \ \isakeywordONE{using}\isamarkupfalse%
\ {\isacartoucheopen}the\ {\isacharparenleft}{\kern0pt}res{\isacharunderscore}{\kern0pt}inst{\isacharunderscore}{\kern0pt}snap{\isacharunderscore}{\kern0pt}action\ {\isasympi}\ Over{\isacharunderscore}{\kern0pt}All{\isacharparenright}{\kern0pt}\ {\isasymin}\ set\ A{\isacartoucheclose}\ \isakeywordONE{by}\isamarkupfalse%
\ {\isacharparenleft}{\kern0pt}metis\ option{\isachardot}{\kern0pt}sel{\isacharparenright}{\kern0pt}\isanewline
\isanewline
\ \ \ \ \isakeywordONE{have}\isamarkupfalse%
\ {\isachardoublequoteopen}{\isacharparenleft}{\kern0pt}t{\isacharcomma}{\kern0pt}A{\isacharparenright}{\kern0pt}\ {\isasymin}\ set\ {\isacharquery}{\kern0pt}hs{\isacharprime}{\kern0pt}{\isachardoublequoteclose}\isanewline
\ \ \ \ \ \ \isakeywordONE{using}\isamarkupfalse%
\ {\isacartoucheopen}t\isactrlsub i\ {\isacharless}{\kern0pt}\ t\ {\isasymand}\ t\ {\isacharless}{\kern0pt}\ t\isactrlsub j{\isacartoucheclose}\ \ {\isacartoucheopen}{\isacharparenleft}{\kern0pt}t{\isacharcomma}{\kern0pt}\ A{\isacharparenright}{\kern0pt}\ {\isasymin}\ set\ hs{\isacartoucheclose}\ \isanewline
\ \ \ \ \ \ \ \ \ \ \ \ strict{\isacharunderscore}{\kern0pt}sorted{\isacharunderscore}{\kern0pt}drop{\isacharunderscore}{\kern0pt}leq{\isacharunderscore}{\kern0pt}take{\isacharunderscore}{\kern0pt}le{\isacharprime}{\kern0pt}{\isacharbrackleft}{\kern0pt}OF\ {\isacartoucheopen}strict{\isacharunderscore}{\kern0pt}sorted\ {\isacharparenleft}{\kern0pt}map\ fst\ hs{\isacharparenright}{\kern0pt}{\isacartoucheclose}{\isacharbrackright}{\kern0pt}\isanewline
\ \ \ \ \ \ \isakeywordONE{by}\isamarkupfalse%
\ auto\isanewline
\ \ \ \ \isakeywordONE{have}\isamarkupfalse%
\ {\isachardoublequoteopen}wf{\isacharunderscore}{\kern0pt}plan{\isacharunderscore}{\kern0pt}action\ {\isasympi}{\isachardoublequoteclose}\isanewline
\ \ \ \ \ \ \isakeywordONE{using}\isamarkupfalse%
\ assms\ {\isacartoucheopen}{\isacharparenleft}{\kern0pt}t\isactrlsub {\isasympi}{\isacharcomma}{\kern0pt}{\isasympi}{\isacharparenright}{\kern0pt}\ {\isasymin}\ durative{\isacharunderscore}{\kern0pt}acts\ {\isasympi}s{\isacartoucheclose}\ \isakeywordONE{unfolding}\isamarkupfalse%
\ wf{\isacharunderscore}{\kern0pt}plan{\isacharunderscore}{\kern0pt}def\ \isanewline
\ \ \ \ \ \ \isakeywordONE{by}\isamarkupfalse%
\ {\isacharparenleft}{\kern0pt}auto\ simp{\isacharcolon}{\kern0pt}\ durative{\isacharunderscore}{\kern0pt}acts{\isacharunderscore}{\kern0pt}def{\isacharparenright}{\kern0pt}\isanewline
\ \ \ \ \isakeywordONE{then}\isamarkupfalse%
\ \isakeywordONE{have}\isamarkupfalse%
\ {\isachardoublequoteopen}Some\ {\isacharparenleft}{\kern0pt}precondition\ a{\isacharparenright}{\kern0pt}\ {\isacharequal}{\kern0pt}\ res{\isacharunderscore}{\kern0pt}inst{\isacharunderscore}{\kern0pt}inv\ {\isasympi}{\isachardoublequoteclose}\isanewline
\ \ \ \ \ \ \isakeywordONE{using}\isamarkupfalse%
\ assms\ {\isacartoucheopen}Some\ a\ {\isacharequal}{\kern0pt}\ res{\isacharunderscore}{\kern0pt}inst{\isacharunderscore}{\kern0pt}snap{\isacharunderscore}{\kern0pt}action\ {\isasympi}\ Over{\isacharunderscore}{\kern0pt}All{\isacartoucheclose}\ res{\isacharunderscore}{\kern0pt}inst{\isacharunderscore}{\kern0pt}inv{\isacharunderscore}{\kern0pt}refine\ \isakeywordONE{by}\isamarkupfalse%
\ auto\isanewline
\ \ \ \ \isakeywordONE{then}\isamarkupfalse%
\ \isakeywordONE{have}\isamarkupfalse%
\ {\isachardoublequoteopen}i\ {\isacharequal}{\kern0pt}\ precondition\ a{\isachardoublequoteclose}\isanewline
\ \ \ \ \ \ \isakeywordONE{using}\isamarkupfalse%
\ {\isacartoucheopen}Some\ i\ {\isacharequal}{\kern0pt}\ res{\isacharunderscore}{\kern0pt}inst{\isacharunderscore}{\kern0pt}inv\ {\isasympi}{\isacartoucheclose}\isanewline
\ \ \ \ \ \ \isakeywordONE{by}\isamarkupfalse%
\ {\isacharparenleft}{\kern0pt}metis\ option{\isachardot}{\kern0pt}sel{\isacharparenright}{\kern0pt}\isanewline
\ \ \ \ \isakeywordTHREE{show}\isamarkupfalse%
\ {\isacharquery}{\kern0pt}thesis\isanewline
\ \ \ \ \ \ \isakeywordONE{using}\isamarkupfalse%
\ {\isacartoucheopen}{\isacharparenleft}{\kern0pt}t{\isacharcomma}{\kern0pt}A{\isacharparenright}{\kern0pt}\ {\isasymin}\ set\ {\isacharquery}{\kern0pt}hs{\isacharprime}{\kern0pt}{\isacartoucheclose}\ {\isacartoucheopen}a\ {\isasymin}\ set\ A{\isacartoucheclose}\ {\isacartoucheopen}i\ {\isacharequal}{\kern0pt}\ precondition\ a{\isacartoucheclose}\ \isakeywordONE{by}\isamarkupfalse%
\ auto\isanewline
\ \ \isakeywordONE{qed}\isamarkupfalse%
%
\endisatagproof
{\isafoldproof}%
%
\isadelimproof
%
\endisadelimproof
%
\begin{isamarkuptext}%
In an induced happening sequence all ground actions between two consecutive happening 
  time points have empty effect. Those ground actions are the instantiated durative action 
  for the \texttt{(over all ...)} time specifier.%
\end{isamarkuptext}\isamarkuptrue%
\ \ \isakeywordONE{lemma}\isamarkupfalse%
\ ind{\isacharunderscore}{\kern0pt}happ{\isacharunderscore}{\kern0pt}seq{\isacharunderscore}{\kern0pt}leq\isactrlsub i{\isacharunderscore}{\kern0pt}le\isactrlsub j{\isacharunderscore}{\kern0pt}empty{\isacharunderscore}{\kern0pt}eff{\isacharcolon}{\kern0pt}\isanewline
\ \ \ \ \isakeywordTWO{assumes}\ {\isachardoublequoteopen}consec{\isacharunderscore}{\kern0pt}htps\ {\isasympi}s\ t\isactrlsub i\ t\isactrlsub j{\isachardoublequoteclose}\isanewline
\ \ \ \ \ \ \ \ \isakeywordTWO{and}\ {\isachardoublequoteopen}wf{\isacharunderscore}{\kern0pt}plan\ {\isasympi}s{\isachardoublequoteclose}\isanewline
\ \ \ \ \ \ \ \ \isakeywordTWO{and}\ {\isachardoublequoteopen}ind{\isacharunderscore}{\kern0pt}happ{\isacharunderscore}{\kern0pt}seq\ {\isasympi}s\ hs{\isachardoublequoteclose}\isanewline
\ \ \ \ \ \ \ \ \isakeywordTWO{and}\ {\isachardoublequoteopen}{\isacharparenleft}{\kern0pt}t\isactrlsub a{\isacharcomma}{\kern0pt}A{\isacharparenright}{\kern0pt}\ {\isasymin}\ set\ {\isacharparenleft}{\kern0pt}dropWhile\ {\isacharparenleft}{\kern0pt}leq\ t\isactrlsub i{\isacharparenright}{\kern0pt}\ {\isacharparenleft}{\kern0pt}takeWhile\ {\isacharparenleft}{\kern0pt}le\ t\isactrlsub j{\isacharparenright}{\kern0pt}\ hs{\isacharparenright}{\kern0pt}{\isacharparenright}{\kern0pt}{\isachardoublequoteclose}\isanewline
\ \ \ \ \ \ \ \ \isakeywordTWO{and}\ {\isachardoublequoteopen}a\ {\isasymin}\ set\ A{\isachardoublequoteclose}\isanewline
\ \ \ \ \isakeywordTWO{shows}\ {\isachardoublequoteopen}effect\ a\ {\isacharequal}{\kern0pt}\ Effect\ {\isacharbrackleft}{\kern0pt}{\isacharbrackright}{\kern0pt}\ {\isacharbrackleft}{\kern0pt}{\isacharbrackright}{\kern0pt}{\isachardoublequoteclose}\isanewline
%
\isadelimproof
\ \ %
\endisadelimproof
%
\isatagproof
\isakeywordONE{proof}\isamarkupfalse%
\ {\isacharminus}{\kern0pt}\isanewline
\ \ \ \ \isakeywordONE{let}\isamarkupfalse%
\ {\isacharquery}{\kern0pt}hs{\isacharprime}{\kern0pt}{\isacharequal}{\kern0pt}{\isachardoublequoteopen}dropWhile\ {\isacharparenleft}{\kern0pt}leq\ t\isactrlsub i{\isacharparenright}{\kern0pt}\ {\isacharparenleft}{\kern0pt}takeWhile\ {\isacharparenleft}{\kern0pt}le\ t\isactrlsub j{\isacharparenright}{\kern0pt}\ hs{\isacharparenright}{\kern0pt}{\isachardoublequoteclose}\isanewline
\ \ \ \ \isakeywordONE{have}\isamarkupfalse%
\ {\isachardoublequoteopen}strict{\isacharunderscore}{\kern0pt}sorted\ {\isacharparenleft}{\kern0pt}map\ fst\ hs{\isacharparenright}{\kern0pt}{\isachardoublequoteclose}\isanewline
\ \ \ \ \ \ \isakeywordONE{using}\isamarkupfalse%
\ {\isacartoucheopen}ind{\isacharunderscore}{\kern0pt}happ{\isacharunderscore}{\kern0pt}seq\ {\isasympi}s\ hs{\isacartoucheclose}\ \isakeywordONE{unfolding}\isamarkupfalse%
\ ind{\isacharunderscore}{\kern0pt}happ{\isacharunderscore}{\kern0pt}seq{\isacharunderscore}{\kern0pt}def\ \isakeywordONE{by}\isamarkupfalse%
\ auto\isanewline
\ \ \ \ \isakeywordONE{then}\isamarkupfalse%
\ \isakeywordONE{have}\isamarkupfalse%
\ {\isachardoublequoteopen}{\isacharparenleft}{\kern0pt}t\isactrlsub a{\isacharcomma}{\kern0pt}A{\isacharparenright}{\kern0pt}\ {\isasymin}\ set\ hs{\isachardoublequoteclose}\isanewline
\ \ \ \ \ \ \isakeywordONE{using}\isamarkupfalse%
\ {\isacartoucheopen}{\isacharparenleft}{\kern0pt}t\isactrlsub a{\isacharcomma}{\kern0pt}A{\isacharparenright}{\kern0pt}\ {\isasymin}\ set\ {\isacharquery}{\kern0pt}hs{\isacharprime}{\kern0pt}{\isacartoucheclose}\ \isakeywordONE{by}\isamarkupfalse%
\ {\isacharparenleft}{\kern0pt}meson\ set{\isacharunderscore}{\kern0pt}dropWhileD\ set{\isacharunderscore}{\kern0pt}takeWhileD{\isacharparenright}{\kern0pt}\isanewline
\ \ \ \ \isakeywordONE{have}\isamarkupfalse%
\ {\isachardoublequoteopen}t\isactrlsub i\ {\isacharless}{\kern0pt}\ t\isactrlsub a\ {\isasymand}\ t\isactrlsub a\ {\isacharless}{\kern0pt}\ t\isactrlsub j{\isachardoublequoteclose}\isanewline
\ \ \ \ \ \ \isakeywordONE{using}\isamarkupfalse%
\ {\isacartoucheopen}{\isacharparenleft}{\kern0pt}t\isactrlsub a{\isacharcomma}{\kern0pt}A{\isacharparenright}{\kern0pt}\ {\isasymin}\ set\ {\isacharquery}{\kern0pt}hs{\isacharprime}{\kern0pt}{\isacartoucheclose}\isanewline
\ \ \ \ \ \ \ \ \ \ \ \ strict{\isacharunderscore}{\kern0pt}sorted{\isacharunderscore}{\kern0pt}drop{\isacharunderscore}{\kern0pt}leq{\isacharunderscore}{\kern0pt}take{\isacharunderscore}{\kern0pt}le{\isacharbrackleft}{\kern0pt}OF\ {\isacartoucheopen}strict{\isacharunderscore}{\kern0pt}sorted\ {\isacharparenleft}{\kern0pt}map\ fst\ hs{\isacharparenright}{\kern0pt}{\isacartoucheclose}\ {\isacartoucheopen}{\isacharparenleft}{\kern0pt}t\isactrlsub a{\isacharcomma}{\kern0pt}A{\isacharparenright}{\kern0pt}\ {\isasymin}\ set\ hs{\isacartoucheclose}{\isacharbrackright}{\kern0pt}\isanewline
\ \ \ \ \ \ \isakeywordONE{by}\isamarkupfalse%
\ auto\isanewline
\ \ \ \ \isakeywordONE{have}\isamarkupfalse%
\ {\isachardoublequoteopen}{\isasymnot}is{\isacharunderscore}{\kern0pt}htp\ {\isasympi}s\ t\isactrlsub a{\isachardoublequoteclose}\isanewline
\ \ \ \ \ \ \isakeywordONE{using}\isamarkupfalse%
\ {\isacartoucheopen}t\isactrlsub i\ {\isacharless}{\kern0pt}\ t\isactrlsub a\ {\isasymand}\ t\isactrlsub a\ {\isacharless}{\kern0pt}\ t\isactrlsub j{\isacartoucheclose}\ {\isacartoucheopen}consec{\isacharunderscore}{\kern0pt}htps\ {\isasympi}s\ t\isactrlsub i\ t\isactrlsub j{\isacartoucheclose}\ \isakeywordONE{unfolding}\isamarkupfalse%
\ consec{\isacharunderscore}{\kern0pt}htps{\isacharunderscore}{\kern0pt}def\ \isakeywordONE{by}\isamarkupfalse%
\ auto\isanewline
\ \ \ \ \isakeywordONE{have}\isamarkupfalse%
\ {\isachardoublequoteopen}A\ {\isasymnoteq}\ {\isacharbrackleft}{\kern0pt}{\isacharbrackright}{\kern0pt}{\isachardoublequoteclose}\isanewline
\ \ \ \ \ \ \isakeywordONE{using}\isamarkupfalse%
\ {\isacartoucheopen}ind{\isacharunderscore}{\kern0pt}happ{\isacharunderscore}{\kern0pt}seq\ {\isasympi}s\ hs{\isacartoucheclose}\ {\isacartoucheopen}{\isacharparenleft}{\kern0pt}t\isactrlsub a{\isacharcomma}{\kern0pt}A{\isacharparenright}{\kern0pt}\ {\isasymin}\ set\ hs{\isacartoucheclose}\ \isakeywordONE{unfolding}\isamarkupfalse%
\ ind{\isacharunderscore}{\kern0pt}happ{\isacharunderscore}{\kern0pt}seq{\isacharunderscore}{\kern0pt}def\ \isakeywordONE{by}\isamarkupfalse%
\ auto\isanewline
\ \ \ \ \isakeywordONE{have}\isamarkupfalse%
\ {\isachardoublequoteopen}{\isasymexists}{\isacharparenleft}{\kern0pt}t\isactrlsub {\isasympi}{\isacharcomma}{\kern0pt}{\isasympi}{\isacharparenright}{\kern0pt}\ {\isasymin}\ set\ {\isasympi}s{\isachardot}{\kern0pt}\ inst{\isacharunderscore}{\kern0pt}of{\isacharunderscore}{\kern0pt}plan{\isacharunderscore}{\kern0pt}action\ {\isasympi}s\ {\isacharparenleft}{\kern0pt}t\isactrlsub {\isasympi}{\isacharcomma}{\kern0pt}{\isasympi}{\isacharparenright}{\kern0pt}\ {\isacharparenleft}{\kern0pt}t\isactrlsub a{\isacharcomma}{\kern0pt}a{\isacharparenright}{\kern0pt}{\isachardoublequoteclose}\isanewline
\ \ \ \ \ \ \isakeywordONE{using}\isamarkupfalse%
\ {\isacartoucheopen}ind{\isacharunderscore}{\kern0pt}happ{\isacharunderscore}{\kern0pt}seq\ {\isasympi}s\ hs{\isacartoucheclose}\ {\isacartoucheopen}{\isacharparenleft}{\kern0pt}t\isactrlsub a{\isacharcomma}{\kern0pt}A{\isacharparenright}{\kern0pt}\ {\isasymin}\ set\ hs{\isacartoucheclose}\ {\isacartoucheopen}a\ {\isasymin}\ set\ A{\isacartoucheclose}\ \isakeywordONE{unfolding}\isamarkupfalse%
\ ind{\isacharunderscore}{\kern0pt}happ{\isacharunderscore}{\kern0pt}seq{\isacharunderscore}{\kern0pt}def\ \isakeywordONE{by}\isamarkupfalse%
\ auto\isanewline
\ \ \ \ \isakeywordONE{then}\isamarkupfalse%
\ \isakeywordTHREE{obtain}\isamarkupfalse%
\ t\isactrlsub {\isasympi}\ {\isasympi}\ \isakeywordTWO{where}\ {\isachardoublequoteopen}{\isacharparenleft}{\kern0pt}t\isactrlsub {\isasympi}{\isacharcomma}{\kern0pt}{\isasympi}{\isacharparenright}{\kern0pt}\ {\isasymin}\ set\ {\isasympi}s{\isachardoublequoteclose}\ \isakeywordTWO{and}\ {\isachardoublequoteopen}inst{\isacharunderscore}{\kern0pt}of{\isacharunderscore}{\kern0pt}plan{\isacharunderscore}{\kern0pt}action\ {\isasympi}s\ {\isacharparenleft}{\kern0pt}t\isactrlsub {\isasympi}{\isacharcomma}{\kern0pt}{\isasympi}{\isacharparenright}{\kern0pt}\ {\isacharparenleft}{\kern0pt}t\isactrlsub a{\isacharcomma}{\kern0pt}a{\isacharparenright}{\kern0pt}{\isachardoublequoteclose}\isanewline
\ \ \ \ \ \ \isakeywordONE{by}\isamarkupfalse%
\ auto\isanewline
\ \ \ \ \isakeywordONE{then}\isamarkupfalse%
\ \isakeywordONE{have}\isamarkupfalse%
\ {\isachardoublequoteopen}is{\isacharunderscore}{\kern0pt}htp\ {\isasympi}s\ t\isactrlsub {\isasympi}{\isachardoublequoteclose}\isanewline
\ \ \ \ \ \ \isakeywordONE{unfolding}\isamarkupfalse%
\ is{\isacharunderscore}{\kern0pt}htp{\isacharunderscore}{\kern0pt}def\ \isakeywordONE{by}\isamarkupfalse%
\ auto\isanewline
\ \ \ \ \isakeywordONE{let}\isamarkupfalse%
\ {\isacharquery}{\kern0pt}case{\isadigit{1}}{\isacharequal}{\kern0pt}{\isachardoublequoteopen}res{\isacharunderscore}{\kern0pt}inst\ {\isasympi}\ At{\isacharunderscore}{\kern0pt}Start\ {\isacharequal}{\kern0pt}\ Some\ a\ {\isasymand}\ t\isactrlsub {\isasympi}\ {\isacharequal}{\kern0pt}\ t\isactrlsub a{\isachardoublequoteclose}\isanewline
\ \ \ \ \isakeywordONE{let}\isamarkupfalse%
\ {\isacharquery}{\kern0pt}case{\isadigit{2}}{\isacharequal}{\kern0pt}{\isachardoublequoteopen}res{\isacharunderscore}{\kern0pt}inst{\isacharunderscore}{\kern0pt}snap{\isacharunderscore}{\kern0pt}action\ {\isasympi}\ At{\isacharunderscore}{\kern0pt}Start\ {\isacharequal}{\kern0pt}\ Some\ a\ {\isasymand}\ t\isactrlsub {\isasympi}\ {\isacharequal}{\kern0pt}\ t\isactrlsub a{\isachardoublequoteclose}\isanewline
\ \ \ \ \isakeywordONE{let}\isamarkupfalse%
\ {\isacharquery}{\kern0pt}case{\isadigit{3}}{\isacharequal}{\kern0pt}{\isachardoublequoteopen}res{\isacharunderscore}{\kern0pt}inst{\isacharunderscore}{\kern0pt}snap{\isacharunderscore}{\kern0pt}action\ {\isasympi}\ At{\isacharunderscore}{\kern0pt}End\ {\isacharequal}{\kern0pt}\ Some\ a\ {\isasymand}\ t\isactrlsub {\isasympi}\ {\isacharplus}{\kern0pt}\ {\isacharparenleft}{\kern0pt}duration\ {\isasympi}{\isacharparenright}{\kern0pt}\ {\isacharequal}{\kern0pt}\ t\isactrlsub a{\isachardoublequoteclose}\isanewline
\ \ \ \ \isakeywordONE{let}\isamarkupfalse%
\ {\isacharquery}{\kern0pt}case{\isadigit{4}}{\isacharequal}{\kern0pt}{\isachardoublequoteopen}res{\isacharunderscore}{\kern0pt}inst{\isacharunderscore}{\kern0pt}snap{\isacharunderscore}{\kern0pt}action\ {\isasympi}\ Over{\isacharunderscore}{\kern0pt}All\ {\isacharequal}{\kern0pt}\ Some\ a\ {\isasymand}\ \isanewline
\ \ \ \ \ \ \ \ \ \ {\isacharparenleft}{\kern0pt}{\isasymexists}t\isactrlsub i\ t\isactrlsub j{\isachardot}{\kern0pt}\ consec{\isacharunderscore}{\kern0pt}htps\ {\isasympi}s\ t\isactrlsub i\ t\isactrlsub j\ {\isasymand}\ t\isactrlsub {\isasympi}\ {\isasymle}\ t\isactrlsub i\ {\isasymand}\ t\isactrlsub j\ {\isasymle}\ t\isactrlsub {\isasympi}\ {\isacharplus}{\kern0pt}\ {\isacharparenleft}{\kern0pt}duration\ {\isasympi}{\isacharparenright}{\kern0pt}\ {\isasymand}\ t\isactrlsub i\ {\isacharless}{\kern0pt}\ t\isactrlsub a\ {\isasymand}\ t\isactrlsub a\ {\isacharless}{\kern0pt}\ t\isactrlsub j{\isacharparenright}{\kern0pt}{\isachardoublequoteclose}\isanewline
\ \ \ \ \isakeywordONE{have}\isamarkupfalse%
\ {\isachardoublequoteopen}{\isacharquery}{\kern0pt}case{\isadigit{1}}\ {\isasymor}\ {\isacharquery}{\kern0pt}case{\isadigit{2}}\ {\isasymor}\ {\isacharquery}{\kern0pt}case{\isadigit{3}}\ {\isasymor}\ {\isacharquery}{\kern0pt}case{\isadigit{4}}{\isachardoublequoteclose}\isanewline
\ \ \ \ \ \ \isakeywordONE{using}\isamarkupfalse%
\ {\isacartoucheopen}inst{\isacharunderscore}{\kern0pt}of{\isacharunderscore}{\kern0pt}plan{\isacharunderscore}{\kern0pt}action\ {\isasympi}s\ {\isacharparenleft}{\kern0pt}t\isactrlsub {\isasympi}{\isacharcomma}{\kern0pt}{\isasympi}{\isacharparenright}{\kern0pt}\ {\isacharparenleft}{\kern0pt}t\isactrlsub a{\isacharcomma}{\kern0pt}a{\isacharparenright}{\kern0pt}{\isacartoucheclose}\ \isakeywordONE{by}\isamarkupfalse%
\ simp\isanewline
\ \ \ \ \isakeywordONE{then}\isamarkupfalse%
\ \isakeywordTHREE{show}\isamarkupfalse%
\ {\isacharquery}{\kern0pt}thesis\isanewline
\ \ \ \ \ \ \isakeywordONE{proof}\isamarkupfalse%
\ {\isacharparenleft}{\kern0pt}elim\ disjE{\isacharparenright}{\kern0pt}\isanewline
\ \ \ \ \ \ \ \ \isakeywordTHREE{assume}\isamarkupfalse%
\ {\isacharquery}{\kern0pt}case{\isadigit{1}}\isanewline
\ \ \ \ \ \ \ \ \isakeywordONE{then}\isamarkupfalse%
\ \isakeywordTHREE{show}\isamarkupfalse%
\ {\isacharquery}{\kern0pt}thesis\isanewline
\ \ \ \ \ \ \ \ \ \ \isakeywordONE{using}\isamarkupfalse%
\ {\isacartoucheopen}{\isasymnot}is{\isacharunderscore}{\kern0pt}htp\ {\isasympi}s\ t\isactrlsub a{\isacartoucheclose}\ {\isacartoucheopen}is{\isacharunderscore}{\kern0pt}htp\ {\isasympi}s\ t\isactrlsub {\isasympi}{\isacartoucheclose}\ \isakeywordONE{by}\isamarkupfalse%
\ auto\isanewline
\ \ \ \ \ \ \isakeywordONE{next}\isamarkupfalse%
\isanewline
\ \ \ \ \ \ \ \ \isakeywordTHREE{assume}\isamarkupfalse%
\ {\isacharquery}{\kern0pt}case{\isadigit{2}}\isanewline
\ \ \ \ \ \ \ \ \isakeywordONE{then}\isamarkupfalse%
\ \isakeywordTHREE{show}\isamarkupfalse%
\ {\isacharquery}{\kern0pt}thesis\isanewline
\ \ \ \ \ \ \ \ \ \ \isakeywordONE{using}\isamarkupfalse%
\ {\isacartoucheopen}{\isasymnot}is{\isacharunderscore}{\kern0pt}htp\ {\isasympi}s\ t\isactrlsub a{\isacartoucheclose}\ {\isacartoucheopen}is{\isacharunderscore}{\kern0pt}htp\ {\isasympi}s\ t\isactrlsub {\isasympi}{\isacartoucheclose}\ \isakeywordONE{by}\isamarkupfalse%
\ auto\isanewline
\ \ \ \ \ \ \isakeywordONE{next}\isamarkupfalse%
\isanewline
\ \ \ \ \ \ \ \ \isakeywordTHREE{assume}\isamarkupfalse%
\ {\isacharquery}{\kern0pt}case{\isadigit{3}}\isanewline
\ \ \ \ \ \ \ \ \isakeywordONE{then}\isamarkupfalse%
\ \isakeywordONE{have}\isamarkupfalse%
\ {\isachardoublequoteopen}{\isacharparenleft}{\kern0pt}t\isactrlsub {\isasympi}{\isacharcomma}{\kern0pt}{\isasympi}{\isacharparenright}{\kern0pt}\ {\isasymin}\ durative{\isacharunderscore}{\kern0pt}acts\ {\isasympi}s{\isachardoublequoteclose}\isanewline
\ \ \ \ \ \ \ \ \ \ \isakeywordONE{unfolding}\isamarkupfalse%
\ durative{\isacharunderscore}{\kern0pt}acts{\isacharunderscore}{\kern0pt}def\ is{\isacharunderscore}{\kern0pt}act{\isacharunderscore}{\kern0pt}simple{\isacharunderscore}{\kern0pt}def\isanewline
\ \ \ \ \ \ \ \ \ \ \isakeywordONE{using}\isamarkupfalse%
\ {\isacartoucheopen}{\isacharparenleft}{\kern0pt}t\isactrlsub {\isasympi}{\isacharcomma}{\kern0pt}{\isasympi}{\isacharparenright}{\kern0pt}\ {\isasymin}\ set\ {\isasympi}s{\isacartoucheclose}\ \isakeywordONE{by}\isamarkupfalse%
\ {\isacharparenleft}{\kern0pt}cases\ {\isasympi}{\isacharparenright}{\kern0pt}\ auto\isanewline
\ \ \ \ \ \ \ \ \isakeywordONE{then}\isamarkupfalse%
\ \isakeywordONE{have}\isamarkupfalse%
\ {\isachardoublequoteopen}is{\isacharunderscore}{\kern0pt}htp\ {\isasympi}s\ {\isacharparenleft}{\kern0pt}t\isactrlsub {\isasympi}\ {\isacharplus}{\kern0pt}\ {\isacharparenleft}{\kern0pt}duration\ {\isasympi}{\isacharparenright}{\kern0pt}{\isacharparenright}{\kern0pt}{\isachardoublequoteclose}\isanewline
\ \ \ \ \ \ \ \ \ \ \isakeywordONE{unfolding}\isamarkupfalse%
\ is{\isacharunderscore}{\kern0pt}htp{\isacharunderscore}{\kern0pt}def\ \isakeywordONE{by}\isamarkupfalse%
\ auto\isanewline
\ \ \ \ \ \ \ \ \isakeywordONE{then}\isamarkupfalse%
\ \isakeywordTHREE{show}\isamarkupfalse%
\ {\isacharquery}{\kern0pt}thesis\isanewline
\ \ \ \ \ \ \ \ \ \ \isakeywordONE{using}\isamarkupfalse%
\ {\isacartoucheopen}{\isacharquery}{\kern0pt}case{\isadigit{3}}{\isacartoucheclose}\ {\isacartoucheopen}{\isasymnot}is{\isacharunderscore}{\kern0pt}htp\ {\isasympi}s\ t\isactrlsub a{\isacartoucheclose}\ {\isacartoucheopen}is{\isacharunderscore}{\kern0pt}htp\ {\isasympi}s\ {\isacharparenleft}{\kern0pt}t\isactrlsub {\isasympi}\ {\isacharplus}{\kern0pt}\ {\isacharparenleft}{\kern0pt}duration\ {\isasympi}{\isacharparenright}{\kern0pt}{\isacharparenright}{\kern0pt}{\isacartoucheclose}\ \isakeywordONE{by}\isamarkupfalse%
\ auto\isanewline
\ \ \ \ \ \ \isakeywordONE{next}\isamarkupfalse%
\isanewline
\ \ \ \ \ \ \ \ \isakeywordTHREE{assume}\isamarkupfalse%
\ {\isacharquery}{\kern0pt}case{\isadigit{4}}\isanewline
\ \ \ \ \ \ \ \ \isakeywordONE{let}\isamarkupfalse%
\ {\isacharquery}{\kern0pt}a\isactrlsub i\isactrlsub n\isactrlsub v{\isacharequal}{\kern0pt}{\isachardoublequoteopen}the\ {\isacharparenleft}{\kern0pt}res{\isacharunderscore}{\kern0pt}inst{\isacharunderscore}{\kern0pt}snap{\isacharunderscore}{\kern0pt}action\ {\isasympi}\ Over{\isacharunderscore}{\kern0pt}All{\isacharparenright}{\kern0pt}{\isachardoublequoteclose}\isanewline
\ \ \ \ \ \ \ \ \isakeywordONE{have}\isamarkupfalse%
\ {\isachardoublequoteopen}wf{\isacharunderscore}{\kern0pt}plan{\isacharunderscore}{\kern0pt}action\ {\isasympi}{\isachardoublequoteclose}\isanewline
\ \ \ \ \ \ \ \ \ \ \isakeywordONE{using}\isamarkupfalse%
\ {\isacartoucheopen}wf{\isacharunderscore}{\kern0pt}plan\ {\isasympi}s{\isacartoucheclose}\ {\isacartoucheopen}{\isacharparenleft}{\kern0pt}t\isactrlsub {\isasympi}{\isacharcomma}{\kern0pt}{\isasympi}{\isacharparenright}{\kern0pt}\ {\isasymin}\ set\ {\isasympi}s{\isacartoucheclose}\isanewline
\ \ \ \ \ \ \ \ \ \ \isakeywordONE{unfolding}\isamarkupfalse%
\ wf{\isacharunderscore}{\kern0pt}plan{\isacharunderscore}{\kern0pt}def\ \isakeywordONE{by}\isamarkupfalse%
\ auto\isanewline
\ \ \ \ \ \ \ \ \isakeywordTHREE{obtain}\isamarkupfalse%
\ a\isactrlsub s\isactrlsub c\isactrlsub h\isactrlsub e\isactrlsub m\isactrlsub a\ \isakeywordTWO{where}\ a\isactrlsub s\isactrlsub c\isactrlsub h\isactrlsub e\isactrlsub m\isactrlsub a{\isacharunderscore}{\kern0pt}def{\isacharcolon}{\kern0pt}\ {\isachardoublequoteopen}Some\ a\isactrlsub s\isactrlsub c\isactrlsub h\isactrlsub e\isactrlsub m\isactrlsub a\ {\isacharequal}{\kern0pt}\ resolve{\isacharunderscore}{\kern0pt}action{\isacharunderscore}{\kern0pt}schema\ {\isacharparenleft}{\kern0pt}name\ {\isasympi}{\isacharparenright}{\kern0pt}{\isachardoublequoteclose}\isanewline
\ \ \ \ \ \ \ \ \ \ \isakeywordONE{using}\isamarkupfalse%
\ {\isacartoucheopen}wf{\isacharunderscore}{\kern0pt}plan{\isacharunderscore}{\kern0pt}action\ {\isasympi}{\isacartoucheclose}\ \isakeywordONE{by}\isamarkupfalse%
\ {\isacharparenleft}{\kern0pt}cases\ {\isasympi}{\isacharparenright}{\kern0pt}\ {\isacharparenleft}{\kern0pt}auto\ split{\isacharcolon}{\kern0pt}\ option{\isachardot}{\kern0pt}splits{\isacharparenright}{\kern0pt}\isanewline
\ \ \ \ \ \ \ \ \isakeywordONE{then}\isamarkupfalse%
\ \isakeywordONE{have}\isamarkupfalse%
\ {\isachardoublequoteopen}a\isactrlsub s\isactrlsub c\isactrlsub h\isactrlsub e\isactrlsub m\isactrlsub a\ {\isasymin}\ set\ {\isacharparenleft}{\kern0pt}actions\ D{\isacharparenright}{\kern0pt}{\isachardoublequoteclose}\isanewline
\ \ \ \ \ \ \ \ \ \ \isakeywordONE{unfolding}\isamarkupfalse%
\ resolve{\isacharunderscore}{\kern0pt}action{\isacharunderscore}{\kern0pt}schema{\isacharunderscore}{\kern0pt}def\ \isakeywordONE{by}\isamarkupfalse%
\ {\isacharparenleft}{\kern0pt}metis\ index{\isacharunderscore}{\kern0pt}by{\isacharunderscore}{\kern0pt}eq{\isacharunderscore}{\kern0pt}SomeD{\isacharparenright}{\kern0pt}\isanewline
\ \ \ \ \ \ \ \ \isakeywordONE{then}\isamarkupfalse%
\ \isakeywordONE{have}\isamarkupfalse%
\ {\isachardoublequoteopen}wf{\isacharunderscore}{\kern0pt}action{\isacharunderscore}{\kern0pt}schema\ a\isactrlsub s\isactrlsub c\isactrlsub h\isactrlsub e\isactrlsub m\isactrlsub a{\isachardoublequoteclose}\isanewline
\ \ \ \ \ \ \ \ \ \ \isakeywordONE{using}\isamarkupfalse%
\ wf{\isacharunderscore}{\kern0pt}domain\isanewline
\ \ \ \ \ \ \ \ \ \ \isakeywordONE{unfolding}\isamarkupfalse%
\ wf{\isacharunderscore}{\kern0pt}domain{\isacharunderscore}{\kern0pt}def\ \isakeywordONE{by}\isamarkupfalse%
\ auto\isanewline
\ \ \ \ \ \ \ \ \isakeywordONE{then}\isamarkupfalse%
\ \isakeywordTHREE{show}\isamarkupfalse%
\ {\isacharquery}{\kern0pt}thesis\isanewline
\ \ \ \ \ \ \ \ \ \ \isakeywordONE{using}\isamarkupfalse%
\ {\isacartoucheopen}{\isacharquery}{\kern0pt}case{\isadigit{4}}{\isacartoucheclose}\ wf{\isacharunderscore}{\kern0pt}over{\isacharunderscore}{\kern0pt}all{\isacharunderscore}{\kern0pt}empty{\isacharunderscore}{\kern0pt}eff{\isacharbrackleft}{\kern0pt}OF\ {\isacartoucheopen}wf{\isacharunderscore}{\kern0pt}plan{\isacharunderscore}{\kern0pt}action\ {\isasympi}{\isacartoucheclose}\ \isanewline
\ \ \ \ \ \ \ \ \ \ \ \ \ \ \ \ \ \ \ \ \ \ \ \ \ \ \ \ \ \ \ \ \ \ \ \ \ \ \ \ \ \ \ \ \ \ \ \ \ \ {\isacartoucheopen}wf{\isacharunderscore}{\kern0pt}action{\isacharunderscore}{\kern0pt}schema\ a\isactrlsub s\isactrlsub c\isactrlsub h\isactrlsub e\isactrlsub m\isactrlsub a{\isacartoucheclose}\ a\isactrlsub s\isactrlsub c\isactrlsub h\isactrlsub e\isactrlsub m\isactrlsub a{\isacharunderscore}{\kern0pt}def{\isacharbrackright}{\kern0pt}\isanewline
\ \ \ \ \ \ \ \ \ \ \isakeywordONE{by}\isamarkupfalse%
\ auto\isanewline
\ \ \ \ \ \ \isakeywordONE{qed}\isamarkupfalse%
\isanewline
\ \ \ \ \isakeywordONE{qed}\isamarkupfalse%
%
\endisatagproof
{\isafoldproof}%
%
\isadelimproof
%
\endisadelimproof
%
\begin{isamarkuptext}%
Applying a happening containing only action with empty effect, doesn't 
  change the current state.%
\end{isamarkuptext}\isamarkuptrue%
\ \ \isakeywordONE{lemma}\isamarkupfalse%
\ apply{\isacharunderscore}{\kern0pt}happ{\isacharunderscore}{\kern0pt}empty{\isacharunderscore}{\kern0pt}effect{\isacharcolon}{\kern0pt}\isanewline
\ \ \ \ \isakeywordTWO{assumes}\ {\isachardoublequoteopen}{\isasymforall}{\isacharparenleft}{\kern0pt}t\isactrlsub a{\isacharcomma}{\kern0pt}A{\isacharparenright}{\kern0pt}\ {\isasymin}\ set\ hs{\isachardot}{\kern0pt}\ {\isasymforall}a\ {\isasymin}\ set\ A{\isachardot}{\kern0pt}\ effect\ a\ {\isacharequal}{\kern0pt}\ Effect\ {\isacharbrackleft}{\kern0pt}{\isacharbrackright}{\kern0pt}\ {\isacharbrackleft}{\kern0pt}{\isacharbrackright}{\kern0pt}{\isachardoublequoteclose}\isanewline
\ \ \ \ \isakeywordTWO{shows}\ {\isachardoublequoteopen}{\isasymforall}h\ {\isasymin}\ set\ hs{\isachardot}{\kern0pt}\ apply{\isacharunderscore}{\kern0pt}happ\ h\ M\ {\isacharequal}{\kern0pt}\ M{\isachardoublequoteclose}\isanewline
%
\isadelimproof
\ \ \ \ %
\endisadelimproof
%
\isatagproof
\isakeywordONE{using}\isamarkupfalse%
\ assms\ \isakeywordONE{by}\isamarkupfalse%
\ {\isacharparenleft}{\kern0pt}induction\ hs{\isacharparenright}{\kern0pt}\ auto%
\endisatagproof
{\isafoldproof}%
%
\isadelimproof
%
\endisadelimproof
%
\begin{isamarkuptext}%
Combining \isa{{\isasymlbrakk}\isaconst{wf{\isacharunderscore}{\kern0pt}ast{\isacharunderscore}{\kern0pt}problem}\ \isavar{{\isacharquery}{\kern0pt}P}{\isacharsemicolon}{\kern0pt}\ \isaconst{consec{\isacharunderscore}{\kern0pt}htps}\ \isavar{{\isacharquery}{\kern0pt}{\isasympi}s}\ \isavar{{\isacharquery}{\kern0pt}t\isactrlsub i}\ \isavar{{\isacharquery}{\kern0pt}t\isactrlsub j}{\isacharsemicolon}{\kern0pt}\ \isaconst{ast{\isacharunderscore}{\kern0pt}problem{\isachardot}{\kern0pt}wf{\isacharunderscore}{\kern0pt}plan}\ \isavar{{\isacharquery}{\kern0pt}P}\ \isavar{{\isacharquery}{\kern0pt}{\isasympi}s}{\isacharsemicolon}{\kern0pt}\ \isaconst{ast{\isacharunderscore}{\kern0pt}problem{\isachardot}{\kern0pt}ind{\isacharunderscore}{\kern0pt}happ{\isacharunderscore}{\kern0pt}seq}\ \isavar{{\isacharquery}{\kern0pt}P}\ \isavar{{\isacharquery}{\kern0pt}{\isasympi}s}\ \isavar{{\isacharquery}{\kern0pt}hs}{\isacharsemicolon}{\kern0pt}\ {\isacharparenleft}{\kern0pt}\isavar{{\isacharquery}{\kern0pt}t\isactrlsub a}{\isacharcomma}{\kern0pt}\ \isavar{{\isacharquery}{\kern0pt}A}{\isacharparenright}{\kern0pt}\ {\isasymin}\ \isaconst{set}\ {\isacharparenleft}{\kern0pt}\isaconst{dropWhile}\ {\isacharparenleft}{\kern0pt}\isaconst{leq}\ \isavar{{\isacharquery}{\kern0pt}t\isactrlsub i}{\isacharparenright}{\kern0pt}\ {\isacharparenleft}{\kern0pt}\isaconst{takeWhile}\ {\isacharparenleft}{\kern0pt}\isaconst{le}\ \isavar{{\isacharquery}{\kern0pt}t\isactrlsub j}{\isacharparenright}{\kern0pt}\ \isavar{{\isacharquery}{\kern0pt}hs}{\isacharparenright}{\kern0pt}{\isacharparenright}{\kern0pt}{\isacharsemicolon}{\kern0pt}\ \isavar{{\isacharquery}{\kern0pt}a}\ {\isasymin}\ \isaconst{set}\ \isavar{{\isacharquery}{\kern0pt}A}{\isasymrbrakk}\ {\isasymLongrightarrow}\ \isaconst{ground{\isacharunderscore}{\kern0pt}action{\isachardot}{\kern0pt}effect}\ \isavar{{\isacharquery}{\kern0pt}a}\ {\isacharequal}{\kern0pt}\ \isaconst{Effect}\ {\isacharbrackleft}{\kern0pt}{\isacharbrackright}{\kern0pt}\ {\isacharbrackleft}{\kern0pt}{\isacharbrackright}{\kern0pt}} 
  and \isa{{\isasymforall}{\isacharparenleft}{\kern0pt}\isabound{t\isactrlsub a}{\isacharcomma}{\kern0pt}\ \isabound{A}{\isacharparenright}{\kern0pt}{\isasymin}\isaconst{set}\ \isavar{{\isacharquery}{\kern0pt}hs}{\isachardot}{\kern0pt}\ {\isasymforall}\isabound{a}{\isasymin}\isaconst{set}\ \isabound{A}{\isachardot}{\kern0pt}\ \isaconst{ground{\isacharunderscore}{\kern0pt}action{\isachardot}{\kern0pt}effect}\ \isabound{a}\ {\isacharequal}{\kern0pt}\ \isaconst{Effect}\ {\isacharbrackleft}{\kern0pt}{\isacharbrackright}{\kern0pt}\ {\isacharbrackleft}{\kern0pt}{\isacharbrackright}{\kern0pt}\ {\isasymLongrightarrow}\ {\isasymforall}\isabound{h}{\isasymin}\isaconst{set}\ \isavar{{\isacharquery}{\kern0pt}hs}{\isachardot}{\kern0pt}\ \isaconst{apply{\isacharunderscore}{\kern0pt}happ}\ \isabound{h}\ \isavar{{\isacharquery}{\kern0pt}M}\ {\isacharequal}{\kern0pt}\ \isavar{{\isacharquery}{\kern0pt}M}}: Applying the happening in between two consecutive 
  happening time points doesn't change the current state.%
\end{isamarkuptext}\isamarkuptrue%
\ \ \isakeywordONE{lemma}\isamarkupfalse%
\ ind{\isacharunderscore}{\kern0pt}happ{\isacharunderscore}{\kern0pt}seq{\isacharunderscore}{\kern0pt}leq\isactrlsub i{\isacharunderscore}{\kern0pt}le\isactrlsub j{\isacharcolon}{\kern0pt}\isanewline
\ \ \ \ \isakeywordTWO{assumes}\ {\isachardoublequoteopen}consec{\isacharunderscore}{\kern0pt}htps\ {\isasympi}s\ t\isactrlsub i\ t\isactrlsub j{\isachardoublequoteclose}\isanewline
\ \ \ \ \ \ \ \ \isakeywordTWO{and}\ {\isachardoublequoteopen}wf{\isacharunderscore}{\kern0pt}plan\ {\isasympi}s{\isachardoublequoteclose}\isanewline
\ \ \ \ \ \ \ \ \isakeywordTWO{and}\ {\isachardoublequoteopen}ind{\isacharunderscore}{\kern0pt}happ{\isacharunderscore}{\kern0pt}seq\ {\isasympi}s\ hs{\isachardoublequoteclose}\isanewline
\ \ \ \ \ \ \isakeywordTWO{shows}\ {\isachardoublequoteopen}{\isasymforall}h\ {\isasymin}\ set\ {\isacharparenleft}{\kern0pt}dropWhile\ {\isacharparenleft}{\kern0pt}leq\ t\isactrlsub i{\isacharparenright}{\kern0pt}\ {\isacharparenleft}{\kern0pt}takeWhile\ {\isacharparenleft}{\kern0pt}le\ t\isactrlsub j{\isacharparenright}{\kern0pt}\ hs{\isacharparenright}{\kern0pt}{\isacharparenright}{\kern0pt}{\isachardot}{\kern0pt}\ apply{\isacharunderscore}{\kern0pt}happ\ h\ M\ {\isacharequal}{\kern0pt}\ M{\isachardoublequoteclose}\isanewline
%
\isadelimproof
\ \ \ \ %
\endisadelimproof
%
\isatagproof
\isakeywordONE{using}\isamarkupfalse%
\ assms\ apply{\isacharunderscore}{\kern0pt}happ{\isacharunderscore}{\kern0pt}empty{\isacharunderscore}{\kern0pt}effect\ ind{\isacharunderscore}{\kern0pt}happ{\isacharunderscore}{\kern0pt}seq{\isacharunderscore}{\kern0pt}leq\isactrlsub i{\isacharunderscore}{\kern0pt}le\isactrlsub j{\isacharunderscore}{\kern0pt}empty{\isacharunderscore}{\kern0pt}eff\isanewline
\ \ \ \ \isakeywordONE{by}\isamarkupfalse%
\ auto%
\endisatagproof
{\isafoldproof}%
%
\isadelimproof
\isanewline
%
\endisadelimproof
\isanewline
\ \ \isakeywordONE{lemma}\isamarkupfalse%
\ valid{\isacharunderscore}{\kern0pt}happ{\isacharunderscore}{\kern0pt}seq{\isacharunderscore}{\kern0pt}M\isactrlsub i{\isacharunderscore}{\kern0pt}M\isactrlsub i{\isacharunderscore}{\kern0pt}precond{\isacharcolon}{\kern0pt}\isanewline
\ \ \ \ \isakeywordTWO{assumes}\ {\isachardoublequoteopen}{\isasymforall}h\ {\isasymin}\ set\ hs{\isachardot}{\kern0pt}\ apply{\isacharunderscore}{\kern0pt}happ\ h\ M\isactrlsub i\ {\isacharequal}{\kern0pt}\ M\isactrlsub i{\isachardoublequoteclose}\isanewline
\ \ \ \ \ \ \ \ \isakeywordTWO{and}\ {\isachardoublequoteopen}valid{\isacharunderscore}{\kern0pt}happ{\isacharunderscore}{\kern0pt}seq\ M\isactrlsub i\ hs\ M\isactrlsub i{\isachardoublequoteclose}\isanewline
\ \ \ \ \ \ \isakeywordTWO{shows}\ {\isachardoublequoteopen}{\isasymforall}{\isacharparenleft}{\kern0pt}t\isactrlsub a{\isacharcomma}{\kern0pt}as{\isacharparenright}{\kern0pt}\ {\isasymin}\ set\ hs{\isachardot}{\kern0pt}\ {\isasymforall}a{\isasymin}\ set\ as{\isachardot}{\kern0pt}\ M\isactrlsub i\ \isactrlsup c{\isasymTTurnstile}\isactrlsub {\isacharequal}{\kern0pt}\ precondition\ a{\isachardoublequoteclose}\isanewline
%
\isadelimproof
\ \ \ \ %
\endisadelimproof
%
\isatagproof
\isakeywordONE{using}\isamarkupfalse%
\ assms\ \isakeywordONE{by}\isamarkupfalse%
\ {\isacharparenleft}{\kern0pt}induction\ M\isactrlsub i\ hs\ M\isactrlsub i\ rule{\isacharcolon}{\kern0pt}\ valid{\isacharunderscore}{\kern0pt}happ{\isacharunderscore}{\kern0pt}seq{\isachardot}{\kern0pt}induct{\isacharparenright}{\kern0pt}\ auto%
\endisatagproof
{\isafoldproof}%
%
\isadelimproof
\isanewline
%
\endisadelimproof
\isanewline
\ \ \isakeywordONE{lemma}\isamarkupfalse%
\ valid{\isacharunderscore}{\kern0pt}happ{\isacharunderscore}{\kern0pt}seq{\isacharunderscore}{\kern0pt}leq\isactrlsub i{\isacharunderscore}{\kern0pt}le\isactrlsub j{\isacharunderscore}{\kern0pt}aux{\isadigit{1}}{\isacharcolon}{\kern0pt}\isanewline
\ \ \ \ \isakeywordTWO{assumes}\ {\isachardoublequoteopen}consec{\isacharunderscore}{\kern0pt}htps\ {\isasympi}s\ t\isactrlsub i\ t\isactrlsub j{\isachardoublequoteclose}\isanewline
\ \ \ \ \ \ \ \ \isakeywordTWO{and}\ {\isachardoublequoteopen}wf{\isacharunderscore}{\kern0pt}plan\ {\isasympi}s{\isachardoublequoteclose}\isanewline
\ \ \ \ \ \ \ \ \isakeywordTWO{and}\ {\isachardoublequoteopen}ind{\isacharunderscore}{\kern0pt}happ{\isacharunderscore}{\kern0pt}seq\ {\isasympi}s\ hs{\isachardoublequoteclose}\isanewline
\ \ \ \ \ \ \ \ \isakeywordTWO{and}\ {\isachardoublequoteopen}valid{\isacharunderscore}{\kern0pt}happ{\isacharunderscore}{\kern0pt}seq\ M\isactrlsub i\ {\isacharparenleft}{\kern0pt}dropWhile\ {\isacharparenleft}{\kern0pt}leq\ t\isactrlsub i{\isacharparenright}{\kern0pt}\ {\isacharparenleft}{\kern0pt}takeWhile\ {\isacharparenleft}{\kern0pt}le\ t\isactrlsub j{\isacharparenright}{\kern0pt}\ hs{\isacharparenright}{\kern0pt}{\isacharparenright}{\kern0pt}\ M\isactrlsub j{\isachardoublequoteclose}\isanewline
\ \ \ \ \ \ \isakeywordTWO{shows}\ {\isachardoublequoteopen}{\isasymforall}{\isacharparenleft}{\kern0pt}t\isactrlsub a{\isacharcomma}{\kern0pt}as{\isacharparenright}{\kern0pt}\ {\isasymin}\ set\ {\isacharparenleft}{\kern0pt}dropWhile\ {\isacharparenleft}{\kern0pt}leq\ t\isactrlsub i{\isacharparenright}{\kern0pt}\ {\isacharparenleft}{\kern0pt}takeWhile\ {\isacharparenleft}{\kern0pt}le\ t\isactrlsub j{\isacharparenright}{\kern0pt}\ hs{\isacharparenright}{\kern0pt}{\isacharparenright}{\kern0pt}{\isachardot}{\kern0pt}\ {\isasymforall}a{\isasymin}\ set\ as{\isachardot}{\kern0pt}\ M\isactrlsub i\ \isactrlsup c{\isasymTTurnstile}\isactrlsub {\isacharequal}{\kern0pt}\ precondition\ a{\isachardoublequoteclose}\isanewline
%
\isadelimproof
\ \ %
\endisadelimproof
%
\isatagproof
\isakeywordONE{proof}\isamarkupfalse%
\ {\isacharminus}{\kern0pt}\isanewline
\ \ \ \ \isakeywordONE{let}\isamarkupfalse%
\ {\isacharquery}{\kern0pt}hs{\isacharprime}{\kern0pt}{\isacharequal}{\kern0pt}{\isachardoublequoteopen}dropWhile\ {\isacharparenleft}{\kern0pt}leq\ t\isactrlsub i{\isacharparenright}{\kern0pt}\ {\isacharparenleft}{\kern0pt}takeWhile\ {\isacharparenleft}{\kern0pt}le\ t\isactrlsub j{\isacharparenright}{\kern0pt}\ hs{\isacharparenright}{\kern0pt}{\isachardoublequoteclose}\isanewline
\ \ \ \ \isakeywordONE{have}\isamarkupfalse%
\ {\isachardoublequoteopen}{\isasymforall}h\ {\isasymin}\ set\ {\isacharquery}{\kern0pt}hs{\isacharprime}{\kern0pt}{\isachardot}{\kern0pt}\ apply{\isacharunderscore}{\kern0pt}happ\ h\ M\isactrlsub i\ {\isacharequal}{\kern0pt}\ M\isactrlsub i{\isachardoublequoteclose}\isanewline
\ \ \ \ \ \ \isakeywordONE{using}\isamarkupfalse%
\ assms\ ind{\isacharunderscore}{\kern0pt}happ{\isacharunderscore}{\kern0pt}seq{\isacharunderscore}{\kern0pt}leq\isactrlsub i{\isacharunderscore}{\kern0pt}le\isactrlsub j\ \isakeywordONE{by}\isamarkupfalse%
\ simp\isanewline
\ \ \ \ \isakeywordONE{have}\isamarkupfalse%
\ {\isachardoublequoteopen}valid{\isacharunderscore}{\kern0pt}happ{\isacharunderscore}{\kern0pt}seq\ M\isactrlsub i\ {\isacharquery}{\kern0pt}hs{\isacharprime}{\kern0pt}\ M\isactrlsub i{\isachardoublequoteclose}\isanewline
\ \ \ \ \ \ \isakeywordONE{using}\isamarkupfalse%
\ assms\ valid{\isacharunderscore}{\kern0pt}happ{\isacharunderscore}{\kern0pt}seq{\isacharunderscore}{\kern0pt}M\isactrlsub j{\isacharunderscore}{\kern0pt}is{\isacharunderscore}{\kern0pt}M\isactrlsub i{\isacharbrackleft}{\kern0pt}OF\ {\isacartoucheopen}{\isasymforall}h\ {\isasymin}\ set\ {\isacharquery}{\kern0pt}hs{\isacharprime}{\kern0pt}{\isachardot}{\kern0pt}\ apply{\isacharunderscore}{\kern0pt}happ\ h\ M\isactrlsub i\ {\isacharequal}{\kern0pt}\ M\isactrlsub i{\isacartoucheclose}{\isacharbrackright}{\kern0pt}\ \isakeywordONE{by}\isamarkupfalse%
\ auto\isanewline
\ \ \ \ \isakeywordONE{then}\isamarkupfalse%
\ \isakeywordTHREE{show}\isamarkupfalse%
\ {\isacharquery}{\kern0pt}thesis\isanewline
\ \ \ \ \ \ \isakeywordONE{using}\isamarkupfalse%
\ assms\ {\isacartoucheopen}{\isasymforall}h\ {\isasymin}\ set\ {\isacharquery}{\kern0pt}hs{\isacharprime}{\kern0pt}{\isachardot}{\kern0pt}\ apply{\isacharunderscore}{\kern0pt}happ\ h\ M\isactrlsub i\ {\isacharequal}{\kern0pt}\ M\isactrlsub i{\isacartoucheclose}\ valid{\isacharunderscore}{\kern0pt}happ{\isacharunderscore}{\kern0pt}seq{\isacharunderscore}{\kern0pt}M\isactrlsub i{\isacharunderscore}{\kern0pt}M\isactrlsub i{\isacharunderscore}{\kern0pt}precond\ \isakeywordONE{by}\isamarkupfalse%
\ simp\isanewline
\ \ \isakeywordONE{qed}\isamarkupfalse%
%
\endisatagproof
{\isafoldproof}%
%
\isadelimproof
\isanewline
%
\endisadelimproof
\isanewline
\ \ \isakeywordONE{lemma}\isamarkupfalse%
\ valid{\isacharunderscore}{\kern0pt}happ{\isacharunderscore}{\kern0pt}seq{\isacharunderscore}{\kern0pt}leq\isactrlsub i{\isacharunderscore}{\kern0pt}le\isactrlsub j{\isacharunderscore}{\kern0pt}aux{\isadigit{2}}{\isacharcolon}{\kern0pt}\isanewline
\ \ \ \ \isakeywordTWO{assumes}\ {\isachardoublequoteopen}{\isasymforall}h\ {\isasymin}\ set\ hs{\isachardot}{\kern0pt}\ apply{\isacharunderscore}{\kern0pt}happ\ h\ M\isactrlsub i\ {\isacharequal}{\kern0pt}\ M\isactrlsub i{\isachardoublequoteclose}\isanewline
\ \ \ \ \ \ \ \ \isakeywordTWO{and}\ {\isachardoublequoteopen}{\isasymforall}{\isacharparenleft}{\kern0pt}t\isactrlsub a{\isacharcomma}{\kern0pt}A{\isacharparenright}{\kern0pt}\ {\isasymin}\ set\ hs{\isachardot}{\kern0pt}\ {\isasymforall}a{\isasymin}\ set\ A{\isachardot}{\kern0pt}\ M\isactrlsub i\ \isactrlsup c{\isasymTTurnstile}\isactrlsub {\isacharequal}{\kern0pt}\ precondition\ a{\isachardoublequoteclose}\isanewline
\ \ \ \ \ \ \ \ \isakeywordTWO{and}\ {\isachardoublequoteopen}{\isasymforall}{\isacharparenleft}{\kern0pt}t\isactrlsub a{\isacharcomma}{\kern0pt}A{\isacharparenright}{\kern0pt}\ {\isasymin}\ set\ hs{\isachardot}{\kern0pt}\ {\isasymforall}{\isacharparenleft}{\kern0pt}t\isactrlsub a{\isacharprime}{\kern0pt}{\isacharcomma}{\kern0pt}A{\isacharprime}{\kern0pt}{\isacharparenright}{\kern0pt}\ {\isasymin}\ set\ hs{\isachardot}{\kern0pt}\ \isanewline
\ \ \ \ \ \ \ \ \ \ \ \ \ {\isasymforall}a\ {\isasymin}\ set\ A{\isachardot}{\kern0pt}\ {\isasymforall}a{\isacharprime}{\kern0pt}\ {\isasymin}\ set\ A{\isacharprime}{\kern0pt}{\isachardot}{\kern0pt}\ acts{\isacharunderscore}{\kern0pt}non{\isacharunderscore}{\kern0pt}intrf\ a\ a{\isacharprime}{\kern0pt}{\isachardoublequoteclose}\isanewline
\ \ \ \ \isakeywordTWO{shows}\ {\isachardoublequoteopen}valid{\isacharunderscore}{\kern0pt}happ{\isacharunderscore}{\kern0pt}seq\ M\isactrlsub i\ hs\ M\isactrlsub i{\isachardoublequoteclose}\isanewline
%
\isadelimproof
\ \ \ \ %
\endisadelimproof
%
\isatagproof
\isakeywordONE{using}\isamarkupfalse%
\ assms\isanewline
\ \ \ \ \isakeywordONE{by}\isamarkupfalse%
\ {\isacharparenleft}{\kern0pt}induction\ hs{\isacharparenright}{\kern0pt}\ {\isacharparenleft}{\kern0pt}auto\ split{\isacharcolon}{\kern0pt}\ prod{\isachardot}{\kern0pt}splits\ simp{\isacharcolon}{\kern0pt}\ pairwise{\isacharunderscore}{\kern0pt}def{\isacharparenright}{\kern0pt}%
\endisatagproof
{\isafoldproof}%
%
\isadelimproof
\isanewline
%
\endisadelimproof
\isanewline
\isakeywordTWO{end}\isamarkupfalse%
\ %
\isamarkupcmt{Context \isa{wf{\isacharunderscore}{\kern0pt}ast{\isacharunderscore}{\kern0pt}problem}%
}\ \isanewline
\isanewline
\isakeywordONE{context}\isamarkupfalse%
\ ast{\isacharunderscore}{\kern0pt}problem\isanewline
\isakeywordTWO{begin}%
\begin{isamarkuptext}%
All invariants from \isa{\isaconst{invs{\isacharunderscore}{\kern0pt}of{\isacharunderscore}{\kern0pt}plan{\isacharunderscore}{\kern0pt}at}} are satisfied iff an induced happening 
  sequence is valid in between two consecutive happening time points.%
\end{isamarkuptext}\isamarkuptrue%
\ \ \isakeywordONE{lemma}\isamarkupfalse%
\ {\isacharparenleft}{\kern0pt}\isakeywordTWO{in}\ wf{\isacharunderscore}{\kern0pt}ast{\isacharunderscore}{\kern0pt}problem{\isacharparenright}{\kern0pt}\ valid{\isacharunderscore}{\kern0pt}happ{\isacharunderscore}{\kern0pt}seq{\isacharunderscore}{\kern0pt}leq\isactrlsub i{\isacharunderscore}{\kern0pt}le\isactrlsub j{\isacharunderscore}{\kern0pt}invs{\isacharunderscore}{\kern0pt}of{\isacharunderscore}{\kern0pt}plan{\isacharunderscore}{\kern0pt}at{\isacharcolon}{\kern0pt}\isanewline
\ \ \ \ \isakeywordTWO{assumes}\ {\isachardoublequoteopen}wf{\isacharunderscore}{\kern0pt}plan\ {\isasympi}s{\isachardoublequoteclose}\isanewline
\ \ \ \ \ \ \ \ \isakeywordTWO{and}\ {\isachardoublequoteopen}consec{\isacharunderscore}{\kern0pt}htps\ {\isasympi}s\ t\isactrlsub i\ t\isactrlsub j{\isachardoublequoteclose}\isanewline
\ \ \ \ \ \ \ \ \isakeywordTWO{and}\ {\isachardoublequoteopen}ind{\isacharunderscore}{\kern0pt}happ{\isacharunderscore}{\kern0pt}seq\ {\isasympi}s\ hs{\isachardoublequoteclose}\isanewline
\ \ \ \ \isakeywordTWO{shows}\ {\isachardoublequoteopen}valid{\isacharunderscore}{\kern0pt}happ{\isacharunderscore}{\kern0pt}seq\ M\ {\isacharparenleft}{\kern0pt}dropWhile\ {\isacharparenleft}{\kern0pt}leq\ t\isactrlsub i{\isacharparenright}{\kern0pt}\ {\isacharparenleft}{\kern0pt}takeWhile\ {\isacharparenleft}{\kern0pt}le\ t\isactrlsub j{\isacharparenright}{\kern0pt}\ hs{\isacharparenright}{\kern0pt}{\isacharparenright}{\kern0pt}\ M\isanewline
\ \ \ \ \ \ \ \ \ \ \ \ {\isasymlongleftrightarrow}\ {\isacharparenleft}{\kern0pt}{\isasymforall}i\ {\isasymin}\ invs{\isacharunderscore}{\kern0pt}of{\isacharunderscore}{\kern0pt}plan{\isacharunderscore}{\kern0pt}at\ t\isactrlsub j\ {\isasympi}s{\isachardot}{\kern0pt}\ M\ \isactrlsup c{\isasymTTurnstile}\isactrlsub {\isacharequal}{\kern0pt}\ i{\isacharparenright}{\kern0pt}{\isachardoublequoteclose}\isanewline
%
\isadelimproof
\ \ %
\endisadelimproof
%
\isatagproof
\isakeywordONE{proof}\isamarkupfalse%
\isanewline
\ \ \ \ \isakeywordONE{let}\isamarkupfalse%
\ {\isacharquery}{\kern0pt}hs{\isacharprime}{\kern0pt}{\isacharequal}{\kern0pt}{\isachardoublequoteopen}dropWhile\ {\isacharparenleft}{\kern0pt}leq\ t\isactrlsub i{\isacharparenright}{\kern0pt}\ {\isacharparenleft}{\kern0pt}takeWhile\ {\isacharparenleft}{\kern0pt}le\ t\isactrlsub j{\isacharparenright}{\kern0pt}\ hs{\isacharparenright}{\kern0pt}{\isachardoublequoteclose}\isanewline
\ \ \ \ \isakeywordTHREE{assume}\isamarkupfalse%
\ {\isachardoublequoteopen}valid{\isacharunderscore}{\kern0pt}happ{\isacharunderscore}{\kern0pt}seq\ M\ {\isacharquery}{\kern0pt}hs{\isacharprime}{\kern0pt}\ M{\isachardoublequoteclose}\isanewline
\ \ \ \ \isakeywordONE{then}\isamarkupfalse%
\ \isakeywordTHREE{show}\isamarkupfalse%
\ {\isachardoublequoteopen}{\isasymforall}i\ {\isasymin}\ invs{\isacharunderscore}{\kern0pt}of{\isacharunderscore}{\kern0pt}plan{\isacharunderscore}{\kern0pt}at\ t\isactrlsub j\ {\isasympi}s{\isachardot}{\kern0pt}\ M\ \isactrlsup c{\isasymTTurnstile}\isactrlsub {\isacharequal}{\kern0pt}\ i{\isachardoublequoteclose}\isanewline
\ \ \ \ \ \ \isakeywordONE{using}\isamarkupfalse%
\ assms\ \ invs{\isacharunderscore}{\kern0pt}of{\isacharunderscore}{\kern0pt}plan{\isacharunderscore}{\kern0pt}at{\isacharunderscore}{\kern0pt}ind{\isacharunderscore}{\kern0pt}happ{\isacharunderscore}{\kern0pt}seq{\isacharbrackleft}{\kern0pt}OF\ assms{\isacharbrackright}{\kern0pt}\isanewline
\ \ \ \ \ \ \ \ \ \ \ \ valid{\isacharunderscore}{\kern0pt}happ{\isacharunderscore}{\kern0pt}seq{\isacharunderscore}{\kern0pt}leq\isactrlsub i{\isacharunderscore}{\kern0pt}le\isactrlsub j{\isacharunderscore}{\kern0pt}aux{\isadigit{1}}\ \isanewline
\ \ \ \ \ \ \ \isakeywordONE{by}\isamarkupfalse%
\ fastforce\isanewline
\ \ \isakeywordONE{next}\isamarkupfalse%
\ \isanewline
\ \ \ \ \isakeywordONE{let}\isamarkupfalse%
\ {\isacharquery}{\kern0pt}hs{\isacharprime}{\kern0pt}{\isacharequal}{\kern0pt}{\isachardoublequoteopen}dropWhile\ {\isacharparenleft}{\kern0pt}leq\ t\isactrlsub i{\isacharparenright}{\kern0pt}\ {\isacharparenleft}{\kern0pt}takeWhile\ {\isacharparenleft}{\kern0pt}le\ t\isactrlsub j{\isacharparenright}{\kern0pt}\ hs{\isacharparenright}{\kern0pt}{\isachardoublequoteclose}\isanewline
\ \ \ \ \isakeywordTHREE{assume}\isamarkupfalse%
\ {\isachardoublequoteopen}{\isasymforall}i\ {\isasymin}\ invs{\isacharunderscore}{\kern0pt}of{\isacharunderscore}{\kern0pt}plan{\isacharunderscore}{\kern0pt}at\ t\isactrlsub j\ {\isasympi}s{\isachardot}{\kern0pt}\ M\ \isactrlsup c{\isasymTTurnstile}\isactrlsub {\isacharequal}{\kern0pt}\ i{\isachardoublequoteclose}\isanewline
\ \ \ \ \isakeywordONE{have}\isamarkupfalse%
\ {\isadigit{1}}{\isacharcolon}{\kern0pt}\ {\isachardoublequoteopen}{\isasymforall}h\ {\isasymin}\ set\ {\isacharquery}{\kern0pt}hs{\isacharprime}{\kern0pt}{\isachardot}{\kern0pt}\ apply{\isacharunderscore}{\kern0pt}happ\ h\ M\ {\isacharequal}{\kern0pt}\ M{\isachardoublequoteclose}\isanewline
\ \ \ \ \ \ \isakeywordONE{using}\isamarkupfalse%
\ assms\ ind{\isacharunderscore}{\kern0pt}happ{\isacharunderscore}{\kern0pt}seq{\isacharunderscore}{\kern0pt}leq\isactrlsub i{\isacharunderscore}{\kern0pt}le\isactrlsub j\ \isakeywordONE{by}\isamarkupfalse%
\ simp\isanewline
\ \ \ \ \isakeywordONE{have}\isamarkupfalse%
\ {\isadigit{2}}{\isacharcolon}{\kern0pt}\ {\isachardoublequoteopen}{\isasymforall}{\isacharparenleft}{\kern0pt}t\isactrlsub a{\isacharcomma}{\kern0pt}A{\isacharparenright}{\kern0pt}\ {\isasymin}\ set\ {\isacharquery}{\kern0pt}hs{\isacharprime}{\kern0pt}{\isachardot}{\kern0pt}\ {\isasymforall}a{\isasymin}\ set\ A{\isachardot}{\kern0pt}\ M\ \isactrlsup c{\isasymTTurnstile}\isactrlsub {\isacharequal}{\kern0pt}\ precondition\ a{\isachardoublequoteclose}\isanewline
\ \ \ \ \ \ \isakeywordONE{using}\isamarkupfalse%
\ {\isacartoucheopen}{\isasymforall}i\ {\isasymin}\ invs{\isacharunderscore}{\kern0pt}of{\isacharunderscore}{\kern0pt}plan{\isacharunderscore}{\kern0pt}at\ t\isactrlsub j\ {\isasympi}s{\isachardot}{\kern0pt}\ M\ \isactrlsup c{\isasymTTurnstile}\isactrlsub {\isacharequal}{\kern0pt}\ i{\isacartoucheclose}\ ind{\isacharunderscore}{\kern0pt}happ{\isacharunderscore}{\kern0pt}seq{\isacharunderscore}{\kern0pt}invs{\isacharunderscore}{\kern0pt}of{\isacharunderscore}{\kern0pt}plan{\isacharunderscore}{\kern0pt}at{\isacharbrackleft}{\kern0pt}OF\ assms{\isacharbrackright}{\kern0pt}\ \isakeywordONE{by}\isamarkupfalse%
\ auto\isanewline
\ \ \ \ \isakeywordONE{have}\isamarkupfalse%
\ {\isachardoublequoteopen}{\isasymforall}{\isacharparenleft}{\kern0pt}t\isactrlsub a{\isacharcomma}{\kern0pt}A{\isacharparenright}{\kern0pt}\ {\isasymin}\ set\ {\isacharquery}{\kern0pt}hs{\isacharprime}{\kern0pt}{\isachardot}{\kern0pt}\ {\isasymforall}a\ {\isasymin}\ set\ A{\isachardot}{\kern0pt}\ effect\ a\ {\isacharequal}{\kern0pt}\ Effect\ {\isacharbrackleft}{\kern0pt}{\isacharbrackright}{\kern0pt}\ {\isacharbrackleft}{\kern0pt}{\isacharbrackright}{\kern0pt}{\isachardoublequoteclose}\isanewline
\ \ \ \ \ \ \isakeywordONE{using}\isamarkupfalse%
\ assms\ ind{\isacharunderscore}{\kern0pt}happ{\isacharunderscore}{\kern0pt}seq{\isacharunderscore}{\kern0pt}leq\isactrlsub i{\isacharunderscore}{\kern0pt}le\isactrlsub j{\isacharunderscore}{\kern0pt}empty{\isacharunderscore}{\kern0pt}eff\ \isakeywordONE{by}\isamarkupfalse%
\ auto\isanewline
\ \ \ \ \isakeywordONE{then}\isamarkupfalse%
\ \isakeywordONE{have}\isamarkupfalse%
\ {\isadigit{3}}{\isacharcolon}{\kern0pt}\ {\isachardoublequoteopen}{\isasymforall}{\isacharparenleft}{\kern0pt}t\isactrlsub a{\isacharcomma}{\kern0pt}A{\isacharparenright}{\kern0pt}\ {\isasymin}\ set\ {\isacharquery}{\kern0pt}hs{\isacharprime}{\kern0pt}{\isachardot}{\kern0pt}\ {\isasymforall}{\isacharparenleft}{\kern0pt}t\isactrlsub a{\isacharprime}{\kern0pt}{\isacharcomma}{\kern0pt}A{\isacharprime}{\kern0pt}{\isacharparenright}{\kern0pt}\ {\isasymin}\ set\ {\isacharquery}{\kern0pt}hs{\isacharprime}{\kern0pt}{\isachardot}{\kern0pt}\ \isanewline
\ \ \ \ \ \ \ \ \ \ \ \ \ \ \ \ \ \ \ \ {\isasymforall}a\ {\isasymin}\ set\ A{\isachardot}{\kern0pt}\ {\isasymforall}a{\isacharprime}{\kern0pt}\ {\isasymin}\ set\ A{\isacharprime}{\kern0pt}{\isachardot}{\kern0pt}\ acts{\isacharunderscore}{\kern0pt}non{\isacharunderscore}{\kern0pt}intrf\ a\ a{\isacharprime}{\kern0pt}{\isachardoublequoteclose}\isanewline
\ \ \ \ \ \ \isakeywordONE{unfolding}\isamarkupfalse%
\ acts{\isacharunderscore}{\kern0pt}non{\isacharunderscore}{\kern0pt}intrf{\isacharunderscore}{\kern0pt}def\ \isakeywordONE{by}\isamarkupfalse%
\ {\isacharparenleft}{\kern0pt}auto\ simp{\isacharcolon}{\kern0pt}\ Let{\isacharunderscore}{\kern0pt}def{\isacharparenright}{\kern0pt}\isanewline
\ \ \ \ \isakeywordTHREE{show}\isamarkupfalse%
\ {\isachardoublequoteopen}valid{\isacharunderscore}{\kern0pt}happ{\isacharunderscore}{\kern0pt}seq\ M\ {\isacharquery}{\kern0pt}hs{\isacharprime}{\kern0pt}\ M{\isachardoublequoteclose}\isanewline
\ \ \ \ \ \ \isakeywordONE{using}\isamarkupfalse%
\ valid{\isacharunderscore}{\kern0pt}happ{\isacharunderscore}{\kern0pt}seq{\isacharunderscore}{\kern0pt}leq\isactrlsub i{\isacharunderscore}{\kern0pt}le\isactrlsub j{\isacharunderscore}{\kern0pt}aux{\isadigit{2}}{\isacharbrackleft}{\kern0pt}OF\ {\isadigit{1}}\ {\isadigit{2}}\ {\isadigit{3}}{\isacharbrackright}{\kern0pt}\ \isakeywordONE{by}\isamarkupfalse%
\ auto\isanewline
\ \ \isakeywordONE{qed}\isamarkupfalse%
%
\endisatagproof
{\isafoldproof}%
%
\isadelimproof
\isanewline
%
\endisadelimproof
\isanewline
\ \ \isakeywordONE{lemma}\isamarkupfalse%
\ drop{\isacharunderscore}{\kern0pt}take{\isacharunderscore}{\kern0pt}leq\isactrlsub i{\isacharunderscore}{\kern0pt}le{\isacharunderscore}{\kern0pt}leq\isactrlsub j{\isacharunderscore}{\kern0pt}aux{\isacharcolon}{\kern0pt}\isanewline
\ \ \ \ \isakeywordTWO{assumes}\ {\isachardoublequoteopen}t\isactrlsub k\ {\isasymle}\ t\isactrlsub j{\isachardoublequoteclose}\isanewline
\ \ \ \ \ \ \ \ \isakeywordTWO{and}\ {\isachardoublequoteopen}strict{\isacharunderscore}{\kern0pt}sorted\ {\isacharparenleft}{\kern0pt}map\ fst\ hs{\isacharparenright}{\kern0pt}{\isachardoublequoteclose}\isanewline
\ \ \ \ \isakeywordTWO{shows}\ {\isachardoublequoteopen}takeWhile\ {\isacharparenleft}{\kern0pt}le\ t\isactrlsub k{\isacharparenright}{\kern0pt}\ hs\ {\isacharat}{\kern0pt}\ takeWhile\ {\isacharparenleft}{\kern0pt}leq\ t\isactrlsub j{\isacharparenright}{\kern0pt}\ {\isacharparenleft}{\kern0pt}dropWhile\ {\isacharparenleft}{\kern0pt}le\ t\isactrlsub k{\isacharparenright}{\kern0pt}\ hs{\isacharparenright}{\kern0pt}\ {\isacharequal}{\kern0pt}\ takeWhile\ {\isacharparenleft}{\kern0pt}leq\ t\isactrlsub j{\isacharparenright}{\kern0pt}\ hs{\isachardoublequoteclose}\isanewline
%
\isadelimproof
\ \ \ \ %
\endisadelimproof
%
\isatagproof
\isakeywordONE{using}\isamarkupfalse%
\ assms\ \isakeywordONE{by}\isamarkupfalse%
\ {\isacharparenleft}{\kern0pt}induction\ hs{\isacharparenright}{\kern0pt}\ {\isacharparenleft}{\kern0pt}auto\ simp{\isacharcolon}{\kern0pt}\ le{\isacharunderscore}{\kern0pt}def\ leq{\isacharunderscore}{\kern0pt}def{\isacharparenright}{\kern0pt}%
\endisatagproof
{\isafoldproof}%
%
\isadelimproof
%
\endisadelimproof
%
\begin{isamarkuptext}%
We can split up a happening sequence via \isa{\isaconst{dropWhile}} and \isa{\isaconst{takeWhile}}%
\end{isamarkuptext}\isamarkuptrue%
\ \ \isakeywordONE{lemma}\isamarkupfalse%
\ drop{\isacharunderscore}{\kern0pt}take{\isacharunderscore}{\kern0pt}leq\isactrlsub i{\isacharunderscore}{\kern0pt}le{\isacharunderscore}{\kern0pt}leq\isactrlsub j{\isacharunderscore}{\kern0pt}split{\isacharunderscore}{\kern0pt}app{\isacharcolon}{\kern0pt}\isanewline
\ \ \ \ \isakeywordTWO{assumes}\ {\isachardoublequoteopen}t\isactrlsub i\ {\isacharless}{\kern0pt}\ t\isactrlsub k{\isachardoublequoteclose}\ \isakeywordTWO{and}\ {\isachardoublequoteopen}t\isactrlsub k\ {\isasymle}\ t\isactrlsub j{\isachardoublequoteclose}\isanewline
\ \ \ \ \ \ \ \ \isakeywordTWO{and}\ {\isachardoublequoteopen}strict{\isacharunderscore}{\kern0pt}sorted\ {\isacharparenleft}{\kern0pt}map\ fst\ hs{\isacharparenright}{\kern0pt}{\isachardoublequoteclose}\isanewline
\ \ \ \ \ \ \isakeywordTWO{shows}\ {\isachardoublequoteopen}dropWhile\ {\isacharparenleft}{\kern0pt}leq\ t\isactrlsub i{\isacharparenright}{\kern0pt}\ {\isacharparenleft}{\kern0pt}takeWhile\ {\isacharparenleft}{\kern0pt}le\ t\isactrlsub k{\isacharparenright}{\kern0pt}\ hs{\isacharparenright}{\kern0pt}\ {\isacharat}{\kern0pt}\ dropWhile\ {\isacharparenleft}{\kern0pt}le\ t\isactrlsub k{\isacharparenright}{\kern0pt}\ {\isacharparenleft}{\kern0pt}takeWhile\ {\isacharparenleft}{\kern0pt}leq\ t\isactrlsub j{\isacharparenright}{\kern0pt}\ hs{\isacharparenright}{\kern0pt}\ \isanewline
\ \ \ \ \ \ \ \ \ \ \ \ \ \ {\isacharequal}{\kern0pt}\ dropWhile\ {\isacharparenleft}{\kern0pt}leq\ t\isactrlsub i{\isacharparenright}{\kern0pt}\ {\isacharparenleft}{\kern0pt}takeWhile\ {\isacharparenleft}{\kern0pt}leq\ t\isactrlsub j{\isacharparenright}{\kern0pt}\ hs{\isacharparenright}{\kern0pt}{\isachardoublequoteclose}\isanewline
%
\isadelimproof
\ \ \ \ %
\endisadelimproof
%
\isatagproof
\isakeywordONE{using}\isamarkupfalse%
\ {\isacartoucheopen}t\isactrlsub i\ {\isacharless}{\kern0pt}\ t\isactrlsub k{\isacartoucheclose}\ drop{\isacharunderscore}{\kern0pt}take{\isacharunderscore}{\kern0pt}leq\isactrlsub i{\isacharunderscore}{\kern0pt}le{\isacharunderscore}{\kern0pt}leq\isactrlsub j{\isacharunderscore}{\kern0pt}aux{\isacharbrackleft}{\kern0pt}OF\ {\isacartoucheopen}t\isactrlsub k\ {\isasymle}\ t\isactrlsub j{\isacartoucheclose}\ {\isacartoucheopen}strict{\isacharunderscore}{\kern0pt}sorted\ {\isacharparenleft}{\kern0pt}map\ fst\ hs{\isacharparenright}{\kern0pt}{\isacartoucheclose}{\isacharbrackright}{\kern0pt}\isanewline
\ \ \ \ \ \ \ \ \ \ strict{\isacharunderscore}{\kern0pt}sorted{\isacharunderscore}{\kern0pt}drop{\isacharunderscore}{\kern0pt}le{\isacharunderscore}{\kern0pt}take{\isacharunderscore}{\kern0pt}leq{\isacharunderscore}{\kern0pt}comm{\isacharbrackleft}{\kern0pt}OF\ {\isacartoucheopen}t\isactrlsub k\ {\isasymle}\ t\isactrlsub j{\isacartoucheclose}\ {\isacartoucheopen}strict{\isacharunderscore}{\kern0pt}sorted\ {\isacharparenleft}{\kern0pt}map\ fst\ hs{\isacharparenright}{\kern0pt}{\isacartoucheclose}{\isacharbrackright}{\kern0pt}\isanewline
\ \ \ \ \isakeywordONE{by}\isamarkupfalse%
\ {\isacharparenleft}{\kern0pt}induction\ hs{\isacharparenright}{\kern0pt}\ {\isacharparenleft}{\kern0pt}auto\ simp{\isacharcolon}{\kern0pt}\ le{\isacharunderscore}{\kern0pt}def\ leq{\isacharunderscore}{\kern0pt}def{\isacharparenright}{\kern0pt}%
\endisatagproof
{\isafoldproof}%
%
\isadelimproof
%
\endisadelimproof
%
\begin{isamarkuptext}%
Some properties of \isa{\isaconst{valid{\isacharunderscore}{\kern0pt}happ{\isacharunderscore}{\kern0pt}seq}} that are needed later.%
\end{isamarkuptext}\isamarkuptrue%
\ \ \isakeywordONE{lemma}\isamarkupfalse%
\ valid{\isacharunderscore}{\kern0pt}happ{\isacharunderscore}{\kern0pt}seq{\isacharunderscore}{\kern0pt}app{\isacharunderscore}{\kern0pt}iff{\isacharcolon}{\kern0pt}\isanewline
\ \ \ \ {\isachardoublequoteopen}valid{\isacharunderscore}{\kern0pt}happ{\isacharunderscore}{\kern0pt}seq\ M\isactrlsub {\isadigit{0}}\ {\isacharparenleft}{\kern0pt}hs\isactrlsub {\isadigit{1}}{\isacharat}{\kern0pt}hs\isactrlsub {\isadigit{2}}{\isacharparenright}{\kern0pt}\ M\isactrlsub j\ {\isasymlongleftrightarrow}\ {\isacharparenleft}{\kern0pt}{\isasymexists}M\isactrlsub i{\isachardot}{\kern0pt}\ valid{\isacharunderscore}{\kern0pt}happ{\isacharunderscore}{\kern0pt}seq\ M\isactrlsub {\isadigit{0}}\ hs\isactrlsub {\isadigit{1}}\ M\isactrlsub i\ {\isasymand}\ valid{\isacharunderscore}{\kern0pt}happ{\isacharunderscore}{\kern0pt}seq\ M\isactrlsub i\ hs\isactrlsub {\isadigit{2}}\ M\isactrlsub j{\isacharparenright}{\kern0pt}{\isachardoublequoteclose}\isanewline
%
\isadelimproof
\ \ \ \ %
\endisadelimproof
%
\isatagproof
\isakeywordONE{by}\isamarkupfalse%
\ {\isacharparenleft}{\kern0pt}induction\ hs\isactrlsub {\isadigit{1}}\ arbitrary{\isacharcolon}{\kern0pt}\ hs\isactrlsub {\isadigit{2}}\ M\isactrlsub {\isadigit{0}}\ M\isactrlsub j{\isacharparenright}{\kern0pt}\ auto%
\endisatagproof
{\isafoldproof}%
%
\isadelimproof
\isanewline
%
\endisadelimproof
\isanewline
\ \ \isakeywordONE{lemma}\isamarkupfalse%
\ valid{\isacharunderscore}{\kern0pt}happ{\isacharunderscore}{\kern0pt}seq{\isacharunderscore}{\kern0pt}split{\isacharunderscore}{\kern0pt}iff{\isacharcolon}{\kern0pt}\isanewline
\ \ \ \ {\isachardoublequoteopen}valid{\isacharunderscore}{\kern0pt}happ{\isacharunderscore}{\kern0pt}seq\ M\isactrlsub {\isadigit{0}}\ hs\ M\isactrlsub j\ {\isasymlongleftrightarrow}\ \isanewline
\ \ \ \ \ \ {\isacharparenleft}{\kern0pt}{\isasymexists}M\isactrlsub i{\isachardot}{\kern0pt}\ valid{\isacharunderscore}{\kern0pt}happ{\isacharunderscore}{\kern0pt}seq\ M\isactrlsub {\isadigit{0}}\ {\isacharparenleft}{\kern0pt}takeWhile\ f\ hs{\isacharparenright}{\kern0pt}\ M\isactrlsub i\ {\isasymand}\ valid{\isacharunderscore}{\kern0pt}happ{\isacharunderscore}{\kern0pt}seq\ M\isactrlsub i\ {\isacharparenleft}{\kern0pt}dropWhile\ f\ hs{\isacharparenright}{\kern0pt}\ M\isactrlsub j{\isacharparenright}{\kern0pt}{\isachardoublequoteclose}\isanewline
%
\isadelimproof
\ \ \ \ %
\endisadelimproof
%
\isatagproof
\isakeywordONE{using}\isamarkupfalse%
\ valid{\isacharunderscore}{\kern0pt}happ{\isacharunderscore}{\kern0pt}seq{\isacharunderscore}{\kern0pt}app{\isacharunderscore}{\kern0pt}iff\ \isakeywordONE{by}\isamarkupfalse%
\ force%
\endisatagproof
{\isafoldproof}%
%
\isadelimproof
%
\endisadelimproof
%
\begin{isamarkuptext}%
Every happening sequence contains a happening at every happening time point.%
\end{isamarkuptext}\isamarkuptrue%
\ \ \isakeywordONE{lemma}\isamarkupfalse%
\ ind{\isacharunderscore}{\kern0pt}happ{\isacharunderscore}{\kern0pt}seq{\isacharunderscore}{\kern0pt}happ{\isacharunderscore}{\kern0pt}at{\isacharunderscore}{\kern0pt}htps{\isacharcolon}{\kern0pt}\isanewline
\ \ \ \ \isakeywordTWO{assumes}\ {\isachardoublequoteopen}is{\isacharunderscore}{\kern0pt}htp\ {\isasympi}s\ t\isactrlsub i{\isachardoublequoteclose}\isanewline
\ \ \ \ \ \ \isakeywordTWO{and}\ {\isachardoublequoteopen}ind{\isacharunderscore}{\kern0pt}happ{\isacharunderscore}{\kern0pt}seq\ {\isasympi}s\ hs{\isachardoublequoteclose}\isanewline
\ \ \ \ \isakeywordTWO{shows}\ {\isachardoublequoteopen}{\isasymexists}A\isactrlsub i{\isachardot}{\kern0pt}\ {\isacharparenleft}{\kern0pt}t\isactrlsub i{\isacharcomma}{\kern0pt}A\isactrlsub i{\isacharparenright}{\kern0pt}\ {\isasymin}\ set\ hs{\isachardoublequoteclose}\isanewline
%
\isadelimproof
\ \ %
\endisadelimproof
%
\isatagproof
\isakeywordONE{proof}\isamarkupfalse%
\ {\isacharminus}{\kern0pt}\isanewline
\ \ \ \ \isakeywordONE{have}\isamarkupfalse%
\ {\isachardoublequoteopen}{\isacharparenleft}{\kern0pt}{\isasymexists}{\isasympi}{\isachardot}{\kern0pt}\ {\isacharparenleft}{\kern0pt}t\isactrlsub i{\isacharcomma}{\kern0pt}{\isasympi}{\isacharparenright}{\kern0pt}\ {\isasymin}\ set\ {\isasympi}s{\isacharparenright}{\kern0pt}\ {\isasymor}\ {\isacharparenleft}{\kern0pt}{\isasymexists}{\isacharparenleft}{\kern0pt}t\isactrlsub {\isasympi}{\isacharcomma}{\kern0pt}{\isasympi}{\isacharparenright}{\kern0pt}\ {\isasymin}\ durative{\isacharunderscore}{\kern0pt}acts\ {\isasympi}s{\isachardot}{\kern0pt}\ t\isactrlsub i\ {\isacharequal}{\kern0pt}\ t\isactrlsub {\isasympi}\ {\isacharplus}{\kern0pt}\ {\isacharparenleft}{\kern0pt}duration\ {\isasympi}{\isacharparenright}{\kern0pt}{\isacharparenright}{\kern0pt}{\isachardoublequoteclose}\isanewline
\ \ \ \ \ \ \isakeywordONE{using}\isamarkupfalse%
\ {\isacartoucheopen}is{\isacharunderscore}{\kern0pt}htp\ {\isasympi}s\ t\isactrlsub i{\isacartoucheclose}\ \isakeywordONE{unfolding}\isamarkupfalse%
\ is{\isacharunderscore}{\kern0pt}htp{\isacharunderscore}{\kern0pt}def\ \isakeywordONE{by}\isamarkupfalse%
\ simp\isanewline
\ \ \ \ \isakeywordONE{then}\isamarkupfalse%
\ \isakeywordTHREE{show}\isamarkupfalse%
\ {\isacharquery}{\kern0pt}thesis\isanewline
\ \ \ \ \isakeywordONE{proof}\isamarkupfalse%
\isanewline
\ \ \ \ \ \ \isakeywordTHREE{assume}\isamarkupfalse%
\ {\isachardoublequoteopen}{\isasymexists}{\isasympi}{\isachardot}{\kern0pt}\ {\isacharparenleft}{\kern0pt}t\isactrlsub i{\isacharcomma}{\kern0pt}{\isasympi}{\isacharparenright}{\kern0pt}\ {\isasymin}\ set\ {\isasympi}s{\isachardoublequoteclose}\isanewline
\ \ \ \ \ \ \isakeywordONE{then}\isamarkupfalse%
\ \isakeywordTHREE{obtain}\isamarkupfalse%
\ {\isasympi}\ \isakeywordTWO{where}\ {\isachardoublequoteopen}{\isacharparenleft}{\kern0pt}t\isactrlsub i{\isacharcomma}{\kern0pt}{\isasympi}{\isacharparenright}{\kern0pt}\ {\isasymin}\ set\ {\isasympi}s{\isachardoublequoteclose}\isanewline
\ \ \ \ \ \ \ \ \isakeywordONE{by}\isamarkupfalse%
\ auto\isanewline
\ \ \ \ \ \ \isakeywordONE{then}\isamarkupfalse%
\ \isakeywordTHREE{show}\isamarkupfalse%
\ {\isacharquery}{\kern0pt}thesis\isanewline
\ \ \ \ \ \ \ \ \isakeywordONE{using}\isamarkupfalse%
\ {\isacartoucheopen}ind{\isacharunderscore}{\kern0pt}happ{\isacharunderscore}{\kern0pt}seq\ {\isasympi}s\ hs{\isacartoucheclose}\isanewline
\ \ \ \ \ \ \ \ \isakeywordONE{unfolding}\isamarkupfalse%
\ ind{\isacharunderscore}{\kern0pt}happ{\isacharunderscore}{\kern0pt}seq{\isacharunderscore}{\kern0pt}def\ durative{\isacharunderscore}{\kern0pt}acts{\isacharunderscore}{\kern0pt}def\ simple{\isacharunderscore}{\kern0pt}acts{\isacharunderscore}{\kern0pt}def\ \isanewline
\ \ \ \ \ \ \ \ \isakeywordONE{apply}\isamarkupfalse%
\ {\isacharparenleft}{\kern0pt}cases\ {\isasympi}{\isacharparenright}{\kern0pt}\ \isakeywordONE{by}\isamarkupfalse%
\ auto\ blast{\isacharplus}{\kern0pt}\isanewline
\ \ \ \ \isakeywordONE{next}\isamarkupfalse%
\isanewline
\ \ \ \ \ \ \isakeywordTHREE{assume}\isamarkupfalse%
\ {\isachardoublequoteopen}{\isasymexists}{\isacharparenleft}{\kern0pt}t\isactrlsub {\isasympi}{\isacharcomma}{\kern0pt}{\isasympi}{\isacharparenright}{\kern0pt}\ {\isasymin}\ durative{\isacharunderscore}{\kern0pt}acts\ {\isasympi}s{\isachardot}{\kern0pt}\ t\isactrlsub i\ {\isacharequal}{\kern0pt}\ t\isactrlsub {\isasympi}\ {\isacharplus}{\kern0pt}\ {\isacharparenleft}{\kern0pt}duration\ {\isasympi}{\isacharparenright}{\kern0pt}{\isachardoublequoteclose}\isanewline
\ \ \ \ \ \ \isakeywordONE{then}\isamarkupfalse%
\ \isakeywordTHREE{obtain}\isamarkupfalse%
\ t\isactrlsub {\isasympi}\ {\isasympi}\ \isakeywordTWO{where}\ {\isachardoublequoteopen}{\isacharparenleft}{\kern0pt}t\isactrlsub {\isasympi}{\isacharcomma}{\kern0pt}{\isasympi}{\isacharparenright}{\kern0pt}\ {\isasymin}\ durative{\isacharunderscore}{\kern0pt}acts\ {\isasympi}s{\isachardoublequoteclose}\ \isakeywordTWO{and}\ {\isachardoublequoteopen}t\isactrlsub i\ {\isacharequal}{\kern0pt}\ t\isactrlsub {\isasympi}\ {\isacharplus}{\kern0pt}\ {\isacharparenleft}{\kern0pt}duration\ {\isasympi}{\isacharparenright}{\kern0pt}{\isachardoublequoteclose}\isanewline
\ \ \ \ \ \ \ \ \isakeywordONE{by}\isamarkupfalse%
\ auto\isanewline
\ \ \ \ \ \ \isakeywordONE{then}\isamarkupfalse%
\ \isakeywordTHREE{show}\isamarkupfalse%
\ {\isacharquery}{\kern0pt}thesis\isanewline
\ \ \ \ \ \ \ \ \isakeywordONE{using}\isamarkupfalse%
\ {\isacartoucheopen}ind{\isacharunderscore}{\kern0pt}happ{\isacharunderscore}{\kern0pt}seq\ {\isasympi}s\ hs{\isacartoucheclose}\isanewline
\ \ \ \ \ \ \ \ \isakeywordONE{unfolding}\isamarkupfalse%
\ ind{\isacharunderscore}{\kern0pt}happ{\isacharunderscore}{\kern0pt}seq{\isacharunderscore}{\kern0pt}def\ durative{\isacharunderscore}{\kern0pt}acts{\isacharunderscore}{\kern0pt}def\ simple{\isacharunderscore}{\kern0pt}acts{\isacharunderscore}{\kern0pt}def\ \isanewline
\ \ \ \ \ \ \ \ \isakeywordONE{apply}\isamarkupfalse%
\ {\isacharparenleft}{\kern0pt}cases\ {\isasympi}{\isacharparenright}{\kern0pt}\ \isakeywordONE{by}\isamarkupfalse%
\ auto\ force{\isacharplus}{\kern0pt}\isanewline
\ \ \ \ \isakeywordONE{qed}\isamarkupfalse%
\isanewline
\ \ \isakeywordONE{qed}\isamarkupfalse%
%
\endisatagproof
{\isafoldproof}%
%
\isadelimproof
\isanewline
%
\endisadelimproof
\isanewline
\ \ \isakeywordONE{lemma}\isamarkupfalse%
\ distinct{\isacharunderscore}{\kern0pt}single{\isacharunderscore}{\kern0pt}set{\isacharunderscore}{\kern0pt}list{\isacharcolon}{\kern0pt}\ {\isachardoublequoteopen}set\ xs\ {\isacharequal}{\kern0pt}\ {\isacharbraceleft}{\kern0pt}x{\isacharbraceright}{\kern0pt}\ {\isasymLongrightarrow}\ distinct\ xs\ {\isasymLongrightarrow}\ xs\ {\isacharequal}{\kern0pt}\ {\isacharbrackleft}{\kern0pt}x{\isacharbrackright}{\kern0pt}{\isachardoublequoteclose}\isanewline
%
\isadelimproof
\ \ \ \ %
\endisadelimproof
%
\isatagproof
\isakeywordONE{by}\isamarkupfalse%
\ {\isacharparenleft}{\kern0pt}induction\ xs{\isacharparenright}{\kern0pt}\ auto%
\endisatagproof
{\isafoldproof}%
%
\isadelimproof
%
\endisadelimproof
%
\begin{isamarkuptext}%
An induced happening sequence contain exactly one happening for 
  every happening time point.%
\end{isamarkuptext}\isamarkuptrue%
\ \ \isakeywordONE{lemma}\isamarkupfalse%
\ ind{\isacharunderscore}{\kern0pt}happ{\isacharunderscore}{\kern0pt}seq{\isacharunderscore}{\kern0pt}le\isactrlsub i{\isacharunderscore}{\kern0pt}leq\isactrlsub i{\isacharcolon}{\kern0pt}\isanewline
\ \ \ \ \isakeywordTWO{assumes}\ {\isachardoublequoteopen}is{\isacharunderscore}{\kern0pt}htp\ {\isasympi}s\ t\isactrlsub i{\isachardoublequoteclose}\isanewline
\ \ \ \ \ \ \ \ \isakeywordTWO{and}\ {\isachardoublequoteopen}ind{\isacharunderscore}{\kern0pt}happ{\isacharunderscore}{\kern0pt}seq\ {\isasympi}s\ hs{\isachardoublequoteclose}\isanewline
\ \ \ \ \isakeywordTWO{shows}\ {\isachardoublequoteopen}{\isasymexists}A\isactrlsub i{\isachardot}{\kern0pt}\ dropWhile\ {\isacharparenleft}{\kern0pt}le\ t\isactrlsub i{\isacharparenright}{\kern0pt}\ {\isacharparenleft}{\kern0pt}takeWhile\ {\isacharparenleft}{\kern0pt}leq\ t\isactrlsub i{\isacharparenright}{\kern0pt}\ hs{\isacharparenright}{\kern0pt}\ {\isacharequal}{\kern0pt}\ {\isacharbrackleft}{\kern0pt}{\isacharparenleft}{\kern0pt}t\isactrlsub i{\isacharcomma}{\kern0pt}A\isactrlsub i{\isacharparenright}{\kern0pt}{\isacharbrackright}{\kern0pt}{\isachardoublequoteclose}\isanewline
%
\isadelimproof
\ \ %
\endisadelimproof
%
\isatagproof
\isakeywordONE{proof}\isamarkupfalse%
\ {\isacharminus}{\kern0pt}\isanewline
\ \ \ \ \isakeywordONE{have}\isamarkupfalse%
\ {\isachardoublequoteopen}strict{\isacharunderscore}{\kern0pt}sorted\ {\isacharparenleft}{\kern0pt}map\ fst\ hs{\isacharparenright}{\kern0pt}{\isachardoublequoteclose}\isanewline
\ \ \ \ \ \ \isakeywordONE{using}\isamarkupfalse%
\ {\isacartoucheopen}ind{\isacharunderscore}{\kern0pt}happ{\isacharunderscore}{\kern0pt}seq\ {\isasympi}s\ hs{\isacartoucheclose}\isanewline
\ \ \ \ \ \ \isakeywordONE{unfolding}\isamarkupfalse%
\ ind{\isacharunderscore}{\kern0pt}happ{\isacharunderscore}{\kern0pt}seq{\isacharunderscore}{\kern0pt}def\ \isakeywordONE{by}\isamarkupfalse%
\ auto\isanewline
\ \ \ \ \isakeywordONE{let}\isamarkupfalse%
\ {\isacharquery}{\kern0pt}hs{\isacharprime}{\kern0pt}{\isacharequal}{\kern0pt}{\isachardoublequoteopen}dropWhile\ {\isacharparenleft}{\kern0pt}le\ t\isactrlsub i{\isacharparenright}{\kern0pt}\ {\isacharparenleft}{\kern0pt}takeWhile\ {\isacharparenleft}{\kern0pt}leq\ t\isactrlsub i{\isacharparenright}{\kern0pt}\ hs{\isacharparenright}{\kern0pt}{\isachardoublequoteclose}\isanewline
\ \ \ \ \isakeywordONE{have}\isamarkupfalse%
\ {\isachardoublequoteopen}set\ {\isacharquery}{\kern0pt}hs{\isacharprime}{\kern0pt}\ {\isasymsubseteq}\ set\ hs{\isachardoublequoteclose}\isanewline
\ \ \ \ \ \ \isakeywordONE{using}\isamarkupfalse%
\ set{\isacharunderscore}{\kern0pt}dropWhileD\ set{\isacharunderscore}{\kern0pt}takeWhileD\ \isakeywordONE{by}\isamarkupfalse%
\ fastforce\isanewline
\ \ \ \ \isakeywordONE{have}\isamarkupfalse%
\ {\isachardoublequoteopen}distinct\ hs{\isachardoublequoteclose}\isanewline
\ \ \ \ \ \ \isakeywordONE{using}\isamarkupfalse%
\ {\isacartoucheopen}strict{\isacharunderscore}{\kern0pt}sorted\ {\isacharparenleft}{\kern0pt}map\ fst\ hs{\isacharparenright}{\kern0pt}{\isacartoucheclose}\ distinct{\isacharunderscore}{\kern0pt}mapI\ \isakeywordONE{by}\isamarkupfalse%
\ {\isacharparenleft}{\kern0pt}auto\ simp{\isacharcolon}{\kern0pt}\ strict{\isacharunderscore}{\kern0pt}sorted{\isacharunderscore}{\kern0pt}iff{\isacharparenright}{\kern0pt}\isanewline
\ \ \ \ \isakeywordONE{then}\isamarkupfalse%
\ \isakeywordONE{have}\isamarkupfalse%
\ {\isachardoublequoteopen}distinct\ {\isacharquery}{\kern0pt}hs{\isacharprime}{\kern0pt}{\isachardoublequoteclose}\isanewline
\ \ \ \ \ \ \isakeywordONE{using}\isamarkupfalse%
\ distinct{\isacharunderscore}{\kern0pt}dropWhile\ distinct{\isacharunderscore}{\kern0pt}takeWhile\ \isakeywordONE{by}\isamarkupfalse%
\ auto\isanewline
\ \ \ \ \isakeywordTHREE{obtain}\isamarkupfalse%
\ A\isactrlsub i\ \isakeywordTWO{where}\ {\isachardoublequoteopen}{\isacharparenleft}{\kern0pt}t\isactrlsub i{\isacharcomma}{\kern0pt}A\isactrlsub i{\isacharparenright}{\kern0pt}\ {\isasymin}\ set\ hs{\isachardoublequoteclose}\isanewline
\ \ \ \ \ \ \isakeywordONE{using}\isamarkupfalse%
\ ind{\isacharunderscore}{\kern0pt}happ{\isacharunderscore}{\kern0pt}seq{\isacharunderscore}{\kern0pt}happ{\isacharunderscore}{\kern0pt}at{\isacharunderscore}{\kern0pt}htps{\isacharbrackleft}{\kern0pt}OF\ {\isacartoucheopen}is{\isacharunderscore}{\kern0pt}htp\ {\isasympi}s\ t\isactrlsub i{\isacartoucheclose}\ {\isacartoucheopen}ind{\isacharunderscore}{\kern0pt}happ{\isacharunderscore}{\kern0pt}seq\ {\isasympi}s\ hs{\isacartoucheclose}{\isacharbrackright}{\kern0pt}\ \isakeywordONE{by}\isamarkupfalse%
\ auto\isanewline
\ \ \ \ \isakeywordONE{then}\isamarkupfalse%
\ \isakeywordONE{have}\isamarkupfalse%
\ {\isachardoublequoteopen}{\isacharparenleft}{\kern0pt}t\isactrlsub i{\isacharcomma}{\kern0pt}A\isactrlsub i{\isacharparenright}{\kern0pt}\ {\isasymin}\ set\ {\isacharquery}{\kern0pt}hs{\isacharprime}{\kern0pt}{\isachardoublequoteclose}\isanewline
\ \ \ \ \ \ \isakeywordONE{unfolding}\isamarkupfalse%
\ le{\isacharunderscore}{\kern0pt}def\ leq{\isacharunderscore}{\kern0pt}def\isanewline
\ \ \ \ \ \ \isakeywordONE{using}\isamarkupfalse%
\ {\isacartoucheopen}strict{\isacharunderscore}{\kern0pt}sorted\ {\isacharparenleft}{\kern0pt}map\ fst\ hs{\isacharparenright}{\kern0pt}{\isacartoucheclose}\ \isakeywordONE{by}\isamarkupfalse%
\ {\isacharparenleft}{\kern0pt}induction\ hs{\isacharparenright}{\kern0pt}\ auto\isanewline
\ \ \ \ \isakeywordONE{have}\isamarkupfalse%
\ {\isachardoublequoteopen}{\isasymforall}{\isacharparenleft}{\kern0pt}t\isactrlsub a{\isacharcomma}{\kern0pt}A{\isacharparenright}{\kern0pt}\ {\isasymin}\ set\ {\isacharquery}{\kern0pt}hs{\isacharprime}{\kern0pt}{\isachardot}{\kern0pt}\ t\isactrlsub a\ {\isacharequal}{\kern0pt}\ t\isactrlsub i{\isachardoublequoteclose}\isanewline
\ \ \ \ \ \ \isakeywordONE{using}\isamarkupfalse%
\ {\isacartoucheopen}set\ {\isacharquery}{\kern0pt}hs{\isacharprime}{\kern0pt}\ {\isasymsubseteq}\ set\ hs{\isacartoucheclose}\ strict{\isacharunderscore}{\kern0pt}sorted{\isacharunderscore}{\kern0pt}drop{\isacharunderscore}{\kern0pt}le{\isacharunderscore}{\kern0pt}take{\isacharunderscore}{\kern0pt}leq{\isacharbrackleft}{\kern0pt}OF\ {\isacartoucheopen}strict{\isacharunderscore}{\kern0pt}sorted\ {\isacharparenleft}{\kern0pt}map\ fst\ hs{\isacharparenright}{\kern0pt}{\isacartoucheclose}{\isacharbrackright}{\kern0pt}\isanewline
\ \ \ \ \ \ \isakeywordONE{by}\isamarkupfalse%
\ auto\isanewline
\ \ \ \ \isakeywordONE{have}\isamarkupfalse%
\ {\isachardoublequoteopen}{\isasymforall}{\isacharparenleft}{\kern0pt}t\isactrlsub {\isadigit{1}}{\isacharcomma}{\kern0pt}A\isactrlsub {\isadigit{1}}{\isacharparenright}{\kern0pt}\ {\isasymin}\ set\ hs{\isachardot}{\kern0pt}\ {\isasymforall}{\isacharparenleft}{\kern0pt}t\isactrlsub {\isadigit{2}}{\isacharcomma}{\kern0pt}A\isactrlsub {\isadigit{2}}{\isacharparenright}{\kern0pt}\ {\isasymin}\ set\ hs{\isachardot}{\kern0pt}\ t\isactrlsub {\isadigit{1}}\ {\isacharequal}{\kern0pt}\ t\isactrlsub {\isadigit{2}}\ {\isasymlongrightarrow}\ A\isactrlsub {\isadigit{1}}\ {\isacharequal}{\kern0pt}\ A\isactrlsub {\isadigit{2}}{\isachardoublequoteclose}\isanewline
\ \ \ \ \ \ \isakeywordONE{using}\isamarkupfalse%
\ {\isacartoucheopen}strict{\isacharunderscore}{\kern0pt}sorted\ {\isacharparenleft}{\kern0pt}map\ fst\ hs{\isacharparenright}{\kern0pt}{\isacartoucheclose}\ \isakeywordONE{by}\isamarkupfalse%
\ {\isacharparenleft}{\kern0pt}auto\ simp{\isacharcolon}{\kern0pt}\ strict{\isacharunderscore}{\kern0pt}sorted{\isacharunderscore}{\kern0pt}iff\ distinct{\isacharunderscore}{\kern0pt}map{\isacharunderscore}{\kern0pt}fstD{\isacharparenright}{\kern0pt}\isanewline
\ \ \ \ \isakeywordONE{then}\isamarkupfalse%
\ \isakeywordONE{have}\isamarkupfalse%
\ {\isachardoublequoteopen}{\isasymforall}{\isacharparenleft}{\kern0pt}t\isactrlsub {\isadigit{1}}{\isacharcomma}{\kern0pt}A\isactrlsub {\isadigit{1}}{\isacharparenright}{\kern0pt}\ {\isasymin}\ set\ {\isacharquery}{\kern0pt}hs{\isacharprime}{\kern0pt}{\isachardot}{\kern0pt}\ {\isasymforall}{\isacharparenleft}{\kern0pt}t\isactrlsub {\isadigit{2}}{\isacharcomma}{\kern0pt}A\isactrlsub {\isadigit{2}}{\isacharparenright}{\kern0pt}\ {\isasymin}\ set\ {\isacharquery}{\kern0pt}hs{\isacharprime}{\kern0pt}{\isachardot}{\kern0pt}\ t\isactrlsub {\isadigit{1}}\ {\isacharequal}{\kern0pt}\ t\isactrlsub {\isadigit{2}}\ {\isasymlongrightarrow}\ A\isactrlsub {\isadigit{1}}\ {\isacharequal}{\kern0pt}\ A\isactrlsub {\isadigit{2}}{\isachardoublequoteclose}\isanewline
\ \ \ \ \ \ \isakeywordONE{using}\isamarkupfalse%
\ {\isacartoucheopen}set\ {\isacharquery}{\kern0pt}hs{\isacharprime}{\kern0pt}\ {\isasymsubseteq}\ set\ hs{\isacartoucheclose}\ \isakeywordONE{by}\isamarkupfalse%
\ blast\isanewline
\ \ \ \ \isakeywordONE{then}\isamarkupfalse%
\ \isakeywordONE{have}\isamarkupfalse%
\ {\isachardoublequoteopen}{\isasymforall}{\isacharparenleft}{\kern0pt}t\isactrlsub a{\isacharcomma}{\kern0pt}as{\isacharprime}{\kern0pt}{\isacharparenright}{\kern0pt}\ {\isasymin}\ set\ {\isacharquery}{\kern0pt}hs{\isacharprime}{\kern0pt}{\isachardot}{\kern0pt}\ t\isactrlsub a\ {\isacharequal}{\kern0pt}\ t\isactrlsub i\ {\isasymand}\ as{\isacharprime}{\kern0pt}\ {\isacharequal}{\kern0pt}\ A\isactrlsub i{\isachardoublequoteclose}\isanewline
\ \ \ \ \ \ \isakeywordONE{using}\isamarkupfalse%
\ {\isacartoucheopen}{\isacharparenleft}{\kern0pt}t\isactrlsub i{\isacharcomma}{\kern0pt}A\isactrlsub i{\isacharparenright}{\kern0pt}\ {\isasymin}\ set\ {\isacharquery}{\kern0pt}hs{\isacharprime}{\kern0pt}{\isacartoucheclose}\ {\isacartoucheopen}{\isasymforall}{\isacharparenleft}{\kern0pt}t\isactrlsub a{\isacharcomma}{\kern0pt}as{\isacharparenright}{\kern0pt}\ {\isasymin}\ set\ {\isacharquery}{\kern0pt}hs{\isacharprime}{\kern0pt}{\isachardot}{\kern0pt}\ t\isactrlsub a\ {\isacharequal}{\kern0pt}\ t\isactrlsub i{\isacartoucheclose}\ \isakeywordONE{by}\isamarkupfalse%
\ auto\ force\isanewline
\ \ \ \ \isakeywordONE{then}\isamarkupfalse%
\ \isakeywordONE{have}\isamarkupfalse%
\ {\isachardoublequoteopen}set\ {\isacharquery}{\kern0pt}hs{\isacharprime}{\kern0pt}\ {\isacharequal}{\kern0pt}\ {\isacharbraceleft}{\kern0pt}{\isacharparenleft}{\kern0pt}t\isactrlsub i{\isacharcomma}{\kern0pt}A\isactrlsub i{\isacharparenright}{\kern0pt}{\isacharbraceright}{\kern0pt}{\isachardoublequoteclose}\isanewline
\ \ \ \ \ \ \isakeywordONE{using}\isamarkupfalse%
\ {\isacartoucheopen}{\isacharparenleft}{\kern0pt}t\isactrlsub i{\isacharcomma}{\kern0pt}A\isactrlsub i{\isacharparenright}{\kern0pt}\ {\isasymin}\ set\ {\isacharquery}{\kern0pt}hs{\isacharprime}{\kern0pt}{\isacartoucheclose}\ \isakeywordONE{by}\isamarkupfalse%
\ blast\isanewline
\ \ \ \ \isakeywordONE{then}\isamarkupfalse%
\ \isakeywordTHREE{show}\isamarkupfalse%
\ {\isacharquery}{\kern0pt}thesis\isanewline
\ \ \ \ \ \ \isakeywordONE{using}\isamarkupfalse%
\ distinct{\isacharunderscore}{\kern0pt}single{\isacharunderscore}{\kern0pt}set{\isacharunderscore}{\kern0pt}list{\isacharbrackleft}{\kern0pt}OF\ {\isacartoucheopen}set\ {\isacharquery}{\kern0pt}hs{\isacharprime}{\kern0pt}\ {\isacharequal}{\kern0pt}\ {\isacharbraceleft}{\kern0pt}{\isacharparenleft}{\kern0pt}t\isactrlsub i{\isacharcomma}{\kern0pt}A\isactrlsub i{\isacharparenright}{\kern0pt}{\isacharbraceright}{\kern0pt}{\isacartoucheclose}\ {\isacartoucheopen}distinct\ {\isacharquery}{\kern0pt}hs{\isacharprime}{\kern0pt}{\isacartoucheclose}{\isacharbrackright}{\kern0pt}\ \isakeywordONE{by}\isamarkupfalse%
\ simp\isanewline
\ \ \isakeywordONE{qed}\isamarkupfalse%
%
\endisatagproof
{\isafoldproof}%
%
\isadelimproof
%
\endisadelimproof
%
\begin{isamarkuptext}%
\isa{\isaconst{invs{\isacharunderscore}{\kern0pt}of{\isacharunderscore}{\kern0pt}plan{\isacharunderscore}{\kern0pt}at}} is valid for a single happening time point iff any 
  induced happening sequence is valid from the previous happening time point to the 
  happening time point.%
\end{isamarkuptext}\isamarkuptrue%
\ \ \isakeywordONE{lemma}\isamarkupfalse%
\ {\isacharparenleft}{\kern0pt}\isakeywordTWO{in}\ wf{\isacharunderscore}{\kern0pt}ast{\isacharunderscore}{\kern0pt}problem{\isacharparenright}{\kern0pt}\ valid{\isacharunderscore}{\kern0pt}happ{\isacharunderscore}{\kern0pt}seq{\isacharunderscore}{\kern0pt}valid{\isacharunderscore}{\kern0pt}state{\isacharunderscore}{\kern0pt}seq{\isacharunderscore}{\kern0pt}consec{\isacharunderscore}{\kern0pt}htps{\isacharcolon}{\kern0pt}\isanewline
\ \ \ \ \isakeywordTWO{assumes}\ {\isachardoublequoteopen}wf{\isacharunderscore}{\kern0pt}plan\ {\isasympi}s{\isachardoublequoteclose}\isanewline
\ \ \ \ \ \ \ \ \isakeywordTWO{and}\ {\isachardoublequoteopen}consec{\isacharunderscore}{\kern0pt}htps\ {\isasympi}s\ t\isactrlsub i\ t\isactrlsub j{\isachardoublequoteclose}\isanewline
\ \ \ \ \ \ \ \ \isakeywordTWO{and}\ {\isachardoublequoteopen}ind{\isacharunderscore}{\kern0pt}happ{\isacharunderscore}{\kern0pt}seq\ {\isasympi}s\ hs{\isachardoublequoteclose}\isanewline
\ \ \ \ \isakeywordTWO{shows}\ {\isachardoublequoteopen}valid{\isacharunderscore}{\kern0pt}happ{\isacharunderscore}{\kern0pt}seq\ M\ {\isacharparenleft}{\kern0pt}dropWhile\ {\isacharparenleft}{\kern0pt}leq\ t\isactrlsub i{\isacharparenright}{\kern0pt}\ {\isacharparenleft}{\kern0pt}takeWhile\ {\isacharparenleft}{\kern0pt}leq\ t\isactrlsub j{\isacharparenright}{\kern0pt}\ hs{\isacharparenright}{\kern0pt}{\isacharparenright}{\kern0pt}\ M{\isacharprime}{\kern0pt}\ \isanewline
\ \ \ \ \ \ \ \ \ \ \ \ {\isasymlongleftrightarrow}\ valid{\isacharunderscore}{\kern0pt}state{\isacharunderscore}{\kern0pt}seq\ M\ {\isacharbrackleft}{\kern0pt}t\isactrlsub j{\isacharbrackright}{\kern0pt}\ {\isasympi}s\ M{\isacharprime}{\kern0pt}{\isachardoublequoteclose}\isanewline
%
\isadelimproof
\ \ %
\endisadelimproof
%
\isatagproof
\isakeywordONE{proof}\isamarkupfalse%
\isanewline
\ \ \ \ \isakeywordONE{let}\isamarkupfalse%
\ {\isacharquery}{\kern0pt}hs{\isacharprime}{\kern0pt}{\isacharequal}{\kern0pt}{\isachardoublequoteopen}dropWhile\ {\isacharparenleft}{\kern0pt}leq\ t\isactrlsub i{\isacharparenright}{\kern0pt}\ {\isacharparenleft}{\kern0pt}takeWhile\ {\isacharparenleft}{\kern0pt}leq\ t\isactrlsub j{\isacharparenright}{\kern0pt}\ hs{\isacharparenright}{\kern0pt}{\isachardoublequoteclose}\isanewline
\ \ \ \ \isakeywordTHREE{assume}\isamarkupfalse%
\ {\isachardoublequoteopen}valid{\isacharunderscore}{\kern0pt}happ{\isacharunderscore}{\kern0pt}seq\ M\ {\isacharquery}{\kern0pt}hs{\isacharprime}{\kern0pt}\ M{\isacharprime}{\kern0pt}{\isachardoublequoteclose}\isanewline
\isanewline
\ \ \ \ \isakeywordONE{have}\isamarkupfalse%
\ {\isachardoublequoteopen}is{\isacharunderscore}{\kern0pt}htp\ {\isasympi}s\ t\isactrlsub j{\isachardoublequoteclose}\ \isakeywordTWO{and}\ {\isachardoublequoteopen}t\isactrlsub i\ {\isacharless}{\kern0pt}\ t\isactrlsub j{\isachardoublequoteclose}\ \isakeywordTWO{and}\ {\isachardoublequoteopen}t\isactrlsub j\ {\isasymle}\ t\isactrlsub j{\isachardoublequoteclose}\isanewline
\ \ \ \ \ \ \isakeywordONE{using}\isamarkupfalse%
\ {\isacartoucheopen}consec{\isacharunderscore}{\kern0pt}htps\ {\isasympi}s\ t\isactrlsub i\ t\isactrlsub j{\isacartoucheclose}\ \isakeywordONE{unfolding}\isamarkupfalse%
\ consec{\isacharunderscore}{\kern0pt}htps{\isacharunderscore}{\kern0pt}def\ \isakeywordONE{by}\isamarkupfalse%
\ auto\isanewline
\ \ \ \ \isakeywordONE{have}\isamarkupfalse%
\ {\isachardoublequoteopen}strict{\isacharunderscore}{\kern0pt}sorted\ {\isacharparenleft}{\kern0pt}map\ fst\ hs{\isacharparenright}{\kern0pt}{\isachardoublequoteclose}\isanewline
\ \ \ \ \ \ \isakeywordONE{using}\isamarkupfalse%
\ {\isacartoucheopen}ind{\isacharunderscore}{\kern0pt}happ{\isacharunderscore}{\kern0pt}seq\ {\isasympi}s\ hs{\isacartoucheclose}\ \isakeywordONE{unfolding}\isamarkupfalse%
\ ind{\isacharunderscore}{\kern0pt}happ{\isacharunderscore}{\kern0pt}seq{\isacharunderscore}{\kern0pt}def\ \isakeywordONE{by}\isamarkupfalse%
\ auto\isanewline
\ \ \ \ \isakeywordONE{then}\isamarkupfalse%
\ \isakeywordONE{have}\isamarkupfalse%
\ {\isachardoublequoteopen}{\isacharquery}{\kern0pt}hs{\isacharprime}{\kern0pt}\ \isanewline
\ \ \ \ \ \ \ \ \ \ \ {\isacharequal}{\kern0pt}\ dropWhile\ {\isacharparenleft}{\kern0pt}leq\ t\isactrlsub i{\isacharparenright}{\kern0pt}\ {\isacharparenleft}{\kern0pt}takeWhile\ {\isacharparenleft}{\kern0pt}le\ t\isactrlsub j{\isacharparenright}{\kern0pt}\ hs{\isacharparenright}{\kern0pt}\ {\isacharat}{\kern0pt}\ dropWhile\ {\isacharparenleft}{\kern0pt}le\ t\isactrlsub j{\isacharparenright}{\kern0pt}\ {\isacharparenleft}{\kern0pt}takeWhile\ {\isacharparenleft}{\kern0pt}leq\ t\isactrlsub j{\isacharparenright}{\kern0pt}\ hs{\isacharparenright}{\kern0pt}{\isachardoublequoteclose}\isanewline
\ \ \ \ \ \ \isakeywordONE{using}\isamarkupfalse%
\ drop{\isacharunderscore}{\kern0pt}take{\isacharunderscore}{\kern0pt}leq\isactrlsub i{\isacharunderscore}{\kern0pt}le{\isacharunderscore}{\kern0pt}leq\isactrlsub j{\isacharunderscore}{\kern0pt}split{\isacharunderscore}{\kern0pt}app{\isacharbrackleft}{\kern0pt}OF\ {\isacartoucheopen}t\isactrlsub i\ {\isacharless}{\kern0pt}\ t\isactrlsub j{\isacartoucheclose}\ {\isacartoucheopen}t\isactrlsub j\ {\isasymle}\ t\isactrlsub j{\isacartoucheclose}{\isacharbrackright}{\kern0pt}\ \isakeywordONE{by}\isamarkupfalse%
\ fastforce\isanewline
\ \ \ \ \isakeywordONE{then}\isamarkupfalse%
\ \isakeywordTHREE{obtain}\isamarkupfalse%
\ M{\isacharprime}{\kern0pt}{\isacharprime}{\kern0pt}\ \isakeywordTWO{where}\ {\isachardoublequoteopen}valid{\isacharunderscore}{\kern0pt}happ{\isacharunderscore}{\kern0pt}seq\ M\ {\isacharparenleft}{\kern0pt}dropWhile\ {\isacharparenleft}{\kern0pt}leq\ t\isactrlsub i{\isacharparenright}{\kern0pt}\ {\isacharparenleft}{\kern0pt}takeWhile\ {\isacharparenleft}{\kern0pt}le\ t\isactrlsub j{\isacharparenright}{\kern0pt}\ hs{\isacharparenright}{\kern0pt}{\isacharparenright}{\kern0pt}\ M{\isacharprime}{\kern0pt}{\isacharprime}{\kern0pt}{\isachardoublequoteclose}\ \isanewline
\ \ \ \ \ \ \ \ \ \ \ \ \ \ \ \ \ \isakeywordTWO{and}\ {\isachardoublequoteopen}valid{\isacharunderscore}{\kern0pt}happ{\isacharunderscore}{\kern0pt}seq\ M{\isacharprime}{\kern0pt}{\isacharprime}{\kern0pt}\ {\isacharparenleft}{\kern0pt}dropWhile\ {\isacharparenleft}{\kern0pt}le\ t\isactrlsub j{\isacharparenright}{\kern0pt}\ {\isacharparenleft}{\kern0pt}takeWhile\ {\isacharparenleft}{\kern0pt}leq\ t\isactrlsub j{\isacharparenright}{\kern0pt}\ hs{\isacharparenright}{\kern0pt}{\isacharparenright}{\kern0pt}\ M{\isacharprime}{\kern0pt}{\isachardoublequoteclose}\isanewline
\ \ \ \ \ \ \isakeywordONE{using}\isamarkupfalse%
\ valid{\isacharunderscore}{\kern0pt}happ{\isacharunderscore}{\kern0pt}seq{\isacharunderscore}{\kern0pt}app{\isacharunderscore}{\kern0pt}iff\ {\isacartoucheopen}valid{\isacharunderscore}{\kern0pt}happ{\isacharunderscore}{\kern0pt}seq\ M\ {\isacharquery}{\kern0pt}hs{\isacharprime}{\kern0pt}\ M{\isacharprime}{\kern0pt}{\isacartoucheclose}\ \isakeywordONE{by}\isamarkupfalse%
\ auto\isanewline
\ \ \ \ \isakeywordONE{then}\isamarkupfalse%
\ \isakeywordONE{have}\isamarkupfalse%
\ {\isachardoublequoteopen}M\ {\isacharequal}{\kern0pt}\ M{\isacharprime}{\kern0pt}{\isacharprime}{\kern0pt}{\isachardoublequoteclose}\ \isanewline
\ \ \ \ \ \ \isakeywordONE{using}\isamarkupfalse%
\ assms\ valid{\isacharunderscore}{\kern0pt}happ{\isacharunderscore}{\kern0pt}seq{\isacharunderscore}{\kern0pt}M\isactrlsub j{\isacharunderscore}{\kern0pt}is{\isacharunderscore}{\kern0pt}M\isactrlsub i{\isacharbrackleft}{\kern0pt}OF\ ind{\isacharunderscore}{\kern0pt}happ{\isacharunderscore}{\kern0pt}seq{\isacharunderscore}{\kern0pt}leq\isactrlsub i{\isacharunderscore}{\kern0pt}le\isactrlsub j{\isacharbrackright}{\kern0pt}\ \isakeywordONE{by}\isamarkupfalse%
\ auto\isanewline
\isanewline
\ \ \ \ \isakeywordONE{have}\isamarkupfalse%
\ {\isachardoublequoteopen}is{\isacharunderscore}{\kern0pt}htp\ {\isasympi}s\ t\isactrlsub j{\isachardoublequoteclose}\isanewline
\ \ \ \ \ \ \isakeywordONE{using}\isamarkupfalse%
\ assms\ \isakeywordONE{unfolding}\isamarkupfalse%
\ consec{\isacharunderscore}{\kern0pt}htps{\isacharunderscore}{\kern0pt}def\ \isakeywordONE{by}\isamarkupfalse%
\ auto\isanewline
\ \ \ \ \isakeywordTHREE{obtain}\isamarkupfalse%
\ A\isactrlsub j\ \isakeywordTWO{where}\ {\isachardoublequoteopen}dropWhile\ {\isacharparenleft}{\kern0pt}le\ t\isactrlsub j{\isacharparenright}{\kern0pt}\ {\isacharparenleft}{\kern0pt}takeWhile\ {\isacharparenleft}{\kern0pt}leq\ t\isactrlsub j{\isacharparenright}{\kern0pt}\ hs{\isacharparenright}{\kern0pt}\ {\isacharequal}{\kern0pt}\ {\isacharbrackleft}{\kern0pt}{\isacharparenleft}{\kern0pt}t\isactrlsub j{\isacharcomma}{\kern0pt}A\isactrlsub j{\isacharparenright}{\kern0pt}{\isacharbrackright}{\kern0pt}{\isachardoublequoteclose}\isanewline
\ \ \ \ \ \ \isakeywordONE{using}\isamarkupfalse%
\ assms\ ind{\isacharunderscore}{\kern0pt}happ{\isacharunderscore}{\kern0pt}seq{\isacharunderscore}{\kern0pt}le\isactrlsub i{\isacharunderscore}{\kern0pt}leq\isactrlsub i{\isacharbrackleft}{\kern0pt}OF\ {\isacartoucheopen}is{\isacharunderscore}{\kern0pt}htp\ {\isasympi}s\ t\isactrlsub j{\isacartoucheclose}{\isacharbrackright}{\kern0pt}\ \isakeywordONE{by}\isamarkupfalse%
\ auto\isanewline
\ \ \ \ \isakeywordONE{then}\isamarkupfalse%
\ \isakeywordONE{have}\isamarkupfalse%
\ {\isachardoublequoteopen}{\isacharparenleft}{\kern0pt}t\isactrlsub j{\isacharcomma}{\kern0pt}A\isactrlsub j{\isacharparenright}{\kern0pt}\ {\isasymin}\ set\ hs{\isachardoublequoteclose}\isanewline
\ \ \ \ \ \ \isakeywordONE{by}\isamarkupfalse%
\ {\isacharparenleft}{\kern0pt}metis\ list{\isachardot}{\kern0pt}set{\isacharunderscore}{\kern0pt}intros{\isacharparenleft}{\kern0pt}{\isadigit{1}}{\isacharparenright}{\kern0pt}\ set{\isacharunderscore}{\kern0pt}dropWhileD\ set{\isacharunderscore}{\kern0pt}takeWhileD{\isacharparenright}{\kern0pt}\isanewline
\ \ \ \ \isakeywordONE{then}\isamarkupfalse%
\ \isakeywordONE{have}\isamarkupfalse%
\ {\isachardoublequoteopen}set\ A\isactrlsub j\ {\isacharequal}{\kern0pt}\ acts{\isacharunderscore}{\kern0pt}of{\isacharunderscore}{\kern0pt}plan{\isacharunderscore}{\kern0pt}at\ t\isactrlsub j\ {\isasympi}s{\isachardoublequoteclose}\isanewline
\ \ \ \ \ \ \isakeywordONE{using}\isamarkupfalse%
\ assms\ {\isacartoucheopen}is{\isacharunderscore}{\kern0pt}htp\ {\isasympi}s\ t\isactrlsub j{\isacartoucheclose}\ ind{\isacharunderscore}{\kern0pt}happ{\isacharunderscore}{\kern0pt}seq{\isacharunderscore}{\kern0pt}htp{\isacharunderscore}{\kern0pt}acts{\isacharunderscore}{\kern0pt}of{\isacharunderscore}{\kern0pt}plan{\isacharunderscore}{\kern0pt}at\ \isakeywordONE{by}\isamarkupfalse%
\ auto\isanewline
\ \ \ \ \isakeywordONE{let}\isamarkupfalse%
\ {\isacharquery}{\kern0pt}as{\isacharequal}{\kern0pt}{\isachardoublequoteopen}acts{\isacharunderscore}{\kern0pt}of{\isacharunderscore}{\kern0pt}plan{\isacharunderscore}{\kern0pt}at\ t\isactrlsub j\ {\isasympi}s{\isachardoublequoteclose}\isanewline
\ \ \ \ \isakeywordONE{have}\isamarkupfalse%
\ {\isadigit{1}}{\isacharcolon}{\kern0pt}\ {\isachardoublequoteopen}{\isasymforall}i\ {\isasymin}\ invs{\isacharunderscore}{\kern0pt}of{\isacharunderscore}{\kern0pt}plan{\isacharunderscore}{\kern0pt}at\ t\isactrlsub j\ {\isasympi}s{\isachardot}{\kern0pt}\ M\ \isactrlsup c{\isasymTTurnstile}\isactrlsub {\isacharequal}{\kern0pt}\ i{\isachardoublequoteclose}\isanewline
\ \ \ \ \ \isakeywordTWO{and}\ {\isadigit{2}}{\isacharcolon}{\kern0pt}\ {\isachardoublequoteopen}{\isasymforall}a{\isasymin}\ {\isacharquery}{\kern0pt}as{\isachardot}{\kern0pt}\ M\ \isactrlsup c{\isasymTTurnstile}\isactrlsub {\isacharequal}{\kern0pt}\ precondition\ a{\isachardoublequoteclose}\ \isanewline
\ \ \ \ \ \isakeywordTWO{and}\ {\isadigit{3}}{\isacharcolon}{\kern0pt}\ {\isachardoublequoteopen}{\isasymforall}a\ {\isasymin}\ {\isacharquery}{\kern0pt}as{\isachardot}{\kern0pt}\ {\isasymforall}b\ {\isasymin}\ {\isacharquery}{\kern0pt}as{\isachardot}{\kern0pt}\ a\ {\isasymnoteq}\ b\ {\isasymlongrightarrow}\ acts{\isacharunderscore}{\kern0pt}non{\isacharunderscore}{\kern0pt}intrf\ a\ b{\isachardoublequoteclose}\isanewline
\ \ \ \ \ \isakeywordTWO{and}\ {\isadigit{4}}{\isacharcolon}{\kern0pt}\ {\isachardoublequoteopen}apply{\isacharunderscore}{\kern0pt}eff\ {\isacharquery}{\kern0pt}as\ M\ {\isacharequal}{\kern0pt}\ M{\isacharprime}{\kern0pt}{\isachardoublequoteclose}\isanewline
\ \ \ \ \ \ \isakeywordONE{using}\isamarkupfalse%
\ valid{\isacharunderscore}{\kern0pt}happ{\isacharunderscore}{\kern0pt}seq{\isacharunderscore}{\kern0pt}leq\isactrlsub i{\isacharunderscore}{\kern0pt}le\isactrlsub j{\isacharunderscore}{\kern0pt}invs{\isacharunderscore}{\kern0pt}of{\isacharunderscore}{\kern0pt}plan{\isacharunderscore}{\kern0pt}at{\isacharbrackleft}{\kern0pt}OF\ assms{\isacharbrackright}{\kern0pt}\isanewline
\ \ \ \ \ \ \ \ \ \ \ \ {\isacartoucheopen}valid{\isacharunderscore}{\kern0pt}happ{\isacharunderscore}{\kern0pt}seq\ M\ {\isacharparenleft}{\kern0pt}dropWhile\ {\isacharparenleft}{\kern0pt}leq\ t\isactrlsub i{\isacharparenright}{\kern0pt}\ {\isacharparenleft}{\kern0pt}takeWhile\ {\isacharparenleft}{\kern0pt}le\ t\isactrlsub j{\isacharparenright}{\kern0pt}\ hs{\isacharparenright}{\kern0pt}{\isacharparenright}{\kern0pt}\ M{\isacharprime}{\kern0pt}{\isacharprime}{\kern0pt}{\isacartoucheclose}\isanewline
\ \ \ \ \ \ \ \ \ \ \ \ {\isacartoucheopen}valid{\isacharunderscore}{\kern0pt}happ{\isacharunderscore}{\kern0pt}seq\ M{\isacharprime}{\kern0pt}{\isacharprime}{\kern0pt}\ {\isacharparenleft}{\kern0pt}dropWhile\ {\isacharparenleft}{\kern0pt}le\ t\isactrlsub j{\isacharparenright}{\kern0pt}\ {\isacharparenleft}{\kern0pt}takeWhile\ {\isacharparenleft}{\kern0pt}leq\ t\isactrlsub j{\isacharparenright}{\kern0pt}\ hs{\isacharparenright}{\kern0pt}{\isacharparenright}{\kern0pt}\ M{\isacharprime}{\kern0pt}{\isacartoucheclose}\ \isanewline
\ \ \ \ \ \ \ \ \ \ \ \ {\isacartoucheopen}M\ {\isacharequal}{\kern0pt}\ M{\isacharprime}{\kern0pt}{\isacharprime}{\kern0pt}{\isacartoucheclose}\isanewline
\ \ \ \ \ \ \ \ \ \ \ \ {\isacartoucheopen}dropWhile\ {\isacharparenleft}{\kern0pt}le\ t\isactrlsub j{\isacharparenright}{\kern0pt}\ {\isacharparenleft}{\kern0pt}takeWhile\ {\isacharparenleft}{\kern0pt}leq\ t\isactrlsub j{\isacharparenright}{\kern0pt}\ hs{\isacharparenright}{\kern0pt}\ {\isacharequal}{\kern0pt}\ {\isacharbrackleft}{\kern0pt}{\isacharparenleft}{\kern0pt}t\isactrlsub j{\isacharcomma}{\kern0pt}A\isactrlsub j{\isacharparenright}{\kern0pt}{\isacharbrackright}{\kern0pt}{\isacartoucheclose}\ \isanewline
\ \ \ \ \ \ \ \ \ \ \ \ {\isacartoucheopen}set\ A\isactrlsub j\ {\isacharequal}{\kern0pt}\ acts{\isacharunderscore}{\kern0pt}of{\isacharunderscore}{\kern0pt}plan{\isacharunderscore}{\kern0pt}at\ t\isactrlsub j\ {\isasympi}s{\isacartoucheclose}\isanewline
\ \ \ \ \ \ \isakeywordONE{by}\isamarkupfalse%
\ {\isacharparenleft}{\kern0pt}auto\ simp{\isacharcolon}{\kern0pt}\ pairwise{\isacharunderscore}{\kern0pt}def{\isacharparenright}{\kern0pt}\isanewline
\ \ \ \ \isakeywordONE{then}\isamarkupfalse%
\ \isakeywordTHREE{show}\isamarkupfalse%
\ {\isachardoublequoteopen}valid{\isacharunderscore}{\kern0pt}state{\isacharunderscore}{\kern0pt}seq\ M\ {\isacharbrackleft}{\kern0pt}t\isactrlsub j{\isacharbrackright}{\kern0pt}\ {\isasympi}s\ M{\isacharprime}{\kern0pt}{\isachardoublequoteclose}\isanewline
\ \ \ \ \ \ \isakeywordONE{by}\isamarkupfalse%
\ auto\isanewline
\ \ \isakeywordONE{next}\isamarkupfalse%
\isanewline
\ \ \ \ \isakeywordONE{let}\isamarkupfalse%
\ {\isacharquery}{\kern0pt}hs{\isacharprime}{\kern0pt}{\isacharequal}{\kern0pt}{\isachardoublequoteopen}dropWhile\ {\isacharparenleft}{\kern0pt}leq\ t\isactrlsub i{\isacharparenright}{\kern0pt}\ {\isacharparenleft}{\kern0pt}takeWhile\ {\isacharparenleft}{\kern0pt}leq\ t\isactrlsub j{\isacharparenright}{\kern0pt}\ hs{\isacharparenright}{\kern0pt}{\isachardoublequoteclose}\isanewline
\ \ \ \ \isakeywordTHREE{assume}\isamarkupfalse%
\ {\isachardoublequoteopen}valid{\isacharunderscore}{\kern0pt}state{\isacharunderscore}{\kern0pt}seq\ M\ {\isacharbrackleft}{\kern0pt}t\isactrlsub j{\isacharbrackright}{\kern0pt}\ {\isasympi}s\ M{\isacharprime}{\kern0pt}{\isachardoublequoteclose}\isanewline
\ \ \ \ \isakeywordONE{let}\isamarkupfalse%
\ {\isacharquery}{\kern0pt}A\isactrlsub j{\isacharequal}{\kern0pt}{\isachardoublequoteopen}acts{\isacharunderscore}{\kern0pt}of{\isacharunderscore}{\kern0pt}plan{\isacharunderscore}{\kern0pt}at\ t\isactrlsub j\ {\isasympi}s{\isachardoublequoteclose}\isanewline
\ \ \ \ \isakeywordONE{have}\isamarkupfalse%
\ {\isadigit{1}}{\isacharcolon}{\kern0pt}\ {\isachardoublequoteopen}{\isasymforall}i\ {\isasymin}\ invs{\isacharunderscore}{\kern0pt}of{\isacharunderscore}{\kern0pt}plan{\isacharunderscore}{\kern0pt}at\ t\isactrlsub j\ {\isasympi}s{\isachardot}{\kern0pt}\ M\ \isactrlsup c{\isasymTTurnstile}\isactrlsub {\isacharequal}{\kern0pt}\ i{\isachardoublequoteclose}\isanewline
\ \ \ \ \ \isakeywordTWO{and}\ {\isadigit{2}}{\isacharcolon}{\kern0pt}\ {\isachardoublequoteopen}{\isasymforall}a{\isasymin}\ {\isacharquery}{\kern0pt}A\isactrlsub j{\isachardot}{\kern0pt}\ M\ \isactrlsup c{\isasymTTurnstile}\isactrlsub {\isacharequal}{\kern0pt}\ precondition\ a{\isachardoublequoteclose}\ \isanewline
\ \ \ \ \ \isakeywordTWO{and}\ {\isadigit{3}}{\isacharcolon}{\kern0pt}\ {\isachardoublequoteopen}{\isasymforall}a\ {\isasymin}\ {\isacharquery}{\kern0pt}A\isactrlsub j{\isachardot}{\kern0pt}\ {\isasymforall}b\ {\isasymin}\ {\isacharquery}{\kern0pt}A\isactrlsub j{\isachardot}{\kern0pt}\ a\ {\isasymnoteq}\ b\ {\isasymlongrightarrow}\ acts{\isacharunderscore}{\kern0pt}non{\isacharunderscore}{\kern0pt}intrf\ a\ b{\isachardoublequoteclose}\isanewline
\ \ \ \ \ \isakeywordTWO{and}\ {\isadigit{4}}{\isacharcolon}{\kern0pt}\ {\isachardoublequoteopen}apply{\isacharunderscore}{\kern0pt}eff\ {\isacharquery}{\kern0pt}A\isactrlsub j\ M\ {\isacharequal}{\kern0pt}\ M{\isacharprime}{\kern0pt}{\isachardoublequoteclose}\isanewline
\ \ \ \ \ \ \isakeywordONE{using}\isamarkupfalse%
\ {\isacartoucheopen}valid{\isacharunderscore}{\kern0pt}state{\isacharunderscore}{\kern0pt}seq\ M\ {\isacharbrackleft}{\kern0pt}t\isactrlsub j{\isacharbrackright}{\kern0pt}\ {\isasympi}s\ M{\isacharprime}{\kern0pt}{\isacartoucheclose}\isanewline
\ \ \ \ \ \ \isakeywordONE{by}\isamarkupfalse%
\ {\isacharparenleft}{\kern0pt}induction\ M\ {\isachardoublequoteopen}{\isacharbrackleft}{\kern0pt}t\isactrlsub j{\isacharbrackright}{\kern0pt}{\isachardoublequoteclose}\ {\isasympi}s\ M{\isacharprime}{\kern0pt}\ rule{\isacharcolon}{\kern0pt}\ valid{\isacharunderscore}{\kern0pt}state{\isacharunderscore}{\kern0pt}seq{\isachardot}{\kern0pt}induct{\isacharparenright}{\kern0pt}\ {\isacharparenleft}{\kern0pt}auto\ simp{\isacharcolon}{\kern0pt}\ Let{\isacharunderscore}{\kern0pt}def{\isacharparenright}{\kern0pt}\isanewline
\isanewline
\ \ \ \ \isakeywordONE{have}\isamarkupfalse%
\ {\isachardoublequoteopen}valid{\isacharunderscore}{\kern0pt}happ{\isacharunderscore}{\kern0pt}seq\ M\ {\isacharparenleft}{\kern0pt}dropWhile\ {\isacharparenleft}{\kern0pt}leq\ t\isactrlsub i{\isacharparenright}{\kern0pt}\ {\isacharparenleft}{\kern0pt}takeWhile\ {\isacharparenleft}{\kern0pt}le\ t\isactrlsub j{\isacharparenright}{\kern0pt}\ hs{\isacharparenright}{\kern0pt}{\isacharparenright}{\kern0pt}\ M{\isachardoublequoteclose}\isanewline
\ \ \ \ \ \ \isakeywordONE{using}\isamarkupfalse%
\ {\isadigit{1}}\ valid{\isacharunderscore}{\kern0pt}happ{\isacharunderscore}{\kern0pt}seq{\isacharunderscore}{\kern0pt}leq\isactrlsub i{\isacharunderscore}{\kern0pt}le\isactrlsub j{\isacharunderscore}{\kern0pt}invs{\isacharunderscore}{\kern0pt}of{\isacharunderscore}{\kern0pt}plan{\isacharunderscore}{\kern0pt}at{\isacharbrackleft}{\kern0pt}OF\ assms{\isacharbrackright}{\kern0pt}\ \isakeywordONE{by}\isamarkupfalse%
\ auto\isanewline
\isanewline
\ \ \ \ \isakeywordONE{have}\isamarkupfalse%
\ {\isachardoublequoteopen}is{\isacharunderscore}{\kern0pt}htp\ {\isasympi}s\ t\isactrlsub j{\isachardoublequoteclose}\isanewline
\ \ \ \ \ \ \isakeywordONE{using}\isamarkupfalse%
\ assms\ \isakeywordONE{unfolding}\isamarkupfalse%
\ consec{\isacharunderscore}{\kern0pt}htps{\isacharunderscore}{\kern0pt}def\ \isakeywordONE{by}\isamarkupfalse%
\ auto\isanewline
\ \ \ \ \isakeywordTHREE{obtain}\isamarkupfalse%
\ A\isactrlsub j\ \isakeywordTWO{where}\ {\isachardoublequoteopen}dropWhile\ {\isacharparenleft}{\kern0pt}le\ t\isactrlsub j{\isacharparenright}{\kern0pt}\ {\isacharparenleft}{\kern0pt}takeWhile\ {\isacharparenleft}{\kern0pt}leq\ t\isactrlsub j{\isacharparenright}{\kern0pt}\ hs{\isacharparenright}{\kern0pt}\ {\isacharequal}{\kern0pt}\ {\isacharbrackleft}{\kern0pt}{\isacharparenleft}{\kern0pt}t\isactrlsub j{\isacharcomma}{\kern0pt}A\isactrlsub j{\isacharparenright}{\kern0pt}{\isacharbrackright}{\kern0pt}{\isachardoublequoteclose}\isanewline
\ \ \ \ \ \ \isakeywordONE{using}\isamarkupfalse%
\ assms\ ind{\isacharunderscore}{\kern0pt}happ{\isacharunderscore}{\kern0pt}seq{\isacharunderscore}{\kern0pt}le\isactrlsub i{\isacharunderscore}{\kern0pt}leq\isactrlsub i{\isacharbrackleft}{\kern0pt}OF\ {\isacartoucheopen}is{\isacharunderscore}{\kern0pt}htp\ {\isasympi}s\ t\isactrlsub j{\isacartoucheclose}{\isacharbrackright}{\kern0pt}\ \isakeywordONE{by}\isamarkupfalse%
\ auto\isanewline
\ \ \ \ \isakeywordONE{then}\isamarkupfalse%
\ \isakeywordONE{have}\isamarkupfalse%
\ {\isachardoublequoteopen}{\isacharparenleft}{\kern0pt}t\isactrlsub j{\isacharcomma}{\kern0pt}A\isactrlsub j{\isacharparenright}{\kern0pt}\ {\isasymin}\ set\ hs{\isachardoublequoteclose}\isanewline
\ \ \ \ \ \ \isakeywordONE{by}\isamarkupfalse%
\ {\isacharparenleft}{\kern0pt}metis\ list{\isachardot}{\kern0pt}set{\isacharunderscore}{\kern0pt}intros{\isacharparenleft}{\kern0pt}{\isadigit{1}}{\isacharparenright}{\kern0pt}\ set{\isacharunderscore}{\kern0pt}dropWhileD\ set{\isacharunderscore}{\kern0pt}takeWhileD{\isacharparenright}{\kern0pt}\isanewline
\ \ \ \ \isakeywordONE{then}\isamarkupfalse%
\ \isakeywordONE{have}\isamarkupfalse%
\ {\isachardoublequoteopen}set\ A\isactrlsub j\ {\isacharequal}{\kern0pt}\ acts{\isacharunderscore}{\kern0pt}of{\isacharunderscore}{\kern0pt}plan{\isacharunderscore}{\kern0pt}at\ t\isactrlsub j\ {\isasympi}s{\isachardoublequoteclose}\isanewline
\ \ \ \ \ \ \isakeywordONE{using}\isamarkupfalse%
\ assms\ {\isacartoucheopen}is{\isacharunderscore}{\kern0pt}htp\ {\isasympi}s\ t\isactrlsub j{\isacartoucheclose}\ ind{\isacharunderscore}{\kern0pt}happ{\isacharunderscore}{\kern0pt}seq{\isacharunderscore}{\kern0pt}htp{\isacharunderscore}{\kern0pt}acts{\isacharunderscore}{\kern0pt}of{\isacharunderscore}{\kern0pt}plan{\isacharunderscore}{\kern0pt}at\ \isakeywordONE{by}\isamarkupfalse%
\ auto\isanewline
\ \ \ \ \isakeywordONE{then}\isamarkupfalse%
\ \isakeywordONE{have}\isamarkupfalse%
\ {\isachardoublequoteopen}valid{\isacharunderscore}{\kern0pt}happ{\isacharunderscore}{\kern0pt}seq\ M\ {\isacharparenleft}{\kern0pt}dropWhile\ {\isacharparenleft}{\kern0pt}le\ t\isactrlsub j{\isacharparenright}{\kern0pt}\ {\isacharparenleft}{\kern0pt}takeWhile\ {\isacharparenleft}{\kern0pt}leq\ t\isactrlsub j{\isacharparenright}{\kern0pt}\ hs{\isacharparenright}{\kern0pt}{\isacharparenright}{\kern0pt}\ M{\isacharprime}{\kern0pt}{\isachardoublequoteclose}\isanewline
\ \ \ \ \ \ \isakeywordONE{using}\isamarkupfalse%
\ {\isadigit{2}}\ {\isadigit{3}}\ {\isadigit{4}}\ {\isacartoucheopen}dropWhile\ {\isacharparenleft}{\kern0pt}le\ t\isactrlsub j{\isacharparenright}{\kern0pt}\ {\isacharparenleft}{\kern0pt}takeWhile\ {\isacharparenleft}{\kern0pt}leq\ t\isactrlsub j{\isacharparenright}{\kern0pt}\ hs{\isacharparenright}{\kern0pt}\ {\isacharequal}{\kern0pt}\ {\isacharbrackleft}{\kern0pt}{\isacharparenleft}{\kern0pt}t\isactrlsub j{\isacharcomma}{\kern0pt}A\isactrlsub j{\isacharparenright}{\kern0pt}{\isacharbrackright}{\kern0pt}{\isacartoucheclose}\ \isakeywordONE{by}\isamarkupfalse%
\ {\isacharparenleft}{\kern0pt}auto\ simp{\isacharcolon}{\kern0pt}\ pairwise{\isacharunderscore}{\kern0pt}def{\isacharparenright}{\kern0pt}\isanewline
\isanewline
\ \ \ \ \isakeywordONE{have}\isamarkupfalse%
\ {\isachardoublequoteopen}is{\isacharunderscore}{\kern0pt}htp\ {\isasympi}s\ t\isactrlsub j{\isachardoublequoteclose}\ \isakeywordTWO{and}\ {\isachardoublequoteopen}t\isactrlsub i\ {\isacharless}{\kern0pt}\ t\isactrlsub j{\isachardoublequoteclose}\ \isakeywordTWO{and}\ {\isachardoublequoteopen}t\isactrlsub j\ {\isasymle}\ t\isactrlsub j{\isachardoublequoteclose}\isanewline
\ \ \ \ \ \ \isakeywordONE{using}\isamarkupfalse%
\ {\isacartoucheopen}consec{\isacharunderscore}{\kern0pt}htps\ {\isasympi}s\ t\isactrlsub i\ t\isactrlsub j{\isacartoucheclose}\ \isakeywordONE{unfolding}\isamarkupfalse%
\ consec{\isacharunderscore}{\kern0pt}htps{\isacharunderscore}{\kern0pt}def\ \isakeywordONE{by}\isamarkupfalse%
\ auto\isanewline
\ \ \ \ \isakeywordONE{have}\isamarkupfalse%
\ {\isachardoublequoteopen}strict{\isacharunderscore}{\kern0pt}sorted\ {\isacharparenleft}{\kern0pt}map\ fst\ hs{\isacharparenright}{\kern0pt}{\isachardoublequoteclose}\isanewline
\ \ \ \ \ \ \isakeywordONE{using}\isamarkupfalse%
\ {\isacartoucheopen}ind{\isacharunderscore}{\kern0pt}happ{\isacharunderscore}{\kern0pt}seq\ {\isasympi}s\ hs{\isacartoucheclose}\ \isakeywordONE{unfolding}\isamarkupfalse%
\ ind{\isacharunderscore}{\kern0pt}happ{\isacharunderscore}{\kern0pt}seq{\isacharunderscore}{\kern0pt}def\ \isakeywordONE{by}\isamarkupfalse%
\ auto\isanewline
\ \ \ \ \isakeywordONE{then}\isamarkupfalse%
\ \isakeywordONE{have}\isamarkupfalse%
\ {\isachardoublequoteopen}{\isacharquery}{\kern0pt}hs{\isacharprime}{\kern0pt}\ \isanewline
\ \ \ \ \ \ \ \ \ \ \ {\isacharequal}{\kern0pt}\ dropWhile\ {\isacharparenleft}{\kern0pt}leq\ t\isactrlsub i{\isacharparenright}{\kern0pt}\ {\isacharparenleft}{\kern0pt}takeWhile\ {\isacharparenleft}{\kern0pt}le\ t\isactrlsub j{\isacharparenright}{\kern0pt}\ hs{\isacharparenright}{\kern0pt}\ {\isacharat}{\kern0pt}\ dropWhile\ {\isacharparenleft}{\kern0pt}le\ t\isactrlsub j{\isacharparenright}{\kern0pt}\ {\isacharparenleft}{\kern0pt}takeWhile\ {\isacharparenleft}{\kern0pt}leq\ t\isactrlsub j{\isacharparenright}{\kern0pt}\ hs{\isacharparenright}{\kern0pt}{\isachardoublequoteclose}\isanewline
\ \ \ \ \ \ \isakeywordONE{using}\isamarkupfalse%
\ drop{\isacharunderscore}{\kern0pt}take{\isacharunderscore}{\kern0pt}leq\isactrlsub i{\isacharunderscore}{\kern0pt}le{\isacharunderscore}{\kern0pt}leq\isactrlsub j{\isacharunderscore}{\kern0pt}split{\isacharunderscore}{\kern0pt}app{\isacharbrackleft}{\kern0pt}OF\ {\isacartoucheopen}t\isactrlsub i\ {\isacharless}{\kern0pt}\ t\isactrlsub j{\isacartoucheclose}\ {\isacartoucheopen}t\isactrlsub j\ {\isasymle}\ t\isactrlsub j{\isacartoucheclose}{\isacharbrackright}{\kern0pt}\ \isakeywordONE{by}\isamarkupfalse%
\ fastforce\isanewline
\ \ \ \ \isakeywordONE{then}\isamarkupfalse%
\ \isakeywordTHREE{show}\isamarkupfalse%
\ {\isachardoublequoteopen}valid{\isacharunderscore}{\kern0pt}happ{\isacharunderscore}{\kern0pt}seq\ M\ {\isacharquery}{\kern0pt}hs{\isacharprime}{\kern0pt}\ M{\isacharprime}{\kern0pt}{\isachardoublequoteclose}\isanewline
\ \ \ \ \ \ \isakeywordONE{using}\isamarkupfalse%
\ {\isacartoucheopen}valid{\isacharunderscore}{\kern0pt}happ{\isacharunderscore}{\kern0pt}seq\ M\ {\isacharparenleft}{\kern0pt}dropWhile\ {\isacharparenleft}{\kern0pt}leq\ t\isactrlsub i{\isacharparenright}{\kern0pt}\ {\isacharparenleft}{\kern0pt}takeWhile\ {\isacharparenleft}{\kern0pt}le\ t\isactrlsub j{\isacharparenright}{\kern0pt}\ hs{\isacharparenright}{\kern0pt}{\isacharparenright}{\kern0pt}\ M{\isacartoucheclose}\ \isanewline
\ \ \ \ \ \ \ \ \ \ \ \ {\isacartoucheopen}valid{\isacharunderscore}{\kern0pt}happ{\isacharunderscore}{\kern0pt}seq\ M\ {\isacharparenleft}{\kern0pt}dropWhile\ {\isacharparenleft}{\kern0pt}le\ t\isactrlsub j{\isacharparenright}{\kern0pt}\ {\isacharparenleft}{\kern0pt}takeWhile\ {\isacharparenleft}{\kern0pt}leq\ t\isactrlsub j{\isacharparenright}{\kern0pt}\ hs{\isacharparenright}{\kern0pt}{\isacharparenright}{\kern0pt}\ M{\isacharprime}{\kern0pt}{\isacartoucheclose}\isanewline
\ \ \ \ \ \ \ \ \ \ \ \ valid{\isacharunderscore}{\kern0pt}happ{\isacharunderscore}{\kern0pt}seq{\isacharunderscore}{\kern0pt}app{\isacharunderscore}{\kern0pt}iff\isanewline
\ \ \ \ \ \ \isakeywordONE{by}\isamarkupfalse%
\ auto\isanewline
\ \ \isakeywordONE{qed}\isamarkupfalse%
%
\endisatagproof
{\isafoldproof}%
%
\isadelimproof
\isanewline
%
\endisadelimproof
\isanewline
\isanewline
\ \ \isakeywordONE{lemma}\isamarkupfalse%
\ cons{\isacharunderscore}{\kern0pt}consec{\isacharunderscore}{\kern0pt}htps{\isacharcolon}{\kern0pt}\isanewline
\ \ \ \ \isakeywordTWO{assumes}\ {\isachardoublequoteopen}htps{\isacharunderscore}{\kern0pt}seq\ {\isasympi}s\ htps{\isachardoublequoteclose}\ \isanewline
\ \ \ \ \ \ \ \ \isakeywordTWO{and}\ {\isachardoublequoteopen}htps\ {\isacharequal}{\kern0pt}\ ts\isactrlsub {\isadigit{1}}\ {\isacharat}{\kern0pt}\ t\isactrlsub i\ {\isacharhash}{\kern0pt}\ t\isactrlsub j\ {\isacharhash}{\kern0pt}\ ts\isactrlsub {\isadigit{2}}{\isachardoublequoteclose}\isanewline
\ \ \ \ \ \ \isakeywordTWO{shows}\ {\isachardoublequoteopen}consec{\isacharunderscore}{\kern0pt}htps\ {\isasympi}s\ t\isactrlsub i\ t\isactrlsub j{\isachardoublequoteclose}\isanewline
%
\isadelimproof
\ \ %
\endisadelimproof
%
\isatagproof
\isakeywordONE{proof}\isamarkupfalse%
\ {\isacharminus}{\kern0pt}\isanewline
\ \ \ \ \isakeywordTHREE{obtain}\isamarkupfalse%
\ ts{\isacharprime}{\kern0pt}\ \isakeywordTWO{where}\ {\isachardoublequoteopen}htps\ {\isacharequal}{\kern0pt}\ ts\isactrlsub {\isadigit{1}}\ {\isacharat}{\kern0pt}\ t\isactrlsub i\ {\isacharhash}{\kern0pt}\ ts{\isacharprime}{\kern0pt}\ {\isacharat}{\kern0pt}\ t\isactrlsub j\ {\isacharhash}{\kern0pt}\ ts\isactrlsub {\isadigit{2}}{\isachardoublequoteclose}\ \isakeywordTWO{and}\ {\isachardoublequoteopen}ts{\isacharprime}{\kern0pt}\ {\isacharequal}{\kern0pt}\ {\isacharbrackleft}{\kern0pt}{\isacharbrackright}{\kern0pt}{\isachardoublequoteclose}\isanewline
\ \ \ \ \ \ \isakeywordONE{using}\isamarkupfalse%
\ assms\ \isakeywordONE{by}\isamarkupfalse%
\ auto\isanewline
\ \ \ \ \isakeywordONE{then}\isamarkupfalse%
\ \isakeywordONE{have}\isamarkupfalse%
\ {\isachardoublequoteopen}{\isasymforall}t\ {\isasymin}\ set\ htps{\isachardot}{\kern0pt}\ t\isactrlsub i\ {\isacharless}{\kern0pt}\ t\ {\isasymand}\ t\ {\isacharless}{\kern0pt}\ t\isactrlsub j\ {\isasymlongrightarrow}\ t\ {\isasymin}\ set\ ts{\isacharprime}{\kern0pt}{\isachardoublequoteclose}\isanewline
\ \ \ \ \ \ \isakeywordONE{using}\isamarkupfalse%
\ assms\ \isakeywordONE{unfolding}\isamarkupfalse%
\ htps{\isacharunderscore}{\kern0pt}seq{\isacharunderscore}{\kern0pt}def\ \isakeywordONE{by}\isamarkupfalse%
\ {\isacharparenleft}{\kern0pt}auto\ simp{\isacharcolon}{\kern0pt}\ sorted{\isacharunderscore}{\kern0pt}append\ strict{\isacharunderscore}{\kern0pt}sorted{\isacharunderscore}{\kern0pt}iff{\isacharparenright}{\kern0pt}\isanewline
\ \ \ \ \isakeywordONE{then}\isamarkupfalse%
\ \isakeywordONE{have}\isamarkupfalse%
\ {\isachardoublequoteopen}{\isasymnot}{\isacharparenleft}{\kern0pt}{\isasymexists}t\ {\isasymin}\ set\ htps{\isachardot}{\kern0pt}\ t\isactrlsub i\ {\isacharless}{\kern0pt}\ t\ {\isasymand}\ t\ {\isacharless}{\kern0pt}\ t\isactrlsub j{\isacharparenright}{\kern0pt}{\isachardoublequoteclose}\isanewline
\ \ \ \ \ \ \isakeywordONE{using}\isamarkupfalse%
\ {\isacartoucheopen}ts{\isacharprime}{\kern0pt}\ {\isacharequal}{\kern0pt}\ {\isacharbrackleft}{\kern0pt}{\isacharbrackright}{\kern0pt}{\isacartoucheclose}\ \isakeywordONE{by}\isamarkupfalse%
\ auto\isanewline
\ \ \ \ \isakeywordONE{then}\isamarkupfalse%
\ \isakeywordONE{have}\isamarkupfalse%
\ {\isachardoublequoteopen}{\isasymforall}t{\isachardot}{\kern0pt}\ t\ {\isasymin}\ set\ htps\ {\isasymlongrightarrow}\ t\ {\isasymle}\ t\isactrlsub i\ {\isasymor}\ t\isactrlsub j\ {\isasymle}\ t{\isachardoublequoteclose}\isanewline
\ \ \ \ \ \ \isakeywordONE{by}\isamarkupfalse%
\ auto\isanewline
\ \ \ \ \isakeywordONE{then}\isamarkupfalse%
\ \isakeywordTHREE{show}\isamarkupfalse%
\ {\isacharquery}{\kern0pt}thesis\isanewline
\ \ \ \ \ \ \isakeywordONE{using}\isamarkupfalse%
\ assms\ \isakeywordONE{unfolding}\isamarkupfalse%
\ consec{\isacharunderscore}{\kern0pt}htps{\isacharunderscore}{\kern0pt}def\ htps{\isacharunderscore}{\kern0pt}seq{\isacharunderscore}{\kern0pt}def\ \isanewline
\ \ \ \ \ \ \isakeywordONE{by}\isamarkupfalse%
\ {\isacharparenleft}{\kern0pt}auto\ simp{\isacharcolon}{\kern0pt}\ sorted{\isacharunderscore}{\kern0pt}append\ strict{\isacharunderscore}{\kern0pt}sorted{\isacharunderscore}{\kern0pt}iff{\isacharparenright}{\kern0pt}\isanewline
\ \ \isakeywordONE{qed}\isamarkupfalse%
%
\endisatagproof
{\isafoldproof}%
%
\isadelimproof
\isanewline
%
\endisadelimproof
\isanewline
\ \ \isakeywordONE{lemma}\isamarkupfalse%
\ {\isacharparenleft}{\kern0pt}\isakeywordTWO{in}\ wf{\isacharunderscore}{\kern0pt}ast{\isacharunderscore}{\kern0pt}problem{\isacharparenright}{\kern0pt}\ valid{\isacharunderscore}{\kern0pt}happ{\isacharunderscore}{\kern0pt}seq{\isacharunderscore}{\kern0pt}valid{\isacharunderscore}{\kern0pt}state{\isacharunderscore}{\kern0pt}seq{\isacharunderscore}{\kern0pt}i{\isacharunderscore}{\kern0pt}to{\isacharunderscore}{\kern0pt}j{\isacharunderscore}{\kern0pt}aux{\isacharcolon}{\kern0pt}\isanewline
\ \ \ \ \isakeywordTWO{fixes}\ {\isasympi}s\isanewline
\ \ \ \ \isakeywordTWO{assumes}\ {\isachardoublequoteopen}htps{\isacharunderscore}{\kern0pt}seq\ {\isasympi}s\ htps{\isachardoublequoteclose}\isanewline
\ \ \ \ \ \ \ \ \isakeywordTWO{and}\ {\isachardoublequoteopen}htps\ {\isacharequal}{\kern0pt}\ ts\isactrlsub {\isadigit{1}}\ {\isacharat}{\kern0pt}\ t\isactrlsub i\ {\isacharhash}{\kern0pt}\ ts\isactrlsub {\isadigit{2}}\ {\isacharat}{\kern0pt}\ t\isactrlsub j\ {\isacharhash}{\kern0pt}\ ts\isactrlsub {\isadigit{3}}{\isachardoublequoteclose}\ \isanewline
\ \ \ \ \ \ \ \ \isakeywordTWO{and}\ {\isachardoublequoteopen}wf{\isacharunderscore}{\kern0pt}plan\ {\isasympi}s{\isachardoublequoteclose}\isanewline
\ \ \ \ \ \ \ \ \isakeywordTWO{and}\ {\isachardoublequoteopen}ind{\isacharunderscore}{\kern0pt}happ{\isacharunderscore}{\kern0pt}seq\ {\isasympi}s\ hs{\isachardoublequoteclose}\isanewline
\ \ \ \ \isakeywordTWO{shows}\ {\isachardoublequoteopen}valid{\isacharunderscore}{\kern0pt}happ{\isacharunderscore}{\kern0pt}seq\ M\ {\isacharparenleft}{\kern0pt}dropWhile\ {\isacharparenleft}{\kern0pt}leq\ t\isactrlsub i{\isacharparenright}{\kern0pt}\ {\isacharparenleft}{\kern0pt}takeWhile\ {\isacharparenleft}{\kern0pt}leq\ t\isactrlsub j{\isacharparenright}{\kern0pt}\ hs{\isacharparenright}{\kern0pt}{\isacharparenright}{\kern0pt}\ M{\isacharprime}{\kern0pt}\ \isanewline
\ \ \ \ \ \ \ \ \ \ \ \ {\isasymlongleftrightarrow}\ valid{\isacharunderscore}{\kern0pt}state{\isacharunderscore}{\kern0pt}seq\ M\ {\isacharparenleft}{\kern0pt}ts\isactrlsub {\isadigit{2}}\ {\isacharat}{\kern0pt}\ {\isacharbrackleft}{\kern0pt}t\isactrlsub j{\isacharbrackright}{\kern0pt}{\isacharparenright}{\kern0pt}\ {\isasympi}s\ M{\isacharprime}{\kern0pt}{\isachardoublequoteclose}\isanewline
%
\isadelimproof
\ \ \ \ %
\endisadelimproof
%
\isatagproof
\isakeywordONE{using}\isamarkupfalse%
\ assms\ \isanewline
\ \ \isakeywordONE{proof}\isamarkupfalse%
\ {\isacharparenleft}{\kern0pt}induction\ ts\isactrlsub {\isadigit{2}}\ arbitrary{\isacharcolon}{\kern0pt}\ ts\isactrlsub {\isadigit{1}}\ t\isactrlsub i\ ts\isactrlsub {\isadigit{3}}\ M{\isacharparenright}{\kern0pt}\isanewline
\ \ \ \ \isakeywordTHREE{case}\isamarkupfalse%
\ Nil\isanewline
\ \ \ \ \isakeywordONE{let}\isamarkupfalse%
\ {\isacharquery}{\kern0pt}hs{\isacharprime}{\kern0pt}{\isacharequal}{\kern0pt}{\isachardoublequoteopen}dropWhile\ {\isacharparenleft}{\kern0pt}leq\ t\isactrlsub i{\isacharparenright}{\kern0pt}\ {\isacharparenleft}{\kern0pt}takeWhile\ {\isacharparenleft}{\kern0pt}leq\ t\isactrlsub j{\isacharparenright}{\kern0pt}\ hs{\isacharparenright}{\kern0pt}{\isachardoublequoteclose}\isanewline
\ \ \ \ \isakeywordONE{have}\isamarkupfalse%
\ {\isachardoublequoteopen}consec{\isacharunderscore}{\kern0pt}htps\ {\isasympi}s\ t\isactrlsub i\ t\isactrlsub j{\isachardoublequoteclose}\isanewline
\ \ \ \ \ \ \isakeywordONE{using}\isamarkupfalse%
\ Nil{\isachardot}{\kern0pt}prems\ cons{\isacharunderscore}{\kern0pt}consec{\isacharunderscore}{\kern0pt}htps\ \isakeywordONE{by}\isamarkupfalse%
\ auto\isanewline
\ \ \ \ \isakeywordONE{then}\isamarkupfalse%
\ \isakeywordTHREE{show}\isamarkupfalse%
\ {\isacharquery}{\kern0pt}case\ \isanewline
\ \ \ \ \ \ \isakeywordONE{using}\isamarkupfalse%
\ Nil{\isachardot}{\kern0pt}prems\ valid{\isacharunderscore}{\kern0pt}happ{\isacharunderscore}{\kern0pt}seq{\isacharunderscore}{\kern0pt}valid{\isacharunderscore}{\kern0pt}state{\isacharunderscore}{\kern0pt}seq{\isacharunderscore}{\kern0pt}consec{\isacharunderscore}{\kern0pt}htps\ \isakeywordONE{by}\isamarkupfalse%
\ auto\isanewline
\ \ \isakeywordONE{next}\isamarkupfalse%
\isanewline
\ \ \ \ \isakeywordTHREE{case}\isamarkupfalse%
\ {\isacharparenleft}{\kern0pt}Cons\ t\isactrlsub i{\isacharprime}{\kern0pt}\ ts\isactrlsub {\isadigit{2}}{\isacharparenright}{\kern0pt}\isanewline
\ \ \ \ \isakeywordONE{have}\isamarkupfalse%
\ {\isachardoublequoteopen}strict{\isacharunderscore}{\kern0pt}sorted\ htps{\isachardoublequoteclose}\isanewline
\ \ \ \ \ \ \isakeywordONE{using}\isamarkupfalse%
\ assms\ \isakeywordONE{unfolding}\isamarkupfalse%
\ htps{\isacharunderscore}{\kern0pt}seq{\isacharunderscore}{\kern0pt}def\ \isakeywordONE{by}\isamarkupfalse%
\ auto\isanewline
\ \ \ \ \isakeywordONE{then}\isamarkupfalse%
\ \isakeywordONE{have}\isamarkupfalse%
\ {\isachardoublequoteopen}consec{\isacharunderscore}{\kern0pt}htps\ {\isasympi}s\ t\isactrlsub i\ t\isactrlsub i{\isacharprime}{\kern0pt}{\isachardoublequoteclose}\isanewline
\ \ \ \ \ \ \isakeywordONE{using}\isamarkupfalse%
\ Cons{\isachardot}{\kern0pt}prems\ cons{\isacharunderscore}{\kern0pt}consec{\isacharunderscore}{\kern0pt}htps\ \isakeywordONE{by}\isamarkupfalse%
\ auto\isanewline
\ \ \ \ \isakeywordONE{have}\isamarkupfalse%
\ {\isachardoublequoteopen}t\isactrlsub i\ {\isasymin}\ set\ htps{\isachardoublequoteclose}\ \isakeywordTWO{and}\ {\isachardoublequoteopen}t\isactrlsub i{\isacharprime}{\kern0pt}\ {\isasymin}\ set\ htps{\isachardoublequoteclose}\ \isakeywordTWO{and}\ {\isachardoublequoteopen}t\isactrlsub j\ {\isasymin}\ set\ htps{\isachardoublequoteclose}\isanewline
\ \ \ \ \ \ \isakeywordONE{using}\isamarkupfalse%
\ Cons{\isachardot}{\kern0pt}prems\ \isakeywordONE{by}\isamarkupfalse%
\ auto\isanewline
\ \ \ \ \isakeywordONE{have}\isamarkupfalse%
\ {\isachardoublequoteopen}t\isactrlsub i\ {\isacharless}{\kern0pt}\ t\isactrlsub j{\isachardoublequoteclose}\isanewline
\ \ \ \ \ \ \isakeywordONE{using}\isamarkupfalse%
\ Cons{\isachardot}{\kern0pt}prems\ split{\isacharunderscore}{\kern0pt}list{\isadigit{2}}{\isacharunderscore}{\kern0pt}le{\isacharbrackleft}{\kern0pt}OF\ {\isacartoucheopen}strict{\isacharunderscore}{\kern0pt}sorted\ htps{\isacartoucheclose}\ {\isacartoucheopen}t\isactrlsub i\ {\isasymin}\ set\ htps{\isacartoucheclose}\ {\isacartoucheopen}t\isactrlsub j\ {\isasymin}\ set\ htps{\isacartoucheclose}{\isacharbrackright}{\kern0pt}\isanewline
\ \ \ \ \ \ \isakeywordONE{by}\isamarkupfalse%
\ blast\isanewline
\ \ \ \ \isakeywordONE{have}\isamarkupfalse%
\ {\isachardoublequoteopen}t\isactrlsub i\ {\isasymle}\ t\isactrlsub i{\isacharprime}{\kern0pt}{\isachardoublequoteclose}\isanewline
\ \ \ \ \ \ \isakeywordONE{using}\isamarkupfalse%
\ {\isacartoucheopen}consec{\isacharunderscore}{\kern0pt}htps\ {\isasympi}s\ t\isactrlsub i\ t\isactrlsub i{\isacharprime}{\kern0pt}{\isacartoucheclose}\ \isakeywordONE{unfolding}\isamarkupfalse%
\ consec{\isacharunderscore}{\kern0pt}htps{\isacharunderscore}{\kern0pt}def\ \isakeywordONE{by}\isamarkupfalse%
\ auto\isanewline
\ \ \ \ \isakeywordONE{have}\isamarkupfalse%
\ {\isachardoublequoteopen}htps\ {\isacharequal}{\kern0pt}\ {\isacharparenleft}{\kern0pt}ts\isactrlsub {\isadigit{1}}\ {\isacharat}{\kern0pt}\ {\isacharbrackleft}{\kern0pt}t\isactrlsub i{\isacharbrackright}{\kern0pt}{\isacharparenright}{\kern0pt}\ {\isacharat}{\kern0pt}\ t\isactrlsub i{\isacharprime}{\kern0pt}\ {\isacharhash}{\kern0pt}\ ts\isactrlsub {\isadigit{2}}\ {\isacharat}{\kern0pt}\ t\isactrlsub j\ {\isacharhash}{\kern0pt}\ ts\isactrlsub {\isadigit{3}}{\isachardoublequoteclose}\isanewline
\ \ \ \ \ \ \isakeywordONE{using}\isamarkupfalse%
\ Cons{\isachardot}{\kern0pt}prems\ \isakeywordONE{by}\isamarkupfalse%
\ auto\isanewline
\ \ \ \ \isakeywordONE{then}\isamarkupfalse%
\ \isakeywordONE{have}\isamarkupfalse%
\ {\isachardoublequoteopen}t\isactrlsub i{\isacharprime}{\kern0pt}\ {\isasymle}\ t\isactrlsub j{\isachardoublequoteclose}\isanewline
\ \ \ \ \ \ \isakeywordONE{using}\isamarkupfalse%
\ Cons{\isachardot}{\kern0pt}prems\ split{\isacharunderscore}{\kern0pt}list{\isadigit{2}}{\isacharunderscore}{\kern0pt}le{\isacharbrackleft}{\kern0pt}OF\ {\isacartoucheopen}strict{\isacharunderscore}{\kern0pt}sorted\ htps{\isacartoucheclose}\ {\isacartoucheopen}t\isactrlsub i{\isacharprime}{\kern0pt}\ {\isasymin}\ set\ htps{\isacartoucheclose}\ {\isacartoucheopen}t\isactrlsub j\ {\isasymin}\ set\ htps{\isacartoucheclose}{\isacharbrackright}{\kern0pt}\isanewline
\ \ \ \ \ \ \ \ \ \ \ \ less{\isacharunderscore}{\kern0pt}eq{\isacharunderscore}{\kern0pt}rat{\isacharunderscore}{\kern0pt}def\ \isanewline
\ \ \ \ \ \ \isakeywordONE{by}\isamarkupfalse%
\ blast\isanewline
\isanewline
\ \ \ \ \isakeywordONE{let}\isamarkupfalse%
\ {\isacharquery}{\kern0pt}hs{\isacharprime}{\kern0pt}{\isacharequal}{\kern0pt}{\isachardoublequoteopen}dropWhile\ {\isacharparenleft}{\kern0pt}leq\ t\isactrlsub i{\isacharparenright}{\kern0pt}\ {\isacharparenleft}{\kern0pt}takeWhile\ {\isacharparenleft}{\kern0pt}leq\ t\isactrlsub j{\isacharparenright}{\kern0pt}\ hs{\isacharparenright}{\kern0pt}{\isachardoublequoteclose}\isanewline
\ \ \ \ \isakeywordONE{have}\isamarkupfalse%
\ {\isachardoublequoteopen}strict{\isacharunderscore}{\kern0pt}sorted\ {\isacharparenleft}{\kern0pt}map\ fst\ hs{\isacharparenright}{\kern0pt}{\isachardoublequoteclose}\isanewline
\ \ \ \ \ \ \isakeywordONE{using}\isamarkupfalse%
\ {\isacartoucheopen}ind{\isacharunderscore}{\kern0pt}happ{\isacharunderscore}{\kern0pt}seq\ {\isasympi}s\ hs{\isacartoucheclose}\ \isakeywordONE{unfolding}\isamarkupfalse%
\ ind{\isacharunderscore}{\kern0pt}happ{\isacharunderscore}{\kern0pt}seq{\isacharunderscore}{\kern0pt}def\ \isakeywordONE{by}\isamarkupfalse%
\ auto\isanewline
\ \ \ \ \isakeywordONE{then}\isamarkupfalse%
\ \isakeywordONE{have}\isamarkupfalse%
\ hs{\isacharprime}{\kern0pt}{\isacharunderscore}{\kern0pt}app{\isacharcolon}{\kern0pt}\ {\isachardoublequoteopen}{\isacharquery}{\kern0pt}hs{\isacharprime}{\kern0pt}\ {\isacharequal}{\kern0pt}\ dropWhile\ {\isacharparenleft}{\kern0pt}leq\ t\isactrlsub i{\isacharparenright}{\kern0pt}\ {\isacharparenleft}{\kern0pt}takeWhile\ {\isacharparenleft}{\kern0pt}leq\ t\isactrlsub i{\isacharprime}{\kern0pt}{\isacharparenright}{\kern0pt}\ hs{\isacharparenright}{\kern0pt}\ \isanewline
\ \ \ \ \ \ \ \ \ \ \ \ \ \ \ \ \ \ \ \ \ \ \ \ {\isacharat}{\kern0pt}\ dropWhile\ {\isacharparenleft}{\kern0pt}leq\ t\isactrlsub i{\isacharprime}{\kern0pt}{\isacharparenright}{\kern0pt}\ {\isacharparenleft}{\kern0pt}takeWhile\ {\isacharparenleft}{\kern0pt}leq\ t\isactrlsub j{\isacharparenright}{\kern0pt}\ hs{\isacharparenright}{\kern0pt}{\isachardoublequoteclose}\isanewline
\ \ \ \ \ \ \isakeywordONE{using}\isamarkupfalse%
\ drop{\isacharunderscore}{\kern0pt}take{\isacharunderscore}{\kern0pt}leq\isactrlsub i{\isacharunderscore}{\kern0pt}leq{\isacharunderscore}{\kern0pt}leq\isactrlsub j{\isacharunderscore}{\kern0pt}simp{\isacharbrackleft}{\kern0pt}OF\ {\isacartoucheopen}t\isactrlsub i\ {\isasymle}\ t\isactrlsub i{\isacharprime}{\kern0pt}{\isacartoucheclose}\ {\isacartoucheopen}t\isactrlsub i{\isacharprime}{\kern0pt}\ {\isasymle}\ t\isactrlsub j{\isacartoucheclose}{\isacharbrackright}{\kern0pt}\ \isakeywordONE{by}\isamarkupfalse%
\ fastforce\isanewline
\isanewline
\ \ \ \ \isakeywordTHREE{show}\isamarkupfalse%
\ {\isachardoublequoteopen}valid{\isacharunderscore}{\kern0pt}happ{\isacharunderscore}{\kern0pt}seq\ M\ {\isacharquery}{\kern0pt}hs{\isacharprime}{\kern0pt}\ M{\isacharprime}{\kern0pt}\ {\isasymlongleftrightarrow}\ valid{\isacharunderscore}{\kern0pt}state{\isacharunderscore}{\kern0pt}seq\ M\ {\isacharparenleft}{\kern0pt}{\isacharparenleft}{\kern0pt}t\isactrlsub i{\isacharprime}{\kern0pt}\ {\isacharhash}{\kern0pt}\ ts\isactrlsub {\isadigit{2}}{\isacharparenright}{\kern0pt}\ {\isacharat}{\kern0pt}\ {\isacharbrackleft}{\kern0pt}t\isactrlsub j{\isacharbrackright}{\kern0pt}{\isacharparenright}{\kern0pt}\ {\isasympi}s\ M{\isacharprime}{\kern0pt}{\isachardoublequoteclose}\isanewline
\ \ \ \ \isakeywordONE{proof}\isamarkupfalse%
\isanewline
\ \ \ \ \ \ \isakeywordTHREE{assume}\isamarkupfalse%
\ {\isachardoublequoteopen}valid{\isacharunderscore}{\kern0pt}happ{\isacharunderscore}{\kern0pt}seq\ M\ {\isacharquery}{\kern0pt}hs{\isacharprime}{\kern0pt}\ M{\isacharprime}{\kern0pt}{\isachardoublequoteclose}\isanewline
\ \ \ \ \ \ \isakeywordONE{then}\isamarkupfalse%
\ \isakeywordTHREE{obtain}\isamarkupfalse%
\ M{\isacharprime}{\kern0pt}{\isacharprime}{\kern0pt}\ \isakeywordTWO{where}\ ii{\isacharprime}{\kern0pt}{\isacharcolon}{\kern0pt}\ {\isachardoublequoteopen}valid{\isacharunderscore}{\kern0pt}happ{\isacharunderscore}{\kern0pt}seq\ M\ {\isacharparenleft}{\kern0pt}dropWhile\ {\isacharparenleft}{\kern0pt}leq\ t\isactrlsub i{\isacharparenright}{\kern0pt}\ {\isacharparenleft}{\kern0pt}takeWhile\ {\isacharparenleft}{\kern0pt}leq\ t\isactrlsub i{\isacharprime}{\kern0pt}{\isacharparenright}{\kern0pt}\ hs{\isacharparenright}{\kern0pt}{\isacharparenright}{\kern0pt}\ M{\isacharprime}{\kern0pt}{\isacharprime}{\kern0pt}{\isachardoublequoteclose}\ \isanewline
\ \ \ \ \ \ \ \ \ \ \ \ \ \ \ \ \ \ \ \ \ \ \ \ \isakeywordTWO{and}\ i{\isacharprime}{\kern0pt}j{\isacharcolon}{\kern0pt}\ {\isachardoublequoteopen}valid{\isacharunderscore}{\kern0pt}happ{\isacharunderscore}{\kern0pt}seq\ M{\isacharprime}{\kern0pt}{\isacharprime}{\kern0pt}\ {\isacharparenleft}{\kern0pt}dropWhile\ {\isacharparenleft}{\kern0pt}leq\ t\isactrlsub i{\isacharprime}{\kern0pt}{\isacharparenright}{\kern0pt}\ {\isacharparenleft}{\kern0pt}takeWhile\ {\isacharparenleft}{\kern0pt}leq\ t\isactrlsub j{\isacharparenright}{\kern0pt}\ hs{\isacharparenright}{\kern0pt}{\isacharparenright}{\kern0pt}\ M{\isacharprime}{\kern0pt}{\isachardoublequoteclose}\isanewline
\ \ \ \ \ \ \ \ \isakeywordONE{using}\isamarkupfalse%
\ hs{\isacharprime}{\kern0pt}{\isacharunderscore}{\kern0pt}app\ valid{\isacharunderscore}{\kern0pt}happ{\isacharunderscore}{\kern0pt}seq{\isacharunderscore}{\kern0pt}app{\isacharunderscore}{\kern0pt}iff\ \isakeywordONE{by}\isamarkupfalse%
\ auto\isanewline
\ \ \ \ \ \ \isakeywordONE{then}\isamarkupfalse%
\ \isakeywordONE{have}\isamarkupfalse%
\ {\isadigit{1}}{\isacharcolon}{\kern0pt}\ {\isachardoublequoteopen}valid{\isacharunderscore}{\kern0pt}state{\isacharunderscore}{\kern0pt}seq\ M\ {\isacharbrackleft}{\kern0pt}t\isactrlsub i{\isacharprime}{\kern0pt}{\isacharbrackright}{\kern0pt}\ {\isasympi}s\ M{\isacharprime}{\kern0pt}{\isacharprime}{\kern0pt}{\isachardoublequoteclose}\isanewline
\ \ \ \ \ \ \ \ \isakeywordONE{using}\isamarkupfalse%
\ Cons{\isachardot}{\kern0pt}prems\ {\isacartoucheopen}consec{\isacharunderscore}{\kern0pt}htps\ {\isasympi}s\ t\isactrlsub i\ t\isactrlsub i{\isacharprime}{\kern0pt}{\isacartoucheclose}\ valid{\isacharunderscore}{\kern0pt}happ{\isacharunderscore}{\kern0pt}seq{\isacharunderscore}{\kern0pt}valid{\isacharunderscore}{\kern0pt}state{\isacharunderscore}{\kern0pt}seq{\isacharunderscore}{\kern0pt}consec{\isacharunderscore}{\kern0pt}htps\ \isakeywordONE{by}\isamarkupfalse%
\ blast\isanewline
\ \ \ \ \ \ \isakeywordONE{have}\isamarkupfalse%
\ {\isadigit{2}}{\isacharcolon}{\kern0pt}\ {\isachardoublequoteopen}valid{\isacharunderscore}{\kern0pt}state{\isacharunderscore}{\kern0pt}seq\ M{\isacharprime}{\kern0pt}{\isacharprime}{\kern0pt}\ {\isacharparenleft}{\kern0pt}ts\isactrlsub {\isadigit{2}}\ {\isacharat}{\kern0pt}\ {\isacharbrackleft}{\kern0pt}t\isactrlsub j{\isacharbrackright}{\kern0pt}{\isacharparenright}{\kern0pt}\ {\isasympi}s\ M{\isacharprime}{\kern0pt}{\isachardoublequoteclose}\isanewline
\ \ \ \ \ \ \ \ \isakeywordONE{using}\isamarkupfalse%
\ assms\ Cons{\isachardot}{\kern0pt}IH\ {\isacartoucheopen}htps\ {\isacharequal}{\kern0pt}\ {\isacharparenleft}{\kern0pt}ts\isactrlsub {\isadigit{1}}\ {\isacharat}{\kern0pt}\ {\isacharbrackleft}{\kern0pt}t\isactrlsub i{\isacharbrackright}{\kern0pt}{\isacharparenright}{\kern0pt}\ {\isacharat}{\kern0pt}\ t\isactrlsub i{\isacharprime}{\kern0pt}\ {\isacharhash}{\kern0pt}\ ts\isactrlsub {\isadigit{2}}\ {\isacharat}{\kern0pt}\ t\isactrlsub j\ {\isacharhash}{\kern0pt}\ ts\isactrlsub {\isadigit{3}}{\isacartoucheclose}\ i{\isacharprime}{\kern0pt}j\ \isakeywordONE{by}\isamarkupfalse%
\ blast\isanewline
\ \ \ \ \ \ \isakeywordONE{then}\isamarkupfalse%
\ \isakeywordTHREE{show}\isamarkupfalse%
\ {\isachardoublequoteopen}valid{\isacharunderscore}{\kern0pt}state{\isacharunderscore}{\kern0pt}seq\ M\ {\isacharparenleft}{\kern0pt}{\isacharparenleft}{\kern0pt}t\isactrlsub i{\isacharprime}{\kern0pt}\ {\isacharhash}{\kern0pt}\ ts\isactrlsub {\isadigit{2}}{\isacharparenright}{\kern0pt}\ {\isacharat}{\kern0pt}\ {\isacharbrackleft}{\kern0pt}t\isactrlsub j{\isacharbrackright}{\kern0pt}{\isacharparenright}{\kern0pt}\ {\isasympi}s\ M{\isacharprime}{\kern0pt}{\isachardoublequoteclose}\isanewline
\ \ \ \ \ \ \ \ \isakeywordONE{using}\isamarkupfalse%
\ {\isadigit{1}}\ {\isadigit{2}}\ valid{\isacharunderscore}{\kern0pt}state{\isacharunderscore}{\kern0pt}seq{\isacharunderscore}{\kern0pt}app{\isacharunderscore}{\kern0pt}iff{\isacharbrackleft}{\kern0pt}\isakeywordTWO{where}\ ts\isactrlsub {\isadigit{1}}{\isacharequal}{\kern0pt}{\isachardoublequoteopen}{\isacharbrackleft}{\kern0pt}t\isactrlsub i{\isacharprime}{\kern0pt}{\isacharbrackright}{\kern0pt}{\isachardoublequoteclose}{\isacharbrackright}{\kern0pt}\ \isakeywordONE{by}\isamarkupfalse%
\ auto\ \ \ \ \ \ \ \ \ \ \ \ \ \ \isanewline
\ \ \ \ \isakeywordONE{next}\isamarkupfalse%
\isanewline
\ \ \ \ \ \ \isakeywordTHREE{assume}\isamarkupfalse%
\ {\isachardoublequoteopen}valid{\isacharunderscore}{\kern0pt}state{\isacharunderscore}{\kern0pt}seq\ M\ {\isacharparenleft}{\kern0pt}{\isacharparenleft}{\kern0pt}t\isactrlsub i{\isacharprime}{\kern0pt}\ {\isacharhash}{\kern0pt}\ ts\isactrlsub {\isadigit{2}}{\isacharparenright}{\kern0pt}\ {\isacharat}{\kern0pt}\ {\isacharbrackleft}{\kern0pt}t\isactrlsub j{\isacharbrackright}{\kern0pt}{\isacharparenright}{\kern0pt}\ {\isasympi}s\ M{\isacharprime}{\kern0pt}{\isachardoublequoteclose}\isanewline
\ \ \ \ \ \ \isakeywordONE{then}\isamarkupfalse%
\ \isakeywordTHREE{obtain}\isamarkupfalse%
\ M{\isacharprime}{\kern0pt}{\isacharprime}{\kern0pt}\ \isakeywordTWO{where}\ i{\isacharprime}{\kern0pt}{\isacharcolon}{\kern0pt}\ {\isachardoublequoteopen}valid{\isacharunderscore}{\kern0pt}state{\isacharunderscore}{\kern0pt}seq\ M\ {\isacharbrackleft}{\kern0pt}t\isactrlsub i{\isacharprime}{\kern0pt}{\isacharbrackright}{\kern0pt}\ {\isasympi}s\ M{\isacharprime}{\kern0pt}{\isacharprime}{\kern0pt}{\isachardoublequoteclose}\ \isanewline
\ \ \ \ \ \ \ \ \ \ \ \ \ \ \ \ \ \ \ \ \ \ \ \ \isakeywordTWO{and}\ ts\isactrlsub {\isadigit{2}}j{\isacharcolon}{\kern0pt}\ {\isachardoublequoteopen}valid{\isacharunderscore}{\kern0pt}state{\isacharunderscore}{\kern0pt}seq\ M{\isacharprime}{\kern0pt}{\isacharprime}{\kern0pt}\ {\isacharparenleft}{\kern0pt}ts\isactrlsub {\isadigit{2}}\ {\isacharat}{\kern0pt}\ {\isacharbrackleft}{\kern0pt}t\isactrlsub j{\isacharbrackright}{\kern0pt}{\isacharparenright}{\kern0pt}\ {\isasympi}s\ M{\isacharprime}{\kern0pt}{\isachardoublequoteclose}\isanewline
\ \ \ \ \ \ \ \ \isakeywordONE{using}\isamarkupfalse%
\ valid{\isacharunderscore}{\kern0pt}state{\isacharunderscore}{\kern0pt}seq{\isacharunderscore}{\kern0pt}app{\isacharunderscore}{\kern0pt}iff{\isacharbrackleft}{\kern0pt}\isakeywordTWO{where}\ ts\isactrlsub {\isadigit{1}}{\isacharequal}{\kern0pt}{\isachardoublequoteopen}{\isacharbrackleft}{\kern0pt}t\isactrlsub i{\isacharprime}{\kern0pt}{\isacharbrackright}{\kern0pt}{\isachardoublequoteclose}{\isacharbrackright}{\kern0pt}\ \isakeywordONE{by}\isamarkupfalse%
\ auto\isanewline
\ \ \ \ \ \ \isakeywordONE{then}\isamarkupfalse%
\ \isakeywordONE{have}\isamarkupfalse%
\ {\isadigit{1}}{\isacharcolon}{\kern0pt}\ {\isachardoublequoteopen}valid{\isacharunderscore}{\kern0pt}happ{\isacharunderscore}{\kern0pt}seq\ M\ {\isacharparenleft}{\kern0pt}dropWhile\ {\isacharparenleft}{\kern0pt}leq\ t\isactrlsub i{\isacharparenright}{\kern0pt}\ {\isacharparenleft}{\kern0pt}takeWhile\ {\isacharparenleft}{\kern0pt}leq\ t\isactrlsub i{\isacharprime}{\kern0pt}{\isacharparenright}{\kern0pt}\ hs{\isacharparenright}{\kern0pt}{\isacharparenright}{\kern0pt}\ M{\isacharprime}{\kern0pt}{\isacharprime}{\kern0pt}{\isachardoublequoteclose}\isanewline
\ \ \ \ \ \ \ \ \isakeywordONE{using}\isamarkupfalse%
\ Cons{\isachardot}{\kern0pt}prems\ {\isacartoucheopen}consec{\isacharunderscore}{\kern0pt}htps\ {\isasympi}s\ t\isactrlsub i\ t\isactrlsub i{\isacharprime}{\kern0pt}{\isacartoucheclose}\ valid{\isacharunderscore}{\kern0pt}happ{\isacharunderscore}{\kern0pt}seq{\isacharunderscore}{\kern0pt}valid{\isacharunderscore}{\kern0pt}state{\isacharunderscore}{\kern0pt}seq{\isacharunderscore}{\kern0pt}consec{\isacharunderscore}{\kern0pt}htps\isanewline
\ \ \ \ \ \ \ \ \isakeywordONE{by}\isamarkupfalse%
\ blast\isanewline
\ \ \ \ \ \ \isakeywordONE{have}\isamarkupfalse%
\ {\isadigit{2}}{\isacharcolon}{\kern0pt}\ {\isachardoublequoteopen}valid{\isacharunderscore}{\kern0pt}happ{\isacharunderscore}{\kern0pt}seq\ M{\isacharprime}{\kern0pt}{\isacharprime}{\kern0pt}\ {\isacharparenleft}{\kern0pt}dropWhile\ {\isacharparenleft}{\kern0pt}leq\ t\isactrlsub i{\isacharprime}{\kern0pt}{\isacharparenright}{\kern0pt}\ {\isacharparenleft}{\kern0pt}takeWhile\ {\isacharparenleft}{\kern0pt}leq\ t\isactrlsub j{\isacharparenright}{\kern0pt}\ hs{\isacharparenright}{\kern0pt}{\isacharparenright}{\kern0pt}\ M{\isacharprime}{\kern0pt}{\isachardoublequoteclose}\isanewline
\ \ \ \ \ \ \ \ \isakeywordONE{using}\isamarkupfalse%
\ assms\ Cons{\isachardot}{\kern0pt}IH\ {\isacartoucheopen}htps\ {\isacharequal}{\kern0pt}\ {\isacharparenleft}{\kern0pt}ts\isactrlsub {\isadigit{1}}\ {\isacharat}{\kern0pt}\ {\isacharbrackleft}{\kern0pt}t\isactrlsub i{\isacharbrackright}{\kern0pt}{\isacharparenright}{\kern0pt}\ {\isacharat}{\kern0pt}\ t\isactrlsub i{\isacharprime}{\kern0pt}\ {\isacharhash}{\kern0pt}\ ts\isactrlsub {\isadigit{2}}\ {\isacharat}{\kern0pt}\ t\isactrlsub j\ {\isacharhash}{\kern0pt}\ ts\isactrlsub {\isadigit{3}}{\isacartoucheclose}\ ts\isactrlsub {\isadigit{2}}j\ \isakeywordONE{by}\isamarkupfalse%
\ blast\isanewline
\ \ \ \ \ \ \isakeywordONE{then}\isamarkupfalse%
\ \isakeywordTHREE{show}\isamarkupfalse%
\ {\isachardoublequoteopen}valid{\isacharunderscore}{\kern0pt}happ{\isacharunderscore}{\kern0pt}seq\ M\ {\isacharquery}{\kern0pt}hs{\isacharprime}{\kern0pt}\ M{\isacharprime}{\kern0pt}{\isachardoublequoteclose}\isanewline
\ \ \ \ \ \ \ \ \isakeywordONE{using}\isamarkupfalse%
\ hs{\isacharprime}{\kern0pt}{\isacharunderscore}{\kern0pt}app\ {\isadigit{1}}\ {\isadigit{2}}\ valid{\isacharunderscore}{\kern0pt}happ{\isacharunderscore}{\kern0pt}seq{\isacharunderscore}{\kern0pt}app{\isacharunderscore}{\kern0pt}iff\ \isakeywordONE{by}\isamarkupfalse%
\ auto\isanewline
\ \ \ \ \isakeywordONE{qed}\isamarkupfalse%
\isanewline
\ \ \isakeywordONE{qed}\isamarkupfalse%
%
\endisatagproof
{\isafoldproof}%
%
\isadelimproof
\isanewline
%
\endisadelimproof
\isanewline
\ \ \isakeywordONE{lemma}\isamarkupfalse%
\ prod{\isacharunderscore}{\kern0pt}list{\isacharunderscore}{\kern0pt}Nil{\isacharcolon}{\kern0pt}\ {\isachardoublequoteopen}{\isasymforall}x\isactrlsub {\isadigit{1}}\ x\isactrlsub {\isadigit{2}}{\isachardot}{\kern0pt}\ {\isacharparenleft}{\kern0pt}x\isactrlsub {\isadigit{1}}{\isacharcomma}{\kern0pt}x\isactrlsub {\isadigit{2}}{\isacharparenright}{\kern0pt}\ {\isasymnotin}\ set\ xs\ {\isasymLongrightarrow}\ xs\ {\isacharequal}{\kern0pt}\ {\isacharbrackleft}{\kern0pt}{\isacharbrackright}{\kern0pt}{\isachardoublequoteclose}\isanewline
%
\isadelimproof
\ \ \ \ %
\endisadelimproof
%
\isatagproof
\isakeywordONE{by}\isamarkupfalse%
\ {\isacharparenleft}{\kern0pt}induction\ xs{\isacharparenright}{\kern0pt}\ auto%
\endisatagproof
{\isafoldproof}%
%
\isadelimproof
%
\endisadelimproof
%
\begin{isamarkuptext}%
An induced happening sequence is only Nil iff the orginal plan was Nil%
\end{isamarkuptext}\isamarkuptrue%
\ \ \isakeywordONE{lemma}\isamarkupfalse%
\ ind{\isacharunderscore}{\kern0pt}happ{\isacharunderscore}{\kern0pt}seq{\isacharunderscore}{\kern0pt}Nil{\isacharcolon}{\kern0pt}\ {\isachardoublequoteopen}ind{\isacharunderscore}{\kern0pt}happ{\isacharunderscore}{\kern0pt}seq\ {\isasympi}s\ hs\ {\isasymLongrightarrow}\ {\isasympi}s\ {\isacharequal}{\kern0pt}\ {\isacharbrackleft}{\kern0pt}{\isacharbrackright}{\kern0pt}\ {\isasymlongleftrightarrow}\ hs\ {\isacharequal}{\kern0pt}\ {\isacharbrackleft}{\kern0pt}{\isacharbrackright}{\kern0pt}{\isachardoublequoteclose}\isanewline
%
\isadelimproof
\ \ \ \ %
\endisadelimproof
%
\isatagproof
\isakeywordONE{unfolding}\isamarkupfalse%
\ ind{\isacharunderscore}{\kern0pt}happ{\isacharunderscore}{\kern0pt}seq{\isacharunderscore}{\kern0pt}def\ simple{\isacharunderscore}{\kern0pt}acts{\isacharunderscore}{\kern0pt}def\ durative{\isacharunderscore}{\kern0pt}acts{\isacharunderscore}{\kern0pt}def\ is{\isacharunderscore}{\kern0pt}act{\isacharunderscore}{\kern0pt}simple{\isacharunderscore}{\kern0pt}def\ \isanewline
\ \ \ \ \isakeywordONE{by}\isamarkupfalse%
\ {\isacharparenleft}{\kern0pt}auto\ simp{\isacharcolon}{\kern0pt}\ prod{\isacharunderscore}{\kern0pt}list{\isacharunderscore}{\kern0pt}Nil{\isacharparenright}{\kern0pt}%
\endisatagproof
{\isafoldproof}%
%
\isadelimproof
%
\endisadelimproof
\ \isanewline
\isanewline
\ \ \isakeywordONE{lemma}\isamarkupfalse%
\ htps{\isacharunderscore}{\kern0pt}seq{\isacharunderscore}{\kern0pt}Nil{\isacharunderscore}{\kern0pt}iff{\isacharcolon}{\kern0pt}\isanewline
\ \ \ \ \isakeywordTWO{assumes}\ {\isachardoublequoteopen}htps{\isacharunderscore}{\kern0pt}seq\ {\isasympi}s\ htps{\isachardoublequoteclose}\isanewline
\ \ \ \ \isakeywordTWO{shows}\ {\isachardoublequoteopen}{\isasympi}s\ {\isacharequal}{\kern0pt}\ {\isacharbrackleft}{\kern0pt}{\isacharbrackright}{\kern0pt}\ {\isasymlongleftrightarrow}\ htps\ {\isacharequal}{\kern0pt}\ {\isacharbrackleft}{\kern0pt}{\isacharbrackright}{\kern0pt}{\isachardoublequoteclose}\isanewline
%
\isadelimproof
\ \ %
\endisadelimproof
%
\isatagproof
\isakeywordONE{proof}\isamarkupfalse%
\isanewline
\ \ \ \ \isakeywordTHREE{assume}\isamarkupfalse%
\ {\isachardoublequoteopen}{\isasympi}s\ {\isacharequal}{\kern0pt}\ {\isacharbrackleft}{\kern0pt}{\isacharbrackright}{\kern0pt}{\isachardoublequoteclose}\isanewline
\ \ \ \ \isakeywordONE{then}\isamarkupfalse%
\ \isakeywordTHREE{show}\isamarkupfalse%
\ {\isachardoublequoteopen}htps\ {\isacharequal}{\kern0pt}\ {\isacharbrackleft}{\kern0pt}{\isacharbrackright}{\kern0pt}{\isachardoublequoteclose}\isanewline
\ \ \ \ \ \ \isakeywordONE{using}\isamarkupfalse%
\ assms\ dur{\isacharunderscore}{\kern0pt}acts{\isacharunderscore}{\kern0pt}subset\ \isakeywordONE{unfolding}\isamarkupfalse%
\ htps{\isacharunderscore}{\kern0pt}seq{\isacharunderscore}{\kern0pt}def\ is{\isacharunderscore}{\kern0pt}htp{\isacharunderscore}{\kern0pt}def\ \isakeywordONE{by}\isamarkupfalse%
\ fastforce\isanewline
\ \ \isakeywordONE{next}\isamarkupfalse%
\isanewline
\ \ \ \ \isakeywordTHREE{assume}\isamarkupfalse%
\ {\isachardoublequoteopen}htps\ {\isacharequal}{\kern0pt}\ {\isacharbrackleft}{\kern0pt}{\isacharbrackright}{\kern0pt}{\isachardoublequoteclose}\isanewline
\ \ \ \ \isakeywordTHREE{show}\isamarkupfalse%
\ {\isachardoublequoteopen}{\isasympi}s\ {\isacharequal}{\kern0pt}\ {\isacharbrackleft}{\kern0pt}{\isacharbrackright}{\kern0pt}{\isachardoublequoteclose}\isanewline
\ \ \ \ \isakeywordONE{proof}\isamarkupfalse%
\ {\isacharparenleft}{\kern0pt}rule\ ccontr{\isacharparenright}{\kern0pt}\isanewline
\ \ \ \ \ \ \isakeywordTHREE{assume}\isamarkupfalse%
\ {\isachardoublequoteopen}{\isasympi}s\ {\isasymnoteq}\ {\isacharbrackleft}{\kern0pt}{\isacharbrackright}{\kern0pt}{\isachardoublequoteclose}\isanewline
\ \ \ \ \ \ \isakeywordONE{then}\isamarkupfalse%
\ \isakeywordTHREE{obtain}\isamarkupfalse%
\ t\isactrlsub {\isasympi}\ {\isasympi}\ \isakeywordTWO{where}\ {\isachardoublequoteopen}{\isacharparenleft}{\kern0pt}t\isactrlsub {\isasympi}{\isacharcomma}{\kern0pt}{\isasympi}{\isacharparenright}{\kern0pt}\ {\isasymin}\ set\ {\isasympi}s{\isachardoublequoteclose}\isanewline
\ \ \ \ \ \ \ \ \isakeywordONE{by}\isamarkupfalse%
\ {\isacharparenleft}{\kern0pt}meson\ prod{\isacharunderscore}{\kern0pt}list{\isacharunderscore}{\kern0pt}Nil{\isacharparenright}{\kern0pt}\isanewline
\ \ \ \ \ \ \isakeywordONE{then}\isamarkupfalse%
\ \isakeywordONE{have}\isamarkupfalse%
\ {\isachardoublequoteopen}t\isactrlsub {\isasympi}\ {\isasymin}\ set\ htps{\isachardoublequoteclose}\isanewline
\ \ \ \ \ \ \ \ \isakeywordONE{using}\isamarkupfalse%
\ assms\ \isakeywordONE{unfolding}\isamarkupfalse%
\ htps{\isacharunderscore}{\kern0pt}seq{\isacharunderscore}{\kern0pt}def\ is{\isacharunderscore}{\kern0pt}htp{\isacharunderscore}{\kern0pt}def\ \isakeywordONE{by}\isamarkupfalse%
\ auto\isanewline
\ \ \ \ \ \ \isakeywordONE{then}\isamarkupfalse%
\ \isakeywordTHREE{show}\isamarkupfalse%
\ {\isachardoublequoteopen}False{\isachardoublequoteclose}\isanewline
\ \ \ \ \ \ \ \ \isakeywordONE{using}\isamarkupfalse%
\ {\isacartoucheopen}htps\ {\isacharequal}{\kern0pt}\ {\isacharbrackleft}{\kern0pt}{\isacharbrackright}{\kern0pt}{\isacartoucheclose}\ \isakeywordONE{by}\isamarkupfalse%
\ simp\isanewline
\ \ \ \ \isakeywordONE{qed}\isamarkupfalse%
\isanewline
\ \ \isakeywordONE{qed}\isamarkupfalse%
%
\endisatagproof
{\isafoldproof}%
%
\isadelimproof
\isanewline
%
\endisadelimproof
\isanewline
\ \ \isakeywordONE{lemma}\isamarkupfalse%
\ non{\isacharunderscore}{\kern0pt}Nil{\isacharunderscore}{\kern0pt}sorted{\isacharunderscore}{\kern0pt}last{\isacharunderscore}{\kern0pt}max{\isacharcolon}{\kern0pt}\ \isanewline
\ \ \ \ \isakeywordTWO{assumes}\ {\isachardoublequoteopen}{\isacharparenleft}{\kern0pt}xs{\isacharcolon}{\kern0pt}{\isacharcolon}{\kern0pt}{\isacharparenleft}{\kern0pt}{\isacharprime}{\kern0pt}a{\isacharcolon}{\kern0pt}{\isacharcolon}{\kern0pt}linorder{\isacharparenright}{\kern0pt}\ list{\isacharparenright}{\kern0pt}\ {\isasymnoteq}\ {\isacharbrackleft}{\kern0pt}{\isacharbrackright}{\kern0pt}{\isachardoublequoteclose}\ \isanewline
\ \ \ \ \ \ \ \ \isakeywordTWO{and}\ {\isachardoublequoteopen}sorted\ xs{\isachardoublequoteclose}\ \isanewline
\ \ \ \ \ \ \ \ \isakeywordTWO{and}\ {\isachardoublequoteopen}x\isactrlsub m\isactrlsub a\isactrlsub x\ {\isacharequal}{\kern0pt}\ last\ xs{\isachardoublequoteclose}\ \isanewline
\ \ \ \ \isakeywordTWO{shows}\ {\isachardoublequoteopen}{\isasymforall}x\ {\isasymin}\ set\ xs{\isachardot}{\kern0pt}\ x\ {\isasymle}\ x\isactrlsub m\isactrlsub a\isactrlsub x{\isachardoublequoteclose}\isanewline
%
\isadelimproof
\ \ \ \ %
\endisadelimproof
%
\isatagproof
\isakeywordONE{using}\isamarkupfalse%
\ assms\ \isakeywordONE{by}\isamarkupfalse%
\ {\isacharparenleft}{\kern0pt}induction\ xs{\isacharparenright}{\kern0pt}\ auto%
\endisatagproof
{\isafoldproof}%
%
\isadelimproof
%
\endisadelimproof
%
\begin{isamarkuptext}%
Proof for the existence of a sequence of happening time point for every plan. 
  Needed later in the proof for the completeness of the semantics.%
\end{isamarkuptext}\isamarkuptrue%
\ \ \isakeywordONE{lemma}\isamarkupfalse%
\ htps{\isacharunderscore}{\kern0pt}seq{\isacharunderscore}{\kern0pt}exists{\isacharcolon}{\kern0pt}\ {\isachardoublequoteopen}{\isasymexists}htps{\isachardot}{\kern0pt}\ htps{\isacharunderscore}{\kern0pt}seq\ {\isasympi}s\ htps{\isachardoublequoteclose}\isanewline
%
\isadelimproof
\ \ %
\endisadelimproof
%
\isatagproof
\isakeywordONE{proof}\isamarkupfalse%
\ {\isacharparenleft}{\kern0pt}induction\ {\isasympi}s{\isacharparenright}{\kern0pt}\isanewline
\ \ \ \ \isakeywordTHREE{case}\isamarkupfalse%
\ Nil\isanewline
\ \ \ \ \isakeywordONE{then}\isamarkupfalse%
\ \isakeywordONE{have}\isamarkupfalse%
\ {\isachardoublequoteopen}htps{\isacharunderscore}{\kern0pt}seq\ {\isacharbrackleft}{\kern0pt}{\isacharbrackright}{\kern0pt}\ {\isacharbrackleft}{\kern0pt}{\isacharbrackright}{\kern0pt}{\isachardoublequoteclose}\isanewline
\ \ \ \ \ \ \isakeywordONE{unfolding}\isamarkupfalse%
\ htps{\isacharunderscore}{\kern0pt}seq{\isacharunderscore}{\kern0pt}def\ is{\isacharunderscore}{\kern0pt}htp{\isacharunderscore}{\kern0pt}def\ \isakeywordONE{using}\isamarkupfalse%
\ dur{\isacharunderscore}{\kern0pt}acts{\isacharunderscore}{\kern0pt}subset\ \isakeywordONE{by}\isamarkupfalse%
\ fastforce\isanewline
\ \ \ \ \isakeywordONE{then}\isamarkupfalse%
\ \isakeywordTHREE{show}\isamarkupfalse%
\ {\isacharquery}{\kern0pt}case\ \isanewline
\ \ \ \ \ \ \isakeywordONE{by}\isamarkupfalse%
\ auto\isanewline
\ \ \isakeywordONE{next}\isamarkupfalse%
\isanewline
\ \ \ \ \isakeywordTHREE{case}\isamarkupfalse%
\ {\isacharparenleft}{\kern0pt}Cons\ {\isasympi}\ {\isasympi}s{\isacharparenright}{\kern0pt}\isanewline
\ \ \ \ \isakeywordONE{then}\isamarkupfalse%
\ \isakeywordTHREE{show}\isamarkupfalse%
\ {\isacharquery}{\kern0pt}case\isanewline
\ \ \ \ \isakeywordONE{proof}\isamarkupfalse%
\ {\isacharparenleft}{\kern0pt}cases\ {\isasympi}{\isacharparenright}{\kern0pt}\isanewline
\ \ \ \ \ \ \isakeywordTHREE{case}\isamarkupfalse%
\ {\isacharparenleft}{\kern0pt}Pair\ t\isactrlsub {\isasympi}\ {\isasympi}{\isacharprime}{\kern0pt}{\isacharparenright}{\kern0pt}\isanewline
\ \ \ \ \ \ \isakeywordTHREE{obtain}\isamarkupfalse%
\ htps{\isacharprime}{\kern0pt}\ \isakeywordTWO{where}\ {\isachardoublequoteopen}htps{\isacharunderscore}{\kern0pt}seq\ {\isasympi}s\ htps{\isacharprime}{\kern0pt}{\isachardoublequoteclose}\isanewline
\ \ \ \ \ \ \ \ \isakeywordONE{using}\isamarkupfalse%
\ Cons{\isachardot}{\kern0pt}IH\ \isakeywordONE{by}\isamarkupfalse%
\ auto\isanewline
\ \ \ \ \ \ \isakeywordONE{then}\isamarkupfalse%
\ \isakeywordONE{have}\isamarkupfalse%
\ {\isachardoublequoteopen}sorted\ htps{\isacharprime}{\kern0pt}{\isachardoublequoteclose}\ \isakeywordTWO{and}\ {\isachardoublequoteopen}distinct\ htps{\isacharprime}{\kern0pt}{\isachardoublequoteclose}\isanewline
\ \ \ \ \ \ \ \ \isakeywordONE{unfolding}\isamarkupfalse%
\ htps{\isacharunderscore}{\kern0pt}seq{\isacharunderscore}{\kern0pt}def\ \isakeywordONE{by}\isamarkupfalse%
\ {\isacharparenleft}{\kern0pt}auto\ simp{\isacharcolon}{\kern0pt}\ strict{\isacharunderscore}{\kern0pt}sorted{\isacharunderscore}{\kern0pt}iff{\isacharparenright}{\kern0pt}\isanewline
\isanewline
\ \ \ \ \ \ \isakeywordONE{have}\isamarkupfalse%
\ {\isachardoublequoteopen}{\isasymforall}t{\isachardot}{\kern0pt}\ t\ {\isasymin}\ set\ htps{\isacharprime}{\kern0pt}\ {\isasymlongleftrightarrow}\ is{\isacharunderscore}{\kern0pt}htp\ {\isasympi}s\ t{\isachardoublequoteclose}\isanewline
\ \ \ \ \ \ \ \ \isakeywordONE{using}\isamarkupfalse%
\ {\isacartoucheopen}htps{\isacharunderscore}{\kern0pt}seq\ {\isasympi}s\ htps{\isacharprime}{\kern0pt}{\isacartoucheclose}\ \isakeywordONE{unfolding}\isamarkupfalse%
\ htps{\isacharunderscore}{\kern0pt}seq{\isacharunderscore}{\kern0pt}def\ \isakeywordONE{by}\isamarkupfalse%
\ auto\isanewline
\ \ \ \ \ \ \isakeywordONE{then}\isamarkupfalse%
\ \isakeywordONE{have}\isamarkupfalse%
\ {\isachardoublequoteopen}{\isasymforall}t{\isachardot}{\kern0pt}\ t\ {\isasymin}\ set\ htps{\isacharprime}{\kern0pt}\ {\isasymlongrightarrow}\ is{\isacharunderscore}{\kern0pt}htp\ {\isacharparenleft}{\kern0pt}{\isasympi}{\isacharhash}{\kern0pt}{\isasympi}s{\isacharparenright}{\kern0pt}\ t{\isachardoublequoteclose}\isanewline
\ \ \ \ \ \ \ \ \isakeywordONE{unfolding}\isamarkupfalse%
\ is{\isacharunderscore}{\kern0pt}htp{\isacharunderscore}{\kern0pt}def\ durative{\isacharunderscore}{\kern0pt}acts{\isacharunderscore}{\kern0pt}def\ is{\isacharunderscore}{\kern0pt}act{\isacharunderscore}{\kern0pt}simple{\isacharunderscore}{\kern0pt}def\ \isakeywordONE{by}\isamarkupfalse%
\ auto\isanewline
\ \ \ \ \ \ \isakeywordONE{have}\isamarkupfalse%
\ {\isachardoublequoteopen}is{\isacharunderscore}{\kern0pt}htp\ {\isacharparenleft}{\kern0pt}{\isasympi}{\isacharhash}{\kern0pt}{\isasympi}s{\isacharparenright}{\kern0pt}\ t\isactrlsub {\isasympi}{\isachardoublequoteclose}\isanewline
\ \ \ \ \ \ \ \ \isakeywordONE{unfolding}\isamarkupfalse%
\ is{\isacharunderscore}{\kern0pt}htp{\isacharunderscore}{\kern0pt}def\ \isakeywordONE{using}\isamarkupfalse%
\ Pair\ \isakeywordONE{by}\isamarkupfalse%
\ auto\isanewline
\ \ \ \ \ \ \isakeywordTHREE{show}\isamarkupfalse%
\ {\isacharquery}{\kern0pt}thesis\isanewline
\ \ \ \ \ \ \isakeywordONE{proof}\isamarkupfalse%
\ {\isacharparenleft}{\kern0pt}cases\ {\isasympi}{\isacharprime}{\kern0pt}{\isacharparenright}{\kern0pt}\isanewline
\ \ \ \ \ \ \ \ \isakeywordTHREE{case}\isamarkupfalse%
\ {\isacharparenleft}{\kern0pt}Simple{\isacharunderscore}{\kern0pt}Plan{\isacharunderscore}{\kern0pt}Action\ n\ args{\isacharparenright}{\kern0pt}\isanewline
\ \ \ \ \ \ \ \ \isakeywordONE{then}\isamarkupfalse%
\ \isakeywordTHREE{obtain}\isamarkupfalse%
\ htps\ \isakeywordTWO{where}\ {\isachardoublequoteopen}htps\ {\isacharequal}{\kern0pt}\ insort{\isacharunderscore}{\kern0pt}insert\ t\isactrlsub {\isasympi}\ htps{\isacharprime}{\kern0pt}{\isachardoublequoteclose}\isanewline
\ \ \ \ \ \ \ \ \ \ \isakeywordONE{by}\isamarkupfalse%
\ auto\isanewline
\ \ \ \ \ \ \ \ \isakeywordONE{then}\isamarkupfalse%
\ \isakeywordONE{have}\isamarkupfalse%
\ {\isachardoublequoteopen}strict{\isacharunderscore}{\kern0pt}sorted\ htps{\isachardoublequoteclose}\isanewline
\ \ \ \ \ \ \ \ \ \ \isakeywordONE{using}\isamarkupfalse%
\ sorted{\isacharunderscore}{\kern0pt}insort{\isacharunderscore}{\kern0pt}insert{\isacharbrackleft}{\kern0pt}OF\ {\isacartoucheopen}sorted\ htps{\isacharprime}{\kern0pt}{\isacartoucheclose}{\isacharbrackright}{\kern0pt}\ \isanewline
\ \ \ \ \ \ \ \ \ \ \ \ \ \ \ \ distinct{\isacharunderscore}{\kern0pt}insort{\isacharunderscore}{\kern0pt}insert{\isacharbrackleft}{\kern0pt}OF\ {\isacartoucheopen}distinct\ htps{\isacharprime}{\kern0pt}{\isacartoucheclose}{\isacharbrackright}{\kern0pt}\isanewline
\ \ \ \ \ \ \ \ \ \ \isakeywordONE{by}\isamarkupfalse%
\ {\isacharparenleft}{\kern0pt}auto\ simp{\isacharcolon}{\kern0pt}\ strict{\isacharunderscore}{\kern0pt}sorted{\isacharunderscore}{\kern0pt}iff{\isacharparenright}{\kern0pt}\isanewline
\ \ \ \ \ \ \ \ \isakeywordONE{have}\isamarkupfalse%
\ {\isachardoublequoteopen}{\isasymforall}t{\isachardot}{\kern0pt}\ is{\isacharunderscore}{\kern0pt}htp\ {\isasympi}s\ t\ {\isasymlongrightarrow}\ t\ {\isasymin}\ set\ htps{\isachardoublequoteclose}\ \isakeywordTWO{and}\ {\isachardoublequoteopen}t\isactrlsub {\isasympi}\ {\isasymin}\ set\ htps{\isachardoublequoteclose}\isanewline
\ \ \ \ \ \ \ \ \ \ \isakeywordONE{using}\isamarkupfalse%
\ {\isacartoucheopen}{\isasymforall}t{\isachardot}{\kern0pt}\ t\ {\isasymin}\ set\ htps{\isacharprime}{\kern0pt}\ {\isasymlongleftrightarrow}\ is{\isacharunderscore}{\kern0pt}htp\ {\isasympi}s\ t{\isacartoucheclose}\ \isanewline
\ \ \ \ \ \ \ \ \ \ \ \ \ \ \ \ {\isacartoucheopen}htps\ {\isacharequal}{\kern0pt}\ insort{\isacharunderscore}{\kern0pt}insert\ t\isactrlsub {\isasympi}\ htps{\isacharprime}{\kern0pt}{\isacartoucheclose}\ \isanewline
\ \ \ \ \ \ \ \ \ \ \isakeywordONE{by}\isamarkupfalse%
\ {\isacharparenleft}{\kern0pt}auto\ simp{\isacharcolon}{\kern0pt}\ set{\isacharunderscore}{\kern0pt}insort{\isacharunderscore}{\kern0pt}insert{\isacharparenright}{\kern0pt}\isanewline
\ \ \ \ \ \ \ \ \isakeywordONE{then}\isamarkupfalse%
\ \isakeywordONE{have}\isamarkupfalse%
\ {\isachardoublequoteopen}{\isasymforall}t{\isachardot}{\kern0pt}\ is{\isacharunderscore}{\kern0pt}htp\ {\isacharparenleft}{\kern0pt}{\isasympi}{\isacharhash}{\kern0pt}{\isasympi}s{\isacharparenright}{\kern0pt}\ t\ {\isasymlongrightarrow}\ t\ {\isasymin}\ set\ htps{\isachardoublequoteclose}\isanewline
\ \ \ \ \ \ \ \ \ \ \isakeywordONE{unfolding}\isamarkupfalse%
\ is{\isacharunderscore}{\kern0pt}htp{\isacharunderscore}{\kern0pt}def\ durative{\isacharunderscore}{\kern0pt}acts{\isacharunderscore}{\kern0pt}def\ is{\isacharunderscore}{\kern0pt}act{\isacharunderscore}{\kern0pt}simple{\isacharunderscore}{\kern0pt}def\isanewline
\ \ \ \ \ \ \ \ \ \ \isakeywordONE{using}\isamarkupfalse%
\ Pair\ Simple{\isacharunderscore}{\kern0pt}Plan{\isacharunderscore}{\kern0pt}Action\ \isakeywordONE{by}\isamarkupfalse%
\ auto\isanewline
\ \ \ \ \ \ \ \ \isakeywordONE{have}\isamarkupfalse%
\ {\isachardoublequoteopen}{\isasymforall}t{\isachardot}{\kern0pt}\ t\ {\isasymin}\ set\ htps\ {\isasymlongrightarrow}\ is{\isacharunderscore}{\kern0pt}htp\ {\isacharparenleft}{\kern0pt}{\isasympi}{\isacharhash}{\kern0pt}{\isasympi}s{\isacharparenright}{\kern0pt}\ t{\isachardoublequoteclose}\isanewline
\ \ \ \ \ \ \ \ \ \ \isakeywordONE{using}\isamarkupfalse%
\ {\isacartoucheopen}{\isasymforall}t{\isachardot}{\kern0pt}\ t\ {\isasymin}\ set\ htps{\isacharprime}{\kern0pt}\ {\isasymlongrightarrow}\ is{\isacharunderscore}{\kern0pt}htp\ {\isacharparenleft}{\kern0pt}{\isasympi}{\isacharhash}{\kern0pt}{\isasympi}s{\isacharparenright}{\kern0pt}\ t{\isacartoucheclose}\ \isanewline
\ \ \ \ \ \ \ \ \ \ \ \ \ \ \ \ {\isacartoucheopen}htps\ {\isacharequal}{\kern0pt}\ insort{\isacharunderscore}{\kern0pt}insert\ t\isactrlsub {\isasympi}\ htps{\isacharprime}{\kern0pt}{\isacartoucheclose}\ \isanewline
\ \ \ \ \ \ \ \ \ \ \ \ \ \ \ \ {\isacartoucheopen}is{\isacharunderscore}{\kern0pt}htp\ {\isacharparenleft}{\kern0pt}{\isasympi}{\isacharhash}{\kern0pt}{\isasympi}s{\isacharparenright}{\kern0pt}\ t\isactrlsub {\isasympi}{\isacartoucheclose}\isanewline
\ \ \ \ \ \ \ \ \ \ \isakeywordONE{by}\isamarkupfalse%
\ {\isacharparenleft}{\kern0pt}auto\ simp{\isacharcolon}{\kern0pt}\ set{\isacharunderscore}{\kern0pt}insort{\isacharunderscore}{\kern0pt}insert{\isacharparenright}{\kern0pt}\isanewline
\ \ \ \ \ \ \ \ \isakeywordTHREE{show}\isamarkupfalse%
\ {\isacharquery}{\kern0pt}thesis\ \isanewline
\ \ \ \ \ \ \ \ \ \ \isakeywordONE{using}\isamarkupfalse%
\ {\isacartoucheopen}strict{\isacharunderscore}{\kern0pt}sorted\ htps{\isacartoucheclose}\ \isanewline
\ \ \ \ \ \ \ \ \ \ \ \ \ \ \ \ {\isacartoucheopen}{\isasymforall}t{\isachardot}{\kern0pt}\ is{\isacharunderscore}{\kern0pt}htp\ {\isacharparenleft}{\kern0pt}{\isasympi}{\isacharhash}{\kern0pt}{\isasympi}s{\isacharparenright}{\kern0pt}\ t\ {\isasymlongrightarrow}\ t\ {\isasymin}\ set\ htps{\isacartoucheclose}\ \isanewline
\ \ \ \ \ \ \ \ \ \ \ \ \ \ \ \ {\isacartoucheopen}{\isasymforall}t{\isachardot}{\kern0pt}\ t\ {\isasymin}\ set\ htps\ {\isasymlongrightarrow}\ is{\isacharunderscore}{\kern0pt}htp\ {\isacharparenleft}{\kern0pt}{\isasympi}{\isacharhash}{\kern0pt}{\isasympi}s{\isacharparenright}{\kern0pt}\ t{\isacartoucheclose}\isanewline
\ \ \ \ \ \ \ \ \ \ \isakeywordONE{by}\isamarkupfalse%
\ {\isacharparenleft}{\kern0pt}auto\ simp{\isacharcolon}{\kern0pt}\ htps{\isacharunderscore}{\kern0pt}seq{\isacharunderscore}{\kern0pt}def{\isacharparenright}{\kern0pt}\isanewline
\ \ \ \ \ \ \isakeywordONE{next}\isamarkupfalse%
\isanewline
\ \ \ \ \ \ \ \ \isakeywordTHREE{case}\isamarkupfalse%
\ {\isacharparenleft}{\kern0pt}Durative{\isacharunderscore}{\kern0pt}Plan{\isacharunderscore}{\kern0pt}Action\ n\ args\ d{\isacharparenright}{\kern0pt}\isanewline
\ \ \ \ \ \ \ \ \isakeywordONE{then}\isamarkupfalse%
\ \isakeywordONE{have}\isamarkupfalse%
\ {\isachardoublequoteopen}is{\isacharunderscore}{\kern0pt}htp\ {\isacharparenleft}{\kern0pt}{\isasympi}{\isacharhash}{\kern0pt}{\isasympi}s{\isacharparenright}{\kern0pt}\ {\isacharparenleft}{\kern0pt}t\isactrlsub {\isasympi}\ {\isacharplus}{\kern0pt}\ d{\isacharparenright}{\kern0pt}{\isachardoublequoteclose}\isanewline
\ \ \ \ \ \ \ \ \ \ \isakeywordONE{unfolding}\isamarkupfalse%
\ is{\isacharunderscore}{\kern0pt}htp{\isacharunderscore}{\kern0pt}def\ durative{\isacharunderscore}{\kern0pt}acts{\isacharunderscore}{\kern0pt}def\ is{\isacharunderscore}{\kern0pt}act{\isacharunderscore}{\kern0pt}simple{\isacharunderscore}{\kern0pt}def\ \isakeywordONE{using}\isamarkupfalse%
\ Pair\ \isakeywordONE{by}\isamarkupfalse%
\ auto\isanewline
\ \ \ \ \ \ \ \ \isakeywordTHREE{obtain}\isamarkupfalse%
\ htps\ \isakeywordTWO{where}\ {\isachardoublequoteopen}htps\ {\isacharequal}{\kern0pt}\ insort{\isacharunderscore}{\kern0pt}insert\ t\isactrlsub {\isasympi}\ {\isacharparenleft}{\kern0pt}insort{\isacharunderscore}{\kern0pt}insert\ {\isacharparenleft}{\kern0pt}t\isactrlsub {\isasympi}\ {\isacharplus}{\kern0pt}\ d{\isacharparenright}{\kern0pt}\ htps{\isacharprime}{\kern0pt}{\isacharparenright}{\kern0pt}{\isachardoublequoteclose}\isanewline
\ \ \ \ \ \ \ \ \ \ \isakeywordONE{by}\isamarkupfalse%
\ auto\isanewline
\ \ \ \ \ \ \ \ \isakeywordONE{then}\isamarkupfalse%
\ \isakeywordONE{have}\isamarkupfalse%
\ {\isachardoublequoteopen}strict{\isacharunderscore}{\kern0pt}sorted\ htps{\isachardoublequoteclose}\isanewline
\ \ \ \ \ \ \ \ \ \ \isakeywordONE{using}\isamarkupfalse%
\ {\isacartoucheopen}sorted\ htps{\isacharprime}{\kern0pt}{\isacartoucheclose}\ {\isacartoucheopen}distinct\ htps{\isacharprime}{\kern0pt}{\isacartoucheclose}\isanewline
\ \ \ \ \ \ \ \ \ \ \isakeywordONE{by}\isamarkupfalse%
\ {\isacharparenleft}{\kern0pt}auto\ simp{\isacharcolon}{\kern0pt}\ sorted{\isacharunderscore}{\kern0pt}insort{\isacharunderscore}{\kern0pt}insert\ distinct{\isacharunderscore}{\kern0pt}insort{\isacharunderscore}{\kern0pt}insert\ strict{\isacharunderscore}{\kern0pt}sorted{\isacharunderscore}{\kern0pt}iff{\isacharparenright}{\kern0pt}\isanewline
\ \ \ \ \ \ \ \ \isakeywordONE{have}\isamarkupfalse%
\ {\isachardoublequoteopen}{\isasymforall}t{\isachardot}{\kern0pt}\ is{\isacharunderscore}{\kern0pt}htp\ {\isasympi}s\ t\ {\isasymlongrightarrow}\ t\ {\isasymin}\ set\ htps{\isachardoublequoteclose}\ \isakeywordTWO{and}\ {\isachardoublequoteopen}t\isactrlsub {\isasympi}\ {\isasymin}\ set\ htps{\isachardoublequoteclose}\ \isakeywordTWO{and}\ {\isachardoublequoteopen}{\isacharparenleft}{\kern0pt}t\isactrlsub {\isasympi}\ {\isacharplus}{\kern0pt}\ d{\isacharparenright}{\kern0pt}\ {\isasymin}\ set\ htps{\isachardoublequoteclose}\isanewline
\ \ \ \ \ \ \ \ \ \ \isakeywordONE{using}\isamarkupfalse%
\ {\isacartoucheopen}{\isasymforall}t{\isachardot}{\kern0pt}\ t\ {\isasymin}\ set\ htps{\isacharprime}{\kern0pt}\ {\isasymlongleftrightarrow}\ is{\isacharunderscore}{\kern0pt}htp\ {\isasympi}s\ t{\isacartoucheclose}\ \isanewline
\ \ \ \ \ \ \ \ \ \ \ \ \ \ \ \ {\isacartoucheopen}htps\ {\isacharequal}{\kern0pt}\ insort{\isacharunderscore}{\kern0pt}insert\ t\isactrlsub {\isasympi}\ {\isacharparenleft}{\kern0pt}insort{\isacharunderscore}{\kern0pt}insert\ {\isacharparenleft}{\kern0pt}t\isactrlsub {\isasympi}\ {\isacharplus}{\kern0pt}\ d{\isacharparenright}{\kern0pt}\ htps{\isacharprime}{\kern0pt}{\isacharparenright}{\kern0pt}{\isacartoucheclose}\ \isanewline
\ \ \ \ \ \ \ \ \ \ \isakeywordONE{by}\isamarkupfalse%
\ {\isacharparenleft}{\kern0pt}auto\ simp{\isacharcolon}{\kern0pt}\ set{\isacharunderscore}{\kern0pt}insort{\isacharunderscore}{\kern0pt}insert{\isacharparenright}{\kern0pt}\isanewline
\ \ \ \ \ \ \ \ \isakeywordONE{then}\isamarkupfalse%
\ \isakeywordONE{have}\isamarkupfalse%
\ {\isachardoublequoteopen}{\isasymforall}t{\isachardot}{\kern0pt}\ is{\isacharunderscore}{\kern0pt}htp\ {\isacharparenleft}{\kern0pt}{\isasympi}{\isacharhash}{\kern0pt}{\isasympi}s{\isacharparenright}{\kern0pt}\ t\ {\isasymlongrightarrow}\ t\ {\isasymin}\ set\ htps{\isachardoublequoteclose}\isanewline
\ \ \ \ \ \ \ \ \ \ \isakeywordONE{unfolding}\isamarkupfalse%
\ is{\isacharunderscore}{\kern0pt}htp{\isacharunderscore}{\kern0pt}def\ durative{\isacharunderscore}{\kern0pt}acts{\isacharunderscore}{\kern0pt}def\ is{\isacharunderscore}{\kern0pt}act{\isacharunderscore}{\kern0pt}simple{\isacharunderscore}{\kern0pt}def\isanewline
\ \ \ \ \ \ \ \ \ \ \isakeywordONE{using}\isamarkupfalse%
\ Pair\ Durative{\isacharunderscore}{\kern0pt}Plan{\isacharunderscore}{\kern0pt}Action\ \isakeywordONE{by}\isamarkupfalse%
\ auto\isanewline
\ \ \ \ \ \ \ \ \isakeywordONE{have}\isamarkupfalse%
\ {\isachardoublequoteopen}{\isasymforall}t{\isachardot}{\kern0pt}\ t\ {\isasymin}\ set\ htps\ {\isasymlongrightarrow}\ is{\isacharunderscore}{\kern0pt}htp\ {\isacharparenleft}{\kern0pt}{\isasympi}{\isacharhash}{\kern0pt}{\isasympi}s{\isacharparenright}{\kern0pt}\ t{\isachardoublequoteclose}\isanewline
\ \ \ \ \ \ \ \ \ \ \isakeywordONE{using}\isamarkupfalse%
\ {\isacartoucheopen}{\isasymforall}t{\isachardot}{\kern0pt}\ t\ {\isasymin}\ set\ htps{\isacharprime}{\kern0pt}\ {\isasymlongrightarrow}\ is{\isacharunderscore}{\kern0pt}htp\ {\isacharparenleft}{\kern0pt}{\isasympi}{\isacharhash}{\kern0pt}{\isasympi}s{\isacharparenright}{\kern0pt}\ t{\isacartoucheclose}\ \isanewline
\ \ \ \ \ \ \ \ \ \ \ \ \ \ \ \ {\isacartoucheopen}htps\ {\isacharequal}{\kern0pt}\ insort{\isacharunderscore}{\kern0pt}insert\ t\isactrlsub {\isasympi}\ {\isacharparenleft}{\kern0pt}insort{\isacharunderscore}{\kern0pt}insert\ {\isacharparenleft}{\kern0pt}t\isactrlsub {\isasympi}\ {\isacharplus}{\kern0pt}\ d{\isacharparenright}{\kern0pt}\ htps{\isacharprime}{\kern0pt}{\isacharparenright}{\kern0pt}{\isacartoucheclose}\ \isanewline
\ \ \ \ \ \ \ \ \ \ \ \ \ \ \ \ {\isacartoucheopen}is{\isacharunderscore}{\kern0pt}htp\ {\isacharparenleft}{\kern0pt}{\isasympi}{\isacharhash}{\kern0pt}{\isasympi}s{\isacharparenright}{\kern0pt}\ t\isactrlsub {\isasympi}{\isacartoucheclose}\ {\isacartoucheopen}is{\isacharunderscore}{\kern0pt}htp\ {\isacharparenleft}{\kern0pt}{\isasympi}{\isacharhash}{\kern0pt}{\isasympi}s{\isacharparenright}{\kern0pt}\ {\isacharparenleft}{\kern0pt}t\isactrlsub {\isasympi}\ {\isacharplus}{\kern0pt}\ d{\isacharparenright}{\kern0pt}{\isacartoucheclose}\isanewline
\ \ \ \ \ \ \ \ \ \ \isakeywordONE{by}\isamarkupfalse%
\ {\isacharparenleft}{\kern0pt}auto\ simp{\isacharcolon}{\kern0pt}\ set{\isacharunderscore}{\kern0pt}insort{\isacharunderscore}{\kern0pt}insert{\isacharparenright}{\kern0pt}\isanewline
\ \ \ \ \ \ \ \ \isakeywordTHREE{show}\isamarkupfalse%
\ {\isacharquery}{\kern0pt}thesis\ \isanewline
\ \ \ \ \ \ \ \ \ \ \isakeywordONE{using}\isamarkupfalse%
\ {\isacartoucheopen}strict{\isacharunderscore}{\kern0pt}sorted\ htps{\isacartoucheclose}\ \isanewline
\ \ \ \ \ \ \ \ \ \ \ \ \ \ \ \ {\isacartoucheopen}{\isasymforall}t{\isachardot}{\kern0pt}\ is{\isacharunderscore}{\kern0pt}htp\ {\isacharparenleft}{\kern0pt}{\isasympi}{\isacharhash}{\kern0pt}{\isasympi}s{\isacharparenright}{\kern0pt}\ t\ {\isasymlongrightarrow}\ t\ {\isasymin}\ set\ htps{\isacartoucheclose}\ \isanewline
\ \ \ \ \ \ \ \ \ \ \ \ \ \ \ \ {\isacartoucheopen}{\isasymforall}t{\isachardot}{\kern0pt}\ t\ {\isasymin}\ set\ htps\ {\isasymlongrightarrow}\ is{\isacharunderscore}{\kern0pt}htp\ {\isacharparenleft}{\kern0pt}{\isasympi}{\isacharhash}{\kern0pt}{\isasympi}s{\isacharparenright}{\kern0pt}\ t{\isacartoucheclose}\isanewline
\ \ \ \ \ \ \ \ \ \ \isakeywordONE{by}\isamarkupfalse%
\ {\isacharparenleft}{\kern0pt}auto\ simp{\isacharcolon}{\kern0pt}\ htps{\isacharunderscore}{\kern0pt}seq{\isacharunderscore}{\kern0pt}def{\isacharparenright}{\kern0pt}\isanewline
\ \ \ \ \ \ \isakeywordONE{qed}\isamarkupfalse%
\isanewline
\ \ \ \ \isakeywordONE{qed}\isamarkupfalse%
\isanewline
\ \ \isakeywordONE{qed}\isamarkupfalse%
%
\endisatagproof
{\isafoldproof}%
%
\isadelimproof
\isanewline
%
\endisadelimproof
\isanewline
\isakeywordTWO{end}\isamarkupfalse%
\isanewline
\isanewline
\isakeywordONE{context}\isamarkupfalse%
\ wf{\isacharunderscore}{\kern0pt}ast{\isacharunderscore}{\kern0pt}problem\isanewline
\isakeywordTWO{begin}%
\begin{isamarkuptext}%
Two induced happening sequences contain the same ground actions in between two 
  consecutive happening time points.%
\end{isamarkuptext}\isamarkuptrue%
\ \ \isakeywordONE{lemma}\isamarkupfalse%
\ ind{\isacharunderscore}{\kern0pt}happ{\isacharunderscore}{\kern0pt}seqs{\isacharunderscore}{\kern0pt}same{\isacharunderscore}{\kern0pt}acts{\isacharunderscore}{\kern0pt}leq\isactrlsub i{\isacharunderscore}{\kern0pt}le\isactrlsub j{\isacharcolon}{\kern0pt}\isanewline
\ \ \ \ \isakeywordTWO{assumes}\ {\isachardoublequoteopen}consec{\isacharunderscore}{\kern0pt}htps\ {\isasympi}s\ t\isactrlsub i\ t\isactrlsub j{\isachardoublequoteclose}\ \isanewline
\ \ \ \ \ \ \ \ \isakeywordTWO{and}\ {\isachardoublequoteopen}ind{\isacharunderscore}{\kern0pt}happ{\isacharunderscore}{\kern0pt}seq\ {\isasympi}s\ hs\isactrlsub {\isadigit{1}}{\isachardoublequoteclose}\ \isanewline
\ \ \ \ \ \ \ \ \isakeywordTWO{and}\ {\isachardoublequoteopen}ind{\isacharunderscore}{\kern0pt}happ{\isacharunderscore}{\kern0pt}seq\ {\isasympi}s\ hs\isactrlsub {\isadigit{2}}{\isachardoublequoteclose}\isanewline
\ \ \ \ \ \ \ \ \isakeywordTWO{and}\ {\isachardoublequoteopen}{\isacharparenleft}{\kern0pt}t\isactrlsub a{\isacharcomma}{\kern0pt}A{\isacharparenright}{\kern0pt}\ {\isasymin}\ set\ {\isacharparenleft}{\kern0pt}dropWhile\ {\isacharparenleft}{\kern0pt}leq\ t\isactrlsub i{\isacharparenright}{\kern0pt}\ {\isacharparenleft}{\kern0pt}takeWhile\ {\isacharparenleft}{\kern0pt}le\ t\isactrlsub j{\isacharparenright}{\kern0pt}\ hs\isactrlsub {\isadigit{1}}{\isacharparenright}{\kern0pt}{\isacharparenright}{\kern0pt}{\isachardoublequoteclose}\isanewline
\ \ \ \ \ \ \ \ \isakeywordTWO{and}\ {\isachardoublequoteopen}a\ {\isasymin}\ set\ A{\isachardoublequoteclose}\isanewline
\ \ \ \ \ \ \isakeywordTWO{shows}\ {\isachardoublequoteopen}{\isasymexists}t\isactrlsub a{\isacharprime}{\kern0pt}\ A{\isacharprime}{\kern0pt}{\isachardot}{\kern0pt}\ {\isacharparenleft}{\kern0pt}t\isactrlsub a{\isacharprime}{\kern0pt}{\isacharcomma}{\kern0pt}A{\isacharprime}{\kern0pt}{\isacharparenright}{\kern0pt}\ {\isasymin}\ set\ {\isacharparenleft}{\kern0pt}dropWhile\ {\isacharparenleft}{\kern0pt}leq\ t\isactrlsub i{\isacharparenright}{\kern0pt}\ {\isacharparenleft}{\kern0pt}takeWhile\ {\isacharparenleft}{\kern0pt}le\ t\isactrlsub j{\isacharparenright}{\kern0pt}\ hs\isactrlsub {\isadigit{2}}{\isacharparenright}{\kern0pt}{\isacharparenright}{\kern0pt}\ {\isasymand}\ a\ {\isasymin}\ set\ A{\isacharprime}{\kern0pt}{\isachardoublequoteclose}\isanewline
%
\isadelimproof
\ \ %
\endisadelimproof
%
\isatagproof
\isakeywordONE{proof}\isamarkupfalse%
\ {\isacharminus}{\kern0pt}\isanewline
\ \ \ \ \isakeywordONE{let}\isamarkupfalse%
\ {\isacharquery}{\kern0pt}hs\isactrlsub {\isadigit{1}}{\isacharprime}{\kern0pt}{\isacharequal}{\kern0pt}{\isachardoublequoteopen}dropWhile\ {\isacharparenleft}{\kern0pt}leq\ t\isactrlsub i{\isacharparenright}{\kern0pt}\ {\isacharparenleft}{\kern0pt}takeWhile\ {\isacharparenleft}{\kern0pt}le\ t\isactrlsub j{\isacharparenright}{\kern0pt}\ hs\isactrlsub {\isadigit{1}}{\isacharparenright}{\kern0pt}{\isachardoublequoteclose}\isanewline
\ \ \ \ \isakeywordONE{let}\isamarkupfalse%
\ {\isacharquery}{\kern0pt}hs\isactrlsub {\isadigit{2}}{\isacharprime}{\kern0pt}{\isacharequal}{\kern0pt}{\isachardoublequoteopen}dropWhile\ {\isacharparenleft}{\kern0pt}leq\ t\isactrlsub i{\isacharparenright}{\kern0pt}\ {\isacharparenleft}{\kern0pt}takeWhile\ {\isacharparenleft}{\kern0pt}le\ t\isactrlsub j{\isacharparenright}{\kern0pt}\ hs\isactrlsub {\isadigit{2}}{\isacharparenright}{\kern0pt}{\isachardoublequoteclose}\isanewline
\ \ \ \ \isakeywordONE{have}\isamarkupfalse%
\ {\isachardoublequoteopen}strict{\isacharunderscore}{\kern0pt}sorted\ {\isacharparenleft}{\kern0pt}map\ fst\ hs\isactrlsub {\isadigit{1}}{\isacharparenright}{\kern0pt}{\isachardoublequoteclose}\ \isakeywordTWO{and}\ {\isachardoublequoteopen}strict{\isacharunderscore}{\kern0pt}sorted\ {\isacharparenleft}{\kern0pt}map\ fst\ hs\isactrlsub {\isadigit{2}}{\isacharparenright}{\kern0pt}{\isachardoublequoteclose}\isanewline
\ \ \ \ \ \ \isakeywordONE{using}\isamarkupfalse%
\ {\isacartoucheopen}ind{\isacharunderscore}{\kern0pt}happ{\isacharunderscore}{\kern0pt}seq\ {\isasympi}s\ hs\isactrlsub {\isadigit{1}}{\isacartoucheclose}\ {\isacartoucheopen}ind{\isacharunderscore}{\kern0pt}happ{\isacharunderscore}{\kern0pt}seq\ {\isasympi}s\ hs\isactrlsub {\isadigit{2}}{\isacartoucheclose}\ \isakeywordONE{unfolding}\isamarkupfalse%
\ ind{\isacharunderscore}{\kern0pt}happ{\isacharunderscore}{\kern0pt}seq{\isacharunderscore}{\kern0pt}def\ \isakeywordONE{by}\isamarkupfalse%
\ auto\isanewline
\ \ \ \ \isakeywordONE{have}\isamarkupfalse%
\ {\isachardoublequoteopen}{\isacharparenleft}{\kern0pt}t\isactrlsub a{\isacharcomma}{\kern0pt}A{\isacharparenright}{\kern0pt}\ {\isasymin}\ set\ hs\isactrlsub {\isadigit{1}}{\isachardoublequoteclose}\isanewline
\ \ \ \ \ \ \isakeywordONE{using}\isamarkupfalse%
\ {\isacartoucheopen}{\isacharparenleft}{\kern0pt}t\isactrlsub a{\isacharcomma}{\kern0pt}A{\isacharparenright}{\kern0pt}\ {\isasymin}\ set\ {\isacharquery}{\kern0pt}hs\isactrlsub {\isadigit{1}}{\isacharprime}{\kern0pt}{\isacartoucheclose}\ \isakeywordONE{by}\isamarkupfalse%
\ {\isacharparenleft}{\kern0pt}meson\ set{\isacharunderscore}{\kern0pt}dropWhileD\ set{\isacharunderscore}{\kern0pt}takeWhileD{\isacharparenright}{\kern0pt}\ \isanewline
\ \ \ \ \isakeywordONE{have}\isamarkupfalse%
\ {\isachardoublequoteopen}t\isactrlsub i\ {\isacharless}{\kern0pt}\ t\isactrlsub a\ {\isasymand}\ t\isactrlsub a\ {\isacharless}{\kern0pt}\ t\isactrlsub j{\isachardoublequoteclose}\isanewline
\ \ \ \ \ \ \isakeywordONE{using}\isamarkupfalse%
\ {\isacartoucheopen}{\isacharparenleft}{\kern0pt}t\isactrlsub a{\isacharcomma}{\kern0pt}A{\isacharparenright}{\kern0pt}\ {\isasymin}\ set\ {\isacharquery}{\kern0pt}hs\isactrlsub {\isadigit{1}}{\isacharprime}{\kern0pt}{\isacartoucheclose}\isanewline
\ \ \ \ \ \ \ \ \ \ \ \ strict{\isacharunderscore}{\kern0pt}sorted{\isacharunderscore}{\kern0pt}drop{\isacharunderscore}{\kern0pt}leq{\isacharunderscore}{\kern0pt}take{\isacharunderscore}{\kern0pt}le{\isacharbrackleft}{\kern0pt}OF\ {\isacartoucheopen}strict{\isacharunderscore}{\kern0pt}sorted\ {\isacharparenleft}{\kern0pt}map\ fst\ hs\isactrlsub {\isadigit{1}}{\isacharparenright}{\kern0pt}{\isacartoucheclose}\ {\isacartoucheopen}{\isacharparenleft}{\kern0pt}t\isactrlsub a{\isacharcomma}{\kern0pt}A{\isacharparenright}{\kern0pt}\ {\isasymin}\ set\ hs\isactrlsub {\isadigit{1}}{\isacartoucheclose}{\isacharbrackright}{\kern0pt}\isanewline
\ \ \ \ \ \ \isakeywordONE{by}\isamarkupfalse%
\ auto\isanewline
\ \ \ \ \isakeywordONE{have}\isamarkupfalse%
\ {\isachardoublequoteopen}{\isasymexists}{\isacharparenleft}{\kern0pt}t\isactrlsub {\isasympi}{\isacharcomma}{\kern0pt}{\isasympi}{\isacharparenright}{\kern0pt}\ {\isasymin}\ set\ {\isasympi}s{\isachardot}{\kern0pt}\ inst{\isacharunderscore}{\kern0pt}of{\isacharunderscore}{\kern0pt}plan{\isacharunderscore}{\kern0pt}action\ {\isasympi}s\ {\isacharparenleft}{\kern0pt}t\isactrlsub {\isasympi}{\isacharcomma}{\kern0pt}{\isasympi}{\isacharparenright}{\kern0pt}\ {\isacharparenleft}{\kern0pt}t\isactrlsub a{\isacharcomma}{\kern0pt}a{\isacharparenright}{\kern0pt}{\isachardoublequoteclose}\isanewline
\ \ \ \ \ \ \isakeywordONE{using}\isamarkupfalse%
\ {\isacartoucheopen}ind{\isacharunderscore}{\kern0pt}happ{\isacharunderscore}{\kern0pt}seq\ {\isasympi}s\ hs\isactrlsub {\isadigit{1}}{\isacartoucheclose}\ {\isacartoucheopen}{\isacharparenleft}{\kern0pt}t\isactrlsub a{\isacharcomma}{\kern0pt}A{\isacharparenright}{\kern0pt}\ {\isasymin}\ set\ hs\isactrlsub {\isadigit{1}}{\isacartoucheclose}\ {\isacartoucheopen}a\ {\isasymin}\ set\ A{\isacartoucheclose}\ \isanewline
\ \ \ \ \ \ \isakeywordONE{unfolding}\isamarkupfalse%
\ ind{\isacharunderscore}{\kern0pt}happ{\isacharunderscore}{\kern0pt}seq{\isacharunderscore}{\kern0pt}def\ \isakeywordONE{by}\isamarkupfalse%
\ auto\isanewline
\ \ \ \ \isakeywordONE{then}\isamarkupfalse%
\ \isakeywordTHREE{obtain}\isamarkupfalse%
\ t\isactrlsub {\isasympi}\ {\isasympi}\ \isakeywordTWO{where}\ {\isachardoublequoteopen}{\isacharparenleft}{\kern0pt}t\isactrlsub {\isasympi}{\isacharcomma}{\kern0pt}{\isasympi}{\isacharparenright}{\kern0pt}\ {\isasymin}\ set\ {\isasympi}s{\isachardoublequoteclose}\ \isakeywordTWO{and}\ {\isachardoublequoteopen}inst{\isacharunderscore}{\kern0pt}of{\isacharunderscore}{\kern0pt}plan{\isacharunderscore}{\kern0pt}action\ {\isasympi}s\ {\isacharparenleft}{\kern0pt}t\isactrlsub {\isasympi}{\isacharcomma}{\kern0pt}{\isasympi}{\isacharparenright}{\kern0pt}\ {\isacharparenleft}{\kern0pt}t\isactrlsub a{\isacharcomma}{\kern0pt}a{\isacharparenright}{\kern0pt}{\isachardoublequoteclose}\isanewline
\ \ \ \ \ \ \isakeywordONE{by}\isamarkupfalse%
\ auto\isanewline
\ \ \ \ \isakeywordONE{let}\isamarkupfalse%
\ {\isacharquery}{\kern0pt}case{\isadigit{1}}{\isacharequal}{\kern0pt}{\isachardoublequoteopen}res{\isacharunderscore}{\kern0pt}inst\ {\isasympi}\ At{\isacharunderscore}{\kern0pt}Start\ {\isacharequal}{\kern0pt}\ Some\ a\ {\isasymand}\ t\isactrlsub {\isasympi}\ {\isacharequal}{\kern0pt}\ t\isactrlsub a{\isachardoublequoteclose}\isanewline
\ \ \ \ \isakeywordONE{let}\isamarkupfalse%
\ {\isacharquery}{\kern0pt}case{\isadigit{2}}{\isacharequal}{\kern0pt}{\isachardoublequoteopen}res{\isacharunderscore}{\kern0pt}inst{\isacharunderscore}{\kern0pt}snap{\isacharunderscore}{\kern0pt}action\ {\isasympi}\ At{\isacharunderscore}{\kern0pt}Start\ {\isacharequal}{\kern0pt}\ Some\ a\ {\isasymand}\ t\isactrlsub {\isasympi}\ {\isacharequal}{\kern0pt}\ t\isactrlsub a{\isachardoublequoteclose}\isanewline
\ \ \ \ \isakeywordONE{let}\isamarkupfalse%
\ {\isacharquery}{\kern0pt}case{\isadigit{3}}{\isacharequal}{\kern0pt}{\isachardoublequoteopen}res{\isacharunderscore}{\kern0pt}inst{\isacharunderscore}{\kern0pt}snap{\isacharunderscore}{\kern0pt}action\ {\isasympi}\ At{\isacharunderscore}{\kern0pt}End\ {\isacharequal}{\kern0pt}\ Some\ a\ {\isasymand}\ t\isactrlsub {\isasympi}\ {\isacharplus}{\kern0pt}\ {\isacharparenleft}{\kern0pt}duration\ {\isasympi}{\isacharparenright}{\kern0pt}\ {\isacharequal}{\kern0pt}\ t\isactrlsub a{\isachardoublequoteclose}\isanewline
\ \ \ \ \isakeywordONE{let}\isamarkupfalse%
\ {\isacharquery}{\kern0pt}case{\isadigit{4}}{\isacharequal}{\kern0pt}{\isachardoublequoteopen}res{\isacharunderscore}{\kern0pt}inst{\isacharunderscore}{\kern0pt}snap{\isacharunderscore}{\kern0pt}action\ {\isasympi}\ Over{\isacharunderscore}{\kern0pt}All\ {\isacharequal}{\kern0pt}\ Some\ a\ {\isasymand}\ \isanewline
\ \ \ \ \ \ \ \ \ \ {\isacharparenleft}{\kern0pt}{\isasymexists}t\isactrlsub i\ t\isactrlsub j{\isachardot}{\kern0pt}\ consec{\isacharunderscore}{\kern0pt}htps\ {\isasympi}s\ t\isactrlsub i\ t\isactrlsub j\ {\isasymand}\ t\isactrlsub {\isasympi}\ {\isasymle}\ t\isactrlsub i\ {\isasymand}\ t\isactrlsub j\ {\isasymle}\ t\isactrlsub {\isasympi}\ {\isacharplus}{\kern0pt}\ {\isacharparenleft}{\kern0pt}duration\ {\isasympi}{\isacharparenright}{\kern0pt}\ {\isasymand}\ t\isactrlsub i\ {\isacharless}{\kern0pt}\ t\isactrlsub a\ {\isasymand}\ t\isactrlsub a\ {\isacharless}{\kern0pt}\ t\isactrlsub j{\isacharparenright}{\kern0pt}{\isachardoublequoteclose}\isanewline
\ \ \ \ \isakeywordONE{have}\isamarkupfalse%
\ {\isachardoublequoteopen}{\isacharquery}{\kern0pt}case{\isadigit{1}}\ {\isasymor}\ {\isacharquery}{\kern0pt}case{\isadigit{2}}\ {\isasymor}\ {\isacharquery}{\kern0pt}case{\isadigit{3}}\ {\isasymor}\ {\isacharquery}{\kern0pt}case{\isadigit{4}}{\isachardoublequoteclose}\isanewline
\ \ \ \ \ \ \isakeywordONE{using}\isamarkupfalse%
\ {\isacartoucheopen}inst{\isacharunderscore}{\kern0pt}of{\isacharunderscore}{\kern0pt}plan{\isacharunderscore}{\kern0pt}action\ {\isasympi}s\ {\isacharparenleft}{\kern0pt}t\isactrlsub {\isasympi}{\isacharcomma}{\kern0pt}{\isasympi}{\isacharparenright}{\kern0pt}\ {\isacharparenleft}{\kern0pt}t\isactrlsub a{\isacharcomma}{\kern0pt}a{\isacharparenright}{\kern0pt}{\isacartoucheclose}\ \isakeywordONE{by}\isamarkupfalse%
\ simp\isanewline
\ \ \ \ \isakeywordONE{then}\isamarkupfalse%
\ \isakeywordTHREE{show}\isamarkupfalse%
\ {\isacharquery}{\kern0pt}thesis\isanewline
\ \ \ \ \isakeywordONE{proof}\isamarkupfalse%
\ {\isacharparenleft}{\kern0pt}elim\ disjE{\isacharparenright}{\kern0pt}\isanewline
\ \ \ \ \ \ \isakeywordTHREE{assume}\isamarkupfalse%
\ {\isacharquery}{\kern0pt}case{\isadigit{1}}\isanewline
\ \ \ \ \ \ \isakeywordONE{then}\isamarkupfalse%
\ \isakeywordONE{have}\isamarkupfalse%
\ {\isachardoublequoteopen}{\isacharparenleft}{\kern0pt}t\isactrlsub {\isasympi}{\isacharcomma}{\kern0pt}{\isasympi}{\isacharparenright}{\kern0pt}\ {\isasymin}\ simple{\isacharunderscore}{\kern0pt}acts\ {\isasympi}s{\isachardoublequoteclose}\isanewline
\ \ \ \ \ \ \ \ \isakeywordONE{unfolding}\isamarkupfalse%
\ simple{\isacharunderscore}{\kern0pt}acts{\isacharunderscore}{\kern0pt}def\ is{\isacharunderscore}{\kern0pt}act{\isacharunderscore}{\kern0pt}simple{\isacharunderscore}{\kern0pt}def\ \isakeywordONE{using}\isamarkupfalse%
\ {\isacartoucheopen}{\isacharparenleft}{\kern0pt}t\isactrlsub {\isasympi}{\isacharcomma}{\kern0pt}{\isasympi}{\isacharparenright}{\kern0pt}\ {\isasymin}\ set\ {\isasympi}s{\isacartoucheclose}\ \isakeywordONE{by}\isamarkupfalse%
\ {\isacharparenleft}{\kern0pt}cases\ {\isasympi}{\isacharparenright}{\kern0pt}\ auto\isanewline
\ \ \ \ \ \ \isakeywordONE{then}\isamarkupfalse%
\ \isakeywordONE{have}\isamarkupfalse%
\ {\isachardoublequoteopen}{\isasymexists}as{\isachardot}{\kern0pt}\ {\isacharparenleft}{\kern0pt}t\isactrlsub {\isasympi}{\isacharcomma}{\kern0pt}as{\isacharparenright}{\kern0pt}\ {\isasymin}\ set\ hs\isactrlsub {\isadigit{2}}\ {\isasymand}\ the\ {\isacharparenleft}{\kern0pt}res{\isacharunderscore}{\kern0pt}inst\ {\isasympi}\ At{\isacharunderscore}{\kern0pt}Start{\isacharparenright}{\kern0pt}\ {\isasymin}\ set\ as{\isachardoublequoteclose}\isanewline
\ \ \ \ \ \ \ \ \isakeywordONE{using}\isamarkupfalse%
\ {\isacartoucheopen}ind{\isacharunderscore}{\kern0pt}happ{\isacharunderscore}{\kern0pt}seq\ {\isasympi}s\ hs\isactrlsub {\isadigit{2}}{\isacartoucheclose}\ \isakeywordONE{unfolding}\isamarkupfalse%
\ ind{\isacharunderscore}{\kern0pt}happ{\isacharunderscore}{\kern0pt}seq{\isacharunderscore}{\kern0pt}def\ \isakeywordONE{by}\isamarkupfalse%
\ auto\isanewline
\ \ \ \ \ \ \isakeywordONE{then}\isamarkupfalse%
\ \isakeywordTHREE{obtain}\isamarkupfalse%
\ A{\isacharprime}{\kern0pt}\ \isakeywordTWO{where}\ {\isachardoublequoteopen}{\isacharparenleft}{\kern0pt}t\isactrlsub {\isasympi}{\isacharcomma}{\kern0pt}A{\isacharprime}{\kern0pt}{\isacharparenright}{\kern0pt}\ {\isasymin}\ set\ hs\isactrlsub {\isadigit{2}}{\isachardoublequoteclose}\ \isakeywordTWO{and}\ {\isachardoublequoteopen}the\ {\isacharparenleft}{\kern0pt}res{\isacharunderscore}{\kern0pt}inst\ {\isasympi}\ At{\isacharunderscore}{\kern0pt}Start{\isacharparenright}{\kern0pt}\ {\isasymin}\ set\ A{\isacharprime}{\kern0pt}{\isachardoublequoteclose}\ \isanewline
\ \ \ \ \ \ \ \ \isakeywordONE{by}\isamarkupfalse%
\ auto\isanewline
\ \ \ \ \ \ \isakeywordONE{then}\isamarkupfalse%
\ \isakeywordONE{have}\isamarkupfalse%
\ {\isachardoublequoteopen}{\isacharparenleft}{\kern0pt}t\isactrlsub a{\isacharcomma}{\kern0pt}A{\isacharprime}{\kern0pt}{\isacharparenright}{\kern0pt}\ {\isasymin}\ set\ hs\isactrlsub {\isadigit{2}}{\isachardoublequoteclose}\isanewline
\ \ \ \ \ \ \ \ \isakeywordONE{using}\isamarkupfalse%
\ {\isacartoucheopen}{\isacharquery}{\kern0pt}case{\isadigit{1}}{\isacartoucheclose}\ \isakeywordONE{by}\isamarkupfalse%
\ auto\isanewline
\ \ \ \ \ \ \isakeywordTHREE{show}\isamarkupfalse%
\ {\isacharquery}{\kern0pt}thesis\isanewline
\ \ \ \ \ \ \ \ \isakeywordONE{using}\isamarkupfalse%
\ {\isacartoucheopen}{\isacharquery}{\kern0pt}case{\isadigit{1}}{\isacartoucheclose}\ {\isacartoucheopen}t\isactrlsub i\ {\isacharless}{\kern0pt}\ t\isactrlsub a\ {\isasymand}\ t\isactrlsub a\ {\isacharless}{\kern0pt}\ t\isactrlsub j{\isacartoucheclose}\ {\isacartoucheopen}the\ {\isacharparenleft}{\kern0pt}res{\isacharunderscore}{\kern0pt}inst\ {\isasympi}\ At{\isacharunderscore}{\kern0pt}Start{\isacharparenright}{\kern0pt}\ {\isasymin}\ set\ A{\isacharprime}{\kern0pt}{\isacartoucheclose}\isanewline
\ \ \ \ \ \ \ \ \ \ \ \ \ \ strict{\isacharunderscore}{\kern0pt}sorted{\isacharunderscore}{\kern0pt}drop{\isacharunderscore}{\kern0pt}leq{\isacharunderscore}{\kern0pt}take{\isacharunderscore}{\kern0pt}le{\isacharbrackleft}{\kern0pt}OF\ {\isacartoucheopen}strict{\isacharunderscore}{\kern0pt}sorted\ {\isacharparenleft}{\kern0pt}map\ fst\ hs\isactrlsub {\isadigit{2}}{\isacharparenright}{\kern0pt}{\isacartoucheclose}\ {\isacartoucheopen}{\isacharparenleft}{\kern0pt}t\isactrlsub a{\isacharcomma}{\kern0pt}A{\isacharprime}{\kern0pt}{\isacharparenright}{\kern0pt}\ {\isasymin}\ set\ hs\isactrlsub {\isadigit{2}}{\isacartoucheclose}{\isacharbrackright}{\kern0pt}\isanewline
\ \ \ \ \ \ \ \ \isakeywordONE{by}\isamarkupfalse%
\ auto\ \ \ \isanewline
\ \ \ \ \isakeywordONE{next}\isamarkupfalse%
\isanewline
\ \ \ \ \ \ \isakeywordTHREE{assume}\isamarkupfalse%
\ {\isacharquery}{\kern0pt}case{\isadigit{2}}\isanewline
\ \ \ \ \ \ \isakeywordONE{let}\isamarkupfalse%
\ {\isacharquery}{\kern0pt}a\isactrlsub s\isactrlsub t\isactrlsub a\isactrlsub r\isactrlsub t{\isacharequal}{\kern0pt}{\isachardoublequoteopen}the\ {\isacharparenleft}{\kern0pt}res{\isacharunderscore}{\kern0pt}inst{\isacharunderscore}{\kern0pt}snap{\isacharunderscore}{\kern0pt}action\ {\isasympi}\ At{\isacharunderscore}{\kern0pt}Start{\isacharparenright}{\kern0pt}{\isachardoublequoteclose}\isanewline
\ \ \ \ \ \ \isakeywordONE{have}\isamarkupfalse%
\ {\isachardoublequoteopen}{\isacharparenleft}{\kern0pt}t\isactrlsub {\isasympi}{\isacharcomma}{\kern0pt}{\isasympi}{\isacharparenright}{\kern0pt}\ {\isasymin}\ durative{\isacharunderscore}{\kern0pt}acts\ {\isasympi}s{\isachardoublequoteclose}\isanewline
\ \ \ \ \ \ \ \ \isakeywordONE{unfolding}\isamarkupfalse%
\ durative{\isacharunderscore}{\kern0pt}acts{\isacharunderscore}{\kern0pt}def\ is{\isacharunderscore}{\kern0pt}act{\isacharunderscore}{\kern0pt}simple{\isacharunderscore}{\kern0pt}def\isanewline
\ \ \ \ \ \ \ \ \isakeywordONE{using}\isamarkupfalse%
\ {\isacartoucheopen}{\isacharquery}{\kern0pt}case{\isadigit{2}}{\isacartoucheclose}\ {\isacartoucheopen}{\isacharparenleft}{\kern0pt}t\isactrlsub {\isasympi}{\isacharcomma}{\kern0pt}{\isasympi}{\isacharparenright}{\kern0pt}\ {\isasymin}\ set\ {\isasympi}s{\isacartoucheclose}\ \isakeywordONE{by}\isamarkupfalse%
\ {\isacharparenleft}{\kern0pt}cases\ {\isasympi}{\isacharparenright}{\kern0pt}\ auto\isanewline
\ \ \ \ \ \ \isakeywordONE{then}\isamarkupfalse%
\ \isakeywordONE{have}\isamarkupfalse%
\ {\isachardoublequoteopen}{\isasymexists}as{\isachardot}{\kern0pt}\ {\isacharparenleft}{\kern0pt}t\isactrlsub {\isasympi}{\isacharcomma}{\kern0pt}as{\isacharparenright}{\kern0pt}\ {\isasymin}\ set\ hs\isactrlsub {\isadigit{2}}\ {\isasymand}\ {\isacharquery}{\kern0pt}a\isactrlsub s\isactrlsub t\isactrlsub a\isactrlsub r\isactrlsub t\ {\isasymin}\ set\ as{\isachardoublequoteclose}\isanewline
\ \ \ \ \ \ \ \ \isakeywordONE{using}\isamarkupfalse%
\ {\isacartoucheopen}ind{\isacharunderscore}{\kern0pt}happ{\isacharunderscore}{\kern0pt}seq\ {\isasympi}s\ hs\isactrlsub {\isadigit{2}}{\isacartoucheclose}\ \isakeywordONE{unfolding}\isamarkupfalse%
\ ind{\isacharunderscore}{\kern0pt}happ{\isacharunderscore}{\kern0pt}seq{\isacharunderscore}{\kern0pt}def\ \isakeywordONE{by}\isamarkupfalse%
\ auto\isanewline
\ \ \ \ \ \ \isakeywordONE{then}\isamarkupfalse%
\ \isakeywordTHREE{obtain}\isamarkupfalse%
\ A{\isacharprime}{\kern0pt}\ \isakeywordTWO{where}\ {\isachardoublequoteopen}{\isacharparenleft}{\kern0pt}t\isactrlsub {\isasympi}{\isacharcomma}{\kern0pt}A{\isacharprime}{\kern0pt}{\isacharparenright}{\kern0pt}\ {\isasymin}\ set\ hs\isactrlsub {\isadigit{2}}{\isachardoublequoteclose}\ \isakeywordTWO{and}\ {\isachardoublequoteopen}{\isacharquery}{\kern0pt}a\isactrlsub s\isactrlsub t\isactrlsub a\isactrlsub r\isactrlsub t\ {\isasymin}\ set\ A{\isacharprime}{\kern0pt}{\isachardoublequoteclose}\ \isanewline
\ \ \ \ \ \ \ \ \isakeywordONE{by}\isamarkupfalse%
\ auto\isanewline
\ \ \ \ \ \ \isakeywordONE{then}\isamarkupfalse%
\ \isakeywordONE{have}\isamarkupfalse%
\ {\isachardoublequoteopen}{\isacharparenleft}{\kern0pt}t\isactrlsub a{\isacharcomma}{\kern0pt}A{\isacharprime}{\kern0pt}{\isacharparenright}{\kern0pt}\ {\isasymin}\ set\ hs\isactrlsub {\isadigit{2}}{\isachardoublequoteclose}\isanewline
\ \ \ \ \ \ \ \ \isakeywordONE{using}\isamarkupfalse%
\ {\isacartoucheopen}{\isacharquery}{\kern0pt}case{\isadigit{2}}{\isacartoucheclose}\ \isakeywordONE{by}\isamarkupfalse%
\ auto\isanewline
\ \ \ \ \ \ \isakeywordTHREE{show}\isamarkupfalse%
\ {\isacharquery}{\kern0pt}thesis\isanewline
\ \ \ \ \ \ \ \ \isakeywordONE{using}\isamarkupfalse%
\ {\isacartoucheopen}{\isacharquery}{\kern0pt}case{\isadigit{2}}{\isacartoucheclose}\ {\isacartoucheopen}t\isactrlsub i\ {\isacharless}{\kern0pt}\ t\isactrlsub a\ {\isasymand}\ t\isactrlsub a\ {\isacharless}{\kern0pt}\ t\isactrlsub j{\isacartoucheclose}\ {\isacartoucheopen}{\isacharquery}{\kern0pt}a\isactrlsub s\isactrlsub t\isactrlsub a\isactrlsub r\isactrlsub t\ {\isasymin}\ set\ A{\isacharprime}{\kern0pt}{\isacartoucheclose}\isanewline
\ \ \ \ \ \ \ \ \ \ \ \ \ \ strict{\isacharunderscore}{\kern0pt}sorted{\isacharunderscore}{\kern0pt}drop{\isacharunderscore}{\kern0pt}leq{\isacharunderscore}{\kern0pt}take{\isacharunderscore}{\kern0pt}le{\isacharbrackleft}{\kern0pt}OF\ {\isacartoucheopen}strict{\isacharunderscore}{\kern0pt}sorted\ {\isacharparenleft}{\kern0pt}map\ fst\ hs\isactrlsub {\isadigit{2}}{\isacharparenright}{\kern0pt}{\isacartoucheclose}\ {\isacartoucheopen}{\isacharparenleft}{\kern0pt}t\isactrlsub a{\isacharcomma}{\kern0pt}A{\isacharprime}{\kern0pt}{\isacharparenright}{\kern0pt}\ {\isasymin}\ set\ hs\isactrlsub {\isadigit{2}}{\isacartoucheclose}{\isacharbrackright}{\kern0pt}\isanewline
\ \ \ \ \ \ \ \ \isakeywordONE{by}\isamarkupfalse%
\ auto\isanewline
\ \ \ \ \isakeywordONE{next}\isamarkupfalse%
\isanewline
\ \ \ \ \ \ \isakeywordTHREE{assume}\isamarkupfalse%
\ {\isacharquery}{\kern0pt}case{\isadigit{3}}\isanewline
\ \ \ \ \ \ \isakeywordONE{let}\isamarkupfalse%
\ {\isacharquery}{\kern0pt}a\isactrlsub e\isactrlsub n\isactrlsub d{\isacharequal}{\kern0pt}{\isachardoublequoteopen}the\ {\isacharparenleft}{\kern0pt}res{\isacharunderscore}{\kern0pt}inst{\isacharunderscore}{\kern0pt}snap{\isacharunderscore}{\kern0pt}action\ {\isasympi}\ At{\isacharunderscore}{\kern0pt}End{\isacharparenright}{\kern0pt}{\isachardoublequoteclose}\isanewline
\ \ \ \ \ \ \isakeywordONE{have}\isamarkupfalse%
\ {\isachardoublequoteopen}{\isacharparenleft}{\kern0pt}t\isactrlsub {\isasympi}{\isacharcomma}{\kern0pt}{\isasympi}{\isacharparenright}{\kern0pt}\ {\isasymin}\ durative{\isacharunderscore}{\kern0pt}acts\ {\isasympi}s{\isachardoublequoteclose}\isanewline
\ \ \ \ \ \ \ \ \isakeywordONE{unfolding}\isamarkupfalse%
\ durative{\isacharunderscore}{\kern0pt}acts{\isacharunderscore}{\kern0pt}def\ is{\isacharunderscore}{\kern0pt}act{\isacharunderscore}{\kern0pt}simple{\isacharunderscore}{\kern0pt}def\isanewline
\ \ \ \ \ \ \ \ \isakeywordONE{using}\isamarkupfalse%
\ {\isacartoucheopen}{\isacharquery}{\kern0pt}case{\isadigit{3}}{\isacartoucheclose}\ {\isacartoucheopen}{\isacharparenleft}{\kern0pt}t\isactrlsub {\isasympi}{\isacharcomma}{\kern0pt}{\isasympi}{\isacharparenright}{\kern0pt}\ {\isasymin}\ set\ {\isasympi}s{\isacartoucheclose}\ \isakeywordONE{by}\isamarkupfalse%
\ {\isacharparenleft}{\kern0pt}cases\ {\isasympi}{\isacharparenright}{\kern0pt}\ auto\isanewline
\ \ \ \ \ \ \isakeywordONE{then}\isamarkupfalse%
\ \isakeywordONE{have}\isamarkupfalse%
\ {\isachardoublequoteopen}{\isasymexists}as{\isachardot}{\kern0pt}\ {\isacharparenleft}{\kern0pt}t\isactrlsub {\isasympi}\ {\isacharplus}{\kern0pt}\ {\isacharparenleft}{\kern0pt}duration\ {\isasympi}{\isacharparenright}{\kern0pt}{\isacharcomma}{\kern0pt}as{\isacharparenright}{\kern0pt}\ {\isasymin}\ set\ hs\isactrlsub {\isadigit{2}}\ {\isasymand}\ {\isacharquery}{\kern0pt}a\isactrlsub e\isactrlsub n\isactrlsub d\ {\isasymin}\ set\ as{\isachardoublequoteclose}\isanewline
\ \ \ \ \ \ \ \ \isakeywordONE{using}\isamarkupfalse%
\ {\isacartoucheopen}ind{\isacharunderscore}{\kern0pt}happ{\isacharunderscore}{\kern0pt}seq\ {\isasympi}s\ hs\isactrlsub {\isadigit{2}}{\isacartoucheclose}\ \isakeywordONE{unfolding}\isamarkupfalse%
\ ind{\isacharunderscore}{\kern0pt}happ{\isacharunderscore}{\kern0pt}seq{\isacharunderscore}{\kern0pt}def\ \isakeywordONE{by}\isamarkupfalse%
\ auto\isanewline
\ \ \ \ \ \ \isakeywordONE{then}\isamarkupfalse%
\ \isakeywordTHREE{obtain}\isamarkupfalse%
\ A{\isacharprime}{\kern0pt}\ \isakeywordTWO{where}\ {\isachardoublequoteopen}{\isacharparenleft}{\kern0pt}t\isactrlsub {\isasympi}\ {\isacharplus}{\kern0pt}\ {\isacharparenleft}{\kern0pt}duration\ {\isasympi}{\isacharparenright}{\kern0pt}{\isacharcomma}{\kern0pt}A{\isacharprime}{\kern0pt}{\isacharparenright}{\kern0pt}\ {\isasymin}\ set\ hs\isactrlsub {\isadigit{2}}{\isachardoublequoteclose}\ \isakeywordTWO{and}\ {\isachardoublequoteopen}{\isacharquery}{\kern0pt}a\isactrlsub e\isactrlsub n\isactrlsub d\ {\isasymin}\ set\ A{\isacharprime}{\kern0pt}{\isachardoublequoteclose}\ \isanewline
\ \ \ \ \ \ \ \ \isakeywordONE{by}\isamarkupfalse%
\ auto\isanewline
\ \ \ \ \ \ \isakeywordONE{then}\isamarkupfalse%
\ \isakeywordONE{have}\isamarkupfalse%
\ {\isachardoublequoteopen}{\isacharparenleft}{\kern0pt}t\isactrlsub a{\isacharcomma}{\kern0pt}A{\isacharprime}{\kern0pt}{\isacharparenright}{\kern0pt}\ {\isasymin}\ set\ hs\isactrlsub {\isadigit{2}}{\isachardoublequoteclose}\isanewline
\ \ \ \ \ \ \ \ \isakeywordONE{using}\isamarkupfalse%
\ {\isacartoucheopen}{\isacharquery}{\kern0pt}case{\isadigit{3}}{\isacartoucheclose}\ \isakeywordONE{by}\isamarkupfalse%
\ auto\isanewline
\ \ \ \ \ \ \isakeywordTHREE{show}\isamarkupfalse%
\ {\isacharquery}{\kern0pt}thesis\isanewline
\ \ \ \ \ \ \ \ \isakeywordONE{using}\isamarkupfalse%
\ {\isacartoucheopen}{\isacharquery}{\kern0pt}case{\isadigit{3}}{\isacartoucheclose}\ {\isacartoucheopen}t\isactrlsub i\ {\isacharless}{\kern0pt}\ t\isactrlsub a\ {\isasymand}\ t\isactrlsub a\ {\isacharless}{\kern0pt}\ t\isactrlsub j{\isacartoucheclose}\ {\isacartoucheopen}{\isacharquery}{\kern0pt}a\isactrlsub e\isactrlsub n\isactrlsub d\ {\isasymin}\ set\ A{\isacharprime}{\kern0pt}{\isacartoucheclose}\isanewline
\ \ \ \ \ \ \ \ \ \ \ \ \ \ strict{\isacharunderscore}{\kern0pt}sorted{\isacharunderscore}{\kern0pt}drop{\isacharunderscore}{\kern0pt}leq{\isacharunderscore}{\kern0pt}take{\isacharunderscore}{\kern0pt}le{\isacharbrackleft}{\kern0pt}OF\ {\isacartoucheopen}strict{\isacharunderscore}{\kern0pt}sorted\ {\isacharparenleft}{\kern0pt}map\ fst\ hs\isactrlsub {\isadigit{2}}{\isacharparenright}{\kern0pt}{\isacartoucheclose}\ {\isacartoucheopen}{\isacharparenleft}{\kern0pt}t\isactrlsub a{\isacharcomma}{\kern0pt}A{\isacharprime}{\kern0pt}{\isacharparenright}{\kern0pt}\ {\isasymin}\ set\ hs\isactrlsub {\isadigit{2}}{\isacartoucheclose}{\isacharbrackright}{\kern0pt}\isanewline
\ \ \ \ \ \ \ \ \isakeywordONE{by}\isamarkupfalse%
\ auto\isanewline
\ \ \ \ \isakeywordONE{next}\isamarkupfalse%
\isanewline
\ \ \ \ \ \ \isakeywordTHREE{assume}\isamarkupfalse%
\ {\isacharquery}{\kern0pt}case{\isadigit{4}}\isanewline
\ \ \ \ \ \ \isakeywordONE{let}\isamarkupfalse%
\ {\isacharquery}{\kern0pt}a\isactrlsub i\isactrlsub n\isactrlsub v{\isacharequal}{\kern0pt}{\isachardoublequoteopen}the\ {\isacharparenleft}{\kern0pt}res{\isacharunderscore}{\kern0pt}inst{\isacharunderscore}{\kern0pt}snap{\isacharunderscore}{\kern0pt}action\ {\isasympi}\ Over{\isacharunderscore}{\kern0pt}All{\isacharparenright}{\kern0pt}{\isachardoublequoteclose}\isanewline
\ \ \ \ \ \ \isakeywordTHREE{obtain}\isamarkupfalse%
\ t\isactrlsub i{\isacharprime}{\kern0pt}\ t\isactrlsub j{\isacharprime}{\kern0pt}\ \isakeywordTWO{where}\ {\isachardoublequoteopen}consec{\isacharunderscore}{\kern0pt}htps\ {\isasympi}s\ t\isactrlsub i{\isacharprime}{\kern0pt}\ t\isactrlsub j{\isacharprime}{\kern0pt}{\isachardoublequoteclose}\ \isanewline
\ \ \ \ \ \ \ \ \ \ \ \ \ \ \ \ \ \ \ \ \ \ \isakeywordTWO{and}\ {\isachardoublequoteopen}t\isactrlsub {\isasympi}\ {\isasymle}\ t\isactrlsub i{\isacharprime}{\kern0pt}\ {\isasymand}\ t\isactrlsub j{\isacharprime}{\kern0pt}\ {\isasymle}\ t\isactrlsub {\isasympi}\ {\isacharplus}{\kern0pt}\ duration\ {\isasympi}{\isachardoublequoteclose}\ \isanewline
\ \ \ \ \ \ \ \ \ \ \ \ \ \ \ \ \ \ \ \ \ \ \isakeywordTWO{and}\ {\isachardoublequoteopen}t\isactrlsub i{\isacharprime}{\kern0pt}\ {\isacharless}{\kern0pt}\ t\isactrlsub a\ {\isasymand}\ t\isactrlsub a\ {\isacharless}{\kern0pt}\ t\isactrlsub j{\isacharprime}{\kern0pt}{\isachardoublequoteclose}\isanewline
\ \ \ \ \ \ \ \ \isakeywordONE{using}\isamarkupfalse%
\ {\isacartoucheopen}{\isacharquery}{\kern0pt}case{\isadigit{4}}{\isacartoucheclose}\ \isakeywordONE{by}\isamarkupfalse%
\ auto\isanewline
\ \ \ \ \ \ \isakeywordONE{have}\isamarkupfalse%
\ {\isachardoublequoteopen}{\isacharparenleft}{\kern0pt}t\isactrlsub {\isasympi}{\isacharcomma}{\kern0pt}{\isasympi}{\isacharparenright}{\kern0pt}\ {\isasymin}\ durative{\isacharunderscore}{\kern0pt}acts\ {\isasympi}s{\isachardoublequoteclose}\isanewline
\ \ \ \ \ \ \ \ \isakeywordONE{unfolding}\isamarkupfalse%
\ durative{\isacharunderscore}{\kern0pt}acts{\isacharunderscore}{\kern0pt}def\ is{\isacharunderscore}{\kern0pt}act{\isacharunderscore}{\kern0pt}simple{\isacharunderscore}{\kern0pt}def\isanewline
\ \ \ \ \ \ \ \ \isakeywordONE{using}\isamarkupfalse%
\ {\isacartoucheopen}{\isacharquery}{\kern0pt}case{\isadigit{4}}{\isacartoucheclose}\ {\isacartoucheopen}{\isacharparenleft}{\kern0pt}t\isactrlsub {\isasympi}{\isacharcomma}{\kern0pt}{\isasympi}{\isacharparenright}{\kern0pt}\ {\isasymin}\ set\ {\isasympi}s{\isacartoucheclose}\ \isakeywordONE{by}\isamarkupfalse%
\ {\isacharparenleft}{\kern0pt}cases\ {\isasympi}{\isacharparenright}{\kern0pt}\ auto\isanewline
\ \ \ \ \ \ \isanewline
\ \ \ \ \ \ \isakeywordONE{then}\isamarkupfalse%
\ \isakeywordONE{have}\isamarkupfalse%
\ {\isachardoublequoteopen}{\isasymexists}{\isacharparenleft}{\kern0pt}t{\isacharprime}{\kern0pt}{\isacharcomma}{\kern0pt}as{\isacharparenright}{\kern0pt}\ {\isasymin}\ set\ hs\isactrlsub {\isadigit{2}}{\isachardot}{\kern0pt}\ t\isactrlsub i{\isacharprime}{\kern0pt}\ {\isacharless}{\kern0pt}\ t{\isacharprime}{\kern0pt}\ {\isasymand}\ t{\isacharprime}{\kern0pt}\ {\isacharless}{\kern0pt}\ t\isactrlsub j{\isacharprime}{\kern0pt}\ {\isasymand}\ {\isacharquery}{\kern0pt}a\isactrlsub i\isactrlsub n\isactrlsub v\ {\isasymin}\ set\ as{\isachardoublequoteclose}\isanewline
\ \ \ \ \ \ \ \ \isakeywordONE{using}\isamarkupfalse%
\ {\isacartoucheopen}ind{\isacharunderscore}{\kern0pt}happ{\isacharunderscore}{\kern0pt}seq\ {\isasympi}s\ hs\isactrlsub {\isadigit{2}}{\isacartoucheclose}\ \isakeywordONE{unfolding}\isamarkupfalse%
\ ind{\isacharunderscore}{\kern0pt}happ{\isacharunderscore}{\kern0pt}seq{\isacharunderscore}{\kern0pt}def\isanewline
\ \ \ \ \ \ \ \ \isakeywordONE{using}\isamarkupfalse%
\ {\isacartoucheopen}consec{\isacharunderscore}{\kern0pt}htps\ {\isasympi}s\ t\isactrlsub i{\isacharprime}{\kern0pt}\ t\isactrlsub j{\isacharprime}{\kern0pt}{\isacartoucheclose}\ {\isacartoucheopen}t\isactrlsub {\isasympi}\ {\isasymle}\ t\isactrlsub i{\isacharprime}{\kern0pt}\ {\isasymand}\ t\isactrlsub j{\isacharprime}{\kern0pt}\ {\isasymle}\ t\isactrlsub {\isasympi}\ {\isacharplus}{\kern0pt}\ duration\ {\isasympi}{\isacartoucheclose}\ \isakeywordONE{by}\isamarkupfalse%
\ auto\isanewline
\ \ \ \ \ \ \isakeywordONE{then}\isamarkupfalse%
\ \isakeywordTHREE{obtain}\isamarkupfalse%
\ t\isactrlsub a{\isacharprime}{\kern0pt}\ A{\isacharprime}{\kern0pt}\ \isakeywordTWO{where}\ {\isachardoublequoteopen}{\isacharparenleft}{\kern0pt}t\isactrlsub a{\isacharprime}{\kern0pt}{\isacharcomma}{\kern0pt}A{\isacharprime}{\kern0pt}{\isacharparenright}{\kern0pt}\ {\isasymin}\ set\ hs\isactrlsub {\isadigit{2}}{\isachardoublequoteclose}\ \isakeywordTWO{and}\ {\isachardoublequoteopen}t\isactrlsub i{\isacharprime}{\kern0pt}\ {\isacharless}{\kern0pt}\ t\isactrlsub a{\isacharprime}{\kern0pt}\ {\isasymand}\ t\isactrlsub a{\isacharprime}{\kern0pt}\ {\isacharless}{\kern0pt}\ t\isactrlsub j{\isacharprime}{\kern0pt}{\isachardoublequoteclose}\ \isakeywordTWO{and}\ {\isachardoublequoteopen}{\isacharquery}{\kern0pt}a\isactrlsub i\isactrlsub n\isactrlsub v\ {\isasymin}\ set\ A{\isacharprime}{\kern0pt}{\isachardoublequoteclose}\isanewline
\ \ \ \ \ \ \ \ \isakeywordONE{by}\isamarkupfalse%
\ auto\isanewline
\ \ \ \ \ \ \isakeywordONE{then}\isamarkupfalse%
\ \isakeywordONE{have}\isamarkupfalse%
\ {\isachardoublequoteopen}{\isacharparenleft}{\kern0pt}t\isactrlsub a{\isacharprime}{\kern0pt}{\isacharcomma}{\kern0pt}A{\isacharprime}{\kern0pt}{\isacharparenright}{\kern0pt}\ {\isasymin}\ set\ hs\isactrlsub {\isadigit{2}}{\isachardoublequoteclose}\isanewline
\ \ \ \ \ \ \ \ \isakeywordONE{using}\isamarkupfalse%
\ {\isacartoucheopen}{\isacharquery}{\kern0pt}case{\isadigit{4}}{\isacartoucheclose}\ \isakeywordONE{by}\isamarkupfalse%
\ auto\isanewline
\ \ \ \ \ \ \isakeywordONE{have}\isamarkupfalse%
\ {\isachardoublequoteopen}t\isactrlsub i\ {\isacharless}{\kern0pt}\ t\isactrlsub a{\isacharprime}{\kern0pt}\ {\isasymand}\ t\isactrlsub a{\isacharprime}{\kern0pt}\ {\isacharless}{\kern0pt}\ t\isactrlsub j{\isachardoublequoteclose}\isanewline
\ \ \ \ \ \ \ \ \isakeywordONE{using}\isamarkupfalse%
\ assms\ {\isacartoucheopen}consec{\isacharunderscore}{\kern0pt}htps\ {\isasympi}s\ t\isactrlsub i{\isacharprime}{\kern0pt}\ t\isactrlsub j{\isacharprime}{\kern0pt}{\isacartoucheclose}\ {\isacartoucheopen}t\isactrlsub i\ {\isacharless}{\kern0pt}\ t\isactrlsub a\ {\isasymand}\ t\isactrlsub a\ {\isacharless}{\kern0pt}\ t\isactrlsub j{\isacartoucheclose}\ {\isacartoucheopen}t\isactrlsub i{\isacharprime}{\kern0pt}\ {\isacharless}{\kern0pt}\ t\isactrlsub a\ {\isasymand}\ t\isactrlsub a\ {\isacharless}{\kern0pt}\ t\isactrlsub j{\isacharprime}{\kern0pt}{\isacartoucheclose}\ \isanewline
\ \ \ \ \ \ \ \ \ \ \ \ \ \ {\isacartoucheopen}t\isactrlsub i{\isacharprime}{\kern0pt}\ {\isacharless}{\kern0pt}\ t\isactrlsub a{\isacharprime}{\kern0pt}\ {\isasymand}\ t\isactrlsub a{\isacharprime}{\kern0pt}\ {\isacharless}{\kern0pt}\ t\isactrlsub j{\isacharprime}{\kern0pt}{\isacartoucheclose}\ \isakeywordONE{unfolding}\isamarkupfalse%
\ consec{\isacharunderscore}{\kern0pt}htps{\isacharunderscore}{\kern0pt}def\ \isakeywordONE{by}\isamarkupfalse%
\ auto\isanewline
\ \ \ \ \ \ \isakeywordTHREE{show}\isamarkupfalse%
\ {\isacharquery}{\kern0pt}thesis\isanewline
\ \ \ \ \ \ \ \ \isakeywordONE{using}\isamarkupfalse%
\ {\isacartoucheopen}{\isacharquery}{\kern0pt}case{\isadigit{4}}{\isacartoucheclose}\ {\isacartoucheopen}{\isacharquery}{\kern0pt}a\isactrlsub i\isactrlsub n\isactrlsub v\ {\isasymin}\ set\ A{\isacharprime}{\kern0pt}{\isacartoucheclose}\ {\isacartoucheopen}t\isactrlsub i\ {\isacharless}{\kern0pt}\ t\isactrlsub a{\isacharprime}{\kern0pt}\ {\isasymand}\ t\isactrlsub a{\isacharprime}{\kern0pt}\ {\isacharless}{\kern0pt}\ t\isactrlsub j{\isacartoucheclose}\isanewline
\ \ \ \ \ \ \ \ \ \ \ \ \ \ strict{\isacharunderscore}{\kern0pt}sorted{\isacharunderscore}{\kern0pt}drop{\isacharunderscore}{\kern0pt}leq{\isacharunderscore}{\kern0pt}take{\isacharunderscore}{\kern0pt}le{\isacharbrackleft}{\kern0pt}OF\ {\isacartoucheopen}strict{\isacharunderscore}{\kern0pt}sorted\ {\isacharparenleft}{\kern0pt}map\ fst\ hs\isactrlsub {\isadigit{2}}{\isacharparenright}{\kern0pt}{\isacartoucheclose}\ {\isacartoucheopen}{\isacharparenleft}{\kern0pt}t\isactrlsub a{\isacharprime}{\kern0pt}{\isacharcomma}{\kern0pt}A{\isacharprime}{\kern0pt}{\isacharparenright}{\kern0pt}\ {\isasymin}\ set\ hs\isactrlsub {\isadigit{2}}{\isacartoucheclose}{\isacharbrackright}{\kern0pt}\isanewline
\ \ \ \ \ \ \ \ \isakeywordONE{by}\isamarkupfalse%
\ auto\isanewline
\ \ \ \ \isakeywordONE{qed}\isamarkupfalse%
\isanewline
\ \ \isakeywordONE{qed}\isamarkupfalse%
%
\endisatagproof
{\isafoldproof}%
%
\isadelimproof
%
\endisadelimproof
%
\begin{isamarkuptext}%
If one induced happening sequence is a valid happening sequence in between two 
  consecutive happening time points, then every other induced happening sequence is 
  also valid in between those consecutive happening time point.%
\end{isamarkuptext}\isamarkuptrue%
\ \ \isakeywordONE{lemma}\isamarkupfalse%
\ valid{\isacharunderscore}{\kern0pt}happ{\isacharunderscore}{\kern0pt}seq{\isacharunderscore}{\kern0pt}leq\isactrlsub i{\isacharunderscore}{\kern0pt}le\isactrlsub j{\isacharcolon}{\kern0pt}\isanewline
\ \ \ \ \isakeywordTWO{assumes}\ {\isachardoublequoteopen}consec{\isacharunderscore}{\kern0pt}htps\ {\isasympi}s\ t\isactrlsub i\ t\isactrlsub j{\isachardoublequoteclose}\ \isanewline
\ \ \ \ \ \ \ \ \isakeywordTWO{and}\ {\isachardoublequoteopen}wf{\isacharunderscore}{\kern0pt}plan\ {\isasympi}s{\isachardoublequoteclose}\isanewline
\ \ \ \ \ \ \ \ \isakeywordTWO{and}\ {\isachardoublequoteopen}ind{\isacharunderscore}{\kern0pt}happ{\isacharunderscore}{\kern0pt}seq\ {\isasympi}s\ hs\isactrlsub {\isadigit{1}}{\isachardoublequoteclose}\ \isanewline
\ \ \ \ \ \ \ \ \isakeywordTWO{and}\ {\isachardoublequoteopen}ind{\isacharunderscore}{\kern0pt}happ{\isacharunderscore}{\kern0pt}seq\ {\isasympi}s\ hs\isactrlsub {\isadigit{2}}{\isachardoublequoteclose}\ \isanewline
\ \ \ \ \ \ \ \ \isakeywordTWO{and}\ {\isachardoublequoteopen}valid{\isacharunderscore}{\kern0pt}happ{\isacharunderscore}{\kern0pt}seq\ M\isactrlsub i\ {\isacharparenleft}{\kern0pt}dropWhile\ {\isacharparenleft}{\kern0pt}leq\ t\isactrlsub i{\isacharparenright}{\kern0pt}\ {\isacharparenleft}{\kern0pt}takeWhile\ {\isacharparenleft}{\kern0pt}le\ t\isactrlsub j{\isacharparenright}{\kern0pt}\ hs\isactrlsub {\isadigit{1}}{\isacharparenright}{\kern0pt}{\isacharparenright}{\kern0pt}\ M\isactrlsub j{\isachardoublequoteclose}\isanewline
\ \ \ \ \isakeywordTWO{shows}\ {\isachardoublequoteopen}valid{\isacharunderscore}{\kern0pt}happ{\isacharunderscore}{\kern0pt}seq\ M\isactrlsub i\ {\isacharparenleft}{\kern0pt}dropWhile\ {\isacharparenleft}{\kern0pt}leq\ t\isactrlsub i{\isacharparenright}{\kern0pt}\ {\isacharparenleft}{\kern0pt}takeWhile\ {\isacharparenleft}{\kern0pt}le\ t\isactrlsub j{\isacharparenright}{\kern0pt}\ hs\isactrlsub {\isadigit{2}}{\isacharparenright}{\kern0pt}{\isacharparenright}{\kern0pt}\ M\isactrlsub j{\isachardoublequoteclose}\isanewline
%
\isadelimproof
\ \ %
\endisadelimproof
%
\isatagproof
\isakeywordONE{proof}\isamarkupfalse%
\ {\isacharminus}{\kern0pt}\isanewline
\ \ \ \ \isakeywordONE{let}\isamarkupfalse%
\ {\isacharquery}{\kern0pt}hs\isactrlsub {\isadigit{1}}{\isacharprime}{\kern0pt}{\isacharequal}{\kern0pt}{\isachardoublequoteopen}dropWhile\ {\isacharparenleft}{\kern0pt}leq\ t\isactrlsub i{\isacharparenright}{\kern0pt}\ {\isacharparenleft}{\kern0pt}takeWhile\ {\isacharparenleft}{\kern0pt}le\ t\isactrlsub j{\isacharparenright}{\kern0pt}\ hs\isactrlsub {\isadigit{1}}{\isacharparenright}{\kern0pt}{\isachardoublequoteclose}\isanewline
\ \ \ \ \isakeywordONE{let}\isamarkupfalse%
\ {\isacharquery}{\kern0pt}hs\isactrlsub {\isadigit{2}}{\isacharprime}{\kern0pt}{\isacharequal}{\kern0pt}{\isachardoublequoteopen}dropWhile\ {\isacharparenleft}{\kern0pt}leq\ t\isactrlsub i{\isacharparenright}{\kern0pt}\ {\isacharparenleft}{\kern0pt}takeWhile\ {\isacharparenleft}{\kern0pt}le\ t\isactrlsub j{\isacharparenright}{\kern0pt}\ hs\isactrlsub {\isadigit{2}}{\isacharparenright}{\kern0pt}{\isachardoublequoteclose}\isanewline
\ \ \ \ \isakeywordONE{have}\isamarkupfalse%
\ {\isachardoublequoteopen}{\isasymforall}h\ {\isasymin}\ set\ {\isacharquery}{\kern0pt}hs\isactrlsub {\isadigit{1}}{\isacharprime}{\kern0pt}{\isachardot}{\kern0pt}\ apply{\isacharunderscore}{\kern0pt}happ\ h\ M\isactrlsub i\ {\isacharequal}{\kern0pt}\ M\isactrlsub i{\isachardoublequoteclose}\isanewline
\ \ \ \ \ \ \isakeywordONE{using}\isamarkupfalse%
\ assms\ ind{\isacharunderscore}{\kern0pt}happ{\isacharunderscore}{\kern0pt}seq{\isacharunderscore}{\kern0pt}leq\isactrlsub i{\isacharunderscore}{\kern0pt}le\isactrlsub j\ \isakeywordONE{by}\isamarkupfalse%
\ simp\isanewline
\ \ \ \ \isakeywordONE{have}\isamarkupfalse%
\ {\isachardoublequoteopen}M\isactrlsub i\ {\isacharequal}{\kern0pt}\ M\isactrlsub j{\isachardoublequoteclose}\isanewline
\ \ \ \ \ \ \isakeywordONE{using}\isamarkupfalse%
\ valid{\isacharunderscore}{\kern0pt}happ{\isacharunderscore}{\kern0pt}seq{\isacharunderscore}{\kern0pt}M\isactrlsub j{\isacharunderscore}{\kern0pt}is{\isacharunderscore}{\kern0pt}M\isactrlsub i{\isacharbrackleft}{\kern0pt}OF\ {\isacartoucheopen}{\isasymforall}h\ {\isasymin}\ set\ {\isacharquery}{\kern0pt}hs\isactrlsub {\isadigit{1}}{\isacharprime}{\kern0pt}{\isachardot}{\kern0pt}\ apply{\isacharunderscore}{\kern0pt}happ\ h\ M\isactrlsub i\ {\isacharequal}{\kern0pt}\ M\isactrlsub i{\isacartoucheclose}{\isacharbrackright}{\kern0pt}\ \isanewline
\ \ \ \ \ \ \ \ \ \ \ \ {\isacartoucheopen}valid{\isacharunderscore}{\kern0pt}happ{\isacharunderscore}{\kern0pt}seq\ M\isactrlsub i\ {\isacharquery}{\kern0pt}hs\isactrlsub {\isadigit{1}}{\isacharprime}{\kern0pt}\ M\isactrlsub j{\isacartoucheclose}\ \isakeywordONE{by}\isamarkupfalse%
\ simp\isanewline
\ \ \ \ \isakeywordONE{have}\isamarkupfalse%
\ {\isachardoublequoteopen}{\isasymforall}{\isacharparenleft}{\kern0pt}t\isactrlsub a{\isacharcomma}{\kern0pt}A{\isacharparenright}{\kern0pt}\ {\isasymin}\ set\ {\isacharquery}{\kern0pt}hs\isactrlsub {\isadigit{1}}{\isacharprime}{\kern0pt}{\isachardot}{\kern0pt}\ {\isasymforall}a{\isasymin}\ set\ A{\isachardot}{\kern0pt}\ M\isactrlsub i\ \isactrlsup c{\isasymTTurnstile}\isactrlsub {\isacharequal}{\kern0pt}\ precondition\ a{\isachardoublequoteclose}\isanewline
\ \ \ \ \ \ \isakeywordONE{using}\isamarkupfalse%
\ assms\ valid{\isacharunderscore}{\kern0pt}happ{\isacharunderscore}{\kern0pt}seq{\isacharunderscore}{\kern0pt}leq\isactrlsub i{\isacharunderscore}{\kern0pt}le\isactrlsub j{\isacharunderscore}{\kern0pt}aux{\isadigit{1}}\ {\isacartoucheopen}valid{\isacharunderscore}{\kern0pt}happ{\isacharunderscore}{\kern0pt}seq\ M\isactrlsub i\ {\isacharquery}{\kern0pt}hs\isactrlsub {\isadigit{1}}{\isacharprime}{\kern0pt}\ M\isactrlsub j{\isacartoucheclose}\ \isakeywordONE{by}\isamarkupfalse%
\ simp\isanewline
\ \ \ \ \isakeywordONE{have}\isamarkupfalse%
\ {\isadigit{1}}{\isacharcolon}{\kern0pt}\ {\isachardoublequoteopen}{\isasymforall}h\ {\isasymin}\ set\ {\isacharquery}{\kern0pt}hs\isactrlsub {\isadigit{2}}{\isacharprime}{\kern0pt}{\isachardot}{\kern0pt}\ apply{\isacharunderscore}{\kern0pt}happ\ h\ M\isactrlsub i\ {\isacharequal}{\kern0pt}\ M\isactrlsub i{\isachardoublequoteclose}\isanewline
\ \ \ \ \ \ \isakeywordONE{using}\isamarkupfalse%
\ assms\ ind{\isacharunderscore}{\kern0pt}happ{\isacharunderscore}{\kern0pt}seq{\isacharunderscore}{\kern0pt}leq\isactrlsub i{\isacharunderscore}{\kern0pt}le\isactrlsub j\ \isakeywordONE{by}\isamarkupfalse%
\ simp\isanewline
\ \ \ \ \isakeywordONE{have}\isamarkupfalse%
\ {\isadigit{2}}{\isacharcolon}{\kern0pt}\ {\isachardoublequoteopen}{\isasymforall}{\isacharparenleft}{\kern0pt}t\isactrlsub a{\isacharcomma}{\kern0pt}A{\isacharparenright}{\kern0pt}\ {\isasymin}\ set\ {\isacharquery}{\kern0pt}hs\isactrlsub {\isadigit{2}}{\isacharprime}{\kern0pt}{\isachardot}{\kern0pt}\ {\isasymforall}a{\isasymin}\ set\ A{\isachardot}{\kern0pt}\ M\isactrlsub i\ \isactrlsup c{\isasymTTurnstile}\isactrlsub {\isacharequal}{\kern0pt}\ precondition\ a{\isachardoublequoteclose}\isanewline
\ \ \ \ \ \ \isakeywordONE{using}\isamarkupfalse%
\ assms\ ind{\isacharunderscore}{\kern0pt}happ{\isacharunderscore}{\kern0pt}seqs{\isacharunderscore}{\kern0pt}same{\isacharunderscore}{\kern0pt}acts{\isacharunderscore}{\kern0pt}leq\isactrlsub i{\isacharunderscore}{\kern0pt}le\isactrlsub j\isanewline
\ \ \ \ \ \ \ \ \ \ \ \ {\isacartoucheopen}{\isasymforall}{\isacharparenleft}{\kern0pt}t\isactrlsub a{\isacharcomma}{\kern0pt}A{\isacharparenright}{\kern0pt}\ {\isasymin}\ set\ {\isacharquery}{\kern0pt}hs\isactrlsub {\isadigit{1}}{\isacharprime}{\kern0pt}{\isachardot}{\kern0pt}\ {\isasymforall}a{\isasymin}\ set\ A{\isachardot}{\kern0pt}\ M\isactrlsub i\ \isactrlsup c{\isasymTTurnstile}\isactrlsub {\isacharequal}{\kern0pt}\ precondition\ a{\isacartoucheclose}\ \isakeywordONE{by}\isamarkupfalse%
\ blast\isanewline
\ \ \ \ \isakeywordONE{have}\isamarkupfalse%
\ {\isachardoublequoteopen}{\isasymforall}{\isacharparenleft}{\kern0pt}t\isactrlsub a{\isacharcomma}{\kern0pt}A{\isacharparenright}{\kern0pt}\ {\isasymin}\ set\ {\isacharquery}{\kern0pt}hs\isactrlsub {\isadigit{2}}{\isacharprime}{\kern0pt}{\isachardot}{\kern0pt}\ {\isasymforall}a\ {\isasymin}\ set\ A{\isachardot}{\kern0pt}\ effect\ a\ {\isacharequal}{\kern0pt}\ Effect\ {\isacharbrackleft}{\kern0pt}{\isacharbrackright}{\kern0pt}\ {\isacharbrackleft}{\kern0pt}{\isacharbrackright}{\kern0pt}{\isachardoublequoteclose}\isanewline
\ \ \ \ \ \ \isakeywordONE{using}\isamarkupfalse%
\ assms\ ind{\isacharunderscore}{\kern0pt}happ{\isacharunderscore}{\kern0pt}seq{\isacharunderscore}{\kern0pt}leq\isactrlsub i{\isacharunderscore}{\kern0pt}le\isactrlsub j{\isacharunderscore}{\kern0pt}empty{\isacharunderscore}{\kern0pt}eff\ \isakeywordONE{by}\isamarkupfalse%
\ auto\isanewline
\ \ \ \ \isakeywordONE{then}\isamarkupfalse%
\ \isakeywordONE{have}\isamarkupfalse%
\ {\isadigit{3}}{\isacharcolon}{\kern0pt}\ {\isachardoublequoteopen}{\isasymforall}{\isacharparenleft}{\kern0pt}t\isactrlsub a{\isacharcomma}{\kern0pt}A{\isacharparenright}{\kern0pt}\ {\isasymin}\ set\ {\isacharquery}{\kern0pt}hs\isactrlsub {\isadigit{2}}{\isacharprime}{\kern0pt}{\isachardot}{\kern0pt}\ {\isasymforall}{\isacharparenleft}{\kern0pt}t\isactrlsub a{\isacharprime}{\kern0pt}{\isacharcomma}{\kern0pt}A{\isacharprime}{\kern0pt}{\isacharparenright}{\kern0pt}\ {\isasymin}\ set\ {\isacharquery}{\kern0pt}hs\isactrlsub {\isadigit{2}}{\isacharprime}{\kern0pt}{\isachardot}{\kern0pt}\ \isanewline
\ \ \ \ \ \ \ \ \ \ \ \ \ \ \ \ \ \ \ \ {\isasymforall}a\ {\isasymin}\ set\ A{\isachardot}{\kern0pt}\ {\isasymforall}a{\isacharprime}{\kern0pt}\ {\isasymin}\ set\ A{\isacharprime}{\kern0pt}{\isachardot}{\kern0pt}\ acts{\isacharunderscore}{\kern0pt}non{\isacharunderscore}{\kern0pt}intrf\ a\ a{\isacharprime}{\kern0pt}{\isachardoublequoteclose}\isanewline
\ \ \ \ \ \ \isakeywordONE{unfolding}\isamarkupfalse%
\ acts{\isacharunderscore}{\kern0pt}non{\isacharunderscore}{\kern0pt}intrf{\isacharunderscore}{\kern0pt}def\ \isakeywordONE{by}\isamarkupfalse%
\ {\isacharparenleft}{\kern0pt}auto\ simp{\isacharcolon}{\kern0pt}\ Let{\isacharunderscore}{\kern0pt}def{\isacharparenright}{\kern0pt}\isanewline
\ \ \ \ \isakeywordONE{then}\isamarkupfalse%
\ \isakeywordTHREE{show}\isamarkupfalse%
\ {\isacharquery}{\kern0pt}thesis\isanewline
\ \ \ \ \ \ \isakeywordONE{using}\isamarkupfalse%
\ valid{\isacharunderscore}{\kern0pt}happ{\isacharunderscore}{\kern0pt}seq{\isacharunderscore}{\kern0pt}leq\isactrlsub i{\isacharunderscore}{\kern0pt}le\isactrlsub j{\isacharunderscore}{\kern0pt}aux{\isadigit{2}}{\isacharbrackleft}{\kern0pt}OF\ {\isadigit{1}}\ {\isadigit{2}}\ {\isadigit{3}}{\isacharbrackright}{\kern0pt}\ {\isacartoucheopen}M\isactrlsub i\ {\isacharequal}{\kern0pt}\ M\isactrlsub j{\isacartoucheclose}\ \isakeywordONE{by}\isamarkupfalse%
\ auto\isanewline
\ \ \isakeywordONE{qed}\isamarkupfalse%
%
\endisatagproof
{\isafoldproof}%
%
\isadelimproof
%
\endisadelimproof
%
\begin{isamarkuptext}%
An induced happening sequence contain exactly one happening for 
  every happening time point.%
\end{isamarkuptext}\isamarkuptrue%
\ \ \isakeywordONE{lemma}\isamarkupfalse%
\ ind{\isacharunderscore}{\kern0pt}happ{\isacharunderscore}{\kern0pt}seq{\isacharunderscore}{\kern0pt}le\isactrlsub i{\isacharunderscore}{\kern0pt}leq\isactrlsub i{\isacharcolon}{\kern0pt}\isanewline
\ \ \ \ \isakeywordTWO{assumes}\ {\isachardoublequoteopen}is{\isacharunderscore}{\kern0pt}htp\ {\isasympi}s\ t\isactrlsub i{\isachardoublequoteclose}\isanewline
\ \ \ \ \ \ \ \ \isakeywordTWO{and}\ {\isachardoublequoteopen}ind{\isacharunderscore}{\kern0pt}happ{\isacharunderscore}{\kern0pt}seq\ {\isasympi}s\ hs{\isachardoublequoteclose}\isanewline
\ \ \ \ \isakeywordTWO{shows}\ {\isachardoublequoteopen}{\isasymexists}A\isactrlsub i{\isachardot}{\kern0pt}\ dropWhile\ {\isacharparenleft}{\kern0pt}le\ t\isactrlsub i{\isacharparenright}{\kern0pt}\ {\isacharparenleft}{\kern0pt}takeWhile\ {\isacharparenleft}{\kern0pt}leq\ t\isactrlsub i{\isacharparenright}{\kern0pt}\ hs{\isacharparenright}{\kern0pt}\ {\isacharequal}{\kern0pt}\ {\isacharbrackleft}{\kern0pt}{\isacharparenleft}{\kern0pt}t\isactrlsub i{\isacharcomma}{\kern0pt}A\isactrlsub i{\isacharparenright}{\kern0pt}{\isacharbrackright}{\kern0pt}{\isachardoublequoteclose}\isanewline
%
\isadelimproof
\ \ %
\endisadelimproof
%
\isatagproof
\isakeywordONE{proof}\isamarkupfalse%
\ {\isacharminus}{\kern0pt}\isanewline
\ \ \ \ \isakeywordONE{have}\isamarkupfalse%
\ {\isachardoublequoteopen}strict{\isacharunderscore}{\kern0pt}sorted\ {\isacharparenleft}{\kern0pt}map\ fst\ hs{\isacharparenright}{\kern0pt}{\isachardoublequoteclose}\isanewline
\ \ \ \ \ \ \isakeywordONE{using}\isamarkupfalse%
\ {\isacartoucheopen}ind{\isacharunderscore}{\kern0pt}happ{\isacharunderscore}{\kern0pt}seq\ {\isasympi}s\ hs{\isacartoucheclose}\isanewline
\ \ \ \ \ \ \isakeywordONE{unfolding}\isamarkupfalse%
\ ind{\isacharunderscore}{\kern0pt}happ{\isacharunderscore}{\kern0pt}seq{\isacharunderscore}{\kern0pt}def\ \isakeywordONE{by}\isamarkupfalse%
\ auto\isanewline
\ \ \ \ \isakeywordONE{let}\isamarkupfalse%
\ {\isacharquery}{\kern0pt}hs{\isacharprime}{\kern0pt}{\isacharequal}{\kern0pt}{\isachardoublequoteopen}dropWhile\ {\isacharparenleft}{\kern0pt}le\ t\isactrlsub i{\isacharparenright}{\kern0pt}\ {\isacharparenleft}{\kern0pt}takeWhile\ {\isacharparenleft}{\kern0pt}leq\ t\isactrlsub i{\isacharparenright}{\kern0pt}\ hs{\isacharparenright}{\kern0pt}{\isachardoublequoteclose}\isanewline
\ \ \ \ \isakeywordONE{have}\isamarkupfalse%
\ {\isachardoublequoteopen}set\ {\isacharquery}{\kern0pt}hs{\isacharprime}{\kern0pt}\ {\isasymsubseteq}\ set\ hs{\isachardoublequoteclose}\isanewline
\ \ \ \ \ \ \isakeywordONE{using}\isamarkupfalse%
\ set{\isacharunderscore}{\kern0pt}dropWhileD\ set{\isacharunderscore}{\kern0pt}takeWhileD\ \isakeywordONE{by}\isamarkupfalse%
\ fastforce\isanewline
\ \ \ \ \isakeywordONE{have}\isamarkupfalse%
\ {\isachardoublequoteopen}distinct\ hs{\isachardoublequoteclose}\isanewline
\ \ \ \ \ \ \isakeywordONE{using}\isamarkupfalse%
\ {\isacartoucheopen}strict{\isacharunderscore}{\kern0pt}sorted\ {\isacharparenleft}{\kern0pt}map\ fst\ hs{\isacharparenright}{\kern0pt}{\isacartoucheclose}\ distinct{\isacharunderscore}{\kern0pt}mapI\ \isakeywordONE{by}\isamarkupfalse%
\ {\isacharparenleft}{\kern0pt}auto\ simp{\isacharcolon}{\kern0pt}\ strict{\isacharunderscore}{\kern0pt}sorted{\isacharunderscore}{\kern0pt}iff{\isacharparenright}{\kern0pt}\isanewline
\ \ \ \ \isakeywordONE{then}\isamarkupfalse%
\ \isakeywordONE{have}\isamarkupfalse%
\ {\isachardoublequoteopen}distinct\ {\isacharquery}{\kern0pt}hs{\isacharprime}{\kern0pt}{\isachardoublequoteclose}\isanewline
\ \ \ \ \ \ \isakeywordONE{using}\isamarkupfalse%
\ distinct{\isacharunderscore}{\kern0pt}dropWhile\ distinct{\isacharunderscore}{\kern0pt}takeWhile\ \isakeywordONE{by}\isamarkupfalse%
\ auto\isanewline
\ \ \ \ \isakeywordTHREE{obtain}\isamarkupfalse%
\ A\isactrlsub i\ \isakeywordTWO{where}\ {\isachardoublequoteopen}{\isacharparenleft}{\kern0pt}t\isactrlsub i{\isacharcomma}{\kern0pt}A\isactrlsub i{\isacharparenright}{\kern0pt}\ {\isasymin}\ set\ hs{\isachardoublequoteclose}\isanewline
\ \ \ \ \ \ \isakeywordONE{using}\isamarkupfalse%
\ assms{\isacharparenleft}{\kern0pt}{\isadigit{1}}{\isacharparenright}{\kern0pt}\ assms{\isacharparenleft}{\kern0pt}{\isadigit{2}}{\isacharparenright}{\kern0pt}\ ast{\isacharunderscore}{\kern0pt}problem{\isachardot}{\kern0pt}ind{\isacharunderscore}{\kern0pt}happ{\isacharunderscore}{\kern0pt}seq{\isacharunderscore}{\kern0pt}happ{\isacharunderscore}{\kern0pt}at{\isacharunderscore}{\kern0pt}htps\isanewline
\ \ \ \ \ \ \isakeywordONE{by}\isamarkupfalse%
\ blast\isanewline
\ \ \ \ \isakeywordONE{then}\isamarkupfalse%
\ \isakeywordONE{have}\isamarkupfalse%
\ {\isachardoublequoteopen}{\isacharparenleft}{\kern0pt}t\isactrlsub i{\isacharcomma}{\kern0pt}A\isactrlsub i{\isacharparenright}{\kern0pt}\ {\isasymin}\ set\ {\isacharquery}{\kern0pt}hs{\isacharprime}{\kern0pt}{\isachardoublequoteclose}\isanewline
\ \ \ \ \ \ \isakeywordONE{unfolding}\isamarkupfalse%
\ le{\isacharunderscore}{\kern0pt}def\ leq{\isacharunderscore}{\kern0pt}def\isanewline
\ \ \ \ \ \ \isakeywordONE{using}\isamarkupfalse%
\ {\isacartoucheopen}strict{\isacharunderscore}{\kern0pt}sorted\ {\isacharparenleft}{\kern0pt}map\ fst\ hs{\isacharparenright}{\kern0pt}{\isacartoucheclose}\ \isakeywordONE{by}\isamarkupfalse%
\ {\isacharparenleft}{\kern0pt}induction\ hs{\isacharparenright}{\kern0pt}\ auto\isanewline
\ \ \ \ \isakeywordONE{have}\isamarkupfalse%
\ {\isachardoublequoteopen}{\isasymforall}{\isacharparenleft}{\kern0pt}t\isactrlsub a{\isacharcomma}{\kern0pt}A{\isacharparenright}{\kern0pt}\ {\isasymin}\ set\ {\isacharquery}{\kern0pt}hs{\isacharprime}{\kern0pt}{\isachardot}{\kern0pt}\ t\isactrlsub a\ {\isacharequal}{\kern0pt}\ t\isactrlsub i{\isachardoublequoteclose}\isanewline
\ \ \ \ \ \ \isakeywordONE{using}\isamarkupfalse%
\ {\isacartoucheopen}set\ {\isacharquery}{\kern0pt}hs{\isacharprime}{\kern0pt}\ {\isasymsubseteq}\ set\ hs{\isacartoucheclose}\ strict{\isacharunderscore}{\kern0pt}sorted{\isacharunderscore}{\kern0pt}drop{\isacharunderscore}{\kern0pt}le{\isacharunderscore}{\kern0pt}take{\isacharunderscore}{\kern0pt}leq{\isacharbrackleft}{\kern0pt}OF\ {\isacartoucheopen}strict{\isacharunderscore}{\kern0pt}sorted\ {\isacharparenleft}{\kern0pt}map\ fst\ hs{\isacharparenright}{\kern0pt}{\isacartoucheclose}{\isacharbrackright}{\kern0pt}\isanewline
\ \ \ \ \ \ \isakeywordONE{by}\isamarkupfalse%
\ auto\isanewline
\ \ \ \ \isakeywordONE{have}\isamarkupfalse%
\ {\isachardoublequoteopen}{\isasymforall}{\isacharparenleft}{\kern0pt}t\isactrlsub {\isadigit{1}}{\isacharcomma}{\kern0pt}A\isactrlsub {\isadigit{1}}{\isacharparenright}{\kern0pt}\ {\isasymin}\ set\ hs{\isachardot}{\kern0pt}\ {\isasymforall}{\isacharparenleft}{\kern0pt}t\isactrlsub {\isadigit{2}}{\isacharcomma}{\kern0pt}A\isactrlsub {\isadigit{2}}{\isacharparenright}{\kern0pt}\ {\isasymin}\ set\ hs{\isachardot}{\kern0pt}\ t\isactrlsub {\isadigit{1}}\ {\isacharequal}{\kern0pt}\ t\isactrlsub {\isadigit{2}}\ {\isasymlongrightarrow}\ A\isactrlsub {\isadigit{1}}\ {\isacharequal}{\kern0pt}\ A\isactrlsub {\isadigit{2}}{\isachardoublequoteclose}\isanewline
\ \ \ \ \ \ \isakeywordONE{using}\isamarkupfalse%
\ {\isacartoucheopen}strict{\isacharunderscore}{\kern0pt}sorted\ {\isacharparenleft}{\kern0pt}map\ fst\ hs{\isacharparenright}{\kern0pt}{\isacartoucheclose}\ \isakeywordONE{by}\isamarkupfalse%
\ {\isacharparenleft}{\kern0pt}auto\ simp{\isacharcolon}{\kern0pt}\ strict{\isacharunderscore}{\kern0pt}sorted{\isacharunderscore}{\kern0pt}iff\ distinct{\isacharunderscore}{\kern0pt}map{\isacharunderscore}{\kern0pt}fstD{\isacharparenright}{\kern0pt}\isanewline
\ \ \ \ \isakeywordONE{then}\isamarkupfalse%
\ \isakeywordONE{have}\isamarkupfalse%
\ {\isachardoublequoteopen}{\isasymforall}{\isacharparenleft}{\kern0pt}t\isactrlsub {\isadigit{1}}{\isacharcomma}{\kern0pt}A\isactrlsub {\isadigit{1}}{\isacharparenright}{\kern0pt}\ {\isasymin}\ set\ {\isacharquery}{\kern0pt}hs{\isacharprime}{\kern0pt}{\isachardot}{\kern0pt}\ {\isasymforall}{\isacharparenleft}{\kern0pt}t\isactrlsub {\isadigit{2}}{\isacharcomma}{\kern0pt}A\isactrlsub {\isadigit{2}}{\isacharparenright}{\kern0pt}\ {\isasymin}\ set\ {\isacharquery}{\kern0pt}hs{\isacharprime}{\kern0pt}{\isachardot}{\kern0pt}\ t\isactrlsub {\isadigit{1}}\ {\isacharequal}{\kern0pt}\ t\isactrlsub {\isadigit{2}}\ {\isasymlongrightarrow}\ A\isactrlsub {\isadigit{1}}\ {\isacharequal}{\kern0pt}\ A\isactrlsub {\isadigit{2}}{\isachardoublequoteclose}\isanewline
\ \ \ \ \ \ \isakeywordONE{using}\isamarkupfalse%
\ {\isacartoucheopen}set\ {\isacharquery}{\kern0pt}hs{\isacharprime}{\kern0pt}\ {\isasymsubseteq}\ set\ hs{\isacartoucheclose}\ \isakeywordONE{by}\isamarkupfalse%
\ blast\isanewline
\ \ \ \ \isakeywordONE{then}\isamarkupfalse%
\ \isakeywordONE{have}\isamarkupfalse%
\ {\isachardoublequoteopen}{\isasymforall}{\isacharparenleft}{\kern0pt}t\isactrlsub a{\isacharcomma}{\kern0pt}as{\isacharprime}{\kern0pt}{\isacharparenright}{\kern0pt}\ {\isasymin}\ set\ {\isacharquery}{\kern0pt}hs{\isacharprime}{\kern0pt}{\isachardot}{\kern0pt}\ t\isactrlsub a\ {\isacharequal}{\kern0pt}\ t\isactrlsub i\ {\isasymand}\ as{\isacharprime}{\kern0pt}\ {\isacharequal}{\kern0pt}\ A\isactrlsub i{\isachardoublequoteclose}\isanewline
\ \ \ \ \ \ \isakeywordONE{using}\isamarkupfalse%
\ {\isacartoucheopen}{\isacharparenleft}{\kern0pt}t\isactrlsub i{\isacharcomma}{\kern0pt}A\isactrlsub i{\isacharparenright}{\kern0pt}\ {\isasymin}\ set\ {\isacharquery}{\kern0pt}hs{\isacharprime}{\kern0pt}{\isacartoucheclose}\ {\isacartoucheopen}{\isasymforall}{\isacharparenleft}{\kern0pt}t\isactrlsub a{\isacharcomma}{\kern0pt}as{\isacharparenright}{\kern0pt}\ {\isasymin}\ set\ {\isacharquery}{\kern0pt}hs{\isacharprime}{\kern0pt}{\isachardot}{\kern0pt}\ t\isactrlsub a\ {\isacharequal}{\kern0pt}\ t\isactrlsub i{\isacartoucheclose}\ \isakeywordONE{by}\isamarkupfalse%
\ auto\ force\isanewline
\ \ \ \ \isakeywordONE{then}\isamarkupfalse%
\ \isakeywordONE{have}\isamarkupfalse%
\ {\isachardoublequoteopen}set\ {\isacharquery}{\kern0pt}hs{\isacharprime}{\kern0pt}\ {\isacharequal}{\kern0pt}\ {\isacharbraceleft}{\kern0pt}{\isacharparenleft}{\kern0pt}t\isactrlsub i{\isacharcomma}{\kern0pt}A\isactrlsub i{\isacharparenright}{\kern0pt}{\isacharbraceright}{\kern0pt}{\isachardoublequoteclose}\isanewline
\ \ \ \ \ \ \isakeywordONE{using}\isamarkupfalse%
\ {\isacartoucheopen}{\isacharparenleft}{\kern0pt}t\isactrlsub i{\isacharcomma}{\kern0pt}A\isactrlsub i{\isacharparenright}{\kern0pt}\ {\isasymin}\ set\ {\isacharquery}{\kern0pt}hs{\isacharprime}{\kern0pt}{\isacartoucheclose}\ \isakeywordONE{by}\isamarkupfalse%
\ blast\isanewline
\ \ \ \ \isakeywordONE{then}\isamarkupfalse%
\ \isakeywordTHREE{show}\isamarkupfalse%
\ {\isacharquery}{\kern0pt}thesis\isanewline
\ \ \ \ \ \ \isakeywordONE{using}\isamarkupfalse%
\ distinct{\isacharunderscore}{\kern0pt}single{\isacharunderscore}{\kern0pt}set{\isacharunderscore}{\kern0pt}list{\isacharbrackleft}{\kern0pt}OF\ {\isacartoucheopen}set\ {\isacharquery}{\kern0pt}hs{\isacharprime}{\kern0pt}\ {\isacharequal}{\kern0pt}\ {\isacharbraceleft}{\kern0pt}{\isacharparenleft}{\kern0pt}t\isactrlsub i{\isacharcomma}{\kern0pt}A\isactrlsub i{\isacharparenright}{\kern0pt}{\isacharbraceright}{\kern0pt}{\isacartoucheclose}\ {\isacartoucheopen}distinct\ {\isacharquery}{\kern0pt}hs{\isacharprime}{\kern0pt}{\isacartoucheclose}{\isacharbrackright}{\kern0pt}\ \isakeywordONE{by}\isamarkupfalse%
\ simp\isanewline
\ \ \isakeywordONE{qed}\isamarkupfalse%
%
\endisatagproof
{\isafoldproof}%
%
\isadelimproof
%
\endisadelimproof
%
\begin{isamarkuptext}%
Two induced happening sequences contain the same ground actions for 
  every happening time point.%
\end{isamarkuptext}\isamarkuptrue%
\ \ \isakeywordONE{lemma}\isamarkupfalse%
\ ind{\isacharunderscore}{\kern0pt}happ{\isacharunderscore}{\kern0pt}seqs{\isacharunderscore}{\kern0pt}same{\isacharunderscore}{\kern0pt}acts{\isacharunderscore}{\kern0pt}le\isactrlsub i{\isacharunderscore}{\kern0pt}leq\isactrlsub i{\isacharcolon}{\kern0pt}\isanewline
\ \ \ \ \isakeywordTWO{assumes}\ {\isachardoublequoteopen}is{\isacharunderscore}{\kern0pt}htp\ {\isasympi}s\ t\isactrlsub i{\isachardoublequoteclose}\ \isanewline
\ \ \ \ \ \ \ \ \isakeywordTWO{and}\ {\isachardoublequoteopen}ind{\isacharunderscore}{\kern0pt}happ{\isacharunderscore}{\kern0pt}seq\ {\isasympi}s\ hs\isactrlsub {\isadigit{1}}{\isachardoublequoteclose}\ \isanewline
\ \ \ \ \ \ \ \ \isakeywordTWO{and}\ {\isachardoublequoteopen}ind{\isacharunderscore}{\kern0pt}happ{\isacharunderscore}{\kern0pt}seq\ {\isasympi}s\ hs\isactrlsub {\isadigit{2}}{\isachardoublequoteclose}\isanewline
\ \ \ \ \ \ \ \ \isakeywordTWO{and}\ {\isachardoublequoteopen}{\isacharparenleft}{\kern0pt}t\isactrlsub a{\isacharcomma}{\kern0pt}A{\isacharparenright}{\kern0pt}\ {\isasymin}\ set\ {\isacharparenleft}{\kern0pt}dropWhile\ {\isacharparenleft}{\kern0pt}le\ t\isactrlsub i{\isacharparenright}{\kern0pt}\ {\isacharparenleft}{\kern0pt}takeWhile\ {\isacharparenleft}{\kern0pt}leq\ t\isactrlsub i{\isacharparenright}{\kern0pt}\ hs\isactrlsub {\isadigit{1}}{\isacharparenright}{\kern0pt}{\isacharparenright}{\kern0pt}{\isachardoublequoteclose}\isanewline
\ \ \ \ \ \ \ \ \isakeywordTWO{and}\ {\isachardoublequoteopen}a\ {\isasymin}\ set\ A{\isachardoublequoteclose}\isanewline
\ \ \ \ \isakeywordTWO{shows}\ {\isachardoublequoteopen}{\isasymexists}t\isactrlsub a{\isacharprime}{\kern0pt}\ A{\isacharprime}{\kern0pt}{\isachardot}{\kern0pt}\ {\isacharparenleft}{\kern0pt}t\isactrlsub a{\isacharprime}{\kern0pt}{\isacharcomma}{\kern0pt}A{\isacharprime}{\kern0pt}{\isacharparenright}{\kern0pt}\ {\isasymin}\ set\ {\isacharparenleft}{\kern0pt}dropWhile\ {\isacharparenleft}{\kern0pt}le\ t\isactrlsub i{\isacharparenright}{\kern0pt}\ {\isacharparenleft}{\kern0pt}takeWhile\ {\isacharparenleft}{\kern0pt}leq\ t\isactrlsub i{\isacharparenright}{\kern0pt}\ hs\isactrlsub {\isadigit{2}}{\isacharparenright}{\kern0pt}{\isacharparenright}{\kern0pt}\ {\isasymand}\ a\ {\isasymin}\ set\ A{\isacharprime}{\kern0pt}{\isachardoublequoteclose}\isanewline
%
\isadelimproof
\ \ %
\endisadelimproof
%
\isatagproof
\isakeywordONE{proof}\isamarkupfalse%
\ {\isacharminus}{\kern0pt}\isanewline
\ \ \ \ \isakeywordONE{have}\isamarkupfalse%
\ {\isachardoublequoteopen}strict{\isacharunderscore}{\kern0pt}sorted\ {\isacharparenleft}{\kern0pt}map\ fst\ hs\isactrlsub {\isadigit{1}}{\isacharparenright}{\kern0pt}{\isachardoublequoteclose}\ \isakeywordTWO{and}\ {\isachardoublequoteopen}strict{\isacharunderscore}{\kern0pt}sorted\ {\isacharparenleft}{\kern0pt}map\ fst\ hs\isactrlsub {\isadigit{2}}{\isacharparenright}{\kern0pt}{\isachardoublequoteclose}\isanewline
\ \ \ \ \ \ \isakeywordONE{using}\isamarkupfalse%
\ {\isacartoucheopen}ind{\isacharunderscore}{\kern0pt}happ{\isacharunderscore}{\kern0pt}seq\ {\isasympi}s\ hs\isactrlsub {\isadigit{1}}{\isacartoucheclose}\ {\isacartoucheopen}ind{\isacharunderscore}{\kern0pt}happ{\isacharunderscore}{\kern0pt}seq\ {\isasympi}s\ hs\isactrlsub {\isadigit{2}}{\isacartoucheclose}\isanewline
\ \ \ \ \ \ \isakeywordONE{unfolding}\isamarkupfalse%
\ ind{\isacharunderscore}{\kern0pt}happ{\isacharunderscore}{\kern0pt}seq{\isacharunderscore}{\kern0pt}def\ \isakeywordONE{by}\isamarkupfalse%
\ auto\isanewline
\ \ \ \ \isakeywordONE{let}\isamarkupfalse%
\ {\isacharquery}{\kern0pt}hs\isactrlsub {\isadigit{1}}{\isacharprime}{\kern0pt}{\isacharequal}{\kern0pt}{\isachardoublequoteopen}dropWhile\ {\isacharparenleft}{\kern0pt}le\ t\isactrlsub i{\isacharparenright}{\kern0pt}\ {\isacharparenleft}{\kern0pt}takeWhile\ {\isacharparenleft}{\kern0pt}leq\ t\isactrlsub i{\isacharparenright}{\kern0pt}\ hs\isactrlsub {\isadigit{1}}{\isacharparenright}{\kern0pt}{\isachardoublequoteclose}\isanewline
\ \ \ \ \isakeywordONE{let}\isamarkupfalse%
\ {\isacharquery}{\kern0pt}hs\isactrlsub {\isadigit{2}}{\isacharprime}{\kern0pt}{\isacharequal}{\kern0pt}{\isachardoublequoteopen}dropWhile\ {\isacharparenleft}{\kern0pt}le\ t\isactrlsub i{\isacharparenright}{\kern0pt}\ {\isacharparenleft}{\kern0pt}takeWhile\ {\isacharparenleft}{\kern0pt}leq\ t\isactrlsub i{\isacharparenright}{\kern0pt}\ hs\isactrlsub {\isadigit{2}}{\isacharparenright}{\kern0pt}{\isachardoublequoteclose}\isanewline
\ \ \ \ \isakeywordONE{have}\isamarkupfalse%
\ {\isachardoublequoteopen}{\isacharparenleft}{\kern0pt}t\isactrlsub a{\isacharcomma}{\kern0pt}A{\isacharparenright}{\kern0pt}\ {\isasymin}\ set\ hs\isactrlsub {\isadigit{1}}{\isachardoublequoteclose}\isanewline
\ \ \ \ \ \ \isakeywordONE{using}\isamarkupfalse%
\ {\isacartoucheopen}{\isacharparenleft}{\kern0pt}t\isactrlsub a{\isacharcomma}{\kern0pt}A{\isacharparenright}{\kern0pt}\ {\isasymin}\ set\ {\isacharquery}{\kern0pt}hs\isactrlsub {\isadigit{1}}{\isacharprime}{\kern0pt}{\isacartoucheclose}\ \isakeywordONE{by}\isamarkupfalse%
\ {\isacharparenleft}{\kern0pt}meson\ set{\isacharunderscore}{\kern0pt}dropWhileD\ set{\isacharunderscore}{\kern0pt}takeWhileD{\isacharparenright}{\kern0pt}\isanewline
\ \ \ \ \isakeywordONE{have}\isamarkupfalse%
\ {\isachardoublequoteopen}t\isactrlsub i\ {\isacharequal}{\kern0pt}\ t\isactrlsub a{\isachardoublequoteclose}\isanewline
\ \ \ \ \ \ \isakeywordONE{using}\isamarkupfalse%
\ {\isacartoucheopen}{\isacharparenleft}{\kern0pt}t\isactrlsub a{\isacharcomma}{\kern0pt}A{\isacharparenright}{\kern0pt}\ {\isasymin}\ set\ {\isacharquery}{\kern0pt}hs\isactrlsub {\isadigit{1}}{\isacharprime}{\kern0pt}{\isacartoucheclose}\isanewline
\ \ \ \ \ \ \ \ \ \ \ \ strict{\isacharunderscore}{\kern0pt}sorted{\isacharunderscore}{\kern0pt}drop{\isacharunderscore}{\kern0pt}le{\isacharunderscore}{\kern0pt}take{\isacharunderscore}{\kern0pt}leq{\isacharbrackleft}{\kern0pt}OF\ {\isacartoucheopen}strict{\isacharunderscore}{\kern0pt}sorted\ {\isacharparenleft}{\kern0pt}map\ fst\ hs\isactrlsub {\isadigit{1}}{\isacharparenright}{\kern0pt}{\isacartoucheclose}\ {\isacartoucheopen}{\isacharparenleft}{\kern0pt}t\isactrlsub a{\isacharcomma}{\kern0pt}A{\isacharparenright}{\kern0pt}\ {\isasymin}\ set\ hs\isactrlsub {\isadigit{1}}{\isacartoucheclose}{\isacharbrackright}{\kern0pt}\isanewline
\ \ \ \ \ \ \isakeywordONE{by}\isamarkupfalse%
\ auto\isanewline
\ \ \ \ \isakeywordONE{have}\isamarkupfalse%
\ {\isachardoublequoteopen}{\isasymexists}{\isacharparenleft}{\kern0pt}t\isactrlsub {\isasympi}{\isacharcomma}{\kern0pt}{\isasympi}{\isacharparenright}{\kern0pt}\ {\isasymin}\ set\ {\isasympi}s{\isachardot}{\kern0pt}\ inst{\isacharunderscore}{\kern0pt}of{\isacharunderscore}{\kern0pt}plan{\isacharunderscore}{\kern0pt}action\ {\isasympi}s\ {\isacharparenleft}{\kern0pt}t\isactrlsub {\isasympi}{\isacharcomma}{\kern0pt}{\isasympi}{\isacharparenright}{\kern0pt}\ {\isacharparenleft}{\kern0pt}t\isactrlsub a{\isacharcomma}{\kern0pt}a{\isacharparenright}{\kern0pt}{\isachardoublequoteclose}\isanewline
\ \ \ \ \ \ \isakeywordONE{using}\isamarkupfalse%
\ {\isacartoucheopen}ind{\isacharunderscore}{\kern0pt}happ{\isacharunderscore}{\kern0pt}seq\ {\isasympi}s\ hs\isactrlsub {\isadigit{1}}{\isacartoucheclose}\ {\isacartoucheopen}{\isacharparenleft}{\kern0pt}t\isactrlsub a{\isacharcomma}{\kern0pt}A{\isacharparenright}{\kern0pt}\ {\isasymin}\ set\ hs\isactrlsub {\isadigit{1}}{\isacartoucheclose}\ {\isacartoucheopen}a\ {\isasymin}\ set\ A{\isacartoucheclose}\isanewline
\ \ \ \ \ \ \isakeywordONE{unfolding}\isamarkupfalse%
\ ind{\isacharunderscore}{\kern0pt}happ{\isacharunderscore}{\kern0pt}seq{\isacharunderscore}{\kern0pt}def\ \isakeywordONE{by}\isamarkupfalse%
\ auto\isanewline
\ \ \ \ \isakeywordONE{then}\isamarkupfalse%
\ \isakeywordTHREE{obtain}\isamarkupfalse%
\ t\isactrlsub {\isasympi}\ {\isasympi}\ \isakeywordTWO{where}\ {\isachardoublequoteopen}{\isacharparenleft}{\kern0pt}t\isactrlsub {\isasympi}{\isacharcomma}{\kern0pt}{\isasympi}{\isacharparenright}{\kern0pt}\ {\isasymin}\ set\ {\isasympi}s{\isachardoublequoteclose}\ \isakeywordTWO{and}\ {\isachardoublequoteopen}inst{\isacharunderscore}{\kern0pt}of{\isacharunderscore}{\kern0pt}plan{\isacharunderscore}{\kern0pt}action\ {\isasympi}s\ {\isacharparenleft}{\kern0pt}t\isactrlsub {\isasympi}{\isacharcomma}{\kern0pt}{\isasympi}{\isacharparenright}{\kern0pt}\ {\isacharparenleft}{\kern0pt}t\isactrlsub a{\isacharcomma}{\kern0pt}a{\isacharparenright}{\kern0pt}{\isachardoublequoteclose}\isanewline
\ \ \ \ \ \ \isakeywordONE{by}\isamarkupfalse%
\ auto\isanewline
\ \ \ \ \isakeywordONE{let}\isamarkupfalse%
\ {\isacharquery}{\kern0pt}case{\isadigit{1}}{\isacharequal}{\kern0pt}{\isachardoublequoteopen}res{\isacharunderscore}{\kern0pt}inst\ {\isasympi}\ At{\isacharunderscore}{\kern0pt}Start\ {\isacharequal}{\kern0pt}\ Some\ a\ {\isasymand}\ t\isactrlsub {\isasympi}\ {\isacharequal}{\kern0pt}\ t\isactrlsub a{\isachardoublequoteclose}\isanewline
\ \ \ \ \isakeywordONE{let}\isamarkupfalse%
\ {\isacharquery}{\kern0pt}case{\isadigit{2}}{\isacharequal}{\kern0pt}{\isachardoublequoteopen}res{\isacharunderscore}{\kern0pt}inst{\isacharunderscore}{\kern0pt}snap{\isacharunderscore}{\kern0pt}action\ {\isasympi}\ At{\isacharunderscore}{\kern0pt}Start\ {\isacharequal}{\kern0pt}\ Some\ a\ {\isasymand}\ t\isactrlsub {\isasympi}\ {\isacharequal}{\kern0pt}\ t\isactrlsub a{\isachardoublequoteclose}\isanewline
\ \ \ \ \isakeywordONE{let}\isamarkupfalse%
\ {\isacharquery}{\kern0pt}case{\isadigit{3}}{\isacharequal}{\kern0pt}{\isachardoublequoteopen}res{\isacharunderscore}{\kern0pt}inst{\isacharunderscore}{\kern0pt}snap{\isacharunderscore}{\kern0pt}action\ {\isasympi}\ At{\isacharunderscore}{\kern0pt}End\ {\isacharequal}{\kern0pt}\ Some\ a\ {\isasymand}\ t\isactrlsub {\isasympi}\ {\isacharplus}{\kern0pt}\ {\isacharparenleft}{\kern0pt}duration\ {\isasympi}{\isacharparenright}{\kern0pt}\ {\isacharequal}{\kern0pt}\ t\isactrlsub a{\isachardoublequoteclose}\isanewline
\ \ \ \ \isakeywordONE{let}\isamarkupfalse%
\ {\isacharquery}{\kern0pt}case{\isadigit{4}}{\isacharequal}{\kern0pt}{\isachardoublequoteopen}res{\isacharunderscore}{\kern0pt}inst{\isacharunderscore}{\kern0pt}snap{\isacharunderscore}{\kern0pt}action\ {\isasympi}\ Over{\isacharunderscore}{\kern0pt}All\ {\isacharequal}{\kern0pt}\ Some\ a\ {\isasymand}\ \isanewline
\ \ \ \ \ \ \ \ \ \ {\isacharparenleft}{\kern0pt}{\isasymexists}t\isactrlsub i\ t\isactrlsub j{\isachardot}{\kern0pt}\ consec{\isacharunderscore}{\kern0pt}htps\ {\isasympi}s\ t\isactrlsub i\ t\isactrlsub j\ {\isasymand}\ t\isactrlsub {\isasympi}\ {\isasymle}\ t\isactrlsub i\ {\isasymand}\ t\isactrlsub j\ {\isasymle}\ t\isactrlsub {\isasympi}\ {\isacharplus}{\kern0pt}\ {\isacharparenleft}{\kern0pt}duration\ {\isasympi}{\isacharparenright}{\kern0pt}\ {\isasymand}\ t\isactrlsub i\ {\isacharless}{\kern0pt}\ t\isactrlsub a\ {\isasymand}\ t\isactrlsub a\ {\isacharless}{\kern0pt}\ t\isactrlsub j{\isacharparenright}{\kern0pt}{\isachardoublequoteclose}\isanewline
\ \ \ \ \isakeywordONE{have}\isamarkupfalse%
\ {\isachardoublequoteopen}{\isacharquery}{\kern0pt}case{\isadigit{1}}\ {\isasymor}\ {\isacharquery}{\kern0pt}case{\isadigit{2}}\ {\isasymor}\ {\isacharquery}{\kern0pt}case{\isadigit{3}}\ {\isasymor}\ {\isacharquery}{\kern0pt}case{\isadigit{4}}{\isachardoublequoteclose}\isanewline
\ \ \ \ \ \ \isakeywordONE{using}\isamarkupfalse%
\ {\isacartoucheopen}inst{\isacharunderscore}{\kern0pt}of{\isacharunderscore}{\kern0pt}plan{\isacharunderscore}{\kern0pt}action\ {\isasympi}s\ {\isacharparenleft}{\kern0pt}t\isactrlsub {\isasympi}{\isacharcomma}{\kern0pt}{\isasympi}{\isacharparenright}{\kern0pt}\ {\isacharparenleft}{\kern0pt}t\isactrlsub a{\isacharcomma}{\kern0pt}a{\isacharparenright}{\kern0pt}{\isacartoucheclose}\ \isakeywordONE{by}\isamarkupfalse%
\ simp\isanewline
\ \ \ \ \isakeywordONE{then}\isamarkupfalse%
\ \isakeywordTHREE{show}\isamarkupfalse%
\ {\isacharquery}{\kern0pt}thesis\isanewline
\ \ \ \ \ \ \isakeywordONE{proof}\isamarkupfalse%
\ {\isacharparenleft}{\kern0pt}elim\ disjE{\isacharparenright}{\kern0pt}\isanewline
\ \ \ \ \ \ \ \ \isakeywordTHREE{assume}\isamarkupfalse%
\ {\isacharquery}{\kern0pt}case{\isadigit{1}}\isanewline
\ \ \ \ \ \ \ \ \isakeywordONE{then}\isamarkupfalse%
\ \isakeywordONE{have}\isamarkupfalse%
\ {\isachardoublequoteopen}{\isacharparenleft}{\kern0pt}t\isactrlsub {\isasympi}{\isacharcomma}{\kern0pt}{\isasympi}{\isacharparenright}{\kern0pt}\ {\isasymin}\ simple{\isacharunderscore}{\kern0pt}acts\ {\isasympi}s{\isachardoublequoteclose}\isanewline
\ \ \ \ \ \ \ \ \ \ \isakeywordONE{unfolding}\isamarkupfalse%
\ simple{\isacharunderscore}{\kern0pt}acts{\isacharunderscore}{\kern0pt}def\ is{\isacharunderscore}{\kern0pt}act{\isacharunderscore}{\kern0pt}simple{\isacharunderscore}{\kern0pt}def\ \isakeywordONE{using}\isamarkupfalse%
\ {\isacartoucheopen}{\isacharparenleft}{\kern0pt}t\isactrlsub {\isasympi}{\isacharcomma}{\kern0pt}{\isasympi}{\isacharparenright}{\kern0pt}\ {\isasymin}\ set\ {\isasympi}s{\isacartoucheclose}\ \isakeywordONE{by}\isamarkupfalse%
\ {\isacharparenleft}{\kern0pt}cases\ {\isasympi}{\isacharparenright}{\kern0pt}\ auto\isanewline
\ \ \ \ \ \ \ \ \isakeywordONE{then}\isamarkupfalse%
\ \isakeywordONE{have}\isamarkupfalse%
\ {\isachardoublequoteopen}{\isasymexists}as{\isachardot}{\kern0pt}\ {\isacharparenleft}{\kern0pt}t\isactrlsub {\isasympi}{\isacharcomma}{\kern0pt}as{\isacharparenright}{\kern0pt}\ {\isasymin}\ set\ hs\isactrlsub {\isadigit{2}}\ {\isasymand}\ the\ {\isacharparenleft}{\kern0pt}res{\isacharunderscore}{\kern0pt}inst\ {\isasympi}\ At{\isacharunderscore}{\kern0pt}Start{\isacharparenright}{\kern0pt}\ {\isasymin}\ set\ as{\isachardoublequoteclose}\isanewline
\ \ \ \ \ \ \ \ \ \ \isakeywordONE{using}\isamarkupfalse%
\ {\isacartoucheopen}ind{\isacharunderscore}{\kern0pt}happ{\isacharunderscore}{\kern0pt}seq\ {\isasympi}s\ hs\isactrlsub {\isadigit{2}}{\isacartoucheclose}\ \isakeywordONE{unfolding}\isamarkupfalse%
\ ind{\isacharunderscore}{\kern0pt}happ{\isacharunderscore}{\kern0pt}seq{\isacharunderscore}{\kern0pt}def\ \isakeywordONE{by}\isamarkupfalse%
\ auto\isanewline
\ \ \ \ \ \ \ \ \isakeywordONE{then}\isamarkupfalse%
\ \isakeywordTHREE{obtain}\isamarkupfalse%
\ A{\isacharprime}{\kern0pt}\ \isakeywordTWO{where}\ {\isachardoublequoteopen}{\isacharparenleft}{\kern0pt}t\isactrlsub {\isasympi}{\isacharcomma}{\kern0pt}A{\isacharprime}{\kern0pt}{\isacharparenright}{\kern0pt}\ {\isasymin}\ set\ hs\isactrlsub {\isadigit{2}}{\isachardoublequoteclose}\ \isakeywordTWO{and}\ {\isachardoublequoteopen}the\ {\isacharparenleft}{\kern0pt}res{\isacharunderscore}{\kern0pt}inst\ {\isasympi}\ At{\isacharunderscore}{\kern0pt}Start{\isacharparenright}{\kern0pt}\ {\isasymin}\ set\ A{\isacharprime}{\kern0pt}{\isachardoublequoteclose}\ \isanewline
\ \ \ \ \ \ \ \ \ \ \isakeywordONE{by}\isamarkupfalse%
\ auto\isanewline
\ \ \ \ \ \ \ \ \isakeywordONE{then}\isamarkupfalse%
\ \isakeywordONE{have}\isamarkupfalse%
\ {\isachardoublequoteopen}{\isacharparenleft}{\kern0pt}t\isactrlsub a{\isacharcomma}{\kern0pt}A{\isacharprime}{\kern0pt}{\isacharparenright}{\kern0pt}\ {\isasymin}\ set\ hs\isactrlsub {\isadigit{2}}{\isachardoublequoteclose}\isanewline
\ \ \ \ \ \ \ \ \ \ \isakeywordONE{using}\isamarkupfalse%
\ {\isacartoucheopen}{\isacharquery}{\kern0pt}case{\isadigit{1}}{\isacartoucheclose}\ \isakeywordONE{by}\isamarkupfalse%
\ auto\isanewline
\ \ \ \ \ \ \ \ \isakeywordTHREE{show}\isamarkupfalse%
\ {\isacharquery}{\kern0pt}thesis\isanewline
\ \ \ \ \ \ \ \ \ \ \isakeywordONE{using}\isamarkupfalse%
\ {\isacartoucheopen}t\isactrlsub i\ {\isacharequal}{\kern0pt}\ t\isactrlsub a{\isacartoucheclose}\ {\isacartoucheopen}{\isacharquery}{\kern0pt}case{\isadigit{1}}{\isacartoucheclose}\ {\isacartoucheopen}the\ {\isacharparenleft}{\kern0pt}res{\isacharunderscore}{\kern0pt}inst\ {\isasympi}\ At{\isacharunderscore}{\kern0pt}Start{\isacharparenright}{\kern0pt}\ {\isasymin}\ set\ A{\isacharprime}{\kern0pt}{\isacartoucheclose}\isanewline
\ \ \ \ \ \ \ \ \ \ \ \ \ \ strict{\isacharunderscore}{\kern0pt}sorted{\isacharunderscore}{\kern0pt}drop{\isacharunderscore}{\kern0pt}le{\isacharunderscore}{\kern0pt}take{\isacharunderscore}{\kern0pt}leq{\isacharbrackleft}{\kern0pt}OF\ {\isacartoucheopen}strict{\isacharunderscore}{\kern0pt}sorted\ {\isacharparenleft}{\kern0pt}map\ fst\ hs\isactrlsub {\isadigit{2}}{\isacharparenright}{\kern0pt}{\isacartoucheclose}\ {\isacartoucheopen}{\isacharparenleft}{\kern0pt}t\isactrlsub a{\isacharcomma}{\kern0pt}A{\isacharprime}{\kern0pt}{\isacharparenright}{\kern0pt}\ {\isasymin}\ set\ hs\isactrlsub {\isadigit{2}}{\isacartoucheclose}{\isacharbrackright}{\kern0pt}\isanewline
\ \ \ \ \ \ \ \ \ \ \isakeywordONE{by}\isamarkupfalse%
\ auto\isanewline
\ \ \ \ \ \ \isakeywordONE{next}\isamarkupfalse%
\isanewline
\ \ \ \ \ \ \ \ \isakeywordTHREE{assume}\isamarkupfalse%
\ {\isacharquery}{\kern0pt}case{\isadigit{2}}\isanewline
\ \ \ \ \ \ \ \ \isakeywordONE{let}\isamarkupfalse%
\ {\isacharquery}{\kern0pt}a\isactrlsub s\isactrlsub t\isactrlsub a\isactrlsub r\isactrlsub t{\isacharequal}{\kern0pt}{\isachardoublequoteopen}the\ {\isacharparenleft}{\kern0pt}res{\isacharunderscore}{\kern0pt}inst{\isacharunderscore}{\kern0pt}snap{\isacharunderscore}{\kern0pt}action\ {\isasympi}\ At{\isacharunderscore}{\kern0pt}Start{\isacharparenright}{\kern0pt}{\isachardoublequoteclose}\isanewline
\ \ \ \ \ \ \ \ \isakeywordONE{have}\isamarkupfalse%
\ {\isachardoublequoteopen}{\isacharparenleft}{\kern0pt}t\isactrlsub {\isasympi}{\isacharcomma}{\kern0pt}{\isasympi}{\isacharparenright}{\kern0pt}\ {\isasymin}\ durative{\isacharunderscore}{\kern0pt}acts\ {\isasympi}s{\isachardoublequoteclose}\isanewline
\ \ \ \ \ \ \ \ \ \ \isakeywordONE{unfolding}\isamarkupfalse%
\ durative{\isacharunderscore}{\kern0pt}acts{\isacharunderscore}{\kern0pt}def\ is{\isacharunderscore}{\kern0pt}act{\isacharunderscore}{\kern0pt}simple{\isacharunderscore}{\kern0pt}def\isanewline
\ \ \ \ \ \ \ \ \ \ \isakeywordONE{using}\isamarkupfalse%
\ {\isacartoucheopen}{\isacharquery}{\kern0pt}case{\isadigit{2}}{\isacartoucheclose}\ {\isacartoucheopen}{\isacharparenleft}{\kern0pt}t\isactrlsub {\isasympi}{\isacharcomma}{\kern0pt}{\isasympi}{\isacharparenright}{\kern0pt}\ {\isasymin}\ set\ {\isasympi}s{\isacartoucheclose}\ \isakeywordONE{by}\isamarkupfalse%
\ {\isacharparenleft}{\kern0pt}cases\ {\isasympi}{\isacharparenright}{\kern0pt}\ auto\isanewline
\ \ \ \ \ \ \ \ \isakeywordONE{then}\isamarkupfalse%
\ \isakeywordONE{have}\isamarkupfalse%
\ {\isachardoublequoteopen}{\isasymexists}as{\isachardot}{\kern0pt}\ {\isacharparenleft}{\kern0pt}t\isactrlsub {\isasympi}{\isacharcomma}{\kern0pt}as{\isacharparenright}{\kern0pt}\ {\isasymin}\ set\ hs\isactrlsub {\isadigit{2}}\ {\isasymand}\ {\isacharquery}{\kern0pt}a\isactrlsub s\isactrlsub t\isactrlsub a\isactrlsub r\isactrlsub t\ {\isasymin}\ set\ as{\isachardoublequoteclose}\isanewline
\ \ \ \ \ \ \ \ \ \ \isakeywordONE{using}\isamarkupfalse%
\ {\isacartoucheopen}ind{\isacharunderscore}{\kern0pt}happ{\isacharunderscore}{\kern0pt}seq\ {\isasympi}s\ hs\isactrlsub {\isadigit{2}}{\isacartoucheclose}\ \isakeywordONE{unfolding}\isamarkupfalse%
\ ind{\isacharunderscore}{\kern0pt}happ{\isacharunderscore}{\kern0pt}seq{\isacharunderscore}{\kern0pt}def\ \isakeywordONE{by}\isamarkupfalse%
\ auto\isanewline
\ \ \ \ \ \ \ \ \isakeywordONE{then}\isamarkupfalse%
\ \isakeywordTHREE{obtain}\isamarkupfalse%
\ A{\isacharprime}{\kern0pt}\ \isakeywordTWO{where}\ {\isachardoublequoteopen}{\isacharparenleft}{\kern0pt}t\isactrlsub {\isasympi}{\isacharcomma}{\kern0pt}A{\isacharprime}{\kern0pt}{\isacharparenright}{\kern0pt}\ {\isasymin}\ set\ hs\isactrlsub {\isadigit{2}}{\isachardoublequoteclose}\ \isakeywordTWO{and}\ {\isachardoublequoteopen}{\isacharquery}{\kern0pt}a\isactrlsub s\isactrlsub t\isactrlsub a\isactrlsub r\isactrlsub t\ {\isasymin}\ set\ A{\isacharprime}{\kern0pt}{\isachardoublequoteclose}\ \isanewline
\ \ \ \ \ \ \ \ \ \ \isakeywordONE{by}\isamarkupfalse%
\ auto\isanewline
\ \ \ \ \ \ \ \ \isakeywordONE{then}\isamarkupfalse%
\ \isakeywordONE{have}\isamarkupfalse%
\ {\isachardoublequoteopen}{\isacharparenleft}{\kern0pt}t\isactrlsub a{\isacharcomma}{\kern0pt}A{\isacharprime}{\kern0pt}{\isacharparenright}{\kern0pt}\ {\isasymin}\ set\ hs\isactrlsub {\isadigit{2}}{\isachardoublequoteclose}\isanewline
\ \ \ \ \ \ \ \ \ \ \isakeywordONE{using}\isamarkupfalse%
\ {\isacartoucheopen}{\isacharquery}{\kern0pt}case{\isadigit{2}}{\isacartoucheclose}\ \isakeywordONE{by}\isamarkupfalse%
\ auto\isanewline
\ \ \ \ \ \ \ \ \isakeywordTHREE{show}\isamarkupfalse%
\ {\isacharquery}{\kern0pt}thesis\isanewline
\ \ \ \ \ \ \ \ \ \ \isakeywordONE{using}\isamarkupfalse%
\ {\isacartoucheopen}t\isactrlsub i\ {\isacharequal}{\kern0pt}\ t\isactrlsub a{\isacartoucheclose}\ {\isacartoucheopen}{\isacharquery}{\kern0pt}case{\isadigit{2}}{\isacartoucheclose}\ {\isacartoucheopen}{\isacharquery}{\kern0pt}a\isactrlsub s\isactrlsub t\isactrlsub a\isactrlsub r\isactrlsub t\ {\isasymin}\ set\ A{\isacharprime}{\kern0pt}{\isacartoucheclose}\isanewline
\ \ \ \ \ \ \ \ \ \ \ \ \ \ strict{\isacharunderscore}{\kern0pt}sorted{\isacharunderscore}{\kern0pt}drop{\isacharunderscore}{\kern0pt}le{\isacharunderscore}{\kern0pt}take{\isacharunderscore}{\kern0pt}leq{\isacharbrackleft}{\kern0pt}OF\ {\isacartoucheopen}strict{\isacharunderscore}{\kern0pt}sorted\ {\isacharparenleft}{\kern0pt}map\ fst\ hs\isactrlsub {\isadigit{2}}{\isacharparenright}{\kern0pt}{\isacartoucheclose}\ {\isacartoucheopen}{\isacharparenleft}{\kern0pt}t\isactrlsub a{\isacharcomma}{\kern0pt}A{\isacharprime}{\kern0pt}{\isacharparenright}{\kern0pt}\ {\isasymin}\ set\ hs\isactrlsub {\isadigit{2}}{\isacartoucheclose}{\isacharbrackright}{\kern0pt}\isanewline
\ \ \ \ \ \ \ \ \ \ \isakeywordONE{by}\isamarkupfalse%
\ auto\isanewline
\ \ \ \ \isakeywordONE{next}\isamarkupfalse%
\isanewline
\ \ \ \ \ \ \isakeywordTHREE{assume}\isamarkupfalse%
\ {\isacharquery}{\kern0pt}case{\isadigit{3}}\isanewline
\ \ \ \ \ \ \ \ \isakeywordONE{let}\isamarkupfalse%
\ {\isacharquery}{\kern0pt}a\isactrlsub e\isactrlsub n\isactrlsub d{\isacharequal}{\kern0pt}{\isachardoublequoteopen}the\ {\isacharparenleft}{\kern0pt}res{\isacharunderscore}{\kern0pt}inst{\isacharunderscore}{\kern0pt}snap{\isacharunderscore}{\kern0pt}action\ {\isasympi}\ At{\isacharunderscore}{\kern0pt}End{\isacharparenright}{\kern0pt}{\isachardoublequoteclose}\isanewline
\ \ \ \ \ \ \ \ \isakeywordONE{have}\isamarkupfalse%
\ {\isachardoublequoteopen}{\isacharparenleft}{\kern0pt}t\isactrlsub {\isasympi}{\isacharcomma}{\kern0pt}{\isasympi}{\isacharparenright}{\kern0pt}\ {\isasymin}\ durative{\isacharunderscore}{\kern0pt}acts\ {\isasympi}s{\isachardoublequoteclose}\isanewline
\ \ \ \ \ \ \ \ \ \ \isakeywordONE{unfolding}\isamarkupfalse%
\ durative{\isacharunderscore}{\kern0pt}acts{\isacharunderscore}{\kern0pt}def\ is{\isacharunderscore}{\kern0pt}act{\isacharunderscore}{\kern0pt}simple{\isacharunderscore}{\kern0pt}def\isanewline
\ \ \ \ \ \ \ \ \ \ \isakeywordONE{using}\isamarkupfalse%
\ {\isacartoucheopen}{\isacharquery}{\kern0pt}case{\isadigit{3}}{\isacartoucheclose}\ {\isacartoucheopen}{\isacharparenleft}{\kern0pt}t\isactrlsub {\isasympi}{\isacharcomma}{\kern0pt}{\isasympi}{\isacharparenright}{\kern0pt}\ {\isasymin}\ set\ {\isasympi}s{\isacartoucheclose}\ \isakeywordONE{by}\isamarkupfalse%
\ {\isacharparenleft}{\kern0pt}cases\ {\isasympi}{\isacharparenright}{\kern0pt}\ auto\isanewline
\ \ \ \ \ \ \ \ \isakeywordONE{then}\isamarkupfalse%
\ \isakeywordONE{have}\isamarkupfalse%
\ {\isachardoublequoteopen}{\isasymexists}as{\isachardot}{\kern0pt}\ {\isacharparenleft}{\kern0pt}t\isactrlsub {\isasympi}\ {\isacharplus}{\kern0pt}\ {\isacharparenleft}{\kern0pt}duration\ {\isasympi}{\isacharparenright}{\kern0pt}{\isacharcomma}{\kern0pt}as{\isacharparenright}{\kern0pt}\ {\isasymin}\ set\ hs\isactrlsub {\isadigit{2}}\ {\isasymand}\ {\isacharquery}{\kern0pt}a\isactrlsub e\isactrlsub n\isactrlsub d\ {\isasymin}\ set\ as{\isachardoublequoteclose}\isanewline
\ \ \ \ \ \ \ \ \ \ \isakeywordONE{using}\isamarkupfalse%
\ {\isacartoucheopen}ind{\isacharunderscore}{\kern0pt}happ{\isacharunderscore}{\kern0pt}seq\ {\isasympi}s\ hs\isactrlsub {\isadigit{2}}{\isacartoucheclose}\ \isakeywordONE{unfolding}\isamarkupfalse%
\ ind{\isacharunderscore}{\kern0pt}happ{\isacharunderscore}{\kern0pt}seq{\isacharunderscore}{\kern0pt}def\ \isakeywordONE{by}\isamarkupfalse%
\ auto\isanewline
\ \ \ \ \ \ \ \ \isakeywordONE{then}\isamarkupfalse%
\ \isakeywordTHREE{obtain}\isamarkupfalse%
\ A{\isacharprime}{\kern0pt}\ \isakeywordTWO{where}\ {\isachardoublequoteopen}{\isacharparenleft}{\kern0pt}t\isactrlsub {\isasympi}\ {\isacharplus}{\kern0pt}\ {\isacharparenleft}{\kern0pt}duration\ {\isasympi}{\isacharparenright}{\kern0pt}{\isacharcomma}{\kern0pt}A{\isacharprime}{\kern0pt}{\isacharparenright}{\kern0pt}\ {\isasymin}\ set\ hs\isactrlsub {\isadigit{2}}{\isachardoublequoteclose}\ \isakeywordTWO{and}\ {\isachardoublequoteopen}{\isacharquery}{\kern0pt}a\isactrlsub e\isactrlsub n\isactrlsub d\ {\isasymin}\ set\ A{\isacharprime}{\kern0pt}{\isachardoublequoteclose}\ \isanewline
\ \ \ \ \ \ \ \ \ \ \isakeywordONE{by}\isamarkupfalse%
\ auto\isanewline
\ \ \ \ \ \ \ \ \isakeywordONE{then}\isamarkupfalse%
\ \isakeywordONE{have}\isamarkupfalse%
\ {\isachardoublequoteopen}{\isacharparenleft}{\kern0pt}t\isactrlsub a{\isacharcomma}{\kern0pt}A{\isacharprime}{\kern0pt}{\isacharparenright}{\kern0pt}\ {\isasymin}\ set\ hs\isactrlsub {\isadigit{2}}{\isachardoublequoteclose}\isanewline
\ \ \ \ \ \ \ \ \ \ \isakeywordONE{using}\isamarkupfalse%
\ {\isacartoucheopen}{\isacharquery}{\kern0pt}case{\isadigit{3}}{\isacartoucheclose}\ \isakeywordONE{by}\isamarkupfalse%
\ auto\isanewline
\ \ \ \ \ \ \ \ \isakeywordTHREE{show}\isamarkupfalse%
\ {\isacharquery}{\kern0pt}thesis\isanewline
\ \ \ \ \ \ \ \ \ \ \isakeywordONE{using}\isamarkupfalse%
\ {\isacartoucheopen}t\isactrlsub i\ {\isacharequal}{\kern0pt}\ t\isactrlsub a{\isacartoucheclose}\ {\isacartoucheopen}{\isacharquery}{\kern0pt}case{\isadigit{3}}{\isacartoucheclose}\ {\isacartoucheopen}{\isacharquery}{\kern0pt}a\isactrlsub e\isactrlsub n\isactrlsub d\ {\isasymin}\ set\ A{\isacharprime}{\kern0pt}{\isacartoucheclose}\isanewline
\ \ \ \ \ \ \ \ \ \ \ \ \ \ strict{\isacharunderscore}{\kern0pt}sorted{\isacharunderscore}{\kern0pt}drop{\isacharunderscore}{\kern0pt}le{\isacharunderscore}{\kern0pt}take{\isacharunderscore}{\kern0pt}leq{\isacharbrackleft}{\kern0pt}OF\ {\isacartoucheopen}strict{\isacharunderscore}{\kern0pt}sorted\ {\isacharparenleft}{\kern0pt}map\ fst\ hs\isactrlsub {\isadigit{2}}{\isacharparenright}{\kern0pt}{\isacartoucheclose}\ {\isacartoucheopen}{\isacharparenleft}{\kern0pt}t\isactrlsub a{\isacharcomma}{\kern0pt}A{\isacharprime}{\kern0pt}{\isacharparenright}{\kern0pt}\ {\isasymin}\ set\ hs\isactrlsub {\isadigit{2}}{\isacartoucheclose}{\isacharbrackright}{\kern0pt}\isanewline
\ \ \ \ \ \ \ \ \ \ \isakeywordONE{by}\isamarkupfalse%
\ auto\isanewline
\ \ \ \ \ \ \isakeywordONE{next}\isamarkupfalse%
\isanewline
\ \ \ \ \ \ \ \ \isakeywordTHREE{assume}\isamarkupfalse%
\ {\isacharquery}{\kern0pt}case{\isadigit{4}}\isanewline
\ \ \ \ \ \ \ \ \isakeywordTHREE{obtain}\isamarkupfalse%
\ t\isactrlsub i{\isacharprime}{\kern0pt}\ t\isactrlsub j{\isacharprime}{\kern0pt}\ \isakeywordTWO{where}\ {\isachardoublequoteopen}consec{\isacharunderscore}{\kern0pt}htps\ {\isasympi}s\ t\isactrlsub i{\isacharprime}{\kern0pt}\ t\isactrlsub j{\isacharprime}{\kern0pt}{\isachardoublequoteclose}\ \isanewline
\ \ \ \ \ \ \ \ \ \ \ \ \ \ \ \ \ \ \ \ \ \ \isakeywordTWO{and}\ {\isachardoublequoteopen}t\isactrlsub {\isasympi}\ {\isasymle}\ t\isactrlsub i{\isacharprime}{\kern0pt}\ {\isasymand}\ t\isactrlsub i{\isacharprime}{\kern0pt}\ {\isacharless}{\kern0pt}\ t\isactrlsub {\isasympi}\ {\isacharplus}{\kern0pt}\ duration\ {\isasympi}{\isachardoublequoteclose}\ \isanewline
\ \ \ \ \ \ \ \ \ \ \ \ \ \ \ \ \ \ \ \ \ \ \isakeywordTWO{and}\ {\isachardoublequoteopen}t\isactrlsub i{\isacharprime}{\kern0pt}\ {\isacharless}{\kern0pt}\ t\isactrlsub a\ {\isasymand}\ t\isactrlsub a\ {\isacharless}{\kern0pt}\ t\isactrlsub j{\isacharprime}{\kern0pt}{\isachardoublequoteclose}\isanewline
\ \ \ \ \ \ \ \ \ \ \isakeywordONE{using}\isamarkupfalse%
\ {\isacartoucheopen}{\isacharquery}{\kern0pt}case{\isadigit{4}}{\isacartoucheclose}\ \isakeywordONE{by}\isamarkupfalse%
\ auto\isanewline
\ \ \ \ \ \ \ \ \isakeywordONE{have}\isamarkupfalse%
\ {\isachardoublequoteopen}is{\isacharunderscore}{\kern0pt}htp\ {\isasympi}s\ t\isactrlsub a{\isachardoublequoteclose}\isanewline
\ \ \ \ \ \ \ \ \ \ \isakeywordONE{using}\isamarkupfalse%
\ {\isacartoucheopen}is{\isacharunderscore}{\kern0pt}htp\ {\isasympi}s\ t\isactrlsub i{\isacartoucheclose}\ {\isacartoucheopen}t\isactrlsub i\ {\isacharequal}{\kern0pt}\ t\isactrlsub a{\isacartoucheclose}\ \isakeywordONE{by}\isamarkupfalse%
\ simp\isanewline
\ \ \ \ \ \ \ \ \isakeywordTHREE{show}\isamarkupfalse%
\ {\isacharquery}{\kern0pt}thesis\isanewline
\ \ \ \ \ \ \ \ \ \ \isakeywordONE{using}\isamarkupfalse%
\ {\isacartoucheopen}is{\isacharunderscore}{\kern0pt}htp\ {\isasympi}s\ t\isactrlsub a{\isacartoucheclose}\ {\isacartoucheopen}t\isactrlsub i{\isacharprime}{\kern0pt}\ {\isacharless}{\kern0pt}\ t\isactrlsub a\ {\isasymand}\ t\isactrlsub a\ {\isacharless}{\kern0pt}\ t\isactrlsub j{\isacharprime}{\kern0pt}{\isacartoucheclose}\ {\isacartoucheopen}consec{\isacharunderscore}{\kern0pt}htps\ {\isasympi}s\ t\isactrlsub i{\isacharprime}{\kern0pt}\ t\isactrlsub j{\isacharprime}{\kern0pt}{\isacartoucheclose}\isanewline
\ \ \ \ \ \ \ \ \ \ \isakeywordONE{unfolding}\isamarkupfalse%
\ consec{\isacharunderscore}{\kern0pt}htps{\isacharunderscore}{\kern0pt}def\ \isakeywordONE{by}\isamarkupfalse%
\ auto\isanewline
\ \ \ \ \ \ \isakeywordONE{qed}\isamarkupfalse%
\isanewline
\ \ \ \ \isakeywordONE{qed}\isamarkupfalse%
%
\endisatagproof
{\isafoldproof}%
%
\isadelimproof
%
\endisadelimproof
%
\begin{isamarkuptext}%
With \isa{{\isasymlbrakk}\isaconst{is{\isacharunderscore}{\kern0pt}htp}\ \isavar{{\isacharquery}{\kern0pt}{\isasympi}s}\ \isavar{{\isacharquery}{\kern0pt}t\isactrlsub i}{\isacharsemicolon}{\kern0pt}\ \isaconst{ind{\isacharunderscore}{\kern0pt}happ{\isacharunderscore}{\kern0pt}seq}\ \isavar{{\isacharquery}{\kern0pt}{\isasympi}s}\ \isavar{{\isacharquery}{\kern0pt}hs\isactrlsub {\isadigit{1}}}{\isacharsemicolon}{\kern0pt}\ \isaconst{ind{\isacharunderscore}{\kern0pt}happ{\isacharunderscore}{\kern0pt}seq}\ \isavar{{\isacharquery}{\kern0pt}{\isasympi}s}\ \isavar{{\isacharquery}{\kern0pt}hs\isactrlsub {\isadigit{2}}}{\isacharsemicolon}{\kern0pt}\ {\isacharparenleft}{\kern0pt}\isavar{{\isacharquery}{\kern0pt}t\isactrlsub a}{\isacharcomma}{\kern0pt}\ \isavar{{\isacharquery}{\kern0pt}A}{\isacharparenright}{\kern0pt}\ {\isasymin}\ \isaconst{set}\ {\isacharparenleft}{\kern0pt}\isaconst{dropWhile}\ {\isacharparenleft}{\kern0pt}\isaconst{le}\ \isavar{{\isacharquery}{\kern0pt}t\isactrlsub i}{\isacharparenright}{\kern0pt}\ {\isacharparenleft}{\kern0pt}\isaconst{takeWhile}\ {\isacharparenleft}{\kern0pt}\isaconst{leq}\ \isavar{{\isacharquery}{\kern0pt}t\isactrlsub i}{\isacharparenright}{\kern0pt}\ \isavar{{\isacharquery}{\kern0pt}hs\isactrlsub {\isadigit{1}}}{\isacharparenright}{\kern0pt}{\isacharparenright}{\kern0pt}{\isacharsemicolon}{\kern0pt}\ \isavar{{\isacharquery}{\kern0pt}a}\ {\isasymin}\ \isaconst{set}\ \isavar{{\isacharquery}{\kern0pt}A}{\isasymrbrakk}\ {\isasymLongrightarrow}\ {\isasymexists}\isabound{t\isactrlsub a{\isacharprime}{\kern0pt}}\ \isabound{A{\isacharprime}{\kern0pt}}{\isachardot}{\kern0pt}\ {\isacharparenleft}{\kern0pt}\isabound{t\isactrlsub a{\isacharprime}{\kern0pt}}{\isacharcomma}{\kern0pt}\ \isabound{A{\isacharprime}{\kern0pt}}{\isacharparenright}{\kern0pt}\ {\isasymin}\ \isaconst{set}\ {\isacharparenleft}{\kern0pt}\isaconst{dropWhile}\ {\isacharparenleft}{\kern0pt}\isaconst{le}\ \isavar{{\isacharquery}{\kern0pt}t\isactrlsub i}{\isacharparenright}{\kern0pt}\ {\isacharparenleft}{\kern0pt}\isaconst{takeWhile}\ {\isacharparenleft}{\kern0pt}\isaconst{leq}\ \isavar{{\isacharquery}{\kern0pt}t\isactrlsub i}{\isacharparenright}{\kern0pt}\ \isavar{{\isacharquery}{\kern0pt}hs\isactrlsub {\isadigit{2}}}{\isacharparenright}{\kern0pt}{\isacharparenright}{\kern0pt}\ {\isasymand}\ \isavar{{\isacharquery}{\kern0pt}a}\ {\isasymin}\ \isaconst{set}\ \isabound{A{\isacharprime}{\kern0pt}}} we can prove: 
  If we have one induced happening sequence that is a valid happening sequence for a 
  single happening time point, then every other induced happening sequence is also valid 
  for that happening time point.%
\end{isamarkuptext}\isamarkuptrue%
\ \ \isakeywordONE{lemma}\isamarkupfalse%
\ valid{\isacharunderscore}{\kern0pt}happ{\isacharunderscore}{\kern0pt}seq{\isacharunderscore}{\kern0pt}le\isactrlsub i{\isacharunderscore}{\kern0pt}leq\isactrlsub i{\isacharcolon}{\kern0pt}\isanewline
\ \ \ \ \isakeywordTWO{assumes}\ {\isachardoublequoteopen}is{\isacharunderscore}{\kern0pt}htp\ {\isasympi}s\ t\isactrlsub i{\isachardoublequoteclose}\ \isanewline
\ \ \ \ \ \ \ \ \isakeywordTWO{and}\ {\isachardoublequoteopen}ind{\isacharunderscore}{\kern0pt}happ{\isacharunderscore}{\kern0pt}seq\ {\isasympi}s\ hs\isactrlsub {\isadigit{1}}{\isachardoublequoteclose}\ \isanewline
\ \ \ \ \ \ \ \ \isakeywordTWO{and}\ {\isachardoublequoteopen}ind{\isacharunderscore}{\kern0pt}happ{\isacharunderscore}{\kern0pt}seq\ {\isasympi}s\ hs\isactrlsub {\isadigit{2}}{\isachardoublequoteclose}\ \isanewline
\ \ \ \ \ \ \ \ \isakeywordTWO{and}\ {\isachardoublequoteopen}valid{\isacharunderscore}{\kern0pt}happ{\isacharunderscore}{\kern0pt}seq\ M\isactrlsub i\ {\isacharparenleft}{\kern0pt}dropWhile\ {\isacharparenleft}{\kern0pt}le\ t\isactrlsub i{\isacharparenright}{\kern0pt}\ {\isacharparenleft}{\kern0pt}takeWhile\ {\isacharparenleft}{\kern0pt}leq\ t\isactrlsub i{\isacharparenright}{\kern0pt}\ hs\isactrlsub {\isadigit{1}}{\isacharparenright}{\kern0pt}{\isacharparenright}{\kern0pt}\ M\isactrlsub j{\isachardoublequoteclose}\isanewline
\ \ \ \ \isakeywordTWO{shows}\ {\isachardoublequoteopen}valid{\isacharunderscore}{\kern0pt}happ{\isacharunderscore}{\kern0pt}seq\ M\isactrlsub i\ {\isacharparenleft}{\kern0pt}dropWhile\ {\isacharparenleft}{\kern0pt}le\ t\isactrlsub i{\isacharparenright}{\kern0pt}\ {\isacharparenleft}{\kern0pt}takeWhile\ {\isacharparenleft}{\kern0pt}leq\ t\isactrlsub i{\isacharparenright}{\kern0pt}\ hs\isactrlsub {\isadigit{2}}{\isacharparenright}{\kern0pt}{\isacharparenright}{\kern0pt}\ M\isactrlsub j{\isachardoublequoteclose}\isanewline
%
\isadelimproof
\ \ %
\endisadelimproof
%
\isatagproof
\isakeywordONE{proof}\isamarkupfalse%
\ {\isacharminus}{\kern0pt}\isanewline
\ \ \ \ \isakeywordTHREE{obtain}\isamarkupfalse%
\ A\isactrlsub {\isadigit{1}}\ \isakeywordTWO{where}\ {\isachardoublequoteopen}dropWhile\ {\isacharparenleft}{\kern0pt}le\ t\isactrlsub i{\isacharparenright}{\kern0pt}\ {\isacharparenleft}{\kern0pt}takeWhile\ {\isacharparenleft}{\kern0pt}leq\ t\isactrlsub i{\isacharparenright}{\kern0pt}\ hs\isactrlsub {\isadigit{1}}{\isacharparenright}{\kern0pt}\ {\isacharequal}{\kern0pt}\ {\isacharbrackleft}{\kern0pt}{\isacharparenleft}{\kern0pt}t\isactrlsub i{\isacharcomma}{\kern0pt}\ A\isactrlsub {\isadigit{1}}{\isacharparenright}{\kern0pt}{\isacharbrackright}{\kern0pt}{\isachardoublequoteclose}\isanewline
\ \ \ \ \ \ \isakeywordONE{using}\isamarkupfalse%
\ assms\ ind{\isacharunderscore}{\kern0pt}happ{\isacharunderscore}{\kern0pt}seq{\isacharunderscore}{\kern0pt}le\isactrlsub i{\isacharunderscore}{\kern0pt}leq\isactrlsub i{\isacharbrackleft}{\kern0pt}OF\ {\isacartoucheopen}is{\isacharunderscore}{\kern0pt}htp\ {\isasympi}s\ t\isactrlsub i{\isacartoucheclose}\ {\isacartoucheopen}ind{\isacharunderscore}{\kern0pt}happ{\isacharunderscore}{\kern0pt}seq\ {\isasympi}s\ hs\isactrlsub {\isadigit{1}}{\isacartoucheclose}{\isacharbrackright}{\kern0pt}\ \isakeywordONE{by}\isamarkupfalse%
\ auto\isanewline
\ \ \ \ \isakeywordTHREE{obtain}\isamarkupfalse%
\ A\isactrlsub {\isadigit{2}}\ \isakeywordTWO{where}\ {\isachardoublequoteopen}dropWhile\ {\isacharparenleft}{\kern0pt}le\ t\isactrlsub i{\isacharparenright}{\kern0pt}\ {\isacharparenleft}{\kern0pt}takeWhile\ {\isacharparenleft}{\kern0pt}leq\ t\isactrlsub i{\isacharparenright}{\kern0pt}\ hs\isactrlsub {\isadigit{2}}{\isacharparenright}{\kern0pt}\ {\isacharequal}{\kern0pt}\ {\isacharbrackleft}{\kern0pt}{\isacharparenleft}{\kern0pt}t\isactrlsub i{\isacharcomma}{\kern0pt}\ A\isactrlsub {\isadigit{2}}{\isacharparenright}{\kern0pt}{\isacharbrackright}{\kern0pt}{\isachardoublequoteclose}\isanewline
\ \ \ \ \ \ \isakeywordONE{using}\isamarkupfalse%
\ assms\ ind{\isacharunderscore}{\kern0pt}happ{\isacharunderscore}{\kern0pt}seq{\isacharunderscore}{\kern0pt}le\isactrlsub i{\isacharunderscore}{\kern0pt}leq\isactrlsub i{\isacharbrackleft}{\kern0pt}OF\ {\isacartoucheopen}is{\isacharunderscore}{\kern0pt}htp\ {\isasympi}s\ t\isactrlsub i{\isacartoucheclose}\ {\isacartoucheopen}ind{\isacharunderscore}{\kern0pt}happ{\isacharunderscore}{\kern0pt}seq\ {\isasympi}s\ hs\isactrlsub {\isadigit{2}}{\isacartoucheclose}{\isacharbrackright}{\kern0pt}\ \isakeywordONE{by}\isamarkupfalse%
\ auto\isanewline
\ \ \ \ \isakeywordONE{have}\isamarkupfalse%
\ {\isachardoublequoteopen}set\ A\isactrlsub {\isadigit{1}}\ {\isacharequal}{\kern0pt}\ set\ A\isactrlsub {\isadigit{2}}{\isachardoublequoteclose}\isanewline
\ \ \ \ \ \ \isakeywordONE{using}\isamarkupfalse%
\ ind{\isacharunderscore}{\kern0pt}happ{\isacharunderscore}{\kern0pt}seqs{\isacharunderscore}{\kern0pt}same{\isacharunderscore}{\kern0pt}acts{\isacharunderscore}{\kern0pt}le\isactrlsub i{\isacharunderscore}{\kern0pt}leq\isactrlsub i{\isacharbrackleft}{\kern0pt}OF\ {\isacartoucheopen}is{\isacharunderscore}{\kern0pt}htp\ {\isasympi}s\ t\isactrlsub i{\isacartoucheclose}\ {\isacartoucheopen}ind{\isacharunderscore}{\kern0pt}happ{\isacharunderscore}{\kern0pt}seq\ {\isasympi}s\ hs\isactrlsub {\isadigit{1}}{\isacartoucheclose}\isanewline
\ \ \ \ \ \ \ \ \ \ \ \ \ \ \ \ \ \ \ \ \ \ \ \ \ \ \ \ \ \ \ \ \ \ \ \ \ \ \ \ \ \ \ \ \ \ \ \ \ \ \ {\isacartoucheopen}ind{\isacharunderscore}{\kern0pt}happ{\isacharunderscore}{\kern0pt}seq\ {\isasympi}s\ hs\isactrlsub {\isadigit{2}}{\isacartoucheclose}{\isacharbrackright}{\kern0pt}\isanewline
\ \ \ \ \ \ \ \ \ \ \ \ \ {\isacartoucheopen}dropWhile\ {\isacharparenleft}{\kern0pt}le\ t\isactrlsub i{\isacharparenright}{\kern0pt}\ {\isacharparenleft}{\kern0pt}takeWhile\ {\isacharparenleft}{\kern0pt}leq\ t\isactrlsub i{\isacharparenright}{\kern0pt}\ hs\isactrlsub {\isadigit{1}}{\isacharparenright}{\kern0pt}\ {\isacharequal}{\kern0pt}\ {\isacharbrackleft}{\kern0pt}{\isacharparenleft}{\kern0pt}t\isactrlsub i{\isacharcomma}{\kern0pt}\ A\isactrlsub {\isadigit{1}}{\isacharparenright}{\kern0pt}{\isacharbrackright}{\kern0pt}{\isacartoucheclose}\isanewline
\ \ \ \ \ \ \ \ \ \ \ \ ind{\isacharunderscore}{\kern0pt}happ{\isacharunderscore}{\kern0pt}seqs{\isacharunderscore}{\kern0pt}same{\isacharunderscore}{\kern0pt}acts{\isacharunderscore}{\kern0pt}le\isactrlsub i{\isacharunderscore}{\kern0pt}leq\isactrlsub i{\isacharbrackleft}{\kern0pt}OF\ {\isacartoucheopen}is{\isacharunderscore}{\kern0pt}htp\ {\isasympi}s\ t\isactrlsub i{\isacartoucheclose}\ {\isacartoucheopen}ind{\isacharunderscore}{\kern0pt}happ{\isacharunderscore}{\kern0pt}seq\ {\isasympi}s\ hs\isactrlsub {\isadigit{2}}{\isacartoucheclose}\isanewline
\ \ \ \ \ \ \ \ \ \ \ \ \ \ \ \ \ \ \ \ \ \ \ \ \ \ \ \ \ \ \ \ \ \ \ \ \ \ \ \ \ \ \ \ \ \ \ \ \ \ \ {\isacartoucheopen}ind{\isacharunderscore}{\kern0pt}happ{\isacharunderscore}{\kern0pt}seq\ {\isasympi}s\ hs\isactrlsub {\isadigit{1}}{\isacartoucheclose}{\isacharbrackright}{\kern0pt}\isanewline
\ \ \ \ \ \ \ \ \ \ \ \ {\isacartoucheopen}dropWhile\ {\isacharparenleft}{\kern0pt}le\ t\isactrlsub i{\isacharparenright}{\kern0pt}\ {\isacharparenleft}{\kern0pt}takeWhile\ {\isacharparenleft}{\kern0pt}leq\ t\isactrlsub i{\isacharparenright}{\kern0pt}\ hs\isactrlsub {\isadigit{2}}{\isacharparenright}{\kern0pt}\ {\isacharequal}{\kern0pt}\ {\isacharbrackleft}{\kern0pt}{\isacharparenleft}{\kern0pt}t\isactrlsub i{\isacharcomma}{\kern0pt}\ A\isactrlsub {\isadigit{2}}{\isacharparenright}{\kern0pt}{\isacharbrackright}{\kern0pt}{\isacartoucheclose}\isanewline
\ \ \ \ \ \ \isakeywordONE{by}\isamarkupfalse%
\ auto\isanewline
\ \ \ \ \isakeywordONE{have}\isamarkupfalse%
\ {\isachardoublequoteopen}valid{\isacharunderscore}{\kern0pt}happ{\isacharunderscore}{\kern0pt}seq\ M\isactrlsub i\ {\isacharbrackleft}{\kern0pt}{\isacharparenleft}{\kern0pt}t\isactrlsub i{\isacharcomma}{\kern0pt}\ A\isactrlsub {\isadigit{1}}{\isacharparenright}{\kern0pt}{\isacharbrackright}{\kern0pt}\ M\isactrlsub j{\isachardoublequoteclose}\isanewline
\ \ \ \ \ \ \isakeywordONE{using}\isamarkupfalse%
\ {\isacartoucheopen}valid{\isacharunderscore}{\kern0pt}happ{\isacharunderscore}{\kern0pt}seq\ M\isactrlsub i\ {\isacharparenleft}{\kern0pt}dropWhile\ {\isacharparenleft}{\kern0pt}le\ t\isactrlsub i{\isacharparenright}{\kern0pt}\ {\isacharparenleft}{\kern0pt}takeWhile\ {\isacharparenleft}{\kern0pt}leq\ t\isactrlsub i{\isacharparenright}{\kern0pt}\ hs\isactrlsub {\isadigit{1}}{\isacharparenright}{\kern0pt}{\isacharparenright}{\kern0pt}\ M\isactrlsub j{\isacartoucheclose}\ \isanewline
\ \ \ \ \ \ \ \ \ \ \ \ {\isacartoucheopen}dropWhile\ {\isacharparenleft}{\kern0pt}le\ t\isactrlsub i{\isacharparenright}{\kern0pt}\ {\isacharparenleft}{\kern0pt}takeWhile\ {\isacharparenleft}{\kern0pt}leq\ t\isactrlsub i{\isacharparenright}{\kern0pt}\ hs\isactrlsub {\isadigit{1}}{\isacharparenright}{\kern0pt}\ {\isacharequal}{\kern0pt}\ {\isacharbrackleft}{\kern0pt}{\isacharparenleft}{\kern0pt}t\isactrlsub i{\isacharcomma}{\kern0pt}\ A\isactrlsub {\isadigit{1}}{\isacharparenright}{\kern0pt}{\isacharbrackright}{\kern0pt}{\isacartoucheclose}\ \isakeywordONE{by}\isamarkupfalse%
\ simp\isanewline
\ \ \ \ \isakeywordONE{then}\isamarkupfalse%
\ \isakeywordONE{have}\isamarkupfalse%
\ {\isachardoublequoteopen}valid{\isacharunderscore}{\kern0pt}happ{\isacharunderscore}{\kern0pt}seq\ M\isactrlsub i\ {\isacharbrackleft}{\kern0pt}{\isacharparenleft}{\kern0pt}t\isactrlsub i{\isacharcomma}{\kern0pt}A\isactrlsub {\isadigit{2}}{\isacharparenright}{\kern0pt}{\isacharbrackright}{\kern0pt}\ M\isactrlsub j{\isachardoublequoteclose}\isanewline
\ \ \ \ \ \ \isakeywordONE{using}\isamarkupfalse%
\ {\isacartoucheopen}set\ A\isactrlsub {\isadigit{1}}\ {\isacharequal}{\kern0pt}\ set\ A\isactrlsub {\isadigit{2}}{\isacartoucheclose}\ \isakeywordONE{by}\isamarkupfalse%
\ auto\isanewline
\ \ \ \ \isakeywordONE{then}\isamarkupfalse%
\ \isakeywordTHREE{show}\isamarkupfalse%
\ {\isacharquery}{\kern0pt}thesis\isanewline
\ \ \ \ \ \ \isakeywordONE{using}\isamarkupfalse%
\ {\isacartoucheopen}dropWhile\ {\isacharparenleft}{\kern0pt}le\ t\isactrlsub i{\isacharparenright}{\kern0pt}\ {\isacharparenleft}{\kern0pt}takeWhile\ {\isacharparenleft}{\kern0pt}leq\ t\isactrlsub i{\isacharparenright}{\kern0pt}\ hs\isactrlsub {\isadigit{2}}{\isacharparenright}{\kern0pt}\ {\isacharequal}{\kern0pt}\ {\isacharbrackleft}{\kern0pt}{\isacharparenleft}{\kern0pt}t\isactrlsub i{\isacharcomma}{\kern0pt}\ A\isactrlsub {\isadigit{2}}{\isacharparenright}{\kern0pt}{\isacharbrackright}{\kern0pt}{\isacartoucheclose}\ \isakeywordONE{by}\isamarkupfalse%
\ auto\isanewline
\ \ \isakeywordONE{qed}\isamarkupfalse%
%
\endisatagproof
{\isafoldproof}%
%
\isadelimproof
%
\endisadelimproof
%
\begin{isamarkuptext}%
Putting together \isa{{\isasymlbrakk}\isaconst{wf{\isacharunderscore}{\kern0pt}ast{\isacharunderscore}{\kern0pt}problem}\ \isavar{{\isacharquery}{\kern0pt}P}{\isacharsemicolon}{\kern0pt}\ \isaconst{consec{\isacharunderscore}{\kern0pt}htps}\ \isavar{{\isacharquery}{\kern0pt}{\isasympi}s}\ \isavar{{\isacharquery}{\kern0pt}t\isactrlsub i}\ \isavar{{\isacharquery}{\kern0pt}t\isactrlsub j}{\isacharsemicolon}{\kern0pt}\ \isaconst{ast{\isacharunderscore}{\kern0pt}problem{\isachardot}{\kern0pt}wf{\isacharunderscore}{\kern0pt}plan}\ \isavar{{\isacharquery}{\kern0pt}P}\ \isavar{{\isacharquery}{\kern0pt}{\isasympi}s}{\isacharsemicolon}{\kern0pt}\ \isaconst{ast{\isacharunderscore}{\kern0pt}problem{\isachardot}{\kern0pt}ind{\isacharunderscore}{\kern0pt}happ{\isacharunderscore}{\kern0pt}seq}\ \isavar{{\isacharquery}{\kern0pt}P}\ \isavar{{\isacharquery}{\kern0pt}{\isasympi}s}\ \isavar{{\isacharquery}{\kern0pt}hs\isactrlsub {\isadigit{1}}}{\isacharsemicolon}{\kern0pt}\ \isaconst{ast{\isacharunderscore}{\kern0pt}problem{\isachardot}{\kern0pt}ind{\isacharunderscore}{\kern0pt}happ{\isacharunderscore}{\kern0pt}seq}\ \isavar{{\isacharquery}{\kern0pt}P}\ \isavar{{\isacharquery}{\kern0pt}{\isasympi}s}\ \isavar{{\isacharquery}{\kern0pt}hs\isactrlsub {\isadigit{2}}}{\isacharsemicolon}{\kern0pt}\ \isaconst{valid{\isacharunderscore}{\kern0pt}happ{\isacharunderscore}{\kern0pt}seq}\ \isavar{{\isacharquery}{\kern0pt}M\isactrlsub i}\ {\isacharparenleft}{\kern0pt}\isaconst{dropWhile}\ {\isacharparenleft}{\kern0pt}\isaconst{leq}\ \isavar{{\isacharquery}{\kern0pt}t\isactrlsub i}{\isacharparenright}{\kern0pt}\ {\isacharparenleft}{\kern0pt}\isaconst{takeWhile}\ {\isacharparenleft}{\kern0pt}\isaconst{le}\ \isavar{{\isacharquery}{\kern0pt}t\isactrlsub j}{\isacharparenright}{\kern0pt}\ \isavar{{\isacharquery}{\kern0pt}hs\isactrlsub {\isadigit{1}}}{\isacharparenright}{\kern0pt}{\isacharparenright}{\kern0pt}\ \isavar{{\isacharquery}{\kern0pt}M\isactrlsub j}{\isasymrbrakk}\ {\isasymLongrightarrow}\ \isaconst{valid{\isacharunderscore}{\kern0pt}happ{\isacharunderscore}{\kern0pt}seq}\ \isavar{{\isacharquery}{\kern0pt}M\isactrlsub i}\ {\isacharparenleft}{\kern0pt}\isaconst{dropWhile}\ {\isacharparenleft}{\kern0pt}\isaconst{leq}\ \isavar{{\isacharquery}{\kern0pt}t\isactrlsub i}{\isacharparenright}{\kern0pt}\ {\isacharparenleft}{\kern0pt}\isaconst{takeWhile}\ {\isacharparenleft}{\kern0pt}\isaconst{le}\ \isavar{{\isacharquery}{\kern0pt}t\isactrlsub j}{\isacharparenright}{\kern0pt}\ \isavar{{\isacharquery}{\kern0pt}hs\isactrlsub {\isadigit{2}}}{\isacharparenright}{\kern0pt}{\isacharparenright}{\kern0pt}\ \isavar{{\isacharquery}{\kern0pt}M\isactrlsub j}} and 
  \isa{{\isasymlbrakk}\isaconst{is{\isacharunderscore}{\kern0pt}htp}\ \isavar{{\isacharquery}{\kern0pt}{\isasympi}s}\ \isavar{{\isacharquery}{\kern0pt}t\isactrlsub i}{\isacharsemicolon}{\kern0pt}\ \isaconst{ind{\isacharunderscore}{\kern0pt}happ{\isacharunderscore}{\kern0pt}seq}\ \isavar{{\isacharquery}{\kern0pt}{\isasympi}s}\ \isavar{{\isacharquery}{\kern0pt}hs\isactrlsub {\isadigit{1}}}{\isacharsemicolon}{\kern0pt}\ \isaconst{ind{\isacharunderscore}{\kern0pt}happ{\isacharunderscore}{\kern0pt}seq}\ \isavar{{\isacharquery}{\kern0pt}{\isasympi}s}\ \isavar{{\isacharquery}{\kern0pt}hs\isactrlsub {\isadigit{2}}}{\isacharsemicolon}{\kern0pt}\ \isaconst{valid{\isacharunderscore}{\kern0pt}happ{\isacharunderscore}{\kern0pt}seq}\ \isavar{{\isacharquery}{\kern0pt}M\isactrlsub i}\ {\isacharparenleft}{\kern0pt}\isaconst{dropWhile}\ {\isacharparenleft}{\kern0pt}\isaconst{le}\ \isavar{{\isacharquery}{\kern0pt}t\isactrlsub i}{\isacharparenright}{\kern0pt}\ {\isacharparenleft}{\kern0pt}\isaconst{takeWhile}\ {\isacharparenleft}{\kern0pt}\isaconst{leq}\ \isavar{{\isacharquery}{\kern0pt}t\isactrlsub i}{\isacharparenright}{\kern0pt}\ \isavar{{\isacharquery}{\kern0pt}hs\isactrlsub {\isadigit{1}}}{\isacharparenright}{\kern0pt}{\isacharparenright}{\kern0pt}\ \isavar{{\isacharquery}{\kern0pt}M\isactrlsub j}{\isasymrbrakk}\ {\isasymLongrightarrow}\ \isaconst{valid{\isacharunderscore}{\kern0pt}happ{\isacharunderscore}{\kern0pt}seq}\ \isavar{{\isacharquery}{\kern0pt}M\isactrlsub i}\ {\isacharparenleft}{\kern0pt}\isaconst{dropWhile}\ {\isacharparenleft}{\kern0pt}\isaconst{le}\ \isavar{{\isacharquery}{\kern0pt}t\isactrlsub i}{\isacharparenright}{\kern0pt}\ {\isacharparenleft}{\kern0pt}\isaconst{takeWhile}\ {\isacharparenleft}{\kern0pt}\isaconst{leq}\ \isavar{{\isacharquery}{\kern0pt}t\isactrlsub i}{\isacharparenright}{\kern0pt}\ \isavar{{\isacharquery}{\kern0pt}hs\isactrlsub {\isadigit{2}}}{\isacharparenright}{\kern0pt}{\isacharparenright}{\kern0pt}\ \isavar{{\isacharquery}{\kern0pt}M\isactrlsub j}}. If we have one induced happening sequence that is 
  valid from a happening time point to the next consecutive happening time point, 
  that every other induced happening sequence is also valid for that interval.%
\end{isamarkuptext}\isamarkuptrue%
\ \ \isakeywordONE{lemma}\isamarkupfalse%
\ valid{\isacharunderscore}{\kern0pt}happ{\isacharunderscore}{\kern0pt}seq{\isacharunderscore}{\kern0pt}consec{\isacharunderscore}{\kern0pt}htps{\isacharcolon}{\kern0pt}\isanewline
\ \ \ \ \isakeywordTWO{assumes}\ {\isachardoublequoteopen}consec{\isacharunderscore}{\kern0pt}htps\ {\isasympi}s\ t\isactrlsub i\ t\isactrlsub j{\isachardoublequoteclose}\isanewline
\ \ \ \ \ \ \ \ \isakeywordTWO{and}\ {\isachardoublequoteopen}wf{\isacharunderscore}{\kern0pt}plan\ {\isasympi}s{\isachardoublequoteclose}\isanewline
\ \ \ \ \ \ \ \ \isakeywordTWO{and}\ {\isachardoublequoteopen}ind{\isacharunderscore}{\kern0pt}happ{\isacharunderscore}{\kern0pt}seq\ {\isasympi}s\ hs\isactrlsub {\isadigit{1}}{\isachardoublequoteclose}\ \isanewline
\ \ \ \ \ \ \ \ \isakeywordTWO{and}\ {\isachardoublequoteopen}ind{\isacharunderscore}{\kern0pt}happ{\isacharunderscore}{\kern0pt}seq\ {\isasympi}s\ hs\isactrlsub {\isadigit{2}}{\isachardoublequoteclose}\ \isanewline
\ \ \ \ \ \ \ \ \isakeywordTWO{and}\ {\isachardoublequoteopen}valid{\isacharunderscore}{\kern0pt}happ{\isacharunderscore}{\kern0pt}seq\ M\isactrlsub i\ {\isacharparenleft}{\kern0pt}dropWhile\ {\isacharparenleft}{\kern0pt}leq\ t\isactrlsub i{\isacharparenright}{\kern0pt}\ {\isacharparenleft}{\kern0pt}takeWhile\ {\isacharparenleft}{\kern0pt}leq\ t\isactrlsub j{\isacharparenright}{\kern0pt}\ hs\isactrlsub {\isadigit{1}}{\isacharparenright}{\kern0pt}{\isacharparenright}{\kern0pt}\ M\isactrlsub j{\isachardoublequoteclose}\isanewline
\ \ \ \ \ \ \isakeywordTWO{shows}\ {\isachardoublequoteopen}valid{\isacharunderscore}{\kern0pt}happ{\isacharunderscore}{\kern0pt}seq\ M\isactrlsub i\ {\isacharparenleft}{\kern0pt}dropWhile\ {\isacharparenleft}{\kern0pt}leq\ t\isactrlsub i{\isacharparenright}{\kern0pt}\ {\isacharparenleft}{\kern0pt}takeWhile\ {\isacharparenleft}{\kern0pt}leq\ t\isactrlsub j{\isacharparenright}{\kern0pt}\ hs\isactrlsub {\isadigit{2}}{\isacharparenright}{\kern0pt}{\isacharparenright}{\kern0pt}\ M\isactrlsub j{\isachardoublequoteclose}\isanewline
%
\isadelimproof
\ \ \ \ %
\endisadelimproof
%
\isatagproof
\isakeywordONE{using}\isamarkupfalse%
\ assms\isanewline
\ \ \isakeywordONE{proof}\isamarkupfalse%
\ {\isacharminus}{\kern0pt}\isanewline
\ \ \ \ \isakeywordONE{have}\isamarkupfalse%
\ {\isachardoublequoteopen}is{\isacharunderscore}{\kern0pt}htp\ {\isasympi}s\ t\isactrlsub j{\isachardoublequoteclose}\ \isakeywordTWO{and}\ {\isachardoublequoteopen}t\isactrlsub i\ {\isacharless}{\kern0pt}\ t\isactrlsub j{\isachardoublequoteclose}\ \isakeywordTWO{and}\ {\isachardoublequoteopen}t\isactrlsub j\ {\isasymle}\ t\isactrlsub j{\isachardoublequoteclose}\isanewline
\ \ \ \ \ \ \isakeywordONE{using}\isamarkupfalse%
\ {\isacartoucheopen}consec{\isacharunderscore}{\kern0pt}htps\ {\isasympi}s\ t\isactrlsub i\ t\isactrlsub j{\isacartoucheclose}\ \isakeywordONE{unfolding}\isamarkupfalse%
\ consec{\isacharunderscore}{\kern0pt}htps{\isacharunderscore}{\kern0pt}def\ \isakeywordONE{by}\isamarkupfalse%
\ auto\isanewline
\ \ \ \ \isakeywordONE{have}\isamarkupfalse%
\ {\isachardoublequoteopen}strict{\isacharunderscore}{\kern0pt}sorted\ {\isacharparenleft}{\kern0pt}map\ fst\ hs\isactrlsub {\isadigit{1}}{\isacharparenright}{\kern0pt}{\isachardoublequoteclose}\ \isakeywordTWO{and}\ {\isachardoublequoteopen}strict{\isacharunderscore}{\kern0pt}sorted\ {\isacharparenleft}{\kern0pt}map\ fst\ hs\isactrlsub {\isadigit{2}}{\isacharparenright}{\kern0pt}{\isachardoublequoteclose}\isanewline
\ \ \ \ \ \ \isakeywordONE{using}\isamarkupfalse%
\ {\isacartoucheopen}ind{\isacharunderscore}{\kern0pt}happ{\isacharunderscore}{\kern0pt}seq\ {\isasympi}s\ hs\isactrlsub {\isadigit{1}}{\isacartoucheclose}\ {\isacartoucheopen}ind{\isacharunderscore}{\kern0pt}happ{\isacharunderscore}{\kern0pt}seq\ {\isasympi}s\ hs\isactrlsub {\isadigit{2}}{\isacartoucheclose}\ \isakeywordONE{unfolding}\isamarkupfalse%
\ ind{\isacharunderscore}{\kern0pt}happ{\isacharunderscore}{\kern0pt}seq{\isacharunderscore}{\kern0pt}def\ \isakeywordONE{by}\isamarkupfalse%
\ auto\isanewline
\ \ \ \ \isakeywordONE{have}\isamarkupfalse%
\ {\isachardoublequoteopen}dropWhile\ {\isacharparenleft}{\kern0pt}leq\ t\isactrlsub i{\isacharparenright}{\kern0pt}\ {\isacharparenleft}{\kern0pt}takeWhile\ {\isacharparenleft}{\kern0pt}leq\ t\isactrlsub j{\isacharparenright}{\kern0pt}\ hs\isactrlsub {\isadigit{1}}{\isacharparenright}{\kern0pt}\ \isanewline
\ \ \ \ \ \ \ \ \ \ \ {\isacharequal}{\kern0pt}\ dropWhile\ {\isacharparenleft}{\kern0pt}leq\ t\isactrlsub i{\isacharparenright}{\kern0pt}\ {\isacharparenleft}{\kern0pt}takeWhile\ {\isacharparenleft}{\kern0pt}le\ t\isactrlsub j{\isacharparenright}{\kern0pt}\ hs\isactrlsub {\isadigit{1}}{\isacharparenright}{\kern0pt}\ {\isacharat}{\kern0pt}\ dropWhile\ {\isacharparenleft}{\kern0pt}le\ t\isactrlsub j{\isacharparenright}{\kern0pt}\ {\isacharparenleft}{\kern0pt}takeWhile\ {\isacharparenleft}{\kern0pt}leq\ t\isactrlsub j{\isacharparenright}{\kern0pt}\ hs\isactrlsub {\isadigit{1}}{\isacharparenright}{\kern0pt}{\isachardoublequoteclose}\isanewline
\ \ \ \ \ \ \isakeywordONE{using}\isamarkupfalse%
\ drop{\isacharunderscore}{\kern0pt}take{\isacharunderscore}{\kern0pt}leq\isactrlsub i{\isacharunderscore}{\kern0pt}le{\isacharunderscore}{\kern0pt}leq\isactrlsub j{\isacharunderscore}{\kern0pt}split{\isacharunderscore}{\kern0pt}app{\isacharbrackleft}{\kern0pt}OF\ {\isacartoucheopen}t\isactrlsub i\ {\isacharless}{\kern0pt}\ t\isactrlsub j{\isacartoucheclose}\ {\isacartoucheopen}t\isactrlsub j\ {\isasymle}\ t\isactrlsub j{\isacartoucheclose}\ {\isacartoucheopen}strict{\isacharunderscore}{\kern0pt}sorted\ {\isacharparenleft}{\kern0pt}map\ fst\ hs\isactrlsub {\isadigit{1}}{\isacharparenright}{\kern0pt}{\isacartoucheclose}{\isacharbrackright}{\kern0pt}\ \isanewline
\ \ \ \ \ \ \isakeywordONE{by}\isamarkupfalse%
\ simp\isanewline
\ \ \ \ \isakeywordONE{then}\isamarkupfalse%
\ \isakeywordTHREE{obtain}\isamarkupfalse%
\ M\isactrlsub i{\isacharprime}{\kern0pt}\ \isakeywordTWO{where}\ {\isachardoublequoteopen}valid{\isacharunderscore}{\kern0pt}happ{\isacharunderscore}{\kern0pt}seq\ M\isactrlsub i\ {\isacharparenleft}{\kern0pt}dropWhile\ {\isacharparenleft}{\kern0pt}leq\ t\isactrlsub i{\isacharparenright}{\kern0pt}\ {\isacharparenleft}{\kern0pt}takeWhile\ {\isacharparenleft}{\kern0pt}le\ t\isactrlsub j{\isacharparenright}{\kern0pt}\ hs\isactrlsub {\isadigit{1}}{\isacharparenright}{\kern0pt}{\isacharparenright}{\kern0pt}\ M\isactrlsub i{\isacharprime}{\kern0pt}{\isachardoublequoteclose}\ \isanewline
\ \ \ \ \ \ \ \ \ \ \ \ \ \ \ \ \ \isakeywordTWO{and}\ {\isachardoublequoteopen}valid{\isacharunderscore}{\kern0pt}happ{\isacharunderscore}{\kern0pt}seq\ M\isactrlsub i{\isacharprime}{\kern0pt}\ {\isacharparenleft}{\kern0pt}dropWhile\ {\isacharparenleft}{\kern0pt}le\ t\isactrlsub j{\isacharparenright}{\kern0pt}\ {\isacharparenleft}{\kern0pt}takeWhile\ {\isacharparenleft}{\kern0pt}leq\ t\isactrlsub j{\isacharparenright}{\kern0pt}\ hs\isactrlsub {\isadigit{1}}{\isacharparenright}{\kern0pt}{\isacharparenright}{\kern0pt}\ M\isactrlsub j{\isachardoublequoteclose}\isanewline
\ \ \ \ \ \ \isakeywordONE{using}\isamarkupfalse%
\ valid{\isacharunderscore}{\kern0pt}happ{\isacharunderscore}{\kern0pt}seq{\isacharunderscore}{\kern0pt}app{\isacharunderscore}{\kern0pt}iff\isanewline
\ \ \ \ \ \ \ \ \ \ \ \ {\isacartoucheopen}valid{\isacharunderscore}{\kern0pt}happ{\isacharunderscore}{\kern0pt}seq\ M\isactrlsub i\ {\isacharparenleft}{\kern0pt}dropWhile\ {\isacharparenleft}{\kern0pt}leq\ t\isactrlsub i{\isacharparenright}{\kern0pt}\ {\isacharparenleft}{\kern0pt}takeWhile\ {\isacharparenleft}{\kern0pt}leq\ t\isactrlsub j{\isacharparenright}{\kern0pt}\ hs\isactrlsub {\isadigit{1}}{\isacharparenright}{\kern0pt}{\isacharparenright}{\kern0pt}\ M\isactrlsub j{\isacartoucheclose}\ \isakeywordONE{by}\isamarkupfalse%
\ auto\isanewline
\ \ \ \ \isakeywordONE{have}\isamarkupfalse%
\ {\isadigit{1}}{\isacharcolon}{\kern0pt}\ {\isachardoublequoteopen}valid{\isacharunderscore}{\kern0pt}happ{\isacharunderscore}{\kern0pt}seq\ M\isactrlsub i\ {\isacharparenleft}{\kern0pt}dropWhile\ {\isacharparenleft}{\kern0pt}leq\ t\isactrlsub i{\isacharparenright}{\kern0pt}\ {\isacharparenleft}{\kern0pt}takeWhile\ {\isacharparenleft}{\kern0pt}le\ t\isactrlsub j{\isacharparenright}{\kern0pt}\ hs\isactrlsub {\isadigit{2}}{\isacharparenright}{\kern0pt}{\isacharparenright}{\kern0pt}\ M\isactrlsub i{\isacharprime}{\kern0pt}{\isachardoublequoteclose}\isanewline
\ \ \ \ \ \ \isakeywordONE{using}\isamarkupfalse%
\ valid{\isacharunderscore}{\kern0pt}happ{\isacharunderscore}{\kern0pt}seq{\isacharunderscore}{\kern0pt}leq\isactrlsub i{\isacharunderscore}{\kern0pt}le\isactrlsub j\isanewline
\ \ \ \ \ \ \ \ \ \ \ \ \ \ {\isacharbrackleft}{\kern0pt}OF\ {\isacartoucheopen}consec{\isacharunderscore}{\kern0pt}htps\ {\isasympi}s\ t\isactrlsub i\ t\isactrlsub j{\isacartoucheclose}\ {\isacartoucheopen}wf{\isacharunderscore}{\kern0pt}plan\ {\isasympi}s{\isacartoucheclose}\ {\isacartoucheopen}ind{\isacharunderscore}{\kern0pt}happ{\isacharunderscore}{\kern0pt}seq\ {\isasympi}s\ hs\isactrlsub {\isadigit{1}}{\isacartoucheclose}\ {\isacartoucheopen}ind{\isacharunderscore}{\kern0pt}happ{\isacharunderscore}{\kern0pt}seq\ {\isasympi}s\ hs\isactrlsub {\isadigit{2}}{\isacartoucheclose}{\isacharbrackright}{\kern0pt}\isanewline
\ \ \ \ \ \ \ \ \ \ \ \ {\isacartoucheopen}valid{\isacharunderscore}{\kern0pt}happ{\isacharunderscore}{\kern0pt}seq\ M\isactrlsub i\ {\isacharparenleft}{\kern0pt}dropWhile\ {\isacharparenleft}{\kern0pt}leq\ t\isactrlsub i{\isacharparenright}{\kern0pt}\ {\isacharparenleft}{\kern0pt}takeWhile\ {\isacharparenleft}{\kern0pt}le\ t\isactrlsub j{\isacharparenright}{\kern0pt}\ hs\isactrlsub {\isadigit{1}}{\isacharparenright}{\kern0pt}{\isacharparenright}{\kern0pt}\ M\isactrlsub i{\isacharprime}{\kern0pt}{\isacartoucheclose}\ \isakeywordONE{by}\isamarkupfalse%
\ auto\isanewline
\ \ \ \ \isakeywordONE{have}\isamarkupfalse%
\ {\isadigit{2}}{\isacharcolon}{\kern0pt}\ {\isachardoublequoteopen}valid{\isacharunderscore}{\kern0pt}happ{\isacharunderscore}{\kern0pt}seq\ M\isactrlsub i{\isacharprime}{\kern0pt}\ {\isacharparenleft}{\kern0pt}dropWhile\ {\isacharparenleft}{\kern0pt}le\ t\isactrlsub j{\isacharparenright}{\kern0pt}\ {\isacharparenleft}{\kern0pt}takeWhile\ {\isacharparenleft}{\kern0pt}leq\ t\isactrlsub j{\isacharparenright}{\kern0pt}\ hs\isactrlsub {\isadigit{2}}{\isacharparenright}{\kern0pt}{\isacharparenright}{\kern0pt}\ M\isactrlsub j{\isachardoublequoteclose}\isanewline
\ \ \ \ \ \ \isakeywordONE{using}\isamarkupfalse%
\ valid{\isacharunderscore}{\kern0pt}happ{\isacharunderscore}{\kern0pt}seq{\isacharunderscore}{\kern0pt}le\isactrlsub i{\isacharunderscore}{\kern0pt}leq\isactrlsub i{\isacharbrackleft}{\kern0pt}OF\ {\isacartoucheopen}is{\isacharunderscore}{\kern0pt}htp\ {\isasympi}s\ t\isactrlsub j{\isacartoucheclose}\ {\isacartoucheopen}ind{\isacharunderscore}{\kern0pt}happ{\isacharunderscore}{\kern0pt}seq\ {\isasympi}s\ hs\isactrlsub {\isadigit{1}}{\isacartoucheclose}\ {\isacartoucheopen}ind{\isacharunderscore}{\kern0pt}happ{\isacharunderscore}{\kern0pt}seq\ {\isasympi}s\ hs\isactrlsub {\isadigit{2}}{\isacartoucheclose}{\isacharbrackright}{\kern0pt}\isanewline
\ \ \ \ \ \ \ \ \ \ \ \ {\isacartoucheopen}valid{\isacharunderscore}{\kern0pt}happ{\isacharunderscore}{\kern0pt}seq\ M\isactrlsub i{\isacharprime}{\kern0pt}\ {\isacharparenleft}{\kern0pt}dropWhile\ {\isacharparenleft}{\kern0pt}le\ t\isactrlsub j{\isacharparenright}{\kern0pt}\ {\isacharparenleft}{\kern0pt}takeWhile\ {\isacharparenleft}{\kern0pt}leq\ t\isactrlsub j{\isacharparenright}{\kern0pt}\ hs\isactrlsub {\isadigit{1}}{\isacharparenright}{\kern0pt}{\isacharparenright}{\kern0pt}\ M\isactrlsub j{\isacartoucheclose}\ \isakeywordONE{by}\isamarkupfalse%
\ auto\isanewline
\ \ \ \ \isakeywordONE{have}\isamarkupfalse%
\ {\isachardoublequoteopen}dropWhile\ {\isacharparenleft}{\kern0pt}leq\ t\isactrlsub i{\isacharparenright}{\kern0pt}\ {\isacharparenleft}{\kern0pt}takeWhile\ {\isacharparenleft}{\kern0pt}leq\ t\isactrlsub j{\isacharparenright}{\kern0pt}\ hs\isactrlsub {\isadigit{2}}{\isacharparenright}{\kern0pt}\ \isanewline
\ \ \ \ \ \ \ \ \ \ \ {\isacharequal}{\kern0pt}\ dropWhile\ {\isacharparenleft}{\kern0pt}leq\ t\isactrlsub i{\isacharparenright}{\kern0pt}\ {\isacharparenleft}{\kern0pt}takeWhile\ {\isacharparenleft}{\kern0pt}le\ t\isactrlsub j{\isacharparenright}{\kern0pt}\ hs\isactrlsub {\isadigit{2}}{\isacharparenright}{\kern0pt}\ {\isacharat}{\kern0pt}\ dropWhile\ {\isacharparenleft}{\kern0pt}le\ t\isactrlsub j{\isacharparenright}{\kern0pt}\ {\isacharparenleft}{\kern0pt}takeWhile\ {\isacharparenleft}{\kern0pt}leq\ t\isactrlsub j{\isacharparenright}{\kern0pt}\ hs\isactrlsub {\isadigit{2}}{\isacharparenright}{\kern0pt}{\isachardoublequoteclose}\isanewline
\ \ \ \ \ \ \isakeywordONE{using}\isamarkupfalse%
\ drop{\isacharunderscore}{\kern0pt}take{\isacharunderscore}{\kern0pt}leq\isactrlsub i{\isacharunderscore}{\kern0pt}le{\isacharunderscore}{\kern0pt}leq\isactrlsub j{\isacharunderscore}{\kern0pt}split{\isacharunderscore}{\kern0pt}app{\isacharbrackleft}{\kern0pt}OF\ {\isacartoucheopen}t\isactrlsub i\ {\isacharless}{\kern0pt}\ t\isactrlsub j{\isacartoucheclose}\ {\isacartoucheopen}t\isactrlsub j\ {\isasymle}\ t\isactrlsub j{\isacartoucheclose}\ {\isacartoucheopen}strict{\isacharunderscore}{\kern0pt}sorted\ {\isacharparenleft}{\kern0pt}map\ fst\ hs\isactrlsub {\isadigit{2}}{\isacharparenright}{\kern0pt}{\isacartoucheclose}{\isacharbrackright}{\kern0pt}\ \isanewline
\ \ \ \ \ \ \isakeywordONE{by}\isamarkupfalse%
\ simp\isanewline
\ \ \ \ \isakeywordONE{then}\isamarkupfalse%
\ \isakeywordTHREE{show}\isamarkupfalse%
\ {\isacharquery}{\kern0pt}thesis\isanewline
\ \ \ \ \ \ \isakeywordONE{using}\isamarkupfalse%
\ valid{\isacharunderscore}{\kern0pt}happ{\isacharunderscore}{\kern0pt}seq{\isacharunderscore}{\kern0pt}app{\isacharunderscore}{\kern0pt}iff\ {\isadigit{1}}\ {\isadigit{2}}\ \isakeywordONE{by}\isamarkupfalse%
\ auto\isanewline
\ \ \isakeywordONE{qed}\isamarkupfalse%
%
\endisatagproof
{\isafoldproof}%
%
\isadelimproof
\isanewline
%
\endisadelimproof
\isanewline
\isanewline
\ \ \isakeywordONE{lemma}\isamarkupfalse%
\ htps{\isacharunderscore}{\kern0pt}state{\isacharunderscore}{\kern0pt}trace{\isacharunderscore}{\kern0pt}unique{\isacharunderscore}{\kern0pt}i{\isacharunderscore}{\kern0pt}to{\isacharunderscore}{\kern0pt}j{\isacharunderscore}{\kern0pt}aux{\isacharcolon}{\kern0pt}\isanewline
\ \ \ \ \isakeywordTWO{fixes}\ {\isasympi}s\isanewline
\ \ \ \ \isakeywordTWO{assumes}\ {\isachardoublequoteopen}htps{\isacharunderscore}{\kern0pt}seq\ {\isasympi}s\ htps{\isachardoublequoteclose}\ \isanewline
\ \ \ \ \ \ \ \ \isakeywordTWO{and}\ {\isachardoublequoteopen}htps\ {\isacharequal}{\kern0pt}\ ts\isactrlsub {\isadigit{1}}\ {\isacharat}{\kern0pt}\ t\isactrlsub i\ {\isacharhash}{\kern0pt}\ ts\isactrlsub {\isadigit{2}}\ {\isacharat}{\kern0pt}\ t\isactrlsub j\ {\isacharhash}{\kern0pt}\ ts\isactrlsub {\isadigit{3}}{\isachardoublequoteclose}\ \isanewline
\ \ \ \ \ \ \ \ \isakeywordTWO{and}\ {\isachardoublequoteopen}wf{\isacharunderscore}{\kern0pt}plan\ {\isasympi}s{\isachardoublequoteclose}\isanewline
\ \ \ \ \ \ \ \ \isakeywordTWO{and}\ {\isachardoublequoteopen}ind{\isacharunderscore}{\kern0pt}happ{\isacharunderscore}{\kern0pt}seq\ {\isasympi}s\ hs\isactrlsub {\isadigit{1}}{\isachardoublequoteclose}\ \isanewline
\ \ \ \ \ \ \ \ \isakeywordTWO{and}\ {\isachardoublequoteopen}ind{\isacharunderscore}{\kern0pt}happ{\isacharunderscore}{\kern0pt}seq\ {\isasympi}s\ hs\isactrlsub {\isadigit{2}}{\isachardoublequoteclose}\ \isanewline
\ \ \ \ \ \ \ \ \isakeywordTWO{and}\ {\isachardoublequoteopen}valid{\isacharunderscore}{\kern0pt}happ{\isacharunderscore}{\kern0pt}seq\ M\isactrlsub i\ {\isacharparenleft}{\kern0pt}dropWhile\ {\isacharparenleft}{\kern0pt}leq\ t\isactrlsub i{\isacharparenright}{\kern0pt}\ {\isacharparenleft}{\kern0pt}takeWhile\ {\isacharparenleft}{\kern0pt}leq\ t\isactrlsub j{\isacharparenright}{\kern0pt}\ hs\isactrlsub {\isadigit{1}}{\isacharparenright}{\kern0pt}{\isacharparenright}{\kern0pt}\ M\isactrlsub j{\isachardoublequoteclose}\isanewline
\ \ \ \ \isakeywordTWO{shows}\ {\isachardoublequoteopen}valid{\isacharunderscore}{\kern0pt}happ{\isacharunderscore}{\kern0pt}seq\ M\isactrlsub i\ {\isacharparenleft}{\kern0pt}dropWhile\ {\isacharparenleft}{\kern0pt}leq\ t\isactrlsub i{\isacharparenright}{\kern0pt}\ {\isacharparenleft}{\kern0pt}takeWhile\ {\isacharparenleft}{\kern0pt}leq\ t\isactrlsub j{\isacharparenright}{\kern0pt}\ hs\isactrlsub {\isadigit{2}}{\isacharparenright}{\kern0pt}{\isacharparenright}{\kern0pt}\ M\isactrlsub j{\isachardoublequoteclose}\isanewline
%
\isadelimproof
\ \ \ \ %
\endisadelimproof
%
\isatagproof
\isakeywordONE{using}\isamarkupfalse%
\ assms\isanewline
\ \ \isakeywordONE{proof}\isamarkupfalse%
\ {\isacharparenleft}{\kern0pt}induction\ ts\isactrlsub {\isadigit{2}}\ arbitrary{\isacharcolon}{\kern0pt}\ ts\isactrlsub {\isadigit{1}}\ t\isactrlsub i\ ts\isactrlsub {\isadigit{3}}\ M\isactrlsub i{\isacharparenright}{\kern0pt}\isanewline
\ \ \ \ \isakeywordTHREE{case}\isamarkupfalse%
\ Nil\isanewline
\ \ \ \ \isakeywordONE{then}\isamarkupfalse%
\ \isakeywordONE{have}\isamarkupfalse%
\ {\isachardoublequoteopen}consec{\isacharunderscore}{\kern0pt}htps\ {\isasympi}s\ t\isactrlsub i\ t\isactrlsub j{\isachardoublequoteclose}\isanewline
\ \ \ \ \ \ \ \isakeywordONE{using}\isamarkupfalse%
\ cons{\isacharunderscore}{\kern0pt}consec{\isacharunderscore}{\kern0pt}htps\ \isakeywordONE{by}\isamarkupfalse%
\ auto\isanewline
\ \ \ \ \isakeywordONE{then}\isamarkupfalse%
\ \isakeywordTHREE{show}\isamarkupfalse%
\ {\isacharquery}{\kern0pt}case\isanewline
\ \ \ \ \ \ \isakeywordONE{using}\isamarkupfalse%
\ valid{\isacharunderscore}{\kern0pt}happ{\isacharunderscore}{\kern0pt}seq{\isacharunderscore}{\kern0pt}consec{\isacharunderscore}{\kern0pt}htps\isanewline
\ \ \ \ \ \ \ \ \ \ \ \ \ \ {\isacharbrackleft}{\kern0pt}OF\ {\isacartoucheopen}consec{\isacharunderscore}{\kern0pt}htps\ {\isasympi}s\ t\isactrlsub i\ t\isactrlsub j{\isacartoucheclose}\ {\isacartoucheopen}wf{\isacharunderscore}{\kern0pt}plan\ {\isasympi}s{\isacartoucheclose}\ {\isacartoucheopen}ind{\isacharunderscore}{\kern0pt}happ{\isacharunderscore}{\kern0pt}seq\ {\isasympi}s\ hs\isactrlsub {\isadigit{1}}{\isacartoucheclose}\ {\isacartoucheopen}ind{\isacharunderscore}{\kern0pt}happ{\isacharunderscore}{\kern0pt}seq\ {\isasympi}s\ hs\isactrlsub {\isadigit{2}}{\isacartoucheclose}{\isacharbrackright}{\kern0pt}\isanewline
\ \ \ \ \ \ \ \ \ \ \ \ {\isacartoucheopen}valid{\isacharunderscore}{\kern0pt}happ{\isacharunderscore}{\kern0pt}seq\ M\isactrlsub i\ {\isacharparenleft}{\kern0pt}dropWhile\ {\isacharparenleft}{\kern0pt}leq\ t\isactrlsub i{\isacharparenright}{\kern0pt}\ {\isacharparenleft}{\kern0pt}takeWhile\ {\isacharparenleft}{\kern0pt}leq\ t\isactrlsub j{\isacharparenright}{\kern0pt}\ hs\isactrlsub {\isadigit{1}}{\isacharparenright}{\kern0pt}{\isacharparenright}{\kern0pt}\ M\isactrlsub j{\isacartoucheclose}\isanewline
\ \ \ \ \ \ \isakeywordONE{by}\isamarkupfalse%
\ {\isacharparenleft}{\kern0pt}simp\ add{\isacharcolon}{\kern0pt}\ leq{\isacharunderscore}{\kern0pt}def{\isacharparenright}{\kern0pt}\isanewline
\ \ \isakeywordONE{next}\isamarkupfalse%
\isanewline
\ \ \ \ \isakeywordTHREE{case}\isamarkupfalse%
\ {\isacharparenleft}{\kern0pt}Cons\ t\isactrlsub i{\isacharprime}{\kern0pt}\ ts\isactrlsub {\isadigit{2}}{\isacharparenright}{\kern0pt}\isanewline
\ \ \ \ \isakeywordONE{then}\isamarkupfalse%
\ \isakeywordONE{have}\isamarkupfalse%
\ {\isachardoublequoteopen}strict{\isacharunderscore}{\kern0pt}sorted\ htps{\isachardoublequoteclose}\ \isakeywordTWO{and}\ {\isachardoublequoteopen}distinct\ htps{\isachardoublequoteclose}\isanewline
\ \ \ \ \ \ \isakeywordONE{using}\isamarkupfalse%
\ assms\ \isakeywordONE{unfolding}\isamarkupfalse%
\ htps{\isacharunderscore}{\kern0pt}seq{\isacharunderscore}{\kern0pt}def\ \isakeywordONE{by}\isamarkupfalse%
\ {\isacharparenleft}{\kern0pt}auto\ simp{\isacharcolon}{\kern0pt}\ strict{\isacharunderscore}{\kern0pt}sorted{\isacharunderscore}{\kern0pt}iff{\isacharparenright}{\kern0pt}\isanewline
\ \ \ \ \isakeywordONE{have}\isamarkupfalse%
\ {\isachardoublequoteopen}consec{\isacharunderscore}{\kern0pt}htps\ {\isasympi}s\ t\isactrlsub i\ t\isactrlsub i{\isacharprime}{\kern0pt}{\isachardoublequoteclose}\isanewline
\ \ \ \ \ \ \isakeywordONE{using}\isamarkupfalse%
\ Cons{\isachardot}{\kern0pt}prems\ cons{\isacharunderscore}{\kern0pt}consec{\isacharunderscore}{\kern0pt}htps\ \isakeywordONE{by}\isamarkupfalse%
\ auto\isanewline
\ \ \ \ \isakeywordONE{have}\isamarkupfalse%
\ {\isachardoublequoteopen}t\isactrlsub i\ {\isasymin}\ set\ htps{\isachardoublequoteclose}\ \isakeywordTWO{and}\ {\isachardoublequoteopen}t\isactrlsub i{\isacharprime}{\kern0pt}\ {\isasymin}\ set\ htps{\isachardoublequoteclose}\ \isakeywordTWO{and}\ {\isachardoublequoteopen}t\isactrlsub j\ {\isasymin}\ set\ htps{\isachardoublequoteclose}\isanewline
\ \ \ \ \ \ \isakeywordONE{using}\isamarkupfalse%
\ {\isacartoucheopen}htps\ {\isacharequal}{\kern0pt}\ ts\isactrlsub {\isadigit{1}}\ {\isacharat}{\kern0pt}\ t\isactrlsub i\ {\isacharhash}{\kern0pt}\ {\isacharparenleft}{\kern0pt}t\isactrlsub i{\isacharprime}{\kern0pt}\ {\isacharhash}{\kern0pt}\ ts\isactrlsub {\isadigit{2}}{\isacharparenright}{\kern0pt}\ {\isacharat}{\kern0pt}\ t\isactrlsub j\ {\isacharhash}{\kern0pt}\ ts\isactrlsub {\isadigit{3}}{\isacartoucheclose}\ \isakeywordONE{by}\isamarkupfalse%
\ auto\isanewline
\ \ \ \ \isakeywordONE{have}\isamarkupfalse%
\ {\isachardoublequoteopen}t\isactrlsub i\ {\isacharless}{\kern0pt}\ t\isactrlsub j{\isachardoublequoteclose}\isanewline
\ \ \ \ \ \ \isakeywordONE{using}\isamarkupfalse%
\ split{\isacharunderscore}{\kern0pt}list{\isadigit{2}}{\isacharunderscore}{\kern0pt}le{\isacharbrackleft}{\kern0pt}OF\ {\isacartoucheopen}strict{\isacharunderscore}{\kern0pt}sorted\ htps{\isacartoucheclose}\ {\isacartoucheopen}t\isactrlsub i\ {\isasymin}\ set\ htps{\isacartoucheclose}\ {\isacartoucheopen}t\isactrlsub j\ {\isasymin}\ set\ htps{\isacartoucheclose}{\isacharbrackright}{\kern0pt}\isanewline
\ \ \ \ \ \ \ \ \ \ \ \ {\isacartoucheopen}htps\ {\isacharequal}{\kern0pt}\ ts\isactrlsub {\isadigit{1}}\ {\isacharat}{\kern0pt}\ t\isactrlsub i\ {\isacharhash}{\kern0pt}\ {\isacharparenleft}{\kern0pt}t\isactrlsub i{\isacharprime}{\kern0pt}\ {\isacharhash}{\kern0pt}\ ts\isactrlsub {\isadigit{2}}{\isacharparenright}{\kern0pt}\ {\isacharat}{\kern0pt}\ t\isactrlsub j\ {\isacharhash}{\kern0pt}\ ts\isactrlsub {\isadigit{3}}{\isacartoucheclose}\ \isakeywordONE{by}\isamarkupfalse%
\ blast\isanewline
\ \ \ \ \isakeywordONE{have}\isamarkupfalse%
\ {\isachardoublequoteopen}t\isactrlsub i\ {\isasymle}\ t\isactrlsub i{\isacharprime}{\kern0pt}{\isachardoublequoteclose}\isanewline
\ \ \ \ \ \ \isakeywordONE{using}\isamarkupfalse%
\ {\isacartoucheopen}consec{\isacharunderscore}{\kern0pt}htps\ {\isasympi}s\ t\isactrlsub i\ t\isactrlsub i{\isacharprime}{\kern0pt}{\isacartoucheclose}\ \isakeywordONE{unfolding}\isamarkupfalse%
\ consec{\isacharunderscore}{\kern0pt}htps{\isacharunderscore}{\kern0pt}def\ \isakeywordONE{by}\isamarkupfalse%
\ auto\isanewline
\ \ \ \ \isakeywordONE{have}\isamarkupfalse%
\ {\isachardoublequoteopen}htps\ {\isacharequal}{\kern0pt}\ {\isacharparenleft}{\kern0pt}ts\isactrlsub {\isadigit{1}}\ {\isacharat}{\kern0pt}\ {\isacharbrackleft}{\kern0pt}t\isactrlsub i{\isacharbrackright}{\kern0pt}{\isacharparenright}{\kern0pt}\ {\isacharat}{\kern0pt}\ t\isactrlsub i{\isacharprime}{\kern0pt}\ {\isacharhash}{\kern0pt}\ ts\isactrlsub {\isadigit{2}}\ {\isacharat}{\kern0pt}\ t\isactrlsub j\ {\isacharhash}{\kern0pt}\ ts\isactrlsub {\isadigit{3}}{\isachardoublequoteclose}\isanewline
\ \ \ \ \ \ \isakeywordONE{using}\isamarkupfalse%
\ Cons{\isachardot}{\kern0pt}prems\ \isakeywordONE{by}\isamarkupfalse%
\ auto\isanewline
\ \ \ \ \isakeywordONE{then}\isamarkupfalse%
\ \isakeywordONE{have}\isamarkupfalse%
\ {\isachardoublequoteopen}t\isactrlsub i{\isacharprime}{\kern0pt}\ {\isasymle}\ t\isactrlsub j{\isachardoublequoteclose}\isanewline
\ \ \ \ \ \ \isakeywordONE{using}\isamarkupfalse%
\ Cons{\isachardot}{\kern0pt}prems\ split{\isacharunderscore}{\kern0pt}list{\isadigit{2}}{\isacharunderscore}{\kern0pt}le{\isacharbrackleft}{\kern0pt}OF\ {\isacartoucheopen}strict{\isacharunderscore}{\kern0pt}sorted\ htps{\isacartoucheclose}\ {\isacartoucheopen}t\isactrlsub i{\isacharprime}{\kern0pt}\ {\isasymin}\ set\ htps{\isacartoucheclose}\ {\isacartoucheopen}t\isactrlsub j\ {\isasymin}\ set\ htps{\isacartoucheclose}{\isacharbrackright}{\kern0pt}\isanewline
\ \ \ \ \ \ \ \ \ \ \ \ less{\isacharunderscore}{\kern0pt}eq{\isacharunderscore}{\kern0pt}rat{\isacharunderscore}{\kern0pt}def\ \isakeywordONE{by}\isamarkupfalse%
\ blast\isanewline
\ \ \ \ \isakeywordONE{have}\isamarkupfalse%
\ {\isachardoublequoteopen}strict{\isacharunderscore}{\kern0pt}sorted\ {\isacharparenleft}{\kern0pt}map\ fst\ hs\isactrlsub {\isadigit{1}}{\isacharparenright}{\kern0pt}{\isachardoublequoteclose}\ \isakeywordTWO{and}\ {\isachardoublequoteopen}strict{\isacharunderscore}{\kern0pt}sorted\ {\isacharparenleft}{\kern0pt}map\ fst\ hs\isactrlsub {\isadigit{2}}{\isacharparenright}{\kern0pt}{\isachardoublequoteclose}\isanewline
\ \ \ \ \ \ \isakeywordONE{using}\isamarkupfalse%
\ {\isacartoucheopen}ind{\isacharunderscore}{\kern0pt}happ{\isacharunderscore}{\kern0pt}seq\ {\isasympi}s\ hs\isactrlsub {\isadigit{1}}{\isacartoucheclose}\ {\isacartoucheopen}ind{\isacharunderscore}{\kern0pt}happ{\isacharunderscore}{\kern0pt}seq\ {\isasympi}s\ hs\isactrlsub {\isadigit{2}}{\isacartoucheclose}\ \isakeywordONE{unfolding}\isamarkupfalse%
\ ind{\isacharunderscore}{\kern0pt}happ{\isacharunderscore}{\kern0pt}seq{\isacharunderscore}{\kern0pt}def\ \isakeywordONE{by}\isamarkupfalse%
\ auto\isanewline
\ \ \ \ \isakeywordONE{have}\isamarkupfalse%
\ {\isachardoublequoteopen}{\isasymexists}M\isactrlsub i{\isacharprime}{\kern0pt}{\isachardot}{\kern0pt}\ valid{\isacharunderscore}{\kern0pt}happ{\isacharunderscore}{\kern0pt}seq\ M\isactrlsub i\ {\isacharparenleft}{\kern0pt}dropWhile\ {\isacharparenleft}{\kern0pt}leq\ t\isactrlsub i{\isacharparenright}{\kern0pt}\ {\isacharparenleft}{\kern0pt}takeWhile\ {\isacharparenleft}{\kern0pt}leq\ t\isactrlsub i{\isacharprime}{\kern0pt}{\isacharparenright}{\kern0pt}\ hs\isactrlsub {\isadigit{1}}{\isacharparenright}{\kern0pt}{\isacharparenright}{\kern0pt}\ M\isactrlsub i{\isacharprime}{\kern0pt}\ \isanewline
\ \ \ \ \ \ \ \ \ \ \ \ {\isasymand}\ valid{\isacharunderscore}{\kern0pt}happ{\isacharunderscore}{\kern0pt}seq\ M\isactrlsub i{\isacharprime}{\kern0pt}\ {\isacharparenleft}{\kern0pt}dropWhile\ {\isacharparenleft}{\kern0pt}leq\ t\isactrlsub i{\isacharprime}{\kern0pt}{\isacharparenright}{\kern0pt}\ {\isacharparenleft}{\kern0pt}takeWhile\ {\isacharparenleft}{\kern0pt}leq\ t\isactrlsub j{\isacharparenright}{\kern0pt}\ hs\isactrlsub {\isadigit{1}}{\isacharparenright}{\kern0pt}{\isacharparenright}{\kern0pt}\ M\isactrlsub j{\isachardoublequoteclose}\isanewline
\ \ \ \ \ \ \isakeywordONE{using}\isamarkupfalse%
\ drop{\isacharunderscore}{\kern0pt}take{\isacharunderscore}{\kern0pt}leq\isactrlsub i{\isacharunderscore}{\kern0pt}leq{\isacharunderscore}{\kern0pt}leq\isactrlsub j{\isacharunderscore}{\kern0pt}simp{\isacharbrackleft}{\kern0pt}OF\ {\isacartoucheopen}t\isactrlsub i\ {\isasymle}\ t\isactrlsub i{\isacharprime}{\kern0pt}{\isacartoucheclose}\ {\isacartoucheopen}t\isactrlsub i{\isacharprime}{\kern0pt}\ {\isasymle}\ t\isactrlsub j{\isacartoucheclose}\ {\isacartoucheopen}strict{\isacharunderscore}{\kern0pt}sorted\ {\isacharparenleft}{\kern0pt}map\ fst\ hs\isactrlsub {\isadigit{1}}{\isacharparenright}{\kern0pt}{\isacartoucheclose}{\isacharbrackright}{\kern0pt}\isanewline
\ \ \ \ \ \ \ \ \ \ \ \ valid{\isacharunderscore}{\kern0pt}happ{\isacharunderscore}{\kern0pt}seq{\isacharunderscore}{\kern0pt}app{\isacharunderscore}{\kern0pt}iff\isanewline
\ \ \ \ \ \ \ \ \ \ \ \ {\isacartoucheopen}valid{\isacharunderscore}{\kern0pt}happ{\isacharunderscore}{\kern0pt}seq\ M\isactrlsub i\ {\isacharparenleft}{\kern0pt}dropWhile\ {\isacharparenleft}{\kern0pt}leq\ t\isactrlsub i{\isacharparenright}{\kern0pt}\ {\isacharparenleft}{\kern0pt}takeWhile\ {\isacharparenleft}{\kern0pt}leq\ t\isactrlsub j{\isacharparenright}{\kern0pt}\ hs\isactrlsub {\isadigit{1}}{\isacharparenright}{\kern0pt}{\isacharparenright}{\kern0pt}\ M\isactrlsub j{\isacartoucheclose}\ \isanewline
\ \ \ \ \ \ \isakeywordONE{by}\isamarkupfalse%
\ {\isacharparenleft}{\kern0pt}metis\ {\isacharparenleft}{\kern0pt}no{\isacharunderscore}{\kern0pt}types{\isacharcomma}{\kern0pt}\ lifting{\isacharparenright}{\kern0pt}{\isacharparenright}{\kern0pt}\isanewline
\ \ \ \ \isakeywordONE{then}\isamarkupfalse%
\ \isakeywordTHREE{obtain}\isamarkupfalse%
\ M\isactrlsub i{\isacharprime}{\kern0pt}\ \isakeywordTWO{where}\ {\isachardoublequoteopen}valid{\isacharunderscore}{\kern0pt}happ{\isacharunderscore}{\kern0pt}seq\ M\isactrlsub i\ {\isacharparenleft}{\kern0pt}dropWhile\ {\isacharparenleft}{\kern0pt}leq\ t\isactrlsub i{\isacharparenright}{\kern0pt}\ {\isacharparenleft}{\kern0pt}takeWhile\ {\isacharparenleft}{\kern0pt}leq\ t\isactrlsub i{\isacharprime}{\kern0pt}{\isacharparenright}{\kern0pt}\ hs\isactrlsub {\isadigit{1}}{\isacharparenright}{\kern0pt}{\isacharparenright}{\kern0pt}\ M\isactrlsub i{\isacharprime}{\kern0pt}{\isachardoublequoteclose}\isanewline
\ \ \ \ \ \ \ \ \ \ \ \ \ \ \ \ \isakeywordTWO{and}\ {\isachardoublequoteopen}valid{\isacharunderscore}{\kern0pt}happ{\isacharunderscore}{\kern0pt}seq\ M\isactrlsub i{\isacharprime}{\kern0pt}\ {\isacharparenleft}{\kern0pt}dropWhile\ {\isacharparenleft}{\kern0pt}leq\ t\isactrlsub i{\isacharprime}{\kern0pt}{\isacharparenright}{\kern0pt}\ {\isacharparenleft}{\kern0pt}takeWhile\ {\isacharparenleft}{\kern0pt}leq\ t\isactrlsub j{\isacharparenright}{\kern0pt}\ hs\isactrlsub {\isadigit{1}}{\isacharparenright}{\kern0pt}{\isacharparenright}{\kern0pt}\ M\isactrlsub j{\isachardoublequoteclose}\isanewline
\ \ \ \ \ \ \isakeywordONE{by}\isamarkupfalse%
\ auto\isanewline
\ \ \ \ \isakeywordONE{then}\isamarkupfalse%
\ \isakeywordONE{have}\isamarkupfalse%
\ {\isachardoublequoteopen}valid{\isacharunderscore}{\kern0pt}happ{\isacharunderscore}{\kern0pt}seq\ M\isactrlsub i\ {\isacharparenleft}{\kern0pt}dropWhile\ {\isacharparenleft}{\kern0pt}leq\ t\isactrlsub i{\isacharparenright}{\kern0pt}\ {\isacharparenleft}{\kern0pt}takeWhile\ {\isacharparenleft}{\kern0pt}leq\ t\isactrlsub i{\isacharprime}{\kern0pt}{\isacharparenright}{\kern0pt}\ hs\isactrlsub {\isadigit{2}}{\isacharparenright}{\kern0pt}{\isacharparenright}{\kern0pt}\ M\isactrlsub i{\isacharprime}{\kern0pt}{\isachardoublequoteclose}\isanewline
\ \ \ \ \ \ \isakeywordONE{using}\isamarkupfalse%
\ valid{\isacharunderscore}{\kern0pt}happ{\isacharunderscore}{\kern0pt}seq{\isacharunderscore}{\kern0pt}consec{\isacharunderscore}{\kern0pt}htps\isanewline
\ \ \ \ \ \ \ \ \ \ \ \ \ \ {\isacharbrackleft}{\kern0pt}OF\ {\isacartoucheopen}consec{\isacharunderscore}{\kern0pt}htps\ {\isasympi}s\ t\isactrlsub i\ t\isactrlsub i{\isacharprime}{\kern0pt}{\isacartoucheclose}\ {\isacartoucheopen}wf{\isacharunderscore}{\kern0pt}plan\ {\isasympi}s{\isacartoucheclose}\ {\isacartoucheopen}ind{\isacharunderscore}{\kern0pt}happ{\isacharunderscore}{\kern0pt}seq\ {\isasympi}s\ hs\isactrlsub {\isadigit{1}}{\isacartoucheclose}\ {\isacartoucheopen}ind{\isacharunderscore}{\kern0pt}happ{\isacharunderscore}{\kern0pt}seq\ {\isasympi}s\ hs\isactrlsub {\isadigit{2}}{\isacartoucheclose}{\isacharbrackright}{\kern0pt}\isanewline
\ \ \ \ \ \ \ \ \ \ \ \ {\isacartoucheopen}valid{\isacharunderscore}{\kern0pt}happ{\isacharunderscore}{\kern0pt}seq\ M\isactrlsub i\ {\isacharparenleft}{\kern0pt}dropWhile\ {\isacharparenleft}{\kern0pt}leq\ t\isactrlsub i{\isacharparenright}{\kern0pt}\ {\isacharparenleft}{\kern0pt}takeWhile\ {\isacharparenleft}{\kern0pt}leq\ t\isactrlsub i{\isacharprime}{\kern0pt}{\isacharparenright}{\kern0pt}\ hs\isactrlsub {\isadigit{1}}{\isacharparenright}{\kern0pt}{\isacharparenright}{\kern0pt}\ M\isactrlsub i{\isacharprime}{\kern0pt}{\isacartoucheclose}\ \isakeywordONE{by}\isamarkupfalse%
\ simp\isanewline
\ \ \ \ \isakeywordONE{have}\isamarkupfalse%
\ {\isachardoublequoteopen}htps\ {\isacharequal}{\kern0pt}\ {\isacharparenleft}{\kern0pt}ts\isactrlsub {\isadigit{1}}\ {\isacharat}{\kern0pt}\ {\isacharbrackleft}{\kern0pt}t\isactrlsub i{\isacharbrackright}{\kern0pt}{\isacharparenright}{\kern0pt}\ {\isacharat}{\kern0pt}\ t\isactrlsub i{\isacharprime}{\kern0pt}\ {\isacharhash}{\kern0pt}\ ts\isactrlsub {\isadigit{2}}\ {\isacharat}{\kern0pt}\ t\isactrlsub j\ {\isacharhash}{\kern0pt}\ ts\isactrlsub {\isadigit{3}}{\isachardoublequoteclose}\isanewline
\ \ \ \ \ \ \isakeywordONE{using}\isamarkupfalse%
\ Cons{\isachardot}{\kern0pt}prems\ \isakeywordONE{by}\isamarkupfalse%
\ simp\isanewline
\ \ \ \ \isakeywordONE{then}\isamarkupfalse%
\ \isakeywordONE{have}\isamarkupfalse%
\ {\isachardoublequoteopen}valid{\isacharunderscore}{\kern0pt}happ{\isacharunderscore}{\kern0pt}seq\ M\isactrlsub i{\isacharprime}{\kern0pt}\ {\isacharparenleft}{\kern0pt}dropWhile\ {\isacharparenleft}{\kern0pt}leq\ t\isactrlsub i{\isacharprime}{\kern0pt}{\isacharparenright}{\kern0pt}\ {\isacharparenleft}{\kern0pt}takeWhile\ {\isacharparenleft}{\kern0pt}leq\ t\isactrlsub j{\isacharparenright}{\kern0pt}\ hs\isactrlsub {\isadigit{2}}{\isacharparenright}{\kern0pt}{\isacharparenright}{\kern0pt}\ M\isactrlsub j{\isachardoublequoteclose}\isanewline
\ \ \ \ \ \ \isakeywordONE{using}\isamarkupfalse%
\ assms\ Cons{\isachardot}{\kern0pt}IH\isanewline
\ \ \ \ \ \ \ \ \ \ \ \ {\isacartoucheopen}valid{\isacharunderscore}{\kern0pt}happ{\isacharunderscore}{\kern0pt}seq\ M\isactrlsub i{\isacharprime}{\kern0pt}\ {\isacharparenleft}{\kern0pt}dropWhile\ {\isacharparenleft}{\kern0pt}leq\ t\isactrlsub i{\isacharprime}{\kern0pt}{\isacharparenright}{\kern0pt}\ {\isacharparenleft}{\kern0pt}takeWhile\ {\isacharparenleft}{\kern0pt}leq\ t\isactrlsub j{\isacharparenright}{\kern0pt}\ hs\isactrlsub {\isadigit{1}}{\isacharparenright}{\kern0pt}{\isacharparenright}{\kern0pt}\ M\isactrlsub j{\isacartoucheclose}\ \isakeywordONE{by}\isamarkupfalse%
\ fastforce\isanewline
\ \ \ \ \isakeywordONE{then}\isamarkupfalse%
\ \isakeywordTHREE{show}\isamarkupfalse%
\ {\isacharquery}{\kern0pt}case\ \isanewline
\ \ \ \ \ \ \isakeywordONE{using}\isamarkupfalse%
\ drop{\isacharunderscore}{\kern0pt}take{\isacharunderscore}{\kern0pt}leq\isactrlsub i{\isacharunderscore}{\kern0pt}leq{\isacharunderscore}{\kern0pt}leq\isactrlsub j{\isacharunderscore}{\kern0pt}simp{\isacharbrackleft}{\kern0pt}OF\ {\isacartoucheopen}t\isactrlsub i\ {\isasymle}\ t\isactrlsub i{\isacharprime}{\kern0pt}{\isacartoucheclose}\ {\isacartoucheopen}t\isactrlsub i{\isacharprime}{\kern0pt}\ {\isasymle}\ t\isactrlsub j{\isacartoucheclose}\ {\isacartoucheopen}strict{\isacharunderscore}{\kern0pt}sorted\ {\isacharparenleft}{\kern0pt}map\ fst\ hs\isactrlsub {\isadigit{2}}{\isacharparenright}{\kern0pt}{\isacartoucheclose}{\isacharbrackright}{\kern0pt}\isanewline
\ \ \ \ \ \ \ \ \ \ \ \ valid{\isacharunderscore}{\kern0pt}happ{\isacharunderscore}{\kern0pt}seq{\isacharunderscore}{\kern0pt}app{\isacharunderscore}{\kern0pt}iff\isanewline
\ \ \ \ \ \ \ \ \ \ \ \ {\isacartoucheopen}valid{\isacharunderscore}{\kern0pt}happ{\isacharunderscore}{\kern0pt}seq\ M\isactrlsub i\ {\isacharparenleft}{\kern0pt}dropWhile\ {\isacharparenleft}{\kern0pt}leq\ t\isactrlsub i{\isacharparenright}{\kern0pt}\ {\isacharparenleft}{\kern0pt}takeWhile\ {\isacharparenleft}{\kern0pt}leq\ t\isactrlsub i{\isacharprime}{\kern0pt}{\isacharparenright}{\kern0pt}\ hs\isactrlsub {\isadigit{2}}{\isacharparenright}{\kern0pt}{\isacharparenright}{\kern0pt}\ M\isactrlsub i{\isacharprime}{\kern0pt}{\isacartoucheclose}\isanewline
\ \ \ \ \ \ \ \ \ \ \ \ {\isacartoucheopen}valid{\isacharunderscore}{\kern0pt}happ{\isacharunderscore}{\kern0pt}seq\ M\isactrlsub i{\isacharprime}{\kern0pt}\ {\isacharparenleft}{\kern0pt}dropWhile\ {\isacharparenleft}{\kern0pt}leq\ t\isactrlsub i{\isacharprime}{\kern0pt}{\isacharparenright}{\kern0pt}\ {\isacharparenleft}{\kern0pt}takeWhile\ {\isacharparenleft}{\kern0pt}leq\ t\isactrlsub j{\isacharparenright}{\kern0pt}\ hs\isactrlsub {\isadigit{2}}{\isacharparenright}{\kern0pt}{\isacharparenright}{\kern0pt}\ M\isactrlsub j{\isacartoucheclose}\isanewline
\ \ \ \ \ \ \isakeywordONE{by}\isamarkupfalse%
\ {\isacharparenleft}{\kern0pt}metis\ {\isacharparenleft}{\kern0pt}no{\isacharunderscore}{\kern0pt}types{\isacharcomma}{\kern0pt}\ lifting{\isacharparenright}{\kern0pt}{\isacharparenright}{\kern0pt}\isanewline
\ \ \isakeywordONE{qed}\isamarkupfalse%
%
\endisatagproof
{\isafoldproof}%
%
\isadelimproof
%
\endisadelimproof
%
\begin{isamarkuptext}%
With \isa{{\isasymlbrakk}\isaconst{wf{\isacharunderscore}{\kern0pt}ast{\isacharunderscore}{\kern0pt}problem}\ \isavar{{\isacharquery}{\kern0pt}P}{\isacharsemicolon}{\kern0pt}\ \isaconst{consec{\isacharunderscore}{\kern0pt}htps}\ \isavar{{\isacharquery}{\kern0pt}{\isasympi}s}\ \isavar{{\isacharquery}{\kern0pt}t\isactrlsub i}\ \isavar{{\isacharquery}{\kern0pt}t\isactrlsub j}{\isacharsemicolon}{\kern0pt}\ \isaconst{ast{\isacharunderscore}{\kern0pt}problem{\isachardot}{\kern0pt}wf{\isacharunderscore}{\kern0pt}plan}\ \isavar{{\isacharquery}{\kern0pt}P}\ \isavar{{\isacharquery}{\kern0pt}{\isasympi}s}{\isacharsemicolon}{\kern0pt}\ \isaconst{ast{\isacharunderscore}{\kern0pt}problem{\isachardot}{\kern0pt}ind{\isacharunderscore}{\kern0pt}happ{\isacharunderscore}{\kern0pt}seq}\ \isavar{{\isacharquery}{\kern0pt}P}\ \isavar{{\isacharquery}{\kern0pt}{\isasympi}s}\ \isavar{{\isacharquery}{\kern0pt}hs\isactrlsub {\isadigit{1}}}{\isacharsemicolon}{\kern0pt}\ \isaconst{ast{\isacharunderscore}{\kern0pt}problem{\isachardot}{\kern0pt}ind{\isacharunderscore}{\kern0pt}happ{\isacharunderscore}{\kern0pt}seq}\ \isavar{{\isacharquery}{\kern0pt}P}\ \isavar{{\isacharquery}{\kern0pt}{\isasympi}s}\ \isavar{{\isacharquery}{\kern0pt}hs\isactrlsub {\isadigit{2}}}{\isacharsemicolon}{\kern0pt}\ \isaconst{valid{\isacharunderscore}{\kern0pt}happ{\isacharunderscore}{\kern0pt}seq}\ \isavar{{\isacharquery}{\kern0pt}M\isactrlsub i}\ {\isacharparenleft}{\kern0pt}\isaconst{dropWhile}\ {\isacharparenleft}{\kern0pt}\isaconst{leq}\ \isavar{{\isacharquery}{\kern0pt}t\isactrlsub i}{\isacharparenright}{\kern0pt}\ {\isacharparenleft}{\kern0pt}\isaconst{takeWhile}\ {\isacharparenleft}{\kern0pt}\isaconst{leq}\ \isavar{{\isacharquery}{\kern0pt}t\isactrlsub j}{\isacharparenright}{\kern0pt}\ \isavar{{\isacharquery}{\kern0pt}hs\isactrlsub {\isadigit{1}}}{\isacharparenright}{\kern0pt}{\isacharparenright}{\kern0pt}\ \isavar{{\isacharquery}{\kern0pt}M\isactrlsub j}{\isasymrbrakk}\ {\isasymLongrightarrow}\ \isaconst{valid{\isacharunderscore}{\kern0pt}happ{\isacharunderscore}{\kern0pt}seq}\ \isavar{{\isacharquery}{\kern0pt}M\isactrlsub i}\ {\isacharparenleft}{\kern0pt}\isaconst{dropWhile}\ {\isacharparenleft}{\kern0pt}\isaconst{leq}\ \isavar{{\isacharquery}{\kern0pt}t\isactrlsub i}{\isacharparenright}{\kern0pt}\ {\isacharparenleft}{\kern0pt}\isaconst{takeWhile}\ {\isacharparenleft}{\kern0pt}\isaconst{leq}\ \isavar{{\isacharquery}{\kern0pt}t\isactrlsub j}{\isacharparenright}{\kern0pt}\ \isavar{{\isacharquery}{\kern0pt}hs\isactrlsub {\isadigit{2}}}{\isacharparenright}{\kern0pt}{\isacharparenright}{\kern0pt}\ \isavar{{\isacharquery}{\kern0pt}M\isactrlsub j}} we can now proof that if one 
  induced happening sequence is valid form a happening time point to another happening time 
  point (not just consecutive), than every other induced happening sequence is valid for 
  between those two happening time points.%
\end{isamarkuptext}\isamarkuptrue%
\ \ \isakeywordONE{lemma}\isamarkupfalse%
\ valid{\isacharunderscore}{\kern0pt}happ{\isacharunderscore}{\kern0pt}seq{\isacharunderscore}{\kern0pt}htp{\isacharunderscore}{\kern0pt}i{\isacharunderscore}{\kern0pt}to{\isacharunderscore}{\kern0pt}htp{\isacharunderscore}{\kern0pt}j{\isacharcolon}{\kern0pt}\isanewline
\ \ \ \ \isakeywordTWO{assumes}\ {\isachardoublequoteopen}t\isactrlsub i\ {\isacharless}{\kern0pt}\ t\isactrlsub j{\isachardoublequoteclose}\ \isanewline
\ \ \ \ \ \ \ \ \isakeywordTWO{and}\ {\isachardoublequoteopen}is{\isacharunderscore}{\kern0pt}htp\ {\isasympi}s\ t\isactrlsub i{\isachardoublequoteclose}\ \isanewline
\ \ \ \ \ \ \ \ \isakeywordTWO{and}\ {\isachardoublequoteopen}is{\isacharunderscore}{\kern0pt}htp\ {\isasympi}s\ t\isactrlsub j{\isachardoublequoteclose}\isanewline
\ \ \ \ \ \ \ \ \isakeywordTWO{and}\ {\isachardoublequoteopen}wf{\isacharunderscore}{\kern0pt}plan\ {\isasympi}s{\isachardoublequoteclose}\isanewline
\ \ \ \ \ \ \ \ \isakeywordTWO{and}\ {\isachardoublequoteopen}ind{\isacharunderscore}{\kern0pt}happ{\isacharunderscore}{\kern0pt}seq\ {\isasympi}s\ hs\isactrlsub {\isadigit{1}}{\isachardoublequoteclose}\ \isanewline
\ \ \ \ \ \ \ \ \isakeywordTWO{and}\ {\isachardoublequoteopen}ind{\isacharunderscore}{\kern0pt}happ{\isacharunderscore}{\kern0pt}seq\ {\isasympi}s\ hs\isactrlsub {\isadigit{2}}{\isachardoublequoteclose}\ \isanewline
\ \ \ \ \ \ \ \ \isakeywordTWO{and}\ {\isachardoublequoteopen}valid{\isacharunderscore}{\kern0pt}happ{\isacharunderscore}{\kern0pt}seq\ M\isactrlsub i\ {\isacharparenleft}{\kern0pt}dropWhile\ {\isacharparenleft}{\kern0pt}leq\ t\isactrlsub i{\isacharparenright}{\kern0pt}\ {\isacharparenleft}{\kern0pt}takeWhile\ {\isacharparenleft}{\kern0pt}leq\ t\isactrlsub j{\isacharparenright}{\kern0pt}\ hs\isactrlsub {\isadigit{1}}{\isacharparenright}{\kern0pt}{\isacharparenright}{\kern0pt}\ M\isactrlsub j{\isachardoublequoteclose}\isanewline
\ \ \ \ \isakeywordTWO{shows}\ {\isachardoublequoteopen}valid{\isacharunderscore}{\kern0pt}happ{\isacharunderscore}{\kern0pt}seq\ M\isactrlsub i\ {\isacharparenleft}{\kern0pt}dropWhile\ {\isacharparenleft}{\kern0pt}leq\ t\isactrlsub i{\isacharparenright}{\kern0pt}\ {\isacharparenleft}{\kern0pt}takeWhile\ {\isacharparenleft}{\kern0pt}leq\ t\isactrlsub j{\isacharparenright}{\kern0pt}\ hs\isactrlsub {\isadigit{2}}{\isacharparenright}{\kern0pt}{\isacharparenright}{\kern0pt}\ M\isactrlsub j{\isachardoublequoteclose}\isanewline
%
\isadelimproof
\ \ %
\endisadelimproof
%
\isatagproof
\isakeywordONE{proof}\isamarkupfalse%
\ {\isacharminus}{\kern0pt}\isanewline
\ \ \ \ \isakeywordTHREE{obtain}\isamarkupfalse%
\ htps\ \isakeywordTWO{where}\ {\isachardoublequoteopen}htps{\isacharunderscore}{\kern0pt}seq\ {\isasympi}s\ htps{\isachardoublequoteclose}\isanewline
\ \ \ \ \ \ \isakeywordONE{using}\isamarkupfalse%
\ htps{\isacharunderscore}{\kern0pt}seq{\isacharunderscore}{\kern0pt}exists\ \isakeywordONE{by}\isamarkupfalse%
\ auto\isanewline
\ \ \ \ \isakeywordONE{then}\isamarkupfalse%
\ \isakeywordONE{have}\isamarkupfalse%
\ {\isachardoublequoteopen}strict{\isacharunderscore}{\kern0pt}sorted\ htps{\isachardoublequoteclose}\ \isakeywordTWO{and}\ {\isachardoublequoteopen}t\isactrlsub i\ {\isasymin}\ set\ htps{\isachardoublequoteclose}\ \isakeywordTWO{and}\ {\isachardoublequoteopen}t\isactrlsub j\ {\isasymin}\ set\ htps{\isachardoublequoteclose}\isanewline
\ \ \ \ \ \ \isakeywordONE{unfolding}\isamarkupfalse%
\ htps{\isacharunderscore}{\kern0pt}seq{\isacharunderscore}{\kern0pt}def\ \isakeywordONE{using}\isamarkupfalse%
\ assms\ \isakeywordONE{by}\isamarkupfalse%
\ {\isacharparenleft}{\kern0pt}auto\ simp{\isacharcolon}{\kern0pt}\ strict{\isacharunderscore}{\kern0pt}sorted{\isacharunderscore}{\kern0pt}iff{\isacharparenright}{\kern0pt}\isanewline
\ \ \ \ \isakeywordONE{then}\isamarkupfalse%
\ \isakeywordTHREE{obtain}\isamarkupfalse%
\ ts\isactrlsub {\isadigit{1}}\ ts\isactrlsub {\isadigit{2}}\ ts\isactrlsub {\isadigit{3}}\ \isakeywordTWO{where}\ {\isachardoublequoteopen}htps\ {\isacharequal}{\kern0pt}\ ts\isactrlsub {\isadigit{1}}\ {\isacharat}{\kern0pt}\ t\isactrlsub i\ {\isacharhash}{\kern0pt}\ ts\isactrlsub {\isadigit{2}}\ {\isacharat}{\kern0pt}\ t\isactrlsub j\ {\isacharhash}{\kern0pt}\ ts\isactrlsub {\isadigit{3}}{\isachardoublequoteclose}\isanewline
\ \ \ \ \ \ \isakeywordONE{using}\isamarkupfalse%
\ {\isacartoucheopen}t\isactrlsub i\ {\isacharless}{\kern0pt}\ t\isactrlsub j{\isacartoucheclose}\ split{\isacharunderscore}{\kern0pt}list{\isadigit{2}}{\isacharunderscore}{\kern0pt}le{\isacharbrackleft}{\kern0pt}OF\ {\isacartoucheopen}strict{\isacharunderscore}{\kern0pt}sorted\ htps{\isacartoucheclose}\ {\isacartoucheopen}t\isactrlsub i\ {\isasymin}\ set\ htps{\isacartoucheclose}\ {\isacartoucheopen}t\isactrlsub j\ {\isasymin}\ set\ htps{\isacartoucheclose}{\isacharbrackright}{\kern0pt}\ \isakeywordONE{by}\isamarkupfalse%
\ auto\isanewline
\ \ \ \ \isakeywordONE{then}\isamarkupfalse%
\ \isakeywordTHREE{show}\isamarkupfalse%
\ {\isacharquery}{\kern0pt}thesis\isanewline
\ \ \ \ \ \ \isakeywordONE{using}\isamarkupfalse%
\ assms\ {\isacartoucheopen}htps{\isacharunderscore}{\kern0pt}seq\ {\isasympi}s\ htps{\isacartoucheclose}\ htps{\isacharunderscore}{\kern0pt}state{\isacharunderscore}{\kern0pt}trace{\isacharunderscore}{\kern0pt}unique{\isacharunderscore}{\kern0pt}i{\isacharunderscore}{\kern0pt}to{\isacharunderscore}{\kern0pt}j{\isacharunderscore}{\kern0pt}aux\ \isakeywordONE{by}\isamarkupfalse%
\ metis\isanewline
\ \ \isakeywordONE{qed}\isamarkupfalse%
%
\endisatagproof
{\isafoldproof}%
%
\isadelimproof
\isanewline
%
\endisadelimproof
\isanewline
\ \ \isakeywordONE{lemma}\isamarkupfalse%
\ happ{\isacharunderscore}{\kern0pt}tp{\isacharunderscore}{\kern0pt}leq{\isacharunderscore}{\kern0pt}t\isactrlsub m\isactrlsub i\isactrlsub n{\isacharunderscore}{\kern0pt}t\isactrlsub m\isactrlsub a\isactrlsub x{\isacharcolon}{\kern0pt}\isanewline
\ \ \ \ \isakeywordTWO{assumes}\ {\isachardoublequoteopen}{\isacharparenleft}{\kern0pt}t\isactrlsub {\isasympi}{\isacharcomma}{\kern0pt}{\isasympi}{\isacharparenright}{\kern0pt}\ {\isasymin}\ set\ {\isasympi}s{\isachardoublequoteclose}\isanewline
\ \ \ \ \ \ \ \ \isakeywordTWO{and}\ {\isachardoublequoteopen}inst{\isacharunderscore}{\kern0pt}of{\isacharunderscore}{\kern0pt}plan{\isacharunderscore}{\kern0pt}action\ {\isasympi}s\ {\isacharparenleft}{\kern0pt}t\isactrlsub {\isasympi}{\isacharcomma}{\kern0pt}{\isasympi}{\isacharparenright}{\kern0pt}\ {\isacharparenleft}{\kern0pt}t\isactrlsub a{\isacharcomma}{\kern0pt}a{\isacharparenright}{\kern0pt}{\isachardoublequoteclose}\ \isanewline
\ \ \ \ \isakeywordTWO{shows}\ {\isachardoublequoteopen}{\isacharparenleft}{\kern0pt}{\isasymexists}t\isactrlsub j{\isachardot}{\kern0pt}\ is{\isacharunderscore}{\kern0pt}htp\ {\isasympi}s\ t\isactrlsub j\ {\isasymand}\ t\isactrlsub j\ {\isasymle}\ t\isactrlsub a{\isacharparenright}{\kern0pt}\ {\isasymand}\ {\isacharparenleft}{\kern0pt}{\isasymexists}t\isactrlsub j{\isachardot}{\kern0pt}\ is{\isacharunderscore}{\kern0pt}htp\ {\isasympi}s\ t\isactrlsub j\ {\isasymand}\ t\isactrlsub a\ {\isasymle}\ t\isactrlsub j{\isacharparenright}{\kern0pt}{\isachardoublequoteclose}\isanewline
%
\isadelimproof
\ \ %
\endisadelimproof
%
\isatagproof
\isakeywordONE{proof}\isamarkupfalse%
\ {\isacharminus}{\kern0pt}\isanewline
\ \ \ \ \isakeywordONE{let}\isamarkupfalse%
\ {\isacharquery}{\kern0pt}case{\isadigit{1}}{\isacharequal}{\kern0pt}{\isachardoublequoteopen}res{\isacharunderscore}{\kern0pt}inst\ {\isasympi}\ At{\isacharunderscore}{\kern0pt}Start\ {\isacharequal}{\kern0pt}\ Some\ a\ {\isasymand}\ t\isactrlsub {\isasympi}\ {\isacharequal}{\kern0pt}\ t\isactrlsub a{\isachardoublequoteclose}\isanewline
\ \ \ \ \isakeywordONE{let}\isamarkupfalse%
\ {\isacharquery}{\kern0pt}case{\isadigit{2}}{\isacharequal}{\kern0pt}{\isachardoublequoteopen}res{\isacharunderscore}{\kern0pt}inst{\isacharunderscore}{\kern0pt}snap{\isacharunderscore}{\kern0pt}action\ {\isasympi}\ At{\isacharunderscore}{\kern0pt}Start\ {\isacharequal}{\kern0pt}\ Some\ a\ {\isasymand}\ t\isactrlsub {\isasympi}\ {\isacharequal}{\kern0pt}\ t\isactrlsub a{\isachardoublequoteclose}\isanewline
\ \ \ \ \isakeywordONE{let}\isamarkupfalse%
\ {\isacharquery}{\kern0pt}case{\isadigit{3}}{\isacharequal}{\kern0pt}{\isachardoublequoteopen}res{\isacharunderscore}{\kern0pt}inst{\isacharunderscore}{\kern0pt}snap{\isacharunderscore}{\kern0pt}action\ {\isasympi}\ At{\isacharunderscore}{\kern0pt}End\ {\isacharequal}{\kern0pt}\ Some\ a\ {\isasymand}\ t\isactrlsub {\isasympi}\ {\isacharplus}{\kern0pt}\ {\isacharparenleft}{\kern0pt}duration\ {\isasympi}{\isacharparenright}{\kern0pt}\ {\isacharequal}{\kern0pt}\ t\isactrlsub a{\isachardoublequoteclose}\isanewline
\ \ \ \ \isakeywordONE{let}\isamarkupfalse%
\ {\isacharquery}{\kern0pt}case{\isadigit{4}}{\isacharequal}{\kern0pt}{\isachardoublequoteopen}res{\isacharunderscore}{\kern0pt}inst{\isacharunderscore}{\kern0pt}snap{\isacharunderscore}{\kern0pt}action\ {\isasympi}\ Over{\isacharunderscore}{\kern0pt}All\ {\isacharequal}{\kern0pt}\ Some\ a\ {\isasymand}\ \isanewline
\ \ \ \ \ \ \ \ \ \ {\isacharparenleft}{\kern0pt}{\isasymexists}t\isactrlsub i\ t\isactrlsub j{\isachardot}{\kern0pt}\ consec{\isacharunderscore}{\kern0pt}htps\ {\isasympi}s\ t\isactrlsub i\ t\isactrlsub j\ {\isasymand}\ t\isactrlsub {\isasympi}\ {\isasymle}\ t\isactrlsub i\ {\isasymand}\ t\isactrlsub j\ {\isasymle}\ t\isactrlsub {\isasympi}\ {\isacharplus}{\kern0pt}\ {\isacharparenleft}{\kern0pt}duration\ {\isasympi}{\isacharparenright}{\kern0pt}\ {\isasymand}\ t\isactrlsub i\ {\isacharless}{\kern0pt}\ t\isactrlsub a\ {\isasymand}\ t\isactrlsub a\ {\isacharless}{\kern0pt}\ t\isactrlsub j{\isacharparenright}{\kern0pt}{\isachardoublequoteclose}\isanewline
\ \ \ \ \isakeywordONE{have}\isamarkupfalse%
\ {\isachardoublequoteopen}{\isacharquery}{\kern0pt}case{\isadigit{1}}\ {\isasymor}\ {\isacharquery}{\kern0pt}case{\isadigit{2}}\ {\isasymor}\ {\isacharquery}{\kern0pt}case{\isadigit{3}}\ {\isasymor}\ {\isacharquery}{\kern0pt}case{\isadigit{4}}{\isachardoublequoteclose}\isanewline
\ \ \ \ \ \ \isakeywordONE{using}\isamarkupfalse%
\ {\isacartoucheopen}inst{\isacharunderscore}{\kern0pt}of{\isacharunderscore}{\kern0pt}plan{\isacharunderscore}{\kern0pt}action\ {\isasympi}s\ {\isacharparenleft}{\kern0pt}t\isactrlsub {\isasympi}{\isacharcomma}{\kern0pt}{\isasympi}{\isacharparenright}{\kern0pt}\ {\isacharparenleft}{\kern0pt}t\isactrlsub a{\isacharcomma}{\kern0pt}a{\isacharparenright}{\kern0pt}{\isacartoucheclose}\ \isakeywordONE{by}\isamarkupfalse%
\ simp\isanewline
\ \ \ \ \isakeywordONE{then}\isamarkupfalse%
\ \isakeywordTHREE{show}\isamarkupfalse%
\ {\isacharquery}{\kern0pt}thesis\isanewline
\ \ \ \ \isakeywordONE{proof}\isamarkupfalse%
\ {\isacharparenleft}{\kern0pt}elim\ disjE{\isacharparenright}{\kern0pt}\isanewline
\ \ \ \ \ \ \isakeywordTHREE{assume}\isamarkupfalse%
\ {\isacharquery}{\kern0pt}case{\isadigit{1}}\isanewline
\ \ \ \ \ \ \isakeywordONE{then}\isamarkupfalse%
\ \isakeywordTHREE{show}\isamarkupfalse%
\ {\isacharquery}{\kern0pt}thesis\isanewline
\ \ \ \ \ \ \ \ \isakeywordONE{using}\isamarkupfalse%
\ {\isacartoucheopen}{\isacharparenleft}{\kern0pt}t\isactrlsub {\isasympi}{\isacharcomma}{\kern0pt}{\isasympi}{\isacharparenright}{\kern0pt}\ {\isasymin}\ set\ {\isasympi}s{\isacartoucheclose}\ \isakeywordONE{by}\isamarkupfalse%
\ {\isacharparenleft}{\kern0pt}auto\ simp{\isacharcolon}{\kern0pt}\ is{\isacharunderscore}{\kern0pt}htp{\isacharunderscore}{\kern0pt}def{\isacharparenright}{\kern0pt}\isanewline
\ \ \ \ \isakeywordONE{next}\isamarkupfalse%
\isanewline
\ \ \ \ \ \ \isakeywordTHREE{assume}\isamarkupfalse%
\ {\isacharquery}{\kern0pt}case{\isadigit{2}}\isanewline
\ \ \ \ \ \ \isakeywordONE{then}\isamarkupfalse%
\ \isakeywordTHREE{show}\isamarkupfalse%
\ {\isacharquery}{\kern0pt}thesis\isanewline
\ \ \ \ \ \ \ \ \isakeywordONE{using}\isamarkupfalse%
\ {\isacartoucheopen}{\isacharparenleft}{\kern0pt}t\isactrlsub {\isasympi}{\isacharcomma}{\kern0pt}{\isasympi}{\isacharparenright}{\kern0pt}\ {\isasymin}\ set\ {\isasympi}s{\isacartoucheclose}\ \isakeywordONE{by}\isamarkupfalse%
\ {\isacharparenleft}{\kern0pt}auto\ simp{\isacharcolon}{\kern0pt}\ is{\isacharunderscore}{\kern0pt}htp{\isacharunderscore}{\kern0pt}def{\isacharparenright}{\kern0pt}\isanewline
\ \ \ \ \isakeywordONE{next}\isamarkupfalse%
\isanewline
\ \ \ \ \ \ \isakeywordTHREE{assume}\isamarkupfalse%
\ {\isacharquery}{\kern0pt}case{\isadigit{3}}\isanewline
\ \ \ \ \ \ \isakeywordONE{then}\isamarkupfalse%
\ \isakeywordTHREE{show}\isamarkupfalse%
\ {\isacharquery}{\kern0pt}thesis\isanewline
\ \ \ \ \ \ \ \ \isakeywordONE{unfolding}\isamarkupfalse%
\ is{\isacharunderscore}{\kern0pt}htp{\isacharunderscore}{\kern0pt}def\ durative{\isacharunderscore}{\kern0pt}acts{\isacharunderscore}{\kern0pt}def\ is{\isacharunderscore}{\kern0pt}act{\isacharunderscore}{\kern0pt}simple{\isacharunderscore}{\kern0pt}def\isanewline
\ \ \ \ \ \ \ \ \isakeywordONE{using}\isamarkupfalse%
\ {\isacartoucheopen}{\isacharparenleft}{\kern0pt}t\isactrlsub {\isasympi}{\isacharcomma}{\kern0pt}{\isasympi}{\isacharparenright}{\kern0pt}\ {\isasymin}\ set\ {\isasympi}s{\isacartoucheclose}\ {\isacartoucheopen}{\isacharquery}{\kern0pt}case{\isadigit{3}}{\isacartoucheclose}\ \isakeywordONE{by}\isamarkupfalse%
\ {\isacharparenleft}{\kern0pt}cases\ {\isasympi}{\isacharparenright}{\kern0pt}\ fastforce{\isacharplus}{\kern0pt}\isanewline
\ \ \ \ \isakeywordONE{next}\isamarkupfalse%
\isanewline
\ \ \ \ \ \ \isakeywordTHREE{assume}\isamarkupfalse%
\ {\isacharquery}{\kern0pt}case{\isadigit{4}}\isanewline
\ \ \ \ \ \ \isakeywordONE{then}\isamarkupfalse%
\ \isakeywordTHREE{obtain}\isamarkupfalse%
\ t\isactrlsub i\ t\isactrlsub j\ \isakeywordTWO{where}\ {\isachardoublequoteopen}consec{\isacharunderscore}{\kern0pt}htps\ {\isasympi}s\ t\isactrlsub i\ t\isactrlsub j{\isachardoublequoteclose}\ \isakeywordTWO{and}\ {\isachardoublequoteopen}t\isactrlsub {\isasympi}\ {\isasymle}\ t\isactrlsub i\ {\isasymand}\ t\isactrlsub i\ {\isacharless}{\kern0pt}\ t\isactrlsub {\isasympi}\ {\isacharplus}{\kern0pt}\ {\isacharparenleft}{\kern0pt}duration\ {\isasympi}{\isacharparenright}{\kern0pt}{\isachardoublequoteclose}\ \isanewline
\ \ \ \ \ \ \ \ \ \ \ \ \ \ \ \ \ \ \ \ \ \ \ \ \ \isakeywordTWO{and}\ {\isachardoublequoteopen}t\isactrlsub i\ {\isacharless}{\kern0pt}\ t\isactrlsub a\ {\isasymand}\ t\isactrlsub a\ {\isacharless}{\kern0pt}\ t\isactrlsub j{\isachardoublequoteclose}\isanewline
\ \ \ \ \ \ \ \ \isakeywordONE{by}\isamarkupfalse%
\ auto\isanewline
\ \ \ \ \ \ \isakeywordONE{then}\isamarkupfalse%
\ \isakeywordTHREE{show}\isamarkupfalse%
\ {\isacharquery}{\kern0pt}thesis\isanewline
\ \ \ \ \ \ \ \ \isakeywordONE{using}\isamarkupfalse%
\ {\isacartoucheopen}consec{\isacharunderscore}{\kern0pt}htps\ {\isasympi}s\ t\isactrlsub i\ t\isactrlsub j{\isacartoucheclose}\ \isakeywordONE{unfolding}\isamarkupfalse%
\ consec{\isacharunderscore}{\kern0pt}htps{\isacharunderscore}{\kern0pt}def\ \isakeywordONE{by}\isamarkupfalse%
\ auto\isanewline
\ \ \ \ \isakeywordONE{qed}\isamarkupfalse%
\isanewline
\ \ \isakeywordONE{qed}\isamarkupfalse%
%
\endisatagproof
{\isafoldproof}%
%
\isadelimproof
%
\endisadelimproof
%
\begin{isamarkuptext}%
Proof that every happening of a induced happening sequence is in between 
  the smallest and largest happening time point.%
\end{isamarkuptext}\isamarkuptrue%
\ \ \isakeywordONE{lemma}\isamarkupfalse%
\ ind{\isacharunderscore}{\kern0pt}happ{\isacharunderscore}{\kern0pt}seq{\isacharunderscore}{\kern0pt}htps{\isacharunderscore}{\kern0pt}leq{\isacharunderscore}{\kern0pt}t\isactrlsub m\isactrlsub i\isactrlsub n{\isacharunderscore}{\kern0pt}t\isactrlsub m\isactrlsub a\isactrlsub x{\isacharcolon}{\kern0pt}\isanewline
\ \ \ \ \isakeywordTWO{fixes}\ {\isasympi}s\isanewline
\ \ \ \ \isakeywordTWO{assumes}\ {\isachardoublequoteopen}ind{\isacharunderscore}{\kern0pt}happ{\isacharunderscore}{\kern0pt}seq\ {\isasympi}s\ {\isacharparenleft}{\kern0pt}hs{\isacharprime}{\kern0pt}{\isacharat}{\kern0pt}hs{\isacharparenright}{\kern0pt}{\isachardoublequoteclose}\isanewline
\ \ \ \ \ \ \ \ \isakeywordTWO{and}\ {\isachardoublequoteopen}{\isasymforall}t\isactrlsub i{\isachardot}{\kern0pt}\ is{\isacharunderscore}{\kern0pt}htp\ {\isasympi}s\ t\isactrlsub i\ {\isasymlongrightarrow}\ t\isactrlsub m\isactrlsub i\isactrlsub n\ {\isasymle}\ t\isactrlsub i{\isachardoublequoteclose}\isanewline
\ \ \ \ \ \ \ \ \isakeywordTWO{and}\ {\isachardoublequoteopen}{\isasymforall}t\isactrlsub i{\isachardot}{\kern0pt}\ is{\isacharunderscore}{\kern0pt}htp\ {\isasympi}s\ t\isactrlsub i\ {\isasymlongrightarrow}\ t\isactrlsub i\ {\isasymle}\ t\isactrlsub m\isactrlsub a\isactrlsub x{\isachardoublequoteclose}\isanewline
\ \ \ \ \isakeywordTWO{shows}\ {\isachardoublequoteopen}{\isasymforall}{\isacharparenleft}{\kern0pt}t\isactrlsub a{\isacharcomma}{\kern0pt}A{\isacharparenright}{\kern0pt}\ {\isasymin}\ set\ hs{\isachardot}{\kern0pt}\ t\isactrlsub m\isactrlsub i\isactrlsub n\ {\isasymle}\ t\isactrlsub a\ {\isasymand}\ t\isactrlsub a\ {\isasymle}\ t\isactrlsub m\isactrlsub a\isactrlsub x{\isachardoublequoteclose}\isanewline
%
\isadelimproof
\ \ \ \ %
\endisadelimproof
%
\isatagproof
\isakeywordONE{using}\isamarkupfalse%
\ assms\isanewline
\ \ \isakeywordONE{proof}\isamarkupfalse%
\ {\isacharparenleft}{\kern0pt}induction\ hs\ arbitrary{\isacharcolon}{\kern0pt}\ hs{\isacharprime}{\kern0pt}{\isacharparenright}{\kern0pt}\isanewline
\ \ \ \ \isakeywordTHREE{case}\isamarkupfalse%
\ Nil\isanewline
\ \ \ \ \isakeywordONE{then}\isamarkupfalse%
\ \isakeywordTHREE{show}\isamarkupfalse%
\ {\isacharquery}{\kern0pt}case\ \isanewline
\ \ \ \ \ \ \isakeywordONE{by}\isamarkupfalse%
\ simp\isanewline
\ \ \isakeywordONE{next}\isamarkupfalse%
\isanewline
\ \ \ \ \isakeywordTHREE{case}\isamarkupfalse%
\ {\isacharparenleft}{\kern0pt}Cons\ h\ hs{\isacharparenright}{\kern0pt}\isanewline
\ \ \ \ \isakeywordONE{then}\isamarkupfalse%
\ \isakeywordTHREE{show}\isamarkupfalse%
\ {\isacharquery}{\kern0pt}case\isanewline
\ \ \ \ \isakeywordONE{proof}\isamarkupfalse%
\ {\isacharparenleft}{\kern0pt}cases\ h{\isacharparenright}{\kern0pt}\isanewline
\ \ \ \ \ \ \isakeywordTHREE{case}\isamarkupfalse%
\ {\isacharparenleft}{\kern0pt}Pair\ t\isactrlsub a\ A{\isacharparenright}{\kern0pt}\isanewline
\ \ \ \ \ \ \ \ \isakeywordONE{then}\isamarkupfalse%
\ \isakeywordTHREE{obtain}\isamarkupfalse%
\ a\ \isakeywordTWO{where}\ {\isachardoublequoteopen}a\ {\isasymin}\ set\ A{\isachardoublequoteclose}\isanewline
\ \ \ \ \ \ \ \ \ \ \isakeywordONE{using}\isamarkupfalse%
\ {\isacartoucheopen}ind{\isacharunderscore}{\kern0pt}happ{\isacharunderscore}{\kern0pt}seq\ {\isasympi}s\ {\isacharparenleft}{\kern0pt}hs{\isacharprime}{\kern0pt}{\isacharat}{\kern0pt}h{\isacharhash}{\kern0pt}hs{\isacharparenright}{\kern0pt}{\isacartoucheclose}\ \isakeywordONE{unfolding}\isamarkupfalse%
\ ind{\isacharunderscore}{\kern0pt}happ{\isacharunderscore}{\kern0pt}seq{\isacharunderscore}{\kern0pt}def\ \isakeywordONE{by}\isamarkupfalse%
\ {\isacharparenleft}{\kern0pt}cases\ A{\isacharparenright}{\kern0pt}\ auto\isanewline
\ \ \ \ \ \ \ \ \isakeywordONE{then}\isamarkupfalse%
\ \isakeywordONE{have}\isamarkupfalse%
\ {\isachardoublequoteopen}{\isasymexists}{\isacharparenleft}{\kern0pt}t\isactrlsub {\isasympi}{\isacharcomma}{\kern0pt}{\isasympi}{\isacharparenright}{\kern0pt}\ {\isasymin}\ set\ {\isasympi}s{\isachardot}{\kern0pt}\ inst{\isacharunderscore}{\kern0pt}of{\isacharunderscore}{\kern0pt}plan{\isacharunderscore}{\kern0pt}action\ {\isasympi}s\ {\isacharparenleft}{\kern0pt}t\isactrlsub {\isasympi}{\isacharcomma}{\kern0pt}{\isasympi}{\isacharparenright}{\kern0pt}\ {\isacharparenleft}{\kern0pt}t\isactrlsub a{\isacharcomma}{\kern0pt}a{\isacharparenright}{\kern0pt}{\isachardoublequoteclose}\isanewline
\ \ \ \ \ \ \ \ \ \ \isakeywordONE{using}\isamarkupfalse%
\ {\isacartoucheopen}ind{\isacharunderscore}{\kern0pt}happ{\isacharunderscore}{\kern0pt}seq\ {\isasympi}s\ {\isacharparenleft}{\kern0pt}hs{\isacharprime}{\kern0pt}{\isacharat}{\kern0pt}h{\isacharhash}{\kern0pt}hs{\isacharparenright}{\kern0pt}{\isacartoucheclose}\ Pair\ \isakeywordONE{unfolding}\isamarkupfalse%
\ ind{\isacharunderscore}{\kern0pt}happ{\isacharunderscore}{\kern0pt}seq{\isacharunderscore}{\kern0pt}def\ \isakeywordONE{by}\isamarkupfalse%
\ auto\isanewline
\ \ \ \ \ \ \ \ \isakeywordONE{then}\isamarkupfalse%
\ \isakeywordTHREE{obtain}\isamarkupfalse%
\ t\isactrlsub {\isasympi}\ {\isasympi}\ \isakeywordTWO{where}\ {\isachardoublequoteopen}{\isacharparenleft}{\kern0pt}t\isactrlsub {\isasympi}{\isacharcomma}{\kern0pt}{\isasympi}{\isacharparenright}{\kern0pt}\ {\isasymin}\ set\ {\isasympi}s{\isachardoublequoteclose}\ \isakeywordTWO{and}\ {\isachardoublequoteopen}inst{\isacharunderscore}{\kern0pt}of{\isacharunderscore}{\kern0pt}plan{\isacharunderscore}{\kern0pt}action\ {\isasympi}s\ {\isacharparenleft}{\kern0pt}t\isactrlsub {\isasympi}{\isacharcomma}{\kern0pt}{\isasympi}{\isacharparenright}{\kern0pt}\ {\isacharparenleft}{\kern0pt}t\isactrlsub a{\isacharcomma}{\kern0pt}a{\isacharparenright}{\kern0pt}{\isachardoublequoteclose}\isanewline
\ \ \ \ \ \ \ \ \ \ \isakeywordONE{by}\isamarkupfalse%
\ auto\isanewline
\ \ \ \ \ \ \ \ \isakeywordTHREE{obtain}\isamarkupfalse%
\ hs{\isacharprime}{\kern0pt}{\isacharprime}{\kern0pt}\ \isakeywordTWO{where}\ {\isachardoublequoteopen}hs{\isacharprime}{\kern0pt}{\isacharprime}{\kern0pt}{\isacharequal}{\kern0pt}hs{\isacharprime}{\kern0pt}{\isacharat}{\kern0pt}{\isacharbrackleft}{\kern0pt}h{\isacharbrackright}{\kern0pt}{\isachardoublequoteclose}\isanewline
\ \ \ \ \ \ \ \ \ \ \isakeywordONE{by}\isamarkupfalse%
\ simp\isanewline
\ \ \ \ \ \ \ \ \isakeywordONE{then}\isamarkupfalse%
\ \isakeywordONE{have}\isamarkupfalse%
\ {\isachardoublequoteopen}ind{\isacharunderscore}{\kern0pt}happ{\isacharunderscore}{\kern0pt}seq\ {\isasympi}s\ {\isacharparenleft}{\kern0pt}hs{\isacharprime}{\kern0pt}{\isacharprime}{\kern0pt}{\isacharat}{\kern0pt}hs{\isacharparenright}{\kern0pt}{\isachardoublequoteclose}\isanewline
\ \ \ \ \ \ \ \ \ \ \isakeywordONE{using}\isamarkupfalse%
\ {\isacartoucheopen}ind{\isacharunderscore}{\kern0pt}happ{\isacharunderscore}{\kern0pt}seq\ {\isasympi}s\ {\isacharparenleft}{\kern0pt}hs{\isacharprime}{\kern0pt}{\isacharat}{\kern0pt}h{\isacharhash}{\kern0pt}hs{\isacharparenright}{\kern0pt}{\isacartoucheclose}\ \isakeywordONE{by}\isamarkupfalse%
\ simp\isanewline
\ \ \ \ \ \ \ \ \ \isakeywordTHREE{show}\isamarkupfalse%
\ {\isacharquery}{\kern0pt}thesis\isanewline
\ \ \ \ \ \ \ \ \ \ \ \isakeywordONE{using}\isamarkupfalse%
\ assms\ Pair\ Cons{\isachardot}{\kern0pt}IH{\isacharbrackleft}{\kern0pt}OF\ {\isacartoucheopen}ind{\isacharunderscore}{\kern0pt}happ{\isacharunderscore}{\kern0pt}seq\ {\isasympi}s\ {\isacharparenleft}{\kern0pt}hs{\isacharprime}{\kern0pt}{\isacharprime}{\kern0pt}{\isacharat}{\kern0pt}hs{\isacharparenright}{\kern0pt}{\isacartoucheclose}{\isacharbrackright}{\kern0pt}\isanewline
\ \ \ \ \ \ \ \ \ \ \ \ \ \ \ \ \ happ{\isacharunderscore}{\kern0pt}tp{\isacharunderscore}{\kern0pt}leq{\isacharunderscore}{\kern0pt}t\isactrlsub m\isactrlsub i\isactrlsub n{\isacharunderscore}{\kern0pt}t\isactrlsub m\isactrlsub a\isactrlsub x{\isacharbrackleft}{\kern0pt}OF\ {\isacartoucheopen}{\isacharparenleft}{\kern0pt}t\isactrlsub {\isasympi}{\isacharcomma}{\kern0pt}{\isasympi}{\isacharparenright}{\kern0pt}\ {\isasymin}\ set\ {\isasympi}s{\isacartoucheclose}\ {\isacartoucheopen}inst{\isacharunderscore}{\kern0pt}of{\isacharunderscore}{\kern0pt}plan{\isacharunderscore}{\kern0pt}action\ {\isasympi}s\ {\isacharparenleft}{\kern0pt}t\isactrlsub {\isasympi}{\isacharcomma}{\kern0pt}{\isasympi}{\isacharparenright}{\kern0pt}\ {\isacharparenleft}{\kern0pt}t\isactrlsub a{\isacharcomma}{\kern0pt}a{\isacharparenright}{\kern0pt}{\isacartoucheclose}{\isacharbrackright}{\kern0pt}\isanewline
\ \ \ \ \ \ \ \ \ \ \ \isakeywordONE{by}\isamarkupfalse%
\ auto\isanewline
\ \ \ \ \isakeywordONE{qed}\isamarkupfalse%
\isanewline
\ \ \isakeywordONE{qed}\isamarkupfalse%
%
\endisatagproof
{\isafoldproof}%
%
\isadelimproof
\isanewline
%
\endisadelimproof
\isanewline
\ \ \isakeywordONE{lemma}\isamarkupfalse%
\ valid{\isacharunderscore}{\kern0pt}happ{\isacharunderscore}{\kern0pt}seq{\isacharunderscore}{\kern0pt}state{\isacharunderscore}{\kern0pt}unique{\isacharcolon}{\kern0pt}\isanewline
\ \ \ \ \isakeywordTWO{assumes}\ {\isachardoublequoteopen}valid{\isacharunderscore}{\kern0pt}happ{\isacharunderscore}{\kern0pt}seq\ M\isactrlsub {\isadigit{0}}\ hs\ M\isactrlsub i{\isachardoublequoteclose}\ \isakeywordTWO{and}\ {\isachardoublequoteopen}valid{\isacharunderscore}{\kern0pt}happ{\isacharunderscore}{\kern0pt}seq\ M\isactrlsub {\isadigit{0}}\ hs\ M\isactrlsub i{\isacharprime}{\kern0pt}{\isachardoublequoteclose}\isanewline
\ \ \ \ \isakeywordTWO{shows}\ {\isachardoublequoteopen}M\isactrlsub i\ {\isacharequal}{\kern0pt}\ M\isactrlsub i{\isacharprime}{\kern0pt}{\isachardoublequoteclose}\isanewline
%
\isadelimproof
\ \ \ \ %
\endisadelimproof
%
\isatagproof
\isakeywordONE{using}\isamarkupfalse%
\ assms\ \isakeywordONE{by}\isamarkupfalse%
\ {\isacharparenleft}{\kern0pt}induction\ M\isactrlsub {\isadigit{0}}\ hs\ M\isactrlsub i\ rule{\isacharcolon}{\kern0pt}\ valid{\isacharunderscore}{\kern0pt}happ{\isacharunderscore}{\kern0pt}seq{\isachardot}{\kern0pt}induct{\isacharparenright}{\kern0pt}\ auto%
\endisatagproof
{\isafoldproof}%
%
\isadelimproof
\isanewline
%
\endisadelimproof
\isanewline
\ \ \isakeywordONE{lemma}\isamarkupfalse%
\ \ valid{\isacharunderscore}{\kern0pt}happ{\isacharunderscore}{\kern0pt}seq{\isacharunderscore}{\kern0pt}iff{\isacharunderscore}{\kern0pt}valid{\isacharunderscore}{\kern0pt}state{\isacharunderscore}{\kern0pt}seq{\isacharcolon}{\kern0pt}\isanewline
\ \ \ \ \isakeywordTWO{fixes}\ {\isasympi}s\isanewline
\ \ \ \ \isakeywordTWO{assumes}\ {\isachardoublequoteopen}htps{\isacharunderscore}{\kern0pt}seq\ {\isasympi}s\ htps{\isachardoublequoteclose}\isanewline
\ \ \ \ \ \ \ \ \isakeywordTWO{and}\ {\isachardoublequoteopen}wf{\isacharunderscore}{\kern0pt}plan\ {\isasympi}s{\isachardoublequoteclose}\isanewline
\ \ \ \ \ \ \ \ \isakeywordTWO{and}\ {\isachardoublequoteopen}ind{\isacharunderscore}{\kern0pt}happ{\isacharunderscore}{\kern0pt}seq\ {\isasympi}s\ hs{\isachardoublequoteclose}\isanewline
\ \ \ \ \isakeywordTWO{shows}\ {\isachardoublequoteopen}valid{\isacharunderscore}{\kern0pt}happ{\isacharunderscore}{\kern0pt}seq\ M\isactrlsub {\isadigit{0}}\ hs\ M\isactrlsub n\ {\isasymlongleftrightarrow}\ valid{\isacharunderscore}{\kern0pt}state{\isacharunderscore}{\kern0pt}seq\ M\isactrlsub {\isadigit{0}}\ htps\ {\isasympi}s\ M\isactrlsub n{\isachardoublequoteclose}\isanewline
%
\isadelimproof
\ \ %
\endisadelimproof
%
\isatagproof
\isakeywordONE{proof}\isamarkupfalse%
\ {\isacharparenleft}{\kern0pt}cases\ {\isachardoublequoteopen}{\isasympi}s\ {\isacharequal}{\kern0pt}\ {\isacharbrackleft}{\kern0pt}{\isacharbrackright}{\kern0pt}{\isachardoublequoteclose}{\isacharparenright}{\kern0pt}\isanewline
\ \ \ \ \isakeywordTHREE{case}\isamarkupfalse%
\ True\isanewline
\ \ \ \ \ \ \isakeywordONE{then}\isamarkupfalse%
\ \isakeywordONE{have}\isamarkupfalse%
\ {\isachardoublequoteopen}hs\ {\isacharequal}{\kern0pt}\ {\isacharbrackleft}{\kern0pt}{\isacharbrackright}{\kern0pt}{\isachardoublequoteclose}\ \isakeywordTWO{and}\ {\isachardoublequoteopen}htps\ {\isacharequal}{\kern0pt}\ {\isacharbrackleft}{\kern0pt}{\isacharbrackright}{\kern0pt}{\isachardoublequoteclose}\isanewline
\ \ \ \ \ \ \ \ \isakeywordONE{using}\isamarkupfalse%
\ assms\ ind{\isacharunderscore}{\kern0pt}happ{\isacharunderscore}{\kern0pt}seq{\isacharunderscore}{\kern0pt}Nil\ htps{\isacharunderscore}{\kern0pt}seq{\isacharunderscore}{\kern0pt}Nil{\isacharunderscore}{\kern0pt}iff\ \isakeywordONE{by}\isamarkupfalse%
\ auto\isanewline
\ \ \ \ \ \ \isakeywordONE{then}\isamarkupfalse%
\ \isakeywordTHREE{show}\isamarkupfalse%
\ {\isacharquery}{\kern0pt}thesis\isanewline
\ \ \ \ \ \ \ \ \isakeywordONE{by}\isamarkupfalse%
\ auto\isanewline
\ \ \isakeywordONE{next}\isamarkupfalse%
\isanewline
\ \ \ \ \isakeywordTHREE{case}\isamarkupfalse%
\ False\isanewline
\isanewline
\ \ \ \ \isakeywordONE{have}\isamarkupfalse%
\ {\isachardoublequoteopen}strict{\isacharunderscore}{\kern0pt}sorted\ {\isacharparenleft}{\kern0pt}map\ fst\ hs{\isacharparenright}{\kern0pt}{\isachardoublequoteclose}\isanewline
\ \ \ \ \ \ \isakeywordONE{using}\isamarkupfalse%
\ {\isacartoucheopen}ind{\isacharunderscore}{\kern0pt}happ{\isacharunderscore}{\kern0pt}seq\ {\isasympi}s\ hs{\isacartoucheclose}\ \isakeywordONE{unfolding}\isamarkupfalse%
\ ind{\isacharunderscore}{\kern0pt}happ{\isacharunderscore}{\kern0pt}seq{\isacharunderscore}{\kern0pt}def\ \isakeywordONE{by}\isamarkupfalse%
\ auto\isanewline
\ \ \ \ \isakeywordONE{have}\isamarkupfalse%
\ {\isachardoublequoteopen}htps\ {\isasymnoteq}\ {\isacharbrackleft}{\kern0pt}{\isacharbrackright}{\kern0pt}{\isachardoublequoteclose}\ \isakeywordTWO{and}\ {\isachardoublequoteopen}sorted\ htps{\isachardoublequoteclose}\ \isakeywordTWO{and}\ {\isachardoublequoteopen}distinct\ htps{\isachardoublequoteclose}\ \isakeywordTWO{and}\ {\isachardoublequoteopen}strict{\isacharunderscore}{\kern0pt}sorted\ htps{\isachardoublequoteclose}\isanewline
\ \ \ \ \ \ \isakeywordONE{using}\isamarkupfalse%
\ assms\ {\isacartoucheopen}{\isasympi}s\ {\isasymnoteq}\ {\isacharbrackleft}{\kern0pt}{\isacharbrackright}{\kern0pt}{\isacartoucheclose}\ htps{\isacharunderscore}{\kern0pt}seq{\isacharunderscore}{\kern0pt}Nil{\isacharunderscore}{\kern0pt}iff\ \isakeywordONE{unfolding}\isamarkupfalse%
\ htps{\isacharunderscore}{\kern0pt}seq{\isacharunderscore}{\kern0pt}def\isanewline
\ \ \ \ \ \ \isakeywordONE{by}\isamarkupfalse%
\ {\isacharparenleft}{\kern0pt}auto\ simp{\isacharcolon}{\kern0pt}\ strict{\isacharunderscore}{\kern0pt}sorted{\isacharunderscore}{\kern0pt}iff{\isacharparenright}{\kern0pt}%
\begin{isamarkuptext}%
Obtain the smallest happening time point.%
\end{isamarkuptext}\isamarkuptrue%
\ \ \ \ \isakeywordTHREE{obtain}\isamarkupfalse%
\ t\isactrlsub m\isactrlsub i\isactrlsub n\ \isakeywordTWO{where}\ {\isachardoublequoteopen}t\isactrlsub m\isactrlsub i\isactrlsub n\ {\isacharequal}{\kern0pt}\ hd\ htps{\isachardoublequoteclose}\isanewline
\ \ \ \ \ \ \isakeywordONE{by}\isamarkupfalse%
\ auto\isanewline
\ \ \ \ \isakeywordONE{then}\isamarkupfalse%
\ \isakeywordONE{have}\isamarkupfalse%
\ {\isachardoublequoteopen}t\isactrlsub m\isactrlsub i\isactrlsub n\ {\isasymin}\ set\ htps{\isachardoublequoteclose}\isanewline
\ \ \ \ \ \ \isakeywordONE{using}\isamarkupfalse%
\ {\isacartoucheopen}htps\ {\isasymnoteq}\ {\isacharbrackleft}{\kern0pt}{\isacharbrackright}{\kern0pt}{\isacartoucheclose}\ \isakeywordONE{by}\isamarkupfalse%
\ auto\isanewline
\ \ \ \ \isakeywordONE{have}\isamarkupfalse%
\ {\isachardoublequoteopen}is{\isacharunderscore}{\kern0pt}htp\ {\isasympi}s\ t\isactrlsub m\isactrlsub i\isactrlsub n{\isachardoublequoteclose}\ \isakeywordTWO{and}\ {\isachardoublequoteopen}{\isasymforall}t\isactrlsub k{\isachardot}{\kern0pt}\ is{\isacharunderscore}{\kern0pt}htp\ {\isasympi}s\ t\isactrlsub k\ {\isasymlongrightarrow}\ t\isactrlsub m\isactrlsub i\isactrlsub n\ {\isasymle}\ t\isactrlsub k{\isachardoublequoteclose}\isanewline
\ \ \ \ \ \ \isakeywordONE{using}\isamarkupfalse%
\ assms\ {\isacartoucheopen}htps\ {\isasymnoteq}\ {\isacharbrackleft}{\kern0pt}{\isacharbrackright}{\kern0pt}{\isacartoucheclose}\ {\isacartoucheopen}sorted\ htps{\isacartoucheclose}\ {\isacartoucheopen}t\isactrlsub m\isactrlsub i\isactrlsub n\ {\isacharequal}{\kern0pt}\ hd\ htps{\isacartoucheclose}\ sorted{\isacharunderscore}{\kern0pt}hd{\isacharunderscore}{\kern0pt}min\ \isanewline
\ \ \ \ \ \ \isakeywordONE{by}\isamarkupfalse%
\ {\isacharparenleft}{\kern0pt}auto\ simp{\isacharcolon}{\kern0pt}\ htps{\isacharunderscore}{\kern0pt}seq{\isacharunderscore}{\kern0pt}def{\isacharparenright}{\kern0pt}\isanewline
\ \ \ \ \isakeywordONE{have}\isamarkupfalse%
\ {\isachardoublequoteopen}{\isasymforall}{\isacharparenleft}{\kern0pt}t\isactrlsub {\isasympi}{\isacharcomma}{\kern0pt}{\isasympi}{\isacharparenright}{\kern0pt}\ {\isasymin}\ set\ {\isasympi}s{\isachardot}{\kern0pt}\ t\isactrlsub m\isactrlsub i\isactrlsub n\ {\isasymle}\ t\isactrlsub {\isasympi}{\isachardoublequoteclose}\isanewline
\ \ \ \ \ \ \isakeywordONE{using}\isamarkupfalse%
\ {\isacartoucheopen}{\isasymforall}t\isactrlsub k{\isachardot}{\kern0pt}\ is{\isacharunderscore}{\kern0pt}htp\ {\isasympi}s\ t\isactrlsub k\ {\isasymlongrightarrow}\ t\isactrlsub m\isactrlsub i\isactrlsub n\ {\isasymle}\ t\isactrlsub k{\isacartoucheclose}\ \isakeywordONE{by}\isamarkupfalse%
\ {\isacharparenleft}{\kern0pt}auto\ simp{\isacharcolon}{\kern0pt}\ is{\isacharunderscore}{\kern0pt}htp{\isacharunderscore}{\kern0pt}def{\isacharparenright}{\kern0pt}%
\begin{isamarkuptext}%
Obtain the largest happening time point.%
\end{isamarkuptext}\isamarkuptrue%
\ \ \ \ \isakeywordTHREE{obtain}\isamarkupfalse%
\ t\isactrlsub m\isactrlsub a\isactrlsub x\ \isakeywordTWO{where}\ {\isachardoublequoteopen}t\isactrlsub m\isactrlsub a\isactrlsub x\ {\isacharequal}{\kern0pt}\ last\ htps{\isachardoublequoteclose}\isanewline
\ \ \ \ \ \ \isakeywordONE{by}\isamarkupfalse%
\ auto\isanewline
\ \ \ \ \isakeywordONE{then}\isamarkupfalse%
\ \isakeywordONE{have}\isamarkupfalse%
\ {\isachardoublequoteopen}t\isactrlsub m\isactrlsub a\isactrlsub x\ {\isasymin}\ set\ htps{\isachardoublequoteclose}\isanewline
\ \ \ \ \ \ \isakeywordONE{using}\isamarkupfalse%
\ {\isacartoucheopen}htps\ {\isasymnoteq}\ {\isacharbrackleft}{\kern0pt}{\isacharbrackright}{\kern0pt}{\isacartoucheclose}\ \isakeywordONE{by}\isamarkupfalse%
\ auto\isanewline
\ \ \ \ \isakeywordONE{then}\isamarkupfalse%
\ \isakeywordONE{have}\isamarkupfalse%
\ {\isachardoublequoteopen}is{\isacharunderscore}{\kern0pt}htp\ {\isasympi}s\ t\isactrlsub m\isactrlsub a\isactrlsub x{\isachardoublequoteclose}\ \isakeywordTWO{and}\ {\isachardoublequoteopen}{\isasymforall}t\isactrlsub k{\isachardot}{\kern0pt}\ is{\isacharunderscore}{\kern0pt}htp\ {\isasympi}s\ t\isactrlsub k\ {\isasymlongrightarrow}\ t\isactrlsub k\ {\isasymle}\ t\isactrlsub m\isactrlsub a\isactrlsub x{\isachardoublequoteclose}\isanewline
\ \ \ \ \ \ \isakeywordONE{using}\isamarkupfalse%
\ assms\ non{\isacharunderscore}{\kern0pt}Nil{\isacharunderscore}{\kern0pt}sorted{\isacharunderscore}{\kern0pt}last{\isacharunderscore}{\kern0pt}max{\isacharbrackleft}{\kern0pt}OF\ {\isacartoucheopen}htps\ {\isasymnoteq}\ {\isacharbrackleft}{\kern0pt}{\isacharbrackright}{\kern0pt}{\isacartoucheclose}\ {\isacartoucheopen}sorted\ htps{\isacartoucheclose}\ {\isacartoucheopen}t\isactrlsub m\isactrlsub a\isactrlsub x\ {\isacharequal}{\kern0pt}\ last\ htps{\isacartoucheclose}{\isacharbrackright}{\kern0pt}\ \isanewline
\ \ \ \ \ \ \isakeywordONE{by}\isamarkupfalse%
\ {\isacharparenleft}{\kern0pt}auto\ simp{\isacharcolon}{\kern0pt}\ htps{\isacharunderscore}{\kern0pt}seq{\isacharunderscore}{\kern0pt}def{\isacharparenright}{\kern0pt}\isanewline
\isanewline
\ \ \ \ \isakeywordONE{have}\isamarkupfalse%
\ {\isachardoublequoteopen}{\isasymforall}{\isacharparenleft}{\kern0pt}t\isactrlsub a{\isacharcomma}{\kern0pt}A{\isacharparenright}{\kern0pt}{\isasymin}set\ hs{\isachardot}{\kern0pt}\ t\isactrlsub m\isactrlsub i\isactrlsub n\ {\isasymle}\ t\isactrlsub a\ {\isasymand}\ t\isactrlsub a\ {\isasymle}\ t\isactrlsub m\isactrlsub a\isactrlsub x{\isachardoublequoteclose}\isanewline
\ \ \ \ \ \ \isakeywordONE{using}\isamarkupfalse%
\ assms\ {\isacartoucheopen}{\isasymforall}t\isactrlsub k{\isachardot}{\kern0pt}\ is{\isacharunderscore}{\kern0pt}htp\ {\isasympi}s\ t\isactrlsub k\ {\isasymlongrightarrow}\ t\isactrlsub k\ {\isasymle}\ t\isactrlsub m\isactrlsub a\isactrlsub x{\isacartoucheclose}\ {\isacartoucheopen}{\isasymforall}t\isactrlsub k{\isachardot}{\kern0pt}\ is{\isacharunderscore}{\kern0pt}htp\ {\isasympi}s\ t\isactrlsub k\ {\isasymlongrightarrow}\ t\isactrlsub m\isactrlsub i\isactrlsub n\ {\isasymle}\ t\isactrlsub k{\isacartoucheclose}\ \isanewline
\ \ \ \ \ \ \ \ \ \ \ \ ind{\isacharunderscore}{\kern0pt}happ{\isacharunderscore}{\kern0pt}seq{\isacharunderscore}{\kern0pt}htps{\isacharunderscore}{\kern0pt}leq{\isacharunderscore}{\kern0pt}t\isactrlsub m\isactrlsub i\isactrlsub n{\isacharunderscore}{\kern0pt}t\isactrlsub m\isactrlsub a\isactrlsub x{\isacharbrackleft}{\kern0pt}\isakeywordTWO{where}\ hs{\isacharprime}{\kern0pt}{\isacharequal}{\kern0pt}{\isachardoublequoteopen}{\isacharbrackleft}{\kern0pt}{\isacharbrackright}{\kern0pt}{\isachardoublequoteclose}{\isacharbrackright}{\kern0pt}\isanewline
\ \ \ \ \ \ \isakeywordONE{by}\isamarkupfalse%
\ auto\isanewline
\ \ \ \ \isanewline
\ \ \ \ \isakeywordONE{have}\isamarkupfalse%
\ {\isachardoublequoteopen}takeWhile\ {\isacharparenleft}{\kern0pt}leq\ t\isactrlsub m\isactrlsub a\isactrlsub x{\isacharparenright}{\kern0pt}\ hs\ {\isacharequal}{\kern0pt}\ hs{\isachardoublequoteclose}\isanewline
\ \ \ \ \ \ \isakeywordONE{using}\isamarkupfalse%
\ {\isacartoucheopen}{\isasymforall}{\isacharparenleft}{\kern0pt}t\isactrlsub a{\isacharcomma}{\kern0pt}A{\isacharparenright}{\kern0pt}{\isasymin}set\ hs{\isachardot}{\kern0pt}\ t\isactrlsub m\isactrlsub i\isactrlsub n\ {\isasymle}\ t\isactrlsub a\ {\isasymand}\ t\isactrlsub a\ {\isasymle}\ t\isactrlsub m\isactrlsub a\isactrlsub x{\isacartoucheclose}\ \isakeywordONE{by}\isamarkupfalse%
\ {\isacharparenleft}{\kern0pt}auto\ simp{\isacharcolon}{\kern0pt}\ leq{\isacharunderscore}{\kern0pt}def{\isacharparenright}{\kern0pt}\isanewline
\ \ \ \ \isakeywordONE{have}\isamarkupfalse%
\ {\isachardoublequoteopen}dropWhile\ {\isacharparenleft}{\kern0pt}le\ t\isactrlsub m\isactrlsub i\isactrlsub n{\isacharparenright}{\kern0pt}\ {\isacharparenleft}{\kern0pt}takeWhile\ {\isacharparenleft}{\kern0pt}leq\ t\isactrlsub m\isactrlsub i\isactrlsub n{\isacharparenright}{\kern0pt}\ hs{\isacharparenright}{\kern0pt}\ {\isacharequal}{\kern0pt}\ takeWhile\ {\isacharparenleft}{\kern0pt}leq\ t\isactrlsub m\isactrlsub i\isactrlsub n{\isacharparenright}{\kern0pt}\ hs{\isachardoublequoteclose}\isanewline
\ \ \ \ \ \ \isakeywordONE{using}\isamarkupfalse%
\ {\isacartoucheopen}{\isasymforall}{\isacharparenleft}{\kern0pt}t\isactrlsub a{\isacharcomma}{\kern0pt}A{\isacharparenright}{\kern0pt}{\isasymin}set\ hs{\isachardot}{\kern0pt}\ t\isactrlsub m\isactrlsub i\isactrlsub n\ {\isasymle}\ t\isactrlsub a\ {\isasymand}\ t\isactrlsub a\ {\isasymle}\ t\isactrlsub m\isactrlsub a\isactrlsub x{\isacartoucheclose}\ {\isacartoucheopen}strict{\isacharunderscore}{\kern0pt}sorted\ {\isacharparenleft}{\kern0pt}map\ fst\ hs{\isacharparenright}{\kern0pt}{\isacartoucheclose}\isanewline
\ \ \ \ \ \ \isakeywordONE{by}\isamarkupfalse%
\ {\isacharparenleft}{\kern0pt}induction\ hs{\isacharparenright}{\kern0pt}\ {\isacharparenleft}{\kern0pt}auto\ simp{\isacharcolon}{\kern0pt}\ le{\isacharunderscore}{\kern0pt}def{\isacharparenright}{\kern0pt}\isanewline
\isanewline
\ \ \ \ \isakeywordTHREE{obtain}\isamarkupfalse%
\ as\isactrlsub m\isactrlsub i\isactrlsub n\ \isakeywordTWO{where}\ {\isachardoublequoteopen}dropWhile\ {\isacharparenleft}{\kern0pt}le\ t\isactrlsub m\isactrlsub i\isactrlsub n{\isacharparenright}{\kern0pt}\ {\isacharparenleft}{\kern0pt}takeWhile\ {\isacharparenleft}{\kern0pt}leq\ t\isactrlsub m\isactrlsub i\isactrlsub n{\isacharparenright}{\kern0pt}\ hs{\isacharparenright}{\kern0pt}\ {\isacharequal}{\kern0pt}\ {\isacharbrackleft}{\kern0pt}{\isacharparenleft}{\kern0pt}t\isactrlsub m\isactrlsub i\isactrlsub n{\isacharcomma}{\kern0pt}as\isactrlsub m\isactrlsub i\isactrlsub n{\isacharparenright}{\kern0pt}{\isacharbrackright}{\kern0pt}{\isachardoublequoteclose}\isanewline
\ \ \ \ \ \ \isakeywordONE{using}\isamarkupfalse%
\ assms\ ind{\isacharunderscore}{\kern0pt}happ{\isacharunderscore}{\kern0pt}seq{\isacharunderscore}{\kern0pt}le\isactrlsub i{\isacharunderscore}{\kern0pt}leq\isactrlsub i{\isacharbrackleft}{\kern0pt}OF\ {\isacartoucheopen}is{\isacharunderscore}{\kern0pt}htp\ {\isasympi}s\ t\isactrlsub m\isactrlsub i\isactrlsub n{\isacartoucheclose}{\isacharbrackright}{\kern0pt}\ \isakeywordONE{by}\isamarkupfalse%
\ auto\isanewline
\ \ \ \ \isakeywordONE{then}\isamarkupfalse%
\ \isakeywordONE{have}\isamarkupfalse%
\ {\isachardoublequoteopen}{\isacharparenleft}{\kern0pt}t\isactrlsub m\isactrlsub i\isactrlsub n{\isacharcomma}{\kern0pt}as\isactrlsub m\isactrlsub i\isactrlsub n{\isacharparenright}{\kern0pt}\ {\isasymin}\ set\ hs{\isachardoublequoteclose}\isanewline
\ \ \ \ \ \ \isakeywordONE{by}\isamarkupfalse%
\ {\isacharparenleft}{\kern0pt}metis\ list{\isachardot}{\kern0pt}set{\isacharunderscore}{\kern0pt}intros{\isacharparenleft}{\kern0pt}{\isadigit{1}}{\isacharparenright}{\kern0pt}\ set{\isacharunderscore}{\kern0pt}dropWhileD\ set{\isacharunderscore}{\kern0pt}takeWhileD{\isacharparenright}{\kern0pt}\isanewline
\ \ \ \ \isakeywordONE{then}\isamarkupfalse%
\ \isakeywordONE{have}\isamarkupfalse%
\ {\isachardoublequoteopen}set\ as\isactrlsub m\isactrlsub i\isactrlsub n\ {\isacharequal}{\kern0pt}\ acts{\isacharunderscore}{\kern0pt}of{\isacharunderscore}{\kern0pt}plan{\isacharunderscore}{\kern0pt}at\ t\isactrlsub m\isactrlsub i\isactrlsub n\ {\isasympi}s{\isachardoublequoteclose}\isanewline
\ \ \ \ \ \ \isakeywordONE{using}\isamarkupfalse%
\ assms\ {\isacartoucheopen}is{\isacharunderscore}{\kern0pt}htp\ {\isasympi}s\ t\isactrlsub m\isactrlsub i\isactrlsub n{\isacartoucheclose}\ ind{\isacharunderscore}{\kern0pt}happ{\isacharunderscore}{\kern0pt}seq{\isacharunderscore}{\kern0pt}htp{\isacharunderscore}{\kern0pt}acts{\isacharunderscore}{\kern0pt}of{\isacharunderscore}{\kern0pt}plan{\isacharunderscore}{\kern0pt}at\ \isakeywordONE{by}\isamarkupfalse%
\ auto\isanewline
\ \ \ \ \isakeywordONE{let}\isamarkupfalse%
\ {\isacharquery}{\kern0pt}as{\isacharequal}{\kern0pt}{\isachardoublequoteopen}acts{\isacharunderscore}{\kern0pt}of{\isacharunderscore}{\kern0pt}plan{\isacharunderscore}{\kern0pt}at\ t\isactrlsub m\isactrlsub i\isactrlsub n\ {\isasympi}s{\isachardoublequoteclose}\isanewline
\isanewline
\ \ \ \ \isakeywordONE{have}\isamarkupfalse%
\ {\isachardoublequoteopen}t\isactrlsub m\isactrlsub i\isactrlsub n\ {\isasymle}\ t\isactrlsub m\isactrlsub a\isactrlsub x{\isachardoublequoteclose}\isanewline
\ \ \ \ \ \ \isakeywordONE{using}\isamarkupfalse%
\ {\isacartoucheopen}is{\isacharunderscore}{\kern0pt}htp\ {\isasympi}s\ t\isactrlsub m\isactrlsub i\isactrlsub n{\isacartoucheclose}\ \isakeywordONE{by}\isamarkupfalse%
\ {\isacharparenleft}{\kern0pt}simp\ add{\isacharcolon}{\kern0pt}\ {\isacartoucheopen}{\isasymforall}t\isactrlsub k{\isachardot}{\kern0pt}\ is{\isacharunderscore}{\kern0pt}htp\ {\isasympi}s\ t\isactrlsub k\ {\isasymlongrightarrow}\ t\isactrlsub k\ {\isasymle}\ t\isactrlsub m\isactrlsub a\isactrlsub x{\isacartoucheclose}{\isacharparenright}{\kern0pt}\isanewline
\isanewline
\ \ \ \ \isakeywordTHREE{show}\isamarkupfalse%
\ {\isachardoublequoteopen}valid{\isacharunderscore}{\kern0pt}happ{\isacharunderscore}{\kern0pt}seq\ M\isactrlsub {\isadigit{0}}\ hs\ M\isactrlsub n\ {\isasymlongleftrightarrow}\ valid{\isacharunderscore}{\kern0pt}state{\isacharunderscore}{\kern0pt}seq\ M\isactrlsub {\isadigit{0}}\ htps\ {\isasympi}s\ M\isactrlsub n{\isachardoublequoteclose}\isanewline
\ \ \ \ \isakeywordONE{proof}\isamarkupfalse%
\isanewline
\ \ \ \ \ \ \isakeywordTHREE{assume}\isamarkupfalse%
\ {\isachardoublequoteopen}valid{\isacharunderscore}{\kern0pt}happ{\isacharunderscore}{\kern0pt}seq\ M\isactrlsub {\isadigit{0}}\ hs\ M\isactrlsub n{\isachardoublequoteclose}%
\begin{isamarkuptext}%
Remove first happening from the happening sequence.%
\end{isamarkuptext}\isamarkuptrue%
\ \ \ \ \ \ \isakeywordONE{then}\isamarkupfalse%
\ \isakeywordONE{have}\isamarkupfalse%
\ {\isachardoublequoteopen}{\isasymexists}M\isactrlsub {\isadigit{1}}{\isachardot}{\kern0pt}\ valid{\isacharunderscore}{\kern0pt}happ{\isacharunderscore}{\kern0pt}seq\ M\isactrlsub {\isadigit{0}}\ {\isacharparenleft}{\kern0pt}takeWhile\ {\isacharparenleft}{\kern0pt}leq\ t\isactrlsub m\isactrlsub i\isactrlsub n{\isacharparenright}{\kern0pt}\ hs{\isacharparenright}{\kern0pt}\ M\isactrlsub {\isadigit{1}}\ \isanewline
\ \ \ \ \ \ \ \ \ \ \ \ \ \ \ \ \ \ \ \ {\isasymand}\ valid{\isacharunderscore}{\kern0pt}happ{\isacharunderscore}{\kern0pt}seq\ M\isactrlsub {\isadigit{1}}\ {\isacharparenleft}{\kern0pt}dropWhile\ {\isacharparenleft}{\kern0pt}leq\ t\isactrlsub m\isactrlsub i\isactrlsub n{\isacharparenright}{\kern0pt}\ hs{\isacharparenright}{\kern0pt}\ M\isactrlsub n{\isachardoublequoteclose}\isanewline
\ \ \ \ \ \ \ \ \isakeywordONE{using}\isamarkupfalse%
\ valid{\isacharunderscore}{\kern0pt}happ{\isacharunderscore}{\kern0pt}seq{\isacharunderscore}{\kern0pt}app{\isacharunderscore}{\kern0pt}iff{\isacharbrackleft}{\kern0pt}\isakeywordTWO{where}\ hs\isactrlsub {\isadigit{1}}{\isacharequal}{\kern0pt}{\isachardoublequoteopen}takeWhile\ {\isacharparenleft}{\kern0pt}leq\ t\isactrlsub m\isactrlsub i\isactrlsub n{\isacharparenright}{\kern0pt}\ hs{\isachardoublequoteclose}\ \isanewline
\ \ \ \ \ \ \ \ \ \ \ \ \ \ \ \ \ \ \ \ \ \ \ \ \ \ \ \ \ \ \ \ \ \ \ \ \ \ \ \isakeywordTWO{and}\ hs\isactrlsub {\isadigit{2}}{\isacharequal}{\kern0pt}{\isachardoublequoteopen}dropWhile\ {\isacharparenleft}{\kern0pt}leq\ t\isactrlsub m\isactrlsub i\isactrlsub n{\isacharparenright}{\kern0pt}\ hs{\isachardoublequoteclose}{\isacharbrackright}{\kern0pt}\ \isanewline
\ \ \ \ \ \ \ \ \isakeywordONE{by}\isamarkupfalse%
\ auto\isanewline
\ \ \ \ \ \ \isakeywordONE{then}\isamarkupfalse%
\ \isakeywordTHREE{obtain}\isamarkupfalse%
\ M\isactrlsub {\isadigit{1}}\ \isakeywordTWO{where}\ {\isachardoublequoteopen}valid{\isacharunderscore}{\kern0pt}happ{\isacharunderscore}{\kern0pt}seq\ M\isactrlsub {\isadigit{0}}\ {\isacharparenleft}{\kern0pt}dropWhile\ {\isacharparenleft}{\kern0pt}le\ t\isactrlsub m\isactrlsub i\isactrlsub n{\isacharparenright}{\kern0pt}\ {\isacharparenleft}{\kern0pt}takeWhile\ {\isacharparenleft}{\kern0pt}leq\ t\isactrlsub m\isactrlsub i\isactrlsub n{\isacharparenright}{\kern0pt}\ hs{\isacharparenright}{\kern0pt}{\isacharparenright}{\kern0pt}\ M\isactrlsub {\isadigit{1}}{\isachardoublequoteclose}\isanewline
\ \ \ \ \ \ \ \ \ \ \ \ \ \ \ \ \ \ \ \ \ \ \ \isakeywordTWO{and}\ {\isachardoublequoteopen}valid{\isacharunderscore}{\kern0pt}happ{\isacharunderscore}{\kern0pt}seq\ M\isactrlsub {\isadigit{1}}\ {\isacharparenleft}{\kern0pt}dropWhile\ {\isacharparenleft}{\kern0pt}leq\ t\isactrlsub m\isactrlsub i\isactrlsub n{\isacharparenright}{\kern0pt}\ {\isacharparenleft}{\kern0pt}takeWhile\ {\isacharparenleft}{\kern0pt}leq\ t\isactrlsub m\isactrlsub a\isactrlsub x{\isacharparenright}{\kern0pt}\ hs{\isacharparenright}{\kern0pt}{\isacharparenright}{\kern0pt}\ M\isactrlsub n{\isachardoublequoteclose}\isanewline
\ \ \ \ \ \ \ \ \isakeywordONE{using}\isamarkupfalse%
\ {\isacartoucheopen}dropWhile\ {\isacharparenleft}{\kern0pt}le\ t\isactrlsub m\isactrlsub i\isactrlsub n{\isacharparenright}{\kern0pt}\ {\isacharparenleft}{\kern0pt}takeWhile\ {\isacharparenleft}{\kern0pt}leq\ t\isactrlsub m\isactrlsub i\isactrlsub n{\isacharparenright}{\kern0pt}\ hs{\isacharparenright}{\kern0pt}\ {\isacharequal}{\kern0pt}\ takeWhile\ {\isacharparenleft}{\kern0pt}leq\ t\isactrlsub m\isactrlsub i\isactrlsub n{\isacharparenright}{\kern0pt}\ hs{\isacartoucheclose}\ \isanewline
\ \ \ \ \ \ \ \ \ \ \ \ \ \ {\isacartoucheopen}takeWhile\ {\isacharparenleft}{\kern0pt}leq\ t\isactrlsub m\isactrlsub a\isactrlsub x{\isacharparenright}{\kern0pt}\ hs\ {\isacharequal}{\kern0pt}\ hs{\isacartoucheclose}{\isacharbrackleft}{\kern0pt}symmetric{\isacharbrackright}{\kern0pt}\ \isanewline
\ \ \ \ \ \ \ \ \isakeywordONE{by}\isamarkupfalse%
\ auto\isanewline
\isanewline
\ \ \ \ \ \ \isakeywordONE{have}\isamarkupfalse%
\ {\isachardoublequoteopen}invs{\isacharunderscore}{\kern0pt}of{\isacharunderscore}{\kern0pt}plan{\isacharunderscore}{\kern0pt}at\ t\isactrlsub m\isactrlsub i\isactrlsub n\ {\isasympi}s\ {\isacharequal}{\kern0pt}\ {\isacharbraceleft}{\kern0pt}{\isacharbraceright}{\kern0pt}{\isachardoublequoteclose}\isanewline
\ \ \ \ \ \ \ \ \isakeywordONE{unfolding}\isamarkupfalse%
\ invs{\isacharunderscore}{\kern0pt}of{\isacharunderscore}{\kern0pt}plan{\isacharunderscore}{\kern0pt}at{\isacharunderscore}{\kern0pt}def\ \isakeywordONE{using}\isamarkupfalse%
\ {\isacartoucheopen}{\isasymforall}{\isacharparenleft}{\kern0pt}t\isactrlsub {\isasympi}{\isacharcomma}{\kern0pt}{\isasympi}{\isacharparenright}{\kern0pt}\ {\isasymin}\ set\ {\isasympi}s{\isachardot}{\kern0pt}\ t\isactrlsub m\isactrlsub i\isactrlsub n\ {\isasymle}\ t\isactrlsub {\isasympi}{\isacartoucheclose}\ \isanewline
\ \ \ \ \ \ \ \ \isakeywordONE{by}\isamarkupfalse%
\ {\isacharparenleft}{\kern0pt}auto\ simp{\isacharcolon}{\kern0pt}\ durative{\isacharunderscore}{\kern0pt}acts{\isacharunderscore}{\kern0pt}def{\isacharparenright}{\kern0pt}\isanewline
\ \ \ \ \ \ \isakeywordONE{then}\isamarkupfalse%
\ \isakeywordONE{have}\isamarkupfalse%
\ {\isachardoublequoteopen}valid{\isacharunderscore}{\kern0pt}state{\isacharunderscore}{\kern0pt}seq\ M\isactrlsub {\isadigit{0}}\ {\isacharbrackleft}{\kern0pt}t\isactrlsub m\isactrlsub i\isactrlsub n{\isacharbrackright}{\kern0pt}\ {\isasympi}s\ M\isactrlsub {\isadigit{1}}{\isachardoublequoteclose}\isanewline
\ \ \ \ \ \ \ \ \isakeywordONE{using}\isamarkupfalse%
\ {\isacartoucheopen}valid{\isacharunderscore}{\kern0pt}happ{\isacharunderscore}{\kern0pt}seq\ M\isactrlsub {\isadigit{0}}\ {\isacharparenleft}{\kern0pt}dropWhile\ {\isacharparenleft}{\kern0pt}le\ t\isactrlsub m\isactrlsub i\isactrlsub n{\isacharparenright}{\kern0pt}\ {\isacharparenleft}{\kern0pt}takeWhile\ {\isacharparenleft}{\kern0pt}leq\ t\isactrlsub m\isactrlsub i\isactrlsub n{\isacharparenright}{\kern0pt}\ hs{\isacharparenright}{\kern0pt}{\isacharparenright}{\kern0pt}\ M\isactrlsub {\isadigit{1}}{\isacartoucheclose}\isanewline
\ \ \ \ \ \ \ \ \ \ \ \ \ \ {\isacartoucheopen}dropWhile\ {\isacharparenleft}{\kern0pt}le\ t\isactrlsub m\isactrlsub i\isactrlsub n{\isacharparenright}{\kern0pt}\ {\isacharparenleft}{\kern0pt}takeWhile\ {\isacharparenleft}{\kern0pt}leq\ t\isactrlsub m\isactrlsub i\isactrlsub n{\isacharparenright}{\kern0pt}\ hs{\isacharparenright}{\kern0pt}\ {\isacharequal}{\kern0pt}\ {\isacharbrackleft}{\kern0pt}{\isacharparenleft}{\kern0pt}t\isactrlsub m\isactrlsub i\isactrlsub n{\isacharcomma}{\kern0pt}as\isactrlsub m\isactrlsub i\isactrlsub n{\isacharparenright}{\kern0pt}{\isacharbrackright}{\kern0pt}{\isacartoucheclose}\ \isanewline
\ \ \ \ \ \ \ \ \ \ \ \ \ \ {\isacartoucheopen}set\ as\isactrlsub m\isactrlsub i\isactrlsub n\ {\isacharequal}{\kern0pt}\ acts{\isacharunderscore}{\kern0pt}of{\isacharunderscore}{\kern0pt}plan{\isacharunderscore}{\kern0pt}at\ t\isactrlsub m\isactrlsub i\isactrlsub n\ {\isasympi}s{\isacartoucheclose}\isanewline
\ \ \ \ \ \ \ \ \isakeywordONE{by}\isamarkupfalse%
\ {\isacharparenleft}{\kern0pt}auto\ simp\ add{\isacharcolon}{\kern0pt}\ pairwise{\isacharunderscore}{\kern0pt}def{\isacharparenright}{\kern0pt}\isanewline
\isanewline
\ \ \ \ \ \ \isakeywordONE{have}\isamarkupfalse%
\ {\isachardoublequoteopen}t\isactrlsub m\isactrlsub i\isactrlsub n\ {\isacharequal}{\kern0pt}\ t\isactrlsub m\isactrlsub a\isactrlsub x\ {\isasymor}\ t\isactrlsub m\isactrlsub i\isactrlsub n\ {\isacharless}{\kern0pt}\ t\isactrlsub m\isactrlsub a\isactrlsub x{\isachardoublequoteclose}\isanewline
\ \ \ \ \ \ \ \ \isakeywordONE{using}\isamarkupfalse%
\ {\isacartoucheopen}t\isactrlsub m\isactrlsub i\isactrlsub n\ {\isasymle}\ t\isactrlsub m\isactrlsub a\isactrlsub x{\isacartoucheclose}\ \isakeywordONE{by}\isamarkupfalse%
\ auto\isanewline
\ \ \ \ \ \ \isakeywordONE{then}\isamarkupfalse%
\ \isakeywordTHREE{show}\isamarkupfalse%
\ {\isachardoublequoteopen}valid{\isacharunderscore}{\kern0pt}state{\isacharunderscore}{\kern0pt}seq\ M\isactrlsub {\isadigit{0}}\ htps\ {\isasympi}s\ M\isactrlsub n{\isachardoublequoteclose}\isanewline
\ \ \ \ \ \ \isakeywordONE{proof}\isamarkupfalse%
\isanewline
\ \ \ \ \ \ \ \ \isakeywordTHREE{assume}\isamarkupfalse%
\ {\isachardoublequoteopen}t\isactrlsub m\isactrlsub i\isactrlsub n\ {\isacharequal}{\kern0pt}\ t\isactrlsub m\isactrlsub a\isactrlsub x{\isachardoublequoteclose}\isanewline
\ \ \ \ \ \ \ \ \isakeywordONE{then}\isamarkupfalse%
\ \isakeywordONE{have}\isamarkupfalse%
\ {\isachardoublequoteopen}valid{\isacharunderscore}{\kern0pt}happ{\isacharunderscore}{\kern0pt}seq\ M\isactrlsub {\isadigit{0}}\ hs\ M\isactrlsub {\isadigit{1}}{\isachardoublequoteclose}\isanewline
\ \ \ \ \ \ \ \ \ \ \isakeywordONE{using}\isamarkupfalse%
\ {\isacartoucheopen}dropWhile\ {\isacharparenleft}{\kern0pt}le\ t\isactrlsub m\isactrlsub i\isactrlsub n{\isacharparenright}{\kern0pt}\ {\isacharparenleft}{\kern0pt}takeWhile\ {\isacharparenleft}{\kern0pt}leq\ t\isactrlsub m\isactrlsub i\isactrlsub n{\isacharparenright}{\kern0pt}\ hs{\isacharparenright}{\kern0pt}\ {\isacharequal}{\kern0pt}\ takeWhile\ {\isacharparenleft}{\kern0pt}leq\ t\isactrlsub m\isactrlsub i\isactrlsub n{\isacharparenright}{\kern0pt}\ hs{\isacartoucheclose}\ \isanewline
\ \ \ \ \ \ \ \ \ \ \ \ \ \ \ \ {\isacartoucheopen}valid{\isacharunderscore}{\kern0pt}happ{\isacharunderscore}{\kern0pt}seq\ M\isactrlsub {\isadigit{0}}\ {\isacharparenleft}{\kern0pt}dropWhile\ {\isacharparenleft}{\kern0pt}le\ t\isactrlsub m\isactrlsub i\isactrlsub n{\isacharparenright}{\kern0pt}\ {\isacharparenleft}{\kern0pt}takeWhile\ {\isacharparenleft}{\kern0pt}leq\ t\isactrlsub m\isactrlsub i\isactrlsub n{\isacharparenright}{\kern0pt}\ hs{\isacharparenright}{\kern0pt}{\isacharparenright}{\kern0pt}\ M\isactrlsub {\isadigit{1}}{\isacartoucheclose}\isanewline
\ \ \ \ \ \ \ \ \ \ \isakeywordONE{by}\isamarkupfalse%
\ {\isacharparenleft}{\kern0pt}simp\ add{\isacharcolon}{\kern0pt}\ {\isacartoucheopen}takeWhile\ {\isacharparenleft}{\kern0pt}leq\ t\isactrlsub m\isactrlsub a\isactrlsub x{\isacharparenright}{\kern0pt}\ hs\ {\isacharequal}{\kern0pt}\ hs{\isacartoucheclose}{\isacharparenright}{\kern0pt}\isanewline
\ \ \ \ \ \ \ \ \isakeywordONE{then}\isamarkupfalse%
\ \isakeywordONE{have}\isamarkupfalse%
\ {\isachardoublequoteopen}M\isactrlsub {\isadigit{1}}\ {\isacharequal}{\kern0pt}\ M\isactrlsub n{\isachardoublequoteclose}\isanewline
\ \ \ \ \ \ \ \ \ \ \isakeywordONE{using}\isamarkupfalse%
\ valid{\isacharunderscore}{\kern0pt}happ{\isacharunderscore}{\kern0pt}seq{\isacharunderscore}{\kern0pt}state{\isacharunderscore}{\kern0pt}unique{\isacharbrackleft}{\kern0pt}OF\ {\isacartoucheopen}valid{\isacharunderscore}{\kern0pt}happ{\isacharunderscore}{\kern0pt}seq\ M\isactrlsub {\isadigit{0}}\ hs\ M\isactrlsub n{\isacartoucheclose}\ {\isacartoucheopen}valid{\isacharunderscore}{\kern0pt}happ{\isacharunderscore}{\kern0pt}seq\ M\isactrlsub {\isadigit{0}}\ hs\ M\isactrlsub {\isadigit{1}}{\isacartoucheclose}{\isacharbrackright}{\kern0pt}\isanewline
\ \ \ \ \ \ \ \ \ \ \isakeywordONE{by}\isamarkupfalse%
\ auto\isanewline
\ \ \ \ \ \ \ \ \isakeywordONE{have}\isamarkupfalse%
\ {\isachardoublequoteopen}htps\ {\isacharequal}{\kern0pt}\ {\isacharbrackleft}{\kern0pt}t\isactrlsub m\isactrlsub i\isactrlsub n{\isacharbrackright}{\kern0pt}{\isachardoublequoteclose}\isanewline
\ \ \ \ \ \ \ \ \ \ \isakeywordONE{using}\isamarkupfalse%
\ {\isacartoucheopen}htps\ {\isasymnoteq}\ {\isacharbrackleft}{\kern0pt}{\isacharbrackright}{\kern0pt}{\isacartoucheclose}\ {\isacartoucheopen}strict{\isacharunderscore}{\kern0pt}sorted\ htps{\isacartoucheclose}\ {\isacartoucheopen}t\isactrlsub m\isactrlsub i\isactrlsub n\ {\isacharequal}{\kern0pt}\ hd\ htps{\isacartoucheclose}\ {\isacartoucheopen}t\isactrlsub m\isactrlsub a\isactrlsub x\ {\isacharequal}{\kern0pt}\ last\ htps{\isacartoucheclose}\ {\isacartoucheopen}t\isactrlsub m\isactrlsub i\isactrlsub n\ {\isacharequal}{\kern0pt}\ t\isactrlsub m\isactrlsub a\isactrlsub x{\isacartoucheclose}\isanewline
\ \ \ \ \ \ \ \ \ \ \isakeywordONE{by}\isamarkupfalse%
\ {\isacharparenleft}{\kern0pt}induction\ htps{\isacharparenright}{\kern0pt}\ {\isacharparenleft}{\kern0pt}auto\ split{\isacharcolon}{\kern0pt}\ if{\isacharunderscore}{\kern0pt}splits{\isacharparenright}{\kern0pt}\isanewline
\ \ \ \ \ \ \ \ \isakeywordONE{then}\isamarkupfalse%
\ \isakeywordTHREE{show}\isamarkupfalse%
\ {\isacharquery}{\kern0pt}thesis\isanewline
\ \ \ \ \ \ \ \ \ \ \isakeywordONE{using}\isamarkupfalse%
\ {\isacartoucheopen}valid{\isacharunderscore}{\kern0pt}state{\isacharunderscore}{\kern0pt}seq\ M\isactrlsub {\isadigit{0}}\ {\isacharbrackleft}{\kern0pt}t\isactrlsub m\isactrlsub i\isactrlsub n{\isacharbrackright}{\kern0pt}\ {\isasympi}s\ M\isactrlsub {\isadigit{1}}{\isacartoucheclose}\ {\isacartoucheopen}M\isactrlsub {\isadigit{1}}\ {\isacharequal}{\kern0pt}\ M\isactrlsub n{\isacartoucheclose}\ \isakeywordONE{by}\isamarkupfalse%
\ auto\isanewline
\ \ \ \ \ \ \isakeywordONE{next}\isamarkupfalse%
\isanewline
\ \ \ \ \ \ \ \ \isakeywordTHREE{assume}\isamarkupfalse%
\ {\isachardoublequoteopen}t\isactrlsub m\isactrlsub i\isactrlsub n\ {\isacharless}{\kern0pt}\ t\isactrlsub m\isactrlsub a\isactrlsub x{\isachardoublequoteclose}\isanewline
\ \ \ \ \ \ \ \ \isakeywordONE{then}\isamarkupfalse%
\ \isakeywordONE{have}\isamarkupfalse%
\ {\isachardoublequoteopen}{\isasymexists}htps{\isacharprime}{\kern0pt}{\isachardot}{\kern0pt}\ htps\ {\isacharequal}{\kern0pt}\ {\isacharbrackleft}{\kern0pt}{\isacharbrackright}{\kern0pt}\ {\isacharat}{\kern0pt}\ t\isactrlsub m\isactrlsub i\isactrlsub n\ {\isacharhash}{\kern0pt}\ htps{\isacharprime}{\kern0pt}\ {\isacharat}{\kern0pt}\ t\isactrlsub m\isactrlsub a\isactrlsub x\ {\isacharhash}{\kern0pt}\ {\isacharbrackleft}{\kern0pt}{\isacharbrackright}{\kern0pt}{\isachardoublequoteclose}\isanewline
\ \ \ \ \ \ \ \ \ \ \isakeywordONE{using}\isamarkupfalse%
\ {\isacartoucheopen}htps\ {\isasymnoteq}\ {\isacharbrackleft}{\kern0pt}{\isacharbrackright}{\kern0pt}{\isacartoucheclose}\ {\isacartoucheopen}t\isactrlsub m\isactrlsub i\isactrlsub n\ {\isacharequal}{\kern0pt}\ hd\ htps{\isacartoucheclose}\ {\isacartoucheopen}t\isactrlsub m\isactrlsub a\isactrlsub x\ {\isacharequal}{\kern0pt}\ last\ htps{\isacartoucheclose}\isanewline
\ \ \ \ \ \ \ \ \ \ \isakeywordONE{apply}\isamarkupfalse%
\ {\isacharparenleft}{\kern0pt}induction\ htps{\isacharparenright}{\kern0pt}\isanewline
\ \ \ \ \ \ \ \ \ \ \isakeywordONE{apply}\isamarkupfalse%
\ {\isacharparenleft}{\kern0pt}auto\ split{\isacharcolon}{\kern0pt}\ if{\isacharunderscore}{\kern0pt}splits{\isacharparenright}{\kern0pt}\isanewline
\ \ \ \ \ \ \ \ \ \ \isakeywordONE{apply}\isamarkupfalse%
\ {\isacharparenleft}{\kern0pt}metis\ append{\isacharunderscore}{\kern0pt}butlast{\isacharunderscore}{\kern0pt}last{\isacharunderscore}{\kern0pt}id{\isacharparenright}{\kern0pt}\isanewline
\ \ \ \ \ \ \ \ \ \ \isakeywordONE{done}\isamarkupfalse%
\isanewline
\ \ \ \ \ \ \ \ \isakeywordONE{then}\isamarkupfalse%
\ \isakeywordTHREE{obtain}\isamarkupfalse%
\ htps{\isacharprime}{\kern0pt}\ \isakeywordTWO{where}\ {\isachardoublequoteopen}htps\ {\isacharequal}{\kern0pt}\ {\isacharbrackleft}{\kern0pt}{\isacharbrackright}{\kern0pt}\ {\isacharat}{\kern0pt}\ t\isactrlsub m\isactrlsub i\isactrlsub n\ {\isacharhash}{\kern0pt}\ htps{\isacharprime}{\kern0pt}\ {\isacharat}{\kern0pt}\ t\isactrlsub m\isactrlsub a\isactrlsub x\ {\isacharhash}{\kern0pt}\ {\isacharbrackleft}{\kern0pt}{\isacharbrackright}{\kern0pt}{\isachardoublequoteclose}\isanewline
\ \ \ \ \ \ \ \ \ \ \isakeywordONE{by}\isamarkupfalse%
\ auto\isanewline
\ \ \ \ \ \ \ \ \isakeywordONE{then}\isamarkupfalse%
\ \isakeywordONE{have}\isamarkupfalse%
\ {\isachardoublequoteopen}htps{\isacharunderscore}{\kern0pt}seq\ {\isasympi}s\ {\isacharparenleft}{\kern0pt}{\isacharbrackleft}{\kern0pt}{\isacharbrackright}{\kern0pt}\ {\isacharat}{\kern0pt}\ t\isactrlsub m\isactrlsub i\isactrlsub n\ {\isacharhash}{\kern0pt}\ htps{\isacharprime}{\kern0pt}\ {\isacharat}{\kern0pt}\ t\isactrlsub m\isactrlsub a\isactrlsub x\ {\isacharhash}{\kern0pt}\ {\isacharbrackleft}{\kern0pt}{\isacharbrackright}{\kern0pt}{\isacharparenright}{\kern0pt}{\isachardoublequoteclose}\isanewline
\ \ \ \ \ \ \ \ \ \ \isakeywordONE{using}\isamarkupfalse%
\ assms\ \isakeywordONE{by}\isamarkupfalse%
\ auto\isanewline
\ \ \ \ \ \ \ \ \isakeywordONE{then}\isamarkupfalse%
\ \isakeywordONE{have}\isamarkupfalse%
\ {\isachardoublequoteopen}valid{\isacharunderscore}{\kern0pt}state{\isacharunderscore}{\kern0pt}seq\ M\isactrlsub {\isadigit{1}}\ {\isacharparenleft}{\kern0pt}htps{\isacharprime}{\kern0pt}\ {\isacharat}{\kern0pt}\ {\isacharbrackleft}{\kern0pt}t\isactrlsub m\isactrlsub a\isactrlsub x{\isacharbrackright}{\kern0pt}{\isacharparenright}{\kern0pt}\ {\isasympi}s\ M\isactrlsub n{\isachardoublequoteclose}\isanewline
\ \ \ \ \ \ \ \ \ \ \isakeywordONE{using}\isamarkupfalse%
\ assms\ {\isacartoucheopen}valid{\isacharunderscore}{\kern0pt}happ{\isacharunderscore}{\kern0pt}seq\ M\isactrlsub {\isadigit{1}}\ {\isacharparenleft}{\kern0pt}dropWhile\ {\isacharparenleft}{\kern0pt}leq\ t\isactrlsub m\isactrlsub i\isactrlsub n{\isacharparenright}{\kern0pt}\ {\isacharparenleft}{\kern0pt}takeWhile\ {\isacharparenleft}{\kern0pt}leq\ t\isactrlsub m\isactrlsub a\isactrlsub x{\isacharparenright}{\kern0pt}\ hs{\isacharparenright}{\kern0pt}{\isacharparenright}{\kern0pt}\ M\isactrlsub n{\isacartoucheclose}\isanewline
\ \ \ \ \ \ \ \ \ \ \ \ \ \ \ \ valid{\isacharunderscore}{\kern0pt}happ{\isacharunderscore}{\kern0pt}seq{\isacharunderscore}{\kern0pt}valid{\isacharunderscore}{\kern0pt}state{\isacharunderscore}{\kern0pt}seq{\isacharunderscore}{\kern0pt}i{\isacharunderscore}{\kern0pt}to{\isacharunderscore}{\kern0pt}j{\isacharunderscore}{\kern0pt}aux\isanewline
\ \ \ \ \ \ \ \ \ \ \ \ \ \ \ \ \ \ {\isacharbrackleft}{\kern0pt}OF\ {\isacartoucheopen}htps{\isacharunderscore}{\kern0pt}seq\ {\isasympi}s\ {\isacharparenleft}{\kern0pt}{\isacharbrackleft}{\kern0pt}{\isacharbrackright}{\kern0pt}\ {\isacharat}{\kern0pt}\ t\isactrlsub m\isactrlsub i\isactrlsub n\ {\isacharhash}{\kern0pt}\ htps{\isacharprime}{\kern0pt}\ {\isacharat}{\kern0pt}\ t\isactrlsub m\isactrlsub a\isactrlsub x\ {\isacharhash}{\kern0pt}\ {\isacharbrackleft}{\kern0pt}{\isacharbrackright}{\kern0pt}{\isacharparenright}{\kern0pt}{\isacartoucheclose}{\isacharbrackright}{\kern0pt}\isanewline
\ \ \ \ \ \ \ \ \ \ \isakeywordONE{by}\isamarkupfalse%
\ blast\isanewline
\ \ \ \ \ \ \ \ \isakeywordONE{then}\isamarkupfalse%
\ \isakeywordTHREE{show}\isamarkupfalse%
\ {\isacharquery}{\kern0pt}thesis\isanewline
\ \ \ \ \ \ \ \ \ \ \isakeywordONE{using}\isamarkupfalse%
\ {\isacartoucheopen}valid{\isacharunderscore}{\kern0pt}state{\isacharunderscore}{\kern0pt}seq\ M\isactrlsub {\isadigit{0}}\ {\isacharbrackleft}{\kern0pt}t\isactrlsub m\isactrlsub i\isactrlsub n{\isacharbrackright}{\kern0pt}\ {\isasympi}s\ M\isactrlsub {\isadigit{1}}{\isacartoucheclose}\ {\isacartoucheopen}htps\ {\isacharequal}{\kern0pt}\ {\isacharbrackleft}{\kern0pt}{\isacharbrackright}{\kern0pt}\ {\isacharat}{\kern0pt}\ t\isactrlsub m\isactrlsub i\isactrlsub n\ {\isacharhash}{\kern0pt}\ htps{\isacharprime}{\kern0pt}\ {\isacharat}{\kern0pt}\ t\isactrlsub m\isactrlsub a\isactrlsub x\ {\isacharhash}{\kern0pt}\ {\isacharbrackleft}{\kern0pt}{\isacharbrackright}{\kern0pt}{\isacartoucheclose}\isanewline
\ \ \ \ \ \ \ \ \ \ \ \ \ \ \ valid{\isacharunderscore}{\kern0pt}state{\isacharunderscore}{\kern0pt}seq{\isacharunderscore}{\kern0pt}app{\isacharunderscore}{\kern0pt}iff{\isacharbrackleft}{\kern0pt}\isakeywordTWO{where}\ ts\isactrlsub {\isadigit{1}}{\isacharequal}{\kern0pt}{\isachardoublequoteopen}{\isacharbrackleft}{\kern0pt}t\isactrlsub m\isactrlsub i\isactrlsub n{\isacharbrackright}{\kern0pt}{\isachardoublequoteclose}\ \isakeywordTWO{and}\ ts\isactrlsub {\isadigit{2}}{\isacharequal}{\kern0pt}{\isachardoublequoteopen}htps{\isacharprime}{\kern0pt}\ {\isacharat}{\kern0pt}\ {\isacharbrackleft}{\kern0pt}t\isactrlsub m\isactrlsub a\isactrlsub x{\isacharbrackright}{\kern0pt}{\isachardoublequoteclose}{\isacharbrackright}{\kern0pt}\isanewline
\ \ \ \ \ \ \ \ \ \ \isakeywordONE{by}\isamarkupfalse%
\ auto\isanewline
\ \ \ \ \ \ \isakeywordONE{qed}\isamarkupfalse%
\isanewline
\ \ \ \ \isakeywordONE{next}\isamarkupfalse%
\isanewline
\ \ \ \ \ \ \isakeywordTHREE{assume}\isamarkupfalse%
\ {\isachardoublequoteopen}valid{\isacharunderscore}{\kern0pt}state{\isacharunderscore}{\kern0pt}seq\ M\isactrlsub {\isadigit{0}}\ htps\ {\isasympi}s\ M\isactrlsub n{\isachardoublequoteclose}%
\begin{isamarkuptext}%
Remove first happening from the happening sequence.%
\end{isamarkuptext}\isamarkuptrue%
\ \ \ \ \ \ \isakeywordONE{then}\isamarkupfalse%
\ \isakeywordTHREE{obtain}\isamarkupfalse%
\ M\isactrlsub {\isadigit{1}}\ \isakeywordTWO{where}\ {\isachardoublequoteopen}valid{\isacharunderscore}{\kern0pt}state{\isacharunderscore}{\kern0pt}seq\ M\isactrlsub {\isadigit{0}}\ {\isacharbrackleft}{\kern0pt}t\isactrlsub m\isactrlsub i\isactrlsub n{\isacharbrackright}{\kern0pt}\ {\isasympi}s\ M\isactrlsub {\isadigit{1}}{\isachardoublequoteclose}\ \isakeywordTWO{and}\ {\isachardoublequoteopen}valid{\isacharunderscore}{\kern0pt}state{\isacharunderscore}{\kern0pt}seq\ M\isactrlsub {\isadigit{1}}\ {\isacharparenleft}{\kern0pt}tl\ htps{\isacharparenright}{\kern0pt}\ {\isasympi}s\ M\isactrlsub n{\isachardoublequoteclose}\isanewline
\ \ \ \ \ \ \ \ \isakeywordONE{using}\isamarkupfalse%
\ {\isacartoucheopen}t\isactrlsub m\isactrlsub i\isactrlsub n\ {\isacharequal}{\kern0pt}\ hd\ htps{\isacartoucheclose}\ valid{\isacharunderscore}{\kern0pt}state{\isacharunderscore}{\kern0pt}seq{\isacharunderscore}{\kern0pt}app{\isacharunderscore}{\kern0pt}iff{\isacharbrackleft}{\kern0pt}\isakeywordTWO{where}\ ts\isactrlsub {\isadigit{1}}{\isacharequal}{\kern0pt}{\isachardoublequoteopen}{\isacharbrackleft}{\kern0pt}t\isactrlsub m\isactrlsub i\isactrlsub n{\isacharbrackright}{\kern0pt}{\isachardoublequoteclose}\ \isakeywordTWO{and}\ ts\isactrlsub {\isadigit{2}}{\isacharequal}{\kern0pt}{\isachardoublequoteopen}tl\ htps{\isachardoublequoteclose}{\isacharbrackright}{\kern0pt}\isanewline
\ \ \ \ \ \ \ \ \isakeywordONE{by}\isamarkupfalse%
\ {\isacharparenleft}{\kern0pt}auto\ simp{\isacharcolon}{\kern0pt}\ {\isacartoucheopen}htps\ {\isasymnoteq}\ {\isacharbrackleft}{\kern0pt}{\isacharbrackright}{\kern0pt}{\isacartoucheclose}{\isacharparenright}{\kern0pt}\isanewline
\ \ \ \ \ \ \isakeywordONE{then}\isamarkupfalse%
\ \isakeywordONE{have}\isamarkupfalse%
\ {\isachardoublequoteopen}valid{\isacharunderscore}{\kern0pt}happ{\isacharunderscore}{\kern0pt}seq\ M\isactrlsub {\isadigit{0}}\ {\isacharparenleft}{\kern0pt}dropWhile\ {\isacharparenleft}{\kern0pt}le\ t\isactrlsub m\isactrlsub i\isactrlsub n{\isacharparenright}{\kern0pt}\ {\isacharparenleft}{\kern0pt}takeWhile\ {\isacharparenleft}{\kern0pt}leq\ t\isactrlsub m\isactrlsub i\isactrlsub n{\isacharparenright}{\kern0pt}\ hs{\isacharparenright}{\kern0pt}{\isacharparenright}{\kern0pt}\ M\isactrlsub {\isadigit{1}}{\isachardoublequoteclose}\isanewline
\ \ \ \ \ \ \ \ \isakeywordONE{using}\isamarkupfalse%
\ {\isacartoucheopen}valid{\isacharunderscore}{\kern0pt}state{\isacharunderscore}{\kern0pt}seq\ M\isactrlsub {\isadigit{0}}\ {\isacharbrackleft}{\kern0pt}t\isactrlsub m\isactrlsub i\isactrlsub n{\isacharbrackright}{\kern0pt}\ {\isasympi}s\ M\isactrlsub {\isadigit{1}}{\isacartoucheclose}\ \isanewline
\ \ \ \ \ \ \ \ \ \ \ \ \ \ {\isacartoucheopen}dropWhile\ {\isacharparenleft}{\kern0pt}le\ t\isactrlsub m\isactrlsub i\isactrlsub n{\isacharparenright}{\kern0pt}\ {\isacharparenleft}{\kern0pt}takeWhile\ {\isacharparenleft}{\kern0pt}leq\ t\isactrlsub m\isactrlsub i\isactrlsub n{\isacharparenright}{\kern0pt}\ hs{\isacharparenright}{\kern0pt}\ {\isacharequal}{\kern0pt}\ {\isacharbrackleft}{\kern0pt}{\isacharparenleft}{\kern0pt}t\isactrlsub m\isactrlsub i\isactrlsub n{\isacharcomma}{\kern0pt}as\isactrlsub m\isactrlsub i\isactrlsub n{\isacharparenright}{\kern0pt}{\isacharbrackright}{\kern0pt}{\isacartoucheclose}\ \isanewline
\ \ \ \ \ \ \ \ \ \ \ \ \ \ {\isacartoucheopen}set\ as\isactrlsub m\isactrlsub i\isactrlsub n\ {\isacharequal}{\kern0pt}\ acts{\isacharunderscore}{\kern0pt}of{\isacharunderscore}{\kern0pt}plan{\isacharunderscore}{\kern0pt}at\ t\isactrlsub m\isactrlsub i\isactrlsub n\ {\isasympi}s{\isacartoucheclose}\isanewline
\ \ \ \ \ \ \ \ \isakeywordONE{by}\isamarkupfalse%
\ {\isacharparenleft}{\kern0pt}auto\ simp{\isacharcolon}{\kern0pt}\ Let{\isacharunderscore}{\kern0pt}def\ pairwise{\isacharunderscore}{\kern0pt}def{\isacharparenright}{\kern0pt}\isanewline
\isanewline
\ \ \ \ \ \ \isakeywordONE{have}\isamarkupfalse%
\ {\isachardoublequoteopen}t\isactrlsub m\isactrlsub i\isactrlsub n\ {\isacharequal}{\kern0pt}\ t\isactrlsub m\isactrlsub a\isactrlsub x\ {\isasymor}\ t\isactrlsub m\isactrlsub i\isactrlsub n\ {\isacharless}{\kern0pt}\ t\isactrlsub m\isactrlsub a\isactrlsub x{\isachardoublequoteclose}\isanewline
\ \ \ \ \ \ \ \ \isakeywordONE{using}\isamarkupfalse%
\ {\isacartoucheopen}t\isactrlsub m\isactrlsub i\isactrlsub n\ {\isasymle}\ t\isactrlsub m\isactrlsub a\isactrlsub x{\isacartoucheclose}\ \isakeywordONE{by}\isamarkupfalse%
\ auto\isanewline
\ \ \ \ \ \ \isakeywordONE{then}\isamarkupfalse%
\ \isakeywordTHREE{show}\isamarkupfalse%
\ {\isachardoublequoteopen}valid{\isacharunderscore}{\kern0pt}happ{\isacharunderscore}{\kern0pt}seq\ M\isactrlsub {\isadigit{0}}\ hs\ M\isactrlsub n{\isachardoublequoteclose}\isanewline
\ \ \ \ \ \ \isakeywordONE{proof}\isamarkupfalse%
\isanewline
\ \ \ \ \ \ \ \ \isakeywordTHREE{assume}\isamarkupfalse%
\ {\isachardoublequoteopen}t\isactrlsub m\isactrlsub i\isactrlsub n\ {\isacharequal}{\kern0pt}\ t\isactrlsub m\isactrlsub a\isactrlsub x{\isachardoublequoteclose}\isanewline
\ \ \ \ \ \ \ \ \isakeywordONE{have}\isamarkupfalse%
\ {\isachardoublequoteopen}htps\ {\isacharequal}{\kern0pt}\ {\isacharbrackleft}{\kern0pt}t\isactrlsub m\isactrlsub i\isactrlsub n{\isacharbrackright}{\kern0pt}{\isachardoublequoteclose}\isanewline
\ \ \ \ \ \ \ \ \ \ \isakeywordONE{using}\isamarkupfalse%
\ {\isacartoucheopen}htps\ {\isasymnoteq}\ {\isacharbrackleft}{\kern0pt}{\isacharbrackright}{\kern0pt}{\isacartoucheclose}\ {\isacartoucheopen}strict{\isacharunderscore}{\kern0pt}sorted\ htps{\isacartoucheclose}\ {\isacartoucheopen}t\isactrlsub m\isactrlsub i\isactrlsub n\ {\isacharequal}{\kern0pt}\ hd\ htps{\isacartoucheclose}\ {\isacartoucheopen}t\isactrlsub m\isactrlsub a\isactrlsub x\ {\isacharequal}{\kern0pt}\ last\ htps{\isacartoucheclose}\ {\isacartoucheopen}t\isactrlsub m\isactrlsub i\isactrlsub n\ {\isacharequal}{\kern0pt}\ t\isactrlsub m\isactrlsub a\isactrlsub x{\isacartoucheclose}\isanewline
\ \ \ \ \ \ \ \ \ \ \isakeywordONE{by}\isamarkupfalse%
\ {\isacharparenleft}{\kern0pt}induction\ htps{\isacharparenright}{\kern0pt}\ {\isacharparenleft}{\kern0pt}auto\ split{\isacharcolon}{\kern0pt}\ if{\isacharunderscore}{\kern0pt}splits{\isacharparenright}{\kern0pt}\isanewline
\ \ \ \ \ \ \ \ \isakeywordONE{then}\isamarkupfalse%
\ \isakeywordONE{have}\isamarkupfalse%
\ {\isachardoublequoteopen}M\isactrlsub {\isadigit{1}}\ {\isacharequal}{\kern0pt}\ M\isactrlsub n{\isachardoublequoteclose}\isanewline
\ \ \ \ \ \ \ \ \ \ \isakeywordONE{using}\isamarkupfalse%
\ {\isacartoucheopen}valid{\isacharunderscore}{\kern0pt}state{\isacharunderscore}{\kern0pt}seq\ M\isactrlsub {\isadigit{0}}\ {\isacharbrackleft}{\kern0pt}t\isactrlsub m\isactrlsub i\isactrlsub n{\isacharbrackright}{\kern0pt}\ {\isasympi}s\ M\isactrlsub {\isadigit{1}}{\isacartoucheclose}\ \isanewline
\ \ \ \ \ \ \ \ \ \ \ \ \ \ \ \ valid{\isacharunderscore}{\kern0pt}state{\isacharunderscore}{\kern0pt}seq{\isacharunderscore}{\kern0pt}state{\isacharunderscore}{\kern0pt}unique\ {\isacharbrackleft}{\kern0pt}OF\ {\isacartoucheopen}valid{\isacharunderscore}{\kern0pt}state{\isacharunderscore}{\kern0pt}seq\ M\isactrlsub {\isadigit{0}}\ htps\ {\isasympi}s\ M\isactrlsub n{\isacartoucheclose}{\isacharbrackright}{\kern0pt}\isanewline
\ \ \ \ \ \ \ \ \ \ \isakeywordONE{by}\isamarkupfalse%
\ {\isacharparenleft}{\kern0pt}auto\ simp\ add{\isacharcolon}{\kern0pt}\ {\isacartoucheopen}takeWhile\ {\isacharparenleft}{\kern0pt}leq\ t\isactrlsub m\isactrlsub a\isactrlsub x{\isacharparenright}{\kern0pt}\ hs\ {\isacharequal}{\kern0pt}\ hs{\isacartoucheclose}{\isacharparenright}{\kern0pt}\isanewline
\ \ \ \ \ \ \ \ \isakeywordONE{then}\isamarkupfalse%
\ \isakeywordTHREE{show}\isamarkupfalse%
\ {\isacharquery}{\kern0pt}thesis\isanewline
\ \ \ \ \ \ \ \ \ \ \isakeywordONE{using}\isamarkupfalse%
\ {\isacartoucheopen}dropWhile\ {\isacharparenleft}{\kern0pt}le\ t\isactrlsub m\isactrlsub i\isactrlsub n{\isacharparenright}{\kern0pt}\ {\isacharparenleft}{\kern0pt}takeWhile\ {\isacharparenleft}{\kern0pt}leq\ t\isactrlsub m\isactrlsub i\isactrlsub n{\isacharparenright}{\kern0pt}\ hs{\isacharparenright}{\kern0pt}\ {\isacharequal}{\kern0pt}\ takeWhile\ {\isacharparenleft}{\kern0pt}leq\ t\isactrlsub m\isactrlsub i\isactrlsub n{\isacharparenright}{\kern0pt}\ hs{\isacartoucheclose}\isanewline
\ \ \ \ \ \ \ \ \ \ \ \ \ \ \ \ {\isacartoucheopen}valid{\isacharunderscore}{\kern0pt}happ{\isacharunderscore}{\kern0pt}seq\ M\isactrlsub {\isadigit{0}}\ {\isacharparenleft}{\kern0pt}dropWhile\ {\isacharparenleft}{\kern0pt}le\ t\isactrlsub m\isactrlsub i\isactrlsub n{\isacharparenright}{\kern0pt}\ {\isacharparenleft}{\kern0pt}takeWhile\ {\isacharparenleft}{\kern0pt}leq\ t\isactrlsub m\isactrlsub i\isactrlsub n{\isacharparenright}{\kern0pt}\ hs{\isacharparenright}{\kern0pt}{\isacharparenright}{\kern0pt}\ M\isactrlsub {\isadigit{1}}{\isacartoucheclose}\ \isanewline
\ \ \ \ \ \ \ \ \ \ \ \ \ \ \ \ {\isacartoucheopen}t\isactrlsub m\isactrlsub i\isactrlsub n\ {\isacharequal}{\kern0pt}\ t\isactrlsub m\isactrlsub a\isactrlsub x{\isacartoucheclose}\isanewline
\ \ \ \ \ \ \ \ \ \ \isakeywordONE{by}\isamarkupfalse%
\ {\isacharparenleft}{\kern0pt}simp\ add{\isacharcolon}{\kern0pt}\ {\isacartoucheopen}takeWhile\ {\isacharparenleft}{\kern0pt}leq\ t\isactrlsub m\isactrlsub a\isactrlsub x{\isacharparenright}{\kern0pt}\ hs\ {\isacharequal}{\kern0pt}\ hs{\isacartoucheclose}{\isacharparenright}{\kern0pt}\isanewline
\ \ \ \ \ \ \isakeywordONE{next}\isamarkupfalse%
\isanewline
\ \ \ \ \ \ \ \ \isakeywordTHREE{assume}\isamarkupfalse%
\ {\isachardoublequoteopen}t\isactrlsub m\isactrlsub i\isactrlsub n\ {\isacharless}{\kern0pt}\ t\isactrlsub m\isactrlsub a\isactrlsub x{\isachardoublequoteclose}\isanewline
\ \ \ \ \ \ \ \ \isakeywordONE{then}\isamarkupfalse%
\ \isakeywordONE{have}\isamarkupfalse%
\ {\isachardoublequoteopen}{\isasymexists}htps{\isacharprime}{\kern0pt}{\isachardot}{\kern0pt}\ htps\ {\isacharequal}{\kern0pt}\ {\isacharbrackleft}{\kern0pt}{\isacharbrackright}{\kern0pt}\ {\isacharat}{\kern0pt}\ t\isactrlsub m\isactrlsub i\isactrlsub n\ {\isacharhash}{\kern0pt}\ htps{\isacharprime}{\kern0pt}\ {\isacharat}{\kern0pt}\ t\isactrlsub m\isactrlsub a\isactrlsub x\ {\isacharhash}{\kern0pt}\ {\isacharbrackleft}{\kern0pt}{\isacharbrackright}{\kern0pt}{\isachardoublequoteclose}\isanewline
\ \ \ \ \ \ \ \ \ \ \isakeywordONE{using}\isamarkupfalse%
\ {\isacartoucheopen}htps\ {\isasymnoteq}\ {\isacharbrackleft}{\kern0pt}{\isacharbrackright}{\kern0pt}{\isacartoucheclose}\ {\isacartoucheopen}t\isactrlsub m\isactrlsub i\isactrlsub n\ {\isacharequal}{\kern0pt}\ hd\ htps{\isacartoucheclose}\ {\isacartoucheopen}t\isactrlsub m\isactrlsub a\isactrlsub x\ {\isacharequal}{\kern0pt}\ last\ htps{\isacartoucheclose}\isanewline
\ \ \ \ \ \ \ \ \ \ \isakeywordONE{apply}\isamarkupfalse%
\ {\isacharparenleft}{\kern0pt}induction\ htps{\isacharparenright}{\kern0pt}\isanewline
\ \ \ \ \ \ \ \ \ \ \isakeywordONE{apply}\isamarkupfalse%
\ {\isacharparenleft}{\kern0pt}auto\ split{\isacharcolon}{\kern0pt}\ if{\isacharunderscore}{\kern0pt}splits{\isacharparenright}{\kern0pt}\isanewline
\ \ \ \ \ \ \ \ \ \ \isakeywordONE{apply}\isamarkupfalse%
\ {\isacharparenleft}{\kern0pt}metis\ append{\isacharunderscore}{\kern0pt}butlast{\isacharunderscore}{\kern0pt}last{\isacharunderscore}{\kern0pt}id{\isacharparenright}{\kern0pt}\isanewline
\ \ \ \ \ \ \ \ \ \ \isakeywordONE{done}\isamarkupfalse%
\isanewline
\ \ \ \ \ \ \ \ \isakeywordONE{then}\isamarkupfalse%
\ \isakeywordTHREE{obtain}\isamarkupfalse%
\ htps{\isacharprime}{\kern0pt}\ \isakeywordTWO{where}\ {\isachardoublequoteopen}htps\ {\isacharequal}{\kern0pt}\ {\isacharbrackleft}{\kern0pt}{\isacharbrackright}{\kern0pt}\ {\isacharat}{\kern0pt}\ t\isactrlsub m\isactrlsub i\isactrlsub n\ {\isacharhash}{\kern0pt}\ htps{\isacharprime}{\kern0pt}\ {\isacharat}{\kern0pt}\ t\isactrlsub m\isactrlsub a\isactrlsub x\ {\isacharhash}{\kern0pt}\ {\isacharbrackleft}{\kern0pt}{\isacharbrackright}{\kern0pt}{\isachardoublequoteclose}\isanewline
\ \ \ \ \ \ \ \ \ \ \isakeywordONE{by}\isamarkupfalse%
\ auto\isanewline
\ \ \ \ \ \ \ \ \isakeywordONE{then}\isamarkupfalse%
\ \isakeywordONE{have}\isamarkupfalse%
\ {\isachardoublequoteopen}htps{\isacharunderscore}{\kern0pt}seq\ {\isasympi}s\ {\isacharparenleft}{\kern0pt}{\isacharbrackleft}{\kern0pt}{\isacharbrackright}{\kern0pt}\ {\isacharat}{\kern0pt}\ t\isactrlsub m\isactrlsub i\isactrlsub n\ {\isacharhash}{\kern0pt}\ htps{\isacharprime}{\kern0pt}\ {\isacharat}{\kern0pt}\ t\isactrlsub m\isactrlsub a\isactrlsub x\ {\isacharhash}{\kern0pt}\ {\isacharbrackleft}{\kern0pt}{\isacharbrackright}{\kern0pt}{\isacharparenright}{\kern0pt}{\isachardoublequoteclose}\isanewline
\ \ \ \ \ \ \ \ \ \ \isakeywordONE{using}\isamarkupfalse%
\ assms\ \isakeywordONE{by}\isamarkupfalse%
\ auto\isanewline
\ \ \ \ \ \ \ \ \isakeywordONE{then}\isamarkupfalse%
\ \isakeywordONE{have}\isamarkupfalse%
\ {\isachardoublequoteopen}valid{\isacharunderscore}{\kern0pt}happ{\isacharunderscore}{\kern0pt}seq\ M\isactrlsub {\isadigit{1}}\ {\isacharparenleft}{\kern0pt}dropWhile\ {\isacharparenleft}{\kern0pt}leq\ t\isactrlsub m\isactrlsub i\isactrlsub n{\isacharparenright}{\kern0pt}\ {\isacharparenleft}{\kern0pt}takeWhile\ {\isacharparenleft}{\kern0pt}leq\ t\isactrlsub m\isactrlsub a\isactrlsub x{\isacharparenright}{\kern0pt}\ hs{\isacharparenright}{\kern0pt}{\isacharparenright}{\kern0pt}\ M\isactrlsub n{\isachardoublequoteclose}\isanewline
\ \ \ \ \ \ \ \ \ \ \isakeywordONE{using}\isamarkupfalse%
\ assms\ {\isacartoucheopen}valid{\isacharunderscore}{\kern0pt}state{\isacharunderscore}{\kern0pt}seq\ M\isactrlsub {\isadigit{1}}\ {\isacharparenleft}{\kern0pt}tl\ htps{\isacharparenright}{\kern0pt}\ {\isasympi}s\ M\isactrlsub n{\isacartoucheclose}\ {\isacartoucheopen}htps\ {\isacharequal}{\kern0pt}\ {\isacharbrackleft}{\kern0pt}{\isacharbrackright}{\kern0pt}\ {\isacharat}{\kern0pt}\ t\isactrlsub m\isactrlsub i\isactrlsub n\ {\isacharhash}{\kern0pt}\ htps{\isacharprime}{\kern0pt}\ {\isacharat}{\kern0pt}\ {\isacharbrackleft}{\kern0pt}t\isactrlsub m\isactrlsub a\isactrlsub x{\isacharbrackright}{\kern0pt}{\isacartoucheclose}\isanewline
\ \ \ \ \ \ \ \ \ \ \ \ \ \ \ \ valid{\isacharunderscore}{\kern0pt}happ{\isacharunderscore}{\kern0pt}seq{\isacharunderscore}{\kern0pt}valid{\isacharunderscore}{\kern0pt}state{\isacharunderscore}{\kern0pt}seq{\isacharunderscore}{\kern0pt}i{\isacharunderscore}{\kern0pt}to{\isacharunderscore}{\kern0pt}j{\isacharunderscore}{\kern0pt}aux\isanewline
\ \ \ \ \ \ \ \ \ \ \ \ \ \ \ \ \ \ {\isacharbrackleft}{\kern0pt}OF\ {\isacartoucheopen}htps{\isacharunderscore}{\kern0pt}seq\ {\isasympi}s\ {\isacharparenleft}{\kern0pt}{\isacharbrackleft}{\kern0pt}{\isacharbrackright}{\kern0pt}\ {\isacharat}{\kern0pt}\ t\isactrlsub m\isactrlsub i\isactrlsub n\ {\isacharhash}{\kern0pt}\ htps{\isacharprime}{\kern0pt}\ {\isacharat}{\kern0pt}\ t\isactrlsub m\isactrlsub a\isactrlsub x\ {\isacharhash}{\kern0pt}\ {\isacharbrackleft}{\kern0pt}{\isacharbrackright}{\kern0pt}{\isacharparenright}{\kern0pt}{\isacartoucheclose}{\isacharbrackright}{\kern0pt}\isanewline
\ \ \ \ \ \ \ \ \ \ \isakeywordONE{by}\isamarkupfalse%
\ fastforce\isanewline
\ \ \ \ \ \ \ \ \isakeywordONE{then}\isamarkupfalse%
\ \isakeywordTHREE{show}\isamarkupfalse%
\ {\isacharquery}{\kern0pt}thesis\isanewline
\ \ \ \ \ \ \ \ \ \ \isakeywordONE{using}\isamarkupfalse%
\ valid{\isacharunderscore}{\kern0pt}happ{\isacharunderscore}{\kern0pt}seq{\isacharunderscore}{\kern0pt}app{\isacharunderscore}{\kern0pt}iff{\isacharbrackleft}{\kern0pt}\isakeywordTWO{where}\ hs\isactrlsub {\isadigit{1}}{\isacharequal}{\kern0pt}{\isachardoublequoteopen}takeWhile\ {\isacharparenleft}{\kern0pt}leq\ t\isactrlsub m\isactrlsub i\isactrlsub n{\isacharparenright}{\kern0pt}\ hs{\isachardoublequoteclose}\ \isanewline
\ \ \ \ \ \ \ \ \ \ \ \ \ \ \ \ \ \ \ \ \ \ \ \ \ \ \ \ \ \ \ \ \ \ \ \ \ \isakeywordTWO{and}\ hs\isactrlsub {\isadigit{2}}{\isacharequal}{\kern0pt}{\isachardoublequoteopen}dropWhile\ {\isacharparenleft}{\kern0pt}leq\ t\isactrlsub m\isactrlsub i\isactrlsub n{\isacharparenright}{\kern0pt}\ hs{\isachardoublequoteclose}{\isacharbrackright}{\kern0pt}\isanewline
\ \ \ \ \ \ \ \ \ \ \ \ \ \ {\isacartoucheopen}valid{\isacharunderscore}{\kern0pt}happ{\isacharunderscore}{\kern0pt}seq\ M\isactrlsub {\isadigit{0}}\ {\isacharparenleft}{\kern0pt}dropWhile\ {\isacharparenleft}{\kern0pt}le\ t\isactrlsub m\isactrlsub i\isactrlsub n{\isacharparenright}{\kern0pt}\ {\isacharparenleft}{\kern0pt}takeWhile\ {\isacharparenleft}{\kern0pt}leq\ t\isactrlsub m\isactrlsub i\isactrlsub n{\isacharparenright}{\kern0pt}\ hs{\isacharparenright}{\kern0pt}{\isacharparenright}{\kern0pt}\ M\isactrlsub {\isadigit{1}}{\isacartoucheclose}\isanewline
\ \ \ \ \ \ \ \ \ \ \ \ \ \ {\isacartoucheopen}takeWhile\ {\isacharparenleft}{\kern0pt}leq\ t\isactrlsub m\isactrlsub a\isactrlsub x{\isacharparenright}{\kern0pt}\ hs\ {\isacharequal}{\kern0pt}\ hs{\isacartoucheclose}\isanewline
\ \ \ \ \ \ \ \ \ \ \ \ \ \ {\isacartoucheopen}dropWhile\ {\isacharparenleft}{\kern0pt}le\ t\isactrlsub m\isactrlsub i\isactrlsub n{\isacharparenright}{\kern0pt}\ {\isacharparenleft}{\kern0pt}takeWhile\ {\isacharparenleft}{\kern0pt}leq\ t\isactrlsub m\isactrlsub i\isactrlsub n{\isacharparenright}{\kern0pt}\ hs{\isacharparenright}{\kern0pt}\ {\isacharequal}{\kern0pt}\ takeWhile\ {\isacharparenleft}{\kern0pt}leq\ t\isactrlsub m\isactrlsub i\isactrlsub n{\isacharparenright}{\kern0pt}\ hs{\isacartoucheclose}\isanewline
\ \ \ \ \ \ \ \ \ \ \isakeywordONE{by}\isamarkupfalse%
\ auto\isanewline
\ \ \ \ \ \ \isakeywordONE{qed}\isamarkupfalse%
\isanewline
\ \ \ \ \isakeywordONE{qed}\isamarkupfalse%
\isanewline
\ \ \isakeywordONE{qed}\isamarkupfalse%
%
\endisatagproof
{\isafoldproof}%
%
\isadelimproof
\isanewline
%
\endisadelimproof
\isanewline
\ \ \isakeywordONE{lemma}\isamarkupfalse%
\ {\isacharparenleft}{\kern0pt}\isakeywordTWO{in}\ wf{\isacharunderscore}{\kern0pt}ast{\isacharunderscore}{\kern0pt}problem{\isacharparenright}{\kern0pt}\ plan{\isacharunderscore}{\kern0pt}happ{\isacharunderscore}{\kern0pt}path{\isacharunderscore}{\kern0pt}iff{\isacharunderscore}{\kern0pt}valid{\isacharunderscore}{\kern0pt}state{\isacharunderscore}{\kern0pt}seq{\isacharcolon}{\kern0pt}\isanewline
\ \ \ \ \isakeywordTWO{assumes}\ {\isachardoublequoteopen}wf{\isacharunderscore}{\kern0pt}plan\ {\isasympi}s{\isachardoublequoteclose}\ \isakeywordTWO{and}\ {\isachardoublequoteopen}htps{\isacharunderscore}{\kern0pt}seq\ {\isasympi}s\ htps{\isachardoublequoteclose}\isanewline
\ \ \ \ \isakeywordTWO{shows}\ {\isachardoublequoteopen}plan{\isacharunderscore}{\kern0pt}happ{\isacharunderscore}{\kern0pt}path\ M\ {\isasympi}s\ M{\isacharprime}{\kern0pt}\ {\isasymlongleftrightarrow}\ valid{\isacharunderscore}{\kern0pt}state{\isacharunderscore}{\kern0pt}seq\ M\ htps\ {\isasympi}s\ M{\isacharprime}{\kern0pt}{\isachardoublequoteclose}\isanewline
%
\isadelimproof
\ \ \ \ %
\endisadelimproof
%
\isatagproof
\isakeywordONE{unfolding}\isamarkupfalse%
\ plan{\isacharunderscore}{\kern0pt}happ{\isacharunderscore}{\kern0pt}path{\isacharunderscore}{\kern0pt}def\isanewline
\ \ \ \ \isakeywordONE{using}\isamarkupfalse%
\ assms\ ind{\isacharunderscore}{\kern0pt}happ{\isacharunderscore}{\kern0pt}seq{\isacharunderscore}{\kern0pt}exists\ valid{\isacharunderscore}{\kern0pt}happ{\isacharunderscore}{\kern0pt}seq{\isacharunderscore}{\kern0pt}iff{\isacharunderscore}{\kern0pt}valid{\isacharunderscore}{\kern0pt}state{\isacharunderscore}{\kern0pt}seq\ \isakeywordONE{by}\isamarkupfalse%
\ blast%
\endisatagproof
{\isafoldproof}%
%
\isadelimproof
\isanewline
%
\endisadelimproof
\isanewline
\ \ \isakeywordONE{lemmas}\isamarkupfalse%
\ plan{\isacharunderscore}{\kern0pt}happ{\isacharunderscore}{\kern0pt}path{\isacharunderscore}{\kern0pt}iff{\isacharunderscore}{\kern0pt}valid{\isacharunderscore}{\kern0pt}state{\isacharunderscore}{\kern0pt}seq{\isacharunderscore}{\kern0pt}iff{\isacharprime}{\kern0pt}\isanewline
\ \ \ \ {\isacharequal}{\kern0pt}\ wf{\isacharunderscore}{\kern0pt}ast{\isacharunderscore}{\kern0pt}problem{\isachardot}{\kern0pt}plan{\isacharunderscore}{\kern0pt}happ{\isacharunderscore}{\kern0pt}path{\isacharunderscore}{\kern0pt}iff{\isacharunderscore}{\kern0pt}valid{\isacharunderscore}{\kern0pt}state{\isacharunderscore}{\kern0pt}seq{\isacharbrackleft}{\kern0pt}of\ P{\isacharcomma}{\kern0pt}\ OF\ wf{\isacharunderscore}{\kern0pt}ast{\isacharunderscore}{\kern0pt}problem{\isachardot}{\kern0pt}intro{\isacharbrackright}{\kern0pt}\isanewline
\isanewline
\ \ \isakeywordONE{lemma}\isamarkupfalse%
\ {\isachardoublequoteopen}strict{\isacharunderscore}{\kern0pt}sorted\ xs\ {\isasymLongrightarrow}\ strict{\isacharunderscore}{\kern0pt}sorted\ ys\ {\isasymLongrightarrow}\ set\ xs\ {\isacharequal}{\kern0pt}\ set\ ys\ {\isasymLongrightarrow}\ xs\ {\isacharequal}{\kern0pt}\ ys{\isachardoublequoteclose}\isanewline
%
\isadelimproof
\ \ \ \ %
\endisadelimproof
%
\isatagproof
\isakeywordONE{by}\isamarkupfalse%
\ {\isacharparenleft}{\kern0pt}simp\ add{\isacharcolon}{\kern0pt}\ strict{\isacharunderscore}{\kern0pt}sorted{\isacharunderscore}{\kern0pt}iff\ sorted{\isacharunderscore}{\kern0pt}distinct{\isacharunderscore}{\kern0pt}set{\isacharunderscore}{\kern0pt}unique{\isacharparenright}{\kern0pt}%
\endisatagproof
{\isafoldproof}%
%
\isadelimproof
%
\endisadelimproof
%
\begin{isamarkuptext}%
Proof that a sequence of happening time points is unique for every plan%
\end{isamarkuptext}\isamarkuptrue%
\ \ \isakeywordONE{lemma}\isamarkupfalse%
\ htps{\isacharunderscore}{\kern0pt}seq{\isacharunderscore}{\kern0pt}unique{\isacharcolon}{\kern0pt}\isanewline
\ \ \ \ \isakeywordTWO{assumes}\ {\isachardoublequoteopen}htps{\isacharunderscore}{\kern0pt}seq\ {\isasympi}s\ htps\isactrlsub {\isadigit{1}}{\isachardoublequoteclose}\isanewline
\ \ \ \ \ \ \ \ \isakeywordTWO{and}\ {\isachardoublequoteopen}htps{\isacharunderscore}{\kern0pt}seq\ {\isasympi}s\ htps\isactrlsub {\isadigit{2}}{\isachardoublequoteclose}\isanewline
\ \ \ \ \isakeywordTWO{shows}\ {\isachardoublequoteopen}htps\isactrlsub {\isadigit{1}}\ {\isacharequal}{\kern0pt}\ htps\isactrlsub {\isadigit{2}}{\isachardoublequoteclose}\isanewline
%
\isadelimproof
\ \ %
\endisadelimproof
%
\isatagproof
\isakeywordONE{proof}\isamarkupfalse%
\ {\isacharminus}{\kern0pt}\isanewline
\ \ \ \ \isakeywordONE{have}\isamarkupfalse%
\ {\isachardoublequoteopen}set\ htps\isactrlsub {\isadigit{1}}\ {\isacharequal}{\kern0pt}\ set\ htps\isactrlsub {\isadigit{2}}{\isachardoublequoteclose}\isanewline
\ \ \ \ \ \ \isakeywordONE{using}\isamarkupfalse%
\ assms\ \isakeywordONE{unfolding}\isamarkupfalse%
\ htps{\isacharunderscore}{\kern0pt}seq{\isacharunderscore}{\kern0pt}def\ \isakeywordONE{by}\isamarkupfalse%
\ auto\isanewline
\ \ \ \ \isakeywordONE{then}\isamarkupfalse%
\ \isakeywordTHREE{show}\isamarkupfalse%
\ {\isacharquery}{\kern0pt}thesis\isanewline
\ \ \ \ \ \ \isakeywordONE{using}\isamarkupfalse%
\ assms\ \isakeywordONE{unfolding}\isamarkupfalse%
\ htps{\isacharunderscore}{\kern0pt}seq{\isacharunderscore}{\kern0pt}def\ \isakeywordONE{by}\isamarkupfalse%
\ {\isacharparenleft}{\kern0pt}simp\ add{\isacharcolon}{\kern0pt}\ strict{\isacharunderscore}{\kern0pt}sorted{\isacharunderscore}{\kern0pt}iff\ sorted{\isacharunderscore}{\kern0pt}distinct{\isacharunderscore}{\kern0pt}set{\isacharunderscore}{\kern0pt}unique{\isacharparenright}{\kern0pt}\isanewline
\ \ \isakeywordONE{qed}\isamarkupfalse%
%
\endisatagproof
{\isafoldproof}%
%
\isadelimproof
\isanewline
%
\endisadelimproof
\ \ \isanewline
\isanewline
\ \ \isakeywordONE{lemma}\isamarkupfalse%
\ valid{\isacharunderscore}{\kern0pt}plan{\isacharunderscore}{\kern0pt}from{\isadigit{2}}{\isacharunderscore}{\kern0pt}iff{\isacharcolon}{\kern0pt}\isanewline
\ \ \ \ \isakeywordTWO{assumes}\ {\isachardoublequoteopen}wf{\isacharunderscore}{\kern0pt}problem{\isachardoublequoteclose}\isanewline
\ \ \ \ \isakeywordTWO{shows}\ {\isachardoublequoteopen}valid{\isacharunderscore}{\kern0pt}plan{\isacharunderscore}{\kern0pt}from{\isadigit{2}}\ I\ {\isasympi}s\ {\isasymlongleftrightarrow}\ valid{\isacharunderscore}{\kern0pt}plan{\isacharunderscore}{\kern0pt}from\ I\ {\isasympi}s{\isachardoublequoteclose}\isanewline
%
\isadelimproof
\ \ \ \ %
\endisadelimproof
%
\isatagproof
\isakeywordONE{unfolding}\isamarkupfalse%
\ valid{\isacharunderscore}{\kern0pt}plan{\isacharunderscore}{\kern0pt}from{\isacharunderscore}{\kern0pt}def\ valid{\isacharunderscore}{\kern0pt}plan{\isacharunderscore}{\kern0pt}from{\isadigit{2}}{\isacharunderscore}{\kern0pt}def\isanewline
\ \ \ \ \isakeywordONE{using}\isamarkupfalse%
\ htps{\isacharunderscore}{\kern0pt}seq{\isacharunderscore}{\kern0pt}exists\ htps{\isacharunderscore}{\kern0pt}seq{\isacharunderscore}{\kern0pt}unique\isanewline
\ \ \ \ \ \ \ \ \ \ plan{\isacharunderscore}{\kern0pt}happ{\isacharunderscore}{\kern0pt}path{\isacharunderscore}{\kern0pt}iff{\isacharunderscore}{\kern0pt}valid{\isacharunderscore}{\kern0pt}state{\isacharunderscore}{\kern0pt}seq{\isacharunderscore}{\kern0pt}iff{\isacharprime}{\kern0pt}{\isacharbrackleft}{\kern0pt}OF\ assms{\isacharbrackright}{\kern0pt}\ \isanewline
\ \ \ \ \isakeywordONE{by}\isamarkupfalse%
\ blast%
\endisatagproof
{\isafoldproof}%
%
\isadelimproof
\isanewline
%
\endisadelimproof
\isanewline
\ \ \isakeywordONE{lemma}\isamarkupfalse%
\ \isanewline
\ \ \ \ \isakeywordTWO{assumes}\ {\isachardoublequoteopen}wf{\isacharunderscore}{\kern0pt}problem{\isachardoublequoteclose}\isanewline
\ \ \ \ \isakeywordTWO{shows}\ {\isachardoublequoteopen}valid{\isacharunderscore}{\kern0pt}plan{\isadigit{2}}\ {\isasympi}s\ {\isasymlongleftrightarrow}\ valid{\isacharunderscore}{\kern0pt}plan\ {\isasympi}s{\isachardoublequoteclose}\isanewline
%
\isadelimproof
\ \ \ \ %
\endisadelimproof
%
\isatagproof
\isakeywordONE{unfolding}\isamarkupfalse%
\ valid{\isacharunderscore}{\kern0pt}plan{\isacharunderscore}{\kern0pt}def\ valid{\isacharunderscore}{\kern0pt}plan{\isadigit{2}}{\isacharunderscore}{\kern0pt}def\ \isanewline
\ \ \ \ \isakeywordONE{using}\isamarkupfalse%
\ valid{\isacharunderscore}{\kern0pt}plan{\isacharunderscore}{\kern0pt}from{\isadigit{2}}{\isacharunderscore}{\kern0pt}iff{\isacharbrackleft}{\kern0pt}OF\ assms{\isacharbrackright}{\kern0pt}\ \isakeywordONE{by}\isamarkupfalse%
\ blast%
\endisatagproof
{\isafoldproof}%
%
\isadelimproof
\isanewline
%
\endisadelimproof
\ \ \isakeywordONE{find{\isacharunderscore}{\kern0pt}theorems}\isamarkupfalse%
\ {\isachardoublequoteopen}acts{\isacharunderscore}{\kern0pt}of{\isacharunderscore}{\kern0pt}plan{\isacharunderscore}{\kern0pt}at{\isachardoublequoteclose}\isanewline
\isakeywordTWO{end}\isamarkupfalse%
\ %
\isamarkupcmt{Context of \isa{ast{\isacharunderscore}{\kern0pt}problem}%
}\isanewline
%
\isadelimtheory
\isanewline
%
\endisadelimtheory
%
\isatagtheory
\isakeywordTWO{end}\isamarkupfalse%
\ %
\isamarkupcmt{Theory%
}%
\endisatagtheory
{\isafoldtheory}%
%
\isadelimtheory
%
\endisadelimtheory
%
\end{isabellebody}%
\endinput
%:%file=~/work/temp_planning_certification/temporal-planning-certification/lib/temporal-pddl-semantics/TEMPORAL_PDDL_Semantics_Alt.thy%:%
%:%11=1%:%
%:%27=2%:%
%:%28=2%:%
%:%29=3%:%
%:%30=4%:%
%:%31=5%:%
%:%32=6%:%
%:%41=8%:%
%:%42=9%:%
%:%44=11%:%
%:%45=11%:%
%:%46=12%:%
%:%47=13%:%
%:%48=14%:%
%:%51=15%:%
%:%55=15%:%
%:%56=15%:%
%:%57=16%:%
%:%58=16%:%
%:%63=16%:%
%:%66=17%:%
%:%67=18%:%
%:%68=18%:%
%:%69=19%:%
%:%70=20%:%
%:%71=21%:%
%:%74=22%:%
%:%78=22%:%
%:%79=22%:%
%:%80=23%:%
%:%81=23%:%
%:%86=23%:%
%:%89=24%:%
%:%90=25%:%
%:%91=25%:%
%:%92=26%:%
%:%93=27%:%
%:%94=28%:%
%:%97=29%:%
%:%101=29%:%
%:%102=29%:%
%:%103=29%:%
%:%108=29%:%
%:%111=30%:%
%:%112=31%:%
%:%113=31%:%
%:%114=32%:%
%:%115=33%:%
%:%116=34%:%
%:%119=35%:%
%:%123=35%:%
%:%124=35%:%
%:%125=36%:%
%:%126=36%:%
%:%127=37%:%
%:%128=37%:%
%:%129=38%:%
%:%130=38%:%
%:%131=39%:%
%:%132=39%:%
%:%133=40%:%
%:%134=40%:%
%:%135=41%:%
%:%136=41%:%
%:%137=41%:%
%:%138=42%:%
%:%139=42%:%
%:%140=43%:%
%:%148=43%:%
%:%149=44%:%
%:%150=45%:%
%:%151=45%:%
%:%152=46%:%
%:%153=47%:%
%:%154=48%:%
%:%157=49%:%
%:%161=49%:%
%:%162=49%:%
%:%163=50%:%
%:%164=50%:%
%:%165=51%:%
%:%166=51%:%
%:%167=52%:%
%:%168=52%:%
%:%169=53%:%
%:%170=53%:%
%:%171=54%:%
%:%172=54%:%
%:%173=55%:%
%:%174=55%:%
%:%175=55%:%
%:%176=56%:%
%:%177=56%:%
%:%178=57%:%
%:%186=57%:%
%:%187=58%:%
%:%188=59%:%
%:%189=59%:%
%:%190=60%:%
%:%191=61%:%
%:%192=62%:%
%:%195=63%:%
%:%199=63%:%
%:%200=63%:%
%:%201=64%:%
%:%202=64%:%
%:%203=65%:%
%:%204=65%:%
%:%205=66%:%
%:%206=66%:%
%:%207=67%:%
%:%208=67%:%
%:%209=68%:%
%:%210=68%:%
%:%211=69%:%
%:%212=69%:%
%:%213=69%:%
%:%214=70%:%
%:%215=70%:%
%:%216=71%:%
%:%224=71%:%
%:%225=72%:%
%:%226=73%:%
%:%227=74%:%
%:%228=74%:%
%:%229=75%:%
%:%230=76%:%
%:%231=77%:%
%:%232=78%:%
%:%235=79%:%
%:%239=79%:%
%:%240=79%:%
%:%241=80%:%
%:%242=80%:%
%:%243=81%:%
%:%244=81%:%
%:%245=82%:%
%:%246=82%:%
%:%247=83%:%
%:%248=83%:%
%:%249=84%:%
%:%250=84%:%
%:%251=84%:%
%:%252=85%:%
%:%253=85%:%
%:%254=86%:%
%:%255=86%:%
%:%256=87%:%
%:%257=87%:%
%:%258=87%:%
%:%259=88%:%
%:%260=88%:%
%:%261=88%:%
%:%262=89%:%
%:%263=89%:%
%:%264=89%:%
%:%265=90%:%
%:%266=90%:%
%:%267=91%:%
%:%268=91%:%
%:%269=91%:%
%:%270=92%:%
%:%271=92%:%
%:%272=92%:%
%:%273=93%:%
%:%274=93%:%
%:%275=94%:%
%:%276=94%:%
%:%277=95%:%
%:%278=95%:%
%:%279=95%:%
%:%280=96%:%
%:%281=96%:%
%:%282=97%:%
%:%283=97%:%
%:%284=97%:%
%:%285=98%:%
%:%286=98%:%
%:%287=99%:%
%:%288=100%:%
%:%289=100%:%
%:%290=101%:%
%:%296=101%:%
%:%299=102%:%
%:%300=103%:%
%:%301=103%:%
%:%302=104%:%
%:%303=105%:%
%:%304=106%:%
%:%307=107%:%
%:%311=107%:%
%:%312=107%:%
%:%313=107%:%
%:%318=107%:%
%:%321=108%:%
%:%322=109%:%
%:%323=109%:%
%:%324=110%:%
%:%325=111%:%
%:%326=112%:%
%:%327=113%:%
%:%330=114%:%
%:%334=114%:%
%:%335=114%:%
%:%336=115%:%
%:%337=116%:%
%:%338=116%:%
%:%352=118%:%
%:%362=120%:%
%:%363=120%:%
%:%364=121%:%
%:%365=122%:%
%:%366=122%:%
%:%367=123%:%
%:%369=125%:%
%:%370=126%:%
%:%371=127%:%
%:%372=127%:%
%:%373=128%:%
%:%374=129%:%
%:%375=130%:%
%:%376=131%:%
%:%377=132%:%
%:%378=132%:%
%:%379=133%:%
%:%380=134%:%
%:%381=135%:%
%:%384=136%:%
%:%388=136%:%
%:%389=136%:%
%:%390=137%:%
%:%391=137%:%
%:%400=139%:%
%:%401=140%:%
%:%403=141%:%
%:%404=141%:%
%:%405=142%:%
%:%410=148%:%
%:%412=149%:%
%:%413=149%:%
%:%414=150%:%
%:%417=154%:%
%:%419=158%:%
%:%420=158%:%
%:%421=159%:%
%:%423=161%:%
%:%424=162%:%
%:%425=163%:%
%:%427=164%:%
%:%428=164%:%
%:%429=165%:%
%:%430=166%:%
%:%435=171%:%
%:%436=172%:%
%:%437=173%:%
%:%438=173%:%
%:%439=174%:%
%:%440=175%:%
%:%441=176%:%
%:%442=176%:%
%:%443=177%:%
%:%445=179%:%
%:%447=180%:%
%:%448=180%:%
%:%451=181%:%
%:%455=181%:%
%:%456=181%:%
%:%457=181%:%
%:%471=183%:%
%:%483=185%:%
%:%485=187%:%
%:%486=187%:%
%:%489=188%:%
%:%493=188%:%
%:%494=188%:%
%:%499=188%:%
%:%502=189%:%
%:%503=190%:%
%:%504=190%:%
%:%505=191%:%
%:%506=192%:%
%:%509=193%:%
%:%510=194%:%
%:%514=194%:%
%:%515=194%:%
%:%520=194%:%
%:%523=195%:%
%:%524=196%:%
%:%525=196%:%
%:%526=197%:%
%:%527=198%:%
%:%530=199%:%
%:%534=199%:%
%:%535=199%:%
%:%536=199%:%
%:%545=201%:%
%:%546=202%:%
%:%548=203%:%
%:%549=203%:%
%:%550=204%:%
%:%551=205%:%
%:%552=206%:%
%:%553=207%:%
%:%556=208%:%
%:%560=208%:%
%:%561=208%:%
%:%562=209%:%
%:%563=209%:%
%:%564=210%:%
%:%565=210%:%
%:%566=211%:%
%:%567=211%:%
%:%568=212%:%
%:%569=212%:%
%:%570=213%:%
%:%571=213%:%
%:%572=213%:%
%:%573=214%:%
%:%574=214%:%
%:%575=214%:%
%:%576=215%:%
%:%577=215%:%
%:%578=215%:%
%:%579=216%:%
%:%580=216%:%
%:%581=217%:%
%:%582=217%:%
%:%583=218%:%
%:%584=218%:%
%:%585=219%:%
%:%586=219%:%
%:%587=220%:%
%:%588=220%:%
%:%589=221%:%
%:%590=222%:%
%:%591=222%:%
%:%592=223%:%
%:%593=223%:%
%:%594=223%:%
%:%595=224%:%
%:%596=224%:%
%:%597=224%:%
%:%598=225%:%
%:%599=225%:%
%:%600=226%:%
%:%601=226%:%
%:%602=227%:%
%:%603=227%:%
%:%604=227%:%
%:%605=228%:%
%:%606=228%:%
%:%607=229%:%
%:%608=229%:%
%:%609=229%:%
%:%610=230%:%
%:%611=230%:%
%:%612=231%:%
%:%613=231%:%
%:%614=232%:%
%:%615=232%:%
%:%616=232%:%
%:%617=233%:%
%:%618=233%:%
%:%619=234%:%
%:%620=234%:%
%:%621=234%:%
%:%622=235%:%
%:%623=235%:%
%:%624=236%:%
%:%625=236%:%
%:%626=237%:%
%:%627=237%:%
%:%628=237%:%
%:%629=238%:%
%:%630=238%:%
%:%631=239%:%
%:%632=239%:%
%:%633=239%:%
%:%634=240%:%
%:%635=240%:%
%:%636=241%:%
%:%637=241%:%
%:%638=242%:%
%:%639=242%:%
%:%640=242%:%
%:%641=243%:%
%:%642=243%:%
%:%643=243%:%
%:%644=244%:%
%:%645=244%:%
%:%646=245%:%
%:%647=245%:%
%:%648=246%:%
%:%649=246%:%
%:%650=247%:%
%:%651=247%:%
%:%652=248%:%
%:%653=248%:%
%:%654=249%:%
%:%655=249%:%
%:%656=250%:%
%:%657=250%:%
%:%658=251%:%
%:%659=251%:%
%:%660=252%:%
%:%661=252%:%
%:%662=253%:%
%:%663=253%:%
%:%664=254%:%
%:%665=254%:%
%:%666=254%:%
%:%667=255%:%
%:%668=255%:%
%:%669=255%:%
%:%670=256%:%
%:%671=256%:%
%:%672=256%:%
%:%673=257%:%
%:%674=257%:%
%:%675=258%:%
%:%676=258%:%
%:%677=259%:%
%:%678=259%:%
%:%679=259%:%
%:%680=260%:%
%:%681=260%:%
%:%682=261%:%
%:%683=261%:%
%:%684=261%:%
%:%685=262%:%
%:%686=262%:%
%:%687=262%:%
%:%688=262%:%
%:%689=263%:%
%:%690=263%:%
%:%691=263%:%
%:%692=264%:%
%:%693=264%:%
%:%694=265%:%
%:%695=265%:%
%:%696=266%:%
%:%697=266%:%
%:%698=266%:%
%:%699=267%:%
%:%700=267%:%
%:%701=268%:%
%:%702=268%:%
%:%703=268%:%
%:%704=269%:%
%:%705=269%:%
%:%706=269%:%
%:%707=270%:%
%:%708=270%:%
%:%709=271%:%
%:%710=271%:%
%:%711=272%:%
%:%712=272%:%
%:%713=272%:%
%:%714=273%:%
%:%715=274%:%
%:%716=274%:%
%:%717=275%:%
%:%718=275%:%
%:%719=275%:%
%:%720=276%:%
%:%721=276%:%
%:%722=276%:%
%:%723=276%:%
%:%724=277%:%
%:%725=277%:%
%:%726=277%:%
%:%727=278%:%
%:%728=278%:%
%:%729=279%:%
%:%730=279%:%
%:%731=280%:%
%:%732=280%:%
%:%733=280%:%
%:%734=281%:%
%:%735=281%:%
%:%736=282%:%
%:%737=282%:%
%:%738=282%:%
%:%739=283%:%
%:%740=283%:%
%:%741=283%:%
%:%742=284%:%
%:%743=284%:%
%:%744=285%:%
%:%745=285%:%
%:%746=286%:%
%:%747=286%:%
%:%748=286%:%
%:%749=287%:%
%:%750=288%:%
%:%751=288%:%
%:%752=289%:%
%:%753=289%:%
%:%754=289%:%
%:%755=290%:%
%:%756=290%:%
%:%757=290%:%
%:%758=290%:%
%:%759=291%:%
%:%760=291%:%
%:%761=291%:%
%:%762=292%:%
%:%763=292%:%
%:%764=293%:%
%:%765=293%:%
%:%766=294%:%
%:%767=294%:%
%:%768=294%:%
%:%769=295%:%
%:%770=295%:%
%:%771=296%:%
%:%772=296%:%
%:%773=296%:%
%:%774=297%:%
%:%775=297%:%
%:%776=297%:%
%:%777=298%:%
%:%778=298%:%
%:%779=299%:%
%:%780=299%:%
%:%781=300%:%
%:%787=300%:%
%:%790=301%:%
%:%791=301%:%
%:%792=302%:%
%:%793=303%:%
%:%794=303%:%
%:%795=304%:%
%:%797=306%:%
%:%798=307%:%
%:%800=308%:%
%:%801=308%:%
%:%802=309%:%
%:%803=310%:%
%:%804=311%:%
%:%805=312%:%
%:%806=313%:%
%:%807=314%:%
%:%810=315%:%
%:%814=315%:%
%:%815=315%:%
%:%816=316%:%
%:%817=316%:%
%:%818=317%:%
%:%819=317%:%
%:%820=318%:%
%:%821=318%:%
%:%822=318%:%
%:%823=318%:%
%:%824=319%:%
%:%825=319%:%
%:%826=320%:%
%:%827=320%:%
%:%828=320%:%
%:%829=321%:%
%:%830=321%:%
%:%831=321%:%
%:%832=322%:%
%:%833=322%:%
%:%834=322%:%
%:%835=323%:%
%:%836=323%:%
%:%837=323%:%
%:%838=324%:%
%:%839=324%:%
%:%840=325%:%
%:%841=325%:%
%:%842=325%:%
%:%843=326%:%
%:%844=326%:%
%:%845=326%:%
%:%846=327%:%
%:%847=327%:%
%:%848=328%:%
%:%849=328%:%
%:%850=329%:%
%:%851=329%:%
%:%852=329%:%
%:%853=330%:%
%:%854=330%:%
%:%855=331%:%
%:%856=331%:%
%:%857=332%:%
%:%858=332%:%
%:%859=333%:%
%:%860=333%:%
%:%861=334%:%
%:%862=335%:%
%:%863=335%:%
%:%864=336%:%
%:%865=336%:%
%:%866=336%:%
%:%867=337%:%
%:%868=337%:%
%:%869=337%:%
%:%870=338%:%
%:%871=338%:%
%:%872=339%:%
%:%873=339%:%
%:%874=340%:%
%:%875=340%:%
%:%876=340%:%
%:%877=341%:%
%:%878=341%:%
%:%879=341%:%
%:%880=342%:%
%:%881=342%:%
%:%882=343%:%
%:%883=343%:%
%:%884=344%:%
%:%885=344%:%
%:%886=344%:%
%:%887=345%:%
%:%888=345%:%
%:%889=345%:%
%:%890=346%:%
%:%891=346%:%
%:%892=347%:%
%:%893=347%:%
%:%894=348%:%
%:%895=348%:%
%:%896=348%:%
%:%897=349%:%
%:%898=349%:%
%:%899=350%:%
%:%900=350%:%
%:%901=350%:%
%:%902=351%:%
%:%903=351%:%
%:%904=351%:%
%:%905=352%:%
%:%906=352%:%
%:%907=352%:%
%:%908=353%:%
%:%909=353%:%
%:%910=353%:%
%:%911=354%:%
%:%912=354%:%
%:%913=354%:%
%:%914=355%:%
%:%915=355%:%
%:%916=356%:%
%:%917=356%:%
%:%918=357%:%
%:%919=357%:%
%:%920=358%:%
%:%921=359%:%
%:%922=360%:%
%:%923=360%:%
%:%924=360%:%
%:%925=361%:%
%:%926=361%:%
%:%927=362%:%
%:%928=362%:%
%:%929=363%:%
%:%930=363%:%
%:%931=363%:%
%:%932=364%:%
%:%933=364%:%
%:%934=364%:%
%:%935=365%:%
%:%936=365%:%
%:%937=365%:%
%:%938=366%:%
%:%939=366%:%
%:%940=367%:%
%:%941=367%:%
%:%942=368%:%
%:%943=369%:%
%:%944=369%:%
%:%945=370%:%
%:%946=370%:%
%:%947=370%:%
%:%948=371%:%
%:%949=371%:%
%:%950=371%:%
%:%951=372%:%
%:%952=372%:%
%:%953=372%:%
%:%954=373%:%
%:%955=373%:%
%:%956=374%:%
%:%957=374%:%
%:%958=374%:%
%:%959=375%:%
%:%960=375%:%
%:%961=375%:%
%:%962=376%:%
%:%963=376%:%
%:%964=376%:%
%:%965=377%:%
%:%966=377%:%
%:%967=377%:%
%:%968=378%:%
%:%969=378%:%
%:%970=378%:%
%:%971=378%:%
%:%972=379%:%
%:%973=380%:%
%:%974=380%:%
%:%975=381%:%
%:%976=381%:%
%:%977=381%:%
%:%978=381%:%
%:%979=382%:%
%:%980=383%:%
%:%981=383%:%
%:%982=383%:%
%:%983=384%:%
%:%984=384%:%
%:%985=385%:%
%:%986=385%:%
%:%987=386%:%
%:%988=387%:%
%:%989=388%:%
%:%990=388%:%
%:%991=389%:%
%:%992=389%:%
%:%993=390%:%
%:%1003=392%:%
%:%1004=393%:%
%:%1005=394%:%
%:%1007=395%:%
%:%1008=395%:%
%:%1009=396%:%
%:%1010=397%:%
%:%1011=398%:%
%:%1012=399%:%
%:%1013=400%:%
%:%1016=401%:%
%:%1020=401%:%
%:%1021=401%:%
%:%1022=402%:%
%:%1023=402%:%
%:%1024=403%:%
%:%1025=403%:%
%:%1026=404%:%
%:%1027=404%:%
%:%1028=404%:%
%:%1029=404%:%
%:%1030=405%:%
%:%1031=405%:%
%:%1032=406%:%
%:%1033=407%:%
%:%1034=407%:%
%:%1035=407%:%
%:%1036=407%:%
%:%1037=408%:%
%:%1038=408%:%
%:%1039=408%:%
%:%1040=409%:%
%:%1041=409%:%
%:%1042=409%:%
%:%1043=410%:%
%:%1044=410%:%
%:%1045=410%:%
%:%1046=411%:%
%:%1047=411%:%
%:%1048=411%:%
%:%1049=411%:%
%:%1050=412%:%
%:%1051=412%:%
%:%1052=412%:%
%:%1053=413%:%
%:%1054=413%:%
%:%1055=413%:%
%:%1056=413%:%
%:%1057=414%:%
%:%1057=415%:%
%:%1058=416%:%
%:%1059=416%:%
%:%1060=416%:%
%:%1061=417%:%
%:%1062=418%:%
%:%1063=418%:%
%:%1064=419%:%
%:%1065=419%:%
%:%1066=420%:%
%:%1067=420%:%
%:%1068=421%:%
%:%1069=421%:%
%:%1070=422%:%
%:%1071=422%:%
%:%1072=422%:%
%:%1073=423%:%
%:%1074=423%:%
%:%1075=423%:%
%:%1076=424%:%
%:%1077=425%:%
%:%1078=425%:%
%:%1079=426%:%
%:%1080=426%:%
%:%1081=427%:%
%:%1082=428%:%
%:%1083=428%:%
%:%1084=429%:%
%:%1085=429%:%
%:%1086=430%:%
%:%1087=430%:%
%:%1088=430%:%
%:%1089=431%:%
%:%1090=431%:%
%:%1091=432%:%
%:%1092=432%:%
%:%1093=432%:%
%:%1094=433%:%
%:%1095=433%:%
%:%1096=433%:%
%:%1097=434%:%
%:%1098=434%:%
%:%1099=434%:%
%:%1100=435%:%
%:%1101=435%:%
%:%1102=436%:%
%:%1103=436%:%
%:%1104=437%:%
%:%1105=437%:%
%:%1106=438%:%
%:%1107=438%:%
%:%1108=438%:%
%:%1109=439%:%
%:%1119=443%:%
%:%1120=444%:%
%:%1121=445%:%
%:%1123=446%:%
%:%1124=446%:%
%:%1125=447%:%
%:%1126=448%:%
%:%1127=449%:%
%:%1128=450%:%
%:%1129=451%:%
%:%1130=452%:%
%:%1133=453%:%
%:%1137=453%:%
%:%1138=453%:%
%:%1139=454%:%
%:%1140=454%:%
%:%1141=455%:%
%:%1142=455%:%
%:%1143=456%:%
%:%1144=456%:%
%:%1145=456%:%
%:%1146=456%:%
%:%1147=457%:%
%:%1148=457%:%
%:%1149=457%:%
%:%1150=458%:%
%:%1151=458%:%
%:%1152=458%:%
%:%1153=459%:%
%:%1154=459%:%
%:%1155=460%:%
%:%1156=460%:%
%:%1157=461%:%
%:%1158=462%:%
%:%1159=462%:%
%:%1160=463%:%
%:%1161=463%:%
%:%1162=464%:%
%:%1163=464%:%
%:%1164=464%:%
%:%1165=464%:%
%:%1166=465%:%
%:%1167=465%:%
%:%1168=466%:%
%:%1169=466%:%
%:%1170=466%:%
%:%1171=466%:%
%:%1172=467%:%
%:%1173=467%:%
%:%1174=468%:%
%:%1175=468%:%
%:%1176=468%:%
%:%1177=468%:%
%:%1178=469%:%
%:%1179=469%:%
%:%1180=469%:%
%:%1181=470%:%
%:%1182=470%:%
%:%1183=471%:%
%:%1184=471%:%
%:%1185=471%:%
%:%1186=472%:%
%:%1187=472%:%
%:%1188=472%:%
%:%1189=473%:%
%:%1190=473%:%
%:%1191=474%:%
%:%1192=474%:%
%:%1193=475%:%
%:%1194=475%:%
%:%1195=476%:%
%:%1196=476%:%
%:%1197=477%:%
%:%1198=478%:%
%:%1199=478%:%
%:%1200=479%:%
%:%1201=479%:%
%:%1202=479%:%
%:%1203=480%:%
%:%1204=480%:%
%:%1205=480%:%
%:%1206=481%:%
%:%1207=481%:%
%:%1208=482%:%
%:%1209=482%:%
%:%1210=483%:%
%:%1211=483%:%
%:%1212=483%:%
%:%1213=484%:%
%:%1214=484%:%
%:%1215=484%:%
%:%1216=485%:%
%:%1217=485%:%
%:%1218=486%:%
%:%1219=486%:%
%:%1220=487%:%
%:%1221=487%:%
%:%1222=487%:%
%:%1223=488%:%
%:%1224=488%:%
%:%1225=488%:%
%:%1226=489%:%
%:%1227=489%:%
%:%1228=490%:%
%:%1229=490%:%
%:%1230=491%:%
%:%1231=491%:%
%:%1232=491%:%
%:%1233=492%:%
%:%1234=492%:%
%:%1235=493%:%
%:%1236=493%:%
%:%1237=493%:%
%:%1238=494%:%
%:%1239=494%:%
%:%1240=494%:%
%:%1241=495%:%
%:%1242=495%:%
%:%1243=495%:%
%:%1244=496%:%
%:%1245=496%:%
%:%1246=496%:%
%:%1247=497%:%
%:%1248=497%:%
%:%1249=497%:%
%:%1250=498%:%
%:%1251=498%:%
%:%1252=499%:%
%:%1253=499%:%
%:%1254=500%:%
%:%1255=500%:%
%:%1256=501%:%
%:%1257=501%:%
%:%1258=502%:%
%:%1259=502%:%
%:%1260=503%:%
%:%1261=503%:%
%:%1262=503%:%
%:%1263=504%:%
%:%1264=504%:%
%:%1265=505%:%
%:%1266=505%:%
%:%1267=505%:%
%:%1268=506%:%
%:%1269=506%:%
%:%1270=506%:%
%:%1271=507%:%
%:%1272=507%:%
%:%1273=507%:%
%:%1274=508%:%
%:%1275=508%:%
%:%1276=508%:%
%:%1277=509%:%
%:%1278=509%:%
%:%1279=510%:%
%:%1280=510%:%
%:%1281=510%:%
%:%1282=511%:%
%:%1283=511%:%
%:%1284=511%:%
%:%1285=512%:%
%:%1286=512%:%
%:%1287=513%:%
%:%1288=514%:%
%:%1289=514%:%
%:%1290=515%:%
%:%1291=515%:%
%:%1292=516%:%
%:%1302=518%:%
%:%1303=519%:%
%:%1305=520%:%
%:%1306=520%:%
%:%1307=521%:%
%:%1308=522%:%
%:%1311=523%:%
%:%1315=523%:%
%:%1316=523%:%
%:%1317=523%:%
%:%1326=525%:%
%:%1327=526%:%
%:%1328=527%:%
%:%1330=528%:%
%:%1331=528%:%
%:%1332=529%:%
%:%1333=530%:%
%:%1334=531%:%
%:%1335=532%:%
%:%1338=533%:%
%:%1342=533%:%
%:%1343=533%:%
%:%1344=534%:%
%:%1345=534%:%
%:%1350=534%:%
%:%1353=535%:%
%:%1354=536%:%
%:%1355=536%:%
%:%1356=537%:%
%:%1357=538%:%
%:%1358=539%:%
%:%1361=540%:%
%:%1365=540%:%
%:%1366=540%:%
%:%1367=540%:%
%:%1372=540%:%
%:%1375=541%:%
%:%1376=542%:%
%:%1377=542%:%
%:%1378=543%:%
%:%1379=544%:%
%:%1380=545%:%
%:%1381=546%:%
%:%1382=547%:%
%:%1385=548%:%
%:%1389=548%:%
%:%1390=548%:%
%:%1391=549%:%
%:%1392=549%:%
%:%1393=550%:%
%:%1394=550%:%
%:%1395=551%:%
%:%1396=551%:%
%:%1397=551%:%
%:%1398=552%:%
%:%1399=552%:%
%:%1400=553%:%
%:%1401=553%:%
%:%1402=553%:%
%:%1403=554%:%
%:%1404=554%:%
%:%1405=554%:%
%:%1406=555%:%
%:%1407=555%:%
%:%1408=555%:%
%:%1409=556%:%
%:%1415=556%:%
%:%1418=557%:%
%:%1419=558%:%
%:%1420=558%:%
%:%1421=559%:%
%:%1422=560%:%
%:%1423=561%:%
%:%1424=562%:%
%:%1425=563%:%
%:%1428=564%:%
%:%1432=564%:%
%:%1433=564%:%
%:%1434=565%:%
%:%1435=565%:%
%:%1440=565%:%
%:%1443=566%:%
%:%1444=567%:%
%:%1445=567%:%
%:%1446=567%:%
%:%1447=567%:%
%:%1448=568%:%
%:%1449=569%:%
%:%1450=569%:%
%:%1451=570%:%
%:%1453=572%:%
%:%1454=573%:%
%:%1456=574%:%
%:%1457=574%:%
%:%1458=575%:%
%:%1459=576%:%
%:%1460=577%:%
%:%1461=578%:%
%:%1462=579%:%
%:%1465=580%:%
%:%1469=580%:%
%:%1470=580%:%
%:%1471=581%:%
%:%1472=581%:%
%:%1473=582%:%
%:%1474=582%:%
%:%1475=583%:%
%:%1476=583%:%
%:%1477=583%:%
%:%1478=584%:%
%:%1479=584%:%
%:%1480=585%:%
%:%1481=586%:%
%:%1482=586%:%
%:%1483=587%:%
%:%1484=587%:%
%:%1485=588%:%
%:%1486=588%:%
%:%1487=589%:%
%:%1488=589%:%
%:%1489=590%:%
%:%1490=590%:%
%:%1491=591%:%
%:%1492=591%:%
%:%1493=591%:%
%:%1494=592%:%
%:%1495=592%:%
%:%1496=593%:%
%:%1497=593%:%
%:%1498=593%:%
%:%1499=594%:%
%:%1500=594%:%
%:%1501=595%:%
%:%1502=595%:%
%:%1503=595%:%
%:%1504=596%:%
%:%1505=596%:%
%:%1506=596%:%
%:%1507=597%:%
%:%1508=598%:%
%:%1509=598%:%
%:%1510=598%:%
%:%1511=599%:%
%:%1512=599%:%
%:%1513=600%:%
%:%1514=600%:%
%:%1515=600%:%
%:%1516=601%:%
%:%1522=601%:%
%:%1525=602%:%
%:%1526=603%:%
%:%1527=603%:%
%:%1528=604%:%
%:%1529=605%:%
%:%1530=606%:%
%:%1533=607%:%
%:%1537=607%:%
%:%1538=607%:%
%:%1539=607%:%
%:%1548=609%:%
%:%1550=610%:%
%:%1551=610%:%
%:%1552=611%:%
%:%1553=612%:%
%:%1554=613%:%
%:%1555=614%:%
%:%1558=615%:%
%:%1562=615%:%
%:%1563=615%:%
%:%1564=616%:%
%:%1565=617%:%
%:%1566=617%:%
%:%1575=619%:%
%:%1577=621%:%
%:%1578=621%:%
%:%1579=622%:%
%:%1582=623%:%
%:%1586=623%:%
%:%1587=623%:%
%:%1592=623%:%
%:%1595=624%:%
%:%1596=625%:%
%:%1597=625%:%
%:%1598=626%:%
%:%1599=627%:%
%:%1602=628%:%
%:%1606=628%:%
%:%1607=628%:%
%:%1608=628%:%
%:%1617=630%:%
%:%1619=631%:%
%:%1620=631%:%
%:%1621=632%:%
%:%1622=633%:%
%:%1623=634%:%
%:%1626=635%:%
%:%1630=635%:%
%:%1631=635%:%
%:%1632=636%:%
%:%1633=636%:%
%:%1634=637%:%
%:%1635=637%:%
%:%1636=637%:%
%:%1637=637%:%
%:%1638=638%:%
%:%1639=638%:%
%:%1640=638%:%
%:%1641=639%:%
%:%1642=639%:%
%:%1643=640%:%
%:%1644=640%:%
%:%1645=641%:%
%:%1646=641%:%
%:%1647=641%:%
%:%1648=642%:%
%:%1649=642%:%
%:%1650=643%:%
%:%1651=643%:%
%:%1652=643%:%
%:%1653=644%:%
%:%1654=644%:%
%:%1655=645%:%
%:%1656=645%:%
%:%1657=646%:%
%:%1658=646%:%
%:%1659=646%:%
%:%1660=647%:%
%:%1661=647%:%
%:%1662=648%:%
%:%1663=648%:%
%:%1664=649%:%
%:%1665=649%:%
%:%1666=649%:%
%:%1667=650%:%
%:%1668=650%:%
%:%1669=651%:%
%:%1670=651%:%
%:%1671=651%:%
%:%1672=652%:%
%:%1673=652%:%
%:%1674=653%:%
%:%1675=653%:%
%:%1676=654%:%
%:%1677=654%:%
%:%1678=654%:%
%:%1679=655%:%
%:%1680=655%:%
%:%1681=656%:%
%:%1687=656%:%
%:%1690=657%:%
%:%1691=658%:%
%:%1692=658%:%
%:%1695=659%:%
%:%1699=659%:%
%:%1700=659%:%
%:%1709=661%:%
%:%1710=662%:%
%:%1712=663%:%
%:%1713=663%:%
%:%1714=664%:%
%:%1715=665%:%
%:%1716=666%:%
%:%1719=667%:%
%:%1723=667%:%
%:%1724=667%:%
%:%1725=668%:%
%:%1726=668%:%
%:%1727=669%:%
%:%1728=669%:%
%:%1729=670%:%
%:%1730=670%:%
%:%1731=670%:%
%:%1732=671%:%
%:%1733=671%:%
%:%1734=672%:%
%:%1735=672%:%
%:%1736=673%:%
%:%1737=673%:%
%:%1738=673%:%
%:%1739=674%:%
%:%1740=674%:%
%:%1741=675%:%
%:%1742=675%:%
%:%1743=675%:%
%:%1744=676%:%
%:%1745=676%:%
%:%1746=676%:%
%:%1747=677%:%
%:%1748=677%:%
%:%1749=677%:%
%:%1750=678%:%
%:%1751=678%:%
%:%1752=679%:%
%:%1753=679%:%
%:%1754=679%:%
%:%1755=680%:%
%:%1756=680%:%
%:%1757=680%:%
%:%1758=681%:%
%:%1759=681%:%
%:%1760=682%:%
%:%1761=682%:%
%:%1762=682%:%
%:%1763=683%:%
%:%1764=683%:%
%:%1765=684%:%
%:%1766=684%:%
%:%1767=685%:%
%:%1768=685%:%
%:%1769=686%:%
%:%1770=686%:%
%:%1771=687%:%
%:%1772=687%:%
%:%1773=687%:%
%:%1774=688%:%
%:%1775=688%:%
%:%1776=688%:%
%:%1777=689%:%
%:%1778=689%:%
%:%1779=689%:%
%:%1780=690%:%
%:%1781=690%:%
%:%1782=690%:%
%:%1783=691%:%
%:%1784=691%:%
%:%1785=691%:%
%:%1786=692%:%
%:%1787=692%:%
%:%1788=692%:%
%:%1789=693%:%
%:%1790=693%:%
%:%1791=693%:%
%:%1792=694%:%
%:%1793=694%:%
%:%1794=694%:%
%:%1795=695%:%
%:%1796=695%:%
%:%1797=695%:%
%:%1798=696%:%
%:%1808=698%:%
%:%1809=699%:%
%:%1810=700%:%
%:%1812=701%:%
%:%1813=701%:%
%:%1814=702%:%
%:%1815=703%:%
%:%1816=704%:%
%:%1817=705%:%
%:%1818=706%:%
%:%1821=707%:%
%:%1825=707%:%
%:%1826=707%:%
%:%1827=708%:%
%:%1828=708%:%
%:%1829=709%:%
%:%1830=709%:%
%:%1831=710%:%
%:%1832=711%:%
%:%1833=711%:%
%:%1834=712%:%
%:%1835=712%:%
%:%1836=712%:%
%:%1837=712%:%
%:%1838=713%:%
%:%1839=713%:%
%:%1840=714%:%
%:%1841=714%:%
%:%1842=714%:%
%:%1843=714%:%
%:%1844=715%:%
%:%1845=715%:%
%:%1846=715%:%
%:%1847=716%:%
%:%1848=717%:%
%:%1849=717%:%
%:%1850=717%:%
%:%1851=718%:%
%:%1852=718%:%
%:%1853=718%:%
%:%1854=719%:%
%:%1855=720%:%
%:%1856=720%:%
%:%1857=720%:%
%:%1858=721%:%
%:%1859=721%:%
%:%1860=721%:%
%:%1861=722%:%
%:%1862=722%:%
%:%1863=722%:%
%:%1864=723%:%
%:%1865=724%:%
%:%1866=724%:%
%:%1867=725%:%
%:%1868=725%:%
%:%1869=725%:%
%:%1870=725%:%
%:%1871=726%:%
%:%1872=726%:%
%:%1873=727%:%
%:%1874=727%:%
%:%1875=727%:%
%:%1876=728%:%
%:%1877=728%:%
%:%1878=728%:%
%:%1879=729%:%
%:%1880=729%:%
%:%1881=730%:%
%:%1882=730%:%
%:%1883=730%:%
%:%1884=731%:%
%:%1885=731%:%
%:%1886=731%:%
%:%1887=732%:%
%:%1888=732%:%
%:%1889=733%:%
%:%1890=733%:%
%:%1891=734%:%
%:%1892=735%:%
%:%1893=736%:%
%:%1894=737%:%
%:%1895=737%:%
%:%1896=738%:%
%:%1897=739%:%
%:%1898=740%:%
%:%1899=741%:%
%:%1900=742%:%
%:%1901=743%:%
%:%1902=743%:%
%:%1903=744%:%
%:%1904=744%:%
%:%1905=744%:%
%:%1906=745%:%
%:%1907=745%:%
%:%1908=746%:%
%:%1909=746%:%
%:%1910=747%:%
%:%1911=747%:%
%:%1912=748%:%
%:%1913=748%:%
%:%1914=749%:%
%:%1915=749%:%
%:%1916=750%:%
%:%1917=750%:%
%:%1918=751%:%
%:%1919=752%:%
%:%1920=753%:%
%:%1921=754%:%
%:%1922=754%:%
%:%1923=755%:%
%:%1924=755%:%
%:%1925=756%:%
%:%1926=757%:%
%:%1927=757%:%
%:%1928=758%:%
%:%1929=758%:%
%:%1930=758%:%
%:%1931=759%:%
%:%1932=760%:%
%:%1933=760%:%
%:%1934=761%:%
%:%1935=761%:%
%:%1936=761%:%
%:%1937=761%:%
%:%1938=762%:%
%:%1939=762%:%
%:%1940=763%:%
%:%1941=763%:%
%:%1942=763%:%
%:%1943=764%:%
%:%1944=764%:%
%:%1945=764%:%
%:%1946=765%:%
%:%1947=765%:%
%:%1948=766%:%
%:%1949=766%:%
%:%1950=766%:%
%:%1951=767%:%
%:%1952=767%:%
%:%1953=767%:%
%:%1954=768%:%
%:%1955=768%:%
%:%1956=768%:%
%:%1957=769%:%
%:%1958=769%:%
%:%1959=769%:%
%:%1960=770%:%
%:%1961=771%:%
%:%1962=771%:%
%:%1963=772%:%
%:%1964=772%:%
%:%1965=772%:%
%:%1966=772%:%
%:%1967=773%:%
%:%1968=773%:%
%:%1969=774%:%
%:%1970=774%:%
%:%1971=774%:%
%:%1972=774%:%
%:%1973=775%:%
%:%1974=775%:%
%:%1975=775%:%
%:%1976=776%:%
%:%1977=777%:%
%:%1978=777%:%
%:%1979=777%:%
%:%1980=778%:%
%:%1981=778%:%
%:%1982=778%:%
%:%1983=779%:%
%:%1984=779%:%
%:%1985=780%:%
%:%1986=781%:%
%:%1987=782%:%
%:%1988=782%:%
%:%1989=783%:%
%:%1995=783%:%
%:%1998=784%:%
%:%1999=785%:%
%:%2000=786%:%
%:%2001=786%:%
%:%2002=787%:%
%:%2003=788%:%
%:%2004=789%:%
%:%2007=790%:%
%:%2011=790%:%
%:%2012=790%:%
%:%2013=791%:%
%:%2014=791%:%
%:%2015=792%:%
%:%2016=792%:%
%:%2017=792%:%
%:%2018=793%:%
%:%2019=793%:%
%:%2020=793%:%
%:%2021=794%:%
%:%2022=794%:%
%:%2023=794%:%
%:%2024=794%:%
%:%2025=795%:%
%:%2026=795%:%
%:%2027=795%:%
%:%2028=796%:%
%:%2029=796%:%
%:%2030=796%:%
%:%2031=797%:%
%:%2032=797%:%
%:%2033=797%:%
%:%2034=798%:%
%:%2035=798%:%
%:%2036=799%:%
%:%2037=799%:%
%:%2038=799%:%
%:%2039=800%:%
%:%2040=800%:%
%:%2041=800%:%
%:%2042=801%:%
%:%2043=801%:%
%:%2044=802%:%
%:%2050=802%:%
%:%2053=803%:%
%:%2054=804%:%
%:%2055=804%:%
%:%2056=805%:%
%:%2057=806%:%
%:%2058=807%:%
%:%2059=808%:%
%:%2060=809%:%
%:%2061=810%:%
%:%2062=811%:%
%:%2065=812%:%
%:%2069=812%:%
%:%2070=812%:%
%:%2071=813%:%
%:%2072=813%:%
%:%2073=814%:%
%:%2074=814%:%
%:%2075=815%:%
%:%2076=815%:%
%:%2077=816%:%
%:%2078=816%:%
%:%2079=817%:%
%:%2080=817%:%
%:%2081=817%:%
%:%2082=818%:%
%:%2083=818%:%
%:%2084=818%:%
%:%2085=819%:%
%:%2086=819%:%
%:%2087=819%:%
%:%2088=820%:%
%:%2089=820%:%
%:%2090=821%:%
%:%2091=821%:%
%:%2092=822%:%
%:%2093=822%:%
%:%2094=823%:%
%:%2095=823%:%
%:%2096=823%:%
%:%2097=823%:%
%:%2098=824%:%
%:%2099=824%:%
%:%2100=824%:%
%:%2101=825%:%
%:%2102=825%:%
%:%2103=825%:%
%:%2104=826%:%
%:%2105=826%:%
%:%2106=827%:%
%:%2107=827%:%
%:%2108=827%:%
%:%2109=828%:%
%:%2110=828%:%
%:%2111=829%:%
%:%2112=829%:%
%:%2113=830%:%
%:%2114=830%:%
%:%2115=831%:%
%:%2116=831%:%
%:%2117=832%:%
%:%2118=832%:%
%:%2119=832%:%
%:%2120=832%:%
%:%2121=833%:%
%:%2122=833%:%
%:%2123=834%:%
%:%2124=834%:%
%:%2125=834%:%
%:%2126=835%:%
%:%2127=835%:%
%:%2128=835%:%
%:%2129=836%:%
%:%2130=836%:%
%:%2131=837%:%
%:%2132=838%:%
%:%2133=838%:%
%:%2134=839%:%
%:%2135=840%:%
%:%2136=840%:%
%:%2137=841%:%
%:%2138=841%:%
%:%2139=842%:%
%:%2140=842%:%
%:%2141=842%:%
%:%2142=842%:%
%:%2143=843%:%
%:%2144=843%:%
%:%2145=843%:%
%:%2146=844%:%
%:%2147=845%:%
%:%2148=845%:%
%:%2149=845%:%
%:%2150=846%:%
%:%2151=847%:%
%:%2152=847%:%
%:%2153=848%:%
%:%2154=848%:%
%:%2155=849%:%
%:%2156=849%:%
%:%2157=850%:%
%:%2158=850%:%
%:%2159=850%:%
%:%2160=851%:%
%:%2161=852%:%
%:%2162=852%:%
%:%2163=852%:%
%:%2164=853%:%
%:%2165=853%:%
%:%2166=853%:%
%:%2167=854%:%
%:%2168=854%:%
%:%2169=854%:%
%:%2170=855%:%
%:%2171=855%:%
%:%2172=856%:%
%:%2173=856%:%
%:%2174=856%:%
%:%2175=857%:%
%:%2176=857%:%
%:%2177=857%:%
%:%2178=858%:%
%:%2179=858%:%
%:%2180=858%:%
%:%2181=859%:%
%:%2182=859%:%
%:%2183=860%:%
%:%2184=860%:%
%:%2185=861%:%
%:%2186=861%:%
%:%2187=861%:%
%:%2188=862%:%
%:%2189=863%:%
%:%2190=863%:%
%:%2191=863%:%
%:%2192=864%:%
%:%2193=864%:%
%:%2194=864%:%
%:%2195=865%:%
%:%2196=865%:%
%:%2197=866%:%
%:%2198=866%:%
%:%2199=867%:%
%:%2200=867%:%
%:%2201=868%:%
%:%2202=868%:%
%:%2203=868%:%
%:%2204=869%:%
%:%2205=869%:%
%:%2206=869%:%
%:%2207=870%:%
%:%2208=870%:%
%:%2209=870%:%
%:%2210=871%:%
%:%2211=871%:%
%:%2212=872%:%
%:%2218=872%:%
%:%2221=873%:%
%:%2222=874%:%
%:%2223=874%:%
%:%2226=875%:%
%:%2230=875%:%
%:%2231=875%:%
%:%2240=877%:%
%:%2242=878%:%
%:%2243=878%:%
%:%2246=879%:%
%:%2250=879%:%
%:%2251=879%:%
%:%2252=880%:%
%:%2253=880%:%
%:%2260=880%:%
%:%2261=881%:%
%:%2262=882%:%
%:%2263=882%:%
%:%2264=883%:%
%:%2265=884%:%
%:%2268=885%:%
%:%2272=885%:%
%:%2273=885%:%
%:%2274=886%:%
%:%2275=886%:%
%:%2276=887%:%
%:%2277=887%:%
%:%2278=887%:%
%:%2279=888%:%
%:%2280=888%:%
%:%2281=888%:%
%:%2282=888%:%
%:%2283=889%:%
%:%2284=889%:%
%:%2285=890%:%
%:%2286=890%:%
%:%2287=891%:%
%:%2288=891%:%
%:%2289=892%:%
%:%2290=892%:%
%:%2291=893%:%
%:%2292=893%:%
%:%2293=894%:%
%:%2294=894%:%
%:%2295=894%:%
%:%2296=895%:%
%:%2297=895%:%
%:%2298=896%:%
%:%2299=896%:%
%:%2300=896%:%
%:%2301=897%:%
%:%2302=897%:%
%:%2303=897%:%
%:%2304=897%:%
%:%2305=898%:%
%:%2306=898%:%
%:%2307=898%:%
%:%2308=899%:%
%:%2309=899%:%
%:%2310=899%:%
%:%2311=900%:%
%:%2312=900%:%
%:%2313=901%:%
%:%2319=901%:%
%:%2322=902%:%
%:%2323=903%:%
%:%2324=903%:%
%:%2325=904%:%
%:%2326=905%:%
%:%2327=906%:%
%:%2328=907%:%
%:%2331=908%:%
%:%2335=908%:%
%:%2336=908%:%
%:%2337=908%:%
%:%2346=910%:%
%:%2347=911%:%
%:%2349=912%:%
%:%2350=912%:%
%:%2353=913%:%
%:%2357=913%:%
%:%2358=913%:%
%:%2359=914%:%
%:%2360=914%:%
%:%2361=915%:%
%:%2362=915%:%
%:%2363=915%:%
%:%2364=916%:%
%:%2365=916%:%
%:%2366=916%:%
%:%2367=916%:%
%:%2368=917%:%
%:%2369=917%:%
%:%2370=917%:%
%:%2371=918%:%
%:%2372=918%:%
%:%2373=919%:%
%:%2374=919%:%
%:%2375=920%:%
%:%2376=920%:%
%:%2377=921%:%
%:%2378=921%:%
%:%2379=921%:%
%:%2380=922%:%
%:%2381=922%:%
%:%2382=923%:%
%:%2383=923%:%
%:%2384=924%:%
%:%2385=924%:%
%:%2386=925%:%
%:%2387=925%:%
%:%2388=925%:%
%:%2389=926%:%
%:%2390=926%:%
%:%2391=926%:%
%:%2392=927%:%
%:%2393=927%:%
%:%2394=927%:%
%:%2395=928%:%
%:%2396=929%:%
%:%2397=929%:%
%:%2398=930%:%
%:%2399=930%:%
%:%2400=930%:%
%:%2401=930%:%
%:%2402=931%:%
%:%2403=931%:%
%:%2404=931%:%
%:%2405=932%:%
%:%2406=932%:%
%:%2407=932%:%
%:%2408=933%:%
%:%2409=933%:%
%:%2410=934%:%
%:%2411=934%:%
%:%2412=934%:%
%:%2413=934%:%
%:%2414=935%:%
%:%2415=935%:%
%:%2416=936%:%
%:%2417=936%:%
%:%2418=937%:%
%:%2419=937%:%
%:%2420=938%:%
%:%2421=938%:%
%:%2422=938%:%
%:%2423=939%:%
%:%2424=939%:%
%:%2425=940%:%
%:%2426=940%:%
%:%2427=940%:%
%:%2428=941%:%
%:%2429=941%:%
%:%2430=942%:%
%:%2431=943%:%
%:%2432=943%:%
%:%2433=944%:%
%:%2434=944%:%
%:%2435=945%:%
%:%2436=945%:%
%:%2437=946%:%
%:%2438=947%:%
%:%2439=947%:%
%:%2440=948%:%
%:%2441=948%:%
%:%2442=948%:%
%:%2443=949%:%
%:%2444=949%:%
%:%2445=950%:%
%:%2446=950%:%
%:%2447=950%:%
%:%2448=951%:%
%:%2449=951%:%
%:%2450=952%:%
%:%2451=952%:%
%:%2452=953%:%
%:%2453=954%:%
%:%2454=955%:%
%:%2455=955%:%
%:%2456=956%:%
%:%2457=956%:%
%:%2458=957%:%
%:%2459=957%:%
%:%2460=958%:%
%:%2461=959%:%
%:%2462=960%:%
%:%2463=960%:%
%:%2464=961%:%
%:%2465=961%:%
%:%2466=962%:%
%:%2467=962%:%
%:%2468=963%:%
%:%2469=963%:%
%:%2470=963%:%
%:%2471=964%:%
%:%2472=964%:%
%:%2473=964%:%
%:%2474=964%:%
%:%2475=965%:%
%:%2476=965%:%
%:%2477=966%:%
%:%2478=966%:%
%:%2479=967%:%
%:%2480=967%:%
%:%2481=967%:%
%:%2482=968%:%
%:%2483=968%:%
%:%2484=969%:%
%:%2485=969%:%
%:%2486=970%:%
%:%2487=970%:%
%:%2488=971%:%
%:%2489=971%:%
%:%2490=972%:%
%:%2491=973%:%
%:%2492=973%:%
%:%2493=974%:%
%:%2494=974%:%
%:%2495=974%:%
%:%2496=975%:%
%:%2497=975%:%
%:%2498=976%:%
%:%2499=976%:%
%:%2500=976%:%
%:%2501=977%:%
%:%2502=977%:%
%:%2503=978%:%
%:%2504=978%:%
%:%2505=979%:%
%:%2506=980%:%
%:%2507=981%:%
%:%2508=981%:%
%:%2509=982%:%
%:%2510=982%:%
%:%2511=983%:%
%:%2512=983%:%
%:%2513=984%:%
%:%2514=985%:%
%:%2515=986%:%
%:%2516=986%:%
%:%2517=987%:%
%:%2518=987%:%
%:%2519=988%:%
%:%2520=988%:%
%:%2521=989%:%
%:%2527=989%:%
%:%2530=990%:%
%:%2531=991%:%
%:%2532=991%:%
%:%2533=992%:%
%:%2534=993%:%
%:%2535=993%:%
%:%2536=994%:%
%:%2538=997%:%
%:%2539=998%:%
%:%2541=999%:%
%:%2542=999%:%
%:%2543=1000%:%
%:%2544=1001%:%
%:%2545=1002%:%
%:%2546=1003%:%
%:%2547=1004%:%
%:%2548=1005%:%
%:%2551=1006%:%
%:%2555=1006%:%
%:%2556=1006%:%
%:%2557=1007%:%
%:%2558=1007%:%
%:%2559=1008%:%
%:%2560=1008%:%
%:%2561=1009%:%
%:%2562=1009%:%
%:%2563=1010%:%
%:%2564=1010%:%
%:%2565=1010%:%
%:%2566=1010%:%
%:%2567=1011%:%
%:%2568=1011%:%
%:%2569=1012%:%
%:%2570=1012%:%
%:%2571=1012%:%
%:%2572=1013%:%
%:%2573=1013%:%
%:%2574=1014%:%
%:%2575=1014%:%
%:%2576=1015%:%
%:%2577=1016%:%
%:%2578=1016%:%
%:%2579=1017%:%
%:%2580=1017%:%
%:%2581=1018%:%
%:%2582=1018%:%
%:%2583=1019%:%
%:%2584=1019%:%
%:%2585=1019%:%
%:%2586=1020%:%
%:%2587=1020%:%
%:%2588=1020%:%
%:%2589=1021%:%
%:%2590=1021%:%
%:%2591=1022%:%
%:%2592=1022%:%
%:%2593=1023%:%
%:%2594=1023%:%
%:%2595=1024%:%
%:%2596=1024%:%
%:%2597=1025%:%
%:%2598=1025%:%
%:%2599=1026%:%
%:%2600=1027%:%
%:%2601=1027%:%
%:%2602=1028%:%
%:%2603=1028%:%
%:%2604=1028%:%
%:%2605=1029%:%
%:%2606=1029%:%
%:%2607=1029%:%
%:%2608=1030%:%
%:%2609=1030%:%
%:%2610=1031%:%
%:%2611=1031%:%
%:%2612=1032%:%
%:%2613=1032%:%
%:%2614=1032%:%
%:%2615=1033%:%
%:%2616=1033%:%
%:%2617=1033%:%
%:%2618=1033%:%
%:%2619=1034%:%
%:%2620=1034%:%
%:%2621=1034%:%
%:%2622=1035%:%
%:%2623=1035%:%
%:%2624=1035%:%
%:%2625=1035%:%
%:%2626=1036%:%
%:%2627=1036%:%
%:%2628=1036%:%
%:%2629=1037%:%
%:%2630=1037%:%
%:%2631=1038%:%
%:%2632=1038%:%
%:%2633=1038%:%
%:%2634=1039%:%
%:%2635=1039%:%
%:%2636=1039%:%
%:%2637=1040%:%
%:%2638=1040%:%
%:%2639=1041%:%
%:%2640=1041%:%
%:%2641=1042%:%
%:%2642=1043%:%
%:%2643=1043%:%
%:%2644=1044%:%
%:%2645=1044%:%
%:%2646=1045%:%
%:%2647=1045%:%
%:%2648=1046%:%
%:%2649=1046%:%
%:%2650=1047%:%
%:%2651=1047%:%
%:%2652=1048%:%
%:%2653=1048%:%
%:%2654=1049%:%
%:%2655=1049%:%
%:%2656=1049%:%
%:%2657=1050%:%
%:%2658=1050%:%
%:%2659=1050%:%
%:%2660=1051%:%
%:%2661=1051%:%
%:%2662=1051%:%
%:%2663=1051%:%
%:%2664=1052%:%
%:%2665=1052%:%
%:%2666=1052%:%
%:%2667=1053%:%
%:%2668=1053%:%
%:%2669=1054%:%
%:%2670=1054%:%
%:%2671=1054%:%
%:%2672=1055%:%
%:%2673=1055%:%
%:%2674=1055%:%
%:%2675=1056%:%
%:%2676=1056%:%
%:%2677=1057%:%
%:%2678=1057%:%
%:%2679=1058%:%
%:%2680=1059%:%
%:%2681=1059%:%
%:%2682=1060%:%
%:%2683=1060%:%
%:%2684=1061%:%
%:%2685=1061%:%
%:%2686=1062%:%
%:%2687=1062%:%
%:%2688=1063%:%
%:%2689=1063%:%
%:%2690=1064%:%
%:%2691=1064%:%
%:%2692=1065%:%
%:%2693=1065%:%
%:%2694=1065%:%
%:%2695=1066%:%
%:%2696=1066%:%
%:%2697=1066%:%
%:%2698=1067%:%
%:%2699=1067%:%
%:%2700=1067%:%
%:%2701=1067%:%
%:%2702=1068%:%
%:%2703=1068%:%
%:%2704=1068%:%
%:%2705=1069%:%
%:%2706=1069%:%
%:%2707=1070%:%
%:%2708=1070%:%
%:%2709=1070%:%
%:%2710=1071%:%
%:%2711=1071%:%
%:%2712=1071%:%
%:%2713=1072%:%
%:%2714=1072%:%
%:%2715=1073%:%
%:%2716=1073%:%
%:%2717=1074%:%
%:%2718=1075%:%
%:%2719=1075%:%
%:%2720=1076%:%
%:%2721=1076%:%
%:%2722=1077%:%
%:%2723=1077%:%
%:%2724=1078%:%
%:%2725=1078%:%
%:%2726=1079%:%
%:%2727=1079%:%
%:%2728=1080%:%
%:%2729=1081%:%
%:%2730=1082%:%
%:%2731=1082%:%
%:%2732=1082%:%
%:%2733=1083%:%
%:%2734=1083%:%
%:%2735=1084%:%
%:%2736=1084%:%
%:%2737=1085%:%
%:%2738=1085%:%
%:%2739=1085%:%
%:%2740=1086%:%
%:%2740=1088%:%
%:%2741=1089%:%
%:%2742=1089%:%
%:%2743=1089%:%
%:%2744=1090%:%
%:%2745=1090%:%
%:%2746=1090%:%
%:%2747=1091%:%
%:%2748=1091%:%
%:%2749=1091%:%
%:%2750=1092%:%
%:%2751=1092%:%
%:%2752=1092%:%
%:%2753=1093%:%
%:%2754=1093%:%
%:%2755=1094%:%
%:%2756=1094%:%
%:%2757=1094%:%
%:%2758=1095%:%
%:%2759=1095%:%
%:%2760=1095%:%
%:%2761=1096%:%
%:%2762=1096%:%
%:%2763=1097%:%
%:%2764=1097%:%
%:%2765=1098%:%
%:%2766=1098%:%
%:%2767=1098%:%
%:%2768=1099%:%
%:%2769=1099%:%
%:%2770=1100%:%
%:%2771=1100%:%
%:%2772=1101%:%
%:%2773=1102%:%
%:%2774=1102%:%
%:%2775=1103%:%
%:%2776=1103%:%
%:%2777=1104%:%
%:%2787=1106%:%
%:%2788=1107%:%
%:%2789=1108%:%
%:%2791=1109%:%
%:%2792=1109%:%
%:%2793=1110%:%
%:%2794=1111%:%
%:%2795=1112%:%
%:%2796=1113%:%
%:%2797=1114%:%
%:%2798=1115%:%
%:%2801=1116%:%
%:%2805=1116%:%
%:%2806=1116%:%
%:%2807=1117%:%
%:%2808=1117%:%
%:%2809=1118%:%
%:%2810=1118%:%
%:%2811=1119%:%
%:%2812=1119%:%
%:%2813=1120%:%
%:%2814=1120%:%
%:%2815=1120%:%
%:%2816=1121%:%
%:%2817=1121%:%
%:%2818=1122%:%
%:%2819=1122%:%
%:%2820=1123%:%
%:%2821=1123%:%
%:%2822=1124%:%
%:%2823=1124%:%
%:%2824=1125%:%
%:%2825=1125%:%
%:%2826=1125%:%
%:%2827=1126%:%
%:%2828=1126%:%
%:%2829=1127%:%
%:%2830=1127%:%
%:%2831=1127%:%
%:%2832=1128%:%
%:%2833=1128%:%
%:%2834=1129%:%
%:%2835=1129%:%
%:%2836=1130%:%
%:%2837=1130%:%
%:%2838=1131%:%
%:%2839=1131%:%
%:%2840=1132%:%
%:%2841=1132%:%
%:%2842=1132%:%
%:%2843=1133%:%
%:%2844=1133%:%
%:%2845=1133%:%
%:%2846=1134%:%
%:%2847=1135%:%
%:%2848=1135%:%
%:%2849=1135%:%
%:%2850=1136%:%
%:%2851=1136%:%
%:%2852=1136%:%
%:%2853=1137%:%
%:%2854=1137%:%
%:%2855=1137%:%
%:%2856=1138%:%
%:%2866=1141%:%
%:%2867=1142%:%
%:%2869=1143%:%
%:%2870=1143%:%
%:%2871=1144%:%
%:%2872=1145%:%
%:%2873=1146%:%
%:%2876=1147%:%
%:%2880=1147%:%
%:%2881=1147%:%
%:%2882=1148%:%
%:%2883=1148%:%
%:%2884=1149%:%
%:%2885=1149%:%
%:%2886=1150%:%
%:%2887=1150%:%
%:%2888=1150%:%
%:%2889=1151%:%
%:%2890=1151%:%
%:%2891=1152%:%
%:%2892=1152%:%
%:%2893=1153%:%
%:%2894=1153%:%
%:%2895=1153%:%
%:%2896=1154%:%
%:%2897=1154%:%
%:%2898=1155%:%
%:%2899=1155%:%
%:%2900=1155%:%
%:%2901=1156%:%
%:%2902=1156%:%
%:%2903=1156%:%
%:%2904=1157%:%
%:%2905=1157%:%
%:%2906=1157%:%
%:%2907=1158%:%
%:%2908=1158%:%
%:%2909=1159%:%
%:%2910=1159%:%
%:%2911=1160%:%
%:%2912=1160%:%
%:%2913=1161%:%
%:%2914=1161%:%
%:%2915=1161%:%
%:%2916=1162%:%
%:%2917=1162%:%
%:%2918=1163%:%
%:%2919=1163%:%
%:%2920=1163%:%
%:%2921=1164%:%
%:%2922=1164%:%
%:%2923=1165%:%
%:%2924=1165%:%
%:%2925=1166%:%
%:%2926=1166%:%
%:%2927=1167%:%
%:%2928=1167%:%
%:%2929=1168%:%
%:%2930=1168%:%
%:%2931=1168%:%
%:%2932=1169%:%
%:%2933=1169%:%
%:%2934=1169%:%
%:%2935=1170%:%
%:%2936=1170%:%
%:%2937=1170%:%
%:%2938=1171%:%
%:%2939=1171%:%
%:%2940=1171%:%
%:%2941=1172%:%
%:%2942=1172%:%
%:%2943=1172%:%
%:%2944=1173%:%
%:%2945=1173%:%
%:%2946=1173%:%
%:%2947=1174%:%
%:%2948=1174%:%
%:%2949=1174%:%
%:%2950=1175%:%
%:%2951=1175%:%
%:%2952=1175%:%
%:%2953=1176%:%
%:%2954=1176%:%
%:%2955=1176%:%
%:%2956=1177%:%
%:%2966=1179%:%
%:%2967=1180%:%
%:%2969=1181%:%
%:%2970=1181%:%
%:%2971=1182%:%
%:%2972=1183%:%
%:%2973=1184%:%
%:%2974=1185%:%
%:%2975=1186%:%
%:%2976=1187%:%
%:%2979=1188%:%
%:%2983=1188%:%
%:%2984=1188%:%
%:%2985=1189%:%
%:%2986=1189%:%
%:%2987=1190%:%
%:%2988=1190%:%
%:%2989=1191%:%
%:%2990=1191%:%
%:%2991=1191%:%
%:%2992=1192%:%
%:%2993=1192%:%
%:%2994=1193%:%
%:%2995=1193%:%
%:%2996=1194%:%
%:%2997=1194%:%
%:%2998=1195%:%
%:%2999=1195%:%
%:%3000=1195%:%
%:%3001=1196%:%
%:%3002=1196%:%
%:%3003=1197%:%
%:%3004=1197%:%
%:%3005=1198%:%
%:%3006=1199%:%
%:%3007=1199%:%
%:%3008=1200%:%
%:%3009=1200%:%
%:%3010=1201%:%
%:%3011=1201%:%
%:%3012=1202%:%
%:%3013=1202%:%
%:%3014=1202%:%
%:%3015=1203%:%
%:%3016=1203%:%
%:%3017=1203%:%
%:%3018=1204%:%
%:%3019=1204%:%
%:%3020=1205%:%
%:%3021=1205%:%
%:%3022=1206%:%
%:%3023=1206%:%
%:%3024=1207%:%
%:%3025=1207%:%
%:%3026=1208%:%
%:%3027=1208%:%
%:%3028=1209%:%
%:%3029=1210%:%
%:%3030=1210%:%
%:%3031=1211%:%
%:%3032=1211%:%
%:%3033=1211%:%
%:%3034=1212%:%
%:%3035=1212%:%
%:%3036=1212%:%
%:%3037=1213%:%
%:%3038=1213%:%
%:%3039=1214%:%
%:%3040=1214%:%
%:%3041=1215%:%
%:%3042=1215%:%
%:%3043=1215%:%
%:%3044=1216%:%
%:%3045=1216%:%
%:%3046=1216%:%
%:%3047=1216%:%
%:%3048=1217%:%
%:%3049=1217%:%
%:%3050=1217%:%
%:%3051=1218%:%
%:%3052=1218%:%
%:%3053=1218%:%
%:%3054=1218%:%
%:%3055=1219%:%
%:%3056=1219%:%
%:%3057=1219%:%
%:%3058=1220%:%
%:%3059=1220%:%
%:%3060=1221%:%
%:%3061=1221%:%
%:%3062=1221%:%
%:%3063=1222%:%
%:%3064=1222%:%
%:%3065=1222%:%
%:%3066=1223%:%
%:%3067=1223%:%
%:%3068=1224%:%
%:%3069=1224%:%
%:%3070=1225%:%
%:%3071=1226%:%
%:%3072=1226%:%
%:%3073=1227%:%
%:%3074=1227%:%
%:%3075=1228%:%
%:%3076=1228%:%
%:%3077=1229%:%
%:%3078=1229%:%
%:%3079=1230%:%
%:%3080=1230%:%
%:%3081=1231%:%
%:%3082=1231%:%
%:%3083=1232%:%
%:%3084=1232%:%
%:%3085=1232%:%
%:%3086=1233%:%
%:%3087=1233%:%
%:%3088=1233%:%
%:%3089=1234%:%
%:%3090=1234%:%
%:%3091=1234%:%
%:%3092=1234%:%
%:%3093=1235%:%
%:%3094=1235%:%
%:%3095=1235%:%
%:%3096=1236%:%
%:%3097=1236%:%
%:%3098=1237%:%
%:%3099=1237%:%
%:%3100=1237%:%
%:%3101=1238%:%
%:%3102=1238%:%
%:%3103=1238%:%
%:%3104=1239%:%
%:%3105=1239%:%
%:%3106=1240%:%
%:%3107=1240%:%
%:%3108=1241%:%
%:%3109=1242%:%
%:%3110=1242%:%
%:%3111=1243%:%
%:%3112=1243%:%
%:%3113=1244%:%
%:%3114=1244%:%
%:%3115=1245%:%
%:%3116=1245%:%
%:%3117=1246%:%
%:%3118=1246%:%
%:%3119=1247%:%
%:%3120=1247%:%
%:%3121=1248%:%
%:%3122=1248%:%
%:%3123=1248%:%
%:%3124=1249%:%
%:%3125=1249%:%
%:%3126=1249%:%
%:%3127=1250%:%
%:%3128=1250%:%
%:%3129=1250%:%
%:%3130=1250%:%
%:%3131=1251%:%
%:%3132=1251%:%
%:%3133=1251%:%
%:%3134=1252%:%
%:%3135=1252%:%
%:%3136=1253%:%
%:%3137=1253%:%
%:%3138=1253%:%
%:%3139=1254%:%
%:%3140=1254%:%
%:%3141=1254%:%
%:%3142=1255%:%
%:%3143=1255%:%
%:%3144=1256%:%
%:%3145=1256%:%
%:%3146=1257%:%
%:%3147=1258%:%
%:%3148=1258%:%
%:%3149=1259%:%
%:%3150=1259%:%
%:%3151=1260%:%
%:%3152=1260%:%
%:%3153=1261%:%
%:%3154=1261%:%
%:%3155=1262%:%
%:%3156=1263%:%
%:%3157=1264%:%
%:%3158=1264%:%
%:%3159=1264%:%
%:%3160=1265%:%
%:%3161=1265%:%
%:%3162=1266%:%
%:%3163=1266%:%
%:%3164=1266%:%
%:%3165=1267%:%
%:%3166=1267%:%
%:%3167=1268%:%
%:%3168=1268%:%
%:%3169=1269%:%
%:%3170=1269%:%
%:%3171=1269%:%
%:%3172=1270%:%
%:%3173=1270%:%
%:%3174=1271%:%
%:%3184=1273%:%
%:%3185=1274%:%
%:%3186=1275%:%
%:%3187=1276%:%
%:%3189=1277%:%
%:%3190=1277%:%
%:%3191=1278%:%
%:%3192=1279%:%
%:%3193=1280%:%
%:%3194=1281%:%
%:%3195=1282%:%
%:%3198=1283%:%
%:%3202=1283%:%
%:%3203=1283%:%
%:%3204=1284%:%
%:%3205=1284%:%
%:%3206=1285%:%
%:%3207=1285%:%
%:%3208=1285%:%
%:%3209=1286%:%
%:%3210=1286%:%
%:%3211=1287%:%
%:%3212=1287%:%
%:%3213=1287%:%
%:%3214=1288%:%
%:%3215=1288%:%
%:%3216=1289%:%
%:%3217=1289%:%
%:%3218=1290%:%
%:%3219=1291%:%
%:%3220=1292%:%
%:%3221=1293%:%
%:%3222=1294%:%
%:%3223=1295%:%
%:%3224=1295%:%
%:%3225=1296%:%
%:%3226=1296%:%
%:%3227=1297%:%
%:%3228=1297%:%
%:%3229=1298%:%
%:%3230=1298%:%
%:%3231=1299%:%
%:%3232=1299%:%
%:%3233=1299%:%
%:%3234=1300%:%
%:%3235=1300%:%
%:%3236=1300%:%
%:%3237=1301%:%
%:%3238=1301%:%
%:%3239=1301%:%
%:%3240=1302%:%
%:%3241=1302%:%
%:%3242=1302%:%
%:%3243=1303%:%
%:%3253=1306%:%
%:%3254=1307%:%
%:%3255=1308%:%
%:%3256=1309%:%
%:%3258=1310%:%
%:%3259=1310%:%
%:%3260=1311%:%
%:%3261=1312%:%
%:%3262=1313%:%
%:%3263=1314%:%
%:%3264=1315%:%
%:%3265=1316%:%
%:%3268=1317%:%
%:%3272=1317%:%
%:%3273=1317%:%
%:%3274=1318%:%
%:%3275=1318%:%
%:%3276=1319%:%
%:%3277=1319%:%
%:%3278=1320%:%
%:%3279=1320%:%
%:%3280=1320%:%
%:%3281=1320%:%
%:%3282=1321%:%
%:%3283=1321%:%
%:%3284=1322%:%
%:%3285=1322%:%
%:%3286=1322%:%
%:%3287=1322%:%
%:%3288=1323%:%
%:%3289=1323%:%
%:%3290=1324%:%
%:%3291=1325%:%
%:%3292=1325%:%
%:%3293=1326%:%
%:%3294=1326%:%
%:%3295=1327%:%
%:%3296=1327%:%
%:%3297=1327%:%
%:%3298=1328%:%
%:%3299=1329%:%
%:%3300=1329%:%
%:%3301=1330%:%
%:%3302=1330%:%
%:%3303=1331%:%
%:%3304=1331%:%
%:%3305=1332%:%
%:%3306=1332%:%
%:%3307=1333%:%
%:%3308=1334%:%
%:%3309=1334%:%
%:%3310=1335%:%
%:%3311=1335%:%
%:%3312=1336%:%
%:%3313=1336%:%
%:%3314=1337%:%
%:%3315=1337%:%
%:%3316=1338%:%
%:%3317=1338%:%
%:%3318=1339%:%
%:%3319=1340%:%
%:%3320=1340%:%
%:%3321=1341%:%
%:%3322=1341%:%
%:%3323=1342%:%
%:%3324=1342%:%
%:%3325=1342%:%
%:%3326=1343%:%
%:%3327=1343%:%
%:%3328=1343%:%
%:%3329=1344%:%
%:%3335=1344%:%
%:%3338=1345%:%
%:%3339=1346%:%
%:%3340=1347%:%
%:%3341=1347%:%
%:%3342=1348%:%
%:%3343=1349%:%
%:%3344=1350%:%
%:%3345=1351%:%
%:%3346=1352%:%
%:%3347=1353%:%
%:%3348=1354%:%
%:%3349=1355%:%
%:%3352=1356%:%
%:%3356=1356%:%
%:%3357=1356%:%
%:%3358=1357%:%
%:%3359=1357%:%
%:%3360=1358%:%
%:%3361=1358%:%
%:%3362=1359%:%
%:%3363=1359%:%
%:%3364=1359%:%
%:%3365=1360%:%
%:%3366=1360%:%
%:%3367=1360%:%
%:%3368=1361%:%
%:%3369=1361%:%
%:%3370=1361%:%
%:%3371=1362%:%
%:%3372=1362%:%
%:%3373=1363%:%
%:%3374=1364%:%
%:%3375=1365%:%
%:%3376=1365%:%
%:%3377=1366%:%
%:%3378=1366%:%
%:%3379=1367%:%
%:%3380=1367%:%
%:%3381=1368%:%
%:%3382=1368%:%
%:%3383=1368%:%
%:%3384=1369%:%
%:%3385=1369%:%
%:%3386=1369%:%
%:%3387=1369%:%
%:%3388=1370%:%
%:%3389=1370%:%
%:%3390=1371%:%
%:%3391=1371%:%
%:%3392=1371%:%
%:%3393=1372%:%
%:%3394=1372%:%
%:%3395=1373%:%
%:%3396=1373%:%
%:%3397=1373%:%
%:%3398=1374%:%
%:%3399=1374%:%
%:%3400=1375%:%
%:%3401=1375%:%
%:%3402=1376%:%
%:%3403=1376%:%
%:%3404=1377%:%
%:%3405=1377%:%
%:%3406=1378%:%
%:%3407=1378%:%
%:%3408=1378%:%
%:%3409=1378%:%
%:%3410=1379%:%
%:%3411=1379%:%
%:%3412=1380%:%
%:%3413=1380%:%
%:%3414=1380%:%
%:%3415=1381%:%
%:%3416=1381%:%
%:%3417=1381%:%
%:%3418=1382%:%
%:%3419=1382%:%
%:%3420=1383%:%
%:%3421=1383%:%
%:%3422=1384%:%
%:%3423=1384%:%
%:%3424=1385%:%
%:%3425=1385%:%
%:%3426=1385%:%
%:%3427=1385%:%
%:%3428=1386%:%
%:%3429=1386%:%
%:%3430=1387%:%
%:%3431=1388%:%
%:%3432=1388%:%
%:%3433=1389%:%
%:%3434=1390%:%
%:%3435=1391%:%
%:%3436=1391%:%
%:%3437=1392%:%
%:%3438=1392%:%
%:%3439=1392%:%
%:%3440=1393%:%
%:%3441=1394%:%
%:%3442=1394%:%
%:%3443=1395%:%
%:%3444=1395%:%
%:%3445=1395%:%
%:%3446=1396%:%
%:%3447=1396%:%
%:%3448=1397%:%
%:%3449=1398%:%
%:%3450=1398%:%
%:%3451=1399%:%
%:%3452=1399%:%
%:%3453=1400%:%
%:%3454=1400%:%
%:%3455=1400%:%
%:%3456=1401%:%
%:%3457=1401%:%
%:%3458=1401%:%
%:%3459=1402%:%
%:%3460=1402%:%
%:%3461=1403%:%
%:%3462=1403%:%
%:%3463=1404%:%
%:%3464=1404%:%
%:%3465=1404%:%
%:%3466=1405%:%
%:%3467=1405%:%
%:%3468=1406%:%
%:%3469=1407%:%
%:%3470=1408%:%
%:%3471=1409%:%
%:%3472=1409%:%
%:%3473=1410%:%
%:%3483=1412%:%
%:%3484=1413%:%
%:%3485=1414%:%
%:%3486=1415%:%
%:%3488=1416%:%
%:%3489=1416%:%
%:%3490=1417%:%
%:%3491=1418%:%
%:%3492=1419%:%
%:%3493=1420%:%
%:%3494=1421%:%
%:%3495=1422%:%
%:%3496=1423%:%
%:%3497=1424%:%
%:%3500=1425%:%
%:%3504=1425%:%
%:%3505=1425%:%
%:%3506=1426%:%
%:%3507=1426%:%
%:%3508=1427%:%
%:%3509=1427%:%
%:%3510=1427%:%
%:%3511=1428%:%
%:%3512=1428%:%
%:%3513=1428%:%
%:%3514=1429%:%
%:%3515=1429%:%
%:%3516=1429%:%
%:%3517=1429%:%
%:%3518=1430%:%
%:%3519=1430%:%
%:%3520=1430%:%
%:%3521=1431%:%
%:%3522=1431%:%
%:%3523=1431%:%
%:%3524=1432%:%
%:%3525=1432%:%
%:%3526=1432%:%
%:%3527=1433%:%
%:%3528=1433%:%
%:%3529=1433%:%
%:%3530=1434%:%
%:%3536=1434%:%
%:%3539=1435%:%
%:%3540=1436%:%
%:%3541=1436%:%
%:%3542=1437%:%
%:%3543=1438%:%
%:%3544=1439%:%
%:%3547=1440%:%
%:%3551=1440%:%
%:%3552=1440%:%
%:%3553=1441%:%
%:%3554=1441%:%
%:%3555=1442%:%
%:%3556=1442%:%
%:%3557=1443%:%
%:%3558=1443%:%
%:%3559=1444%:%
%:%3560=1444%:%
%:%3561=1445%:%
%:%3562=1446%:%
%:%3563=1446%:%
%:%3564=1447%:%
%:%3565=1447%:%
%:%3566=1447%:%
%:%3567=1448%:%
%:%3568=1448%:%
%:%3569=1448%:%
%:%3570=1449%:%
%:%3571=1449%:%
%:%3572=1450%:%
%:%3573=1450%:%
%:%3574=1451%:%
%:%3575=1451%:%
%:%3576=1451%:%
%:%3577=1452%:%
%:%3578=1452%:%
%:%3579=1452%:%
%:%3580=1453%:%
%:%3581=1453%:%
%:%3582=1454%:%
%:%3583=1454%:%
%:%3584=1455%:%
%:%3585=1455%:%
%:%3586=1455%:%
%:%3587=1456%:%
%:%3588=1456%:%
%:%3589=1456%:%
%:%3590=1457%:%
%:%3591=1457%:%
%:%3592=1458%:%
%:%3593=1458%:%
%:%3594=1459%:%
%:%3595=1459%:%
%:%3596=1459%:%
%:%3597=1460%:%
%:%3598=1460%:%
%:%3599=1461%:%
%:%3600=1461%:%
%:%3601=1461%:%
%:%3602=1462%:%
%:%3603=1462%:%
%:%3604=1463%:%
%:%3605=1463%:%
%:%3606=1464%:%
%:%3607=1464%:%
%:%3608=1464%:%
%:%3609=1465%:%
%:%3610=1466%:%
%:%3611=1466%:%
%:%3612=1467%:%
%:%3613=1467%:%
%:%3614=1467%:%
%:%3615=1468%:%
%:%3616=1468%:%
%:%3617=1468%:%
%:%3618=1468%:%
%:%3619=1469%:%
%:%3620=1469%:%
%:%3621=1470%:%
%:%3631=1472%:%
%:%3632=1473%:%
%:%3634=1474%:%
%:%3635=1474%:%
%:%3636=1475%:%
%:%3637=1476%:%
%:%3638=1477%:%
%:%3639=1478%:%
%:%3640=1479%:%
%:%3643=1480%:%
%:%3647=1480%:%
%:%3648=1480%:%
%:%3649=1481%:%
%:%3650=1481%:%
%:%3651=1482%:%
%:%3652=1482%:%
%:%3653=1483%:%
%:%3654=1483%:%
%:%3655=1483%:%
%:%3656=1484%:%
%:%3657=1484%:%
%:%3658=1485%:%
%:%3659=1485%:%
%:%3660=1486%:%
%:%3661=1486%:%
%:%3662=1487%:%
%:%3663=1487%:%
%:%3664=1487%:%
%:%3665=1488%:%
%:%3666=1488%:%
%:%3667=1489%:%
%:%3668=1489%:%
%:%3669=1490%:%
%:%3670=1490%:%
%:%3671=1490%:%
%:%3672=1491%:%
%:%3673=1491%:%
%:%3674=1491%:%
%:%3675=1491%:%
%:%3676=1492%:%
%:%3677=1492%:%
%:%3678=1492%:%
%:%3679=1493%:%
%:%3680=1493%:%
%:%3681=1493%:%
%:%3682=1493%:%
%:%3683=1494%:%
%:%3684=1494%:%
%:%3685=1494%:%
%:%3686=1495%:%
%:%3687=1495%:%
%:%3688=1496%:%
%:%3689=1496%:%
%:%3690=1497%:%
%:%3691=1497%:%
%:%3692=1498%:%
%:%3693=1498%:%
%:%3694=1498%:%
%:%3695=1499%:%
%:%3696=1499%:%
%:%3697=1499%:%
%:%3698=1500%:%
%:%3699=1500%:%
%:%3700=1501%:%
%:%3701=1501%:%
%:%3702=1502%:%
%:%3703=1503%:%
%:%3704=1503%:%
%:%3705=1504%:%
%:%3706=1504%:%
%:%3707=1505%:%
%:%3713=1505%:%
%:%3716=1506%:%
%:%3717=1507%:%
%:%3718=1507%:%
%:%3719=1508%:%
%:%3720=1509%:%
%:%3723=1510%:%
%:%3727=1510%:%
%:%3728=1510%:%
%:%3729=1510%:%
%:%3734=1510%:%
%:%3737=1511%:%
%:%3738=1512%:%
%:%3739=1512%:%
%:%3740=1513%:%
%:%3741=1514%:%
%:%3742=1515%:%
%:%3743=1516%:%
%:%3744=1517%:%
%:%3747=1518%:%
%:%3751=1518%:%
%:%3752=1518%:%
%:%3753=1519%:%
%:%3754=1519%:%
%:%3755=1520%:%
%:%3756=1520%:%
%:%3757=1520%:%
%:%3758=1521%:%
%:%3759=1521%:%
%:%3760=1521%:%
%:%3761=1522%:%
%:%3762=1522%:%
%:%3763=1522%:%
%:%3764=1523%:%
%:%3765=1523%:%
%:%3766=1524%:%
%:%3767=1524%:%
%:%3768=1525%:%
%:%3769=1525%:%
%:%3770=1526%:%
%:%3771=1527%:%
%:%3772=1527%:%
%:%3773=1528%:%
%:%3774=1528%:%
%:%3775=1528%:%
%:%3776=1528%:%
%:%3777=1529%:%
%:%3778=1529%:%
%:%3779=1530%:%
%:%3780=1530%:%
%:%3781=1530%:%
%:%3782=1531%:%
%:%3783=1531%:%
%:%3785=1533%:%
%:%3787=1534%:%
%:%3788=1534%:%
%:%3789=1535%:%
%:%3790=1535%:%
%:%3791=1536%:%
%:%3792=1536%:%
%:%3793=1536%:%
%:%3794=1537%:%
%:%3795=1537%:%
%:%3796=1537%:%
%:%3797=1538%:%
%:%3798=1538%:%
%:%3799=1539%:%
%:%3800=1539%:%
%:%3801=1540%:%
%:%3802=1540%:%
%:%3803=1541%:%
%:%3804=1541%:%
%:%3805=1542%:%
%:%3806=1542%:%
%:%3807=1542%:%
%:%3809=1544%:%
%:%3811=1545%:%
%:%3812=1545%:%
%:%3813=1546%:%
%:%3814=1546%:%
%:%3815=1547%:%
%:%3816=1547%:%
%:%3817=1547%:%
%:%3818=1548%:%
%:%3819=1548%:%
%:%3820=1548%:%
%:%3821=1549%:%
%:%3822=1549%:%
%:%3823=1549%:%
%:%3824=1550%:%
%:%3825=1550%:%
%:%3826=1551%:%
%:%3827=1551%:%
%:%3828=1552%:%
%:%3829=1553%:%
%:%3830=1553%:%
%:%3831=1554%:%
%:%3832=1554%:%
%:%3833=1555%:%
%:%3834=1556%:%
%:%3835=1556%:%
%:%3836=1557%:%
%:%3837=1558%:%
%:%3838=1558%:%
%:%3839=1559%:%
%:%3840=1559%:%
%:%3841=1559%:%
%:%3842=1560%:%
%:%3843=1560%:%
%:%3844=1561%:%
%:%3845=1561%:%
%:%3846=1562%:%
%:%3847=1562%:%
%:%3848=1563%:%
%:%3849=1564%:%
%:%3850=1564%:%
%:%3851=1565%:%
%:%3852=1565%:%
%:%3853=1565%:%
%:%3854=1566%:%
%:%3855=1566%:%
%:%3856=1566%:%
%:%3857=1567%:%
%:%3858=1567%:%
%:%3859=1568%:%
%:%3860=1568%:%
%:%3861=1568%:%
%:%3862=1569%:%
%:%3863=1569%:%
%:%3864=1569%:%
%:%3865=1570%:%
%:%3866=1570%:%
%:%3867=1571%:%
%:%3868=1572%:%
%:%3869=1572%:%
%:%3870=1573%:%
%:%3871=1573%:%
%:%3872=1573%:%
%:%3873=1574%:%
%:%3874=1575%:%
%:%3875=1575%:%
%:%3876=1576%:%
%:%3877=1576%:%
%:%3878=1577%:%
%:%3879=1577%:%
%:%3881=1578%:%
%:%3883=1579%:%
%:%3884=1579%:%
%:%3885=1579%:%
%:%3886=1580%:%
%:%3887=1581%:%
%:%3888=1581%:%
%:%3889=1582%:%
%:%3890=1583%:%
%:%3891=1583%:%
%:%3892=1584%:%
%:%3893=1584%:%
%:%3894=1584%:%
%:%3895=1585%:%
%:%3896=1586%:%
%:%3897=1586%:%
%:%3898=1587%:%
%:%3899=1588%:%
%:%3900=1588%:%
%:%3901=1589%:%
%:%3902=1590%:%
%:%3903=1590%:%
%:%3904=1591%:%
%:%3905=1591%:%
%:%3906=1591%:%
%:%3907=1592%:%
%:%3908=1592%:%
%:%3909=1593%:%
%:%3910=1593%:%
%:%3911=1593%:%
%:%3912=1594%:%
%:%3913=1594%:%
%:%3914=1595%:%
%:%3915=1596%:%
%:%3916=1597%:%
%:%3917=1597%:%
%:%3918=1598%:%
%:%3919=1599%:%
%:%3920=1599%:%
%:%3921=1600%:%
%:%3922=1600%:%
%:%3923=1600%:%
%:%3924=1601%:%
%:%3925=1601%:%
%:%3926=1601%:%
%:%3927=1602%:%
%:%3928=1602%:%
%:%3929=1603%:%
%:%3930=1603%:%
%:%3931=1604%:%
%:%3932=1604%:%
%:%3933=1604%:%
%:%3934=1605%:%
%:%3935=1605%:%
%:%3936=1606%:%
%:%3937=1607%:%
%:%3938=1607%:%
%:%3939=1608%:%
%:%3940=1608%:%
%:%3941=1608%:%
%:%3942=1609%:%
%:%3943=1609%:%
%:%3944=1610%:%
%:%3945=1610%:%
%:%3946=1611%:%
%:%3947=1611%:%
%:%3948=1612%:%
%:%3949=1612%:%
%:%3950=1613%:%
%:%3951=1613%:%
%:%3952=1614%:%
%:%3953=1614%:%
%:%3954=1614%:%
%:%3955=1615%:%
%:%3956=1615%:%
%:%3957=1615%:%
%:%3958=1616%:%
%:%3959=1616%:%
%:%3960=1617%:%
%:%3961=1617%:%
%:%3962=1618%:%
%:%3963=1618%:%
%:%3964=1618%:%
%:%3965=1619%:%
%:%3966=1619%:%
%:%3967=1620%:%
%:%3968=1620%:%
%:%3969=1621%:%
%:%3970=1621%:%
%:%3971=1622%:%
%:%3972=1622%:%
%:%3973=1623%:%
%:%3974=1623%:%
%:%3975=1624%:%
%:%3976=1624%:%
%:%3977=1624%:%
%:%3978=1625%:%
%:%3979=1625%:%
%:%3980=1626%:%
%:%3981=1626%:%
%:%3982=1626%:%
%:%3983=1627%:%
%:%3984=1627%:%
%:%3985=1627%:%
%:%3986=1628%:%
%:%3987=1628%:%
%:%3988=1628%:%
%:%3989=1629%:%
%:%3990=1629%:%
%:%3991=1630%:%
%:%3992=1631%:%
%:%3993=1632%:%
%:%3994=1632%:%
%:%3995=1633%:%
%:%3996=1633%:%
%:%3997=1633%:%
%:%3998=1634%:%
%:%3999=1634%:%
%:%4000=1635%:%
%:%4001=1636%:%
%:%4002=1636%:%
%:%4003=1637%:%
%:%4004=1637%:%
%:%4005=1638%:%
%:%4006=1638%:%
%:%4007=1639%:%
%:%4008=1639%:%
%:%4010=1640%:%
%:%4012=1641%:%
%:%4013=1641%:%
%:%4014=1641%:%
%:%4015=1642%:%
%:%4016=1642%:%
%:%4017=1643%:%
%:%4018=1643%:%
%:%4019=1644%:%
%:%4020=1644%:%
%:%4021=1644%:%
%:%4022=1645%:%
%:%4023=1645%:%
%:%4024=1646%:%
%:%4025=1647%:%
%:%4026=1648%:%
%:%4027=1648%:%
%:%4028=1649%:%
%:%4029=1650%:%
%:%4030=1650%:%
%:%4031=1651%:%
%:%4032=1651%:%
%:%4033=1651%:%
%:%4034=1652%:%
%:%4035=1652%:%
%:%4036=1652%:%
%:%4037=1653%:%
%:%4038=1653%:%
%:%4039=1654%:%
%:%4040=1654%:%
%:%4041=1655%:%
%:%4042=1655%:%
%:%4043=1656%:%
%:%4044=1656%:%
%:%4045=1657%:%
%:%4046=1657%:%
%:%4047=1658%:%
%:%4048=1658%:%
%:%4049=1658%:%
%:%4050=1659%:%
%:%4051=1659%:%
%:%4052=1660%:%
%:%4053=1661%:%
%:%4054=1661%:%
%:%4055=1662%:%
%:%4056=1662%:%
%:%4057=1662%:%
%:%4058=1663%:%
%:%4059=1663%:%
%:%4060=1664%:%
%:%4061=1665%:%
%:%4062=1666%:%
%:%4063=1666%:%
%:%4064=1667%:%
%:%4065=1667%:%
%:%4066=1668%:%
%:%4067=1668%:%
%:%4068=1669%:%
%:%4069=1669%:%
%:%4070=1669%:%
%:%4071=1670%:%
%:%4072=1670%:%
%:%4073=1671%:%
%:%4074=1671%:%
%:%4075=1672%:%
%:%4076=1672%:%
%:%4077=1673%:%
%:%4078=1673%:%
%:%4079=1674%:%
%:%4080=1674%:%
%:%4081=1675%:%
%:%4082=1675%:%
%:%4083=1675%:%
%:%4084=1676%:%
%:%4085=1676%:%
%:%4086=1677%:%
%:%4087=1677%:%
%:%4088=1677%:%
%:%4089=1678%:%
%:%4090=1678%:%
%:%4091=1678%:%
%:%4092=1679%:%
%:%4093=1679%:%
%:%4094=1679%:%
%:%4095=1680%:%
%:%4096=1680%:%
%:%4097=1681%:%
%:%4098=1682%:%
%:%4099=1683%:%
%:%4100=1683%:%
%:%4101=1684%:%
%:%4102=1684%:%
%:%4103=1684%:%
%:%4104=1685%:%
%:%4105=1685%:%
%:%4106=1686%:%
%:%4107=1687%:%
%:%4108=1688%:%
%:%4109=1689%:%
%:%4110=1690%:%
%:%4111=1690%:%
%:%4112=1691%:%
%:%4113=1691%:%
%:%4114=1692%:%
%:%4115=1692%:%
%:%4116=1693%:%
%:%4122=1693%:%
%:%4125=1694%:%
%:%4126=1695%:%
%:%4127=1695%:%
%:%4128=1696%:%
%:%4129=1697%:%
%:%4132=1698%:%
%:%4136=1698%:%
%:%4137=1698%:%
%:%4138=1699%:%
%:%4139=1699%:%
%:%4140=1699%:%
%:%4145=1699%:%
%:%4148=1700%:%
%:%4149=1701%:%
%:%4150=1701%:%
%:%4151=1702%:%
%:%4152=1703%:%
%:%4153=1704%:%
%:%4154=1704%:%
%:%4157=1705%:%
%:%4161=1705%:%
%:%4162=1705%:%
%:%4171=1708%:%
%:%4173=1709%:%
%:%4174=1709%:%
%:%4175=1710%:%
%:%4176=1711%:%
%:%4177=1712%:%
%:%4180=1713%:%
%:%4184=1713%:%
%:%4185=1713%:%
%:%4186=1714%:%
%:%4187=1714%:%
%:%4188=1715%:%
%:%4189=1715%:%
%:%4190=1715%:%
%:%4191=1715%:%
%:%4192=1716%:%
%:%4193=1716%:%
%:%4194=1716%:%
%:%4195=1717%:%
%:%4196=1717%:%
%:%4197=1717%:%
%:%4198=1717%:%
%:%4199=1718%:%
%:%4205=1718%:%
%:%4208=1719%:%
%:%4209=1720%:%
%:%4210=1721%:%
%:%4211=1721%:%
%:%4212=1722%:%
%:%4213=1723%:%
%:%4216=1724%:%
%:%4220=1724%:%
%:%4221=1724%:%
%:%4222=1725%:%
%:%4223=1725%:%
%:%4224=1726%:%
%:%4225=1727%:%
%:%4226=1727%:%
%:%4231=1727%:%
%:%4234=1728%:%
%:%4235=1729%:%
%:%4236=1729%:%
%:%4237=1730%:%
%:%4238=1731%:%
%:%4241=1732%:%
%:%4245=1732%:%
%:%4246=1732%:%
%:%4247=1733%:%
%:%4248=1733%:%
%:%4249=1733%:%
%:%4254=1733%:%
%:%4257=1734%:%
%:%4258=1734%:%
%:%4259=1735%:%
%:%4260=1735%:%
%:%4261=1735%:%
%:%4262=1735%:%
%:%4265=1736%:%
%:%4270=1737%:%
%:%4271=1737%:%
%:%4272=1737%:%



% optional bibliography
\bibliographystyle{abbrv}
\bibliography{root}

\end{document}

%%% Local Variables:
%%% mode: latex
%%% TeX-master: t
%%% End:
\endinput
%:%file=~/work/temp_planning_certification/temporal-planning-certification/lib/temporal-pddl-semantics/document/root.tex%:%
\endinput
%:%file=~/work/temp_planning_certification/temporal-planning-certification/lib/temporal-pddl-semantics/document/root.tex%:%
